\section{Overview of the Formalisation}

This section provides an overview of the formalisation in Agda.
We show how propositions in \cref{sec:name-embedding} and \cref{sec:value-embedding} are formalised in Agda.
We also give the main helper lemmas that are needed to prove these propositions.
The formalisation is available on GitHub\footnote{\url{www.thisurlwillbeavailablelater.com}}


\subsection{File Structure}
\label{sec:file-structure}
The Agda formalisation is structured into a number of files. Here, we go over the file structure.
The project has the following three main directories: examples, formalisation and proofs.
The \texttt{examples} directory contains some examples of \lterms, substitution, embeddings and reductions.
The \texttt{formalisation} directory is the most important directory and contains the formalisation of \lterms, substitution, embeddings and reductions.
Here, we list all files in this directory and give a short description of their contents:
\begin{itemize}
  \item \texttt{boxCalc.agda}: Contains the definition of \lterms in \lab, the definition of substitution and the definition of values in \lab.
  \item \texttt{boxCalcRed.agda}: Contains all the reduction relations in \lab.
  \item \texttt{calc.agda}: Contains the definition of \lterms in the standard \lc, the definition of substitution and the definition of values.
  \item \texttt{embedCBNIntoCBB.agda}: Contains the definition of Girard's embedding.
  \item \texttt{embedCBVIntoCBB.agda}: Contains the definition of Gödel's embedding.
  \item \texttt{nameCalcRed.agda}: Contains the reduction relations for \lan.
  \item \texttt{valueCalcRed.agda}: Contains the reduction relations for \lav.
\end{itemize}

The \texttt{proofs} directory contains the proofs of the propositions in \cref{sec:name-embedding} and \cref{sec:value-embedding} in the files \texttt{nameProofs.agda} and \texttt{valueProofs.agda} respectively.

\subsection{Helper Lemmas}
\label{sec:helper-lemmas}
This subsection goes over some helper lemmas needed to prove the main reduction properties of \lab. 
The first lemma we discuss is the lemma \textsf{pres-subst}, which states that the order of embedding and substitution does not matter.
So substitution is preserved and reflected under embedding. 
We prove this lemma by using another lemma, \textsf{pres-simultaneous-subst}, which states that simultaneous substitution is preserved and reflected under embedding.
This can be proven by induction on the structure of the \lterm in which substitution is done. However, in \cref{def:substitution:lab}, we defined substitution in terms of \textsf{extsb},
which is defined in terms of \textsf{renameb} and \textsf{extb}. 
Therefore, we also make the following helper lemmas: \textsf{pres-ext}, \textsf{pres-rename} and \textsf{pres-exts}.

We only show the lemmas \textsf{pres-subst} and \textsf{pres-simultaneous-subst} for Girard's embedding, since the idea for Gödel's translation is the same.
\Cref{lem:pres-subst} shows proof of the \textsf{pres-subst} lemma, and \cref{lem:pres-simultaneous-subst} shows the proof of the \textsf{pres-simultaneous-subst} lemma.
Lemma \textsf{pres-subst-func} is a helper lemma in \cref{lem:pres-subst} that proofs assumption $h$ in \cref{lem:pres-simultaneous-subst}.
This assumption shows how functions fb' and f' are related to each other, which is needed to prove \textsf{pres-simultaneous-subst}.

\begin{lemma}
  \label{lem:pres-subst}
  \textbf{Preservation of single substitution}:
  \small
  \begin{code}[hide]%
\>[0]\AgdaKeyword{module}\AgdaSpace{}%
\AgdaModule{nameProofs1}\AgdaSpace{}%
\AgdaKeyword{where}\<%
\\
%
\\[\AgdaEmptyExtraSkip]%
\>[0]\AgdaKeyword{open}\AgdaSpace{}%
\AgdaKeyword{import}\AgdaSpace{}%
\AgdaModule{boxCalc}\<%
\\
\>[0]\AgdaKeyword{open}\AgdaSpace{}%
\AgdaKeyword{import}\AgdaSpace{}%
\AgdaModule{calc}\<%
\\
\>[0]\AgdaKeyword{open}\AgdaSpace{}%
\AgdaKeyword{import}\AgdaSpace{}%
\AgdaModule{embedCBNIntoCBB}\<%
\\
\>[0]\AgdaKeyword{open}\AgdaSpace{}%
\AgdaKeyword{import}\AgdaSpace{}%
\AgdaModule{nameCalcRed}\<%
\\
\>[0]\AgdaKeyword{open}\AgdaSpace{}%
\AgdaKeyword{import}\AgdaSpace{}%
\AgdaModule{boxCalcRed}\<%
\\
%
\\[\AgdaEmptyExtraSkip]%
\>[0]\AgdaKeyword{open}\AgdaSpace{}%
\AgdaKeyword{import}\AgdaSpace{}%
\AgdaModule{Data.Empty}\AgdaSpace{}%
\AgdaKeyword{using}\AgdaSpace{}%
\AgdaSymbol{(}\AgdaFunction{⊥}\AgdaSymbol{;}\AgdaSpace{}%
\AgdaFunction{⊥-elim}\AgdaSymbol{)}\<%
\\
\>[0]\AgdaKeyword{import}\AgdaSpace{}%
\AgdaModule{Relation.Binary.PropositionalEquality}\AgdaSpace{}%
\AgdaSymbol{as}\AgdaSpace{}%
\AgdaModule{Eq}\<%
\\
\>[0]\AgdaKeyword{open}\AgdaSpace{}%
\AgdaModule{Eq}\AgdaSpace{}%
\AgdaKeyword{using}\AgdaSpace{}%
\AgdaSymbol{(}\AgdaOperator{\AgdaDatatype{\AgdaUnderscore{}≡\AgdaUnderscore{}}}\AgdaSymbol{;}\AgdaSpace{}%
\AgdaInductiveConstructor{refl}\AgdaSymbol{;}\AgdaSpace{}%
\AgdaFunction{cong}\AgdaSymbol{;}\AgdaSpace{}%
\AgdaFunction{cong₂}\AgdaSymbol{;}\AgdaSpace{}%
\AgdaFunction{sym}\AgdaSymbol{)}\<%
\\
%
\\[\AgdaEmptyExtraSkip]%
\>[0]\AgdaKeyword{open}\AgdaSpace{}%
\AgdaModule{Eq.≡-Reasoning}\AgdaSpace{}%
\AgdaKeyword{using}\AgdaSpace{}%
\AgdaSymbol{(}\AgdaOperator{\AgdaFunction{begin\AgdaUnderscore{}}}\AgdaSymbol{;}\AgdaSpace{}%
\AgdaFunction{step-≡-∣}\AgdaSymbol{;}\AgdaSpace{}%
\AgdaFunction{step-≡-⟩}\AgdaSymbol{;}\AgdaSpace{}%
\AgdaOperator{\AgdaFunction{\AgdaUnderscore{}∎}}\AgdaSymbol{)}\<%
\\
\>[0]\AgdaKeyword{open}\AgdaSpace{}%
\AgdaModule{Eq.≡-Reasoning}\AgdaSpace{}%
\AgdaKeyword{renaming}\AgdaSpace{}%
\AgdaSymbol{(}\AgdaOperator{\AgdaFunction{begin\AgdaUnderscore{}}}\AgdaSpace{}%
\AgdaSymbol{to}\AgdaSpace{}%
\AgdaOperator{\AgdaFunction{beginEq\AgdaUnderscore{}}}\AgdaSymbol{;}\AgdaSpace{}%
\AgdaOperator{\AgdaFunction{\AgdaUnderscore{}∎}}\AgdaSpace{}%
\AgdaSymbol{to}\AgdaSpace{}%
\AgdaOperator{\AgdaFunction{\AgdaUnderscore{}∎eq}}\AgdaSymbol{)}\<%
\\
\>[0]\AgdaKeyword{open}\AgdaSpace{}%
\AgdaKeyword{import}\AgdaSpace{}%
\AgdaModule{Function.Bundles}\AgdaSpace{}%
\AgdaKeyword{using}\AgdaSpace{}%
\AgdaSymbol{(}\AgdaOperator{\AgdaFunction{\AgdaUnderscore{}⇔\AgdaUnderscore{}}}\AgdaSymbol{;}\AgdaSpace{}%
\AgdaFunction{mk⇔}\AgdaSymbol{)}\<%
\\
%
\\[\AgdaEmptyExtraSkip]%
\>[0]\AgdaComment{--\ GIRARD'S\ embedding}\<%
\\
%
\\[\AgdaEmptyExtraSkip]%
\>[0]\AgdaComment{--\ Preservation\ of\ typing\ by\ Girard's\ translation}\<%
\\
\>[0]\AgdaFunction{pres-typing}\AgdaSpace{}%
\AgdaSymbol{:}\AgdaSpace{}%
\AgdaSymbol{∀}\AgdaSpace{}%
\AgdaSymbol{\{}\AgdaBound{Γ}\AgdaSpace{}%
\AgdaBound{A}\AgdaSymbol{\}}\AgdaSpace{}%
\AgdaSymbol{→}\AgdaSpace{}%
\AgdaBound{Γ}\AgdaSpace{}%
\AgdaOperator{\AgdaDatatype{⊢}}\AgdaSpace{}%
\AgdaBound{A}\AgdaSpace{}%
\AgdaSymbol{→}\AgdaSpace{}%
\AgdaFunction{embedContext}\AgdaSpace{}%
\AgdaBound{Γ}\AgdaSpace{}%
\AgdaOperator{\AgdaDatatype{⊢b}}\AgdaSpace{}%
\AgdaFunction{embedType}\AgdaSpace{}%
\AgdaBound{A}\<%
\\
\>[0]\AgdaFunction{pres-typing}\AgdaSpace{}%
\AgdaBound{x}\AgdaSpace{}%
\AgdaSymbol{=}\AgdaSpace{}%
\AgdaFunction{embedTerm}\AgdaSpace{}%
\AgdaBound{x}\<%
\\
%
\\[\AgdaEmptyExtraSkip]%
\>[0]\AgdaFunction{pres-ext}\AgdaSpace{}%
\AgdaSymbol{:}\AgdaSpace{}%
\AgdaSymbol{∀}\AgdaSpace{}%
\AgdaSymbol{\{}\AgdaBound{Γ}\AgdaSpace{}%
\AgdaBound{Δ}\AgdaSpace{}%
\AgdaBound{A'}\AgdaSpace{}%
\AgdaBound{A}\AgdaSymbol{\}}\<%
\\
\>[0][@{}l@{\AgdaIndent{0}}]%
\>[2]\AgdaSymbol{→}\AgdaSpace{}%
\AgdaSymbol{(}\AgdaBound{f'}\AgdaSpace{}%
\AgdaSymbol{:}\AgdaSpace{}%
\AgdaSymbol{∀}\AgdaSpace{}%
\AgdaSymbol{\{}\AgdaBound{A}\AgdaSymbol{\}}\AgdaSpace{}%
\AgdaSymbol{→}\AgdaSpace{}%
\AgdaBound{A}\AgdaSpace{}%
\AgdaOperator{\AgdaDatatype{∈}}\AgdaSpace{}%
\AgdaBound{Γ}\AgdaSpace{}%
\AgdaSymbol{→}\AgdaSpace{}%
\AgdaBound{A}\AgdaSpace{}%
\AgdaOperator{\AgdaDatatype{∈}}\AgdaSpace{}%
\AgdaBound{Δ}\AgdaSymbol{)}\<%
\\
%
\>[2]\AgdaSymbol{→}\AgdaSpace{}%
\AgdaSymbol{(}\AgdaBound{fb'}\AgdaSpace{}%
\AgdaSymbol{:}\AgdaSpace{}%
\AgdaSymbol{∀}\AgdaSpace{}%
\AgdaSymbol{\{}\AgdaBound{Ab}\AgdaSymbol{\}}\AgdaSpace{}%
\AgdaSymbol{→}\AgdaSpace{}%
\AgdaBound{Ab}\AgdaSpace{}%
\AgdaOperator{\AgdaDatatype{∈b}}\AgdaSpace{}%
\AgdaFunction{embedContext}\AgdaSpace{}%
\AgdaBound{Γ}\AgdaSpace{}%
\AgdaSymbol{→}\AgdaSpace{}%
\AgdaBound{Ab}\AgdaSpace{}%
\AgdaOperator{\AgdaDatatype{∈b}}\AgdaSpace{}%
\AgdaFunction{embedContext}\AgdaSpace{}%
\AgdaBound{Δ}\AgdaSymbol{)}\<%
\\
%
\>[2]\AgdaSymbol{→}\AgdaSpace{}%
\AgdaSymbol{(\{}\AgdaBound{A'}\AgdaSpace{}%
\AgdaSymbol{:}\AgdaSpace{}%
\AgdaDatatype{Type}\AgdaSymbol{\}}\AgdaSpace{}%
\AgdaSymbol{(}\AgdaBound{x}\AgdaSpace{}%
\AgdaSymbol{:}\AgdaSpace{}%
\AgdaBound{A'}\AgdaSpace{}%
\AgdaOperator{\AgdaDatatype{∈}}\AgdaSpace{}%
\AgdaBound{Γ}\AgdaSymbol{)}\AgdaSpace{}%
\AgdaSymbol{→}\AgdaSpace{}%
\AgdaBound{fb'}\AgdaSpace{}%
\AgdaSymbol{(}\AgdaFunction{embedLookup}\AgdaSpace{}%
\AgdaBound{x}\AgdaSymbol{)}\AgdaSpace{}%
\AgdaOperator{\AgdaDatatype{≡}}\AgdaSpace{}%
\AgdaFunction{embedLookup}\AgdaSpace{}%
\AgdaSymbol{(}\AgdaBound{f'}\AgdaSpace{}%
\AgdaBound{x}\AgdaSymbol{))}\<%
\\
%
\>[2]\AgdaSymbol{→}\AgdaSpace{}%
\AgdaSymbol{(}\AgdaBound{x}\AgdaSpace{}%
\AgdaSymbol{:}\AgdaSpace{}%
\AgdaBound{A'}\AgdaSpace{}%
\AgdaOperator{\AgdaDatatype{∈}}\AgdaSpace{}%
\AgdaBound{Γ}\AgdaSpace{}%
\AgdaOperator{\AgdaInductiveConstructor{,}}\AgdaSpace{}%
\AgdaBound{A}\AgdaSymbol{)}\<%
\\
%
\>[2]\AgdaSymbol{→}\AgdaSpace{}%
\AgdaFunction{extb}\AgdaSpace{}%
\AgdaBound{fb'}\AgdaSpace{}%
\AgdaSymbol{(}\AgdaFunction{embedLookup}\AgdaSpace{}%
\AgdaBound{x}\AgdaSymbol{)}\AgdaSpace{}%
\AgdaOperator{\AgdaDatatype{≡}}\AgdaSpace{}%
\AgdaFunction{embedLookup}\AgdaSpace{}%
\AgdaSymbol{(}\AgdaFunction{ext}\AgdaSpace{}%
\AgdaBound{f'}\AgdaSpace{}%
\AgdaBound{x}\AgdaSymbol{)}\<%
\\
\>[0]\AgdaFunction{pres-ext}\AgdaSpace{}%
\AgdaBound{f'}\AgdaSpace{}%
\AgdaBound{fb'}\AgdaSpace{}%
\AgdaBound{h}\AgdaSpace{}%
\AgdaInductiveConstructor{Z}\AgdaSpace{}%
\AgdaSymbol{=}\AgdaSpace{}%
\AgdaInductiveConstructor{refl}\<%
\\
\>[0]\AgdaFunction{pres-ext}\AgdaSpace{}%
\AgdaBound{f'}\AgdaSpace{}%
\AgdaBound{fb'}\AgdaSpace{}%
\AgdaBound{h}\AgdaSpace{}%
\AgdaSymbol{(}\AgdaInductiveConstructor{S}\AgdaSpace{}%
\AgdaBound{x}\AgdaSymbol{)}\AgdaSpace{}%
\AgdaSymbol{=}\AgdaSpace{}%
\AgdaFunction{cong}\AgdaSpace{}%
\AgdaInductiveConstructor{Sb}\AgdaSpace{}%
\AgdaSymbol{(}\AgdaBound{h}\AgdaSpace{}%
\AgdaBound{x}\AgdaSymbol{)}\<%
\\
%
\\[\AgdaEmptyExtraSkip]%
\>[0]\AgdaFunction{pres-rename}\AgdaSpace{}%
\AgdaSymbol{:}\AgdaSpace{}%
\AgdaSymbol{∀}\AgdaSpace{}%
\AgdaSymbol{\{}\AgdaBound{Γ}\AgdaSpace{}%
\AgdaBound{A}\AgdaSpace{}%
\AgdaBound{Δ}\AgdaSymbol{\}}\<%
\\
\>[0][@{}l@{\AgdaIndent{0}}]%
\>[2]\AgdaSymbol{→}\AgdaSpace{}%
\AgdaSymbol{(}\AgdaBound{f'}\AgdaSpace{}%
\AgdaSymbol{:}\AgdaSpace{}%
\AgdaSymbol{∀}\AgdaSpace{}%
\AgdaSymbol{\{}\AgdaBound{A}\AgdaSymbol{\}}\AgdaSpace{}%
\AgdaSymbol{→}\AgdaSpace{}%
\AgdaBound{A}\AgdaSpace{}%
\AgdaOperator{\AgdaDatatype{∈}}\AgdaSpace{}%
\AgdaBound{Γ}\AgdaSpace{}%
\AgdaSymbol{→}\AgdaSpace{}%
\AgdaBound{A}\AgdaSpace{}%
\AgdaOperator{\AgdaDatatype{∈}}\AgdaSpace{}%
\AgdaBound{Δ}\AgdaSymbol{)}\<%
\\
%
\>[2]\AgdaSymbol{→}\AgdaSpace{}%
\AgdaSymbol{(}\AgdaBound{fb'}\AgdaSpace{}%
\AgdaSymbol{:}\AgdaSpace{}%
\AgdaSymbol{∀}\AgdaSpace{}%
\AgdaSymbol{\{}\AgdaBound{Ab}\AgdaSymbol{\}}\AgdaSpace{}%
\AgdaSymbol{→}\AgdaSpace{}%
\AgdaBound{Ab}\AgdaSpace{}%
\AgdaOperator{\AgdaDatatype{∈b}}\AgdaSpace{}%
\AgdaFunction{embedContext}\AgdaSpace{}%
\AgdaBound{Γ}\AgdaSpace{}%
\AgdaSymbol{→}\AgdaSpace{}%
\AgdaBound{Ab}\AgdaSpace{}%
\AgdaOperator{\AgdaDatatype{∈b}}\AgdaSpace{}%
\AgdaFunction{embedContext}\AgdaSpace{}%
\AgdaBound{Δ}\AgdaSymbol{)}\<%
\\
%
\>[2]\AgdaSymbol{→}\AgdaSpace{}%
\AgdaSymbol{(\{}\AgdaBound{A'}\AgdaSpace{}%
\AgdaSymbol{:}\AgdaSpace{}%
\AgdaDatatype{Type}\AgdaSymbol{\}}\AgdaSpace{}%
\AgdaSymbol{(}\AgdaBound{x}\AgdaSpace{}%
\AgdaSymbol{:}\AgdaSpace{}%
\AgdaBound{A'}\AgdaSpace{}%
\AgdaOperator{\AgdaDatatype{∈}}\AgdaSpace{}%
\AgdaBound{Γ}\AgdaSymbol{)}\AgdaSpace{}%
\AgdaSymbol{→}\AgdaSpace{}%
\AgdaBound{fb'}\AgdaSpace{}%
\AgdaSymbol{(}\AgdaFunction{embedLookup}\AgdaSpace{}%
\AgdaBound{x}\AgdaSymbol{)}\AgdaSpace{}%
\AgdaOperator{\AgdaDatatype{≡}}\AgdaSpace{}%
\AgdaFunction{embedLookup}\AgdaSpace{}%
\AgdaSymbol{(}\AgdaBound{f'}\AgdaSpace{}%
\AgdaBound{x}\AgdaSymbol{))}\<%
\\
%
\>[2]\AgdaSymbol{→}\AgdaSpace{}%
\AgdaSymbol{(}\AgdaBound{t}\AgdaSpace{}%
\AgdaSymbol{:}\AgdaSpace{}%
\AgdaBound{Γ}\AgdaSpace{}%
\AgdaOperator{\AgdaDatatype{⊢}}\AgdaSpace{}%
\AgdaBound{A}\AgdaSymbol{)}\<%
\\
%
\>[2]\AgdaSymbol{→}\AgdaSpace{}%
\AgdaFunction{renameb}\AgdaSpace{}%
\AgdaBound{fb'}\AgdaSpace{}%
\AgdaSymbol{(}\AgdaFunction{embedTerm}\AgdaSpace{}%
\AgdaBound{t}\AgdaSymbol{)}\AgdaSpace{}%
\AgdaOperator{\AgdaDatatype{≡}}\AgdaSpace{}%
\AgdaFunction{embedTerm}\AgdaSpace{}%
\AgdaSymbol{(}\AgdaFunction{rename}\AgdaSpace{}%
\AgdaBound{f'}\AgdaSpace{}%
\AgdaBound{t}\AgdaSymbol{)}\<%
\\
\>[0]\AgdaFunction{pres-rename}\AgdaSpace{}%
\AgdaBound{f'}\AgdaSpace{}%
\AgdaBound{fb'}\AgdaSpace{}%
\AgdaBound{h}\AgdaSpace{}%
\AgdaSymbol{(}\AgdaInductiveConstructor{`}\AgdaSpace{}%
\AgdaBound{x}\AgdaSymbol{)}\AgdaSpace{}%
\AgdaSymbol{=}\<%
\\
\>[0][@{}l@{\AgdaIndent{0}}]%
\>[2]\AgdaFunction{cong}\AgdaSpace{}%
\AgdaInductiveConstructor{ε}\AgdaSpace{}%
\AgdaSymbol{(}\AgdaBound{h}\AgdaSpace{}%
\AgdaBound{x}\AgdaSymbol{)}\<%
\\
\>[0]\AgdaFunction{pres-rename}\AgdaSpace{}%
\AgdaBound{f'}\AgdaSpace{}%
\AgdaBound{fb'}\AgdaSpace{}%
\AgdaBound{h}\AgdaSpace{}%
\AgdaSymbol{(}\AgdaInductiveConstructor{ƛ}\AgdaSpace{}%
\AgdaBound{t}\AgdaSymbol{)}\AgdaSpace{}%
\AgdaSymbol{=}\<%
\\
\>[0][@{}l@{\AgdaIndent{0}}]%
\>[2]\AgdaFunction{cong}\AgdaSpace{}%
\AgdaInductiveConstructor{ƛb}\AgdaSpace{}%
\AgdaSymbol{(}\AgdaFunction{pres-rename}\AgdaSpace{}%
\AgdaSymbol{(}\AgdaFunction{calc.ext}\AgdaSpace{}%
\AgdaBound{f'}\AgdaSymbol{)}\AgdaSpace{}%
\AgdaSymbol{(}\AgdaFunction{extb}\AgdaSpace{}%
\AgdaBound{fb'}\AgdaSymbol{)}\AgdaSpace{}%
\AgdaSymbol{(}\AgdaFunction{pres-ext}\AgdaSpace{}%
\AgdaBound{f'}\AgdaSpace{}%
\AgdaBound{fb'}\AgdaSpace{}%
\AgdaBound{h}\AgdaSymbol{)}\AgdaSpace{}%
\AgdaBound{t}\AgdaSymbol{)}\<%
\\
\>[0]\AgdaFunction{pres-rename}\AgdaSpace{}%
\AgdaBound{f'}\AgdaSpace{}%
\AgdaBound{fb'}\AgdaSpace{}%
\AgdaBound{h}\AgdaSpace{}%
\AgdaSymbol{(}\AgdaBound{t}\AgdaSpace{}%
\AgdaOperator{\AgdaInductiveConstructor{·}}\AgdaSpace{}%
\AgdaBound{t₁}\AgdaSymbol{)}\AgdaSpace{}%
\AgdaSymbol{=}\<%
\\
\>[0][@{}l@{\AgdaIndent{0}}]%
\>[2]\AgdaFunction{cong₂}\AgdaSpace{}%
\AgdaOperator{\AgdaInductiveConstructor{\AgdaUnderscore{}·b\AgdaUnderscore{}}}\AgdaSpace{}%
\AgdaSymbol{(}\AgdaFunction{pres-rename}\AgdaSpace{}%
\AgdaBound{f'}\AgdaSpace{}%
\AgdaBound{fb'}\AgdaSpace{}%
\AgdaBound{h}\AgdaSpace{}%
\AgdaBound{t}\AgdaSymbol{)}\AgdaSpace{}%
\AgdaSymbol{(}\AgdaFunction{cong}\AgdaSpace{}%
\AgdaInductiveConstructor{box}\AgdaSpace{}%
\AgdaSymbol{(}\AgdaFunction{pres-rename}\AgdaSpace{}%
\AgdaBound{f'}\AgdaSpace{}%
\AgdaBound{fb'}\AgdaSpace{}%
\AgdaBound{h}\AgdaSpace{}%
\AgdaBound{t₁}\AgdaSymbol{))}\<%
\\
%
\\[\AgdaEmptyExtraSkip]%
%
\\[\AgdaEmptyExtraSkip]%
\>[0]\AgdaFunction{pres-exts}\AgdaSpace{}%
\AgdaSymbol{:}\AgdaSpace{}%
\AgdaSymbol{∀}\AgdaSpace{}%
\AgdaSymbol{\{}\AgdaBound{Γ}\AgdaSpace{}%
\AgdaBound{Δ}\AgdaSymbol{\}}\<%
\\
\>[0][@{}l@{\AgdaIndent{0}}]%
\>[2]\AgdaSymbol{→}\AgdaSpace{}%
\AgdaSymbol{(}\AgdaBound{fb'}\AgdaSpace{}%
\AgdaSymbol{:}\AgdaSpace{}%
\AgdaSymbol{∀}\AgdaSpace{}%
\AgdaSymbol{\{}\AgdaBound{Ab}\AgdaSymbol{\}}\AgdaSpace{}%
\AgdaSymbol{→}\AgdaSpace{}%
\AgdaInductiveConstructor{□}\AgdaSpace{}%
\AgdaBound{Ab}\AgdaSpace{}%
\AgdaOperator{\AgdaDatatype{∈b}}\AgdaSpace{}%
\AgdaFunction{embedContext}\AgdaSpace{}%
\AgdaBound{Γ}\AgdaSpace{}%
\AgdaSymbol{→}\AgdaSpace{}%
\AgdaFunction{embedContext}\AgdaSpace{}%
\AgdaBound{Δ}\AgdaSpace{}%
\AgdaOperator{\AgdaDatatype{⊢b}}\AgdaSpace{}%
\AgdaBound{Ab}\AgdaSymbol{)}\<%
\\
%
\>[2]\AgdaSymbol{→}\AgdaSpace{}%
\AgdaSymbol{(}\AgdaBound{f'}\AgdaSpace{}%
\AgdaSymbol{:}\AgdaSpace{}%
\AgdaSymbol{∀}\AgdaSpace{}%
\AgdaSymbol{\{}\AgdaBound{A}\AgdaSymbol{\}}\AgdaSpace{}%
\AgdaSymbol{→}\AgdaSpace{}%
\AgdaBound{A}\AgdaSpace{}%
\AgdaOperator{\AgdaDatatype{∈}}\AgdaSpace{}%
\AgdaBound{Γ}\AgdaSpace{}%
\AgdaSymbol{→}\AgdaSpace{}%
\AgdaBound{Δ}\AgdaSpace{}%
\AgdaOperator{\AgdaDatatype{⊢}}\AgdaSpace{}%
\AgdaBound{A}\AgdaSymbol{)}\<%
\\
%
\>[2]\AgdaSymbol{→}\AgdaSpace{}%
\AgdaSymbol{(\{}\AgdaBound{A'}\AgdaSpace{}%
\AgdaSymbol{:}\AgdaSpace{}%
\AgdaDatatype{Type}\AgdaSymbol{\}}\AgdaSpace{}%
\AgdaSymbol{(}\AgdaBound{x}\AgdaSpace{}%
\AgdaSymbol{:}\AgdaSpace{}%
\AgdaBound{A'}\AgdaSpace{}%
\AgdaOperator{\AgdaDatatype{∈}}\AgdaSpace{}%
\AgdaBound{Γ}\AgdaSymbol{)}\AgdaSpace{}%
\AgdaSymbol{→}\AgdaSpace{}%
\AgdaBound{fb'}\AgdaSpace{}%
\AgdaSymbol{(}\AgdaFunction{embedLookup}\AgdaSpace{}%
\AgdaBound{x}\AgdaSymbol{)}\AgdaSpace{}%
\AgdaOperator{\AgdaDatatype{≡}}\AgdaSpace{}%
\AgdaFunction{embedTerm}\AgdaSpace{}%
\AgdaSymbol{(}\AgdaBound{f'}\AgdaSpace{}%
\AgdaBound{x}\AgdaSymbol{))}\<%
\\
%
\>[2]\AgdaSymbol{→}\AgdaSpace{}%
\AgdaSymbol{\{}\AgdaBound{A'}\AgdaSpace{}%
\AgdaBound{A}\AgdaSpace{}%
\AgdaSymbol{:}\AgdaSpace{}%
\AgdaDatatype{Type}\AgdaSymbol{\}}\AgdaSpace{}%
\AgdaSymbol{(}\AgdaBound{x}\AgdaSpace{}%
\AgdaSymbol{:}\AgdaSpace{}%
\AgdaBound{A'}\AgdaSpace{}%
\AgdaOperator{\AgdaDatatype{∈}}\AgdaSpace{}%
\AgdaBound{Γ}\AgdaSpace{}%
\AgdaOperator{\AgdaInductiveConstructor{,}}\AgdaSpace{}%
\AgdaBound{A}\AgdaSymbol{)}\AgdaSpace{}%
\AgdaSymbol{→}\AgdaSpace{}%
\AgdaFunction{extsb}\AgdaSpace{}%
\AgdaBound{fb'}\AgdaSpace{}%
\AgdaSymbol{(}\AgdaFunction{embedLookup}\AgdaSpace{}%
\AgdaBound{x}\AgdaSymbol{)}\AgdaSpace{}%
\AgdaOperator{\AgdaDatatype{≡}}\AgdaSpace{}%
\AgdaFunction{embedTerm}\AgdaSpace{}%
\AgdaSymbol{(}\AgdaFunction{exts}\AgdaSpace{}%
\AgdaBound{f'}\AgdaSpace{}%
\AgdaBound{x}\AgdaSymbol{)}\<%
\\
\>[0]\AgdaFunction{pres-exts}\AgdaSpace{}%
\AgdaBound{fb'}\AgdaSpace{}%
\AgdaBound{f'}\AgdaSpace{}%
\AgdaBound{h}\AgdaSpace{}%
\AgdaInductiveConstructor{Z}\AgdaSpace{}%
\AgdaSymbol{=}\AgdaSpace{}%
\AgdaInductiveConstructor{refl}\<%
\\
\>[0]\AgdaFunction{pres-exts}\AgdaSpace{}%
\AgdaBound{fb'}\AgdaSpace{}%
\AgdaBound{f'}\AgdaSpace{}%
\AgdaBound{h}\AgdaSpace{}%
\AgdaSymbol{(}\AgdaInductiveConstructor{S}\AgdaSpace{}%
\AgdaBound{x}\AgdaSymbol{)}\AgdaSpace{}%
\AgdaSymbol{=}\<%
\\
\>[0][@{}l@{\AgdaIndent{0}}]%
\>[2]\AgdaOperator{\AgdaFunction{beginEq}}\<%
\\
\>[2][@{}l@{\AgdaIndent{0}}]%
\>[4]\AgdaFunction{renameb}\AgdaSpace{}%
\AgdaInductiveConstructor{Sb}\AgdaSpace{}%
\AgdaSymbol{(}\AgdaBound{fb'}\AgdaSpace{}%
\AgdaSymbol{(}\AgdaFunction{embedLookup}\AgdaSpace{}%
\AgdaBound{x}\AgdaSymbol{))}\<%
\\
%
\>[2]\AgdaFunction{≡⟨}\AgdaSpace{}%
\AgdaFunction{cong}\AgdaSpace{}%
\AgdaSymbol{(}\AgdaFunction{renameb}\AgdaSpace{}%
\AgdaInductiveConstructor{Sb}\AgdaSymbol{)}\AgdaSpace{}%
\AgdaSymbol{(}\AgdaBound{h}\AgdaSpace{}%
\AgdaBound{x}\AgdaSymbol{)}\AgdaSpace{}%
\AgdaFunction{⟩}\<%
\\
\>[2][@{}l@{\AgdaIndent{0}}]%
\>[4]\AgdaFunction{renameb}\AgdaSpace{}%
\AgdaInductiveConstructor{Sb}\AgdaSpace{}%
\AgdaSymbol{(}\AgdaFunction{embedTerm}\AgdaSpace{}%
\AgdaSymbol{(}\AgdaBound{f'}\AgdaSpace{}%
\AgdaBound{x}\AgdaSymbol{))}\<%
\\
%
\>[2]\AgdaFunction{≡⟨}\AgdaSpace{}%
\AgdaFunction{pres-rename}\AgdaSpace{}%
\AgdaInductiveConstructor{S}\AgdaSpace{}%
\AgdaInductiveConstructor{Sb}\AgdaSpace{}%
\AgdaSymbol{(λ}\AgdaSpace{}%
\AgdaBound{x₁}\AgdaSpace{}%
\AgdaSymbol{→}\AgdaSpace{}%
\AgdaInductiveConstructor{refl}\AgdaSymbol{)}\AgdaSpace{}%
\AgdaSymbol{(}\AgdaBound{f'}\AgdaSpace{}%
\AgdaBound{x}\AgdaSymbol{)}\AgdaSpace{}%
\AgdaFunction{⟩}\<%
\\
\>[2][@{}l@{\AgdaIndent{0}}]%
\>[4]\AgdaFunction{embedTerm}\AgdaSpace{}%
\AgdaSymbol{(}\AgdaFunction{rename}\AgdaSpace{}%
\AgdaInductiveConstructor{S}\AgdaSpace{}%
\AgdaSymbol{(}\AgdaBound{f'}\AgdaSpace{}%
\AgdaBound{x}\AgdaSymbol{))}\<%
\\
%
\>[2]\AgdaOperator{\AgdaFunction{∎eq}}\<%
\\
%
\\[\AgdaEmptyExtraSkip]%
\>[0]\AgdaFunction{pres-simultaneous-subst}\AgdaSpace{}%
\AgdaSymbol{:}\AgdaSpace{}%
\AgdaSymbol{∀}\AgdaSpace{}%
\AgdaSymbol{\{}\AgdaBound{Γ}\AgdaSpace{}%
\AgdaBound{Δ}\AgdaSpace{}%
\AgdaBound{A}\AgdaSymbol{\}}\<%
\\
\>[0][@{}l@{\AgdaIndent{0}}]%
\>[2]\AgdaSymbol{→}\AgdaSpace{}%
\AgdaSymbol{(}\AgdaBound{fb'}\AgdaSpace{}%
\AgdaSymbol{:}\AgdaSpace{}%
\AgdaSymbol{∀}\AgdaSpace{}%
\AgdaSymbol{\{}\AgdaBound{Ab}\AgdaSymbol{\}}\AgdaSpace{}%
\AgdaSymbol{→}\AgdaSpace{}%
\AgdaInductiveConstructor{□}\AgdaSpace{}%
\AgdaBound{Ab}\AgdaSpace{}%
\AgdaOperator{\AgdaDatatype{∈b}}\AgdaSpace{}%
\AgdaFunction{embedContext}\AgdaSpace{}%
\AgdaBound{Γ}\AgdaSpace{}%
\AgdaSymbol{→}\AgdaSpace{}%
\AgdaFunction{embedContext}\AgdaSpace{}%
\AgdaBound{Δ}\AgdaSpace{}%
\AgdaOperator{\AgdaDatatype{⊢b}}\AgdaSpace{}%
\AgdaBound{Ab}\AgdaSymbol{)}\<%
\\
%
\>[2]\AgdaSymbol{→}\AgdaSpace{}%
\AgdaSymbol{(}\AgdaBound{f'}\AgdaSpace{}%
\AgdaSymbol{:}\AgdaSpace{}%
\AgdaSymbol{∀}\AgdaSpace{}%
\AgdaSymbol{\{}\AgdaBound{A}\AgdaSymbol{\}}\AgdaSpace{}%
\AgdaSymbol{→}\AgdaSpace{}%
\AgdaBound{A}\AgdaSpace{}%
\AgdaOperator{\AgdaDatatype{∈}}\AgdaSpace{}%
\AgdaBound{Γ}\AgdaSpace{}%
\AgdaSymbol{→}\AgdaSpace{}%
\AgdaBound{Δ}\AgdaSpace{}%
\AgdaOperator{\AgdaDatatype{⊢}}\AgdaSpace{}%
\AgdaBound{A}\AgdaSymbol{)}\<%
\\
%
\>[2]\AgdaSymbol{→}\AgdaSpace{}%
\AgdaSymbol{(}\AgdaBound{h}\AgdaSpace{}%
\AgdaSymbol{:}\AgdaSpace{}%
\AgdaSymbol{∀}\AgdaSpace{}%
\AgdaSymbol{\{}\AgdaBound{A'}\AgdaSymbol{\}}\AgdaSpace{}%
\AgdaSymbol{→}\AgdaSpace{}%
\AgdaSymbol{(}\AgdaBound{x}\AgdaSpace{}%
\AgdaSymbol{:}\AgdaSpace{}%
\AgdaBound{A'}\AgdaSpace{}%
\AgdaOperator{\AgdaDatatype{∈}}\AgdaSpace{}%
\AgdaBound{Γ}\AgdaSymbol{)}\AgdaSpace{}%
\AgdaSymbol{→}\AgdaSpace{}%
\AgdaBound{fb'}\AgdaSpace{}%
\AgdaSymbol{(}\AgdaFunction{embedLookup}\AgdaSpace{}%
\AgdaBound{x}\AgdaSymbol{)}\AgdaSpace{}%
\AgdaOperator{\AgdaDatatype{≡}}\AgdaSpace{}%
\AgdaFunction{embedTerm}\AgdaSpace{}%
\AgdaSymbol{(}\AgdaBound{f'}\AgdaSpace{}%
\AgdaBound{x}\AgdaSymbol{)}\AgdaSpace{}%
\AgdaSymbol{)}\<%
\\
%
\>[2]\AgdaSymbol{→}\AgdaSpace{}%
\AgdaSymbol{(}\AgdaBound{M}\AgdaSpace{}%
\AgdaSymbol{:}\AgdaSpace{}%
\AgdaBound{Γ}\AgdaSpace{}%
\AgdaOperator{\AgdaDatatype{⊢}}\AgdaSpace{}%
\AgdaBound{A}\AgdaSymbol{)}\<%
\\
%
\>[2]\AgdaSymbol{→}\AgdaSpace{}%
\AgdaFunction{substb}\AgdaSpace{}%
\AgdaBound{fb'}\AgdaSpace{}%
\AgdaSymbol{(}\AgdaFunction{embedTerm}\AgdaSpace{}%
\AgdaBound{M}\AgdaSymbol{)}\AgdaSpace{}%
\AgdaOperator{\AgdaDatatype{≡}}\AgdaSpace{}%
\AgdaFunction{embedTerm}\AgdaSpace{}%
\AgdaSymbol{(}\AgdaFunction{subst}\AgdaSpace{}%
\AgdaBound{f'}\AgdaSpace{}%
\AgdaBound{M}\AgdaSymbol{)}\<%
\\
\>[0]\AgdaFunction{pres-simultaneous-subst}\AgdaSpace{}%
\AgdaBound{fb'}\AgdaSpace{}%
\AgdaBound{f'}\AgdaSpace{}%
\AgdaBound{h}\AgdaSpace{}%
\AgdaSymbol{(}\AgdaInductiveConstructor{`}\AgdaSpace{}%
\AgdaBound{x}\AgdaSymbol{)}\AgdaSpace{}%
\AgdaSymbol{=}\<%
\\
\>[0][@{}l@{\AgdaIndent{0}}]%
\>[2]\AgdaBound{h}\AgdaSpace{}%
\AgdaBound{x}\<%
\\
\>[0]\AgdaFunction{pres-simultaneous-subst}\AgdaSpace{}%
\AgdaBound{fb'}\AgdaSpace{}%
\AgdaBound{f'}\AgdaSpace{}%
\AgdaBound{h}\AgdaSpace{}%
\AgdaSymbol{(}\AgdaInductiveConstructor{ƛ}\AgdaSpace{}%
\AgdaBound{M}\AgdaSymbol{)}\AgdaSpace{}%
\AgdaSymbol{=}\<%
\\
\>[0][@{}l@{\AgdaIndent{0}}]%
\>[2]\AgdaFunction{cong}\AgdaSpace{}%
\AgdaInductiveConstructor{ƛb}\AgdaSpace{}%
\AgdaSymbol{(}\AgdaSpace{}%
\AgdaFunction{pres-simultaneous-subst}\AgdaSpace{}%
\AgdaSymbol{(}\AgdaFunction{extsb}\AgdaSpace{}%
\AgdaBound{fb'}\AgdaSymbol{)}\AgdaSpace{}%
\AgdaSymbol{(}\AgdaFunction{exts}\AgdaSpace{}%
\AgdaBound{f'}\AgdaSymbol{)}\AgdaSpace{}%
\AgdaSymbol{(}\AgdaFunction{pres-exts}\AgdaSpace{}%
\AgdaBound{fb'}\AgdaSpace{}%
\AgdaBound{f'}\AgdaSpace{}%
\AgdaBound{h}\AgdaSymbol{)}\AgdaSpace{}%
\AgdaBound{M}\AgdaSpace{}%
\AgdaSymbol{)}\<%
\\
\>[0]\AgdaFunction{pres-simultaneous-subst}\AgdaSpace{}%
\AgdaBound{fb'}\AgdaSpace{}%
\AgdaBound{f'}\AgdaSpace{}%
\AgdaBound{h}\AgdaSpace{}%
\AgdaSymbol{(}\AgdaBound{M}\AgdaSpace{}%
\AgdaOperator{\AgdaInductiveConstructor{·}}\AgdaSpace{}%
\AgdaBound{M₁}\AgdaSymbol{)}\AgdaSpace{}%
\AgdaSymbol{=}\<%
\\
\>[0][@{}l@{\AgdaIndent{0}}]%
\>[2]\AgdaFunction{cong₂}\AgdaSpace{}%
\AgdaOperator{\AgdaInductiveConstructor{\AgdaUnderscore{}·b\AgdaUnderscore{}}}\AgdaSpace{}%
\AgdaSymbol{(}\AgdaFunction{pres-simultaneous-subst}\AgdaSpace{}%
\AgdaBound{fb'}\AgdaSpace{}%
\AgdaBound{f'}\AgdaSpace{}%
\AgdaBound{h}\AgdaSpace{}%
\AgdaBound{M}\AgdaSymbol{)}\AgdaSpace{}%
\AgdaSymbol{(}\AgdaFunction{cong}\AgdaSpace{}%
\AgdaInductiveConstructor{box}\AgdaSpace{}%
\AgdaSymbol{(}\AgdaFunction{pres-simultaneous-subst}\AgdaSpace{}%
\AgdaBound{fb'}\AgdaSpace{}%
\AgdaBound{f'}\AgdaSpace{}%
\AgdaBound{h}\AgdaSpace{}%
\AgdaBound{M₁}\AgdaSymbol{))}\<%
\\
%
\\[\AgdaEmptyExtraSkip]%
%
\\[\AgdaEmptyExtraSkip]%
\>[0]\AgdaComment{--\ Proof\ of\ h}\<%
\\
\>[0]\AgdaFunction{pres-subst-func}\AgdaSpace{}%
\AgdaSymbol{:}\AgdaSpace{}%
\AgdaSymbol{∀}\AgdaSpace{}%
\AgdaSymbol{\{}\AgdaBound{A}\AgdaSpace{}%
\AgdaBound{C}\AgdaSpace{}%
\AgdaBound{Γ}\AgdaSymbol{\}}\<%
\\
\>[0][@{}l@{\AgdaIndent{0}}]%
\>[2]\AgdaSymbol{→}\AgdaSpace{}%
\AgdaSymbol{(}\AgdaBound{N}\AgdaSpace{}%
\AgdaSymbol{:}\AgdaSpace{}%
\AgdaBound{Γ}\AgdaSpace{}%
\AgdaOperator{\AgdaDatatype{⊢}}\AgdaSpace{}%
\AgdaBound{C}\AgdaSymbol{)}\<%
\\
%
\>[2]\AgdaSymbol{→}\AgdaSpace{}%
\AgdaSymbol{(}\AgdaBound{x}\AgdaSpace{}%
\AgdaSymbol{:}\AgdaSpace{}%
\AgdaBound{A}\AgdaSpace{}%
\AgdaOperator{\AgdaDatatype{∈}}\AgdaSpace{}%
\AgdaBound{Γ}\AgdaSpace{}%
\AgdaOperator{\AgdaInductiveConstructor{,}}\AgdaSpace{}%
\AgdaBound{C}\AgdaSymbol{)}\<%
\\
%
\>[2]\AgdaSymbol{→}\AgdaSpace{}%
\AgdaFunction{fb}\AgdaSpace{}%
\AgdaSymbol{(}\AgdaFunction{embedTerm}\AgdaSpace{}%
\AgdaBound{N}\AgdaSymbol{)}\AgdaSpace{}%
\AgdaSymbol{(}\AgdaFunction{embedLookup}\AgdaSpace{}%
\AgdaBound{x}\AgdaSymbol{)}\AgdaSpace{}%
\AgdaOperator{\AgdaDatatype{≡}}\AgdaSpace{}%
\AgdaFunction{embedTerm}\AgdaSpace{}%
\AgdaSymbol{((}\AgdaFunction{f}\AgdaSpace{}%
\AgdaBound{N}\AgdaSymbol{)}\AgdaSpace{}%
\AgdaBound{x}\AgdaSymbol{)}\<%
\\
\>[0]\AgdaFunction{pres-subst-func}\AgdaSpace{}%
\AgdaBound{N}\AgdaSpace{}%
\AgdaInductiveConstructor{Z}\AgdaSpace{}%
\AgdaSymbol{=}\AgdaSpace{}%
\AgdaInductiveConstructor{refl}\<%
\\
\>[0]\AgdaFunction{pres-subst-func}\AgdaSpace{}%
\AgdaBound{N}\AgdaSpace{}%
\AgdaSymbol{(}\AgdaInductiveConstructor{S}\AgdaSpace{}%
\AgdaBound{x}\AgdaSymbol{)}\AgdaSpace{}%
\AgdaSymbol{=}\AgdaSpace{}%
\AgdaInductiveConstructor{refl}\<%
\end{code}
\begin{code}%
\>[0]\AgdaComment{--\ Preservation\ and\ Reflection\ of\ }\<%
\\
\>[0]\AgdaComment{--\ single\ substitution\ of\ Girard's\ translation}\<%
\\
\>[0]\AgdaFunction{pres-subst}\AgdaSpace{}%
\AgdaSymbol{:}\AgdaSpace{}%
\AgdaSymbol{∀}\AgdaSpace{}%
\AgdaSymbol{\{}\AgdaBound{Γ}\AgdaSpace{}%
\AgdaBound{A}\AgdaSpace{}%
\AgdaBound{B}\AgdaSymbol{\}}\<%
\\
\>[0][@{}l@{\AgdaIndent{0}}]%
\>[2]\AgdaSymbol{→}\AgdaSpace{}%
\AgdaSymbol{(}\AgdaBound{M}\AgdaSpace{}%
\AgdaSymbol{:}\AgdaSpace{}%
\AgdaBound{Γ}\AgdaSpace{}%
\AgdaOperator{\AgdaInductiveConstructor{,}}\AgdaSpace{}%
\AgdaBound{B}\AgdaSpace{}%
\AgdaOperator{\AgdaDatatype{⊢}}\AgdaSpace{}%
\AgdaBound{A}\AgdaSymbol{)}\<%
\\
%
\>[2]\AgdaSymbol{→}\AgdaSpace{}%
\AgdaSymbol{(}\AgdaBound{N}\AgdaSpace{}%
\AgdaSymbol{:}\AgdaSpace{}%
\AgdaBound{Γ}\AgdaSpace{}%
\AgdaOperator{\AgdaDatatype{⊢}}\AgdaSpace{}%
\AgdaBound{B}\AgdaSymbol{)}\<%
\\
%
\>[2]\AgdaSymbol{→}\AgdaSpace{}%
\AgdaSymbol{(}\AgdaFunction{embedTerm}\AgdaSpace{}%
\AgdaBound{M}\AgdaSymbol{)}\AgdaSpace{}%
\AgdaOperator{\AgdaFunction{[}}\AgdaSpace{}%
\AgdaFunction{embedTerm}\AgdaSpace{}%
\AgdaBound{N}\AgdaSpace{}%
\AgdaOperator{\AgdaFunction{]b}}\AgdaSpace{}%
\AgdaOperator{\AgdaDatatype{≡}}\AgdaSpace{}%
\AgdaFunction{embedTerm}\AgdaSpace{}%
\AgdaSymbol{(}\AgdaBound{M}\AgdaSpace{}%
\AgdaOperator{\AgdaFunction{[}}\AgdaSpace{}%
\AgdaBound{N}\AgdaSpace{}%
\AgdaOperator{\AgdaFunction{]}}\AgdaSymbol{)}\<%
\\
\>[0]\AgdaFunction{pres-subst}\AgdaSpace{}%
\AgdaBound{M}\AgdaSpace{}%
\AgdaBound{N}\AgdaSpace{}%
\AgdaSymbol{=}\<%
\\
\>[0][@{}l@{\AgdaIndent{0}}]%
\>[2]\AgdaFunction{pres-simultaneous-subst}\<%
\\
\>[2][@{}l@{\AgdaIndent{0}}]%
\>[4]\AgdaSymbol{(}\AgdaFunction{fb}\AgdaSpace{}%
\AgdaSymbol{(}\AgdaFunction{embedTerm}\AgdaSpace{}%
\AgdaBound{N}\AgdaSymbol{))}\<%
\\
%
\>[4]\AgdaSymbol{(}\AgdaFunction{f}\AgdaSpace{}%
\AgdaBound{N}\AgdaSymbol{)}\<%
\\
%
\>[4]\AgdaSymbol{(}\AgdaFunction{pres-subst-func}\AgdaSpace{}%
\AgdaBound{N}\AgdaSymbol{)}\AgdaSpace{}%
\AgdaBound{M}\<%
\end{code}
\begin{code}[hide]%
\>[0]\<%
\\
\>[0]\AgdaComment{--\ Preservation\ of\ reduction\ of\ Girard's\ translation}\<%
\\
%
\\[\AgdaEmptyExtraSkip]%
\>[0]\AgdaFunction{pres-red}\AgdaSpace{}%
\AgdaSymbol{:}\AgdaSpace{}%
\AgdaSymbol{∀}\AgdaSpace{}%
\AgdaSymbol{\{}\AgdaBound{Γ}\AgdaSpace{}%
\AgdaBound{A}\AgdaSymbol{\}}\<%
\\
\>[0][@{}l@{\AgdaIndent{0}}]%
\>[2]\AgdaSymbol{→}\AgdaSpace{}%
\AgdaSymbol{(}\AgdaBound{M}\AgdaSpace{}%
\AgdaSymbol{:}\AgdaSpace{}%
\AgdaBound{Γ}\AgdaSpace{}%
\AgdaOperator{\AgdaDatatype{⊢}}\AgdaSpace{}%
\AgdaBound{A}\AgdaSymbol{)}\<%
\\
%
\>[2]\AgdaSymbol{→}\AgdaSpace{}%
\AgdaSymbol{(}\AgdaBound{N}\AgdaSpace{}%
\AgdaSymbol{:}\AgdaSpace{}%
\AgdaBound{Γ}\AgdaSpace{}%
\AgdaOperator{\AgdaDatatype{⊢}}\AgdaSpace{}%
\AgdaBound{A}\AgdaSymbol{)}\<%
\\
%
\>[2]\AgdaSymbol{→}\AgdaSpace{}%
\AgdaSymbol{(}\AgdaBound{M}\AgdaSpace{}%
\AgdaOperator{\AgdaDatatype{↝βₙ}}\AgdaSpace{}%
\AgdaBound{N}\AgdaSymbol{)}\<%
\\
%
\>[2]\AgdaSymbol{→}\AgdaSpace{}%
\AgdaSymbol{(}\AgdaFunction{embedTerm}\AgdaSpace{}%
\AgdaBound{M}\AgdaSpace{}%
\AgdaOperator{\AgdaDatatype{↝βb}}\AgdaSpace{}%
\AgdaFunction{embedTerm}\AgdaSpace{}%
\AgdaBound{N}\AgdaSymbol{)}\<%
\\
\>[0]\AgdaFunction{pres-red}\AgdaSpace{}%
\AgdaSymbol{(}\AgdaInductiveConstructor{ƛ}\AgdaSpace{}%
\AgdaBound{M₁}\AgdaSpace{}%
\AgdaOperator{\AgdaInductiveConstructor{·}}\AgdaSpace{}%
\AgdaBound{M₂}\AgdaSymbol{)}\AgdaSpace{}%
\AgdaBound{N}\AgdaSpace{}%
\AgdaInductiveConstructor{βₙ}\AgdaSpace{}%
\AgdaSymbol{=}\<%
\\
\>[0][@{}l@{\AgdaIndent{0}}]%
\>[2]\AgdaFunction{Eq.subst}\AgdaSpace{}%
\AgdaSymbol{(λ}\AgdaSpace{}%
\AgdaBound{t}\AgdaSpace{}%
\AgdaSymbol{→}\AgdaSpace{}%
\AgdaInductiveConstructor{ƛb}\AgdaSpace{}%
\AgdaSymbol{(}\AgdaFunction{embedTerm}\AgdaSpace{}%
\AgdaBound{M₁}\AgdaSymbol{)}\AgdaSpace{}%
\AgdaOperator{\AgdaInductiveConstructor{·b}}\AgdaSpace{}%
\AgdaInductiveConstructor{box}\AgdaSpace{}%
\AgdaSymbol{(}\AgdaFunction{embedTerm}\AgdaSpace{}%
\AgdaBound{M₂}\AgdaSymbol{)}\AgdaSpace{}%
\AgdaOperator{\AgdaDatatype{↝βb}}\AgdaSpace{}%
\AgdaBound{t}\AgdaSymbol{)}\<%
\\
%
\>[2]\AgdaSymbol{(}\AgdaFunction{pres-subst}\AgdaSpace{}%
\AgdaBound{M₁}\AgdaSpace{}%
\AgdaBound{M₂}\AgdaSymbol{)}\<%
\\
%
\>[2]\AgdaInductiveConstructor{βb}\<%
\\
\>[0]\AgdaFunction{pres-red}\AgdaSpace{}%
\AgdaBound{M}\AgdaSpace{}%
\AgdaBound{N}\AgdaSpace{}%
\AgdaSymbol{(}\AgdaInductiveConstructor{μ}\AgdaSpace{}%
\AgdaSymbol{\{}\AgdaArgument{M}\AgdaSpace{}%
\AgdaSymbol{=}\AgdaSpace{}%
\AgdaBound{M₁}\AgdaSymbol{\}}\AgdaSpace{}%
\AgdaSymbol{\{}\AgdaArgument{M'}\AgdaSpace{}%
\AgdaSymbol{=}\AgdaSpace{}%
\AgdaBound{M'}\AgdaSymbol{\}}\AgdaSpace{}%
\AgdaBound{red}\AgdaSymbol{)}\AgdaSpace{}%
\AgdaSymbol{=}\<%
\\
\>[0][@{}l@{\AgdaIndent{0}}]%
\>[2]\AgdaInductiveConstructor{μ}\AgdaSpace{}%
\AgdaSymbol{(}\AgdaFunction{pres-red}\AgdaSpace{}%
\AgdaBound{M₁}\AgdaSpace{}%
\AgdaBound{M'}\AgdaSpace{}%
\AgdaBound{red}\AgdaSymbol{)}\<%
\\
\>[0]\AgdaFunction{pres-red}\AgdaSpace{}%
\AgdaBound{M}\AgdaSpace{}%
\AgdaBound{N}\AgdaSpace{}%
\AgdaSymbol{(}\AgdaInductiveConstructor{ν<}\AgdaSpace{}%
\AgdaSymbol{\{}\AgdaArgument{M}\AgdaSpace{}%
\AgdaSymbol{=}\AgdaSpace{}%
\AgdaBound{M₁}\AgdaSymbol{\}}\AgdaSpace{}%
\AgdaSymbol{\{}\AgdaArgument{N}\AgdaSpace{}%
\AgdaSymbol{=}\AgdaSpace{}%
\AgdaBound{N₁}\AgdaSymbol{\}}\AgdaSpace{}%
\AgdaSymbol{\{}\AgdaArgument{N'}\AgdaSpace{}%
\AgdaSymbol{=}\AgdaSpace{}%
\AgdaBound{N'}\AgdaSymbol{\}}\AgdaSpace{}%
\AgdaBound{x}\AgdaSpace{}%
\AgdaBound{red}\AgdaSymbol{)}\AgdaSpace{}%
\AgdaSymbol{=}\<%
\\
\>[0][@{}l@{\AgdaIndent{0}}]%
\>[2]\AgdaInductiveConstructor{ν}\AgdaSpace{}%
\AgdaSymbol{(}\AgdaInductiveConstructor{ζ}\AgdaSpace{}%
\AgdaSymbol{(}\AgdaFunction{pres-red}\AgdaSpace{}%
\AgdaBound{N₁}\AgdaSpace{}%
\AgdaBound{N'}\AgdaSpace{}%
\AgdaBound{red}\AgdaSymbol{))}\<%
\\
\>[0]\AgdaFunction{pres-red}\AgdaSpace{}%
\AgdaBound{M}\AgdaSpace{}%
\AgdaBound{N}\AgdaSpace{}%
\AgdaSymbol{(}\AgdaInductiveConstructor{ν}\AgdaSpace{}%
\AgdaSymbol{\{}\AgdaArgument{M}\AgdaSpace{}%
\AgdaSymbol{=}\AgdaSpace{}%
\AgdaBound{M₁}\AgdaSymbol{\}}\AgdaSpace{}%
\AgdaSymbol{\{}\AgdaArgument{N}\AgdaSpace{}%
\AgdaSymbol{=}\AgdaSpace{}%
\AgdaBound{N₁}\AgdaSymbol{\}}\AgdaSpace{}%
\AgdaSymbol{\{}\AgdaArgument{N'}\AgdaSpace{}%
\AgdaSymbol{=}\AgdaSpace{}%
\AgdaBound{N'}\AgdaSymbol{\}}\AgdaSpace{}%
\AgdaBound{red}\AgdaSymbol{)}\AgdaSpace{}%
\AgdaSymbol{=}\<%
\\
\>[0][@{}l@{\AgdaIndent{0}}]%
\>[2]\AgdaInductiveConstructor{ν}\AgdaSpace{}%
\AgdaSymbol{(}\AgdaInductiveConstructor{ζ}\AgdaSpace{}%
\AgdaSymbol{(}\AgdaFunction{pres-red}\AgdaSpace{}%
\AgdaBound{N₁}\AgdaSpace{}%
\AgdaBound{N'}\AgdaSpace{}%
\AgdaBound{red}\AgdaSymbol{))}\<%
\\
\>[0]\AgdaFunction{pres-red}\AgdaSpace{}%
\AgdaBound{M}\AgdaSpace{}%
\AgdaBound{N}\AgdaSpace{}%
\AgdaSymbol{(}\AgdaInductiveConstructor{ξ}\AgdaSpace{}%
\AgdaSymbol{\{}\AgdaArgument{M}\AgdaSpace{}%
\AgdaSymbol{=}\AgdaSpace{}%
\AgdaBound{M₁}\AgdaSymbol{\}}\AgdaSpace{}%
\AgdaSymbol{\{}\AgdaArgument{M'}\AgdaSpace{}%
\AgdaSymbol{=}\AgdaSpace{}%
\AgdaBound{M'}\AgdaSymbol{\}}\AgdaSpace{}%
\AgdaBound{red}\AgdaSymbol{)}\AgdaSpace{}%
\AgdaSymbol{=}\<%
\\
\>[0][@{}l@{\AgdaIndent{0}}]%
\>[2]\AgdaInductiveConstructor{ξ}\AgdaSpace{}%
\AgdaSymbol{(}\AgdaFunction{pres-red}\AgdaSpace{}%
\AgdaBound{M₁}\AgdaSpace{}%
\AgdaBound{M'}\AgdaSpace{}%
\AgdaBound{red}\AgdaSymbol{)}\<%
\\
%
\\[\AgdaEmptyExtraSkip]%
\>[0]\AgdaFunction{pres-eval}\AgdaSpace{}%
\AgdaSymbol{:}\AgdaSpace{}%
\AgdaSymbol{∀}\AgdaSpace{}%
\AgdaSymbol{\{}\AgdaBound{Γ}\AgdaSpace{}%
\AgdaBound{A}\AgdaSymbol{\}}\<%
\\
\>[0][@{}l@{\AgdaIndent{0}}]%
\>[2]\AgdaSymbol{→}\AgdaSpace{}%
\AgdaSymbol{(}\AgdaBound{M}\AgdaSpace{}%
\AgdaSymbol{:}\AgdaSpace{}%
\AgdaBound{Γ}\AgdaSpace{}%
\AgdaOperator{\AgdaDatatype{⊢}}\AgdaSpace{}%
\AgdaBound{A}\AgdaSymbol{)}\<%
\\
%
\>[2]\AgdaSymbol{→}\AgdaSpace{}%
\AgdaSymbol{(}\AgdaBound{N}\AgdaSpace{}%
\AgdaSymbol{:}\AgdaSpace{}%
\AgdaBound{Γ}\AgdaSpace{}%
\AgdaOperator{\AgdaDatatype{⊢}}\AgdaSpace{}%
\AgdaBound{A}\AgdaSymbol{)}\<%
\\
%
\>[2]\AgdaSymbol{→}\AgdaSpace{}%
\AgdaSymbol{(}\AgdaBound{M}\AgdaSpace{}%
\AgdaOperator{\AgdaDatatype{↝ₙ}}\AgdaSpace{}%
\AgdaBound{N}\AgdaSymbol{)}\<%
\\
%
\>[2]\AgdaSymbol{→}\AgdaSpace{}%
\AgdaSymbol{(}\AgdaFunction{embedTerm}\AgdaSpace{}%
\AgdaBound{M}\AgdaSpace{}%
\AgdaOperator{\AgdaDatatype{↝bₙ}}\AgdaSpace{}%
\AgdaFunction{embedTerm}\AgdaSpace{}%
\AgdaBound{N}\AgdaSymbol{)}\<%
\\
\>[0]\AgdaFunction{pres-eval}\AgdaSpace{}%
\AgdaSymbol{(}\AgdaInductiveConstructor{ƛ}\AgdaSpace{}%
\AgdaBound{M₁}\AgdaSpace{}%
\AgdaOperator{\AgdaInductiveConstructor{·}}\AgdaSpace{}%
\AgdaBound{M₂}\AgdaSymbol{)}\AgdaSpace{}%
\AgdaBound{N}\AgdaSpace{}%
\AgdaInductiveConstructor{βₙ}\AgdaSpace{}%
\AgdaSymbol{=}\<%
\\
\>[0][@{}l@{\AgdaIndent{0}}]%
\>[2]\AgdaFunction{Eq.subst}\<%
\\
\>[2][@{}l@{\AgdaIndent{0}}]%
\>[4]\AgdaSymbol{(λ}\AgdaSpace{}%
\AgdaBound{t}\AgdaSpace{}%
\AgdaSymbol{→}\AgdaSpace{}%
\AgdaInductiveConstructor{ƛb}\AgdaSpace{}%
\AgdaSymbol{(}\AgdaFunction{embedTerm}\AgdaSpace{}%
\AgdaBound{M₁}\AgdaSymbol{)}\AgdaSpace{}%
\AgdaOperator{\AgdaInductiveConstructor{·b}}\AgdaSpace{}%
\AgdaInductiveConstructor{box}\AgdaSpace{}%
\AgdaSymbol{(}\AgdaFunction{embedTerm}\AgdaSpace{}%
\AgdaBound{M₂}\AgdaSymbol{)}\AgdaSpace{}%
\AgdaOperator{\AgdaDatatype{↝bₙ}}\AgdaSpace{}%
\AgdaBound{t}\AgdaSymbol{)}\<%
\\
%
\>[4]\AgdaSymbol{(}\AgdaFunction{pres-subst}\AgdaSpace{}%
\AgdaBound{M₁}\AgdaSpace{}%
\AgdaBound{M₂}\AgdaSymbol{)}\<%
\\
%
\>[4]\AgdaInductiveConstructor{βb}\<%
\\
\>[0]\AgdaFunction{pres-eval}\AgdaSpace{}%
\AgdaBound{M}\AgdaSpace{}%
\AgdaBound{N}\AgdaSpace{}%
\AgdaSymbol{(}\AgdaInductiveConstructor{μ}\AgdaSpace{}%
\AgdaSymbol{\{}\AgdaArgument{M}\AgdaSpace{}%
\AgdaSymbol{=}\AgdaSpace{}%
\AgdaBound{M₁}\AgdaSymbol{\}}\AgdaSpace{}%
\AgdaSymbol{\{}\AgdaArgument{M'}\AgdaSpace{}%
\AgdaSymbol{=}\AgdaSpace{}%
\AgdaBound{M'}\AgdaSymbol{\}}\AgdaSpace{}%
\AgdaBound{red}\AgdaSymbol{)}\AgdaSpace{}%
\AgdaSymbol{=}\<%
\\
\>[0][@{}l@{\AgdaIndent{0}}]%
\>[2]\AgdaInductiveConstructor{μ}\AgdaSpace{}%
\AgdaSymbol{(}\AgdaFunction{pres-eval}\AgdaSpace{}%
\AgdaBound{M₁}\AgdaSpace{}%
\AgdaBound{M'}\AgdaSpace{}%
\AgdaBound{red}\AgdaSymbol{)}\<%
\end{code}
\end{lemma}

\begin{lemma}
  \label{lem:pres-simultaneous-subst}
  \textbf{Preservation of simultaneous substitution}:
  \small
  \begin{code}[hide]%
\>[0]\AgdaKeyword{module}\AgdaSpace{}%
\AgdaModule{nameProofs2}\AgdaSpace{}%
\AgdaKeyword{where}\<%
\\
%
\\[\AgdaEmptyExtraSkip]%
\>[0]\AgdaKeyword{open}\AgdaSpace{}%
\AgdaKeyword{import}\AgdaSpace{}%
\AgdaModule{boxCalc}\<%
\\
\>[0]\AgdaKeyword{open}\AgdaSpace{}%
\AgdaKeyword{import}\AgdaSpace{}%
\AgdaModule{calc}\<%
\\
\>[0]\AgdaKeyword{open}\AgdaSpace{}%
\AgdaKeyword{import}\AgdaSpace{}%
\AgdaModule{embedCBNIntoCBB}\<%
\\
\>[0]\AgdaKeyword{open}\AgdaSpace{}%
\AgdaKeyword{import}\AgdaSpace{}%
\AgdaModule{nameCalcRed}\<%
\\
\>[0]\AgdaKeyword{open}\AgdaSpace{}%
\AgdaKeyword{import}\AgdaSpace{}%
\AgdaModule{boxCalcRed}\<%
\\
%
\\[\AgdaEmptyExtraSkip]%
\>[0]\AgdaKeyword{open}\AgdaSpace{}%
\AgdaKeyword{import}\AgdaSpace{}%
\AgdaModule{Data.Empty}\AgdaSpace{}%
\AgdaKeyword{using}\AgdaSpace{}%
\AgdaSymbol{(}\AgdaFunction{⊥}\AgdaSymbol{;}\AgdaSpace{}%
\AgdaFunction{⊥-elim}\AgdaSymbol{)}\<%
\\
\>[0]\AgdaKeyword{import}\AgdaSpace{}%
\AgdaModule{Relation.Binary.PropositionalEquality}\AgdaSpace{}%
\AgdaSymbol{as}\AgdaSpace{}%
\AgdaModule{Eq}\<%
\\
\>[0]\AgdaKeyword{open}\AgdaSpace{}%
\AgdaModule{Eq}\AgdaSpace{}%
\AgdaKeyword{using}\AgdaSpace{}%
\AgdaSymbol{(}\AgdaOperator{\AgdaDatatype{\AgdaUnderscore{}≡\AgdaUnderscore{}}}\AgdaSymbol{;}\AgdaSpace{}%
\AgdaInductiveConstructor{refl}\AgdaSymbol{;}\AgdaSpace{}%
\AgdaFunction{cong}\AgdaSymbol{;}\AgdaSpace{}%
\AgdaFunction{cong₂}\AgdaSymbol{;}\AgdaSpace{}%
\AgdaFunction{sym}\AgdaSymbol{)}\<%
\\
%
\\[\AgdaEmptyExtraSkip]%
\>[0]\AgdaKeyword{open}\AgdaSpace{}%
\AgdaModule{Eq.≡-Reasoning}\AgdaSpace{}%
\AgdaKeyword{using}\AgdaSpace{}%
\AgdaSymbol{(}\AgdaOperator{\AgdaFunction{begin\AgdaUnderscore{}}}\AgdaSymbol{;}\AgdaSpace{}%
\AgdaFunction{step-≡-∣}\AgdaSymbol{;}\AgdaSpace{}%
\AgdaFunction{step-≡-⟩}\AgdaSymbol{;}\AgdaSpace{}%
\AgdaOperator{\AgdaFunction{\AgdaUnderscore{}∎}}\AgdaSymbol{)}\<%
\\
\>[0]\AgdaKeyword{open}\AgdaSpace{}%
\AgdaModule{Eq.≡-Reasoning}\AgdaSpace{}%
\AgdaKeyword{renaming}\AgdaSpace{}%
\AgdaSymbol{(}\AgdaOperator{\AgdaFunction{begin\AgdaUnderscore{}}}\AgdaSpace{}%
\AgdaSymbol{to}\AgdaSpace{}%
\AgdaOperator{\AgdaFunction{beginEq\AgdaUnderscore{}}}\AgdaSymbol{;}\AgdaSpace{}%
\AgdaOperator{\AgdaFunction{\AgdaUnderscore{}∎}}\AgdaSpace{}%
\AgdaSymbol{to}\AgdaSpace{}%
\AgdaOperator{\AgdaFunction{\AgdaUnderscore{}∎eq}}\AgdaSymbol{)}\<%
\\
\>[0]\AgdaKeyword{open}\AgdaSpace{}%
\AgdaKeyword{import}\AgdaSpace{}%
\AgdaModule{Function.Bundles}\AgdaSpace{}%
\AgdaKeyword{using}\AgdaSpace{}%
\AgdaSymbol{(}\AgdaOperator{\AgdaFunction{\AgdaUnderscore{}⇔\AgdaUnderscore{}}}\AgdaSymbol{;}\AgdaSpace{}%
\AgdaFunction{mk⇔}\AgdaSymbol{)}\<%
\\
%
\\[\AgdaEmptyExtraSkip]%
\>[0]\AgdaComment{--\ GIRARD'S\ embedding}\<%
\\
%
\\[\AgdaEmptyExtraSkip]%
\>[0]\AgdaComment{--\ Preservation\ of\ typing\ by\ Girard's\ translation}\<%
\\
\>[0]\AgdaFunction{pres-typing}\AgdaSpace{}%
\AgdaSymbol{:}\AgdaSpace{}%
\AgdaSymbol{∀}\AgdaSpace{}%
\AgdaSymbol{\{}\AgdaBound{Γ}\AgdaSpace{}%
\AgdaBound{A}\AgdaSymbol{\}}\AgdaSpace{}%
\AgdaSymbol{→}\AgdaSpace{}%
\AgdaBound{Γ}\AgdaSpace{}%
\AgdaOperator{\AgdaDatatype{⊢}}\AgdaSpace{}%
\AgdaBound{A}\AgdaSpace{}%
\AgdaSymbol{→}\AgdaSpace{}%
\AgdaFunction{embedContext}\AgdaSpace{}%
\AgdaBound{Γ}\AgdaSpace{}%
\AgdaOperator{\AgdaDatatype{⊢b}}\AgdaSpace{}%
\AgdaFunction{embedType}\AgdaSpace{}%
\AgdaBound{A}\<%
\\
\>[0]\AgdaFunction{pres-typing}\AgdaSpace{}%
\AgdaBound{x}\AgdaSpace{}%
\AgdaSymbol{=}\AgdaSpace{}%
\AgdaFunction{embedTerm}\AgdaSpace{}%
\AgdaBound{x}\<%
\\
%
\\[\AgdaEmptyExtraSkip]%
\>[0]\AgdaFunction{pres-ext}\AgdaSpace{}%
\AgdaSymbol{:}\AgdaSpace{}%
\AgdaSymbol{∀}\AgdaSpace{}%
\AgdaSymbol{\{}\AgdaBound{Γ}\AgdaSpace{}%
\AgdaBound{Δ}\AgdaSpace{}%
\AgdaBound{A'}\AgdaSpace{}%
\AgdaBound{A}\AgdaSymbol{\}}\<%
\\
\>[0][@{}l@{\AgdaIndent{0}}]%
\>[2]\AgdaSymbol{→}\AgdaSpace{}%
\AgdaSymbol{(}\AgdaBound{f'}\AgdaSpace{}%
\AgdaSymbol{:}\AgdaSpace{}%
\AgdaSymbol{∀}\AgdaSpace{}%
\AgdaSymbol{\{}\AgdaBound{A}\AgdaSymbol{\}}\AgdaSpace{}%
\AgdaSymbol{→}\AgdaSpace{}%
\AgdaBound{A}\AgdaSpace{}%
\AgdaOperator{\AgdaDatatype{∈}}\AgdaSpace{}%
\AgdaBound{Γ}\AgdaSpace{}%
\AgdaSymbol{→}\AgdaSpace{}%
\AgdaBound{A}\AgdaSpace{}%
\AgdaOperator{\AgdaDatatype{∈}}\AgdaSpace{}%
\AgdaBound{Δ}\AgdaSymbol{)}\<%
\\
%
\>[2]\AgdaSymbol{→}\AgdaSpace{}%
\AgdaSymbol{(}\AgdaBound{fb'}\AgdaSpace{}%
\AgdaSymbol{:}\AgdaSpace{}%
\AgdaSymbol{∀}\AgdaSpace{}%
\AgdaSymbol{\{}\AgdaBound{Ab}\AgdaSymbol{\}}\AgdaSpace{}%
\AgdaSymbol{→}\AgdaSpace{}%
\AgdaBound{Ab}\AgdaSpace{}%
\AgdaOperator{\AgdaDatatype{∈b}}\AgdaSpace{}%
\AgdaFunction{embedContext}\AgdaSpace{}%
\AgdaBound{Γ}\AgdaSpace{}%
\AgdaSymbol{→}\AgdaSpace{}%
\AgdaBound{Ab}\AgdaSpace{}%
\AgdaOperator{\AgdaDatatype{∈b}}\AgdaSpace{}%
\AgdaFunction{embedContext}\AgdaSpace{}%
\AgdaBound{Δ}\AgdaSymbol{)}\<%
\\
%
\>[2]\AgdaSymbol{→}\AgdaSpace{}%
\AgdaSymbol{(\{}\AgdaBound{A'}\AgdaSpace{}%
\AgdaSymbol{:}\AgdaSpace{}%
\AgdaDatatype{Type}\AgdaSymbol{\}}\AgdaSpace{}%
\AgdaSymbol{(}\AgdaBound{x}\AgdaSpace{}%
\AgdaSymbol{:}\AgdaSpace{}%
\AgdaBound{A'}\AgdaSpace{}%
\AgdaOperator{\AgdaDatatype{∈}}\AgdaSpace{}%
\AgdaBound{Γ}\AgdaSymbol{)}\AgdaSpace{}%
\AgdaSymbol{→}\AgdaSpace{}%
\AgdaBound{fb'}\AgdaSpace{}%
\AgdaSymbol{(}\AgdaFunction{embedLookup}\AgdaSpace{}%
\AgdaBound{x}\AgdaSymbol{)}\AgdaSpace{}%
\AgdaOperator{\AgdaDatatype{≡}}\AgdaSpace{}%
\AgdaFunction{embedLookup}\AgdaSpace{}%
\AgdaSymbol{(}\AgdaBound{f'}\AgdaSpace{}%
\AgdaBound{x}\AgdaSymbol{))}\<%
\\
%
\>[2]\AgdaSymbol{→}\AgdaSpace{}%
\AgdaSymbol{(}\AgdaBound{x}\AgdaSpace{}%
\AgdaSymbol{:}\AgdaSpace{}%
\AgdaBound{A'}\AgdaSpace{}%
\AgdaOperator{\AgdaDatatype{∈}}\AgdaSpace{}%
\AgdaBound{Γ}\AgdaSpace{}%
\AgdaOperator{\AgdaInductiveConstructor{,}}\AgdaSpace{}%
\AgdaBound{A}\AgdaSymbol{)}\<%
\\
%
\>[2]\AgdaSymbol{→}\AgdaSpace{}%
\AgdaFunction{extb}\AgdaSpace{}%
\AgdaBound{fb'}\AgdaSpace{}%
\AgdaSymbol{(}\AgdaFunction{embedLookup}\AgdaSpace{}%
\AgdaBound{x}\AgdaSymbol{)}\AgdaSpace{}%
\AgdaOperator{\AgdaDatatype{≡}}\AgdaSpace{}%
\AgdaFunction{embedLookup}\AgdaSpace{}%
\AgdaSymbol{(}\AgdaFunction{ext}\AgdaSpace{}%
\AgdaBound{f'}\AgdaSpace{}%
\AgdaBound{x}\AgdaSymbol{)}\<%
\\
\>[0]\AgdaFunction{pres-ext}\AgdaSpace{}%
\AgdaBound{f'}\AgdaSpace{}%
\AgdaBound{fb'}\AgdaSpace{}%
\AgdaBound{h}\AgdaSpace{}%
\AgdaInductiveConstructor{Z}\AgdaSpace{}%
\AgdaSymbol{=}\AgdaSpace{}%
\AgdaInductiveConstructor{refl}\<%
\\
\>[0]\AgdaFunction{pres-ext}\AgdaSpace{}%
\AgdaBound{f'}\AgdaSpace{}%
\AgdaBound{fb'}\AgdaSpace{}%
\AgdaBound{h}\AgdaSpace{}%
\AgdaSymbol{(}\AgdaInductiveConstructor{S}\AgdaSpace{}%
\AgdaBound{x}\AgdaSymbol{)}\AgdaSpace{}%
\AgdaSymbol{=}\AgdaSpace{}%
\AgdaFunction{cong}\AgdaSpace{}%
\AgdaInductiveConstructor{Sb}\AgdaSpace{}%
\AgdaSymbol{(}\AgdaBound{h}\AgdaSpace{}%
\AgdaBound{x}\AgdaSymbol{)}\<%
\\
%
\\[\AgdaEmptyExtraSkip]%
\>[0]\AgdaFunction{pres-rename}\AgdaSpace{}%
\AgdaSymbol{:}\AgdaSpace{}%
\AgdaSymbol{∀}\AgdaSpace{}%
\AgdaSymbol{\{}\AgdaBound{Γ}\AgdaSpace{}%
\AgdaBound{A}\AgdaSpace{}%
\AgdaBound{Δ}\AgdaSymbol{\}}\<%
\\
\>[0][@{}l@{\AgdaIndent{0}}]%
\>[2]\AgdaSymbol{→}\AgdaSpace{}%
\AgdaSymbol{(}\AgdaBound{f'}\AgdaSpace{}%
\AgdaSymbol{:}\AgdaSpace{}%
\AgdaSymbol{∀}\AgdaSpace{}%
\AgdaSymbol{\{}\AgdaBound{A}\AgdaSymbol{\}}\AgdaSpace{}%
\AgdaSymbol{→}\AgdaSpace{}%
\AgdaBound{A}\AgdaSpace{}%
\AgdaOperator{\AgdaDatatype{∈}}\AgdaSpace{}%
\AgdaBound{Γ}\AgdaSpace{}%
\AgdaSymbol{→}\AgdaSpace{}%
\AgdaBound{A}\AgdaSpace{}%
\AgdaOperator{\AgdaDatatype{∈}}\AgdaSpace{}%
\AgdaBound{Δ}\AgdaSymbol{)}\<%
\\
%
\>[2]\AgdaSymbol{→}\AgdaSpace{}%
\AgdaSymbol{(}\AgdaBound{fb'}\AgdaSpace{}%
\AgdaSymbol{:}\AgdaSpace{}%
\AgdaSymbol{∀}\AgdaSpace{}%
\AgdaSymbol{\{}\AgdaBound{Ab}\AgdaSymbol{\}}\AgdaSpace{}%
\AgdaSymbol{→}\AgdaSpace{}%
\AgdaBound{Ab}\AgdaSpace{}%
\AgdaOperator{\AgdaDatatype{∈b}}\AgdaSpace{}%
\AgdaFunction{embedContext}\AgdaSpace{}%
\AgdaBound{Γ}\AgdaSpace{}%
\AgdaSymbol{→}\AgdaSpace{}%
\AgdaBound{Ab}\AgdaSpace{}%
\AgdaOperator{\AgdaDatatype{∈b}}\AgdaSpace{}%
\AgdaFunction{embedContext}\AgdaSpace{}%
\AgdaBound{Δ}\AgdaSymbol{)}\<%
\\
%
\>[2]\AgdaSymbol{→}\AgdaSpace{}%
\AgdaSymbol{(\{}\AgdaBound{A'}\AgdaSpace{}%
\AgdaSymbol{:}\AgdaSpace{}%
\AgdaDatatype{Type}\AgdaSymbol{\}}\AgdaSpace{}%
\AgdaSymbol{(}\AgdaBound{x}\AgdaSpace{}%
\AgdaSymbol{:}\AgdaSpace{}%
\AgdaBound{A'}\AgdaSpace{}%
\AgdaOperator{\AgdaDatatype{∈}}\AgdaSpace{}%
\AgdaBound{Γ}\AgdaSymbol{)}\AgdaSpace{}%
\AgdaSymbol{→}\AgdaSpace{}%
\AgdaBound{fb'}\AgdaSpace{}%
\AgdaSymbol{(}\AgdaFunction{embedLookup}\AgdaSpace{}%
\AgdaBound{x}\AgdaSymbol{)}\AgdaSpace{}%
\AgdaOperator{\AgdaDatatype{≡}}\AgdaSpace{}%
\AgdaFunction{embedLookup}\AgdaSpace{}%
\AgdaSymbol{(}\AgdaBound{f'}\AgdaSpace{}%
\AgdaBound{x}\AgdaSymbol{))}\<%
\\
%
\>[2]\AgdaSymbol{→}\AgdaSpace{}%
\AgdaSymbol{(}\AgdaBound{t}\AgdaSpace{}%
\AgdaSymbol{:}\AgdaSpace{}%
\AgdaBound{Γ}\AgdaSpace{}%
\AgdaOperator{\AgdaDatatype{⊢}}\AgdaSpace{}%
\AgdaBound{A}\AgdaSymbol{)}\<%
\\
%
\>[2]\AgdaSymbol{→}\AgdaSpace{}%
\AgdaFunction{renameb}\AgdaSpace{}%
\AgdaBound{fb'}\AgdaSpace{}%
\AgdaSymbol{(}\AgdaFunction{embedTerm}\AgdaSpace{}%
\AgdaBound{t}\AgdaSymbol{)}\AgdaSpace{}%
\AgdaOperator{\AgdaDatatype{≡}}\AgdaSpace{}%
\AgdaFunction{embedTerm}\AgdaSpace{}%
\AgdaSymbol{(}\AgdaFunction{rename}\AgdaSpace{}%
\AgdaBound{f'}\AgdaSpace{}%
\AgdaBound{t}\AgdaSymbol{)}\<%
\\
\>[0]\AgdaFunction{pres-rename}\AgdaSpace{}%
\AgdaBound{f'}\AgdaSpace{}%
\AgdaBound{fb'}\AgdaSpace{}%
\AgdaBound{h}\AgdaSpace{}%
\AgdaSymbol{(}\AgdaInductiveConstructor{`}\AgdaSpace{}%
\AgdaBound{x}\AgdaSymbol{)}\AgdaSpace{}%
\AgdaSymbol{=}\<%
\\
\>[0][@{}l@{\AgdaIndent{0}}]%
\>[2]\AgdaFunction{cong}\AgdaSpace{}%
\AgdaInductiveConstructor{ε}\AgdaSpace{}%
\AgdaSymbol{(}\AgdaBound{h}\AgdaSpace{}%
\AgdaBound{x}\AgdaSymbol{)}\<%
\\
\>[0]\AgdaFunction{pres-rename}\AgdaSpace{}%
\AgdaBound{f'}\AgdaSpace{}%
\AgdaBound{fb'}\AgdaSpace{}%
\AgdaBound{h}\AgdaSpace{}%
\AgdaSymbol{(}\AgdaInductiveConstructor{ƛ}\AgdaSpace{}%
\AgdaBound{t}\AgdaSymbol{)}\AgdaSpace{}%
\AgdaSymbol{=}\<%
\\
\>[0][@{}l@{\AgdaIndent{0}}]%
\>[2]\AgdaFunction{cong}\AgdaSpace{}%
\AgdaInductiveConstructor{ƛb}\AgdaSpace{}%
\AgdaSymbol{(}\AgdaFunction{pres-rename}\AgdaSpace{}%
\AgdaSymbol{(}\AgdaFunction{calc.ext}\AgdaSpace{}%
\AgdaBound{f'}\AgdaSymbol{)}\AgdaSpace{}%
\AgdaSymbol{(}\AgdaFunction{extb}\AgdaSpace{}%
\AgdaBound{fb'}\AgdaSymbol{)}\AgdaSpace{}%
\AgdaSymbol{(}\AgdaFunction{pres-ext}\AgdaSpace{}%
\AgdaBound{f'}\AgdaSpace{}%
\AgdaBound{fb'}\AgdaSpace{}%
\AgdaBound{h}\AgdaSymbol{)}\AgdaSpace{}%
\AgdaBound{t}\AgdaSymbol{)}\<%
\\
\>[0]\AgdaFunction{pres-rename}\AgdaSpace{}%
\AgdaBound{f'}\AgdaSpace{}%
\AgdaBound{fb'}\AgdaSpace{}%
\AgdaBound{h}\AgdaSpace{}%
\AgdaSymbol{(}\AgdaBound{t}\AgdaSpace{}%
\AgdaOperator{\AgdaInductiveConstructor{·}}\AgdaSpace{}%
\AgdaBound{t₁}\AgdaSymbol{)}\AgdaSpace{}%
\AgdaSymbol{=}\<%
\\
\>[0][@{}l@{\AgdaIndent{0}}]%
\>[2]\AgdaFunction{cong₂}\AgdaSpace{}%
\AgdaOperator{\AgdaInductiveConstructor{\AgdaUnderscore{}·b\AgdaUnderscore{}}}\AgdaSpace{}%
\AgdaSymbol{(}\AgdaFunction{pres-rename}\AgdaSpace{}%
\AgdaBound{f'}\AgdaSpace{}%
\AgdaBound{fb'}\AgdaSpace{}%
\AgdaBound{h}\AgdaSpace{}%
\AgdaBound{t}\AgdaSymbol{)}\AgdaSpace{}%
\AgdaSymbol{(}\AgdaFunction{cong}\AgdaSpace{}%
\AgdaInductiveConstructor{box}\AgdaSpace{}%
\AgdaSymbol{(}\AgdaFunction{pres-rename}\AgdaSpace{}%
\AgdaBound{f'}\AgdaSpace{}%
\AgdaBound{fb'}\AgdaSpace{}%
\AgdaBound{h}\AgdaSpace{}%
\AgdaBound{t₁}\AgdaSymbol{))}\<%
\\
%
\\[\AgdaEmptyExtraSkip]%
%
\\[\AgdaEmptyExtraSkip]%
\>[0]\AgdaFunction{pres-exts}\AgdaSpace{}%
\AgdaSymbol{:}\AgdaSpace{}%
\AgdaSymbol{∀}\AgdaSpace{}%
\AgdaSymbol{\{}\AgdaBound{Γ}\AgdaSpace{}%
\AgdaBound{Δ}\AgdaSymbol{\}}\<%
\\
\>[0][@{}l@{\AgdaIndent{0}}]%
\>[2]\AgdaSymbol{→}\AgdaSpace{}%
\AgdaSymbol{(}\AgdaBound{fb'}\AgdaSpace{}%
\AgdaSymbol{:}\AgdaSpace{}%
\AgdaSymbol{∀}\AgdaSpace{}%
\AgdaSymbol{\{}\AgdaBound{Ab}\AgdaSymbol{\}}\AgdaSpace{}%
\AgdaSymbol{→}\AgdaSpace{}%
\AgdaInductiveConstructor{□}\AgdaSpace{}%
\AgdaBound{Ab}\AgdaSpace{}%
\AgdaOperator{\AgdaDatatype{∈b}}\AgdaSpace{}%
\AgdaFunction{embedContext}\AgdaSpace{}%
\AgdaBound{Γ}\AgdaSpace{}%
\AgdaSymbol{→}\AgdaSpace{}%
\AgdaFunction{embedContext}\AgdaSpace{}%
\AgdaBound{Δ}\AgdaSpace{}%
\AgdaOperator{\AgdaDatatype{⊢b}}\AgdaSpace{}%
\AgdaBound{Ab}\AgdaSymbol{)}\<%
\\
%
\>[2]\AgdaSymbol{→}\AgdaSpace{}%
\AgdaSymbol{(}\AgdaBound{f'}\AgdaSpace{}%
\AgdaSymbol{:}\AgdaSpace{}%
\AgdaSymbol{∀}\AgdaSpace{}%
\AgdaSymbol{\{}\AgdaBound{A}\AgdaSymbol{\}}\AgdaSpace{}%
\AgdaSymbol{→}\AgdaSpace{}%
\AgdaBound{A}\AgdaSpace{}%
\AgdaOperator{\AgdaDatatype{∈}}\AgdaSpace{}%
\AgdaBound{Γ}\AgdaSpace{}%
\AgdaSymbol{→}\AgdaSpace{}%
\AgdaBound{Δ}\AgdaSpace{}%
\AgdaOperator{\AgdaDatatype{⊢}}\AgdaSpace{}%
\AgdaBound{A}\AgdaSymbol{)}\<%
\\
%
\>[2]\AgdaSymbol{→}\AgdaSpace{}%
\AgdaSymbol{(\{}\AgdaBound{A'}\AgdaSpace{}%
\AgdaSymbol{:}\AgdaSpace{}%
\AgdaDatatype{Type}\AgdaSymbol{\}}\AgdaSpace{}%
\AgdaSymbol{(}\AgdaBound{x}\AgdaSpace{}%
\AgdaSymbol{:}\AgdaSpace{}%
\AgdaBound{A'}\AgdaSpace{}%
\AgdaOperator{\AgdaDatatype{∈}}\AgdaSpace{}%
\AgdaBound{Γ}\AgdaSymbol{)}\AgdaSpace{}%
\AgdaSymbol{→}\AgdaSpace{}%
\AgdaBound{fb'}\AgdaSpace{}%
\AgdaSymbol{(}\AgdaFunction{embedLookup}\AgdaSpace{}%
\AgdaBound{x}\AgdaSymbol{)}\AgdaSpace{}%
\AgdaOperator{\AgdaDatatype{≡}}\AgdaSpace{}%
\AgdaFunction{embedTerm}\AgdaSpace{}%
\AgdaSymbol{(}\AgdaBound{f'}\AgdaSpace{}%
\AgdaBound{x}\AgdaSymbol{))}\<%
\\
%
\>[2]\AgdaSymbol{→}\AgdaSpace{}%
\AgdaSymbol{\{}\AgdaBound{A'}\AgdaSpace{}%
\AgdaBound{A}\AgdaSpace{}%
\AgdaSymbol{:}\AgdaSpace{}%
\AgdaDatatype{Type}\AgdaSymbol{\}}\AgdaSpace{}%
\AgdaSymbol{(}\AgdaBound{x}\AgdaSpace{}%
\AgdaSymbol{:}\AgdaSpace{}%
\AgdaBound{A'}\AgdaSpace{}%
\AgdaOperator{\AgdaDatatype{∈}}\AgdaSpace{}%
\AgdaBound{Γ}\AgdaSpace{}%
\AgdaOperator{\AgdaInductiveConstructor{,}}\AgdaSpace{}%
\AgdaBound{A}\AgdaSymbol{)}\AgdaSpace{}%
\AgdaSymbol{→}\AgdaSpace{}%
\AgdaFunction{extsb}\AgdaSpace{}%
\AgdaBound{fb'}\AgdaSpace{}%
\AgdaSymbol{(}\AgdaFunction{embedLookup}\AgdaSpace{}%
\AgdaBound{x}\AgdaSymbol{)}\AgdaSpace{}%
\AgdaOperator{\AgdaDatatype{≡}}\AgdaSpace{}%
\AgdaFunction{embedTerm}\AgdaSpace{}%
\AgdaSymbol{(}\AgdaFunction{exts}\AgdaSpace{}%
\AgdaBound{f'}\AgdaSpace{}%
\AgdaBound{x}\AgdaSymbol{)}\<%
\\
\>[0]\AgdaFunction{pres-exts}\AgdaSpace{}%
\AgdaBound{fb'}\AgdaSpace{}%
\AgdaBound{f'}\AgdaSpace{}%
\AgdaBound{h}\AgdaSpace{}%
\AgdaInductiveConstructor{Z}\AgdaSpace{}%
\AgdaSymbol{=}\AgdaSpace{}%
\AgdaInductiveConstructor{refl}\<%
\\
\>[0]\AgdaFunction{pres-exts}\AgdaSpace{}%
\AgdaBound{fb'}\AgdaSpace{}%
\AgdaBound{f'}\AgdaSpace{}%
\AgdaBound{h}\AgdaSpace{}%
\AgdaSymbol{(}\AgdaInductiveConstructor{S}\AgdaSpace{}%
\AgdaBound{x}\AgdaSymbol{)}\AgdaSpace{}%
\AgdaSymbol{=}\<%
\\
\>[0][@{}l@{\AgdaIndent{0}}]%
\>[2]\AgdaOperator{\AgdaFunction{beginEq}}\<%
\\
\>[2][@{}l@{\AgdaIndent{0}}]%
\>[4]\AgdaFunction{renameb}\AgdaSpace{}%
\AgdaInductiveConstructor{Sb}\AgdaSpace{}%
\AgdaSymbol{(}\AgdaBound{fb'}\AgdaSpace{}%
\AgdaSymbol{(}\AgdaFunction{embedLookup}\AgdaSpace{}%
\AgdaBound{x}\AgdaSymbol{))}\<%
\\
%
\>[2]\AgdaFunction{≡⟨}\AgdaSpace{}%
\AgdaFunction{cong}\AgdaSpace{}%
\AgdaSymbol{(}\AgdaFunction{renameb}\AgdaSpace{}%
\AgdaInductiveConstructor{Sb}\AgdaSymbol{)}\AgdaSpace{}%
\AgdaSymbol{(}\AgdaBound{h}\AgdaSpace{}%
\AgdaBound{x}\AgdaSymbol{)}\AgdaSpace{}%
\AgdaFunction{⟩}\<%
\\
\>[2][@{}l@{\AgdaIndent{0}}]%
\>[4]\AgdaFunction{renameb}\AgdaSpace{}%
\AgdaInductiveConstructor{Sb}\AgdaSpace{}%
\AgdaSymbol{(}\AgdaFunction{embedTerm}\AgdaSpace{}%
\AgdaSymbol{(}\AgdaBound{f'}\AgdaSpace{}%
\AgdaBound{x}\AgdaSymbol{))}\<%
\\
%
\>[2]\AgdaFunction{≡⟨}\AgdaSpace{}%
\AgdaFunction{pres-rename}\AgdaSpace{}%
\AgdaInductiveConstructor{S}\AgdaSpace{}%
\AgdaInductiveConstructor{Sb}\AgdaSpace{}%
\AgdaSymbol{(λ}\AgdaSpace{}%
\AgdaBound{x₁}\AgdaSpace{}%
\AgdaSymbol{→}\AgdaSpace{}%
\AgdaInductiveConstructor{refl}\AgdaSymbol{)}\AgdaSpace{}%
\AgdaSymbol{(}\AgdaBound{f'}\AgdaSpace{}%
\AgdaBound{x}\AgdaSymbol{)}\AgdaSpace{}%
\AgdaFunction{⟩}\<%
\\
\>[2][@{}l@{\AgdaIndent{0}}]%
\>[4]\AgdaFunction{embedTerm}\AgdaSpace{}%
\AgdaSymbol{(}\AgdaFunction{rename}\AgdaSpace{}%
\AgdaInductiveConstructor{S}\AgdaSpace{}%
\AgdaSymbol{(}\AgdaBound{f'}\AgdaSpace{}%
\AgdaBound{x}\AgdaSymbol{))}\<%
\\
%
\>[2]\AgdaOperator{\AgdaFunction{∎eq}}\<%
\\
\>[0]\<%
\end{code}
\begin{code}%
\>[0]\AgdaComment{--\ Preservation\ and\ Reflection\ of\ }\<%
\\
\>[0]\AgdaComment{--\ simultaneous\ substitution\ of\ Girard's\ translation}\<%
\\
\>[0]\AgdaFunction{pres-simultaneous-subst}\AgdaSpace{}%
\AgdaSymbol{:}\AgdaSpace{}%
\AgdaSymbol{∀}\AgdaSpace{}%
\AgdaSymbol{\{}\AgdaBound{Γ}\AgdaSpace{}%
\AgdaBound{Δ}\AgdaSpace{}%
\AgdaBound{A}\AgdaSymbol{\}}\<%
\\
\>[0][@{}l@{\AgdaIndent{0}}]%
\>[2]\AgdaSymbol{→}\AgdaSpace{}%
\AgdaSymbol{(}\AgdaBound{fb'}\AgdaSpace{}%
\AgdaSymbol{:}\AgdaSpace{}%
\AgdaSymbol{∀}\AgdaSpace{}%
\AgdaSymbol{\{}\AgdaBound{Ab}\AgdaSymbol{\}}\AgdaSpace{}%
\AgdaSymbol{→}\AgdaSpace{}%
\AgdaInductiveConstructor{□}\AgdaSpace{}%
\AgdaBound{Ab}\AgdaSpace{}%
\AgdaOperator{\AgdaDatatype{∈b}}\AgdaSpace{}%
\AgdaFunction{embedContext}\AgdaSpace{}%
\AgdaBound{Γ}\AgdaSpace{}%
\AgdaSymbol{→}\AgdaSpace{}%
\AgdaFunction{embedContext}\AgdaSpace{}%
\AgdaBound{Δ}\AgdaSpace{}%
\AgdaOperator{\AgdaDatatype{⊢b}}\AgdaSpace{}%
\AgdaBound{Ab}\AgdaSymbol{)}\<%
\\
%
\>[2]\AgdaSymbol{→}\AgdaSpace{}%
\AgdaSymbol{(}\AgdaBound{f'}\AgdaSpace{}%
\AgdaSymbol{:}\AgdaSpace{}%
\AgdaSymbol{∀}\AgdaSpace{}%
\AgdaSymbol{\{}\AgdaBound{A}\AgdaSymbol{\}}\AgdaSpace{}%
\AgdaSymbol{→}\AgdaSpace{}%
\AgdaBound{A}\AgdaSpace{}%
\AgdaOperator{\AgdaDatatype{∈}}\AgdaSpace{}%
\AgdaBound{Γ}\AgdaSpace{}%
\AgdaSymbol{→}\AgdaSpace{}%
\AgdaBound{Δ}\AgdaSpace{}%
\AgdaOperator{\AgdaDatatype{⊢}}\AgdaSpace{}%
\AgdaBound{A}\AgdaSymbol{)}\<%
\\
%
\>[2]\AgdaSymbol{→}\AgdaSpace{}%
\AgdaSymbol{(}\AgdaBound{h}\AgdaSpace{}%
\AgdaSymbol{:}\AgdaSpace{}%
\AgdaSymbol{∀}\AgdaSpace{}%
\AgdaSymbol{\{}\AgdaBound{A'}\AgdaSymbol{\}}\AgdaSpace{}%
\AgdaSymbol{→}\AgdaSpace{}%
\AgdaSymbol{(}\AgdaBound{x}\AgdaSpace{}%
\AgdaSymbol{:}\AgdaSpace{}%
\AgdaBound{A'}\AgdaSpace{}%
\AgdaOperator{\AgdaDatatype{∈}}\AgdaSpace{}%
\AgdaBound{Γ}\AgdaSymbol{)}\AgdaSpace{}%
\AgdaSymbol{→}\AgdaSpace{}%
\AgdaBound{fb'}\AgdaSpace{}%
\AgdaSymbol{(}\AgdaFunction{embedLookup}\AgdaSpace{}%
\AgdaBound{x}\AgdaSymbol{)}\AgdaSpace{}%
\AgdaOperator{\AgdaDatatype{≡}}\AgdaSpace{}%
\AgdaFunction{embedTerm}\AgdaSpace{}%
\AgdaSymbol{(}\AgdaBound{f'}\AgdaSpace{}%
\AgdaBound{x}\AgdaSymbol{)}\AgdaSpace{}%
\AgdaSymbol{)}\<%
\\
%
\>[2]\AgdaSymbol{→}\AgdaSpace{}%
\AgdaSymbol{(}\AgdaBound{M}\AgdaSpace{}%
\AgdaSymbol{:}\AgdaSpace{}%
\AgdaBound{Γ}\AgdaSpace{}%
\AgdaOperator{\AgdaDatatype{⊢}}\AgdaSpace{}%
\AgdaBound{A}\AgdaSymbol{)}\<%
\\
%
\>[2]\AgdaSymbol{→}\AgdaSpace{}%
\AgdaFunction{substb}\AgdaSpace{}%
\AgdaBound{fb'}\AgdaSpace{}%
\AgdaSymbol{(}\AgdaFunction{embedTerm}\AgdaSpace{}%
\AgdaBound{M}\AgdaSymbol{)}\AgdaSpace{}%
\AgdaOperator{\AgdaDatatype{≡}}\AgdaSpace{}%
\AgdaFunction{embedTerm}\AgdaSpace{}%
\AgdaSymbol{(}\AgdaFunction{subst}\AgdaSpace{}%
\AgdaBound{f'}\AgdaSpace{}%
\AgdaBound{M}\AgdaSymbol{)}\<%
\\
\>[0]\AgdaFunction{pres-simultaneous-subst}\AgdaSpace{}%
\AgdaBound{fb'}\AgdaSpace{}%
\AgdaBound{f'}\AgdaSpace{}%
\AgdaBound{h}\AgdaSpace{}%
\AgdaSymbol{(}\AgdaInductiveConstructor{`}\AgdaSpace{}%
\AgdaBound{x}\AgdaSymbol{)}\AgdaSpace{}%
\AgdaSymbol{=}\<%
\\
\>[0][@{}l@{\AgdaIndent{0}}]%
\>[2]\AgdaBound{h}\AgdaSpace{}%
\AgdaBound{x}\<%
\\
\>[0]\AgdaFunction{pres-simultaneous-subst}\AgdaSpace{}%
\AgdaBound{fb'}\AgdaSpace{}%
\AgdaBound{f'}\AgdaSpace{}%
\AgdaBound{h}\AgdaSpace{}%
\AgdaSymbol{(}\AgdaInductiveConstructor{ƛ}\AgdaSpace{}%
\AgdaBound{M}\AgdaSymbol{)}\AgdaSpace{}%
\AgdaSymbol{=}\<%
\\
\>[0][@{}l@{\AgdaIndent{0}}]%
\>[2]\AgdaFunction{cong}\<%
\\
\>[2][@{}l@{\AgdaIndent{0}}]%
\>[4]\AgdaInductiveConstructor{ƛb}\<%
\\
%
\>[4]\AgdaSymbol{(}\AgdaSpace{}%
\AgdaFunction{pres-simultaneous-subst}\AgdaSpace{}%
\AgdaSymbol{(}\AgdaFunction{extsb}\AgdaSpace{}%
\AgdaBound{fb'}\AgdaSymbol{)}\AgdaSpace{}%
\AgdaSymbol{(}\AgdaFunction{exts}\AgdaSpace{}%
\AgdaBound{f'}\AgdaSymbol{)}\AgdaSpace{}%
\AgdaSymbol{(}\AgdaFunction{pres-exts}\AgdaSpace{}%
\AgdaBound{fb'}\AgdaSpace{}%
\AgdaBound{f'}\AgdaSpace{}%
\AgdaBound{h}\AgdaSymbol{)}\AgdaSpace{}%
\AgdaBound{M}\AgdaSpace{}%
\AgdaSymbol{)}\<%
\\
\>[0]\AgdaFunction{pres-simultaneous-subst}\AgdaSpace{}%
\AgdaBound{fb'}\AgdaSpace{}%
\AgdaBound{f'}\AgdaSpace{}%
\AgdaBound{h}\AgdaSpace{}%
\AgdaSymbol{(}\AgdaBound{M}\AgdaSpace{}%
\AgdaOperator{\AgdaInductiveConstructor{·}}\AgdaSpace{}%
\AgdaBound{M₁}\AgdaSymbol{)}\AgdaSpace{}%
\AgdaSymbol{=}\<%
\\
\>[0][@{}l@{\AgdaIndent{0}}]%
\>[2]\AgdaFunction{cong₂}\<%
\\
\>[2][@{}l@{\AgdaIndent{0}}]%
\>[4]\AgdaOperator{\AgdaInductiveConstructor{\AgdaUnderscore{}·b\AgdaUnderscore{}}}\<%
\\
%
\>[4]\AgdaSymbol{(}\AgdaFunction{pres-simultaneous-subst}\AgdaSpace{}%
\AgdaBound{fb'}\AgdaSpace{}%
\AgdaBound{f'}\AgdaSpace{}%
\AgdaBound{h}\AgdaSpace{}%
\AgdaBound{M}\AgdaSymbol{)}\<%
\\
%
\>[4]\AgdaSymbol{(}\AgdaFunction{cong}\AgdaSpace{}%
\AgdaInductiveConstructor{box}\AgdaSpace{}%
\AgdaSymbol{(}\AgdaFunction{pres-simultaneous-subst}\AgdaSpace{}%
\AgdaBound{fb'}\AgdaSpace{}%
\AgdaBound{f'}\AgdaSpace{}%
\AgdaBound{h}\AgdaSpace{}%
\AgdaBound{M₁}\AgdaSymbol{))}\<%
\end{code}
\begin{code}[hide]%
\>[0]\<%
\\
%
\\[\AgdaEmptyExtraSkip]%
\>[0]\AgdaComment{--\ Proof\ of\ h}\<%
\\
\>[0]\AgdaFunction{pres-subst-func}\AgdaSpace{}%
\AgdaSymbol{:}\AgdaSpace{}%
\AgdaSymbol{∀}\AgdaSpace{}%
\AgdaSymbol{\{}\AgdaBound{A}\AgdaSpace{}%
\AgdaBound{C}\AgdaSpace{}%
\AgdaBound{Γ}\AgdaSymbol{\}}\<%
\\
\>[0][@{}l@{\AgdaIndent{0}}]%
\>[2]\AgdaSymbol{→}\AgdaSpace{}%
\AgdaSymbol{(}\AgdaBound{N}\AgdaSpace{}%
\AgdaSymbol{:}\AgdaSpace{}%
\AgdaBound{Γ}\AgdaSpace{}%
\AgdaOperator{\AgdaDatatype{⊢}}\AgdaSpace{}%
\AgdaBound{C}\AgdaSymbol{)}\<%
\\
%
\>[2]\AgdaSymbol{→}\AgdaSpace{}%
\AgdaSymbol{(}\AgdaBound{x}\AgdaSpace{}%
\AgdaSymbol{:}\AgdaSpace{}%
\AgdaBound{A}\AgdaSpace{}%
\AgdaOperator{\AgdaDatatype{∈}}\AgdaSpace{}%
\AgdaBound{Γ}\AgdaSpace{}%
\AgdaOperator{\AgdaInductiveConstructor{,}}\AgdaSpace{}%
\AgdaBound{C}\AgdaSymbol{)}\<%
\\
%
\>[2]\AgdaSymbol{→}\AgdaSpace{}%
\AgdaFunction{fb}\AgdaSpace{}%
\AgdaSymbol{(}\AgdaFunction{embedTerm}\AgdaSpace{}%
\AgdaBound{N}\AgdaSymbol{)}\AgdaSpace{}%
\AgdaSymbol{(}\AgdaFunction{embedLookup}\AgdaSpace{}%
\AgdaBound{x}\AgdaSymbol{)}\AgdaSpace{}%
\AgdaOperator{\AgdaDatatype{≡}}\AgdaSpace{}%
\AgdaFunction{embedTerm}\AgdaSpace{}%
\AgdaSymbol{((}\AgdaFunction{f}\AgdaSpace{}%
\AgdaBound{N}\AgdaSymbol{)}\AgdaSpace{}%
\AgdaBound{x}\AgdaSymbol{)}\<%
\\
\>[0]\AgdaFunction{pres-subst-func}\AgdaSpace{}%
\AgdaBound{N}\AgdaSpace{}%
\AgdaInductiveConstructor{Z}\AgdaSpace{}%
\AgdaSymbol{=}\AgdaSpace{}%
\AgdaInductiveConstructor{refl}\<%
\\
\>[0]\AgdaFunction{pres-subst-func}\AgdaSpace{}%
\AgdaBound{N}\AgdaSpace{}%
\AgdaSymbol{(}\AgdaInductiveConstructor{S}\AgdaSpace{}%
\AgdaBound{x}\AgdaSymbol{)}\AgdaSpace{}%
\AgdaSymbol{=}\AgdaSpace{}%
\AgdaInductiveConstructor{refl}\<%
\\
\>[0]\AgdaComment{--\ Reflection\ and\ preservation\ of\ single\ substitution\ }\<%
\\
\>[0]\AgdaComment{--\ of\ refined\ Girard's\ translation\ }\<%
\\
\>[0]\AgdaFunction{pres-subst}\AgdaSpace{}%
\AgdaSymbol{:}\AgdaSpace{}%
\AgdaSymbol{∀}\AgdaSpace{}%
\AgdaSymbol{\{}\AgdaBound{Γ}\AgdaSpace{}%
\AgdaBound{A}\AgdaSpace{}%
\AgdaBound{B}\AgdaSymbol{\}}\<%
\\
\>[0][@{}l@{\AgdaIndent{0}}]%
\>[2]\AgdaSymbol{→}\AgdaSpace{}%
\AgdaSymbol{(}\AgdaBound{M}\AgdaSpace{}%
\AgdaSymbol{:}\AgdaSpace{}%
\AgdaBound{Γ}\AgdaSpace{}%
\AgdaOperator{\AgdaInductiveConstructor{,}}\AgdaSpace{}%
\AgdaBound{B}\AgdaSpace{}%
\AgdaOperator{\AgdaDatatype{⊢}}\AgdaSpace{}%
\AgdaBound{A}\AgdaSymbol{)}\<%
\\
%
\>[2]\AgdaSymbol{→}\AgdaSpace{}%
\AgdaSymbol{(}\AgdaBound{N}\AgdaSpace{}%
\AgdaSymbol{:}\AgdaSpace{}%
\AgdaBound{Γ}\AgdaSpace{}%
\AgdaOperator{\AgdaDatatype{⊢}}\AgdaSpace{}%
\AgdaBound{B}\AgdaSymbol{)}\<%
\\
%
\>[2]\AgdaSymbol{→}\AgdaSpace{}%
\AgdaSymbol{(}\AgdaFunction{embedTerm}\AgdaSpace{}%
\AgdaBound{M}\AgdaSymbol{)}\AgdaSpace{}%
\AgdaOperator{\AgdaFunction{[}}\AgdaSpace{}%
\AgdaFunction{embedTerm}\AgdaSpace{}%
\AgdaBound{N}\AgdaSpace{}%
\AgdaOperator{\AgdaFunction{]b}}\AgdaSpace{}%
\AgdaOperator{\AgdaDatatype{≡}}\AgdaSpace{}%
\AgdaFunction{embedTerm}\AgdaSpace{}%
\AgdaSymbol{(}\AgdaBound{M}\AgdaSpace{}%
\AgdaOperator{\AgdaFunction{[}}\AgdaSpace{}%
\AgdaBound{N}\AgdaSpace{}%
\AgdaOperator{\AgdaFunction{]}}\AgdaSymbol{)}\<%
\\
\>[0]\AgdaFunction{pres-subst}\AgdaSpace{}%
\AgdaBound{M}\AgdaSpace{}%
\AgdaBound{N}\AgdaSpace{}%
\AgdaSymbol{=}\<%
\\
\>[0][@{}l@{\AgdaIndent{0}}]%
\>[2]\AgdaFunction{pres-simultaneous-subst}\AgdaSpace{}%
\AgdaSymbol{(}\AgdaFunction{fb}\AgdaSpace{}%
\AgdaSymbol{(}\AgdaFunction{embedTerm}\AgdaSpace{}%
\AgdaBound{N}\AgdaSymbol{))}\AgdaSpace{}%
\AgdaSymbol{(}\AgdaFunction{f}\AgdaSpace{}%
\AgdaBound{N}\AgdaSymbol{)}\AgdaSpace{}%
\AgdaSymbol{(}\AgdaFunction{pres-subst-func}\AgdaSpace{}%
\AgdaBound{N}\AgdaSymbol{)}\AgdaSpace{}%
\AgdaBound{M}\<%
\\
%
\\[\AgdaEmptyExtraSkip]%
\>[0]\AgdaComment{--\ Preservation\ of\ reduction\ of\ Girard's\ translation}\<%
\\
%
\\[\AgdaEmptyExtraSkip]%
\>[0]\AgdaFunction{pres-red}\AgdaSpace{}%
\AgdaSymbol{:}\AgdaSpace{}%
\AgdaSymbol{∀}\AgdaSpace{}%
\AgdaSymbol{\{}\AgdaBound{Γ}\AgdaSpace{}%
\AgdaBound{A}\AgdaSymbol{\}}\<%
\\
\>[0][@{}l@{\AgdaIndent{0}}]%
\>[2]\AgdaSymbol{→}\AgdaSpace{}%
\AgdaSymbol{(}\AgdaBound{M}\AgdaSpace{}%
\AgdaSymbol{:}\AgdaSpace{}%
\AgdaBound{Γ}\AgdaSpace{}%
\AgdaOperator{\AgdaDatatype{⊢}}\AgdaSpace{}%
\AgdaBound{A}\AgdaSymbol{)}\<%
\\
%
\>[2]\AgdaSymbol{→}\AgdaSpace{}%
\AgdaSymbol{(}\AgdaBound{N}\AgdaSpace{}%
\AgdaSymbol{:}\AgdaSpace{}%
\AgdaBound{Γ}\AgdaSpace{}%
\AgdaOperator{\AgdaDatatype{⊢}}\AgdaSpace{}%
\AgdaBound{A}\AgdaSymbol{)}\<%
\\
%
\>[2]\AgdaSymbol{→}\AgdaSpace{}%
\AgdaSymbol{(}\AgdaBound{M}\AgdaSpace{}%
\AgdaOperator{\AgdaDatatype{↝βₙ}}\AgdaSpace{}%
\AgdaBound{N}\AgdaSymbol{)}\<%
\\
%
\>[2]\AgdaSymbol{→}\AgdaSpace{}%
\AgdaSymbol{(}\AgdaFunction{embedTerm}\AgdaSpace{}%
\AgdaBound{M}\AgdaSpace{}%
\AgdaOperator{\AgdaDatatype{↝βb}}\AgdaSpace{}%
\AgdaFunction{embedTerm}\AgdaSpace{}%
\AgdaBound{N}\AgdaSymbol{)}\<%
\\
\>[0]\AgdaFunction{pres-red}\AgdaSpace{}%
\AgdaSymbol{(}\AgdaInductiveConstructor{ƛ}\AgdaSpace{}%
\AgdaBound{M₁}\AgdaSpace{}%
\AgdaOperator{\AgdaInductiveConstructor{·}}\AgdaSpace{}%
\AgdaBound{M₂}\AgdaSymbol{)}\AgdaSpace{}%
\AgdaBound{N}\AgdaSpace{}%
\AgdaInductiveConstructor{βₙ}\AgdaSpace{}%
\AgdaSymbol{=}\<%
\\
\>[0][@{}l@{\AgdaIndent{0}}]%
\>[2]\AgdaFunction{Eq.subst}\AgdaSpace{}%
\AgdaSymbol{(λ}\AgdaSpace{}%
\AgdaBound{t}\AgdaSpace{}%
\AgdaSymbol{→}\AgdaSpace{}%
\AgdaInductiveConstructor{ƛb}\AgdaSpace{}%
\AgdaSymbol{(}\AgdaFunction{embedTerm}\AgdaSpace{}%
\AgdaBound{M₁}\AgdaSymbol{)}\AgdaSpace{}%
\AgdaOperator{\AgdaInductiveConstructor{·b}}\AgdaSpace{}%
\AgdaInductiveConstructor{box}\AgdaSpace{}%
\AgdaSymbol{(}\AgdaFunction{embedTerm}\AgdaSpace{}%
\AgdaBound{M₂}\AgdaSymbol{)}\AgdaSpace{}%
\AgdaOperator{\AgdaDatatype{↝βb}}\AgdaSpace{}%
\AgdaBound{t}\AgdaSymbol{)}\<%
\\
%
\>[2]\AgdaSymbol{(}\AgdaFunction{pres-subst}\AgdaSpace{}%
\AgdaBound{M₁}\AgdaSpace{}%
\AgdaBound{M₂}\AgdaSymbol{)}\<%
\\
%
\>[2]\AgdaInductiveConstructor{βb}\<%
\\
\>[0]\AgdaFunction{pres-red}\AgdaSpace{}%
\AgdaBound{M}\AgdaSpace{}%
\AgdaBound{N}\AgdaSpace{}%
\AgdaSymbol{(}\AgdaInductiveConstructor{μ}\AgdaSpace{}%
\AgdaSymbol{\{}\AgdaArgument{M}\AgdaSpace{}%
\AgdaSymbol{=}\AgdaSpace{}%
\AgdaBound{M₁}\AgdaSymbol{\}}\AgdaSpace{}%
\AgdaSymbol{\{}\AgdaArgument{M'}\AgdaSpace{}%
\AgdaSymbol{=}\AgdaSpace{}%
\AgdaBound{M'}\AgdaSymbol{\}}\AgdaSpace{}%
\AgdaBound{red}\AgdaSymbol{)}\AgdaSpace{}%
\AgdaSymbol{=}\<%
\\
\>[0][@{}l@{\AgdaIndent{0}}]%
\>[2]\AgdaInductiveConstructor{μ}\AgdaSpace{}%
\AgdaSymbol{(}\AgdaFunction{pres-red}\AgdaSpace{}%
\AgdaBound{M₁}\AgdaSpace{}%
\AgdaBound{M'}\AgdaSpace{}%
\AgdaBound{red}\AgdaSymbol{)}\<%
\\
\>[0]\AgdaFunction{pres-red}\AgdaSpace{}%
\AgdaBound{M}\AgdaSpace{}%
\AgdaBound{N}\AgdaSpace{}%
\AgdaSymbol{(}\AgdaInductiveConstructor{ν<}\AgdaSpace{}%
\AgdaSymbol{\{}\AgdaArgument{M}\AgdaSpace{}%
\AgdaSymbol{=}\AgdaSpace{}%
\AgdaBound{M₁}\AgdaSymbol{\}}\AgdaSpace{}%
\AgdaSymbol{\{}\AgdaArgument{N}\AgdaSpace{}%
\AgdaSymbol{=}\AgdaSpace{}%
\AgdaBound{N₁}\AgdaSymbol{\}}\AgdaSpace{}%
\AgdaSymbol{\{}\AgdaArgument{N'}\AgdaSpace{}%
\AgdaSymbol{=}\AgdaSpace{}%
\AgdaBound{N'}\AgdaSymbol{\}}\AgdaSpace{}%
\AgdaBound{x}\AgdaSpace{}%
\AgdaBound{red}\AgdaSymbol{)}\AgdaSpace{}%
\AgdaSymbol{=}\<%
\\
\>[0][@{}l@{\AgdaIndent{0}}]%
\>[2]\AgdaInductiveConstructor{ν}\AgdaSpace{}%
\AgdaSymbol{(}\AgdaInductiveConstructor{ζ}\AgdaSpace{}%
\AgdaSymbol{(}\AgdaFunction{pres-red}\AgdaSpace{}%
\AgdaBound{N₁}\AgdaSpace{}%
\AgdaBound{N'}\AgdaSpace{}%
\AgdaBound{red}\AgdaSymbol{))}\<%
\\
\>[0]\AgdaFunction{pres-red}\AgdaSpace{}%
\AgdaBound{M}\AgdaSpace{}%
\AgdaBound{N}\AgdaSpace{}%
\AgdaSymbol{(}\AgdaInductiveConstructor{ν}\AgdaSpace{}%
\AgdaSymbol{\{}\AgdaArgument{M}\AgdaSpace{}%
\AgdaSymbol{=}\AgdaSpace{}%
\AgdaBound{M₁}\AgdaSymbol{\}}\AgdaSpace{}%
\AgdaSymbol{\{}\AgdaArgument{N}\AgdaSpace{}%
\AgdaSymbol{=}\AgdaSpace{}%
\AgdaBound{N₁}\AgdaSymbol{\}}\AgdaSpace{}%
\AgdaSymbol{\{}\AgdaArgument{N'}\AgdaSpace{}%
\AgdaSymbol{=}\AgdaSpace{}%
\AgdaBound{N'}\AgdaSymbol{\}}\AgdaSpace{}%
\AgdaBound{red}\AgdaSymbol{)}\AgdaSpace{}%
\AgdaSymbol{=}\<%
\\
\>[0][@{}l@{\AgdaIndent{0}}]%
\>[2]\AgdaInductiveConstructor{ν}\AgdaSpace{}%
\AgdaSymbol{(}\AgdaInductiveConstructor{ζ}\AgdaSpace{}%
\AgdaSymbol{(}\AgdaFunction{pres-red}\AgdaSpace{}%
\AgdaBound{N₁}\AgdaSpace{}%
\AgdaBound{N'}\AgdaSpace{}%
\AgdaBound{red}\AgdaSymbol{))}\<%
\\
\>[0]\AgdaFunction{pres-red}\AgdaSpace{}%
\AgdaBound{M}\AgdaSpace{}%
\AgdaBound{N}\AgdaSpace{}%
\AgdaSymbol{(}\AgdaInductiveConstructor{ξ}\AgdaSpace{}%
\AgdaSymbol{\{}\AgdaArgument{M}\AgdaSpace{}%
\AgdaSymbol{=}\AgdaSpace{}%
\AgdaBound{M₁}\AgdaSymbol{\}}\AgdaSpace{}%
\AgdaSymbol{\{}\AgdaArgument{M'}\AgdaSpace{}%
\AgdaSymbol{=}\AgdaSpace{}%
\AgdaBound{M'}\AgdaSymbol{\}}\AgdaSpace{}%
\AgdaBound{red}\AgdaSymbol{)}\AgdaSpace{}%
\AgdaSymbol{=}\<%
\\
\>[0][@{}l@{\AgdaIndent{0}}]%
\>[2]\AgdaInductiveConstructor{ξ}\AgdaSpace{}%
\AgdaSymbol{(}\AgdaFunction{pres-red}\AgdaSpace{}%
\AgdaBound{M₁}\AgdaSpace{}%
\AgdaBound{M'}\AgdaSpace{}%
\AgdaBound{red}\AgdaSymbol{)}\<%
\\
%
\\[\AgdaEmptyExtraSkip]%
\>[0]\AgdaFunction{pres-eval}\AgdaSpace{}%
\AgdaSymbol{:}\AgdaSpace{}%
\AgdaSymbol{∀}\AgdaSpace{}%
\AgdaSymbol{\{}\AgdaBound{Γ}\AgdaSpace{}%
\AgdaBound{A}\AgdaSymbol{\}}\<%
\\
\>[0][@{}l@{\AgdaIndent{0}}]%
\>[2]\AgdaSymbol{→}\AgdaSpace{}%
\AgdaSymbol{(}\AgdaBound{M}\AgdaSpace{}%
\AgdaSymbol{:}\AgdaSpace{}%
\AgdaBound{Γ}\AgdaSpace{}%
\AgdaOperator{\AgdaDatatype{⊢}}\AgdaSpace{}%
\AgdaBound{A}\AgdaSymbol{)}\<%
\\
%
\>[2]\AgdaSymbol{→}\AgdaSpace{}%
\AgdaSymbol{(}\AgdaBound{N}\AgdaSpace{}%
\AgdaSymbol{:}\AgdaSpace{}%
\AgdaBound{Γ}\AgdaSpace{}%
\AgdaOperator{\AgdaDatatype{⊢}}\AgdaSpace{}%
\AgdaBound{A}\AgdaSymbol{)}\<%
\\
%
\>[2]\AgdaSymbol{→}\AgdaSpace{}%
\AgdaSymbol{(}\AgdaBound{M}\AgdaSpace{}%
\AgdaOperator{\AgdaDatatype{↝ₙ}}\AgdaSpace{}%
\AgdaBound{N}\AgdaSymbol{)}\<%
\\
%
\>[2]\AgdaSymbol{→}\AgdaSpace{}%
\AgdaSymbol{(}\AgdaFunction{embedTerm}\AgdaSpace{}%
\AgdaBound{M}\AgdaSpace{}%
\AgdaOperator{\AgdaDatatype{↝bₙ}}\AgdaSpace{}%
\AgdaFunction{embedTerm}\AgdaSpace{}%
\AgdaBound{N}\AgdaSymbol{)}\<%
\\
\>[0]\AgdaFunction{pres-eval}\AgdaSpace{}%
\AgdaSymbol{(}\AgdaInductiveConstructor{ƛ}\AgdaSpace{}%
\AgdaBound{M₁}\AgdaSpace{}%
\AgdaOperator{\AgdaInductiveConstructor{·}}\AgdaSpace{}%
\AgdaBound{M₂}\AgdaSymbol{)}\AgdaSpace{}%
\AgdaBound{N}\AgdaSpace{}%
\AgdaInductiveConstructor{βₙ}\AgdaSpace{}%
\AgdaSymbol{=}\<%
\\
\>[0][@{}l@{\AgdaIndent{0}}]%
\>[2]\AgdaFunction{Eq.subst}\<%
\\
\>[2][@{}l@{\AgdaIndent{0}}]%
\>[4]\AgdaSymbol{(λ}\AgdaSpace{}%
\AgdaBound{t}\AgdaSpace{}%
\AgdaSymbol{→}\AgdaSpace{}%
\AgdaInductiveConstructor{ƛb}\AgdaSpace{}%
\AgdaSymbol{(}\AgdaFunction{embedTerm}\AgdaSpace{}%
\AgdaBound{M₁}\AgdaSymbol{)}\AgdaSpace{}%
\AgdaOperator{\AgdaInductiveConstructor{·b}}\AgdaSpace{}%
\AgdaInductiveConstructor{box}\AgdaSpace{}%
\AgdaSymbol{(}\AgdaFunction{embedTerm}\AgdaSpace{}%
\AgdaBound{M₂}\AgdaSymbol{)}\AgdaSpace{}%
\AgdaOperator{\AgdaDatatype{↝bₙ}}\AgdaSpace{}%
\AgdaBound{t}\AgdaSymbol{)}\<%
\\
%
\>[4]\AgdaSymbol{(}\AgdaFunction{pres-subst}\AgdaSpace{}%
\AgdaBound{M₁}\AgdaSpace{}%
\AgdaBound{M₂}\AgdaSymbol{)}\<%
\\
%
\>[4]\AgdaInductiveConstructor{βb}\<%
\\
\>[0]\AgdaFunction{pres-eval}\AgdaSpace{}%
\AgdaBound{M}\AgdaSpace{}%
\AgdaBound{N}\AgdaSpace{}%
\AgdaSymbol{(}\AgdaInductiveConstructor{μ}\AgdaSpace{}%
\AgdaSymbol{\{}\AgdaArgument{M}\AgdaSpace{}%
\AgdaSymbol{=}\AgdaSpace{}%
\AgdaBound{M₁}\AgdaSymbol{\}}\AgdaSpace{}%
\AgdaSymbol{\{}\AgdaArgument{M'}\AgdaSpace{}%
\AgdaSymbol{=}\AgdaSpace{}%
\AgdaBound{M'}\AgdaSymbol{\}}\AgdaSpace{}%
\AgdaBound{red}\AgdaSymbol{)}\AgdaSpace{}%
\AgdaSymbol{=}\<%
\\
\>[0][@{}l@{\AgdaIndent{0}}]%
\>[2]\AgdaInductiveConstructor{μ}\AgdaSpace{}%
\AgdaSymbol{(}\AgdaFunction{pres-eval}\AgdaSpace{}%
\AgdaBound{M₁}\AgdaSpace{}%
\AgdaBound{M'}\AgdaSpace{}%
\AgdaBound{red}\AgdaSymbol{)}\<%
\end{code}
\end{lemma}

Some other interesting helper lemmas are the congruence, compose and commute lemmas, which were needed in the poofs of Gödel's embedding.
This refined embedding makes use of the abbrevation \raiseembn, which is formally defined in \cref{def:raise:agda} with the use of \textsf{renameb}.
This use of \textsf{renameb} makes the proofs a lot more complex and requires the need for these helper lemmas. 
We only give congruence, compose and commute helper lemma's for simultaneous substitution because these lemmas are similar for \textsf{extsb}, \textsf{renameb} and \textsf{extb}.
\Cref{lem:subst:commute} uses helper lemmas \textsf{subst-renameb-comp} and substb-rename-commute, of which \textsf{subst-renameb-comp} is given in \cref{lemm:subst:comp}.
The lemma \textsf{substb-cong} is used in \cref{lem:subst:commute} and given in \cref{lem:cong:subst}.

\begin{lemma}
  \label{lem:subst:commute}
  \textbf{Commutativity} of renaming followed by substitution under an extended context and substitution followed by renaming:
  \small
  \begin{code}[hide]%
\>[0]\AgdaKeyword{module}\AgdaSpace{}%
\AgdaModule{valueProofs1}\AgdaSpace{}%
\AgdaKeyword{where}\<%
\\
\>[0]\<%
\\
\>[0]\AgdaKeyword{open}\AgdaSpace{}%
\AgdaKeyword{import}\AgdaSpace{}%
\AgdaModule{calc}\<%
\\
\>[0]\AgdaKeyword{open}\AgdaSpace{}%
\AgdaKeyword{import}\AgdaSpace{}%
\AgdaModule{boxCalc}\<%
\\
\>[0]\AgdaKeyword{open}\AgdaSpace{}%
\AgdaKeyword{import}\AgdaSpace{}%
\AgdaModule{embedCBVIntoCBB}\<%
\\
\>[0]\AgdaKeyword{open}\AgdaSpace{}%
\AgdaKeyword{import}\AgdaSpace{}%
\AgdaModule{valueCalcRed}\<%
\\
\>[0]\AgdaKeyword{open}\AgdaSpace{}%
\AgdaKeyword{import}\AgdaSpace{}%
\AgdaModule{boxCalcRed}\<%
\\
%
\\[\AgdaEmptyExtraSkip]%
\>[0]\AgdaKeyword{open}\AgdaSpace{}%
\AgdaKeyword{import}\AgdaSpace{}%
\AgdaModule{Data.Empty}\AgdaSpace{}%
\AgdaKeyword{using}\AgdaSpace{}%
\AgdaSymbol{(}\AgdaFunction{⊥}\AgdaSymbol{;}\AgdaSpace{}%
\AgdaFunction{⊥-elim}\AgdaSymbol{)}\<%
\\
\>[0]\AgdaKeyword{import}\AgdaSpace{}%
\AgdaModule{Relation.Binary.PropositionalEquality}\AgdaSpace{}%
\AgdaSymbol{as}\AgdaSpace{}%
\AgdaModule{Eq}\<%
\\
\>[0]\AgdaKeyword{open}\AgdaSpace{}%
\AgdaModule{Eq}\AgdaSpace{}%
\AgdaKeyword{using}\AgdaSpace{}%
\AgdaSymbol{(}\AgdaOperator{\AgdaDatatype{\AgdaUnderscore{}≡\AgdaUnderscore{}}}\AgdaSymbol{;}\AgdaSpace{}%
\AgdaInductiveConstructor{refl}\AgdaSymbol{;}\AgdaSpace{}%
\AgdaFunction{cong}\AgdaSymbol{;}\AgdaSpace{}%
\AgdaFunction{cong₂}\AgdaSymbol{;}\AgdaSpace{}%
\AgdaFunction{sym}\AgdaSymbol{)}\<%
\\
%
\\[\AgdaEmptyExtraSkip]%
\>[0]\AgdaKeyword{open}\AgdaSpace{}%
\AgdaModule{Eq.≡-Reasoning}\AgdaSpace{}%
\AgdaKeyword{using}\AgdaSpace{}%
\AgdaSymbol{(}\AgdaOperator{\AgdaFunction{begin\AgdaUnderscore{}}}\AgdaSymbol{;}\AgdaSpace{}%
\AgdaFunction{step-≡-∣}\AgdaSymbol{;}\AgdaSpace{}%
\AgdaFunction{step-≡-⟩}\AgdaSymbol{;}\AgdaSpace{}%
\AgdaOperator{\AgdaFunction{\AgdaUnderscore{}∎}}\AgdaSymbol{)}\<%
\\
\>[0]\AgdaKeyword{open}\AgdaSpace{}%
\AgdaModule{Eq.≡-Reasoning}\AgdaSpace{}%
\AgdaKeyword{renaming}\AgdaSpace{}%
\AgdaSymbol{(}\AgdaOperator{\AgdaFunction{begin\AgdaUnderscore{}}}\AgdaSpace{}%
\AgdaSymbol{to}\AgdaSpace{}%
\AgdaOperator{\AgdaFunction{beginEq\AgdaUnderscore{}}}\AgdaSymbol{;}\AgdaSpace{}%
\AgdaOperator{\AgdaFunction{\AgdaUnderscore{}∎}}\AgdaSpace{}%
\AgdaSymbol{to}\AgdaSpace{}%
\AgdaOperator{\AgdaFunction{\AgdaUnderscore{}∎eq}}\AgdaSymbol{)}\<%
\\
%
\\[\AgdaEmptyExtraSkip]%
\>[0]\AgdaComment{--\ Preservation\ of\ typing\ by\ Gödel's\ translation}\<%
\\
\>[0]\AgdaFunction{pres-typing}\AgdaSpace{}%
\AgdaSymbol{:}\AgdaSpace{}%
\AgdaSymbol{∀}\AgdaSpace{}%
\AgdaSymbol{\{}\AgdaBound{Γ}\AgdaSpace{}%
\AgdaBound{A}\AgdaSymbol{\}}\AgdaSpace{}%
\AgdaSymbol{→}\AgdaSpace{}%
\AgdaBound{Γ}\AgdaSpace{}%
\AgdaOperator{\AgdaDatatype{⊢}}\AgdaSpace{}%
\AgdaBound{A}\AgdaSpace{}%
\AgdaSymbol{→}\AgdaSpace{}%
\AgdaFunction{embedContext}\AgdaSpace{}%
\AgdaBound{Γ}\AgdaSpace{}%
\AgdaOperator{\AgdaDatatype{⊢b}}\AgdaSpace{}%
\AgdaFunction{embedType}\AgdaSpace{}%
\AgdaBound{A}\<%
\\
\>[0]\AgdaFunction{pres-typing}\AgdaSpace{}%
\AgdaBound{x}\AgdaSpace{}%
\AgdaSymbol{=}\AgdaSpace{}%
\AgdaFunction{embedTerm}\AgdaSpace{}%
\AgdaBound{x}\<%
\\
%
\\[\AgdaEmptyExtraSkip]%
\>[0]\AgdaComment{--\ Preservation\ of\ typing\ of\ Values\ by\ Gödel's\ translation}\<%
\\
\>[0]\AgdaFunction{pres-typing-val}\AgdaSpace{}%
\AgdaSymbol{:}\AgdaSpace{}%
\AgdaSymbol{∀}\AgdaSpace{}%
\AgdaSymbol{\{}\AgdaBound{Γ}\AgdaSpace{}%
\AgdaBound{A}\AgdaSymbol{\}}\AgdaSpace{}%
\AgdaSymbol{→}\AgdaSpace{}%
\AgdaSymbol{(}\AgdaBound{V}\AgdaSpace{}%
\AgdaSymbol{:}\AgdaSpace{}%
\AgdaBound{Γ}\AgdaSpace{}%
\AgdaOperator{\AgdaDatatype{⊢}}\AgdaSpace{}%
\AgdaBound{A}\AgdaSymbol{)}\AgdaSpace{}%
\AgdaSymbol{→}\AgdaSpace{}%
\AgdaSymbol{(}\AgdaDatatype{Value}\AgdaSpace{}%
\AgdaBound{V}\AgdaSymbol{)}\AgdaSpace{}%
\AgdaSymbol{→}\AgdaSpace{}%
\AgdaFunction{embedContext}\AgdaSpace{}%
\AgdaBound{Γ}\AgdaSpace{}%
\AgdaOperator{\AgdaDatatype{⊢b}}\AgdaSpace{}%
\AgdaFunction{embedTypeValue}\AgdaSpace{}%
\AgdaBound{A}\<%
\\
\>[0]\AgdaFunction{pres-typing-val}\AgdaSpace{}%
\AgdaBound{V}\AgdaSpace{}%
\AgdaBound{x}\AgdaSpace{}%
\AgdaSymbol{=}\AgdaSpace{}%
\AgdaFunction{embedValue}\AgdaSpace{}%
\AgdaBound{V}\AgdaSpace{}%
\AgdaBound{x}\<%
\\
%
\\[\AgdaEmptyExtraSkip]%
\>[0]\AgdaComment{--\ Helper\ lemma:\ applying\ embedTerm\ to\ a\ Value\ is\ equivalent\ to\ wrapping\ the\ result\ of\ embedValue\ in\ a\ box}\<%
\\
\>[0]\AgdaFunction{box-embedValue}\AgdaSpace{}%
\AgdaSymbol{:}\AgdaSpace{}%
\AgdaSymbol{∀}\AgdaSpace{}%
\AgdaSymbol{\{}\AgdaBound{Γ}\AgdaSpace{}%
\AgdaBound{A}\AgdaSymbol{\}}\AgdaSpace{}%
\AgdaSymbol{→}\AgdaSpace{}%
\AgdaSymbol{\{}\AgdaBound{V}\AgdaSpace{}%
\AgdaSymbol{:}\AgdaSpace{}%
\AgdaBound{Γ}\AgdaSpace{}%
\AgdaOperator{\AgdaDatatype{⊢}}\AgdaSpace{}%
\AgdaBound{A}\AgdaSymbol{\}}\AgdaSpace{}%
\AgdaSymbol{→}\AgdaSpace{}%
\AgdaSymbol{(}\AgdaBound{x}\AgdaSpace{}%
\AgdaSymbol{:}\AgdaSpace{}%
\AgdaDatatype{Value}\AgdaSpace{}%
\AgdaBound{V}\AgdaSymbol{)}\AgdaSpace{}%
\AgdaSymbol{→}\AgdaSpace{}%
\AgdaInductiveConstructor{box}\AgdaSpace{}%
\AgdaSymbol{(}\AgdaFunction{embedValue}\AgdaSpace{}%
\AgdaBound{V}\AgdaSpace{}%
\AgdaBound{x}\AgdaSymbol{)}\AgdaSpace{}%
\AgdaOperator{\AgdaDatatype{≡}}\AgdaSpace{}%
\AgdaFunction{embedTerm}\AgdaSpace{}%
\AgdaBound{V}\<%
\\
\>[0]\AgdaFunction{box-embedValue}\AgdaSpace{}%
\AgdaInductiveConstructor{V-ƛ}\AgdaSpace{}%
\AgdaSymbol{=}\AgdaSpace{}%
\AgdaInductiveConstructor{refl}\<%
\\
\>[0]\AgdaFunction{box-embedValue}\AgdaSpace{}%
\AgdaInductiveConstructor{V-`}\AgdaSpace{}%
\AgdaSymbol{=}\AgdaSpace{}%
\AgdaInductiveConstructor{refl}\<%
\\
%
\\[\AgdaEmptyExtraSkip]%
\>[0]\AgdaComment{--\ HELPER\ LEMMA'S\ SIMULTANIOUS\ SUBSTITUTION}\<%
\\
%
\\[\AgdaEmptyExtraSkip]%
\>[0]\AgdaComment{--\ Preservation\ and\ Reflection\ of\ }\<%
\\
\>[0]\AgdaComment{--\ ext\ of\ Gödel's\ translation}\<%
\\
\>[0]\AgdaFunction{pres-ext}\AgdaSpace{}%
\AgdaSymbol{:}\AgdaSpace{}%
\AgdaSymbol{∀}\AgdaSpace{}%
\AgdaSymbol{\{}\AgdaBound{Γ}\AgdaSpace{}%
\AgdaBound{Δ}\AgdaSpace{}%
\AgdaBound{A'}\AgdaSpace{}%
\AgdaBound{A}\AgdaSymbol{\}}\<%
\\
\>[0][@{}l@{\AgdaIndent{0}}]%
\>[2]\AgdaSymbol{→}\AgdaSpace{}%
\AgdaSymbol{(}\AgdaBound{f'}\AgdaSpace{}%
\AgdaSymbol{:}\AgdaSpace{}%
\AgdaSymbol{∀}\AgdaSpace{}%
\AgdaSymbol{\{}\AgdaBound{A}\AgdaSymbol{\}}\AgdaSpace{}%
\AgdaSymbol{→}\AgdaSpace{}%
\AgdaBound{A}\AgdaSpace{}%
\AgdaOperator{\AgdaDatatype{∈}}\AgdaSpace{}%
\AgdaBound{Γ}\AgdaSpace{}%
\AgdaSymbol{→}\AgdaSpace{}%
\AgdaBound{A}\AgdaSpace{}%
\AgdaOperator{\AgdaDatatype{∈}}\AgdaSpace{}%
\AgdaBound{Δ}\AgdaSymbol{)}\<%
\\
%
\>[2]\AgdaSymbol{→}\AgdaSpace{}%
\AgdaSymbol{(}\AgdaBound{fb'}\AgdaSpace{}%
\AgdaSymbol{:}\AgdaSpace{}%
\AgdaSymbol{∀}\AgdaSpace{}%
\AgdaSymbol{\{}\AgdaBound{Ab}\AgdaSymbol{\}}\AgdaSpace{}%
\AgdaSymbol{→}\AgdaSpace{}%
\AgdaBound{Ab}\AgdaSpace{}%
\AgdaOperator{\AgdaDatatype{∈b}}\AgdaSpace{}%
\AgdaFunction{embedContext}\AgdaSpace{}%
\AgdaBound{Γ}\AgdaSpace{}%
\AgdaSymbol{→}\AgdaSpace{}%
\AgdaBound{Ab}\AgdaSpace{}%
\AgdaOperator{\AgdaDatatype{∈b}}\AgdaSpace{}%
\AgdaFunction{embedContext}\AgdaSpace{}%
\AgdaBound{Δ}\AgdaSymbol{)}\<%
\\
%
\>[2]\AgdaSymbol{→}\AgdaSpace{}%
\AgdaSymbol{(\{}\AgdaBound{A'}\AgdaSpace{}%
\AgdaSymbol{:}\AgdaSpace{}%
\AgdaDatatype{Type}\AgdaSymbol{\}}\AgdaSpace{}%
\AgdaSymbol{(}\AgdaBound{x}\AgdaSpace{}%
\AgdaSymbol{:}\AgdaSpace{}%
\AgdaBound{A'}\AgdaSpace{}%
\AgdaOperator{\AgdaDatatype{∈}}\AgdaSpace{}%
\AgdaBound{Γ}\AgdaSymbol{)}\AgdaSpace{}%
\AgdaSymbol{→}\AgdaSpace{}%
\AgdaBound{fb'}\AgdaSpace{}%
\AgdaSymbol{(}\AgdaFunction{embedLookup}\AgdaSpace{}%
\AgdaBound{x}\AgdaSymbol{)}\AgdaSpace{}%
\AgdaOperator{\AgdaDatatype{≡}}\AgdaSpace{}%
\AgdaFunction{embedLookup}\AgdaSpace{}%
\AgdaSymbol{(}\AgdaBound{f'}\AgdaSpace{}%
\AgdaBound{x}\AgdaSymbol{))}\<%
\\
%
\>[2]\AgdaSymbol{→}\AgdaSpace{}%
\AgdaSymbol{(}\AgdaBound{x}\AgdaSpace{}%
\AgdaSymbol{:}\AgdaSpace{}%
\AgdaBound{A'}\AgdaSpace{}%
\AgdaOperator{\AgdaDatatype{∈}}\AgdaSpace{}%
\AgdaBound{Γ}\AgdaSpace{}%
\AgdaOperator{\AgdaInductiveConstructor{,}}\AgdaSpace{}%
\AgdaBound{A}\AgdaSymbol{)}\<%
\\
%
\>[2]\AgdaSymbol{→}\AgdaSpace{}%
\AgdaFunction{extb}\AgdaSpace{}%
\AgdaBound{fb'}\AgdaSpace{}%
\AgdaSymbol{(}\AgdaFunction{embedLookup}\AgdaSpace{}%
\AgdaBound{x}\AgdaSymbol{)}\AgdaSpace{}%
\AgdaOperator{\AgdaDatatype{≡}}\AgdaSpace{}%
\AgdaFunction{embedLookup}\AgdaSpace{}%
\AgdaSymbol{(}\AgdaFunction{ext}\AgdaSpace{}%
\AgdaBound{f'}\AgdaSpace{}%
\AgdaBound{x}\AgdaSymbol{)}\<%
\\
\>[0]\AgdaFunction{pres-ext}\AgdaSpace{}%
\AgdaBound{f'}\AgdaSpace{}%
\AgdaBound{fb'}\AgdaSpace{}%
\AgdaBound{h}\AgdaSpace{}%
\AgdaInductiveConstructor{Z}\AgdaSpace{}%
\AgdaSymbol{=}\AgdaSpace{}%
\AgdaInductiveConstructor{refl}\<%
\\
\>[0]\AgdaFunction{pres-ext}\AgdaSpace{}%
\AgdaBound{f'}\AgdaSpace{}%
\AgdaBound{fb'}\AgdaSpace{}%
\AgdaBound{h}\AgdaSpace{}%
\AgdaSymbol{(}\AgdaInductiveConstructor{S}\AgdaSpace{}%
\AgdaBound{x}\AgdaSymbol{)}\AgdaSpace{}%
\AgdaSymbol{=}\AgdaSpace{}%
\AgdaFunction{cong}\AgdaSpace{}%
\AgdaInductiveConstructor{Sb}\AgdaSpace{}%
\AgdaSymbol{(}\AgdaBound{h}\AgdaSpace{}%
\AgdaBound{x}\AgdaSymbol{)}\<%
\\
%
\\[\AgdaEmptyExtraSkip]%
\>[0]\AgdaComment{--\ Extb\ preserves\ pointwise\ equality\ }\<%
\\
\>[0]\AgdaComment{--\ of\ renaming\ functions}\<%
\\
\>[0]\AgdaFunction{extb-cong}\AgdaSpace{}%
\AgdaSymbol{:}\AgdaSpace{}%
\AgdaSymbol{∀}\AgdaSpace{}%
\AgdaSymbol{\{}\AgdaBound{Γ₁}\AgdaSpace{}%
\AgdaBound{Γ₂}\AgdaSpace{}%
\AgdaBound{A}\AgdaSpace{}%
\AgdaBound{A'}\AgdaSymbol{\}}\<%
\\
\>[0][@{}l@{\AgdaIndent{0}}]%
\>[2]\AgdaSymbol{→}\AgdaSpace{}%
\AgdaSymbol{\{}\AgdaBound{fb'}\AgdaSpace{}%
\AgdaSymbol{:}\AgdaSpace{}%
\AgdaSymbol{∀}\AgdaSpace{}%
\AgdaSymbol{\{}\AgdaBound{Ab}\AgdaSymbol{\}}\AgdaSpace{}%
\AgdaSymbol{→}\AgdaSpace{}%
\AgdaBound{Ab}\AgdaSpace{}%
\AgdaOperator{\AgdaDatatype{∈b}}\AgdaSpace{}%
\AgdaBound{Γ₁}\AgdaSpace{}%
\AgdaSymbol{→}\AgdaSpace{}%
\AgdaBound{Ab}\AgdaSpace{}%
\AgdaOperator{\AgdaDatatype{∈b}}\AgdaSpace{}%
\AgdaBound{Γ₂}\AgdaSymbol{\}}\<%
\\
%
\>[2]\AgdaSymbol{→}\AgdaSpace{}%
\AgdaSymbol{\{}\AgdaBound{fb''}\AgdaSpace{}%
\AgdaSymbol{:}\AgdaSpace{}%
\AgdaSymbol{∀}\AgdaSpace{}%
\AgdaSymbol{\{}\AgdaBound{Ab}\AgdaSymbol{\}}\AgdaSpace{}%
\AgdaSymbol{→}\AgdaSpace{}%
\AgdaBound{Ab}\AgdaSpace{}%
\AgdaOperator{\AgdaDatatype{∈b}}\AgdaSpace{}%
\AgdaBound{Γ₁}\AgdaSpace{}%
\AgdaSymbol{→}\AgdaSpace{}%
\AgdaBound{Ab}\AgdaSpace{}%
\AgdaOperator{\AgdaDatatype{∈b}}\AgdaSpace{}%
\AgdaBound{Γ₂}\AgdaSymbol{\}}\<%
\\
%
\>[2]\AgdaSymbol{→}\AgdaSpace{}%
\AgdaSymbol{(}\AgdaBound{h}\AgdaSpace{}%
\AgdaSymbol{:}\AgdaSpace{}%
\AgdaSymbol{∀}\AgdaSpace{}%
\AgdaSymbol{\{}\AgdaBound{Ab}\AgdaSymbol{\}}\AgdaSpace{}%
\AgdaSymbol{→}\AgdaSpace{}%
\AgdaSymbol{(}\AgdaBound{x}\AgdaSpace{}%
\AgdaSymbol{:}\AgdaSpace{}%
\AgdaBound{Ab}\AgdaSpace{}%
\AgdaOperator{\AgdaDatatype{∈b}}\AgdaSpace{}%
\AgdaBound{Γ₁}\AgdaSymbol{)}\AgdaSpace{}%
\AgdaSymbol{→}\AgdaSpace{}%
\AgdaBound{fb'}\AgdaSpace{}%
\AgdaBound{x}\AgdaSpace{}%
\AgdaOperator{\AgdaDatatype{≡}}\AgdaSpace{}%
\AgdaBound{fb''}\AgdaSpace{}%
\AgdaBound{x}\AgdaSymbol{)}\<%
\\
%
\>[2]\AgdaSymbol{→}\AgdaSpace{}%
\AgdaSymbol{(}\AgdaBound{x}\AgdaSpace{}%
\AgdaSymbol{:}\AgdaSpace{}%
\AgdaBound{A}\AgdaSpace{}%
\AgdaOperator{\AgdaDatatype{∈b}}\AgdaSpace{}%
\AgdaBound{Γ₁}\AgdaSpace{}%
\AgdaOperator{\AgdaInductiveConstructor{,b}}\AgdaSpace{}%
\AgdaInductiveConstructor{□}\AgdaSpace{}%
\AgdaBound{A'}\AgdaSymbol{)}\<%
\\
%
\>[2]\AgdaSymbol{→}\AgdaSpace{}%
\AgdaFunction{extb}\AgdaSpace{}%
\AgdaBound{fb'}\AgdaSpace{}%
\AgdaBound{x}\AgdaSpace{}%
\AgdaOperator{\AgdaDatatype{≡}}\AgdaSpace{}%
\AgdaFunction{extb}\AgdaSpace{}%
\AgdaBound{fb''}\AgdaSpace{}%
\AgdaBound{x}\<%
\\
\>[0]\AgdaFunction{extb-cong}\AgdaSpace{}%
\AgdaBound{h}\AgdaSpace{}%
\AgdaInductiveConstructor{Zb}\AgdaSpace{}%
\AgdaSymbol{=}\AgdaSpace{}%
\AgdaInductiveConstructor{refl}\<%
\\
\>[0]\AgdaFunction{extb-cong}\AgdaSpace{}%
\AgdaBound{h}\AgdaSpace{}%
\AgdaSymbol{(}\AgdaInductiveConstructor{Sb}\AgdaSpace{}%
\AgdaBound{x}\AgdaSymbol{)}\AgdaSpace{}%
\AgdaSymbol{=}\AgdaSpace{}%
\AgdaFunction{cong}\AgdaSpace{}%
\AgdaInductiveConstructor{Sb}\AgdaSpace{}%
\AgdaSymbol{(}\AgdaBound{h}\AgdaSpace{}%
\AgdaBound{x}\AgdaSymbol{)}\<%
\\
%
\\[\AgdaEmptyExtraSkip]%
\>[0]\AgdaComment{--\ Renameb\ preserves\ pointwise\ equality}\<%
\\
\>[0]\AgdaComment{--\ of\ renaming\ functions}\<%
\\
\>[0]\AgdaFunction{renameb-cong}\AgdaSpace{}%
\AgdaSymbol{:}\AgdaSpace{}%
\AgdaSymbol{∀}\AgdaSpace{}%
\AgdaSymbol{\{}\AgdaBound{Γ₁}\AgdaSpace{}%
\AgdaBound{Γ₂}\AgdaSpace{}%
\AgdaBound{A}\AgdaSymbol{\}}\<%
\\
\>[0][@{}l@{\AgdaIndent{0}}]%
\>[2]\AgdaSymbol{→}\AgdaSpace{}%
\AgdaSymbol{\{}\AgdaBound{fb'}\AgdaSpace{}%
\AgdaSymbol{:}\AgdaSpace{}%
\AgdaSymbol{∀}\AgdaSpace{}%
\AgdaSymbol{\{}\AgdaBound{Ab}\AgdaSymbol{\}}\AgdaSpace{}%
\AgdaSymbol{→}\AgdaSpace{}%
\AgdaBound{Ab}\AgdaSpace{}%
\AgdaOperator{\AgdaDatatype{∈b}}\AgdaSpace{}%
\AgdaBound{Γ₁}\AgdaSpace{}%
\AgdaSymbol{→}\AgdaSpace{}%
\AgdaBound{Ab}\AgdaSpace{}%
\AgdaOperator{\AgdaDatatype{∈b}}\AgdaSpace{}%
\AgdaBound{Γ₂}\AgdaSymbol{\}}\<%
\\
%
\>[2]\AgdaSymbol{→}\AgdaSpace{}%
\AgdaSymbol{\{}\AgdaBound{fb''}\AgdaSpace{}%
\AgdaSymbol{:}\AgdaSpace{}%
\AgdaSymbol{∀}\AgdaSpace{}%
\AgdaSymbol{\{}\AgdaBound{Ab}\AgdaSymbol{\}}\AgdaSpace{}%
\AgdaSymbol{→}\AgdaSpace{}%
\AgdaBound{Ab}\AgdaSpace{}%
\AgdaOperator{\AgdaDatatype{∈b}}\AgdaSpace{}%
\AgdaBound{Γ₁}\AgdaSpace{}%
\AgdaSymbol{→}\AgdaSpace{}%
\AgdaBound{Ab}\AgdaSpace{}%
\AgdaOperator{\AgdaDatatype{∈b}}\AgdaSpace{}%
\AgdaBound{Γ₂}\AgdaSymbol{\}}\<%
\\
%
\>[2]\AgdaSymbol{→}\AgdaSpace{}%
\AgdaSymbol{(}\AgdaBound{h}\AgdaSpace{}%
\AgdaSymbol{:}\AgdaSpace{}%
\AgdaSymbol{∀}\AgdaSpace{}%
\AgdaSymbol{\{}\AgdaBound{Ab}\AgdaSymbol{\}}\AgdaSpace{}%
\AgdaSymbol{→}\AgdaSpace{}%
\AgdaSymbol{(}\AgdaBound{x}\AgdaSpace{}%
\AgdaSymbol{:}\AgdaSpace{}%
\AgdaBound{Ab}\AgdaSpace{}%
\AgdaOperator{\AgdaDatatype{∈b}}\AgdaSpace{}%
\AgdaBound{Γ₁}\AgdaSymbol{)}\AgdaSpace{}%
\AgdaSymbol{→}\AgdaSpace{}%
\AgdaBound{fb'}\AgdaSpace{}%
\AgdaBound{x}\AgdaSpace{}%
\AgdaOperator{\AgdaDatatype{≡}}\AgdaSpace{}%
\AgdaBound{fb''}\AgdaSpace{}%
\AgdaBound{x}\AgdaSymbol{)}\<%
\\
%
\>[2]\AgdaSymbol{→}\AgdaSpace{}%
\AgdaSymbol{(}\AgdaBound{t}\AgdaSpace{}%
\AgdaSymbol{:}\AgdaSpace{}%
\AgdaBound{Γ₁}\AgdaSpace{}%
\AgdaOperator{\AgdaDatatype{⊢b}}\AgdaSpace{}%
\AgdaBound{A}\AgdaSymbol{)}\<%
\\
%
\>[2]\AgdaSymbol{→}\AgdaSpace{}%
\AgdaFunction{renameb}\AgdaSpace{}%
\AgdaBound{fb'}\AgdaSpace{}%
\AgdaBound{t}\AgdaSpace{}%
\AgdaOperator{\AgdaDatatype{≡}}\AgdaSpace{}%
\AgdaFunction{renameb}\AgdaSpace{}%
\AgdaBound{fb''}\AgdaSpace{}%
\AgdaBound{t}\<%
\\
\>[0]\AgdaFunction{renameb-cong}\AgdaSpace{}%
\AgdaBound{h}\AgdaSpace{}%
\AgdaSymbol{(}\AgdaInductiveConstructor{ε}\AgdaSpace{}%
\AgdaBound{x}\AgdaSymbol{)}\AgdaSpace{}%
\AgdaSymbol{=}\AgdaSpace{}%
\AgdaFunction{cong}\AgdaSpace{}%
\AgdaInductiveConstructor{ε}\AgdaSpace{}%
\AgdaSymbol{(}\AgdaBound{h}\AgdaSpace{}%
\AgdaBound{x}\AgdaSymbol{)}\<%
\\
\>[0]\AgdaFunction{renameb-cong}\AgdaSpace{}%
\AgdaBound{h}\AgdaSpace{}%
\AgdaSymbol{(}\AgdaInductiveConstructor{box}\AgdaSpace{}%
\AgdaBound{t}\AgdaSymbol{)}\AgdaSpace{}%
\AgdaSymbol{=}\AgdaSpace{}%
\AgdaFunction{cong}\AgdaSpace{}%
\AgdaInductiveConstructor{box}\AgdaSpace{}%
\AgdaSymbol{(}\AgdaFunction{renameb-cong}\AgdaSpace{}%
\AgdaBound{h}\AgdaSpace{}%
\AgdaBound{t}\AgdaSymbol{)}\<%
\\
\>[0]\AgdaFunction{renameb-cong}\AgdaSpace{}%
\AgdaBound{h}\AgdaSpace{}%
\AgdaSymbol{(}\AgdaInductiveConstructor{ƛb}\AgdaSpace{}%
\AgdaBound{t}\AgdaSymbol{)}\AgdaSpace{}%
\AgdaSymbol{=}\AgdaSpace{}%
\AgdaFunction{cong}\AgdaSpace{}%
\AgdaInductiveConstructor{ƛb}\AgdaSpace{}%
\AgdaSymbol{(}\AgdaFunction{renameb-cong}\AgdaSpace{}%
\AgdaSymbol{(}\AgdaFunction{extb-cong}\AgdaSpace{}%
\AgdaBound{h}\AgdaSymbol{)}\AgdaSpace{}%
\AgdaBound{t}\AgdaSymbol{)}\<%
\\
\>[0]\AgdaFunction{renameb-cong}\AgdaSpace{}%
\AgdaBound{h}\AgdaSpace{}%
\AgdaSymbol{(}\AgdaBound{t}\AgdaSpace{}%
\AgdaOperator{\AgdaInductiveConstructor{·b}}\AgdaSpace{}%
\AgdaBound{t₁}\AgdaSymbol{)}\AgdaSpace{}%
\AgdaSymbol{=}\AgdaSpace{}%
\AgdaFunction{cong₂}\AgdaSpace{}%
\AgdaOperator{\AgdaInductiveConstructor{\AgdaUnderscore{}·b\AgdaUnderscore{}}}\AgdaSpace{}%
\AgdaSymbol{(}\AgdaFunction{renameb-cong}\AgdaSpace{}%
\AgdaBound{h}\AgdaSpace{}%
\AgdaBound{t}\AgdaSymbol{)}\AgdaSpace{}%
\AgdaSymbol{(}\AgdaFunction{renameb-cong}\AgdaSpace{}%
\AgdaBound{h}\AgdaSpace{}%
\AgdaBound{t₁}\AgdaSymbol{)}\<%
\\
%
\\[\AgdaEmptyExtraSkip]%
\>[0]\AgdaComment{--\ extb\ distributes\ over\ composition:}\<%
\\
\>[0]\AgdaComment{--\ Extending\ first\ with\ fb',\ then\ with\ fb'',\ }\<%
\\
\>[0]\AgdaComment{--\ is\ the\ same\ as\ extending\ once\ with\ fb''\ ∘\ fb'.}\<%
\\
\>[0]\AgdaFunction{extb-comp}\AgdaSpace{}%
\AgdaSymbol{:}\AgdaSpace{}%
\AgdaSymbol{∀}\AgdaSpace{}%
\AgdaSymbol{\{}\AgdaBound{Γ₁}\AgdaSpace{}%
\AgdaBound{Γ₂}\AgdaSpace{}%
\AgdaBound{Γ₃}\AgdaSpace{}%
\AgdaBound{A}\AgdaSpace{}%
\AgdaBound{B}\AgdaSymbol{\}}\<%
\\
\>[0][@{}l@{\AgdaIndent{0}}]%
\>[2]\AgdaSymbol{→}\AgdaSpace{}%
\AgdaSymbol{(}\AgdaBound{fb'}\AgdaSpace{}%
\AgdaSymbol{:}\AgdaSpace{}%
\AgdaSymbol{∀}\AgdaSpace{}%
\AgdaSymbol{\{}\AgdaBound{Ab}\AgdaSymbol{\}}\AgdaSpace{}%
\AgdaSymbol{→}\AgdaSpace{}%
\AgdaBound{Ab}\AgdaSpace{}%
\AgdaOperator{\AgdaDatatype{∈b}}\AgdaSpace{}%
\AgdaBound{Γ₁}\AgdaSpace{}%
\AgdaSymbol{→}\AgdaSpace{}%
\AgdaBound{Ab}\AgdaSpace{}%
\AgdaOperator{\AgdaDatatype{∈b}}\AgdaSpace{}%
\AgdaBound{Γ₂}\AgdaSymbol{)}\<%
\\
%
\>[2]\AgdaSymbol{→}\AgdaSpace{}%
\AgdaSymbol{(}\AgdaBound{fb''}\AgdaSpace{}%
\AgdaSymbol{:}\AgdaSpace{}%
\AgdaSymbol{∀}\AgdaSpace{}%
\AgdaSymbol{\{}\AgdaBound{Ab}\AgdaSymbol{\}}\AgdaSpace{}%
\AgdaSymbol{→}\AgdaSpace{}%
\AgdaBound{Ab}\AgdaSpace{}%
\AgdaOperator{\AgdaDatatype{∈b}}\AgdaSpace{}%
\AgdaBound{Γ₂}\AgdaSpace{}%
\AgdaSymbol{→}\AgdaSpace{}%
\AgdaBound{Ab}\AgdaSpace{}%
\AgdaOperator{\AgdaDatatype{∈b}}\AgdaSpace{}%
\AgdaBound{Γ₃}\AgdaSymbol{)}\<%
\\
%
\>[2]\AgdaSymbol{→}\AgdaSpace{}%
\AgdaSymbol{(}\AgdaBound{x}\AgdaSpace{}%
\AgdaSymbol{:}\AgdaSpace{}%
\AgdaBound{A}\AgdaSpace{}%
\AgdaOperator{\AgdaDatatype{∈b}}\AgdaSpace{}%
\AgdaBound{Γ₁}\AgdaSpace{}%
\AgdaOperator{\AgdaInductiveConstructor{,b}}\AgdaSpace{}%
\AgdaBound{B}\AgdaSymbol{)}\<%
\\
%
\>[2]\AgdaSymbol{→}\AgdaSpace{}%
\AgdaFunction{extb}\AgdaSpace{}%
\AgdaBound{fb''}\AgdaSpace{}%
\AgdaSymbol{(}\AgdaFunction{extb}\AgdaSpace{}%
\AgdaBound{fb'}\AgdaSpace{}%
\AgdaBound{x}\AgdaSymbol{)}\AgdaSpace{}%
\AgdaOperator{\AgdaDatatype{≡}}\AgdaSpace{}%
\AgdaFunction{extb}\AgdaSpace{}%
\AgdaSymbol{(λ}\AgdaSpace{}%
\AgdaBound{x₁}\AgdaSpace{}%
\AgdaSymbol{→}\AgdaSpace{}%
\AgdaBound{fb''}\AgdaSpace{}%
\AgdaSymbol{(}\AgdaBound{fb'}\AgdaSpace{}%
\AgdaBound{x₁}\AgdaSymbol{))}\AgdaSpace{}%
\AgdaBound{x}\<%
\\
\>[0]\AgdaFunction{extb-comp}\AgdaSpace{}%
\AgdaBound{fb'}\AgdaSpace{}%
\AgdaBound{fb''}\AgdaSpace{}%
\AgdaInductiveConstructor{Zb}\AgdaSpace{}%
\AgdaSymbol{=}\AgdaSpace{}%
\AgdaInductiveConstructor{refl}\<%
\\
\>[0]\AgdaFunction{extb-comp}\AgdaSpace{}%
\AgdaBound{fb'}\AgdaSpace{}%
\AgdaBound{fb''}\AgdaSpace{}%
\AgdaSymbol{(}\AgdaInductiveConstructor{Sb}\AgdaSpace{}%
\AgdaBound{x}\AgdaSymbol{)}\AgdaSpace{}%
\AgdaSymbol{=}\AgdaSpace{}%
\AgdaInductiveConstructor{refl}\<%
\\
%
\\[\AgdaEmptyExtraSkip]%
\>[0]\AgdaComment{--\ renameb\ distributes\ over\ composition:}\<%
\\
\>[0]\AgdaComment{--\ Renaming\ first\ with\ fb',\ then\ with\ fb''}\<%
\\
\>[0]\AgdaComment{--\ is\ the\ same\ as\ renaming\ once\ with\ fb''\ ∘\ fb'.}\<%
\\
\>[0]\AgdaFunction{renameb-comp}\AgdaSpace{}%
\AgdaSymbol{:}\AgdaSpace{}%
\AgdaSymbol{∀}\AgdaSpace{}%
\AgdaSymbol{\{}\AgdaBound{Γ₁}\AgdaSpace{}%
\AgdaBound{Γ₂}\AgdaSpace{}%
\AgdaBound{Γ₃}\AgdaSpace{}%
\AgdaBound{A}\AgdaSymbol{\}}\<%
\\
\>[0][@{}l@{\AgdaIndent{0}}]%
\>[2]\AgdaSymbol{→}\AgdaSpace{}%
\AgdaSymbol{(}\AgdaBound{fb'}\AgdaSpace{}%
\AgdaSymbol{:}\AgdaSpace{}%
\AgdaSymbol{∀}\AgdaSpace{}%
\AgdaSymbol{\{}\AgdaBound{Ab}\AgdaSymbol{\}}\AgdaSpace{}%
\AgdaSymbol{→}\AgdaSpace{}%
\AgdaBound{Ab}\AgdaSpace{}%
\AgdaOperator{\AgdaDatatype{∈b}}\AgdaSpace{}%
\AgdaBound{Γ₁}\AgdaSpace{}%
\AgdaSymbol{→}\AgdaSpace{}%
\AgdaBound{Ab}\AgdaSpace{}%
\AgdaOperator{\AgdaDatatype{∈b}}\AgdaSpace{}%
\AgdaBound{Γ₂}\AgdaSymbol{)}\<%
\\
%
\>[2]\AgdaSymbol{→}\AgdaSpace{}%
\AgdaSymbol{(}\AgdaBound{fb''}\AgdaSpace{}%
\AgdaSymbol{:}\AgdaSpace{}%
\AgdaSymbol{∀}\AgdaSpace{}%
\AgdaSymbol{\{}\AgdaBound{Ab}\AgdaSymbol{\}}\AgdaSpace{}%
\AgdaSymbol{→}\AgdaSpace{}%
\AgdaBound{Ab}\AgdaSpace{}%
\AgdaOperator{\AgdaDatatype{∈b}}\AgdaSpace{}%
\AgdaBound{Γ₂}\AgdaSpace{}%
\AgdaSymbol{→}\AgdaSpace{}%
\AgdaBound{Ab}\AgdaSpace{}%
\AgdaOperator{\AgdaDatatype{∈b}}\AgdaSpace{}%
\AgdaBound{Γ₃}\AgdaSymbol{)}\<%
\\
%
\>[2]\AgdaSymbol{→}\AgdaSpace{}%
\AgdaSymbol{(}\AgdaBound{t}\AgdaSpace{}%
\AgdaSymbol{:}\AgdaSpace{}%
\AgdaBound{Γ₁}\AgdaSpace{}%
\AgdaOperator{\AgdaDatatype{⊢b}}\AgdaSpace{}%
\AgdaBound{A}\AgdaSymbol{)}\<%
\\
%
\>[2]\AgdaSymbol{→}\AgdaSpace{}%
\AgdaFunction{renameb}\AgdaSpace{}%
\AgdaBound{fb''}\AgdaSpace{}%
\AgdaSymbol{(}\AgdaFunction{renameb}\AgdaSpace{}%
\AgdaBound{fb'}\AgdaSpace{}%
\AgdaBound{t}\AgdaSymbol{)}\AgdaSpace{}%
\AgdaOperator{\AgdaDatatype{≡}}\AgdaSpace{}%
\AgdaFunction{renameb}\AgdaSpace{}%
\AgdaSymbol{(λ}\AgdaSpace{}%
\AgdaBound{x}\AgdaSpace{}%
\AgdaSymbol{→}\AgdaSpace{}%
\AgdaBound{fb''}\AgdaSpace{}%
\AgdaSymbol{(}\AgdaBound{fb'}\AgdaSpace{}%
\AgdaBound{x}\AgdaSymbol{))}\AgdaSpace{}%
\AgdaBound{t}\<%
\\
\>[0]\AgdaFunction{renameb-comp}\AgdaSpace{}%
\AgdaBound{fb'}\AgdaSpace{}%
\AgdaBound{fb''}\AgdaSpace{}%
\AgdaSymbol{(}\AgdaInductiveConstructor{ε}\AgdaSpace{}%
\AgdaBound{x}\AgdaSymbol{)}\AgdaSpace{}%
\AgdaSymbol{=}\AgdaSpace{}%
\AgdaInductiveConstructor{refl}\<%
\\
\>[0]\AgdaFunction{renameb-comp}\AgdaSpace{}%
\AgdaBound{fb'}\AgdaSpace{}%
\AgdaBound{fb''}\AgdaSpace{}%
\AgdaSymbol{(}\AgdaInductiveConstructor{box}\AgdaSpace{}%
\AgdaBound{t}\AgdaSymbol{)}\AgdaSpace{}%
\AgdaSymbol{=}\AgdaSpace{}%
\AgdaFunction{cong}\AgdaSpace{}%
\AgdaInductiveConstructor{box}\AgdaSpace{}%
\AgdaSymbol{(}\AgdaFunction{renameb-comp}\AgdaSpace{}%
\AgdaBound{fb'}\AgdaSpace{}%
\AgdaBound{fb''}\AgdaSpace{}%
\AgdaBound{t}\AgdaSymbol{)}\<%
\\
\>[0]\AgdaFunction{renameb-comp}\AgdaSpace{}%
\AgdaBound{fb'}\AgdaSpace{}%
\AgdaBound{fb''}\AgdaSpace{}%
\AgdaSymbol{(}\AgdaInductiveConstructor{ƛb}\AgdaSpace{}%
\AgdaBound{t}\AgdaSymbol{)}\AgdaSpace{}%
\AgdaSymbol{=}\<%
\\
\>[0][@{}l@{\AgdaIndent{0}}]%
\>[2]\AgdaOperator{\AgdaFunction{beginEq}}\<%
\\
\>[2][@{}l@{\AgdaIndent{0}}]%
\>[4]\AgdaInductiveConstructor{ƛb}\AgdaSpace{}%
\AgdaSymbol{(}\AgdaFunction{renameb}\AgdaSpace{}%
\AgdaSymbol{(}\AgdaFunction{extb}\AgdaSpace{}%
\AgdaBound{fb''}\AgdaSymbol{)}\AgdaSpace{}%
\AgdaSymbol{(}\AgdaFunction{renameb}\AgdaSpace{}%
\AgdaSymbol{(}\AgdaFunction{extb}\AgdaSpace{}%
\AgdaBound{fb'}\AgdaSymbol{)}\AgdaSpace{}%
\AgdaBound{t}\AgdaSymbol{))}\<%
\\
%
\>[2]\AgdaFunction{≡⟨}\AgdaSpace{}%
\AgdaFunction{cong}\AgdaSpace{}%
\AgdaInductiveConstructor{ƛb}\AgdaSpace{}%
\AgdaSymbol{(}\AgdaFunction{renameb-comp}\AgdaSpace{}%
\AgdaSymbol{(}\AgdaFunction{extb}\AgdaSpace{}%
\AgdaBound{fb'}\AgdaSymbol{)}\AgdaSpace{}%
\AgdaSymbol{(}\AgdaFunction{extb}\AgdaSpace{}%
\AgdaBound{fb''}\AgdaSymbol{)}\AgdaSpace{}%
\AgdaBound{t}\AgdaSymbol{)}\AgdaSpace{}%
\AgdaFunction{⟩}\<%
\\
\>[2][@{}l@{\AgdaIndent{0}}]%
\>[4]\AgdaInductiveConstructor{ƛb}\AgdaSpace{}%
\AgdaSymbol{(}\AgdaFunction{renameb}\AgdaSpace{}%
\AgdaSymbol{(λ}\AgdaSpace{}%
\AgdaBound{x}\AgdaSpace{}%
\AgdaSymbol{→}\AgdaSpace{}%
\AgdaFunction{extb}\AgdaSpace{}%
\AgdaBound{fb''}\AgdaSpace{}%
\AgdaSymbol{(}\AgdaFunction{extb}\AgdaSpace{}%
\AgdaBound{fb'}\AgdaSpace{}%
\AgdaBound{x}\AgdaSymbol{))}\AgdaSpace{}%
\AgdaBound{t}\AgdaSymbol{)}\<%
\\
%
\>[2]\AgdaFunction{≡⟨}\AgdaSpace{}%
\AgdaFunction{cong}\AgdaSpace{}%
\AgdaInductiveConstructor{ƛb}\AgdaSpace{}%
\AgdaSymbol{(}\AgdaFunction{renameb-cong}\AgdaSpace{}%
\AgdaSymbol{(}\AgdaFunction{extb-comp}\AgdaSpace{}%
\AgdaBound{fb'}\AgdaSpace{}%
\AgdaBound{fb''}\AgdaSymbol{)}\AgdaSpace{}%
\AgdaBound{t}\AgdaSymbol{)}\AgdaSpace{}%
\AgdaFunction{⟩}\<%
\\
\>[2][@{}l@{\AgdaIndent{0}}]%
\>[4]\AgdaInductiveConstructor{ƛb}\AgdaSpace{}%
\AgdaSymbol{(}\AgdaFunction{renameb}\AgdaSpace{}%
\AgdaSymbol{(}\AgdaFunction{extb}\AgdaSpace{}%
\AgdaSymbol{(λ}\AgdaSpace{}%
\AgdaBound{x}\AgdaSpace{}%
\AgdaSymbol{→}\AgdaSpace{}%
\AgdaBound{fb''}\AgdaSpace{}%
\AgdaSymbol{(}\AgdaBound{fb'}\AgdaSpace{}%
\AgdaBound{x}\AgdaSymbol{)))}\AgdaSpace{}%
\AgdaBound{t}\AgdaSymbol{)}\<%
\\
%
\>[2]\AgdaOperator{\AgdaFunction{∎eq}}\<%
\\
\>[0]\AgdaFunction{renameb-comp}\AgdaSpace{}%
\AgdaBound{fb'}\AgdaSpace{}%
\AgdaBound{fb''}\AgdaSpace{}%
\AgdaSymbol{(}\AgdaBound{t}\AgdaSpace{}%
\AgdaOperator{\AgdaInductiveConstructor{·b}}\AgdaSpace{}%
\AgdaBound{t₁}\AgdaSymbol{)}\AgdaSpace{}%
\AgdaSymbol{=}\AgdaSpace{}%
\AgdaFunction{cong₂}\AgdaSpace{}%
\AgdaOperator{\AgdaInductiveConstructor{\AgdaUnderscore{}·b\AgdaUnderscore{}}}\AgdaSpace{}%
\AgdaSymbol{(}\AgdaFunction{renameb-comp}\AgdaSpace{}%
\AgdaBound{fb'}\AgdaSpace{}%
\AgdaBound{fb''}\AgdaSpace{}%
\AgdaBound{t}\AgdaSymbol{)}\AgdaSpace{}%
\AgdaSymbol{(}\AgdaFunction{renameb-comp}\AgdaSpace{}%
\AgdaBound{fb'}\AgdaSpace{}%
\AgdaBound{fb''}\AgdaSpace{}%
\AgdaBound{t₁}\AgdaSymbol{)}\<%
\\
%
\\[\AgdaEmptyExtraSkip]%
\>[0]\AgdaComment{--\ extb\ and\ Sb\ commute\ under\ renameb.}\<%
\\
\>[0]\AgdaFunction{rename-extb-Sb-commute}\AgdaSpace{}%
\AgdaSymbol{:}\AgdaSpace{}%
\AgdaSymbol{∀}\AgdaSpace{}%
\AgdaSymbol{\{}\AgdaBound{Γ}\AgdaSpace{}%
\AgdaBound{A}\AgdaSpace{}%
\AgdaBound{B}\AgdaSpace{}%
\AgdaBound{Δ}\AgdaSymbol{\}}\<%
\\
\>[0][@{}l@{\AgdaIndent{0}}]%
\>[2]\AgdaSymbol{→}\AgdaSpace{}%
\AgdaSymbol{(}\AgdaBound{fb'}\AgdaSpace{}%
\AgdaSymbol{:}\AgdaSpace{}%
\AgdaSymbol{∀}\AgdaSpace{}%
\AgdaSymbol{\{}\AgdaBound{Ab}\AgdaSymbol{\}}\AgdaSpace{}%
\AgdaSymbol{→}\AgdaSpace{}%
\AgdaBound{Ab}\AgdaSpace{}%
\AgdaOperator{\AgdaDatatype{∈b}}\AgdaSpace{}%
\AgdaBound{Γ}\AgdaSpace{}%
\AgdaSymbol{→}\AgdaSpace{}%
\AgdaBound{Ab}\AgdaSpace{}%
\AgdaOperator{\AgdaDatatype{∈b}}\AgdaSpace{}%
\AgdaBound{Δ}\AgdaSymbol{)}\<%
\\
%
\>[2]\AgdaSymbol{→}\AgdaSpace{}%
\AgdaSymbol{(}\AgdaBound{t}\AgdaSpace{}%
\AgdaSymbol{:}\AgdaSpace{}%
\AgdaBound{Γ}\AgdaSpace{}%
\AgdaOperator{\AgdaDatatype{⊢b}}\AgdaSpace{}%
\AgdaBound{A}\AgdaSymbol{)}\<%
\\
%
\>[2]\AgdaSymbol{→}%
\>[579I]\AgdaFunction{renameb}\AgdaSpace{}%
\AgdaSymbol{(}\AgdaFunction{extb}\AgdaSpace{}%
\AgdaBound{fb'}\AgdaSpace{}%
\AgdaSymbol{\{}\AgdaArgument{B}\AgdaSpace{}%
\AgdaSymbol{=}\AgdaSpace{}%
\AgdaBound{B}\AgdaSymbol{\})}\AgdaSpace{}%
\AgdaSymbol{(}\AgdaFunction{renameb}\AgdaSpace{}%
\AgdaInductiveConstructor{Sb}\AgdaSpace{}%
\AgdaBound{t}\AgdaSymbol{)}\<%
\\
\>[.][@{}l@{}]\<[579I]%
\>[4]\AgdaOperator{\AgdaDatatype{≡}}\AgdaSpace{}%
\AgdaFunction{renameb}\AgdaSpace{}%
\AgdaInductiveConstructor{Sb}\AgdaSpace{}%
\AgdaSymbol{(}\AgdaFunction{renameb}\AgdaSpace{}%
\AgdaBound{fb'}\AgdaSpace{}%
\AgdaBound{t}\AgdaSymbol{)}\<%
\\
\>[0]\AgdaFunction{rename-extb-Sb-commute}\AgdaSpace{}%
\AgdaBound{fb'}\AgdaSpace{}%
\AgdaBound{t}\AgdaSpace{}%
\AgdaSymbol{=}\<%
\\
\>[0][@{}l@{\AgdaIndent{0}}]%
\>[2]\AgdaOperator{\AgdaFunction{beginEq}}\<%
\\
\>[2][@{}l@{\AgdaIndent{0}}]%
\>[4]\AgdaFunction{renameb}\AgdaSpace{}%
\AgdaSymbol{(}\AgdaFunction{extb}\AgdaSpace{}%
\AgdaBound{fb'}\AgdaSymbol{)}\AgdaSpace{}%
\AgdaSymbol{(}\AgdaFunction{renameb}\AgdaSpace{}%
\AgdaInductiveConstructor{Sb}\AgdaSpace{}%
\AgdaBound{t}\AgdaSymbol{)}\<%
\\
%
\>[2]\AgdaFunction{≡⟨}\AgdaSpace{}%
\AgdaFunction{renameb-comp}\AgdaSpace{}%
\AgdaInductiveConstructor{Sb}\AgdaSpace{}%
\AgdaSymbol{(}\AgdaFunction{extb}\AgdaSpace{}%
\AgdaBound{fb'}\AgdaSymbol{)}\AgdaSpace{}%
\AgdaBound{t}\AgdaSpace{}%
\AgdaFunction{⟩}\<%
\\
\>[2][@{}l@{\AgdaIndent{0}}]%
\>[4]\AgdaFunction{renameb}\AgdaSpace{}%
\AgdaSymbol{(λ}\AgdaSpace{}%
\AgdaBound{x}\AgdaSpace{}%
\AgdaSymbol{→}\AgdaSpace{}%
\AgdaFunction{extb}\AgdaSpace{}%
\AgdaBound{fb'}\AgdaSpace{}%
\AgdaSymbol{(}\AgdaInductiveConstructor{Sb}\AgdaSpace{}%
\AgdaBound{x}\AgdaSymbol{))}\AgdaSpace{}%
\AgdaBound{t}\<%
\\
%
\>[2]\AgdaFunction{≡⟨}\AgdaSpace{}%
\AgdaFunction{renameb-cong}\AgdaSpace{}%
\AgdaSymbol{(λ}\AgdaSpace{}%
\AgdaBound{x}\AgdaSpace{}%
\AgdaSymbol{→}\AgdaSpace{}%
\AgdaInductiveConstructor{refl}\AgdaSymbol{)}\AgdaSpace{}%
\AgdaBound{t}\AgdaSpace{}%
\AgdaFunction{⟩}\<%
\\
\>[2][@{}l@{\AgdaIndent{0}}]%
\>[4]\AgdaFunction{renameb}\AgdaSpace{}%
\AgdaSymbol{(λ}\AgdaSpace{}%
\AgdaBound{x}\AgdaSpace{}%
\AgdaSymbol{→}\AgdaSpace{}%
\AgdaInductiveConstructor{Sb}\AgdaSpace{}%
\AgdaSymbol{(}\AgdaBound{fb'}\AgdaSpace{}%
\AgdaBound{x}\AgdaSymbol{))}\AgdaSpace{}%
\AgdaBound{t}\<%
\\
%
\>[2]\AgdaFunction{≡⟨}\AgdaSpace{}%
\AgdaFunction{sym}\AgdaSpace{}%
\AgdaSymbol{(}\AgdaFunction{renameb-comp}\AgdaSpace{}%
\AgdaBound{fb'}\AgdaSpace{}%
\AgdaInductiveConstructor{Sb}\AgdaSpace{}%
\AgdaBound{t}\AgdaSymbol{)}\AgdaSpace{}%
\AgdaFunction{⟩}\<%
\\
\>[2][@{}l@{\AgdaIndent{0}}]%
\>[4]\AgdaFunction{renameb}\AgdaSpace{}%
\AgdaInductiveConstructor{Sb}\AgdaSpace{}%
\AgdaSymbol{(}\AgdaFunction{renameb}\AgdaSpace{}%
\AgdaBound{fb'}\AgdaSpace{}%
\AgdaBound{t}\AgdaSymbol{)}\<%
\\
%
\>[2]\AgdaOperator{\AgdaFunction{∎eq}}\<%
\\
%
\\[\AgdaEmptyExtraSkip]%
\>[0]\AgdaComment{--\ Preservation\ and\ Reflection\ of\ }\<%
\\
\>[0]\AgdaComment{--\ rename\ of\ Gödel's\ translation}\<%
\\
\>[0]\AgdaFunction{pres-rename}\AgdaSpace{}%
\AgdaSymbol{:}\AgdaSpace{}%
\AgdaSymbol{∀}\AgdaSpace{}%
\AgdaSymbol{\{}\AgdaBound{Γ}\AgdaSpace{}%
\AgdaBound{A}\AgdaSpace{}%
\AgdaBound{Δ}\AgdaSymbol{\}}\<%
\\
\>[0][@{}l@{\AgdaIndent{0}}]%
\>[2]\AgdaSymbol{→}\AgdaSpace{}%
\AgdaSymbol{(}\AgdaBound{f'}\AgdaSpace{}%
\AgdaSymbol{:}\AgdaSpace{}%
\AgdaSymbol{∀}\AgdaSpace{}%
\AgdaSymbol{\{}\AgdaBound{A}\AgdaSymbol{\}}\AgdaSpace{}%
\AgdaSymbol{→}\AgdaSpace{}%
\AgdaBound{A}\AgdaSpace{}%
\AgdaOperator{\AgdaDatatype{∈}}\AgdaSpace{}%
\AgdaBound{Γ}\AgdaSpace{}%
\AgdaSymbol{→}\AgdaSpace{}%
\AgdaBound{A}\AgdaSpace{}%
\AgdaOperator{\AgdaDatatype{∈}}\AgdaSpace{}%
\AgdaBound{Δ}\AgdaSymbol{)}\<%
\\
%
\>[2]\AgdaSymbol{→}\AgdaSpace{}%
\AgdaSymbol{(}\AgdaBound{fb'}\AgdaSpace{}%
\AgdaSymbol{:}\AgdaSpace{}%
\AgdaSymbol{∀}\AgdaSpace{}%
\AgdaSymbol{\{}\AgdaBound{Ab}\AgdaSymbol{\}}\AgdaSpace{}%
\AgdaSymbol{→}\AgdaSpace{}%
\AgdaBound{Ab}\AgdaSpace{}%
\AgdaOperator{\AgdaDatatype{∈b}}\AgdaSpace{}%
\AgdaFunction{embedContext}\AgdaSpace{}%
\AgdaBound{Γ}\AgdaSpace{}%
\AgdaSymbol{→}\AgdaSpace{}%
\AgdaBound{Ab}\AgdaSpace{}%
\AgdaOperator{\AgdaDatatype{∈b}}\AgdaSpace{}%
\AgdaFunction{embedContext}\AgdaSpace{}%
\AgdaBound{Δ}\AgdaSymbol{)}\<%
\\
%
\>[2]\AgdaSymbol{→}\AgdaSpace{}%
\AgdaSymbol{(\{}\AgdaBound{A'}\AgdaSpace{}%
\AgdaSymbol{:}\AgdaSpace{}%
\AgdaDatatype{Type}\AgdaSymbol{\}}\AgdaSpace{}%
\AgdaSymbol{(}\AgdaBound{x}\AgdaSpace{}%
\AgdaSymbol{:}\AgdaSpace{}%
\AgdaBound{A'}\AgdaSpace{}%
\AgdaOperator{\AgdaDatatype{∈}}\AgdaSpace{}%
\AgdaBound{Γ}\AgdaSymbol{)}\AgdaSpace{}%
\AgdaSymbol{→}\AgdaSpace{}%
\AgdaBound{fb'}\AgdaSpace{}%
\AgdaSymbol{(}\AgdaFunction{embedLookup}\AgdaSpace{}%
\AgdaBound{x}\AgdaSymbol{)}\AgdaSpace{}%
\AgdaOperator{\AgdaDatatype{≡}}\AgdaSpace{}%
\AgdaFunction{embedLookup}\AgdaSpace{}%
\AgdaSymbol{(}\AgdaBound{f'}\AgdaSpace{}%
\AgdaBound{x}\AgdaSymbol{))}\<%
\\
%
\>[2]\AgdaSymbol{→}\AgdaSpace{}%
\AgdaSymbol{(}\AgdaBound{t}\AgdaSpace{}%
\AgdaSymbol{:}\AgdaSpace{}%
\AgdaBound{Γ}\AgdaSpace{}%
\AgdaOperator{\AgdaDatatype{⊢}}\AgdaSpace{}%
\AgdaBound{A}\AgdaSymbol{)}\<%
\\
%
\>[2]\AgdaSymbol{→}\AgdaSpace{}%
\AgdaFunction{renameb}\AgdaSpace{}%
\AgdaBound{fb'}\AgdaSpace{}%
\AgdaSymbol{(}\AgdaFunction{embedTerm}\AgdaSpace{}%
\AgdaBound{t}\AgdaSymbol{)}\AgdaSpace{}%
\AgdaOperator{\AgdaDatatype{≡}}\AgdaSpace{}%
\AgdaFunction{embedTerm}\AgdaSpace{}%
\AgdaSymbol{(}\AgdaFunction{rename}\AgdaSpace{}%
\AgdaBound{f'}\AgdaSpace{}%
\AgdaBound{t}\AgdaSymbol{)}\<%
\\
\>[0]\AgdaFunction{pres-rename}\AgdaSpace{}%
\AgdaBound{f'}\AgdaSpace{}%
\AgdaBound{fb'}\AgdaSpace{}%
\AgdaBound{h}\AgdaSpace{}%
\AgdaSymbol{(}\AgdaInductiveConstructor{`}\AgdaSpace{}%
\AgdaBound{x}\AgdaSymbol{)}\AgdaSpace{}%
\AgdaSymbol{=}\AgdaSpace{}%
\AgdaFunction{cong}\AgdaSpace{}%
\AgdaInductiveConstructor{box}\AgdaSpace{}%
\AgdaSymbol{(}\AgdaFunction{cong}\AgdaSpace{}%
\AgdaInductiveConstructor{ε}\AgdaSpace{}%
\AgdaSymbol{(}\AgdaBound{h}\AgdaSpace{}%
\AgdaBound{x}\AgdaSymbol{))}\<%
\\
\>[0]\AgdaFunction{pres-rename}\AgdaSpace{}%
\AgdaBound{f'}\AgdaSpace{}%
\AgdaBound{fb'}\AgdaSpace{}%
\AgdaBound{h}\AgdaSpace{}%
\AgdaSymbol{(}\AgdaInductiveConstructor{ƛ}\AgdaSpace{}%
\AgdaBound{t}\AgdaSymbol{)}\AgdaSpace{}%
\AgdaSymbol{=}\AgdaSpace{}%
\AgdaFunction{cong}\AgdaSpace{}%
\AgdaInductiveConstructor{box}\AgdaSpace{}%
\AgdaSymbol{(}\AgdaFunction{cong}\AgdaSpace{}%
\AgdaInductiveConstructor{ƛb}\AgdaSpace{}%
\AgdaSymbol{(}\AgdaFunction{pres-rename}\AgdaSpace{}%
\AgdaSymbol{(}\AgdaFunction{calc.ext}\AgdaSpace{}%
\AgdaBound{f'}\AgdaSymbol{)}\AgdaSpace{}%
\AgdaSymbol{(}\AgdaFunction{boxCalc.extb}\AgdaSpace{}%
\AgdaBound{fb'}\AgdaSymbol{)}\AgdaSpace{}%
\AgdaSymbol{(}\AgdaFunction{pres-ext}\AgdaSpace{}%
\AgdaBound{f'}\AgdaSpace{}%
\AgdaBound{fb'}\AgdaSpace{}%
\AgdaBound{h}\AgdaSymbol{)}\AgdaSpace{}%
\AgdaBound{t}\AgdaSymbol{))}\<%
\\
\>[0]\AgdaFunction{pres-rename}\AgdaSpace{}%
\AgdaBound{f'}\AgdaSpace{}%
\AgdaBound{fb'}\AgdaSpace{}%
\AgdaBound{h}\AgdaSpace{}%
\AgdaSymbol{(}\AgdaBound{t}\AgdaSpace{}%
\AgdaOperator{\AgdaInductiveConstructor{·}}\AgdaSpace{}%
\AgdaBound{t₁}\AgdaSymbol{)}\AgdaSpace{}%
\AgdaSymbol{=}\<%
\\
\>[0][@{}l@{\AgdaIndent{0}}]%
\>[2]\AgdaOperator{\AgdaFunction{beginEq}}\<%
\\
\>[2][@{}l@{\AgdaIndent{0}}]%
\>[4]\AgdaInductiveConstructor{ƛb}\AgdaSpace{}%
\AgdaSymbol{(}\AgdaInductiveConstructor{ε}\AgdaSpace{}%
\AgdaInductiveConstructor{Zb}\AgdaSpace{}%
\AgdaOperator{\AgdaInductiveConstructor{·b}}\AgdaSpace{}%
\AgdaFunction{renameb}\AgdaSpace{}%
\AgdaSymbol{(}\AgdaFunction{extb}\AgdaSpace{}%
\AgdaBound{fb'}\AgdaSymbol{)}\AgdaSpace{}%
\AgdaSymbol{(}\AgdaFunction{renameb}\AgdaSpace{}%
\AgdaInductiveConstructor{Sb}\AgdaSpace{}%
\AgdaSymbol{(}\AgdaFunction{embedTerm}\AgdaSpace{}%
\AgdaBound{t₁}\AgdaSymbol{)))}\AgdaSpace{}%
\AgdaOperator{\AgdaInductiveConstructor{·b}}\AgdaSpace{}%
\AgdaFunction{renameb}\AgdaSpace{}%
\AgdaBound{fb'}\AgdaSpace{}%
\AgdaSymbol{(}\AgdaFunction{embedTerm}\AgdaSpace{}%
\AgdaBound{t}\AgdaSymbol{)}\<%
\\
%
\>[2]\AgdaFunction{≡⟨}\AgdaSpace{}%
\AgdaFunction{cong₂}\AgdaSpace{}%
\AgdaOperator{\AgdaInductiveConstructor{\AgdaUnderscore{}·b\AgdaUnderscore{}}}\AgdaSpace{}%
\AgdaInductiveConstructor{refl}\AgdaSpace{}%
\AgdaSymbol{(}\AgdaFunction{pres-rename}\AgdaSpace{}%
\AgdaBound{f'}\AgdaSpace{}%
\AgdaBound{fb'}\AgdaSpace{}%
\AgdaBound{h}\AgdaSpace{}%
\AgdaBound{t}\AgdaSymbol{)}\AgdaSpace{}%
\AgdaFunction{⟩}\<%
\\
\>[2][@{}l@{\AgdaIndent{0}}]%
\>[4]\AgdaInductiveConstructor{ƛb}\AgdaSpace{}%
\AgdaSymbol{(}\AgdaInductiveConstructor{ε}\AgdaSpace{}%
\AgdaInductiveConstructor{Zb}\AgdaSpace{}%
\AgdaOperator{\AgdaInductiveConstructor{·b}}\AgdaSpace{}%
\AgdaFunction{renameb}\AgdaSpace{}%
\AgdaSymbol{(}\AgdaFunction{extb}\AgdaSpace{}%
\AgdaBound{fb'}\AgdaSymbol{)}\AgdaSpace{}%
\AgdaSymbol{(}\AgdaFunction{renameb}\AgdaSpace{}%
\AgdaInductiveConstructor{Sb}\AgdaSpace{}%
\AgdaSymbol{(}\AgdaFunction{embedTerm}\AgdaSpace{}%
\AgdaBound{t₁}\AgdaSymbol{)))}\AgdaSpace{}%
\AgdaOperator{\AgdaInductiveConstructor{·b}}\AgdaSpace{}%
\AgdaFunction{embedTerm}\AgdaSpace{}%
\AgdaSymbol{(}\AgdaFunction{rename}\AgdaSpace{}%
\AgdaBound{f'}\AgdaSpace{}%
\AgdaBound{t}\AgdaSymbol{)}\<%
\\
%
\>[2]\AgdaFunction{≡⟨}\AgdaSpace{}%
\AgdaFunction{cong₂}\AgdaSpace{}%
\AgdaOperator{\AgdaInductiveConstructor{\AgdaUnderscore{}·b\AgdaUnderscore{}}}\AgdaSpace{}%
\AgdaSymbol{(}\AgdaFunction{cong}\AgdaSpace{}%
\AgdaInductiveConstructor{ƛb}\AgdaSpace{}%
\AgdaSymbol{(}\AgdaFunction{cong₂}\AgdaSpace{}%
\AgdaOperator{\AgdaInductiveConstructor{\AgdaUnderscore{}·b\AgdaUnderscore{}}}\AgdaSpace{}%
\AgdaInductiveConstructor{refl}\AgdaSpace{}%
\AgdaSymbol{(}\AgdaFunction{rename-extb-Sb-commute}\AgdaSpace{}%
\AgdaBound{fb'}\AgdaSpace{}%
\AgdaSymbol{(}\AgdaFunction{embedTerm}\AgdaSpace{}%
\AgdaBound{t₁}\AgdaSymbol{))))}\AgdaSpace{}%
\AgdaInductiveConstructor{refl}\AgdaSpace{}%
\AgdaFunction{⟩}\<%
\\
\>[2][@{}l@{\AgdaIndent{0}}]%
\>[4]\AgdaInductiveConstructor{ƛb}\AgdaSpace{}%
\AgdaSymbol{(}\AgdaInductiveConstructor{ε}\AgdaSpace{}%
\AgdaInductiveConstructor{Zb}\AgdaSpace{}%
\AgdaOperator{\AgdaInductiveConstructor{·b}}\AgdaSpace{}%
\AgdaFunction{renameb}\AgdaSpace{}%
\AgdaInductiveConstructor{Sb}\AgdaSpace{}%
\AgdaSymbol{(}\AgdaFunction{renameb}\AgdaSpace{}%
\AgdaBound{fb'}\AgdaSpace{}%
\AgdaSymbol{(}\AgdaFunction{embedTerm}\AgdaSpace{}%
\AgdaBound{t₁}\AgdaSymbol{)))}\AgdaSpace{}%
\AgdaOperator{\AgdaInductiveConstructor{·b}}\AgdaSpace{}%
\AgdaFunction{embedTerm}\AgdaSpace{}%
\AgdaSymbol{(}\AgdaFunction{rename}\AgdaSpace{}%
\AgdaBound{f'}\AgdaSpace{}%
\AgdaBound{t}\AgdaSymbol{)}\<%
\\
%
\>[2]\AgdaFunction{≡⟨}\AgdaSpace{}%
\AgdaFunction{cong₂}\AgdaSpace{}%
\AgdaOperator{\AgdaInductiveConstructor{\AgdaUnderscore{}·b\AgdaUnderscore{}}}\AgdaSpace{}%
\AgdaSymbol{(}\AgdaFunction{cong}\AgdaSpace{}%
\AgdaInductiveConstructor{ƛb}\AgdaSpace{}%
\AgdaSymbol{(}\AgdaFunction{cong₂}\AgdaSpace{}%
\AgdaOperator{\AgdaInductiveConstructor{\AgdaUnderscore{}·b\AgdaUnderscore{}}}\AgdaSpace{}%
\AgdaInductiveConstructor{refl}\AgdaSpace{}%
\AgdaSymbol{(}\AgdaFunction{cong}\AgdaSpace{}%
\AgdaSymbol{(}\AgdaFunction{renameb}\AgdaSpace{}%
\AgdaInductiveConstructor{Sb}\AgdaSymbol{)}\AgdaSpace{}%
\AgdaSymbol{(}\AgdaFunction{pres-rename}\AgdaSpace{}%
\AgdaBound{f'}\AgdaSpace{}%
\AgdaBound{fb'}\AgdaSpace{}%
\AgdaBound{h}\AgdaSpace{}%
\AgdaBound{t₁}\AgdaSymbol{))))}\AgdaSpace{}%
\AgdaInductiveConstructor{refl}\AgdaSpace{}%
\AgdaFunction{⟩}\<%
\\
\>[2][@{}l@{\AgdaIndent{0}}]%
\>[4]\AgdaInductiveConstructor{ƛb}\AgdaSpace{}%
\AgdaSymbol{(}\AgdaInductiveConstructor{ε}\AgdaSpace{}%
\AgdaInductiveConstructor{Zb}\AgdaSpace{}%
\AgdaOperator{\AgdaInductiveConstructor{·b}}\AgdaSpace{}%
\AgdaFunction{renameb}\AgdaSpace{}%
\AgdaInductiveConstructor{Sb}\AgdaSpace{}%
\AgdaSymbol{(}\AgdaFunction{embedTerm}\AgdaSpace{}%
\AgdaSymbol{(}\AgdaFunction{rename}\AgdaSpace{}%
\AgdaBound{f'}\AgdaSpace{}%
\AgdaBound{t₁}\AgdaSymbol{)))}\AgdaSpace{}%
\AgdaOperator{\AgdaInductiveConstructor{·b}}\AgdaSpace{}%
\AgdaFunction{embedTerm}\AgdaSpace{}%
\AgdaSymbol{(}\AgdaFunction{rename}\AgdaSpace{}%
\AgdaBound{f'}\AgdaSpace{}%
\AgdaBound{t}\AgdaSymbol{)}\<%
\\
%
\>[2]\AgdaOperator{\AgdaFunction{∎eq}}\<%
\\
%
\\[\AgdaEmptyExtraSkip]%
\>[0]\AgdaComment{--\ Preservation\ and\ Reflection\ of\ }\<%
\\
\>[0]\AgdaComment{--\ exts\ of\ Gödel's\ translation}\<%
\\
\>[0]\AgdaFunction{pres-exts}\AgdaSpace{}%
\AgdaSymbol{:}\AgdaSpace{}%
\AgdaSymbol{∀}\AgdaSpace{}%
\AgdaSymbol{\{}\AgdaBound{Γ}\AgdaSpace{}%
\AgdaBound{Δ}\AgdaSymbol{\}}\<%
\\
\>[0][@{}l@{\AgdaIndent{0}}]%
\>[2]\AgdaSymbol{→}\AgdaSpace{}%
\AgdaSymbol{(}\AgdaBound{fb'}\AgdaSpace{}%
\AgdaSymbol{:}\AgdaSpace{}%
\AgdaSymbol{∀}\AgdaSpace{}%
\AgdaSymbol{\{}\AgdaBound{Ab}\AgdaSymbol{\}}\AgdaSpace{}%
\AgdaSymbol{→}\AgdaSpace{}%
\AgdaInductiveConstructor{□}\AgdaSpace{}%
\AgdaBound{Ab}\AgdaSpace{}%
\AgdaOperator{\AgdaDatatype{∈b}}\AgdaSpace{}%
\AgdaFunction{embedContext}\AgdaSpace{}%
\AgdaBound{Γ}\AgdaSpace{}%
\AgdaSymbol{→}\AgdaSpace{}%
\AgdaFunction{embedContext}\AgdaSpace{}%
\AgdaBound{Δ}\AgdaSpace{}%
\AgdaOperator{\AgdaDatatype{⊢b}}\AgdaSpace{}%
\AgdaBound{Ab}\AgdaSymbol{)}\<%
\\
%
\>[2]\AgdaSymbol{→}\AgdaSpace{}%
\AgdaSymbol{(}\AgdaBound{f'}\AgdaSpace{}%
\AgdaSymbol{:}\AgdaSpace{}%
\AgdaSymbol{∀}\AgdaSpace{}%
\AgdaSymbol{\{}\AgdaBound{A}\AgdaSymbol{\}}\AgdaSpace{}%
\AgdaSymbol{→}\AgdaSpace{}%
\AgdaBound{A}\AgdaSpace{}%
\AgdaOperator{\AgdaDatatype{∈}}\AgdaSpace{}%
\AgdaBound{Γ}\AgdaSpace{}%
\AgdaSymbol{→}\AgdaSpace{}%
\AgdaBound{Δ}\AgdaSpace{}%
\AgdaOperator{\AgdaDatatype{⊢}}\AgdaSpace{}%
\AgdaBound{A}\AgdaSymbol{)}\<%
\\
%
\>[2]\AgdaSymbol{→}\AgdaSpace{}%
\AgdaSymbol{(\{}\AgdaBound{A'}\AgdaSpace{}%
\AgdaSymbol{:}\AgdaSpace{}%
\AgdaDatatype{Type}\AgdaSymbol{\}}\AgdaSpace{}%
\AgdaSymbol{(}\AgdaBound{x}\AgdaSpace{}%
\AgdaSymbol{:}\AgdaSpace{}%
\AgdaBound{A'}\AgdaSpace{}%
\AgdaOperator{\AgdaDatatype{∈}}\AgdaSpace{}%
\AgdaBound{Γ}\AgdaSymbol{)}\AgdaSpace{}%
\AgdaSymbol{→}\AgdaSpace{}%
\AgdaInductiveConstructor{box}\AgdaSpace{}%
\AgdaSymbol{(}\AgdaBound{fb'}\AgdaSpace{}%
\AgdaSymbol{(}\AgdaFunction{embedLookup}\AgdaSpace{}%
\AgdaBound{x}\AgdaSymbol{))}\AgdaSpace{}%
\AgdaOperator{\AgdaDatatype{≡}}\AgdaSpace{}%
\AgdaFunction{embedTerm}\AgdaSpace{}%
\AgdaSymbol{(}\AgdaBound{f'}\AgdaSpace{}%
\AgdaBound{x}\AgdaSymbol{))}\<%
\\
%
\>[2]\AgdaSymbol{→}\AgdaSpace{}%
\AgdaSymbol{\{}\AgdaBound{A'}\AgdaSpace{}%
\AgdaBound{A}\AgdaSpace{}%
\AgdaSymbol{:}\AgdaSpace{}%
\AgdaDatatype{Type}\AgdaSymbol{\}}\AgdaSpace{}%
\AgdaSymbol{(}\AgdaBound{x}\AgdaSpace{}%
\AgdaSymbol{:}\AgdaSpace{}%
\AgdaBound{A'}\AgdaSpace{}%
\AgdaOperator{\AgdaDatatype{∈}}\AgdaSpace{}%
\AgdaBound{Γ}\AgdaSpace{}%
\AgdaOperator{\AgdaInductiveConstructor{,}}\AgdaSpace{}%
\AgdaBound{A}\AgdaSymbol{)}\AgdaSpace{}%
\AgdaSymbol{→}\AgdaSpace{}%
\AgdaInductiveConstructor{box}\AgdaSpace{}%
\AgdaSymbol{(}\AgdaFunction{extsb}\AgdaSpace{}%
\AgdaBound{fb'}\AgdaSpace{}%
\AgdaSymbol{(}\AgdaFunction{embedLookup}\AgdaSpace{}%
\AgdaBound{x}\AgdaSymbol{))}\AgdaSpace{}%
\AgdaOperator{\AgdaDatatype{≡}}\AgdaSpace{}%
\AgdaFunction{embedTerm}\AgdaSpace{}%
\AgdaSymbol{(}\AgdaFunction{exts}\AgdaSpace{}%
\AgdaBound{f'}\AgdaSpace{}%
\AgdaBound{x}\AgdaSymbol{)}\<%
\\
\>[0]\AgdaFunction{pres-exts}\AgdaSpace{}%
\AgdaBound{fb'}\AgdaSpace{}%
\AgdaBound{f'}\AgdaSpace{}%
\AgdaBound{h}\AgdaSpace{}%
\AgdaInductiveConstructor{Z}\AgdaSpace{}%
\AgdaSymbol{=}\AgdaSpace{}%
\AgdaInductiveConstructor{refl}\<%
\\
\>[0]\AgdaFunction{pres-exts}\AgdaSpace{}%
\AgdaBound{fb'}\AgdaSpace{}%
\AgdaBound{f'}\AgdaSpace{}%
\AgdaBound{h}\AgdaSpace{}%
\AgdaSymbol{(}\AgdaInductiveConstructor{S}\AgdaSpace{}%
\AgdaBound{x}\AgdaSymbol{)}\AgdaSpace{}%
\AgdaSymbol{=}\<%
\\
\>[0][@{}l@{\AgdaIndent{0}}]%
\>[2]\AgdaOperator{\AgdaFunction{beginEq}}\<%
\\
\>[2][@{}l@{\AgdaIndent{0}}]%
\>[4]\AgdaInductiveConstructor{box}\AgdaSpace{}%
\AgdaSymbol{(}\AgdaFunction{renameb}\AgdaSpace{}%
\AgdaInductiveConstructor{Sb}\AgdaSpace{}%
\AgdaSymbol{(}\AgdaBound{fb'}\AgdaSpace{}%
\AgdaSymbol{(}\AgdaFunction{embedLookup}\AgdaSpace{}%
\AgdaBound{x}\AgdaSymbol{)))}\<%
\\
%
\>[2]\AgdaFunction{≡⟨}\AgdaSpace{}%
\AgdaInductiveConstructor{refl}\AgdaSpace{}%
\AgdaFunction{⟩}\<%
\\
\>[2][@{}l@{\AgdaIndent{0}}]%
\>[4]\AgdaFunction{renameb}\AgdaSpace{}%
\AgdaInductiveConstructor{Sb}\AgdaSpace{}%
\AgdaSymbol{(}\AgdaInductiveConstructor{box}\AgdaSpace{}%
\AgdaSymbol{(}\AgdaBound{fb'}\AgdaSpace{}%
\AgdaSymbol{(}\AgdaFunction{embedLookup}\AgdaSpace{}%
\AgdaBound{x}\AgdaSymbol{)))}\<%
\\
%
\>[2]\AgdaFunction{≡⟨}\AgdaSpace{}%
\AgdaFunction{cong}\AgdaSpace{}%
\AgdaSymbol{(}\AgdaFunction{renameb}\AgdaSpace{}%
\AgdaInductiveConstructor{Sb}\AgdaSymbol{)}\AgdaSpace{}%
\AgdaSymbol{(}\AgdaBound{h}\AgdaSpace{}%
\AgdaBound{x}\AgdaSymbol{)}\AgdaSpace{}%
\AgdaFunction{⟩}\<%
\\
\>[2][@{}l@{\AgdaIndent{0}}]%
\>[4]\AgdaFunction{renameb}\AgdaSpace{}%
\AgdaInductiveConstructor{Sb}\AgdaSpace{}%
\AgdaSymbol{(}\AgdaFunction{embedTerm}\AgdaSpace{}%
\AgdaSymbol{(}\AgdaBound{f'}\AgdaSpace{}%
\AgdaBound{x}\AgdaSymbol{))}\<%
\\
%
\>[2]\AgdaFunction{≡⟨}\AgdaSpace{}%
\AgdaFunction{pres-rename}\AgdaSpace{}%
\AgdaInductiveConstructor{S}\AgdaSpace{}%
\AgdaInductiveConstructor{Sb}\AgdaSpace{}%
\AgdaSymbol{(λ}\AgdaSpace{}%
\AgdaBound{x₁}\AgdaSpace{}%
\AgdaSymbol{→}\AgdaSpace{}%
\AgdaInductiveConstructor{refl}\AgdaSymbol{)}\AgdaSpace{}%
\AgdaSymbol{(}\AgdaBound{f'}\AgdaSpace{}%
\AgdaBound{x}\AgdaSymbol{)}\AgdaSpace{}%
\AgdaFunction{⟩}\<%
\\
\>[2][@{}l@{\AgdaIndent{0}}]%
\>[4]\AgdaFunction{embedTerm}\AgdaSpace{}%
\AgdaSymbol{(}\AgdaFunction{rename}\AgdaSpace{}%
\AgdaInductiveConstructor{S}\AgdaSpace{}%
\AgdaSymbol{(}\AgdaBound{f'}\AgdaSpace{}%
\AgdaBound{x}\AgdaSymbol{))}\<%
\\
%
\>[2]\AgdaOperator{\AgdaFunction{∎eq}}\<%
\\
%
\\[\AgdaEmptyExtraSkip]%
\>[0]\AgdaComment{--\ Extsb\ preserves\ pointwise\ equality}\<%
\\
\>[0]\AgdaComment{--\ of\ renaming\ functions}\<%
\\
\>[0]\AgdaFunction{extsb-cong}\AgdaSpace{}%
\AgdaSymbol{:}\AgdaSpace{}%
\AgdaSymbol{∀}\AgdaSpace{}%
\AgdaSymbol{\{}\AgdaBound{Γ₁}\AgdaSpace{}%
\AgdaBound{Γ₂}\AgdaSpace{}%
\AgdaBound{A}\AgdaSpace{}%
\AgdaBound{A'}\AgdaSymbol{\}}\<%
\\
\>[0][@{}l@{\AgdaIndent{0}}]%
\>[2]\AgdaSymbol{→}\AgdaSpace{}%
\AgdaSymbol{\{}\AgdaBound{fb'}\AgdaSpace{}%
\AgdaSymbol{:}\AgdaSpace{}%
\AgdaSymbol{∀}\AgdaSpace{}%
\AgdaSymbol{\{}\AgdaBound{Ab}\AgdaSymbol{\}}\AgdaSpace{}%
\AgdaSymbol{→}\AgdaSpace{}%
\AgdaInductiveConstructor{□}\AgdaSpace{}%
\AgdaBound{Ab}\AgdaSpace{}%
\AgdaOperator{\AgdaDatatype{∈b}}\AgdaSpace{}%
\AgdaBound{Γ₁}\AgdaSpace{}%
\AgdaSymbol{→}\AgdaSpace{}%
\AgdaBound{Γ₂}\AgdaSpace{}%
\AgdaOperator{\AgdaDatatype{⊢b}}\AgdaSpace{}%
\AgdaBound{Ab}\AgdaSymbol{\}}\<%
\\
%
\>[2]\AgdaSymbol{→}\AgdaSpace{}%
\AgdaSymbol{\{}\AgdaBound{fb''}\AgdaSpace{}%
\AgdaSymbol{:}\AgdaSpace{}%
\AgdaSymbol{∀}\AgdaSpace{}%
\AgdaSymbol{\{}\AgdaBound{Ab}\AgdaSymbol{\}}\AgdaSpace{}%
\AgdaSymbol{→}\AgdaSpace{}%
\AgdaInductiveConstructor{□}\AgdaSpace{}%
\AgdaBound{Ab}\AgdaSpace{}%
\AgdaOperator{\AgdaDatatype{∈b}}\AgdaSpace{}%
\AgdaBound{Γ₁}\AgdaSpace{}%
\AgdaSymbol{→}\AgdaSpace{}%
\AgdaBound{Γ₂}\AgdaSpace{}%
\AgdaOperator{\AgdaDatatype{⊢b}}\AgdaSpace{}%
\AgdaBound{Ab}\AgdaSymbol{\}}\<%
\\
%
\>[2]\AgdaSymbol{→}\AgdaSpace{}%
\AgdaSymbol{(}\AgdaBound{h}\AgdaSpace{}%
\AgdaSymbol{:}\AgdaSpace{}%
\AgdaSymbol{∀}\AgdaSpace{}%
\AgdaSymbol{\{}\AgdaBound{Ab}\AgdaSymbol{\}}\AgdaSpace{}%
\AgdaSymbol{→}\AgdaSpace{}%
\AgdaSymbol{(}\AgdaBound{x}\AgdaSpace{}%
\AgdaSymbol{:}\AgdaSpace{}%
\AgdaInductiveConstructor{□}\AgdaSpace{}%
\AgdaBound{Ab}\AgdaSpace{}%
\AgdaOperator{\AgdaDatatype{∈b}}\AgdaSpace{}%
\AgdaBound{Γ₁}\AgdaSymbol{)}\AgdaSpace{}%
\AgdaSymbol{→}\AgdaSpace{}%
\AgdaBound{fb'}\AgdaSpace{}%
\AgdaBound{x}\AgdaSpace{}%
\AgdaOperator{\AgdaDatatype{≡}}\AgdaSpace{}%
\AgdaBound{fb''}\AgdaSpace{}%
\AgdaBound{x}\AgdaSymbol{)}\<%
\\
%
\>[2]\AgdaSymbol{→}\AgdaSpace{}%
\AgdaSymbol{(}\AgdaBound{x}\AgdaSpace{}%
\AgdaSymbol{:}\AgdaSpace{}%
\AgdaInductiveConstructor{□}\AgdaSpace{}%
\AgdaBound{A}\AgdaSpace{}%
\AgdaOperator{\AgdaDatatype{∈b}}\AgdaSpace{}%
\AgdaBound{Γ₁}\AgdaSpace{}%
\AgdaOperator{\AgdaInductiveConstructor{,b}}\AgdaSpace{}%
\AgdaInductiveConstructor{□}\AgdaSpace{}%
\AgdaBound{A'}\AgdaSymbol{)}\<%
\\
%
\>[2]\AgdaSymbol{→}\AgdaSpace{}%
\AgdaFunction{extsb}\AgdaSpace{}%
\AgdaBound{fb'}\AgdaSpace{}%
\AgdaBound{x}\AgdaSpace{}%
\AgdaOperator{\AgdaDatatype{≡}}\AgdaSpace{}%
\AgdaFunction{extsb}\AgdaSpace{}%
\AgdaBound{fb''}\AgdaSpace{}%
\AgdaBound{x}\<%
\\
\>[0]\AgdaFunction{extsb-cong}\AgdaSpace{}%
\AgdaBound{h}\AgdaSpace{}%
\AgdaInductiveConstructor{Zb}\AgdaSpace{}%
\AgdaSymbol{=}\AgdaSpace{}%
\AgdaInductiveConstructor{refl}\<%
\\
\>[0]\AgdaFunction{extsb-cong}\AgdaSpace{}%
\AgdaBound{h}\AgdaSpace{}%
\AgdaSymbol{(}\AgdaInductiveConstructor{Sb}\AgdaSpace{}%
\AgdaBound{x}\AgdaSymbol{)}\AgdaSpace{}%
\AgdaSymbol{=}\AgdaSpace{}%
\AgdaFunction{cong}\AgdaSpace{}%
\AgdaSymbol{(}\AgdaFunction{renameb}\AgdaSpace{}%
\AgdaInductiveConstructor{Sb}\AgdaSymbol{)}\AgdaSpace{}%
\AgdaSymbol{(}\AgdaBound{h}\AgdaSpace{}%
\AgdaBound{x}\AgdaSymbol{)}\<%
\\
%
\\[\AgdaEmptyExtraSkip]%
\>[0]\AgdaComment{--\ substb\ preserves\ pointwise\ equality}\<%
\\
\>[0]\AgdaComment{--\ of\ renaming\ functions}\<%
\\
\>[0]\AgdaFunction{substb-cong}\AgdaSpace{}%
\AgdaSymbol{:}\AgdaSpace{}%
\AgdaSymbol{∀}\AgdaSpace{}%
\AgdaSymbol{\{}\AgdaBound{Γ₁}\AgdaSpace{}%
\AgdaBound{Γ₂}\AgdaSpace{}%
\AgdaBound{A}\AgdaSymbol{\}}\<%
\\
\>[0][@{}l@{\AgdaIndent{0}}]%
\>[2]\AgdaSymbol{→}\AgdaSpace{}%
\AgdaSymbol{\{}\AgdaBound{fb'}\AgdaSpace{}%
\AgdaSymbol{:}\AgdaSpace{}%
\AgdaSymbol{∀}\AgdaSpace{}%
\AgdaSymbol{\{}\AgdaBound{Ab}\AgdaSymbol{\}}\AgdaSpace{}%
\AgdaSymbol{→}\AgdaSpace{}%
\AgdaInductiveConstructor{□}\AgdaSpace{}%
\AgdaBound{Ab}\AgdaSpace{}%
\AgdaOperator{\AgdaDatatype{∈b}}\AgdaSpace{}%
\AgdaBound{Γ₁}\AgdaSpace{}%
\AgdaSymbol{→}\AgdaSpace{}%
\AgdaBound{Γ₂}\AgdaSpace{}%
\AgdaOperator{\AgdaDatatype{⊢b}}\AgdaSpace{}%
\AgdaBound{Ab}\AgdaSymbol{\}}\<%
\\
%
\>[2]\AgdaSymbol{→}\AgdaSpace{}%
\AgdaSymbol{\{}\AgdaBound{fb''}\AgdaSpace{}%
\AgdaSymbol{:}\AgdaSpace{}%
\AgdaSymbol{∀}\AgdaSpace{}%
\AgdaSymbol{\{}\AgdaBound{Ab}\AgdaSymbol{\}}\AgdaSpace{}%
\AgdaSymbol{→}\AgdaSpace{}%
\AgdaInductiveConstructor{□}\AgdaSpace{}%
\AgdaBound{Ab}\AgdaSpace{}%
\AgdaOperator{\AgdaDatatype{∈b}}\AgdaSpace{}%
\AgdaBound{Γ₁}\AgdaSpace{}%
\AgdaSymbol{→}\AgdaSpace{}%
\AgdaBound{Γ₂}\AgdaSpace{}%
\AgdaOperator{\AgdaDatatype{⊢b}}\AgdaSpace{}%
\AgdaBound{Ab}\AgdaSymbol{\}}\<%
\\
%
\>[2]\AgdaSymbol{→}\AgdaSpace{}%
\AgdaSymbol{(}\AgdaBound{h}\AgdaSpace{}%
\AgdaSymbol{:}\AgdaSpace{}%
\AgdaSymbol{∀}\AgdaSpace{}%
\AgdaSymbol{\{}\AgdaBound{Ab}\AgdaSymbol{\}}\AgdaSpace{}%
\AgdaSymbol{→}\AgdaSpace{}%
\AgdaSymbol{(}\AgdaBound{x}\AgdaSpace{}%
\AgdaSymbol{:}\AgdaSpace{}%
\AgdaInductiveConstructor{□}\AgdaSpace{}%
\AgdaBound{Ab}\AgdaSpace{}%
\AgdaOperator{\AgdaDatatype{∈b}}\AgdaSpace{}%
\AgdaBound{Γ₁}\AgdaSymbol{)}\AgdaSpace{}%
\AgdaSymbol{→}\AgdaSpace{}%
\AgdaBound{fb'}\AgdaSpace{}%
\AgdaBound{x}\AgdaSpace{}%
\AgdaOperator{\AgdaDatatype{≡}}\AgdaSpace{}%
\AgdaBound{fb''}\AgdaSpace{}%
\AgdaBound{x}\AgdaSymbol{)}\<%
\\
%
\>[2]\AgdaSymbol{→}\AgdaSpace{}%
\AgdaSymbol{(}\AgdaBound{t}\AgdaSpace{}%
\AgdaSymbol{:}\AgdaSpace{}%
\AgdaBound{Γ₁}\AgdaSpace{}%
\AgdaOperator{\AgdaDatatype{⊢b}}\AgdaSpace{}%
\AgdaBound{A}\AgdaSymbol{)}\<%
\\
%
\>[2]\AgdaSymbol{→}\AgdaSpace{}%
\AgdaFunction{substb}\AgdaSpace{}%
\AgdaBound{fb'}\AgdaSpace{}%
\AgdaBound{t}\AgdaSpace{}%
\AgdaOperator{\AgdaDatatype{≡}}\AgdaSpace{}%
\AgdaFunction{substb}\AgdaSpace{}%
\AgdaBound{fb''}\AgdaSpace{}%
\AgdaBound{t}\<%
\\
\>[0]\AgdaFunction{substb-cong}\AgdaSpace{}%
\AgdaBound{h}\AgdaSpace{}%
\AgdaSymbol{(}\AgdaInductiveConstructor{ε}\AgdaSpace{}%
\AgdaBound{x}\AgdaSymbol{)}\AgdaSpace{}%
\AgdaSymbol{=}\AgdaSpace{}%
\AgdaBound{h}\AgdaSpace{}%
\AgdaBound{x}\<%
\\
\>[0]\AgdaFunction{substb-cong}\AgdaSpace{}%
\AgdaBound{h}\AgdaSpace{}%
\AgdaSymbol{(}\AgdaInductiveConstructor{box}\AgdaSpace{}%
\AgdaBound{t}\AgdaSymbol{)}\AgdaSpace{}%
\AgdaSymbol{=}\AgdaSpace{}%
\AgdaFunction{cong}\AgdaSpace{}%
\AgdaInductiveConstructor{box}\AgdaSpace{}%
\AgdaSymbol{(}\AgdaFunction{substb-cong}\AgdaSpace{}%
\AgdaBound{h}\AgdaSpace{}%
\AgdaBound{t}\AgdaSymbol{)}\<%
\\
\>[0]\AgdaFunction{substb-cong}\AgdaSpace{}%
\AgdaBound{h}\AgdaSpace{}%
\AgdaSymbol{(}\AgdaInductiveConstructor{ƛb}\AgdaSpace{}%
\AgdaBound{t}\AgdaSymbol{)}\AgdaSpace{}%
\AgdaSymbol{=}\AgdaSpace{}%
\AgdaFunction{cong}\AgdaSpace{}%
\AgdaInductiveConstructor{ƛb}\AgdaSpace{}%
\AgdaSymbol{(}\AgdaFunction{substb-cong}\AgdaSpace{}%
\AgdaSymbol{(}\AgdaFunction{extsb-cong}\AgdaSpace{}%
\AgdaBound{h}\AgdaSymbol{)}\AgdaSpace{}%
\AgdaBound{t}\AgdaSymbol{)}\<%
\\
\>[0]\AgdaFunction{substb-cong}\AgdaSpace{}%
\AgdaBound{h}\AgdaSpace{}%
\AgdaSymbol{(}\AgdaBound{t}\AgdaSpace{}%
\AgdaOperator{\AgdaInductiveConstructor{·b}}\AgdaSpace{}%
\AgdaBound{t₁}\AgdaSymbol{)}\AgdaSpace{}%
\AgdaSymbol{=}\AgdaSpace{}%
\AgdaFunction{cong₂}\AgdaSpace{}%
\AgdaOperator{\AgdaInductiveConstructor{\AgdaUnderscore{}·b\AgdaUnderscore{}}}\AgdaSpace{}%
\AgdaSymbol{(}\AgdaFunction{substb-cong}\AgdaSpace{}%
\AgdaBound{h}\AgdaSpace{}%
\AgdaBound{t}\AgdaSymbol{)}\AgdaSpace{}%
\AgdaSymbol{(}\AgdaFunction{substb-cong}\AgdaSpace{}%
\AgdaBound{h}\AgdaSpace{}%
\AgdaBound{t₁}\AgdaSymbol{)}\<%
\\
%
\\[\AgdaEmptyExtraSkip]%
\>[0]\AgdaComment{--\ Commutativity\ of\ renameb\ and\ extsb:}\<%
\\
\>[0]\AgdaFunction{extsb-rename-commute}\AgdaSpace{}%
\AgdaSymbol{:}\AgdaSpace{}%
\AgdaSymbol{∀}\AgdaSpace{}%
\AgdaSymbol{\{}\AgdaBound{Γ₁}\AgdaSpace{}%
\AgdaBound{Γ₂}\AgdaSpace{}%
\AgdaBound{Γ₃}\AgdaSpace{}%
\AgdaBound{A}\AgdaSpace{}%
\AgdaBound{B}\AgdaSymbol{\}}\<%
\\
\>[0][@{}l@{\AgdaIndent{0}}]%
\>[2]\AgdaSymbol{→}\AgdaSpace{}%
\AgdaSymbol{(}\AgdaBound{fb'}\AgdaSpace{}%
\AgdaSymbol{:}\AgdaSpace{}%
\AgdaSymbol{∀}\AgdaSpace{}%
\AgdaSymbol{\{}\AgdaBound{Ab}\AgdaSymbol{\}}\AgdaSpace{}%
\AgdaSymbol{→}\AgdaSpace{}%
\AgdaBound{Ab}\AgdaSpace{}%
\AgdaOperator{\AgdaDatatype{∈b}}\AgdaSpace{}%
\AgdaBound{Γ₂}\AgdaSpace{}%
\AgdaSymbol{→}\AgdaSpace{}%
\AgdaBound{Ab}\AgdaSpace{}%
\AgdaOperator{\AgdaDatatype{∈b}}\AgdaSpace{}%
\AgdaBound{Γ₃}\AgdaSymbol{)}\<%
\\
%
\>[2]\AgdaSymbol{→}\AgdaSpace{}%
\AgdaSymbol{(}\AgdaBound{fb''}\AgdaSpace{}%
\AgdaSymbol{:}\AgdaSpace{}%
\AgdaSymbol{∀}\AgdaSpace{}%
\AgdaSymbol{\{}\AgdaBound{Ab}\AgdaSymbol{\}}\AgdaSpace{}%
\AgdaSymbol{→}\AgdaSpace{}%
\AgdaInductiveConstructor{□}\AgdaSpace{}%
\AgdaBound{Ab}\AgdaSpace{}%
\AgdaOperator{\AgdaDatatype{∈b}}\AgdaSpace{}%
\AgdaBound{Γ₁}\AgdaSpace{}%
\AgdaSymbol{→}\AgdaSpace{}%
\AgdaBound{Γ₂}\AgdaSpace{}%
\AgdaOperator{\AgdaDatatype{⊢b}}\AgdaSpace{}%
\AgdaBound{Ab}\AgdaSymbol{)}\<%
\\
%
\>[2]\AgdaSymbol{→}\AgdaSpace{}%
\AgdaSymbol{(}\AgdaBound{x}\AgdaSpace{}%
\AgdaSymbol{:}\AgdaSpace{}%
\AgdaInductiveConstructor{□}\AgdaSpace{}%
\AgdaBound{A}\AgdaSpace{}%
\AgdaOperator{\AgdaDatatype{∈b}}\AgdaSpace{}%
\AgdaBound{Γ₁}\AgdaSpace{}%
\AgdaOperator{\AgdaInductiveConstructor{,b}}\AgdaSpace{}%
\AgdaBound{B}\AgdaSymbol{)}\<%
\\
%
\>[2]\AgdaSymbol{→}%
\>[1170I]\AgdaFunction{extsb}\AgdaSpace{}%
\AgdaSymbol{(λ}\AgdaSpace{}%
\AgdaBound{x₁}\AgdaSpace{}%
\AgdaSymbol{→}\AgdaSpace{}%
\AgdaFunction{renameb}\AgdaSpace{}%
\AgdaBound{fb'}\AgdaSpace{}%
\AgdaSymbol{(}\AgdaBound{fb''}\AgdaSpace{}%
\AgdaBound{x₁}\AgdaSymbol{))}\AgdaSpace{}%
\AgdaBound{x}\AgdaSpace{}%
\AgdaOperator{\AgdaDatatype{≡}}\<%
\\
\>[.][@{}l@{}]\<[1170I]%
\>[4]\AgdaFunction{renameb}\AgdaSpace{}%
\AgdaSymbol{(}\AgdaFunction{extb}\AgdaSpace{}%
\AgdaBound{fb'}\AgdaSymbol{)}\AgdaSpace{}%
\AgdaSymbol{(}\AgdaFunction{extsb}\AgdaSpace{}%
\AgdaBound{fb''}\AgdaSpace{}%
\AgdaBound{x}\AgdaSymbol{)}\<%
\\
\>[0]\AgdaFunction{extsb-rename-commute}\AgdaSpace{}%
\AgdaBound{fb'}\AgdaSpace{}%
\AgdaBound{fb''}\AgdaSpace{}%
\AgdaInductiveConstructor{Zb}\AgdaSpace{}%
\AgdaSymbol{=}\AgdaSpace{}%
\AgdaInductiveConstructor{refl}\<%
\\
\>[0]\AgdaFunction{extsb-rename-commute}\AgdaSpace{}%
\AgdaBound{fb'}\AgdaSpace{}%
\AgdaBound{fb''}\AgdaSpace{}%
\AgdaSymbol{(}\AgdaInductiveConstructor{Sb}\AgdaSpace{}%
\AgdaBound{x}\AgdaSymbol{)}\AgdaSpace{}%
\AgdaSymbol{=}\AgdaSpace{}%
\AgdaFunction{sym}\AgdaSpace{}%
\AgdaSymbol{(}\AgdaFunction{rename-extb-Sb-commute}\AgdaSpace{}%
\AgdaBound{fb'}\AgdaSpace{}%
\AgdaSymbol{(}\AgdaBound{fb''}\AgdaSpace{}%
\AgdaBound{x}\AgdaSymbol{))}\<%
\\
%
\\[\AgdaEmptyExtraSkip]%
\>[0]\AgdaComment{--\ Commutativity\ of\ substb\ and\ renameb}\<%
\\
\>[0]\AgdaFunction{substb-rename-commute}\AgdaSpace{}%
\AgdaSymbol{:}\AgdaSpace{}%
\AgdaSymbol{∀}\AgdaSpace{}%
\AgdaSymbol{\{}\AgdaBound{Γ₁}\AgdaSpace{}%
\AgdaBound{Γ₂}\AgdaSpace{}%
\AgdaBound{Γ₃}\AgdaSpace{}%
\AgdaBound{A}\AgdaSymbol{\}}\<%
\\
\>[0][@{}l@{\AgdaIndent{0}}]%
\>[2]\AgdaSymbol{→}\AgdaSpace{}%
\AgdaSymbol{(}\AgdaBound{fb'}\AgdaSpace{}%
\AgdaSymbol{:}\AgdaSpace{}%
\AgdaSymbol{∀}\AgdaSpace{}%
\AgdaSymbol{\{}\AgdaBound{Ab}\AgdaSymbol{\}}\AgdaSpace{}%
\AgdaSymbol{→}\AgdaSpace{}%
\AgdaBound{Ab}\AgdaSpace{}%
\AgdaOperator{\AgdaDatatype{∈b}}\AgdaSpace{}%
\AgdaBound{Γ₂}\AgdaSpace{}%
\AgdaSymbol{→}\AgdaSpace{}%
\AgdaBound{Ab}\AgdaSpace{}%
\AgdaOperator{\AgdaDatatype{∈b}}\AgdaSpace{}%
\AgdaBound{Γ₃}\AgdaSymbol{)}\<%
\\
%
\>[2]\AgdaSymbol{→}\AgdaSpace{}%
\AgdaSymbol{(}\AgdaBound{fb''}\AgdaSpace{}%
\AgdaSymbol{:}\AgdaSpace{}%
\AgdaSymbol{∀}\AgdaSpace{}%
\AgdaSymbol{\{}\AgdaBound{Ab}\AgdaSymbol{\}}\AgdaSpace{}%
\AgdaSymbol{→}\AgdaSpace{}%
\AgdaInductiveConstructor{□}\AgdaSpace{}%
\AgdaBound{Ab}\AgdaSpace{}%
\AgdaOperator{\AgdaDatatype{∈b}}\AgdaSpace{}%
\AgdaBound{Γ₁}\AgdaSpace{}%
\AgdaSymbol{→}\AgdaSpace{}%
\AgdaBound{Γ₂}\AgdaSpace{}%
\AgdaOperator{\AgdaDatatype{⊢b}}\AgdaSpace{}%
\AgdaBound{Ab}\AgdaSymbol{)}\<%
\\
%
\>[2]\AgdaSymbol{→}\AgdaSpace{}%
\AgdaSymbol{(}\AgdaBound{t}\AgdaSpace{}%
\AgdaSymbol{:}\AgdaSpace{}%
\AgdaBound{Γ₁}\AgdaSpace{}%
\AgdaOperator{\AgdaDatatype{⊢b}}\AgdaSpace{}%
\AgdaBound{A}\AgdaSymbol{)}\<%
\\
%
\>[2]\AgdaSymbol{→}\AgdaSpace{}%
\AgdaFunction{substb}\AgdaSpace{}%
\AgdaSymbol{(λ}\AgdaSpace{}%
\AgdaBound{x}\AgdaSpace{}%
\AgdaSymbol{→}\AgdaSpace{}%
\AgdaFunction{renameb}\AgdaSpace{}%
\AgdaBound{fb'}\AgdaSpace{}%
\AgdaSymbol{(}\AgdaBound{fb''}\AgdaSpace{}%
\AgdaBound{x}\AgdaSymbol{))}\AgdaSpace{}%
\AgdaBound{t}\AgdaSpace{}%
\AgdaOperator{\AgdaDatatype{≡}}\AgdaSpace{}%
\AgdaFunction{renameb}\AgdaSpace{}%
\AgdaBound{fb'}\AgdaSpace{}%
\AgdaSymbol{(}\AgdaFunction{substb}\AgdaSpace{}%
\AgdaBound{fb''}\AgdaSpace{}%
\AgdaBound{t}\AgdaSymbol{)}\<%
\\
\>[0]\AgdaFunction{substb-rename-commute}\AgdaSpace{}%
\AgdaBound{fb'}\AgdaSpace{}%
\AgdaBound{fb''}\AgdaSpace{}%
\AgdaSymbol{(}\AgdaInductiveConstructor{ε}\AgdaSpace{}%
\AgdaBound{x}\AgdaSymbol{)}\AgdaSpace{}%
\AgdaSymbol{=}\AgdaSpace{}%
\AgdaInductiveConstructor{refl}\<%
\\
\>[0]\AgdaFunction{substb-rename-commute}\AgdaSpace{}%
\AgdaBound{fb'}\AgdaSpace{}%
\AgdaBound{fb''}\AgdaSpace{}%
\AgdaSymbol{(}\AgdaInductiveConstructor{box}\AgdaSpace{}%
\AgdaBound{t}\AgdaSymbol{)}\AgdaSpace{}%
\AgdaSymbol{=}\<%
\\
\>[0][@{}l@{\AgdaIndent{0}}]%
\>[2]\AgdaFunction{cong}\<%
\\
\>[2][@{}l@{\AgdaIndent{0}}]%
\>[4]\AgdaInductiveConstructor{box}\<%
\\
%
\>[4]\AgdaSymbol{(}\AgdaFunction{substb-rename-commute}\AgdaSpace{}%
\AgdaBound{fb'}\AgdaSpace{}%
\AgdaBound{fb''}\AgdaSpace{}%
\AgdaBound{t}\AgdaSymbol{)}\<%
\\
\>[0]\AgdaFunction{substb-rename-commute}\AgdaSpace{}%
\AgdaBound{fb'}\AgdaSpace{}%
\AgdaBound{fb''}\AgdaSpace{}%
\AgdaSymbol{(}\AgdaInductiveConstructor{ƛb}\AgdaSpace{}%
\AgdaBound{t}\AgdaSymbol{)}\AgdaSpace{}%
\AgdaSymbol{=}\<%
\\
\>[0][@{}l@{\AgdaIndent{0}}]%
\>[2]\AgdaOperator{\AgdaFunction{beginEq}}\<%
\\
\>[2][@{}l@{\AgdaIndent{0}}]%
\>[4]\AgdaInductiveConstructor{ƛb}\AgdaSpace{}%
\AgdaSymbol{(}\AgdaFunction{substb}\AgdaSpace{}%
\AgdaSymbol{(}\AgdaFunction{extsb}\AgdaSpace{}%
\AgdaSymbol{(λ}\AgdaSpace{}%
\AgdaBound{x}\AgdaSpace{}%
\AgdaSymbol{→}\AgdaSpace{}%
\AgdaFunction{renameb}\AgdaSpace{}%
\AgdaBound{fb'}\AgdaSpace{}%
\AgdaSymbol{(}\AgdaBound{fb''}\AgdaSpace{}%
\AgdaBound{x}\AgdaSymbol{)))}\AgdaSpace{}%
\AgdaBound{t}\AgdaSymbol{)}\<%
\\
%
\>[2]\AgdaFunction{≡⟨}\AgdaSpace{}%
\AgdaFunction{cong}\AgdaSpace{}%
\AgdaInductiveConstructor{ƛb}\AgdaSpace{}%
\AgdaSymbol{(}\AgdaFunction{substb-cong}\AgdaSpace{}%
\AgdaSymbol{(λ}\AgdaSpace{}%
\AgdaBound{x}\AgdaSpace{}%
\AgdaSymbol{→}\AgdaSpace{}%
\AgdaFunction{extsb-rename-commute}\AgdaSpace{}%
\AgdaBound{fb'}\AgdaSpace{}%
\AgdaBound{fb''}\AgdaSpace{}%
\AgdaBound{x}\AgdaSymbol{)}\AgdaSpace{}%
\AgdaBound{t}\AgdaSymbol{)}\AgdaSpace{}%
\AgdaFunction{⟩}\<%
\\
\>[2][@{}l@{\AgdaIndent{0}}]%
\>[4]\AgdaInductiveConstructor{ƛb}\AgdaSpace{}%
\AgdaSymbol{(}\AgdaFunction{substb}\AgdaSpace{}%
\AgdaSymbol{(λ}\AgdaSpace{}%
\AgdaBound{x}\AgdaSpace{}%
\AgdaSymbol{→}\AgdaSpace{}%
\AgdaFunction{renameb}\AgdaSpace{}%
\AgdaSymbol{(}\AgdaFunction{extb}\AgdaSpace{}%
\AgdaBound{fb'}\AgdaSymbol{)}\AgdaSpace{}%
\AgdaSymbol{((}\AgdaFunction{extsb}\AgdaSpace{}%
\AgdaBound{fb''}\AgdaSymbol{)}\AgdaSpace{}%
\AgdaBound{x}\AgdaSymbol{))}\AgdaSpace{}%
\AgdaBound{t}\AgdaSymbol{)}\<%
\\
%
\>[2]\AgdaFunction{≡⟨}\AgdaSpace{}%
\AgdaFunction{cong}\AgdaSpace{}%
\AgdaInductiveConstructor{ƛb}\AgdaSpace{}%
\AgdaSymbol{(}\AgdaFunction{substb-rename-commute}\AgdaSpace{}%
\AgdaSymbol{(}\AgdaFunction{extb}\AgdaSpace{}%
\AgdaBound{fb'}\AgdaSymbol{)}\AgdaSpace{}%
\AgdaSymbol{(}\AgdaFunction{extsb}\AgdaSpace{}%
\AgdaBound{fb''}\AgdaSymbol{)}\AgdaSpace{}%
\AgdaBound{t}\AgdaSymbol{)}\AgdaSpace{}%
\AgdaFunction{⟩}\<%
\\
\>[2][@{}l@{\AgdaIndent{0}}]%
\>[4]\AgdaInductiveConstructor{ƛb}\AgdaSpace{}%
\AgdaSymbol{(}\AgdaFunction{renameb}\AgdaSpace{}%
\AgdaSymbol{(}\AgdaFunction{extb}\AgdaSpace{}%
\AgdaBound{fb'}\AgdaSymbol{)}\AgdaSpace{}%
\AgdaSymbol{(}\AgdaFunction{substb}\AgdaSpace{}%
\AgdaSymbol{(}\AgdaFunction{extsb}\AgdaSpace{}%
\AgdaBound{fb''}\AgdaSymbol{)}\AgdaSpace{}%
\AgdaBound{t}\AgdaSymbol{))}\<%
\\
%
\>[2]\AgdaOperator{\AgdaFunction{∎eq}}\<%
\\
\>[0]\AgdaFunction{substb-rename-commute}\AgdaSpace{}%
\AgdaBound{fb'}\AgdaSpace{}%
\AgdaBound{fb''}\AgdaSpace{}%
\AgdaSymbol{(}\AgdaBound{t}\AgdaSpace{}%
\AgdaOperator{\AgdaInductiveConstructor{·b}}\AgdaSpace{}%
\AgdaBound{t₁}\AgdaSymbol{)}\AgdaSpace{}%
\AgdaSymbol{=}\<%
\\
\>[0][@{}l@{\AgdaIndent{0}}]%
\>[2]\AgdaFunction{cong₂}\<%
\\
\>[2][@{}l@{\AgdaIndent{0}}]%
\>[4]\AgdaOperator{\AgdaInductiveConstructor{\AgdaUnderscore{}·b\AgdaUnderscore{}}}\<%
\\
%
\>[4]\AgdaSymbol{(}\AgdaFunction{substb-rename-commute}\AgdaSpace{}%
\AgdaBound{fb'}\AgdaSpace{}%
\AgdaBound{fb''}\AgdaSpace{}%
\AgdaBound{t}\AgdaSymbol{)}\<%
\\
%
\>[4]\AgdaSymbol{(}\AgdaFunction{substb-rename-commute}\AgdaSpace{}%
\AgdaBound{fb'}\AgdaSpace{}%
\AgdaBound{fb''}\AgdaSpace{}%
\AgdaBound{t₁}\AgdaSymbol{)}\<%
\\
%
\\[\AgdaEmptyExtraSkip]%
\>[0]\AgdaComment{--\ extsb\ distributes\ over\ composition:}\<%
\\
\>[0]\AgdaComment{--\ Executing\ extsb\ first\ with\ fb',\ then\ with\ fb''}\<%
\\
\>[0]\AgdaComment{--\ is\ the\ same\ as\ doing\ extsb\ once\ with\ fb''\ ∘\ fb'.}\<%
\\
\>[0]\AgdaFunction{extsb-comp}\AgdaSpace{}%
\AgdaSymbol{:}\AgdaSpace{}%
\AgdaSymbol{∀}\AgdaSpace{}%
\AgdaSymbol{\{}\AgdaBound{Γ₁}\AgdaSpace{}%
\AgdaBound{Γ₂}\AgdaSpace{}%
\AgdaBound{Γ₃}\AgdaSpace{}%
\AgdaBound{A}\AgdaSpace{}%
\AgdaBound{B}\AgdaSymbol{\}}\<%
\\
\>[0][@{}l@{\AgdaIndent{0}}]%
\>[2]\AgdaSymbol{→}\AgdaSpace{}%
\AgdaSymbol{(}\AgdaBound{fb'}\AgdaSpace{}%
\AgdaSymbol{:}\AgdaSpace{}%
\AgdaSymbol{∀}\AgdaSpace{}%
\AgdaSymbol{\{}\AgdaBound{Ab}\AgdaSymbol{\}}\AgdaSpace{}%
\AgdaSymbol{→}\AgdaSpace{}%
\AgdaBound{Ab}\AgdaSpace{}%
\AgdaOperator{\AgdaDatatype{∈b}}\AgdaSpace{}%
\AgdaBound{Γ₁}\AgdaSpace{}%
\AgdaSymbol{→}\AgdaSpace{}%
\AgdaBound{Ab}\AgdaSpace{}%
\AgdaOperator{\AgdaDatatype{∈b}}\AgdaSpace{}%
\AgdaBound{Γ₂}\AgdaSymbol{)}\<%
\\
%
\>[2]\AgdaSymbol{→}\AgdaSpace{}%
\AgdaSymbol{(}\AgdaBound{fb''}\AgdaSpace{}%
\AgdaSymbol{:}\AgdaSpace{}%
\AgdaSymbol{∀}\AgdaSpace{}%
\AgdaSymbol{\{}\AgdaBound{Ab}\AgdaSymbol{\}}\AgdaSpace{}%
\AgdaSymbol{→}\AgdaSpace{}%
\AgdaInductiveConstructor{□}\AgdaSpace{}%
\AgdaBound{Ab}\AgdaSpace{}%
\AgdaOperator{\AgdaDatatype{∈b}}\AgdaSpace{}%
\AgdaBound{Γ₂}\AgdaSpace{}%
\AgdaSymbol{→}\AgdaSpace{}%
\AgdaBound{Γ₃}\AgdaSpace{}%
\AgdaOperator{\AgdaDatatype{⊢b}}\AgdaSpace{}%
\AgdaBound{Ab}\AgdaSymbol{)}\<%
\\
%
\>[2]\AgdaSymbol{→}\AgdaSpace{}%
\AgdaSymbol{(}\AgdaBound{x}\AgdaSpace{}%
\AgdaSymbol{:}\AgdaSpace{}%
\AgdaInductiveConstructor{□}\AgdaSpace{}%
\AgdaBound{A}\AgdaSpace{}%
\AgdaOperator{\AgdaDatatype{∈b}}\AgdaSpace{}%
\AgdaBound{Γ₁}\AgdaSpace{}%
\AgdaOperator{\AgdaInductiveConstructor{,b}}\AgdaSpace{}%
\AgdaBound{B}\AgdaSymbol{)}\<%
\\
%
\>[2]\AgdaSymbol{→}\AgdaSpace{}%
\AgdaFunction{extsb}\AgdaSpace{}%
\AgdaBound{fb''}\AgdaSpace{}%
\AgdaSymbol{(}\AgdaFunction{extb}\AgdaSpace{}%
\AgdaBound{fb'}\AgdaSpace{}%
\AgdaBound{x}\AgdaSymbol{)}\AgdaSpace{}%
\AgdaOperator{\AgdaDatatype{≡}}\AgdaSpace{}%
\AgdaFunction{extsb}\AgdaSpace{}%
\AgdaSymbol{(λ}\AgdaSpace{}%
\AgdaBound{x₁}\AgdaSpace{}%
\AgdaSymbol{→}\AgdaSpace{}%
\AgdaBound{fb''}\AgdaSpace{}%
\AgdaSymbol{(}\AgdaBound{fb'}\AgdaSpace{}%
\AgdaBound{x₁}\AgdaSymbol{))}\AgdaSpace{}%
\AgdaBound{x}\<%
\\
\>[0]\AgdaFunction{extsb-comp}\AgdaSpace{}%
\AgdaBound{fb'}\AgdaSpace{}%
\AgdaBound{fb''}\AgdaSpace{}%
\AgdaInductiveConstructor{Zb}\AgdaSpace{}%
\AgdaSymbol{=}\AgdaSpace{}%
\AgdaInductiveConstructor{refl}\<%
\\
\>[0]\AgdaFunction{extsb-comp}\AgdaSpace{}%
\AgdaBound{fb'}\AgdaSpace{}%
\AgdaBound{fb''}\AgdaSpace{}%
\AgdaSymbol{(}\AgdaInductiveConstructor{Sb}\AgdaSpace{}%
\AgdaBound{x}\AgdaSymbol{)}\AgdaSpace{}%
\AgdaSymbol{=}\AgdaSpace{}%
\AgdaInductiveConstructor{refl}\<%
\\
%
\\[\AgdaEmptyExtraSkip]%
\>[0]\AgdaComment{--\ substb\ and\ rename\ distribute\ over\ composition:}\<%
\\
\>[0]\AgdaComment{--\ Renaming\ first\ with\ fb'\ and\ substituting\ with\ fb''}\<%
\\
\>[0]\AgdaComment{--\ is\ the\ same\ as\ substituting\ with\ fb''\ ∘\ fb'}\<%
\\
\>[0]\AgdaFunction{subst-renameb-comp}\AgdaSpace{}%
\AgdaSymbol{:}\AgdaSpace{}%
\AgdaSymbol{∀}\AgdaSpace{}%
\AgdaSymbol{\{}\AgdaBound{Γ₁}\AgdaSpace{}%
\AgdaBound{Γ₂}\AgdaSpace{}%
\AgdaBound{Γ₃}\AgdaSpace{}%
\AgdaBound{A}\AgdaSymbol{\}}\<%
\\
\>[0][@{}l@{\AgdaIndent{0}}]%
\>[2]\AgdaSymbol{→}\AgdaSpace{}%
\AgdaSymbol{(}\AgdaBound{fb'}\AgdaSpace{}%
\AgdaSymbol{:}\AgdaSpace{}%
\AgdaSymbol{∀}\AgdaSpace{}%
\AgdaSymbol{\{}\AgdaBound{Ab}\AgdaSymbol{\}}\AgdaSpace{}%
\AgdaSymbol{→}\AgdaSpace{}%
\AgdaBound{Ab}\AgdaSpace{}%
\AgdaOperator{\AgdaDatatype{∈b}}\AgdaSpace{}%
\AgdaBound{Γ₁}\AgdaSpace{}%
\AgdaSymbol{→}\AgdaSpace{}%
\AgdaBound{Ab}\AgdaSpace{}%
\AgdaOperator{\AgdaDatatype{∈b}}\AgdaSpace{}%
\AgdaBound{Γ₂}\AgdaSymbol{)}\<%
\\
%
\>[2]\AgdaSymbol{→}\AgdaSpace{}%
\AgdaSymbol{(}\AgdaBound{fb''}\AgdaSpace{}%
\AgdaSymbol{:}\AgdaSpace{}%
\AgdaSymbol{∀}\AgdaSpace{}%
\AgdaSymbol{\{}\AgdaBound{Ab}\AgdaSymbol{\}}\AgdaSpace{}%
\AgdaSymbol{→}\AgdaSpace{}%
\AgdaInductiveConstructor{□}\AgdaSpace{}%
\AgdaBound{Ab}\AgdaSpace{}%
\AgdaOperator{\AgdaDatatype{∈b}}\AgdaSpace{}%
\AgdaBound{Γ₂}\AgdaSpace{}%
\AgdaSymbol{→}\AgdaSpace{}%
\AgdaBound{Γ₃}\AgdaSpace{}%
\AgdaOperator{\AgdaDatatype{⊢b}}\AgdaSpace{}%
\AgdaBound{Ab}\AgdaSymbol{)}\<%
\\
%
\>[2]\AgdaSymbol{→}\AgdaSpace{}%
\AgdaSymbol{(}\AgdaBound{t}\AgdaSpace{}%
\AgdaSymbol{:}\AgdaSpace{}%
\AgdaBound{Γ₁}\AgdaSpace{}%
\AgdaOperator{\AgdaDatatype{⊢b}}\AgdaSpace{}%
\AgdaBound{A}\AgdaSymbol{)}\<%
\\
%
\>[2]\AgdaSymbol{→}\AgdaSpace{}%
\AgdaFunction{substb}\AgdaSpace{}%
\AgdaBound{fb''}\AgdaSpace{}%
\AgdaSymbol{(}\AgdaFunction{renameb}\AgdaSpace{}%
\AgdaBound{fb'}\AgdaSpace{}%
\AgdaBound{t}\AgdaSymbol{)}\AgdaSpace{}%
\AgdaOperator{\AgdaDatatype{≡}}\AgdaSpace{}%
\AgdaFunction{substb}\AgdaSpace{}%
\AgdaSymbol{(λ}\AgdaSpace{}%
\AgdaBound{x}\AgdaSpace{}%
\AgdaSymbol{→}\AgdaSpace{}%
\AgdaBound{fb''}\AgdaSpace{}%
\AgdaSymbol{(}\AgdaBound{fb'}\AgdaSpace{}%
\AgdaBound{x}\AgdaSymbol{))}\AgdaSpace{}%
\AgdaBound{t}\<%
\\
\>[0]\AgdaFunction{subst-renameb-comp}\AgdaSpace{}%
\AgdaBound{fb'}\AgdaSpace{}%
\AgdaBound{fb''}\AgdaSpace{}%
\AgdaSymbol{(}\AgdaInductiveConstructor{ε}\AgdaSpace{}%
\AgdaBound{x}\AgdaSymbol{)}\AgdaSpace{}%
\AgdaSymbol{=}\AgdaSpace{}%
\AgdaInductiveConstructor{refl}\<%
\\
\>[0]\AgdaFunction{subst-renameb-comp}\AgdaSpace{}%
\AgdaBound{fb'}\AgdaSpace{}%
\AgdaBound{fb''}\AgdaSpace{}%
\AgdaSymbol{(}\AgdaInductiveConstructor{box}\AgdaSpace{}%
\AgdaBound{t}\AgdaSymbol{)}\AgdaSpace{}%
\AgdaSymbol{=}\AgdaSpace{}%
\AgdaFunction{cong}\AgdaSpace{}%
\AgdaInductiveConstructor{box}\AgdaSpace{}%
\AgdaSymbol{(}\AgdaFunction{subst-renameb-comp}\AgdaSpace{}%
\AgdaBound{fb'}\AgdaSpace{}%
\AgdaBound{fb''}\AgdaSpace{}%
\AgdaBound{t}\AgdaSymbol{)}\<%
\\
\>[0]\AgdaFunction{subst-renameb-comp}\AgdaSpace{}%
\AgdaBound{fb'}\AgdaSpace{}%
\AgdaBound{fb''}\AgdaSpace{}%
\AgdaSymbol{(}\AgdaInductiveConstructor{ƛb}\AgdaSpace{}%
\AgdaBound{t}\AgdaSymbol{)}\AgdaSpace{}%
\AgdaSymbol{=}\<%
\\
\>[0][@{}l@{\AgdaIndent{0}}]%
\>[2]\AgdaOperator{\AgdaFunction{beginEq}}\<%
\\
\>[2][@{}l@{\AgdaIndent{0}}]%
\>[4]\AgdaInductiveConstructor{ƛb}\AgdaSpace{}%
\AgdaSymbol{(}\AgdaFunction{substb}\AgdaSpace{}%
\AgdaSymbol{(}\AgdaFunction{extsb}\AgdaSpace{}%
\AgdaBound{fb''}\AgdaSymbol{)}\AgdaSpace{}%
\AgdaSymbol{(}\AgdaFunction{renameb}\AgdaSpace{}%
\AgdaSymbol{(}\AgdaFunction{extb}\AgdaSpace{}%
\AgdaBound{fb'}\AgdaSymbol{)}\AgdaSpace{}%
\AgdaBound{t}\AgdaSymbol{))}\<%
\\
%
\>[2]\AgdaFunction{≡⟨}\AgdaSpace{}%
\AgdaFunction{cong}\AgdaSpace{}%
\AgdaInductiveConstructor{ƛb}\AgdaSpace{}%
\AgdaSymbol{(}\AgdaFunction{subst-renameb-comp}\AgdaSpace{}%
\AgdaSymbol{(}\AgdaFunction{extb}\AgdaSpace{}%
\AgdaBound{fb'}\AgdaSymbol{)}\AgdaSpace{}%
\AgdaSymbol{(}\AgdaFunction{extsb}\AgdaSpace{}%
\AgdaBound{fb''}\AgdaSymbol{)}\AgdaSpace{}%
\AgdaBound{t}\AgdaSymbol{)}\AgdaSpace{}%
\AgdaFunction{⟩}\<%
\\
\>[2][@{}l@{\AgdaIndent{0}}]%
\>[4]\AgdaInductiveConstructor{ƛb}\AgdaSpace{}%
\AgdaSymbol{(}\AgdaFunction{substb}\AgdaSpace{}%
\AgdaSymbol{(λ}\AgdaSpace{}%
\AgdaBound{x}\AgdaSpace{}%
\AgdaSymbol{→}\AgdaSpace{}%
\AgdaSymbol{(}\AgdaFunction{extsb}\AgdaSpace{}%
\AgdaBound{fb''}\AgdaSymbol{)}\AgdaSpace{}%
\AgdaSymbol{((}\AgdaFunction{extb}\AgdaSpace{}%
\AgdaBound{fb'}\AgdaSymbol{)}\AgdaSpace{}%
\AgdaBound{x}\AgdaSymbol{))}\AgdaSpace{}%
\AgdaBound{t}\AgdaSymbol{)}\<%
\\
%
\>[2]\AgdaFunction{≡⟨}\AgdaSpace{}%
\AgdaFunction{cong}\AgdaSpace{}%
\AgdaInductiveConstructor{ƛb}\AgdaSpace{}%
\AgdaSymbol{(}\AgdaFunction{substb-cong}\AgdaSpace{}%
\AgdaSymbol{(λ}\AgdaSpace{}%
\AgdaBound{x}\AgdaSpace{}%
\AgdaSymbol{→}\AgdaSpace{}%
\AgdaFunction{extsb-comp}\AgdaSpace{}%
\AgdaBound{fb'}\AgdaSpace{}%
\AgdaBound{fb''}\AgdaSpace{}%
\AgdaBound{x}\AgdaSymbol{)}\AgdaSpace{}%
\AgdaBound{t}\AgdaSymbol{)}\AgdaSpace{}%
\AgdaFunction{⟩}\<%
\\
\>[2][@{}l@{\AgdaIndent{0}}]%
\>[4]\AgdaInductiveConstructor{ƛb}\AgdaSpace{}%
\AgdaSymbol{(}\AgdaFunction{substb}\AgdaSpace{}%
\AgdaSymbol{(}\AgdaFunction{extsb}\AgdaSpace{}%
\AgdaSymbol{(λ}\AgdaSpace{}%
\AgdaBound{x}\AgdaSpace{}%
\AgdaSymbol{→}\AgdaSpace{}%
\AgdaBound{fb''}\AgdaSpace{}%
\AgdaSymbol{(}\AgdaBound{fb'}\AgdaSpace{}%
\AgdaBound{x}\AgdaSymbol{)))}\AgdaSpace{}%
\AgdaBound{t}\AgdaSymbol{)}\<%
\\
%
\>[2]\AgdaOperator{\AgdaFunction{∎eq}}\<%
\\
\>[0]\AgdaFunction{subst-renameb-comp}\AgdaSpace{}%
\AgdaBound{fb'}\AgdaSpace{}%
\AgdaBound{fb''}\AgdaSpace{}%
\AgdaSymbol{(}\AgdaBound{t}\AgdaSpace{}%
\AgdaOperator{\AgdaInductiveConstructor{·b}}\AgdaSpace{}%
\AgdaBound{t₁}\AgdaSymbol{)}\AgdaSpace{}%
\AgdaSymbol{=}\AgdaSpace{}%
\AgdaFunction{cong₂}\AgdaSpace{}%
\AgdaOperator{\AgdaInductiveConstructor{\AgdaUnderscore{}·b\AgdaUnderscore{}}}\AgdaSpace{}%
\AgdaSymbol{(}\AgdaFunction{subst-renameb-comp}\AgdaSpace{}%
\AgdaBound{fb'}\AgdaSpace{}%
\AgdaBound{fb''}\AgdaSpace{}%
\AgdaBound{t}\AgdaSymbol{)}\AgdaSpace{}%
\AgdaSymbol{(}\AgdaFunction{subst-renameb-comp}\AgdaSpace{}%
\AgdaBound{fb'}\AgdaSpace{}%
\AgdaBound{fb''}\AgdaSpace{}%
\AgdaBound{t₁}\AgdaSymbol{)}\<%
\end{code}
\begin{code}%
\>[0]\AgdaComment{--\ Renaming\ followed\ by\ substitution\ with\ an\ extended\ context\ (extsb)}\<%
\\
\>[0]\AgdaComment{--\ commutes\ with\ substitution\ followed\ by\ renaming.}\<%
\\
\>[0]\AgdaFunction{substb-exts-rename-Sb-commute}\AgdaSpace{}%
\AgdaSymbol{:}\AgdaSpace{}%
\AgdaSymbol{∀}\AgdaSpace{}%
\AgdaSymbol{\{}\AgdaBound{Γ}\AgdaSpace{}%
\AgdaBound{A}\AgdaSpace{}%
\AgdaBound{B}\AgdaSpace{}%
\AgdaBound{Δ}\AgdaSymbol{\}}\<%
\\
\>[0][@{}l@{\AgdaIndent{0}}]%
\>[2]\AgdaSymbol{→}\AgdaSpace{}%
\AgdaSymbol{(}\AgdaBound{fb'}\AgdaSpace{}%
\AgdaSymbol{:}\AgdaSpace{}%
\AgdaSymbol{∀}\AgdaSpace{}%
\AgdaSymbol{\{}\AgdaBound{Ab}\AgdaSymbol{\}}\AgdaSpace{}%
\AgdaSymbol{→}\AgdaSpace{}%
\AgdaInductiveConstructor{□}\AgdaSpace{}%
\AgdaBound{Ab}\AgdaSpace{}%
\AgdaOperator{\AgdaDatatype{∈b}}\AgdaSpace{}%
\AgdaBound{Γ}\AgdaSpace{}%
\AgdaSymbol{→}\AgdaSpace{}%
\AgdaBound{Δ}\AgdaSpace{}%
\AgdaOperator{\AgdaDatatype{⊢b}}\AgdaSpace{}%
\AgdaBound{Ab}\AgdaSymbol{)}\<%
\\
%
\>[2]\AgdaSymbol{→}\AgdaSpace{}%
\AgdaSymbol{(}\AgdaBound{t}\AgdaSpace{}%
\AgdaSymbol{:}\AgdaSpace{}%
\AgdaBound{Γ}\AgdaSpace{}%
\AgdaOperator{\AgdaDatatype{⊢b}}\AgdaSpace{}%
\AgdaBound{A}\AgdaSymbol{)}\<%
\\
%
\>[2]\AgdaSymbol{→}%
\>[1555I]\AgdaFunction{substb}\AgdaSpace{}%
\AgdaSymbol{(}\AgdaFunction{extsb}\AgdaSpace{}%
\AgdaBound{fb'}\AgdaSpace{}%
\AgdaSymbol{\{}\AgdaArgument{B}\AgdaSpace{}%
\AgdaSymbol{=}\AgdaSpace{}%
\AgdaBound{B}\AgdaSymbol{\})}\AgdaSpace{}%
\AgdaSymbol{(}\AgdaFunction{renameb}\AgdaSpace{}%
\AgdaInductiveConstructor{Sb}\AgdaSpace{}%
\AgdaBound{t}\AgdaSymbol{)}\<%
\\
\>[.][@{}l@{}]\<[1555I]%
\>[4]\AgdaOperator{\AgdaDatatype{≡}}\AgdaSpace{}%
\AgdaFunction{renameb}\AgdaSpace{}%
\AgdaInductiveConstructor{Sb}\AgdaSpace{}%
\AgdaSymbol{(}\AgdaFunction{substb}\AgdaSpace{}%
\AgdaBound{fb'}\AgdaSpace{}%
\AgdaBound{t}\AgdaSymbol{)}\<%
\\
\>[0]\AgdaFunction{substb-exts-rename-Sb-commute}\AgdaSpace{}%
\AgdaBound{fb'}\AgdaSpace{}%
\AgdaBound{t}\AgdaSpace{}%
\AgdaSymbol{=}\<%
\\
\>[0][@{}l@{\AgdaIndent{0}}]%
\>[2]\AgdaOperator{\AgdaFunction{beginEq}}\<%
\\
\>[2][@{}l@{\AgdaIndent{0}}]%
\>[4]\AgdaFunction{substb}\AgdaSpace{}%
\AgdaSymbol{(}\AgdaFunction{extsb}\AgdaSpace{}%
\AgdaBound{fb'}\AgdaSymbol{)}\AgdaSpace{}%
\AgdaSymbol{(}\AgdaFunction{renameb}\AgdaSpace{}%
\AgdaInductiveConstructor{Sb}\AgdaSpace{}%
\AgdaBound{t}\AgdaSymbol{)}\<%
\\
%
\>[2]\AgdaFunction{≡⟨}\AgdaSpace{}%
\AgdaFunction{subst-renameb-comp}\AgdaSpace{}%
\AgdaInductiveConstructor{Sb}\AgdaSpace{}%
\AgdaSymbol{(}\AgdaFunction{extsb}\AgdaSpace{}%
\AgdaBound{fb'}\AgdaSymbol{)}\AgdaSpace{}%
\AgdaBound{t}\AgdaSpace{}%
\AgdaFunction{⟩}\<%
\\
\>[2][@{}l@{\AgdaIndent{0}}]%
\>[4]\AgdaFunction{substb}\AgdaSpace{}%
\AgdaSymbol{(λ}\AgdaSpace{}%
\AgdaBound{x}\AgdaSpace{}%
\AgdaSymbol{→}\AgdaSpace{}%
\AgdaFunction{renameb}\AgdaSpace{}%
\AgdaInductiveConstructor{Sb}\AgdaSpace{}%
\AgdaSymbol{(}\AgdaBound{fb'}\AgdaSpace{}%
\AgdaBound{x}\AgdaSymbol{))}\AgdaSpace{}%
\AgdaBound{t}\<%
\\
%
\>[2]\AgdaFunction{≡⟨}\AgdaSpace{}%
\AgdaFunction{substb-rename-commute}\AgdaSpace{}%
\AgdaInductiveConstructor{Sb}\AgdaSpace{}%
\AgdaBound{fb'}\AgdaSpace{}%
\AgdaBound{t}\AgdaSpace{}%
\AgdaFunction{⟩}\<%
\\
\>[2][@{}l@{\AgdaIndent{0}}]%
\>[4]\AgdaFunction{renameb}\AgdaSpace{}%
\AgdaInductiveConstructor{Sb}\AgdaSpace{}%
\AgdaSymbol{(}\AgdaFunction{substb}\AgdaSpace{}%
\AgdaBound{fb'}\AgdaSpace{}%
\AgdaBound{t}\AgdaSymbol{)}\<%
\\
%
\>[2]\AgdaOperator{\AgdaFunction{∎eq}}\<%
\end{code}
\begin{code}[hide]%
\>[0]\AgdaComment{--\ Preservation\ and\ Reflection\ of\ }\<%
\\
\>[0]\AgdaComment{--\ simultaneous\ substitution\ of\ Gödel's\ translation}\<%
\\
\>[0]\AgdaFunction{pres-simultaneous-subst}\AgdaSpace{}%
\AgdaSymbol{:}\AgdaSpace{}%
\AgdaSymbol{∀}\AgdaSpace{}%
\AgdaSymbol{\{}\AgdaBound{Γ}\AgdaSpace{}%
\AgdaBound{Δ}\AgdaSpace{}%
\AgdaBound{A}\AgdaSymbol{\}}\<%
\\
\>[0][@{}l@{\AgdaIndent{0}}]%
\>[2]\AgdaSymbol{→}\AgdaSpace{}%
\AgdaSymbol{(}\AgdaBound{fb'}\AgdaSpace{}%
\AgdaSymbol{:}\AgdaSpace{}%
\AgdaSymbol{∀}\AgdaSpace{}%
\AgdaSymbol{\{}\AgdaBound{Ab}\AgdaSymbol{\}}\AgdaSpace{}%
\AgdaSymbol{→}\AgdaSpace{}%
\AgdaInductiveConstructor{□}\AgdaSpace{}%
\AgdaBound{Ab}\AgdaSpace{}%
\AgdaOperator{\AgdaDatatype{∈b}}\AgdaSpace{}%
\AgdaFunction{embedContext}\AgdaSpace{}%
\AgdaBound{Γ}\AgdaSpace{}%
\AgdaSymbol{→}\AgdaSpace{}%
\AgdaFunction{embedContext}\AgdaSpace{}%
\AgdaBound{Δ}\AgdaSpace{}%
\AgdaOperator{\AgdaDatatype{⊢b}}\AgdaSpace{}%
\AgdaBound{Ab}\AgdaSymbol{)}\<%
\\
%
\>[2]\AgdaSymbol{→}\AgdaSpace{}%
\AgdaSymbol{(}\AgdaBound{f'}\AgdaSpace{}%
\AgdaSymbol{:}\AgdaSpace{}%
\AgdaSymbol{∀}\AgdaSpace{}%
\AgdaSymbol{\{}\AgdaBound{A}\AgdaSymbol{\}}\AgdaSpace{}%
\AgdaSymbol{→}\AgdaSpace{}%
\AgdaBound{A}\AgdaSpace{}%
\AgdaOperator{\AgdaDatatype{∈}}\AgdaSpace{}%
\AgdaBound{Γ}\AgdaSpace{}%
\AgdaSymbol{→}\AgdaSpace{}%
\AgdaBound{Δ}\AgdaSpace{}%
\AgdaOperator{\AgdaDatatype{⊢}}\AgdaSpace{}%
\AgdaBound{A}\AgdaSymbol{)}\<%
\\
%
\>[2]\AgdaSymbol{→}\AgdaSpace{}%
\AgdaSymbol{(}\AgdaBound{h}\AgdaSpace{}%
\AgdaSymbol{:}\AgdaSpace{}%
\AgdaSymbol{∀}\AgdaSpace{}%
\AgdaSymbol{\{}\AgdaBound{A'}\AgdaSymbol{\}}\AgdaSpace{}%
\AgdaSymbol{→}\AgdaSpace{}%
\AgdaSymbol{(}\AgdaBound{x}\AgdaSpace{}%
\AgdaSymbol{:}\AgdaSpace{}%
\AgdaBound{A'}\AgdaSpace{}%
\AgdaOperator{\AgdaDatatype{∈}}\AgdaSpace{}%
\AgdaBound{Γ}\AgdaSymbol{)}\AgdaSpace{}%
\AgdaSymbol{→}\AgdaSpace{}%
\AgdaInductiveConstructor{box}\AgdaSpace{}%
\AgdaSymbol{(}\AgdaBound{fb'}\AgdaSpace{}%
\AgdaSymbol{(}\AgdaFunction{embedLookup}\AgdaSpace{}%
\AgdaBound{x}\AgdaSymbol{))}\AgdaSpace{}%
\AgdaOperator{\AgdaDatatype{≡}}\AgdaSpace{}%
\AgdaFunction{embedTerm}\AgdaSpace{}%
\AgdaSymbol{(}\AgdaBound{f'}\AgdaSpace{}%
\AgdaBound{x}\AgdaSymbol{)}\AgdaSpace{}%
\AgdaSymbol{)}\<%
\\
%
\>[2]\AgdaSymbol{→}\AgdaSpace{}%
\AgdaSymbol{(}\AgdaBound{M}\AgdaSpace{}%
\AgdaSymbol{:}\AgdaSpace{}%
\AgdaBound{Γ}\AgdaSpace{}%
\AgdaOperator{\AgdaDatatype{⊢}}\AgdaSpace{}%
\AgdaBound{A}\AgdaSymbol{)}\<%
\\
%
\>[2]\AgdaSymbol{→}\AgdaSpace{}%
\AgdaFunction{substb}\AgdaSpace{}%
\AgdaBound{fb'}\AgdaSpace{}%
\AgdaSymbol{(}\AgdaFunction{embedTerm}\AgdaSpace{}%
\AgdaBound{M}\AgdaSymbol{)}\AgdaSpace{}%
\AgdaOperator{\AgdaDatatype{≡}}\AgdaSpace{}%
\AgdaFunction{embedTerm}\AgdaSpace{}%
\AgdaSymbol{(}\AgdaFunction{subst}\AgdaSpace{}%
\AgdaBound{f'}\AgdaSpace{}%
\AgdaBound{M}\AgdaSymbol{)}\<%
\\
\>[0]\AgdaFunction{pres-simultaneous-subst}\AgdaSpace{}%
\AgdaBound{fb'}\AgdaSpace{}%
\AgdaBound{f'}\AgdaSpace{}%
\AgdaBound{h}\AgdaSpace{}%
\AgdaSymbol{(}\AgdaInductiveConstructor{`}\AgdaSpace{}%
\AgdaBound{x}\AgdaSymbol{)}\AgdaSpace{}%
\AgdaSymbol{=}\<%
\\
\>[0][@{}l@{\AgdaIndent{0}}]%
\>[2]\AgdaBound{h}\AgdaSpace{}%
\AgdaBound{x}\<%
\\
\>[0]\AgdaFunction{pres-simultaneous-subst}\AgdaSpace{}%
\AgdaBound{fb'}\AgdaSpace{}%
\AgdaBound{f'}\AgdaSpace{}%
\AgdaBound{h}\AgdaSpace{}%
\AgdaSymbol{(}\AgdaInductiveConstructor{ƛ}\AgdaSpace{}%
\AgdaBound{M}\AgdaSymbol{)}\AgdaSpace{}%
\AgdaSymbol{=}\<%
\\
\>[0][@{}l@{\AgdaIndent{0}}]%
\>[2]\AgdaFunction{cong}\AgdaSpace{}%
\AgdaInductiveConstructor{box}\AgdaSpace{}%
\AgdaSymbol{(}\AgdaFunction{cong}\AgdaSpace{}%
\AgdaInductiveConstructor{ƛb}\AgdaSpace{}%
\AgdaSymbol{(}\AgdaFunction{pres-simultaneous-subst}\AgdaSpace{}%
\AgdaSymbol{(}\AgdaFunction{extsb}\AgdaSpace{}%
\AgdaBound{fb'}\AgdaSymbol{)}\AgdaSpace{}%
\AgdaSymbol{(}\AgdaFunction{exts}\AgdaSpace{}%
\AgdaBound{f'}\AgdaSymbol{)}\AgdaSpace{}%
\AgdaSymbol{(}\AgdaFunction{pres-exts}\AgdaSpace{}%
\AgdaBound{fb'}\AgdaSpace{}%
\AgdaBound{f'}\AgdaSpace{}%
\AgdaBound{h}\AgdaSymbol{)}\AgdaSpace{}%
\AgdaBound{M}\AgdaSymbol{))}\<%
\\
\>[0]\AgdaFunction{pres-simultaneous-subst}\AgdaSpace{}%
\AgdaBound{fb'}\AgdaSpace{}%
\AgdaBound{f'}\AgdaSpace{}%
\AgdaBound{h}\AgdaSpace{}%
\AgdaSymbol{(}\AgdaBound{M}\AgdaSpace{}%
\AgdaOperator{\AgdaInductiveConstructor{·}}\AgdaSpace{}%
\AgdaBound{M₁}\AgdaSymbol{)}\AgdaSpace{}%
\AgdaSymbol{=}\<%
\\
\>[0][@{}l@{\AgdaIndent{0}}]%
\>[2]\AgdaOperator{\AgdaFunction{beginEq}}\<%
\\
\>[2][@{}l@{\AgdaIndent{0}}]%
\>[4]\AgdaInductiveConstructor{ƛb}\AgdaSpace{}%
\AgdaSymbol{(}\AgdaInductiveConstructor{ε}\AgdaSpace{}%
\AgdaInductiveConstructor{Zb}\AgdaSpace{}%
\AgdaOperator{\AgdaInductiveConstructor{·b}}\AgdaSpace{}%
\AgdaFunction{substb}\AgdaSpace{}%
\AgdaSymbol{(}\AgdaFunction{extsb}\AgdaSpace{}%
\AgdaBound{fb'}\AgdaSymbol{)}\AgdaSpace{}%
\AgdaSymbol{(}\AgdaFunction{renameb}\AgdaSpace{}%
\AgdaInductiveConstructor{Sb}\AgdaSpace{}%
\AgdaSymbol{(}\AgdaFunction{embedTerm}\AgdaSpace{}%
\AgdaBound{M₁}\AgdaSymbol{)))}\AgdaSpace{}%
\AgdaOperator{\AgdaInductiveConstructor{·b}}\AgdaSpace{}%
\AgdaFunction{substb}\AgdaSpace{}%
\AgdaBound{fb'}\AgdaSpace{}%
\AgdaSymbol{(}\AgdaFunction{embedTerm}\AgdaSpace{}%
\AgdaBound{M}\AgdaSymbol{)}\<%
\\
%
\>[2]\AgdaFunction{≡⟨}\AgdaSpace{}%
\AgdaFunction{cong₂}\AgdaSpace{}%
\AgdaOperator{\AgdaInductiveConstructor{\AgdaUnderscore{}·b\AgdaUnderscore{}}}\AgdaSpace{}%
\AgdaInductiveConstructor{refl}\AgdaSpace{}%
\AgdaSymbol{(}\AgdaFunction{pres-simultaneous-subst}\AgdaSpace{}%
\AgdaBound{fb'}\AgdaSpace{}%
\AgdaBound{f'}\AgdaSpace{}%
\AgdaBound{h}\AgdaSpace{}%
\AgdaBound{M}\AgdaSymbol{)}\AgdaSpace{}%
\AgdaFunction{⟩}\<%
\\
\>[2][@{}l@{\AgdaIndent{0}}]%
\>[4]\AgdaInductiveConstructor{ƛb}\AgdaSpace{}%
\AgdaSymbol{(}\AgdaInductiveConstructor{ε}\AgdaSpace{}%
\AgdaInductiveConstructor{Zb}\AgdaSpace{}%
\AgdaOperator{\AgdaInductiveConstructor{·b}}\AgdaSpace{}%
\AgdaFunction{substb}\AgdaSpace{}%
\AgdaSymbol{(}\AgdaFunction{extsb}\AgdaSpace{}%
\AgdaBound{fb'}\AgdaSymbol{)}\AgdaSpace{}%
\AgdaSymbol{(}\AgdaFunction{renameb}\AgdaSpace{}%
\AgdaInductiveConstructor{Sb}\AgdaSpace{}%
\AgdaSymbol{(}\AgdaFunction{embedTerm}\AgdaSpace{}%
\AgdaBound{M₁}\AgdaSymbol{)))}\AgdaSpace{}%
\AgdaOperator{\AgdaInductiveConstructor{·b}}\AgdaSpace{}%
\AgdaFunction{embedTerm}\AgdaSpace{}%
\AgdaSymbol{(}\AgdaFunction{subst}\AgdaSpace{}%
\AgdaBound{f'}\AgdaSpace{}%
\AgdaBound{M}\AgdaSymbol{)}\<%
\\
%
\>[2]\AgdaFunction{≡⟨}\AgdaSpace{}%
\AgdaFunction{cong₂}\AgdaSpace{}%
\AgdaOperator{\AgdaInductiveConstructor{\AgdaUnderscore{}·b\AgdaUnderscore{}}}\AgdaSpace{}%
\AgdaSymbol{(}\AgdaFunction{cong}\AgdaSpace{}%
\AgdaInductiveConstructor{ƛb}\AgdaSpace{}%
\AgdaSymbol{(}\AgdaFunction{cong₂}\AgdaSpace{}%
\AgdaOperator{\AgdaInductiveConstructor{\AgdaUnderscore{}·b\AgdaUnderscore{}}}\AgdaSpace{}%
\AgdaInductiveConstructor{refl}\AgdaSpace{}%
\AgdaSymbol{(}\AgdaFunction{substb-exts-rename-Sb-commute}\AgdaSpace{}%
\AgdaBound{fb'}\AgdaSpace{}%
\AgdaSymbol{(}\AgdaFunction{embedTerm}\AgdaSpace{}%
\AgdaBound{M₁}\AgdaSymbol{))))}\AgdaSpace{}%
\AgdaInductiveConstructor{refl}\AgdaSpace{}%
\AgdaFunction{⟩}\<%
\\
\>[2][@{}l@{\AgdaIndent{0}}]%
\>[4]\AgdaInductiveConstructor{ƛb}\AgdaSpace{}%
\AgdaSymbol{(}\AgdaInductiveConstructor{ε}\AgdaSpace{}%
\AgdaInductiveConstructor{Zb}\AgdaSpace{}%
\AgdaOperator{\AgdaInductiveConstructor{·b}}\AgdaSpace{}%
\AgdaFunction{renameb}\AgdaSpace{}%
\AgdaInductiveConstructor{Sb}\AgdaSpace{}%
\AgdaSymbol{(}\AgdaFunction{substb}\AgdaSpace{}%
\AgdaBound{fb'}\AgdaSpace{}%
\AgdaSymbol{(}\AgdaFunction{embedTerm}\AgdaSpace{}%
\AgdaBound{M₁}\AgdaSymbol{)))}\AgdaSpace{}%
\AgdaOperator{\AgdaInductiveConstructor{·b}}\AgdaSpace{}%
\AgdaFunction{embedTerm}\AgdaSpace{}%
\AgdaSymbol{(}\AgdaFunction{subst}\AgdaSpace{}%
\AgdaBound{f'}\AgdaSpace{}%
\AgdaBound{M}\AgdaSymbol{)}\<%
\\
%
\>[2]\AgdaFunction{≡⟨}\AgdaSpace{}%
\AgdaFunction{cong₂}\AgdaSpace{}%
\AgdaOperator{\AgdaInductiveConstructor{\AgdaUnderscore{}·b\AgdaUnderscore{}}}\AgdaSpace{}%
\AgdaSymbol{(}\AgdaFunction{cong}\AgdaSpace{}%
\AgdaInductiveConstructor{ƛb}\AgdaSpace{}%
\AgdaSymbol{(}\AgdaFunction{cong₂}\AgdaSpace{}%
\AgdaOperator{\AgdaInductiveConstructor{\AgdaUnderscore{}·b\AgdaUnderscore{}}}\AgdaSpace{}%
\AgdaInductiveConstructor{refl}\AgdaSpace{}%
\AgdaSymbol{(}\AgdaFunction{cong}\AgdaSpace{}%
\AgdaSymbol{(}\AgdaFunction{renameb}\AgdaSpace{}%
\AgdaInductiveConstructor{Sb}\AgdaSymbol{)}\AgdaSpace{}%
\AgdaSymbol{(}\AgdaFunction{pres-simultaneous-subst}\AgdaSpace{}%
\AgdaBound{fb'}\AgdaSpace{}%
\AgdaBound{f'}\AgdaSpace{}%
\AgdaBound{h}\AgdaSpace{}%
\AgdaBound{M₁}\AgdaSymbol{))))}\AgdaSpace{}%
\AgdaInductiveConstructor{refl}\AgdaSpace{}%
\AgdaFunction{⟩}\<%
\\
\>[2][@{}l@{\AgdaIndent{0}}]%
\>[4]\AgdaInductiveConstructor{ƛb}\AgdaSpace{}%
\AgdaSymbol{(}\AgdaInductiveConstructor{ε}\AgdaSpace{}%
\AgdaInductiveConstructor{Zb}\AgdaSpace{}%
\AgdaOperator{\AgdaInductiveConstructor{·b}}\AgdaSpace{}%
\AgdaFunction{renameb}\AgdaSpace{}%
\AgdaInductiveConstructor{Sb}\AgdaSpace{}%
\AgdaSymbol{(}\AgdaFunction{embedTerm}\AgdaSpace{}%
\AgdaSymbol{(}\AgdaFunction{subst}\AgdaSpace{}%
\AgdaBound{f'}\AgdaSpace{}%
\AgdaBound{M₁}\AgdaSymbol{)))}\AgdaSpace{}%
\AgdaOperator{\AgdaInductiveConstructor{·b}}\AgdaSpace{}%
\AgdaFunction{embedTerm}\AgdaSpace{}%
\AgdaSymbol{(}\AgdaFunction{subst}\AgdaSpace{}%
\AgdaBound{f'}\AgdaSpace{}%
\AgdaBound{M}\AgdaSymbol{)}\<%
\\
%
\>[2]\AgdaOperator{\AgdaFunction{∎eq}}\<%
\\
%
\\[\AgdaEmptyExtraSkip]%
\>[0]\AgdaComment{--\ HELPER\ LEMMA'S\ TO\ PROVE\ h}\<%
\\
%
\\[\AgdaEmptyExtraSkip]%
\>[0]\AgdaComment{--\ Lemma\ for\ assumption\ h\ in\ lemma}\<%
\\
\>[0]\AgdaComment{--\ pres-simultaneous-subst}\<%
\\
\>[0]\AgdaFunction{pres-subst-func}\AgdaSpace{}%
\AgdaSymbol{:}\AgdaSpace{}%
\AgdaSymbol{∀}\AgdaSpace{}%
\AgdaSymbol{\{}\AgdaBound{A}\AgdaSpace{}%
\AgdaBound{C}\AgdaSpace{}%
\AgdaBound{Γ}\AgdaSymbol{\}}\<%
\\
\>[0][@{}l@{\AgdaIndent{0}}]%
\>[2]\AgdaSymbol{→}\AgdaSpace{}%
\AgdaSymbol{(}\AgdaBound{V}\AgdaSpace{}%
\AgdaSymbol{:}\AgdaSpace{}%
\AgdaBound{Γ}\AgdaSpace{}%
\AgdaOperator{\AgdaDatatype{⊢}}\AgdaSpace{}%
\AgdaBound{C}\AgdaSymbol{)}\<%
\\
%
\>[2]\AgdaSymbol{→}\AgdaSpace{}%
\AgdaSymbol{(}\AgdaBound{y}\AgdaSpace{}%
\AgdaSymbol{:}\AgdaSpace{}%
\AgdaDatatype{Value}\AgdaSpace{}%
\AgdaBound{V}\AgdaSymbol{)}\<%
\\
%
\>[2]\AgdaSymbol{→}\AgdaSpace{}%
\AgdaSymbol{(}\AgdaBound{x}\AgdaSpace{}%
\AgdaSymbol{:}\AgdaSpace{}%
\AgdaBound{A}\AgdaSpace{}%
\AgdaOperator{\AgdaDatatype{∈}}\AgdaSpace{}%
\AgdaBound{Γ}\AgdaSpace{}%
\AgdaOperator{\AgdaInductiveConstructor{,}}\AgdaSpace{}%
\AgdaBound{C}\AgdaSymbol{)}\<%
\\
%
\>[2]\AgdaSymbol{→}\AgdaSpace{}%
\AgdaInductiveConstructor{box}\AgdaSpace{}%
\AgdaSymbol{(}\AgdaFunction{fb}\AgdaSpace{}%
\AgdaSymbol{(}\AgdaFunction{embedValue}\AgdaSpace{}%
\AgdaBound{V}\AgdaSpace{}%
\AgdaBound{y}\AgdaSymbol{)}\AgdaSpace{}%
\AgdaSymbol{(}\AgdaFunction{embedLookup}\AgdaSpace{}%
\AgdaBound{x}\AgdaSymbol{))}\AgdaSpace{}%
\AgdaOperator{\AgdaDatatype{≡}}\AgdaSpace{}%
\AgdaFunction{embedTerm}\AgdaSpace{}%
\AgdaSymbol{((}\AgdaFunction{f}\AgdaSpace{}%
\AgdaBound{V}\AgdaSymbol{)}\AgdaSpace{}%
\AgdaBound{x}\AgdaSymbol{)}\<%
\\
\>[0]\AgdaFunction{pres-subst-func}\AgdaSpace{}%
\AgdaBound{V}\AgdaSpace{}%
\AgdaBound{y}\AgdaSpace{}%
\AgdaInductiveConstructor{Z}\AgdaSpace{}%
\AgdaSymbol{=}\AgdaSpace{}%
\AgdaFunction{box-embedValue}\AgdaSpace{}%
\AgdaBound{y}\<%
\\
\>[0]\AgdaFunction{pres-subst-func}\AgdaSpace{}%
\AgdaBound{V}\AgdaSpace{}%
\AgdaBound{y}\AgdaSpace{}%
\AgdaSymbol{(}\AgdaInductiveConstructor{S}\AgdaSpace{}%
\AgdaBound{x}\AgdaSymbol{)}\AgdaSpace{}%
\AgdaSymbol{=}\AgdaSpace{}%
\AgdaInductiveConstructor{refl}\<%
\\
%
\\[\AgdaEmptyExtraSkip]%
\>[0]\AgdaComment{--\ exts\ lemma\ for\ assumption\ h\ in\ lemma}\<%
\\
\>[0]\AgdaFunction{exts-pres-subst-func}\AgdaSpace{}%
\AgdaSymbol{:}\AgdaSpace{}%
\AgdaSymbol{∀}\AgdaSpace{}%
\AgdaSymbol{\{}\AgdaBound{A}\AgdaSpace{}%
\AgdaBound{C}\AgdaSpace{}%
\AgdaBound{Γ}\AgdaSpace{}%
\AgdaBound{B}\AgdaSymbol{\}}\<%
\\
\>[0][@{}l@{\AgdaIndent{0}}]%
\>[2]\AgdaSymbol{→}\AgdaSpace{}%
\AgdaSymbol{(}\AgdaBound{V}\AgdaSpace{}%
\AgdaSymbol{:}\AgdaSpace{}%
\AgdaBound{Γ}\AgdaSpace{}%
\AgdaOperator{\AgdaDatatype{⊢}}\AgdaSpace{}%
\AgdaBound{C}\AgdaSymbol{)}\<%
\\
%
\>[2]\AgdaSymbol{→}\AgdaSpace{}%
\AgdaSymbol{(}\AgdaBound{y}\AgdaSpace{}%
\AgdaSymbol{:}\AgdaSpace{}%
\AgdaDatatype{Value}\AgdaSpace{}%
\AgdaBound{V}\AgdaSymbol{)}\<%
\\
%
\>[2]\AgdaSymbol{→}\AgdaSpace{}%
\AgdaSymbol{(}\AgdaBound{x}\AgdaSpace{}%
\AgdaSymbol{:}\AgdaSpace{}%
\AgdaBound{A}\AgdaSpace{}%
\AgdaOperator{\AgdaDatatype{∈}}\AgdaSpace{}%
\AgdaBound{Γ}\AgdaSpace{}%
\AgdaOperator{\AgdaInductiveConstructor{,}}\AgdaSpace{}%
\AgdaBound{C}\AgdaSpace{}%
\AgdaOperator{\AgdaInductiveConstructor{,}}\AgdaSpace{}%
\AgdaBound{B}\AgdaSymbol{)}\<%
\\
%
\>[2]\AgdaSymbol{→}\AgdaSpace{}%
\AgdaInductiveConstructor{box}\AgdaSpace{}%
\AgdaSymbol{(}\AgdaFunction{extsb}\AgdaSpace{}%
\AgdaSymbol{(}\AgdaFunction{fb}\AgdaSpace{}%
\AgdaSymbol{(}\AgdaFunction{embedValue}\AgdaSpace{}%
\AgdaBound{V}\AgdaSpace{}%
\AgdaBound{y}\AgdaSymbol{))}\AgdaSpace{}%
\AgdaSymbol{(}\AgdaFunction{embedLookup}\AgdaSpace{}%
\AgdaBound{x}\AgdaSymbol{))}\AgdaSpace{}%
\AgdaOperator{\AgdaDatatype{≡}}\AgdaSpace{}%
\AgdaFunction{embedTerm}\AgdaSpace{}%
\AgdaSymbol{(}\AgdaFunction{exts}\AgdaSpace{}%
\AgdaSymbol{(}\AgdaFunction{f}\AgdaSpace{}%
\AgdaBound{V}\AgdaSymbol{)}\AgdaSpace{}%
\AgdaBound{x}\AgdaSymbol{)}\<%
\\
\>[0]\AgdaFunction{exts-pres-subst-func}\AgdaSpace{}%
\AgdaBound{V}\AgdaSpace{}%
\AgdaBound{y}\AgdaSpace{}%
\AgdaInductiveConstructor{Z}\AgdaSpace{}%
\AgdaSymbol{=}\AgdaSpace{}%
\AgdaInductiveConstructor{refl}\<%
\\
\>[0]\AgdaFunction{exts-pres-subst-func}\AgdaSpace{}%
\AgdaBound{V}\AgdaSpace{}%
\AgdaBound{y}\AgdaSpace{}%
\AgdaSymbol{(}\AgdaInductiveConstructor{S}\AgdaSpace{}%
\AgdaBound{x}\AgdaSymbol{)}\AgdaSpace{}%
\AgdaSymbol{=}\<%
\\
\>[0][@{}l@{\AgdaIndent{0}}]%
\>[2]\AgdaOperator{\AgdaFunction{beginEq}}\<%
\\
\>[2][@{}l@{\AgdaIndent{0}}]%
\>[4]\AgdaInductiveConstructor{box}\AgdaSpace{}%
\AgdaSymbol{(}\AgdaFunction{renameb}\AgdaSpace{}%
\AgdaInductiveConstructor{Sb}\AgdaSpace{}%
\AgdaSymbol{(}\AgdaFunction{fb}\AgdaSpace{}%
\AgdaSymbol{(}\AgdaFunction{embedValue}\AgdaSpace{}%
\AgdaBound{V}\AgdaSpace{}%
\AgdaBound{y}\AgdaSymbol{)}\AgdaSpace{}%
\AgdaSymbol{(}\AgdaFunction{embedLookup}\AgdaSpace{}%
\AgdaBound{x}\AgdaSymbol{)))}\<%
\\
%
\>[2]\AgdaFunction{≡⟨}\AgdaSpace{}%
\AgdaInductiveConstructor{refl}\AgdaSpace{}%
\AgdaFunction{⟩}\<%
\\
\>[2][@{}l@{\AgdaIndent{0}}]%
\>[4]\AgdaFunction{renameb}\AgdaSpace{}%
\AgdaInductiveConstructor{Sb}\AgdaSpace{}%
\AgdaSymbol{(}\AgdaInductiveConstructor{box}\AgdaSpace{}%
\AgdaSymbol{(}\AgdaFunction{fb}\AgdaSpace{}%
\AgdaSymbol{(}\AgdaFunction{embedValue}\AgdaSpace{}%
\AgdaBound{V}\AgdaSpace{}%
\AgdaBound{y}\AgdaSymbol{)}\AgdaSpace{}%
\AgdaSymbol{(}\AgdaFunction{embedLookup}\AgdaSpace{}%
\AgdaBound{x}\AgdaSymbol{)))}\<%
\\
%
\>[2]\AgdaFunction{≡⟨}\AgdaSpace{}%
\AgdaFunction{cong}\AgdaSpace{}%
\AgdaSymbol{(}\AgdaFunction{renameb}\AgdaSpace{}%
\AgdaInductiveConstructor{Sb}\AgdaSymbol{)}\AgdaSpace{}%
\AgdaSymbol{(}\AgdaFunction{pres-subst-func}\AgdaSpace{}%
\AgdaBound{V}\AgdaSpace{}%
\AgdaBound{y}\AgdaSpace{}%
\AgdaBound{x}\AgdaSymbol{)}\AgdaSpace{}%
\AgdaFunction{⟩}\<%
\\
\>[2][@{}l@{\AgdaIndent{0}}]%
\>[4]\AgdaFunction{renameb}\AgdaSpace{}%
\AgdaInductiveConstructor{Sb}\AgdaSpace{}%
\AgdaSymbol{(}\AgdaFunction{embedTerm}\AgdaSpace{}%
\AgdaSymbol{(}\AgdaSpace{}%
\AgdaSymbol{(}\AgdaFunction{f}\AgdaSpace{}%
\AgdaBound{V}\AgdaSymbol{)}\AgdaSpace{}%
\AgdaBound{x}\AgdaSpace{}%
\AgdaSymbol{))}\<%
\\
%
\>[2]\AgdaFunction{≡⟨}\AgdaSpace{}%
\AgdaFunction{pres-rename}\AgdaSpace{}%
\AgdaInductiveConstructor{S}\AgdaSpace{}%
\AgdaInductiveConstructor{Sb}\AgdaSpace{}%
\AgdaSymbol{(λ}\AgdaSpace{}%
\AgdaBound{x₁}\AgdaSpace{}%
\AgdaSymbol{→}\AgdaSpace{}%
\AgdaInductiveConstructor{refl}\AgdaSymbol{)}\AgdaSpace{}%
\AgdaSymbol{(}\AgdaFunction{f}\AgdaSpace{}%
\AgdaBound{V}\AgdaSpace{}%
\AgdaBound{x}\AgdaSymbol{)}\AgdaSpace{}%
\AgdaFunction{⟩}\<%
\\
\>[2][@{}l@{\AgdaIndent{0}}]%
\>[4]\AgdaFunction{embedTerm}\AgdaSpace{}%
\AgdaSymbol{(}\AgdaFunction{rename}\AgdaSpace{}%
\AgdaInductiveConstructor{S}\AgdaSpace{}%
\AgdaSymbol{(}\AgdaFunction{f}\AgdaSpace{}%
\AgdaBound{V}\AgdaSpace{}%
\AgdaBound{x}\AgdaSymbol{))}\<%
\\
%
\>[2]\AgdaOperator{\AgdaFunction{∎eq}}\<%
\\
%
\\[\AgdaEmptyExtraSkip]%
\>[0]\AgdaComment{--\ Preservation\ and\ Reflection\ of\ }\<%
\\
\>[0]\AgdaComment{--\ single\ substitution\ of\ Gödel's\ translation}\<%
\\
\>[0]\AgdaFunction{pres-subst}\AgdaSpace{}%
\AgdaSymbol{:}\AgdaSpace{}%
\AgdaSymbol{∀}\AgdaSpace{}%
\AgdaSymbol{\{}\AgdaBound{Γ}\AgdaSpace{}%
\AgdaBound{A}\AgdaSpace{}%
\AgdaBound{B}\AgdaSymbol{\}}\AgdaSpace{}%
\AgdaSymbol{→}\AgdaSpace{}%
\AgdaSymbol{(}\AgdaBound{M}\AgdaSpace{}%
\AgdaSymbol{:}\AgdaSpace{}%
\AgdaBound{Γ}\AgdaSpace{}%
\AgdaOperator{\AgdaInductiveConstructor{,}}\AgdaSpace{}%
\AgdaBound{B}\AgdaSpace{}%
\AgdaOperator{\AgdaDatatype{⊢}}\AgdaSpace{}%
\AgdaBound{A}\AgdaSymbol{)}\AgdaSpace{}%
\AgdaSymbol{→}\AgdaSpace{}%
\AgdaSymbol{(}\AgdaBound{V}\AgdaSpace{}%
\AgdaSymbol{:}\AgdaSpace{}%
\AgdaBound{Γ}\AgdaSpace{}%
\AgdaOperator{\AgdaDatatype{⊢}}\AgdaSpace{}%
\AgdaBound{B}\AgdaSymbol{)}\AgdaSpace{}%
\AgdaSymbol{→}\AgdaSpace{}%
\AgdaSymbol{(}\AgdaBound{x}\AgdaSpace{}%
\AgdaSymbol{:}\AgdaSpace{}%
\AgdaDatatype{Value}\AgdaSpace{}%
\AgdaBound{V}\AgdaSymbol{)}\AgdaSpace{}%
\AgdaSymbol{→}\AgdaSpace{}%
\AgdaFunction{embedTerm}\AgdaSpace{}%
\AgdaBound{M}\AgdaSpace{}%
\AgdaOperator{\AgdaFunction{[}}\AgdaSpace{}%
\AgdaFunction{embedValue}\AgdaSpace{}%
\AgdaBound{V}\AgdaSpace{}%
\AgdaBound{x}\AgdaSpace{}%
\AgdaOperator{\AgdaFunction{]b}}\AgdaSpace{}%
\AgdaOperator{\AgdaDatatype{≡}}\AgdaSpace{}%
\AgdaFunction{embedTerm}\AgdaSpace{}%
\AgdaSymbol{(}\AgdaBound{M}\AgdaSpace{}%
\AgdaOperator{\AgdaFunction{[}}\AgdaSpace{}%
\AgdaBound{V}\AgdaSpace{}%
\AgdaOperator{\AgdaFunction{]}}\AgdaSymbol{)}\<%
\\
\>[0]\AgdaFunction{pres-subst}\AgdaSpace{}%
\AgdaBound{M}\AgdaSpace{}%
\AgdaBound{V}\AgdaSpace{}%
\AgdaBound{x}\AgdaSpace{}%
\AgdaSymbol{=}\AgdaSpace{}%
\AgdaFunction{pres-simultaneous-subst}\AgdaSpace{}%
\AgdaSymbol{(}\AgdaFunction{fb}\AgdaSpace{}%
\AgdaSymbol{(}\AgdaFunction{embedValue}\AgdaSpace{}%
\AgdaBound{V}\AgdaSpace{}%
\AgdaBound{x}\AgdaSymbol{))}\AgdaSpace{}%
\AgdaSymbol{(}\AgdaFunction{f}\AgdaSpace{}%
\AgdaBound{V}\AgdaSymbol{)}\AgdaSpace{}%
\AgdaSymbol{(}\AgdaFunction{pres-subst-func}\AgdaSpace{}%
\AgdaBound{V}\AgdaSpace{}%
\AgdaBound{x}\AgdaSymbol{)}\AgdaSpace{}%
\AgdaBound{M}\<%
\\
%
\\[\AgdaEmptyExtraSkip]%
\>[0]\AgdaComment{--\ Helper\ lemma:\ if\ a\ value\ is\ substituted\ in\ a\ value,\ the\ resulting\ lambda\ term\ will\ also\ be\ a\ value}\<%
\\
\>[0]\AgdaFunction{val-after-subst}\AgdaSpace{}%
\AgdaSymbol{:}\AgdaSpace{}%
\AgdaSymbol{∀}\AgdaSpace{}%
\AgdaSymbol{\{}\AgdaBound{Γ}\AgdaSpace{}%
\AgdaBound{A}\AgdaSpace{}%
\AgdaBound{B}\AgdaSymbol{\}}\AgdaSpace{}%
\AgdaSymbol{→}\AgdaSpace{}%
\AgdaSymbol{(}\AgdaBound{M}\AgdaSpace{}%
\AgdaSymbol{:}\AgdaSpace{}%
\AgdaBound{Γ}\AgdaSpace{}%
\AgdaOperator{\AgdaInductiveConstructor{,}}\AgdaSpace{}%
\AgdaBound{B}\AgdaSpace{}%
\AgdaOperator{\AgdaDatatype{⊢}}\AgdaSpace{}%
\AgdaBound{A}\AgdaSymbol{)}\AgdaSpace{}%
\AgdaSymbol{→}\AgdaSpace{}%
\AgdaSymbol{(}\AgdaBound{N}\AgdaSpace{}%
\AgdaSymbol{:}\AgdaSpace{}%
\AgdaBound{Γ}\AgdaSpace{}%
\AgdaOperator{\AgdaDatatype{⊢}}\AgdaSpace{}%
\AgdaBound{B}\AgdaSymbol{)}\AgdaSpace{}%
\AgdaSymbol{→}\AgdaSpace{}%
\AgdaSymbol{(}\AgdaBound{x}\AgdaSpace{}%
\AgdaSymbol{:}\AgdaSpace{}%
\AgdaDatatype{Value}\AgdaSpace{}%
\AgdaBound{M}\AgdaSymbol{)}\AgdaSpace{}%
\AgdaSymbol{→}\AgdaSpace{}%
\AgdaSymbol{(}\AgdaBound{y}\AgdaSpace{}%
\AgdaSymbol{:}\AgdaSpace{}%
\AgdaDatatype{Value}\AgdaSpace{}%
\AgdaBound{N}\AgdaSymbol{)}\AgdaSpace{}%
\AgdaSymbol{→}\AgdaSpace{}%
\AgdaDatatype{Value}\AgdaSpace{}%
\AgdaSymbol{(}\AgdaBound{M}\AgdaSpace{}%
\AgdaOperator{\AgdaFunction{[}}\AgdaSpace{}%
\AgdaBound{N}\AgdaSpace{}%
\AgdaOperator{\AgdaFunction{]}}\AgdaSymbol{)}\<%
\\
\>[0]\AgdaFunction{val-after-subst}\AgdaSpace{}%
\AgdaSymbol{(}\AgdaInductiveConstructor{`}\AgdaSpace{}%
\AgdaInductiveConstructor{Z}\AgdaSymbol{)}\AgdaSpace{}%
\AgdaBound{N}\AgdaSpace{}%
\AgdaBound{x}\AgdaSpace{}%
\AgdaBound{y}\AgdaSpace{}%
\AgdaSymbol{=}\AgdaSpace{}%
\AgdaBound{y}\<%
\\
\>[0]\AgdaFunction{val-after-subst}\AgdaSpace{}%
\AgdaSymbol{(}\AgdaInductiveConstructor{`}\AgdaSpace{}%
\AgdaSymbol{(}\AgdaInductiveConstructor{S}\AgdaSpace{}%
\AgdaBound{x₁}\AgdaSymbol{))}\AgdaSpace{}%
\AgdaBound{N}\AgdaSpace{}%
\AgdaBound{x}\AgdaSpace{}%
\AgdaBound{y}\AgdaSpace{}%
\AgdaSymbol{=}\AgdaSpace{}%
\AgdaInductiveConstructor{V-`}\<%
\\
\>[0]\AgdaFunction{val-after-subst}\AgdaSpace{}%
\AgdaSymbol{(}\AgdaInductiveConstructor{ƛ}\AgdaSpace{}%
\AgdaBound{M}\AgdaSymbol{)}\AgdaSpace{}%
\AgdaBound{N}\AgdaSpace{}%
\AgdaBound{x}\AgdaSpace{}%
\AgdaBound{y}\AgdaSpace{}%
\AgdaSymbol{=}\AgdaSpace{}%
\AgdaInductiveConstructor{V-ƛ}\<%
\\
%
\\[\AgdaEmptyExtraSkip]%
\>[0]\AgdaComment{--\ Preservation\ and\ Reflection\ of\ }\<%
\\
\>[0]\AgdaComment{--\ single\ substitution\ in\ values\ of\ Gödel's\ translation\ }\<%
\\
\>[0]\AgdaFunction{pres-subst-val}\AgdaSpace{}%
\AgdaSymbol{:}\AgdaSpace{}%
\AgdaSymbol{∀}\AgdaSpace{}%
\AgdaSymbol{\{}\AgdaBound{Γ}\AgdaSpace{}%
\AgdaBound{A}\AgdaSpace{}%
\AgdaBound{B}\AgdaSymbol{\}}\<%
\\
\>[0][@{}l@{\AgdaIndent{0}}]%
\>[2]\AgdaSymbol{→}\AgdaSpace{}%
\AgdaSymbol{(}\AgdaBound{M}\AgdaSpace{}%
\AgdaSymbol{:}\AgdaSpace{}%
\AgdaBound{Γ}\AgdaSpace{}%
\AgdaOperator{\AgdaInductiveConstructor{,}}\AgdaSpace{}%
\AgdaBound{B}\AgdaSpace{}%
\AgdaOperator{\AgdaDatatype{⊢}}\AgdaSpace{}%
\AgdaBound{A}\AgdaSymbol{)}\<%
\\
%
\>[2]\AgdaSymbol{→}\AgdaSpace{}%
\AgdaSymbol{(}\AgdaBound{V}\AgdaSpace{}%
\AgdaSymbol{:}\AgdaSpace{}%
\AgdaBound{Γ}\AgdaSpace{}%
\AgdaOperator{\AgdaDatatype{⊢}}\AgdaSpace{}%
\AgdaBound{B}\AgdaSymbol{)}\<%
\\
%
\>[2]\AgdaSymbol{→}\AgdaSpace{}%
\AgdaSymbol{(}\AgdaBound{x}\AgdaSpace{}%
\AgdaSymbol{:}\AgdaSpace{}%
\AgdaDatatype{Value}\AgdaSpace{}%
\AgdaBound{M}\AgdaSymbol{)}\<%
\\
%
\>[2]\AgdaSymbol{→}\AgdaSpace{}%
\AgdaSymbol{(}\AgdaBound{y}\AgdaSpace{}%
\AgdaSymbol{:}\AgdaSpace{}%
\AgdaDatatype{Value}\AgdaSpace{}%
\AgdaBound{V}\AgdaSymbol{)}\<%
\\
%
\>[2]\AgdaSymbol{→}\AgdaSpace{}%
\AgdaFunction{embedValue}\AgdaSpace{}%
\AgdaBound{M}\AgdaSpace{}%
\AgdaBound{x}\AgdaSpace{}%
\AgdaOperator{\AgdaFunction{[}}\AgdaSpace{}%
\AgdaFunction{embedValue}\AgdaSpace{}%
\AgdaBound{V}\AgdaSpace{}%
\AgdaBound{y}\AgdaSpace{}%
\AgdaOperator{\AgdaFunction{]b}}\AgdaSpace{}%
\AgdaOperator{\AgdaDatatype{≡}}\AgdaSpace{}%
\AgdaFunction{embedValue}\AgdaSpace{}%
\AgdaSymbol{(}\AgdaBound{M}\AgdaSpace{}%
\AgdaOperator{\AgdaFunction{[}}\AgdaSpace{}%
\AgdaBound{V}\AgdaSpace{}%
\AgdaOperator{\AgdaFunction{]}}\AgdaSymbol{)}\AgdaSpace{}%
\AgdaSymbol{(}\AgdaFunction{val-after-subst}\AgdaSpace{}%
\AgdaBound{M}\AgdaSpace{}%
\AgdaBound{V}\AgdaSpace{}%
\AgdaBound{x}\AgdaSpace{}%
\AgdaBound{y}\AgdaSymbol{)}\<%
\\
\>[0]\AgdaFunction{pres-subst-val}\AgdaSpace{}%
\AgdaSymbol{(}\AgdaInductiveConstructor{`}\AgdaSpace{}%
\AgdaInductiveConstructor{Z}\AgdaSymbol{)}\AgdaSpace{}%
\AgdaBound{V}\AgdaSpace{}%
\AgdaBound{x}\AgdaSpace{}%
\AgdaBound{y}\AgdaSpace{}%
\AgdaSymbol{=}\AgdaSpace{}%
\AgdaInductiveConstructor{refl}\<%
\\
\>[0]\AgdaFunction{pres-subst-val}\AgdaSpace{}%
\AgdaSymbol{(}\AgdaInductiveConstructor{`}\AgdaSpace{}%
\AgdaSymbol{(}\AgdaInductiveConstructor{S}\AgdaSpace{}%
\AgdaBound{x₁}\AgdaSymbol{))}\AgdaSpace{}%
\AgdaBound{V}\AgdaSpace{}%
\AgdaBound{x}\AgdaSpace{}%
\AgdaBound{y}\AgdaSpace{}%
\AgdaSymbol{=}\AgdaSpace{}%
\AgdaInductiveConstructor{refl}\<%
\\
\>[0]\AgdaFunction{pres-subst-val}\AgdaSpace{}%
\AgdaSymbol{(}\AgdaInductiveConstructor{ƛ}\AgdaSpace{}%
\AgdaBound{M}\AgdaSymbol{)}\AgdaSpace{}%
\AgdaBound{V}\AgdaSpace{}%
\AgdaBound{x}\AgdaSpace{}%
\AgdaBound{y}\AgdaSpace{}%
\AgdaSymbol{=}\<%
\\
\>[0][@{}l@{\AgdaIndent{0}}]%
\>[2]\AgdaOperator{\AgdaFunction{beginEq}}\<%
\\
\>[2][@{}l@{\AgdaIndent{0}}]%
\>[4]\AgdaFunction{embedValue}\AgdaSpace{}%
\AgdaSymbol{(}\AgdaInductiveConstructor{ƛ}\AgdaSpace{}%
\AgdaBound{M}\AgdaSymbol{)}\AgdaSpace{}%
\AgdaBound{x}\AgdaSpace{}%
\AgdaOperator{\AgdaFunction{[}}\AgdaSpace{}%
\AgdaFunction{embedValue}\AgdaSpace{}%
\AgdaBound{V}\AgdaSpace{}%
\AgdaBound{y}\AgdaSpace{}%
\AgdaOperator{\AgdaFunction{]b}}\<%
\\
%
\>[2]\AgdaFunction{≡⟨}\AgdaSpace{}%
\AgdaInductiveConstructor{refl}\AgdaSpace{}%
\AgdaFunction{⟩}\<%
\\
\>[2][@{}l@{\AgdaIndent{0}}]%
\>[4]\AgdaInductiveConstructor{ƛb}\AgdaSpace{}%
\AgdaSymbol{(}\AgdaFunction{substb}\AgdaSpace{}%
\AgdaSymbol{(}\AgdaFunction{extsb}\AgdaSpace{}%
\AgdaSymbol{(}\AgdaFunction{fb}\AgdaSpace{}%
\AgdaSymbol{(}\AgdaFunction{embedValue}\AgdaSpace{}%
\AgdaBound{V}\AgdaSpace{}%
\AgdaBound{y}\AgdaSymbol{)))}\AgdaSpace{}%
\AgdaSymbol{(}\AgdaFunction{embedTerm}\AgdaSpace{}%
\AgdaBound{M}\AgdaSymbol{))}\<%
\\
%
\>[2]\AgdaFunction{≡⟨}\AgdaSpace{}%
\AgdaFunction{cong}\AgdaSpace{}%
\AgdaInductiveConstructor{ƛb}\AgdaSpace{}%
\AgdaSymbol{(}\AgdaFunction{pres-simultaneous-subst}\AgdaSpace{}%
\AgdaSymbol{(}\AgdaFunction{extsb}\AgdaSpace{}%
\AgdaSymbol{(}\AgdaFunction{fb}\AgdaSpace{}%
\AgdaSymbol{(}\AgdaFunction{embedValue}\AgdaSpace{}%
\AgdaBound{V}\AgdaSpace{}%
\AgdaBound{y}\AgdaSymbol{)))}\AgdaSpace{}%
\AgdaSymbol{(}\AgdaFunction{exts}\AgdaSpace{}%
\AgdaSymbol{(}\AgdaFunction{f}\AgdaSpace{}%
\AgdaBound{V}\AgdaSymbol{))}\AgdaSpace{}%
\AgdaSymbol{(}\AgdaFunction{exts-pres-subst-func}\AgdaSpace{}%
\AgdaBound{V}\AgdaSpace{}%
\AgdaBound{y}\AgdaSymbol{)}\AgdaSpace{}%
\AgdaBound{M}\AgdaSymbol{)}\AgdaSpace{}%
\AgdaFunction{⟩}\<%
\\
\>[2][@{}l@{\AgdaIndent{0}}]%
\>[4]\AgdaInductiveConstructor{ƛb}\AgdaSpace{}%
\AgdaSymbol{(}\AgdaFunction{embedTerm}\AgdaSpace{}%
\AgdaSymbol{(}\AgdaFunction{subst}\AgdaSpace{}%
\AgdaSymbol{(}\AgdaFunction{exts}\AgdaSpace{}%
\AgdaSymbol{(}\AgdaFunction{f}\AgdaSpace{}%
\AgdaBound{V}\AgdaSymbol{))}\AgdaSpace{}%
\AgdaBound{M}\AgdaSymbol{))}\<%
\\
%
\>[2]\AgdaFunction{≡⟨}\AgdaSpace{}%
\AgdaInductiveConstructor{refl}\AgdaSpace{}%
\AgdaFunction{⟩}\<%
\\
\>[2][@{}l@{\AgdaIndent{0}}]%
\>[4]\AgdaFunction{embedValue}\AgdaSpace{}%
\AgdaSymbol{((}\AgdaInductiveConstructor{ƛ}\AgdaSpace{}%
\AgdaBound{M}\AgdaSymbol{)}\AgdaSpace{}%
\AgdaOperator{\AgdaFunction{[}}\AgdaSpace{}%
\AgdaBound{V}\AgdaSpace{}%
\AgdaOperator{\AgdaFunction{]}}\AgdaSymbol{)}\AgdaSpace{}%
\AgdaSymbol{(}\AgdaFunction{val-after-subst}\AgdaSpace{}%
\AgdaSymbol{(}\AgdaInductiveConstructor{ƛ}\AgdaSpace{}%
\AgdaBound{M}\AgdaSymbol{)}\AgdaSpace{}%
\AgdaBound{V}\AgdaSpace{}%
\AgdaBound{x}\AgdaSpace{}%
\AgdaBound{y}\AgdaSymbol{)}\<%
\\
%
\>[2]\AgdaOperator{\AgdaFunction{∎eq}}\<%
\\
%
\\[\AgdaEmptyExtraSkip]%
\>[0]\AgdaComment{--\ fb\ and\ renameb\ commute\ under\ extended\ context.}\<%
\\
\>[0]\AgdaFunction{fb-renameb-extb-commute}\AgdaSpace{}%
\AgdaSymbol{:}\AgdaSpace{}%
\AgdaSymbol{∀}\AgdaSpace{}%
\AgdaSymbol{\{}\AgdaBound{Γ}\AgdaSpace{}%
\AgdaBound{A}\AgdaSpace{}%
\AgdaBound{B}\AgdaSpace{}%
\AgdaBound{Δ}\AgdaSymbol{\}}\<%
\\
\>[0][@{}l@{\AgdaIndent{0}}]%
\>[2]\AgdaSymbol{→}\AgdaSpace{}%
\AgdaSymbol{(}\AgdaBound{fb'}\AgdaSpace{}%
\AgdaSymbol{:}\AgdaSpace{}%
\AgdaSymbol{∀}\AgdaSpace{}%
\AgdaSymbol{\{}\AgdaBound{Ab}\AgdaSymbol{\}}\AgdaSpace{}%
\AgdaSymbol{→}\AgdaSpace{}%
\AgdaBound{Ab}\AgdaSpace{}%
\AgdaOperator{\AgdaDatatype{∈b}}\AgdaSpace{}%
\AgdaBound{Γ}\AgdaSpace{}%
\AgdaSymbol{→}\AgdaSpace{}%
\AgdaBound{Ab}\AgdaSpace{}%
\AgdaOperator{\AgdaDatatype{∈b}}\AgdaSpace{}%
\AgdaBound{Δ}\AgdaSymbol{)}\<%
\\
%
\>[2]\AgdaSymbol{→}\AgdaSpace{}%
\AgdaSymbol{(}\AgdaBound{W}\AgdaSpace{}%
\AgdaSymbol{:}\AgdaSpace{}%
\AgdaBound{Γ}\AgdaSpace{}%
\AgdaOperator{\AgdaDatatype{⊢b}}\AgdaSpace{}%
\AgdaBound{B}\AgdaSymbol{)}\<%
\\
%
\>[2]\AgdaSymbol{→}\AgdaSpace{}%
\AgdaSymbol{(}\AgdaBound{x}\AgdaSpace{}%
\AgdaSymbol{:}\AgdaSpace{}%
\AgdaInductiveConstructor{□}\AgdaSpace{}%
\AgdaBound{A}\AgdaSpace{}%
\AgdaOperator{\AgdaDatatype{∈b}}\AgdaSpace{}%
\AgdaBound{Γ}\AgdaSpace{}%
\AgdaOperator{\AgdaInductiveConstructor{,b}}\AgdaSpace{}%
\AgdaInductiveConstructor{□}\AgdaSpace{}%
\AgdaBound{B}\AgdaSymbol{)}\<%
\\
%
\>[2]\AgdaSymbol{→}\AgdaSpace{}%
\AgdaFunction{fb}\AgdaSpace{}%
\AgdaSymbol{(}\AgdaFunction{renameb}\AgdaSpace{}%
\AgdaBound{fb'}\AgdaSpace{}%
\AgdaBound{W}\AgdaSymbol{)}\AgdaSpace{}%
\AgdaSymbol{(}\AgdaFunction{extb}\AgdaSpace{}%
\AgdaBound{fb'}\AgdaSpace{}%
\AgdaBound{x}\AgdaSymbol{)}\AgdaSpace{}%
\AgdaOperator{\AgdaDatatype{≡}}\AgdaSpace{}%
\AgdaFunction{renameb}\AgdaSpace{}%
\AgdaBound{fb'}\AgdaSpace{}%
\AgdaSymbol{(}\AgdaFunction{fb}\AgdaSpace{}%
\AgdaBound{W}\AgdaSpace{}%
\AgdaBound{x}\AgdaSymbol{)}\<%
\\
\>[0]\AgdaFunction{fb-renameb-extb-commute}\AgdaSpace{}%
\AgdaBound{fb'}\AgdaSpace{}%
\AgdaBound{W}\AgdaSpace{}%
\AgdaInductiveConstructor{Zb}\AgdaSpace{}%
\AgdaSymbol{=}\AgdaSpace{}%
\AgdaInductiveConstructor{refl}\<%
\\
\>[0]\AgdaFunction{fb-renameb-extb-commute}\AgdaSpace{}%
\AgdaBound{fb'}\AgdaSpace{}%
\AgdaBound{W}\AgdaSpace{}%
\AgdaSymbol{(}\AgdaInductiveConstructor{Sb}\AgdaSpace{}%
\AgdaBound{x}\AgdaSymbol{)}\AgdaSpace{}%
\AgdaSymbol{=}\AgdaSpace{}%
\AgdaInductiveConstructor{refl}\<%
\\
%
\\[\AgdaEmptyExtraSkip]%
\>[0]\AgdaComment{--\ Preservation\ and\ reflection\ of\ }\<%
\\
\>[0]\AgdaComment{--\ single\ substitution\ under\ renaming}\<%
\\
\>[0]\AgdaFunction{pres-subst-func-rename}\AgdaSpace{}%
\AgdaSymbol{:}\AgdaSpace{}%
\AgdaSymbol{∀}\AgdaSpace{}%
\AgdaSymbol{\{}\AgdaBound{Γ}\AgdaSpace{}%
\AgdaBound{B}\AgdaSpace{}%
\AgdaBound{A}\AgdaSpace{}%
\AgdaBound{Δ}\AgdaSymbol{\}}\<%
\\
\>[0][@{}l@{\AgdaIndent{0}}]%
\>[2]\AgdaSymbol{→}\AgdaSpace{}%
\AgdaSymbol{(}\AgdaBound{fb'}\AgdaSpace{}%
\AgdaSymbol{:}\AgdaSpace{}%
\AgdaSymbol{∀}\AgdaSpace{}%
\AgdaSymbol{\{}\AgdaBound{Ab}\AgdaSymbol{\}}\AgdaSpace{}%
\AgdaSymbol{→}\AgdaSpace{}%
\AgdaBound{Ab}\AgdaSpace{}%
\AgdaOperator{\AgdaDatatype{∈b}}\AgdaSpace{}%
\AgdaBound{Γ}\AgdaSpace{}%
\AgdaSymbol{→}\AgdaSpace{}%
\AgdaBound{Ab}\AgdaSpace{}%
\AgdaOperator{\AgdaDatatype{∈b}}\AgdaSpace{}%
\AgdaBound{Δ}\AgdaSymbol{)}\<%
\\
%
\>[2]\AgdaSymbol{→}\AgdaSpace{}%
\AgdaSymbol{(}\AgdaBound{M}\AgdaSpace{}%
\AgdaSymbol{:}\AgdaSpace{}%
\AgdaBound{Γ}\AgdaSpace{}%
\AgdaOperator{\AgdaInductiveConstructor{,b}}\AgdaSpace{}%
\AgdaInductiveConstructor{□}\AgdaSpace{}%
\AgdaBound{B}\AgdaSpace{}%
\AgdaOperator{\AgdaDatatype{⊢b}}\AgdaSpace{}%
\AgdaBound{A}\AgdaSymbol{)}\<%
\\
%
\>[2]\AgdaSymbol{→}\AgdaSpace{}%
\AgdaSymbol{(}\AgdaBound{W}\AgdaSpace{}%
\AgdaSymbol{:}\AgdaSpace{}%
\AgdaBound{Γ}\AgdaSpace{}%
\AgdaOperator{\AgdaDatatype{⊢b}}\AgdaSpace{}%
\AgdaBound{B}\AgdaSymbol{)}\<%
\\
%
\>[2]\AgdaSymbol{→}%
\>[2253I]\AgdaFunction{renameb}\AgdaSpace{}%
\AgdaSymbol{(}\AgdaFunction{extb}\AgdaSpace{}%
\AgdaBound{fb'}\AgdaSymbol{)}\AgdaSpace{}%
\AgdaBound{M}\AgdaSpace{}%
\AgdaOperator{\AgdaFunction{[}}\AgdaSpace{}%
\AgdaFunction{renameb}\AgdaSpace{}%
\AgdaBound{fb'}\AgdaSpace{}%
\AgdaBound{W}\AgdaSpace{}%
\AgdaOperator{\AgdaFunction{]b}}\<%
\\
\>[.][@{}l@{}]\<[2253I]%
\>[4]\AgdaOperator{\AgdaDatatype{≡}}\AgdaSpace{}%
\AgdaFunction{renameb}\AgdaSpace{}%
\AgdaBound{fb'}\AgdaSpace{}%
\AgdaSymbol{(}\AgdaBound{M}\AgdaSpace{}%
\AgdaOperator{\AgdaFunction{[}}\AgdaSpace{}%
\AgdaBound{W}\AgdaSpace{}%
\AgdaOperator{\AgdaFunction{]b}}\AgdaSymbol{)}\<%
\\
\>[0]\AgdaFunction{pres-subst-func-rename}\AgdaSpace{}%
\AgdaBound{fb'}\AgdaSpace{}%
\AgdaBound{M}\AgdaSpace{}%
\AgdaBound{W}\AgdaSpace{}%
\AgdaSymbol{=}\<%
\\
\>[0][@{}l@{\AgdaIndent{0}}]%
\>[2]\AgdaOperator{\AgdaFunction{beginEq}}\<%
\\
\>[2][@{}l@{\AgdaIndent{0}}]%
\>[4]\AgdaFunction{substb}\AgdaSpace{}%
\AgdaSymbol{(}\AgdaFunction{fb}\AgdaSpace{}%
\AgdaSymbol{(}\AgdaFunction{renameb}\AgdaSpace{}%
\AgdaBound{fb'}\AgdaSpace{}%
\AgdaBound{W}\AgdaSymbol{))}\AgdaSpace{}%
\AgdaSymbol{(}\AgdaFunction{renameb}\AgdaSpace{}%
\AgdaSymbol{(}\AgdaFunction{extb}\AgdaSpace{}%
\AgdaBound{fb'}\AgdaSymbol{)}\AgdaSpace{}%
\AgdaBound{M}\AgdaSymbol{)}\<%
\\
%
\>[2]\AgdaFunction{≡⟨}\AgdaSpace{}%
\AgdaFunction{subst-renameb-comp}\AgdaSpace{}%
\AgdaSymbol{(}\AgdaFunction{extb}\AgdaSpace{}%
\AgdaBound{fb'}\AgdaSymbol{)}\AgdaSpace{}%
\AgdaSymbol{(}\AgdaFunction{fb}\AgdaSpace{}%
\AgdaSymbol{(}\AgdaFunction{renameb}\AgdaSpace{}%
\AgdaBound{fb'}\AgdaSpace{}%
\AgdaBound{W}\AgdaSymbol{))}\AgdaSpace{}%
\AgdaBound{M}\AgdaSpace{}%
\AgdaFunction{⟩}\<%
\\
\>[2][@{}l@{\AgdaIndent{0}}]%
\>[4]\AgdaFunction{substb}\AgdaSpace{}%
\AgdaSymbol{(λ}\AgdaSpace{}%
\AgdaBound{x}\AgdaSpace{}%
\AgdaSymbol{→}\AgdaSpace{}%
\AgdaFunction{fb}\AgdaSpace{}%
\AgdaSymbol{(}\AgdaFunction{renameb}\AgdaSpace{}%
\AgdaBound{fb'}\AgdaSpace{}%
\AgdaBound{W}\AgdaSymbol{)}\AgdaSpace{}%
\AgdaSymbol{(}\AgdaFunction{extb}\AgdaSpace{}%
\AgdaBound{fb'}\AgdaSpace{}%
\AgdaBound{x}\AgdaSymbol{))}\AgdaSpace{}%
\AgdaBound{M}\<%
\\
%
\>[2]\AgdaFunction{≡⟨}\AgdaSpace{}%
\AgdaFunction{substb-cong}\AgdaSpace{}%
\AgdaSymbol{(}\AgdaFunction{fb-renameb-extb-commute}\AgdaSpace{}%
\AgdaBound{fb'}\AgdaSpace{}%
\AgdaBound{W}\AgdaSymbol{)}\AgdaSpace{}%
\AgdaBound{M}\AgdaSpace{}%
\AgdaFunction{⟩}\<%
\\
\>[2][@{}l@{\AgdaIndent{0}}]%
\>[4]\AgdaFunction{substb}\AgdaSpace{}%
\AgdaSymbol{(λ}\AgdaSpace{}%
\AgdaBound{x}\AgdaSpace{}%
\AgdaSymbol{→}\AgdaSpace{}%
\AgdaFunction{renameb}\AgdaSpace{}%
\AgdaBound{fb'}\AgdaSpace{}%
\AgdaSymbol{(}\AgdaFunction{fb}\AgdaSpace{}%
\AgdaBound{W}\AgdaSpace{}%
\AgdaBound{x}\AgdaSymbol{))}\AgdaSpace{}%
\AgdaBound{M}\<%
\\
%
\>[2]\AgdaFunction{≡⟨}\AgdaSpace{}%
\AgdaFunction{substb-rename-commute}\AgdaSpace{}%
\AgdaBound{fb'}\AgdaSpace{}%
\AgdaSymbol{(}\AgdaFunction{fb}\AgdaSpace{}%
\AgdaBound{W}\AgdaSymbol{)}\AgdaSpace{}%
\AgdaBound{M}\AgdaSpace{}%
\AgdaFunction{⟩}\<%
\\
\>[2][@{}l@{\AgdaIndent{0}}]%
\>[4]\AgdaFunction{renameb}\AgdaSpace{}%
\AgdaBound{fb'}\AgdaSpace{}%
\AgdaSymbol{(}\AgdaFunction{substb}\AgdaSpace{}%
\AgdaSymbol{(}\AgdaFunction{fb}\AgdaSpace{}%
\AgdaBound{W}\AgdaSymbol{)}\AgdaSpace{}%
\AgdaBound{M}\AgdaSymbol{)}\<%
\\
%
\>[2]\AgdaOperator{\AgdaFunction{∎eq}}\<%
\\
%
\\[\AgdaEmptyExtraSkip]%
\>[0]\AgdaComment{--\ Preservation\ of\ reduction\ under\ renaming\ \ \ }\<%
\\
\>[0]\AgdaFunction{pres-rename-red}\AgdaSpace{}%
\AgdaSymbol{:}\AgdaSpace{}%
\AgdaSymbol{∀}\AgdaSpace{}%
\AgdaSymbol{\{}\AgdaBound{Γ}\AgdaSpace{}%
\AgdaBound{A}\AgdaSpace{}%
\AgdaBound{Δ}\AgdaSymbol{\}}\<%
\\
\>[0][@{}l@{\AgdaIndent{0}}]%
\>[2]\AgdaSymbol{→}\AgdaSpace{}%
\AgdaSymbol{(}\AgdaBound{fb'}\AgdaSpace{}%
\AgdaSymbol{:}\AgdaSpace{}%
\AgdaSymbol{∀}\AgdaSpace{}%
\AgdaSymbol{\{}\AgdaBound{Ab}\AgdaSymbol{\}}\AgdaSpace{}%
\AgdaSymbol{→}\AgdaSpace{}%
\AgdaBound{Ab}\AgdaSpace{}%
\AgdaOperator{\AgdaDatatype{∈b}}\AgdaSpace{}%
\AgdaBound{Γ}\AgdaSpace{}%
\AgdaSymbol{→}\AgdaSpace{}%
\AgdaBound{Ab}\AgdaSpace{}%
\AgdaOperator{\AgdaDatatype{∈b}}\AgdaSpace{}%
\AgdaBound{Δ}\AgdaSymbol{)}\<%
\\
%
\>[2]\AgdaSymbol{→}\AgdaSpace{}%
\AgdaSymbol{(}\AgdaBound{M}\AgdaSpace{}%
\AgdaSymbol{:}\AgdaSpace{}%
\AgdaBound{Γ}\AgdaSpace{}%
\AgdaOperator{\AgdaDatatype{⊢b}}\AgdaSpace{}%
\AgdaBound{A}\AgdaSymbol{)}\<%
\\
%
\>[2]\AgdaSymbol{→}\AgdaSpace{}%
\AgdaSymbol{(}\AgdaBound{N}\AgdaSpace{}%
\AgdaSymbol{:}\AgdaSpace{}%
\AgdaBound{Γ}\AgdaSpace{}%
\AgdaOperator{\AgdaDatatype{⊢b}}\AgdaSpace{}%
\AgdaBound{A}\AgdaSymbol{)}\<%
\\
%
\>[2]\AgdaSymbol{→}\AgdaSpace{}%
\AgdaSymbol{(}\AgdaBound{M}\AgdaSpace{}%
\AgdaOperator{\AgdaDatatype{↝βb}}\AgdaSpace{}%
\AgdaBound{N}\AgdaSymbol{)}\<%
\\
%
\>[2]\AgdaSymbol{→}\AgdaSpace{}%
\AgdaFunction{renameb}\AgdaSpace{}%
\AgdaBound{fb'}\AgdaSpace{}%
\AgdaBound{M}\AgdaSpace{}%
\AgdaOperator{\AgdaDatatype{↝βb}}\AgdaSpace{}%
\AgdaFunction{renameb}\AgdaSpace{}%
\AgdaBound{fb'}\AgdaSpace{}%
\AgdaBound{N}\<%
\\
\>[0]\AgdaFunction{pres-rename-red}\AgdaSpace{}%
\AgdaBound{fb'}\AgdaSpace{}%
\AgdaSymbol{(}\AgdaInductiveConstructor{ƛb}\AgdaSpace{}%
\AgdaBound{M₁}\AgdaSpace{}%
\AgdaOperator{\AgdaInductiveConstructor{·b}}\AgdaSpace{}%
\AgdaInductiveConstructor{box}\AgdaSpace{}%
\AgdaBound{M₂}\AgdaSymbol{)}\AgdaSpace{}%
\AgdaBound{N}\AgdaSpace{}%
\AgdaInductiveConstructor{βb}\AgdaSpace{}%
\AgdaSymbol{=}\AgdaSpace{}%
\AgdaFunction{Eq.subst}\AgdaSpace{}%
\AgdaSymbol{(λ}\AgdaSpace{}%
\AgdaBound{x}\AgdaSpace{}%
\AgdaSymbol{→}\AgdaSpace{}%
\AgdaInductiveConstructor{ƛb}\AgdaSpace{}%
\AgdaSymbol{(}\AgdaFunction{renameb}\AgdaSpace{}%
\AgdaSymbol{(}\AgdaFunction{extb}\AgdaSpace{}%
\AgdaBound{fb'}\AgdaSymbol{)}\AgdaSpace{}%
\AgdaBound{M₁}\AgdaSymbol{)}\AgdaSpace{}%
\AgdaOperator{\AgdaInductiveConstructor{·b}}\AgdaSpace{}%
\AgdaInductiveConstructor{box}\AgdaSpace{}%
\AgdaSymbol{(}\AgdaFunction{renameb}\AgdaSpace{}%
\AgdaBound{fb'}\AgdaSpace{}%
\AgdaBound{M₂}\AgdaSymbol{)}\AgdaSpace{}%
\AgdaOperator{\AgdaDatatype{↝βb}}\AgdaSpace{}%
\AgdaBound{x}\AgdaSymbol{)}\AgdaSpace{}%
\AgdaSymbol{(}\AgdaFunction{pres-subst-func-rename}\AgdaSpace{}%
\AgdaBound{fb'}\AgdaSpace{}%
\AgdaBound{M₁}\AgdaSpace{}%
\AgdaBound{M₂}\AgdaSymbol{)}\AgdaSpace{}%
\AgdaInductiveConstructor{βb}\<%
\\
\>[0]\AgdaFunction{pres-rename-red}\AgdaSpace{}%
\AgdaBound{fb'}\AgdaSpace{}%
\AgdaBound{M}\AgdaSpace{}%
\AgdaBound{N}\AgdaSpace{}%
\AgdaSymbol{(}\AgdaInductiveConstructor{μ}\AgdaSpace{}%
\AgdaSymbol{\{}\AgdaArgument{M}\AgdaSpace{}%
\AgdaSymbol{=}\AgdaSpace{}%
\AgdaBound{M₁}\AgdaSymbol{\}}\AgdaSpace{}%
\AgdaSymbol{\{}\AgdaArgument{M'}\AgdaSpace{}%
\AgdaSymbol{=}\AgdaSpace{}%
\AgdaBound{M'}\AgdaSymbol{\}}\AgdaSpace{}%
\AgdaBound{red}\AgdaSymbol{)}\AgdaSpace{}%
\AgdaSymbol{=}\AgdaSpace{}%
\AgdaInductiveConstructor{μ}\AgdaSpace{}%
\AgdaSymbol{(}\AgdaFunction{pres-rename-red}\AgdaSpace{}%
\AgdaBound{fb'}\AgdaSpace{}%
\AgdaBound{M₁}\AgdaSpace{}%
\AgdaBound{M'}\AgdaSpace{}%
\AgdaBound{red}\AgdaSymbol{)}\<%
\\
\>[0]\AgdaFunction{pres-rename-red}\AgdaSpace{}%
\AgdaBound{fb'}\AgdaSpace{}%
\AgdaBound{M}\AgdaSpace{}%
\AgdaBound{N}\AgdaSpace{}%
\AgdaSymbol{(}\AgdaInductiveConstructor{ν}\AgdaSpace{}%
\AgdaSymbol{\{}\AgdaArgument{M}\AgdaSpace{}%
\AgdaSymbol{=}\AgdaSpace{}%
\AgdaBound{M₁}\AgdaSymbol{\}}\AgdaSpace{}%
\AgdaSymbol{\{}\AgdaArgument{N}\AgdaSpace{}%
\AgdaSymbol{=}\AgdaSpace{}%
\AgdaBound{N₁}\AgdaSymbol{\}}\AgdaSpace{}%
\AgdaSymbol{\{}\AgdaArgument{N'}\AgdaSpace{}%
\AgdaSymbol{=}\AgdaSpace{}%
\AgdaBound{N'}\AgdaSymbol{\}}\AgdaSpace{}%
\AgdaBound{red}\AgdaSymbol{)}\AgdaSpace{}%
\AgdaSymbol{=}\AgdaSpace{}%
\AgdaInductiveConstructor{ν}\AgdaSpace{}%
\AgdaSymbol{(}\AgdaFunction{pres-rename-red}\AgdaSpace{}%
\AgdaBound{fb'}\AgdaSpace{}%
\AgdaBound{N₁}\AgdaSpace{}%
\AgdaBound{N'}\AgdaSpace{}%
\AgdaBound{red}\AgdaSymbol{)}\<%
\\
\>[0]\AgdaFunction{pres-rename-red}\AgdaSpace{}%
\AgdaBound{fb'}\AgdaSpace{}%
\AgdaBound{M}\AgdaSpace{}%
\AgdaBound{N}\AgdaSpace{}%
\AgdaSymbol{(}\AgdaInductiveConstructor{ξ}\AgdaSpace{}%
\AgdaSymbol{\{}\AgdaArgument{M}\AgdaSpace{}%
\AgdaSymbol{=}\AgdaSpace{}%
\AgdaBound{M₁}\AgdaSymbol{\}}\AgdaSpace{}%
\AgdaSymbol{\{}\AgdaArgument{M'}\AgdaSpace{}%
\AgdaSymbol{=}\AgdaSpace{}%
\AgdaBound{M'}\AgdaSymbol{\}}\AgdaSpace{}%
\AgdaBound{red}\AgdaSymbol{)}\AgdaSpace{}%
\AgdaSymbol{=}\AgdaSpace{}%
\AgdaInductiveConstructor{ξ}\AgdaSpace{}%
\AgdaSymbol{(}\AgdaFunction{pres-rename-red}\AgdaSpace{}%
\AgdaSymbol{(}\AgdaFunction{extb}\AgdaSpace{}%
\AgdaBound{fb'}\AgdaSymbol{)}\AgdaSpace{}%
\AgdaBound{M₁}\AgdaSpace{}%
\AgdaBound{M'}\AgdaSpace{}%
\AgdaBound{red}\AgdaSymbol{)}\<%
\\
\>[0]\AgdaFunction{pres-rename-red}\AgdaSpace{}%
\AgdaBound{fb'}\AgdaSpace{}%
\AgdaBound{M}\AgdaSpace{}%
\AgdaBound{N}\AgdaSpace{}%
\AgdaSymbol{(}\AgdaInductiveConstructor{ζ}\AgdaSpace{}%
\AgdaSymbol{\{}\AgdaArgument{M}\AgdaSpace{}%
\AgdaSymbol{=}\AgdaSpace{}%
\AgdaBound{M₁}\AgdaSymbol{\}}\AgdaSpace{}%
\AgdaSymbol{\{}\AgdaArgument{M'}\AgdaSpace{}%
\AgdaSymbol{=}\AgdaSpace{}%
\AgdaBound{M'}\AgdaSymbol{\}}\AgdaSpace{}%
\AgdaBound{red}\AgdaSymbol{)}\AgdaSpace{}%
\AgdaSymbol{=}\AgdaSpace{}%
\AgdaInductiveConstructor{ζ}\AgdaSpace{}%
\AgdaSymbol{(}\AgdaFunction{pres-rename-red}\AgdaSpace{}%
\AgdaBound{fb'}\AgdaSpace{}%
\AgdaBound{M₁}\AgdaSpace{}%
\AgdaBound{M'}\AgdaSpace{}%
\AgdaBound{red}\AgdaSymbol{)}\<%
\\
%
\\[\AgdaEmptyExtraSkip]%
\>[0]\AgdaComment{--\ Preservation\ of\ 2\ step\ reduction}\<%
\\
\>[0]\AgdaComment{--\ under\ renaming}\<%
\\
\>[0]\AgdaFunction{pres-rename-red-2}\AgdaSpace{}%
\AgdaSymbol{:}\AgdaSpace{}%
\AgdaSymbol{∀}\AgdaSpace{}%
\AgdaSymbol{\{}\AgdaBound{Γ}\AgdaSpace{}%
\AgdaBound{A}\AgdaSpace{}%
\AgdaBound{Δ}\AgdaSymbol{\}}\<%
\\
\>[0][@{}l@{\AgdaIndent{0}}]%
\>[2]\AgdaSymbol{→}\AgdaSpace{}%
\AgdaSymbol{(}\AgdaBound{fb'}\AgdaSpace{}%
\AgdaSymbol{:}\AgdaSpace{}%
\AgdaSymbol{∀}\AgdaSpace{}%
\AgdaSymbol{\{}\AgdaBound{Ab}\AgdaSymbol{\}}\AgdaSpace{}%
\AgdaSymbol{→}\AgdaSpace{}%
\AgdaBound{Ab}\AgdaSpace{}%
\AgdaOperator{\AgdaDatatype{∈b}}\AgdaSpace{}%
\AgdaBound{Γ}\AgdaSpace{}%
\AgdaSymbol{→}\AgdaSpace{}%
\AgdaBound{Ab}\AgdaSpace{}%
\AgdaOperator{\AgdaDatatype{∈b}}\AgdaSpace{}%
\AgdaBound{Δ}\AgdaSymbol{)}\<%
\\
%
\>[2]\AgdaSymbol{→}\AgdaSpace{}%
\AgdaSymbol{(}\AgdaBound{M}\AgdaSpace{}%
\AgdaSymbol{:}\AgdaSpace{}%
\AgdaBound{Γ}\AgdaSpace{}%
\AgdaOperator{\AgdaDatatype{⊢b}}\AgdaSpace{}%
\AgdaBound{A}\AgdaSymbol{)}\<%
\\
%
\>[2]\AgdaSymbol{→}\AgdaSpace{}%
\AgdaSymbol{(}\AgdaBound{N}\AgdaSpace{}%
\AgdaSymbol{:}\AgdaSpace{}%
\AgdaBound{Γ}\AgdaSpace{}%
\AgdaOperator{\AgdaDatatype{⊢b}}\AgdaSpace{}%
\AgdaBound{A}\AgdaSymbol{)}\<%
\\
%
\>[2]\AgdaSymbol{→}\AgdaSpace{}%
\AgdaSymbol{(}\AgdaBound{M}\AgdaSpace{}%
\AgdaOperator{\AgdaDatatype{↝βb²}}\AgdaSpace{}%
\AgdaBound{N}\AgdaSymbol{)}\<%
\\
%
\>[2]\AgdaSymbol{→}\AgdaSpace{}%
\AgdaFunction{renameb}\AgdaSpace{}%
\AgdaBound{fb'}\AgdaSpace{}%
\AgdaBound{M}\AgdaSpace{}%
\AgdaOperator{\AgdaDatatype{↝βb²}}\AgdaSpace{}%
\AgdaFunction{renameb}\AgdaSpace{}%
\AgdaBound{fb'}\AgdaSpace{}%
\AgdaBound{N}\<%
\\
\>[0]\AgdaFunction{pres-rename-red-2}\AgdaSpace{}%
\AgdaBound{fb'}\AgdaSpace{}%
\AgdaBound{M}\AgdaSpace{}%
\AgdaBound{N}\AgdaSpace{}%
\AgdaSymbol{(}\AgdaInductiveConstructor{βb²}\AgdaSpace{}%
\AgdaDottedPattern{\AgdaSymbol{.}}\AgdaDottedPattern{\AgdaBound{M}}\AgdaSpace{}%
\AgdaBound{t₂}\AgdaSpace{}%
\AgdaDottedPattern{\AgdaSymbol{.}}\AgdaDottedPattern{\AgdaBound{N}}\AgdaSpace{}%
\AgdaBound{x}\AgdaSpace{}%
\AgdaBound{x₁}\AgdaSymbol{)}\AgdaSpace{}%
\AgdaSymbol{=}\AgdaSpace{}%
\AgdaInductiveConstructor{βb²}\AgdaSpace{}%
\AgdaSymbol{(}\AgdaFunction{renameb}\AgdaSpace{}%
\AgdaBound{fb'}\AgdaSpace{}%
\AgdaBound{M}\AgdaSymbol{)}\AgdaSpace{}%
\AgdaSymbol{(}\AgdaFunction{renameb}\AgdaSpace{}%
\AgdaBound{fb'}\AgdaSpace{}%
\AgdaBound{t₂}\AgdaSymbol{)}\AgdaSpace{}%
\AgdaSymbol{(}\AgdaFunction{renameb}\AgdaSpace{}%
\AgdaBound{fb'}\AgdaSpace{}%
\AgdaBound{N}\AgdaSymbol{)}\AgdaSpace{}%
\AgdaSymbol{(}\AgdaFunction{pres-rename-red}\AgdaSpace{}%
\AgdaBound{fb'}\AgdaSpace{}%
\AgdaBound{M}\AgdaSpace{}%
\AgdaBound{t₂}\AgdaSpace{}%
\AgdaBound{x}\AgdaSymbol{)}\AgdaSpace{}%
\AgdaSymbol{(}\AgdaFunction{pres-rename-red}\AgdaSpace{}%
\AgdaBound{fb'}\AgdaSpace{}%
\AgdaBound{t₂}\AgdaSpace{}%
\AgdaBound{N}\AgdaSpace{}%
\AgdaBound{x₁}\AgdaSymbol{)}\<%
\\
\>[0]\AgdaFunction{pres-rename-red-2}\AgdaSpace{}%
\AgdaBound{fb'}\AgdaSpace{}%
\AgdaBound{M}\AgdaSpace{}%
\AgdaBound{N}\AgdaSpace{}%
\AgdaSymbol{(}\AgdaInductiveConstructor{μ}\AgdaSpace{}%
\AgdaSymbol{\{}\AgdaArgument{M}\AgdaSpace{}%
\AgdaSymbol{=}\AgdaSpace{}%
\AgdaBound{M₁}\AgdaSymbol{\}}\AgdaSpace{}%
\AgdaSymbol{\{}\AgdaArgument{M'}\AgdaSpace{}%
\AgdaSymbol{=}\AgdaSpace{}%
\AgdaBound{M'}\AgdaSymbol{\}}\AgdaSpace{}%
\AgdaBound{red}\AgdaSymbol{)}\AgdaSpace{}%
\AgdaSymbol{=}\AgdaSpace{}%
\AgdaInductiveConstructor{μ}\AgdaSpace{}%
\AgdaSymbol{(}\AgdaFunction{pres-rename-red-2}\AgdaSpace{}%
\AgdaBound{fb'}\AgdaSpace{}%
\AgdaBound{M₁}\AgdaSpace{}%
\AgdaBound{M'}\AgdaSpace{}%
\AgdaBound{red}\AgdaSymbol{)}\<%
\\
\>[0]\AgdaFunction{pres-rename-red-2}\AgdaSpace{}%
\AgdaBound{fb'}\AgdaSpace{}%
\AgdaBound{M}\AgdaSpace{}%
\AgdaBound{N}\AgdaSpace{}%
\AgdaSymbol{(}\AgdaInductiveConstructor{ν}\AgdaSpace{}%
\AgdaSymbol{\{}\AgdaArgument{M}\AgdaSpace{}%
\AgdaSymbol{=}\AgdaSpace{}%
\AgdaBound{M₁}\AgdaSymbol{\}}\AgdaSpace{}%
\AgdaSymbol{\{}\AgdaArgument{N}\AgdaSpace{}%
\AgdaSymbol{=}\AgdaSpace{}%
\AgdaBound{N₁}\AgdaSymbol{\}}\AgdaSpace{}%
\AgdaSymbol{\{}\AgdaArgument{N'}\AgdaSpace{}%
\AgdaSymbol{=}\AgdaSpace{}%
\AgdaBound{N'}\AgdaSymbol{\}}\AgdaSpace{}%
\AgdaBound{red}\AgdaSymbol{)}\AgdaSpace{}%
\AgdaSymbol{=}\AgdaSpace{}%
\AgdaInductiveConstructor{ν}\AgdaSpace{}%
\AgdaSymbol{(}\AgdaFunction{pres-rename-red-2}\AgdaSpace{}%
\AgdaBound{fb'}\AgdaSpace{}%
\AgdaBound{N₁}\AgdaSpace{}%
\AgdaBound{N'}\AgdaSpace{}%
\AgdaBound{red}\AgdaSymbol{)}\<%
\\
\>[0]\AgdaFunction{pres-rename-red-2}\AgdaSpace{}%
\AgdaBound{fb'}\AgdaSpace{}%
\AgdaBound{M}\AgdaSpace{}%
\AgdaBound{N}\AgdaSpace{}%
\AgdaSymbol{(}\AgdaInductiveConstructor{ξ}\AgdaSpace{}%
\AgdaSymbol{\{}\AgdaArgument{M}\AgdaSpace{}%
\AgdaSymbol{=}\AgdaSpace{}%
\AgdaBound{M₁}\AgdaSymbol{\}}\AgdaSpace{}%
\AgdaSymbol{\{}\AgdaArgument{M'}\AgdaSpace{}%
\AgdaSymbol{=}\AgdaSpace{}%
\AgdaBound{M'}\AgdaSymbol{\}}\AgdaSpace{}%
\AgdaBound{red}\AgdaSymbol{)}\AgdaSpace{}%
\AgdaSymbol{=}\AgdaSpace{}%
\AgdaInductiveConstructor{ξ}\AgdaSpace{}%
\AgdaSymbol{(}\AgdaFunction{pres-rename-red-2}\AgdaSpace{}%
\AgdaSymbol{(}\AgdaFunction{extb}\AgdaSpace{}%
\AgdaBound{fb'}\AgdaSymbol{)}\AgdaSpace{}%
\AgdaBound{M₁}\AgdaSpace{}%
\AgdaBound{M'}\AgdaSpace{}%
\AgdaBound{red}\AgdaSymbol{)}\<%
\\
\>[0]\AgdaFunction{pres-rename-red-2}\AgdaSpace{}%
\AgdaBound{fb'}\AgdaSpace{}%
\AgdaBound{M}\AgdaSpace{}%
\AgdaBound{N}\AgdaSpace{}%
\AgdaSymbol{(}\AgdaInductiveConstructor{ζ}\AgdaSpace{}%
\AgdaSymbol{\{}\AgdaArgument{M}\AgdaSpace{}%
\AgdaSymbol{=}\AgdaSpace{}%
\AgdaBound{M₁}\AgdaSymbol{\}}\AgdaSpace{}%
\AgdaSymbol{\{}\AgdaArgument{M'}\AgdaSpace{}%
\AgdaSymbol{=}\AgdaSpace{}%
\AgdaBound{M'}\AgdaSymbol{\}}\AgdaSpace{}%
\AgdaBound{red}\AgdaSymbol{)}\AgdaSpace{}%
\AgdaSymbol{=}\AgdaSpace{}%
\AgdaInductiveConstructor{ζ}\AgdaSpace{}%
\AgdaSymbol{(}\AgdaFunction{pres-rename-red-2}\AgdaSpace{}%
\AgdaBound{fb'}\AgdaSpace{}%
\AgdaBound{M₁}\AgdaSpace{}%
\AgdaBound{M'}\AgdaSpace{}%
\AgdaBound{red}\AgdaSymbol{)}\<%
\\
%
\\[\AgdaEmptyExtraSkip]%
\>[0]\AgdaComment{--\ extsb\ of\ ε\ is\ ε}\<%
\\
\>[0]\AgdaFunction{ε-exts}\AgdaSpace{}%
\AgdaSymbol{:}\AgdaSpace{}%
\AgdaSymbol{∀}\AgdaSpace{}%
\AgdaSymbol{\{}\AgdaBound{Γ}\AgdaSpace{}%
\AgdaBound{A}\AgdaSpace{}%
\AgdaBound{B}\AgdaSymbol{\}}\<%
\\
\>[0][@{}l@{\AgdaIndent{0}}]%
\>[2]\AgdaSymbol{→}\AgdaSpace{}%
\AgdaSymbol{(}\AgdaBound{x}\AgdaSpace{}%
\AgdaSymbol{:}\AgdaSpace{}%
\AgdaInductiveConstructor{□}\AgdaSpace{}%
\AgdaBound{A}\AgdaSpace{}%
\AgdaOperator{\AgdaDatatype{∈b}}\AgdaSpace{}%
\AgdaBound{Γ}\AgdaSpace{}%
\AgdaOperator{\AgdaInductiveConstructor{,b}}\AgdaSpace{}%
\AgdaBound{B}\AgdaSymbol{)}\<%
\\
%
\>[2]\AgdaSymbol{→}\AgdaSpace{}%
\AgdaFunction{extsb}\AgdaSpace{}%
\AgdaInductiveConstructor{ε}\AgdaSpace{}%
\AgdaBound{x}\AgdaSpace{}%
\AgdaOperator{\AgdaDatatype{≡}}\AgdaSpace{}%
\AgdaInductiveConstructor{ε}\AgdaSpace{}%
\AgdaBound{x}\<%
\\
\>[0]\AgdaFunction{ε-exts}\AgdaSpace{}%
\AgdaInductiveConstructor{Zb}\AgdaSpace{}%
\AgdaSymbol{=}\AgdaSpace{}%
\AgdaInductiveConstructor{refl}\<%
\\
\>[0]\AgdaFunction{ε-exts}\AgdaSpace{}%
\AgdaSymbol{(}\AgdaInductiveConstructor{Sb}\AgdaSpace{}%
\AgdaBound{x}\AgdaSymbol{)}\AgdaSpace{}%
\AgdaSymbol{=}\AgdaSpace{}%
\AgdaInductiveConstructor{refl}\<%
\\
%
\\[\AgdaEmptyExtraSkip]%
\>[0]\AgdaComment{--\ Simultaneous\ substitution\ with\ ε\ }\<%
\\
\>[0]\AgdaComment{--\ leaves\ the\ term\ unchanged\ \ }\<%
\\
\>[0]\AgdaFunction{subst-with-id}\AgdaSpace{}%
\AgdaSymbol{:}\AgdaSpace{}%
\AgdaSymbol{∀}\AgdaSpace{}%
\AgdaSymbol{\{}\AgdaBound{Γ}\AgdaSpace{}%
\AgdaBound{A}\AgdaSymbol{\}}\<%
\\
\>[0][@{}l@{\AgdaIndent{0}}]%
\>[2]\AgdaSymbol{→}\AgdaSpace{}%
\AgdaSymbol{(}\AgdaBound{M}\AgdaSpace{}%
\AgdaSymbol{:}\AgdaSpace{}%
\AgdaBound{Γ}\AgdaSpace{}%
\AgdaOperator{\AgdaDatatype{⊢b}}\AgdaSpace{}%
\AgdaBound{A}\AgdaSymbol{)}\<%
\\
%
\>[2]\AgdaSymbol{→}\AgdaSpace{}%
\AgdaFunction{substb}\AgdaSpace{}%
\AgdaSymbol{(λ}\AgdaSpace{}%
\AgdaBound{x₁}\AgdaSpace{}%
\AgdaSymbol{→}\AgdaSpace{}%
\AgdaInductiveConstructor{ε}\AgdaSpace{}%
\AgdaBound{x₁}\AgdaSymbol{)}\AgdaSpace{}%
\AgdaBound{M}\AgdaSpace{}%
\AgdaOperator{\AgdaDatatype{≡}}\AgdaSpace{}%
\AgdaBound{M}\<%
\\
\>[0]\AgdaFunction{subst-with-id}\AgdaSpace{}%
\AgdaSymbol{(}\AgdaInductiveConstructor{ε}\AgdaSpace{}%
\AgdaBound{x}\AgdaSymbol{)}\AgdaSpace{}%
\AgdaSymbol{=}\AgdaSpace{}%
\AgdaInductiveConstructor{refl}\<%
\\
\>[0]\AgdaFunction{subst-with-id}\AgdaSpace{}%
\AgdaSymbol{(}\AgdaInductiveConstructor{box}\AgdaSpace{}%
\AgdaBound{M}\AgdaSymbol{)}\AgdaSpace{}%
\AgdaSymbol{=}\AgdaSpace{}%
\AgdaFunction{cong}\AgdaSpace{}%
\AgdaInductiveConstructor{box}\AgdaSpace{}%
\AgdaSymbol{(}\AgdaFunction{subst-with-id}\AgdaSpace{}%
\AgdaBound{M}\AgdaSymbol{)}\<%
\\
\>[0]\AgdaFunction{subst-with-id}\AgdaSpace{}%
\AgdaSymbol{(}\AgdaInductiveConstructor{ƛb}\AgdaSpace{}%
\AgdaBound{M}\AgdaSymbol{)}\AgdaSpace{}%
\AgdaSymbol{=}\<%
\\
\>[0][@{}l@{\AgdaIndent{0}}]%
\>[2]\AgdaOperator{\AgdaFunction{beginEq}}\<%
\\
\>[2][@{}l@{\AgdaIndent{0}}]%
\>[4]\AgdaInductiveConstructor{ƛb}\AgdaSpace{}%
\AgdaSymbol{(}\AgdaFunction{substb}\AgdaSpace{}%
\AgdaSymbol{(}\AgdaFunction{extsb}\AgdaSpace{}%
\AgdaInductiveConstructor{ε}\AgdaSymbol{)}\AgdaSpace{}%
\AgdaBound{M}\AgdaSymbol{)}\<%
\\
%
\>[2]\AgdaFunction{≡⟨}\AgdaSpace{}%
\AgdaFunction{cong}\AgdaSpace{}%
\AgdaInductiveConstructor{ƛb}\AgdaSpace{}%
\AgdaSymbol{(}\AgdaFunction{substb-cong}\AgdaSpace{}%
\AgdaSymbol{(λ}\AgdaSpace{}%
\AgdaBound{x}\AgdaSpace{}%
\AgdaSymbol{→}\AgdaSpace{}%
\AgdaFunction{ε-exts}\AgdaSpace{}%
\AgdaBound{x}\AgdaSymbol{)}\AgdaSpace{}%
\AgdaBound{M}\AgdaSymbol{)}\AgdaSpace{}%
\AgdaFunction{⟩}\<%
\\
\>[2][@{}l@{\AgdaIndent{0}}]%
\>[4]\AgdaInductiveConstructor{ƛb}\AgdaSpace{}%
\AgdaSymbol{(}\AgdaFunction{substb}\AgdaSpace{}%
\AgdaInductiveConstructor{ε}\AgdaSpace{}%
\AgdaBound{M}\AgdaSymbol{)}\<%
\\
%
\>[2]\AgdaFunction{≡⟨}\AgdaSpace{}%
\AgdaFunction{cong}\AgdaSpace{}%
\AgdaInductiveConstructor{ƛb}\AgdaSpace{}%
\AgdaSymbol{(}\AgdaFunction{subst-with-id}\AgdaSpace{}%
\AgdaBound{M}\AgdaSymbol{)}\AgdaSpace{}%
\AgdaFunction{⟩}\<%
\\
\>[2][@{}l@{\AgdaIndent{0}}]%
\>[4]\AgdaInductiveConstructor{ƛb}\AgdaSpace{}%
\AgdaBound{M}\<%
\\
%
\>[2]\AgdaOperator{\AgdaFunction{∎eq}}\<%
\\
\>[0]\AgdaFunction{subst-with-id}\AgdaSpace{}%
\AgdaSymbol{(}\AgdaBound{M}\AgdaSpace{}%
\AgdaOperator{\AgdaInductiveConstructor{·b}}\AgdaSpace{}%
\AgdaBound{M₁}\AgdaSymbol{)}\AgdaSpace{}%
\AgdaSymbol{=}\AgdaSpace{}%
\AgdaFunction{cong₂}\AgdaSpace{}%
\AgdaOperator{\AgdaInductiveConstructor{\AgdaUnderscore{}·b\AgdaUnderscore{}}}\AgdaSpace{}%
\AgdaSymbol{(}\AgdaFunction{subst-with-id}\AgdaSpace{}%
\AgdaBound{M}\AgdaSymbol{)}\AgdaSpace{}%
\AgdaSymbol{(}\AgdaFunction{subst-with-id}\AgdaSpace{}%
\AgdaBound{M₁}\AgdaSymbol{)}\<%
\\
%
\\[\AgdaEmptyExtraSkip]%
\>[0]\AgdaComment{--\ Substitution\ composed\ with\ renaming\ with\ Sb\ }\<%
\\
\>[0]\AgdaComment{--\ is\ the\ same\ as\ embedding\ that\ term.}\<%
\\
\>[0]\AgdaFunction{subst-rename-fb-comp}\AgdaSpace{}%
\AgdaSymbol{:}\AgdaSpace{}%
\AgdaSymbol{∀}\AgdaSpace{}%
\AgdaSymbol{\{}\AgdaBound{Γ}\AgdaSpace{}%
\AgdaBound{A}\AgdaSpace{}%
\AgdaBound{B}\AgdaSymbol{\}}\<%
\\
\>[0][@{}l@{\AgdaIndent{0}}]%
\>[2]\AgdaSymbol{→}\AgdaSpace{}%
\AgdaSymbol{(}\AgdaBound{M₁}\AgdaSpace{}%
\AgdaSymbol{:}\AgdaSpace{}%
\AgdaBound{Γ}\AgdaSpace{}%
\AgdaOperator{\AgdaInductiveConstructor{,}}\AgdaSpace{}%
\AgdaBound{B}\AgdaSpace{}%
\AgdaOperator{\AgdaDatatype{⊢}}\AgdaSpace{}%
\AgdaBound{A}\AgdaSymbol{)}\<%
\\
%
\>[2]\AgdaSymbol{→}\AgdaSpace{}%
\AgdaSymbol{(}\AgdaBound{M₂}\AgdaSpace{}%
\AgdaSymbol{:}\AgdaSpace{}%
\AgdaBound{Γ}\AgdaSpace{}%
\AgdaOperator{\AgdaDatatype{⊢}}\AgdaSpace{}%
\AgdaBound{B}\AgdaSymbol{)}\<%
\\
%
\>[2]\AgdaSymbol{→}\AgdaSpace{}%
\AgdaSymbol{(}\AgdaBound{x}\AgdaSpace{}%
\AgdaSymbol{:}\AgdaSpace{}%
\AgdaDatatype{Value}\AgdaSpace{}%
\AgdaBound{M₂}\AgdaSpace{}%
\AgdaSymbol{)}\<%
\\
%
\>[2]\AgdaSymbol{→}%
\>[2725I]\AgdaFunction{substb}\AgdaSpace{}%
\AgdaSymbol{(}\AgdaFunction{fb}\AgdaSpace{}%
\AgdaSymbol{(}\AgdaInductiveConstructor{ƛb}\AgdaSpace{}%
\AgdaSymbol{(}\AgdaFunction{embedTerm}\AgdaSpace{}%
\AgdaBound{M₁}\AgdaSymbol{)))}\AgdaSpace{}%
\AgdaSymbol{(}\AgdaFunction{renameb}\AgdaSpace{}%
\AgdaInductiveConstructor{Sb}\AgdaSpace{}%
\AgdaSymbol{(}\AgdaFunction{embedTerm}\AgdaSpace{}%
\AgdaBound{M₂}\AgdaSymbol{))}\<%
\\
\>[.][@{}l@{}]\<[2725I]%
\>[4]\AgdaOperator{\AgdaDatatype{≡}}\AgdaSpace{}%
\AgdaInductiveConstructor{box}\AgdaSpace{}%
\AgdaSymbol{(}\AgdaFunction{embedValue}\AgdaSpace{}%
\AgdaBound{M₂}\AgdaSpace{}%
\AgdaBound{x}\AgdaSymbol{)}\<%
\\
\>[0]\AgdaFunction{subst-rename-fb-comp}\AgdaSpace{}%
\AgdaBound{M₁}\AgdaSpace{}%
\AgdaBound{M₂}\AgdaSpace{}%
\AgdaBound{x}\AgdaSpace{}%
\AgdaSymbol{=}\<%
\\
\>[0][@{}l@{\AgdaIndent{0}}]%
\>[2]\AgdaOperator{\AgdaFunction{beginEq}}\<%
\\
\>[2][@{}l@{\AgdaIndent{0}}]%
\>[4]\AgdaFunction{substb}\AgdaSpace{}%
\AgdaSymbol{(}\AgdaFunction{fb}\AgdaSpace{}%
\AgdaSymbol{(}\AgdaInductiveConstructor{ƛb}\AgdaSpace{}%
\AgdaSymbol{(}\AgdaFunction{embedTerm}\AgdaSpace{}%
\AgdaBound{M₁}\AgdaSymbol{)))}\AgdaSpace{}%
\AgdaSymbol{(}\AgdaFunction{renameb}\AgdaSpace{}%
\AgdaInductiveConstructor{Sb}\AgdaSpace{}%
\AgdaSymbol{(}\AgdaFunction{embedTerm}\AgdaSpace{}%
\AgdaBound{M₂}\AgdaSymbol{))}\<%
\\
%
\>[2]\AgdaFunction{≡⟨}\AgdaSpace{}%
\AgdaFunction{subst-renameb-comp}\AgdaSpace{}%
\AgdaInductiveConstructor{Sb}\AgdaSpace{}%
\AgdaSymbol{(}\AgdaFunction{fb}\AgdaSpace{}%
\AgdaSymbol{(}\AgdaInductiveConstructor{ƛb}\AgdaSpace{}%
\AgdaSymbol{(}\AgdaFunction{embedTerm}\AgdaSpace{}%
\AgdaBound{M₁}\AgdaSymbol{)))}\AgdaSpace{}%
\AgdaSymbol{(}\AgdaFunction{embedTerm}\AgdaSpace{}%
\AgdaBound{M₂}\AgdaSymbol{)}\AgdaSpace{}%
\AgdaFunction{⟩}\<%
\\
\>[2][@{}l@{\AgdaIndent{0}}]%
\>[4]\AgdaFunction{substb}\AgdaSpace{}%
\AgdaSymbol{(λ}\AgdaSpace{}%
\AgdaBound{x₁}\AgdaSpace{}%
\AgdaSymbol{→}\AgdaSpace{}%
\AgdaInductiveConstructor{ε}\AgdaSpace{}%
\AgdaBound{x₁}\AgdaSymbol{)}\AgdaSpace{}%
\AgdaSymbol{(}\AgdaFunction{embedTerm}\AgdaSpace{}%
\AgdaBound{M₂}\AgdaSymbol{)}\<%
\\
%
\>[2]\AgdaFunction{≡⟨}\AgdaSpace{}%
\AgdaFunction{subst-with-id}\AgdaSpace{}%
\AgdaSymbol{(}\AgdaFunction{embedTerm}\AgdaSpace{}%
\AgdaBound{M₂}\AgdaSymbol{)}\AgdaSpace{}%
\AgdaFunction{⟩}\<%
\\
\>[2][@{}l@{\AgdaIndent{0}}]%
\>[4]\AgdaFunction{embedTerm}\AgdaSpace{}%
\AgdaBound{M₂}\<%
\\
%
\>[2]\AgdaFunction{≡⟨}\AgdaSpace{}%
\AgdaFunction{sym}\AgdaSpace{}%
\AgdaSymbol{(}\AgdaFunction{box-embedValue}\AgdaSpace{}%
\AgdaBound{x}\AgdaSymbol{)}\AgdaSpace{}%
\AgdaFunction{⟩}\<%
\\
\>[2][@{}l@{\AgdaIndent{0}}]%
\>[4]\AgdaInductiveConstructor{box}\AgdaSpace{}%
\AgdaSymbol{(}\AgdaFunction{embedValue}\AgdaSpace{}%
\AgdaBound{M₂}\AgdaSpace{}%
\AgdaBound{x}\AgdaSymbol{)}\<%
\\
%
\>[2]\AgdaOperator{\AgdaFunction{∎eq}}\<%
\\
%
\\[\AgdaEmptyExtraSkip]%
\>[0]\AgdaComment{--\ Preservation\ of\ 2\ step\ reduction\ of\ Gödel's\ translation}\<%
\\
\>[0]\AgdaFunction{pres-red}\AgdaSpace{}%
\AgdaSymbol{:}\AgdaSpace{}%
\AgdaSymbol{∀}\AgdaSpace{}%
\AgdaSymbol{\{}\AgdaBound{Γ}\AgdaSpace{}%
\AgdaBound{A}\AgdaSymbol{\}}\<%
\\
\>[0][@{}l@{\AgdaIndent{0}}]%
\>[2]\AgdaSymbol{→}\AgdaSpace{}%
\AgdaSymbol{(}\AgdaBound{M}\AgdaSpace{}%
\AgdaSymbol{:}\AgdaSpace{}%
\AgdaBound{Γ}\AgdaSpace{}%
\AgdaOperator{\AgdaDatatype{⊢}}\AgdaSpace{}%
\AgdaBound{A}\AgdaSymbol{)}\<%
\\
%
\>[2]\AgdaSymbol{→}\AgdaSpace{}%
\AgdaSymbol{(}\AgdaBound{N}\AgdaSpace{}%
\AgdaSymbol{:}\AgdaSpace{}%
\AgdaBound{Γ}\AgdaSpace{}%
\AgdaOperator{\AgdaDatatype{⊢}}\AgdaSpace{}%
\AgdaBound{A}\AgdaSymbol{)}\<%
\\
%
\>[2]\AgdaSymbol{→}\AgdaSpace{}%
\AgdaSymbol{(}\AgdaBound{M}\AgdaSpace{}%
\AgdaOperator{\AgdaDatatype{↝βᵥ}}\AgdaSpace{}%
\AgdaBound{N}\AgdaSymbol{)}\<%
\\
%
\>[2]\AgdaSymbol{→}\AgdaSpace{}%
\AgdaSymbol{(}\AgdaFunction{embedTerm}\AgdaSpace{}%
\AgdaBound{M}\AgdaSpace{}%
\AgdaOperator{\AgdaDatatype{↝βb²}}\AgdaSpace{}%
\AgdaFunction{embedTerm}\AgdaSpace{}%
\AgdaBound{N}\AgdaSymbol{)}\<%
\\
\>[0]\AgdaFunction{pres-red}\AgdaSpace{}%
\AgdaSymbol{(}\AgdaInductiveConstructor{ƛ}\AgdaSpace{}%
\AgdaBound{M₁}\AgdaSpace{}%
\AgdaOperator{\AgdaInductiveConstructor{·}}\AgdaSpace{}%
\AgdaBound{M₂}\AgdaSymbol{)}\AgdaSpace{}%
\AgdaBound{N}\AgdaSpace{}%
\AgdaSymbol{(}\AgdaInductiveConstructor{βᵥ}\AgdaSpace{}%
\AgdaBound{x}\AgdaSymbol{)}\AgdaSpace{}%
\AgdaSymbol{=}\<%
\\
\>[0][@{}l@{\AgdaIndent{0}}]%
\>[2]\AgdaInductiveConstructor{βb²}\<%
\\
%
\>[2]\AgdaSymbol{(}\AgdaFunction{raise}\AgdaSpace{}%
\AgdaSymbol{(}\AgdaFunction{embedTerm}\AgdaSpace{}%
\AgdaBound{M₂}\AgdaSymbol{)}\AgdaSpace{}%
\AgdaOperator{\AgdaInductiveConstructor{·b}}\AgdaSpace{}%
\AgdaInductiveConstructor{box}\AgdaSpace{}%
\AgdaSymbol{(}\AgdaInductiveConstructor{ƛb}\AgdaSpace{}%
\AgdaSymbol{(}\AgdaFunction{embedTerm}\AgdaSpace{}%
\AgdaBound{M₁}\AgdaSymbol{)))}\<%
\\
%
\>[2]\AgdaSymbol{(}\AgdaInductiveConstructor{ƛb}\AgdaSpace{}%
\AgdaSymbol{(}\AgdaFunction{embedTerm}\AgdaSpace{}%
\AgdaBound{M₁}\AgdaSymbol{)}\AgdaSpace{}%
\AgdaOperator{\AgdaInductiveConstructor{·b}}\AgdaSpace{}%
\AgdaInductiveConstructor{box}\AgdaSpace{}%
\AgdaSymbol{(}\AgdaFunction{embedValue}\AgdaSpace{}%
\AgdaBound{M₂}\AgdaSpace{}%
\AgdaBound{x}\AgdaSymbol{))}\<%
\\
%
\>[2]\AgdaSymbol{(}\AgdaFunction{embedTerm}\AgdaSpace{}%
\AgdaSymbol{(}\AgdaBound{M₁}\AgdaSpace{}%
\AgdaOperator{\AgdaFunction{[}}\AgdaSpace{}%
\AgdaBound{M₂}\AgdaSpace{}%
\AgdaOperator{\AgdaFunction{]}}\AgdaSymbol{))}\<%
\\
%
\>[2]\AgdaSymbol{(}\AgdaFunction{Eq.subst}\<%
\\
\>[2][@{}l@{\AgdaIndent{0}}]%
\>[4]\AgdaSymbol{(λ}\AgdaSpace{}%
\AgdaBound{x}\AgdaSpace{}%
\AgdaSymbol{→}\AgdaSpace{}%
\AgdaInductiveConstructor{ƛb}\AgdaSpace{}%
\AgdaSymbol{(}\AgdaInductiveConstructor{ε}\AgdaSpace{}%
\AgdaInductiveConstructor{Zb}\AgdaSpace{}%
\AgdaOperator{\AgdaInductiveConstructor{·b}}\AgdaSpace{}%
\AgdaFunction{renameb}\AgdaSpace{}%
\AgdaInductiveConstructor{Sb}\AgdaSpace{}%
\AgdaSymbol{(}\AgdaFunction{embedTerm}\AgdaSpace{}%
\AgdaBound{M₂}\AgdaSymbol{))}\AgdaSpace{}%
\AgdaOperator{\AgdaInductiveConstructor{·b}}\AgdaSpace{}%
\AgdaInductiveConstructor{box}\AgdaSpace{}%
\AgdaSymbol{(}\AgdaInductiveConstructor{ƛb}\AgdaSpace{}%
\AgdaSymbol{(}\AgdaFunction{embedTerm}\AgdaSpace{}%
\AgdaBound{M₁}\AgdaSymbol{))}\AgdaSpace{}%
\AgdaOperator{\AgdaDatatype{↝βb}}\AgdaSpace{}%
\AgdaBound{x}\AgdaSymbol{)}\<%
\\
%
\>[4]\AgdaSymbol{(}\AgdaFunction{cong₂}\AgdaSpace{}%
\AgdaOperator{\AgdaInductiveConstructor{\AgdaUnderscore{}·b\AgdaUnderscore{}}}\AgdaSpace{}%
\AgdaInductiveConstructor{refl}\AgdaSpace{}%
\AgdaSymbol{(}\AgdaFunction{subst-rename-fb-comp}\AgdaSpace{}%
\AgdaBound{M₁}\AgdaSpace{}%
\AgdaBound{M₂}\AgdaSpace{}%
\AgdaBound{x}\AgdaSymbol{))}\<%
\\
%
\>[4]\AgdaInductiveConstructor{βb}\AgdaSymbol{)}\<%
\\
%
\>[2]\AgdaSymbol{(}\AgdaFunction{Eq.subst}\AgdaSpace{}%
\AgdaSymbol{(λ}\AgdaSpace{}%
\AgdaBound{t}\AgdaSpace{}%
\AgdaSymbol{→}\AgdaSpace{}%
\AgdaInductiveConstructor{ƛb}\AgdaSpace{}%
\AgdaSymbol{(}\AgdaFunction{embedTerm}\AgdaSpace{}%
\AgdaBound{M₁}\AgdaSymbol{)}\AgdaSpace{}%
\AgdaOperator{\AgdaInductiveConstructor{·b}}\AgdaSpace{}%
\AgdaInductiveConstructor{box}\AgdaSpace{}%
\AgdaSymbol{(}\AgdaFunction{embedValue}\AgdaSpace{}%
\AgdaBound{M₂}\AgdaSpace{}%
\AgdaBound{x}\AgdaSymbol{)}\AgdaSpace{}%
\AgdaOperator{\AgdaDatatype{↝βb}}\AgdaSpace{}%
\AgdaBound{t}\AgdaSymbol{)}\AgdaSpace{}%
\AgdaSymbol{(}\AgdaFunction{pres-subst}\AgdaSpace{}%
\AgdaBound{M₁}\AgdaSpace{}%
\AgdaBound{M₂}\AgdaSpace{}%
\AgdaBound{x}\AgdaSymbol{)}\AgdaSpace{}%
\AgdaInductiveConstructor{βb}\AgdaSymbol{)}\<%
\\
\>[0]\AgdaFunction{pres-red}\AgdaSpace{}%
\AgdaBound{M}\AgdaSpace{}%
\AgdaBound{N}\AgdaSpace{}%
\AgdaSymbol{(}\AgdaInductiveConstructor{μ}\AgdaSpace{}%
\AgdaSymbol{\{}\AgdaArgument{M}\AgdaSpace{}%
\AgdaSymbol{=}\AgdaSpace{}%
\AgdaBound{M₁}\AgdaSymbol{\}}\AgdaSpace{}%
\AgdaSymbol{\{}\AgdaArgument{M'}\AgdaSpace{}%
\AgdaSymbol{=}\AgdaSpace{}%
\AgdaBound{M'}\AgdaSymbol{\}}\AgdaSpace{}%
\AgdaBound{red}\AgdaSymbol{)}\AgdaSpace{}%
\AgdaSymbol{=}\AgdaSpace{}%
\AgdaInductiveConstructor{ν}\AgdaSpace{}%
\AgdaSymbol{(}\AgdaFunction{pres-red}\AgdaSpace{}%
\AgdaBound{M₁}\AgdaSpace{}%
\AgdaBound{M'}\AgdaSpace{}%
\AgdaBound{red}\AgdaSymbol{)}\<%
\\
\>[0]\AgdaFunction{pres-red}\AgdaSpace{}%
\AgdaBound{M}\AgdaSpace{}%
\AgdaBound{N}\AgdaSpace{}%
\AgdaSymbol{(}\AgdaInductiveConstructor{ν<}\AgdaSpace{}%
\AgdaSymbol{\{}\AgdaArgument{M}\AgdaSpace{}%
\AgdaSymbol{=}\AgdaSpace{}%
\AgdaBound{M₁}\AgdaSymbol{\}}\AgdaSpace{}%
\AgdaSymbol{\{}\AgdaArgument{N}\AgdaSpace{}%
\AgdaSymbol{=}\AgdaSpace{}%
\AgdaBound{N₁}\AgdaSymbol{\}}\AgdaSpace{}%
\AgdaSymbol{\{}\AgdaArgument{N'}\AgdaSpace{}%
\AgdaSymbol{=}\AgdaSpace{}%
\AgdaBound{N'}\AgdaSymbol{\}}\AgdaSpace{}%
\AgdaBound{x}\AgdaSpace{}%
\AgdaBound{red}\AgdaSymbol{)}\AgdaSpace{}%
\AgdaSymbol{=}\AgdaSpace{}%
\AgdaInductiveConstructor{μ}\AgdaSpace{}%
\AgdaSymbol{(}\AgdaInductiveConstructor{ξ}\AgdaSpace{}%
\AgdaSymbol{(}\AgdaInductiveConstructor{ν}\AgdaSpace{}%
\AgdaSymbol{(}\AgdaFunction{pres-rename-red-2}\AgdaSpace{}%
\AgdaInductiveConstructor{Sb}\AgdaSpace{}%
\AgdaSymbol{(}\AgdaFunction{embedTerm}\AgdaSpace{}%
\AgdaBound{N₁}\AgdaSymbol{)}\AgdaSpace{}%
\AgdaSymbol{(}\AgdaFunction{embedTerm}\AgdaSpace{}%
\AgdaBound{N'}\AgdaSymbol{)}\AgdaSpace{}%
\AgdaSymbol{(}\AgdaFunction{pres-red}\AgdaSpace{}%
\AgdaBound{N₁}\AgdaSpace{}%
\AgdaBound{N'}\AgdaSpace{}%
\AgdaBound{red}\AgdaSymbol{))))}\<%
\\
\>[0]\AgdaFunction{pres-red}\AgdaSpace{}%
\AgdaBound{M}\AgdaSpace{}%
\AgdaBound{N}\AgdaSpace{}%
\AgdaSymbol{(}\AgdaInductiveConstructor{ν}\AgdaSpace{}%
\AgdaSymbol{\{}\AgdaArgument{M}\AgdaSpace{}%
\AgdaSymbol{=}\AgdaSpace{}%
\AgdaBound{M₁}\AgdaSymbol{\}}\AgdaSpace{}%
\AgdaSymbol{\{}\AgdaArgument{N}\AgdaSpace{}%
\AgdaSymbol{=}\AgdaSpace{}%
\AgdaBound{N₁}\AgdaSymbol{\}}\AgdaSpace{}%
\AgdaSymbol{\{}\AgdaArgument{N'}\AgdaSpace{}%
\AgdaSymbol{=}\AgdaSpace{}%
\AgdaBound{N'}\AgdaSymbol{\}}\AgdaSpace{}%
\AgdaBound{red}\AgdaSymbol{)}\AgdaSpace{}%
\AgdaSymbol{=}\AgdaSpace{}%
\AgdaInductiveConstructor{μ}\AgdaSpace{}%
\AgdaSymbol{(}\AgdaInductiveConstructor{ξ}\AgdaSpace{}%
\AgdaSymbol{(}\AgdaInductiveConstructor{ν}\AgdaSpace{}%
\AgdaSymbol{(}\AgdaFunction{pres-rename-red-2}\AgdaSpace{}%
\AgdaInductiveConstructor{Sb}\AgdaSpace{}%
\AgdaSymbol{(}\AgdaFunction{embedTerm}\AgdaSpace{}%
\AgdaBound{N₁}\AgdaSymbol{)}\AgdaSpace{}%
\AgdaSymbol{(}\AgdaFunction{embedTerm}\AgdaSpace{}%
\AgdaBound{N'}\AgdaSymbol{)}\AgdaSpace{}%
\AgdaSymbol{(}\AgdaFunction{pres-red}\AgdaSpace{}%
\AgdaBound{N₁}\AgdaSpace{}%
\AgdaBound{N'}\AgdaSpace{}%
\AgdaBound{red}\AgdaSymbol{))))}\<%
\\
\>[0]\AgdaFunction{pres-red}\AgdaSpace{}%
\AgdaBound{M}\AgdaSpace{}%
\AgdaBound{N}\AgdaSpace{}%
\AgdaSymbol{(}\AgdaInductiveConstructor{ξ}\AgdaSpace{}%
\AgdaSymbol{\{}\AgdaArgument{M}\AgdaSpace{}%
\AgdaSymbol{=}\AgdaSpace{}%
\AgdaBound{M₁}\AgdaSymbol{\}}\AgdaSpace{}%
\AgdaSymbol{\{}\AgdaArgument{M'}\AgdaSpace{}%
\AgdaSymbol{=}\AgdaSpace{}%
\AgdaBound{M'}\AgdaSymbol{\}}\AgdaSpace{}%
\AgdaBound{red}\AgdaSymbol{)}\AgdaSpace{}%
\AgdaSymbol{=}\AgdaSpace{}%
\AgdaInductiveConstructor{ζ}\AgdaSpace{}%
\AgdaSymbol{(}\AgdaInductiveConstructor{ξ}\AgdaSpace{}%
\AgdaSymbol{(}\AgdaFunction{pres-red}\AgdaSpace{}%
\AgdaBound{M₁}\AgdaSpace{}%
\AgdaBound{M'}\AgdaSpace{}%
\AgdaBound{red}\AgdaSymbol{))}\<%
\\
%
\\[\AgdaEmptyExtraSkip]%
\>[0]\AgdaComment{--\ Preservation\ of\ the\ 2\ step\ reduction\ rule\ under\ rename}\<%
\\
\>[0]\AgdaFunction{pres-rename-red-↝βb²}\AgdaSpace{}%
\AgdaSymbol{:}\AgdaSpace{}%
\AgdaSymbol{∀}\AgdaSpace{}%
\AgdaSymbol{\{}\AgdaBound{Γ}\AgdaSpace{}%
\AgdaBound{A}\AgdaSpace{}%
\AgdaBound{Δ}\AgdaSymbol{\}}\<%
\\
\>[0][@{}l@{\AgdaIndent{0}}]%
\>[2]\AgdaSymbol{→}\AgdaSpace{}%
\AgdaSymbol{(}\AgdaBound{fb'}\AgdaSpace{}%
\AgdaSymbol{:}\AgdaSpace{}%
\AgdaSymbol{∀}\AgdaSpace{}%
\AgdaSymbol{\{}\AgdaBound{Ab}\AgdaSymbol{\}}\AgdaSpace{}%
\AgdaSymbol{→}\AgdaSpace{}%
\AgdaBound{Ab}\AgdaSpace{}%
\AgdaOperator{\AgdaDatatype{∈b}}\AgdaSpace{}%
\AgdaBound{Γ}\AgdaSpace{}%
\AgdaSymbol{→}\AgdaSpace{}%
\AgdaBound{Ab}\AgdaSpace{}%
\AgdaOperator{\AgdaDatatype{∈b}}\AgdaSpace{}%
\AgdaBound{Δ}\AgdaSymbol{)}\<%
\\
%
\>[2]\AgdaSymbol{→}\AgdaSpace{}%
\AgdaSymbol{(}\AgdaBound{M}\AgdaSpace{}%
\AgdaSymbol{:}\AgdaSpace{}%
\AgdaBound{Γ}\AgdaSpace{}%
\AgdaOperator{\AgdaDatatype{⊢b}}\AgdaSpace{}%
\AgdaBound{A}\AgdaSymbol{)}\<%
\\
%
\>[2]\AgdaSymbol{→}\AgdaSpace{}%
\AgdaSymbol{(}\AgdaBound{N}\AgdaSpace{}%
\AgdaSymbol{:}\AgdaSpace{}%
\AgdaBound{Γ}\AgdaSpace{}%
\AgdaOperator{\AgdaDatatype{⊢b}}\AgdaSpace{}%
\AgdaBound{A}\AgdaSymbol{)}\<%
\\
%
\>[2]\AgdaSymbol{→}\AgdaSpace{}%
\AgdaSymbol{(}\AgdaBound{M}\AgdaSpace{}%
\AgdaOperator{\AgdaDatatype{↝βb₂}}\AgdaSpace{}%
\AgdaBound{N}\AgdaSymbol{)}\<%
\\
%
\>[2]\AgdaSymbol{→}\AgdaSpace{}%
\AgdaFunction{renameb}\AgdaSpace{}%
\AgdaBound{fb'}\AgdaSpace{}%
\AgdaBound{M}\AgdaSpace{}%
\AgdaOperator{\AgdaDatatype{↝βb₂}}\AgdaSpace{}%
\AgdaFunction{renameb}\AgdaSpace{}%
\AgdaBound{fb'}\AgdaSpace{}%
\AgdaBound{N}\<%
\\
\>[0]\AgdaFunction{pres-rename-red-↝βb²}\AgdaSpace{}%
\AgdaBound{fb'}\AgdaSpace{}%
\AgdaBound{M}\AgdaSpace{}%
\AgdaBound{N}\AgdaSpace{}%
\AgdaSymbol{(}\AgdaInductiveConstructor{βb₂}\AgdaSpace{}%
\AgdaSymbol{\{}\AgdaArgument{M}\AgdaSpace{}%
\AgdaSymbol{=}\AgdaSpace{}%
\AgdaBound{M'}\AgdaSymbol{\}}\AgdaSpace{}%
\AgdaSymbol{\{}\AgdaArgument{W}\AgdaSpace{}%
\AgdaSymbol{=}\AgdaSpace{}%
\AgdaBound{W'}\AgdaSymbol{\})}\AgdaSpace{}%
\AgdaSymbol{=}\<%
\\
\>[0][@{}l@{\AgdaIndent{0}}]%
\>[2]\AgdaFunction{Eq.subst}\<%
\\
\>[2][@{}l@{\AgdaIndent{0}}]%
\>[4]\AgdaSymbol{(λ}\AgdaSpace{}%
\AgdaBound{x}\AgdaSpace{}%
\AgdaSymbol{→}\AgdaSpace{}%
\AgdaBound{x}\AgdaSpace{}%
\AgdaOperator{\AgdaDatatype{↝βb₂}}\AgdaSpace{}%
\AgdaFunction{renameb}\AgdaSpace{}%
\AgdaBound{fb'}\<%
\\
%
\>[4]\AgdaSymbol{(}\AgdaBound{M'}\AgdaSpace{}%
\AgdaOperator{\AgdaFunction{[}}\AgdaSpace{}%
\AgdaBound{W'}\AgdaSpace{}%
\AgdaOperator{\AgdaFunction{]b}}\AgdaSymbol{))}\<%
\\
%
\>[4]\AgdaSymbol{(}\AgdaFunction{cong₂}\AgdaSpace{}%
\AgdaOperator{\AgdaInductiveConstructor{\AgdaUnderscore{}·b\AgdaUnderscore{}}}\AgdaSpace{}%
\AgdaSymbol{(}\AgdaFunction{cong}\AgdaSpace{}%
\AgdaInductiveConstructor{ƛb}\AgdaSpace{}%
\AgdaSymbol{(}\AgdaFunction{cong₂}\AgdaSpace{}%
\AgdaOperator{\AgdaInductiveConstructor{\AgdaUnderscore{}·b\AgdaUnderscore{}}}\AgdaSpace{}%
\AgdaInductiveConstructor{refl}\AgdaSpace{}%
\AgdaSymbol{(}\AgdaFunction{cong}\AgdaSpace{}%
\AgdaInductiveConstructor{box}\AgdaSpace{}%
\AgdaSymbol{(}\AgdaFunction{sym}\AgdaSpace{}%
\AgdaSymbol{(}\AgdaFunction{rename-extb-Sb-commute}\AgdaSpace{}%
\AgdaBound{fb'}\AgdaSpace{}%
\AgdaBound{W'}\AgdaSymbol{)))))}\AgdaSpace{}%
\AgdaInductiveConstructor{refl}\AgdaSymbol{)}\<%
\\
%
\>[4]\AgdaSymbol{(}\AgdaFunction{Eq.subst}\<%
\\
\>[4][@{}l@{\AgdaIndent{0}}]%
\>[6]\AgdaSymbol{(λ}\AgdaSpace{}%
\AgdaBound{x}\AgdaSpace{}%
\AgdaSymbol{→}\AgdaSpace{}%
\AgdaSymbol{(}\AgdaFunction{raise}%
\>[21]\AgdaSymbol{(}\AgdaInductiveConstructor{box}\AgdaSpace{}%
\AgdaSymbol{(}\AgdaFunction{renameb}\AgdaSpace{}%
\AgdaBound{fb'}\AgdaSpace{}%
\AgdaBound{W'}\AgdaSymbol{)))}\AgdaSpace{}%
\AgdaOperator{\AgdaInductiveConstructor{·b}}\AgdaSpace{}%
\AgdaInductiveConstructor{box}\AgdaSpace{}%
\AgdaSymbol{(}\AgdaInductiveConstructor{ƛb}\AgdaSpace{}%
\AgdaSymbol{(}\AgdaFunction{renameb}\AgdaSpace{}%
\AgdaSymbol{(}\AgdaFunction{extb}\AgdaSpace{}%
\AgdaBound{fb'}\AgdaSymbol{)}\AgdaSpace{}%
\AgdaBound{M'}\AgdaSymbol{))}\AgdaSpace{}%
\AgdaOperator{\AgdaDatatype{↝βb₂}}\AgdaSpace{}%
\AgdaBound{x}\AgdaSymbol{)}\<%
\\
%
\>[6]\AgdaSymbol{(}\AgdaFunction{pres-subst-func-rename}\AgdaSpace{}%
\AgdaBound{fb'}\AgdaSpace{}%
\AgdaBound{M'}\AgdaSpace{}%
\AgdaBound{W'}\AgdaSymbol{)}\<%
\\
%
\>[6]\AgdaInductiveConstructor{βb₂}\<%
\\
%
\>[4]\AgdaSymbol{)}\<%
\\
\>[0]\AgdaFunction{pres-rename-red-↝βb²}\AgdaSpace{}%
\AgdaBound{fb'}\AgdaSpace{}%
\AgdaBound{M}\AgdaSpace{}%
\AgdaBound{N}\AgdaSpace{}%
\AgdaSymbol{(}\AgdaInductiveConstructor{μ}\AgdaSpace{}%
\AgdaSymbol{\{}\AgdaArgument{M}\AgdaSpace{}%
\AgdaSymbol{=}\AgdaSpace{}%
\AgdaBound{M₁}\AgdaSymbol{\}}\AgdaSpace{}%
\AgdaSymbol{\{}\AgdaArgument{M'}\AgdaSpace{}%
\AgdaSymbol{=}\AgdaSpace{}%
\AgdaBound{M'}\AgdaSymbol{\}}\AgdaSpace{}%
\AgdaBound{red}\AgdaSymbol{)}\AgdaSpace{}%
\AgdaSymbol{=}\AgdaSpace{}%
\AgdaInductiveConstructor{μ}\AgdaSpace{}%
\AgdaSymbol{(}\AgdaFunction{pres-rename-red-↝βb²}\AgdaSpace{}%
\AgdaBound{fb'}\AgdaSpace{}%
\AgdaBound{M₁}\AgdaSpace{}%
\AgdaBound{M'}\AgdaSpace{}%
\AgdaBound{red}\AgdaSymbol{)}\<%
\\
\>[0]\AgdaFunction{pres-rename-red-↝βb²}\AgdaSpace{}%
\AgdaBound{fb'}\AgdaSpace{}%
\AgdaBound{M}\AgdaSpace{}%
\AgdaBound{N}\AgdaSpace{}%
\AgdaSymbol{(}\AgdaInductiveConstructor{ν}\AgdaSpace{}%
\AgdaSymbol{\{}\AgdaArgument{N}\AgdaSpace{}%
\AgdaSymbol{=}\AgdaSpace{}%
\AgdaBound{N₁}\AgdaSymbol{\}}\AgdaSpace{}%
\AgdaSymbol{\{}\AgdaArgument{N'}\AgdaSpace{}%
\AgdaSymbol{=}\AgdaSpace{}%
\AgdaBound{N'}\AgdaSymbol{\}}\AgdaSpace{}%
\AgdaBound{red}\AgdaSymbol{)}\AgdaSpace{}%
\AgdaSymbol{=}\AgdaSpace{}%
\AgdaInductiveConstructor{ν}\AgdaSpace{}%
\AgdaSymbol{(}\AgdaFunction{pres-rename-red-↝βb²}\AgdaSpace{}%
\AgdaBound{fb'}\AgdaSpace{}%
\AgdaBound{N₁}\AgdaSpace{}%
\AgdaBound{N'}\AgdaSpace{}%
\AgdaBound{red}\AgdaSymbol{)}\<%
\\
\>[0]\AgdaFunction{pres-rename-red-↝βb²}\AgdaSpace{}%
\AgdaBound{fb'}\AgdaSpace{}%
\AgdaBound{M}\AgdaSpace{}%
\AgdaBound{N}\AgdaSpace{}%
\AgdaSymbol{(}\AgdaInductiveConstructor{ξ}\AgdaSpace{}%
\AgdaSymbol{\{}\AgdaArgument{M}\AgdaSpace{}%
\AgdaSymbol{=}\AgdaSpace{}%
\AgdaBound{M₁}\AgdaSymbol{\}}\AgdaSpace{}%
\AgdaSymbol{\{}\AgdaArgument{M'}\AgdaSpace{}%
\AgdaSymbol{=}\AgdaSpace{}%
\AgdaBound{M'}\AgdaSymbol{\}}\AgdaSpace{}%
\AgdaBound{red}\AgdaSymbol{)}\AgdaSpace{}%
\AgdaSymbol{=}\AgdaSpace{}%
\AgdaInductiveConstructor{ξ}\AgdaSpace{}%
\AgdaSymbol{(}\AgdaFunction{pres-rename-red-↝βb²}\AgdaSpace{}%
\AgdaSymbol{(}\AgdaFunction{extb}\AgdaSpace{}%
\AgdaBound{fb'}\AgdaSymbol{)}\AgdaSpace{}%
\AgdaBound{M₁}\AgdaSpace{}%
\AgdaBound{M'}\AgdaSpace{}%
\AgdaBound{red}\AgdaSymbol{)}\<%
\\
\>[0]\AgdaFunction{pres-rename-red-↝βb²}\AgdaSpace{}%
\AgdaBound{fb'}\AgdaSpace{}%
\AgdaBound{M}\AgdaSpace{}%
\AgdaBound{N}\AgdaSpace{}%
\AgdaSymbol{(}\AgdaInductiveConstructor{ζ}\AgdaSpace{}%
\AgdaSymbol{\{}\AgdaArgument{M}\AgdaSpace{}%
\AgdaSymbol{=}\AgdaSpace{}%
\AgdaBound{M₁}\AgdaSymbol{\}}\AgdaSpace{}%
\AgdaSymbol{\{}\AgdaArgument{M'}\AgdaSpace{}%
\AgdaSymbol{=}\AgdaSpace{}%
\AgdaBound{M'}\AgdaSymbol{\}}\AgdaSpace{}%
\AgdaBound{red}\AgdaSymbol{)}\AgdaSpace{}%
\AgdaSymbol{=}\AgdaSpace{}%
\AgdaInductiveConstructor{ζ}\AgdaSpace{}%
\AgdaSymbol{(}\AgdaFunction{pres-rename-red-↝βb²}\AgdaSpace{}%
\AgdaBound{fb'}\AgdaSpace{}%
\AgdaBound{M₁}\AgdaSpace{}%
\AgdaBound{M'}\AgdaSpace{}%
\AgdaBound{red}\AgdaSymbol{)}\<%
\\
%
\\[\AgdaEmptyExtraSkip]%
\>[0]\AgdaComment{--\ Preservation\ of\ reduction\ by\ Gödel's\ translation}\<%
\\
\>[0]\AgdaComment{--\ with\ the\ two\ step\ reduction\ rule}\<%
\\
\>[0]\AgdaFunction{pres-red↝βb²}\AgdaSpace{}%
\AgdaSymbol{:}\AgdaSpace{}%
\AgdaSymbol{∀}\AgdaSpace{}%
\AgdaSymbol{\{}\AgdaBound{Γ}\AgdaSpace{}%
\AgdaBound{A}\AgdaSymbol{\}}\<%
\\
\>[0][@{}l@{\AgdaIndent{0}}]%
\>[2]\AgdaSymbol{→}\AgdaSpace{}%
\AgdaSymbol{(}\AgdaBound{M}\AgdaSpace{}%
\AgdaSymbol{:}\AgdaSpace{}%
\AgdaBound{Γ}\AgdaSpace{}%
\AgdaOperator{\AgdaDatatype{⊢}}\AgdaSpace{}%
\AgdaBound{A}\AgdaSymbol{)}\<%
\\
%
\>[2]\AgdaSymbol{→}\AgdaSpace{}%
\AgdaSymbol{(}\AgdaBound{N}\AgdaSpace{}%
\AgdaSymbol{:}\AgdaSpace{}%
\AgdaBound{Γ}\AgdaSpace{}%
\AgdaOperator{\AgdaDatatype{⊢}}\AgdaSpace{}%
\AgdaBound{A}\AgdaSymbol{)}\<%
\\
%
\>[2]\AgdaSymbol{→}\AgdaSpace{}%
\AgdaSymbol{(}\AgdaBound{M}\AgdaSpace{}%
\AgdaOperator{\AgdaDatatype{↝βᵥ}}\AgdaSpace{}%
\AgdaBound{N}\AgdaSymbol{)}\<%
\\
%
\>[2]\AgdaSymbol{→}\AgdaSpace{}%
\AgdaSymbol{(}\AgdaFunction{embedTerm}\AgdaSpace{}%
\AgdaBound{M}\AgdaSpace{}%
\AgdaOperator{\AgdaDatatype{↝βb₂}}\AgdaSpace{}%
\AgdaFunction{embedTerm}\AgdaSpace{}%
\AgdaBound{N}\AgdaSymbol{)}\<%
\\
\>[0]\AgdaFunction{pres-red↝βb²}\AgdaSpace{}%
\AgdaSymbol{(}\AgdaInductiveConstructor{ƛ}\AgdaSpace{}%
\AgdaBound{M₁}\AgdaSpace{}%
\AgdaOperator{\AgdaInductiveConstructor{·}}\AgdaSpace{}%
\AgdaBound{M₂}\AgdaSymbol{)}\AgdaSpace{}%
\AgdaBound{N}\AgdaSpace{}%
\AgdaSymbol{(}\AgdaInductiveConstructor{βᵥ}\AgdaSpace{}%
\AgdaBound{x}\AgdaSymbol{)}\AgdaSpace{}%
\AgdaSymbol{=}\<%
\\
\>[0][@{}l@{\AgdaIndent{0}}]%
\>[2]\AgdaFunction{Eq.subst}\<%
\\
\>[2][@{}l@{\AgdaIndent{0}}]%
\>[4]\AgdaSymbol{(λ}\AgdaSpace{}%
\AgdaBound{t}\AgdaSpace{}%
\AgdaSymbol{→}\AgdaSpace{}%
\AgdaBound{t}\AgdaSpace{}%
\AgdaOperator{\AgdaDatatype{↝βb₂}}\AgdaSpace{}%
\AgdaFunction{embedTerm}\AgdaSpace{}%
\AgdaSymbol{(}\AgdaBound{M₁}\AgdaSpace{}%
\AgdaOperator{\AgdaFunction{[}}\AgdaSpace{}%
\AgdaBound{M₂}\AgdaSpace{}%
\AgdaOperator{\AgdaFunction{]}}\AgdaSymbol{))}\<%
\\
%
\>[4]\AgdaSymbol{(}\AgdaFunction{cong₂}\AgdaSpace{}%
\AgdaOperator{\AgdaInductiveConstructor{\AgdaUnderscore{}·b\AgdaUnderscore{}}}\AgdaSpace{}%
\AgdaSymbol{(}\AgdaFunction{cong}\AgdaSpace{}%
\AgdaInductiveConstructor{ƛb}\AgdaSpace{}%
\AgdaSymbol{(}\AgdaFunction{cong₂}\AgdaSpace{}%
\AgdaOperator{\AgdaInductiveConstructor{\AgdaUnderscore{}·b\AgdaUnderscore{}}}\AgdaSpace{}%
\AgdaInductiveConstructor{refl}\AgdaSpace{}%
\AgdaSymbol{(}\AgdaFunction{cong}\AgdaSpace{}%
\AgdaSymbol{(}\AgdaFunction{renameb}\AgdaSpace{}%
\AgdaInductiveConstructor{Sb}\AgdaSymbol{)}\AgdaSpace{}%
\AgdaSymbol{(}\AgdaFunction{box-embedValue}\AgdaSpace{}%
\AgdaBound{x}\AgdaSymbol{))))}\AgdaSpace{}%
\AgdaInductiveConstructor{refl}\AgdaSymbol{)}\<%
\\
%
\>[4]\AgdaSymbol{(}\AgdaFunction{Eq.subst}\<%
\\
\>[4][@{}l@{\AgdaIndent{0}}]%
\>[6]\AgdaSymbol{(λ}\AgdaSpace{}%
\AgdaBound{t}\AgdaSpace{}%
\AgdaSymbol{→}\AgdaSpace{}%
\AgdaFunction{raise}\AgdaSpace{}%
\AgdaSymbol{(}\AgdaSpace{}%
\AgdaInductiveConstructor{box}\AgdaSpace{}%
\AgdaSymbol{(}\AgdaFunction{embedValue}\AgdaSpace{}%
\AgdaBound{M₂}\AgdaSpace{}%
\AgdaBound{x}\AgdaSymbol{)}\AgdaSpace{}%
\AgdaSymbol{)}\AgdaSpace{}%
\AgdaOperator{\AgdaInductiveConstructor{·b}}\AgdaSpace{}%
\AgdaInductiveConstructor{box}\AgdaSpace{}%
\AgdaSymbol{(}\AgdaInductiveConstructor{ƛb}\AgdaSpace{}%
\AgdaSymbol{(}\AgdaFunction{embedTerm}\AgdaSpace{}%
\AgdaBound{M₁}\AgdaSymbol{))}\AgdaSpace{}%
\AgdaOperator{\AgdaDatatype{↝βb₂}}\AgdaSpace{}%
\AgdaBound{t}\AgdaSymbol{)}\<%
\\
%
\>[6]\AgdaSymbol{(}\AgdaFunction{pres-subst}\AgdaSpace{}%
\AgdaBound{M₁}\AgdaSpace{}%
\AgdaBound{M₂}\AgdaSpace{}%
\AgdaBound{x}\AgdaSymbol{)}\<%
\\
%
\>[6]\AgdaInductiveConstructor{βb₂}\<%
\\
%
\>[4]\AgdaSymbol{)}\<%
\\
\>[0]\AgdaFunction{pres-red↝βb²}\AgdaSpace{}%
\AgdaBound{M}\AgdaSpace{}%
\AgdaBound{N}\AgdaSpace{}%
\AgdaSymbol{(}\AgdaInductiveConstructor{μ}\AgdaSpace{}%
\AgdaSymbol{\{}\AgdaArgument{M}\AgdaSpace{}%
\AgdaSymbol{=}\AgdaSpace{}%
\AgdaBound{M₁}\AgdaSymbol{\}}\AgdaSpace{}%
\AgdaSymbol{\{}\AgdaArgument{M'}\AgdaSpace{}%
\AgdaSymbol{=}\AgdaSpace{}%
\AgdaBound{M'}\AgdaSymbol{\}}\AgdaSpace{}%
\AgdaBound{red}\AgdaSymbol{)}\AgdaSpace{}%
\AgdaSymbol{=}\AgdaSpace{}%
\AgdaInductiveConstructor{ν}\AgdaSpace{}%
\AgdaSymbol{(}\AgdaFunction{pres-red↝βb²}\AgdaSpace{}%
\AgdaBound{M₁}\AgdaSpace{}%
\AgdaBound{M'}\AgdaSpace{}%
\AgdaBound{red}\AgdaSymbol{)}\<%
\\
\>[0]\AgdaFunction{pres-red↝βb²}\AgdaSpace{}%
\AgdaBound{M}\AgdaSpace{}%
\AgdaBound{N}\AgdaSpace{}%
\AgdaSymbol{(}\AgdaInductiveConstructor{ν<}\AgdaSpace{}%
\AgdaSymbol{\{}\AgdaArgument{N}\AgdaSpace{}%
\AgdaSymbol{=}\AgdaSpace{}%
\AgdaBound{N₁}\AgdaSymbol{\}}\AgdaSpace{}%
\AgdaSymbol{\{}\AgdaArgument{N'}\AgdaSpace{}%
\AgdaSymbol{=}\AgdaSpace{}%
\AgdaBound{N'}\AgdaSymbol{\}}\AgdaSpace{}%
\AgdaBound{x}\AgdaSpace{}%
\AgdaBound{red}\AgdaSymbol{)}\AgdaSpace{}%
\AgdaSymbol{=}\AgdaSpace{}%
\AgdaInductiveConstructor{μ}\AgdaSpace{}%
\AgdaSymbol{(}\AgdaInductiveConstructor{ξ}\AgdaSpace{}%
\AgdaSymbol{(}\AgdaInductiveConstructor{ν}\AgdaSpace{}%
\AgdaSymbol{(}\AgdaFunction{pres-rename-red-↝βb²}\AgdaSpace{}%
\AgdaInductiveConstructor{Sb}\AgdaSpace{}%
\AgdaSymbol{(}\AgdaFunction{embedTerm}\AgdaSpace{}%
\AgdaBound{N₁}\AgdaSymbol{)}\AgdaSpace{}%
\AgdaSymbol{(}\AgdaFunction{embedTerm}\AgdaSpace{}%
\AgdaBound{N'}\AgdaSymbol{)}\AgdaSpace{}%
\AgdaSymbol{(}\AgdaFunction{pres-red↝βb²}\AgdaSpace{}%
\AgdaBound{N₁}\AgdaSpace{}%
\AgdaBound{N'}\AgdaSpace{}%
\AgdaBound{red}\AgdaSymbol{))))}\<%
\\
\>[0]\AgdaFunction{pres-red↝βb²}\AgdaSpace{}%
\AgdaBound{M}\AgdaSpace{}%
\AgdaBound{N}\AgdaSpace{}%
\AgdaSymbol{(}\AgdaInductiveConstructor{ν}\AgdaSpace{}%
\AgdaSymbol{\{}\AgdaArgument{N}\AgdaSpace{}%
\AgdaSymbol{=}\AgdaSpace{}%
\AgdaBound{N₁}\AgdaSymbol{\}}\AgdaSpace{}%
\AgdaSymbol{\{}\AgdaArgument{N'}\AgdaSpace{}%
\AgdaSymbol{=}\AgdaSpace{}%
\AgdaBound{N'}\AgdaSymbol{\}}\AgdaSpace{}%
\AgdaBound{red}\AgdaSymbol{)}\AgdaSpace{}%
\AgdaSymbol{=}\AgdaSpace{}%
\AgdaInductiveConstructor{μ}\AgdaSpace{}%
\AgdaSymbol{(}\AgdaInductiveConstructor{ξ}\AgdaSpace{}%
\AgdaSymbol{(}\AgdaInductiveConstructor{ν}\AgdaSpace{}%
\AgdaSymbol{(}\AgdaFunction{pres-rename-red-↝βb²}\AgdaSpace{}%
\AgdaInductiveConstructor{Sb}\AgdaSpace{}%
\AgdaSymbol{(}\AgdaFunction{embedTerm}\AgdaSpace{}%
\AgdaBound{N₁}\AgdaSymbol{)}\AgdaSpace{}%
\AgdaSymbol{(}\AgdaFunction{embedTerm}\AgdaSpace{}%
\AgdaBound{N'}\AgdaSymbol{)}\AgdaSpace{}%
\AgdaSymbol{(}\AgdaFunction{pres-red↝βb²}\AgdaSpace{}%
\AgdaBound{N₁}\AgdaSpace{}%
\AgdaBound{N'}\AgdaSpace{}%
\AgdaBound{red}\AgdaSymbol{))))}\<%
\\
\>[0]\AgdaFunction{pres-red↝βb²}\AgdaSpace{}%
\AgdaBound{M}\AgdaSpace{}%
\AgdaBound{N}\AgdaSpace{}%
\AgdaSymbol{(}\AgdaInductiveConstructor{ξ}\AgdaSpace{}%
\AgdaSymbol{\{}\AgdaArgument{M}\AgdaSpace{}%
\AgdaSymbol{=}\AgdaSpace{}%
\AgdaBound{M₁}\AgdaSymbol{\}}\AgdaSpace{}%
\AgdaSymbol{\{}\AgdaArgument{M'}\AgdaSpace{}%
\AgdaSymbol{=}\AgdaSpace{}%
\AgdaBound{M'}\AgdaSymbol{\}}\AgdaSpace{}%
\AgdaBound{red}\AgdaSymbol{)}\AgdaSpace{}%
\AgdaSymbol{=}\AgdaSpace{}%
\AgdaInductiveConstructor{ζ}\AgdaSpace{}%
\AgdaSymbol{(}\AgdaInductiveConstructor{ξ}\AgdaSpace{}%
\AgdaSymbol{(}\AgdaFunction{pres-red↝βb²}\AgdaSpace{}%
\AgdaBound{M₁}\AgdaSpace{}%
\AgdaBound{M'}\AgdaSpace{}%
\AgdaBound{red}\AgdaSymbol{))}\<%
\end{code}  
\end{lemma}

\begin{lemma}
  \label{lemm:subst:comp}
  \textbf{Compose} of renaming and substitution:
  \small
  \begin{code}[hide]%
\>[0]\AgdaKeyword{module}\AgdaSpace{}%
\AgdaModule{valueProofs2}\AgdaSpace{}%
\AgdaKeyword{where}\<%
\\
\>[0]\<%
\\
\>[0]\AgdaKeyword{open}\AgdaSpace{}%
\AgdaKeyword{import}\AgdaSpace{}%
\AgdaModule{calc}\<%
\\
\>[0]\AgdaKeyword{open}\AgdaSpace{}%
\AgdaKeyword{import}\AgdaSpace{}%
\AgdaModule{boxCalc}\<%
\\
\>[0]\AgdaKeyword{open}\AgdaSpace{}%
\AgdaKeyword{import}\AgdaSpace{}%
\AgdaModule{embedCBVIntoCBB}\<%
\\
\>[0]\AgdaKeyword{open}\AgdaSpace{}%
\AgdaKeyword{import}\AgdaSpace{}%
\AgdaModule{valueCalcRed}\<%
\\
\>[0]\AgdaKeyword{open}\AgdaSpace{}%
\AgdaKeyword{import}\AgdaSpace{}%
\AgdaModule{boxCalcRed}\<%
\\
%
\\[\AgdaEmptyExtraSkip]%
\>[0]\AgdaKeyword{open}\AgdaSpace{}%
\AgdaKeyword{import}\AgdaSpace{}%
\AgdaModule{Data.Empty}\AgdaSpace{}%
\AgdaKeyword{using}\AgdaSpace{}%
\AgdaSymbol{(}\AgdaFunction{⊥}\AgdaSymbol{;}\AgdaSpace{}%
\AgdaFunction{⊥-elim}\AgdaSymbol{)}\<%
\\
\>[0]\AgdaKeyword{import}\AgdaSpace{}%
\AgdaModule{Relation.Binary.PropositionalEquality}\AgdaSpace{}%
\AgdaSymbol{as}\AgdaSpace{}%
\AgdaModule{Eq}\<%
\\
\>[0]\AgdaKeyword{open}\AgdaSpace{}%
\AgdaModule{Eq}\AgdaSpace{}%
\AgdaKeyword{using}\AgdaSpace{}%
\AgdaSymbol{(}\AgdaOperator{\AgdaDatatype{\AgdaUnderscore{}≡\AgdaUnderscore{}}}\AgdaSymbol{;}\AgdaSpace{}%
\AgdaInductiveConstructor{refl}\AgdaSymbol{;}\AgdaSpace{}%
\AgdaFunction{cong}\AgdaSymbol{;}\AgdaSpace{}%
\AgdaFunction{cong₂}\AgdaSymbol{;}\AgdaSpace{}%
\AgdaFunction{sym}\AgdaSymbol{)}\<%
\\
%
\\[\AgdaEmptyExtraSkip]%
\>[0]\AgdaKeyword{open}\AgdaSpace{}%
\AgdaModule{Eq.≡-Reasoning}\AgdaSpace{}%
\AgdaKeyword{using}\AgdaSpace{}%
\AgdaSymbol{(}\AgdaOperator{\AgdaFunction{begin\AgdaUnderscore{}}}\AgdaSymbol{;}\AgdaSpace{}%
\AgdaFunction{step-≡-∣}\AgdaSymbol{;}\AgdaSpace{}%
\AgdaFunction{step-≡-⟩}\AgdaSymbol{;}\AgdaSpace{}%
\AgdaOperator{\AgdaFunction{\AgdaUnderscore{}∎}}\AgdaSymbol{)}\<%
\\
\>[0]\AgdaKeyword{open}\AgdaSpace{}%
\AgdaModule{Eq.≡-Reasoning}\AgdaSpace{}%
\AgdaKeyword{renaming}\AgdaSpace{}%
\AgdaSymbol{(}\AgdaOperator{\AgdaFunction{begin\AgdaUnderscore{}}}\AgdaSpace{}%
\AgdaSymbol{to}\AgdaSpace{}%
\AgdaOperator{\AgdaFunction{beginEq\AgdaUnderscore{}}}\AgdaSymbol{;}\AgdaSpace{}%
\AgdaOperator{\AgdaFunction{\AgdaUnderscore{}∎}}\AgdaSpace{}%
\AgdaSymbol{to}\AgdaSpace{}%
\AgdaOperator{\AgdaFunction{\AgdaUnderscore{}∎eq}}\AgdaSymbol{)}\<%
\\
%
\\[\AgdaEmptyExtraSkip]%
\>[0]\AgdaComment{--\ Preservation\ of\ typing\ by\ Gödel's\ translation}\<%
\\
\>[0]\AgdaFunction{pres-typing}\AgdaSpace{}%
\AgdaSymbol{:}\AgdaSpace{}%
\AgdaSymbol{∀}\AgdaSpace{}%
\AgdaSymbol{\{}\AgdaBound{Γ}\AgdaSpace{}%
\AgdaBound{A}\AgdaSymbol{\}}\AgdaSpace{}%
\AgdaSymbol{→}\AgdaSpace{}%
\AgdaBound{Γ}\AgdaSpace{}%
\AgdaOperator{\AgdaDatatype{⊢}}\AgdaSpace{}%
\AgdaBound{A}\AgdaSpace{}%
\AgdaSymbol{→}\AgdaSpace{}%
\AgdaFunction{embedContext}\AgdaSpace{}%
\AgdaBound{Γ}\AgdaSpace{}%
\AgdaOperator{\AgdaDatatype{⊢b}}\AgdaSpace{}%
\AgdaFunction{embedType}\AgdaSpace{}%
\AgdaBound{A}\<%
\\
\>[0]\AgdaFunction{pres-typing}\AgdaSpace{}%
\AgdaBound{x}\AgdaSpace{}%
\AgdaSymbol{=}\AgdaSpace{}%
\AgdaFunction{embedTerm}\AgdaSpace{}%
\AgdaBound{x}\<%
\\
%
\\[\AgdaEmptyExtraSkip]%
\>[0]\AgdaComment{--\ Preservation\ of\ typing\ of\ Values\ by\ Gödel's\ translation}\<%
\\
\>[0]\AgdaFunction{pres-typing-val}\AgdaSpace{}%
\AgdaSymbol{:}\AgdaSpace{}%
\AgdaSymbol{∀}\AgdaSpace{}%
\AgdaSymbol{\{}\AgdaBound{Γ}\AgdaSpace{}%
\AgdaBound{A}\AgdaSymbol{\}}\AgdaSpace{}%
\AgdaSymbol{→}\AgdaSpace{}%
\AgdaSymbol{(}\AgdaBound{V}\AgdaSpace{}%
\AgdaSymbol{:}\AgdaSpace{}%
\AgdaBound{Γ}\AgdaSpace{}%
\AgdaOperator{\AgdaDatatype{⊢}}\AgdaSpace{}%
\AgdaBound{A}\AgdaSymbol{)}\AgdaSpace{}%
\AgdaSymbol{→}\AgdaSpace{}%
\AgdaSymbol{(}\AgdaDatatype{Value}\AgdaSpace{}%
\AgdaBound{V}\AgdaSymbol{)}\AgdaSpace{}%
\AgdaSymbol{→}\AgdaSpace{}%
\AgdaFunction{embedContext}\AgdaSpace{}%
\AgdaBound{Γ}\AgdaSpace{}%
\AgdaOperator{\AgdaDatatype{⊢b}}\AgdaSpace{}%
\AgdaFunction{embedTypeValue}\AgdaSpace{}%
\AgdaBound{A}\<%
\\
\>[0]\AgdaFunction{pres-typing-val}\AgdaSpace{}%
\AgdaBound{V}\AgdaSpace{}%
\AgdaBound{x}\AgdaSpace{}%
\AgdaSymbol{=}\AgdaSpace{}%
\AgdaFunction{embedValue}\AgdaSpace{}%
\AgdaBound{V}\AgdaSpace{}%
\AgdaBound{x}\<%
\\
%
\\[\AgdaEmptyExtraSkip]%
\>[0]\AgdaComment{--\ Helper\ lemma:\ applying\ embedTerm\ to\ a\ Value\ is\ equivalent\ to\ wrapping\ the\ result\ of\ embedValue\ in\ a\ box}\<%
\\
\>[0]\AgdaFunction{box-embedValue}\AgdaSpace{}%
\AgdaSymbol{:}\AgdaSpace{}%
\AgdaSymbol{∀}\AgdaSpace{}%
\AgdaSymbol{\{}\AgdaBound{Γ}\AgdaSpace{}%
\AgdaBound{A}\AgdaSymbol{\}}\AgdaSpace{}%
\AgdaSymbol{→}\AgdaSpace{}%
\AgdaSymbol{\{}\AgdaBound{V}\AgdaSpace{}%
\AgdaSymbol{:}\AgdaSpace{}%
\AgdaBound{Γ}\AgdaSpace{}%
\AgdaOperator{\AgdaDatatype{⊢}}\AgdaSpace{}%
\AgdaBound{A}\AgdaSymbol{\}}\AgdaSpace{}%
\AgdaSymbol{→}\AgdaSpace{}%
\AgdaSymbol{(}\AgdaBound{x}\AgdaSpace{}%
\AgdaSymbol{:}\AgdaSpace{}%
\AgdaDatatype{Value}\AgdaSpace{}%
\AgdaBound{V}\AgdaSymbol{)}\AgdaSpace{}%
\AgdaSymbol{→}\AgdaSpace{}%
\AgdaInductiveConstructor{box}\AgdaSpace{}%
\AgdaSymbol{(}\AgdaFunction{embedValue}\AgdaSpace{}%
\AgdaBound{V}\AgdaSpace{}%
\AgdaBound{x}\AgdaSymbol{)}\AgdaSpace{}%
\AgdaOperator{\AgdaDatatype{≡}}\AgdaSpace{}%
\AgdaFunction{embedTerm}\AgdaSpace{}%
\AgdaBound{V}\<%
\\
\>[0]\AgdaFunction{box-embedValue}\AgdaSpace{}%
\AgdaInductiveConstructor{V-ƛ}\AgdaSpace{}%
\AgdaSymbol{=}\AgdaSpace{}%
\AgdaInductiveConstructor{refl}\<%
\\
\>[0]\AgdaFunction{box-embedValue}\AgdaSpace{}%
\AgdaInductiveConstructor{V-`}\AgdaSpace{}%
\AgdaSymbol{=}\AgdaSpace{}%
\AgdaInductiveConstructor{refl}\<%
\\
%
\\[\AgdaEmptyExtraSkip]%
\>[0]\AgdaComment{--\ HELPER\ LEMMA'S\ SIMULTANIOUS\ SUBSTITUTION}\<%
\\
%
\\[\AgdaEmptyExtraSkip]%
\>[0]\AgdaComment{--\ Preservation\ and\ Reflection\ of\ }\<%
\\
\>[0]\AgdaComment{--\ ext\ of\ Gödel's\ translation}\<%
\\
\>[0]\AgdaFunction{pres-ext}\AgdaSpace{}%
\AgdaSymbol{:}\AgdaSpace{}%
\AgdaSymbol{∀}\AgdaSpace{}%
\AgdaSymbol{\{}\AgdaBound{Γ}\AgdaSpace{}%
\AgdaBound{Δ}\AgdaSpace{}%
\AgdaBound{A'}\AgdaSpace{}%
\AgdaBound{A}\AgdaSymbol{\}}\<%
\\
\>[0][@{}l@{\AgdaIndent{0}}]%
\>[2]\AgdaSymbol{→}\AgdaSpace{}%
\AgdaSymbol{(}\AgdaBound{f'}\AgdaSpace{}%
\AgdaSymbol{:}\AgdaSpace{}%
\AgdaSymbol{∀}\AgdaSpace{}%
\AgdaSymbol{\{}\AgdaBound{A}\AgdaSymbol{\}}\AgdaSpace{}%
\AgdaSymbol{→}\AgdaSpace{}%
\AgdaBound{A}\AgdaSpace{}%
\AgdaOperator{\AgdaDatatype{∈}}\AgdaSpace{}%
\AgdaBound{Γ}\AgdaSpace{}%
\AgdaSymbol{→}\AgdaSpace{}%
\AgdaBound{A}\AgdaSpace{}%
\AgdaOperator{\AgdaDatatype{∈}}\AgdaSpace{}%
\AgdaBound{Δ}\AgdaSymbol{)}\<%
\\
%
\>[2]\AgdaSymbol{→}\AgdaSpace{}%
\AgdaSymbol{(}\AgdaBound{fb'}\AgdaSpace{}%
\AgdaSymbol{:}\AgdaSpace{}%
\AgdaSymbol{∀}\AgdaSpace{}%
\AgdaSymbol{\{}\AgdaBound{Ab}\AgdaSymbol{\}}\AgdaSpace{}%
\AgdaSymbol{→}\AgdaSpace{}%
\AgdaBound{Ab}\AgdaSpace{}%
\AgdaOperator{\AgdaDatatype{∈b}}\AgdaSpace{}%
\AgdaFunction{embedContext}\AgdaSpace{}%
\AgdaBound{Γ}\AgdaSpace{}%
\AgdaSymbol{→}\AgdaSpace{}%
\AgdaBound{Ab}\AgdaSpace{}%
\AgdaOperator{\AgdaDatatype{∈b}}\AgdaSpace{}%
\AgdaFunction{embedContext}\AgdaSpace{}%
\AgdaBound{Δ}\AgdaSymbol{)}\<%
\\
%
\>[2]\AgdaSymbol{→}\AgdaSpace{}%
\AgdaSymbol{(\{}\AgdaBound{A'}\AgdaSpace{}%
\AgdaSymbol{:}\AgdaSpace{}%
\AgdaDatatype{Type}\AgdaSymbol{\}}\AgdaSpace{}%
\AgdaSymbol{(}\AgdaBound{x}\AgdaSpace{}%
\AgdaSymbol{:}\AgdaSpace{}%
\AgdaBound{A'}\AgdaSpace{}%
\AgdaOperator{\AgdaDatatype{∈}}\AgdaSpace{}%
\AgdaBound{Γ}\AgdaSymbol{)}\AgdaSpace{}%
\AgdaSymbol{→}\AgdaSpace{}%
\AgdaBound{fb'}\AgdaSpace{}%
\AgdaSymbol{(}\AgdaFunction{embedLookup}\AgdaSpace{}%
\AgdaBound{x}\AgdaSymbol{)}\AgdaSpace{}%
\AgdaOperator{\AgdaDatatype{≡}}\AgdaSpace{}%
\AgdaFunction{embedLookup}\AgdaSpace{}%
\AgdaSymbol{(}\AgdaBound{f'}\AgdaSpace{}%
\AgdaBound{x}\AgdaSymbol{))}\<%
\\
%
\>[2]\AgdaSymbol{→}\AgdaSpace{}%
\AgdaSymbol{(}\AgdaBound{x}\AgdaSpace{}%
\AgdaSymbol{:}\AgdaSpace{}%
\AgdaBound{A'}\AgdaSpace{}%
\AgdaOperator{\AgdaDatatype{∈}}\AgdaSpace{}%
\AgdaBound{Γ}\AgdaSpace{}%
\AgdaOperator{\AgdaInductiveConstructor{,}}\AgdaSpace{}%
\AgdaBound{A}\AgdaSymbol{)}\<%
\\
%
\>[2]\AgdaSymbol{→}\AgdaSpace{}%
\AgdaFunction{extb}\AgdaSpace{}%
\AgdaBound{fb'}\AgdaSpace{}%
\AgdaSymbol{(}\AgdaFunction{embedLookup}\AgdaSpace{}%
\AgdaBound{x}\AgdaSymbol{)}\AgdaSpace{}%
\AgdaOperator{\AgdaDatatype{≡}}\AgdaSpace{}%
\AgdaFunction{embedLookup}\AgdaSpace{}%
\AgdaSymbol{(}\AgdaFunction{ext}\AgdaSpace{}%
\AgdaBound{f'}\AgdaSpace{}%
\AgdaBound{x}\AgdaSymbol{)}\<%
\\
\>[0]\AgdaFunction{pres-ext}\AgdaSpace{}%
\AgdaBound{f'}\AgdaSpace{}%
\AgdaBound{fb'}\AgdaSpace{}%
\AgdaBound{h}\AgdaSpace{}%
\AgdaInductiveConstructor{Z}\AgdaSpace{}%
\AgdaSymbol{=}\AgdaSpace{}%
\AgdaInductiveConstructor{refl}\<%
\\
\>[0]\AgdaFunction{pres-ext}\AgdaSpace{}%
\AgdaBound{f'}\AgdaSpace{}%
\AgdaBound{fb'}\AgdaSpace{}%
\AgdaBound{h}\AgdaSpace{}%
\AgdaSymbol{(}\AgdaInductiveConstructor{S}\AgdaSpace{}%
\AgdaBound{x}\AgdaSymbol{)}\AgdaSpace{}%
\AgdaSymbol{=}\AgdaSpace{}%
\AgdaFunction{cong}\AgdaSpace{}%
\AgdaInductiveConstructor{Sb}\AgdaSpace{}%
\AgdaSymbol{(}\AgdaBound{h}\AgdaSpace{}%
\AgdaBound{x}\AgdaSymbol{)}\<%
\\
%
\\[\AgdaEmptyExtraSkip]%
\>[0]\AgdaComment{--\ Extb\ preserves\ pointwise\ equality\ }\<%
\\
\>[0]\AgdaComment{--\ of\ renaming\ functions}\<%
\\
\>[0]\AgdaFunction{extb-cong}\AgdaSpace{}%
\AgdaSymbol{:}\AgdaSpace{}%
\AgdaSymbol{∀}\AgdaSpace{}%
\AgdaSymbol{\{}\AgdaBound{Γ₁}\AgdaSpace{}%
\AgdaBound{Γ₂}\AgdaSpace{}%
\AgdaBound{A}\AgdaSpace{}%
\AgdaBound{A'}\AgdaSymbol{\}}\<%
\\
\>[0][@{}l@{\AgdaIndent{0}}]%
\>[2]\AgdaSymbol{→}\AgdaSpace{}%
\AgdaSymbol{\{}\AgdaBound{fb'}\AgdaSpace{}%
\AgdaSymbol{:}\AgdaSpace{}%
\AgdaSymbol{∀}\AgdaSpace{}%
\AgdaSymbol{\{}\AgdaBound{Ab}\AgdaSymbol{\}}\AgdaSpace{}%
\AgdaSymbol{→}\AgdaSpace{}%
\AgdaBound{Ab}\AgdaSpace{}%
\AgdaOperator{\AgdaDatatype{∈b}}\AgdaSpace{}%
\AgdaBound{Γ₁}\AgdaSpace{}%
\AgdaSymbol{→}\AgdaSpace{}%
\AgdaBound{Ab}\AgdaSpace{}%
\AgdaOperator{\AgdaDatatype{∈b}}\AgdaSpace{}%
\AgdaBound{Γ₂}\AgdaSymbol{\}}\<%
\\
%
\>[2]\AgdaSymbol{→}\AgdaSpace{}%
\AgdaSymbol{\{}\AgdaBound{fb''}\AgdaSpace{}%
\AgdaSymbol{:}\AgdaSpace{}%
\AgdaSymbol{∀}\AgdaSpace{}%
\AgdaSymbol{\{}\AgdaBound{Ab}\AgdaSymbol{\}}\AgdaSpace{}%
\AgdaSymbol{→}\AgdaSpace{}%
\AgdaBound{Ab}\AgdaSpace{}%
\AgdaOperator{\AgdaDatatype{∈b}}\AgdaSpace{}%
\AgdaBound{Γ₁}\AgdaSpace{}%
\AgdaSymbol{→}\AgdaSpace{}%
\AgdaBound{Ab}\AgdaSpace{}%
\AgdaOperator{\AgdaDatatype{∈b}}\AgdaSpace{}%
\AgdaBound{Γ₂}\AgdaSymbol{\}}\<%
\\
%
\>[2]\AgdaSymbol{→}\AgdaSpace{}%
\AgdaSymbol{(}\AgdaBound{h}\AgdaSpace{}%
\AgdaSymbol{:}\AgdaSpace{}%
\AgdaSymbol{∀}\AgdaSpace{}%
\AgdaSymbol{\{}\AgdaBound{Ab}\AgdaSymbol{\}}\AgdaSpace{}%
\AgdaSymbol{→}\AgdaSpace{}%
\AgdaSymbol{(}\AgdaBound{x}\AgdaSpace{}%
\AgdaSymbol{:}\AgdaSpace{}%
\AgdaBound{Ab}\AgdaSpace{}%
\AgdaOperator{\AgdaDatatype{∈b}}\AgdaSpace{}%
\AgdaBound{Γ₁}\AgdaSymbol{)}\AgdaSpace{}%
\AgdaSymbol{→}\AgdaSpace{}%
\AgdaBound{fb'}\AgdaSpace{}%
\AgdaBound{x}\AgdaSpace{}%
\AgdaOperator{\AgdaDatatype{≡}}\AgdaSpace{}%
\AgdaBound{fb''}\AgdaSpace{}%
\AgdaBound{x}\AgdaSymbol{)}\<%
\\
%
\>[2]\AgdaSymbol{→}\AgdaSpace{}%
\AgdaSymbol{(}\AgdaBound{x}\AgdaSpace{}%
\AgdaSymbol{:}\AgdaSpace{}%
\AgdaBound{A}\AgdaSpace{}%
\AgdaOperator{\AgdaDatatype{∈b}}\AgdaSpace{}%
\AgdaBound{Γ₁}\AgdaSpace{}%
\AgdaOperator{\AgdaInductiveConstructor{,b}}\AgdaSpace{}%
\AgdaInductiveConstructor{□}\AgdaSpace{}%
\AgdaBound{A'}\AgdaSymbol{)}\<%
\\
%
\>[2]\AgdaSymbol{→}\AgdaSpace{}%
\AgdaFunction{extb}\AgdaSpace{}%
\AgdaBound{fb'}\AgdaSpace{}%
\AgdaBound{x}\AgdaSpace{}%
\AgdaOperator{\AgdaDatatype{≡}}\AgdaSpace{}%
\AgdaFunction{extb}\AgdaSpace{}%
\AgdaBound{fb''}\AgdaSpace{}%
\AgdaBound{x}\<%
\\
\>[0]\AgdaFunction{extb-cong}\AgdaSpace{}%
\AgdaBound{h}\AgdaSpace{}%
\AgdaInductiveConstructor{Zb}\AgdaSpace{}%
\AgdaSymbol{=}\AgdaSpace{}%
\AgdaInductiveConstructor{refl}\<%
\\
\>[0]\AgdaFunction{extb-cong}\AgdaSpace{}%
\AgdaBound{h}\AgdaSpace{}%
\AgdaSymbol{(}\AgdaInductiveConstructor{Sb}\AgdaSpace{}%
\AgdaBound{x}\AgdaSymbol{)}\AgdaSpace{}%
\AgdaSymbol{=}\AgdaSpace{}%
\AgdaFunction{cong}\AgdaSpace{}%
\AgdaInductiveConstructor{Sb}\AgdaSpace{}%
\AgdaSymbol{(}\AgdaBound{h}\AgdaSpace{}%
\AgdaBound{x}\AgdaSymbol{)}\<%
\\
%
\\[\AgdaEmptyExtraSkip]%
\>[0]\AgdaComment{--\ Renameb\ preserves\ pointwise\ equality}\<%
\\
\>[0]\AgdaComment{--\ of\ renaming\ functions}\<%
\\
\>[0]\AgdaFunction{renameb-cong}\AgdaSpace{}%
\AgdaSymbol{:}\AgdaSpace{}%
\AgdaSymbol{∀}\AgdaSpace{}%
\AgdaSymbol{\{}\AgdaBound{Γ₁}\AgdaSpace{}%
\AgdaBound{Γ₂}\AgdaSpace{}%
\AgdaBound{A}\AgdaSymbol{\}}\<%
\\
\>[0][@{}l@{\AgdaIndent{0}}]%
\>[2]\AgdaSymbol{→}\AgdaSpace{}%
\AgdaSymbol{\{}\AgdaBound{fb'}\AgdaSpace{}%
\AgdaSymbol{:}\AgdaSpace{}%
\AgdaSymbol{∀}\AgdaSpace{}%
\AgdaSymbol{\{}\AgdaBound{Ab}\AgdaSymbol{\}}\AgdaSpace{}%
\AgdaSymbol{→}\AgdaSpace{}%
\AgdaBound{Ab}\AgdaSpace{}%
\AgdaOperator{\AgdaDatatype{∈b}}\AgdaSpace{}%
\AgdaBound{Γ₁}\AgdaSpace{}%
\AgdaSymbol{→}\AgdaSpace{}%
\AgdaBound{Ab}\AgdaSpace{}%
\AgdaOperator{\AgdaDatatype{∈b}}\AgdaSpace{}%
\AgdaBound{Γ₂}\AgdaSymbol{\}}\<%
\\
%
\>[2]\AgdaSymbol{→}\AgdaSpace{}%
\AgdaSymbol{\{}\AgdaBound{fb''}\AgdaSpace{}%
\AgdaSymbol{:}\AgdaSpace{}%
\AgdaSymbol{∀}\AgdaSpace{}%
\AgdaSymbol{\{}\AgdaBound{Ab}\AgdaSymbol{\}}\AgdaSpace{}%
\AgdaSymbol{→}\AgdaSpace{}%
\AgdaBound{Ab}\AgdaSpace{}%
\AgdaOperator{\AgdaDatatype{∈b}}\AgdaSpace{}%
\AgdaBound{Γ₁}\AgdaSpace{}%
\AgdaSymbol{→}\AgdaSpace{}%
\AgdaBound{Ab}\AgdaSpace{}%
\AgdaOperator{\AgdaDatatype{∈b}}\AgdaSpace{}%
\AgdaBound{Γ₂}\AgdaSymbol{\}}\<%
\\
%
\>[2]\AgdaSymbol{→}\AgdaSpace{}%
\AgdaSymbol{(}\AgdaBound{h}\AgdaSpace{}%
\AgdaSymbol{:}\AgdaSpace{}%
\AgdaSymbol{∀}\AgdaSpace{}%
\AgdaSymbol{\{}\AgdaBound{Ab}\AgdaSymbol{\}}\AgdaSpace{}%
\AgdaSymbol{→}\AgdaSpace{}%
\AgdaSymbol{(}\AgdaBound{x}\AgdaSpace{}%
\AgdaSymbol{:}\AgdaSpace{}%
\AgdaBound{Ab}\AgdaSpace{}%
\AgdaOperator{\AgdaDatatype{∈b}}\AgdaSpace{}%
\AgdaBound{Γ₁}\AgdaSymbol{)}\AgdaSpace{}%
\AgdaSymbol{→}\AgdaSpace{}%
\AgdaBound{fb'}\AgdaSpace{}%
\AgdaBound{x}\AgdaSpace{}%
\AgdaOperator{\AgdaDatatype{≡}}\AgdaSpace{}%
\AgdaBound{fb''}\AgdaSpace{}%
\AgdaBound{x}\AgdaSymbol{)}\<%
\\
%
\>[2]\AgdaSymbol{→}\AgdaSpace{}%
\AgdaSymbol{(}\AgdaBound{t}\AgdaSpace{}%
\AgdaSymbol{:}\AgdaSpace{}%
\AgdaBound{Γ₁}\AgdaSpace{}%
\AgdaOperator{\AgdaDatatype{⊢b}}\AgdaSpace{}%
\AgdaBound{A}\AgdaSymbol{)}\<%
\\
%
\>[2]\AgdaSymbol{→}\AgdaSpace{}%
\AgdaFunction{renameb}\AgdaSpace{}%
\AgdaBound{fb'}\AgdaSpace{}%
\AgdaBound{t}\AgdaSpace{}%
\AgdaOperator{\AgdaDatatype{≡}}\AgdaSpace{}%
\AgdaFunction{renameb}\AgdaSpace{}%
\AgdaBound{fb''}\AgdaSpace{}%
\AgdaBound{t}\<%
\\
\>[0]\AgdaFunction{renameb-cong}\AgdaSpace{}%
\AgdaBound{h}\AgdaSpace{}%
\AgdaSymbol{(}\AgdaInductiveConstructor{ε}\AgdaSpace{}%
\AgdaBound{x}\AgdaSymbol{)}\AgdaSpace{}%
\AgdaSymbol{=}\AgdaSpace{}%
\AgdaFunction{cong}\AgdaSpace{}%
\AgdaInductiveConstructor{ε}\AgdaSpace{}%
\AgdaSymbol{(}\AgdaBound{h}\AgdaSpace{}%
\AgdaBound{x}\AgdaSymbol{)}\<%
\\
\>[0]\AgdaFunction{renameb-cong}\AgdaSpace{}%
\AgdaBound{h}\AgdaSpace{}%
\AgdaSymbol{(}\AgdaInductiveConstructor{box}\AgdaSpace{}%
\AgdaBound{t}\AgdaSymbol{)}\AgdaSpace{}%
\AgdaSymbol{=}\AgdaSpace{}%
\AgdaFunction{cong}\AgdaSpace{}%
\AgdaInductiveConstructor{box}\AgdaSpace{}%
\AgdaSymbol{(}\AgdaFunction{renameb-cong}\AgdaSpace{}%
\AgdaBound{h}\AgdaSpace{}%
\AgdaBound{t}\AgdaSymbol{)}\<%
\\
\>[0]\AgdaFunction{renameb-cong}\AgdaSpace{}%
\AgdaBound{h}\AgdaSpace{}%
\AgdaSymbol{(}\AgdaInductiveConstructor{ƛb}\AgdaSpace{}%
\AgdaBound{t}\AgdaSymbol{)}\AgdaSpace{}%
\AgdaSymbol{=}\AgdaSpace{}%
\AgdaFunction{cong}\AgdaSpace{}%
\AgdaInductiveConstructor{ƛb}\AgdaSpace{}%
\AgdaSymbol{(}\AgdaFunction{renameb-cong}\AgdaSpace{}%
\AgdaSymbol{(}\AgdaFunction{extb-cong}\AgdaSpace{}%
\AgdaBound{h}\AgdaSymbol{)}\AgdaSpace{}%
\AgdaBound{t}\AgdaSymbol{)}\<%
\\
\>[0]\AgdaFunction{renameb-cong}\AgdaSpace{}%
\AgdaBound{h}\AgdaSpace{}%
\AgdaSymbol{(}\AgdaBound{t}\AgdaSpace{}%
\AgdaOperator{\AgdaInductiveConstructor{·b}}\AgdaSpace{}%
\AgdaBound{t₁}\AgdaSymbol{)}\AgdaSpace{}%
\AgdaSymbol{=}\AgdaSpace{}%
\AgdaFunction{cong₂}\AgdaSpace{}%
\AgdaOperator{\AgdaInductiveConstructor{\AgdaUnderscore{}·b\AgdaUnderscore{}}}\AgdaSpace{}%
\AgdaSymbol{(}\AgdaFunction{renameb-cong}\AgdaSpace{}%
\AgdaBound{h}\AgdaSpace{}%
\AgdaBound{t}\AgdaSymbol{)}\AgdaSpace{}%
\AgdaSymbol{(}\AgdaFunction{renameb-cong}\AgdaSpace{}%
\AgdaBound{h}\AgdaSpace{}%
\AgdaBound{t₁}\AgdaSymbol{)}\<%
\\
%
\\[\AgdaEmptyExtraSkip]%
\>[0]\AgdaComment{--\ extb\ distributes\ over\ composition:}\<%
\\
\>[0]\AgdaComment{--\ Extending\ first\ with\ fb',\ then\ with\ fb'',\ }\<%
\\
\>[0]\AgdaComment{--\ is\ the\ same\ as\ extending\ once\ with\ fb''\ ∘\ fb'.}\<%
\\
\>[0]\AgdaFunction{extb-comp}\AgdaSpace{}%
\AgdaSymbol{:}\AgdaSpace{}%
\AgdaSymbol{∀}\AgdaSpace{}%
\AgdaSymbol{\{}\AgdaBound{Γ₁}\AgdaSpace{}%
\AgdaBound{Γ₂}\AgdaSpace{}%
\AgdaBound{Γ₃}\AgdaSpace{}%
\AgdaBound{A}\AgdaSpace{}%
\AgdaBound{B}\AgdaSymbol{\}}\<%
\\
\>[0][@{}l@{\AgdaIndent{0}}]%
\>[2]\AgdaSymbol{→}\AgdaSpace{}%
\AgdaSymbol{(}\AgdaBound{fb'}\AgdaSpace{}%
\AgdaSymbol{:}\AgdaSpace{}%
\AgdaSymbol{∀}\AgdaSpace{}%
\AgdaSymbol{\{}\AgdaBound{Ab}\AgdaSymbol{\}}\AgdaSpace{}%
\AgdaSymbol{→}\AgdaSpace{}%
\AgdaBound{Ab}\AgdaSpace{}%
\AgdaOperator{\AgdaDatatype{∈b}}\AgdaSpace{}%
\AgdaBound{Γ₁}\AgdaSpace{}%
\AgdaSymbol{→}\AgdaSpace{}%
\AgdaBound{Ab}\AgdaSpace{}%
\AgdaOperator{\AgdaDatatype{∈b}}\AgdaSpace{}%
\AgdaBound{Γ₂}\AgdaSymbol{)}\<%
\\
%
\>[2]\AgdaSymbol{→}\AgdaSpace{}%
\AgdaSymbol{(}\AgdaBound{fb''}\AgdaSpace{}%
\AgdaSymbol{:}\AgdaSpace{}%
\AgdaSymbol{∀}\AgdaSpace{}%
\AgdaSymbol{\{}\AgdaBound{Ab}\AgdaSymbol{\}}\AgdaSpace{}%
\AgdaSymbol{→}\AgdaSpace{}%
\AgdaBound{Ab}\AgdaSpace{}%
\AgdaOperator{\AgdaDatatype{∈b}}\AgdaSpace{}%
\AgdaBound{Γ₂}\AgdaSpace{}%
\AgdaSymbol{→}\AgdaSpace{}%
\AgdaBound{Ab}\AgdaSpace{}%
\AgdaOperator{\AgdaDatatype{∈b}}\AgdaSpace{}%
\AgdaBound{Γ₃}\AgdaSymbol{)}\<%
\\
%
\>[2]\AgdaSymbol{→}\AgdaSpace{}%
\AgdaSymbol{(}\AgdaBound{x}\AgdaSpace{}%
\AgdaSymbol{:}\AgdaSpace{}%
\AgdaBound{A}\AgdaSpace{}%
\AgdaOperator{\AgdaDatatype{∈b}}\AgdaSpace{}%
\AgdaBound{Γ₁}\AgdaSpace{}%
\AgdaOperator{\AgdaInductiveConstructor{,b}}\AgdaSpace{}%
\AgdaBound{B}\AgdaSymbol{)}\<%
\\
%
\>[2]\AgdaSymbol{→}\AgdaSpace{}%
\AgdaFunction{extb}\AgdaSpace{}%
\AgdaBound{fb''}\AgdaSpace{}%
\AgdaSymbol{(}\AgdaFunction{extb}\AgdaSpace{}%
\AgdaBound{fb'}\AgdaSpace{}%
\AgdaBound{x}\AgdaSymbol{)}\AgdaSpace{}%
\AgdaOperator{\AgdaDatatype{≡}}\AgdaSpace{}%
\AgdaFunction{extb}\AgdaSpace{}%
\AgdaSymbol{(λ}\AgdaSpace{}%
\AgdaBound{x₁}\AgdaSpace{}%
\AgdaSymbol{→}\AgdaSpace{}%
\AgdaBound{fb''}\AgdaSpace{}%
\AgdaSymbol{(}\AgdaBound{fb'}\AgdaSpace{}%
\AgdaBound{x₁}\AgdaSymbol{))}\AgdaSpace{}%
\AgdaBound{x}\<%
\\
\>[0]\AgdaFunction{extb-comp}\AgdaSpace{}%
\AgdaBound{fb'}\AgdaSpace{}%
\AgdaBound{fb''}\AgdaSpace{}%
\AgdaInductiveConstructor{Zb}\AgdaSpace{}%
\AgdaSymbol{=}\AgdaSpace{}%
\AgdaInductiveConstructor{refl}\<%
\\
\>[0]\AgdaFunction{extb-comp}\AgdaSpace{}%
\AgdaBound{fb'}\AgdaSpace{}%
\AgdaBound{fb''}\AgdaSpace{}%
\AgdaSymbol{(}\AgdaInductiveConstructor{Sb}\AgdaSpace{}%
\AgdaBound{x}\AgdaSymbol{)}\AgdaSpace{}%
\AgdaSymbol{=}\AgdaSpace{}%
\AgdaInductiveConstructor{refl}\<%
\\
%
\\[\AgdaEmptyExtraSkip]%
\>[0]\AgdaComment{--\ renameb\ distributes\ over\ composition:}\<%
\\
\>[0]\AgdaComment{--\ Renaming\ first\ with\ fb',\ then\ with\ fb''}\<%
\\
\>[0]\AgdaComment{--\ is\ the\ same\ as\ renaming\ once\ with\ fb''\ ∘\ fb'.}\<%
\\
\>[0]\AgdaFunction{renameb-comp}\AgdaSpace{}%
\AgdaSymbol{:}\AgdaSpace{}%
\AgdaSymbol{∀}\AgdaSpace{}%
\AgdaSymbol{\{}\AgdaBound{Γ₁}\AgdaSpace{}%
\AgdaBound{Γ₂}\AgdaSpace{}%
\AgdaBound{Γ₃}\AgdaSpace{}%
\AgdaBound{A}\AgdaSymbol{\}}\<%
\\
\>[0][@{}l@{\AgdaIndent{0}}]%
\>[2]\AgdaSymbol{→}\AgdaSpace{}%
\AgdaSymbol{(}\AgdaBound{fb'}\AgdaSpace{}%
\AgdaSymbol{:}\AgdaSpace{}%
\AgdaSymbol{∀}\AgdaSpace{}%
\AgdaSymbol{\{}\AgdaBound{Ab}\AgdaSymbol{\}}\AgdaSpace{}%
\AgdaSymbol{→}\AgdaSpace{}%
\AgdaBound{Ab}\AgdaSpace{}%
\AgdaOperator{\AgdaDatatype{∈b}}\AgdaSpace{}%
\AgdaBound{Γ₁}\AgdaSpace{}%
\AgdaSymbol{→}\AgdaSpace{}%
\AgdaBound{Ab}\AgdaSpace{}%
\AgdaOperator{\AgdaDatatype{∈b}}\AgdaSpace{}%
\AgdaBound{Γ₂}\AgdaSymbol{)}\<%
\\
%
\>[2]\AgdaSymbol{→}\AgdaSpace{}%
\AgdaSymbol{(}\AgdaBound{fb''}\AgdaSpace{}%
\AgdaSymbol{:}\AgdaSpace{}%
\AgdaSymbol{∀}\AgdaSpace{}%
\AgdaSymbol{\{}\AgdaBound{Ab}\AgdaSymbol{\}}\AgdaSpace{}%
\AgdaSymbol{→}\AgdaSpace{}%
\AgdaBound{Ab}\AgdaSpace{}%
\AgdaOperator{\AgdaDatatype{∈b}}\AgdaSpace{}%
\AgdaBound{Γ₂}\AgdaSpace{}%
\AgdaSymbol{→}\AgdaSpace{}%
\AgdaBound{Ab}\AgdaSpace{}%
\AgdaOperator{\AgdaDatatype{∈b}}\AgdaSpace{}%
\AgdaBound{Γ₃}\AgdaSymbol{)}\<%
\\
%
\>[2]\AgdaSymbol{→}\AgdaSpace{}%
\AgdaSymbol{(}\AgdaBound{t}\AgdaSpace{}%
\AgdaSymbol{:}\AgdaSpace{}%
\AgdaBound{Γ₁}\AgdaSpace{}%
\AgdaOperator{\AgdaDatatype{⊢b}}\AgdaSpace{}%
\AgdaBound{A}\AgdaSymbol{)}\<%
\\
%
\>[2]\AgdaSymbol{→}\AgdaSpace{}%
\AgdaFunction{renameb}\AgdaSpace{}%
\AgdaBound{fb''}\AgdaSpace{}%
\AgdaSymbol{(}\AgdaFunction{renameb}\AgdaSpace{}%
\AgdaBound{fb'}\AgdaSpace{}%
\AgdaBound{t}\AgdaSymbol{)}\AgdaSpace{}%
\AgdaOperator{\AgdaDatatype{≡}}\AgdaSpace{}%
\AgdaFunction{renameb}\AgdaSpace{}%
\AgdaSymbol{(λ}\AgdaSpace{}%
\AgdaBound{x}\AgdaSpace{}%
\AgdaSymbol{→}\AgdaSpace{}%
\AgdaBound{fb''}\AgdaSpace{}%
\AgdaSymbol{(}\AgdaBound{fb'}\AgdaSpace{}%
\AgdaBound{x}\AgdaSymbol{))}\AgdaSpace{}%
\AgdaBound{t}\<%
\\
\>[0]\AgdaFunction{renameb-comp}\AgdaSpace{}%
\AgdaBound{fb'}\AgdaSpace{}%
\AgdaBound{fb''}\AgdaSpace{}%
\AgdaSymbol{(}\AgdaInductiveConstructor{ε}\AgdaSpace{}%
\AgdaBound{x}\AgdaSymbol{)}\AgdaSpace{}%
\AgdaSymbol{=}\AgdaSpace{}%
\AgdaInductiveConstructor{refl}\<%
\\
\>[0]\AgdaFunction{renameb-comp}\AgdaSpace{}%
\AgdaBound{fb'}\AgdaSpace{}%
\AgdaBound{fb''}\AgdaSpace{}%
\AgdaSymbol{(}\AgdaInductiveConstructor{box}\AgdaSpace{}%
\AgdaBound{t}\AgdaSymbol{)}\AgdaSpace{}%
\AgdaSymbol{=}\AgdaSpace{}%
\AgdaFunction{cong}\AgdaSpace{}%
\AgdaInductiveConstructor{box}\AgdaSpace{}%
\AgdaSymbol{(}\AgdaFunction{renameb-comp}\AgdaSpace{}%
\AgdaBound{fb'}\AgdaSpace{}%
\AgdaBound{fb''}\AgdaSpace{}%
\AgdaBound{t}\AgdaSymbol{)}\<%
\\
\>[0]\AgdaFunction{renameb-comp}\AgdaSpace{}%
\AgdaBound{fb'}\AgdaSpace{}%
\AgdaBound{fb''}\AgdaSpace{}%
\AgdaSymbol{(}\AgdaInductiveConstructor{ƛb}\AgdaSpace{}%
\AgdaBound{t}\AgdaSymbol{)}\AgdaSpace{}%
\AgdaSymbol{=}\<%
\\
\>[0][@{}l@{\AgdaIndent{0}}]%
\>[2]\AgdaOperator{\AgdaFunction{beginEq}}\<%
\\
\>[2][@{}l@{\AgdaIndent{0}}]%
\>[4]\AgdaInductiveConstructor{ƛb}\AgdaSpace{}%
\AgdaSymbol{(}\AgdaFunction{renameb}\AgdaSpace{}%
\AgdaSymbol{(}\AgdaFunction{extb}\AgdaSpace{}%
\AgdaBound{fb''}\AgdaSymbol{)}\AgdaSpace{}%
\AgdaSymbol{(}\AgdaFunction{renameb}\AgdaSpace{}%
\AgdaSymbol{(}\AgdaFunction{extb}\AgdaSpace{}%
\AgdaBound{fb'}\AgdaSymbol{)}\AgdaSpace{}%
\AgdaBound{t}\AgdaSymbol{))}\<%
\\
%
\>[2]\AgdaFunction{≡⟨}\AgdaSpace{}%
\AgdaFunction{cong}\AgdaSpace{}%
\AgdaInductiveConstructor{ƛb}\AgdaSpace{}%
\AgdaSymbol{(}\AgdaFunction{renameb-comp}\AgdaSpace{}%
\AgdaSymbol{(}\AgdaFunction{extb}\AgdaSpace{}%
\AgdaBound{fb'}\AgdaSymbol{)}\AgdaSpace{}%
\AgdaSymbol{(}\AgdaFunction{extb}\AgdaSpace{}%
\AgdaBound{fb''}\AgdaSymbol{)}\AgdaSpace{}%
\AgdaBound{t}\AgdaSymbol{)}\AgdaSpace{}%
\AgdaFunction{⟩}\<%
\\
\>[2][@{}l@{\AgdaIndent{0}}]%
\>[4]\AgdaInductiveConstructor{ƛb}\AgdaSpace{}%
\AgdaSymbol{(}\AgdaFunction{renameb}\AgdaSpace{}%
\AgdaSymbol{(λ}\AgdaSpace{}%
\AgdaBound{x}\AgdaSpace{}%
\AgdaSymbol{→}\AgdaSpace{}%
\AgdaFunction{extb}\AgdaSpace{}%
\AgdaBound{fb''}\AgdaSpace{}%
\AgdaSymbol{(}\AgdaFunction{extb}\AgdaSpace{}%
\AgdaBound{fb'}\AgdaSpace{}%
\AgdaBound{x}\AgdaSymbol{))}\AgdaSpace{}%
\AgdaBound{t}\AgdaSymbol{)}\<%
\\
%
\>[2]\AgdaFunction{≡⟨}\AgdaSpace{}%
\AgdaFunction{cong}\AgdaSpace{}%
\AgdaInductiveConstructor{ƛb}\AgdaSpace{}%
\AgdaSymbol{(}\AgdaFunction{renameb-cong}\AgdaSpace{}%
\AgdaSymbol{(}\AgdaFunction{extb-comp}\AgdaSpace{}%
\AgdaBound{fb'}\AgdaSpace{}%
\AgdaBound{fb''}\AgdaSymbol{)}\AgdaSpace{}%
\AgdaBound{t}\AgdaSymbol{)}\AgdaSpace{}%
\AgdaFunction{⟩}\<%
\\
\>[2][@{}l@{\AgdaIndent{0}}]%
\>[4]\AgdaInductiveConstructor{ƛb}\AgdaSpace{}%
\AgdaSymbol{(}\AgdaFunction{renameb}\AgdaSpace{}%
\AgdaSymbol{(}\AgdaFunction{extb}\AgdaSpace{}%
\AgdaSymbol{(λ}\AgdaSpace{}%
\AgdaBound{x}\AgdaSpace{}%
\AgdaSymbol{→}\AgdaSpace{}%
\AgdaBound{fb''}\AgdaSpace{}%
\AgdaSymbol{(}\AgdaBound{fb'}\AgdaSpace{}%
\AgdaBound{x}\AgdaSymbol{)))}\AgdaSpace{}%
\AgdaBound{t}\AgdaSymbol{)}\<%
\\
%
\>[2]\AgdaOperator{\AgdaFunction{∎eq}}\<%
\\
\>[0]\AgdaFunction{renameb-comp}\AgdaSpace{}%
\AgdaBound{fb'}\AgdaSpace{}%
\AgdaBound{fb''}\AgdaSpace{}%
\AgdaSymbol{(}\AgdaBound{t}\AgdaSpace{}%
\AgdaOperator{\AgdaInductiveConstructor{·b}}\AgdaSpace{}%
\AgdaBound{t₁}\AgdaSymbol{)}\AgdaSpace{}%
\AgdaSymbol{=}\AgdaSpace{}%
\AgdaFunction{cong₂}\AgdaSpace{}%
\AgdaOperator{\AgdaInductiveConstructor{\AgdaUnderscore{}·b\AgdaUnderscore{}}}\AgdaSpace{}%
\AgdaSymbol{(}\AgdaFunction{renameb-comp}\AgdaSpace{}%
\AgdaBound{fb'}\AgdaSpace{}%
\AgdaBound{fb''}\AgdaSpace{}%
\AgdaBound{t}\AgdaSymbol{)}\AgdaSpace{}%
\AgdaSymbol{(}\AgdaFunction{renameb-comp}\AgdaSpace{}%
\AgdaBound{fb'}\AgdaSpace{}%
\AgdaBound{fb''}\AgdaSpace{}%
\AgdaBound{t₁}\AgdaSymbol{)}\<%
\\
%
\\[\AgdaEmptyExtraSkip]%
\>[0]\AgdaComment{--\ extb\ and\ Sb\ commute\ under\ renameb.}\<%
\\
\>[0]\AgdaFunction{rename-extb-Sb-commute}\AgdaSpace{}%
\AgdaSymbol{:}\AgdaSpace{}%
\AgdaSymbol{∀}\AgdaSpace{}%
\AgdaSymbol{\{}\AgdaBound{Γ}\AgdaSpace{}%
\AgdaBound{A}\AgdaSpace{}%
\AgdaBound{B}\AgdaSpace{}%
\AgdaBound{Δ}\AgdaSymbol{\}}\<%
\\
\>[0][@{}l@{\AgdaIndent{0}}]%
\>[2]\AgdaSymbol{→}\AgdaSpace{}%
\AgdaSymbol{(}\AgdaBound{fb'}\AgdaSpace{}%
\AgdaSymbol{:}\AgdaSpace{}%
\AgdaSymbol{∀}\AgdaSpace{}%
\AgdaSymbol{\{}\AgdaBound{Ab}\AgdaSymbol{\}}\AgdaSpace{}%
\AgdaSymbol{→}\AgdaSpace{}%
\AgdaBound{Ab}\AgdaSpace{}%
\AgdaOperator{\AgdaDatatype{∈b}}\AgdaSpace{}%
\AgdaBound{Γ}\AgdaSpace{}%
\AgdaSymbol{→}\AgdaSpace{}%
\AgdaBound{Ab}\AgdaSpace{}%
\AgdaOperator{\AgdaDatatype{∈b}}\AgdaSpace{}%
\AgdaBound{Δ}\AgdaSymbol{)}\<%
\\
%
\>[2]\AgdaSymbol{→}\AgdaSpace{}%
\AgdaSymbol{(}\AgdaBound{t}\AgdaSpace{}%
\AgdaSymbol{:}\AgdaSpace{}%
\AgdaBound{Γ}\AgdaSpace{}%
\AgdaOperator{\AgdaDatatype{⊢b}}\AgdaSpace{}%
\AgdaBound{A}\AgdaSymbol{)}\<%
\\
%
\>[2]\AgdaSymbol{→}%
\>[579I]\AgdaFunction{renameb}\AgdaSpace{}%
\AgdaSymbol{(}\AgdaFunction{extb}\AgdaSpace{}%
\AgdaBound{fb'}\AgdaSpace{}%
\AgdaSymbol{\{}\AgdaArgument{B}\AgdaSpace{}%
\AgdaSymbol{=}\AgdaSpace{}%
\AgdaBound{B}\AgdaSymbol{\})}\AgdaSpace{}%
\AgdaSymbol{(}\AgdaFunction{renameb}\AgdaSpace{}%
\AgdaInductiveConstructor{Sb}\AgdaSpace{}%
\AgdaBound{t}\AgdaSymbol{)}\<%
\\
\>[.][@{}l@{}]\<[579I]%
\>[4]\AgdaOperator{\AgdaDatatype{≡}}\AgdaSpace{}%
\AgdaFunction{renameb}\AgdaSpace{}%
\AgdaInductiveConstructor{Sb}\AgdaSpace{}%
\AgdaSymbol{(}\AgdaFunction{renameb}\AgdaSpace{}%
\AgdaBound{fb'}\AgdaSpace{}%
\AgdaBound{t}\AgdaSymbol{)}\<%
\\
\>[0]\AgdaFunction{rename-extb-Sb-commute}\AgdaSpace{}%
\AgdaBound{fb'}\AgdaSpace{}%
\AgdaBound{t}\AgdaSpace{}%
\AgdaSymbol{=}\<%
\\
\>[0][@{}l@{\AgdaIndent{0}}]%
\>[2]\AgdaOperator{\AgdaFunction{beginEq}}\<%
\\
\>[2][@{}l@{\AgdaIndent{0}}]%
\>[4]\AgdaFunction{renameb}\AgdaSpace{}%
\AgdaSymbol{(}\AgdaFunction{extb}\AgdaSpace{}%
\AgdaBound{fb'}\AgdaSymbol{)}\AgdaSpace{}%
\AgdaSymbol{(}\AgdaFunction{renameb}\AgdaSpace{}%
\AgdaInductiveConstructor{Sb}\AgdaSpace{}%
\AgdaBound{t}\AgdaSymbol{)}\<%
\\
%
\>[2]\AgdaFunction{≡⟨}\AgdaSpace{}%
\AgdaFunction{renameb-comp}\AgdaSpace{}%
\AgdaInductiveConstructor{Sb}\AgdaSpace{}%
\AgdaSymbol{(}\AgdaFunction{extb}\AgdaSpace{}%
\AgdaBound{fb'}\AgdaSymbol{)}\AgdaSpace{}%
\AgdaBound{t}\AgdaSpace{}%
\AgdaFunction{⟩}\<%
\\
\>[2][@{}l@{\AgdaIndent{0}}]%
\>[4]\AgdaFunction{renameb}\AgdaSpace{}%
\AgdaSymbol{(λ}\AgdaSpace{}%
\AgdaBound{x}\AgdaSpace{}%
\AgdaSymbol{→}\AgdaSpace{}%
\AgdaFunction{extb}\AgdaSpace{}%
\AgdaBound{fb'}\AgdaSpace{}%
\AgdaSymbol{(}\AgdaInductiveConstructor{Sb}\AgdaSpace{}%
\AgdaBound{x}\AgdaSymbol{))}\AgdaSpace{}%
\AgdaBound{t}\<%
\\
%
\>[2]\AgdaFunction{≡⟨}\AgdaSpace{}%
\AgdaFunction{renameb-cong}\AgdaSpace{}%
\AgdaSymbol{(λ}\AgdaSpace{}%
\AgdaBound{x}\AgdaSpace{}%
\AgdaSymbol{→}\AgdaSpace{}%
\AgdaInductiveConstructor{refl}\AgdaSymbol{)}\AgdaSpace{}%
\AgdaBound{t}\AgdaSpace{}%
\AgdaFunction{⟩}\<%
\\
\>[2][@{}l@{\AgdaIndent{0}}]%
\>[4]\AgdaFunction{renameb}\AgdaSpace{}%
\AgdaSymbol{(λ}\AgdaSpace{}%
\AgdaBound{x}\AgdaSpace{}%
\AgdaSymbol{→}\AgdaSpace{}%
\AgdaInductiveConstructor{Sb}\AgdaSpace{}%
\AgdaSymbol{(}\AgdaBound{fb'}\AgdaSpace{}%
\AgdaBound{x}\AgdaSymbol{))}\AgdaSpace{}%
\AgdaBound{t}\<%
\\
%
\>[2]\AgdaFunction{≡⟨}\AgdaSpace{}%
\AgdaFunction{sym}\AgdaSpace{}%
\AgdaSymbol{(}\AgdaFunction{renameb-comp}\AgdaSpace{}%
\AgdaBound{fb'}\AgdaSpace{}%
\AgdaInductiveConstructor{Sb}\AgdaSpace{}%
\AgdaBound{t}\AgdaSymbol{)}\AgdaSpace{}%
\AgdaFunction{⟩}\<%
\\
\>[2][@{}l@{\AgdaIndent{0}}]%
\>[4]\AgdaFunction{renameb}\AgdaSpace{}%
\AgdaInductiveConstructor{Sb}\AgdaSpace{}%
\AgdaSymbol{(}\AgdaFunction{renameb}\AgdaSpace{}%
\AgdaBound{fb'}\AgdaSpace{}%
\AgdaBound{t}\AgdaSymbol{)}\<%
\\
%
\>[2]\AgdaOperator{\AgdaFunction{∎eq}}\<%
\\
%
\\[\AgdaEmptyExtraSkip]%
\>[0]\AgdaComment{--\ Preservation\ and\ Reflection\ of\ }\<%
\\
\>[0]\AgdaComment{--\ rename\ of\ Gödel's\ translation}\<%
\\
\>[0]\AgdaFunction{pres-rename}\AgdaSpace{}%
\AgdaSymbol{:}\AgdaSpace{}%
\AgdaSymbol{∀}\AgdaSpace{}%
\AgdaSymbol{\{}\AgdaBound{Γ}\AgdaSpace{}%
\AgdaBound{A}\AgdaSpace{}%
\AgdaBound{Δ}\AgdaSymbol{\}}\<%
\\
\>[0][@{}l@{\AgdaIndent{0}}]%
\>[2]\AgdaSymbol{→}\AgdaSpace{}%
\AgdaSymbol{(}\AgdaBound{f'}\AgdaSpace{}%
\AgdaSymbol{:}\AgdaSpace{}%
\AgdaSymbol{∀}\AgdaSpace{}%
\AgdaSymbol{\{}\AgdaBound{A}\AgdaSymbol{\}}\AgdaSpace{}%
\AgdaSymbol{→}\AgdaSpace{}%
\AgdaBound{A}\AgdaSpace{}%
\AgdaOperator{\AgdaDatatype{∈}}\AgdaSpace{}%
\AgdaBound{Γ}\AgdaSpace{}%
\AgdaSymbol{→}\AgdaSpace{}%
\AgdaBound{A}\AgdaSpace{}%
\AgdaOperator{\AgdaDatatype{∈}}\AgdaSpace{}%
\AgdaBound{Δ}\AgdaSymbol{)}\<%
\\
%
\>[2]\AgdaSymbol{→}\AgdaSpace{}%
\AgdaSymbol{(}\AgdaBound{fb'}\AgdaSpace{}%
\AgdaSymbol{:}\AgdaSpace{}%
\AgdaSymbol{∀}\AgdaSpace{}%
\AgdaSymbol{\{}\AgdaBound{Ab}\AgdaSymbol{\}}\AgdaSpace{}%
\AgdaSymbol{→}\AgdaSpace{}%
\AgdaBound{Ab}\AgdaSpace{}%
\AgdaOperator{\AgdaDatatype{∈b}}\AgdaSpace{}%
\AgdaFunction{embedContext}\AgdaSpace{}%
\AgdaBound{Γ}\AgdaSpace{}%
\AgdaSymbol{→}\AgdaSpace{}%
\AgdaBound{Ab}\AgdaSpace{}%
\AgdaOperator{\AgdaDatatype{∈b}}\AgdaSpace{}%
\AgdaFunction{embedContext}\AgdaSpace{}%
\AgdaBound{Δ}\AgdaSymbol{)}\<%
\\
%
\>[2]\AgdaSymbol{→}\AgdaSpace{}%
\AgdaSymbol{(\{}\AgdaBound{A'}\AgdaSpace{}%
\AgdaSymbol{:}\AgdaSpace{}%
\AgdaDatatype{Type}\AgdaSymbol{\}}\AgdaSpace{}%
\AgdaSymbol{(}\AgdaBound{x}\AgdaSpace{}%
\AgdaSymbol{:}\AgdaSpace{}%
\AgdaBound{A'}\AgdaSpace{}%
\AgdaOperator{\AgdaDatatype{∈}}\AgdaSpace{}%
\AgdaBound{Γ}\AgdaSymbol{)}\AgdaSpace{}%
\AgdaSymbol{→}\AgdaSpace{}%
\AgdaBound{fb'}\AgdaSpace{}%
\AgdaSymbol{(}\AgdaFunction{embedLookup}\AgdaSpace{}%
\AgdaBound{x}\AgdaSymbol{)}\AgdaSpace{}%
\AgdaOperator{\AgdaDatatype{≡}}\AgdaSpace{}%
\AgdaFunction{embedLookup}\AgdaSpace{}%
\AgdaSymbol{(}\AgdaBound{f'}\AgdaSpace{}%
\AgdaBound{x}\AgdaSymbol{))}\<%
\\
%
\>[2]\AgdaSymbol{→}\AgdaSpace{}%
\AgdaSymbol{(}\AgdaBound{t}\AgdaSpace{}%
\AgdaSymbol{:}\AgdaSpace{}%
\AgdaBound{Γ}\AgdaSpace{}%
\AgdaOperator{\AgdaDatatype{⊢}}\AgdaSpace{}%
\AgdaBound{A}\AgdaSymbol{)}\<%
\\
%
\>[2]\AgdaSymbol{→}\AgdaSpace{}%
\AgdaFunction{renameb}\AgdaSpace{}%
\AgdaBound{fb'}\AgdaSpace{}%
\AgdaSymbol{(}\AgdaFunction{embedTerm}\AgdaSpace{}%
\AgdaBound{t}\AgdaSymbol{)}\AgdaSpace{}%
\AgdaOperator{\AgdaDatatype{≡}}\AgdaSpace{}%
\AgdaFunction{embedTerm}\AgdaSpace{}%
\AgdaSymbol{(}\AgdaFunction{rename}\AgdaSpace{}%
\AgdaBound{f'}\AgdaSpace{}%
\AgdaBound{t}\AgdaSymbol{)}\<%
\\
\>[0]\AgdaFunction{pres-rename}\AgdaSpace{}%
\AgdaBound{f'}\AgdaSpace{}%
\AgdaBound{fb'}\AgdaSpace{}%
\AgdaBound{h}\AgdaSpace{}%
\AgdaSymbol{(}\AgdaInductiveConstructor{`}\AgdaSpace{}%
\AgdaBound{x}\AgdaSymbol{)}\AgdaSpace{}%
\AgdaSymbol{=}\AgdaSpace{}%
\AgdaFunction{cong}\AgdaSpace{}%
\AgdaInductiveConstructor{box}\AgdaSpace{}%
\AgdaSymbol{(}\AgdaFunction{cong}\AgdaSpace{}%
\AgdaInductiveConstructor{ε}\AgdaSpace{}%
\AgdaSymbol{(}\AgdaBound{h}\AgdaSpace{}%
\AgdaBound{x}\AgdaSymbol{))}\<%
\\
\>[0]\AgdaFunction{pres-rename}\AgdaSpace{}%
\AgdaBound{f'}\AgdaSpace{}%
\AgdaBound{fb'}\AgdaSpace{}%
\AgdaBound{h}\AgdaSpace{}%
\AgdaSymbol{(}\AgdaInductiveConstructor{ƛ}\AgdaSpace{}%
\AgdaBound{t}\AgdaSymbol{)}\AgdaSpace{}%
\AgdaSymbol{=}\AgdaSpace{}%
\AgdaFunction{cong}\AgdaSpace{}%
\AgdaInductiveConstructor{box}\AgdaSpace{}%
\AgdaSymbol{(}\AgdaFunction{cong}\AgdaSpace{}%
\AgdaInductiveConstructor{ƛb}\AgdaSpace{}%
\AgdaSymbol{(}\AgdaFunction{pres-rename}\AgdaSpace{}%
\AgdaSymbol{(}\AgdaFunction{calc.ext}\AgdaSpace{}%
\AgdaBound{f'}\AgdaSymbol{)}\AgdaSpace{}%
\AgdaSymbol{(}\AgdaFunction{boxCalc.extb}\AgdaSpace{}%
\AgdaBound{fb'}\AgdaSymbol{)}\AgdaSpace{}%
\AgdaSymbol{(}\AgdaFunction{pres-ext}\AgdaSpace{}%
\AgdaBound{f'}\AgdaSpace{}%
\AgdaBound{fb'}\AgdaSpace{}%
\AgdaBound{h}\AgdaSymbol{)}\AgdaSpace{}%
\AgdaBound{t}\AgdaSymbol{))}\<%
\\
\>[0]\AgdaFunction{pres-rename}\AgdaSpace{}%
\AgdaBound{f'}\AgdaSpace{}%
\AgdaBound{fb'}\AgdaSpace{}%
\AgdaBound{h}\AgdaSpace{}%
\AgdaSymbol{(}\AgdaBound{t}\AgdaSpace{}%
\AgdaOperator{\AgdaInductiveConstructor{·}}\AgdaSpace{}%
\AgdaBound{t₁}\AgdaSymbol{)}\AgdaSpace{}%
\AgdaSymbol{=}\<%
\\
\>[0][@{}l@{\AgdaIndent{0}}]%
\>[2]\AgdaOperator{\AgdaFunction{beginEq}}\<%
\\
\>[2][@{}l@{\AgdaIndent{0}}]%
\>[4]\AgdaInductiveConstructor{ƛb}\AgdaSpace{}%
\AgdaSymbol{(}\AgdaInductiveConstructor{ε}\AgdaSpace{}%
\AgdaInductiveConstructor{Zb}\AgdaSpace{}%
\AgdaOperator{\AgdaInductiveConstructor{·b}}\AgdaSpace{}%
\AgdaFunction{renameb}\AgdaSpace{}%
\AgdaSymbol{(}\AgdaFunction{extb}\AgdaSpace{}%
\AgdaBound{fb'}\AgdaSymbol{)}\AgdaSpace{}%
\AgdaSymbol{(}\AgdaFunction{renameb}\AgdaSpace{}%
\AgdaInductiveConstructor{Sb}\AgdaSpace{}%
\AgdaSymbol{(}\AgdaFunction{embedTerm}\AgdaSpace{}%
\AgdaBound{t₁}\AgdaSymbol{)))}\AgdaSpace{}%
\AgdaOperator{\AgdaInductiveConstructor{·b}}\AgdaSpace{}%
\AgdaFunction{renameb}\AgdaSpace{}%
\AgdaBound{fb'}\AgdaSpace{}%
\AgdaSymbol{(}\AgdaFunction{embedTerm}\AgdaSpace{}%
\AgdaBound{t}\AgdaSymbol{)}\<%
\\
%
\>[2]\AgdaFunction{≡⟨}\AgdaSpace{}%
\AgdaFunction{cong₂}\AgdaSpace{}%
\AgdaOperator{\AgdaInductiveConstructor{\AgdaUnderscore{}·b\AgdaUnderscore{}}}\AgdaSpace{}%
\AgdaInductiveConstructor{refl}\AgdaSpace{}%
\AgdaSymbol{(}\AgdaFunction{pres-rename}\AgdaSpace{}%
\AgdaBound{f'}\AgdaSpace{}%
\AgdaBound{fb'}\AgdaSpace{}%
\AgdaBound{h}\AgdaSpace{}%
\AgdaBound{t}\AgdaSymbol{)}\AgdaSpace{}%
\AgdaFunction{⟩}\<%
\\
\>[2][@{}l@{\AgdaIndent{0}}]%
\>[4]\AgdaInductiveConstructor{ƛb}\AgdaSpace{}%
\AgdaSymbol{(}\AgdaInductiveConstructor{ε}\AgdaSpace{}%
\AgdaInductiveConstructor{Zb}\AgdaSpace{}%
\AgdaOperator{\AgdaInductiveConstructor{·b}}\AgdaSpace{}%
\AgdaFunction{renameb}\AgdaSpace{}%
\AgdaSymbol{(}\AgdaFunction{extb}\AgdaSpace{}%
\AgdaBound{fb'}\AgdaSymbol{)}\AgdaSpace{}%
\AgdaSymbol{(}\AgdaFunction{renameb}\AgdaSpace{}%
\AgdaInductiveConstructor{Sb}\AgdaSpace{}%
\AgdaSymbol{(}\AgdaFunction{embedTerm}\AgdaSpace{}%
\AgdaBound{t₁}\AgdaSymbol{)))}\AgdaSpace{}%
\AgdaOperator{\AgdaInductiveConstructor{·b}}\AgdaSpace{}%
\AgdaFunction{embedTerm}\AgdaSpace{}%
\AgdaSymbol{(}\AgdaFunction{rename}\AgdaSpace{}%
\AgdaBound{f'}\AgdaSpace{}%
\AgdaBound{t}\AgdaSymbol{)}\<%
\\
%
\>[2]\AgdaFunction{≡⟨}\AgdaSpace{}%
\AgdaFunction{cong₂}\AgdaSpace{}%
\AgdaOperator{\AgdaInductiveConstructor{\AgdaUnderscore{}·b\AgdaUnderscore{}}}\AgdaSpace{}%
\AgdaSymbol{(}\AgdaFunction{cong}\AgdaSpace{}%
\AgdaInductiveConstructor{ƛb}\AgdaSpace{}%
\AgdaSymbol{(}\AgdaFunction{cong₂}\AgdaSpace{}%
\AgdaOperator{\AgdaInductiveConstructor{\AgdaUnderscore{}·b\AgdaUnderscore{}}}\AgdaSpace{}%
\AgdaInductiveConstructor{refl}\AgdaSpace{}%
\AgdaSymbol{(}\AgdaFunction{rename-extb-Sb-commute}\AgdaSpace{}%
\AgdaBound{fb'}\AgdaSpace{}%
\AgdaSymbol{(}\AgdaFunction{embedTerm}\AgdaSpace{}%
\AgdaBound{t₁}\AgdaSymbol{))))}\AgdaSpace{}%
\AgdaInductiveConstructor{refl}\AgdaSpace{}%
\AgdaFunction{⟩}\<%
\\
\>[2][@{}l@{\AgdaIndent{0}}]%
\>[4]\AgdaInductiveConstructor{ƛb}\AgdaSpace{}%
\AgdaSymbol{(}\AgdaInductiveConstructor{ε}\AgdaSpace{}%
\AgdaInductiveConstructor{Zb}\AgdaSpace{}%
\AgdaOperator{\AgdaInductiveConstructor{·b}}\AgdaSpace{}%
\AgdaFunction{renameb}\AgdaSpace{}%
\AgdaInductiveConstructor{Sb}\AgdaSpace{}%
\AgdaSymbol{(}\AgdaFunction{renameb}\AgdaSpace{}%
\AgdaBound{fb'}\AgdaSpace{}%
\AgdaSymbol{(}\AgdaFunction{embedTerm}\AgdaSpace{}%
\AgdaBound{t₁}\AgdaSymbol{)))}\AgdaSpace{}%
\AgdaOperator{\AgdaInductiveConstructor{·b}}\AgdaSpace{}%
\AgdaFunction{embedTerm}\AgdaSpace{}%
\AgdaSymbol{(}\AgdaFunction{rename}\AgdaSpace{}%
\AgdaBound{f'}\AgdaSpace{}%
\AgdaBound{t}\AgdaSymbol{)}\<%
\\
%
\>[2]\AgdaFunction{≡⟨}\AgdaSpace{}%
\AgdaFunction{cong₂}\AgdaSpace{}%
\AgdaOperator{\AgdaInductiveConstructor{\AgdaUnderscore{}·b\AgdaUnderscore{}}}\AgdaSpace{}%
\AgdaSymbol{(}\AgdaFunction{cong}\AgdaSpace{}%
\AgdaInductiveConstructor{ƛb}\AgdaSpace{}%
\AgdaSymbol{(}\AgdaFunction{cong₂}\AgdaSpace{}%
\AgdaOperator{\AgdaInductiveConstructor{\AgdaUnderscore{}·b\AgdaUnderscore{}}}\AgdaSpace{}%
\AgdaInductiveConstructor{refl}\AgdaSpace{}%
\AgdaSymbol{(}\AgdaFunction{cong}\AgdaSpace{}%
\AgdaSymbol{(}\AgdaFunction{renameb}\AgdaSpace{}%
\AgdaInductiveConstructor{Sb}\AgdaSymbol{)}\AgdaSpace{}%
\AgdaSymbol{(}\AgdaFunction{pres-rename}\AgdaSpace{}%
\AgdaBound{f'}\AgdaSpace{}%
\AgdaBound{fb'}\AgdaSpace{}%
\AgdaBound{h}\AgdaSpace{}%
\AgdaBound{t₁}\AgdaSymbol{))))}\AgdaSpace{}%
\AgdaInductiveConstructor{refl}\AgdaSpace{}%
\AgdaFunction{⟩}\<%
\\
\>[2][@{}l@{\AgdaIndent{0}}]%
\>[4]\AgdaInductiveConstructor{ƛb}\AgdaSpace{}%
\AgdaSymbol{(}\AgdaInductiveConstructor{ε}\AgdaSpace{}%
\AgdaInductiveConstructor{Zb}\AgdaSpace{}%
\AgdaOperator{\AgdaInductiveConstructor{·b}}\AgdaSpace{}%
\AgdaFunction{renameb}\AgdaSpace{}%
\AgdaInductiveConstructor{Sb}\AgdaSpace{}%
\AgdaSymbol{(}\AgdaFunction{embedTerm}\AgdaSpace{}%
\AgdaSymbol{(}\AgdaFunction{rename}\AgdaSpace{}%
\AgdaBound{f'}\AgdaSpace{}%
\AgdaBound{t₁}\AgdaSymbol{)))}\AgdaSpace{}%
\AgdaOperator{\AgdaInductiveConstructor{·b}}\AgdaSpace{}%
\AgdaFunction{embedTerm}\AgdaSpace{}%
\AgdaSymbol{(}\AgdaFunction{rename}\AgdaSpace{}%
\AgdaBound{f'}\AgdaSpace{}%
\AgdaBound{t}\AgdaSymbol{)}\<%
\\
%
\>[2]\AgdaOperator{\AgdaFunction{∎eq}}\<%
\\
%
\\[\AgdaEmptyExtraSkip]%
\>[0]\AgdaComment{--\ Preservation\ and\ Reflection\ of\ }\<%
\\
\>[0]\AgdaComment{--\ exts\ of\ Gödel's\ translation}\<%
\\
\>[0]\AgdaFunction{pres-exts}\AgdaSpace{}%
\AgdaSymbol{:}\AgdaSpace{}%
\AgdaSymbol{∀}\AgdaSpace{}%
\AgdaSymbol{\{}\AgdaBound{Γ}\AgdaSpace{}%
\AgdaBound{Δ}\AgdaSymbol{\}}\<%
\\
\>[0][@{}l@{\AgdaIndent{0}}]%
\>[2]\AgdaSymbol{→}\AgdaSpace{}%
\AgdaSymbol{(}\AgdaBound{fb'}\AgdaSpace{}%
\AgdaSymbol{:}\AgdaSpace{}%
\AgdaSymbol{∀}\AgdaSpace{}%
\AgdaSymbol{\{}\AgdaBound{Ab}\AgdaSymbol{\}}\AgdaSpace{}%
\AgdaSymbol{→}\AgdaSpace{}%
\AgdaInductiveConstructor{□}\AgdaSpace{}%
\AgdaBound{Ab}\AgdaSpace{}%
\AgdaOperator{\AgdaDatatype{∈b}}\AgdaSpace{}%
\AgdaFunction{embedContext}\AgdaSpace{}%
\AgdaBound{Γ}\AgdaSpace{}%
\AgdaSymbol{→}\AgdaSpace{}%
\AgdaFunction{embedContext}\AgdaSpace{}%
\AgdaBound{Δ}\AgdaSpace{}%
\AgdaOperator{\AgdaDatatype{⊢b}}\AgdaSpace{}%
\AgdaBound{Ab}\AgdaSymbol{)}\<%
\\
%
\>[2]\AgdaSymbol{→}\AgdaSpace{}%
\AgdaSymbol{(}\AgdaBound{f'}\AgdaSpace{}%
\AgdaSymbol{:}\AgdaSpace{}%
\AgdaSymbol{∀}\AgdaSpace{}%
\AgdaSymbol{\{}\AgdaBound{A}\AgdaSymbol{\}}\AgdaSpace{}%
\AgdaSymbol{→}\AgdaSpace{}%
\AgdaBound{A}\AgdaSpace{}%
\AgdaOperator{\AgdaDatatype{∈}}\AgdaSpace{}%
\AgdaBound{Γ}\AgdaSpace{}%
\AgdaSymbol{→}\AgdaSpace{}%
\AgdaBound{Δ}\AgdaSpace{}%
\AgdaOperator{\AgdaDatatype{⊢}}\AgdaSpace{}%
\AgdaBound{A}\AgdaSymbol{)}\<%
\\
%
\>[2]\AgdaSymbol{→}\AgdaSpace{}%
\AgdaSymbol{(\{}\AgdaBound{A'}\AgdaSpace{}%
\AgdaSymbol{:}\AgdaSpace{}%
\AgdaDatatype{Type}\AgdaSymbol{\}}\AgdaSpace{}%
\AgdaSymbol{(}\AgdaBound{x}\AgdaSpace{}%
\AgdaSymbol{:}\AgdaSpace{}%
\AgdaBound{A'}\AgdaSpace{}%
\AgdaOperator{\AgdaDatatype{∈}}\AgdaSpace{}%
\AgdaBound{Γ}\AgdaSymbol{)}\AgdaSpace{}%
\AgdaSymbol{→}\AgdaSpace{}%
\AgdaInductiveConstructor{box}\AgdaSpace{}%
\AgdaSymbol{(}\AgdaBound{fb'}\AgdaSpace{}%
\AgdaSymbol{(}\AgdaFunction{embedLookup}\AgdaSpace{}%
\AgdaBound{x}\AgdaSymbol{))}\AgdaSpace{}%
\AgdaOperator{\AgdaDatatype{≡}}\AgdaSpace{}%
\AgdaFunction{embedTerm}\AgdaSpace{}%
\AgdaSymbol{(}\AgdaBound{f'}\AgdaSpace{}%
\AgdaBound{x}\AgdaSymbol{))}\<%
\\
%
\>[2]\AgdaSymbol{→}\AgdaSpace{}%
\AgdaSymbol{\{}\AgdaBound{A'}\AgdaSpace{}%
\AgdaBound{A}\AgdaSpace{}%
\AgdaSymbol{:}\AgdaSpace{}%
\AgdaDatatype{Type}\AgdaSymbol{\}}\AgdaSpace{}%
\AgdaSymbol{(}\AgdaBound{x}\AgdaSpace{}%
\AgdaSymbol{:}\AgdaSpace{}%
\AgdaBound{A'}\AgdaSpace{}%
\AgdaOperator{\AgdaDatatype{∈}}\AgdaSpace{}%
\AgdaBound{Γ}\AgdaSpace{}%
\AgdaOperator{\AgdaInductiveConstructor{,}}\AgdaSpace{}%
\AgdaBound{A}\AgdaSymbol{)}\AgdaSpace{}%
\AgdaSymbol{→}\AgdaSpace{}%
\AgdaInductiveConstructor{box}\AgdaSpace{}%
\AgdaSymbol{(}\AgdaFunction{extsb}\AgdaSpace{}%
\AgdaBound{fb'}\AgdaSpace{}%
\AgdaSymbol{(}\AgdaFunction{embedLookup}\AgdaSpace{}%
\AgdaBound{x}\AgdaSymbol{))}\AgdaSpace{}%
\AgdaOperator{\AgdaDatatype{≡}}\AgdaSpace{}%
\AgdaFunction{embedTerm}\AgdaSpace{}%
\AgdaSymbol{(}\AgdaFunction{exts}\AgdaSpace{}%
\AgdaBound{f'}\AgdaSpace{}%
\AgdaBound{x}\AgdaSymbol{)}\<%
\\
\>[0]\AgdaFunction{pres-exts}\AgdaSpace{}%
\AgdaBound{fb'}\AgdaSpace{}%
\AgdaBound{f'}\AgdaSpace{}%
\AgdaBound{h}\AgdaSpace{}%
\AgdaInductiveConstructor{Z}\AgdaSpace{}%
\AgdaSymbol{=}\AgdaSpace{}%
\AgdaInductiveConstructor{refl}\<%
\\
\>[0]\AgdaFunction{pres-exts}\AgdaSpace{}%
\AgdaBound{fb'}\AgdaSpace{}%
\AgdaBound{f'}\AgdaSpace{}%
\AgdaBound{h}\AgdaSpace{}%
\AgdaSymbol{(}\AgdaInductiveConstructor{S}\AgdaSpace{}%
\AgdaBound{x}\AgdaSymbol{)}\AgdaSpace{}%
\AgdaSymbol{=}\<%
\\
\>[0][@{}l@{\AgdaIndent{0}}]%
\>[2]\AgdaOperator{\AgdaFunction{beginEq}}\<%
\\
\>[2][@{}l@{\AgdaIndent{0}}]%
\>[4]\AgdaInductiveConstructor{box}\AgdaSpace{}%
\AgdaSymbol{(}\AgdaFunction{renameb}\AgdaSpace{}%
\AgdaInductiveConstructor{Sb}\AgdaSpace{}%
\AgdaSymbol{(}\AgdaBound{fb'}\AgdaSpace{}%
\AgdaSymbol{(}\AgdaFunction{embedLookup}\AgdaSpace{}%
\AgdaBound{x}\AgdaSymbol{)))}\<%
\\
%
\>[2]\AgdaFunction{≡⟨}\AgdaSpace{}%
\AgdaInductiveConstructor{refl}\AgdaSpace{}%
\AgdaFunction{⟩}\<%
\\
\>[2][@{}l@{\AgdaIndent{0}}]%
\>[4]\AgdaFunction{renameb}\AgdaSpace{}%
\AgdaInductiveConstructor{Sb}\AgdaSpace{}%
\AgdaSymbol{(}\AgdaInductiveConstructor{box}\AgdaSpace{}%
\AgdaSymbol{(}\AgdaBound{fb'}\AgdaSpace{}%
\AgdaSymbol{(}\AgdaFunction{embedLookup}\AgdaSpace{}%
\AgdaBound{x}\AgdaSymbol{)))}\<%
\\
%
\>[2]\AgdaFunction{≡⟨}\AgdaSpace{}%
\AgdaFunction{cong}\AgdaSpace{}%
\AgdaSymbol{(}\AgdaFunction{renameb}\AgdaSpace{}%
\AgdaInductiveConstructor{Sb}\AgdaSymbol{)}\AgdaSpace{}%
\AgdaSymbol{(}\AgdaBound{h}\AgdaSpace{}%
\AgdaBound{x}\AgdaSymbol{)}\AgdaSpace{}%
\AgdaFunction{⟩}\<%
\\
\>[2][@{}l@{\AgdaIndent{0}}]%
\>[4]\AgdaFunction{renameb}\AgdaSpace{}%
\AgdaInductiveConstructor{Sb}\AgdaSpace{}%
\AgdaSymbol{(}\AgdaFunction{embedTerm}\AgdaSpace{}%
\AgdaSymbol{(}\AgdaBound{f'}\AgdaSpace{}%
\AgdaBound{x}\AgdaSymbol{))}\<%
\\
%
\>[2]\AgdaFunction{≡⟨}\AgdaSpace{}%
\AgdaFunction{pres-rename}\AgdaSpace{}%
\AgdaInductiveConstructor{S}\AgdaSpace{}%
\AgdaInductiveConstructor{Sb}\AgdaSpace{}%
\AgdaSymbol{(λ}\AgdaSpace{}%
\AgdaBound{x₁}\AgdaSpace{}%
\AgdaSymbol{→}\AgdaSpace{}%
\AgdaInductiveConstructor{refl}\AgdaSymbol{)}\AgdaSpace{}%
\AgdaSymbol{(}\AgdaBound{f'}\AgdaSpace{}%
\AgdaBound{x}\AgdaSymbol{)}\AgdaSpace{}%
\AgdaFunction{⟩}\<%
\\
\>[2][@{}l@{\AgdaIndent{0}}]%
\>[4]\AgdaFunction{embedTerm}\AgdaSpace{}%
\AgdaSymbol{(}\AgdaFunction{rename}\AgdaSpace{}%
\AgdaInductiveConstructor{S}\AgdaSpace{}%
\AgdaSymbol{(}\AgdaBound{f'}\AgdaSpace{}%
\AgdaBound{x}\AgdaSymbol{))}\<%
\\
%
\>[2]\AgdaOperator{\AgdaFunction{∎eq}}\<%
\\
%
\\[\AgdaEmptyExtraSkip]%
\>[0]\AgdaComment{--\ Extsb\ preserves\ pointwise\ equality}\<%
\\
\>[0]\AgdaComment{--\ of\ renaming\ functions}\<%
\\
\>[0]\AgdaFunction{extsb-cong}\AgdaSpace{}%
\AgdaSymbol{:}\AgdaSpace{}%
\AgdaSymbol{∀}\AgdaSpace{}%
\AgdaSymbol{\{}\AgdaBound{Γ₁}\AgdaSpace{}%
\AgdaBound{Γ₂}\AgdaSpace{}%
\AgdaBound{A}\AgdaSpace{}%
\AgdaBound{A'}\AgdaSymbol{\}}\<%
\\
\>[0][@{}l@{\AgdaIndent{0}}]%
\>[2]\AgdaSymbol{→}\AgdaSpace{}%
\AgdaSymbol{\{}\AgdaBound{fb'}\AgdaSpace{}%
\AgdaSymbol{:}\AgdaSpace{}%
\AgdaSymbol{∀}\AgdaSpace{}%
\AgdaSymbol{\{}\AgdaBound{Ab}\AgdaSymbol{\}}\AgdaSpace{}%
\AgdaSymbol{→}\AgdaSpace{}%
\AgdaInductiveConstructor{□}\AgdaSpace{}%
\AgdaBound{Ab}\AgdaSpace{}%
\AgdaOperator{\AgdaDatatype{∈b}}\AgdaSpace{}%
\AgdaBound{Γ₁}\AgdaSpace{}%
\AgdaSymbol{→}\AgdaSpace{}%
\AgdaBound{Γ₂}\AgdaSpace{}%
\AgdaOperator{\AgdaDatatype{⊢b}}\AgdaSpace{}%
\AgdaBound{Ab}\AgdaSymbol{\}}\<%
\\
%
\>[2]\AgdaSymbol{→}\AgdaSpace{}%
\AgdaSymbol{\{}\AgdaBound{fb''}\AgdaSpace{}%
\AgdaSymbol{:}\AgdaSpace{}%
\AgdaSymbol{∀}\AgdaSpace{}%
\AgdaSymbol{\{}\AgdaBound{Ab}\AgdaSymbol{\}}\AgdaSpace{}%
\AgdaSymbol{→}\AgdaSpace{}%
\AgdaInductiveConstructor{□}\AgdaSpace{}%
\AgdaBound{Ab}\AgdaSpace{}%
\AgdaOperator{\AgdaDatatype{∈b}}\AgdaSpace{}%
\AgdaBound{Γ₁}\AgdaSpace{}%
\AgdaSymbol{→}\AgdaSpace{}%
\AgdaBound{Γ₂}\AgdaSpace{}%
\AgdaOperator{\AgdaDatatype{⊢b}}\AgdaSpace{}%
\AgdaBound{Ab}\AgdaSymbol{\}}\<%
\\
%
\>[2]\AgdaSymbol{→}\AgdaSpace{}%
\AgdaSymbol{(}\AgdaBound{h}\AgdaSpace{}%
\AgdaSymbol{:}\AgdaSpace{}%
\AgdaSymbol{∀}\AgdaSpace{}%
\AgdaSymbol{\{}\AgdaBound{Ab}\AgdaSymbol{\}}\AgdaSpace{}%
\AgdaSymbol{→}\AgdaSpace{}%
\AgdaSymbol{(}\AgdaBound{x}\AgdaSpace{}%
\AgdaSymbol{:}\AgdaSpace{}%
\AgdaInductiveConstructor{□}\AgdaSpace{}%
\AgdaBound{Ab}\AgdaSpace{}%
\AgdaOperator{\AgdaDatatype{∈b}}\AgdaSpace{}%
\AgdaBound{Γ₁}\AgdaSymbol{)}\AgdaSpace{}%
\AgdaSymbol{→}\AgdaSpace{}%
\AgdaBound{fb'}\AgdaSpace{}%
\AgdaBound{x}\AgdaSpace{}%
\AgdaOperator{\AgdaDatatype{≡}}\AgdaSpace{}%
\AgdaBound{fb''}\AgdaSpace{}%
\AgdaBound{x}\AgdaSymbol{)}\<%
\\
%
\>[2]\AgdaSymbol{→}\AgdaSpace{}%
\AgdaSymbol{(}\AgdaBound{x}\AgdaSpace{}%
\AgdaSymbol{:}\AgdaSpace{}%
\AgdaInductiveConstructor{□}\AgdaSpace{}%
\AgdaBound{A}\AgdaSpace{}%
\AgdaOperator{\AgdaDatatype{∈b}}\AgdaSpace{}%
\AgdaBound{Γ₁}\AgdaSpace{}%
\AgdaOperator{\AgdaInductiveConstructor{,b}}\AgdaSpace{}%
\AgdaInductiveConstructor{□}\AgdaSpace{}%
\AgdaBound{A'}\AgdaSymbol{)}\<%
\\
%
\>[2]\AgdaSymbol{→}\AgdaSpace{}%
\AgdaFunction{extsb}\AgdaSpace{}%
\AgdaBound{fb'}\AgdaSpace{}%
\AgdaBound{x}\AgdaSpace{}%
\AgdaOperator{\AgdaDatatype{≡}}\AgdaSpace{}%
\AgdaFunction{extsb}\AgdaSpace{}%
\AgdaBound{fb''}\AgdaSpace{}%
\AgdaBound{x}\<%
\\
\>[0]\AgdaFunction{extsb-cong}\AgdaSpace{}%
\AgdaBound{h}\AgdaSpace{}%
\AgdaInductiveConstructor{Zb}\AgdaSpace{}%
\AgdaSymbol{=}\AgdaSpace{}%
\AgdaInductiveConstructor{refl}\<%
\\
\>[0]\AgdaFunction{extsb-cong}\AgdaSpace{}%
\AgdaBound{h}\AgdaSpace{}%
\AgdaSymbol{(}\AgdaInductiveConstructor{Sb}\AgdaSpace{}%
\AgdaBound{x}\AgdaSymbol{)}\AgdaSpace{}%
\AgdaSymbol{=}\AgdaSpace{}%
\AgdaFunction{cong}\AgdaSpace{}%
\AgdaSymbol{(}\AgdaFunction{renameb}\AgdaSpace{}%
\AgdaInductiveConstructor{Sb}\AgdaSymbol{)}\AgdaSpace{}%
\AgdaSymbol{(}\AgdaBound{h}\AgdaSpace{}%
\AgdaBound{x}\AgdaSymbol{)}\<%
\\
%
\\[\AgdaEmptyExtraSkip]%
\>[0]\AgdaComment{--\ substb\ preserves\ pointwise\ equality}\<%
\\
\>[0]\AgdaComment{--\ of\ renaming\ functions}\<%
\\
\>[0]\AgdaFunction{substb-cong}\AgdaSpace{}%
\AgdaSymbol{:}\AgdaSpace{}%
\AgdaSymbol{∀}\AgdaSpace{}%
\AgdaSymbol{\{}\AgdaBound{Γ₁}\AgdaSpace{}%
\AgdaBound{Γ₂}\AgdaSpace{}%
\AgdaBound{A}\AgdaSymbol{\}}\<%
\\
\>[0][@{}l@{\AgdaIndent{0}}]%
\>[2]\AgdaSymbol{→}\AgdaSpace{}%
\AgdaSymbol{\{}\AgdaBound{fb'}\AgdaSpace{}%
\AgdaSymbol{:}\AgdaSpace{}%
\AgdaSymbol{∀}\AgdaSpace{}%
\AgdaSymbol{\{}\AgdaBound{Ab}\AgdaSymbol{\}}\AgdaSpace{}%
\AgdaSymbol{→}\AgdaSpace{}%
\AgdaInductiveConstructor{□}\AgdaSpace{}%
\AgdaBound{Ab}\AgdaSpace{}%
\AgdaOperator{\AgdaDatatype{∈b}}\AgdaSpace{}%
\AgdaBound{Γ₁}\AgdaSpace{}%
\AgdaSymbol{→}\AgdaSpace{}%
\AgdaBound{Γ₂}\AgdaSpace{}%
\AgdaOperator{\AgdaDatatype{⊢b}}\AgdaSpace{}%
\AgdaBound{Ab}\AgdaSymbol{\}}\<%
\\
%
\>[2]\AgdaSymbol{→}\AgdaSpace{}%
\AgdaSymbol{\{}\AgdaBound{fb''}\AgdaSpace{}%
\AgdaSymbol{:}\AgdaSpace{}%
\AgdaSymbol{∀}\AgdaSpace{}%
\AgdaSymbol{\{}\AgdaBound{Ab}\AgdaSymbol{\}}\AgdaSpace{}%
\AgdaSymbol{→}\AgdaSpace{}%
\AgdaInductiveConstructor{□}\AgdaSpace{}%
\AgdaBound{Ab}\AgdaSpace{}%
\AgdaOperator{\AgdaDatatype{∈b}}\AgdaSpace{}%
\AgdaBound{Γ₁}\AgdaSpace{}%
\AgdaSymbol{→}\AgdaSpace{}%
\AgdaBound{Γ₂}\AgdaSpace{}%
\AgdaOperator{\AgdaDatatype{⊢b}}\AgdaSpace{}%
\AgdaBound{Ab}\AgdaSymbol{\}}\<%
\\
%
\>[2]\AgdaSymbol{→}\AgdaSpace{}%
\AgdaSymbol{(}\AgdaBound{h}\AgdaSpace{}%
\AgdaSymbol{:}\AgdaSpace{}%
\AgdaSymbol{∀}\AgdaSpace{}%
\AgdaSymbol{\{}\AgdaBound{Ab}\AgdaSymbol{\}}\AgdaSpace{}%
\AgdaSymbol{→}\AgdaSpace{}%
\AgdaSymbol{(}\AgdaBound{x}\AgdaSpace{}%
\AgdaSymbol{:}\AgdaSpace{}%
\AgdaInductiveConstructor{□}\AgdaSpace{}%
\AgdaBound{Ab}\AgdaSpace{}%
\AgdaOperator{\AgdaDatatype{∈b}}\AgdaSpace{}%
\AgdaBound{Γ₁}\AgdaSymbol{)}\AgdaSpace{}%
\AgdaSymbol{→}\AgdaSpace{}%
\AgdaBound{fb'}\AgdaSpace{}%
\AgdaBound{x}\AgdaSpace{}%
\AgdaOperator{\AgdaDatatype{≡}}\AgdaSpace{}%
\AgdaBound{fb''}\AgdaSpace{}%
\AgdaBound{x}\AgdaSymbol{)}\<%
\\
%
\>[2]\AgdaSymbol{→}\AgdaSpace{}%
\AgdaSymbol{(}\AgdaBound{t}\AgdaSpace{}%
\AgdaSymbol{:}\AgdaSpace{}%
\AgdaBound{Γ₁}\AgdaSpace{}%
\AgdaOperator{\AgdaDatatype{⊢b}}\AgdaSpace{}%
\AgdaBound{A}\AgdaSymbol{)}\<%
\\
%
\>[2]\AgdaSymbol{→}\AgdaSpace{}%
\AgdaFunction{substb}\AgdaSpace{}%
\AgdaBound{fb'}\AgdaSpace{}%
\AgdaBound{t}\AgdaSpace{}%
\AgdaOperator{\AgdaDatatype{≡}}\AgdaSpace{}%
\AgdaFunction{substb}\AgdaSpace{}%
\AgdaBound{fb''}\AgdaSpace{}%
\AgdaBound{t}\<%
\\
\>[0]\AgdaFunction{substb-cong}\AgdaSpace{}%
\AgdaBound{h}\AgdaSpace{}%
\AgdaSymbol{(}\AgdaInductiveConstructor{ε}\AgdaSpace{}%
\AgdaBound{x}\AgdaSymbol{)}\AgdaSpace{}%
\AgdaSymbol{=}\AgdaSpace{}%
\AgdaBound{h}\AgdaSpace{}%
\AgdaBound{x}\<%
\\
\>[0]\AgdaFunction{substb-cong}\AgdaSpace{}%
\AgdaBound{h}\AgdaSpace{}%
\AgdaSymbol{(}\AgdaInductiveConstructor{box}\AgdaSpace{}%
\AgdaBound{t}\AgdaSymbol{)}\AgdaSpace{}%
\AgdaSymbol{=}\AgdaSpace{}%
\AgdaFunction{cong}\AgdaSpace{}%
\AgdaInductiveConstructor{box}\AgdaSpace{}%
\AgdaSymbol{(}\AgdaFunction{substb-cong}\AgdaSpace{}%
\AgdaBound{h}\AgdaSpace{}%
\AgdaBound{t}\AgdaSymbol{)}\<%
\\
\>[0]\AgdaFunction{substb-cong}\AgdaSpace{}%
\AgdaBound{h}\AgdaSpace{}%
\AgdaSymbol{(}\AgdaInductiveConstructor{ƛb}\AgdaSpace{}%
\AgdaBound{t}\AgdaSymbol{)}\AgdaSpace{}%
\AgdaSymbol{=}\AgdaSpace{}%
\AgdaFunction{cong}\AgdaSpace{}%
\AgdaInductiveConstructor{ƛb}\AgdaSpace{}%
\AgdaSymbol{(}\AgdaFunction{substb-cong}\AgdaSpace{}%
\AgdaSymbol{(}\AgdaFunction{extsb-cong}\AgdaSpace{}%
\AgdaBound{h}\AgdaSymbol{)}\AgdaSpace{}%
\AgdaBound{t}\AgdaSymbol{)}\<%
\\
\>[0]\AgdaFunction{substb-cong}\AgdaSpace{}%
\AgdaBound{h}\AgdaSpace{}%
\AgdaSymbol{(}\AgdaBound{t}\AgdaSpace{}%
\AgdaOperator{\AgdaInductiveConstructor{·b}}\AgdaSpace{}%
\AgdaBound{t₁}\AgdaSymbol{)}\AgdaSpace{}%
\AgdaSymbol{=}\AgdaSpace{}%
\AgdaFunction{cong₂}\AgdaSpace{}%
\AgdaOperator{\AgdaInductiveConstructor{\AgdaUnderscore{}·b\AgdaUnderscore{}}}\AgdaSpace{}%
\AgdaSymbol{(}\AgdaFunction{substb-cong}\AgdaSpace{}%
\AgdaBound{h}\AgdaSpace{}%
\AgdaBound{t}\AgdaSymbol{)}\AgdaSpace{}%
\AgdaSymbol{(}\AgdaFunction{substb-cong}\AgdaSpace{}%
\AgdaBound{h}\AgdaSpace{}%
\AgdaBound{t₁}\AgdaSymbol{)}\<%
\\
%
\\[\AgdaEmptyExtraSkip]%
\>[0]\AgdaComment{--\ Commutativity\ of\ renameb\ and\ extsb:}\<%
\\
\>[0]\AgdaFunction{extsb-rename-commute}\AgdaSpace{}%
\AgdaSymbol{:}\AgdaSpace{}%
\AgdaSymbol{∀}\AgdaSpace{}%
\AgdaSymbol{\{}\AgdaBound{Γ₁}\AgdaSpace{}%
\AgdaBound{Γ₂}\AgdaSpace{}%
\AgdaBound{Γ₃}\AgdaSpace{}%
\AgdaBound{A}\AgdaSpace{}%
\AgdaBound{B}\AgdaSymbol{\}}\<%
\\
\>[0][@{}l@{\AgdaIndent{0}}]%
\>[2]\AgdaSymbol{→}\AgdaSpace{}%
\AgdaSymbol{(}\AgdaBound{fb'}\AgdaSpace{}%
\AgdaSymbol{:}\AgdaSpace{}%
\AgdaSymbol{∀}\AgdaSpace{}%
\AgdaSymbol{\{}\AgdaBound{Ab}\AgdaSymbol{\}}\AgdaSpace{}%
\AgdaSymbol{→}\AgdaSpace{}%
\AgdaBound{Ab}\AgdaSpace{}%
\AgdaOperator{\AgdaDatatype{∈b}}\AgdaSpace{}%
\AgdaBound{Γ₂}\AgdaSpace{}%
\AgdaSymbol{→}\AgdaSpace{}%
\AgdaBound{Ab}\AgdaSpace{}%
\AgdaOperator{\AgdaDatatype{∈b}}\AgdaSpace{}%
\AgdaBound{Γ₃}\AgdaSymbol{)}\<%
\\
%
\>[2]\AgdaSymbol{→}\AgdaSpace{}%
\AgdaSymbol{(}\AgdaBound{fb''}\AgdaSpace{}%
\AgdaSymbol{:}\AgdaSpace{}%
\AgdaSymbol{∀}\AgdaSpace{}%
\AgdaSymbol{\{}\AgdaBound{Ab}\AgdaSymbol{\}}\AgdaSpace{}%
\AgdaSymbol{→}\AgdaSpace{}%
\AgdaInductiveConstructor{□}\AgdaSpace{}%
\AgdaBound{Ab}\AgdaSpace{}%
\AgdaOperator{\AgdaDatatype{∈b}}\AgdaSpace{}%
\AgdaBound{Γ₁}\AgdaSpace{}%
\AgdaSymbol{→}\AgdaSpace{}%
\AgdaBound{Γ₂}\AgdaSpace{}%
\AgdaOperator{\AgdaDatatype{⊢b}}\AgdaSpace{}%
\AgdaBound{Ab}\AgdaSymbol{)}\<%
\\
%
\>[2]\AgdaSymbol{→}\AgdaSpace{}%
\AgdaSymbol{(}\AgdaBound{x}\AgdaSpace{}%
\AgdaSymbol{:}\AgdaSpace{}%
\AgdaInductiveConstructor{□}\AgdaSpace{}%
\AgdaBound{A}\AgdaSpace{}%
\AgdaOperator{\AgdaDatatype{∈b}}\AgdaSpace{}%
\AgdaBound{Γ₁}\AgdaSpace{}%
\AgdaOperator{\AgdaInductiveConstructor{,b}}\AgdaSpace{}%
\AgdaBound{B}\AgdaSymbol{)}\<%
\\
%
\>[2]\AgdaSymbol{→}%
\>[1170I]\AgdaFunction{extsb}\AgdaSpace{}%
\AgdaSymbol{(λ}\AgdaSpace{}%
\AgdaBound{x₁}\AgdaSpace{}%
\AgdaSymbol{→}\AgdaSpace{}%
\AgdaFunction{renameb}\AgdaSpace{}%
\AgdaBound{fb'}\AgdaSpace{}%
\AgdaSymbol{(}\AgdaBound{fb''}\AgdaSpace{}%
\AgdaBound{x₁}\AgdaSymbol{))}\AgdaSpace{}%
\AgdaBound{x}\AgdaSpace{}%
\AgdaOperator{\AgdaDatatype{≡}}\<%
\\
\>[.][@{}l@{}]\<[1170I]%
\>[4]\AgdaFunction{renameb}\AgdaSpace{}%
\AgdaSymbol{(}\AgdaFunction{extb}\AgdaSpace{}%
\AgdaBound{fb'}\AgdaSymbol{)}\AgdaSpace{}%
\AgdaSymbol{(}\AgdaFunction{extsb}\AgdaSpace{}%
\AgdaBound{fb''}\AgdaSpace{}%
\AgdaBound{x}\AgdaSymbol{)}\<%
\\
\>[0]\AgdaFunction{extsb-rename-commute}\AgdaSpace{}%
\AgdaBound{fb'}\AgdaSpace{}%
\AgdaBound{fb''}\AgdaSpace{}%
\AgdaInductiveConstructor{Zb}\AgdaSpace{}%
\AgdaSymbol{=}\AgdaSpace{}%
\AgdaInductiveConstructor{refl}\<%
\\
\>[0]\AgdaFunction{extsb-rename-commute}\AgdaSpace{}%
\AgdaBound{fb'}\AgdaSpace{}%
\AgdaBound{fb''}\AgdaSpace{}%
\AgdaSymbol{(}\AgdaInductiveConstructor{Sb}\AgdaSpace{}%
\AgdaBound{x}\AgdaSymbol{)}\AgdaSpace{}%
\AgdaSymbol{=}\AgdaSpace{}%
\AgdaFunction{sym}\AgdaSpace{}%
\AgdaSymbol{(}\AgdaFunction{rename-extb-Sb-commute}\AgdaSpace{}%
\AgdaBound{fb'}\AgdaSpace{}%
\AgdaSymbol{(}\AgdaBound{fb''}\AgdaSpace{}%
\AgdaBound{x}\AgdaSymbol{))}\<%
\\
%
\\[\AgdaEmptyExtraSkip]%
\>[0]\AgdaComment{--\ Commutativity\ of\ substb\ and\ renameb}\<%
\\
\>[0]\AgdaFunction{substb-rename-commute}\AgdaSpace{}%
\AgdaSymbol{:}\AgdaSpace{}%
\AgdaSymbol{∀}\AgdaSpace{}%
\AgdaSymbol{\{}\AgdaBound{Γ₁}\AgdaSpace{}%
\AgdaBound{Γ₂}\AgdaSpace{}%
\AgdaBound{Γ₃}\AgdaSpace{}%
\AgdaBound{A}\AgdaSymbol{\}}\<%
\\
\>[0][@{}l@{\AgdaIndent{0}}]%
\>[2]\AgdaSymbol{→}\AgdaSpace{}%
\AgdaSymbol{(}\AgdaBound{fb'}\AgdaSpace{}%
\AgdaSymbol{:}\AgdaSpace{}%
\AgdaSymbol{∀}\AgdaSpace{}%
\AgdaSymbol{\{}\AgdaBound{Ab}\AgdaSymbol{\}}\AgdaSpace{}%
\AgdaSymbol{→}\AgdaSpace{}%
\AgdaBound{Ab}\AgdaSpace{}%
\AgdaOperator{\AgdaDatatype{∈b}}\AgdaSpace{}%
\AgdaBound{Γ₂}\AgdaSpace{}%
\AgdaSymbol{→}\AgdaSpace{}%
\AgdaBound{Ab}\AgdaSpace{}%
\AgdaOperator{\AgdaDatatype{∈b}}\AgdaSpace{}%
\AgdaBound{Γ₃}\AgdaSymbol{)}\<%
\\
%
\>[2]\AgdaSymbol{→}\AgdaSpace{}%
\AgdaSymbol{(}\AgdaBound{fb''}\AgdaSpace{}%
\AgdaSymbol{:}\AgdaSpace{}%
\AgdaSymbol{∀}\AgdaSpace{}%
\AgdaSymbol{\{}\AgdaBound{Ab}\AgdaSymbol{\}}\AgdaSpace{}%
\AgdaSymbol{→}\AgdaSpace{}%
\AgdaInductiveConstructor{□}\AgdaSpace{}%
\AgdaBound{Ab}\AgdaSpace{}%
\AgdaOperator{\AgdaDatatype{∈b}}\AgdaSpace{}%
\AgdaBound{Γ₁}\AgdaSpace{}%
\AgdaSymbol{→}\AgdaSpace{}%
\AgdaBound{Γ₂}\AgdaSpace{}%
\AgdaOperator{\AgdaDatatype{⊢b}}\AgdaSpace{}%
\AgdaBound{Ab}\AgdaSymbol{)}\<%
\\
%
\>[2]\AgdaSymbol{→}\AgdaSpace{}%
\AgdaSymbol{(}\AgdaBound{t}\AgdaSpace{}%
\AgdaSymbol{:}\AgdaSpace{}%
\AgdaBound{Γ₁}\AgdaSpace{}%
\AgdaOperator{\AgdaDatatype{⊢b}}\AgdaSpace{}%
\AgdaBound{A}\AgdaSymbol{)}\<%
\\
%
\>[2]\AgdaSymbol{→}\AgdaSpace{}%
\AgdaFunction{substb}\AgdaSpace{}%
\AgdaSymbol{(λ}\AgdaSpace{}%
\AgdaBound{x}\AgdaSpace{}%
\AgdaSymbol{→}\AgdaSpace{}%
\AgdaFunction{renameb}\AgdaSpace{}%
\AgdaBound{fb'}\AgdaSpace{}%
\AgdaSymbol{(}\AgdaBound{fb''}\AgdaSpace{}%
\AgdaBound{x}\AgdaSymbol{))}\AgdaSpace{}%
\AgdaBound{t}\AgdaSpace{}%
\AgdaOperator{\AgdaDatatype{≡}}\AgdaSpace{}%
\AgdaFunction{renameb}\AgdaSpace{}%
\AgdaBound{fb'}\AgdaSpace{}%
\AgdaSymbol{(}\AgdaFunction{substb}\AgdaSpace{}%
\AgdaBound{fb''}\AgdaSpace{}%
\AgdaBound{t}\AgdaSymbol{)}\<%
\\
\>[0]\AgdaFunction{substb-rename-commute}\AgdaSpace{}%
\AgdaBound{fb'}\AgdaSpace{}%
\AgdaBound{fb''}\AgdaSpace{}%
\AgdaSymbol{(}\AgdaInductiveConstructor{ε}\AgdaSpace{}%
\AgdaBound{x}\AgdaSymbol{)}\AgdaSpace{}%
\AgdaSymbol{=}\AgdaSpace{}%
\AgdaInductiveConstructor{refl}\<%
\\
\>[0]\AgdaFunction{substb-rename-commute}\AgdaSpace{}%
\AgdaBound{fb'}\AgdaSpace{}%
\AgdaBound{fb''}\AgdaSpace{}%
\AgdaSymbol{(}\AgdaInductiveConstructor{box}\AgdaSpace{}%
\AgdaBound{t}\AgdaSymbol{)}\AgdaSpace{}%
\AgdaSymbol{=}\<%
\\
\>[0][@{}l@{\AgdaIndent{0}}]%
\>[2]\AgdaFunction{cong}\<%
\\
\>[2][@{}l@{\AgdaIndent{0}}]%
\>[4]\AgdaInductiveConstructor{box}\<%
\\
%
\>[4]\AgdaSymbol{(}\AgdaFunction{substb-rename-commute}\AgdaSpace{}%
\AgdaBound{fb'}\AgdaSpace{}%
\AgdaBound{fb''}\AgdaSpace{}%
\AgdaBound{t}\AgdaSymbol{)}\<%
\\
\>[0]\AgdaFunction{substb-rename-commute}\AgdaSpace{}%
\AgdaBound{fb'}\AgdaSpace{}%
\AgdaBound{fb''}\AgdaSpace{}%
\AgdaSymbol{(}\AgdaInductiveConstructor{ƛb}\AgdaSpace{}%
\AgdaBound{t}\AgdaSymbol{)}\AgdaSpace{}%
\AgdaSymbol{=}\<%
\\
\>[0][@{}l@{\AgdaIndent{0}}]%
\>[2]\AgdaOperator{\AgdaFunction{beginEq}}\<%
\\
\>[2][@{}l@{\AgdaIndent{0}}]%
\>[4]\AgdaInductiveConstructor{ƛb}\AgdaSpace{}%
\AgdaSymbol{(}\AgdaFunction{substb}\AgdaSpace{}%
\AgdaSymbol{(}\AgdaFunction{extsb}\AgdaSpace{}%
\AgdaSymbol{(λ}\AgdaSpace{}%
\AgdaBound{x}\AgdaSpace{}%
\AgdaSymbol{→}\AgdaSpace{}%
\AgdaFunction{renameb}\AgdaSpace{}%
\AgdaBound{fb'}\AgdaSpace{}%
\AgdaSymbol{(}\AgdaBound{fb''}\AgdaSpace{}%
\AgdaBound{x}\AgdaSymbol{)))}\AgdaSpace{}%
\AgdaBound{t}\AgdaSymbol{)}\<%
\\
%
\>[2]\AgdaFunction{≡⟨}\AgdaSpace{}%
\AgdaFunction{cong}\AgdaSpace{}%
\AgdaInductiveConstructor{ƛb}\AgdaSpace{}%
\AgdaSymbol{(}\AgdaFunction{substb-cong}\AgdaSpace{}%
\AgdaSymbol{(λ}\AgdaSpace{}%
\AgdaBound{x}\AgdaSpace{}%
\AgdaSymbol{→}\AgdaSpace{}%
\AgdaFunction{extsb-rename-commute}\AgdaSpace{}%
\AgdaBound{fb'}\AgdaSpace{}%
\AgdaBound{fb''}\AgdaSpace{}%
\AgdaBound{x}\AgdaSymbol{)}\AgdaSpace{}%
\AgdaBound{t}\AgdaSymbol{)}\AgdaSpace{}%
\AgdaFunction{⟩}\<%
\\
\>[2][@{}l@{\AgdaIndent{0}}]%
\>[4]\AgdaInductiveConstructor{ƛb}\AgdaSpace{}%
\AgdaSymbol{(}\AgdaFunction{substb}\AgdaSpace{}%
\AgdaSymbol{(λ}\AgdaSpace{}%
\AgdaBound{x}\AgdaSpace{}%
\AgdaSymbol{→}\AgdaSpace{}%
\AgdaFunction{renameb}\AgdaSpace{}%
\AgdaSymbol{(}\AgdaFunction{extb}\AgdaSpace{}%
\AgdaBound{fb'}\AgdaSymbol{)}\AgdaSpace{}%
\AgdaSymbol{((}\AgdaFunction{extsb}\AgdaSpace{}%
\AgdaBound{fb''}\AgdaSymbol{)}\AgdaSpace{}%
\AgdaBound{x}\AgdaSymbol{))}\AgdaSpace{}%
\AgdaBound{t}\AgdaSymbol{)}\<%
\\
%
\>[2]\AgdaFunction{≡⟨}\AgdaSpace{}%
\AgdaFunction{cong}\AgdaSpace{}%
\AgdaInductiveConstructor{ƛb}\AgdaSpace{}%
\AgdaSymbol{(}\AgdaFunction{substb-rename-commute}\AgdaSpace{}%
\AgdaSymbol{(}\AgdaFunction{extb}\AgdaSpace{}%
\AgdaBound{fb'}\AgdaSymbol{)}\AgdaSpace{}%
\AgdaSymbol{(}\AgdaFunction{extsb}\AgdaSpace{}%
\AgdaBound{fb''}\AgdaSymbol{)}\AgdaSpace{}%
\AgdaBound{t}\AgdaSymbol{)}\AgdaSpace{}%
\AgdaFunction{⟩}\<%
\\
\>[2][@{}l@{\AgdaIndent{0}}]%
\>[4]\AgdaInductiveConstructor{ƛb}\AgdaSpace{}%
\AgdaSymbol{(}\AgdaFunction{renameb}\AgdaSpace{}%
\AgdaSymbol{(}\AgdaFunction{extb}\AgdaSpace{}%
\AgdaBound{fb'}\AgdaSymbol{)}\AgdaSpace{}%
\AgdaSymbol{(}\AgdaFunction{substb}\AgdaSpace{}%
\AgdaSymbol{(}\AgdaFunction{extsb}\AgdaSpace{}%
\AgdaBound{fb''}\AgdaSymbol{)}\AgdaSpace{}%
\AgdaBound{t}\AgdaSymbol{))}\<%
\\
%
\>[2]\AgdaOperator{\AgdaFunction{∎eq}}\<%
\\
\>[0]\AgdaFunction{substb-rename-commute}\AgdaSpace{}%
\AgdaBound{fb'}\AgdaSpace{}%
\AgdaBound{fb''}\AgdaSpace{}%
\AgdaSymbol{(}\AgdaBound{t}\AgdaSpace{}%
\AgdaOperator{\AgdaInductiveConstructor{·b}}\AgdaSpace{}%
\AgdaBound{t₁}\AgdaSymbol{)}\AgdaSpace{}%
\AgdaSymbol{=}\<%
\\
\>[0][@{}l@{\AgdaIndent{0}}]%
\>[2]\AgdaFunction{cong₂}\<%
\\
\>[2][@{}l@{\AgdaIndent{0}}]%
\>[4]\AgdaOperator{\AgdaInductiveConstructor{\AgdaUnderscore{}·b\AgdaUnderscore{}}}\<%
\\
%
\>[4]\AgdaSymbol{(}\AgdaFunction{substb-rename-commute}\AgdaSpace{}%
\AgdaBound{fb'}\AgdaSpace{}%
\AgdaBound{fb''}\AgdaSpace{}%
\AgdaBound{t}\AgdaSymbol{)}\<%
\\
%
\>[4]\AgdaSymbol{(}\AgdaFunction{substb-rename-commute}\AgdaSpace{}%
\AgdaBound{fb'}\AgdaSpace{}%
\AgdaBound{fb''}\AgdaSpace{}%
\AgdaBound{t₁}\AgdaSymbol{)}\<%
\\
%
\\[\AgdaEmptyExtraSkip]%
\>[0]\AgdaComment{--\ extsb\ distributes\ over\ composition:}\<%
\\
\>[0]\AgdaComment{--\ Executing\ extsb\ first\ with\ fb',\ then\ with\ fb''}\<%
\\
\>[0]\AgdaComment{--\ is\ the\ same\ as\ doing\ extsb\ once\ with\ fb''\ ∘\ fb'.}\<%
\\
\>[0]\AgdaFunction{extsb-comp}\AgdaSpace{}%
\AgdaSymbol{:}\AgdaSpace{}%
\AgdaSymbol{∀}\AgdaSpace{}%
\AgdaSymbol{\{}\AgdaBound{Γ₁}\AgdaSpace{}%
\AgdaBound{Γ₂}\AgdaSpace{}%
\AgdaBound{Γ₃}\AgdaSpace{}%
\AgdaBound{A}\AgdaSpace{}%
\AgdaBound{B}\AgdaSymbol{\}}\<%
\\
\>[0][@{}l@{\AgdaIndent{0}}]%
\>[2]\AgdaSymbol{→}\AgdaSpace{}%
\AgdaSymbol{(}\AgdaBound{fb'}\AgdaSpace{}%
\AgdaSymbol{:}\AgdaSpace{}%
\AgdaSymbol{∀}\AgdaSpace{}%
\AgdaSymbol{\{}\AgdaBound{Ab}\AgdaSymbol{\}}\AgdaSpace{}%
\AgdaSymbol{→}\AgdaSpace{}%
\AgdaBound{Ab}\AgdaSpace{}%
\AgdaOperator{\AgdaDatatype{∈b}}\AgdaSpace{}%
\AgdaBound{Γ₁}\AgdaSpace{}%
\AgdaSymbol{→}\AgdaSpace{}%
\AgdaBound{Ab}\AgdaSpace{}%
\AgdaOperator{\AgdaDatatype{∈b}}\AgdaSpace{}%
\AgdaBound{Γ₂}\AgdaSymbol{)}\<%
\\
%
\>[2]\AgdaSymbol{→}\AgdaSpace{}%
\AgdaSymbol{(}\AgdaBound{fb''}\AgdaSpace{}%
\AgdaSymbol{:}\AgdaSpace{}%
\AgdaSymbol{∀}\AgdaSpace{}%
\AgdaSymbol{\{}\AgdaBound{Ab}\AgdaSymbol{\}}\AgdaSpace{}%
\AgdaSymbol{→}\AgdaSpace{}%
\AgdaInductiveConstructor{□}\AgdaSpace{}%
\AgdaBound{Ab}\AgdaSpace{}%
\AgdaOperator{\AgdaDatatype{∈b}}\AgdaSpace{}%
\AgdaBound{Γ₂}\AgdaSpace{}%
\AgdaSymbol{→}\AgdaSpace{}%
\AgdaBound{Γ₃}\AgdaSpace{}%
\AgdaOperator{\AgdaDatatype{⊢b}}\AgdaSpace{}%
\AgdaBound{Ab}\AgdaSymbol{)}\<%
\\
%
\>[2]\AgdaSymbol{→}\AgdaSpace{}%
\AgdaSymbol{(}\AgdaBound{x}\AgdaSpace{}%
\AgdaSymbol{:}\AgdaSpace{}%
\AgdaInductiveConstructor{□}\AgdaSpace{}%
\AgdaBound{A}\AgdaSpace{}%
\AgdaOperator{\AgdaDatatype{∈b}}\AgdaSpace{}%
\AgdaBound{Γ₁}\AgdaSpace{}%
\AgdaOperator{\AgdaInductiveConstructor{,b}}\AgdaSpace{}%
\AgdaBound{B}\AgdaSymbol{)}\<%
\\
%
\>[2]\AgdaSymbol{→}\AgdaSpace{}%
\AgdaFunction{extsb}\AgdaSpace{}%
\AgdaBound{fb''}\AgdaSpace{}%
\AgdaSymbol{(}\AgdaFunction{extb}\AgdaSpace{}%
\AgdaBound{fb'}\AgdaSpace{}%
\AgdaBound{x}\AgdaSymbol{)}\AgdaSpace{}%
\AgdaOperator{\AgdaDatatype{≡}}\AgdaSpace{}%
\AgdaFunction{extsb}\AgdaSpace{}%
\AgdaSymbol{(λ}\AgdaSpace{}%
\AgdaBound{x₁}\AgdaSpace{}%
\AgdaSymbol{→}\AgdaSpace{}%
\AgdaBound{fb''}\AgdaSpace{}%
\AgdaSymbol{(}\AgdaBound{fb'}\AgdaSpace{}%
\AgdaBound{x₁}\AgdaSymbol{))}\AgdaSpace{}%
\AgdaBound{x}\<%
\\
\>[0]\AgdaFunction{extsb-comp}\AgdaSpace{}%
\AgdaBound{fb'}\AgdaSpace{}%
\AgdaBound{fb''}\AgdaSpace{}%
\AgdaInductiveConstructor{Zb}\AgdaSpace{}%
\AgdaSymbol{=}\AgdaSpace{}%
\AgdaInductiveConstructor{refl}\<%
\\
\>[0]\AgdaFunction{extsb-comp}\AgdaSpace{}%
\AgdaBound{fb'}\AgdaSpace{}%
\AgdaBound{fb''}\AgdaSpace{}%
\AgdaSymbol{(}\AgdaInductiveConstructor{Sb}\AgdaSpace{}%
\AgdaBound{x}\AgdaSymbol{)}\AgdaSpace{}%
\AgdaSymbol{=}\AgdaSpace{}%
\AgdaInductiveConstructor{refl}\<%
\\
\>[0]\<%
\end{code}
\begin{code}%
\>[0]\AgdaComment{--\ substb\ and\ rename\ distribute\ over\ composition:}\<%
\\
\>[0]\AgdaComment{--\ Renaming\ first\ with\ fb'\ and\ substituting\ with\ fb''}\<%
\\
\>[0]\AgdaComment{--\ is\ the\ same\ as\ substituting\ with\ fb''\ ∘\ fb'}\<%
\\
\>[0]\AgdaFunction{subst-renameb-comp}\AgdaSpace{}%
\AgdaSymbol{:}\AgdaSpace{}%
\AgdaSymbol{∀}\AgdaSpace{}%
\AgdaSymbol{\{}\AgdaBound{Γ₁}\AgdaSpace{}%
\AgdaBound{Γ₂}\AgdaSpace{}%
\AgdaBound{Γ₃}\AgdaSpace{}%
\AgdaBound{A}\AgdaSymbol{\}}\<%
\\
\>[0][@{}l@{\AgdaIndent{0}}]%
\>[2]\AgdaSymbol{→}\AgdaSpace{}%
\AgdaSymbol{(}\AgdaBound{fb'}\AgdaSpace{}%
\AgdaSymbol{:}\AgdaSpace{}%
\AgdaSymbol{∀}\AgdaSpace{}%
\AgdaSymbol{\{}\AgdaBound{Ab}\AgdaSymbol{\}}\AgdaSpace{}%
\AgdaSymbol{→}\AgdaSpace{}%
\AgdaBound{Ab}\AgdaSpace{}%
\AgdaOperator{\AgdaDatatype{∈b}}\AgdaSpace{}%
\AgdaBound{Γ₁}\AgdaSpace{}%
\AgdaSymbol{→}\AgdaSpace{}%
\AgdaBound{Ab}\AgdaSpace{}%
\AgdaOperator{\AgdaDatatype{∈b}}\AgdaSpace{}%
\AgdaBound{Γ₂}\AgdaSymbol{)}\<%
\\
%
\>[2]\AgdaSymbol{→}\AgdaSpace{}%
\AgdaSymbol{(}\AgdaBound{fb''}\AgdaSpace{}%
\AgdaSymbol{:}\AgdaSpace{}%
\AgdaSymbol{∀}\AgdaSpace{}%
\AgdaSymbol{\{}\AgdaBound{Ab}\AgdaSymbol{\}}\AgdaSpace{}%
\AgdaSymbol{→}\AgdaSpace{}%
\AgdaInductiveConstructor{□}\AgdaSpace{}%
\AgdaBound{Ab}\AgdaSpace{}%
\AgdaOperator{\AgdaDatatype{∈b}}\AgdaSpace{}%
\AgdaBound{Γ₂}\AgdaSpace{}%
\AgdaSymbol{→}\AgdaSpace{}%
\AgdaBound{Γ₃}\AgdaSpace{}%
\AgdaOperator{\AgdaDatatype{⊢b}}\AgdaSpace{}%
\AgdaBound{Ab}\AgdaSymbol{)}\<%
\\
%
\>[2]\AgdaSymbol{→}\AgdaSpace{}%
\AgdaSymbol{(}\AgdaBound{t}\AgdaSpace{}%
\AgdaSymbol{:}\AgdaSpace{}%
\AgdaBound{Γ₁}\AgdaSpace{}%
\AgdaOperator{\AgdaDatatype{⊢b}}\AgdaSpace{}%
\AgdaBound{A}\AgdaSymbol{)}\<%
\\
%
\>[2]\AgdaSymbol{→}\AgdaSpace{}%
\AgdaFunction{substb}\AgdaSpace{}%
\AgdaBound{fb''}\AgdaSpace{}%
\AgdaSymbol{(}\AgdaFunction{renameb}\AgdaSpace{}%
\AgdaBound{fb'}\AgdaSpace{}%
\AgdaBound{t}\AgdaSymbol{)}\AgdaSpace{}%
\AgdaOperator{\AgdaDatatype{≡}}\AgdaSpace{}%
\AgdaFunction{substb}\AgdaSpace{}%
\AgdaSymbol{(λ}\AgdaSpace{}%
\AgdaBound{x}\AgdaSpace{}%
\AgdaSymbol{→}\AgdaSpace{}%
\AgdaBound{fb''}\AgdaSpace{}%
\AgdaSymbol{(}\AgdaBound{fb'}\AgdaSpace{}%
\AgdaBound{x}\AgdaSymbol{))}\AgdaSpace{}%
\AgdaBound{t}\<%
\\
\>[0]\AgdaFunction{subst-renameb-comp}\AgdaSpace{}%
\AgdaBound{fb'}\AgdaSpace{}%
\AgdaBound{fb''}\AgdaSpace{}%
\AgdaSymbol{(}\AgdaInductiveConstructor{ε}\AgdaSpace{}%
\AgdaBound{x}\AgdaSymbol{)}\AgdaSpace{}%
\AgdaSymbol{=}\AgdaSpace{}%
\AgdaInductiveConstructor{refl}\<%
\\
\>[0]\AgdaFunction{subst-renameb-comp}\AgdaSpace{}%
\AgdaBound{fb'}\AgdaSpace{}%
\AgdaBound{fb''}\AgdaSpace{}%
\AgdaSymbol{(}\AgdaInductiveConstructor{box}\AgdaSpace{}%
\AgdaBound{t}\AgdaSymbol{)}\AgdaSpace{}%
\AgdaSymbol{=}\AgdaSpace{}%
\AgdaFunction{cong}\AgdaSpace{}%
\AgdaInductiveConstructor{box}\AgdaSpace{}%
\AgdaSymbol{(}\AgdaFunction{subst-renameb-comp}\AgdaSpace{}%
\AgdaBound{fb'}\AgdaSpace{}%
\AgdaBound{fb''}\AgdaSpace{}%
\AgdaBound{t}\AgdaSymbol{)}\<%
\\
\>[0]\AgdaFunction{subst-renameb-comp}\AgdaSpace{}%
\AgdaBound{fb'}\AgdaSpace{}%
\AgdaBound{fb''}\AgdaSpace{}%
\AgdaSymbol{(}\AgdaInductiveConstructor{ƛb}\AgdaSpace{}%
\AgdaBound{t}\AgdaSymbol{)}\AgdaSpace{}%
\AgdaSymbol{=}\<%
\\
\>[0][@{}l@{\AgdaIndent{0}}]%
\>[2]\AgdaOperator{\AgdaFunction{beginEq}}\<%
\\
\>[2][@{}l@{\AgdaIndent{0}}]%
\>[4]\AgdaInductiveConstructor{ƛb}\AgdaSpace{}%
\AgdaSymbol{(}\AgdaFunction{substb}\AgdaSpace{}%
\AgdaSymbol{(}\AgdaFunction{extsb}\AgdaSpace{}%
\AgdaBound{fb''}\AgdaSymbol{)}\AgdaSpace{}%
\AgdaSymbol{(}\AgdaFunction{renameb}\AgdaSpace{}%
\AgdaSymbol{(}\AgdaFunction{extb}\AgdaSpace{}%
\AgdaBound{fb'}\AgdaSymbol{)}\AgdaSpace{}%
\AgdaBound{t}\AgdaSymbol{))}\<%
\\
%
\>[2]\AgdaFunction{≡⟨}\AgdaSpace{}%
\AgdaFunction{cong}\AgdaSpace{}%
\AgdaInductiveConstructor{ƛb}\AgdaSpace{}%
\AgdaSymbol{(}\AgdaFunction{subst-renameb-comp}\AgdaSpace{}%
\AgdaSymbol{(}\AgdaFunction{extb}\AgdaSpace{}%
\AgdaBound{fb'}\AgdaSymbol{)}\AgdaSpace{}%
\AgdaSymbol{(}\AgdaFunction{extsb}\AgdaSpace{}%
\AgdaBound{fb''}\AgdaSymbol{)}\AgdaSpace{}%
\AgdaBound{t}\AgdaSymbol{)}\AgdaSpace{}%
\AgdaFunction{⟩}\<%
\\
\>[2][@{}l@{\AgdaIndent{0}}]%
\>[4]\AgdaInductiveConstructor{ƛb}\AgdaSpace{}%
\AgdaSymbol{(}\AgdaFunction{substb}\AgdaSpace{}%
\AgdaSymbol{(λ}\AgdaSpace{}%
\AgdaBound{x}\AgdaSpace{}%
\AgdaSymbol{→}\AgdaSpace{}%
\AgdaSymbol{(}\AgdaFunction{extsb}\AgdaSpace{}%
\AgdaBound{fb''}\AgdaSymbol{)}\AgdaSpace{}%
\AgdaSymbol{((}\AgdaFunction{extb}\AgdaSpace{}%
\AgdaBound{fb'}\AgdaSymbol{)}\AgdaSpace{}%
\AgdaBound{x}\AgdaSymbol{))}\AgdaSpace{}%
\AgdaBound{t}\AgdaSymbol{)}\<%
\\
%
\>[2]\AgdaFunction{≡⟨}\AgdaSpace{}%
\AgdaFunction{cong}\AgdaSpace{}%
\AgdaInductiveConstructor{ƛb}\AgdaSpace{}%
\AgdaSymbol{(}\AgdaFunction{substb-cong}\AgdaSpace{}%
\AgdaSymbol{(λ}\AgdaSpace{}%
\AgdaBound{x}\AgdaSpace{}%
\AgdaSymbol{→}\AgdaSpace{}%
\AgdaFunction{extsb-comp}\AgdaSpace{}%
\AgdaBound{fb'}\AgdaSpace{}%
\AgdaBound{fb''}\AgdaSpace{}%
\AgdaBound{x}\AgdaSymbol{)}\AgdaSpace{}%
\AgdaBound{t}\AgdaSymbol{)}\AgdaSpace{}%
\AgdaFunction{⟩}\<%
\\
\>[2][@{}l@{\AgdaIndent{0}}]%
\>[4]\AgdaInductiveConstructor{ƛb}\AgdaSpace{}%
\AgdaSymbol{(}\AgdaFunction{substb}\AgdaSpace{}%
\AgdaSymbol{(}\AgdaFunction{extsb}\AgdaSpace{}%
\AgdaSymbol{(λ}\AgdaSpace{}%
\AgdaBound{x}\AgdaSpace{}%
\AgdaSymbol{→}\AgdaSpace{}%
\AgdaBound{fb''}\AgdaSpace{}%
\AgdaSymbol{(}\AgdaBound{fb'}\AgdaSpace{}%
\AgdaBound{x}\AgdaSymbol{)))}\AgdaSpace{}%
\AgdaBound{t}\AgdaSymbol{)}\<%
\\
%
\>[2]\AgdaOperator{\AgdaFunction{∎eq}}\<%
\end{code}
\begin{code}[hide]%
\>[0]\AgdaFunction{subst-renameb-comp}\AgdaSpace{}%
\AgdaBound{fb'}\AgdaSpace{}%
\AgdaBound{fb''}\AgdaSpace{}%
\AgdaSymbol{(}\AgdaBound{t}\AgdaSpace{}%
\AgdaOperator{\AgdaInductiveConstructor{·b}}\AgdaSpace{}%
\AgdaBound{t₁}\AgdaSymbol{)}\AgdaSpace{}%
\AgdaSymbol{=}\AgdaSpace{}%
\AgdaFunction{cong₂}\AgdaSpace{}%
\AgdaOperator{\AgdaInductiveConstructor{\AgdaUnderscore{}·b\AgdaUnderscore{}}}\AgdaSpace{}%
\AgdaSymbol{(}\AgdaFunction{subst-renameb-comp}\AgdaSpace{}%
\AgdaBound{fb'}\AgdaSpace{}%
\AgdaBound{fb''}\AgdaSpace{}%
\AgdaBound{t}\AgdaSymbol{)}\AgdaSpace{}%
\AgdaSymbol{(}\AgdaFunction{subst-renameb-comp}\AgdaSpace{}%
\AgdaBound{fb'}\AgdaSpace{}%
\AgdaBound{fb''}\AgdaSpace{}%
\AgdaBound{t₁}\AgdaSymbol{)}\<%
\\
\>[0]\AgdaComment{--\ Renaming\ followed\ by\ substitution\ with\ an\ extended\ context\ (extsb)}\<%
\\
\>[0]\AgdaComment{--\ commutes\ with\ substitution\ followed\ by\ renaming.}\<%
\\
\>[0]\AgdaFunction{substb-exts-rename-Sb-commute}\AgdaSpace{}%
\AgdaSymbol{:}\AgdaSpace{}%
\AgdaSymbol{∀}\AgdaSpace{}%
\AgdaSymbol{\{}\AgdaBound{Γ}\AgdaSpace{}%
\AgdaBound{A}\AgdaSpace{}%
\AgdaBound{B}\AgdaSpace{}%
\AgdaBound{Δ}\AgdaSymbol{\}}\<%
\\
\>[0][@{}l@{\AgdaIndent{0}}]%
\>[2]\AgdaSymbol{→}\AgdaSpace{}%
\AgdaSymbol{(}\AgdaBound{fb'}\AgdaSpace{}%
\AgdaSymbol{:}\AgdaSpace{}%
\AgdaSymbol{∀}\AgdaSpace{}%
\AgdaSymbol{\{}\AgdaBound{Ab}\AgdaSymbol{\}}\AgdaSpace{}%
\AgdaSymbol{→}\AgdaSpace{}%
\AgdaInductiveConstructor{□}\AgdaSpace{}%
\AgdaBound{Ab}\AgdaSpace{}%
\AgdaOperator{\AgdaDatatype{∈b}}\AgdaSpace{}%
\AgdaBound{Γ}\AgdaSpace{}%
\AgdaSymbol{→}\AgdaSpace{}%
\AgdaBound{Δ}\AgdaSpace{}%
\AgdaOperator{\AgdaDatatype{⊢b}}\AgdaSpace{}%
\AgdaBound{Ab}\AgdaSymbol{)}\<%
\\
%
\>[2]\AgdaSymbol{→}\AgdaSpace{}%
\AgdaSymbol{(}\AgdaBound{t}\AgdaSpace{}%
\AgdaSymbol{:}\AgdaSpace{}%
\AgdaBound{Γ}\AgdaSpace{}%
\AgdaOperator{\AgdaDatatype{⊢b}}\AgdaSpace{}%
\AgdaBound{A}\AgdaSymbol{)}\<%
\\
%
\>[2]\AgdaSymbol{→}%
\>[1555I]\AgdaFunction{substb}\AgdaSpace{}%
\AgdaSymbol{(}\AgdaFunction{extsb}\AgdaSpace{}%
\AgdaBound{fb'}\AgdaSpace{}%
\AgdaSymbol{\{}\AgdaArgument{B}\AgdaSpace{}%
\AgdaSymbol{=}\AgdaSpace{}%
\AgdaBound{B}\AgdaSymbol{\})}\AgdaSpace{}%
\AgdaSymbol{(}\AgdaFunction{renameb}\AgdaSpace{}%
\AgdaInductiveConstructor{Sb}\AgdaSpace{}%
\AgdaBound{t}\AgdaSymbol{)}\<%
\\
\>[.][@{}l@{}]\<[1555I]%
\>[4]\AgdaOperator{\AgdaDatatype{≡}}\AgdaSpace{}%
\AgdaFunction{renameb}\AgdaSpace{}%
\AgdaInductiveConstructor{Sb}\AgdaSpace{}%
\AgdaSymbol{(}\AgdaFunction{substb}\AgdaSpace{}%
\AgdaBound{fb'}\AgdaSpace{}%
\AgdaBound{t}\AgdaSymbol{)}\<%
\\
\>[0]\AgdaFunction{substb-exts-rename-Sb-commute}\AgdaSpace{}%
\AgdaBound{fb'}\AgdaSpace{}%
\AgdaBound{t}\AgdaSpace{}%
\AgdaSymbol{=}\<%
\\
\>[0][@{}l@{\AgdaIndent{0}}]%
\>[2]\AgdaOperator{\AgdaFunction{beginEq}}\<%
\\
\>[2][@{}l@{\AgdaIndent{0}}]%
\>[4]\AgdaFunction{substb}\AgdaSpace{}%
\AgdaSymbol{(}\AgdaFunction{extsb}\AgdaSpace{}%
\AgdaBound{fb'}\AgdaSymbol{)}\AgdaSpace{}%
\AgdaSymbol{(}\AgdaFunction{renameb}\AgdaSpace{}%
\AgdaInductiveConstructor{Sb}\AgdaSpace{}%
\AgdaBound{t}\AgdaSymbol{)}\<%
\\
%
\>[2]\AgdaFunction{≡⟨}\AgdaSpace{}%
\AgdaFunction{subst-renameb-comp}\AgdaSpace{}%
\AgdaInductiveConstructor{Sb}\AgdaSpace{}%
\AgdaSymbol{(}\AgdaFunction{extsb}\AgdaSpace{}%
\AgdaBound{fb'}\AgdaSymbol{)}\AgdaSpace{}%
\AgdaBound{t}\AgdaSpace{}%
\AgdaFunction{⟩}\<%
\\
\>[2][@{}l@{\AgdaIndent{0}}]%
\>[4]\AgdaFunction{substb}\AgdaSpace{}%
\AgdaSymbol{(λ}\AgdaSpace{}%
\AgdaBound{x}\AgdaSpace{}%
\AgdaSymbol{→}\AgdaSpace{}%
\AgdaFunction{renameb}\AgdaSpace{}%
\AgdaInductiveConstructor{Sb}\AgdaSpace{}%
\AgdaSymbol{(}\AgdaBound{fb'}\AgdaSpace{}%
\AgdaBound{x}\AgdaSymbol{))}\AgdaSpace{}%
\AgdaBound{t}\<%
\\
%
\>[2]\AgdaFunction{≡⟨}\AgdaSpace{}%
\AgdaFunction{substb-rename-commute}\AgdaSpace{}%
\AgdaInductiveConstructor{Sb}\AgdaSpace{}%
\AgdaBound{fb'}\AgdaSpace{}%
\AgdaBound{t}\AgdaSpace{}%
\AgdaFunction{⟩}\<%
\\
\>[2][@{}l@{\AgdaIndent{0}}]%
\>[4]\AgdaFunction{renameb}\AgdaSpace{}%
\AgdaInductiveConstructor{Sb}\AgdaSpace{}%
\AgdaSymbol{(}\AgdaFunction{substb}\AgdaSpace{}%
\AgdaBound{fb'}\AgdaSpace{}%
\AgdaBound{t}\AgdaSymbol{)}\<%
\\
%
\>[2]\AgdaOperator{\AgdaFunction{∎eq}}\<%
\\
\>[0]\AgdaComment{--\ Preservation\ and\ Reflection\ of\ }\<%
\\
\>[0]\AgdaComment{--\ simultaneous\ substitution\ of\ Gödel's\ translation}\<%
\\
\>[0]\AgdaFunction{pres-simultaneous-subst}\AgdaSpace{}%
\AgdaSymbol{:}\AgdaSpace{}%
\AgdaSymbol{∀}\AgdaSpace{}%
\AgdaSymbol{\{}\AgdaBound{Γ}\AgdaSpace{}%
\AgdaBound{Δ}\AgdaSpace{}%
\AgdaBound{A}\AgdaSymbol{\}}\<%
\\
\>[0][@{}l@{\AgdaIndent{0}}]%
\>[2]\AgdaSymbol{→}\AgdaSpace{}%
\AgdaSymbol{(}\AgdaBound{fb'}\AgdaSpace{}%
\AgdaSymbol{:}\AgdaSpace{}%
\AgdaSymbol{∀}\AgdaSpace{}%
\AgdaSymbol{\{}\AgdaBound{Ab}\AgdaSymbol{\}}\AgdaSpace{}%
\AgdaSymbol{→}\AgdaSpace{}%
\AgdaInductiveConstructor{□}\AgdaSpace{}%
\AgdaBound{Ab}\AgdaSpace{}%
\AgdaOperator{\AgdaDatatype{∈b}}\AgdaSpace{}%
\AgdaFunction{embedContext}\AgdaSpace{}%
\AgdaBound{Γ}\AgdaSpace{}%
\AgdaSymbol{→}\AgdaSpace{}%
\AgdaFunction{embedContext}\AgdaSpace{}%
\AgdaBound{Δ}\AgdaSpace{}%
\AgdaOperator{\AgdaDatatype{⊢b}}\AgdaSpace{}%
\AgdaBound{Ab}\AgdaSymbol{)}\<%
\\
%
\>[2]\AgdaSymbol{→}\AgdaSpace{}%
\AgdaSymbol{(}\AgdaBound{f'}\AgdaSpace{}%
\AgdaSymbol{:}\AgdaSpace{}%
\AgdaSymbol{∀}\AgdaSpace{}%
\AgdaSymbol{\{}\AgdaBound{A}\AgdaSymbol{\}}\AgdaSpace{}%
\AgdaSymbol{→}\AgdaSpace{}%
\AgdaBound{A}\AgdaSpace{}%
\AgdaOperator{\AgdaDatatype{∈}}\AgdaSpace{}%
\AgdaBound{Γ}\AgdaSpace{}%
\AgdaSymbol{→}\AgdaSpace{}%
\AgdaBound{Δ}\AgdaSpace{}%
\AgdaOperator{\AgdaDatatype{⊢}}\AgdaSpace{}%
\AgdaBound{A}\AgdaSymbol{)}\<%
\\
%
\>[2]\AgdaSymbol{→}\AgdaSpace{}%
\AgdaSymbol{(}\AgdaBound{h}\AgdaSpace{}%
\AgdaSymbol{:}\AgdaSpace{}%
\AgdaSymbol{∀}\AgdaSpace{}%
\AgdaSymbol{\{}\AgdaBound{A'}\AgdaSymbol{\}}\AgdaSpace{}%
\AgdaSymbol{→}\AgdaSpace{}%
\AgdaSymbol{(}\AgdaBound{x}\AgdaSpace{}%
\AgdaSymbol{:}\AgdaSpace{}%
\AgdaBound{A'}\AgdaSpace{}%
\AgdaOperator{\AgdaDatatype{∈}}\AgdaSpace{}%
\AgdaBound{Γ}\AgdaSymbol{)}\AgdaSpace{}%
\AgdaSymbol{→}\AgdaSpace{}%
\AgdaInductiveConstructor{box}\AgdaSpace{}%
\AgdaSymbol{(}\AgdaBound{fb'}\AgdaSpace{}%
\AgdaSymbol{(}\AgdaFunction{embedLookup}\AgdaSpace{}%
\AgdaBound{x}\AgdaSymbol{))}\AgdaSpace{}%
\AgdaOperator{\AgdaDatatype{≡}}\AgdaSpace{}%
\AgdaFunction{embedTerm}\AgdaSpace{}%
\AgdaSymbol{(}\AgdaBound{f'}\AgdaSpace{}%
\AgdaBound{x}\AgdaSymbol{)}\AgdaSpace{}%
\AgdaSymbol{)}\<%
\\
%
\>[2]\AgdaSymbol{→}\AgdaSpace{}%
\AgdaSymbol{(}\AgdaBound{M}\AgdaSpace{}%
\AgdaSymbol{:}\AgdaSpace{}%
\AgdaBound{Γ}\AgdaSpace{}%
\AgdaOperator{\AgdaDatatype{⊢}}\AgdaSpace{}%
\AgdaBound{A}\AgdaSymbol{)}\<%
\\
%
\>[2]\AgdaSymbol{→}\AgdaSpace{}%
\AgdaFunction{substb}\AgdaSpace{}%
\AgdaBound{fb'}\AgdaSpace{}%
\AgdaSymbol{(}\AgdaFunction{embedTerm}\AgdaSpace{}%
\AgdaBound{M}\AgdaSymbol{)}\AgdaSpace{}%
\AgdaOperator{\AgdaDatatype{≡}}\AgdaSpace{}%
\AgdaFunction{embedTerm}\AgdaSpace{}%
\AgdaSymbol{(}\AgdaFunction{subst}\AgdaSpace{}%
\AgdaBound{f'}\AgdaSpace{}%
\AgdaBound{M}\AgdaSymbol{)}\<%
\\
\>[0]\AgdaFunction{pres-simultaneous-subst}\AgdaSpace{}%
\AgdaBound{fb'}\AgdaSpace{}%
\AgdaBound{f'}\AgdaSpace{}%
\AgdaBound{h}\AgdaSpace{}%
\AgdaSymbol{(}\AgdaInductiveConstructor{`}\AgdaSpace{}%
\AgdaBound{x}\AgdaSymbol{)}\AgdaSpace{}%
\AgdaSymbol{=}\<%
\\
\>[0][@{}l@{\AgdaIndent{0}}]%
\>[2]\AgdaBound{h}\AgdaSpace{}%
\AgdaBound{x}\<%
\\
\>[0]\AgdaFunction{pres-simultaneous-subst}\AgdaSpace{}%
\AgdaBound{fb'}\AgdaSpace{}%
\AgdaBound{f'}\AgdaSpace{}%
\AgdaBound{h}\AgdaSpace{}%
\AgdaSymbol{(}\AgdaInductiveConstructor{ƛ}\AgdaSpace{}%
\AgdaBound{M}\AgdaSymbol{)}\AgdaSpace{}%
\AgdaSymbol{=}\<%
\\
\>[0][@{}l@{\AgdaIndent{0}}]%
\>[2]\AgdaFunction{cong}\AgdaSpace{}%
\AgdaInductiveConstructor{box}\AgdaSpace{}%
\AgdaSymbol{(}\AgdaFunction{cong}\AgdaSpace{}%
\AgdaInductiveConstructor{ƛb}\AgdaSpace{}%
\AgdaSymbol{(}\AgdaFunction{pres-simultaneous-subst}\AgdaSpace{}%
\AgdaSymbol{(}\AgdaFunction{extsb}\AgdaSpace{}%
\AgdaBound{fb'}\AgdaSymbol{)}\AgdaSpace{}%
\AgdaSymbol{(}\AgdaFunction{exts}\AgdaSpace{}%
\AgdaBound{f'}\AgdaSymbol{)}\AgdaSpace{}%
\AgdaSymbol{(}\AgdaFunction{pres-exts}\AgdaSpace{}%
\AgdaBound{fb'}\AgdaSpace{}%
\AgdaBound{f'}\AgdaSpace{}%
\AgdaBound{h}\AgdaSymbol{)}\AgdaSpace{}%
\AgdaBound{M}\AgdaSymbol{))}\<%
\\
\>[0]\AgdaFunction{pres-simultaneous-subst}\AgdaSpace{}%
\AgdaBound{fb'}\AgdaSpace{}%
\AgdaBound{f'}\AgdaSpace{}%
\AgdaBound{h}\AgdaSpace{}%
\AgdaSymbol{(}\AgdaBound{M}\AgdaSpace{}%
\AgdaOperator{\AgdaInductiveConstructor{·}}\AgdaSpace{}%
\AgdaBound{M₁}\AgdaSymbol{)}\AgdaSpace{}%
\AgdaSymbol{=}\<%
\\
\>[0][@{}l@{\AgdaIndent{0}}]%
\>[2]\AgdaOperator{\AgdaFunction{beginEq}}\<%
\\
\>[2][@{}l@{\AgdaIndent{0}}]%
\>[4]\AgdaInductiveConstructor{ƛb}\AgdaSpace{}%
\AgdaSymbol{(}\AgdaInductiveConstructor{ε}\AgdaSpace{}%
\AgdaInductiveConstructor{Zb}\AgdaSpace{}%
\AgdaOperator{\AgdaInductiveConstructor{·b}}\AgdaSpace{}%
\AgdaFunction{substb}\AgdaSpace{}%
\AgdaSymbol{(}\AgdaFunction{extsb}\AgdaSpace{}%
\AgdaBound{fb'}\AgdaSymbol{)}\AgdaSpace{}%
\AgdaSymbol{(}\AgdaFunction{renameb}\AgdaSpace{}%
\AgdaInductiveConstructor{Sb}\AgdaSpace{}%
\AgdaSymbol{(}\AgdaFunction{embedTerm}\AgdaSpace{}%
\AgdaBound{M₁}\AgdaSymbol{)))}\AgdaSpace{}%
\AgdaOperator{\AgdaInductiveConstructor{·b}}\AgdaSpace{}%
\AgdaFunction{substb}\AgdaSpace{}%
\AgdaBound{fb'}\AgdaSpace{}%
\AgdaSymbol{(}\AgdaFunction{embedTerm}\AgdaSpace{}%
\AgdaBound{M}\AgdaSymbol{)}\<%
\\
%
\>[2]\AgdaFunction{≡⟨}\AgdaSpace{}%
\AgdaFunction{cong₂}\AgdaSpace{}%
\AgdaOperator{\AgdaInductiveConstructor{\AgdaUnderscore{}·b\AgdaUnderscore{}}}\AgdaSpace{}%
\AgdaInductiveConstructor{refl}\AgdaSpace{}%
\AgdaSymbol{(}\AgdaFunction{pres-simultaneous-subst}\AgdaSpace{}%
\AgdaBound{fb'}\AgdaSpace{}%
\AgdaBound{f'}\AgdaSpace{}%
\AgdaBound{h}\AgdaSpace{}%
\AgdaBound{M}\AgdaSymbol{)}\AgdaSpace{}%
\AgdaFunction{⟩}\<%
\\
\>[2][@{}l@{\AgdaIndent{0}}]%
\>[4]\AgdaInductiveConstructor{ƛb}\AgdaSpace{}%
\AgdaSymbol{(}\AgdaInductiveConstructor{ε}\AgdaSpace{}%
\AgdaInductiveConstructor{Zb}\AgdaSpace{}%
\AgdaOperator{\AgdaInductiveConstructor{·b}}\AgdaSpace{}%
\AgdaFunction{substb}\AgdaSpace{}%
\AgdaSymbol{(}\AgdaFunction{extsb}\AgdaSpace{}%
\AgdaBound{fb'}\AgdaSymbol{)}\AgdaSpace{}%
\AgdaSymbol{(}\AgdaFunction{renameb}\AgdaSpace{}%
\AgdaInductiveConstructor{Sb}\AgdaSpace{}%
\AgdaSymbol{(}\AgdaFunction{embedTerm}\AgdaSpace{}%
\AgdaBound{M₁}\AgdaSymbol{)))}\AgdaSpace{}%
\AgdaOperator{\AgdaInductiveConstructor{·b}}\AgdaSpace{}%
\AgdaFunction{embedTerm}\AgdaSpace{}%
\AgdaSymbol{(}\AgdaFunction{subst}\AgdaSpace{}%
\AgdaBound{f'}\AgdaSpace{}%
\AgdaBound{M}\AgdaSymbol{)}\<%
\\
%
\>[2]\AgdaFunction{≡⟨}\AgdaSpace{}%
\AgdaFunction{cong₂}\AgdaSpace{}%
\AgdaOperator{\AgdaInductiveConstructor{\AgdaUnderscore{}·b\AgdaUnderscore{}}}\AgdaSpace{}%
\AgdaSymbol{(}\AgdaFunction{cong}\AgdaSpace{}%
\AgdaInductiveConstructor{ƛb}\AgdaSpace{}%
\AgdaSymbol{(}\AgdaFunction{cong₂}\AgdaSpace{}%
\AgdaOperator{\AgdaInductiveConstructor{\AgdaUnderscore{}·b\AgdaUnderscore{}}}\AgdaSpace{}%
\AgdaInductiveConstructor{refl}\AgdaSpace{}%
\AgdaSymbol{(}\AgdaFunction{substb-exts-rename-Sb-commute}\AgdaSpace{}%
\AgdaBound{fb'}\AgdaSpace{}%
\AgdaSymbol{(}\AgdaFunction{embedTerm}\AgdaSpace{}%
\AgdaBound{M₁}\AgdaSymbol{))))}\AgdaSpace{}%
\AgdaInductiveConstructor{refl}\AgdaSpace{}%
\AgdaFunction{⟩}\<%
\\
\>[2][@{}l@{\AgdaIndent{0}}]%
\>[4]\AgdaInductiveConstructor{ƛb}\AgdaSpace{}%
\AgdaSymbol{(}\AgdaInductiveConstructor{ε}\AgdaSpace{}%
\AgdaInductiveConstructor{Zb}\AgdaSpace{}%
\AgdaOperator{\AgdaInductiveConstructor{·b}}\AgdaSpace{}%
\AgdaFunction{renameb}\AgdaSpace{}%
\AgdaInductiveConstructor{Sb}\AgdaSpace{}%
\AgdaSymbol{(}\AgdaFunction{substb}\AgdaSpace{}%
\AgdaBound{fb'}\AgdaSpace{}%
\AgdaSymbol{(}\AgdaFunction{embedTerm}\AgdaSpace{}%
\AgdaBound{M₁}\AgdaSymbol{)))}\AgdaSpace{}%
\AgdaOperator{\AgdaInductiveConstructor{·b}}\AgdaSpace{}%
\AgdaFunction{embedTerm}\AgdaSpace{}%
\AgdaSymbol{(}\AgdaFunction{subst}\AgdaSpace{}%
\AgdaBound{f'}\AgdaSpace{}%
\AgdaBound{M}\AgdaSymbol{)}\<%
\\
%
\>[2]\AgdaFunction{≡⟨}\AgdaSpace{}%
\AgdaFunction{cong₂}\AgdaSpace{}%
\AgdaOperator{\AgdaInductiveConstructor{\AgdaUnderscore{}·b\AgdaUnderscore{}}}\AgdaSpace{}%
\AgdaSymbol{(}\AgdaFunction{cong}\AgdaSpace{}%
\AgdaInductiveConstructor{ƛb}\AgdaSpace{}%
\AgdaSymbol{(}\AgdaFunction{cong₂}\AgdaSpace{}%
\AgdaOperator{\AgdaInductiveConstructor{\AgdaUnderscore{}·b\AgdaUnderscore{}}}\AgdaSpace{}%
\AgdaInductiveConstructor{refl}\AgdaSpace{}%
\AgdaSymbol{(}\AgdaFunction{cong}\AgdaSpace{}%
\AgdaSymbol{(}\AgdaFunction{renameb}\AgdaSpace{}%
\AgdaInductiveConstructor{Sb}\AgdaSymbol{)}\AgdaSpace{}%
\AgdaSymbol{(}\AgdaFunction{pres-simultaneous-subst}\AgdaSpace{}%
\AgdaBound{fb'}\AgdaSpace{}%
\AgdaBound{f'}\AgdaSpace{}%
\AgdaBound{h}\AgdaSpace{}%
\AgdaBound{M₁}\AgdaSymbol{))))}\AgdaSpace{}%
\AgdaInductiveConstructor{refl}\AgdaSpace{}%
\AgdaFunction{⟩}\<%
\\
\>[2][@{}l@{\AgdaIndent{0}}]%
\>[4]\AgdaInductiveConstructor{ƛb}\AgdaSpace{}%
\AgdaSymbol{(}\AgdaInductiveConstructor{ε}\AgdaSpace{}%
\AgdaInductiveConstructor{Zb}\AgdaSpace{}%
\AgdaOperator{\AgdaInductiveConstructor{·b}}\AgdaSpace{}%
\AgdaFunction{renameb}\AgdaSpace{}%
\AgdaInductiveConstructor{Sb}\AgdaSpace{}%
\AgdaSymbol{(}\AgdaFunction{embedTerm}\AgdaSpace{}%
\AgdaSymbol{(}\AgdaFunction{subst}\AgdaSpace{}%
\AgdaBound{f'}\AgdaSpace{}%
\AgdaBound{M₁}\AgdaSymbol{)))}\AgdaSpace{}%
\AgdaOperator{\AgdaInductiveConstructor{·b}}\AgdaSpace{}%
\AgdaFunction{embedTerm}\AgdaSpace{}%
\AgdaSymbol{(}\AgdaFunction{subst}\AgdaSpace{}%
\AgdaBound{f'}\AgdaSpace{}%
\AgdaBound{M}\AgdaSymbol{)}\<%
\\
%
\>[2]\AgdaOperator{\AgdaFunction{∎eq}}\<%
\\
%
\\[\AgdaEmptyExtraSkip]%
\>[0]\AgdaComment{--\ HELPER\ LEMMA'S\ TO\ PROVE\ h}\<%
\\
%
\\[\AgdaEmptyExtraSkip]%
\>[0]\AgdaComment{--\ Lemma\ for\ assumption\ h\ in\ lemma}\<%
\\
\>[0]\AgdaComment{--\ pres-simultaneous-subst}\<%
\\
\>[0]\AgdaFunction{pres-subst-func}\AgdaSpace{}%
\AgdaSymbol{:}\AgdaSpace{}%
\AgdaSymbol{∀}\AgdaSpace{}%
\AgdaSymbol{\{}\AgdaBound{A}\AgdaSpace{}%
\AgdaBound{C}\AgdaSpace{}%
\AgdaBound{Γ}\AgdaSymbol{\}}\<%
\\
\>[0][@{}l@{\AgdaIndent{0}}]%
\>[2]\AgdaSymbol{→}\AgdaSpace{}%
\AgdaSymbol{(}\AgdaBound{V}\AgdaSpace{}%
\AgdaSymbol{:}\AgdaSpace{}%
\AgdaBound{Γ}\AgdaSpace{}%
\AgdaOperator{\AgdaDatatype{⊢}}\AgdaSpace{}%
\AgdaBound{C}\AgdaSymbol{)}\<%
\\
%
\>[2]\AgdaSymbol{→}\AgdaSpace{}%
\AgdaSymbol{(}\AgdaBound{y}\AgdaSpace{}%
\AgdaSymbol{:}\AgdaSpace{}%
\AgdaDatatype{Value}\AgdaSpace{}%
\AgdaBound{V}\AgdaSymbol{)}\<%
\\
%
\>[2]\AgdaSymbol{→}\AgdaSpace{}%
\AgdaSymbol{(}\AgdaBound{x}\AgdaSpace{}%
\AgdaSymbol{:}\AgdaSpace{}%
\AgdaBound{A}\AgdaSpace{}%
\AgdaOperator{\AgdaDatatype{∈}}\AgdaSpace{}%
\AgdaBound{Γ}\AgdaSpace{}%
\AgdaOperator{\AgdaInductiveConstructor{,}}\AgdaSpace{}%
\AgdaBound{C}\AgdaSymbol{)}\<%
\\
%
\>[2]\AgdaSymbol{→}\AgdaSpace{}%
\AgdaInductiveConstructor{box}\AgdaSpace{}%
\AgdaSymbol{(}\AgdaFunction{fb}\AgdaSpace{}%
\AgdaSymbol{(}\AgdaFunction{embedValue}\AgdaSpace{}%
\AgdaBound{V}\AgdaSpace{}%
\AgdaBound{y}\AgdaSymbol{)}\AgdaSpace{}%
\AgdaSymbol{(}\AgdaFunction{embedLookup}\AgdaSpace{}%
\AgdaBound{x}\AgdaSymbol{))}\AgdaSpace{}%
\AgdaOperator{\AgdaDatatype{≡}}\AgdaSpace{}%
\AgdaFunction{embedTerm}\AgdaSpace{}%
\AgdaSymbol{((}\AgdaFunction{f}\AgdaSpace{}%
\AgdaBound{V}\AgdaSymbol{)}\AgdaSpace{}%
\AgdaBound{x}\AgdaSymbol{)}\<%
\\
\>[0]\AgdaFunction{pres-subst-func}\AgdaSpace{}%
\AgdaBound{V}\AgdaSpace{}%
\AgdaBound{y}\AgdaSpace{}%
\AgdaInductiveConstructor{Z}\AgdaSpace{}%
\AgdaSymbol{=}\AgdaSpace{}%
\AgdaFunction{box-embedValue}\AgdaSpace{}%
\AgdaBound{y}\<%
\\
\>[0]\AgdaFunction{pres-subst-func}\AgdaSpace{}%
\AgdaBound{V}\AgdaSpace{}%
\AgdaBound{y}\AgdaSpace{}%
\AgdaSymbol{(}\AgdaInductiveConstructor{S}\AgdaSpace{}%
\AgdaBound{x}\AgdaSymbol{)}\AgdaSpace{}%
\AgdaSymbol{=}\AgdaSpace{}%
\AgdaInductiveConstructor{refl}\<%
\\
%
\\[\AgdaEmptyExtraSkip]%
\>[0]\AgdaComment{--\ exts\ lemma\ for\ assumption\ h\ in\ lemma}\<%
\\
\>[0]\AgdaFunction{exts-pres-subst-func}\AgdaSpace{}%
\AgdaSymbol{:}\AgdaSpace{}%
\AgdaSymbol{∀}\AgdaSpace{}%
\AgdaSymbol{\{}\AgdaBound{A}\AgdaSpace{}%
\AgdaBound{C}\AgdaSpace{}%
\AgdaBound{Γ}\AgdaSpace{}%
\AgdaBound{B}\AgdaSymbol{\}}\<%
\\
\>[0][@{}l@{\AgdaIndent{0}}]%
\>[2]\AgdaSymbol{→}\AgdaSpace{}%
\AgdaSymbol{(}\AgdaBound{V}\AgdaSpace{}%
\AgdaSymbol{:}\AgdaSpace{}%
\AgdaBound{Γ}\AgdaSpace{}%
\AgdaOperator{\AgdaDatatype{⊢}}\AgdaSpace{}%
\AgdaBound{C}\AgdaSymbol{)}\<%
\\
%
\>[2]\AgdaSymbol{→}\AgdaSpace{}%
\AgdaSymbol{(}\AgdaBound{y}\AgdaSpace{}%
\AgdaSymbol{:}\AgdaSpace{}%
\AgdaDatatype{Value}\AgdaSpace{}%
\AgdaBound{V}\AgdaSymbol{)}\<%
\\
%
\>[2]\AgdaSymbol{→}\AgdaSpace{}%
\AgdaSymbol{(}\AgdaBound{x}\AgdaSpace{}%
\AgdaSymbol{:}\AgdaSpace{}%
\AgdaBound{A}\AgdaSpace{}%
\AgdaOperator{\AgdaDatatype{∈}}\AgdaSpace{}%
\AgdaBound{Γ}\AgdaSpace{}%
\AgdaOperator{\AgdaInductiveConstructor{,}}\AgdaSpace{}%
\AgdaBound{C}\AgdaSpace{}%
\AgdaOperator{\AgdaInductiveConstructor{,}}\AgdaSpace{}%
\AgdaBound{B}\AgdaSymbol{)}\<%
\\
%
\>[2]\AgdaSymbol{→}\AgdaSpace{}%
\AgdaInductiveConstructor{box}\AgdaSpace{}%
\AgdaSymbol{(}\AgdaFunction{extsb}\AgdaSpace{}%
\AgdaSymbol{(}\AgdaFunction{fb}\AgdaSpace{}%
\AgdaSymbol{(}\AgdaFunction{embedValue}\AgdaSpace{}%
\AgdaBound{V}\AgdaSpace{}%
\AgdaBound{y}\AgdaSymbol{))}\AgdaSpace{}%
\AgdaSymbol{(}\AgdaFunction{embedLookup}\AgdaSpace{}%
\AgdaBound{x}\AgdaSymbol{))}\AgdaSpace{}%
\AgdaOperator{\AgdaDatatype{≡}}\AgdaSpace{}%
\AgdaFunction{embedTerm}\AgdaSpace{}%
\AgdaSymbol{(}\AgdaFunction{exts}\AgdaSpace{}%
\AgdaSymbol{(}\AgdaFunction{f}\AgdaSpace{}%
\AgdaBound{V}\AgdaSymbol{)}\AgdaSpace{}%
\AgdaBound{x}\AgdaSymbol{)}\<%
\\
\>[0]\AgdaFunction{exts-pres-subst-func}\AgdaSpace{}%
\AgdaBound{V}\AgdaSpace{}%
\AgdaBound{y}\AgdaSpace{}%
\AgdaInductiveConstructor{Z}\AgdaSpace{}%
\AgdaSymbol{=}\AgdaSpace{}%
\AgdaInductiveConstructor{refl}\<%
\\
\>[0]\AgdaFunction{exts-pres-subst-func}\AgdaSpace{}%
\AgdaBound{V}\AgdaSpace{}%
\AgdaBound{y}\AgdaSpace{}%
\AgdaSymbol{(}\AgdaInductiveConstructor{S}\AgdaSpace{}%
\AgdaBound{x}\AgdaSymbol{)}\AgdaSpace{}%
\AgdaSymbol{=}\<%
\\
\>[0][@{}l@{\AgdaIndent{0}}]%
\>[2]\AgdaOperator{\AgdaFunction{beginEq}}\<%
\\
\>[2][@{}l@{\AgdaIndent{0}}]%
\>[4]\AgdaInductiveConstructor{box}\AgdaSpace{}%
\AgdaSymbol{(}\AgdaFunction{renameb}\AgdaSpace{}%
\AgdaInductiveConstructor{Sb}\AgdaSpace{}%
\AgdaSymbol{(}\AgdaFunction{fb}\AgdaSpace{}%
\AgdaSymbol{(}\AgdaFunction{embedValue}\AgdaSpace{}%
\AgdaBound{V}\AgdaSpace{}%
\AgdaBound{y}\AgdaSymbol{)}\AgdaSpace{}%
\AgdaSymbol{(}\AgdaFunction{embedLookup}\AgdaSpace{}%
\AgdaBound{x}\AgdaSymbol{)))}\<%
\\
%
\>[2]\AgdaFunction{≡⟨}\AgdaSpace{}%
\AgdaInductiveConstructor{refl}\AgdaSpace{}%
\AgdaFunction{⟩}\<%
\\
\>[2][@{}l@{\AgdaIndent{0}}]%
\>[4]\AgdaFunction{renameb}\AgdaSpace{}%
\AgdaInductiveConstructor{Sb}\AgdaSpace{}%
\AgdaSymbol{(}\AgdaInductiveConstructor{box}\AgdaSpace{}%
\AgdaSymbol{(}\AgdaFunction{fb}\AgdaSpace{}%
\AgdaSymbol{(}\AgdaFunction{embedValue}\AgdaSpace{}%
\AgdaBound{V}\AgdaSpace{}%
\AgdaBound{y}\AgdaSymbol{)}\AgdaSpace{}%
\AgdaSymbol{(}\AgdaFunction{embedLookup}\AgdaSpace{}%
\AgdaBound{x}\AgdaSymbol{)))}\<%
\\
%
\>[2]\AgdaFunction{≡⟨}\AgdaSpace{}%
\AgdaFunction{cong}\AgdaSpace{}%
\AgdaSymbol{(}\AgdaFunction{renameb}\AgdaSpace{}%
\AgdaInductiveConstructor{Sb}\AgdaSymbol{)}\AgdaSpace{}%
\AgdaSymbol{(}\AgdaFunction{pres-subst-func}\AgdaSpace{}%
\AgdaBound{V}\AgdaSpace{}%
\AgdaBound{y}\AgdaSpace{}%
\AgdaBound{x}\AgdaSymbol{)}\AgdaSpace{}%
\AgdaFunction{⟩}\<%
\\
\>[2][@{}l@{\AgdaIndent{0}}]%
\>[4]\AgdaFunction{renameb}\AgdaSpace{}%
\AgdaInductiveConstructor{Sb}\AgdaSpace{}%
\AgdaSymbol{(}\AgdaFunction{embedTerm}\AgdaSpace{}%
\AgdaSymbol{(}\AgdaSpace{}%
\AgdaSymbol{(}\AgdaFunction{f}\AgdaSpace{}%
\AgdaBound{V}\AgdaSymbol{)}\AgdaSpace{}%
\AgdaBound{x}\AgdaSpace{}%
\AgdaSymbol{))}\<%
\\
%
\>[2]\AgdaFunction{≡⟨}\AgdaSpace{}%
\AgdaFunction{pres-rename}\AgdaSpace{}%
\AgdaInductiveConstructor{S}\AgdaSpace{}%
\AgdaInductiveConstructor{Sb}\AgdaSpace{}%
\AgdaSymbol{(λ}\AgdaSpace{}%
\AgdaBound{x₁}\AgdaSpace{}%
\AgdaSymbol{→}\AgdaSpace{}%
\AgdaInductiveConstructor{refl}\AgdaSymbol{)}\AgdaSpace{}%
\AgdaSymbol{(}\AgdaFunction{f}\AgdaSpace{}%
\AgdaBound{V}\AgdaSpace{}%
\AgdaBound{x}\AgdaSymbol{)}\AgdaSpace{}%
\AgdaFunction{⟩}\<%
\\
\>[2][@{}l@{\AgdaIndent{0}}]%
\>[4]\AgdaFunction{embedTerm}\AgdaSpace{}%
\AgdaSymbol{(}\AgdaFunction{rename}\AgdaSpace{}%
\AgdaInductiveConstructor{S}\AgdaSpace{}%
\AgdaSymbol{(}\AgdaFunction{f}\AgdaSpace{}%
\AgdaBound{V}\AgdaSpace{}%
\AgdaBound{x}\AgdaSymbol{))}\<%
\\
%
\>[2]\AgdaOperator{\AgdaFunction{∎eq}}\<%
\\
%
\\[\AgdaEmptyExtraSkip]%
\>[0]\AgdaComment{--\ Preservation\ and\ Reflection\ of\ }\<%
\\
\>[0]\AgdaComment{--\ single\ substitution\ of\ Gödel's\ translation}\<%
\\
\>[0]\AgdaFunction{pres-subst}\AgdaSpace{}%
\AgdaSymbol{:}\AgdaSpace{}%
\AgdaSymbol{∀}\AgdaSpace{}%
\AgdaSymbol{\{}\AgdaBound{Γ}\AgdaSpace{}%
\AgdaBound{A}\AgdaSpace{}%
\AgdaBound{B}\AgdaSymbol{\}}\AgdaSpace{}%
\AgdaSymbol{→}\AgdaSpace{}%
\AgdaSymbol{(}\AgdaBound{M}\AgdaSpace{}%
\AgdaSymbol{:}\AgdaSpace{}%
\AgdaBound{Γ}\AgdaSpace{}%
\AgdaOperator{\AgdaInductiveConstructor{,}}\AgdaSpace{}%
\AgdaBound{B}\AgdaSpace{}%
\AgdaOperator{\AgdaDatatype{⊢}}\AgdaSpace{}%
\AgdaBound{A}\AgdaSymbol{)}\AgdaSpace{}%
\AgdaSymbol{→}\AgdaSpace{}%
\AgdaSymbol{(}\AgdaBound{V}\AgdaSpace{}%
\AgdaSymbol{:}\AgdaSpace{}%
\AgdaBound{Γ}\AgdaSpace{}%
\AgdaOperator{\AgdaDatatype{⊢}}\AgdaSpace{}%
\AgdaBound{B}\AgdaSymbol{)}\AgdaSpace{}%
\AgdaSymbol{→}\AgdaSpace{}%
\AgdaSymbol{(}\AgdaBound{x}\AgdaSpace{}%
\AgdaSymbol{:}\AgdaSpace{}%
\AgdaDatatype{Value}\AgdaSpace{}%
\AgdaBound{V}\AgdaSymbol{)}\AgdaSpace{}%
\AgdaSymbol{→}\AgdaSpace{}%
\AgdaFunction{embedTerm}\AgdaSpace{}%
\AgdaBound{M}\AgdaSpace{}%
\AgdaOperator{\AgdaFunction{[}}\AgdaSpace{}%
\AgdaFunction{embedValue}\AgdaSpace{}%
\AgdaBound{V}\AgdaSpace{}%
\AgdaBound{x}\AgdaSpace{}%
\AgdaOperator{\AgdaFunction{]b}}\AgdaSpace{}%
\AgdaOperator{\AgdaDatatype{≡}}\AgdaSpace{}%
\AgdaFunction{embedTerm}\AgdaSpace{}%
\AgdaSymbol{(}\AgdaBound{M}\AgdaSpace{}%
\AgdaOperator{\AgdaFunction{[}}\AgdaSpace{}%
\AgdaBound{V}\AgdaSpace{}%
\AgdaOperator{\AgdaFunction{]}}\AgdaSymbol{)}\<%
\\
\>[0]\AgdaFunction{pres-subst}\AgdaSpace{}%
\AgdaBound{M}\AgdaSpace{}%
\AgdaBound{V}\AgdaSpace{}%
\AgdaBound{x}\AgdaSpace{}%
\AgdaSymbol{=}\AgdaSpace{}%
\AgdaFunction{pres-simultaneous-subst}\AgdaSpace{}%
\AgdaSymbol{(}\AgdaFunction{fb}\AgdaSpace{}%
\AgdaSymbol{(}\AgdaFunction{embedValue}\AgdaSpace{}%
\AgdaBound{V}\AgdaSpace{}%
\AgdaBound{x}\AgdaSymbol{))}\AgdaSpace{}%
\AgdaSymbol{(}\AgdaFunction{f}\AgdaSpace{}%
\AgdaBound{V}\AgdaSymbol{)}\AgdaSpace{}%
\AgdaSymbol{(}\AgdaFunction{pres-subst-func}\AgdaSpace{}%
\AgdaBound{V}\AgdaSpace{}%
\AgdaBound{x}\AgdaSymbol{)}\AgdaSpace{}%
\AgdaBound{M}\<%
\\
%
\\[\AgdaEmptyExtraSkip]%
\>[0]\AgdaComment{--\ Helper\ lemma:\ if\ a\ value\ is\ substituted\ in\ a\ value,\ the\ resulting\ lambda\ term\ will\ also\ be\ a\ value}\<%
\\
\>[0]\AgdaFunction{val-after-subst}\AgdaSpace{}%
\AgdaSymbol{:}\AgdaSpace{}%
\AgdaSymbol{∀}\AgdaSpace{}%
\AgdaSymbol{\{}\AgdaBound{Γ}\AgdaSpace{}%
\AgdaBound{A}\AgdaSpace{}%
\AgdaBound{B}\AgdaSymbol{\}}\AgdaSpace{}%
\AgdaSymbol{→}\AgdaSpace{}%
\AgdaSymbol{(}\AgdaBound{M}\AgdaSpace{}%
\AgdaSymbol{:}\AgdaSpace{}%
\AgdaBound{Γ}\AgdaSpace{}%
\AgdaOperator{\AgdaInductiveConstructor{,}}\AgdaSpace{}%
\AgdaBound{B}\AgdaSpace{}%
\AgdaOperator{\AgdaDatatype{⊢}}\AgdaSpace{}%
\AgdaBound{A}\AgdaSymbol{)}\AgdaSpace{}%
\AgdaSymbol{→}\AgdaSpace{}%
\AgdaSymbol{(}\AgdaBound{N}\AgdaSpace{}%
\AgdaSymbol{:}\AgdaSpace{}%
\AgdaBound{Γ}\AgdaSpace{}%
\AgdaOperator{\AgdaDatatype{⊢}}\AgdaSpace{}%
\AgdaBound{B}\AgdaSymbol{)}\AgdaSpace{}%
\AgdaSymbol{→}\AgdaSpace{}%
\AgdaSymbol{(}\AgdaBound{x}\AgdaSpace{}%
\AgdaSymbol{:}\AgdaSpace{}%
\AgdaDatatype{Value}\AgdaSpace{}%
\AgdaBound{M}\AgdaSymbol{)}\AgdaSpace{}%
\AgdaSymbol{→}\AgdaSpace{}%
\AgdaSymbol{(}\AgdaBound{y}\AgdaSpace{}%
\AgdaSymbol{:}\AgdaSpace{}%
\AgdaDatatype{Value}\AgdaSpace{}%
\AgdaBound{N}\AgdaSymbol{)}\AgdaSpace{}%
\AgdaSymbol{→}\AgdaSpace{}%
\AgdaDatatype{Value}\AgdaSpace{}%
\AgdaSymbol{(}\AgdaBound{M}\AgdaSpace{}%
\AgdaOperator{\AgdaFunction{[}}\AgdaSpace{}%
\AgdaBound{N}\AgdaSpace{}%
\AgdaOperator{\AgdaFunction{]}}\AgdaSymbol{)}\<%
\\
\>[0]\AgdaFunction{val-after-subst}\AgdaSpace{}%
\AgdaSymbol{(}\AgdaInductiveConstructor{`}\AgdaSpace{}%
\AgdaInductiveConstructor{Z}\AgdaSymbol{)}\AgdaSpace{}%
\AgdaBound{N}\AgdaSpace{}%
\AgdaBound{x}\AgdaSpace{}%
\AgdaBound{y}\AgdaSpace{}%
\AgdaSymbol{=}\AgdaSpace{}%
\AgdaBound{y}\<%
\\
\>[0]\AgdaFunction{val-after-subst}\AgdaSpace{}%
\AgdaSymbol{(}\AgdaInductiveConstructor{`}\AgdaSpace{}%
\AgdaSymbol{(}\AgdaInductiveConstructor{S}\AgdaSpace{}%
\AgdaBound{x₁}\AgdaSymbol{))}\AgdaSpace{}%
\AgdaBound{N}\AgdaSpace{}%
\AgdaBound{x}\AgdaSpace{}%
\AgdaBound{y}\AgdaSpace{}%
\AgdaSymbol{=}\AgdaSpace{}%
\AgdaInductiveConstructor{V-`}\<%
\\
\>[0]\AgdaFunction{val-after-subst}\AgdaSpace{}%
\AgdaSymbol{(}\AgdaInductiveConstructor{ƛ}\AgdaSpace{}%
\AgdaBound{M}\AgdaSymbol{)}\AgdaSpace{}%
\AgdaBound{N}\AgdaSpace{}%
\AgdaBound{x}\AgdaSpace{}%
\AgdaBound{y}\AgdaSpace{}%
\AgdaSymbol{=}\AgdaSpace{}%
\AgdaInductiveConstructor{V-ƛ}\<%
\\
%
\\[\AgdaEmptyExtraSkip]%
\>[0]\AgdaComment{--\ Preservation\ and\ Reflection\ of\ }\<%
\\
\>[0]\AgdaComment{--\ single\ substitution\ in\ values\ of\ Gödel's\ translation\ }\<%
\\
\>[0]\AgdaFunction{pres-subst-val}\AgdaSpace{}%
\AgdaSymbol{:}\AgdaSpace{}%
\AgdaSymbol{∀}\AgdaSpace{}%
\AgdaSymbol{\{}\AgdaBound{Γ}\AgdaSpace{}%
\AgdaBound{A}\AgdaSpace{}%
\AgdaBound{B}\AgdaSymbol{\}}\<%
\\
\>[0][@{}l@{\AgdaIndent{0}}]%
\>[2]\AgdaSymbol{→}\AgdaSpace{}%
\AgdaSymbol{(}\AgdaBound{M}\AgdaSpace{}%
\AgdaSymbol{:}\AgdaSpace{}%
\AgdaBound{Γ}\AgdaSpace{}%
\AgdaOperator{\AgdaInductiveConstructor{,}}\AgdaSpace{}%
\AgdaBound{B}\AgdaSpace{}%
\AgdaOperator{\AgdaDatatype{⊢}}\AgdaSpace{}%
\AgdaBound{A}\AgdaSymbol{)}\<%
\\
%
\>[2]\AgdaSymbol{→}\AgdaSpace{}%
\AgdaSymbol{(}\AgdaBound{V}\AgdaSpace{}%
\AgdaSymbol{:}\AgdaSpace{}%
\AgdaBound{Γ}\AgdaSpace{}%
\AgdaOperator{\AgdaDatatype{⊢}}\AgdaSpace{}%
\AgdaBound{B}\AgdaSymbol{)}\<%
\\
%
\>[2]\AgdaSymbol{→}\AgdaSpace{}%
\AgdaSymbol{(}\AgdaBound{x}\AgdaSpace{}%
\AgdaSymbol{:}\AgdaSpace{}%
\AgdaDatatype{Value}\AgdaSpace{}%
\AgdaBound{M}\AgdaSymbol{)}\<%
\\
%
\>[2]\AgdaSymbol{→}\AgdaSpace{}%
\AgdaSymbol{(}\AgdaBound{y}\AgdaSpace{}%
\AgdaSymbol{:}\AgdaSpace{}%
\AgdaDatatype{Value}\AgdaSpace{}%
\AgdaBound{V}\AgdaSymbol{)}\<%
\\
%
\>[2]\AgdaSymbol{→}\AgdaSpace{}%
\AgdaFunction{embedValue}\AgdaSpace{}%
\AgdaBound{M}\AgdaSpace{}%
\AgdaBound{x}\AgdaSpace{}%
\AgdaOperator{\AgdaFunction{[}}\AgdaSpace{}%
\AgdaFunction{embedValue}\AgdaSpace{}%
\AgdaBound{V}\AgdaSpace{}%
\AgdaBound{y}\AgdaSpace{}%
\AgdaOperator{\AgdaFunction{]b}}\AgdaSpace{}%
\AgdaOperator{\AgdaDatatype{≡}}\AgdaSpace{}%
\AgdaFunction{embedValue}\AgdaSpace{}%
\AgdaSymbol{(}\AgdaBound{M}\AgdaSpace{}%
\AgdaOperator{\AgdaFunction{[}}\AgdaSpace{}%
\AgdaBound{V}\AgdaSpace{}%
\AgdaOperator{\AgdaFunction{]}}\AgdaSymbol{)}\AgdaSpace{}%
\AgdaSymbol{(}\AgdaFunction{val-after-subst}\AgdaSpace{}%
\AgdaBound{M}\AgdaSpace{}%
\AgdaBound{V}\AgdaSpace{}%
\AgdaBound{x}\AgdaSpace{}%
\AgdaBound{y}\AgdaSymbol{)}\<%
\\
\>[0]\AgdaFunction{pres-subst-val}\AgdaSpace{}%
\AgdaSymbol{(}\AgdaInductiveConstructor{`}\AgdaSpace{}%
\AgdaInductiveConstructor{Z}\AgdaSymbol{)}\AgdaSpace{}%
\AgdaBound{V}\AgdaSpace{}%
\AgdaBound{x}\AgdaSpace{}%
\AgdaBound{y}\AgdaSpace{}%
\AgdaSymbol{=}\AgdaSpace{}%
\AgdaInductiveConstructor{refl}\<%
\\
\>[0]\AgdaFunction{pres-subst-val}\AgdaSpace{}%
\AgdaSymbol{(}\AgdaInductiveConstructor{`}\AgdaSpace{}%
\AgdaSymbol{(}\AgdaInductiveConstructor{S}\AgdaSpace{}%
\AgdaBound{x₁}\AgdaSymbol{))}\AgdaSpace{}%
\AgdaBound{V}\AgdaSpace{}%
\AgdaBound{x}\AgdaSpace{}%
\AgdaBound{y}\AgdaSpace{}%
\AgdaSymbol{=}\AgdaSpace{}%
\AgdaInductiveConstructor{refl}\<%
\\
\>[0]\AgdaFunction{pres-subst-val}\AgdaSpace{}%
\AgdaSymbol{(}\AgdaInductiveConstructor{ƛ}\AgdaSpace{}%
\AgdaBound{M}\AgdaSymbol{)}\AgdaSpace{}%
\AgdaBound{V}\AgdaSpace{}%
\AgdaBound{x}\AgdaSpace{}%
\AgdaBound{y}\AgdaSpace{}%
\AgdaSymbol{=}\<%
\\
\>[0][@{}l@{\AgdaIndent{0}}]%
\>[2]\AgdaOperator{\AgdaFunction{beginEq}}\<%
\\
\>[2][@{}l@{\AgdaIndent{0}}]%
\>[4]\AgdaFunction{embedValue}\AgdaSpace{}%
\AgdaSymbol{(}\AgdaInductiveConstructor{ƛ}\AgdaSpace{}%
\AgdaBound{M}\AgdaSymbol{)}\AgdaSpace{}%
\AgdaBound{x}\AgdaSpace{}%
\AgdaOperator{\AgdaFunction{[}}\AgdaSpace{}%
\AgdaFunction{embedValue}\AgdaSpace{}%
\AgdaBound{V}\AgdaSpace{}%
\AgdaBound{y}\AgdaSpace{}%
\AgdaOperator{\AgdaFunction{]b}}\<%
\\
%
\>[2]\AgdaFunction{≡⟨}\AgdaSpace{}%
\AgdaInductiveConstructor{refl}\AgdaSpace{}%
\AgdaFunction{⟩}\<%
\\
\>[2][@{}l@{\AgdaIndent{0}}]%
\>[4]\AgdaInductiveConstructor{ƛb}\AgdaSpace{}%
\AgdaSymbol{(}\AgdaFunction{substb}\AgdaSpace{}%
\AgdaSymbol{(}\AgdaFunction{extsb}\AgdaSpace{}%
\AgdaSymbol{(}\AgdaFunction{fb}\AgdaSpace{}%
\AgdaSymbol{(}\AgdaFunction{embedValue}\AgdaSpace{}%
\AgdaBound{V}\AgdaSpace{}%
\AgdaBound{y}\AgdaSymbol{)))}\AgdaSpace{}%
\AgdaSymbol{(}\AgdaFunction{embedTerm}\AgdaSpace{}%
\AgdaBound{M}\AgdaSymbol{))}\<%
\\
%
\>[2]\AgdaFunction{≡⟨}\AgdaSpace{}%
\AgdaFunction{cong}\AgdaSpace{}%
\AgdaInductiveConstructor{ƛb}\AgdaSpace{}%
\AgdaSymbol{(}\AgdaFunction{pres-simultaneous-subst}\AgdaSpace{}%
\AgdaSymbol{(}\AgdaFunction{extsb}\AgdaSpace{}%
\AgdaSymbol{(}\AgdaFunction{fb}\AgdaSpace{}%
\AgdaSymbol{(}\AgdaFunction{embedValue}\AgdaSpace{}%
\AgdaBound{V}\AgdaSpace{}%
\AgdaBound{y}\AgdaSymbol{)))}\AgdaSpace{}%
\AgdaSymbol{(}\AgdaFunction{exts}\AgdaSpace{}%
\AgdaSymbol{(}\AgdaFunction{f}\AgdaSpace{}%
\AgdaBound{V}\AgdaSymbol{))}\AgdaSpace{}%
\AgdaSymbol{(}\AgdaFunction{exts-pres-subst-func}\AgdaSpace{}%
\AgdaBound{V}\AgdaSpace{}%
\AgdaBound{y}\AgdaSymbol{)}\AgdaSpace{}%
\AgdaBound{M}\AgdaSymbol{)}\AgdaSpace{}%
\AgdaFunction{⟩}\<%
\\
\>[2][@{}l@{\AgdaIndent{0}}]%
\>[4]\AgdaInductiveConstructor{ƛb}\AgdaSpace{}%
\AgdaSymbol{(}\AgdaFunction{embedTerm}\AgdaSpace{}%
\AgdaSymbol{(}\AgdaFunction{subst}\AgdaSpace{}%
\AgdaSymbol{(}\AgdaFunction{exts}\AgdaSpace{}%
\AgdaSymbol{(}\AgdaFunction{f}\AgdaSpace{}%
\AgdaBound{V}\AgdaSymbol{))}\AgdaSpace{}%
\AgdaBound{M}\AgdaSymbol{))}\<%
\\
%
\>[2]\AgdaFunction{≡⟨}\AgdaSpace{}%
\AgdaInductiveConstructor{refl}\AgdaSpace{}%
\AgdaFunction{⟩}\<%
\\
\>[2][@{}l@{\AgdaIndent{0}}]%
\>[4]\AgdaFunction{embedValue}\AgdaSpace{}%
\AgdaSymbol{((}\AgdaInductiveConstructor{ƛ}\AgdaSpace{}%
\AgdaBound{M}\AgdaSymbol{)}\AgdaSpace{}%
\AgdaOperator{\AgdaFunction{[}}\AgdaSpace{}%
\AgdaBound{V}\AgdaSpace{}%
\AgdaOperator{\AgdaFunction{]}}\AgdaSymbol{)}\AgdaSpace{}%
\AgdaSymbol{(}\AgdaFunction{val-after-subst}\AgdaSpace{}%
\AgdaSymbol{(}\AgdaInductiveConstructor{ƛ}\AgdaSpace{}%
\AgdaBound{M}\AgdaSymbol{)}\AgdaSpace{}%
\AgdaBound{V}\AgdaSpace{}%
\AgdaBound{x}\AgdaSpace{}%
\AgdaBound{y}\AgdaSymbol{)}\<%
\\
%
\>[2]\AgdaOperator{\AgdaFunction{∎eq}}\<%
\\
%
\\[\AgdaEmptyExtraSkip]%
\>[0]\AgdaComment{--\ fb\ and\ renameb\ commute\ under\ extended\ context.}\<%
\\
\>[0]\AgdaFunction{fb-renameb-extb-commute}\AgdaSpace{}%
\AgdaSymbol{:}\AgdaSpace{}%
\AgdaSymbol{∀}\AgdaSpace{}%
\AgdaSymbol{\{}\AgdaBound{Γ}\AgdaSpace{}%
\AgdaBound{A}\AgdaSpace{}%
\AgdaBound{B}\AgdaSpace{}%
\AgdaBound{Δ}\AgdaSymbol{\}}\<%
\\
\>[0][@{}l@{\AgdaIndent{0}}]%
\>[2]\AgdaSymbol{→}\AgdaSpace{}%
\AgdaSymbol{(}\AgdaBound{fb'}\AgdaSpace{}%
\AgdaSymbol{:}\AgdaSpace{}%
\AgdaSymbol{∀}\AgdaSpace{}%
\AgdaSymbol{\{}\AgdaBound{Ab}\AgdaSymbol{\}}\AgdaSpace{}%
\AgdaSymbol{→}\AgdaSpace{}%
\AgdaBound{Ab}\AgdaSpace{}%
\AgdaOperator{\AgdaDatatype{∈b}}\AgdaSpace{}%
\AgdaBound{Γ}\AgdaSpace{}%
\AgdaSymbol{→}\AgdaSpace{}%
\AgdaBound{Ab}\AgdaSpace{}%
\AgdaOperator{\AgdaDatatype{∈b}}\AgdaSpace{}%
\AgdaBound{Δ}\AgdaSymbol{)}\<%
\\
%
\>[2]\AgdaSymbol{→}\AgdaSpace{}%
\AgdaSymbol{(}\AgdaBound{W}\AgdaSpace{}%
\AgdaSymbol{:}\AgdaSpace{}%
\AgdaBound{Γ}\AgdaSpace{}%
\AgdaOperator{\AgdaDatatype{⊢b}}\AgdaSpace{}%
\AgdaBound{B}\AgdaSymbol{)}\<%
\\
%
\>[2]\AgdaSymbol{→}\AgdaSpace{}%
\AgdaSymbol{(}\AgdaBound{x}\AgdaSpace{}%
\AgdaSymbol{:}\AgdaSpace{}%
\AgdaInductiveConstructor{□}\AgdaSpace{}%
\AgdaBound{A}\AgdaSpace{}%
\AgdaOperator{\AgdaDatatype{∈b}}\AgdaSpace{}%
\AgdaBound{Γ}\AgdaSpace{}%
\AgdaOperator{\AgdaInductiveConstructor{,b}}\AgdaSpace{}%
\AgdaInductiveConstructor{□}\AgdaSpace{}%
\AgdaBound{B}\AgdaSymbol{)}\<%
\\
%
\>[2]\AgdaSymbol{→}\AgdaSpace{}%
\AgdaFunction{fb}\AgdaSpace{}%
\AgdaSymbol{(}\AgdaFunction{renameb}\AgdaSpace{}%
\AgdaBound{fb'}\AgdaSpace{}%
\AgdaBound{W}\AgdaSymbol{)}\AgdaSpace{}%
\AgdaSymbol{(}\AgdaFunction{extb}\AgdaSpace{}%
\AgdaBound{fb'}\AgdaSpace{}%
\AgdaBound{x}\AgdaSymbol{)}\AgdaSpace{}%
\AgdaOperator{\AgdaDatatype{≡}}\AgdaSpace{}%
\AgdaFunction{renameb}\AgdaSpace{}%
\AgdaBound{fb'}\AgdaSpace{}%
\AgdaSymbol{(}\AgdaFunction{fb}\AgdaSpace{}%
\AgdaBound{W}\AgdaSpace{}%
\AgdaBound{x}\AgdaSymbol{)}\<%
\\
\>[0]\AgdaFunction{fb-renameb-extb-commute}\AgdaSpace{}%
\AgdaBound{fb'}\AgdaSpace{}%
\AgdaBound{W}\AgdaSpace{}%
\AgdaInductiveConstructor{Zb}\AgdaSpace{}%
\AgdaSymbol{=}\AgdaSpace{}%
\AgdaInductiveConstructor{refl}\<%
\\
\>[0]\AgdaFunction{fb-renameb-extb-commute}\AgdaSpace{}%
\AgdaBound{fb'}\AgdaSpace{}%
\AgdaBound{W}\AgdaSpace{}%
\AgdaSymbol{(}\AgdaInductiveConstructor{Sb}\AgdaSpace{}%
\AgdaBound{x}\AgdaSymbol{)}\AgdaSpace{}%
\AgdaSymbol{=}\AgdaSpace{}%
\AgdaInductiveConstructor{refl}\<%
\\
%
\\[\AgdaEmptyExtraSkip]%
\>[0]\AgdaComment{--\ Preservation\ and\ reflection\ of\ }\<%
\\
\>[0]\AgdaComment{--\ single\ substitution\ under\ renaming}\<%
\\
\>[0]\AgdaFunction{pres-subst-func-rename}\AgdaSpace{}%
\AgdaSymbol{:}\AgdaSpace{}%
\AgdaSymbol{∀}\AgdaSpace{}%
\AgdaSymbol{\{}\AgdaBound{Γ}\AgdaSpace{}%
\AgdaBound{B}\AgdaSpace{}%
\AgdaBound{A}\AgdaSpace{}%
\AgdaBound{Δ}\AgdaSymbol{\}}\<%
\\
\>[0][@{}l@{\AgdaIndent{0}}]%
\>[2]\AgdaSymbol{→}\AgdaSpace{}%
\AgdaSymbol{(}\AgdaBound{fb'}\AgdaSpace{}%
\AgdaSymbol{:}\AgdaSpace{}%
\AgdaSymbol{∀}\AgdaSpace{}%
\AgdaSymbol{\{}\AgdaBound{Ab}\AgdaSymbol{\}}\AgdaSpace{}%
\AgdaSymbol{→}\AgdaSpace{}%
\AgdaBound{Ab}\AgdaSpace{}%
\AgdaOperator{\AgdaDatatype{∈b}}\AgdaSpace{}%
\AgdaBound{Γ}\AgdaSpace{}%
\AgdaSymbol{→}\AgdaSpace{}%
\AgdaBound{Ab}\AgdaSpace{}%
\AgdaOperator{\AgdaDatatype{∈b}}\AgdaSpace{}%
\AgdaBound{Δ}\AgdaSymbol{)}\<%
\\
%
\>[2]\AgdaSymbol{→}\AgdaSpace{}%
\AgdaSymbol{(}\AgdaBound{M}\AgdaSpace{}%
\AgdaSymbol{:}\AgdaSpace{}%
\AgdaBound{Γ}\AgdaSpace{}%
\AgdaOperator{\AgdaInductiveConstructor{,b}}\AgdaSpace{}%
\AgdaInductiveConstructor{□}\AgdaSpace{}%
\AgdaBound{B}\AgdaSpace{}%
\AgdaOperator{\AgdaDatatype{⊢b}}\AgdaSpace{}%
\AgdaBound{A}\AgdaSymbol{)}\<%
\\
%
\>[2]\AgdaSymbol{→}\AgdaSpace{}%
\AgdaSymbol{(}\AgdaBound{W}\AgdaSpace{}%
\AgdaSymbol{:}\AgdaSpace{}%
\AgdaBound{Γ}\AgdaSpace{}%
\AgdaOperator{\AgdaDatatype{⊢b}}\AgdaSpace{}%
\AgdaBound{B}\AgdaSymbol{)}\<%
\\
%
\>[2]\AgdaSymbol{→}%
\>[2253I]\AgdaFunction{renameb}\AgdaSpace{}%
\AgdaSymbol{(}\AgdaFunction{extb}\AgdaSpace{}%
\AgdaBound{fb'}\AgdaSymbol{)}\AgdaSpace{}%
\AgdaBound{M}\AgdaSpace{}%
\AgdaOperator{\AgdaFunction{[}}\AgdaSpace{}%
\AgdaFunction{renameb}\AgdaSpace{}%
\AgdaBound{fb'}\AgdaSpace{}%
\AgdaBound{W}\AgdaSpace{}%
\AgdaOperator{\AgdaFunction{]b}}\<%
\\
\>[.][@{}l@{}]\<[2253I]%
\>[4]\AgdaOperator{\AgdaDatatype{≡}}\AgdaSpace{}%
\AgdaFunction{renameb}\AgdaSpace{}%
\AgdaBound{fb'}\AgdaSpace{}%
\AgdaSymbol{(}\AgdaBound{M}\AgdaSpace{}%
\AgdaOperator{\AgdaFunction{[}}\AgdaSpace{}%
\AgdaBound{W}\AgdaSpace{}%
\AgdaOperator{\AgdaFunction{]b}}\AgdaSymbol{)}\<%
\\
\>[0]\AgdaFunction{pres-subst-func-rename}\AgdaSpace{}%
\AgdaBound{fb'}\AgdaSpace{}%
\AgdaBound{M}\AgdaSpace{}%
\AgdaBound{W}\AgdaSpace{}%
\AgdaSymbol{=}\<%
\\
\>[0][@{}l@{\AgdaIndent{0}}]%
\>[2]\AgdaOperator{\AgdaFunction{beginEq}}\<%
\\
\>[2][@{}l@{\AgdaIndent{0}}]%
\>[4]\AgdaFunction{substb}\AgdaSpace{}%
\AgdaSymbol{(}\AgdaFunction{fb}\AgdaSpace{}%
\AgdaSymbol{(}\AgdaFunction{renameb}\AgdaSpace{}%
\AgdaBound{fb'}\AgdaSpace{}%
\AgdaBound{W}\AgdaSymbol{))}\AgdaSpace{}%
\AgdaSymbol{(}\AgdaFunction{renameb}\AgdaSpace{}%
\AgdaSymbol{(}\AgdaFunction{extb}\AgdaSpace{}%
\AgdaBound{fb'}\AgdaSymbol{)}\AgdaSpace{}%
\AgdaBound{M}\AgdaSymbol{)}\<%
\\
%
\>[2]\AgdaFunction{≡⟨}\AgdaSpace{}%
\AgdaFunction{subst-renameb-comp}\AgdaSpace{}%
\AgdaSymbol{(}\AgdaFunction{extb}\AgdaSpace{}%
\AgdaBound{fb'}\AgdaSymbol{)}\AgdaSpace{}%
\AgdaSymbol{(}\AgdaFunction{fb}\AgdaSpace{}%
\AgdaSymbol{(}\AgdaFunction{renameb}\AgdaSpace{}%
\AgdaBound{fb'}\AgdaSpace{}%
\AgdaBound{W}\AgdaSymbol{))}\AgdaSpace{}%
\AgdaBound{M}\AgdaSpace{}%
\AgdaFunction{⟩}\<%
\\
\>[2][@{}l@{\AgdaIndent{0}}]%
\>[4]\AgdaFunction{substb}\AgdaSpace{}%
\AgdaSymbol{(λ}\AgdaSpace{}%
\AgdaBound{x}\AgdaSpace{}%
\AgdaSymbol{→}\AgdaSpace{}%
\AgdaFunction{fb}\AgdaSpace{}%
\AgdaSymbol{(}\AgdaFunction{renameb}\AgdaSpace{}%
\AgdaBound{fb'}\AgdaSpace{}%
\AgdaBound{W}\AgdaSymbol{)}\AgdaSpace{}%
\AgdaSymbol{(}\AgdaFunction{extb}\AgdaSpace{}%
\AgdaBound{fb'}\AgdaSpace{}%
\AgdaBound{x}\AgdaSymbol{))}\AgdaSpace{}%
\AgdaBound{M}\<%
\\
%
\>[2]\AgdaFunction{≡⟨}\AgdaSpace{}%
\AgdaFunction{substb-cong}\AgdaSpace{}%
\AgdaSymbol{(}\AgdaFunction{fb-renameb-extb-commute}\AgdaSpace{}%
\AgdaBound{fb'}\AgdaSpace{}%
\AgdaBound{W}\AgdaSymbol{)}\AgdaSpace{}%
\AgdaBound{M}\AgdaSpace{}%
\AgdaFunction{⟩}\<%
\\
\>[2][@{}l@{\AgdaIndent{0}}]%
\>[4]\AgdaFunction{substb}\AgdaSpace{}%
\AgdaSymbol{(λ}\AgdaSpace{}%
\AgdaBound{x}\AgdaSpace{}%
\AgdaSymbol{→}\AgdaSpace{}%
\AgdaFunction{renameb}\AgdaSpace{}%
\AgdaBound{fb'}\AgdaSpace{}%
\AgdaSymbol{(}\AgdaFunction{fb}\AgdaSpace{}%
\AgdaBound{W}\AgdaSpace{}%
\AgdaBound{x}\AgdaSymbol{))}\AgdaSpace{}%
\AgdaBound{M}\<%
\\
%
\>[2]\AgdaFunction{≡⟨}\AgdaSpace{}%
\AgdaFunction{substb-rename-commute}\AgdaSpace{}%
\AgdaBound{fb'}\AgdaSpace{}%
\AgdaSymbol{(}\AgdaFunction{fb}\AgdaSpace{}%
\AgdaBound{W}\AgdaSymbol{)}\AgdaSpace{}%
\AgdaBound{M}\AgdaSpace{}%
\AgdaFunction{⟩}\<%
\\
\>[2][@{}l@{\AgdaIndent{0}}]%
\>[4]\AgdaFunction{renameb}\AgdaSpace{}%
\AgdaBound{fb'}\AgdaSpace{}%
\AgdaSymbol{(}\AgdaFunction{substb}\AgdaSpace{}%
\AgdaSymbol{(}\AgdaFunction{fb}\AgdaSpace{}%
\AgdaBound{W}\AgdaSymbol{)}\AgdaSpace{}%
\AgdaBound{M}\AgdaSymbol{)}\<%
\\
%
\>[2]\AgdaOperator{\AgdaFunction{∎eq}}\<%
\\
%
\\[\AgdaEmptyExtraSkip]%
\>[0]\AgdaComment{--\ Preservation\ of\ reduction\ under\ renaming\ \ \ }\<%
\\
\>[0]\AgdaFunction{pres-rename-red}\AgdaSpace{}%
\AgdaSymbol{:}\AgdaSpace{}%
\AgdaSymbol{∀}\AgdaSpace{}%
\AgdaSymbol{\{}\AgdaBound{Γ}\AgdaSpace{}%
\AgdaBound{A}\AgdaSpace{}%
\AgdaBound{Δ}\AgdaSymbol{\}}\<%
\\
\>[0][@{}l@{\AgdaIndent{0}}]%
\>[2]\AgdaSymbol{→}\AgdaSpace{}%
\AgdaSymbol{(}\AgdaBound{fb'}\AgdaSpace{}%
\AgdaSymbol{:}\AgdaSpace{}%
\AgdaSymbol{∀}\AgdaSpace{}%
\AgdaSymbol{\{}\AgdaBound{Ab}\AgdaSymbol{\}}\AgdaSpace{}%
\AgdaSymbol{→}\AgdaSpace{}%
\AgdaBound{Ab}\AgdaSpace{}%
\AgdaOperator{\AgdaDatatype{∈b}}\AgdaSpace{}%
\AgdaBound{Γ}\AgdaSpace{}%
\AgdaSymbol{→}\AgdaSpace{}%
\AgdaBound{Ab}\AgdaSpace{}%
\AgdaOperator{\AgdaDatatype{∈b}}\AgdaSpace{}%
\AgdaBound{Δ}\AgdaSymbol{)}\<%
\\
%
\>[2]\AgdaSymbol{→}\AgdaSpace{}%
\AgdaSymbol{(}\AgdaBound{M}\AgdaSpace{}%
\AgdaSymbol{:}\AgdaSpace{}%
\AgdaBound{Γ}\AgdaSpace{}%
\AgdaOperator{\AgdaDatatype{⊢b}}\AgdaSpace{}%
\AgdaBound{A}\AgdaSymbol{)}\<%
\\
%
\>[2]\AgdaSymbol{→}\AgdaSpace{}%
\AgdaSymbol{(}\AgdaBound{N}\AgdaSpace{}%
\AgdaSymbol{:}\AgdaSpace{}%
\AgdaBound{Γ}\AgdaSpace{}%
\AgdaOperator{\AgdaDatatype{⊢b}}\AgdaSpace{}%
\AgdaBound{A}\AgdaSymbol{)}\<%
\\
%
\>[2]\AgdaSymbol{→}\AgdaSpace{}%
\AgdaSymbol{(}\AgdaBound{M}\AgdaSpace{}%
\AgdaOperator{\AgdaDatatype{↝βb}}\AgdaSpace{}%
\AgdaBound{N}\AgdaSymbol{)}\<%
\\
%
\>[2]\AgdaSymbol{→}\AgdaSpace{}%
\AgdaFunction{renameb}\AgdaSpace{}%
\AgdaBound{fb'}\AgdaSpace{}%
\AgdaBound{M}\AgdaSpace{}%
\AgdaOperator{\AgdaDatatype{↝βb}}\AgdaSpace{}%
\AgdaFunction{renameb}\AgdaSpace{}%
\AgdaBound{fb'}\AgdaSpace{}%
\AgdaBound{N}\<%
\\
\>[0]\AgdaFunction{pres-rename-red}\AgdaSpace{}%
\AgdaBound{fb'}\AgdaSpace{}%
\AgdaSymbol{(}\AgdaInductiveConstructor{ƛb}\AgdaSpace{}%
\AgdaBound{M₁}\AgdaSpace{}%
\AgdaOperator{\AgdaInductiveConstructor{·b}}\AgdaSpace{}%
\AgdaInductiveConstructor{box}\AgdaSpace{}%
\AgdaBound{M₂}\AgdaSymbol{)}\AgdaSpace{}%
\AgdaBound{N}\AgdaSpace{}%
\AgdaInductiveConstructor{βb}\AgdaSpace{}%
\AgdaSymbol{=}\AgdaSpace{}%
\AgdaFunction{Eq.subst}\AgdaSpace{}%
\AgdaSymbol{(λ}\AgdaSpace{}%
\AgdaBound{x}\AgdaSpace{}%
\AgdaSymbol{→}\AgdaSpace{}%
\AgdaInductiveConstructor{ƛb}\AgdaSpace{}%
\AgdaSymbol{(}\AgdaFunction{renameb}\AgdaSpace{}%
\AgdaSymbol{(}\AgdaFunction{extb}\AgdaSpace{}%
\AgdaBound{fb'}\AgdaSymbol{)}\AgdaSpace{}%
\AgdaBound{M₁}\AgdaSymbol{)}\AgdaSpace{}%
\AgdaOperator{\AgdaInductiveConstructor{·b}}\AgdaSpace{}%
\AgdaInductiveConstructor{box}\AgdaSpace{}%
\AgdaSymbol{(}\AgdaFunction{renameb}\AgdaSpace{}%
\AgdaBound{fb'}\AgdaSpace{}%
\AgdaBound{M₂}\AgdaSymbol{)}\AgdaSpace{}%
\AgdaOperator{\AgdaDatatype{↝βb}}\AgdaSpace{}%
\AgdaBound{x}\AgdaSymbol{)}\AgdaSpace{}%
\AgdaSymbol{(}\AgdaFunction{pres-subst-func-rename}\AgdaSpace{}%
\AgdaBound{fb'}\AgdaSpace{}%
\AgdaBound{M₁}\AgdaSpace{}%
\AgdaBound{M₂}\AgdaSymbol{)}\AgdaSpace{}%
\AgdaInductiveConstructor{βb}\<%
\\
\>[0]\AgdaFunction{pres-rename-red}\AgdaSpace{}%
\AgdaBound{fb'}\AgdaSpace{}%
\AgdaBound{M}\AgdaSpace{}%
\AgdaBound{N}\AgdaSpace{}%
\AgdaSymbol{(}\AgdaInductiveConstructor{μ}\AgdaSpace{}%
\AgdaSymbol{\{}\AgdaArgument{M}\AgdaSpace{}%
\AgdaSymbol{=}\AgdaSpace{}%
\AgdaBound{M₁}\AgdaSymbol{\}}\AgdaSpace{}%
\AgdaSymbol{\{}\AgdaArgument{M'}\AgdaSpace{}%
\AgdaSymbol{=}\AgdaSpace{}%
\AgdaBound{M'}\AgdaSymbol{\}}\AgdaSpace{}%
\AgdaBound{red}\AgdaSymbol{)}\AgdaSpace{}%
\AgdaSymbol{=}\AgdaSpace{}%
\AgdaInductiveConstructor{μ}\AgdaSpace{}%
\AgdaSymbol{(}\AgdaFunction{pres-rename-red}\AgdaSpace{}%
\AgdaBound{fb'}\AgdaSpace{}%
\AgdaBound{M₁}\AgdaSpace{}%
\AgdaBound{M'}\AgdaSpace{}%
\AgdaBound{red}\AgdaSymbol{)}\<%
\\
\>[0]\AgdaFunction{pres-rename-red}\AgdaSpace{}%
\AgdaBound{fb'}\AgdaSpace{}%
\AgdaBound{M}\AgdaSpace{}%
\AgdaBound{N}\AgdaSpace{}%
\AgdaSymbol{(}\AgdaInductiveConstructor{ν}\AgdaSpace{}%
\AgdaSymbol{\{}\AgdaArgument{M}\AgdaSpace{}%
\AgdaSymbol{=}\AgdaSpace{}%
\AgdaBound{M₁}\AgdaSymbol{\}}\AgdaSpace{}%
\AgdaSymbol{\{}\AgdaArgument{N}\AgdaSpace{}%
\AgdaSymbol{=}\AgdaSpace{}%
\AgdaBound{N₁}\AgdaSymbol{\}}\AgdaSpace{}%
\AgdaSymbol{\{}\AgdaArgument{N'}\AgdaSpace{}%
\AgdaSymbol{=}\AgdaSpace{}%
\AgdaBound{N'}\AgdaSymbol{\}}\AgdaSpace{}%
\AgdaBound{red}\AgdaSymbol{)}\AgdaSpace{}%
\AgdaSymbol{=}\AgdaSpace{}%
\AgdaInductiveConstructor{ν}\AgdaSpace{}%
\AgdaSymbol{(}\AgdaFunction{pres-rename-red}\AgdaSpace{}%
\AgdaBound{fb'}\AgdaSpace{}%
\AgdaBound{N₁}\AgdaSpace{}%
\AgdaBound{N'}\AgdaSpace{}%
\AgdaBound{red}\AgdaSymbol{)}\<%
\\
\>[0]\AgdaFunction{pres-rename-red}\AgdaSpace{}%
\AgdaBound{fb'}\AgdaSpace{}%
\AgdaBound{M}\AgdaSpace{}%
\AgdaBound{N}\AgdaSpace{}%
\AgdaSymbol{(}\AgdaInductiveConstructor{ξ}\AgdaSpace{}%
\AgdaSymbol{\{}\AgdaArgument{M}\AgdaSpace{}%
\AgdaSymbol{=}\AgdaSpace{}%
\AgdaBound{M₁}\AgdaSymbol{\}}\AgdaSpace{}%
\AgdaSymbol{\{}\AgdaArgument{M'}\AgdaSpace{}%
\AgdaSymbol{=}\AgdaSpace{}%
\AgdaBound{M'}\AgdaSymbol{\}}\AgdaSpace{}%
\AgdaBound{red}\AgdaSymbol{)}\AgdaSpace{}%
\AgdaSymbol{=}\AgdaSpace{}%
\AgdaInductiveConstructor{ξ}\AgdaSpace{}%
\AgdaSymbol{(}\AgdaFunction{pres-rename-red}\AgdaSpace{}%
\AgdaSymbol{(}\AgdaFunction{extb}\AgdaSpace{}%
\AgdaBound{fb'}\AgdaSymbol{)}\AgdaSpace{}%
\AgdaBound{M₁}\AgdaSpace{}%
\AgdaBound{M'}\AgdaSpace{}%
\AgdaBound{red}\AgdaSymbol{)}\<%
\\
\>[0]\AgdaFunction{pres-rename-red}\AgdaSpace{}%
\AgdaBound{fb'}\AgdaSpace{}%
\AgdaBound{M}\AgdaSpace{}%
\AgdaBound{N}\AgdaSpace{}%
\AgdaSymbol{(}\AgdaInductiveConstructor{ζ}\AgdaSpace{}%
\AgdaSymbol{\{}\AgdaArgument{M}\AgdaSpace{}%
\AgdaSymbol{=}\AgdaSpace{}%
\AgdaBound{M₁}\AgdaSymbol{\}}\AgdaSpace{}%
\AgdaSymbol{\{}\AgdaArgument{M'}\AgdaSpace{}%
\AgdaSymbol{=}\AgdaSpace{}%
\AgdaBound{M'}\AgdaSymbol{\}}\AgdaSpace{}%
\AgdaBound{red}\AgdaSymbol{)}\AgdaSpace{}%
\AgdaSymbol{=}\AgdaSpace{}%
\AgdaInductiveConstructor{ζ}\AgdaSpace{}%
\AgdaSymbol{(}\AgdaFunction{pres-rename-red}\AgdaSpace{}%
\AgdaBound{fb'}\AgdaSpace{}%
\AgdaBound{M₁}\AgdaSpace{}%
\AgdaBound{M'}\AgdaSpace{}%
\AgdaBound{red}\AgdaSymbol{)}\<%
\\
%
\\[\AgdaEmptyExtraSkip]%
\>[0]\AgdaComment{--\ Preservation\ of\ 2\ step\ reduction}\<%
\\
\>[0]\AgdaComment{--\ under\ renaming}\<%
\\
\>[0]\AgdaFunction{pres-rename-red-2}\AgdaSpace{}%
\AgdaSymbol{:}\AgdaSpace{}%
\AgdaSymbol{∀}\AgdaSpace{}%
\AgdaSymbol{\{}\AgdaBound{Γ}\AgdaSpace{}%
\AgdaBound{A}\AgdaSpace{}%
\AgdaBound{Δ}\AgdaSymbol{\}}\<%
\\
\>[0][@{}l@{\AgdaIndent{0}}]%
\>[2]\AgdaSymbol{→}\AgdaSpace{}%
\AgdaSymbol{(}\AgdaBound{fb'}\AgdaSpace{}%
\AgdaSymbol{:}\AgdaSpace{}%
\AgdaSymbol{∀}\AgdaSpace{}%
\AgdaSymbol{\{}\AgdaBound{Ab}\AgdaSymbol{\}}\AgdaSpace{}%
\AgdaSymbol{→}\AgdaSpace{}%
\AgdaBound{Ab}\AgdaSpace{}%
\AgdaOperator{\AgdaDatatype{∈b}}\AgdaSpace{}%
\AgdaBound{Γ}\AgdaSpace{}%
\AgdaSymbol{→}\AgdaSpace{}%
\AgdaBound{Ab}\AgdaSpace{}%
\AgdaOperator{\AgdaDatatype{∈b}}\AgdaSpace{}%
\AgdaBound{Δ}\AgdaSymbol{)}\<%
\\
%
\>[2]\AgdaSymbol{→}\AgdaSpace{}%
\AgdaSymbol{(}\AgdaBound{M}\AgdaSpace{}%
\AgdaSymbol{:}\AgdaSpace{}%
\AgdaBound{Γ}\AgdaSpace{}%
\AgdaOperator{\AgdaDatatype{⊢b}}\AgdaSpace{}%
\AgdaBound{A}\AgdaSymbol{)}\<%
\\
%
\>[2]\AgdaSymbol{→}\AgdaSpace{}%
\AgdaSymbol{(}\AgdaBound{N}\AgdaSpace{}%
\AgdaSymbol{:}\AgdaSpace{}%
\AgdaBound{Γ}\AgdaSpace{}%
\AgdaOperator{\AgdaDatatype{⊢b}}\AgdaSpace{}%
\AgdaBound{A}\AgdaSymbol{)}\<%
\\
%
\>[2]\AgdaSymbol{→}\AgdaSpace{}%
\AgdaSymbol{(}\AgdaBound{M}\AgdaSpace{}%
\AgdaOperator{\AgdaDatatype{↝βb²}}\AgdaSpace{}%
\AgdaBound{N}\AgdaSymbol{)}\<%
\\
%
\>[2]\AgdaSymbol{→}\AgdaSpace{}%
\AgdaFunction{renameb}\AgdaSpace{}%
\AgdaBound{fb'}\AgdaSpace{}%
\AgdaBound{M}\AgdaSpace{}%
\AgdaOperator{\AgdaDatatype{↝βb²}}\AgdaSpace{}%
\AgdaFunction{renameb}\AgdaSpace{}%
\AgdaBound{fb'}\AgdaSpace{}%
\AgdaBound{N}\<%
\\
\>[0]\AgdaFunction{pres-rename-red-2}\AgdaSpace{}%
\AgdaBound{fb'}\AgdaSpace{}%
\AgdaBound{M}\AgdaSpace{}%
\AgdaBound{N}\AgdaSpace{}%
\AgdaSymbol{(}\AgdaInductiveConstructor{βb²}\AgdaSpace{}%
\AgdaDottedPattern{\AgdaSymbol{.}}\AgdaDottedPattern{\AgdaBound{M}}\AgdaSpace{}%
\AgdaBound{t₂}\AgdaSpace{}%
\AgdaDottedPattern{\AgdaSymbol{.}}\AgdaDottedPattern{\AgdaBound{N}}\AgdaSpace{}%
\AgdaBound{x}\AgdaSpace{}%
\AgdaBound{x₁}\AgdaSymbol{)}\AgdaSpace{}%
\AgdaSymbol{=}\AgdaSpace{}%
\AgdaInductiveConstructor{βb²}\AgdaSpace{}%
\AgdaSymbol{(}\AgdaFunction{renameb}\AgdaSpace{}%
\AgdaBound{fb'}\AgdaSpace{}%
\AgdaBound{M}\AgdaSymbol{)}\AgdaSpace{}%
\AgdaSymbol{(}\AgdaFunction{renameb}\AgdaSpace{}%
\AgdaBound{fb'}\AgdaSpace{}%
\AgdaBound{t₂}\AgdaSymbol{)}\AgdaSpace{}%
\AgdaSymbol{(}\AgdaFunction{renameb}\AgdaSpace{}%
\AgdaBound{fb'}\AgdaSpace{}%
\AgdaBound{N}\AgdaSymbol{)}\AgdaSpace{}%
\AgdaSymbol{(}\AgdaFunction{pres-rename-red}\AgdaSpace{}%
\AgdaBound{fb'}\AgdaSpace{}%
\AgdaBound{M}\AgdaSpace{}%
\AgdaBound{t₂}\AgdaSpace{}%
\AgdaBound{x}\AgdaSymbol{)}\AgdaSpace{}%
\AgdaSymbol{(}\AgdaFunction{pres-rename-red}\AgdaSpace{}%
\AgdaBound{fb'}\AgdaSpace{}%
\AgdaBound{t₂}\AgdaSpace{}%
\AgdaBound{N}\AgdaSpace{}%
\AgdaBound{x₁}\AgdaSymbol{)}\<%
\\
\>[0]\AgdaFunction{pres-rename-red-2}\AgdaSpace{}%
\AgdaBound{fb'}\AgdaSpace{}%
\AgdaBound{M}\AgdaSpace{}%
\AgdaBound{N}\AgdaSpace{}%
\AgdaSymbol{(}\AgdaInductiveConstructor{μ}\AgdaSpace{}%
\AgdaSymbol{\{}\AgdaArgument{M}\AgdaSpace{}%
\AgdaSymbol{=}\AgdaSpace{}%
\AgdaBound{M₁}\AgdaSymbol{\}}\AgdaSpace{}%
\AgdaSymbol{\{}\AgdaArgument{M'}\AgdaSpace{}%
\AgdaSymbol{=}\AgdaSpace{}%
\AgdaBound{M'}\AgdaSymbol{\}}\AgdaSpace{}%
\AgdaBound{red}\AgdaSymbol{)}\AgdaSpace{}%
\AgdaSymbol{=}\AgdaSpace{}%
\AgdaInductiveConstructor{μ}\AgdaSpace{}%
\AgdaSymbol{(}\AgdaFunction{pres-rename-red-2}\AgdaSpace{}%
\AgdaBound{fb'}\AgdaSpace{}%
\AgdaBound{M₁}\AgdaSpace{}%
\AgdaBound{M'}\AgdaSpace{}%
\AgdaBound{red}\AgdaSymbol{)}\<%
\\
\>[0]\AgdaFunction{pres-rename-red-2}\AgdaSpace{}%
\AgdaBound{fb'}\AgdaSpace{}%
\AgdaBound{M}\AgdaSpace{}%
\AgdaBound{N}\AgdaSpace{}%
\AgdaSymbol{(}\AgdaInductiveConstructor{ν}\AgdaSpace{}%
\AgdaSymbol{\{}\AgdaArgument{M}\AgdaSpace{}%
\AgdaSymbol{=}\AgdaSpace{}%
\AgdaBound{M₁}\AgdaSymbol{\}}\AgdaSpace{}%
\AgdaSymbol{\{}\AgdaArgument{N}\AgdaSpace{}%
\AgdaSymbol{=}\AgdaSpace{}%
\AgdaBound{N₁}\AgdaSymbol{\}}\AgdaSpace{}%
\AgdaSymbol{\{}\AgdaArgument{N'}\AgdaSpace{}%
\AgdaSymbol{=}\AgdaSpace{}%
\AgdaBound{N'}\AgdaSymbol{\}}\AgdaSpace{}%
\AgdaBound{red}\AgdaSymbol{)}\AgdaSpace{}%
\AgdaSymbol{=}\AgdaSpace{}%
\AgdaInductiveConstructor{ν}\AgdaSpace{}%
\AgdaSymbol{(}\AgdaFunction{pres-rename-red-2}\AgdaSpace{}%
\AgdaBound{fb'}\AgdaSpace{}%
\AgdaBound{N₁}\AgdaSpace{}%
\AgdaBound{N'}\AgdaSpace{}%
\AgdaBound{red}\AgdaSymbol{)}\<%
\\
\>[0]\AgdaFunction{pres-rename-red-2}\AgdaSpace{}%
\AgdaBound{fb'}\AgdaSpace{}%
\AgdaBound{M}\AgdaSpace{}%
\AgdaBound{N}\AgdaSpace{}%
\AgdaSymbol{(}\AgdaInductiveConstructor{ξ}\AgdaSpace{}%
\AgdaSymbol{\{}\AgdaArgument{M}\AgdaSpace{}%
\AgdaSymbol{=}\AgdaSpace{}%
\AgdaBound{M₁}\AgdaSymbol{\}}\AgdaSpace{}%
\AgdaSymbol{\{}\AgdaArgument{M'}\AgdaSpace{}%
\AgdaSymbol{=}\AgdaSpace{}%
\AgdaBound{M'}\AgdaSymbol{\}}\AgdaSpace{}%
\AgdaBound{red}\AgdaSymbol{)}\AgdaSpace{}%
\AgdaSymbol{=}\AgdaSpace{}%
\AgdaInductiveConstructor{ξ}\AgdaSpace{}%
\AgdaSymbol{(}\AgdaFunction{pres-rename-red-2}\AgdaSpace{}%
\AgdaSymbol{(}\AgdaFunction{extb}\AgdaSpace{}%
\AgdaBound{fb'}\AgdaSymbol{)}\AgdaSpace{}%
\AgdaBound{M₁}\AgdaSpace{}%
\AgdaBound{M'}\AgdaSpace{}%
\AgdaBound{red}\AgdaSymbol{)}\<%
\\
\>[0]\AgdaFunction{pres-rename-red-2}\AgdaSpace{}%
\AgdaBound{fb'}\AgdaSpace{}%
\AgdaBound{M}\AgdaSpace{}%
\AgdaBound{N}\AgdaSpace{}%
\AgdaSymbol{(}\AgdaInductiveConstructor{ζ}\AgdaSpace{}%
\AgdaSymbol{\{}\AgdaArgument{M}\AgdaSpace{}%
\AgdaSymbol{=}\AgdaSpace{}%
\AgdaBound{M₁}\AgdaSymbol{\}}\AgdaSpace{}%
\AgdaSymbol{\{}\AgdaArgument{M'}\AgdaSpace{}%
\AgdaSymbol{=}\AgdaSpace{}%
\AgdaBound{M'}\AgdaSymbol{\}}\AgdaSpace{}%
\AgdaBound{red}\AgdaSymbol{)}\AgdaSpace{}%
\AgdaSymbol{=}\AgdaSpace{}%
\AgdaInductiveConstructor{ζ}\AgdaSpace{}%
\AgdaSymbol{(}\AgdaFunction{pres-rename-red-2}\AgdaSpace{}%
\AgdaBound{fb'}\AgdaSpace{}%
\AgdaBound{M₁}\AgdaSpace{}%
\AgdaBound{M'}\AgdaSpace{}%
\AgdaBound{red}\AgdaSymbol{)}\<%
\\
%
\\[\AgdaEmptyExtraSkip]%
\>[0]\AgdaComment{--\ extsb\ of\ ε\ is\ ε}\<%
\\
\>[0]\AgdaFunction{ε-exts}\AgdaSpace{}%
\AgdaSymbol{:}\AgdaSpace{}%
\AgdaSymbol{∀}\AgdaSpace{}%
\AgdaSymbol{\{}\AgdaBound{Γ}\AgdaSpace{}%
\AgdaBound{A}\AgdaSpace{}%
\AgdaBound{B}\AgdaSymbol{\}}\<%
\\
\>[0][@{}l@{\AgdaIndent{0}}]%
\>[2]\AgdaSymbol{→}\AgdaSpace{}%
\AgdaSymbol{(}\AgdaBound{x}\AgdaSpace{}%
\AgdaSymbol{:}\AgdaSpace{}%
\AgdaInductiveConstructor{□}\AgdaSpace{}%
\AgdaBound{A}\AgdaSpace{}%
\AgdaOperator{\AgdaDatatype{∈b}}\AgdaSpace{}%
\AgdaBound{Γ}\AgdaSpace{}%
\AgdaOperator{\AgdaInductiveConstructor{,b}}\AgdaSpace{}%
\AgdaBound{B}\AgdaSymbol{)}\<%
\\
%
\>[2]\AgdaSymbol{→}\AgdaSpace{}%
\AgdaFunction{extsb}\AgdaSpace{}%
\AgdaInductiveConstructor{ε}\AgdaSpace{}%
\AgdaBound{x}\AgdaSpace{}%
\AgdaOperator{\AgdaDatatype{≡}}\AgdaSpace{}%
\AgdaInductiveConstructor{ε}\AgdaSpace{}%
\AgdaBound{x}\<%
\\
\>[0]\AgdaFunction{ε-exts}\AgdaSpace{}%
\AgdaInductiveConstructor{Zb}\AgdaSpace{}%
\AgdaSymbol{=}\AgdaSpace{}%
\AgdaInductiveConstructor{refl}\<%
\\
\>[0]\AgdaFunction{ε-exts}\AgdaSpace{}%
\AgdaSymbol{(}\AgdaInductiveConstructor{Sb}\AgdaSpace{}%
\AgdaBound{x}\AgdaSymbol{)}\AgdaSpace{}%
\AgdaSymbol{=}\AgdaSpace{}%
\AgdaInductiveConstructor{refl}\<%
\\
%
\\[\AgdaEmptyExtraSkip]%
\>[0]\AgdaComment{--\ Simultaneous\ substitution\ with\ ε\ }\<%
\\
\>[0]\AgdaComment{--\ leaves\ the\ term\ unchanged\ \ }\<%
\\
\>[0]\AgdaFunction{subst-with-id}\AgdaSpace{}%
\AgdaSymbol{:}\AgdaSpace{}%
\AgdaSymbol{∀}\AgdaSpace{}%
\AgdaSymbol{\{}\AgdaBound{Γ}\AgdaSpace{}%
\AgdaBound{A}\AgdaSymbol{\}}\<%
\\
\>[0][@{}l@{\AgdaIndent{0}}]%
\>[2]\AgdaSymbol{→}\AgdaSpace{}%
\AgdaSymbol{(}\AgdaBound{M}\AgdaSpace{}%
\AgdaSymbol{:}\AgdaSpace{}%
\AgdaBound{Γ}\AgdaSpace{}%
\AgdaOperator{\AgdaDatatype{⊢b}}\AgdaSpace{}%
\AgdaBound{A}\AgdaSymbol{)}\<%
\\
%
\>[2]\AgdaSymbol{→}\AgdaSpace{}%
\AgdaFunction{substb}\AgdaSpace{}%
\AgdaSymbol{(λ}\AgdaSpace{}%
\AgdaBound{x₁}\AgdaSpace{}%
\AgdaSymbol{→}\AgdaSpace{}%
\AgdaInductiveConstructor{ε}\AgdaSpace{}%
\AgdaBound{x₁}\AgdaSymbol{)}\AgdaSpace{}%
\AgdaBound{M}\AgdaSpace{}%
\AgdaOperator{\AgdaDatatype{≡}}\AgdaSpace{}%
\AgdaBound{M}\<%
\\
\>[0]\AgdaFunction{subst-with-id}\AgdaSpace{}%
\AgdaSymbol{(}\AgdaInductiveConstructor{ε}\AgdaSpace{}%
\AgdaBound{x}\AgdaSymbol{)}\AgdaSpace{}%
\AgdaSymbol{=}\AgdaSpace{}%
\AgdaInductiveConstructor{refl}\<%
\\
\>[0]\AgdaFunction{subst-with-id}\AgdaSpace{}%
\AgdaSymbol{(}\AgdaInductiveConstructor{box}\AgdaSpace{}%
\AgdaBound{M}\AgdaSymbol{)}\AgdaSpace{}%
\AgdaSymbol{=}\AgdaSpace{}%
\AgdaFunction{cong}\AgdaSpace{}%
\AgdaInductiveConstructor{box}\AgdaSpace{}%
\AgdaSymbol{(}\AgdaFunction{subst-with-id}\AgdaSpace{}%
\AgdaBound{M}\AgdaSymbol{)}\<%
\\
\>[0]\AgdaFunction{subst-with-id}\AgdaSpace{}%
\AgdaSymbol{(}\AgdaInductiveConstructor{ƛb}\AgdaSpace{}%
\AgdaBound{M}\AgdaSymbol{)}\AgdaSpace{}%
\AgdaSymbol{=}\<%
\\
\>[0][@{}l@{\AgdaIndent{0}}]%
\>[2]\AgdaOperator{\AgdaFunction{beginEq}}\<%
\\
\>[2][@{}l@{\AgdaIndent{0}}]%
\>[4]\AgdaInductiveConstructor{ƛb}\AgdaSpace{}%
\AgdaSymbol{(}\AgdaFunction{substb}\AgdaSpace{}%
\AgdaSymbol{(}\AgdaFunction{extsb}\AgdaSpace{}%
\AgdaInductiveConstructor{ε}\AgdaSymbol{)}\AgdaSpace{}%
\AgdaBound{M}\AgdaSymbol{)}\<%
\\
%
\>[2]\AgdaFunction{≡⟨}\AgdaSpace{}%
\AgdaFunction{cong}\AgdaSpace{}%
\AgdaInductiveConstructor{ƛb}\AgdaSpace{}%
\AgdaSymbol{(}\AgdaFunction{substb-cong}\AgdaSpace{}%
\AgdaSymbol{(λ}\AgdaSpace{}%
\AgdaBound{x}\AgdaSpace{}%
\AgdaSymbol{→}\AgdaSpace{}%
\AgdaFunction{ε-exts}\AgdaSpace{}%
\AgdaBound{x}\AgdaSymbol{)}\AgdaSpace{}%
\AgdaBound{M}\AgdaSymbol{)}\AgdaSpace{}%
\AgdaFunction{⟩}\<%
\\
\>[2][@{}l@{\AgdaIndent{0}}]%
\>[4]\AgdaInductiveConstructor{ƛb}\AgdaSpace{}%
\AgdaSymbol{(}\AgdaFunction{substb}\AgdaSpace{}%
\AgdaInductiveConstructor{ε}\AgdaSpace{}%
\AgdaBound{M}\AgdaSymbol{)}\<%
\\
%
\>[2]\AgdaFunction{≡⟨}\AgdaSpace{}%
\AgdaFunction{cong}\AgdaSpace{}%
\AgdaInductiveConstructor{ƛb}\AgdaSpace{}%
\AgdaSymbol{(}\AgdaFunction{subst-with-id}\AgdaSpace{}%
\AgdaBound{M}\AgdaSymbol{)}\AgdaSpace{}%
\AgdaFunction{⟩}\<%
\\
\>[2][@{}l@{\AgdaIndent{0}}]%
\>[4]\AgdaInductiveConstructor{ƛb}\AgdaSpace{}%
\AgdaBound{M}\<%
\\
%
\>[2]\AgdaOperator{\AgdaFunction{∎eq}}\<%
\\
\>[0]\AgdaFunction{subst-with-id}\AgdaSpace{}%
\AgdaSymbol{(}\AgdaBound{M}\AgdaSpace{}%
\AgdaOperator{\AgdaInductiveConstructor{·b}}\AgdaSpace{}%
\AgdaBound{M₁}\AgdaSymbol{)}\AgdaSpace{}%
\AgdaSymbol{=}\AgdaSpace{}%
\AgdaFunction{cong₂}\AgdaSpace{}%
\AgdaOperator{\AgdaInductiveConstructor{\AgdaUnderscore{}·b\AgdaUnderscore{}}}\AgdaSpace{}%
\AgdaSymbol{(}\AgdaFunction{subst-with-id}\AgdaSpace{}%
\AgdaBound{M}\AgdaSymbol{)}\AgdaSpace{}%
\AgdaSymbol{(}\AgdaFunction{subst-with-id}\AgdaSpace{}%
\AgdaBound{M₁}\AgdaSymbol{)}\<%
\\
%
\\[\AgdaEmptyExtraSkip]%
\>[0]\AgdaComment{--\ Substitution\ composed\ with\ renaming\ with\ Sb\ }\<%
\\
\>[0]\AgdaComment{--\ is\ the\ same\ as\ embedding\ that\ term.}\<%
\\
\>[0]\AgdaFunction{subst-rename-fb-comp}\AgdaSpace{}%
\AgdaSymbol{:}\AgdaSpace{}%
\AgdaSymbol{∀}\AgdaSpace{}%
\AgdaSymbol{\{}\AgdaBound{Γ}\AgdaSpace{}%
\AgdaBound{A}\AgdaSpace{}%
\AgdaBound{B}\AgdaSymbol{\}}\<%
\\
\>[0][@{}l@{\AgdaIndent{0}}]%
\>[2]\AgdaSymbol{→}\AgdaSpace{}%
\AgdaSymbol{(}\AgdaBound{M₁}\AgdaSpace{}%
\AgdaSymbol{:}\AgdaSpace{}%
\AgdaBound{Γ}\AgdaSpace{}%
\AgdaOperator{\AgdaInductiveConstructor{,}}\AgdaSpace{}%
\AgdaBound{B}\AgdaSpace{}%
\AgdaOperator{\AgdaDatatype{⊢}}\AgdaSpace{}%
\AgdaBound{A}\AgdaSymbol{)}\<%
\\
%
\>[2]\AgdaSymbol{→}\AgdaSpace{}%
\AgdaSymbol{(}\AgdaBound{M₂}\AgdaSpace{}%
\AgdaSymbol{:}\AgdaSpace{}%
\AgdaBound{Γ}\AgdaSpace{}%
\AgdaOperator{\AgdaDatatype{⊢}}\AgdaSpace{}%
\AgdaBound{B}\AgdaSymbol{)}\<%
\\
%
\>[2]\AgdaSymbol{→}\AgdaSpace{}%
\AgdaSymbol{(}\AgdaBound{x}\AgdaSpace{}%
\AgdaSymbol{:}\AgdaSpace{}%
\AgdaDatatype{Value}\AgdaSpace{}%
\AgdaBound{M₂}\AgdaSpace{}%
\AgdaSymbol{)}\<%
\\
%
\>[2]\AgdaSymbol{→}%
\>[2725I]\AgdaFunction{substb}\AgdaSpace{}%
\AgdaSymbol{(}\AgdaFunction{fb}\AgdaSpace{}%
\AgdaSymbol{(}\AgdaInductiveConstructor{ƛb}\AgdaSpace{}%
\AgdaSymbol{(}\AgdaFunction{embedTerm}\AgdaSpace{}%
\AgdaBound{M₁}\AgdaSymbol{)))}\AgdaSpace{}%
\AgdaSymbol{(}\AgdaFunction{renameb}\AgdaSpace{}%
\AgdaInductiveConstructor{Sb}\AgdaSpace{}%
\AgdaSymbol{(}\AgdaFunction{embedTerm}\AgdaSpace{}%
\AgdaBound{M₂}\AgdaSymbol{))}\<%
\\
\>[.][@{}l@{}]\<[2725I]%
\>[4]\AgdaOperator{\AgdaDatatype{≡}}\AgdaSpace{}%
\AgdaInductiveConstructor{box}\AgdaSpace{}%
\AgdaSymbol{(}\AgdaFunction{embedValue}\AgdaSpace{}%
\AgdaBound{M₂}\AgdaSpace{}%
\AgdaBound{x}\AgdaSymbol{)}\<%
\\
\>[0]\AgdaFunction{subst-rename-fb-comp}\AgdaSpace{}%
\AgdaBound{M₁}\AgdaSpace{}%
\AgdaBound{M₂}\AgdaSpace{}%
\AgdaBound{x}\AgdaSpace{}%
\AgdaSymbol{=}\<%
\\
\>[0][@{}l@{\AgdaIndent{0}}]%
\>[2]\AgdaOperator{\AgdaFunction{beginEq}}\<%
\\
\>[2][@{}l@{\AgdaIndent{0}}]%
\>[4]\AgdaFunction{substb}\AgdaSpace{}%
\AgdaSymbol{(}\AgdaFunction{fb}\AgdaSpace{}%
\AgdaSymbol{(}\AgdaInductiveConstructor{ƛb}\AgdaSpace{}%
\AgdaSymbol{(}\AgdaFunction{embedTerm}\AgdaSpace{}%
\AgdaBound{M₁}\AgdaSymbol{)))}\AgdaSpace{}%
\AgdaSymbol{(}\AgdaFunction{renameb}\AgdaSpace{}%
\AgdaInductiveConstructor{Sb}\AgdaSpace{}%
\AgdaSymbol{(}\AgdaFunction{embedTerm}\AgdaSpace{}%
\AgdaBound{M₂}\AgdaSymbol{))}\<%
\\
%
\>[2]\AgdaFunction{≡⟨}\AgdaSpace{}%
\AgdaFunction{subst-renameb-comp}\AgdaSpace{}%
\AgdaInductiveConstructor{Sb}\AgdaSpace{}%
\AgdaSymbol{(}\AgdaFunction{fb}\AgdaSpace{}%
\AgdaSymbol{(}\AgdaInductiveConstructor{ƛb}\AgdaSpace{}%
\AgdaSymbol{(}\AgdaFunction{embedTerm}\AgdaSpace{}%
\AgdaBound{M₁}\AgdaSymbol{)))}\AgdaSpace{}%
\AgdaSymbol{(}\AgdaFunction{embedTerm}\AgdaSpace{}%
\AgdaBound{M₂}\AgdaSymbol{)}\AgdaSpace{}%
\AgdaFunction{⟩}\<%
\\
\>[2][@{}l@{\AgdaIndent{0}}]%
\>[4]\AgdaFunction{substb}\AgdaSpace{}%
\AgdaSymbol{(λ}\AgdaSpace{}%
\AgdaBound{x₁}\AgdaSpace{}%
\AgdaSymbol{→}\AgdaSpace{}%
\AgdaInductiveConstructor{ε}\AgdaSpace{}%
\AgdaBound{x₁}\AgdaSymbol{)}\AgdaSpace{}%
\AgdaSymbol{(}\AgdaFunction{embedTerm}\AgdaSpace{}%
\AgdaBound{M₂}\AgdaSymbol{)}\<%
\\
%
\>[2]\AgdaFunction{≡⟨}\AgdaSpace{}%
\AgdaFunction{subst-with-id}\AgdaSpace{}%
\AgdaSymbol{(}\AgdaFunction{embedTerm}\AgdaSpace{}%
\AgdaBound{M₂}\AgdaSymbol{)}\AgdaSpace{}%
\AgdaFunction{⟩}\<%
\\
\>[2][@{}l@{\AgdaIndent{0}}]%
\>[4]\AgdaFunction{embedTerm}\AgdaSpace{}%
\AgdaBound{M₂}\<%
\\
%
\>[2]\AgdaFunction{≡⟨}\AgdaSpace{}%
\AgdaFunction{sym}\AgdaSpace{}%
\AgdaSymbol{(}\AgdaFunction{box-embedValue}\AgdaSpace{}%
\AgdaBound{x}\AgdaSymbol{)}\AgdaSpace{}%
\AgdaFunction{⟩}\<%
\\
\>[2][@{}l@{\AgdaIndent{0}}]%
\>[4]\AgdaInductiveConstructor{box}\AgdaSpace{}%
\AgdaSymbol{(}\AgdaFunction{embedValue}\AgdaSpace{}%
\AgdaBound{M₂}\AgdaSpace{}%
\AgdaBound{x}\AgdaSymbol{)}\<%
\\
%
\>[2]\AgdaOperator{\AgdaFunction{∎eq}}\<%
\\
%
\\[\AgdaEmptyExtraSkip]%
\>[0]\AgdaComment{--\ Preservation\ of\ 2\ step\ reduction\ of\ Gödel's\ translation}\<%
\\
\>[0]\AgdaFunction{pres-red}\AgdaSpace{}%
\AgdaSymbol{:}\AgdaSpace{}%
\AgdaSymbol{∀}\AgdaSpace{}%
\AgdaSymbol{\{}\AgdaBound{Γ}\AgdaSpace{}%
\AgdaBound{A}\AgdaSymbol{\}}\<%
\\
\>[0][@{}l@{\AgdaIndent{0}}]%
\>[2]\AgdaSymbol{→}\AgdaSpace{}%
\AgdaSymbol{(}\AgdaBound{M}\AgdaSpace{}%
\AgdaSymbol{:}\AgdaSpace{}%
\AgdaBound{Γ}\AgdaSpace{}%
\AgdaOperator{\AgdaDatatype{⊢}}\AgdaSpace{}%
\AgdaBound{A}\AgdaSymbol{)}\<%
\\
%
\>[2]\AgdaSymbol{→}\AgdaSpace{}%
\AgdaSymbol{(}\AgdaBound{N}\AgdaSpace{}%
\AgdaSymbol{:}\AgdaSpace{}%
\AgdaBound{Γ}\AgdaSpace{}%
\AgdaOperator{\AgdaDatatype{⊢}}\AgdaSpace{}%
\AgdaBound{A}\AgdaSymbol{)}\<%
\\
%
\>[2]\AgdaSymbol{→}\AgdaSpace{}%
\AgdaSymbol{(}\AgdaBound{M}\AgdaSpace{}%
\AgdaOperator{\AgdaDatatype{↝βᵥ}}\AgdaSpace{}%
\AgdaBound{N}\AgdaSymbol{)}\<%
\\
%
\>[2]\AgdaSymbol{→}\AgdaSpace{}%
\AgdaSymbol{(}\AgdaFunction{embedTerm}\AgdaSpace{}%
\AgdaBound{M}\AgdaSpace{}%
\AgdaOperator{\AgdaDatatype{↝βb²}}\AgdaSpace{}%
\AgdaFunction{embedTerm}\AgdaSpace{}%
\AgdaBound{N}\AgdaSymbol{)}\<%
\\
\>[0]\AgdaFunction{pres-red}\AgdaSpace{}%
\AgdaSymbol{(}\AgdaInductiveConstructor{ƛ}\AgdaSpace{}%
\AgdaBound{M₁}\AgdaSpace{}%
\AgdaOperator{\AgdaInductiveConstructor{·}}\AgdaSpace{}%
\AgdaBound{M₂}\AgdaSymbol{)}\AgdaSpace{}%
\AgdaBound{N}\AgdaSpace{}%
\AgdaSymbol{(}\AgdaInductiveConstructor{βᵥ}\AgdaSpace{}%
\AgdaBound{x}\AgdaSymbol{)}\AgdaSpace{}%
\AgdaSymbol{=}\<%
\\
\>[0][@{}l@{\AgdaIndent{0}}]%
\>[2]\AgdaInductiveConstructor{βb²}\<%
\\
%
\>[2]\AgdaSymbol{(}\AgdaFunction{raise}\AgdaSpace{}%
\AgdaSymbol{(}\AgdaFunction{embedTerm}\AgdaSpace{}%
\AgdaBound{M₂}\AgdaSymbol{)}\AgdaSpace{}%
\AgdaOperator{\AgdaInductiveConstructor{·b}}\AgdaSpace{}%
\AgdaInductiveConstructor{box}\AgdaSpace{}%
\AgdaSymbol{(}\AgdaInductiveConstructor{ƛb}\AgdaSpace{}%
\AgdaSymbol{(}\AgdaFunction{embedTerm}\AgdaSpace{}%
\AgdaBound{M₁}\AgdaSymbol{)))}\<%
\\
%
\>[2]\AgdaSymbol{(}\AgdaInductiveConstructor{ƛb}\AgdaSpace{}%
\AgdaSymbol{(}\AgdaFunction{embedTerm}\AgdaSpace{}%
\AgdaBound{M₁}\AgdaSymbol{)}\AgdaSpace{}%
\AgdaOperator{\AgdaInductiveConstructor{·b}}\AgdaSpace{}%
\AgdaInductiveConstructor{box}\AgdaSpace{}%
\AgdaSymbol{(}\AgdaFunction{embedValue}\AgdaSpace{}%
\AgdaBound{M₂}\AgdaSpace{}%
\AgdaBound{x}\AgdaSymbol{))}\<%
\\
%
\>[2]\AgdaSymbol{(}\AgdaFunction{embedTerm}\AgdaSpace{}%
\AgdaSymbol{(}\AgdaBound{M₁}\AgdaSpace{}%
\AgdaOperator{\AgdaFunction{[}}\AgdaSpace{}%
\AgdaBound{M₂}\AgdaSpace{}%
\AgdaOperator{\AgdaFunction{]}}\AgdaSymbol{))}\<%
\\
%
\>[2]\AgdaSymbol{(}\AgdaFunction{Eq.subst}\<%
\\
\>[2][@{}l@{\AgdaIndent{0}}]%
\>[4]\AgdaSymbol{(λ}\AgdaSpace{}%
\AgdaBound{x}\AgdaSpace{}%
\AgdaSymbol{→}\AgdaSpace{}%
\AgdaInductiveConstructor{ƛb}\AgdaSpace{}%
\AgdaSymbol{(}\AgdaInductiveConstructor{ε}\AgdaSpace{}%
\AgdaInductiveConstructor{Zb}\AgdaSpace{}%
\AgdaOperator{\AgdaInductiveConstructor{·b}}\AgdaSpace{}%
\AgdaFunction{renameb}\AgdaSpace{}%
\AgdaInductiveConstructor{Sb}\AgdaSpace{}%
\AgdaSymbol{(}\AgdaFunction{embedTerm}\AgdaSpace{}%
\AgdaBound{M₂}\AgdaSymbol{))}\AgdaSpace{}%
\AgdaOperator{\AgdaInductiveConstructor{·b}}\AgdaSpace{}%
\AgdaInductiveConstructor{box}\AgdaSpace{}%
\AgdaSymbol{(}\AgdaInductiveConstructor{ƛb}\AgdaSpace{}%
\AgdaSymbol{(}\AgdaFunction{embedTerm}\AgdaSpace{}%
\AgdaBound{M₁}\AgdaSymbol{))}\AgdaSpace{}%
\AgdaOperator{\AgdaDatatype{↝βb}}\AgdaSpace{}%
\AgdaBound{x}\AgdaSymbol{)}\<%
\\
%
\>[4]\AgdaSymbol{(}\AgdaFunction{cong₂}\AgdaSpace{}%
\AgdaOperator{\AgdaInductiveConstructor{\AgdaUnderscore{}·b\AgdaUnderscore{}}}\AgdaSpace{}%
\AgdaInductiveConstructor{refl}\AgdaSpace{}%
\AgdaSymbol{(}\AgdaFunction{subst-rename-fb-comp}\AgdaSpace{}%
\AgdaBound{M₁}\AgdaSpace{}%
\AgdaBound{M₂}\AgdaSpace{}%
\AgdaBound{x}\AgdaSymbol{))}\<%
\\
%
\>[4]\AgdaInductiveConstructor{βb}\AgdaSymbol{)}\<%
\\
%
\>[2]\AgdaSymbol{(}\AgdaFunction{Eq.subst}\AgdaSpace{}%
\AgdaSymbol{(λ}\AgdaSpace{}%
\AgdaBound{t}\AgdaSpace{}%
\AgdaSymbol{→}\AgdaSpace{}%
\AgdaInductiveConstructor{ƛb}\AgdaSpace{}%
\AgdaSymbol{(}\AgdaFunction{embedTerm}\AgdaSpace{}%
\AgdaBound{M₁}\AgdaSymbol{)}\AgdaSpace{}%
\AgdaOperator{\AgdaInductiveConstructor{·b}}\AgdaSpace{}%
\AgdaInductiveConstructor{box}\AgdaSpace{}%
\AgdaSymbol{(}\AgdaFunction{embedValue}\AgdaSpace{}%
\AgdaBound{M₂}\AgdaSpace{}%
\AgdaBound{x}\AgdaSymbol{)}\AgdaSpace{}%
\AgdaOperator{\AgdaDatatype{↝βb}}\AgdaSpace{}%
\AgdaBound{t}\AgdaSymbol{)}\AgdaSpace{}%
\AgdaSymbol{(}\AgdaFunction{pres-subst}\AgdaSpace{}%
\AgdaBound{M₁}\AgdaSpace{}%
\AgdaBound{M₂}\AgdaSpace{}%
\AgdaBound{x}\AgdaSymbol{)}\AgdaSpace{}%
\AgdaInductiveConstructor{βb}\AgdaSymbol{)}\<%
\\
\>[0]\AgdaFunction{pres-red}\AgdaSpace{}%
\AgdaBound{M}\AgdaSpace{}%
\AgdaBound{N}\AgdaSpace{}%
\AgdaSymbol{(}\AgdaInductiveConstructor{μ}\AgdaSpace{}%
\AgdaSymbol{\{}\AgdaArgument{M}\AgdaSpace{}%
\AgdaSymbol{=}\AgdaSpace{}%
\AgdaBound{M₁}\AgdaSymbol{\}}\AgdaSpace{}%
\AgdaSymbol{\{}\AgdaArgument{M'}\AgdaSpace{}%
\AgdaSymbol{=}\AgdaSpace{}%
\AgdaBound{M'}\AgdaSymbol{\}}\AgdaSpace{}%
\AgdaBound{red}\AgdaSymbol{)}\AgdaSpace{}%
\AgdaSymbol{=}\AgdaSpace{}%
\AgdaInductiveConstructor{ν}\AgdaSpace{}%
\AgdaSymbol{(}\AgdaFunction{pres-red}\AgdaSpace{}%
\AgdaBound{M₁}\AgdaSpace{}%
\AgdaBound{M'}\AgdaSpace{}%
\AgdaBound{red}\AgdaSymbol{)}\<%
\\
\>[0]\AgdaFunction{pres-red}\AgdaSpace{}%
\AgdaBound{M}\AgdaSpace{}%
\AgdaBound{N}\AgdaSpace{}%
\AgdaSymbol{(}\AgdaInductiveConstructor{ν<}\AgdaSpace{}%
\AgdaSymbol{\{}\AgdaArgument{M}\AgdaSpace{}%
\AgdaSymbol{=}\AgdaSpace{}%
\AgdaBound{M₁}\AgdaSymbol{\}}\AgdaSpace{}%
\AgdaSymbol{\{}\AgdaArgument{N}\AgdaSpace{}%
\AgdaSymbol{=}\AgdaSpace{}%
\AgdaBound{N₁}\AgdaSymbol{\}}\AgdaSpace{}%
\AgdaSymbol{\{}\AgdaArgument{N'}\AgdaSpace{}%
\AgdaSymbol{=}\AgdaSpace{}%
\AgdaBound{N'}\AgdaSymbol{\}}\AgdaSpace{}%
\AgdaBound{x}\AgdaSpace{}%
\AgdaBound{red}\AgdaSymbol{)}\AgdaSpace{}%
\AgdaSymbol{=}\AgdaSpace{}%
\AgdaInductiveConstructor{μ}\AgdaSpace{}%
\AgdaSymbol{(}\AgdaInductiveConstructor{ξ}\AgdaSpace{}%
\AgdaSymbol{(}\AgdaInductiveConstructor{ν}\AgdaSpace{}%
\AgdaSymbol{(}\AgdaFunction{pres-rename-red-2}\AgdaSpace{}%
\AgdaInductiveConstructor{Sb}\AgdaSpace{}%
\AgdaSymbol{(}\AgdaFunction{embedTerm}\AgdaSpace{}%
\AgdaBound{N₁}\AgdaSymbol{)}\AgdaSpace{}%
\AgdaSymbol{(}\AgdaFunction{embedTerm}\AgdaSpace{}%
\AgdaBound{N'}\AgdaSymbol{)}\AgdaSpace{}%
\AgdaSymbol{(}\AgdaFunction{pres-red}\AgdaSpace{}%
\AgdaBound{N₁}\AgdaSpace{}%
\AgdaBound{N'}\AgdaSpace{}%
\AgdaBound{red}\AgdaSymbol{))))}\<%
\\
\>[0]\AgdaFunction{pres-red}\AgdaSpace{}%
\AgdaBound{M}\AgdaSpace{}%
\AgdaBound{N}\AgdaSpace{}%
\AgdaSymbol{(}\AgdaInductiveConstructor{ν}\AgdaSpace{}%
\AgdaSymbol{\{}\AgdaArgument{M}\AgdaSpace{}%
\AgdaSymbol{=}\AgdaSpace{}%
\AgdaBound{M₁}\AgdaSymbol{\}}\AgdaSpace{}%
\AgdaSymbol{\{}\AgdaArgument{N}\AgdaSpace{}%
\AgdaSymbol{=}\AgdaSpace{}%
\AgdaBound{N₁}\AgdaSymbol{\}}\AgdaSpace{}%
\AgdaSymbol{\{}\AgdaArgument{N'}\AgdaSpace{}%
\AgdaSymbol{=}\AgdaSpace{}%
\AgdaBound{N'}\AgdaSymbol{\}}\AgdaSpace{}%
\AgdaBound{red}\AgdaSymbol{)}\AgdaSpace{}%
\AgdaSymbol{=}\AgdaSpace{}%
\AgdaInductiveConstructor{μ}\AgdaSpace{}%
\AgdaSymbol{(}\AgdaInductiveConstructor{ξ}\AgdaSpace{}%
\AgdaSymbol{(}\AgdaInductiveConstructor{ν}\AgdaSpace{}%
\AgdaSymbol{(}\AgdaFunction{pres-rename-red-2}\AgdaSpace{}%
\AgdaInductiveConstructor{Sb}\AgdaSpace{}%
\AgdaSymbol{(}\AgdaFunction{embedTerm}\AgdaSpace{}%
\AgdaBound{N₁}\AgdaSymbol{)}\AgdaSpace{}%
\AgdaSymbol{(}\AgdaFunction{embedTerm}\AgdaSpace{}%
\AgdaBound{N'}\AgdaSymbol{)}\AgdaSpace{}%
\AgdaSymbol{(}\AgdaFunction{pres-red}\AgdaSpace{}%
\AgdaBound{N₁}\AgdaSpace{}%
\AgdaBound{N'}\AgdaSpace{}%
\AgdaBound{red}\AgdaSymbol{))))}\<%
\\
\>[0]\AgdaFunction{pres-red}\AgdaSpace{}%
\AgdaBound{M}\AgdaSpace{}%
\AgdaBound{N}\AgdaSpace{}%
\AgdaSymbol{(}\AgdaInductiveConstructor{ξ}\AgdaSpace{}%
\AgdaSymbol{\{}\AgdaArgument{M}\AgdaSpace{}%
\AgdaSymbol{=}\AgdaSpace{}%
\AgdaBound{M₁}\AgdaSymbol{\}}\AgdaSpace{}%
\AgdaSymbol{\{}\AgdaArgument{M'}\AgdaSpace{}%
\AgdaSymbol{=}\AgdaSpace{}%
\AgdaBound{M'}\AgdaSymbol{\}}\AgdaSpace{}%
\AgdaBound{red}\AgdaSymbol{)}\AgdaSpace{}%
\AgdaSymbol{=}\AgdaSpace{}%
\AgdaInductiveConstructor{ζ}\AgdaSpace{}%
\AgdaSymbol{(}\AgdaInductiveConstructor{ξ}\AgdaSpace{}%
\AgdaSymbol{(}\AgdaFunction{pres-red}\AgdaSpace{}%
\AgdaBound{M₁}\AgdaSpace{}%
\AgdaBound{M'}\AgdaSpace{}%
\AgdaBound{red}\AgdaSymbol{))}\<%
\\
%
\\[\AgdaEmptyExtraSkip]%
\>[0]\AgdaComment{--\ Preservation\ of\ the\ 2\ step\ reduction\ rule\ under\ rename}\<%
\\
\>[0]\AgdaFunction{pres-rename-red-↝βb²}\AgdaSpace{}%
\AgdaSymbol{:}\AgdaSpace{}%
\AgdaSymbol{∀}\AgdaSpace{}%
\AgdaSymbol{\{}\AgdaBound{Γ}\AgdaSpace{}%
\AgdaBound{A}\AgdaSpace{}%
\AgdaBound{Δ}\AgdaSymbol{\}}\<%
\\
\>[0][@{}l@{\AgdaIndent{0}}]%
\>[2]\AgdaSymbol{→}\AgdaSpace{}%
\AgdaSymbol{(}\AgdaBound{fb'}\AgdaSpace{}%
\AgdaSymbol{:}\AgdaSpace{}%
\AgdaSymbol{∀}\AgdaSpace{}%
\AgdaSymbol{\{}\AgdaBound{Ab}\AgdaSymbol{\}}\AgdaSpace{}%
\AgdaSymbol{→}\AgdaSpace{}%
\AgdaBound{Ab}\AgdaSpace{}%
\AgdaOperator{\AgdaDatatype{∈b}}\AgdaSpace{}%
\AgdaBound{Γ}\AgdaSpace{}%
\AgdaSymbol{→}\AgdaSpace{}%
\AgdaBound{Ab}\AgdaSpace{}%
\AgdaOperator{\AgdaDatatype{∈b}}\AgdaSpace{}%
\AgdaBound{Δ}\AgdaSymbol{)}\<%
\\
%
\>[2]\AgdaSymbol{→}\AgdaSpace{}%
\AgdaSymbol{(}\AgdaBound{M}\AgdaSpace{}%
\AgdaSymbol{:}\AgdaSpace{}%
\AgdaBound{Γ}\AgdaSpace{}%
\AgdaOperator{\AgdaDatatype{⊢b}}\AgdaSpace{}%
\AgdaBound{A}\AgdaSymbol{)}\<%
\\
%
\>[2]\AgdaSymbol{→}\AgdaSpace{}%
\AgdaSymbol{(}\AgdaBound{N}\AgdaSpace{}%
\AgdaSymbol{:}\AgdaSpace{}%
\AgdaBound{Γ}\AgdaSpace{}%
\AgdaOperator{\AgdaDatatype{⊢b}}\AgdaSpace{}%
\AgdaBound{A}\AgdaSymbol{)}\<%
\\
%
\>[2]\AgdaSymbol{→}\AgdaSpace{}%
\AgdaSymbol{(}\AgdaBound{M}\AgdaSpace{}%
\AgdaOperator{\AgdaDatatype{↝βb₂}}\AgdaSpace{}%
\AgdaBound{N}\AgdaSymbol{)}\<%
\\
%
\>[2]\AgdaSymbol{→}\AgdaSpace{}%
\AgdaFunction{renameb}\AgdaSpace{}%
\AgdaBound{fb'}\AgdaSpace{}%
\AgdaBound{M}\AgdaSpace{}%
\AgdaOperator{\AgdaDatatype{↝βb₂}}\AgdaSpace{}%
\AgdaFunction{renameb}\AgdaSpace{}%
\AgdaBound{fb'}\AgdaSpace{}%
\AgdaBound{N}\<%
\\
\>[0]\AgdaFunction{pres-rename-red-↝βb²}\AgdaSpace{}%
\AgdaBound{fb'}\AgdaSpace{}%
\AgdaBound{M}\AgdaSpace{}%
\AgdaBound{N}\AgdaSpace{}%
\AgdaSymbol{(}\AgdaInductiveConstructor{βb₂}\AgdaSpace{}%
\AgdaSymbol{\{}\AgdaArgument{M}\AgdaSpace{}%
\AgdaSymbol{=}\AgdaSpace{}%
\AgdaBound{M'}\AgdaSymbol{\}}\AgdaSpace{}%
\AgdaSymbol{\{}\AgdaArgument{W}\AgdaSpace{}%
\AgdaSymbol{=}\AgdaSpace{}%
\AgdaBound{W'}\AgdaSymbol{\})}\AgdaSpace{}%
\AgdaSymbol{=}\<%
\\
\>[0][@{}l@{\AgdaIndent{0}}]%
\>[2]\AgdaFunction{Eq.subst}\<%
\\
\>[2][@{}l@{\AgdaIndent{0}}]%
\>[4]\AgdaSymbol{(λ}\AgdaSpace{}%
\AgdaBound{x}\AgdaSpace{}%
\AgdaSymbol{→}\AgdaSpace{}%
\AgdaBound{x}\AgdaSpace{}%
\AgdaOperator{\AgdaDatatype{↝βb₂}}\AgdaSpace{}%
\AgdaFunction{renameb}\AgdaSpace{}%
\AgdaBound{fb'}\<%
\\
%
\>[4]\AgdaSymbol{(}\AgdaBound{M'}\AgdaSpace{}%
\AgdaOperator{\AgdaFunction{[}}\AgdaSpace{}%
\AgdaBound{W'}\AgdaSpace{}%
\AgdaOperator{\AgdaFunction{]b}}\AgdaSymbol{))}\<%
\\
%
\>[4]\AgdaSymbol{(}\AgdaFunction{cong₂}\AgdaSpace{}%
\AgdaOperator{\AgdaInductiveConstructor{\AgdaUnderscore{}·b\AgdaUnderscore{}}}\AgdaSpace{}%
\AgdaSymbol{(}\AgdaFunction{cong}\AgdaSpace{}%
\AgdaInductiveConstructor{ƛb}\AgdaSpace{}%
\AgdaSymbol{(}\AgdaFunction{cong₂}\AgdaSpace{}%
\AgdaOperator{\AgdaInductiveConstructor{\AgdaUnderscore{}·b\AgdaUnderscore{}}}\AgdaSpace{}%
\AgdaInductiveConstructor{refl}\AgdaSpace{}%
\AgdaSymbol{(}\AgdaFunction{cong}\AgdaSpace{}%
\AgdaInductiveConstructor{box}\AgdaSpace{}%
\AgdaSymbol{(}\AgdaFunction{sym}\AgdaSpace{}%
\AgdaSymbol{(}\AgdaFunction{rename-extb-Sb-commute}\AgdaSpace{}%
\AgdaBound{fb'}\AgdaSpace{}%
\AgdaBound{W'}\AgdaSymbol{)))))}\AgdaSpace{}%
\AgdaInductiveConstructor{refl}\AgdaSymbol{)}\<%
\\
%
\>[4]\AgdaSymbol{(}\AgdaFunction{Eq.subst}\<%
\\
\>[4][@{}l@{\AgdaIndent{0}}]%
\>[6]\AgdaSymbol{(λ}\AgdaSpace{}%
\AgdaBound{x}\AgdaSpace{}%
\AgdaSymbol{→}\AgdaSpace{}%
\AgdaSymbol{(}\AgdaFunction{raise}%
\>[21]\AgdaSymbol{(}\AgdaInductiveConstructor{box}\AgdaSpace{}%
\AgdaSymbol{(}\AgdaFunction{renameb}\AgdaSpace{}%
\AgdaBound{fb'}\AgdaSpace{}%
\AgdaBound{W'}\AgdaSymbol{)))}\AgdaSpace{}%
\AgdaOperator{\AgdaInductiveConstructor{·b}}\AgdaSpace{}%
\AgdaInductiveConstructor{box}\AgdaSpace{}%
\AgdaSymbol{(}\AgdaInductiveConstructor{ƛb}\AgdaSpace{}%
\AgdaSymbol{(}\AgdaFunction{renameb}\AgdaSpace{}%
\AgdaSymbol{(}\AgdaFunction{extb}\AgdaSpace{}%
\AgdaBound{fb'}\AgdaSymbol{)}\AgdaSpace{}%
\AgdaBound{M'}\AgdaSymbol{))}\AgdaSpace{}%
\AgdaOperator{\AgdaDatatype{↝βb₂}}\AgdaSpace{}%
\AgdaBound{x}\AgdaSymbol{)}\<%
\\
%
\>[6]\AgdaSymbol{(}\AgdaFunction{pres-subst-func-rename}\AgdaSpace{}%
\AgdaBound{fb'}\AgdaSpace{}%
\AgdaBound{M'}\AgdaSpace{}%
\AgdaBound{W'}\AgdaSymbol{)}\<%
\\
%
\>[6]\AgdaInductiveConstructor{βb₂}\<%
\\
%
\>[4]\AgdaSymbol{)}\<%
\\
\>[0]\AgdaFunction{pres-rename-red-↝βb²}\AgdaSpace{}%
\AgdaBound{fb'}\AgdaSpace{}%
\AgdaBound{M}\AgdaSpace{}%
\AgdaBound{N}\AgdaSpace{}%
\AgdaSymbol{(}\AgdaInductiveConstructor{μ}\AgdaSpace{}%
\AgdaSymbol{\{}\AgdaArgument{M}\AgdaSpace{}%
\AgdaSymbol{=}\AgdaSpace{}%
\AgdaBound{M₁}\AgdaSymbol{\}}\AgdaSpace{}%
\AgdaSymbol{\{}\AgdaArgument{M'}\AgdaSpace{}%
\AgdaSymbol{=}\AgdaSpace{}%
\AgdaBound{M'}\AgdaSymbol{\}}\AgdaSpace{}%
\AgdaBound{red}\AgdaSymbol{)}\AgdaSpace{}%
\AgdaSymbol{=}\AgdaSpace{}%
\AgdaInductiveConstructor{μ}\AgdaSpace{}%
\AgdaSymbol{(}\AgdaFunction{pres-rename-red-↝βb²}\AgdaSpace{}%
\AgdaBound{fb'}\AgdaSpace{}%
\AgdaBound{M₁}\AgdaSpace{}%
\AgdaBound{M'}\AgdaSpace{}%
\AgdaBound{red}\AgdaSymbol{)}\<%
\\
\>[0]\AgdaFunction{pres-rename-red-↝βb²}\AgdaSpace{}%
\AgdaBound{fb'}\AgdaSpace{}%
\AgdaBound{M}\AgdaSpace{}%
\AgdaBound{N}\AgdaSpace{}%
\AgdaSymbol{(}\AgdaInductiveConstructor{ν}\AgdaSpace{}%
\AgdaSymbol{\{}\AgdaArgument{N}\AgdaSpace{}%
\AgdaSymbol{=}\AgdaSpace{}%
\AgdaBound{N₁}\AgdaSymbol{\}}\AgdaSpace{}%
\AgdaSymbol{\{}\AgdaArgument{N'}\AgdaSpace{}%
\AgdaSymbol{=}\AgdaSpace{}%
\AgdaBound{N'}\AgdaSymbol{\}}\AgdaSpace{}%
\AgdaBound{red}\AgdaSymbol{)}\AgdaSpace{}%
\AgdaSymbol{=}\AgdaSpace{}%
\AgdaInductiveConstructor{ν}\AgdaSpace{}%
\AgdaSymbol{(}\AgdaFunction{pres-rename-red-↝βb²}\AgdaSpace{}%
\AgdaBound{fb'}\AgdaSpace{}%
\AgdaBound{N₁}\AgdaSpace{}%
\AgdaBound{N'}\AgdaSpace{}%
\AgdaBound{red}\AgdaSymbol{)}\<%
\\
\>[0]\AgdaFunction{pres-rename-red-↝βb²}\AgdaSpace{}%
\AgdaBound{fb'}\AgdaSpace{}%
\AgdaBound{M}\AgdaSpace{}%
\AgdaBound{N}\AgdaSpace{}%
\AgdaSymbol{(}\AgdaInductiveConstructor{ξ}\AgdaSpace{}%
\AgdaSymbol{\{}\AgdaArgument{M}\AgdaSpace{}%
\AgdaSymbol{=}\AgdaSpace{}%
\AgdaBound{M₁}\AgdaSymbol{\}}\AgdaSpace{}%
\AgdaSymbol{\{}\AgdaArgument{M'}\AgdaSpace{}%
\AgdaSymbol{=}\AgdaSpace{}%
\AgdaBound{M'}\AgdaSymbol{\}}\AgdaSpace{}%
\AgdaBound{red}\AgdaSymbol{)}\AgdaSpace{}%
\AgdaSymbol{=}\AgdaSpace{}%
\AgdaInductiveConstructor{ξ}\AgdaSpace{}%
\AgdaSymbol{(}\AgdaFunction{pres-rename-red-↝βb²}\AgdaSpace{}%
\AgdaSymbol{(}\AgdaFunction{extb}\AgdaSpace{}%
\AgdaBound{fb'}\AgdaSymbol{)}\AgdaSpace{}%
\AgdaBound{M₁}\AgdaSpace{}%
\AgdaBound{M'}\AgdaSpace{}%
\AgdaBound{red}\AgdaSymbol{)}\<%
\\
\>[0]\AgdaFunction{pres-rename-red-↝βb²}\AgdaSpace{}%
\AgdaBound{fb'}\AgdaSpace{}%
\AgdaBound{M}\AgdaSpace{}%
\AgdaBound{N}\AgdaSpace{}%
\AgdaSymbol{(}\AgdaInductiveConstructor{ζ}\AgdaSpace{}%
\AgdaSymbol{\{}\AgdaArgument{M}\AgdaSpace{}%
\AgdaSymbol{=}\AgdaSpace{}%
\AgdaBound{M₁}\AgdaSymbol{\}}\AgdaSpace{}%
\AgdaSymbol{\{}\AgdaArgument{M'}\AgdaSpace{}%
\AgdaSymbol{=}\AgdaSpace{}%
\AgdaBound{M'}\AgdaSymbol{\}}\AgdaSpace{}%
\AgdaBound{red}\AgdaSymbol{)}\AgdaSpace{}%
\AgdaSymbol{=}\AgdaSpace{}%
\AgdaInductiveConstructor{ζ}\AgdaSpace{}%
\AgdaSymbol{(}\AgdaFunction{pres-rename-red-↝βb²}\AgdaSpace{}%
\AgdaBound{fb'}\AgdaSpace{}%
\AgdaBound{M₁}\AgdaSpace{}%
\AgdaBound{M'}\AgdaSpace{}%
\AgdaBound{red}\AgdaSymbol{)}\<%
\\
%
\\[\AgdaEmptyExtraSkip]%
\>[0]\AgdaComment{--\ Preservation\ of\ reduction\ by\ Gödel's\ translation}\<%
\\
\>[0]\AgdaComment{--\ with\ the\ two\ step\ reduction\ rule}\<%
\\
\>[0]\AgdaFunction{pres-red↝βb²}\AgdaSpace{}%
\AgdaSymbol{:}\AgdaSpace{}%
\AgdaSymbol{∀}\AgdaSpace{}%
\AgdaSymbol{\{}\AgdaBound{Γ}\AgdaSpace{}%
\AgdaBound{A}\AgdaSymbol{\}}\<%
\\
\>[0][@{}l@{\AgdaIndent{0}}]%
\>[2]\AgdaSymbol{→}\AgdaSpace{}%
\AgdaSymbol{(}\AgdaBound{M}\AgdaSpace{}%
\AgdaSymbol{:}\AgdaSpace{}%
\AgdaBound{Γ}\AgdaSpace{}%
\AgdaOperator{\AgdaDatatype{⊢}}\AgdaSpace{}%
\AgdaBound{A}\AgdaSymbol{)}\<%
\\
%
\>[2]\AgdaSymbol{→}\AgdaSpace{}%
\AgdaSymbol{(}\AgdaBound{N}\AgdaSpace{}%
\AgdaSymbol{:}\AgdaSpace{}%
\AgdaBound{Γ}\AgdaSpace{}%
\AgdaOperator{\AgdaDatatype{⊢}}\AgdaSpace{}%
\AgdaBound{A}\AgdaSymbol{)}\<%
\\
%
\>[2]\AgdaSymbol{→}\AgdaSpace{}%
\AgdaSymbol{(}\AgdaBound{M}\AgdaSpace{}%
\AgdaOperator{\AgdaDatatype{↝βᵥ}}\AgdaSpace{}%
\AgdaBound{N}\AgdaSymbol{)}\<%
\\
%
\>[2]\AgdaSymbol{→}\AgdaSpace{}%
\AgdaSymbol{(}\AgdaFunction{embedTerm}\AgdaSpace{}%
\AgdaBound{M}\AgdaSpace{}%
\AgdaOperator{\AgdaDatatype{↝βb₂}}\AgdaSpace{}%
\AgdaFunction{embedTerm}\AgdaSpace{}%
\AgdaBound{N}\AgdaSymbol{)}\<%
\\
\>[0]\AgdaFunction{pres-red↝βb²}\AgdaSpace{}%
\AgdaSymbol{(}\AgdaInductiveConstructor{ƛ}\AgdaSpace{}%
\AgdaBound{M₁}\AgdaSpace{}%
\AgdaOperator{\AgdaInductiveConstructor{·}}\AgdaSpace{}%
\AgdaBound{M₂}\AgdaSymbol{)}\AgdaSpace{}%
\AgdaBound{N}\AgdaSpace{}%
\AgdaSymbol{(}\AgdaInductiveConstructor{βᵥ}\AgdaSpace{}%
\AgdaBound{x}\AgdaSymbol{)}\AgdaSpace{}%
\AgdaSymbol{=}\<%
\\
\>[0][@{}l@{\AgdaIndent{0}}]%
\>[2]\AgdaFunction{Eq.subst}\<%
\\
\>[2][@{}l@{\AgdaIndent{0}}]%
\>[4]\AgdaSymbol{(λ}\AgdaSpace{}%
\AgdaBound{t}\AgdaSpace{}%
\AgdaSymbol{→}\AgdaSpace{}%
\AgdaBound{t}\AgdaSpace{}%
\AgdaOperator{\AgdaDatatype{↝βb₂}}\AgdaSpace{}%
\AgdaFunction{embedTerm}\AgdaSpace{}%
\AgdaSymbol{(}\AgdaBound{M₁}\AgdaSpace{}%
\AgdaOperator{\AgdaFunction{[}}\AgdaSpace{}%
\AgdaBound{M₂}\AgdaSpace{}%
\AgdaOperator{\AgdaFunction{]}}\AgdaSymbol{))}\<%
\\
%
\>[4]\AgdaSymbol{(}\AgdaFunction{cong₂}\AgdaSpace{}%
\AgdaOperator{\AgdaInductiveConstructor{\AgdaUnderscore{}·b\AgdaUnderscore{}}}\AgdaSpace{}%
\AgdaSymbol{(}\AgdaFunction{cong}\AgdaSpace{}%
\AgdaInductiveConstructor{ƛb}\AgdaSpace{}%
\AgdaSymbol{(}\AgdaFunction{cong₂}\AgdaSpace{}%
\AgdaOperator{\AgdaInductiveConstructor{\AgdaUnderscore{}·b\AgdaUnderscore{}}}\AgdaSpace{}%
\AgdaInductiveConstructor{refl}\AgdaSpace{}%
\AgdaSymbol{(}\AgdaFunction{cong}\AgdaSpace{}%
\AgdaSymbol{(}\AgdaFunction{renameb}\AgdaSpace{}%
\AgdaInductiveConstructor{Sb}\AgdaSymbol{)}\AgdaSpace{}%
\AgdaSymbol{(}\AgdaFunction{box-embedValue}\AgdaSpace{}%
\AgdaBound{x}\AgdaSymbol{))))}\AgdaSpace{}%
\AgdaInductiveConstructor{refl}\AgdaSymbol{)}\<%
\\
%
\>[4]\AgdaSymbol{(}\AgdaFunction{Eq.subst}\<%
\\
\>[4][@{}l@{\AgdaIndent{0}}]%
\>[6]\AgdaSymbol{(λ}\AgdaSpace{}%
\AgdaBound{t}\AgdaSpace{}%
\AgdaSymbol{→}\AgdaSpace{}%
\AgdaFunction{raise}\AgdaSpace{}%
\AgdaSymbol{(}\AgdaSpace{}%
\AgdaInductiveConstructor{box}\AgdaSpace{}%
\AgdaSymbol{(}\AgdaFunction{embedValue}\AgdaSpace{}%
\AgdaBound{M₂}\AgdaSpace{}%
\AgdaBound{x}\AgdaSymbol{)}\AgdaSpace{}%
\AgdaSymbol{)}\AgdaSpace{}%
\AgdaOperator{\AgdaInductiveConstructor{·b}}\AgdaSpace{}%
\AgdaInductiveConstructor{box}\AgdaSpace{}%
\AgdaSymbol{(}\AgdaInductiveConstructor{ƛb}\AgdaSpace{}%
\AgdaSymbol{(}\AgdaFunction{embedTerm}\AgdaSpace{}%
\AgdaBound{M₁}\AgdaSymbol{))}\AgdaSpace{}%
\AgdaOperator{\AgdaDatatype{↝βb₂}}\AgdaSpace{}%
\AgdaBound{t}\AgdaSymbol{)}\<%
\\
%
\>[6]\AgdaSymbol{(}\AgdaFunction{pres-subst}\AgdaSpace{}%
\AgdaBound{M₁}\AgdaSpace{}%
\AgdaBound{M₂}\AgdaSpace{}%
\AgdaBound{x}\AgdaSymbol{)}\<%
\\
%
\>[6]\AgdaInductiveConstructor{βb₂}\<%
\\
%
\>[4]\AgdaSymbol{)}\<%
\\
\>[0]\AgdaFunction{pres-red↝βb²}\AgdaSpace{}%
\AgdaBound{M}\AgdaSpace{}%
\AgdaBound{N}\AgdaSpace{}%
\AgdaSymbol{(}\AgdaInductiveConstructor{μ}\AgdaSpace{}%
\AgdaSymbol{\{}\AgdaArgument{M}\AgdaSpace{}%
\AgdaSymbol{=}\AgdaSpace{}%
\AgdaBound{M₁}\AgdaSymbol{\}}\AgdaSpace{}%
\AgdaSymbol{\{}\AgdaArgument{M'}\AgdaSpace{}%
\AgdaSymbol{=}\AgdaSpace{}%
\AgdaBound{M'}\AgdaSymbol{\}}\AgdaSpace{}%
\AgdaBound{red}\AgdaSymbol{)}\AgdaSpace{}%
\AgdaSymbol{=}\AgdaSpace{}%
\AgdaInductiveConstructor{ν}\AgdaSpace{}%
\AgdaSymbol{(}\AgdaFunction{pres-red↝βb²}\AgdaSpace{}%
\AgdaBound{M₁}\AgdaSpace{}%
\AgdaBound{M'}\AgdaSpace{}%
\AgdaBound{red}\AgdaSymbol{)}\<%
\\
\>[0]\AgdaFunction{pres-red↝βb²}\AgdaSpace{}%
\AgdaBound{M}\AgdaSpace{}%
\AgdaBound{N}\AgdaSpace{}%
\AgdaSymbol{(}\AgdaInductiveConstructor{ν<}\AgdaSpace{}%
\AgdaSymbol{\{}\AgdaArgument{N}\AgdaSpace{}%
\AgdaSymbol{=}\AgdaSpace{}%
\AgdaBound{N₁}\AgdaSymbol{\}}\AgdaSpace{}%
\AgdaSymbol{\{}\AgdaArgument{N'}\AgdaSpace{}%
\AgdaSymbol{=}\AgdaSpace{}%
\AgdaBound{N'}\AgdaSymbol{\}}\AgdaSpace{}%
\AgdaBound{x}\AgdaSpace{}%
\AgdaBound{red}\AgdaSymbol{)}\AgdaSpace{}%
\AgdaSymbol{=}\AgdaSpace{}%
\AgdaInductiveConstructor{μ}\AgdaSpace{}%
\AgdaSymbol{(}\AgdaInductiveConstructor{ξ}\AgdaSpace{}%
\AgdaSymbol{(}\AgdaInductiveConstructor{ν}\AgdaSpace{}%
\AgdaSymbol{(}\AgdaFunction{pres-rename-red-↝βb²}\AgdaSpace{}%
\AgdaInductiveConstructor{Sb}\AgdaSpace{}%
\AgdaSymbol{(}\AgdaFunction{embedTerm}\AgdaSpace{}%
\AgdaBound{N₁}\AgdaSymbol{)}\AgdaSpace{}%
\AgdaSymbol{(}\AgdaFunction{embedTerm}\AgdaSpace{}%
\AgdaBound{N'}\AgdaSymbol{)}\AgdaSpace{}%
\AgdaSymbol{(}\AgdaFunction{pres-red↝βb²}\AgdaSpace{}%
\AgdaBound{N₁}\AgdaSpace{}%
\AgdaBound{N'}\AgdaSpace{}%
\AgdaBound{red}\AgdaSymbol{))))}\<%
\\
\>[0]\AgdaFunction{pres-red↝βb²}\AgdaSpace{}%
\AgdaBound{M}\AgdaSpace{}%
\AgdaBound{N}\AgdaSpace{}%
\AgdaSymbol{(}\AgdaInductiveConstructor{ν}\AgdaSpace{}%
\AgdaSymbol{\{}\AgdaArgument{N}\AgdaSpace{}%
\AgdaSymbol{=}\AgdaSpace{}%
\AgdaBound{N₁}\AgdaSymbol{\}}\AgdaSpace{}%
\AgdaSymbol{\{}\AgdaArgument{N'}\AgdaSpace{}%
\AgdaSymbol{=}\AgdaSpace{}%
\AgdaBound{N'}\AgdaSymbol{\}}\AgdaSpace{}%
\AgdaBound{red}\AgdaSymbol{)}\AgdaSpace{}%
\AgdaSymbol{=}\AgdaSpace{}%
\AgdaInductiveConstructor{μ}\AgdaSpace{}%
\AgdaSymbol{(}\AgdaInductiveConstructor{ξ}\AgdaSpace{}%
\AgdaSymbol{(}\AgdaInductiveConstructor{ν}\AgdaSpace{}%
\AgdaSymbol{(}\AgdaFunction{pres-rename-red-↝βb²}\AgdaSpace{}%
\AgdaInductiveConstructor{Sb}\AgdaSpace{}%
\AgdaSymbol{(}\AgdaFunction{embedTerm}\AgdaSpace{}%
\AgdaBound{N₁}\AgdaSymbol{)}\AgdaSpace{}%
\AgdaSymbol{(}\AgdaFunction{embedTerm}\AgdaSpace{}%
\AgdaBound{N'}\AgdaSymbol{)}\AgdaSpace{}%
\AgdaSymbol{(}\AgdaFunction{pres-red↝βb²}\AgdaSpace{}%
\AgdaBound{N₁}\AgdaSpace{}%
\AgdaBound{N'}\AgdaSpace{}%
\AgdaBound{red}\AgdaSymbol{))))}\<%
\\
\>[0]\AgdaFunction{pres-red↝βb²}\AgdaSpace{}%
\AgdaBound{M}\AgdaSpace{}%
\AgdaBound{N}\AgdaSpace{}%
\AgdaSymbol{(}\AgdaInductiveConstructor{ξ}\AgdaSpace{}%
\AgdaSymbol{\{}\AgdaArgument{M}\AgdaSpace{}%
\AgdaSymbol{=}\AgdaSpace{}%
\AgdaBound{M₁}\AgdaSymbol{\}}\AgdaSpace{}%
\AgdaSymbol{\{}\AgdaArgument{M'}\AgdaSpace{}%
\AgdaSymbol{=}\AgdaSpace{}%
\AgdaBound{M'}\AgdaSymbol{\}}\AgdaSpace{}%
\AgdaBound{red}\AgdaSymbol{)}\AgdaSpace{}%
\AgdaSymbol{=}\AgdaSpace{}%
\AgdaInductiveConstructor{ζ}\AgdaSpace{}%
\AgdaSymbol{(}\AgdaInductiveConstructor{ξ}\AgdaSpace{}%
\AgdaSymbol{(}\AgdaFunction{pres-red↝βb²}\AgdaSpace{}%
\AgdaBound{M₁}\AgdaSpace{}%
\AgdaBound{M'}\AgdaSpace{}%
\AgdaBound{red}\AgdaSymbol{))}\<%
\end{code}  
\end{lemma}

\begin{lemma}
  \label{lem:cong:subst}
  \textbf{Congruence} of substitution:
  \small
  \begin{code}[hide]%
\>[0]\AgdaKeyword{module}\AgdaSpace{}%
\AgdaModule{valueProofs3}\AgdaSpace{}%
\AgdaKeyword{where}\<%
\\
\>[0]\<%
\\
\>[0]\AgdaKeyword{open}\AgdaSpace{}%
\AgdaKeyword{import}\AgdaSpace{}%
\AgdaModule{calc}\<%
\\
\>[0]\AgdaKeyword{open}\AgdaSpace{}%
\AgdaKeyword{import}\AgdaSpace{}%
\AgdaModule{boxCalc}\<%
\\
\>[0]\AgdaKeyword{open}\AgdaSpace{}%
\AgdaKeyword{import}\AgdaSpace{}%
\AgdaModule{embedCBVIntoCBB}\<%
\\
\>[0]\AgdaKeyword{open}\AgdaSpace{}%
\AgdaKeyword{import}\AgdaSpace{}%
\AgdaModule{valueCalcRed}\<%
\\
\>[0]\AgdaKeyword{open}\AgdaSpace{}%
\AgdaKeyword{import}\AgdaSpace{}%
\AgdaModule{boxCalcRed}\<%
\\
%
\\[\AgdaEmptyExtraSkip]%
\>[0]\AgdaKeyword{open}\AgdaSpace{}%
\AgdaKeyword{import}\AgdaSpace{}%
\AgdaModule{Data.Empty}\AgdaSpace{}%
\AgdaKeyword{using}\AgdaSpace{}%
\AgdaSymbol{(}\AgdaFunction{⊥}\AgdaSymbol{;}\AgdaSpace{}%
\AgdaFunction{⊥-elim}\AgdaSymbol{)}\<%
\\
\>[0]\AgdaKeyword{import}\AgdaSpace{}%
\AgdaModule{Relation.Binary.PropositionalEquality}\AgdaSpace{}%
\AgdaSymbol{as}\AgdaSpace{}%
\AgdaModule{Eq}\<%
\\
\>[0]\AgdaKeyword{open}\AgdaSpace{}%
\AgdaModule{Eq}\AgdaSpace{}%
\AgdaKeyword{using}\AgdaSpace{}%
\AgdaSymbol{(}\AgdaOperator{\AgdaDatatype{\AgdaUnderscore{}≡\AgdaUnderscore{}}}\AgdaSymbol{;}\AgdaSpace{}%
\AgdaInductiveConstructor{refl}\AgdaSymbol{;}\AgdaSpace{}%
\AgdaFunction{cong}\AgdaSymbol{;}\AgdaSpace{}%
\AgdaFunction{cong₂}\AgdaSymbol{;}\AgdaSpace{}%
\AgdaFunction{sym}\AgdaSymbol{)}\<%
\\
%
\\[\AgdaEmptyExtraSkip]%
\>[0]\AgdaKeyword{open}\AgdaSpace{}%
\AgdaModule{Eq.≡-Reasoning}\AgdaSpace{}%
\AgdaKeyword{using}\AgdaSpace{}%
\AgdaSymbol{(}\AgdaOperator{\AgdaFunction{begin\AgdaUnderscore{}}}\AgdaSymbol{;}\AgdaSpace{}%
\AgdaFunction{step-≡-∣}\AgdaSymbol{;}\AgdaSpace{}%
\AgdaFunction{step-≡-⟩}\AgdaSymbol{;}\AgdaSpace{}%
\AgdaOperator{\AgdaFunction{\AgdaUnderscore{}∎}}\AgdaSymbol{)}\<%
\\
\>[0]\AgdaKeyword{open}\AgdaSpace{}%
\AgdaModule{Eq.≡-Reasoning}\AgdaSpace{}%
\AgdaKeyword{renaming}\AgdaSpace{}%
\AgdaSymbol{(}\AgdaOperator{\AgdaFunction{begin\AgdaUnderscore{}}}\AgdaSpace{}%
\AgdaSymbol{to}\AgdaSpace{}%
\AgdaOperator{\AgdaFunction{beginEq\AgdaUnderscore{}}}\AgdaSymbol{;}\AgdaSpace{}%
\AgdaOperator{\AgdaFunction{\AgdaUnderscore{}∎}}\AgdaSpace{}%
\AgdaSymbol{to}\AgdaSpace{}%
\AgdaOperator{\AgdaFunction{\AgdaUnderscore{}∎eq}}\AgdaSymbol{)}\<%
\\
%
\\[\AgdaEmptyExtraSkip]%
\>[0]\AgdaComment{--\ Preservation\ of\ typing\ by\ Gödel's\ translation}\<%
\\
\>[0]\AgdaFunction{pres-typing}\AgdaSpace{}%
\AgdaSymbol{:}\AgdaSpace{}%
\AgdaSymbol{∀}\AgdaSpace{}%
\AgdaSymbol{\{}\AgdaBound{Γ}\AgdaSpace{}%
\AgdaBound{A}\AgdaSymbol{\}}\AgdaSpace{}%
\AgdaSymbol{→}\AgdaSpace{}%
\AgdaBound{Γ}\AgdaSpace{}%
\AgdaOperator{\AgdaDatatype{⊢}}\AgdaSpace{}%
\AgdaBound{A}\AgdaSpace{}%
\AgdaSymbol{→}\AgdaSpace{}%
\AgdaFunction{embedContext}\AgdaSpace{}%
\AgdaBound{Γ}\AgdaSpace{}%
\AgdaOperator{\AgdaDatatype{⊢b}}\AgdaSpace{}%
\AgdaFunction{embedType}\AgdaSpace{}%
\AgdaBound{A}\<%
\\
\>[0]\AgdaFunction{pres-typing}\AgdaSpace{}%
\AgdaBound{x}\AgdaSpace{}%
\AgdaSymbol{=}\AgdaSpace{}%
\AgdaFunction{embedTerm}\AgdaSpace{}%
\AgdaBound{x}\<%
\\
%
\\[\AgdaEmptyExtraSkip]%
\>[0]\AgdaComment{--\ Preservation\ of\ typing\ of\ Values\ by\ Gödel's\ translation}\<%
\\
\>[0]\AgdaFunction{pres-typing-val}\AgdaSpace{}%
\AgdaSymbol{:}\AgdaSpace{}%
\AgdaSymbol{∀}\AgdaSpace{}%
\AgdaSymbol{\{}\AgdaBound{Γ}\AgdaSpace{}%
\AgdaBound{A}\AgdaSymbol{\}}\AgdaSpace{}%
\AgdaSymbol{→}\AgdaSpace{}%
\AgdaSymbol{(}\AgdaBound{V}\AgdaSpace{}%
\AgdaSymbol{:}\AgdaSpace{}%
\AgdaBound{Γ}\AgdaSpace{}%
\AgdaOperator{\AgdaDatatype{⊢}}\AgdaSpace{}%
\AgdaBound{A}\AgdaSymbol{)}\AgdaSpace{}%
\AgdaSymbol{→}\AgdaSpace{}%
\AgdaSymbol{(}\AgdaDatatype{Value}\AgdaSpace{}%
\AgdaBound{V}\AgdaSymbol{)}\AgdaSpace{}%
\AgdaSymbol{→}\AgdaSpace{}%
\AgdaFunction{embedContext}\AgdaSpace{}%
\AgdaBound{Γ}\AgdaSpace{}%
\AgdaOperator{\AgdaDatatype{⊢b}}\AgdaSpace{}%
\AgdaFunction{embedTypeValue}\AgdaSpace{}%
\AgdaBound{A}\<%
\\
\>[0]\AgdaFunction{pres-typing-val}\AgdaSpace{}%
\AgdaBound{V}\AgdaSpace{}%
\AgdaBound{x}\AgdaSpace{}%
\AgdaSymbol{=}\AgdaSpace{}%
\AgdaFunction{embedValue}\AgdaSpace{}%
\AgdaBound{V}\AgdaSpace{}%
\AgdaBound{x}\<%
\\
%
\\[\AgdaEmptyExtraSkip]%
\>[0]\AgdaComment{--\ Helper\ lemma:\ applying\ embedTerm\ to\ a\ Value\ is\ equivalent\ to\ wrapping\ the\ result\ of\ embedValue\ in\ a\ box}\<%
\\
\>[0]\AgdaFunction{box-embedValue}\AgdaSpace{}%
\AgdaSymbol{:}\AgdaSpace{}%
\AgdaSymbol{∀}\AgdaSpace{}%
\AgdaSymbol{\{}\AgdaBound{Γ}\AgdaSpace{}%
\AgdaBound{A}\AgdaSymbol{\}}\AgdaSpace{}%
\AgdaSymbol{→}\AgdaSpace{}%
\AgdaSymbol{\{}\AgdaBound{V}\AgdaSpace{}%
\AgdaSymbol{:}\AgdaSpace{}%
\AgdaBound{Γ}\AgdaSpace{}%
\AgdaOperator{\AgdaDatatype{⊢}}\AgdaSpace{}%
\AgdaBound{A}\AgdaSymbol{\}}\AgdaSpace{}%
\AgdaSymbol{→}\AgdaSpace{}%
\AgdaSymbol{(}\AgdaBound{x}\AgdaSpace{}%
\AgdaSymbol{:}\AgdaSpace{}%
\AgdaDatatype{Value}\AgdaSpace{}%
\AgdaBound{V}\AgdaSymbol{)}\AgdaSpace{}%
\AgdaSymbol{→}\AgdaSpace{}%
\AgdaInductiveConstructor{box}\AgdaSpace{}%
\AgdaSymbol{(}\AgdaFunction{embedValue}\AgdaSpace{}%
\AgdaBound{V}\AgdaSpace{}%
\AgdaBound{x}\AgdaSymbol{)}\AgdaSpace{}%
\AgdaOperator{\AgdaDatatype{≡}}\AgdaSpace{}%
\AgdaFunction{embedTerm}\AgdaSpace{}%
\AgdaBound{V}\<%
\\
\>[0]\AgdaFunction{box-embedValue}\AgdaSpace{}%
\AgdaInductiveConstructor{V-ƛ}\AgdaSpace{}%
\AgdaSymbol{=}\AgdaSpace{}%
\AgdaInductiveConstructor{refl}\<%
\\
\>[0]\AgdaFunction{box-embedValue}\AgdaSpace{}%
\AgdaInductiveConstructor{V-`}\AgdaSpace{}%
\AgdaSymbol{=}\AgdaSpace{}%
\AgdaInductiveConstructor{refl}\<%
\\
%
\\[\AgdaEmptyExtraSkip]%
\>[0]\AgdaComment{--\ HELPER\ LEMMA'S\ SIMULTANIOUS\ SUBSTITUTION}\<%
\\
%
\\[\AgdaEmptyExtraSkip]%
\>[0]\AgdaComment{--\ Preservation\ and\ Reflection\ of\ }\<%
\\
\>[0]\AgdaComment{--\ ext\ of\ Gödel's\ translation}\<%
\\
\>[0]\AgdaFunction{pres-ext}\AgdaSpace{}%
\AgdaSymbol{:}\AgdaSpace{}%
\AgdaSymbol{∀}\AgdaSpace{}%
\AgdaSymbol{\{}\AgdaBound{Γ}\AgdaSpace{}%
\AgdaBound{Δ}\AgdaSpace{}%
\AgdaBound{A'}\AgdaSpace{}%
\AgdaBound{A}\AgdaSymbol{\}}\<%
\\
\>[0][@{}l@{\AgdaIndent{0}}]%
\>[2]\AgdaSymbol{→}\AgdaSpace{}%
\AgdaSymbol{(}\AgdaBound{f'}\AgdaSpace{}%
\AgdaSymbol{:}\AgdaSpace{}%
\AgdaSymbol{∀}\AgdaSpace{}%
\AgdaSymbol{\{}\AgdaBound{A}\AgdaSymbol{\}}\AgdaSpace{}%
\AgdaSymbol{→}\AgdaSpace{}%
\AgdaBound{A}\AgdaSpace{}%
\AgdaOperator{\AgdaDatatype{∈}}\AgdaSpace{}%
\AgdaBound{Γ}\AgdaSpace{}%
\AgdaSymbol{→}\AgdaSpace{}%
\AgdaBound{A}\AgdaSpace{}%
\AgdaOperator{\AgdaDatatype{∈}}\AgdaSpace{}%
\AgdaBound{Δ}\AgdaSymbol{)}\<%
\\
%
\>[2]\AgdaSymbol{→}\AgdaSpace{}%
\AgdaSymbol{(}\AgdaBound{fb'}\AgdaSpace{}%
\AgdaSymbol{:}\AgdaSpace{}%
\AgdaSymbol{∀}\AgdaSpace{}%
\AgdaSymbol{\{}\AgdaBound{Ab}\AgdaSymbol{\}}\AgdaSpace{}%
\AgdaSymbol{→}\AgdaSpace{}%
\AgdaBound{Ab}\AgdaSpace{}%
\AgdaOperator{\AgdaDatatype{∈b}}\AgdaSpace{}%
\AgdaFunction{embedContext}\AgdaSpace{}%
\AgdaBound{Γ}\AgdaSpace{}%
\AgdaSymbol{→}\AgdaSpace{}%
\AgdaBound{Ab}\AgdaSpace{}%
\AgdaOperator{\AgdaDatatype{∈b}}\AgdaSpace{}%
\AgdaFunction{embedContext}\AgdaSpace{}%
\AgdaBound{Δ}\AgdaSymbol{)}\<%
\\
%
\>[2]\AgdaSymbol{→}\AgdaSpace{}%
\AgdaSymbol{(\{}\AgdaBound{A'}\AgdaSpace{}%
\AgdaSymbol{:}\AgdaSpace{}%
\AgdaDatatype{Type}\AgdaSymbol{\}}\AgdaSpace{}%
\AgdaSymbol{(}\AgdaBound{x}\AgdaSpace{}%
\AgdaSymbol{:}\AgdaSpace{}%
\AgdaBound{A'}\AgdaSpace{}%
\AgdaOperator{\AgdaDatatype{∈}}\AgdaSpace{}%
\AgdaBound{Γ}\AgdaSymbol{)}\AgdaSpace{}%
\AgdaSymbol{→}\AgdaSpace{}%
\AgdaBound{fb'}\AgdaSpace{}%
\AgdaSymbol{(}\AgdaFunction{embedLookup}\AgdaSpace{}%
\AgdaBound{x}\AgdaSymbol{)}\AgdaSpace{}%
\AgdaOperator{\AgdaDatatype{≡}}\AgdaSpace{}%
\AgdaFunction{embedLookup}\AgdaSpace{}%
\AgdaSymbol{(}\AgdaBound{f'}\AgdaSpace{}%
\AgdaBound{x}\AgdaSymbol{))}\<%
\\
%
\>[2]\AgdaSymbol{→}\AgdaSpace{}%
\AgdaSymbol{(}\AgdaBound{x}\AgdaSpace{}%
\AgdaSymbol{:}\AgdaSpace{}%
\AgdaBound{A'}\AgdaSpace{}%
\AgdaOperator{\AgdaDatatype{∈}}\AgdaSpace{}%
\AgdaBound{Γ}\AgdaSpace{}%
\AgdaOperator{\AgdaInductiveConstructor{,}}\AgdaSpace{}%
\AgdaBound{A}\AgdaSymbol{)}\<%
\\
%
\>[2]\AgdaSymbol{→}\AgdaSpace{}%
\AgdaFunction{extb}\AgdaSpace{}%
\AgdaBound{fb'}\AgdaSpace{}%
\AgdaSymbol{(}\AgdaFunction{embedLookup}\AgdaSpace{}%
\AgdaBound{x}\AgdaSymbol{)}\AgdaSpace{}%
\AgdaOperator{\AgdaDatatype{≡}}\AgdaSpace{}%
\AgdaFunction{embedLookup}\AgdaSpace{}%
\AgdaSymbol{(}\AgdaFunction{ext}\AgdaSpace{}%
\AgdaBound{f'}\AgdaSpace{}%
\AgdaBound{x}\AgdaSymbol{)}\<%
\\
\>[0]\AgdaFunction{pres-ext}\AgdaSpace{}%
\AgdaBound{f'}\AgdaSpace{}%
\AgdaBound{fb'}\AgdaSpace{}%
\AgdaBound{h}\AgdaSpace{}%
\AgdaInductiveConstructor{Z}\AgdaSpace{}%
\AgdaSymbol{=}\AgdaSpace{}%
\AgdaInductiveConstructor{refl}\<%
\\
\>[0]\AgdaFunction{pres-ext}\AgdaSpace{}%
\AgdaBound{f'}\AgdaSpace{}%
\AgdaBound{fb'}\AgdaSpace{}%
\AgdaBound{h}\AgdaSpace{}%
\AgdaSymbol{(}\AgdaInductiveConstructor{S}\AgdaSpace{}%
\AgdaBound{x}\AgdaSymbol{)}\AgdaSpace{}%
\AgdaSymbol{=}\AgdaSpace{}%
\AgdaFunction{cong}\AgdaSpace{}%
\AgdaInductiveConstructor{Sb}\AgdaSpace{}%
\AgdaSymbol{(}\AgdaBound{h}\AgdaSpace{}%
\AgdaBound{x}\AgdaSymbol{)}\<%
\\
%
\\[\AgdaEmptyExtraSkip]%
\>[0]\AgdaComment{--\ Extb\ preserves\ pointwise\ equality\ }\<%
\\
\>[0]\AgdaComment{--\ of\ renaming\ functions}\<%
\\
\>[0]\AgdaFunction{extb-cong}\AgdaSpace{}%
\AgdaSymbol{:}\AgdaSpace{}%
\AgdaSymbol{∀}\AgdaSpace{}%
\AgdaSymbol{\{}\AgdaBound{Γ₁}\AgdaSpace{}%
\AgdaBound{Γ₂}\AgdaSpace{}%
\AgdaBound{A}\AgdaSpace{}%
\AgdaBound{A'}\AgdaSymbol{\}}\<%
\\
\>[0][@{}l@{\AgdaIndent{0}}]%
\>[2]\AgdaSymbol{→}\AgdaSpace{}%
\AgdaSymbol{\{}\AgdaBound{fb'}\AgdaSpace{}%
\AgdaSymbol{:}\AgdaSpace{}%
\AgdaSymbol{∀}\AgdaSpace{}%
\AgdaSymbol{\{}\AgdaBound{Ab}\AgdaSymbol{\}}\AgdaSpace{}%
\AgdaSymbol{→}\AgdaSpace{}%
\AgdaBound{Ab}\AgdaSpace{}%
\AgdaOperator{\AgdaDatatype{∈b}}\AgdaSpace{}%
\AgdaBound{Γ₁}\AgdaSpace{}%
\AgdaSymbol{→}\AgdaSpace{}%
\AgdaBound{Ab}\AgdaSpace{}%
\AgdaOperator{\AgdaDatatype{∈b}}\AgdaSpace{}%
\AgdaBound{Γ₂}\AgdaSymbol{\}}\<%
\\
%
\>[2]\AgdaSymbol{→}\AgdaSpace{}%
\AgdaSymbol{\{}\AgdaBound{fb''}\AgdaSpace{}%
\AgdaSymbol{:}\AgdaSpace{}%
\AgdaSymbol{∀}\AgdaSpace{}%
\AgdaSymbol{\{}\AgdaBound{Ab}\AgdaSymbol{\}}\AgdaSpace{}%
\AgdaSymbol{→}\AgdaSpace{}%
\AgdaBound{Ab}\AgdaSpace{}%
\AgdaOperator{\AgdaDatatype{∈b}}\AgdaSpace{}%
\AgdaBound{Γ₁}\AgdaSpace{}%
\AgdaSymbol{→}\AgdaSpace{}%
\AgdaBound{Ab}\AgdaSpace{}%
\AgdaOperator{\AgdaDatatype{∈b}}\AgdaSpace{}%
\AgdaBound{Γ₂}\AgdaSymbol{\}}\<%
\\
%
\>[2]\AgdaSymbol{→}\AgdaSpace{}%
\AgdaSymbol{(}\AgdaBound{h}\AgdaSpace{}%
\AgdaSymbol{:}\AgdaSpace{}%
\AgdaSymbol{∀}\AgdaSpace{}%
\AgdaSymbol{\{}\AgdaBound{Ab}\AgdaSymbol{\}}\AgdaSpace{}%
\AgdaSymbol{→}\AgdaSpace{}%
\AgdaSymbol{(}\AgdaBound{x}\AgdaSpace{}%
\AgdaSymbol{:}\AgdaSpace{}%
\AgdaBound{Ab}\AgdaSpace{}%
\AgdaOperator{\AgdaDatatype{∈b}}\AgdaSpace{}%
\AgdaBound{Γ₁}\AgdaSymbol{)}\AgdaSpace{}%
\AgdaSymbol{→}\AgdaSpace{}%
\AgdaBound{fb'}\AgdaSpace{}%
\AgdaBound{x}\AgdaSpace{}%
\AgdaOperator{\AgdaDatatype{≡}}\AgdaSpace{}%
\AgdaBound{fb''}\AgdaSpace{}%
\AgdaBound{x}\AgdaSymbol{)}\<%
\\
%
\>[2]\AgdaSymbol{→}\AgdaSpace{}%
\AgdaSymbol{(}\AgdaBound{x}\AgdaSpace{}%
\AgdaSymbol{:}\AgdaSpace{}%
\AgdaBound{A}\AgdaSpace{}%
\AgdaOperator{\AgdaDatatype{∈b}}\AgdaSpace{}%
\AgdaBound{Γ₁}\AgdaSpace{}%
\AgdaOperator{\AgdaInductiveConstructor{,b}}\AgdaSpace{}%
\AgdaInductiveConstructor{□}\AgdaSpace{}%
\AgdaBound{A'}\AgdaSymbol{)}\<%
\\
%
\>[2]\AgdaSymbol{→}\AgdaSpace{}%
\AgdaFunction{extb}\AgdaSpace{}%
\AgdaBound{fb'}\AgdaSpace{}%
\AgdaBound{x}\AgdaSpace{}%
\AgdaOperator{\AgdaDatatype{≡}}\AgdaSpace{}%
\AgdaFunction{extb}\AgdaSpace{}%
\AgdaBound{fb''}\AgdaSpace{}%
\AgdaBound{x}\<%
\\
\>[0]\AgdaFunction{extb-cong}\AgdaSpace{}%
\AgdaBound{h}\AgdaSpace{}%
\AgdaInductiveConstructor{Zb}\AgdaSpace{}%
\AgdaSymbol{=}\AgdaSpace{}%
\AgdaInductiveConstructor{refl}\<%
\\
\>[0]\AgdaFunction{extb-cong}\AgdaSpace{}%
\AgdaBound{h}\AgdaSpace{}%
\AgdaSymbol{(}\AgdaInductiveConstructor{Sb}\AgdaSpace{}%
\AgdaBound{x}\AgdaSymbol{)}\AgdaSpace{}%
\AgdaSymbol{=}\AgdaSpace{}%
\AgdaFunction{cong}\AgdaSpace{}%
\AgdaInductiveConstructor{Sb}\AgdaSpace{}%
\AgdaSymbol{(}\AgdaBound{h}\AgdaSpace{}%
\AgdaBound{x}\AgdaSymbol{)}\<%
\\
%
\\[\AgdaEmptyExtraSkip]%
\>[0]\AgdaComment{--\ Renameb\ preserves\ pointwise\ equality}\<%
\\
\>[0]\AgdaComment{--\ of\ renaming\ functions}\<%
\\
\>[0]\AgdaFunction{renameb-cong}\AgdaSpace{}%
\AgdaSymbol{:}\AgdaSpace{}%
\AgdaSymbol{∀}\AgdaSpace{}%
\AgdaSymbol{\{}\AgdaBound{Γ₁}\AgdaSpace{}%
\AgdaBound{Γ₂}\AgdaSpace{}%
\AgdaBound{A}\AgdaSymbol{\}}\<%
\\
\>[0][@{}l@{\AgdaIndent{0}}]%
\>[2]\AgdaSymbol{→}\AgdaSpace{}%
\AgdaSymbol{\{}\AgdaBound{fb'}\AgdaSpace{}%
\AgdaSymbol{:}\AgdaSpace{}%
\AgdaSymbol{∀}\AgdaSpace{}%
\AgdaSymbol{\{}\AgdaBound{Ab}\AgdaSymbol{\}}\AgdaSpace{}%
\AgdaSymbol{→}\AgdaSpace{}%
\AgdaBound{Ab}\AgdaSpace{}%
\AgdaOperator{\AgdaDatatype{∈b}}\AgdaSpace{}%
\AgdaBound{Γ₁}\AgdaSpace{}%
\AgdaSymbol{→}\AgdaSpace{}%
\AgdaBound{Ab}\AgdaSpace{}%
\AgdaOperator{\AgdaDatatype{∈b}}\AgdaSpace{}%
\AgdaBound{Γ₂}\AgdaSymbol{\}}\<%
\\
%
\>[2]\AgdaSymbol{→}\AgdaSpace{}%
\AgdaSymbol{\{}\AgdaBound{fb''}\AgdaSpace{}%
\AgdaSymbol{:}\AgdaSpace{}%
\AgdaSymbol{∀}\AgdaSpace{}%
\AgdaSymbol{\{}\AgdaBound{Ab}\AgdaSymbol{\}}\AgdaSpace{}%
\AgdaSymbol{→}\AgdaSpace{}%
\AgdaBound{Ab}\AgdaSpace{}%
\AgdaOperator{\AgdaDatatype{∈b}}\AgdaSpace{}%
\AgdaBound{Γ₁}\AgdaSpace{}%
\AgdaSymbol{→}\AgdaSpace{}%
\AgdaBound{Ab}\AgdaSpace{}%
\AgdaOperator{\AgdaDatatype{∈b}}\AgdaSpace{}%
\AgdaBound{Γ₂}\AgdaSymbol{\}}\<%
\\
%
\>[2]\AgdaSymbol{→}\AgdaSpace{}%
\AgdaSymbol{(}\AgdaBound{h}\AgdaSpace{}%
\AgdaSymbol{:}\AgdaSpace{}%
\AgdaSymbol{∀}\AgdaSpace{}%
\AgdaSymbol{\{}\AgdaBound{Ab}\AgdaSymbol{\}}\AgdaSpace{}%
\AgdaSymbol{→}\AgdaSpace{}%
\AgdaSymbol{(}\AgdaBound{x}\AgdaSpace{}%
\AgdaSymbol{:}\AgdaSpace{}%
\AgdaBound{Ab}\AgdaSpace{}%
\AgdaOperator{\AgdaDatatype{∈b}}\AgdaSpace{}%
\AgdaBound{Γ₁}\AgdaSymbol{)}\AgdaSpace{}%
\AgdaSymbol{→}\AgdaSpace{}%
\AgdaBound{fb'}\AgdaSpace{}%
\AgdaBound{x}\AgdaSpace{}%
\AgdaOperator{\AgdaDatatype{≡}}\AgdaSpace{}%
\AgdaBound{fb''}\AgdaSpace{}%
\AgdaBound{x}\AgdaSymbol{)}\<%
\\
%
\>[2]\AgdaSymbol{→}\AgdaSpace{}%
\AgdaSymbol{(}\AgdaBound{t}\AgdaSpace{}%
\AgdaSymbol{:}\AgdaSpace{}%
\AgdaBound{Γ₁}\AgdaSpace{}%
\AgdaOperator{\AgdaDatatype{⊢b}}\AgdaSpace{}%
\AgdaBound{A}\AgdaSymbol{)}\<%
\\
%
\>[2]\AgdaSymbol{→}\AgdaSpace{}%
\AgdaFunction{renameb}\AgdaSpace{}%
\AgdaBound{fb'}\AgdaSpace{}%
\AgdaBound{t}\AgdaSpace{}%
\AgdaOperator{\AgdaDatatype{≡}}\AgdaSpace{}%
\AgdaFunction{renameb}\AgdaSpace{}%
\AgdaBound{fb''}\AgdaSpace{}%
\AgdaBound{t}\<%
\\
\>[0]\AgdaFunction{renameb-cong}\AgdaSpace{}%
\AgdaBound{h}\AgdaSpace{}%
\AgdaSymbol{(}\AgdaInductiveConstructor{ε}\AgdaSpace{}%
\AgdaBound{x}\AgdaSymbol{)}\AgdaSpace{}%
\AgdaSymbol{=}\AgdaSpace{}%
\AgdaFunction{cong}\AgdaSpace{}%
\AgdaInductiveConstructor{ε}\AgdaSpace{}%
\AgdaSymbol{(}\AgdaBound{h}\AgdaSpace{}%
\AgdaBound{x}\AgdaSymbol{)}\<%
\\
\>[0]\AgdaFunction{renameb-cong}\AgdaSpace{}%
\AgdaBound{h}\AgdaSpace{}%
\AgdaSymbol{(}\AgdaInductiveConstructor{box}\AgdaSpace{}%
\AgdaBound{t}\AgdaSymbol{)}\AgdaSpace{}%
\AgdaSymbol{=}\AgdaSpace{}%
\AgdaFunction{cong}\AgdaSpace{}%
\AgdaInductiveConstructor{box}\AgdaSpace{}%
\AgdaSymbol{(}\AgdaFunction{renameb-cong}\AgdaSpace{}%
\AgdaBound{h}\AgdaSpace{}%
\AgdaBound{t}\AgdaSymbol{)}\<%
\\
\>[0]\AgdaFunction{renameb-cong}\AgdaSpace{}%
\AgdaBound{h}\AgdaSpace{}%
\AgdaSymbol{(}\AgdaInductiveConstructor{ƛb}\AgdaSpace{}%
\AgdaBound{t}\AgdaSymbol{)}\AgdaSpace{}%
\AgdaSymbol{=}\AgdaSpace{}%
\AgdaFunction{cong}\AgdaSpace{}%
\AgdaInductiveConstructor{ƛb}\AgdaSpace{}%
\AgdaSymbol{(}\AgdaFunction{renameb-cong}\AgdaSpace{}%
\AgdaSymbol{(}\AgdaFunction{extb-cong}\AgdaSpace{}%
\AgdaBound{h}\AgdaSymbol{)}\AgdaSpace{}%
\AgdaBound{t}\AgdaSymbol{)}\<%
\\
\>[0]\AgdaFunction{renameb-cong}\AgdaSpace{}%
\AgdaBound{h}\AgdaSpace{}%
\AgdaSymbol{(}\AgdaBound{t}\AgdaSpace{}%
\AgdaOperator{\AgdaInductiveConstructor{·b}}\AgdaSpace{}%
\AgdaBound{t₁}\AgdaSymbol{)}\AgdaSpace{}%
\AgdaSymbol{=}\AgdaSpace{}%
\AgdaFunction{cong₂}\AgdaSpace{}%
\AgdaOperator{\AgdaInductiveConstructor{\AgdaUnderscore{}·b\AgdaUnderscore{}}}\AgdaSpace{}%
\AgdaSymbol{(}\AgdaFunction{renameb-cong}\AgdaSpace{}%
\AgdaBound{h}\AgdaSpace{}%
\AgdaBound{t}\AgdaSymbol{)}\AgdaSpace{}%
\AgdaSymbol{(}\AgdaFunction{renameb-cong}\AgdaSpace{}%
\AgdaBound{h}\AgdaSpace{}%
\AgdaBound{t₁}\AgdaSymbol{)}\<%
\\
%
\\[\AgdaEmptyExtraSkip]%
\>[0]\AgdaComment{--\ extb\ distributes\ over\ composition:}\<%
\\
\>[0]\AgdaComment{--\ Extending\ first\ with\ fb',\ then\ with\ fb'',\ }\<%
\\
\>[0]\AgdaComment{--\ is\ the\ same\ as\ extending\ once\ with\ fb''\ ∘\ fb'.}\<%
\\
\>[0]\AgdaFunction{extb-comp}\AgdaSpace{}%
\AgdaSymbol{:}\AgdaSpace{}%
\AgdaSymbol{∀}\AgdaSpace{}%
\AgdaSymbol{\{}\AgdaBound{Γ₁}\AgdaSpace{}%
\AgdaBound{Γ₂}\AgdaSpace{}%
\AgdaBound{Γ₃}\AgdaSpace{}%
\AgdaBound{A}\AgdaSpace{}%
\AgdaBound{B}\AgdaSymbol{\}}\<%
\\
\>[0][@{}l@{\AgdaIndent{0}}]%
\>[2]\AgdaSymbol{→}\AgdaSpace{}%
\AgdaSymbol{(}\AgdaBound{fb'}\AgdaSpace{}%
\AgdaSymbol{:}\AgdaSpace{}%
\AgdaSymbol{∀}\AgdaSpace{}%
\AgdaSymbol{\{}\AgdaBound{Ab}\AgdaSymbol{\}}\AgdaSpace{}%
\AgdaSymbol{→}\AgdaSpace{}%
\AgdaBound{Ab}\AgdaSpace{}%
\AgdaOperator{\AgdaDatatype{∈b}}\AgdaSpace{}%
\AgdaBound{Γ₁}\AgdaSpace{}%
\AgdaSymbol{→}\AgdaSpace{}%
\AgdaBound{Ab}\AgdaSpace{}%
\AgdaOperator{\AgdaDatatype{∈b}}\AgdaSpace{}%
\AgdaBound{Γ₂}\AgdaSymbol{)}\<%
\\
%
\>[2]\AgdaSymbol{→}\AgdaSpace{}%
\AgdaSymbol{(}\AgdaBound{fb''}\AgdaSpace{}%
\AgdaSymbol{:}\AgdaSpace{}%
\AgdaSymbol{∀}\AgdaSpace{}%
\AgdaSymbol{\{}\AgdaBound{Ab}\AgdaSymbol{\}}\AgdaSpace{}%
\AgdaSymbol{→}\AgdaSpace{}%
\AgdaBound{Ab}\AgdaSpace{}%
\AgdaOperator{\AgdaDatatype{∈b}}\AgdaSpace{}%
\AgdaBound{Γ₂}\AgdaSpace{}%
\AgdaSymbol{→}\AgdaSpace{}%
\AgdaBound{Ab}\AgdaSpace{}%
\AgdaOperator{\AgdaDatatype{∈b}}\AgdaSpace{}%
\AgdaBound{Γ₃}\AgdaSymbol{)}\<%
\\
%
\>[2]\AgdaSymbol{→}\AgdaSpace{}%
\AgdaSymbol{(}\AgdaBound{x}\AgdaSpace{}%
\AgdaSymbol{:}\AgdaSpace{}%
\AgdaBound{A}\AgdaSpace{}%
\AgdaOperator{\AgdaDatatype{∈b}}\AgdaSpace{}%
\AgdaBound{Γ₁}\AgdaSpace{}%
\AgdaOperator{\AgdaInductiveConstructor{,b}}\AgdaSpace{}%
\AgdaBound{B}\AgdaSymbol{)}\<%
\\
%
\>[2]\AgdaSymbol{→}\AgdaSpace{}%
\AgdaFunction{extb}\AgdaSpace{}%
\AgdaBound{fb''}\AgdaSpace{}%
\AgdaSymbol{(}\AgdaFunction{extb}\AgdaSpace{}%
\AgdaBound{fb'}\AgdaSpace{}%
\AgdaBound{x}\AgdaSymbol{)}\AgdaSpace{}%
\AgdaOperator{\AgdaDatatype{≡}}\AgdaSpace{}%
\AgdaFunction{extb}\AgdaSpace{}%
\AgdaSymbol{(λ}\AgdaSpace{}%
\AgdaBound{x₁}\AgdaSpace{}%
\AgdaSymbol{→}\AgdaSpace{}%
\AgdaBound{fb''}\AgdaSpace{}%
\AgdaSymbol{(}\AgdaBound{fb'}\AgdaSpace{}%
\AgdaBound{x₁}\AgdaSymbol{))}\AgdaSpace{}%
\AgdaBound{x}\<%
\\
\>[0]\AgdaFunction{extb-comp}\AgdaSpace{}%
\AgdaBound{fb'}\AgdaSpace{}%
\AgdaBound{fb''}\AgdaSpace{}%
\AgdaInductiveConstructor{Zb}\AgdaSpace{}%
\AgdaSymbol{=}\AgdaSpace{}%
\AgdaInductiveConstructor{refl}\<%
\\
\>[0]\AgdaFunction{extb-comp}\AgdaSpace{}%
\AgdaBound{fb'}\AgdaSpace{}%
\AgdaBound{fb''}\AgdaSpace{}%
\AgdaSymbol{(}\AgdaInductiveConstructor{Sb}\AgdaSpace{}%
\AgdaBound{x}\AgdaSymbol{)}\AgdaSpace{}%
\AgdaSymbol{=}\AgdaSpace{}%
\AgdaInductiveConstructor{refl}\<%
\\
%
\\[\AgdaEmptyExtraSkip]%
\>[0]\AgdaComment{--\ renameb\ distributes\ over\ composition:}\<%
\\
\>[0]\AgdaComment{--\ Renaming\ first\ with\ fb',\ then\ with\ fb''}\<%
\\
\>[0]\AgdaComment{--\ is\ the\ same\ as\ renaming\ once\ with\ fb''\ ∘\ fb'.}\<%
\\
\>[0]\AgdaFunction{renameb-comp}\AgdaSpace{}%
\AgdaSymbol{:}\AgdaSpace{}%
\AgdaSymbol{∀}\AgdaSpace{}%
\AgdaSymbol{\{}\AgdaBound{Γ₁}\AgdaSpace{}%
\AgdaBound{Γ₂}\AgdaSpace{}%
\AgdaBound{Γ₃}\AgdaSpace{}%
\AgdaBound{A}\AgdaSymbol{\}}\<%
\\
\>[0][@{}l@{\AgdaIndent{0}}]%
\>[2]\AgdaSymbol{→}\AgdaSpace{}%
\AgdaSymbol{(}\AgdaBound{fb'}\AgdaSpace{}%
\AgdaSymbol{:}\AgdaSpace{}%
\AgdaSymbol{∀}\AgdaSpace{}%
\AgdaSymbol{\{}\AgdaBound{Ab}\AgdaSymbol{\}}\AgdaSpace{}%
\AgdaSymbol{→}\AgdaSpace{}%
\AgdaBound{Ab}\AgdaSpace{}%
\AgdaOperator{\AgdaDatatype{∈b}}\AgdaSpace{}%
\AgdaBound{Γ₁}\AgdaSpace{}%
\AgdaSymbol{→}\AgdaSpace{}%
\AgdaBound{Ab}\AgdaSpace{}%
\AgdaOperator{\AgdaDatatype{∈b}}\AgdaSpace{}%
\AgdaBound{Γ₂}\AgdaSymbol{)}\<%
\\
%
\>[2]\AgdaSymbol{→}\AgdaSpace{}%
\AgdaSymbol{(}\AgdaBound{fb''}\AgdaSpace{}%
\AgdaSymbol{:}\AgdaSpace{}%
\AgdaSymbol{∀}\AgdaSpace{}%
\AgdaSymbol{\{}\AgdaBound{Ab}\AgdaSymbol{\}}\AgdaSpace{}%
\AgdaSymbol{→}\AgdaSpace{}%
\AgdaBound{Ab}\AgdaSpace{}%
\AgdaOperator{\AgdaDatatype{∈b}}\AgdaSpace{}%
\AgdaBound{Γ₂}\AgdaSpace{}%
\AgdaSymbol{→}\AgdaSpace{}%
\AgdaBound{Ab}\AgdaSpace{}%
\AgdaOperator{\AgdaDatatype{∈b}}\AgdaSpace{}%
\AgdaBound{Γ₃}\AgdaSymbol{)}\<%
\\
%
\>[2]\AgdaSymbol{→}\AgdaSpace{}%
\AgdaSymbol{(}\AgdaBound{t}\AgdaSpace{}%
\AgdaSymbol{:}\AgdaSpace{}%
\AgdaBound{Γ₁}\AgdaSpace{}%
\AgdaOperator{\AgdaDatatype{⊢b}}\AgdaSpace{}%
\AgdaBound{A}\AgdaSymbol{)}\<%
\\
%
\>[2]\AgdaSymbol{→}\AgdaSpace{}%
\AgdaFunction{renameb}\AgdaSpace{}%
\AgdaBound{fb''}\AgdaSpace{}%
\AgdaSymbol{(}\AgdaFunction{renameb}\AgdaSpace{}%
\AgdaBound{fb'}\AgdaSpace{}%
\AgdaBound{t}\AgdaSymbol{)}\AgdaSpace{}%
\AgdaOperator{\AgdaDatatype{≡}}\AgdaSpace{}%
\AgdaFunction{renameb}\AgdaSpace{}%
\AgdaSymbol{(λ}\AgdaSpace{}%
\AgdaBound{x}\AgdaSpace{}%
\AgdaSymbol{→}\AgdaSpace{}%
\AgdaBound{fb''}\AgdaSpace{}%
\AgdaSymbol{(}\AgdaBound{fb'}\AgdaSpace{}%
\AgdaBound{x}\AgdaSymbol{))}\AgdaSpace{}%
\AgdaBound{t}\<%
\\
\>[0]\AgdaFunction{renameb-comp}\AgdaSpace{}%
\AgdaBound{fb'}\AgdaSpace{}%
\AgdaBound{fb''}\AgdaSpace{}%
\AgdaSymbol{(}\AgdaInductiveConstructor{ε}\AgdaSpace{}%
\AgdaBound{x}\AgdaSymbol{)}\AgdaSpace{}%
\AgdaSymbol{=}\AgdaSpace{}%
\AgdaInductiveConstructor{refl}\<%
\\
\>[0]\AgdaFunction{renameb-comp}\AgdaSpace{}%
\AgdaBound{fb'}\AgdaSpace{}%
\AgdaBound{fb''}\AgdaSpace{}%
\AgdaSymbol{(}\AgdaInductiveConstructor{box}\AgdaSpace{}%
\AgdaBound{t}\AgdaSymbol{)}\AgdaSpace{}%
\AgdaSymbol{=}\AgdaSpace{}%
\AgdaFunction{cong}\AgdaSpace{}%
\AgdaInductiveConstructor{box}\AgdaSpace{}%
\AgdaSymbol{(}\AgdaFunction{renameb-comp}\AgdaSpace{}%
\AgdaBound{fb'}\AgdaSpace{}%
\AgdaBound{fb''}\AgdaSpace{}%
\AgdaBound{t}\AgdaSymbol{)}\<%
\\
\>[0]\AgdaFunction{renameb-comp}\AgdaSpace{}%
\AgdaBound{fb'}\AgdaSpace{}%
\AgdaBound{fb''}\AgdaSpace{}%
\AgdaSymbol{(}\AgdaInductiveConstructor{ƛb}\AgdaSpace{}%
\AgdaBound{t}\AgdaSymbol{)}\AgdaSpace{}%
\AgdaSymbol{=}\<%
\\
\>[0][@{}l@{\AgdaIndent{0}}]%
\>[2]\AgdaOperator{\AgdaFunction{beginEq}}\<%
\\
\>[2][@{}l@{\AgdaIndent{0}}]%
\>[4]\AgdaInductiveConstructor{ƛb}\AgdaSpace{}%
\AgdaSymbol{(}\AgdaFunction{renameb}\AgdaSpace{}%
\AgdaSymbol{(}\AgdaFunction{extb}\AgdaSpace{}%
\AgdaBound{fb''}\AgdaSymbol{)}\AgdaSpace{}%
\AgdaSymbol{(}\AgdaFunction{renameb}\AgdaSpace{}%
\AgdaSymbol{(}\AgdaFunction{extb}\AgdaSpace{}%
\AgdaBound{fb'}\AgdaSymbol{)}\AgdaSpace{}%
\AgdaBound{t}\AgdaSymbol{))}\<%
\\
%
\>[2]\AgdaFunction{≡⟨}\AgdaSpace{}%
\AgdaFunction{cong}\AgdaSpace{}%
\AgdaInductiveConstructor{ƛb}\AgdaSpace{}%
\AgdaSymbol{(}\AgdaFunction{renameb-comp}\AgdaSpace{}%
\AgdaSymbol{(}\AgdaFunction{extb}\AgdaSpace{}%
\AgdaBound{fb'}\AgdaSymbol{)}\AgdaSpace{}%
\AgdaSymbol{(}\AgdaFunction{extb}\AgdaSpace{}%
\AgdaBound{fb''}\AgdaSymbol{)}\AgdaSpace{}%
\AgdaBound{t}\AgdaSymbol{)}\AgdaSpace{}%
\AgdaFunction{⟩}\<%
\\
\>[2][@{}l@{\AgdaIndent{0}}]%
\>[4]\AgdaInductiveConstructor{ƛb}\AgdaSpace{}%
\AgdaSymbol{(}\AgdaFunction{renameb}\AgdaSpace{}%
\AgdaSymbol{(λ}\AgdaSpace{}%
\AgdaBound{x}\AgdaSpace{}%
\AgdaSymbol{→}\AgdaSpace{}%
\AgdaFunction{extb}\AgdaSpace{}%
\AgdaBound{fb''}\AgdaSpace{}%
\AgdaSymbol{(}\AgdaFunction{extb}\AgdaSpace{}%
\AgdaBound{fb'}\AgdaSpace{}%
\AgdaBound{x}\AgdaSymbol{))}\AgdaSpace{}%
\AgdaBound{t}\AgdaSymbol{)}\<%
\\
%
\>[2]\AgdaFunction{≡⟨}\AgdaSpace{}%
\AgdaFunction{cong}\AgdaSpace{}%
\AgdaInductiveConstructor{ƛb}\AgdaSpace{}%
\AgdaSymbol{(}\AgdaFunction{renameb-cong}\AgdaSpace{}%
\AgdaSymbol{(}\AgdaFunction{extb-comp}\AgdaSpace{}%
\AgdaBound{fb'}\AgdaSpace{}%
\AgdaBound{fb''}\AgdaSymbol{)}\AgdaSpace{}%
\AgdaBound{t}\AgdaSymbol{)}\AgdaSpace{}%
\AgdaFunction{⟩}\<%
\\
\>[2][@{}l@{\AgdaIndent{0}}]%
\>[4]\AgdaInductiveConstructor{ƛb}\AgdaSpace{}%
\AgdaSymbol{(}\AgdaFunction{renameb}\AgdaSpace{}%
\AgdaSymbol{(}\AgdaFunction{extb}\AgdaSpace{}%
\AgdaSymbol{(λ}\AgdaSpace{}%
\AgdaBound{x}\AgdaSpace{}%
\AgdaSymbol{→}\AgdaSpace{}%
\AgdaBound{fb''}\AgdaSpace{}%
\AgdaSymbol{(}\AgdaBound{fb'}\AgdaSpace{}%
\AgdaBound{x}\AgdaSymbol{)))}\AgdaSpace{}%
\AgdaBound{t}\AgdaSymbol{)}\<%
\\
%
\>[2]\AgdaOperator{\AgdaFunction{∎eq}}\<%
\\
\>[0]\AgdaFunction{renameb-comp}\AgdaSpace{}%
\AgdaBound{fb'}\AgdaSpace{}%
\AgdaBound{fb''}\AgdaSpace{}%
\AgdaSymbol{(}\AgdaBound{t}\AgdaSpace{}%
\AgdaOperator{\AgdaInductiveConstructor{·b}}\AgdaSpace{}%
\AgdaBound{t₁}\AgdaSymbol{)}\AgdaSpace{}%
\AgdaSymbol{=}\AgdaSpace{}%
\AgdaFunction{cong₂}\AgdaSpace{}%
\AgdaOperator{\AgdaInductiveConstructor{\AgdaUnderscore{}·b\AgdaUnderscore{}}}\AgdaSpace{}%
\AgdaSymbol{(}\AgdaFunction{renameb-comp}\AgdaSpace{}%
\AgdaBound{fb'}\AgdaSpace{}%
\AgdaBound{fb''}\AgdaSpace{}%
\AgdaBound{t}\AgdaSymbol{)}\AgdaSpace{}%
\AgdaSymbol{(}\AgdaFunction{renameb-comp}\AgdaSpace{}%
\AgdaBound{fb'}\AgdaSpace{}%
\AgdaBound{fb''}\AgdaSpace{}%
\AgdaBound{t₁}\AgdaSymbol{)}\<%
\\
%
\\[\AgdaEmptyExtraSkip]%
\>[0]\AgdaComment{--\ extb\ and\ Sb\ commute\ under\ renameb.}\<%
\\
\>[0]\AgdaFunction{rename-extb-Sb-commute}\AgdaSpace{}%
\AgdaSymbol{:}\AgdaSpace{}%
\AgdaSymbol{∀}\AgdaSpace{}%
\AgdaSymbol{\{}\AgdaBound{Γ}\AgdaSpace{}%
\AgdaBound{A}\AgdaSpace{}%
\AgdaBound{B}\AgdaSpace{}%
\AgdaBound{Δ}\AgdaSymbol{\}}\<%
\\
\>[0][@{}l@{\AgdaIndent{0}}]%
\>[2]\AgdaSymbol{→}\AgdaSpace{}%
\AgdaSymbol{(}\AgdaBound{fb'}\AgdaSpace{}%
\AgdaSymbol{:}\AgdaSpace{}%
\AgdaSymbol{∀}\AgdaSpace{}%
\AgdaSymbol{\{}\AgdaBound{Ab}\AgdaSymbol{\}}\AgdaSpace{}%
\AgdaSymbol{→}\AgdaSpace{}%
\AgdaBound{Ab}\AgdaSpace{}%
\AgdaOperator{\AgdaDatatype{∈b}}\AgdaSpace{}%
\AgdaBound{Γ}\AgdaSpace{}%
\AgdaSymbol{→}\AgdaSpace{}%
\AgdaBound{Ab}\AgdaSpace{}%
\AgdaOperator{\AgdaDatatype{∈b}}\AgdaSpace{}%
\AgdaBound{Δ}\AgdaSymbol{)}\<%
\\
%
\>[2]\AgdaSymbol{→}\AgdaSpace{}%
\AgdaSymbol{(}\AgdaBound{t}\AgdaSpace{}%
\AgdaSymbol{:}\AgdaSpace{}%
\AgdaBound{Γ}\AgdaSpace{}%
\AgdaOperator{\AgdaDatatype{⊢b}}\AgdaSpace{}%
\AgdaBound{A}\AgdaSymbol{)}\<%
\\
%
\>[2]\AgdaSymbol{→}%
\>[579I]\AgdaFunction{renameb}\AgdaSpace{}%
\AgdaSymbol{(}\AgdaFunction{extb}\AgdaSpace{}%
\AgdaBound{fb'}\AgdaSpace{}%
\AgdaSymbol{\{}\AgdaArgument{B}\AgdaSpace{}%
\AgdaSymbol{=}\AgdaSpace{}%
\AgdaBound{B}\AgdaSymbol{\})}\AgdaSpace{}%
\AgdaSymbol{(}\AgdaFunction{renameb}\AgdaSpace{}%
\AgdaInductiveConstructor{Sb}\AgdaSpace{}%
\AgdaBound{t}\AgdaSymbol{)}\<%
\\
\>[.][@{}l@{}]\<[579I]%
\>[4]\AgdaOperator{\AgdaDatatype{≡}}\AgdaSpace{}%
\AgdaFunction{renameb}\AgdaSpace{}%
\AgdaInductiveConstructor{Sb}\AgdaSpace{}%
\AgdaSymbol{(}\AgdaFunction{renameb}\AgdaSpace{}%
\AgdaBound{fb'}\AgdaSpace{}%
\AgdaBound{t}\AgdaSymbol{)}\<%
\\
\>[0]\AgdaFunction{rename-extb-Sb-commute}\AgdaSpace{}%
\AgdaBound{fb'}\AgdaSpace{}%
\AgdaBound{t}\AgdaSpace{}%
\AgdaSymbol{=}\<%
\\
\>[0][@{}l@{\AgdaIndent{0}}]%
\>[2]\AgdaOperator{\AgdaFunction{beginEq}}\<%
\\
\>[2][@{}l@{\AgdaIndent{0}}]%
\>[4]\AgdaFunction{renameb}\AgdaSpace{}%
\AgdaSymbol{(}\AgdaFunction{extb}\AgdaSpace{}%
\AgdaBound{fb'}\AgdaSymbol{)}\AgdaSpace{}%
\AgdaSymbol{(}\AgdaFunction{renameb}\AgdaSpace{}%
\AgdaInductiveConstructor{Sb}\AgdaSpace{}%
\AgdaBound{t}\AgdaSymbol{)}\<%
\\
%
\>[2]\AgdaFunction{≡⟨}\AgdaSpace{}%
\AgdaFunction{renameb-comp}\AgdaSpace{}%
\AgdaInductiveConstructor{Sb}\AgdaSpace{}%
\AgdaSymbol{(}\AgdaFunction{extb}\AgdaSpace{}%
\AgdaBound{fb'}\AgdaSymbol{)}\AgdaSpace{}%
\AgdaBound{t}\AgdaSpace{}%
\AgdaFunction{⟩}\<%
\\
\>[2][@{}l@{\AgdaIndent{0}}]%
\>[4]\AgdaFunction{renameb}\AgdaSpace{}%
\AgdaSymbol{(λ}\AgdaSpace{}%
\AgdaBound{x}\AgdaSpace{}%
\AgdaSymbol{→}\AgdaSpace{}%
\AgdaFunction{extb}\AgdaSpace{}%
\AgdaBound{fb'}\AgdaSpace{}%
\AgdaSymbol{(}\AgdaInductiveConstructor{Sb}\AgdaSpace{}%
\AgdaBound{x}\AgdaSymbol{))}\AgdaSpace{}%
\AgdaBound{t}\<%
\\
%
\>[2]\AgdaFunction{≡⟨}\AgdaSpace{}%
\AgdaFunction{renameb-cong}\AgdaSpace{}%
\AgdaSymbol{(λ}\AgdaSpace{}%
\AgdaBound{x}\AgdaSpace{}%
\AgdaSymbol{→}\AgdaSpace{}%
\AgdaInductiveConstructor{refl}\AgdaSymbol{)}\AgdaSpace{}%
\AgdaBound{t}\AgdaSpace{}%
\AgdaFunction{⟩}\<%
\\
\>[2][@{}l@{\AgdaIndent{0}}]%
\>[4]\AgdaFunction{renameb}\AgdaSpace{}%
\AgdaSymbol{(λ}\AgdaSpace{}%
\AgdaBound{x}\AgdaSpace{}%
\AgdaSymbol{→}\AgdaSpace{}%
\AgdaInductiveConstructor{Sb}\AgdaSpace{}%
\AgdaSymbol{(}\AgdaBound{fb'}\AgdaSpace{}%
\AgdaBound{x}\AgdaSymbol{))}\AgdaSpace{}%
\AgdaBound{t}\<%
\\
%
\>[2]\AgdaFunction{≡⟨}\AgdaSpace{}%
\AgdaFunction{sym}\AgdaSpace{}%
\AgdaSymbol{(}\AgdaFunction{renameb-comp}\AgdaSpace{}%
\AgdaBound{fb'}\AgdaSpace{}%
\AgdaInductiveConstructor{Sb}\AgdaSpace{}%
\AgdaBound{t}\AgdaSymbol{)}\AgdaSpace{}%
\AgdaFunction{⟩}\<%
\\
\>[2][@{}l@{\AgdaIndent{0}}]%
\>[4]\AgdaFunction{renameb}\AgdaSpace{}%
\AgdaInductiveConstructor{Sb}\AgdaSpace{}%
\AgdaSymbol{(}\AgdaFunction{renameb}\AgdaSpace{}%
\AgdaBound{fb'}\AgdaSpace{}%
\AgdaBound{t}\AgdaSymbol{)}\<%
\\
%
\>[2]\AgdaOperator{\AgdaFunction{∎eq}}\<%
\\
%
\\[\AgdaEmptyExtraSkip]%
\>[0]\AgdaComment{--\ Preservation\ and\ Reflection\ of\ }\<%
\\
\>[0]\AgdaComment{--\ rename\ of\ Gödel's\ translation}\<%
\\
\>[0]\AgdaFunction{pres-rename}\AgdaSpace{}%
\AgdaSymbol{:}\AgdaSpace{}%
\AgdaSymbol{∀}\AgdaSpace{}%
\AgdaSymbol{\{}\AgdaBound{Γ}\AgdaSpace{}%
\AgdaBound{A}\AgdaSpace{}%
\AgdaBound{Δ}\AgdaSymbol{\}}\<%
\\
\>[0][@{}l@{\AgdaIndent{0}}]%
\>[2]\AgdaSymbol{→}\AgdaSpace{}%
\AgdaSymbol{(}\AgdaBound{f'}\AgdaSpace{}%
\AgdaSymbol{:}\AgdaSpace{}%
\AgdaSymbol{∀}\AgdaSpace{}%
\AgdaSymbol{\{}\AgdaBound{A}\AgdaSymbol{\}}\AgdaSpace{}%
\AgdaSymbol{→}\AgdaSpace{}%
\AgdaBound{A}\AgdaSpace{}%
\AgdaOperator{\AgdaDatatype{∈}}\AgdaSpace{}%
\AgdaBound{Γ}\AgdaSpace{}%
\AgdaSymbol{→}\AgdaSpace{}%
\AgdaBound{A}\AgdaSpace{}%
\AgdaOperator{\AgdaDatatype{∈}}\AgdaSpace{}%
\AgdaBound{Δ}\AgdaSymbol{)}\<%
\\
%
\>[2]\AgdaSymbol{→}\AgdaSpace{}%
\AgdaSymbol{(}\AgdaBound{fb'}\AgdaSpace{}%
\AgdaSymbol{:}\AgdaSpace{}%
\AgdaSymbol{∀}\AgdaSpace{}%
\AgdaSymbol{\{}\AgdaBound{Ab}\AgdaSymbol{\}}\AgdaSpace{}%
\AgdaSymbol{→}\AgdaSpace{}%
\AgdaBound{Ab}\AgdaSpace{}%
\AgdaOperator{\AgdaDatatype{∈b}}\AgdaSpace{}%
\AgdaFunction{embedContext}\AgdaSpace{}%
\AgdaBound{Γ}\AgdaSpace{}%
\AgdaSymbol{→}\AgdaSpace{}%
\AgdaBound{Ab}\AgdaSpace{}%
\AgdaOperator{\AgdaDatatype{∈b}}\AgdaSpace{}%
\AgdaFunction{embedContext}\AgdaSpace{}%
\AgdaBound{Δ}\AgdaSymbol{)}\<%
\\
%
\>[2]\AgdaSymbol{→}\AgdaSpace{}%
\AgdaSymbol{(\{}\AgdaBound{A'}\AgdaSpace{}%
\AgdaSymbol{:}\AgdaSpace{}%
\AgdaDatatype{Type}\AgdaSymbol{\}}\AgdaSpace{}%
\AgdaSymbol{(}\AgdaBound{x}\AgdaSpace{}%
\AgdaSymbol{:}\AgdaSpace{}%
\AgdaBound{A'}\AgdaSpace{}%
\AgdaOperator{\AgdaDatatype{∈}}\AgdaSpace{}%
\AgdaBound{Γ}\AgdaSymbol{)}\AgdaSpace{}%
\AgdaSymbol{→}\AgdaSpace{}%
\AgdaBound{fb'}\AgdaSpace{}%
\AgdaSymbol{(}\AgdaFunction{embedLookup}\AgdaSpace{}%
\AgdaBound{x}\AgdaSymbol{)}\AgdaSpace{}%
\AgdaOperator{\AgdaDatatype{≡}}\AgdaSpace{}%
\AgdaFunction{embedLookup}\AgdaSpace{}%
\AgdaSymbol{(}\AgdaBound{f'}\AgdaSpace{}%
\AgdaBound{x}\AgdaSymbol{))}\<%
\\
%
\>[2]\AgdaSymbol{→}\AgdaSpace{}%
\AgdaSymbol{(}\AgdaBound{t}\AgdaSpace{}%
\AgdaSymbol{:}\AgdaSpace{}%
\AgdaBound{Γ}\AgdaSpace{}%
\AgdaOperator{\AgdaDatatype{⊢}}\AgdaSpace{}%
\AgdaBound{A}\AgdaSymbol{)}\<%
\\
%
\>[2]\AgdaSymbol{→}\AgdaSpace{}%
\AgdaFunction{renameb}\AgdaSpace{}%
\AgdaBound{fb'}\AgdaSpace{}%
\AgdaSymbol{(}\AgdaFunction{embedTerm}\AgdaSpace{}%
\AgdaBound{t}\AgdaSymbol{)}\AgdaSpace{}%
\AgdaOperator{\AgdaDatatype{≡}}\AgdaSpace{}%
\AgdaFunction{embedTerm}\AgdaSpace{}%
\AgdaSymbol{(}\AgdaFunction{rename}\AgdaSpace{}%
\AgdaBound{f'}\AgdaSpace{}%
\AgdaBound{t}\AgdaSymbol{)}\<%
\\
\>[0]\AgdaFunction{pres-rename}\AgdaSpace{}%
\AgdaBound{f'}\AgdaSpace{}%
\AgdaBound{fb'}\AgdaSpace{}%
\AgdaBound{h}\AgdaSpace{}%
\AgdaSymbol{(}\AgdaInductiveConstructor{`}\AgdaSpace{}%
\AgdaBound{x}\AgdaSymbol{)}\AgdaSpace{}%
\AgdaSymbol{=}\AgdaSpace{}%
\AgdaFunction{cong}\AgdaSpace{}%
\AgdaInductiveConstructor{box}\AgdaSpace{}%
\AgdaSymbol{(}\AgdaFunction{cong}\AgdaSpace{}%
\AgdaInductiveConstructor{ε}\AgdaSpace{}%
\AgdaSymbol{(}\AgdaBound{h}\AgdaSpace{}%
\AgdaBound{x}\AgdaSymbol{))}\<%
\\
\>[0]\AgdaFunction{pres-rename}\AgdaSpace{}%
\AgdaBound{f'}\AgdaSpace{}%
\AgdaBound{fb'}\AgdaSpace{}%
\AgdaBound{h}\AgdaSpace{}%
\AgdaSymbol{(}\AgdaInductiveConstructor{ƛ}\AgdaSpace{}%
\AgdaBound{t}\AgdaSymbol{)}\AgdaSpace{}%
\AgdaSymbol{=}\AgdaSpace{}%
\AgdaFunction{cong}\AgdaSpace{}%
\AgdaInductiveConstructor{box}\AgdaSpace{}%
\AgdaSymbol{(}\AgdaFunction{cong}\AgdaSpace{}%
\AgdaInductiveConstructor{ƛb}\AgdaSpace{}%
\AgdaSymbol{(}\AgdaFunction{pres-rename}\AgdaSpace{}%
\AgdaSymbol{(}\AgdaFunction{calc.ext}\AgdaSpace{}%
\AgdaBound{f'}\AgdaSymbol{)}\AgdaSpace{}%
\AgdaSymbol{(}\AgdaFunction{boxCalc.extb}\AgdaSpace{}%
\AgdaBound{fb'}\AgdaSymbol{)}\AgdaSpace{}%
\AgdaSymbol{(}\AgdaFunction{pres-ext}\AgdaSpace{}%
\AgdaBound{f'}\AgdaSpace{}%
\AgdaBound{fb'}\AgdaSpace{}%
\AgdaBound{h}\AgdaSymbol{)}\AgdaSpace{}%
\AgdaBound{t}\AgdaSymbol{))}\<%
\\
\>[0]\AgdaFunction{pres-rename}\AgdaSpace{}%
\AgdaBound{f'}\AgdaSpace{}%
\AgdaBound{fb'}\AgdaSpace{}%
\AgdaBound{h}\AgdaSpace{}%
\AgdaSymbol{(}\AgdaBound{t}\AgdaSpace{}%
\AgdaOperator{\AgdaInductiveConstructor{·}}\AgdaSpace{}%
\AgdaBound{t₁}\AgdaSymbol{)}\AgdaSpace{}%
\AgdaSymbol{=}\<%
\\
\>[0][@{}l@{\AgdaIndent{0}}]%
\>[2]\AgdaOperator{\AgdaFunction{beginEq}}\<%
\\
\>[2][@{}l@{\AgdaIndent{0}}]%
\>[4]\AgdaInductiveConstructor{ƛb}\AgdaSpace{}%
\AgdaSymbol{(}\AgdaInductiveConstructor{ε}\AgdaSpace{}%
\AgdaInductiveConstructor{Zb}\AgdaSpace{}%
\AgdaOperator{\AgdaInductiveConstructor{·b}}\AgdaSpace{}%
\AgdaFunction{renameb}\AgdaSpace{}%
\AgdaSymbol{(}\AgdaFunction{extb}\AgdaSpace{}%
\AgdaBound{fb'}\AgdaSymbol{)}\AgdaSpace{}%
\AgdaSymbol{(}\AgdaFunction{renameb}\AgdaSpace{}%
\AgdaInductiveConstructor{Sb}\AgdaSpace{}%
\AgdaSymbol{(}\AgdaFunction{embedTerm}\AgdaSpace{}%
\AgdaBound{t₁}\AgdaSymbol{)))}\AgdaSpace{}%
\AgdaOperator{\AgdaInductiveConstructor{·b}}\AgdaSpace{}%
\AgdaFunction{renameb}\AgdaSpace{}%
\AgdaBound{fb'}\AgdaSpace{}%
\AgdaSymbol{(}\AgdaFunction{embedTerm}\AgdaSpace{}%
\AgdaBound{t}\AgdaSymbol{)}\<%
\\
%
\>[2]\AgdaFunction{≡⟨}\AgdaSpace{}%
\AgdaFunction{cong₂}\AgdaSpace{}%
\AgdaOperator{\AgdaInductiveConstructor{\AgdaUnderscore{}·b\AgdaUnderscore{}}}\AgdaSpace{}%
\AgdaInductiveConstructor{refl}\AgdaSpace{}%
\AgdaSymbol{(}\AgdaFunction{pres-rename}\AgdaSpace{}%
\AgdaBound{f'}\AgdaSpace{}%
\AgdaBound{fb'}\AgdaSpace{}%
\AgdaBound{h}\AgdaSpace{}%
\AgdaBound{t}\AgdaSymbol{)}\AgdaSpace{}%
\AgdaFunction{⟩}\<%
\\
\>[2][@{}l@{\AgdaIndent{0}}]%
\>[4]\AgdaInductiveConstructor{ƛb}\AgdaSpace{}%
\AgdaSymbol{(}\AgdaInductiveConstructor{ε}\AgdaSpace{}%
\AgdaInductiveConstructor{Zb}\AgdaSpace{}%
\AgdaOperator{\AgdaInductiveConstructor{·b}}\AgdaSpace{}%
\AgdaFunction{renameb}\AgdaSpace{}%
\AgdaSymbol{(}\AgdaFunction{extb}\AgdaSpace{}%
\AgdaBound{fb'}\AgdaSymbol{)}\AgdaSpace{}%
\AgdaSymbol{(}\AgdaFunction{renameb}\AgdaSpace{}%
\AgdaInductiveConstructor{Sb}\AgdaSpace{}%
\AgdaSymbol{(}\AgdaFunction{embedTerm}\AgdaSpace{}%
\AgdaBound{t₁}\AgdaSymbol{)))}\AgdaSpace{}%
\AgdaOperator{\AgdaInductiveConstructor{·b}}\AgdaSpace{}%
\AgdaFunction{embedTerm}\AgdaSpace{}%
\AgdaSymbol{(}\AgdaFunction{rename}\AgdaSpace{}%
\AgdaBound{f'}\AgdaSpace{}%
\AgdaBound{t}\AgdaSymbol{)}\<%
\\
%
\>[2]\AgdaFunction{≡⟨}\AgdaSpace{}%
\AgdaFunction{cong₂}\AgdaSpace{}%
\AgdaOperator{\AgdaInductiveConstructor{\AgdaUnderscore{}·b\AgdaUnderscore{}}}\AgdaSpace{}%
\AgdaSymbol{(}\AgdaFunction{cong}\AgdaSpace{}%
\AgdaInductiveConstructor{ƛb}\AgdaSpace{}%
\AgdaSymbol{(}\AgdaFunction{cong₂}\AgdaSpace{}%
\AgdaOperator{\AgdaInductiveConstructor{\AgdaUnderscore{}·b\AgdaUnderscore{}}}\AgdaSpace{}%
\AgdaInductiveConstructor{refl}\AgdaSpace{}%
\AgdaSymbol{(}\AgdaFunction{rename-extb-Sb-commute}\AgdaSpace{}%
\AgdaBound{fb'}\AgdaSpace{}%
\AgdaSymbol{(}\AgdaFunction{embedTerm}\AgdaSpace{}%
\AgdaBound{t₁}\AgdaSymbol{))))}\AgdaSpace{}%
\AgdaInductiveConstructor{refl}\AgdaSpace{}%
\AgdaFunction{⟩}\<%
\\
\>[2][@{}l@{\AgdaIndent{0}}]%
\>[4]\AgdaInductiveConstructor{ƛb}\AgdaSpace{}%
\AgdaSymbol{(}\AgdaInductiveConstructor{ε}\AgdaSpace{}%
\AgdaInductiveConstructor{Zb}\AgdaSpace{}%
\AgdaOperator{\AgdaInductiveConstructor{·b}}\AgdaSpace{}%
\AgdaFunction{renameb}\AgdaSpace{}%
\AgdaInductiveConstructor{Sb}\AgdaSpace{}%
\AgdaSymbol{(}\AgdaFunction{renameb}\AgdaSpace{}%
\AgdaBound{fb'}\AgdaSpace{}%
\AgdaSymbol{(}\AgdaFunction{embedTerm}\AgdaSpace{}%
\AgdaBound{t₁}\AgdaSymbol{)))}\AgdaSpace{}%
\AgdaOperator{\AgdaInductiveConstructor{·b}}\AgdaSpace{}%
\AgdaFunction{embedTerm}\AgdaSpace{}%
\AgdaSymbol{(}\AgdaFunction{rename}\AgdaSpace{}%
\AgdaBound{f'}\AgdaSpace{}%
\AgdaBound{t}\AgdaSymbol{)}\<%
\\
%
\>[2]\AgdaFunction{≡⟨}\AgdaSpace{}%
\AgdaFunction{cong₂}\AgdaSpace{}%
\AgdaOperator{\AgdaInductiveConstructor{\AgdaUnderscore{}·b\AgdaUnderscore{}}}\AgdaSpace{}%
\AgdaSymbol{(}\AgdaFunction{cong}\AgdaSpace{}%
\AgdaInductiveConstructor{ƛb}\AgdaSpace{}%
\AgdaSymbol{(}\AgdaFunction{cong₂}\AgdaSpace{}%
\AgdaOperator{\AgdaInductiveConstructor{\AgdaUnderscore{}·b\AgdaUnderscore{}}}\AgdaSpace{}%
\AgdaInductiveConstructor{refl}\AgdaSpace{}%
\AgdaSymbol{(}\AgdaFunction{cong}\AgdaSpace{}%
\AgdaSymbol{(}\AgdaFunction{renameb}\AgdaSpace{}%
\AgdaInductiveConstructor{Sb}\AgdaSymbol{)}\AgdaSpace{}%
\AgdaSymbol{(}\AgdaFunction{pres-rename}\AgdaSpace{}%
\AgdaBound{f'}\AgdaSpace{}%
\AgdaBound{fb'}\AgdaSpace{}%
\AgdaBound{h}\AgdaSpace{}%
\AgdaBound{t₁}\AgdaSymbol{))))}\AgdaSpace{}%
\AgdaInductiveConstructor{refl}\AgdaSpace{}%
\AgdaFunction{⟩}\<%
\\
\>[2][@{}l@{\AgdaIndent{0}}]%
\>[4]\AgdaInductiveConstructor{ƛb}\AgdaSpace{}%
\AgdaSymbol{(}\AgdaInductiveConstructor{ε}\AgdaSpace{}%
\AgdaInductiveConstructor{Zb}\AgdaSpace{}%
\AgdaOperator{\AgdaInductiveConstructor{·b}}\AgdaSpace{}%
\AgdaFunction{renameb}\AgdaSpace{}%
\AgdaInductiveConstructor{Sb}\AgdaSpace{}%
\AgdaSymbol{(}\AgdaFunction{embedTerm}\AgdaSpace{}%
\AgdaSymbol{(}\AgdaFunction{rename}\AgdaSpace{}%
\AgdaBound{f'}\AgdaSpace{}%
\AgdaBound{t₁}\AgdaSymbol{)))}\AgdaSpace{}%
\AgdaOperator{\AgdaInductiveConstructor{·b}}\AgdaSpace{}%
\AgdaFunction{embedTerm}\AgdaSpace{}%
\AgdaSymbol{(}\AgdaFunction{rename}\AgdaSpace{}%
\AgdaBound{f'}\AgdaSpace{}%
\AgdaBound{t}\AgdaSymbol{)}\<%
\\
%
\>[2]\AgdaOperator{\AgdaFunction{∎eq}}\<%
\\
%
\\[\AgdaEmptyExtraSkip]%
\>[0]\AgdaComment{--\ Preservation\ and\ Reflection\ of\ }\<%
\\
\>[0]\AgdaComment{--\ exts\ of\ Gödel's\ translation}\<%
\\
\>[0]\AgdaFunction{pres-exts}\AgdaSpace{}%
\AgdaSymbol{:}\AgdaSpace{}%
\AgdaSymbol{∀}\AgdaSpace{}%
\AgdaSymbol{\{}\AgdaBound{Γ}\AgdaSpace{}%
\AgdaBound{Δ}\AgdaSymbol{\}}\<%
\\
\>[0][@{}l@{\AgdaIndent{0}}]%
\>[2]\AgdaSymbol{→}\AgdaSpace{}%
\AgdaSymbol{(}\AgdaBound{fb'}\AgdaSpace{}%
\AgdaSymbol{:}\AgdaSpace{}%
\AgdaSymbol{∀}\AgdaSpace{}%
\AgdaSymbol{\{}\AgdaBound{Ab}\AgdaSymbol{\}}\AgdaSpace{}%
\AgdaSymbol{→}\AgdaSpace{}%
\AgdaInductiveConstructor{□}\AgdaSpace{}%
\AgdaBound{Ab}\AgdaSpace{}%
\AgdaOperator{\AgdaDatatype{∈b}}\AgdaSpace{}%
\AgdaFunction{embedContext}\AgdaSpace{}%
\AgdaBound{Γ}\AgdaSpace{}%
\AgdaSymbol{→}\AgdaSpace{}%
\AgdaFunction{embedContext}\AgdaSpace{}%
\AgdaBound{Δ}\AgdaSpace{}%
\AgdaOperator{\AgdaDatatype{⊢b}}\AgdaSpace{}%
\AgdaBound{Ab}\AgdaSymbol{)}\<%
\\
%
\>[2]\AgdaSymbol{→}\AgdaSpace{}%
\AgdaSymbol{(}\AgdaBound{f'}\AgdaSpace{}%
\AgdaSymbol{:}\AgdaSpace{}%
\AgdaSymbol{∀}\AgdaSpace{}%
\AgdaSymbol{\{}\AgdaBound{A}\AgdaSymbol{\}}\AgdaSpace{}%
\AgdaSymbol{→}\AgdaSpace{}%
\AgdaBound{A}\AgdaSpace{}%
\AgdaOperator{\AgdaDatatype{∈}}\AgdaSpace{}%
\AgdaBound{Γ}\AgdaSpace{}%
\AgdaSymbol{→}\AgdaSpace{}%
\AgdaBound{Δ}\AgdaSpace{}%
\AgdaOperator{\AgdaDatatype{⊢}}\AgdaSpace{}%
\AgdaBound{A}\AgdaSymbol{)}\<%
\\
%
\>[2]\AgdaSymbol{→}\AgdaSpace{}%
\AgdaSymbol{(\{}\AgdaBound{A'}\AgdaSpace{}%
\AgdaSymbol{:}\AgdaSpace{}%
\AgdaDatatype{Type}\AgdaSymbol{\}}\AgdaSpace{}%
\AgdaSymbol{(}\AgdaBound{x}\AgdaSpace{}%
\AgdaSymbol{:}\AgdaSpace{}%
\AgdaBound{A'}\AgdaSpace{}%
\AgdaOperator{\AgdaDatatype{∈}}\AgdaSpace{}%
\AgdaBound{Γ}\AgdaSymbol{)}\AgdaSpace{}%
\AgdaSymbol{→}\AgdaSpace{}%
\AgdaInductiveConstructor{box}\AgdaSpace{}%
\AgdaSymbol{(}\AgdaBound{fb'}\AgdaSpace{}%
\AgdaSymbol{(}\AgdaFunction{embedLookup}\AgdaSpace{}%
\AgdaBound{x}\AgdaSymbol{))}\AgdaSpace{}%
\AgdaOperator{\AgdaDatatype{≡}}\AgdaSpace{}%
\AgdaFunction{embedTerm}\AgdaSpace{}%
\AgdaSymbol{(}\AgdaBound{f'}\AgdaSpace{}%
\AgdaBound{x}\AgdaSymbol{))}\<%
\\
%
\>[2]\AgdaSymbol{→}\AgdaSpace{}%
\AgdaSymbol{\{}\AgdaBound{A'}\AgdaSpace{}%
\AgdaBound{A}\AgdaSpace{}%
\AgdaSymbol{:}\AgdaSpace{}%
\AgdaDatatype{Type}\AgdaSymbol{\}}\AgdaSpace{}%
\AgdaSymbol{(}\AgdaBound{x}\AgdaSpace{}%
\AgdaSymbol{:}\AgdaSpace{}%
\AgdaBound{A'}\AgdaSpace{}%
\AgdaOperator{\AgdaDatatype{∈}}\AgdaSpace{}%
\AgdaBound{Γ}\AgdaSpace{}%
\AgdaOperator{\AgdaInductiveConstructor{,}}\AgdaSpace{}%
\AgdaBound{A}\AgdaSymbol{)}\AgdaSpace{}%
\AgdaSymbol{→}\AgdaSpace{}%
\AgdaInductiveConstructor{box}\AgdaSpace{}%
\AgdaSymbol{(}\AgdaFunction{extsb}\AgdaSpace{}%
\AgdaBound{fb'}\AgdaSpace{}%
\AgdaSymbol{(}\AgdaFunction{embedLookup}\AgdaSpace{}%
\AgdaBound{x}\AgdaSymbol{))}\AgdaSpace{}%
\AgdaOperator{\AgdaDatatype{≡}}\AgdaSpace{}%
\AgdaFunction{embedTerm}\AgdaSpace{}%
\AgdaSymbol{(}\AgdaFunction{exts}\AgdaSpace{}%
\AgdaBound{f'}\AgdaSpace{}%
\AgdaBound{x}\AgdaSymbol{)}\<%
\\
\>[0]\AgdaFunction{pres-exts}\AgdaSpace{}%
\AgdaBound{fb'}\AgdaSpace{}%
\AgdaBound{f'}\AgdaSpace{}%
\AgdaBound{h}\AgdaSpace{}%
\AgdaInductiveConstructor{Z}\AgdaSpace{}%
\AgdaSymbol{=}\AgdaSpace{}%
\AgdaInductiveConstructor{refl}\<%
\\
\>[0]\AgdaFunction{pres-exts}\AgdaSpace{}%
\AgdaBound{fb'}\AgdaSpace{}%
\AgdaBound{f'}\AgdaSpace{}%
\AgdaBound{h}\AgdaSpace{}%
\AgdaSymbol{(}\AgdaInductiveConstructor{S}\AgdaSpace{}%
\AgdaBound{x}\AgdaSymbol{)}\AgdaSpace{}%
\AgdaSymbol{=}\<%
\\
\>[0][@{}l@{\AgdaIndent{0}}]%
\>[2]\AgdaOperator{\AgdaFunction{beginEq}}\<%
\\
\>[2][@{}l@{\AgdaIndent{0}}]%
\>[4]\AgdaInductiveConstructor{box}\AgdaSpace{}%
\AgdaSymbol{(}\AgdaFunction{renameb}\AgdaSpace{}%
\AgdaInductiveConstructor{Sb}\AgdaSpace{}%
\AgdaSymbol{(}\AgdaBound{fb'}\AgdaSpace{}%
\AgdaSymbol{(}\AgdaFunction{embedLookup}\AgdaSpace{}%
\AgdaBound{x}\AgdaSymbol{)))}\<%
\\
%
\>[2]\AgdaFunction{≡⟨}\AgdaSpace{}%
\AgdaInductiveConstructor{refl}\AgdaSpace{}%
\AgdaFunction{⟩}\<%
\\
\>[2][@{}l@{\AgdaIndent{0}}]%
\>[4]\AgdaFunction{renameb}\AgdaSpace{}%
\AgdaInductiveConstructor{Sb}\AgdaSpace{}%
\AgdaSymbol{(}\AgdaInductiveConstructor{box}\AgdaSpace{}%
\AgdaSymbol{(}\AgdaBound{fb'}\AgdaSpace{}%
\AgdaSymbol{(}\AgdaFunction{embedLookup}\AgdaSpace{}%
\AgdaBound{x}\AgdaSymbol{)))}\<%
\\
%
\>[2]\AgdaFunction{≡⟨}\AgdaSpace{}%
\AgdaFunction{cong}\AgdaSpace{}%
\AgdaSymbol{(}\AgdaFunction{renameb}\AgdaSpace{}%
\AgdaInductiveConstructor{Sb}\AgdaSymbol{)}\AgdaSpace{}%
\AgdaSymbol{(}\AgdaBound{h}\AgdaSpace{}%
\AgdaBound{x}\AgdaSymbol{)}\AgdaSpace{}%
\AgdaFunction{⟩}\<%
\\
\>[2][@{}l@{\AgdaIndent{0}}]%
\>[4]\AgdaFunction{renameb}\AgdaSpace{}%
\AgdaInductiveConstructor{Sb}\AgdaSpace{}%
\AgdaSymbol{(}\AgdaFunction{embedTerm}\AgdaSpace{}%
\AgdaSymbol{(}\AgdaBound{f'}\AgdaSpace{}%
\AgdaBound{x}\AgdaSymbol{))}\<%
\\
%
\>[2]\AgdaFunction{≡⟨}\AgdaSpace{}%
\AgdaFunction{pres-rename}\AgdaSpace{}%
\AgdaInductiveConstructor{S}\AgdaSpace{}%
\AgdaInductiveConstructor{Sb}\AgdaSpace{}%
\AgdaSymbol{(λ}\AgdaSpace{}%
\AgdaBound{x₁}\AgdaSpace{}%
\AgdaSymbol{→}\AgdaSpace{}%
\AgdaInductiveConstructor{refl}\AgdaSymbol{)}\AgdaSpace{}%
\AgdaSymbol{(}\AgdaBound{f'}\AgdaSpace{}%
\AgdaBound{x}\AgdaSymbol{)}\AgdaSpace{}%
\AgdaFunction{⟩}\<%
\\
\>[2][@{}l@{\AgdaIndent{0}}]%
\>[4]\AgdaFunction{embedTerm}\AgdaSpace{}%
\AgdaSymbol{(}\AgdaFunction{rename}\AgdaSpace{}%
\AgdaInductiveConstructor{S}\AgdaSpace{}%
\AgdaSymbol{(}\AgdaBound{f'}\AgdaSpace{}%
\AgdaBound{x}\AgdaSymbol{))}\<%
\\
%
\>[2]\AgdaOperator{\AgdaFunction{∎eq}}\<%
\\
%
\\[\AgdaEmptyExtraSkip]%
\>[0]\AgdaComment{--\ Extsb\ preserves\ pointwise\ equality}\<%
\\
\>[0]\AgdaComment{--\ of\ renaming\ functions}\<%
\\
\>[0]\AgdaFunction{extsb-cong}\AgdaSpace{}%
\AgdaSymbol{:}\AgdaSpace{}%
\AgdaSymbol{∀}\AgdaSpace{}%
\AgdaSymbol{\{}\AgdaBound{Γ₁}\AgdaSpace{}%
\AgdaBound{Γ₂}\AgdaSpace{}%
\AgdaBound{A}\AgdaSpace{}%
\AgdaBound{A'}\AgdaSymbol{\}}\<%
\\
\>[0][@{}l@{\AgdaIndent{0}}]%
\>[2]\AgdaSymbol{→}\AgdaSpace{}%
\AgdaSymbol{\{}\AgdaBound{fb'}\AgdaSpace{}%
\AgdaSymbol{:}\AgdaSpace{}%
\AgdaSymbol{∀}\AgdaSpace{}%
\AgdaSymbol{\{}\AgdaBound{Ab}\AgdaSymbol{\}}\AgdaSpace{}%
\AgdaSymbol{→}\AgdaSpace{}%
\AgdaInductiveConstructor{□}\AgdaSpace{}%
\AgdaBound{Ab}\AgdaSpace{}%
\AgdaOperator{\AgdaDatatype{∈b}}\AgdaSpace{}%
\AgdaBound{Γ₁}\AgdaSpace{}%
\AgdaSymbol{→}\AgdaSpace{}%
\AgdaBound{Γ₂}\AgdaSpace{}%
\AgdaOperator{\AgdaDatatype{⊢b}}\AgdaSpace{}%
\AgdaBound{Ab}\AgdaSymbol{\}}\<%
\\
%
\>[2]\AgdaSymbol{→}\AgdaSpace{}%
\AgdaSymbol{\{}\AgdaBound{fb''}\AgdaSpace{}%
\AgdaSymbol{:}\AgdaSpace{}%
\AgdaSymbol{∀}\AgdaSpace{}%
\AgdaSymbol{\{}\AgdaBound{Ab}\AgdaSymbol{\}}\AgdaSpace{}%
\AgdaSymbol{→}\AgdaSpace{}%
\AgdaInductiveConstructor{□}\AgdaSpace{}%
\AgdaBound{Ab}\AgdaSpace{}%
\AgdaOperator{\AgdaDatatype{∈b}}\AgdaSpace{}%
\AgdaBound{Γ₁}\AgdaSpace{}%
\AgdaSymbol{→}\AgdaSpace{}%
\AgdaBound{Γ₂}\AgdaSpace{}%
\AgdaOperator{\AgdaDatatype{⊢b}}\AgdaSpace{}%
\AgdaBound{Ab}\AgdaSymbol{\}}\<%
\\
%
\>[2]\AgdaSymbol{→}\AgdaSpace{}%
\AgdaSymbol{(}\AgdaBound{h}\AgdaSpace{}%
\AgdaSymbol{:}\AgdaSpace{}%
\AgdaSymbol{∀}\AgdaSpace{}%
\AgdaSymbol{\{}\AgdaBound{Ab}\AgdaSymbol{\}}\AgdaSpace{}%
\AgdaSymbol{→}\AgdaSpace{}%
\AgdaSymbol{(}\AgdaBound{x}\AgdaSpace{}%
\AgdaSymbol{:}\AgdaSpace{}%
\AgdaInductiveConstructor{□}\AgdaSpace{}%
\AgdaBound{Ab}\AgdaSpace{}%
\AgdaOperator{\AgdaDatatype{∈b}}\AgdaSpace{}%
\AgdaBound{Γ₁}\AgdaSymbol{)}\AgdaSpace{}%
\AgdaSymbol{→}\AgdaSpace{}%
\AgdaBound{fb'}\AgdaSpace{}%
\AgdaBound{x}\AgdaSpace{}%
\AgdaOperator{\AgdaDatatype{≡}}\AgdaSpace{}%
\AgdaBound{fb''}\AgdaSpace{}%
\AgdaBound{x}\AgdaSymbol{)}\<%
\\
%
\>[2]\AgdaSymbol{→}\AgdaSpace{}%
\AgdaSymbol{(}\AgdaBound{x}\AgdaSpace{}%
\AgdaSymbol{:}\AgdaSpace{}%
\AgdaInductiveConstructor{□}\AgdaSpace{}%
\AgdaBound{A}\AgdaSpace{}%
\AgdaOperator{\AgdaDatatype{∈b}}\AgdaSpace{}%
\AgdaBound{Γ₁}\AgdaSpace{}%
\AgdaOperator{\AgdaInductiveConstructor{,b}}\AgdaSpace{}%
\AgdaInductiveConstructor{□}\AgdaSpace{}%
\AgdaBound{A'}\AgdaSymbol{)}\<%
\\
%
\>[2]\AgdaSymbol{→}\AgdaSpace{}%
\AgdaFunction{extsb}\AgdaSpace{}%
\AgdaBound{fb'}\AgdaSpace{}%
\AgdaBound{x}\AgdaSpace{}%
\AgdaOperator{\AgdaDatatype{≡}}\AgdaSpace{}%
\AgdaFunction{extsb}\AgdaSpace{}%
\AgdaBound{fb''}\AgdaSpace{}%
\AgdaBound{x}\<%
\\
\>[0]\AgdaFunction{extsb-cong}\AgdaSpace{}%
\AgdaBound{h}\AgdaSpace{}%
\AgdaInductiveConstructor{Zb}\AgdaSpace{}%
\AgdaSymbol{=}\AgdaSpace{}%
\AgdaInductiveConstructor{refl}\<%
\\
\>[0]\AgdaFunction{extsb-cong}\AgdaSpace{}%
\AgdaBound{h}\AgdaSpace{}%
\AgdaSymbol{(}\AgdaInductiveConstructor{Sb}\AgdaSpace{}%
\AgdaBound{x}\AgdaSymbol{)}\AgdaSpace{}%
\AgdaSymbol{=}\AgdaSpace{}%
\AgdaFunction{cong}\AgdaSpace{}%
\AgdaSymbol{(}\AgdaFunction{renameb}\AgdaSpace{}%
\AgdaInductiveConstructor{Sb}\AgdaSymbol{)}\AgdaSpace{}%
\AgdaSymbol{(}\AgdaBound{h}\AgdaSpace{}%
\AgdaBound{x}\AgdaSymbol{)}\<%
\\
\>[0]\<%
\end{code}
\begin{code}%
\>[0]\AgdaComment{--\ substb\ preserves\ pointwise\ equality}\<%
\\
\>[0]\AgdaComment{--\ of\ renaming\ functions}\<%
\\
\>[0]\AgdaFunction{substb-cong}\AgdaSpace{}%
\AgdaSymbol{:}\AgdaSpace{}%
\AgdaSymbol{∀}\AgdaSpace{}%
\AgdaSymbol{\{}\AgdaBound{Γ₁}\AgdaSpace{}%
\AgdaBound{Γ₂}\AgdaSpace{}%
\AgdaBound{A}\AgdaSymbol{\}}\<%
\\
\>[0][@{}l@{\AgdaIndent{0}}]%
\>[2]\AgdaSymbol{→}\AgdaSpace{}%
\AgdaSymbol{\{}\AgdaBound{fb'}\AgdaSpace{}%
\AgdaSymbol{:}\AgdaSpace{}%
\AgdaSymbol{∀}\AgdaSpace{}%
\AgdaSymbol{\{}\AgdaBound{Ab}\AgdaSymbol{\}}\AgdaSpace{}%
\AgdaSymbol{→}\AgdaSpace{}%
\AgdaInductiveConstructor{□}\AgdaSpace{}%
\AgdaBound{Ab}\AgdaSpace{}%
\AgdaOperator{\AgdaDatatype{∈b}}\AgdaSpace{}%
\AgdaBound{Γ₁}\AgdaSpace{}%
\AgdaSymbol{→}\AgdaSpace{}%
\AgdaBound{Γ₂}\AgdaSpace{}%
\AgdaOperator{\AgdaDatatype{⊢b}}\AgdaSpace{}%
\AgdaBound{Ab}\AgdaSymbol{\}}\<%
\\
%
\>[2]\AgdaSymbol{→}\AgdaSpace{}%
\AgdaSymbol{\{}\AgdaBound{fb''}\AgdaSpace{}%
\AgdaSymbol{:}\AgdaSpace{}%
\AgdaSymbol{∀}\AgdaSpace{}%
\AgdaSymbol{\{}\AgdaBound{Ab}\AgdaSymbol{\}}\AgdaSpace{}%
\AgdaSymbol{→}\AgdaSpace{}%
\AgdaInductiveConstructor{□}\AgdaSpace{}%
\AgdaBound{Ab}\AgdaSpace{}%
\AgdaOperator{\AgdaDatatype{∈b}}\AgdaSpace{}%
\AgdaBound{Γ₁}\AgdaSpace{}%
\AgdaSymbol{→}\AgdaSpace{}%
\AgdaBound{Γ₂}\AgdaSpace{}%
\AgdaOperator{\AgdaDatatype{⊢b}}\AgdaSpace{}%
\AgdaBound{Ab}\AgdaSymbol{\}}\<%
\\
%
\>[2]\AgdaSymbol{→}\AgdaSpace{}%
\AgdaSymbol{(}\AgdaBound{h}\AgdaSpace{}%
\AgdaSymbol{:}\AgdaSpace{}%
\AgdaSymbol{∀}\AgdaSpace{}%
\AgdaSymbol{\{}\AgdaBound{Ab}\AgdaSymbol{\}}\AgdaSpace{}%
\AgdaSymbol{→}\AgdaSpace{}%
\AgdaSymbol{(}\AgdaBound{x}\AgdaSpace{}%
\AgdaSymbol{:}\AgdaSpace{}%
\AgdaInductiveConstructor{□}\AgdaSpace{}%
\AgdaBound{Ab}\AgdaSpace{}%
\AgdaOperator{\AgdaDatatype{∈b}}\AgdaSpace{}%
\AgdaBound{Γ₁}\AgdaSymbol{)}\AgdaSpace{}%
\AgdaSymbol{→}\AgdaSpace{}%
\AgdaBound{fb'}\AgdaSpace{}%
\AgdaBound{x}\AgdaSpace{}%
\AgdaOperator{\AgdaDatatype{≡}}\AgdaSpace{}%
\AgdaBound{fb''}\AgdaSpace{}%
\AgdaBound{x}\AgdaSymbol{)}\<%
\\
%
\>[2]\AgdaSymbol{→}\AgdaSpace{}%
\AgdaSymbol{(}\AgdaBound{t}\AgdaSpace{}%
\AgdaSymbol{:}\AgdaSpace{}%
\AgdaBound{Γ₁}\AgdaSpace{}%
\AgdaOperator{\AgdaDatatype{⊢b}}\AgdaSpace{}%
\AgdaBound{A}\AgdaSymbol{)}\<%
\\
%
\>[2]\AgdaSymbol{→}\AgdaSpace{}%
\AgdaFunction{substb}\AgdaSpace{}%
\AgdaBound{fb'}\AgdaSpace{}%
\AgdaBound{t}\AgdaSpace{}%
\AgdaOperator{\AgdaDatatype{≡}}\AgdaSpace{}%
\AgdaFunction{substb}\AgdaSpace{}%
\AgdaBound{fb''}\AgdaSpace{}%
\AgdaBound{t}\<%
\\
\>[0]\AgdaFunction{substb-cong}\AgdaSpace{}%
\AgdaBound{h}\AgdaSpace{}%
\AgdaSymbol{(}\AgdaInductiveConstructor{ε}\AgdaSpace{}%
\AgdaBound{x}\AgdaSymbol{)}\AgdaSpace{}%
\AgdaSymbol{=}\AgdaSpace{}%
\AgdaBound{h}\AgdaSpace{}%
\AgdaBound{x}\<%
\\
\>[0]\AgdaFunction{substb-cong}\AgdaSpace{}%
\AgdaBound{h}\AgdaSpace{}%
\AgdaSymbol{(}\AgdaInductiveConstructor{box}\AgdaSpace{}%
\AgdaBound{t}\AgdaSymbol{)}\AgdaSpace{}%
\AgdaSymbol{=}\AgdaSpace{}%
\AgdaFunction{cong}\AgdaSpace{}%
\AgdaInductiveConstructor{box}\AgdaSpace{}%
\AgdaSymbol{(}\AgdaFunction{substb-cong}\AgdaSpace{}%
\AgdaBound{h}\AgdaSpace{}%
\AgdaBound{t}\AgdaSymbol{)}\<%
\\
\>[0]\AgdaFunction{substb-cong}\AgdaSpace{}%
\AgdaBound{h}\AgdaSpace{}%
\AgdaSymbol{(}\AgdaInductiveConstructor{ƛb}\AgdaSpace{}%
\AgdaBound{t}\AgdaSymbol{)}\AgdaSpace{}%
\AgdaSymbol{=}\AgdaSpace{}%
\AgdaFunction{cong}\AgdaSpace{}%
\AgdaInductiveConstructor{ƛb}\AgdaSpace{}%
\AgdaSymbol{(}\AgdaFunction{substb-cong}\AgdaSpace{}%
\AgdaSymbol{(}\AgdaFunction{extsb-cong}\AgdaSpace{}%
\AgdaBound{h}\AgdaSymbol{)}\AgdaSpace{}%
\AgdaBound{t}\AgdaSymbol{)}\<%
\\
\>[0]\AgdaFunction{substb-cong}\AgdaSpace{}%
\AgdaBound{h}\AgdaSpace{}%
\AgdaSymbol{(}\AgdaBound{t}\AgdaSpace{}%
\AgdaOperator{\AgdaInductiveConstructor{·b}}\AgdaSpace{}%
\AgdaBound{t₁}\AgdaSymbol{)}\AgdaSpace{}%
\AgdaSymbol{=}\AgdaSpace{}%
\AgdaFunction{cong₂}\AgdaSpace{}%
\AgdaOperator{\AgdaInductiveConstructor{\AgdaUnderscore{}·b\AgdaUnderscore{}}}\AgdaSpace{}%
\AgdaSymbol{(}\AgdaFunction{substb-cong}\AgdaSpace{}%
\AgdaBound{h}\AgdaSpace{}%
\AgdaBound{t}\AgdaSymbol{)}\AgdaSpace{}%
\AgdaSymbol{(}\AgdaFunction{substb-cong}\AgdaSpace{}%
\AgdaBound{h}\AgdaSpace{}%
\AgdaBound{t₁}\AgdaSymbol{)}\<%
\end{code}
\begin{code}[hide]%
\>[0]\<%
\\
\>[0]\AgdaComment{--\ Commutativity\ of\ renameb\ and\ extsb:}\<%
\\
\>[0]\AgdaFunction{extsb-rename-commute}\AgdaSpace{}%
\AgdaSymbol{:}\AgdaSpace{}%
\AgdaSymbol{∀}\AgdaSpace{}%
\AgdaSymbol{\{}\AgdaBound{Γ₁}\AgdaSpace{}%
\AgdaBound{Γ₂}\AgdaSpace{}%
\AgdaBound{Γ₃}\AgdaSpace{}%
\AgdaBound{A}\AgdaSpace{}%
\AgdaBound{B}\AgdaSymbol{\}}\<%
\\
\>[0][@{}l@{\AgdaIndent{0}}]%
\>[2]\AgdaSymbol{→}\AgdaSpace{}%
\AgdaSymbol{(}\AgdaBound{fb'}\AgdaSpace{}%
\AgdaSymbol{:}\AgdaSpace{}%
\AgdaSymbol{∀}\AgdaSpace{}%
\AgdaSymbol{\{}\AgdaBound{Ab}\AgdaSymbol{\}}\AgdaSpace{}%
\AgdaSymbol{→}\AgdaSpace{}%
\AgdaBound{Ab}\AgdaSpace{}%
\AgdaOperator{\AgdaDatatype{∈b}}\AgdaSpace{}%
\AgdaBound{Γ₂}\AgdaSpace{}%
\AgdaSymbol{→}\AgdaSpace{}%
\AgdaBound{Ab}\AgdaSpace{}%
\AgdaOperator{\AgdaDatatype{∈b}}\AgdaSpace{}%
\AgdaBound{Γ₃}\AgdaSymbol{)}\<%
\\
%
\>[2]\AgdaSymbol{→}\AgdaSpace{}%
\AgdaSymbol{(}\AgdaBound{fb''}\AgdaSpace{}%
\AgdaSymbol{:}\AgdaSpace{}%
\AgdaSymbol{∀}\AgdaSpace{}%
\AgdaSymbol{\{}\AgdaBound{Ab}\AgdaSymbol{\}}\AgdaSpace{}%
\AgdaSymbol{→}\AgdaSpace{}%
\AgdaInductiveConstructor{□}\AgdaSpace{}%
\AgdaBound{Ab}\AgdaSpace{}%
\AgdaOperator{\AgdaDatatype{∈b}}\AgdaSpace{}%
\AgdaBound{Γ₁}\AgdaSpace{}%
\AgdaSymbol{→}\AgdaSpace{}%
\AgdaBound{Γ₂}\AgdaSpace{}%
\AgdaOperator{\AgdaDatatype{⊢b}}\AgdaSpace{}%
\AgdaBound{Ab}\AgdaSymbol{)}\<%
\\
%
\>[2]\AgdaSymbol{→}\AgdaSpace{}%
\AgdaSymbol{(}\AgdaBound{x}\AgdaSpace{}%
\AgdaSymbol{:}\AgdaSpace{}%
\AgdaInductiveConstructor{□}\AgdaSpace{}%
\AgdaBound{A}\AgdaSpace{}%
\AgdaOperator{\AgdaDatatype{∈b}}\AgdaSpace{}%
\AgdaBound{Γ₁}\AgdaSpace{}%
\AgdaOperator{\AgdaInductiveConstructor{,b}}\AgdaSpace{}%
\AgdaBound{B}\AgdaSymbol{)}\<%
\\
%
\>[2]\AgdaSymbol{→}%
\>[1170I]\AgdaFunction{extsb}\AgdaSpace{}%
\AgdaSymbol{(λ}\AgdaSpace{}%
\AgdaBound{x₁}\AgdaSpace{}%
\AgdaSymbol{→}\AgdaSpace{}%
\AgdaFunction{renameb}\AgdaSpace{}%
\AgdaBound{fb'}\AgdaSpace{}%
\AgdaSymbol{(}\AgdaBound{fb''}\AgdaSpace{}%
\AgdaBound{x₁}\AgdaSymbol{))}\AgdaSpace{}%
\AgdaBound{x}\AgdaSpace{}%
\AgdaOperator{\AgdaDatatype{≡}}\<%
\\
\>[.][@{}l@{}]\<[1170I]%
\>[4]\AgdaFunction{renameb}\AgdaSpace{}%
\AgdaSymbol{(}\AgdaFunction{extb}\AgdaSpace{}%
\AgdaBound{fb'}\AgdaSymbol{)}\AgdaSpace{}%
\AgdaSymbol{(}\AgdaFunction{extsb}\AgdaSpace{}%
\AgdaBound{fb''}\AgdaSpace{}%
\AgdaBound{x}\AgdaSymbol{)}\<%
\\
\>[0]\AgdaFunction{extsb-rename-commute}\AgdaSpace{}%
\AgdaBound{fb'}\AgdaSpace{}%
\AgdaBound{fb''}\AgdaSpace{}%
\AgdaInductiveConstructor{Zb}\AgdaSpace{}%
\AgdaSymbol{=}\AgdaSpace{}%
\AgdaInductiveConstructor{refl}\<%
\\
\>[0]\AgdaFunction{extsb-rename-commute}\AgdaSpace{}%
\AgdaBound{fb'}\AgdaSpace{}%
\AgdaBound{fb''}\AgdaSpace{}%
\AgdaSymbol{(}\AgdaInductiveConstructor{Sb}\AgdaSpace{}%
\AgdaBound{x}\AgdaSymbol{)}\AgdaSpace{}%
\AgdaSymbol{=}\AgdaSpace{}%
\AgdaFunction{sym}\AgdaSpace{}%
\AgdaSymbol{(}\AgdaFunction{rename-extb-Sb-commute}\AgdaSpace{}%
\AgdaBound{fb'}\AgdaSpace{}%
\AgdaSymbol{(}\AgdaBound{fb''}\AgdaSpace{}%
\AgdaBound{x}\AgdaSymbol{))}\<%
\\
%
\\[\AgdaEmptyExtraSkip]%
\>[0]\AgdaComment{--\ Commutativity\ of\ substb\ and\ renameb}\<%
\\
\>[0]\AgdaFunction{substb-rename-commute}\AgdaSpace{}%
\AgdaSymbol{:}\AgdaSpace{}%
\AgdaSymbol{∀}\AgdaSpace{}%
\AgdaSymbol{\{}\AgdaBound{Γ₁}\AgdaSpace{}%
\AgdaBound{Γ₂}\AgdaSpace{}%
\AgdaBound{Γ₃}\AgdaSpace{}%
\AgdaBound{A}\AgdaSymbol{\}}\<%
\\
\>[0][@{}l@{\AgdaIndent{0}}]%
\>[2]\AgdaSymbol{→}\AgdaSpace{}%
\AgdaSymbol{(}\AgdaBound{fb'}\AgdaSpace{}%
\AgdaSymbol{:}\AgdaSpace{}%
\AgdaSymbol{∀}\AgdaSpace{}%
\AgdaSymbol{\{}\AgdaBound{Ab}\AgdaSymbol{\}}\AgdaSpace{}%
\AgdaSymbol{→}\AgdaSpace{}%
\AgdaBound{Ab}\AgdaSpace{}%
\AgdaOperator{\AgdaDatatype{∈b}}\AgdaSpace{}%
\AgdaBound{Γ₂}\AgdaSpace{}%
\AgdaSymbol{→}\AgdaSpace{}%
\AgdaBound{Ab}\AgdaSpace{}%
\AgdaOperator{\AgdaDatatype{∈b}}\AgdaSpace{}%
\AgdaBound{Γ₃}\AgdaSymbol{)}\<%
\\
%
\>[2]\AgdaSymbol{→}\AgdaSpace{}%
\AgdaSymbol{(}\AgdaBound{fb''}\AgdaSpace{}%
\AgdaSymbol{:}\AgdaSpace{}%
\AgdaSymbol{∀}\AgdaSpace{}%
\AgdaSymbol{\{}\AgdaBound{Ab}\AgdaSymbol{\}}\AgdaSpace{}%
\AgdaSymbol{→}\AgdaSpace{}%
\AgdaInductiveConstructor{□}\AgdaSpace{}%
\AgdaBound{Ab}\AgdaSpace{}%
\AgdaOperator{\AgdaDatatype{∈b}}\AgdaSpace{}%
\AgdaBound{Γ₁}\AgdaSpace{}%
\AgdaSymbol{→}\AgdaSpace{}%
\AgdaBound{Γ₂}\AgdaSpace{}%
\AgdaOperator{\AgdaDatatype{⊢b}}\AgdaSpace{}%
\AgdaBound{Ab}\AgdaSymbol{)}\<%
\\
%
\>[2]\AgdaSymbol{→}\AgdaSpace{}%
\AgdaSymbol{(}\AgdaBound{t}\AgdaSpace{}%
\AgdaSymbol{:}\AgdaSpace{}%
\AgdaBound{Γ₁}\AgdaSpace{}%
\AgdaOperator{\AgdaDatatype{⊢b}}\AgdaSpace{}%
\AgdaBound{A}\AgdaSymbol{)}\<%
\\
%
\>[2]\AgdaSymbol{→}\AgdaSpace{}%
\AgdaFunction{substb}\AgdaSpace{}%
\AgdaSymbol{(λ}\AgdaSpace{}%
\AgdaBound{x}\AgdaSpace{}%
\AgdaSymbol{→}\AgdaSpace{}%
\AgdaFunction{renameb}\AgdaSpace{}%
\AgdaBound{fb'}\AgdaSpace{}%
\AgdaSymbol{(}\AgdaBound{fb''}\AgdaSpace{}%
\AgdaBound{x}\AgdaSymbol{))}\AgdaSpace{}%
\AgdaBound{t}\AgdaSpace{}%
\AgdaOperator{\AgdaDatatype{≡}}\AgdaSpace{}%
\AgdaFunction{renameb}\AgdaSpace{}%
\AgdaBound{fb'}\AgdaSpace{}%
\AgdaSymbol{(}\AgdaFunction{substb}\AgdaSpace{}%
\AgdaBound{fb''}\AgdaSpace{}%
\AgdaBound{t}\AgdaSymbol{)}\<%
\\
\>[0]\AgdaFunction{substb-rename-commute}\AgdaSpace{}%
\AgdaBound{fb'}\AgdaSpace{}%
\AgdaBound{fb''}\AgdaSpace{}%
\AgdaSymbol{(}\AgdaInductiveConstructor{ε}\AgdaSpace{}%
\AgdaBound{x}\AgdaSymbol{)}\AgdaSpace{}%
\AgdaSymbol{=}\AgdaSpace{}%
\AgdaInductiveConstructor{refl}\<%
\\
\>[0]\AgdaFunction{substb-rename-commute}\AgdaSpace{}%
\AgdaBound{fb'}\AgdaSpace{}%
\AgdaBound{fb''}\AgdaSpace{}%
\AgdaSymbol{(}\AgdaInductiveConstructor{box}\AgdaSpace{}%
\AgdaBound{t}\AgdaSymbol{)}\AgdaSpace{}%
\AgdaSymbol{=}\<%
\\
\>[0][@{}l@{\AgdaIndent{0}}]%
\>[2]\AgdaFunction{cong}\<%
\\
\>[2][@{}l@{\AgdaIndent{0}}]%
\>[4]\AgdaInductiveConstructor{box}\<%
\\
%
\>[4]\AgdaSymbol{(}\AgdaFunction{substb-rename-commute}\AgdaSpace{}%
\AgdaBound{fb'}\AgdaSpace{}%
\AgdaBound{fb''}\AgdaSpace{}%
\AgdaBound{t}\AgdaSymbol{)}\<%
\\
\>[0]\AgdaFunction{substb-rename-commute}\AgdaSpace{}%
\AgdaBound{fb'}\AgdaSpace{}%
\AgdaBound{fb''}\AgdaSpace{}%
\AgdaSymbol{(}\AgdaInductiveConstructor{ƛb}\AgdaSpace{}%
\AgdaBound{t}\AgdaSymbol{)}\AgdaSpace{}%
\AgdaSymbol{=}\<%
\\
\>[0][@{}l@{\AgdaIndent{0}}]%
\>[2]\AgdaOperator{\AgdaFunction{beginEq}}\<%
\\
\>[2][@{}l@{\AgdaIndent{0}}]%
\>[4]\AgdaInductiveConstructor{ƛb}\AgdaSpace{}%
\AgdaSymbol{(}\AgdaFunction{substb}\AgdaSpace{}%
\AgdaSymbol{(}\AgdaFunction{extsb}\AgdaSpace{}%
\AgdaSymbol{(λ}\AgdaSpace{}%
\AgdaBound{x}\AgdaSpace{}%
\AgdaSymbol{→}\AgdaSpace{}%
\AgdaFunction{renameb}\AgdaSpace{}%
\AgdaBound{fb'}\AgdaSpace{}%
\AgdaSymbol{(}\AgdaBound{fb''}\AgdaSpace{}%
\AgdaBound{x}\AgdaSymbol{)))}\AgdaSpace{}%
\AgdaBound{t}\AgdaSymbol{)}\<%
\\
%
\>[2]\AgdaFunction{≡⟨}\AgdaSpace{}%
\AgdaFunction{cong}\AgdaSpace{}%
\AgdaInductiveConstructor{ƛb}\AgdaSpace{}%
\AgdaSymbol{(}\AgdaFunction{substb-cong}\AgdaSpace{}%
\AgdaSymbol{(λ}\AgdaSpace{}%
\AgdaBound{x}\AgdaSpace{}%
\AgdaSymbol{→}\AgdaSpace{}%
\AgdaFunction{extsb-rename-commute}\AgdaSpace{}%
\AgdaBound{fb'}\AgdaSpace{}%
\AgdaBound{fb''}\AgdaSpace{}%
\AgdaBound{x}\AgdaSymbol{)}\AgdaSpace{}%
\AgdaBound{t}\AgdaSymbol{)}\AgdaSpace{}%
\AgdaFunction{⟩}\<%
\\
\>[2][@{}l@{\AgdaIndent{0}}]%
\>[4]\AgdaInductiveConstructor{ƛb}\AgdaSpace{}%
\AgdaSymbol{(}\AgdaFunction{substb}\AgdaSpace{}%
\AgdaSymbol{(λ}\AgdaSpace{}%
\AgdaBound{x}\AgdaSpace{}%
\AgdaSymbol{→}\AgdaSpace{}%
\AgdaFunction{renameb}\AgdaSpace{}%
\AgdaSymbol{(}\AgdaFunction{extb}\AgdaSpace{}%
\AgdaBound{fb'}\AgdaSymbol{)}\AgdaSpace{}%
\AgdaSymbol{((}\AgdaFunction{extsb}\AgdaSpace{}%
\AgdaBound{fb''}\AgdaSymbol{)}\AgdaSpace{}%
\AgdaBound{x}\AgdaSymbol{))}\AgdaSpace{}%
\AgdaBound{t}\AgdaSymbol{)}\<%
\\
%
\>[2]\AgdaFunction{≡⟨}\AgdaSpace{}%
\AgdaFunction{cong}\AgdaSpace{}%
\AgdaInductiveConstructor{ƛb}\AgdaSpace{}%
\AgdaSymbol{(}\AgdaFunction{substb-rename-commute}\AgdaSpace{}%
\AgdaSymbol{(}\AgdaFunction{extb}\AgdaSpace{}%
\AgdaBound{fb'}\AgdaSymbol{)}\AgdaSpace{}%
\AgdaSymbol{(}\AgdaFunction{extsb}\AgdaSpace{}%
\AgdaBound{fb''}\AgdaSymbol{)}\AgdaSpace{}%
\AgdaBound{t}\AgdaSymbol{)}\AgdaSpace{}%
\AgdaFunction{⟩}\<%
\\
\>[2][@{}l@{\AgdaIndent{0}}]%
\>[4]\AgdaInductiveConstructor{ƛb}\AgdaSpace{}%
\AgdaSymbol{(}\AgdaFunction{renameb}\AgdaSpace{}%
\AgdaSymbol{(}\AgdaFunction{extb}\AgdaSpace{}%
\AgdaBound{fb'}\AgdaSymbol{)}\AgdaSpace{}%
\AgdaSymbol{(}\AgdaFunction{substb}\AgdaSpace{}%
\AgdaSymbol{(}\AgdaFunction{extsb}\AgdaSpace{}%
\AgdaBound{fb''}\AgdaSymbol{)}\AgdaSpace{}%
\AgdaBound{t}\AgdaSymbol{))}\<%
\\
%
\>[2]\AgdaOperator{\AgdaFunction{∎eq}}\<%
\\
\>[0]\AgdaFunction{substb-rename-commute}\AgdaSpace{}%
\AgdaBound{fb'}\AgdaSpace{}%
\AgdaBound{fb''}\AgdaSpace{}%
\AgdaSymbol{(}\AgdaBound{t}\AgdaSpace{}%
\AgdaOperator{\AgdaInductiveConstructor{·b}}\AgdaSpace{}%
\AgdaBound{t₁}\AgdaSymbol{)}\AgdaSpace{}%
\AgdaSymbol{=}\<%
\\
\>[0][@{}l@{\AgdaIndent{0}}]%
\>[2]\AgdaFunction{cong₂}\<%
\\
\>[2][@{}l@{\AgdaIndent{0}}]%
\>[4]\AgdaOperator{\AgdaInductiveConstructor{\AgdaUnderscore{}·b\AgdaUnderscore{}}}\<%
\\
%
\>[4]\AgdaSymbol{(}\AgdaFunction{substb-rename-commute}\AgdaSpace{}%
\AgdaBound{fb'}\AgdaSpace{}%
\AgdaBound{fb''}\AgdaSpace{}%
\AgdaBound{t}\AgdaSymbol{)}\<%
\\
%
\>[4]\AgdaSymbol{(}\AgdaFunction{substb-rename-commute}\AgdaSpace{}%
\AgdaBound{fb'}\AgdaSpace{}%
\AgdaBound{fb''}\AgdaSpace{}%
\AgdaBound{t₁}\AgdaSymbol{)}\<%
\\
%
\\[\AgdaEmptyExtraSkip]%
\>[0]\AgdaComment{--\ extsb\ distributes\ over\ composition:}\<%
\\
\>[0]\AgdaComment{--\ Executing\ extsb\ first\ with\ fb',\ then\ with\ fb''}\<%
\\
\>[0]\AgdaComment{--\ is\ the\ same\ as\ doing\ extsb\ once\ with\ fb''\ ∘\ fb'.}\<%
\\
\>[0]\AgdaFunction{extsb-comp}\AgdaSpace{}%
\AgdaSymbol{:}\AgdaSpace{}%
\AgdaSymbol{∀}\AgdaSpace{}%
\AgdaSymbol{\{}\AgdaBound{Γ₁}\AgdaSpace{}%
\AgdaBound{Γ₂}\AgdaSpace{}%
\AgdaBound{Γ₃}\AgdaSpace{}%
\AgdaBound{A}\AgdaSpace{}%
\AgdaBound{B}\AgdaSymbol{\}}\<%
\\
\>[0][@{}l@{\AgdaIndent{0}}]%
\>[2]\AgdaSymbol{→}\AgdaSpace{}%
\AgdaSymbol{(}\AgdaBound{fb'}\AgdaSpace{}%
\AgdaSymbol{:}\AgdaSpace{}%
\AgdaSymbol{∀}\AgdaSpace{}%
\AgdaSymbol{\{}\AgdaBound{Ab}\AgdaSymbol{\}}\AgdaSpace{}%
\AgdaSymbol{→}\AgdaSpace{}%
\AgdaBound{Ab}\AgdaSpace{}%
\AgdaOperator{\AgdaDatatype{∈b}}\AgdaSpace{}%
\AgdaBound{Γ₁}\AgdaSpace{}%
\AgdaSymbol{→}\AgdaSpace{}%
\AgdaBound{Ab}\AgdaSpace{}%
\AgdaOperator{\AgdaDatatype{∈b}}\AgdaSpace{}%
\AgdaBound{Γ₂}\AgdaSymbol{)}\<%
\\
%
\>[2]\AgdaSymbol{→}\AgdaSpace{}%
\AgdaSymbol{(}\AgdaBound{fb''}\AgdaSpace{}%
\AgdaSymbol{:}\AgdaSpace{}%
\AgdaSymbol{∀}\AgdaSpace{}%
\AgdaSymbol{\{}\AgdaBound{Ab}\AgdaSymbol{\}}\AgdaSpace{}%
\AgdaSymbol{→}\AgdaSpace{}%
\AgdaInductiveConstructor{□}\AgdaSpace{}%
\AgdaBound{Ab}\AgdaSpace{}%
\AgdaOperator{\AgdaDatatype{∈b}}\AgdaSpace{}%
\AgdaBound{Γ₂}\AgdaSpace{}%
\AgdaSymbol{→}\AgdaSpace{}%
\AgdaBound{Γ₃}\AgdaSpace{}%
\AgdaOperator{\AgdaDatatype{⊢b}}\AgdaSpace{}%
\AgdaBound{Ab}\AgdaSymbol{)}\<%
\\
%
\>[2]\AgdaSymbol{→}\AgdaSpace{}%
\AgdaSymbol{(}\AgdaBound{x}\AgdaSpace{}%
\AgdaSymbol{:}\AgdaSpace{}%
\AgdaInductiveConstructor{□}\AgdaSpace{}%
\AgdaBound{A}\AgdaSpace{}%
\AgdaOperator{\AgdaDatatype{∈b}}\AgdaSpace{}%
\AgdaBound{Γ₁}\AgdaSpace{}%
\AgdaOperator{\AgdaInductiveConstructor{,b}}\AgdaSpace{}%
\AgdaBound{B}\AgdaSymbol{)}\<%
\\
%
\>[2]\AgdaSymbol{→}\AgdaSpace{}%
\AgdaFunction{extsb}\AgdaSpace{}%
\AgdaBound{fb''}\AgdaSpace{}%
\AgdaSymbol{(}\AgdaFunction{extb}\AgdaSpace{}%
\AgdaBound{fb'}\AgdaSpace{}%
\AgdaBound{x}\AgdaSymbol{)}\AgdaSpace{}%
\AgdaOperator{\AgdaDatatype{≡}}\AgdaSpace{}%
\AgdaFunction{extsb}\AgdaSpace{}%
\AgdaSymbol{(λ}\AgdaSpace{}%
\AgdaBound{x₁}\AgdaSpace{}%
\AgdaSymbol{→}\AgdaSpace{}%
\AgdaBound{fb''}\AgdaSpace{}%
\AgdaSymbol{(}\AgdaBound{fb'}\AgdaSpace{}%
\AgdaBound{x₁}\AgdaSymbol{))}\AgdaSpace{}%
\AgdaBound{x}\<%
\\
\>[0]\AgdaFunction{extsb-comp}\AgdaSpace{}%
\AgdaBound{fb'}\AgdaSpace{}%
\AgdaBound{fb''}\AgdaSpace{}%
\AgdaInductiveConstructor{Zb}\AgdaSpace{}%
\AgdaSymbol{=}\AgdaSpace{}%
\AgdaInductiveConstructor{refl}\<%
\\
\>[0]\AgdaFunction{extsb-comp}\AgdaSpace{}%
\AgdaBound{fb'}\AgdaSpace{}%
\AgdaBound{fb''}\AgdaSpace{}%
\AgdaSymbol{(}\AgdaInductiveConstructor{Sb}\AgdaSpace{}%
\AgdaBound{x}\AgdaSymbol{)}\AgdaSpace{}%
\AgdaSymbol{=}\AgdaSpace{}%
\AgdaInductiveConstructor{refl}\<%
\\
%
\\[\AgdaEmptyExtraSkip]%
\>[0]\AgdaComment{--\ substb\ and\ rename\ distribute\ over\ composition:}\<%
\\
\>[0]\AgdaComment{--\ Renaming\ first\ with\ fb'\ and\ substituting\ with\ fb''}\<%
\\
\>[0]\AgdaComment{--\ is\ the\ same\ as\ substituting\ with\ fb''\ ∘\ fb'}\<%
\\
\>[0]\AgdaFunction{subst-renameb-comp}\AgdaSpace{}%
\AgdaSymbol{:}\AgdaSpace{}%
\AgdaSymbol{∀}\AgdaSpace{}%
\AgdaSymbol{\{}\AgdaBound{Γ₁}\AgdaSpace{}%
\AgdaBound{Γ₂}\AgdaSpace{}%
\AgdaBound{Γ₃}\AgdaSpace{}%
\AgdaBound{A}\AgdaSymbol{\}}\<%
\\
\>[0][@{}l@{\AgdaIndent{0}}]%
\>[2]\AgdaSymbol{→}\AgdaSpace{}%
\AgdaSymbol{(}\AgdaBound{fb'}\AgdaSpace{}%
\AgdaSymbol{:}\AgdaSpace{}%
\AgdaSymbol{∀}\AgdaSpace{}%
\AgdaSymbol{\{}\AgdaBound{Ab}\AgdaSymbol{\}}\AgdaSpace{}%
\AgdaSymbol{→}\AgdaSpace{}%
\AgdaBound{Ab}\AgdaSpace{}%
\AgdaOperator{\AgdaDatatype{∈b}}\AgdaSpace{}%
\AgdaBound{Γ₁}\AgdaSpace{}%
\AgdaSymbol{→}\AgdaSpace{}%
\AgdaBound{Ab}\AgdaSpace{}%
\AgdaOperator{\AgdaDatatype{∈b}}\AgdaSpace{}%
\AgdaBound{Γ₂}\AgdaSymbol{)}\<%
\\
%
\>[2]\AgdaSymbol{→}\AgdaSpace{}%
\AgdaSymbol{(}\AgdaBound{fb''}\AgdaSpace{}%
\AgdaSymbol{:}\AgdaSpace{}%
\AgdaSymbol{∀}\AgdaSpace{}%
\AgdaSymbol{\{}\AgdaBound{Ab}\AgdaSymbol{\}}\AgdaSpace{}%
\AgdaSymbol{→}\AgdaSpace{}%
\AgdaInductiveConstructor{□}\AgdaSpace{}%
\AgdaBound{Ab}\AgdaSpace{}%
\AgdaOperator{\AgdaDatatype{∈b}}\AgdaSpace{}%
\AgdaBound{Γ₂}\AgdaSpace{}%
\AgdaSymbol{→}\AgdaSpace{}%
\AgdaBound{Γ₃}\AgdaSpace{}%
\AgdaOperator{\AgdaDatatype{⊢b}}\AgdaSpace{}%
\AgdaBound{Ab}\AgdaSymbol{)}\<%
\\
%
\>[2]\AgdaSymbol{→}\AgdaSpace{}%
\AgdaSymbol{(}\AgdaBound{t}\AgdaSpace{}%
\AgdaSymbol{:}\AgdaSpace{}%
\AgdaBound{Γ₁}\AgdaSpace{}%
\AgdaOperator{\AgdaDatatype{⊢b}}\AgdaSpace{}%
\AgdaBound{A}\AgdaSymbol{)}\<%
\\
%
\>[2]\AgdaSymbol{→}\AgdaSpace{}%
\AgdaFunction{substb}\AgdaSpace{}%
\AgdaBound{fb''}\AgdaSpace{}%
\AgdaSymbol{(}\AgdaFunction{renameb}\AgdaSpace{}%
\AgdaBound{fb'}\AgdaSpace{}%
\AgdaBound{t}\AgdaSymbol{)}\AgdaSpace{}%
\AgdaOperator{\AgdaDatatype{≡}}\AgdaSpace{}%
\AgdaFunction{substb}\AgdaSpace{}%
\AgdaSymbol{(λ}\AgdaSpace{}%
\AgdaBound{x}\AgdaSpace{}%
\AgdaSymbol{→}\AgdaSpace{}%
\AgdaBound{fb''}\AgdaSpace{}%
\AgdaSymbol{(}\AgdaBound{fb'}\AgdaSpace{}%
\AgdaBound{x}\AgdaSymbol{))}\AgdaSpace{}%
\AgdaBound{t}\<%
\\
\>[0]\AgdaFunction{subst-renameb-comp}\AgdaSpace{}%
\AgdaBound{fb'}\AgdaSpace{}%
\AgdaBound{fb''}\AgdaSpace{}%
\AgdaSymbol{(}\AgdaInductiveConstructor{ε}\AgdaSpace{}%
\AgdaBound{x}\AgdaSymbol{)}\AgdaSpace{}%
\AgdaSymbol{=}\AgdaSpace{}%
\AgdaInductiveConstructor{refl}\<%
\\
\>[0]\AgdaFunction{subst-renameb-comp}\AgdaSpace{}%
\AgdaBound{fb'}\AgdaSpace{}%
\AgdaBound{fb''}\AgdaSpace{}%
\AgdaSymbol{(}\AgdaInductiveConstructor{box}\AgdaSpace{}%
\AgdaBound{t}\AgdaSymbol{)}\AgdaSpace{}%
\AgdaSymbol{=}\AgdaSpace{}%
\AgdaFunction{cong}\AgdaSpace{}%
\AgdaInductiveConstructor{box}\AgdaSpace{}%
\AgdaSymbol{(}\AgdaFunction{subst-renameb-comp}\AgdaSpace{}%
\AgdaBound{fb'}\AgdaSpace{}%
\AgdaBound{fb''}\AgdaSpace{}%
\AgdaBound{t}\AgdaSymbol{)}\<%
\\
\>[0]\AgdaFunction{subst-renameb-comp}\AgdaSpace{}%
\AgdaBound{fb'}\AgdaSpace{}%
\AgdaBound{fb''}\AgdaSpace{}%
\AgdaSymbol{(}\AgdaInductiveConstructor{ƛb}\AgdaSpace{}%
\AgdaBound{t}\AgdaSymbol{)}\AgdaSpace{}%
\AgdaSymbol{=}\<%
\\
\>[0][@{}l@{\AgdaIndent{0}}]%
\>[2]\AgdaOperator{\AgdaFunction{beginEq}}\<%
\\
\>[2][@{}l@{\AgdaIndent{0}}]%
\>[4]\AgdaInductiveConstructor{ƛb}\AgdaSpace{}%
\AgdaSymbol{(}\AgdaFunction{substb}\AgdaSpace{}%
\AgdaSymbol{(}\AgdaFunction{extsb}\AgdaSpace{}%
\AgdaBound{fb''}\AgdaSymbol{)}\AgdaSpace{}%
\AgdaSymbol{(}\AgdaFunction{renameb}\AgdaSpace{}%
\AgdaSymbol{(}\AgdaFunction{extb}\AgdaSpace{}%
\AgdaBound{fb'}\AgdaSymbol{)}\AgdaSpace{}%
\AgdaBound{t}\AgdaSymbol{))}\<%
\\
%
\>[2]\AgdaFunction{≡⟨}\AgdaSpace{}%
\AgdaFunction{cong}\AgdaSpace{}%
\AgdaInductiveConstructor{ƛb}\AgdaSpace{}%
\AgdaSymbol{(}\AgdaFunction{subst-renameb-comp}\AgdaSpace{}%
\AgdaSymbol{(}\AgdaFunction{extb}\AgdaSpace{}%
\AgdaBound{fb'}\AgdaSymbol{)}\AgdaSpace{}%
\AgdaSymbol{(}\AgdaFunction{extsb}\AgdaSpace{}%
\AgdaBound{fb''}\AgdaSymbol{)}\AgdaSpace{}%
\AgdaBound{t}\AgdaSymbol{)}\AgdaSpace{}%
\AgdaFunction{⟩}\<%
\\
\>[2][@{}l@{\AgdaIndent{0}}]%
\>[4]\AgdaInductiveConstructor{ƛb}\AgdaSpace{}%
\AgdaSymbol{(}\AgdaFunction{substb}\AgdaSpace{}%
\AgdaSymbol{(λ}\AgdaSpace{}%
\AgdaBound{x}\AgdaSpace{}%
\AgdaSymbol{→}\AgdaSpace{}%
\AgdaSymbol{(}\AgdaFunction{extsb}\AgdaSpace{}%
\AgdaBound{fb''}\AgdaSymbol{)}\AgdaSpace{}%
\AgdaSymbol{((}\AgdaFunction{extb}\AgdaSpace{}%
\AgdaBound{fb'}\AgdaSymbol{)}\AgdaSpace{}%
\AgdaBound{x}\AgdaSymbol{))}\AgdaSpace{}%
\AgdaBound{t}\AgdaSymbol{)}\<%
\\
%
\>[2]\AgdaFunction{≡⟨}\AgdaSpace{}%
\AgdaFunction{cong}\AgdaSpace{}%
\AgdaInductiveConstructor{ƛb}\AgdaSpace{}%
\AgdaSymbol{(}\AgdaFunction{substb-cong}\AgdaSpace{}%
\AgdaSymbol{(λ}\AgdaSpace{}%
\AgdaBound{x}\AgdaSpace{}%
\AgdaSymbol{→}\AgdaSpace{}%
\AgdaFunction{extsb-comp}\AgdaSpace{}%
\AgdaBound{fb'}\AgdaSpace{}%
\AgdaBound{fb''}\AgdaSpace{}%
\AgdaBound{x}\AgdaSymbol{)}\AgdaSpace{}%
\AgdaBound{t}\AgdaSymbol{)}\AgdaSpace{}%
\AgdaFunction{⟩}\<%
\\
\>[2][@{}l@{\AgdaIndent{0}}]%
\>[4]\AgdaInductiveConstructor{ƛb}\AgdaSpace{}%
\AgdaSymbol{(}\AgdaFunction{substb}\AgdaSpace{}%
\AgdaSymbol{(}\AgdaFunction{extsb}\AgdaSpace{}%
\AgdaSymbol{(λ}\AgdaSpace{}%
\AgdaBound{x}\AgdaSpace{}%
\AgdaSymbol{→}\AgdaSpace{}%
\AgdaBound{fb''}\AgdaSpace{}%
\AgdaSymbol{(}\AgdaBound{fb'}\AgdaSpace{}%
\AgdaBound{x}\AgdaSymbol{)))}\AgdaSpace{}%
\AgdaBound{t}\AgdaSymbol{)}\<%
\\
%
\>[2]\AgdaOperator{\AgdaFunction{∎eq}}\<%
\\
\>[0]\AgdaFunction{subst-renameb-comp}\AgdaSpace{}%
\AgdaBound{fb'}\AgdaSpace{}%
\AgdaBound{fb''}\AgdaSpace{}%
\AgdaSymbol{(}\AgdaBound{t}\AgdaSpace{}%
\AgdaOperator{\AgdaInductiveConstructor{·b}}\AgdaSpace{}%
\AgdaBound{t₁}\AgdaSymbol{)}\AgdaSpace{}%
\AgdaSymbol{=}\AgdaSpace{}%
\AgdaFunction{cong₂}\AgdaSpace{}%
\AgdaOperator{\AgdaInductiveConstructor{\AgdaUnderscore{}·b\AgdaUnderscore{}}}\AgdaSpace{}%
\AgdaSymbol{(}\AgdaFunction{subst-renameb-comp}\AgdaSpace{}%
\AgdaBound{fb'}\AgdaSpace{}%
\AgdaBound{fb''}\AgdaSpace{}%
\AgdaBound{t}\AgdaSymbol{)}\AgdaSpace{}%
\AgdaSymbol{(}\AgdaFunction{subst-renameb-comp}\AgdaSpace{}%
\AgdaBound{fb'}\AgdaSpace{}%
\AgdaBound{fb''}\AgdaSpace{}%
\AgdaBound{t₁}\AgdaSymbol{)}\<%
\\
\>[0]\AgdaComment{--\ Renaming\ followed\ by\ substitution\ with\ an\ extended\ context\ (extsb)}\<%
\\
\>[0]\AgdaComment{--\ commutes\ with\ substitution\ followed\ by\ renaming.}\<%
\\
\>[0]\AgdaFunction{substb-exts-rename-Sb-commute}\AgdaSpace{}%
\AgdaSymbol{:}\AgdaSpace{}%
\AgdaSymbol{∀}\AgdaSpace{}%
\AgdaSymbol{\{}\AgdaBound{Γ}\AgdaSpace{}%
\AgdaBound{A}\AgdaSpace{}%
\AgdaBound{B}\AgdaSpace{}%
\AgdaBound{Δ}\AgdaSymbol{\}}\<%
\\
\>[0][@{}l@{\AgdaIndent{0}}]%
\>[2]\AgdaSymbol{→}\AgdaSpace{}%
\AgdaSymbol{(}\AgdaBound{fb'}\AgdaSpace{}%
\AgdaSymbol{:}\AgdaSpace{}%
\AgdaSymbol{∀}\AgdaSpace{}%
\AgdaSymbol{\{}\AgdaBound{Ab}\AgdaSymbol{\}}\AgdaSpace{}%
\AgdaSymbol{→}\AgdaSpace{}%
\AgdaInductiveConstructor{□}\AgdaSpace{}%
\AgdaBound{Ab}\AgdaSpace{}%
\AgdaOperator{\AgdaDatatype{∈b}}\AgdaSpace{}%
\AgdaBound{Γ}\AgdaSpace{}%
\AgdaSymbol{→}\AgdaSpace{}%
\AgdaBound{Δ}\AgdaSpace{}%
\AgdaOperator{\AgdaDatatype{⊢b}}\AgdaSpace{}%
\AgdaBound{Ab}\AgdaSymbol{)}\<%
\\
%
\>[2]\AgdaSymbol{→}\AgdaSpace{}%
\AgdaSymbol{(}\AgdaBound{t}\AgdaSpace{}%
\AgdaSymbol{:}\AgdaSpace{}%
\AgdaBound{Γ}\AgdaSpace{}%
\AgdaOperator{\AgdaDatatype{⊢b}}\AgdaSpace{}%
\AgdaBound{A}\AgdaSymbol{)}\<%
\\
%
\>[2]\AgdaSymbol{→}%
\>[1555I]\AgdaFunction{substb}\AgdaSpace{}%
\AgdaSymbol{(}\AgdaFunction{extsb}\AgdaSpace{}%
\AgdaBound{fb'}\AgdaSpace{}%
\AgdaSymbol{\{}\AgdaArgument{B}\AgdaSpace{}%
\AgdaSymbol{=}\AgdaSpace{}%
\AgdaBound{B}\AgdaSymbol{\})}\AgdaSpace{}%
\AgdaSymbol{(}\AgdaFunction{renameb}\AgdaSpace{}%
\AgdaInductiveConstructor{Sb}\AgdaSpace{}%
\AgdaBound{t}\AgdaSymbol{)}\<%
\\
\>[.][@{}l@{}]\<[1555I]%
\>[4]\AgdaOperator{\AgdaDatatype{≡}}\AgdaSpace{}%
\AgdaFunction{renameb}\AgdaSpace{}%
\AgdaInductiveConstructor{Sb}\AgdaSpace{}%
\AgdaSymbol{(}\AgdaFunction{substb}\AgdaSpace{}%
\AgdaBound{fb'}\AgdaSpace{}%
\AgdaBound{t}\AgdaSymbol{)}\<%
\\
\>[0]\AgdaFunction{substb-exts-rename-Sb-commute}\AgdaSpace{}%
\AgdaBound{fb'}\AgdaSpace{}%
\AgdaBound{t}\AgdaSpace{}%
\AgdaSymbol{=}\<%
\\
\>[0][@{}l@{\AgdaIndent{0}}]%
\>[2]\AgdaOperator{\AgdaFunction{beginEq}}\<%
\\
\>[2][@{}l@{\AgdaIndent{0}}]%
\>[4]\AgdaFunction{substb}\AgdaSpace{}%
\AgdaSymbol{(}\AgdaFunction{extsb}\AgdaSpace{}%
\AgdaBound{fb'}\AgdaSymbol{)}\AgdaSpace{}%
\AgdaSymbol{(}\AgdaFunction{renameb}\AgdaSpace{}%
\AgdaInductiveConstructor{Sb}\AgdaSpace{}%
\AgdaBound{t}\AgdaSymbol{)}\<%
\\
%
\>[2]\AgdaFunction{≡⟨}\AgdaSpace{}%
\AgdaFunction{subst-renameb-comp}\AgdaSpace{}%
\AgdaInductiveConstructor{Sb}\AgdaSpace{}%
\AgdaSymbol{(}\AgdaFunction{extsb}\AgdaSpace{}%
\AgdaBound{fb'}\AgdaSymbol{)}\AgdaSpace{}%
\AgdaBound{t}\AgdaSpace{}%
\AgdaFunction{⟩}\<%
\\
\>[2][@{}l@{\AgdaIndent{0}}]%
\>[4]\AgdaFunction{substb}\AgdaSpace{}%
\AgdaSymbol{(λ}\AgdaSpace{}%
\AgdaBound{x}\AgdaSpace{}%
\AgdaSymbol{→}\AgdaSpace{}%
\AgdaFunction{renameb}\AgdaSpace{}%
\AgdaInductiveConstructor{Sb}\AgdaSpace{}%
\AgdaSymbol{(}\AgdaBound{fb'}\AgdaSpace{}%
\AgdaBound{x}\AgdaSymbol{))}\AgdaSpace{}%
\AgdaBound{t}\<%
\\
%
\>[2]\AgdaFunction{≡⟨}\AgdaSpace{}%
\AgdaFunction{substb-rename-commute}\AgdaSpace{}%
\AgdaInductiveConstructor{Sb}\AgdaSpace{}%
\AgdaBound{fb'}\AgdaSpace{}%
\AgdaBound{t}\AgdaSpace{}%
\AgdaFunction{⟩}\<%
\\
\>[2][@{}l@{\AgdaIndent{0}}]%
\>[4]\AgdaFunction{renameb}\AgdaSpace{}%
\AgdaInductiveConstructor{Sb}\AgdaSpace{}%
\AgdaSymbol{(}\AgdaFunction{substb}\AgdaSpace{}%
\AgdaBound{fb'}\AgdaSpace{}%
\AgdaBound{t}\AgdaSymbol{)}\<%
\\
%
\>[2]\AgdaOperator{\AgdaFunction{∎eq}}\<%
\\
\>[0]\AgdaComment{--\ Preservation\ and\ Reflection\ of\ }\<%
\\
\>[0]\AgdaComment{--\ simultaneous\ substitution\ of\ Gödel's\ translation}\<%
\\
\>[0]\AgdaFunction{pres-simultaneous-subst}\AgdaSpace{}%
\AgdaSymbol{:}\AgdaSpace{}%
\AgdaSymbol{∀}\AgdaSpace{}%
\AgdaSymbol{\{}\AgdaBound{Γ}\AgdaSpace{}%
\AgdaBound{Δ}\AgdaSpace{}%
\AgdaBound{A}\AgdaSymbol{\}}\<%
\\
\>[0][@{}l@{\AgdaIndent{0}}]%
\>[2]\AgdaSymbol{→}\AgdaSpace{}%
\AgdaSymbol{(}\AgdaBound{fb'}\AgdaSpace{}%
\AgdaSymbol{:}\AgdaSpace{}%
\AgdaSymbol{∀}\AgdaSpace{}%
\AgdaSymbol{\{}\AgdaBound{Ab}\AgdaSymbol{\}}\AgdaSpace{}%
\AgdaSymbol{→}\AgdaSpace{}%
\AgdaInductiveConstructor{□}\AgdaSpace{}%
\AgdaBound{Ab}\AgdaSpace{}%
\AgdaOperator{\AgdaDatatype{∈b}}\AgdaSpace{}%
\AgdaFunction{embedContext}\AgdaSpace{}%
\AgdaBound{Γ}\AgdaSpace{}%
\AgdaSymbol{→}\AgdaSpace{}%
\AgdaFunction{embedContext}\AgdaSpace{}%
\AgdaBound{Δ}\AgdaSpace{}%
\AgdaOperator{\AgdaDatatype{⊢b}}\AgdaSpace{}%
\AgdaBound{Ab}\AgdaSymbol{)}\<%
\\
%
\>[2]\AgdaSymbol{→}\AgdaSpace{}%
\AgdaSymbol{(}\AgdaBound{f'}\AgdaSpace{}%
\AgdaSymbol{:}\AgdaSpace{}%
\AgdaSymbol{∀}\AgdaSpace{}%
\AgdaSymbol{\{}\AgdaBound{A}\AgdaSymbol{\}}\AgdaSpace{}%
\AgdaSymbol{→}\AgdaSpace{}%
\AgdaBound{A}\AgdaSpace{}%
\AgdaOperator{\AgdaDatatype{∈}}\AgdaSpace{}%
\AgdaBound{Γ}\AgdaSpace{}%
\AgdaSymbol{→}\AgdaSpace{}%
\AgdaBound{Δ}\AgdaSpace{}%
\AgdaOperator{\AgdaDatatype{⊢}}\AgdaSpace{}%
\AgdaBound{A}\AgdaSymbol{)}\<%
\\
%
\>[2]\AgdaSymbol{→}\AgdaSpace{}%
\AgdaSymbol{(}\AgdaBound{h}\AgdaSpace{}%
\AgdaSymbol{:}\AgdaSpace{}%
\AgdaSymbol{∀}\AgdaSpace{}%
\AgdaSymbol{\{}\AgdaBound{A'}\AgdaSymbol{\}}\AgdaSpace{}%
\AgdaSymbol{→}\AgdaSpace{}%
\AgdaSymbol{(}\AgdaBound{x}\AgdaSpace{}%
\AgdaSymbol{:}\AgdaSpace{}%
\AgdaBound{A'}\AgdaSpace{}%
\AgdaOperator{\AgdaDatatype{∈}}\AgdaSpace{}%
\AgdaBound{Γ}\AgdaSymbol{)}\AgdaSpace{}%
\AgdaSymbol{→}\AgdaSpace{}%
\AgdaInductiveConstructor{box}\AgdaSpace{}%
\AgdaSymbol{(}\AgdaBound{fb'}\AgdaSpace{}%
\AgdaSymbol{(}\AgdaFunction{embedLookup}\AgdaSpace{}%
\AgdaBound{x}\AgdaSymbol{))}\AgdaSpace{}%
\AgdaOperator{\AgdaDatatype{≡}}\AgdaSpace{}%
\AgdaFunction{embedTerm}\AgdaSpace{}%
\AgdaSymbol{(}\AgdaBound{f'}\AgdaSpace{}%
\AgdaBound{x}\AgdaSymbol{)}\AgdaSpace{}%
\AgdaSymbol{)}\<%
\\
%
\>[2]\AgdaSymbol{→}\AgdaSpace{}%
\AgdaSymbol{(}\AgdaBound{M}\AgdaSpace{}%
\AgdaSymbol{:}\AgdaSpace{}%
\AgdaBound{Γ}\AgdaSpace{}%
\AgdaOperator{\AgdaDatatype{⊢}}\AgdaSpace{}%
\AgdaBound{A}\AgdaSymbol{)}\<%
\\
%
\>[2]\AgdaSymbol{→}\AgdaSpace{}%
\AgdaFunction{substb}\AgdaSpace{}%
\AgdaBound{fb'}\AgdaSpace{}%
\AgdaSymbol{(}\AgdaFunction{embedTerm}\AgdaSpace{}%
\AgdaBound{M}\AgdaSymbol{)}\AgdaSpace{}%
\AgdaOperator{\AgdaDatatype{≡}}\AgdaSpace{}%
\AgdaFunction{embedTerm}\AgdaSpace{}%
\AgdaSymbol{(}\AgdaFunction{subst}\AgdaSpace{}%
\AgdaBound{f'}\AgdaSpace{}%
\AgdaBound{M}\AgdaSymbol{)}\<%
\\
\>[0]\AgdaFunction{pres-simultaneous-subst}\AgdaSpace{}%
\AgdaBound{fb'}\AgdaSpace{}%
\AgdaBound{f'}\AgdaSpace{}%
\AgdaBound{h}\AgdaSpace{}%
\AgdaSymbol{(}\AgdaInductiveConstructor{`}\AgdaSpace{}%
\AgdaBound{x}\AgdaSymbol{)}\AgdaSpace{}%
\AgdaSymbol{=}\<%
\\
\>[0][@{}l@{\AgdaIndent{0}}]%
\>[2]\AgdaBound{h}\AgdaSpace{}%
\AgdaBound{x}\<%
\\
\>[0]\AgdaFunction{pres-simultaneous-subst}\AgdaSpace{}%
\AgdaBound{fb'}\AgdaSpace{}%
\AgdaBound{f'}\AgdaSpace{}%
\AgdaBound{h}\AgdaSpace{}%
\AgdaSymbol{(}\AgdaInductiveConstructor{ƛ}\AgdaSpace{}%
\AgdaBound{M}\AgdaSymbol{)}\AgdaSpace{}%
\AgdaSymbol{=}\<%
\\
\>[0][@{}l@{\AgdaIndent{0}}]%
\>[2]\AgdaFunction{cong}\AgdaSpace{}%
\AgdaInductiveConstructor{box}\AgdaSpace{}%
\AgdaSymbol{(}\AgdaFunction{cong}\AgdaSpace{}%
\AgdaInductiveConstructor{ƛb}\AgdaSpace{}%
\AgdaSymbol{(}\AgdaFunction{pres-simultaneous-subst}\AgdaSpace{}%
\AgdaSymbol{(}\AgdaFunction{extsb}\AgdaSpace{}%
\AgdaBound{fb'}\AgdaSymbol{)}\AgdaSpace{}%
\AgdaSymbol{(}\AgdaFunction{exts}\AgdaSpace{}%
\AgdaBound{f'}\AgdaSymbol{)}\AgdaSpace{}%
\AgdaSymbol{(}\AgdaFunction{pres-exts}\AgdaSpace{}%
\AgdaBound{fb'}\AgdaSpace{}%
\AgdaBound{f'}\AgdaSpace{}%
\AgdaBound{h}\AgdaSymbol{)}\AgdaSpace{}%
\AgdaBound{M}\AgdaSymbol{))}\<%
\\
\>[0]\AgdaFunction{pres-simultaneous-subst}\AgdaSpace{}%
\AgdaBound{fb'}\AgdaSpace{}%
\AgdaBound{f'}\AgdaSpace{}%
\AgdaBound{h}\AgdaSpace{}%
\AgdaSymbol{(}\AgdaBound{M}\AgdaSpace{}%
\AgdaOperator{\AgdaInductiveConstructor{·}}\AgdaSpace{}%
\AgdaBound{M₁}\AgdaSymbol{)}\AgdaSpace{}%
\AgdaSymbol{=}\<%
\\
\>[0][@{}l@{\AgdaIndent{0}}]%
\>[2]\AgdaOperator{\AgdaFunction{beginEq}}\<%
\\
\>[2][@{}l@{\AgdaIndent{0}}]%
\>[4]\AgdaInductiveConstructor{ƛb}\AgdaSpace{}%
\AgdaSymbol{(}\AgdaInductiveConstructor{ε}\AgdaSpace{}%
\AgdaInductiveConstructor{Zb}\AgdaSpace{}%
\AgdaOperator{\AgdaInductiveConstructor{·b}}\AgdaSpace{}%
\AgdaFunction{substb}\AgdaSpace{}%
\AgdaSymbol{(}\AgdaFunction{extsb}\AgdaSpace{}%
\AgdaBound{fb'}\AgdaSymbol{)}\AgdaSpace{}%
\AgdaSymbol{(}\AgdaFunction{renameb}\AgdaSpace{}%
\AgdaInductiveConstructor{Sb}\AgdaSpace{}%
\AgdaSymbol{(}\AgdaFunction{embedTerm}\AgdaSpace{}%
\AgdaBound{M₁}\AgdaSymbol{)))}\AgdaSpace{}%
\AgdaOperator{\AgdaInductiveConstructor{·b}}\AgdaSpace{}%
\AgdaFunction{substb}\AgdaSpace{}%
\AgdaBound{fb'}\AgdaSpace{}%
\AgdaSymbol{(}\AgdaFunction{embedTerm}\AgdaSpace{}%
\AgdaBound{M}\AgdaSymbol{)}\<%
\\
%
\>[2]\AgdaFunction{≡⟨}\AgdaSpace{}%
\AgdaFunction{cong₂}\AgdaSpace{}%
\AgdaOperator{\AgdaInductiveConstructor{\AgdaUnderscore{}·b\AgdaUnderscore{}}}\AgdaSpace{}%
\AgdaInductiveConstructor{refl}\AgdaSpace{}%
\AgdaSymbol{(}\AgdaFunction{pres-simultaneous-subst}\AgdaSpace{}%
\AgdaBound{fb'}\AgdaSpace{}%
\AgdaBound{f'}\AgdaSpace{}%
\AgdaBound{h}\AgdaSpace{}%
\AgdaBound{M}\AgdaSymbol{)}\AgdaSpace{}%
\AgdaFunction{⟩}\<%
\\
\>[2][@{}l@{\AgdaIndent{0}}]%
\>[4]\AgdaInductiveConstructor{ƛb}\AgdaSpace{}%
\AgdaSymbol{(}\AgdaInductiveConstructor{ε}\AgdaSpace{}%
\AgdaInductiveConstructor{Zb}\AgdaSpace{}%
\AgdaOperator{\AgdaInductiveConstructor{·b}}\AgdaSpace{}%
\AgdaFunction{substb}\AgdaSpace{}%
\AgdaSymbol{(}\AgdaFunction{extsb}\AgdaSpace{}%
\AgdaBound{fb'}\AgdaSymbol{)}\AgdaSpace{}%
\AgdaSymbol{(}\AgdaFunction{renameb}\AgdaSpace{}%
\AgdaInductiveConstructor{Sb}\AgdaSpace{}%
\AgdaSymbol{(}\AgdaFunction{embedTerm}\AgdaSpace{}%
\AgdaBound{M₁}\AgdaSymbol{)))}\AgdaSpace{}%
\AgdaOperator{\AgdaInductiveConstructor{·b}}\AgdaSpace{}%
\AgdaFunction{embedTerm}\AgdaSpace{}%
\AgdaSymbol{(}\AgdaFunction{subst}\AgdaSpace{}%
\AgdaBound{f'}\AgdaSpace{}%
\AgdaBound{M}\AgdaSymbol{)}\<%
\\
%
\>[2]\AgdaFunction{≡⟨}\AgdaSpace{}%
\AgdaFunction{cong₂}\AgdaSpace{}%
\AgdaOperator{\AgdaInductiveConstructor{\AgdaUnderscore{}·b\AgdaUnderscore{}}}\AgdaSpace{}%
\AgdaSymbol{(}\AgdaFunction{cong}\AgdaSpace{}%
\AgdaInductiveConstructor{ƛb}\AgdaSpace{}%
\AgdaSymbol{(}\AgdaFunction{cong₂}\AgdaSpace{}%
\AgdaOperator{\AgdaInductiveConstructor{\AgdaUnderscore{}·b\AgdaUnderscore{}}}\AgdaSpace{}%
\AgdaInductiveConstructor{refl}\AgdaSpace{}%
\AgdaSymbol{(}\AgdaFunction{substb-exts-rename-Sb-commute}\AgdaSpace{}%
\AgdaBound{fb'}\AgdaSpace{}%
\AgdaSymbol{(}\AgdaFunction{embedTerm}\AgdaSpace{}%
\AgdaBound{M₁}\AgdaSymbol{))))}\AgdaSpace{}%
\AgdaInductiveConstructor{refl}\AgdaSpace{}%
\AgdaFunction{⟩}\<%
\\
\>[2][@{}l@{\AgdaIndent{0}}]%
\>[4]\AgdaInductiveConstructor{ƛb}\AgdaSpace{}%
\AgdaSymbol{(}\AgdaInductiveConstructor{ε}\AgdaSpace{}%
\AgdaInductiveConstructor{Zb}\AgdaSpace{}%
\AgdaOperator{\AgdaInductiveConstructor{·b}}\AgdaSpace{}%
\AgdaFunction{renameb}\AgdaSpace{}%
\AgdaInductiveConstructor{Sb}\AgdaSpace{}%
\AgdaSymbol{(}\AgdaFunction{substb}\AgdaSpace{}%
\AgdaBound{fb'}\AgdaSpace{}%
\AgdaSymbol{(}\AgdaFunction{embedTerm}\AgdaSpace{}%
\AgdaBound{M₁}\AgdaSymbol{)))}\AgdaSpace{}%
\AgdaOperator{\AgdaInductiveConstructor{·b}}\AgdaSpace{}%
\AgdaFunction{embedTerm}\AgdaSpace{}%
\AgdaSymbol{(}\AgdaFunction{subst}\AgdaSpace{}%
\AgdaBound{f'}\AgdaSpace{}%
\AgdaBound{M}\AgdaSymbol{)}\<%
\\
%
\>[2]\AgdaFunction{≡⟨}\AgdaSpace{}%
\AgdaFunction{cong₂}\AgdaSpace{}%
\AgdaOperator{\AgdaInductiveConstructor{\AgdaUnderscore{}·b\AgdaUnderscore{}}}\AgdaSpace{}%
\AgdaSymbol{(}\AgdaFunction{cong}\AgdaSpace{}%
\AgdaInductiveConstructor{ƛb}\AgdaSpace{}%
\AgdaSymbol{(}\AgdaFunction{cong₂}\AgdaSpace{}%
\AgdaOperator{\AgdaInductiveConstructor{\AgdaUnderscore{}·b\AgdaUnderscore{}}}\AgdaSpace{}%
\AgdaInductiveConstructor{refl}\AgdaSpace{}%
\AgdaSymbol{(}\AgdaFunction{cong}\AgdaSpace{}%
\AgdaSymbol{(}\AgdaFunction{renameb}\AgdaSpace{}%
\AgdaInductiveConstructor{Sb}\AgdaSymbol{)}\AgdaSpace{}%
\AgdaSymbol{(}\AgdaFunction{pres-simultaneous-subst}\AgdaSpace{}%
\AgdaBound{fb'}\AgdaSpace{}%
\AgdaBound{f'}\AgdaSpace{}%
\AgdaBound{h}\AgdaSpace{}%
\AgdaBound{M₁}\AgdaSymbol{))))}\AgdaSpace{}%
\AgdaInductiveConstructor{refl}\AgdaSpace{}%
\AgdaFunction{⟩}\<%
\\
\>[2][@{}l@{\AgdaIndent{0}}]%
\>[4]\AgdaInductiveConstructor{ƛb}\AgdaSpace{}%
\AgdaSymbol{(}\AgdaInductiveConstructor{ε}\AgdaSpace{}%
\AgdaInductiveConstructor{Zb}\AgdaSpace{}%
\AgdaOperator{\AgdaInductiveConstructor{·b}}\AgdaSpace{}%
\AgdaFunction{renameb}\AgdaSpace{}%
\AgdaInductiveConstructor{Sb}\AgdaSpace{}%
\AgdaSymbol{(}\AgdaFunction{embedTerm}\AgdaSpace{}%
\AgdaSymbol{(}\AgdaFunction{subst}\AgdaSpace{}%
\AgdaBound{f'}\AgdaSpace{}%
\AgdaBound{M₁}\AgdaSymbol{)))}\AgdaSpace{}%
\AgdaOperator{\AgdaInductiveConstructor{·b}}\AgdaSpace{}%
\AgdaFunction{embedTerm}\AgdaSpace{}%
\AgdaSymbol{(}\AgdaFunction{subst}\AgdaSpace{}%
\AgdaBound{f'}\AgdaSpace{}%
\AgdaBound{M}\AgdaSymbol{)}\<%
\\
%
\>[2]\AgdaOperator{\AgdaFunction{∎eq}}\<%
\\
%
\\[\AgdaEmptyExtraSkip]%
\>[0]\AgdaComment{--\ HELPER\ LEMMA'S\ TO\ PROVE\ h}\<%
\\
%
\\[\AgdaEmptyExtraSkip]%
\>[0]\AgdaComment{--\ Lemma\ for\ assumption\ h\ in\ lemma}\<%
\\
\>[0]\AgdaComment{--\ pres-simultaneous-subst}\<%
\\
\>[0]\AgdaFunction{pres-subst-func}\AgdaSpace{}%
\AgdaSymbol{:}\AgdaSpace{}%
\AgdaSymbol{∀}\AgdaSpace{}%
\AgdaSymbol{\{}\AgdaBound{A}\AgdaSpace{}%
\AgdaBound{C}\AgdaSpace{}%
\AgdaBound{Γ}\AgdaSymbol{\}}\<%
\\
\>[0][@{}l@{\AgdaIndent{0}}]%
\>[2]\AgdaSymbol{→}\AgdaSpace{}%
\AgdaSymbol{(}\AgdaBound{V}\AgdaSpace{}%
\AgdaSymbol{:}\AgdaSpace{}%
\AgdaBound{Γ}\AgdaSpace{}%
\AgdaOperator{\AgdaDatatype{⊢}}\AgdaSpace{}%
\AgdaBound{C}\AgdaSymbol{)}\<%
\\
%
\>[2]\AgdaSymbol{→}\AgdaSpace{}%
\AgdaSymbol{(}\AgdaBound{y}\AgdaSpace{}%
\AgdaSymbol{:}\AgdaSpace{}%
\AgdaDatatype{Value}\AgdaSpace{}%
\AgdaBound{V}\AgdaSymbol{)}\<%
\\
%
\>[2]\AgdaSymbol{→}\AgdaSpace{}%
\AgdaSymbol{(}\AgdaBound{x}\AgdaSpace{}%
\AgdaSymbol{:}\AgdaSpace{}%
\AgdaBound{A}\AgdaSpace{}%
\AgdaOperator{\AgdaDatatype{∈}}\AgdaSpace{}%
\AgdaBound{Γ}\AgdaSpace{}%
\AgdaOperator{\AgdaInductiveConstructor{,}}\AgdaSpace{}%
\AgdaBound{C}\AgdaSymbol{)}\<%
\\
%
\>[2]\AgdaSymbol{→}\AgdaSpace{}%
\AgdaInductiveConstructor{box}\AgdaSpace{}%
\AgdaSymbol{(}\AgdaFunction{fb}\AgdaSpace{}%
\AgdaSymbol{(}\AgdaFunction{embedValue}\AgdaSpace{}%
\AgdaBound{V}\AgdaSpace{}%
\AgdaBound{y}\AgdaSymbol{)}\AgdaSpace{}%
\AgdaSymbol{(}\AgdaFunction{embedLookup}\AgdaSpace{}%
\AgdaBound{x}\AgdaSymbol{))}\AgdaSpace{}%
\AgdaOperator{\AgdaDatatype{≡}}\AgdaSpace{}%
\AgdaFunction{embedTerm}\AgdaSpace{}%
\AgdaSymbol{((}\AgdaFunction{f}\AgdaSpace{}%
\AgdaBound{V}\AgdaSymbol{)}\AgdaSpace{}%
\AgdaBound{x}\AgdaSymbol{)}\<%
\\
\>[0]\AgdaFunction{pres-subst-func}\AgdaSpace{}%
\AgdaBound{V}\AgdaSpace{}%
\AgdaBound{y}\AgdaSpace{}%
\AgdaInductiveConstructor{Z}\AgdaSpace{}%
\AgdaSymbol{=}\AgdaSpace{}%
\AgdaFunction{box-embedValue}\AgdaSpace{}%
\AgdaBound{y}\<%
\\
\>[0]\AgdaFunction{pres-subst-func}\AgdaSpace{}%
\AgdaBound{V}\AgdaSpace{}%
\AgdaBound{y}\AgdaSpace{}%
\AgdaSymbol{(}\AgdaInductiveConstructor{S}\AgdaSpace{}%
\AgdaBound{x}\AgdaSymbol{)}\AgdaSpace{}%
\AgdaSymbol{=}\AgdaSpace{}%
\AgdaInductiveConstructor{refl}\<%
\\
%
\\[\AgdaEmptyExtraSkip]%
\>[0]\AgdaComment{--\ exts\ lemma\ for\ assumption\ h\ in\ lemma}\<%
\\
\>[0]\AgdaFunction{exts-pres-subst-func}\AgdaSpace{}%
\AgdaSymbol{:}\AgdaSpace{}%
\AgdaSymbol{∀}\AgdaSpace{}%
\AgdaSymbol{\{}\AgdaBound{A}\AgdaSpace{}%
\AgdaBound{C}\AgdaSpace{}%
\AgdaBound{Γ}\AgdaSpace{}%
\AgdaBound{B}\AgdaSymbol{\}}\<%
\\
\>[0][@{}l@{\AgdaIndent{0}}]%
\>[2]\AgdaSymbol{→}\AgdaSpace{}%
\AgdaSymbol{(}\AgdaBound{V}\AgdaSpace{}%
\AgdaSymbol{:}\AgdaSpace{}%
\AgdaBound{Γ}\AgdaSpace{}%
\AgdaOperator{\AgdaDatatype{⊢}}\AgdaSpace{}%
\AgdaBound{C}\AgdaSymbol{)}\<%
\\
%
\>[2]\AgdaSymbol{→}\AgdaSpace{}%
\AgdaSymbol{(}\AgdaBound{y}\AgdaSpace{}%
\AgdaSymbol{:}\AgdaSpace{}%
\AgdaDatatype{Value}\AgdaSpace{}%
\AgdaBound{V}\AgdaSymbol{)}\<%
\\
%
\>[2]\AgdaSymbol{→}\AgdaSpace{}%
\AgdaSymbol{(}\AgdaBound{x}\AgdaSpace{}%
\AgdaSymbol{:}\AgdaSpace{}%
\AgdaBound{A}\AgdaSpace{}%
\AgdaOperator{\AgdaDatatype{∈}}\AgdaSpace{}%
\AgdaBound{Γ}\AgdaSpace{}%
\AgdaOperator{\AgdaInductiveConstructor{,}}\AgdaSpace{}%
\AgdaBound{C}\AgdaSpace{}%
\AgdaOperator{\AgdaInductiveConstructor{,}}\AgdaSpace{}%
\AgdaBound{B}\AgdaSymbol{)}\<%
\\
%
\>[2]\AgdaSymbol{→}\AgdaSpace{}%
\AgdaInductiveConstructor{box}\AgdaSpace{}%
\AgdaSymbol{(}\AgdaFunction{extsb}\AgdaSpace{}%
\AgdaSymbol{(}\AgdaFunction{fb}\AgdaSpace{}%
\AgdaSymbol{(}\AgdaFunction{embedValue}\AgdaSpace{}%
\AgdaBound{V}\AgdaSpace{}%
\AgdaBound{y}\AgdaSymbol{))}\AgdaSpace{}%
\AgdaSymbol{(}\AgdaFunction{embedLookup}\AgdaSpace{}%
\AgdaBound{x}\AgdaSymbol{))}\AgdaSpace{}%
\AgdaOperator{\AgdaDatatype{≡}}\AgdaSpace{}%
\AgdaFunction{embedTerm}\AgdaSpace{}%
\AgdaSymbol{(}\AgdaFunction{exts}\AgdaSpace{}%
\AgdaSymbol{(}\AgdaFunction{f}\AgdaSpace{}%
\AgdaBound{V}\AgdaSymbol{)}\AgdaSpace{}%
\AgdaBound{x}\AgdaSymbol{)}\<%
\\
\>[0]\AgdaFunction{exts-pres-subst-func}\AgdaSpace{}%
\AgdaBound{V}\AgdaSpace{}%
\AgdaBound{y}\AgdaSpace{}%
\AgdaInductiveConstructor{Z}\AgdaSpace{}%
\AgdaSymbol{=}\AgdaSpace{}%
\AgdaInductiveConstructor{refl}\<%
\\
\>[0]\AgdaFunction{exts-pres-subst-func}\AgdaSpace{}%
\AgdaBound{V}\AgdaSpace{}%
\AgdaBound{y}\AgdaSpace{}%
\AgdaSymbol{(}\AgdaInductiveConstructor{S}\AgdaSpace{}%
\AgdaBound{x}\AgdaSymbol{)}\AgdaSpace{}%
\AgdaSymbol{=}\<%
\\
\>[0][@{}l@{\AgdaIndent{0}}]%
\>[2]\AgdaOperator{\AgdaFunction{beginEq}}\<%
\\
\>[2][@{}l@{\AgdaIndent{0}}]%
\>[4]\AgdaInductiveConstructor{box}\AgdaSpace{}%
\AgdaSymbol{(}\AgdaFunction{renameb}\AgdaSpace{}%
\AgdaInductiveConstructor{Sb}\AgdaSpace{}%
\AgdaSymbol{(}\AgdaFunction{fb}\AgdaSpace{}%
\AgdaSymbol{(}\AgdaFunction{embedValue}\AgdaSpace{}%
\AgdaBound{V}\AgdaSpace{}%
\AgdaBound{y}\AgdaSymbol{)}\AgdaSpace{}%
\AgdaSymbol{(}\AgdaFunction{embedLookup}\AgdaSpace{}%
\AgdaBound{x}\AgdaSymbol{)))}\<%
\\
%
\>[2]\AgdaFunction{≡⟨}\AgdaSpace{}%
\AgdaInductiveConstructor{refl}\AgdaSpace{}%
\AgdaFunction{⟩}\<%
\\
\>[2][@{}l@{\AgdaIndent{0}}]%
\>[4]\AgdaFunction{renameb}\AgdaSpace{}%
\AgdaInductiveConstructor{Sb}\AgdaSpace{}%
\AgdaSymbol{(}\AgdaInductiveConstructor{box}\AgdaSpace{}%
\AgdaSymbol{(}\AgdaFunction{fb}\AgdaSpace{}%
\AgdaSymbol{(}\AgdaFunction{embedValue}\AgdaSpace{}%
\AgdaBound{V}\AgdaSpace{}%
\AgdaBound{y}\AgdaSymbol{)}\AgdaSpace{}%
\AgdaSymbol{(}\AgdaFunction{embedLookup}\AgdaSpace{}%
\AgdaBound{x}\AgdaSymbol{)))}\<%
\\
%
\>[2]\AgdaFunction{≡⟨}\AgdaSpace{}%
\AgdaFunction{cong}\AgdaSpace{}%
\AgdaSymbol{(}\AgdaFunction{renameb}\AgdaSpace{}%
\AgdaInductiveConstructor{Sb}\AgdaSymbol{)}\AgdaSpace{}%
\AgdaSymbol{(}\AgdaFunction{pres-subst-func}\AgdaSpace{}%
\AgdaBound{V}\AgdaSpace{}%
\AgdaBound{y}\AgdaSpace{}%
\AgdaBound{x}\AgdaSymbol{)}\AgdaSpace{}%
\AgdaFunction{⟩}\<%
\\
\>[2][@{}l@{\AgdaIndent{0}}]%
\>[4]\AgdaFunction{renameb}\AgdaSpace{}%
\AgdaInductiveConstructor{Sb}\AgdaSpace{}%
\AgdaSymbol{(}\AgdaFunction{embedTerm}\AgdaSpace{}%
\AgdaSymbol{(}\AgdaSpace{}%
\AgdaSymbol{(}\AgdaFunction{f}\AgdaSpace{}%
\AgdaBound{V}\AgdaSymbol{)}\AgdaSpace{}%
\AgdaBound{x}\AgdaSpace{}%
\AgdaSymbol{))}\<%
\\
%
\>[2]\AgdaFunction{≡⟨}\AgdaSpace{}%
\AgdaFunction{pres-rename}\AgdaSpace{}%
\AgdaInductiveConstructor{S}\AgdaSpace{}%
\AgdaInductiveConstructor{Sb}\AgdaSpace{}%
\AgdaSymbol{(λ}\AgdaSpace{}%
\AgdaBound{x₁}\AgdaSpace{}%
\AgdaSymbol{→}\AgdaSpace{}%
\AgdaInductiveConstructor{refl}\AgdaSymbol{)}\AgdaSpace{}%
\AgdaSymbol{(}\AgdaFunction{f}\AgdaSpace{}%
\AgdaBound{V}\AgdaSpace{}%
\AgdaBound{x}\AgdaSymbol{)}\AgdaSpace{}%
\AgdaFunction{⟩}\<%
\\
\>[2][@{}l@{\AgdaIndent{0}}]%
\>[4]\AgdaFunction{embedTerm}\AgdaSpace{}%
\AgdaSymbol{(}\AgdaFunction{rename}\AgdaSpace{}%
\AgdaInductiveConstructor{S}\AgdaSpace{}%
\AgdaSymbol{(}\AgdaFunction{f}\AgdaSpace{}%
\AgdaBound{V}\AgdaSpace{}%
\AgdaBound{x}\AgdaSymbol{))}\<%
\\
%
\>[2]\AgdaOperator{\AgdaFunction{∎eq}}\<%
\\
%
\\[\AgdaEmptyExtraSkip]%
\>[0]\AgdaComment{--\ Preservation\ and\ Reflection\ of\ }\<%
\\
\>[0]\AgdaComment{--\ single\ substitution\ of\ Gödel's\ translation}\<%
\\
\>[0]\AgdaFunction{pres-subst}\AgdaSpace{}%
\AgdaSymbol{:}\AgdaSpace{}%
\AgdaSymbol{∀}\AgdaSpace{}%
\AgdaSymbol{\{}\AgdaBound{Γ}\AgdaSpace{}%
\AgdaBound{A}\AgdaSpace{}%
\AgdaBound{B}\AgdaSymbol{\}}\AgdaSpace{}%
\AgdaSymbol{→}\AgdaSpace{}%
\AgdaSymbol{(}\AgdaBound{M}\AgdaSpace{}%
\AgdaSymbol{:}\AgdaSpace{}%
\AgdaBound{Γ}\AgdaSpace{}%
\AgdaOperator{\AgdaInductiveConstructor{,}}\AgdaSpace{}%
\AgdaBound{B}\AgdaSpace{}%
\AgdaOperator{\AgdaDatatype{⊢}}\AgdaSpace{}%
\AgdaBound{A}\AgdaSymbol{)}\AgdaSpace{}%
\AgdaSymbol{→}\AgdaSpace{}%
\AgdaSymbol{(}\AgdaBound{V}\AgdaSpace{}%
\AgdaSymbol{:}\AgdaSpace{}%
\AgdaBound{Γ}\AgdaSpace{}%
\AgdaOperator{\AgdaDatatype{⊢}}\AgdaSpace{}%
\AgdaBound{B}\AgdaSymbol{)}\AgdaSpace{}%
\AgdaSymbol{→}\AgdaSpace{}%
\AgdaSymbol{(}\AgdaBound{x}\AgdaSpace{}%
\AgdaSymbol{:}\AgdaSpace{}%
\AgdaDatatype{Value}\AgdaSpace{}%
\AgdaBound{V}\AgdaSymbol{)}\AgdaSpace{}%
\AgdaSymbol{→}\AgdaSpace{}%
\AgdaFunction{embedTerm}\AgdaSpace{}%
\AgdaBound{M}\AgdaSpace{}%
\AgdaOperator{\AgdaFunction{[}}\AgdaSpace{}%
\AgdaFunction{embedValue}\AgdaSpace{}%
\AgdaBound{V}\AgdaSpace{}%
\AgdaBound{x}\AgdaSpace{}%
\AgdaOperator{\AgdaFunction{]b}}\AgdaSpace{}%
\AgdaOperator{\AgdaDatatype{≡}}\AgdaSpace{}%
\AgdaFunction{embedTerm}\AgdaSpace{}%
\AgdaSymbol{(}\AgdaBound{M}\AgdaSpace{}%
\AgdaOperator{\AgdaFunction{[}}\AgdaSpace{}%
\AgdaBound{V}\AgdaSpace{}%
\AgdaOperator{\AgdaFunction{]}}\AgdaSymbol{)}\<%
\\
\>[0]\AgdaFunction{pres-subst}\AgdaSpace{}%
\AgdaBound{M}\AgdaSpace{}%
\AgdaBound{V}\AgdaSpace{}%
\AgdaBound{x}\AgdaSpace{}%
\AgdaSymbol{=}\AgdaSpace{}%
\AgdaFunction{pres-simultaneous-subst}\AgdaSpace{}%
\AgdaSymbol{(}\AgdaFunction{fb}\AgdaSpace{}%
\AgdaSymbol{(}\AgdaFunction{embedValue}\AgdaSpace{}%
\AgdaBound{V}\AgdaSpace{}%
\AgdaBound{x}\AgdaSymbol{))}\AgdaSpace{}%
\AgdaSymbol{(}\AgdaFunction{f}\AgdaSpace{}%
\AgdaBound{V}\AgdaSymbol{)}\AgdaSpace{}%
\AgdaSymbol{(}\AgdaFunction{pres-subst-func}\AgdaSpace{}%
\AgdaBound{V}\AgdaSpace{}%
\AgdaBound{x}\AgdaSymbol{)}\AgdaSpace{}%
\AgdaBound{M}\<%
\\
%
\\[\AgdaEmptyExtraSkip]%
\>[0]\AgdaComment{--\ Helper\ lemma:\ if\ a\ value\ is\ substituted\ in\ a\ value,\ the\ resulting\ lambda\ term\ will\ also\ be\ a\ value}\<%
\\
\>[0]\AgdaFunction{val-after-subst}\AgdaSpace{}%
\AgdaSymbol{:}\AgdaSpace{}%
\AgdaSymbol{∀}\AgdaSpace{}%
\AgdaSymbol{\{}\AgdaBound{Γ}\AgdaSpace{}%
\AgdaBound{A}\AgdaSpace{}%
\AgdaBound{B}\AgdaSymbol{\}}\AgdaSpace{}%
\AgdaSymbol{→}\AgdaSpace{}%
\AgdaSymbol{(}\AgdaBound{M}\AgdaSpace{}%
\AgdaSymbol{:}\AgdaSpace{}%
\AgdaBound{Γ}\AgdaSpace{}%
\AgdaOperator{\AgdaInductiveConstructor{,}}\AgdaSpace{}%
\AgdaBound{B}\AgdaSpace{}%
\AgdaOperator{\AgdaDatatype{⊢}}\AgdaSpace{}%
\AgdaBound{A}\AgdaSymbol{)}\AgdaSpace{}%
\AgdaSymbol{→}\AgdaSpace{}%
\AgdaSymbol{(}\AgdaBound{N}\AgdaSpace{}%
\AgdaSymbol{:}\AgdaSpace{}%
\AgdaBound{Γ}\AgdaSpace{}%
\AgdaOperator{\AgdaDatatype{⊢}}\AgdaSpace{}%
\AgdaBound{B}\AgdaSymbol{)}\AgdaSpace{}%
\AgdaSymbol{→}\AgdaSpace{}%
\AgdaSymbol{(}\AgdaBound{x}\AgdaSpace{}%
\AgdaSymbol{:}\AgdaSpace{}%
\AgdaDatatype{Value}\AgdaSpace{}%
\AgdaBound{M}\AgdaSymbol{)}\AgdaSpace{}%
\AgdaSymbol{→}\AgdaSpace{}%
\AgdaSymbol{(}\AgdaBound{y}\AgdaSpace{}%
\AgdaSymbol{:}\AgdaSpace{}%
\AgdaDatatype{Value}\AgdaSpace{}%
\AgdaBound{N}\AgdaSymbol{)}\AgdaSpace{}%
\AgdaSymbol{→}\AgdaSpace{}%
\AgdaDatatype{Value}\AgdaSpace{}%
\AgdaSymbol{(}\AgdaBound{M}\AgdaSpace{}%
\AgdaOperator{\AgdaFunction{[}}\AgdaSpace{}%
\AgdaBound{N}\AgdaSpace{}%
\AgdaOperator{\AgdaFunction{]}}\AgdaSymbol{)}\<%
\\
\>[0]\AgdaFunction{val-after-subst}\AgdaSpace{}%
\AgdaSymbol{(}\AgdaInductiveConstructor{`}\AgdaSpace{}%
\AgdaInductiveConstructor{Z}\AgdaSymbol{)}\AgdaSpace{}%
\AgdaBound{N}\AgdaSpace{}%
\AgdaBound{x}\AgdaSpace{}%
\AgdaBound{y}\AgdaSpace{}%
\AgdaSymbol{=}\AgdaSpace{}%
\AgdaBound{y}\<%
\\
\>[0]\AgdaFunction{val-after-subst}\AgdaSpace{}%
\AgdaSymbol{(}\AgdaInductiveConstructor{`}\AgdaSpace{}%
\AgdaSymbol{(}\AgdaInductiveConstructor{S}\AgdaSpace{}%
\AgdaBound{x₁}\AgdaSymbol{))}\AgdaSpace{}%
\AgdaBound{N}\AgdaSpace{}%
\AgdaBound{x}\AgdaSpace{}%
\AgdaBound{y}\AgdaSpace{}%
\AgdaSymbol{=}\AgdaSpace{}%
\AgdaInductiveConstructor{V-`}\<%
\\
\>[0]\AgdaFunction{val-after-subst}\AgdaSpace{}%
\AgdaSymbol{(}\AgdaInductiveConstructor{ƛ}\AgdaSpace{}%
\AgdaBound{M}\AgdaSymbol{)}\AgdaSpace{}%
\AgdaBound{N}\AgdaSpace{}%
\AgdaBound{x}\AgdaSpace{}%
\AgdaBound{y}\AgdaSpace{}%
\AgdaSymbol{=}\AgdaSpace{}%
\AgdaInductiveConstructor{V-ƛ}\<%
\\
%
\\[\AgdaEmptyExtraSkip]%
\>[0]\AgdaComment{--\ Preservation\ and\ Reflection\ of\ }\<%
\\
\>[0]\AgdaComment{--\ single\ substitution\ in\ values\ of\ Gödel's\ translation\ }\<%
\\
\>[0]\AgdaFunction{pres-subst-val}\AgdaSpace{}%
\AgdaSymbol{:}\AgdaSpace{}%
\AgdaSymbol{∀}\AgdaSpace{}%
\AgdaSymbol{\{}\AgdaBound{Γ}\AgdaSpace{}%
\AgdaBound{A}\AgdaSpace{}%
\AgdaBound{B}\AgdaSymbol{\}}\<%
\\
\>[0][@{}l@{\AgdaIndent{0}}]%
\>[2]\AgdaSymbol{→}\AgdaSpace{}%
\AgdaSymbol{(}\AgdaBound{M}\AgdaSpace{}%
\AgdaSymbol{:}\AgdaSpace{}%
\AgdaBound{Γ}\AgdaSpace{}%
\AgdaOperator{\AgdaInductiveConstructor{,}}\AgdaSpace{}%
\AgdaBound{B}\AgdaSpace{}%
\AgdaOperator{\AgdaDatatype{⊢}}\AgdaSpace{}%
\AgdaBound{A}\AgdaSymbol{)}\<%
\\
%
\>[2]\AgdaSymbol{→}\AgdaSpace{}%
\AgdaSymbol{(}\AgdaBound{V}\AgdaSpace{}%
\AgdaSymbol{:}\AgdaSpace{}%
\AgdaBound{Γ}\AgdaSpace{}%
\AgdaOperator{\AgdaDatatype{⊢}}\AgdaSpace{}%
\AgdaBound{B}\AgdaSymbol{)}\<%
\\
%
\>[2]\AgdaSymbol{→}\AgdaSpace{}%
\AgdaSymbol{(}\AgdaBound{x}\AgdaSpace{}%
\AgdaSymbol{:}\AgdaSpace{}%
\AgdaDatatype{Value}\AgdaSpace{}%
\AgdaBound{M}\AgdaSymbol{)}\<%
\\
%
\>[2]\AgdaSymbol{→}\AgdaSpace{}%
\AgdaSymbol{(}\AgdaBound{y}\AgdaSpace{}%
\AgdaSymbol{:}\AgdaSpace{}%
\AgdaDatatype{Value}\AgdaSpace{}%
\AgdaBound{V}\AgdaSymbol{)}\<%
\\
%
\>[2]\AgdaSymbol{→}\AgdaSpace{}%
\AgdaFunction{embedValue}\AgdaSpace{}%
\AgdaBound{M}\AgdaSpace{}%
\AgdaBound{x}\AgdaSpace{}%
\AgdaOperator{\AgdaFunction{[}}\AgdaSpace{}%
\AgdaFunction{embedValue}\AgdaSpace{}%
\AgdaBound{V}\AgdaSpace{}%
\AgdaBound{y}\AgdaSpace{}%
\AgdaOperator{\AgdaFunction{]b}}\AgdaSpace{}%
\AgdaOperator{\AgdaDatatype{≡}}\AgdaSpace{}%
\AgdaFunction{embedValue}\AgdaSpace{}%
\AgdaSymbol{(}\AgdaBound{M}\AgdaSpace{}%
\AgdaOperator{\AgdaFunction{[}}\AgdaSpace{}%
\AgdaBound{V}\AgdaSpace{}%
\AgdaOperator{\AgdaFunction{]}}\AgdaSymbol{)}\AgdaSpace{}%
\AgdaSymbol{(}\AgdaFunction{val-after-subst}\AgdaSpace{}%
\AgdaBound{M}\AgdaSpace{}%
\AgdaBound{V}\AgdaSpace{}%
\AgdaBound{x}\AgdaSpace{}%
\AgdaBound{y}\AgdaSymbol{)}\<%
\\
\>[0]\AgdaFunction{pres-subst-val}\AgdaSpace{}%
\AgdaSymbol{(}\AgdaInductiveConstructor{`}\AgdaSpace{}%
\AgdaInductiveConstructor{Z}\AgdaSymbol{)}\AgdaSpace{}%
\AgdaBound{V}\AgdaSpace{}%
\AgdaBound{x}\AgdaSpace{}%
\AgdaBound{y}\AgdaSpace{}%
\AgdaSymbol{=}\AgdaSpace{}%
\AgdaInductiveConstructor{refl}\<%
\\
\>[0]\AgdaFunction{pres-subst-val}\AgdaSpace{}%
\AgdaSymbol{(}\AgdaInductiveConstructor{`}\AgdaSpace{}%
\AgdaSymbol{(}\AgdaInductiveConstructor{S}\AgdaSpace{}%
\AgdaBound{x₁}\AgdaSymbol{))}\AgdaSpace{}%
\AgdaBound{V}\AgdaSpace{}%
\AgdaBound{x}\AgdaSpace{}%
\AgdaBound{y}\AgdaSpace{}%
\AgdaSymbol{=}\AgdaSpace{}%
\AgdaInductiveConstructor{refl}\<%
\\
\>[0]\AgdaFunction{pres-subst-val}\AgdaSpace{}%
\AgdaSymbol{(}\AgdaInductiveConstructor{ƛ}\AgdaSpace{}%
\AgdaBound{M}\AgdaSymbol{)}\AgdaSpace{}%
\AgdaBound{V}\AgdaSpace{}%
\AgdaBound{x}\AgdaSpace{}%
\AgdaBound{y}\AgdaSpace{}%
\AgdaSymbol{=}\<%
\\
\>[0][@{}l@{\AgdaIndent{0}}]%
\>[2]\AgdaOperator{\AgdaFunction{beginEq}}\<%
\\
\>[2][@{}l@{\AgdaIndent{0}}]%
\>[4]\AgdaFunction{embedValue}\AgdaSpace{}%
\AgdaSymbol{(}\AgdaInductiveConstructor{ƛ}\AgdaSpace{}%
\AgdaBound{M}\AgdaSymbol{)}\AgdaSpace{}%
\AgdaBound{x}\AgdaSpace{}%
\AgdaOperator{\AgdaFunction{[}}\AgdaSpace{}%
\AgdaFunction{embedValue}\AgdaSpace{}%
\AgdaBound{V}\AgdaSpace{}%
\AgdaBound{y}\AgdaSpace{}%
\AgdaOperator{\AgdaFunction{]b}}\<%
\\
%
\>[2]\AgdaFunction{≡⟨}\AgdaSpace{}%
\AgdaInductiveConstructor{refl}\AgdaSpace{}%
\AgdaFunction{⟩}\<%
\\
\>[2][@{}l@{\AgdaIndent{0}}]%
\>[4]\AgdaInductiveConstructor{ƛb}\AgdaSpace{}%
\AgdaSymbol{(}\AgdaFunction{substb}\AgdaSpace{}%
\AgdaSymbol{(}\AgdaFunction{extsb}\AgdaSpace{}%
\AgdaSymbol{(}\AgdaFunction{fb}\AgdaSpace{}%
\AgdaSymbol{(}\AgdaFunction{embedValue}\AgdaSpace{}%
\AgdaBound{V}\AgdaSpace{}%
\AgdaBound{y}\AgdaSymbol{)))}\AgdaSpace{}%
\AgdaSymbol{(}\AgdaFunction{embedTerm}\AgdaSpace{}%
\AgdaBound{M}\AgdaSymbol{))}\<%
\\
%
\>[2]\AgdaFunction{≡⟨}\AgdaSpace{}%
\AgdaFunction{cong}\AgdaSpace{}%
\AgdaInductiveConstructor{ƛb}\AgdaSpace{}%
\AgdaSymbol{(}\AgdaFunction{pres-simultaneous-subst}\AgdaSpace{}%
\AgdaSymbol{(}\AgdaFunction{extsb}\AgdaSpace{}%
\AgdaSymbol{(}\AgdaFunction{fb}\AgdaSpace{}%
\AgdaSymbol{(}\AgdaFunction{embedValue}\AgdaSpace{}%
\AgdaBound{V}\AgdaSpace{}%
\AgdaBound{y}\AgdaSymbol{)))}\AgdaSpace{}%
\AgdaSymbol{(}\AgdaFunction{exts}\AgdaSpace{}%
\AgdaSymbol{(}\AgdaFunction{f}\AgdaSpace{}%
\AgdaBound{V}\AgdaSymbol{))}\AgdaSpace{}%
\AgdaSymbol{(}\AgdaFunction{exts-pres-subst-func}\AgdaSpace{}%
\AgdaBound{V}\AgdaSpace{}%
\AgdaBound{y}\AgdaSymbol{)}\AgdaSpace{}%
\AgdaBound{M}\AgdaSymbol{)}\AgdaSpace{}%
\AgdaFunction{⟩}\<%
\\
\>[2][@{}l@{\AgdaIndent{0}}]%
\>[4]\AgdaInductiveConstructor{ƛb}\AgdaSpace{}%
\AgdaSymbol{(}\AgdaFunction{embedTerm}\AgdaSpace{}%
\AgdaSymbol{(}\AgdaFunction{subst}\AgdaSpace{}%
\AgdaSymbol{(}\AgdaFunction{exts}\AgdaSpace{}%
\AgdaSymbol{(}\AgdaFunction{f}\AgdaSpace{}%
\AgdaBound{V}\AgdaSymbol{))}\AgdaSpace{}%
\AgdaBound{M}\AgdaSymbol{))}\<%
\\
%
\>[2]\AgdaFunction{≡⟨}\AgdaSpace{}%
\AgdaInductiveConstructor{refl}\AgdaSpace{}%
\AgdaFunction{⟩}\<%
\\
\>[2][@{}l@{\AgdaIndent{0}}]%
\>[4]\AgdaFunction{embedValue}\AgdaSpace{}%
\AgdaSymbol{((}\AgdaInductiveConstructor{ƛ}\AgdaSpace{}%
\AgdaBound{M}\AgdaSymbol{)}\AgdaSpace{}%
\AgdaOperator{\AgdaFunction{[}}\AgdaSpace{}%
\AgdaBound{V}\AgdaSpace{}%
\AgdaOperator{\AgdaFunction{]}}\AgdaSymbol{)}\AgdaSpace{}%
\AgdaSymbol{(}\AgdaFunction{val-after-subst}\AgdaSpace{}%
\AgdaSymbol{(}\AgdaInductiveConstructor{ƛ}\AgdaSpace{}%
\AgdaBound{M}\AgdaSymbol{)}\AgdaSpace{}%
\AgdaBound{V}\AgdaSpace{}%
\AgdaBound{x}\AgdaSpace{}%
\AgdaBound{y}\AgdaSymbol{)}\<%
\\
%
\>[2]\AgdaOperator{\AgdaFunction{∎eq}}\<%
\\
%
\\[\AgdaEmptyExtraSkip]%
\>[0]\AgdaComment{--\ fb\ and\ renameb\ commute\ under\ extended\ context.}\<%
\\
\>[0]\AgdaFunction{fb-renameb-extb-commute}\AgdaSpace{}%
\AgdaSymbol{:}\AgdaSpace{}%
\AgdaSymbol{∀}\AgdaSpace{}%
\AgdaSymbol{\{}\AgdaBound{Γ}\AgdaSpace{}%
\AgdaBound{A}\AgdaSpace{}%
\AgdaBound{B}\AgdaSpace{}%
\AgdaBound{Δ}\AgdaSymbol{\}}\<%
\\
\>[0][@{}l@{\AgdaIndent{0}}]%
\>[2]\AgdaSymbol{→}\AgdaSpace{}%
\AgdaSymbol{(}\AgdaBound{fb'}\AgdaSpace{}%
\AgdaSymbol{:}\AgdaSpace{}%
\AgdaSymbol{∀}\AgdaSpace{}%
\AgdaSymbol{\{}\AgdaBound{Ab}\AgdaSymbol{\}}\AgdaSpace{}%
\AgdaSymbol{→}\AgdaSpace{}%
\AgdaBound{Ab}\AgdaSpace{}%
\AgdaOperator{\AgdaDatatype{∈b}}\AgdaSpace{}%
\AgdaBound{Γ}\AgdaSpace{}%
\AgdaSymbol{→}\AgdaSpace{}%
\AgdaBound{Ab}\AgdaSpace{}%
\AgdaOperator{\AgdaDatatype{∈b}}\AgdaSpace{}%
\AgdaBound{Δ}\AgdaSymbol{)}\<%
\\
%
\>[2]\AgdaSymbol{→}\AgdaSpace{}%
\AgdaSymbol{(}\AgdaBound{W}\AgdaSpace{}%
\AgdaSymbol{:}\AgdaSpace{}%
\AgdaBound{Γ}\AgdaSpace{}%
\AgdaOperator{\AgdaDatatype{⊢b}}\AgdaSpace{}%
\AgdaBound{B}\AgdaSymbol{)}\<%
\\
%
\>[2]\AgdaSymbol{→}\AgdaSpace{}%
\AgdaSymbol{(}\AgdaBound{x}\AgdaSpace{}%
\AgdaSymbol{:}\AgdaSpace{}%
\AgdaInductiveConstructor{□}\AgdaSpace{}%
\AgdaBound{A}\AgdaSpace{}%
\AgdaOperator{\AgdaDatatype{∈b}}\AgdaSpace{}%
\AgdaBound{Γ}\AgdaSpace{}%
\AgdaOperator{\AgdaInductiveConstructor{,b}}\AgdaSpace{}%
\AgdaInductiveConstructor{□}\AgdaSpace{}%
\AgdaBound{B}\AgdaSymbol{)}\<%
\\
%
\>[2]\AgdaSymbol{→}\AgdaSpace{}%
\AgdaFunction{fb}\AgdaSpace{}%
\AgdaSymbol{(}\AgdaFunction{renameb}\AgdaSpace{}%
\AgdaBound{fb'}\AgdaSpace{}%
\AgdaBound{W}\AgdaSymbol{)}\AgdaSpace{}%
\AgdaSymbol{(}\AgdaFunction{extb}\AgdaSpace{}%
\AgdaBound{fb'}\AgdaSpace{}%
\AgdaBound{x}\AgdaSymbol{)}\AgdaSpace{}%
\AgdaOperator{\AgdaDatatype{≡}}\AgdaSpace{}%
\AgdaFunction{renameb}\AgdaSpace{}%
\AgdaBound{fb'}\AgdaSpace{}%
\AgdaSymbol{(}\AgdaFunction{fb}\AgdaSpace{}%
\AgdaBound{W}\AgdaSpace{}%
\AgdaBound{x}\AgdaSymbol{)}\<%
\\
\>[0]\AgdaFunction{fb-renameb-extb-commute}\AgdaSpace{}%
\AgdaBound{fb'}\AgdaSpace{}%
\AgdaBound{W}\AgdaSpace{}%
\AgdaInductiveConstructor{Zb}\AgdaSpace{}%
\AgdaSymbol{=}\AgdaSpace{}%
\AgdaInductiveConstructor{refl}\<%
\\
\>[0]\AgdaFunction{fb-renameb-extb-commute}\AgdaSpace{}%
\AgdaBound{fb'}\AgdaSpace{}%
\AgdaBound{W}\AgdaSpace{}%
\AgdaSymbol{(}\AgdaInductiveConstructor{Sb}\AgdaSpace{}%
\AgdaBound{x}\AgdaSymbol{)}\AgdaSpace{}%
\AgdaSymbol{=}\AgdaSpace{}%
\AgdaInductiveConstructor{refl}\<%
\\
%
\\[\AgdaEmptyExtraSkip]%
\>[0]\AgdaComment{--\ Preservation\ and\ reflection\ of\ }\<%
\\
\>[0]\AgdaComment{--\ single\ substitution\ under\ renaming}\<%
\\
\>[0]\AgdaFunction{pres-subst-func-rename}\AgdaSpace{}%
\AgdaSymbol{:}\AgdaSpace{}%
\AgdaSymbol{∀}\AgdaSpace{}%
\AgdaSymbol{\{}\AgdaBound{Γ}\AgdaSpace{}%
\AgdaBound{B}\AgdaSpace{}%
\AgdaBound{A}\AgdaSpace{}%
\AgdaBound{Δ}\AgdaSymbol{\}}\<%
\\
\>[0][@{}l@{\AgdaIndent{0}}]%
\>[2]\AgdaSymbol{→}\AgdaSpace{}%
\AgdaSymbol{(}\AgdaBound{fb'}\AgdaSpace{}%
\AgdaSymbol{:}\AgdaSpace{}%
\AgdaSymbol{∀}\AgdaSpace{}%
\AgdaSymbol{\{}\AgdaBound{Ab}\AgdaSymbol{\}}\AgdaSpace{}%
\AgdaSymbol{→}\AgdaSpace{}%
\AgdaBound{Ab}\AgdaSpace{}%
\AgdaOperator{\AgdaDatatype{∈b}}\AgdaSpace{}%
\AgdaBound{Γ}\AgdaSpace{}%
\AgdaSymbol{→}\AgdaSpace{}%
\AgdaBound{Ab}\AgdaSpace{}%
\AgdaOperator{\AgdaDatatype{∈b}}\AgdaSpace{}%
\AgdaBound{Δ}\AgdaSymbol{)}\<%
\\
%
\>[2]\AgdaSymbol{→}\AgdaSpace{}%
\AgdaSymbol{(}\AgdaBound{M}\AgdaSpace{}%
\AgdaSymbol{:}\AgdaSpace{}%
\AgdaBound{Γ}\AgdaSpace{}%
\AgdaOperator{\AgdaInductiveConstructor{,b}}\AgdaSpace{}%
\AgdaInductiveConstructor{□}\AgdaSpace{}%
\AgdaBound{B}\AgdaSpace{}%
\AgdaOperator{\AgdaDatatype{⊢b}}\AgdaSpace{}%
\AgdaBound{A}\AgdaSymbol{)}\<%
\\
%
\>[2]\AgdaSymbol{→}\AgdaSpace{}%
\AgdaSymbol{(}\AgdaBound{W}\AgdaSpace{}%
\AgdaSymbol{:}\AgdaSpace{}%
\AgdaBound{Γ}\AgdaSpace{}%
\AgdaOperator{\AgdaDatatype{⊢b}}\AgdaSpace{}%
\AgdaBound{B}\AgdaSymbol{)}\<%
\\
%
\>[2]\AgdaSymbol{→}%
\>[2253I]\AgdaFunction{renameb}\AgdaSpace{}%
\AgdaSymbol{(}\AgdaFunction{extb}\AgdaSpace{}%
\AgdaBound{fb'}\AgdaSymbol{)}\AgdaSpace{}%
\AgdaBound{M}\AgdaSpace{}%
\AgdaOperator{\AgdaFunction{[}}\AgdaSpace{}%
\AgdaFunction{renameb}\AgdaSpace{}%
\AgdaBound{fb'}\AgdaSpace{}%
\AgdaBound{W}\AgdaSpace{}%
\AgdaOperator{\AgdaFunction{]b}}\<%
\\
\>[.][@{}l@{}]\<[2253I]%
\>[4]\AgdaOperator{\AgdaDatatype{≡}}\AgdaSpace{}%
\AgdaFunction{renameb}\AgdaSpace{}%
\AgdaBound{fb'}\AgdaSpace{}%
\AgdaSymbol{(}\AgdaBound{M}\AgdaSpace{}%
\AgdaOperator{\AgdaFunction{[}}\AgdaSpace{}%
\AgdaBound{W}\AgdaSpace{}%
\AgdaOperator{\AgdaFunction{]b}}\AgdaSymbol{)}\<%
\\
\>[0]\AgdaFunction{pres-subst-func-rename}\AgdaSpace{}%
\AgdaBound{fb'}\AgdaSpace{}%
\AgdaBound{M}\AgdaSpace{}%
\AgdaBound{W}\AgdaSpace{}%
\AgdaSymbol{=}\<%
\\
\>[0][@{}l@{\AgdaIndent{0}}]%
\>[2]\AgdaOperator{\AgdaFunction{beginEq}}\<%
\\
\>[2][@{}l@{\AgdaIndent{0}}]%
\>[4]\AgdaFunction{substb}\AgdaSpace{}%
\AgdaSymbol{(}\AgdaFunction{fb}\AgdaSpace{}%
\AgdaSymbol{(}\AgdaFunction{renameb}\AgdaSpace{}%
\AgdaBound{fb'}\AgdaSpace{}%
\AgdaBound{W}\AgdaSymbol{))}\AgdaSpace{}%
\AgdaSymbol{(}\AgdaFunction{renameb}\AgdaSpace{}%
\AgdaSymbol{(}\AgdaFunction{extb}\AgdaSpace{}%
\AgdaBound{fb'}\AgdaSymbol{)}\AgdaSpace{}%
\AgdaBound{M}\AgdaSymbol{)}\<%
\\
%
\>[2]\AgdaFunction{≡⟨}\AgdaSpace{}%
\AgdaFunction{subst-renameb-comp}\AgdaSpace{}%
\AgdaSymbol{(}\AgdaFunction{extb}\AgdaSpace{}%
\AgdaBound{fb'}\AgdaSymbol{)}\AgdaSpace{}%
\AgdaSymbol{(}\AgdaFunction{fb}\AgdaSpace{}%
\AgdaSymbol{(}\AgdaFunction{renameb}\AgdaSpace{}%
\AgdaBound{fb'}\AgdaSpace{}%
\AgdaBound{W}\AgdaSymbol{))}\AgdaSpace{}%
\AgdaBound{M}\AgdaSpace{}%
\AgdaFunction{⟩}\<%
\\
\>[2][@{}l@{\AgdaIndent{0}}]%
\>[4]\AgdaFunction{substb}\AgdaSpace{}%
\AgdaSymbol{(λ}\AgdaSpace{}%
\AgdaBound{x}\AgdaSpace{}%
\AgdaSymbol{→}\AgdaSpace{}%
\AgdaFunction{fb}\AgdaSpace{}%
\AgdaSymbol{(}\AgdaFunction{renameb}\AgdaSpace{}%
\AgdaBound{fb'}\AgdaSpace{}%
\AgdaBound{W}\AgdaSymbol{)}\AgdaSpace{}%
\AgdaSymbol{(}\AgdaFunction{extb}\AgdaSpace{}%
\AgdaBound{fb'}\AgdaSpace{}%
\AgdaBound{x}\AgdaSymbol{))}\AgdaSpace{}%
\AgdaBound{M}\<%
\\
%
\>[2]\AgdaFunction{≡⟨}\AgdaSpace{}%
\AgdaFunction{substb-cong}\AgdaSpace{}%
\AgdaSymbol{(}\AgdaFunction{fb-renameb-extb-commute}\AgdaSpace{}%
\AgdaBound{fb'}\AgdaSpace{}%
\AgdaBound{W}\AgdaSymbol{)}\AgdaSpace{}%
\AgdaBound{M}\AgdaSpace{}%
\AgdaFunction{⟩}\<%
\\
\>[2][@{}l@{\AgdaIndent{0}}]%
\>[4]\AgdaFunction{substb}\AgdaSpace{}%
\AgdaSymbol{(λ}\AgdaSpace{}%
\AgdaBound{x}\AgdaSpace{}%
\AgdaSymbol{→}\AgdaSpace{}%
\AgdaFunction{renameb}\AgdaSpace{}%
\AgdaBound{fb'}\AgdaSpace{}%
\AgdaSymbol{(}\AgdaFunction{fb}\AgdaSpace{}%
\AgdaBound{W}\AgdaSpace{}%
\AgdaBound{x}\AgdaSymbol{))}\AgdaSpace{}%
\AgdaBound{M}\<%
\\
%
\>[2]\AgdaFunction{≡⟨}\AgdaSpace{}%
\AgdaFunction{substb-rename-commute}\AgdaSpace{}%
\AgdaBound{fb'}\AgdaSpace{}%
\AgdaSymbol{(}\AgdaFunction{fb}\AgdaSpace{}%
\AgdaBound{W}\AgdaSymbol{)}\AgdaSpace{}%
\AgdaBound{M}\AgdaSpace{}%
\AgdaFunction{⟩}\<%
\\
\>[2][@{}l@{\AgdaIndent{0}}]%
\>[4]\AgdaFunction{renameb}\AgdaSpace{}%
\AgdaBound{fb'}\AgdaSpace{}%
\AgdaSymbol{(}\AgdaFunction{substb}\AgdaSpace{}%
\AgdaSymbol{(}\AgdaFunction{fb}\AgdaSpace{}%
\AgdaBound{W}\AgdaSymbol{)}\AgdaSpace{}%
\AgdaBound{M}\AgdaSymbol{)}\<%
\\
%
\>[2]\AgdaOperator{\AgdaFunction{∎eq}}\<%
\\
%
\\[\AgdaEmptyExtraSkip]%
\>[0]\AgdaComment{--\ Preservation\ of\ reduction\ under\ renaming\ \ \ }\<%
\\
\>[0]\AgdaFunction{pres-rename-red}\AgdaSpace{}%
\AgdaSymbol{:}\AgdaSpace{}%
\AgdaSymbol{∀}\AgdaSpace{}%
\AgdaSymbol{\{}\AgdaBound{Γ}\AgdaSpace{}%
\AgdaBound{A}\AgdaSpace{}%
\AgdaBound{Δ}\AgdaSymbol{\}}\<%
\\
\>[0][@{}l@{\AgdaIndent{0}}]%
\>[2]\AgdaSymbol{→}\AgdaSpace{}%
\AgdaSymbol{(}\AgdaBound{fb'}\AgdaSpace{}%
\AgdaSymbol{:}\AgdaSpace{}%
\AgdaSymbol{∀}\AgdaSpace{}%
\AgdaSymbol{\{}\AgdaBound{Ab}\AgdaSymbol{\}}\AgdaSpace{}%
\AgdaSymbol{→}\AgdaSpace{}%
\AgdaBound{Ab}\AgdaSpace{}%
\AgdaOperator{\AgdaDatatype{∈b}}\AgdaSpace{}%
\AgdaBound{Γ}\AgdaSpace{}%
\AgdaSymbol{→}\AgdaSpace{}%
\AgdaBound{Ab}\AgdaSpace{}%
\AgdaOperator{\AgdaDatatype{∈b}}\AgdaSpace{}%
\AgdaBound{Δ}\AgdaSymbol{)}\<%
\\
%
\>[2]\AgdaSymbol{→}\AgdaSpace{}%
\AgdaSymbol{(}\AgdaBound{M}\AgdaSpace{}%
\AgdaSymbol{:}\AgdaSpace{}%
\AgdaBound{Γ}\AgdaSpace{}%
\AgdaOperator{\AgdaDatatype{⊢b}}\AgdaSpace{}%
\AgdaBound{A}\AgdaSymbol{)}\<%
\\
%
\>[2]\AgdaSymbol{→}\AgdaSpace{}%
\AgdaSymbol{(}\AgdaBound{N}\AgdaSpace{}%
\AgdaSymbol{:}\AgdaSpace{}%
\AgdaBound{Γ}\AgdaSpace{}%
\AgdaOperator{\AgdaDatatype{⊢b}}\AgdaSpace{}%
\AgdaBound{A}\AgdaSymbol{)}\<%
\\
%
\>[2]\AgdaSymbol{→}\AgdaSpace{}%
\AgdaSymbol{(}\AgdaBound{M}\AgdaSpace{}%
\AgdaOperator{\AgdaDatatype{↝βb}}\AgdaSpace{}%
\AgdaBound{N}\AgdaSymbol{)}\<%
\\
%
\>[2]\AgdaSymbol{→}\AgdaSpace{}%
\AgdaFunction{renameb}\AgdaSpace{}%
\AgdaBound{fb'}\AgdaSpace{}%
\AgdaBound{M}\AgdaSpace{}%
\AgdaOperator{\AgdaDatatype{↝βb}}\AgdaSpace{}%
\AgdaFunction{renameb}\AgdaSpace{}%
\AgdaBound{fb'}\AgdaSpace{}%
\AgdaBound{N}\<%
\\
\>[0]\AgdaFunction{pres-rename-red}\AgdaSpace{}%
\AgdaBound{fb'}\AgdaSpace{}%
\AgdaSymbol{(}\AgdaInductiveConstructor{ƛb}\AgdaSpace{}%
\AgdaBound{M₁}\AgdaSpace{}%
\AgdaOperator{\AgdaInductiveConstructor{·b}}\AgdaSpace{}%
\AgdaInductiveConstructor{box}\AgdaSpace{}%
\AgdaBound{M₂}\AgdaSymbol{)}\AgdaSpace{}%
\AgdaBound{N}\AgdaSpace{}%
\AgdaInductiveConstructor{βb}\AgdaSpace{}%
\AgdaSymbol{=}\AgdaSpace{}%
\AgdaFunction{Eq.subst}\AgdaSpace{}%
\AgdaSymbol{(λ}\AgdaSpace{}%
\AgdaBound{x}\AgdaSpace{}%
\AgdaSymbol{→}\AgdaSpace{}%
\AgdaInductiveConstructor{ƛb}\AgdaSpace{}%
\AgdaSymbol{(}\AgdaFunction{renameb}\AgdaSpace{}%
\AgdaSymbol{(}\AgdaFunction{extb}\AgdaSpace{}%
\AgdaBound{fb'}\AgdaSymbol{)}\AgdaSpace{}%
\AgdaBound{M₁}\AgdaSymbol{)}\AgdaSpace{}%
\AgdaOperator{\AgdaInductiveConstructor{·b}}\AgdaSpace{}%
\AgdaInductiveConstructor{box}\AgdaSpace{}%
\AgdaSymbol{(}\AgdaFunction{renameb}\AgdaSpace{}%
\AgdaBound{fb'}\AgdaSpace{}%
\AgdaBound{M₂}\AgdaSymbol{)}\AgdaSpace{}%
\AgdaOperator{\AgdaDatatype{↝βb}}\AgdaSpace{}%
\AgdaBound{x}\AgdaSymbol{)}\AgdaSpace{}%
\AgdaSymbol{(}\AgdaFunction{pres-subst-func-rename}\AgdaSpace{}%
\AgdaBound{fb'}\AgdaSpace{}%
\AgdaBound{M₁}\AgdaSpace{}%
\AgdaBound{M₂}\AgdaSymbol{)}\AgdaSpace{}%
\AgdaInductiveConstructor{βb}\<%
\\
\>[0]\AgdaFunction{pres-rename-red}\AgdaSpace{}%
\AgdaBound{fb'}\AgdaSpace{}%
\AgdaBound{M}\AgdaSpace{}%
\AgdaBound{N}\AgdaSpace{}%
\AgdaSymbol{(}\AgdaInductiveConstructor{μ}\AgdaSpace{}%
\AgdaSymbol{\{}\AgdaArgument{M}\AgdaSpace{}%
\AgdaSymbol{=}\AgdaSpace{}%
\AgdaBound{M₁}\AgdaSymbol{\}}\AgdaSpace{}%
\AgdaSymbol{\{}\AgdaArgument{M'}\AgdaSpace{}%
\AgdaSymbol{=}\AgdaSpace{}%
\AgdaBound{M'}\AgdaSymbol{\}}\AgdaSpace{}%
\AgdaBound{red}\AgdaSymbol{)}\AgdaSpace{}%
\AgdaSymbol{=}\AgdaSpace{}%
\AgdaInductiveConstructor{μ}\AgdaSpace{}%
\AgdaSymbol{(}\AgdaFunction{pres-rename-red}\AgdaSpace{}%
\AgdaBound{fb'}\AgdaSpace{}%
\AgdaBound{M₁}\AgdaSpace{}%
\AgdaBound{M'}\AgdaSpace{}%
\AgdaBound{red}\AgdaSymbol{)}\<%
\\
\>[0]\AgdaFunction{pres-rename-red}\AgdaSpace{}%
\AgdaBound{fb'}\AgdaSpace{}%
\AgdaBound{M}\AgdaSpace{}%
\AgdaBound{N}\AgdaSpace{}%
\AgdaSymbol{(}\AgdaInductiveConstructor{ν}\AgdaSpace{}%
\AgdaSymbol{\{}\AgdaArgument{M}\AgdaSpace{}%
\AgdaSymbol{=}\AgdaSpace{}%
\AgdaBound{M₁}\AgdaSymbol{\}}\AgdaSpace{}%
\AgdaSymbol{\{}\AgdaArgument{N}\AgdaSpace{}%
\AgdaSymbol{=}\AgdaSpace{}%
\AgdaBound{N₁}\AgdaSymbol{\}}\AgdaSpace{}%
\AgdaSymbol{\{}\AgdaArgument{N'}\AgdaSpace{}%
\AgdaSymbol{=}\AgdaSpace{}%
\AgdaBound{N'}\AgdaSymbol{\}}\AgdaSpace{}%
\AgdaBound{red}\AgdaSymbol{)}\AgdaSpace{}%
\AgdaSymbol{=}\AgdaSpace{}%
\AgdaInductiveConstructor{ν}\AgdaSpace{}%
\AgdaSymbol{(}\AgdaFunction{pres-rename-red}\AgdaSpace{}%
\AgdaBound{fb'}\AgdaSpace{}%
\AgdaBound{N₁}\AgdaSpace{}%
\AgdaBound{N'}\AgdaSpace{}%
\AgdaBound{red}\AgdaSymbol{)}\<%
\\
\>[0]\AgdaFunction{pres-rename-red}\AgdaSpace{}%
\AgdaBound{fb'}\AgdaSpace{}%
\AgdaBound{M}\AgdaSpace{}%
\AgdaBound{N}\AgdaSpace{}%
\AgdaSymbol{(}\AgdaInductiveConstructor{ξ}\AgdaSpace{}%
\AgdaSymbol{\{}\AgdaArgument{M}\AgdaSpace{}%
\AgdaSymbol{=}\AgdaSpace{}%
\AgdaBound{M₁}\AgdaSymbol{\}}\AgdaSpace{}%
\AgdaSymbol{\{}\AgdaArgument{M'}\AgdaSpace{}%
\AgdaSymbol{=}\AgdaSpace{}%
\AgdaBound{M'}\AgdaSymbol{\}}\AgdaSpace{}%
\AgdaBound{red}\AgdaSymbol{)}\AgdaSpace{}%
\AgdaSymbol{=}\AgdaSpace{}%
\AgdaInductiveConstructor{ξ}\AgdaSpace{}%
\AgdaSymbol{(}\AgdaFunction{pres-rename-red}\AgdaSpace{}%
\AgdaSymbol{(}\AgdaFunction{extb}\AgdaSpace{}%
\AgdaBound{fb'}\AgdaSymbol{)}\AgdaSpace{}%
\AgdaBound{M₁}\AgdaSpace{}%
\AgdaBound{M'}\AgdaSpace{}%
\AgdaBound{red}\AgdaSymbol{)}\<%
\\
\>[0]\AgdaFunction{pres-rename-red}\AgdaSpace{}%
\AgdaBound{fb'}\AgdaSpace{}%
\AgdaBound{M}\AgdaSpace{}%
\AgdaBound{N}\AgdaSpace{}%
\AgdaSymbol{(}\AgdaInductiveConstructor{ζ}\AgdaSpace{}%
\AgdaSymbol{\{}\AgdaArgument{M}\AgdaSpace{}%
\AgdaSymbol{=}\AgdaSpace{}%
\AgdaBound{M₁}\AgdaSymbol{\}}\AgdaSpace{}%
\AgdaSymbol{\{}\AgdaArgument{M'}\AgdaSpace{}%
\AgdaSymbol{=}\AgdaSpace{}%
\AgdaBound{M'}\AgdaSymbol{\}}\AgdaSpace{}%
\AgdaBound{red}\AgdaSymbol{)}\AgdaSpace{}%
\AgdaSymbol{=}\AgdaSpace{}%
\AgdaInductiveConstructor{ζ}\AgdaSpace{}%
\AgdaSymbol{(}\AgdaFunction{pres-rename-red}\AgdaSpace{}%
\AgdaBound{fb'}\AgdaSpace{}%
\AgdaBound{M₁}\AgdaSpace{}%
\AgdaBound{M'}\AgdaSpace{}%
\AgdaBound{red}\AgdaSymbol{)}\<%
\\
%
\\[\AgdaEmptyExtraSkip]%
\>[0]\AgdaComment{--\ Preservation\ of\ 2\ step\ reduction}\<%
\\
\>[0]\AgdaComment{--\ under\ renaming}\<%
\\
\>[0]\AgdaFunction{pres-rename-red-2}\AgdaSpace{}%
\AgdaSymbol{:}\AgdaSpace{}%
\AgdaSymbol{∀}\AgdaSpace{}%
\AgdaSymbol{\{}\AgdaBound{Γ}\AgdaSpace{}%
\AgdaBound{A}\AgdaSpace{}%
\AgdaBound{Δ}\AgdaSymbol{\}}\<%
\\
\>[0][@{}l@{\AgdaIndent{0}}]%
\>[2]\AgdaSymbol{→}\AgdaSpace{}%
\AgdaSymbol{(}\AgdaBound{fb'}\AgdaSpace{}%
\AgdaSymbol{:}\AgdaSpace{}%
\AgdaSymbol{∀}\AgdaSpace{}%
\AgdaSymbol{\{}\AgdaBound{Ab}\AgdaSymbol{\}}\AgdaSpace{}%
\AgdaSymbol{→}\AgdaSpace{}%
\AgdaBound{Ab}\AgdaSpace{}%
\AgdaOperator{\AgdaDatatype{∈b}}\AgdaSpace{}%
\AgdaBound{Γ}\AgdaSpace{}%
\AgdaSymbol{→}\AgdaSpace{}%
\AgdaBound{Ab}\AgdaSpace{}%
\AgdaOperator{\AgdaDatatype{∈b}}\AgdaSpace{}%
\AgdaBound{Δ}\AgdaSymbol{)}\<%
\\
%
\>[2]\AgdaSymbol{→}\AgdaSpace{}%
\AgdaSymbol{(}\AgdaBound{M}\AgdaSpace{}%
\AgdaSymbol{:}\AgdaSpace{}%
\AgdaBound{Γ}\AgdaSpace{}%
\AgdaOperator{\AgdaDatatype{⊢b}}\AgdaSpace{}%
\AgdaBound{A}\AgdaSymbol{)}\<%
\\
%
\>[2]\AgdaSymbol{→}\AgdaSpace{}%
\AgdaSymbol{(}\AgdaBound{N}\AgdaSpace{}%
\AgdaSymbol{:}\AgdaSpace{}%
\AgdaBound{Γ}\AgdaSpace{}%
\AgdaOperator{\AgdaDatatype{⊢b}}\AgdaSpace{}%
\AgdaBound{A}\AgdaSymbol{)}\<%
\\
%
\>[2]\AgdaSymbol{→}\AgdaSpace{}%
\AgdaSymbol{(}\AgdaBound{M}\AgdaSpace{}%
\AgdaOperator{\AgdaDatatype{↝βb²}}\AgdaSpace{}%
\AgdaBound{N}\AgdaSymbol{)}\<%
\\
%
\>[2]\AgdaSymbol{→}\AgdaSpace{}%
\AgdaFunction{renameb}\AgdaSpace{}%
\AgdaBound{fb'}\AgdaSpace{}%
\AgdaBound{M}\AgdaSpace{}%
\AgdaOperator{\AgdaDatatype{↝βb²}}\AgdaSpace{}%
\AgdaFunction{renameb}\AgdaSpace{}%
\AgdaBound{fb'}\AgdaSpace{}%
\AgdaBound{N}\<%
\\
\>[0]\AgdaFunction{pres-rename-red-2}\AgdaSpace{}%
\AgdaBound{fb'}\AgdaSpace{}%
\AgdaBound{M}\AgdaSpace{}%
\AgdaBound{N}\AgdaSpace{}%
\AgdaSymbol{(}\AgdaInductiveConstructor{βb²}\AgdaSpace{}%
\AgdaDottedPattern{\AgdaSymbol{.}}\AgdaDottedPattern{\AgdaBound{M}}\AgdaSpace{}%
\AgdaBound{t₂}\AgdaSpace{}%
\AgdaDottedPattern{\AgdaSymbol{.}}\AgdaDottedPattern{\AgdaBound{N}}\AgdaSpace{}%
\AgdaBound{x}\AgdaSpace{}%
\AgdaBound{x₁}\AgdaSymbol{)}\AgdaSpace{}%
\AgdaSymbol{=}\AgdaSpace{}%
\AgdaInductiveConstructor{βb²}\AgdaSpace{}%
\AgdaSymbol{(}\AgdaFunction{renameb}\AgdaSpace{}%
\AgdaBound{fb'}\AgdaSpace{}%
\AgdaBound{M}\AgdaSymbol{)}\AgdaSpace{}%
\AgdaSymbol{(}\AgdaFunction{renameb}\AgdaSpace{}%
\AgdaBound{fb'}\AgdaSpace{}%
\AgdaBound{t₂}\AgdaSymbol{)}\AgdaSpace{}%
\AgdaSymbol{(}\AgdaFunction{renameb}\AgdaSpace{}%
\AgdaBound{fb'}\AgdaSpace{}%
\AgdaBound{N}\AgdaSymbol{)}\AgdaSpace{}%
\AgdaSymbol{(}\AgdaFunction{pres-rename-red}\AgdaSpace{}%
\AgdaBound{fb'}\AgdaSpace{}%
\AgdaBound{M}\AgdaSpace{}%
\AgdaBound{t₂}\AgdaSpace{}%
\AgdaBound{x}\AgdaSymbol{)}\AgdaSpace{}%
\AgdaSymbol{(}\AgdaFunction{pres-rename-red}\AgdaSpace{}%
\AgdaBound{fb'}\AgdaSpace{}%
\AgdaBound{t₂}\AgdaSpace{}%
\AgdaBound{N}\AgdaSpace{}%
\AgdaBound{x₁}\AgdaSymbol{)}\<%
\\
\>[0]\AgdaFunction{pres-rename-red-2}\AgdaSpace{}%
\AgdaBound{fb'}\AgdaSpace{}%
\AgdaBound{M}\AgdaSpace{}%
\AgdaBound{N}\AgdaSpace{}%
\AgdaSymbol{(}\AgdaInductiveConstructor{μ}\AgdaSpace{}%
\AgdaSymbol{\{}\AgdaArgument{M}\AgdaSpace{}%
\AgdaSymbol{=}\AgdaSpace{}%
\AgdaBound{M₁}\AgdaSymbol{\}}\AgdaSpace{}%
\AgdaSymbol{\{}\AgdaArgument{M'}\AgdaSpace{}%
\AgdaSymbol{=}\AgdaSpace{}%
\AgdaBound{M'}\AgdaSymbol{\}}\AgdaSpace{}%
\AgdaBound{red}\AgdaSymbol{)}\AgdaSpace{}%
\AgdaSymbol{=}\AgdaSpace{}%
\AgdaInductiveConstructor{μ}\AgdaSpace{}%
\AgdaSymbol{(}\AgdaFunction{pres-rename-red-2}\AgdaSpace{}%
\AgdaBound{fb'}\AgdaSpace{}%
\AgdaBound{M₁}\AgdaSpace{}%
\AgdaBound{M'}\AgdaSpace{}%
\AgdaBound{red}\AgdaSymbol{)}\<%
\\
\>[0]\AgdaFunction{pres-rename-red-2}\AgdaSpace{}%
\AgdaBound{fb'}\AgdaSpace{}%
\AgdaBound{M}\AgdaSpace{}%
\AgdaBound{N}\AgdaSpace{}%
\AgdaSymbol{(}\AgdaInductiveConstructor{ν}\AgdaSpace{}%
\AgdaSymbol{\{}\AgdaArgument{M}\AgdaSpace{}%
\AgdaSymbol{=}\AgdaSpace{}%
\AgdaBound{M₁}\AgdaSymbol{\}}\AgdaSpace{}%
\AgdaSymbol{\{}\AgdaArgument{N}\AgdaSpace{}%
\AgdaSymbol{=}\AgdaSpace{}%
\AgdaBound{N₁}\AgdaSymbol{\}}\AgdaSpace{}%
\AgdaSymbol{\{}\AgdaArgument{N'}\AgdaSpace{}%
\AgdaSymbol{=}\AgdaSpace{}%
\AgdaBound{N'}\AgdaSymbol{\}}\AgdaSpace{}%
\AgdaBound{red}\AgdaSymbol{)}\AgdaSpace{}%
\AgdaSymbol{=}\AgdaSpace{}%
\AgdaInductiveConstructor{ν}\AgdaSpace{}%
\AgdaSymbol{(}\AgdaFunction{pres-rename-red-2}\AgdaSpace{}%
\AgdaBound{fb'}\AgdaSpace{}%
\AgdaBound{N₁}\AgdaSpace{}%
\AgdaBound{N'}\AgdaSpace{}%
\AgdaBound{red}\AgdaSymbol{)}\<%
\\
\>[0]\AgdaFunction{pres-rename-red-2}\AgdaSpace{}%
\AgdaBound{fb'}\AgdaSpace{}%
\AgdaBound{M}\AgdaSpace{}%
\AgdaBound{N}\AgdaSpace{}%
\AgdaSymbol{(}\AgdaInductiveConstructor{ξ}\AgdaSpace{}%
\AgdaSymbol{\{}\AgdaArgument{M}\AgdaSpace{}%
\AgdaSymbol{=}\AgdaSpace{}%
\AgdaBound{M₁}\AgdaSymbol{\}}\AgdaSpace{}%
\AgdaSymbol{\{}\AgdaArgument{M'}\AgdaSpace{}%
\AgdaSymbol{=}\AgdaSpace{}%
\AgdaBound{M'}\AgdaSymbol{\}}\AgdaSpace{}%
\AgdaBound{red}\AgdaSymbol{)}\AgdaSpace{}%
\AgdaSymbol{=}\AgdaSpace{}%
\AgdaInductiveConstructor{ξ}\AgdaSpace{}%
\AgdaSymbol{(}\AgdaFunction{pres-rename-red-2}\AgdaSpace{}%
\AgdaSymbol{(}\AgdaFunction{extb}\AgdaSpace{}%
\AgdaBound{fb'}\AgdaSymbol{)}\AgdaSpace{}%
\AgdaBound{M₁}\AgdaSpace{}%
\AgdaBound{M'}\AgdaSpace{}%
\AgdaBound{red}\AgdaSymbol{)}\<%
\\
\>[0]\AgdaFunction{pres-rename-red-2}\AgdaSpace{}%
\AgdaBound{fb'}\AgdaSpace{}%
\AgdaBound{M}\AgdaSpace{}%
\AgdaBound{N}\AgdaSpace{}%
\AgdaSymbol{(}\AgdaInductiveConstructor{ζ}\AgdaSpace{}%
\AgdaSymbol{\{}\AgdaArgument{M}\AgdaSpace{}%
\AgdaSymbol{=}\AgdaSpace{}%
\AgdaBound{M₁}\AgdaSymbol{\}}\AgdaSpace{}%
\AgdaSymbol{\{}\AgdaArgument{M'}\AgdaSpace{}%
\AgdaSymbol{=}\AgdaSpace{}%
\AgdaBound{M'}\AgdaSymbol{\}}\AgdaSpace{}%
\AgdaBound{red}\AgdaSymbol{)}\AgdaSpace{}%
\AgdaSymbol{=}\AgdaSpace{}%
\AgdaInductiveConstructor{ζ}\AgdaSpace{}%
\AgdaSymbol{(}\AgdaFunction{pres-rename-red-2}\AgdaSpace{}%
\AgdaBound{fb'}\AgdaSpace{}%
\AgdaBound{M₁}\AgdaSpace{}%
\AgdaBound{M'}\AgdaSpace{}%
\AgdaBound{red}\AgdaSymbol{)}\<%
\\
%
\\[\AgdaEmptyExtraSkip]%
\>[0]\AgdaComment{--\ extsb\ of\ ε\ is\ ε}\<%
\\
\>[0]\AgdaFunction{ε-exts}\AgdaSpace{}%
\AgdaSymbol{:}\AgdaSpace{}%
\AgdaSymbol{∀}\AgdaSpace{}%
\AgdaSymbol{\{}\AgdaBound{Γ}\AgdaSpace{}%
\AgdaBound{A}\AgdaSpace{}%
\AgdaBound{B}\AgdaSymbol{\}}\<%
\\
\>[0][@{}l@{\AgdaIndent{0}}]%
\>[2]\AgdaSymbol{→}\AgdaSpace{}%
\AgdaSymbol{(}\AgdaBound{x}\AgdaSpace{}%
\AgdaSymbol{:}\AgdaSpace{}%
\AgdaInductiveConstructor{□}\AgdaSpace{}%
\AgdaBound{A}\AgdaSpace{}%
\AgdaOperator{\AgdaDatatype{∈b}}\AgdaSpace{}%
\AgdaBound{Γ}\AgdaSpace{}%
\AgdaOperator{\AgdaInductiveConstructor{,b}}\AgdaSpace{}%
\AgdaBound{B}\AgdaSymbol{)}\<%
\\
%
\>[2]\AgdaSymbol{→}\AgdaSpace{}%
\AgdaFunction{extsb}\AgdaSpace{}%
\AgdaInductiveConstructor{ε}\AgdaSpace{}%
\AgdaBound{x}\AgdaSpace{}%
\AgdaOperator{\AgdaDatatype{≡}}\AgdaSpace{}%
\AgdaInductiveConstructor{ε}\AgdaSpace{}%
\AgdaBound{x}\<%
\\
\>[0]\AgdaFunction{ε-exts}\AgdaSpace{}%
\AgdaInductiveConstructor{Zb}\AgdaSpace{}%
\AgdaSymbol{=}\AgdaSpace{}%
\AgdaInductiveConstructor{refl}\<%
\\
\>[0]\AgdaFunction{ε-exts}\AgdaSpace{}%
\AgdaSymbol{(}\AgdaInductiveConstructor{Sb}\AgdaSpace{}%
\AgdaBound{x}\AgdaSymbol{)}\AgdaSpace{}%
\AgdaSymbol{=}\AgdaSpace{}%
\AgdaInductiveConstructor{refl}\<%
\\
%
\\[\AgdaEmptyExtraSkip]%
\>[0]\AgdaComment{--\ Simultaneous\ substitution\ with\ ε\ }\<%
\\
\>[0]\AgdaComment{--\ leaves\ the\ term\ unchanged\ \ }\<%
\\
\>[0]\AgdaFunction{subst-with-id}\AgdaSpace{}%
\AgdaSymbol{:}\AgdaSpace{}%
\AgdaSymbol{∀}\AgdaSpace{}%
\AgdaSymbol{\{}\AgdaBound{Γ}\AgdaSpace{}%
\AgdaBound{A}\AgdaSymbol{\}}\<%
\\
\>[0][@{}l@{\AgdaIndent{0}}]%
\>[2]\AgdaSymbol{→}\AgdaSpace{}%
\AgdaSymbol{(}\AgdaBound{M}\AgdaSpace{}%
\AgdaSymbol{:}\AgdaSpace{}%
\AgdaBound{Γ}\AgdaSpace{}%
\AgdaOperator{\AgdaDatatype{⊢b}}\AgdaSpace{}%
\AgdaBound{A}\AgdaSymbol{)}\<%
\\
%
\>[2]\AgdaSymbol{→}\AgdaSpace{}%
\AgdaFunction{substb}\AgdaSpace{}%
\AgdaSymbol{(λ}\AgdaSpace{}%
\AgdaBound{x₁}\AgdaSpace{}%
\AgdaSymbol{→}\AgdaSpace{}%
\AgdaInductiveConstructor{ε}\AgdaSpace{}%
\AgdaBound{x₁}\AgdaSymbol{)}\AgdaSpace{}%
\AgdaBound{M}\AgdaSpace{}%
\AgdaOperator{\AgdaDatatype{≡}}\AgdaSpace{}%
\AgdaBound{M}\<%
\\
\>[0]\AgdaFunction{subst-with-id}\AgdaSpace{}%
\AgdaSymbol{(}\AgdaInductiveConstructor{ε}\AgdaSpace{}%
\AgdaBound{x}\AgdaSymbol{)}\AgdaSpace{}%
\AgdaSymbol{=}\AgdaSpace{}%
\AgdaInductiveConstructor{refl}\<%
\\
\>[0]\AgdaFunction{subst-with-id}\AgdaSpace{}%
\AgdaSymbol{(}\AgdaInductiveConstructor{box}\AgdaSpace{}%
\AgdaBound{M}\AgdaSymbol{)}\AgdaSpace{}%
\AgdaSymbol{=}\AgdaSpace{}%
\AgdaFunction{cong}\AgdaSpace{}%
\AgdaInductiveConstructor{box}\AgdaSpace{}%
\AgdaSymbol{(}\AgdaFunction{subst-with-id}\AgdaSpace{}%
\AgdaBound{M}\AgdaSymbol{)}\<%
\\
\>[0]\AgdaFunction{subst-with-id}\AgdaSpace{}%
\AgdaSymbol{(}\AgdaInductiveConstructor{ƛb}\AgdaSpace{}%
\AgdaBound{M}\AgdaSymbol{)}\AgdaSpace{}%
\AgdaSymbol{=}\<%
\\
\>[0][@{}l@{\AgdaIndent{0}}]%
\>[2]\AgdaOperator{\AgdaFunction{beginEq}}\<%
\\
\>[2][@{}l@{\AgdaIndent{0}}]%
\>[4]\AgdaInductiveConstructor{ƛb}\AgdaSpace{}%
\AgdaSymbol{(}\AgdaFunction{substb}\AgdaSpace{}%
\AgdaSymbol{(}\AgdaFunction{extsb}\AgdaSpace{}%
\AgdaInductiveConstructor{ε}\AgdaSymbol{)}\AgdaSpace{}%
\AgdaBound{M}\AgdaSymbol{)}\<%
\\
%
\>[2]\AgdaFunction{≡⟨}\AgdaSpace{}%
\AgdaFunction{cong}\AgdaSpace{}%
\AgdaInductiveConstructor{ƛb}\AgdaSpace{}%
\AgdaSymbol{(}\AgdaFunction{substb-cong}\AgdaSpace{}%
\AgdaSymbol{(λ}\AgdaSpace{}%
\AgdaBound{x}\AgdaSpace{}%
\AgdaSymbol{→}\AgdaSpace{}%
\AgdaFunction{ε-exts}\AgdaSpace{}%
\AgdaBound{x}\AgdaSymbol{)}\AgdaSpace{}%
\AgdaBound{M}\AgdaSymbol{)}\AgdaSpace{}%
\AgdaFunction{⟩}\<%
\\
\>[2][@{}l@{\AgdaIndent{0}}]%
\>[4]\AgdaInductiveConstructor{ƛb}\AgdaSpace{}%
\AgdaSymbol{(}\AgdaFunction{substb}\AgdaSpace{}%
\AgdaInductiveConstructor{ε}\AgdaSpace{}%
\AgdaBound{M}\AgdaSymbol{)}\<%
\\
%
\>[2]\AgdaFunction{≡⟨}\AgdaSpace{}%
\AgdaFunction{cong}\AgdaSpace{}%
\AgdaInductiveConstructor{ƛb}\AgdaSpace{}%
\AgdaSymbol{(}\AgdaFunction{subst-with-id}\AgdaSpace{}%
\AgdaBound{M}\AgdaSymbol{)}\AgdaSpace{}%
\AgdaFunction{⟩}\<%
\\
\>[2][@{}l@{\AgdaIndent{0}}]%
\>[4]\AgdaInductiveConstructor{ƛb}\AgdaSpace{}%
\AgdaBound{M}\<%
\\
%
\>[2]\AgdaOperator{\AgdaFunction{∎eq}}\<%
\\
\>[0]\AgdaFunction{subst-with-id}\AgdaSpace{}%
\AgdaSymbol{(}\AgdaBound{M}\AgdaSpace{}%
\AgdaOperator{\AgdaInductiveConstructor{·b}}\AgdaSpace{}%
\AgdaBound{M₁}\AgdaSymbol{)}\AgdaSpace{}%
\AgdaSymbol{=}\AgdaSpace{}%
\AgdaFunction{cong₂}\AgdaSpace{}%
\AgdaOperator{\AgdaInductiveConstructor{\AgdaUnderscore{}·b\AgdaUnderscore{}}}\AgdaSpace{}%
\AgdaSymbol{(}\AgdaFunction{subst-with-id}\AgdaSpace{}%
\AgdaBound{M}\AgdaSymbol{)}\AgdaSpace{}%
\AgdaSymbol{(}\AgdaFunction{subst-with-id}\AgdaSpace{}%
\AgdaBound{M₁}\AgdaSymbol{)}\<%
\\
%
\\[\AgdaEmptyExtraSkip]%
\>[0]\AgdaComment{--\ Substitution\ composed\ with\ renaming\ with\ Sb\ }\<%
\\
\>[0]\AgdaComment{--\ is\ the\ same\ as\ embedding\ that\ term.}\<%
\\
\>[0]\AgdaFunction{subst-rename-fb-comp}\AgdaSpace{}%
\AgdaSymbol{:}\AgdaSpace{}%
\AgdaSymbol{∀}\AgdaSpace{}%
\AgdaSymbol{\{}\AgdaBound{Γ}\AgdaSpace{}%
\AgdaBound{A}\AgdaSpace{}%
\AgdaBound{B}\AgdaSymbol{\}}\<%
\\
\>[0][@{}l@{\AgdaIndent{0}}]%
\>[2]\AgdaSymbol{→}\AgdaSpace{}%
\AgdaSymbol{(}\AgdaBound{M₁}\AgdaSpace{}%
\AgdaSymbol{:}\AgdaSpace{}%
\AgdaBound{Γ}\AgdaSpace{}%
\AgdaOperator{\AgdaInductiveConstructor{,}}\AgdaSpace{}%
\AgdaBound{B}\AgdaSpace{}%
\AgdaOperator{\AgdaDatatype{⊢}}\AgdaSpace{}%
\AgdaBound{A}\AgdaSymbol{)}\<%
\\
%
\>[2]\AgdaSymbol{→}\AgdaSpace{}%
\AgdaSymbol{(}\AgdaBound{M₂}\AgdaSpace{}%
\AgdaSymbol{:}\AgdaSpace{}%
\AgdaBound{Γ}\AgdaSpace{}%
\AgdaOperator{\AgdaDatatype{⊢}}\AgdaSpace{}%
\AgdaBound{B}\AgdaSymbol{)}\<%
\\
%
\>[2]\AgdaSymbol{→}\AgdaSpace{}%
\AgdaSymbol{(}\AgdaBound{x}\AgdaSpace{}%
\AgdaSymbol{:}\AgdaSpace{}%
\AgdaDatatype{Value}\AgdaSpace{}%
\AgdaBound{M₂}\AgdaSpace{}%
\AgdaSymbol{)}\<%
\\
%
\>[2]\AgdaSymbol{→}%
\>[2725I]\AgdaFunction{substb}\AgdaSpace{}%
\AgdaSymbol{(}\AgdaFunction{fb}\AgdaSpace{}%
\AgdaSymbol{(}\AgdaInductiveConstructor{ƛb}\AgdaSpace{}%
\AgdaSymbol{(}\AgdaFunction{embedTerm}\AgdaSpace{}%
\AgdaBound{M₁}\AgdaSymbol{)))}\AgdaSpace{}%
\AgdaSymbol{(}\AgdaFunction{renameb}\AgdaSpace{}%
\AgdaInductiveConstructor{Sb}\AgdaSpace{}%
\AgdaSymbol{(}\AgdaFunction{embedTerm}\AgdaSpace{}%
\AgdaBound{M₂}\AgdaSymbol{))}\<%
\\
\>[.][@{}l@{}]\<[2725I]%
\>[4]\AgdaOperator{\AgdaDatatype{≡}}\AgdaSpace{}%
\AgdaInductiveConstructor{box}\AgdaSpace{}%
\AgdaSymbol{(}\AgdaFunction{embedValue}\AgdaSpace{}%
\AgdaBound{M₂}\AgdaSpace{}%
\AgdaBound{x}\AgdaSymbol{)}\<%
\\
\>[0]\AgdaFunction{subst-rename-fb-comp}\AgdaSpace{}%
\AgdaBound{M₁}\AgdaSpace{}%
\AgdaBound{M₂}\AgdaSpace{}%
\AgdaBound{x}\AgdaSpace{}%
\AgdaSymbol{=}\<%
\\
\>[0][@{}l@{\AgdaIndent{0}}]%
\>[2]\AgdaOperator{\AgdaFunction{beginEq}}\<%
\\
\>[2][@{}l@{\AgdaIndent{0}}]%
\>[4]\AgdaFunction{substb}\AgdaSpace{}%
\AgdaSymbol{(}\AgdaFunction{fb}\AgdaSpace{}%
\AgdaSymbol{(}\AgdaInductiveConstructor{ƛb}\AgdaSpace{}%
\AgdaSymbol{(}\AgdaFunction{embedTerm}\AgdaSpace{}%
\AgdaBound{M₁}\AgdaSymbol{)))}\AgdaSpace{}%
\AgdaSymbol{(}\AgdaFunction{renameb}\AgdaSpace{}%
\AgdaInductiveConstructor{Sb}\AgdaSpace{}%
\AgdaSymbol{(}\AgdaFunction{embedTerm}\AgdaSpace{}%
\AgdaBound{M₂}\AgdaSymbol{))}\<%
\\
%
\>[2]\AgdaFunction{≡⟨}\AgdaSpace{}%
\AgdaFunction{subst-renameb-comp}\AgdaSpace{}%
\AgdaInductiveConstructor{Sb}\AgdaSpace{}%
\AgdaSymbol{(}\AgdaFunction{fb}\AgdaSpace{}%
\AgdaSymbol{(}\AgdaInductiveConstructor{ƛb}\AgdaSpace{}%
\AgdaSymbol{(}\AgdaFunction{embedTerm}\AgdaSpace{}%
\AgdaBound{M₁}\AgdaSymbol{)))}\AgdaSpace{}%
\AgdaSymbol{(}\AgdaFunction{embedTerm}\AgdaSpace{}%
\AgdaBound{M₂}\AgdaSymbol{)}\AgdaSpace{}%
\AgdaFunction{⟩}\<%
\\
\>[2][@{}l@{\AgdaIndent{0}}]%
\>[4]\AgdaFunction{substb}\AgdaSpace{}%
\AgdaSymbol{(λ}\AgdaSpace{}%
\AgdaBound{x₁}\AgdaSpace{}%
\AgdaSymbol{→}\AgdaSpace{}%
\AgdaInductiveConstructor{ε}\AgdaSpace{}%
\AgdaBound{x₁}\AgdaSymbol{)}\AgdaSpace{}%
\AgdaSymbol{(}\AgdaFunction{embedTerm}\AgdaSpace{}%
\AgdaBound{M₂}\AgdaSymbol{)}\<%
\\
%
\>[2]\AgdaFunction{≡⟨}\AgdaSpace{}%
\AgdaFunction{subst-with-id}\AgdaSpace{}%
\AgdaSymbol{(}\AgdaFunction{embedTerm}\AgdaSpace{}%
\AgdaBound{M₂}\AgdaSymbol{)}\AgdaSpace{}%
\AgdaFunction{⟩}\<%
\\
\>[2][@{}l@{\AgdaIndent{0}}]%
\>[4]\AgdaFunction{embedTerm}\AgdaSpace{}%
\AgdaBound{M₂}\<%
\\
%
\>[2]\AgdaFunction{≡⟨}\AgdaSpace{}%
\AgdaFunction{sym}\AgdaSpace{}%
\AgdaSymbol{(}\AgdaFunction{box-embedValue}\AgdaSpace{}%
\AgdaBound{x}\AgdaSymbol{)}\AgdaSpace{}%
\AgdaFunction{⟩}\<%
\\
\>[2][@{}l@{\AgdaIndent{0}}]%
\>[4]\AgdaInductiveConstructor{box}\AgdaSpace{}%
\AgdaSymbol{(}\AgdaFunction{embedValue}\AgdaSpace{}%
\AgdaBound{M₂}\AgdaSpace{}%
\AgdaBound{x}\AgdaSymbol{)}\<%
\\
%
\>[2]\AgdaOperator{\AgdaFunction{∎eq}}\<%
\\
%
\\[\AgdaEmptyExtraSkip]%
\>[0]\AgdaComment{--\ Preservation\ of\ 2\ step\ reduction\ of\ Gödel's\ translation}\<%
\\
\>[0]\AgdaFunction{pres-red}\AgdaSpace{}%
\AgdaSymbol{:}\AgdaSpace{}%
\AgdaSymbol{∀}\AgdaSpace{}%
\AgdaSymbol{\{}\AgdaBound{Γ}\AgdaSpace{}%
\AgdaBound{A}\AgdaSymbol{\}}\<%
\\
\>[0][@{}l@{\AgdaIndent{0}}]%
\>[2]\AgdaSymbol{→}\AgdaSpace{}%
\AgdaSymbol{(}\AgdaBound{M}\AgdaSpace{}%
\AgdaSymbol{:}\AgdaSpace{}%
\AgdaBound{Γ}\AgdaSpace{}%
\AgdaOperator{\AgdaDatatype{⊢}}\AgdaSpace{}%
\AgdaBound{A}\AgdaSymbol{)}\<%
\\
%
\>[2]\AgdaSymbol{→}\AgdaSpace{}%
\AgdaSymbol{(}\AgdaBound{N}\AgdaSpace{}%
\AgdaSymbol{:}\AgdaSpace{}%
\AgdaBound{Γ}\AgdaSpace{}%
\AgdaOperator{\AgdaDatatype{⊢}}\AgdaSpace{}%
\AgdaBound{A}\AgdaSymbol{)}\<%
\\
%
\>[2]\AgdaSymbol{→}\AgdaSpace{}%
\AgdaSymbol{(}\AgdaBound{M}\AgdaSpace{}%
\AgdaOperator{\AgdaDatatype{↝βᵥ}}\AgdaSpace{}%
\AgdaBound{N}\AgdaSymbol{)}\<%
\\
%
\>[2]\AgdaSymbol{→}\AgdaSpace{}%
\AgdaSymbol{(}\AgdaFunction{embedTerm}\AgdaSpace{}%
\AgdaBound{M}\AgdaSpace{}%
\AgdaOperator{\AgdaDatatype{↝βb²}}\AgdaSpace{}%
\AgdaFunction{embedTerm}\AgdaSpace{}%
\AgdaBound{N}\AgdaSymbol{)}\<%
\\
\>[0]\AgdaFunction{pres-red}\AgdaSpace{}%
\AgdaSymbol{(}\AgdaInductiveConstructor{ƛ}\AgdaSpace{}%
\AgdaBound{M₁}\AgdaSpace{}%
\AgdaOperator{\AgdaInductiveConstructor{·}}\AgdaSpace{}%
\AgdaBound{M₂}\AgdaSymbol{)}\AgdaSpace{}%
\AgdaBound{N}\AgdaSpace{}%
\AgdaSymbol{(}\AgdaInductiveConstructor{βᵥ}\AgdaSpace{}%
\AgdaBound{x}\AgdaSymbol{)}\AgdaSpace{}%
\AgdaSymbol{=}\<%
\\
\>[0][@{}l@{\AgdaIndent{0}}]%
\>[2]\AgdaInductiveConstructor{βb²}\<%
\\
%
\>[2]\AgdaSymbol{(}\AgdaFunction{raise}\AgdaSpace{}%
\AgdaSymbol{(}\AgdaFunction{embedTerm}\AgdaSpace{}%
\AgdaBound{M₂}\AgdaSymbol{)}\AgdaSpace{}%
\AgdaOperator{\AgdaInductiveConstructor{·b}}\AgdaSpace{}%
\AgdaInductiveConstructor{box}\AgdaSpace{}%
\AgdaSymbol{(}\AgdaInductiveConstructor{ƛb}\AgdaSpace{}%
\AgdaSymbol{(}\AgdaFunction{embedTerm}\AgdaSpace{}%
\AgdaBound{M₁}\AgdaSymbol{)))}\<%
\\
%
\>[2]\AgdaSymbol{(}\AgdaInductiveConstructor{ƛb}\AgdaSpace{}%
\AgdaSymbol{(}\AgdaFunction{embedTerm}\AgdaSpace{}%
\AgdaBound{M₁}\AgdaSymbol{)}\AgdaSpace{}%
\AgdaOperator{\AgdaInductiveConstructor{·b}}\AgdaSpace{}%
\AgdaInductiveConstructor{box}\AgdaSpace{}%
\AgdaSymbol{(}\AgdaFunction{embedValue}\AgdaSpace{}%
\AgdaBound{M₂}\AgdaSpace{}%
\AgdaBound{x}\AgdaSymbol{))}\<%
\\
%
\>[2]\AgdaSymbol{(}\AgdaFunction{embedTerm}\AgdaSpace{}%
\AgdaSymbol{(}\AgdaBound{M₁}\AgdaSpace{}%
\AgdaOperator{\AgdaFunction{[}}\AgdaSpace{}%
\AgdaBound{M₂}\AgdaSpace{}%
\AgdaOperator{\AgdaFunction{]}}\AgdaSymbol{))}\<%
\\
%
\>[2]\AgdaSymbol{(}\AgdaFunction{Eq.subst}\<%
\\
\>[2][@{}l@{\AgdaIndent{0}}]%
\>[4]\AgdaSymbol{(λ}\AgdaSpace{}%
\AgdaBound{x}\AgdaSpace{}%
\AgdaSymbol{→}\AgdaSpace{}%
\AgdaInductiveConstructor{ƛb}\AgdaSpace{}%
\AgdaSymbol{(}\AgdaInductiveConstructor{ε}\AgdaSpace{}%
\AgdaInductiveConstructor{Zb}\AgdaSpace{}%
\AgdaOperator{\AgdaInductiveConstructor{·b}}\AgdaSpace{}%
\AgdaFunction{renameb}\AgdaSpace{}%
\AgdaInductiveConstructor{Sb}\AgdaSpace{}%
\AgdaSymbol{(}\AgdaFunction{embedTerm}\AgdaSpace{}%
\AgdaBound{M₂}\AgdaSymbol{))}\AgdaSpace{}%
\AgdaOperator{\AgdaInductiveConstructor{·b}}\AgdaSpace{}%
\AgdaInductiveConstructor{box}\AgdaSpace{}%
\AgdaSymbol{(}\AgdaInductiveConstructor{ƛb}\AgdaSpace{}%
\AgdaSymbol{(}\AgdaFunction{embedTerm}\AgdaSpace{}%
\AgdaBound{M₁}\AgdaSymbol{))}\AgdaSpace{}%
\AgdaOperator{\AgdaDatatype{↝βb}}\AgdaSpace{}%
\AgdaBound{x}\AgdaSymbol{)}\<%
\\
%
\>[4]\AgdaSymbol{(}\AgdaFunction{cong₂}\AgdaSpace{}%
\AgdaOperator{\AgdaInductiveConstructor{\AgdaUnderscore{}·b\AgdaUnderscore{}}}\AgdaSpace{}%
\AgdaInductiveConstructor{refl}\AgdaSpace{}%
\AgdaSymbol{(}\AgdaFunction{subst-rename-fb-comp}\AgdaSpace{}%
\AgdaBound{M₁}\AgdaSpace{}%
\AgdaBound{M₂}\AgdaSpace{}%
\AgdaBound{x}\AgdaSymbol{))}\<%
\\
%
\>[4]\AgdaInductiveConstructor{βb}\AgdaSymbol{)}\<%
\\
%
\>[2]\AgdaSymbol{(}\AgdaFunction{Eq.subst}\AgdaSpace{}%
\AgdaSymbol{(λ}\AgdaSpace{}%
\AgdaBound{t}\AgdaSpace{}%
\AgdaSymbol{→}\AgdaSpace{}%
\AgdaInductiveConstructor{ƛb}\AgdaSpace{}%
\AgdaSymbol{(}\AgdaFunction{embedTerm}\AgdaSpace{}%
\AgdaBound{M₁}\AgdaSymbol{)}\AgdaSpace{}%
\AgdaOperator{\AgdaInductiveConstructor{·b}}\AgdaSpace{}%
\AgdaInductiveConstructor{box}\AgdaSpace{}%
\AgdaSymbol{(}\AgdaFunction{embedValue}\AgdaSpace{}%
\AgdaBound{M₂}\AgdaSpace{}%
\AgdaBound{x}\AgdaSymbol{)}\AgdaSpace{}%
\AgdaOperator{\AgdaDatatype{↝βb}}\AgdaSpace{}%
\AgdaBound{t}\AgdaSymbol{)}\AgdaSpace{}%
\AgdaSymbol{(}\AgdaFunction{pres-subst}\AgdaSpace{}%
\AgdaBound{M₁}\AgdaSpace{}%
\AgdaBound{M₂}\AgdaSpace{}%
\AgdaBound{x}\AgdaSymbol{)}\AgdaSpace{}%
\AgdaInductiveConstructor{βb}\AgdaSymbol{)}\<%
\\
\>[0]\AgdaFunction{pres-red}\AgdaSpace{}%
\AgdaBound{M}\AgdaSpace{}%
\AgdaBound{N}\AgdaSpace{}%
\AgdaSymbol{(}\AgdaInductiveConstructor{μ}\AgdaSpace{}%
\AgdaSymbol{\{}\AgdaArgument{M}\AgdaSpace{}%
\AgdaSymbol{=}\AgdaSpace{}%
\AgdaBound{M₁}\AgdaSymbol{\}}\AgdaSpace{}%
\AgdaSymbol{\{}\AgdaArgument{M'}\AgdaSpace{}%
\AgdaSymbol{=}\AgdaSpace{}%
\AgdaBound{M'}\AgdaSymbol{\}}\AgdaSpace{}%
\AgdaBound{red}\AgdaSymbol{)}\AgdaSpace{}%
\AgdaSymbol{=}\AgdaSpace{}%
\AgdaInductiveConstructor{ν}\AgdaSpace{}%
\AgdaSymbol{(}\AgdaFunction{pres-red}\AgdaSpace{}%
\AgdaBound{M₁}\AgdaSpace{}%
\AgdaBound{M'}\AgdaSpace{}%
\AgdaBound{red}\AgdaSymbol{)}\<%
\\
\>[0]\AgdaFunction{pres-red}\AgdaSpace{}%
\AgdaBound{M}\AgdaSpace{}%
\AgdaBound{N}\AgdaSpace{}%
\AgdaSymbol{(}\AgdaInductiveConstructor{ν<}\AgdaSpace{}%
\AgdaSymbol{\{}\AgdaArgument{M}\AgdaSpace{}%
\AgdaSymbol{=}\AgdaSpace{}%
\AgdaBound{M₁}\AgdaSymbol{\}}\AgdaSpace{}%
\AgdaSymbol{\{}\AgdaArgument{N}\AgdaSpace{}%
\AgdaSymbol{=}\AgdaSpace{}%
\AgdaBound{N₁}\AgdaSymbol{\}}\AgdaSpace{}%
\AgdaSymbol{\{}\AgdaArgument{N'}\AgdaSpace{}%
\AgdaSymbol{=}\AgdaSpace{}%
\AgdaBound{N'}\AgdaSymbol{\}}\AgdaSpace{}%
\AgdaBound{x}\AgdaSpace{}%
\AgdaBound{red}\AgdaSymbol{)}\AgdaSpace{}%
\AgdaSymbol{=}\AgdaSpace{}%
\AgdaInductiveConstructor{μ}\AgdaSpace{}%
\AgdaSymbol{(}\AgdaInductiveConstructor{ξ}\AgdaSpace{}%
\AgdaSymbol{(}\AgdaInductiveConstructor{ν}\AgdaSpace{}%
\AgdaSymbol{(}\AgdaFunction{pres-rename-red-2}\AgdaSpace{}%
\AgdaInductiveConstructor{Sb}\AgdaSpace{}%
\AgdaSymbol{(}\AgdaFunction{embedTerm}\AgdaSpace{}%
\AgdaBound{N₁}\AgdaSymbol{)}\AgdaSpace{}%
\AgdaSymbol{(}\AgdaFunction{embedTerm}\AgdaSpace{}%
\AgdaBound{N'}\AgdaSymbol{)}\AgdaSpace{}%
\AgdaSymbol{(}\AgdaFunction{pres-red}\AgdaSpace{}%
\AgdaBound{N₁}\AgdaSpace{}%
\AgdaBound{N'}\AgdaSpace{}%
\AgdaBound{red}\AgdaSymbol{))))}\<%
\\
\>[0]\AgdaFunction{pres-red}\AgdaSpace{}%
\AgdaBound{M}\AgdaSpace{}%
\AgdaBound{N}\AgdaSpace{}%
\AgdaSymbol{(}\AgdaInductiveConstructor{ν}\AgdaSpace{}%
\AgdaSymbol{\{}\AgdaArgument{M}\AgdaSpace{}%
\AgdaSymbol{=}\AgdaSpace{}%
\AgdaBound{M₁}\AgdaSymbol{\}}\AgdaSpace{}%
\AgdaSymbol{\{}\AgdaArgument{N}\AgdaSpace{}%
\AgdaSymbol{=}\AgdaSpace{}%
\AgdaBound{N₁}\AgdaSymbol{\}}\AgdaSpace{}%
\AgdaSymbol{\{}\AgdaArgument{N'}\AgdaSpace{}%
\AgdaSymbol{=}\AgdaSpace{}%
\AgdaBound{N'}\AgdaSymbol{\}}\AgdaSpace{}%
\AgdaBound{red}\AgdaSymbol{)}\AgdaSpace{}%
\AgdaSymbol{=}\AgdaSpace{}%
\AgdaInductiveConstructor{μ}\AgdaSpace{}%
\AgdaSymbol{(}\AgdaInductiveConstructor{ξ}\AgdaSpace{}%
\AgdaSymbol{(}\AgdaInductiveConstructor{ν}\AgdaSpace{}%
\AgdaSymbol{(}\AgdaFunction{pres-rename-red-2}\AgdaSpace{}%
\AgdaInductiveConstructor{Sb}\AgdaSpace{}%
\AgdaSymbol{(}\AgdaFunction{embedTerm}\AgdaSpace{}%
\AgdaBound{N₁}\AgdaSymbol{)}\AgdaSpace{}%
\AgdaSymbol{(}\AgdaFunction{embedTerm}\AgdaSpace{}%
\AgdaBound{N'}\AgdaSymbol{)}\AgdaSpace{}%
\AgdaSymbol{(}\AgdaFunction{pres-red}\AgdaSpace{}%
\AgdaBound{N₁}\AgdaSpace{}%
\AgdaBound{N'}\AgdaSpace{}%
\AgdaBound{red}\AgdaSymbol{))))}\<%
\\
\>[0]\AgdaFunction{pres-red}\AgdaSpace{}%
\AgdaBound{M}\AgdaSpace{}%
\AgdaBound{N}\AgdaSpace{}%
\AgdaSymbol{(}\AgdaInductiveConstructor{ξ}\AgdaSpace{}%
\AgdaSymbol{\{}\AgdaArgument{M}\AgdaSpace{}%
\AgdaSymbol{=}\AgdaSpace{}%
\AgdaBound{M₁}\AgdaSymbol{\}}\AgdaSpace{}%
\AgdaSymbol{\{}\AgdaArgument{M'}\AgdaSpace{}%
\AgdaSymbol{=}\AgdaSpace{}%
\AgdaBound{M'}\AgdaSymbol{\}}\AgdaSpace{}%
\AgdaBound{red}\AgdaSymbol{)}\AgdaSpace{}%
\AgdaSymbol{=}\AgdaSpace{}%
\AgdaInductiveConstructor{ζ}\AgdaSpace{}%
\AgdaSymbol{(}\AgdaInductiveConstructor{ξ}\AgdaSpace{}%
\AgdaSymbol{(}\AgdaFunction{pres-red}\AgdaSpace{}%
\AgdaBound{M₁}\AgdaSpace{}%
\AgdaBound{M'}\AgdaSpace{}%
\AgdaBound{red}\AgdaSymbol{))}\<%
\\
%
\\[\AgdaEmptyExtraSkip]%
\>[0]\AgdaComment{--\ Preservation\ of\ the\ 2\ step\ reduction\ rule\ under\ rename}\<%
\\
\>[0]\AgdaFunction{pres-rename-red-↝βb²}\AgdaSpace{}%
\AgdaSymbol{:}\AgdaSpace{}%
\AgdaSymbol{∀}\AgdaSpace{}%
\AgdaSymbol{\{}\AgdaBound{Γ}\AgdaSpace{}%
\AgdaBound{A}\AgdaSpace{}%
\AgdaBound{Δ}\AgdaSymbol{\}}\<%
\\
\>[0][@{}l@{\AgdaIndent{0}}]%
\>[2]\AgdaSymbol{→}\AgdaSpace{}%
\AgdaSymbol{(}\AgdaBound{fb'}\AgdaSpace{}%
\AgdaSymbol{:}\AgdaSpace{}%
\AgdaSymbol{∀}\AgdaSpace{}%
\AgdaSymbol{\{}\AgdaBound{Ab}\AgdaSymbol{\}}\AgdaSpace{}%
\AgdaSymbol{→}\AgdaSpace{}%
\AgdaBound{Ab}\AgdaSpace{}%
\AgdaOperator{\AgdaDatatype{∈b}}\AgdaSpace{}%
\AgdaBound{Γ}\AgdaSpace{}%
\AgdaSymbol{→}\AgdaSpace{}%
\AgdaBound{Ab}\AgdaSpace{}%
\AgdaOperator{\AgdaDatatype{∈b}}\AgdaSpace{}%
\AgdaBound{Δ}\AgdaSymbol{)}\<%
\\
%
\>[2]\AgdaSymbol{→}\AgdaSpace{}%
\AgdaSymbol{(}\AgdaBound{M}\AgdaSpace{}%
\AgdaSymbol{:}\AgdaSpace{}%
\AgdaBound{Γ}\AgdaSpace{}%
\AgdaOperator{\AgdaDatatype{⊢b}}\AgdaSpace{}%
\AgdaBound{A}\AgdaSymbol{)}\<%
\\
%
\>[2]\AgdaSymbol{→}\AgdaSpace{}%
\AgdaSymbol{(}\AgdaBound{N}\AgdaSpace{}%
\AgdaSymbol{:}\AgdaSpace{}%
\AgdaBound{Γ}\AgdaSpace{}%
\AgdaOperator{\AgdaDatatype{⊢b}}\AgdaSpace{}%
\AgdaBound{A}\AgdaSymbol{)}\<%
\\
%
\>[2]\AgdaSymbol{→}\AgdaSpace{}%
\AgdaSymbol{(}\AgdaBound{M}\AgdaSpace{}%
\AgdaOperator{\AgdaDatatype{↝βb₂}}\AgdaSpace{}%
\AgdaBound{N}\AgdaSymbol{)}\<%
\\
%
\>[2]\AgdaSymbol{→}\AgdaSpace{}%
\AgdaFunction{renameb}\AgdaSpace{}%
\AgdaBound{fb'}\AgdaSpace{}%
\AgdaBound{M}\AgdaSpace{}%
\AgdaOperator{\AgdaDatatype{↝βb₂}}\AgdaSpace{}%
\AgdaFunction{renameb}\AgdaSpace{}%
\AgdaBound{fb'}\AgdaSpace{}%
\AgdaBound{N}\<%
\\
\>[0]\AgdaFunction{pres-rename-red-↝βb²}\AgdaSpace{}%
\AgdaBound{fb'}\AgdaSpace{}%
\AgdaBound{M}\AgdaSpace{}%
\AgdaBound{N}\AgdaSpace{}%
\AgdaSymbol{(}\AgdaInductiveConstructor{βb₂}\AgdaSpace{}%
\AgdaSymbol{\{}\AgdaArgument{M}\AgdaSpace{}%
\AgdaSymbol{=}\AgdaSpace{}%
\AgdaBound{M'}\AgdaSymbol{\}}\AgdaSpace{}%
\AgdaSymbol{\{}\AgdaArgument{W}\AgdaSpace{}%
\AgdaSymbol{=}\AgdaSpace{}%
\AgdaBound{W'}\AgdaSymbol{\})}\AgdaSpace{}%
\AgdaSymbol{=}\<%
\\
\>[0][@{}l@{\AgdaIndent{0}}]%
\>[2]\AgdaFunction{Eq.subst}\<%
\\
\>[2][@{}l@{\AgdaIndent{0}}]%
\>[4]\AgdaSymbol{(λ}\AgdaSpace{}%
\AgdaBound{x}\AgdaSpace{}%
\AgdaSymbol{→}\AgdaSpace{}%
\AgdaBound{x}\AgdaSpace{}%
\AgdaOperator{\AgdaDatatype{↝βb₂}}\AgdaSpace{}%
\AgdaFunction{renameb}\AgdaSpace{}%
\AgdaBound{fb'}\<%
\\
%
\>[4]\AgdaSymbol{(}\AgdaBound{M'}\AgdaSpace{}%
\AgdaOperator{\AgdaFunction{[}}\AgdaSpace{}%
\AgdaBound{W'}\AgdaSpace{}%
\AgdaOperator{\AgdaFunction{]b}}\AgdaSymbol{))}\<%
\\
%
\>[4]\AgdaSymbol{(}\AgdaFunction{cong₂}\AgdaSpace{}%
\AgdaOperator{\AgdaInductiveConstructor{\AgdaUnderscore{}·b\AgdaUnderscore{}}}\AgdaSpace{}%
\AgdaSymbol{(}\AgdaFunction{cong}\AgdaSpace{}%
\AgdaInductiveConstructor{ƛb}\AgdaSpace{}%
\AgdaSymbol{(}\AgdaFunction{cong₂}\AgdaSpace{}%
\AgdaOperator{\AgdaInductiveConstructor{\AgdaUnderscore{}·b\AgdaUnderscore{}}}\AgdaSpace{}%
\AgdaInductiveConstructor{refl}\AgdaSpace{}%
\AgdaSymbol{(}\AgdaFunction{cong}\AgdaSpace{}%
\AgdaInductiveConstructor{box}\AgdaSpace{}%
\AgdaSymbol{(}\AgdaFunction{sym}\AgdaSpace{}%
\AgdaSymbol{(}\AgdaFunction{rename-extb-Sb-commute}\AgdaSpace{}%
\AgdaBound{fb'}\AgdaSpace{}%
\AgdaBound{W'}\AgdaSymbol{)))))}\AgdaSpace{}%
\AgdaInductiveConstructor{refl}\AgdaSymbol{)}\<%
\\
%
\>[4]\AgdaSymbol{(}\AgdaFunction{Eq.subst}\<%
\\
\>[4][@{}l@{\AgdaIndent{0}}]%
\>[6]\AgdaSymbol{(λ}\AgdaSpace{}%
\AgdaBound{x}\AgdaSpace{}%
\AgdaSymbol{→}\AgdaSpace{}%
\AgdaSymbol{(}\AgdaFunction{raise}%
\>[21]\AgdaSymbol{(}\AgdaInductiveConstructor{box}\AgdaSpace{}%
\AgdaSymbol{(}\AgdaFunction{renameb}\AgdaSpace{}%
\AgdaBound{fb'}\AgdaSpace{}%
\AgdaBound{W'}\AgdaSymbol{)))}\AgdaSpace{}%
\AgdaOperator{\AgdaInductiveConstructor{·b}}\AgdaSpace{}%
\AgdaInductiveConstructor{box}\AgdaSpace{}%
\AgdaSymbol{(}\AgdaInductiveConstructor{ƛb}\AgdaSpace{}%
\AgdaSymbol{(}\AgdaFunction{renameb}\AgdaSpace{}%
\AgdaSymbol{(}\AgdaFunction{extb}\AgdaSpace{}%
\AgdaBound{fb'}\AgdaSymbol{)}\AgdaSpace{}%
\AgdaBound{M'}\AgdaSymbol{))}\AgdaSpace{}%
\AgdaOperator{\AgdaDatatype{↝βb₂}}\AgdaSpace{}%
\AgdaBound{x}\AgdaSymbol{)}\<%
\\
%
\>[6]\AgdaSymbol{(}\AgdaFunction{pres-subst-func-rename}\AgdaSpace{}%
\AgdaBound{fb'}\AgdaSpace{}%
\AgdaBound{M'}\AgdaSpace{}%
\AgdaBound{W'}\AgdaSymbol{)}\<%
\\
%
\>[6]\AgdaInductiveConstructor{βb₂}\<%
\\
%
\>[4]\AgdaSymbol{)}\<%
\\
\>[0]\AgdaFunction{pres-rename-red-↝βb²}\AgdaSpace{}%
\AgdaBound{fb'}\AgdaSpace{}%
\AgdaBound{M}\AgdaSpace{}%
\AgdaBound{N}\AgdaSpace{}%
\AgdaSymbol{(}\AgdaInductiveConstructor{μ}\AgdaSpace{}%
\AgdaSymbol{\{}\AgdaArgument{M}\AgdaSpace{}%
\AgdaSymbol{=}\AgdaSpace{}%
\AgdaBound{M₁}\AgdaSymbol{\}}\AgdaSpace{}%
\AgdaSymbol{\{}\AgdaArgument{M'}\AgdaSpace{}%
\AgdaSymbol{=}\AgdaSpace{}%
\AgdaBound{M'}\AgdaSymbol{\}}\AgdaSpace{}%
\AgdaBound{red}\AgdaSymbol{)}\AgdaSpace{}%
\AgdaSymbol{=}\AgdaSpace{}%
\AgdaInductiveConstructor{μ}\AgdaSpace{}%
\AgdaSymbol{(}\AgdaFunction{pres-rename-red-↝βb²}\AgdaSpace{}%
\AgdaBound{fb'}\AgdaSpace{}%
\AgdaBound{M₁}\AgdaSpace{}%
\AgdaBound{M'}\AgdaSpace{}%
\AgdaBound{red}\AgdaSymbol{)}\<%
\\
\>[0]\AgdaFunction{pres-rename-red-↝βb²}\AgdaSpace{}%
\AgdaBound{fb'}\AgdaSpace{}%
\AgdaBound{M}\AgdaSpace{}%
\AgdaBound{N}\AgdaSpace{}%
\AgdaSymbol{(}\AgdaInductiveConstructor{ν}\AgdaSpace{}%
\AgdaSymbol{\{}\AgdaArgument{N}\AgdaSpace{}%
\AgdaSymbol{=}\AgdaSpace{}%
\AgdaBound{N₁}\AgdaSymbol{\}}\AgdaSpace{}%
\AgdaSymbol{\{}\AgdaArgument{N'}\AgdaSpace{}%
\AgdaSymbol{=}\AgdaSpace{}%
\AgdaBound{N'}\AgdaSymbol{\}}\AgdaSpace{}%
\AgdaBound{red}\AgdaSymbol{)}\AgdaSpace{}%
\AgdaSymbol{=}\AgdaSpace{}%
\AgdaInductiveConstructor{ν}\AgdaSpace{}%
\AgdaSymbol{(}\AgdaFunction{pres-rename-red-↝βb²}\AgdaSpace{}%
\AgdaBound{fb'}\AgdaSpace{}%
\AgdaBound{N₁}\AgdaSpace{}%
\AgdaBound{N'}\AgdaSpace{}%
\AgdaBound{red}\AgdaSymbol{)}\<%
\\
\>[0]\AgdaFunction{pres-rename-red-↝βb²}\AgdaSpace{}%
\AgdaBound{fb'}\AgdaSpace{}%
\AgdaBound{M}\AgdaSpace{}%
\AgdaBound{N}\AgdaSpace{}%
\AgdaSymbol{(}\AgdaInductiveConstructor{ξ}\AgdaSpace{}%
\AgdaSymbol{\{}\AgdaArgument{M}\AgdaSpace{}%
\AgdaSymbol{=}\AgdaSpace{}%
\AgdaBound{M₁}\AgdaSymbol{\}}\AgdaSpace{}%
\AgdaSymbol{\{}\AgdaArgument{M'}\AgdaSpace{}%
\AgdaSymbol{=}\AgdaSpace{}%
\AgdaBound{M'}\AgdaSymbol{\}}\AgdaSpace{}%
\AgdaBound{red}\AgdaSymbol{)}\AgdaSpace{}%
\AgdaSymbol{=}\AgdaSpace{}%
\AgdaInductiveConstructor{ξ}\AgdaSpace{}%
\AgdaSymbol{(}\AgdaFunction{pres-rename-red-↝βb²}\AgdaSpace{}%
\AgdaSymbol{(}\AgdaFunction{extb}\AgdaSpace{}%
\AgdaBound{fb'}\AgdaSymbol{)}\AgdaSpace{}%
\AgdaBound{M₁}\AgdaSpace{}%
\AgdaBound{M'}\AgdaSpace{}%
\AgdaBound{red}\AgdaSymbol{)}\<%
\\
\>[0]\AgdaFunction{pres-rename-red-↝βb²}\AgdaSpace{}%
\AgdaBound{fb'}\AgdaSpace{}%
\AgdaBound{M}\AgdaSpace{}%
\AgdaBound{N}\AgdaSpace{}%
\AgdaSymbol{(}\AgdaInductiveConstructor{ζ}\AgdaSpace{}%
\AgdaSymbol{\{}\AgdaArgument{M}\AgdaSpace{}%
\AgdaSymbol{=}\AgdaSpace{}%
\AgdaBound{M₁}\AgdaSymbol{\}}\AgdaSpace{}%
\AgdaSymbol{\{}\AgdaArgument{M'}\AgdaSpace{}%
\AgdaSymbol{=}\AgdaSpace{}%
\AgdaBound{M'}\AgdaSymbol{\}}\AgdaSpace{}%
\AgdaBound{red}\AgdaSymbol{)}\AgdaSpace{}%
\AgdaSymbol{=}\AgdaSpace{}%
\AgdaInductiveConstructor{ζ}\AgdaSpace{}%
\AgdaSymbol{(}\AgdaFunction{pres-rename-red-↝βb²}\AgdaSpace{}%
\AgdaBound{fb'}\AgdaSpace{}%
\AgdaBound{M₁}\AgdaSpace{}%
\AgdaBound{M'}\AgdaSpace{}%
\AgdaBound{red}\AgdaSymbol{)}\<%
\\
%
\\[\AgdaEmptyExtraSkip]%
\>[0]\AgdaComment{--\ Preservation\ of\ reduction\ by\ Gödel's\ translation}\<%
\\
\>[0]\AgdaComment{--\ with\ the\ two\ step\ reduction\ rule}\<%
\\
\>[0]\AgdaFunction{pres-red↝βb²}\AgdaSpace{}%
\AgdaSymbol{:}\AgdaSpace{}%
\AgdaSymbol{∀}\AgdaSpace{}%
\AgdaSymbol{\{}\AgdaBound{Γ}\AgdaSpace{}%
\AgdaBound{A}\AgdaSymbol{\}}\<%
\\
\>[0][@{}l@{\AgdaIndent{0}}]%
\>[2]\AgdaSymbol{→}\AgdaSpace{}%
\AgdaSymbol{(}\AgdaBound{M}\AgdaSpace{}%
\AgdaSymbol{:}\AgdaSpace{}%
\AgdaBound{Γ}\AgdaSpace{}%
\AgdaOperator{\AgdaDatatype{⊢}}\AgdaSpace{}%
\AgdaBound{A}\AgdaSymbol{)}\<%
\\
%
\>[2]\AgdaSymbol{→}\AgdaSpace{}%
\AgdaSymbol{(}\AgdaBound{N}\AgdaSpace{}%
\AgdaSymbol{:}\AgdaSpace{}%
\AgdaBound{Γ}\AgdaSpace{}%
\AgdaOperator{\AgdaDatatype{⊢}}\AgdaSpace{}%
\AgdaBound{A}\AgdaSymbol{)}\<%
\\
%
\>[2]\AgdaSymbol{→}\AgdaSpace{}%
\AgdaSymbol{(}\AgdaBound{M}\AgdaSpace{}%
\AgdaOperator{\AgdaDatatype{↝βᵥ}}\AgdaSpace{}%
\AgdaBound{N}\AgdaSymbol{)}\<%
\\
%
\>[2]\AgdaSymbol{→}\AgdaSpace{}%
\AgdaSymbol{(}\AgdaFunction{embedTerm}\AgdaSpace{}%
\AgdaBound{M}\AgdaSpace{}%
\AgdaOperator{\AgdaDatatype{↝βb₂}}\AgdaSpace{}%
\AgdaFunction{embedTerm}\AgdaSpace{}%
\AgdaBound{N}\AgdaSymbol{)}\<%
\\
\>[0]\AgdaFunction{pres-red↝βb²}\AgdaSpace{}%
\AgdaSymbol{(}\AgdaInductiveConstructor{ƛ}\AgdaSpace{}%
\AgdaBound{M₁}\AgdaSpace{}%
\AgdaOperator{\AgdaInductiveConstructor{·}}\AgdaSpace{}%
\AgdaBound{M₂}\AgdaSymbol{)}\AgdaSpace{}%
\AgdaBound{N}\AgdaSpace{}%
\AgdaSymbol{(}\AgdaInductiveConstructor{βᵥ}\AgdaSpace{}%
\AgdaBound{x}\AgdaSymbol{)}\AgdaSpace{}%
\AgdaSymbol{=}\<%
\\
\>[0][@{}l@{\AgdaIndent{0}}]%
\>[2]\AgdaFunction{Eq.subst}\<%
\\
\>[2][@{}l@{\AgdaIndent{0}}]%
\>[4]\AgdaSymbol{(λ}\AgdaSpace{}%
\AgdaBound{t}\AgdaSpace{}%
\AgdaSymbol{→}\AgdaSpace{}%
\AgdaBound{t}\AgdaSpace{}%
\AgdaOperator{\AgdaDatatype{↝βb₂}}\AgdaSpace{}%
\AgdaFunction{embedTerm}\AgdaSpace{}%
\AgdaSymbol{(}\AgdaBound{M₁}\AgdaSpace{}%
\AgdaOperator{\AgdaFunction{[}}\AgdaSpace{}%
\AgdaBound{M₂}\AgdaSpace{}%
\AgdaOperator{\AgdaFunction{]}}\AgdaSymbol{))}\<%
\\
%
\>[4]\AgdaSymbol{(}\AgdaFunction{cong₂}\AgdaSpace{}%
\AgdaOperator{\AgdaInductiveConstructor{\AgdaUnderscore{}·b\AgdaUnderscore{}}}\AgdaSpace{}%
\AgdaSymbol{(}\AgdaFunction{cong}\AgdaSpace{}%
\AgdaInductiveConstructor{ƛb}\AgdaSpace{}%
\AgdaSymbol{(}\AgdaFunction{cong₂}\AgdaSpace{}%
\AgdaOperator{\AgdaInductiveConstructor{\AgdaUnderscore{}·b\AgdaUnderscore{}}}\AgdaSpace{}%
\AgdaInductiveConstructor{refl}\AgdaSpace{}%
\AgdaSymbol{(}\AgdaFunction{cong}\AgdaSpace{}%
\AgdaSymbol{(}\AgdaFunction{renameb}\AgdaSpace{}%
\AgdaInductiveConstructor{Sb}\AgdaSymbol{)}\AgdaSpace{}%
\AgdaSymbol{(}\AgdaFunction{box-embedValue}\AgdaSpace{}%
\AgdaBound{x}\AgdaSymbol{))))}\AgdaSpace{}%
\AgdaInductiveConstructor{refl}\AgdaSymbol{)}\<%
\\
%
\>[4]\AgdaSymbol{(}\AgdaFunction{Eq.subst}\<%
\\
\>[4][@{}l@{\AgdaIndent{0}}]%
\>[6]\AgdaSymbol{(λ}\AgdaSpace{}%
\AgdaBound{t}\AgdaSpace{}%
\AgdaSymbol{→}\AgdaSpace{}%
\AgdaFunction{raise}\AgdaSpace{}%
\AgdaSymbol{(}\AgdaSpace{}%
\AgdaInductiveConstructor{box}\AgdaSpace{}%
\AgdaSymbol{(}\AgdaFunction{embedValue}\AgdaSpace{}%
\AgdaBound{M₂}\AgdaSpace{}%
\AgdaBound{x}\AgdaSymbol{)}\AgdaSpace{}%
\AgdaSymbol{)}\AgdaSpace{}%
\AgdaOperator{\AgdaInductiveConstructor{·b}}\AgdaSpace{}%
\AgdaInductiveConstructor{box}\AgdaSpace{}%
\AgdaSymbol{(}\AgdaInductiveConstructor{ƛb}\AgdaSpace{}%
\AgdaSymbol{(}\AgdaFunction{embedTerm}\AgdaSpace{}%
\AgdaBound{M₁}\AgdaSymbol{))}\AgdaSpace{}%
\AgdaOperator{\AgdaDatatype{↝βb₂}}\AgdaSpace{}%
\AgdaBound{t}\AgdaSymbol{)}\<%
\\
%
\>[6]\AgdaSymbol{(}\AgdaFunction{pres-subst}\AgdaSpace{}%
\AgdaBound{M₁}\AgdaSpace{}%
\AgdaBound{M₂}\AgdaSpace{}%
\AgdaBound{x}\AgdaSymbol{)}\<%
\\
%
\>[6]\AgdaInductiveConstructor{βb₂}\<%
\\
%
\>[4]\AgdaSymbol{)}\<%
\\
\>[0]\AgdaFunction{pres-red↝βb²}\AgdaSpace{}%
\AgdaBound{M}\AgdaSpace{}%
\AgdaBound{N}\AgdaSpace{}%
\AgdaSymbol{(}\AgdaInductiveConstructor{μ}\AgdaSpace{}%
\AgdaSymbol{\{}\AgdaArgument{M}\AgdaSpace{}%
\AgdaSymbol{=}\AgdaSpace{}%
\AgdaBound{M₁}\AgdaSymbol{\}}\AgdaSpace{}%
\AgdaSymbol{\{}\AgdaArgument{M'}\AgdaSpace{}%
\AgdaSymbol{=}\AgdaSpace{}%
\AgdaBound{M'}\AgdaSymbol{\}}\AgdaSpace{}%
\AgdaBound{red}\AgdaSymbol{)}\AgdaSpace{}%
\AgdaSymbol{=}\AgdaSpace{}%
\AgdaInductiveConstructor{ν}\AgdaSpace{}%
\AgdaSymbol{(}\AgdaFunction{pres-red↝βb²}\AgdaSpace{}%
\AgdaBound{M₁}\AgdaSpace{}%
\AgdaBound{M'}\AgdaSpace{}%
\AgdaBound{red}\AgdaSymbol{)}\<%
\\
\>[0]\AgdaFunction{pres-red↝βb²}\AgdaSpace{}%
\AgdaBound{M}\AgdaSpace{}%
\AgdaBound{N}\AgdaSpace{}%
\AgdaSymbol{(}\AgdaInductiveConstructor{ν<}\AgdaSpace{}%
\AgdaSymbol{\{}\AgdaArgument{N}\AgdaSpace{}%
\AgdaSymbol{=}\AgdaSpace{}%
\AgdaBound{N₁}\AgdaSymbol{\}}\AgdaSpace{}%
\AgdaSymbol{\{}\AgdaArgument{N'}\AgdaSpace{}%
\AgdaSymbol{=}\AgdaSpace{}%
\AgdaBound{N'}\AgdaSymbol{\}}\AgdaSpace{}%
\AgdaBound{x}\AgdaSpace{}%
\AgdaBound{red}\AgdaSymbol{)}\AgdaSpace{}%
\AgdaSymbol{=}\AgdaSpace{}%
\AgdaInductiveConstructor{μ}\AgdaSpace{}%
\AgdaSymbol{(}\AgdaInductiveConstructor{ξ}\AgdaSpace{}%
\AgdaSymbol{(}\AgdaInductiveConstructor{ν}\AgdaSpace{}%
\AgdaSymbol{(}\AgdaFunction{pres-rename-red-↝βb²}\AgdaSpace{}%
\AgdaInductiveConstructor{Sb}\AgdaSpace{}%
\AgdaSymbol{(}\AgdaFunction{embedTerm}\AgdaSpace{}%
\AgdaBound{N₁}\AgdaSymbol{)}\AgdaSpace{}%
\AgdaSymbol{(}\AgdaFunction{embedTerm}\AgdaSpace{}%
\AgdaBound{N'}\AgdaSymbol{)}\AgdaSpace{}%
\AgdaSymbol{(}\AgdaFunction{pres-red↝βb²}\AgdaSpace{}%
\AgdaBound{N₁}\AgdaSpace{}%
\AgdaBound{N'}\AgdaSpace{}%
\AgdaBound{red}\AgdaSymbol{))))}\<%
\\
\>[0]\AgdaFunction{pres-red↝βb²}\AgdaSpace{}%
\AgdaBound{M}\AgdaSpace{}%
\AgdaBound{N}\AgdaSpace{}%
\AgdaSymbol{(}\AgdaInductiveConstructor{ν}\AgdaSpace{}%
\AgdaSymbol{\{}\AgdaArgument{N}\AgdaSpace{}%
\AgdaSymbol{=}\AgdaSpace{}%
\AgdaBound{N₁}\AgdaSymbol{\}}\AgdaSpace{}%
\AgdaSymbol{\{}\AgdaArgument{N'}\AgdaSpace{}%
\AgdaSymbol{=}\AgdaSpace{}%
\AgdaBound{N'}\AgdaSymbol{\}}\AgdaSpace{}%
\AgdaBound{red}\AgdaSymbol{)}\AgdaSpace{}%
\AgdaSymbol{=}\AgdaSpace{}%
\AgdaInductiveConstructor{μ}\AgdaSpace{}%
\AgdaSymbol{(}\AgdaInductiveConstructor{ξ}\AgdaSpace{}%
\AgdaSymbol{(}\AgdaInductiveConstructor{ν}\AgdaSpace{}%
\AgdaSymbol{(}\AgdaFunction{pres-rename-red-↝βb²}\AgdaSpace{}%
\AgdaInductiveConstructor{Sb}\AgdaSpace{}%
\AgdaSymbol{(}\AgdaFunction{embedTerm}\AgdaSpace{}%
\AgdaBound{N₁}\AgdaSymbol{)}\AgdaSpace{}%
\AgdaSymbol{(}\AgdaFunction{embedTerm}\AgdaSpace{}%
\AgdaBound{N'}\AgdaSymbol{)}\AgdaSpace{}%
\AgdaSymbol{(}\AgdaFunction{pres-red↝βb²}\AgdaSpace{}%
\AgdaBound{N₁}\AgdaSpace{}%
\AgdaBound{N'}\AgdaSpace{}%
\AgdaBound{red}\AgdaSymbol{))))}\<%
\\
\>[0]\AgdaFunction{pres-red↝βb²}\AgdaSpace{}%
\AgdaBound{M}\AgdaSpace{}%
\AgdaBound{N}\AgdaSpace{}%
\AgdaSymbol{(}\AgdaInductiveConstructor{ξ}\AgdaSpace{}%
\AgdaSymbol{\{}\AgdaArgument{M}\AgdaSpace{}%
\AgdaSymbol{=}\AgdaSpace{}%
\AgdaBound{M₁}\AgdaSymbol{\}}\AgdaSpace{}%
\AgdaSymbol{\{}\AgdaArgument{M'}\AgdaSpace{}%
\AgdaSymbol{=}\AgdaSpace{}%
\AgdaBound{M'}\AgdaSymbol{\}}\AgdaSpace{}%
\AgdaBound{red}\AgdaSymbol{)}\AgdaSpace{}%
\AgdaSymbol{=}\AgdaSpace{}%
\AgdaInductiveConstructor{ζ}\AgdaSpace{}%
\AgdaSymbol{(}\AgdaInductiveConstructor{ξ}\AgdaSpace{}%
\AgdaSymbol{(}\AgdaFunction{pres-red↝βb²}\AgdaSpace{}%
\AgdaBound{M₁}\AgdaSpace{}%
\AgdaBound{M'}\AgdaSpace{}%
\AgdaBound{red}\AgdaSymbol{))}\<%
\end{code}  
\end{lemma}

\subsection{Proofs of the Main Propositions}
Here we give the proofs of \cref{prop:girard-preservation-red} and \cref{prop:girard-preservation-eval}.
These proofs are not very difficult and depend on the preservation of substitution, which is given in \cref{lem:girard-pres-subst}.
This proof was a lot more complicated. 

\begin{proofthm}
  \label{proof:girard-preservation-red}
  \textbf{Preservation of reduction} under Girard's translation.
  \small
  \begin{code}[hide]%
\>[0]\AgdaKeyword{module}\AgdaSpace{}%
\AgdaModule{nameProofs3}\AgdaSpace{}%
\AgdaKeyword{where}\<%
\\
%
\\[\AgdaEmptyExtraSkip]%
\>[0]\AgdaKeyword{open}\AgdaSpace{}%
\AgdaKeyword{import}\AgdaSpace{}%
\AgdaModule{boxCalc}\<%
\\
\>[0]\AgdaKeyword{open}\AgdaSpace{}%
\AgdaKeyword{import}\AgdaSpace{}%
\AgdaModule{calc}\<%
\\
\>[0]\AgdaKeyword{open}\AgdaSpace{}%
\AgdaKeyword{import}\AgdaSpace{}%
\AgdaModule{embedCBNIntoCBB}\<%
\\
\>[0]\AgdaKeyword{open}\AgdaSpace{}%
\AgdaKeyword{import}\AgdaSpace{}%
\AgdaModule{nameCalcRed}\<%
\\
\>[0]\AgdaKeyword{open}\AgdaSpace{}%
\AgdaKeyword{import}\AgdaSpace{}%
\AgdaModule{boxCalcRed}\<%
\\
%
\\[\AgdaEmptyExtraSkip]%
\>[0]\AgdaKeyword{open}\AgdaSpace{}%
\AgdaKeyword{import}\AgdaSpace{}%
\AgdaModule{Data.Empty}\AgdaSpace{}%
\AgdaKeyword{using}\AgdaSpace{}%
\AgdaSymbol{(}\AgdaFunction{⊥}\AgdaSymbol{;}\AgdaSpace{}%
\AgdaFunction{⊥-elim}\AgdaSymbol{)}\<%
\\
\>[0]\AgdaKeyword{import}\AgdaSpace{}%
\AgdaModule{Relation.Binary.PropositionalEquality}\AgdaSpace{}%
\AgdaSymbol{as}\AgdaSpace{}%
\AgdaModule{Eq}\<%
\\
\>[0]\AgdaKeyword{open}\AgdaSpace{}%
\AgdaModule{Eq}\AgdaSpace{}%
\AgdaKeyword{using}\AgdaSpace{}%
\AgdaSymbol{(}\AgdaOperator{\AgdaDatatype{\AgdaUnderscore{}≡\AgdaUnderscore{}}}\AgdaSymbol{;}\AgdaSpace{}%
\AgdaInductiveConstructor{refl}\AgdaSymbol{;}\AgdaSpace{}%
\AgdaFunction{cong}\AgdaSymbol{;}\AgdaSpace{}%
\AgdaFunction{cong₂}\AgdaSymbol{;}\AgdaSpace{}%
\AgdaFunction{sym}\AgdaSymbol{)}\<%
\\
%
\\[\AgdaEmptyExtraSkip]%
\>[0]\AgdaKeyword{open}\AgdaSpace{}%
\AgdaModule{Eq.≡-Reasoning}\AgdaSpace{}%
\AgdaKeyword{using}\AgdaSpace{}%
\AgdaSymbol{(}\AgdaOperator{\AgdaFunction{begin\AgdaUnderscore{}}}\AgdaSymbol{;}\AgdaSpace{}%
\AgdaFunction{step-≡-∣}\AgdaSymbol{;}\AgdaSpace{}%
\AgdaFunction{step-≡-⟩}\AgdaSymbol{;}\AgdaSpace{}%
\AgdaOperator{\AgdaFunction{\AgdaUnderscore{}∎}}\AgdaSymbol{)}\<%
\\
\>[0]\AgdaKeyword{open}\AgdaSpace{}%
\AgdaModule{Eq.≡-Reasoning}\AgdaSpace{}%
\AgdaKeyword{renaming}\AgdaSpace{}%
\AgdaSymbol{(}\AgdaOperator{\AgdaFunction{begin\AgdaUnderscore{}}}\AgdaSpace{}%
\AgdaSymbol{to}\AgdaSpace{}%
\AgdaOperator{\AgdaFunction{beginEq\AgdaUnderscore{}}}\AgdaSymbol{;}\AgdaSpace{}%
\AgdaOperator{\AgdaFunction{\AgdaUnderscore{}∎}}\AgdaSpace{}%
\AgdaSymbol{to}\AgdaSpace{}%
\AgdaOperator{\AgdaFunction{\AgdaUnderscore{}∎eq}}\AgdaSymbol{)}\<%
\\
\>[0]\AgdaKeyword{open}\AgdaSpace{}%
\AgdaKeyword{import}\AgdaSpace{}%
\AgdaModule{Function.Bundles}\AgdaSpace{}%
\AgdaKeyword{using}\AgdaSpace{}%
\AgdaSymbol{(}\AgdaOperator{\AgdaFunction{\AgdaUnderscore{}⇔\AgdaUnderscore{}}}\AgdaSymbol{;}\AgdaSpace{}%
\AgdaFunction{mk⇔}\AgdaSymbol{)}\<%
\\
%
\\[\AgdaEmptyExtraSkip]%
\>[0]\AgdaComment{--\ GIRARD'S\ embedding}\<%
\\
%
\\[\AgdaEmptyExtraSkip]%
\>[0]\AgdaComment{--\ Preservation\ of\ typing\ by\ Girard's\ translation}\<%
\\
\>[0]\AgdaFunction{pres-typing}\AgdaSpace{}%
\AgdaSymbol{:}\AgdaSpace{}%
\AgdaSymbol{∀}\AgdaSpace{}%
\AgdaSymbol{\{}\AgdaBound{Γ}\AgdaSpace{}%
\AgdaBound{A}\AgdaSymbol{\}}\AgdaSpace{}%
\AgdaSymbol{→}\AgdaSpace{}%
\AgdaBound{Γ}\AgdaSpace{}%
\AgdaOperator{\AgdaDatatype{⊢}}\AgdaSpace{}%
\AgdaBound{A}\AgdaSpace{}%
\AgdaSymbol{→}\AgdaSpace{}%
\AgdaFunction{embedContext}\AgdaSpace{}%
\AgdaBound{Γ}\AgdaSpace{}%
\AgdaOperator{\AgdaDatatype{⊢b}}\AgdaSpace{}%
\AgdaFunction{embedType}\AgdaSpace{}%
\AgdaBound{A}\<%
\\
\>[0]\AgdaFunction{pres-typing}\AgdaSpace{}%
\AgdaBound{x}\AgdaSpace{}%
\AgdaSymbol{=}\AgdaSpace{}%
\AgdaFunction{embedTerm}\AgdaSpace{}%
\AgdaBound{x}\<%
\\
%
\\[\AgdaEmptyExtraSkip]%
\>[0]\AgdaFunction{pres-ext}\AgdaSpace{}%
\AgdaSymbol{:}\AgdaSpace{}%
\AgdaSymbol{∀}\AgdaSpace{}%
\AgdaSymbol{\{}\AgdaBound{Γ}\AgdaSpace{}%
\AgdaBound{Δ}\AgdaSpace{}%
\AgdaBound{A'}\AgdaSpace{}%
\AgdaBound{A}\AgdaSymbol{\}}\<%
\\
\>[0][@{}l@{\AgdaIndent{0}}]%
\>[2]\AgdaSymbol{→}\AgdaSpace{}%
\AgdaSymbol{(}\AgdaBound{f'}\AgdaSpace{}%
\AgdaSymbol{:}\AgdaSpace{}%
\AgdaSymbol{∀}\AgdaSpace{}%
\AgdaSymbol{\{}\AgdaBound{A}\AgdaSymbol{\}}\AgdaSpace{}%
\AgdaSymbol{→}\AgdaSpace{}%
\AgdaBound{A}\AgdaSpace{}%
\AgdaOperator{\AgdaDatatype{∈}}\AgdaSpace{}%
\AgdaBound{Γ}\AgdaSpace{}%
\AgdaSymbol{→}\AgdaSpace{}%
\AgdaBound{A}\AgdaSpace{}%
\AgdaOperator{\AgdaDatatype{∈}}\AgdaSpace{}%
\AgdaBound{Δ}\AgdaSymbol{)}\<%
\\
%
\>[2]\AgdaSymbol{→}\AgdaSpace{}%
\AgdaSymbol{(}\AgdaBound{fb'}\AgdaSpace{}%
\AgdaSymbol{:}\AgdaSpace{}%
\AgdaSymbol{∀}\AgdaSpace{}%
\AgdaSymbol{\{}\AgdaBound{Ab}\AgdaSymbol{\}}\AgdaSpace{}%
\AgdaSymbol{→}\AgdaSpace{}%
\AgdaBound{Ab}\AgdaSpace{}%
\AgdaOperator{\AgdaDatatype{∈b}}\AgdaSpace{}%
\AgdaFunction{embedContext}\AgdaSpace{}%
\AgdaBound{Γ}\AgdaSpace{}%
\AgdaSymbol{→}\AgdaSpace{}%
\AgdaBound{Ab}\AgdaSpace{}%
\AgdaOperator{\AgdaDatatype{∈b}}\AgdaSpace{}%
\AgdaFunction{embedContext}\AgdaSpace{}%
\AgdaBound{Δ}\AgdaSymbol{)}\<%
\\
%
\>[2]\AgdaSymbol{→}\AgdaSpace{}%
\AgdaSymbol{(\{}\AgdaBound{A'}\AgdaSpace{}%
\AgdaSymbol{:}\AgdaSpace{}%
\AgdaDatatype{Type}\AgdaSymbol{\}}\AgdaSpace{}%
\AgdaSymbol{(}\AgdaBound{x}\AgdaSpace{}%
\AgdaSymbol{:}\AgdaSpace{}%
\AgdaBound{A'}\AgdaSpace{}%
\AgdaOperator{\AgdaDatatype{∈}}\AgdaSpace{}%
\AgdaBound{Γ}\AgdaSymbol{)}\AgdaSpace{}%
\AgdaSymbol{→}\AgdaSpace{}%
\AgdaBound{fb'}\AgdaSpace{}%
\AgdaSymbol{(}\AgdaFunction{embedLookup}\AgdaSpace{}%
\AgdaBound{x}\AgdaSymbol{)}\AgdaSpace{}%
\AgdaOperator{\AgdaDatatype{≡}}\AgdaSpace{}%
\AgdaFunction{embedLookup}\AgdaSpace{}%
\AgdaSymbol{(}\AgdaBound{f'}\AgdaSpace{}%
\AgdaBound{x}\AgdaSymbol{))}\<%
\\
%
\>[2]\AgdaSymbol{→}\AgdaSpace{}%
\AgdaSymbol{(}\AgdaBound{x}\AgdaSpace{}%
\AgdaSymbol{:}\AgdaSpace{}%
\AgdaBound{A'}\AgdaSpace{}%
\AgdaOperator{\AgdaDatatype{∈}}\AgdaSpace{}%
\AgdaBound{Γ}\AgdaSpace{}%
\AgdaOperator{\AgdaInductiveConstructor{,}}\AgdaSpace{}%
\AgdaBound{A}\AgdaSymbol{)}\<%
\\
%
\>[2]\AgdaSymbol{→}\AgdaSpace{}%
\AgdaFunction{extb}\AgdaSpace{}%
\AgdaBound{fb'}\AgdaSpace{}%
\AgdaSymbol{(}\AgdaFunction{embedLookup}\AgdaSpace{}%
\AgdaBound{x}\AgdaSymbol{)}\AgdaSpace{}%
\AgdaOperator{\AgdaDatatype{≡}}\AgdaSpace{}%
\AgdaFunction{embedLookup}\AgdaSpace{}%
\AgdaSymbol{(}\AgdaFunction{ext}\AgdaSpace{}%
\AgdaBound{f'}\AgdaSpace{}%
\AgdaBound{x}\AgdaSymbol{)}\<%
\\
\>[0]\AgdaFunction{pres-ext}\AgdaSpace{}%
\AgdaBound{f'}\AgdaSpace{}%
\AgdaBound{fb'}\AgdaSpace{}%
\AgdaBound{h}\AgdaSpace{}%
\AgdaInductiveConstructor{Z}\AgdaSpace{}%
\AgdaSymbol{=}\AgdaSpace{}%
\AgdaInductiveConstructor{refl}\<%
\\
\>[0]\AgdaFunction{pres-ext}\AgdaSpace{}%
\AgdaBound{f'}\AgdaSpace{}%
\AgdaBound{fb'}\AgdaSpace{}%
\AgdaBound{h}\AgdaSpace{}%
\AgdaSymbol{(}\AgdaInductiveConstructor{S}\AgdaSpace{}%
\AgdaBound{x}\AgdaSymbol{)}\AgdaSpace{}%
\AgdaSymbol{=}\AgdaSpace{}%
\AgdaFunction{cong}\AgdaSpace{}%
\AgdaInductiveConstructor{Sb}\AgdaSpace{}%
\AgdaSymbol{(}\AgdaBound{h}\AgdaSpace{}%
\AgdaBound{x}\AgdaSymbol{)}\<%
\\
%
\\[\AgdaEmptyExtraSkip]%
\>[0]\AgdaFunction{pres-rename}\AgdaSpace{}%
\AgdaSymbol{:}\AgdaSpace{}%
\AgdaSymbol{∀}\AgdaSpace{}%
\AgdaSymbol{\{}\AgdaBound{Γ}\AgdaSpace{}%
\AgdaBound{A}\AgdaSpace{}%
\AgdaBound{Δ}\AgdaSymbol{\}}\<%
\\
\>[0][@{}l@{\AgdaIndent{0}}]%
\>[2]\AgdaSymbol{→}\AgdaSpace{}%
\AgdaSymbol{(}\AgdaBound{f'}\AgdaSpace{}%
\AgdaSymbol{:}\AgdaSpace{}%
\AgdaSymbol{∀}\AgdaSpace{}%
\AgdaSymbol{\{}\AgdaBound{A}\AgdaSymbol{\}}\AgdaSpace{}%
\AgdaSymbol{→}\AgdaSpace{}%
\AgdaBound{A}\AgdaSpace{}%
\AgdaOperator{\AgdaDatatype{∈}}\AgdaSpace{}%
\AgdaBound{Γ}\AgdaSpace{}%
\AgdaSymbol{→}\AgdaSpace{}%
\AgdaBound{A}\AgdaSpace{}%
\AgdaOperator{\AgdaDatatype{∈}}\AgdaSpace{}%
\AgdaBound{Δ}\AgdaSymbol{)}\<%
\\
%
\>[2]\AgdaSymbol{→}\AgdaSpace{}%
\AgdaSymbol{(}\AgdaBound{fb'}\AgdaSpace{}%
\AgdaSymbol{:}\AgdaSpace{}%
\AgdaSymbol{∀}\AgdaSpace{}%
\AgdaSymbol{\{}\AgdaBound{Ab}\AgdaSymbol{\}}\AgdaSpace{}%
\AgdaSymbol{→}\AgdaSpace{}%
\AgdaBound{Ab}\AgdaSpace{}%
\AgdaOperator{\AgdaDatatype{∈b}}\AgdaSpace{}%
\AgdaFunction{embedContext}\AgdaSpace{}%
\AgdaBound{Γ}\AgdaSpace{}%
\AgdaSymbol{→}\AgdaSpace{}%
\AgdaBound{Ab}\AgdaSpace{}%
\AgdaOperator{\AgdaDatatype{∈b}}\AgdaSpace{}%
\AgdaFunction{embedContext}\AgdaSpace{}%
\AgdaBound{Δ}\AgdaSymbol{)}\<%
\\
%
\>[2]\AgdaSymbol{→}\AgdaSpace{}%
\AgdaSymbol{(\{}\AgdaBound{A'}\AgdaSpace{}%
\AgdaSymbol{:}\AgdaSpace{}%
\AgdaDatatype{Type}\AgdaSymbol{\}}\AgdaSpace{}%
\AgdaSymbol{(}\AgdaBound{x}\AgdaSpace{}%
\AgdaSymbol{:}\AgdaSpace{}%
\AgdaBound{A'}\AgdaSpace{}%
\AgdaOperator{\AgdaDatatype{∈}}\AgdaSpace{}%
\AgdaBound{Γ}\AgdaSymbol{)}\AgdaSpace{}%
\AgdaSymbol{→}\AgdaSpace{}%
\AgdaBound{fb'}\AgdaSpace{}%
\AgdaSymbol{(}\AgdaFunction{embedLookup}\AgdaSpace{}%
\AgdaBound{x}\AgdaSymbol{)}\AgdaSpace{}%
\AgdaOperator{\AgdaDatatype{≡}}\AgdaSpace{}%
\AgdaFunction{embedLookup}\AgdaSpace{}%
\AgdaSymbol{(}\AgdaBound{f'}\AgdaSpace{}%
\AgdaBound{x}\AgdaSymbol{))}\<%
\\
%
\>[2]\AgdaSymbol{→}\AgdaSpace{}%
\AgdaSymbol{(}\AgdaBound{t}\AgdaSpace{}%
\AgdaSymbol{:}\AgdaSpace{}%
\AgdaBound{Γ}\AgdaSpace{}%
\AgdaOperator{\AgdaDatatype{⊢}}\AgdaSpace{}%
\AgdaBound{A}\AgdaSymbol{)}\<%
\\
%
\>[2]\AgdaSymbol{→}\AgdaSpace{}%
\AgdaFunction{renameb}\AgdaSpace{}%
\AgdaBound{fb'}\AgdaSpace{}%
\AgdaSymbol{(}\AgdaFunction{embedTerm}\AgdaSpace{}%
\AgdaBound{t}\AgdaSymbol{)}\AgdaSpace{}%
\AgdaOperator{\AgdaDatatype{≡}}\AgdaSpace{}%
\AgdaFunction{embedTerm}\AgdaSpace{}%
\AgdaSymbol{(}\AgdaFunction{rename}\AgdaSpace{}%
\AgdaBound{f'}\AgdaSpace{}%
\AgdaBound{t}\AgdaSymbol{)}\<%
\\
\>[0]\AgdaFunction{pres-rename}\AgdaSpace{}%
\AgdaBound{f'}\AgdaSpace{}%
\AgdaBound{fb'}\AgdaSpace{}%
\AgdaBound{h}\AgdaSpace{}%
\AgdaSymbol{(}\AgdaInductiveConstructor{`}\AgdaSpace{}%
\AgdaBound{x}\AgdaSymbol{)}\AgdaSpace{}%
\AgdaSymbol{=}\<%
\\
\>[0][@{}l@{\AgdaIndent{0}}]%
\>[2]\AgdaFunction{cong}\AgdaSpace{}%
\AgdaInductiveConstructor{ε}\AgdaSpace{}%
\AgdaSymbol{(}\AgdaBound{h}\AgdaSpace{}%
\AgdaBound{x}\AgdaSymbol{)}\<%
\\
\>[0]\AgdaFunction{pres-rename}\AgdaSpace{}%
\AgdaBound{f'}\AgdaSpace{}%
\AgdaBound{fb'}\AgdaSpace{}%
\AgdaBound{h}\AgdaSpace{}%
\AgdaSymbol{(}\AgdaInductiveConstructor{ƛ}\AgdaSpace{}%
\AgdaBound{t}\AgdaSymbol{)}\AgdaSpace{}%
\AgdaSymbol{=}\<%
\\
\>[0][@{}l@{\AgdaIndent{0}}]%
\>[2]\AgdaFunction{cong}\AgdaSpace{}%
\AgdaInductiveConstructor{ƛb}\AgdaSpace{}%
\AgdaSymbol{(}\AgdaFunction{pres-rename}\AgdaSpace{}%
\AgdaSymbol{(}\AgdaFunction{calc.ext}\AgdaSpace{}%
\AgdaBound{f'}\AgdaSymbol{)}\AgdaSpace{}%
\AgdaSymbol{(}\AgdaFunction{extb}\AgdaSpace{}%
\AgdaBound{fb'}\AgdaSymbol{)}\AgdaSpace{}%
\AgdaSymbol{(}\AgdaFunction{pres-ext}\AgdaSpace{}%
\AgdaBound{f'}\AgdaSpace{}%
\AgdaBound{fb'}\AgdaSpace{}%
\AgdaBound{h}\AgdaSymbol{)}\AgdaSpace{}%
\AgdaBound{t}\AgdaSymbol{)}\<%
\\
\>[0]\AgdaFunction{pres-rename}\AgdaSpace{}%
\AgdaBound{f'}\AgdaSpace{}%
\AgdaBound{fb'}\AgdaSpace{}%
\AgdaBound{h}\AgdaSpace{}%
\AgdaSymbol{(}\AgdaBound{t}\AgdaSpace{}%
\AgdaOperator{\AgdaInductiveConstructor{·}}\AgdaSpace{}%
\AgdaBound{t₁}\AgdaSymbol{)}\AgdaSpace{}%
\AgdaSymbol{=}\<%
\\
\>[0][@{}l@{\AgdaIndent{0}}]%
\>[2]\AgdaFunction{cong₂}\AgdaSpace{}%
\AgdaOperator{\AgdaInductiveConstructor{\AgdaUnderscore{}·b\AgdaUnderscore{}}}\AgdaSpace{}%
\AgdaSymbol{(}\AgdaFunction{pres-rename}\AgdaSpace{}%
\AgdaBound{f'}\AgdaSpace{}%
\AgdaBound{fb'}\AgdaSpace{}%
\AgdaBound{h}\AgdaSpace{}%
\AgdaBound{t}\AgdaSymbol{)}\AgdaSpace{}%
\AgdaSymbol{(}\AgdaFunction{cong}\AgdaSpace{}%
\AgdaInductiveConstructor{box}\AgdaSpace{}%
\AgdaSymbol{(}\AgdaFunction{pres-rename}\AgdaSpace{}%
\AgdaBound{f'}\AgdaSpace{}%
\AgdaBound{fb'}\AgdaSpace{}%
\AgdaBound{h}\AgdaSpace{}%
\AgdaBound{t₁}\AgdaSymbol{))}\<%
\\
%
\\[\AgdaEmptyExtraSkip]%
%
\\[\AgdaEmptyExtraSkip]%
\>[0]\AgdaFunction{pres-exts}\AgdaSpace{}%
\AgdaSymbol{:}\AgdaSpace{}%
\AgdaSymbol{∀}\AgdaSpace{}%
\AgdaSymbol{\{}\AgdaBound{Γ}\AgdaSpace{}%
\AgdaBound{Δ}\AgdaSymbol{\}}\<%
\\
\>[0][@{}l@{\AgdaIndent{0}}]%
\>[2]\AgdaSymbol{→}\AgdaSpace{}%
\AgdaSymbol{(}\AgdaBound{fb'}\AgdaSpace{}%
\AgdaSymbol{:}\AgdaSpace{}%
\AgdaSymbol{∀}\AgdaSpace{}%
\AgdaSymbol{\{}\AgdaBound{Ab}\AgdaSymbol{\}}\AgdaSpace{}%
\AgdaSymbol{→}\AgdaSpace{}%
\AgdaInductiveConstructor{□}\AgdaSpace{}%
\AgdaBound{Ab}\AgdaSpace{}%
\AgdaOperator{\AgdaDatatype{∈b}}\AgdaSpace{}%
\AgdaFunction{embedContext}\AgdaSpace{}%
\AgdaBound{Γ}\AgdaSpace{}%
\AgdaSymbol{→}\AgdaSpace{}%
\AgdaFunction{embedContext}\AgdaSpace{}%
\AgdaBound{Δ}\AgdaSpace{}%
\AgdaOperator{\AgdaDatatype{⊢b}}\AgdaSpace{}%
\AgdaBound{Ab}\AgdaSymbol{)}\<%
\\
%
\>[2]\AgdaSymbol{→}\AgdaSpace{}%
\AgdaSymbol{(}\AgdaBound{f'}\AgdaSpace{}%
\AgdaSymbol{:}\AgdaSpace{}%
\AgdaSymbol{∀}\AgdaSpace{}%
\AgdaSymbol{\{}\AgdaBound{A}\AgdaSymbol{\}}\AgdaSpace{}%
\AgdaSymbol{→}\AgdaSpace{}%
\AgdaBound{A}\AgdaSpace{}%
\AgdaOperator{\AgdaDatatype{∈}}\AgdaSpace{}%
\AgdaBound{Γ}\AgdaSpace{}%
\AgdaSymbol{→}\AgdaSpace{}%
\AgdaBound{Δ}\AgdaSpace{}%
\AgdaOperator{\AgdaDatatype{⊢}}\AgdaSpace{}%
\AgdaBound{A}\AgdaSymbol{)}\<%
\\
%
\>[2]\AgdaSymbol{→}\AgdaSpace{}%
\AgdaSymbol{(\{}\AgdaBound{A'}\AgdaSpace{}%
\AgdaSymbol{:}\AgdaSpace{}%
\AgdaDatatype{Type}\AgdaSymbol{\}}\AgdaSpace{}%
\AgdaSymbol{(}\AgdaBound{x}\AgdaSpace{}%
\AgdaSymbol{:}\AgdaSpace{}%
\AgdaBound{A'}\AgdaSpace{}%
\AgdaOperator{\AgdaDatatype{∈}}\AgdaSpace{}%
\AgdaBound{Γ}\AgdaSymbol{)}\AgdaSpace{}%
\AgdaSymbol{→}\AgdaSpace{}%
\AgdaBound{fb'}\AgdaSpace{}%
\AgdaSymbol{(}\AgdaFunction{embedLookup}\AgdaSpace{}%
\AgdaBound{x}\AgdaSymbol{)}\AgdaSpace{}%
\AgdaOperator{\AgdaDatatype{≡}}\AgdaSpace{}%
\AgdaFunction{embedTerm}\AgdaSpace{}%
\AgdaSymbol{(}\AgdaBound{f'}\AgdaSpace{}%
\AgdaBound{x}\AgdaSymbol{))}\<%
\\
%
\>[2]\AgdaSymbol{→}\AgdaSpace{}%
\AgdaSymbol{\{}\AgdaBound{A'}\AgdaSpace{}%
\AgdaBound{A}\AgdaSpace{}%
\AgdaSymbol{:}\AgdaSpace{}%
\AgdaDatatype{Type}\AgdaSymbol{\}}\AgdaSpace{}%
\AgdaSymbol{(}\AgdaBound{x}\AgdaSpace{}%
\AgdaSymbol{:}\AgdaSpace{}%
\AgdaBound{A'}\AgdaSpace{}%
\AgdaOperator{\AgdaDatatype{∈}}\AgdaSpace{}%
\AgdaBound{Γ}\AgdaSpace{}%
\AgdaOperator{\AgdaInductiveConstructor{,}}\AgdaSpace{}%
\AgdaBound{A}\AgdaSymbol{)}\AgdaSpace{}%
\AgdaSymbol{→}\AgdaSpace{}%
\AgdaFunction{extsb}\AgdaSpace{}%
\AgdaBound{fb'}\AgdaSpace{}%
\AgdaSymbol{(}\AgdaFunction{embedLookup}\AgdaSpace{}%
\AgdaBound{x}\AgdaSymbol{)}\AgdaSpace{}%
\AgdaOperator{\AgdaDatatype{≡}}\AgdaSpace{}%
\AgdaFunction{embedTerm}\AgdaSpace{}%
\AgdaSymbol{(}\AgdaFunction{exts}\AgdaSpace{}%
\AgdaBound{f'}\AgdaSpace{}%
\AgdaBound{x}\AgdaSymbol{)}\<%
\\
\>[0]\AgdaFunction{pres-exts}\AgdaSpace{}%
\AgdaBound{fb'}\AgdaSpace{}%
\AgdaBound{f'}\AgdaSpace{}%
\AgdaBound{h}\AgdaSpace{}%
\AgdaInductiveConstructor{Z}\AgdaSpace{}%
\AgdaSymbol{=}\AgdaSpace{}%
\AgdaInductiveConstructor{refl}\<%
\\
\>[0]\AgdaFunction{pres-exts}\AgdaSpace{}%
\AgdaBound{fb'}\AgdaSpace{}%
\AgdaBound{f'}\AgdaSpace{}%
\AgdaBound{h}\AgdaSpace{}%
\AgdaSymbol{(}\AgdaInductiveConstructor{S}\AgdaSpace{}%
\AgdaBound{x}\AgdaSymbol{)}\AgdaSpace{}%
\AgdaSymbol{=}\<%
\\
\>[0][@{}l@{\AgdaIndent{0}}]%
\>[2]\AgdaOperator{\AgdaFunction{beginEq}}\<%
\\
\>[2][@{}l@{\AgdaIndent{0}}]%
\>[4]\AgdaFunction{renameb}\AgdaSpace{}%
\AgdaInductiveConstructor{Sb}\AgdaSpace{}%
\AgdaSymbol{(}\AgdaBound{fb'}\AgdaSpace{}%
\AgdaSymbol{(}\AgdaFunction{embedLookup}\AgdaSpace{}%
\AgdaBound{x}\AgdaSymbol{))}\<%
\\
%
\>[2]\AgdaFunction{≡⟨}\AgdaSpace{}%
\AgdaFunction{cong}\AgdaSpace{}%
\AgdaSymbol{(}\AgdaFunction{renameb}\AgdaSpace{}%
\AgdaInductiveConstructor{Sb}\AgdaSymbol{)}\AgdaSpace{}%
\AgdaSymbol{(}\AgdaBound{h}\AgdaSpace{}%
\AgdaBound{x}\AgdaSymbol{)}\AgdaSpace{}%
\AgdaFunction{⟩}\<%
\\
\>[2][@{}l@{\AgdaIndent{0}}]%
\>[4]\AgdaFunction{renameb}\AgdaSpace{}%
\AgdaInductiveConstructor{Sb}\AgdaSpace{}%
\AgdaSymbol{(}\AgdaFunction{embedTerm}\AgdaSpace{}%
\AgdaSymbol{(}\AgdaBound{f'}\AgdaSpace{}%
\AgdaBound{x}\AgdaSymbol{))}\<%
\\
%
\>[2]\AgdaFunction{≡⟨}\AgdaSpace{}%
\AgdaFunction{pres-rename}\AgdaSpace{}%
\AgdaInductiveConstructor{S}\AgdaSpace{}%
\AgdaInductiveConstructor{Sb}\AgdaSpace{}%
\AgdaSymbol{(λ}\AgdaSpace{}%
\AgdaBound{x₁}\AgdaSpace{}%
\AgdaSymbol{→}\AgdaSpace{}%
\AgdaInductiveConstructor{refl}\AgdaSymbol{)}\AgdaSpace{}%
\AgdaSymbol{(}\AgdaBound{f'}\AgdaSpace{}%
\AgdaBound{x}\AgdaSymbol{)}\AgdaSpace{}%
\AgdaFunction{⟩}\<%
\\
\>[2][@{}l@{\AgdaIndent{0}}]%
\>[4]\AgdaFunction{embedTerm}\AgdaSpace{}%
\AgdaSymbol{(}\AgdaFunction{rename}\AgdaSpace{}%
\AgdaInductiveConstructor{S}\AgdaSpace{}%
\AgdaSymbol{(}\AgdaBound{f'}\AgdaSpace{}%
\AgdaBound{x}\AgdaSymbol{))}\<%
\\
%
\>[2]\AgdaOperator{\AgdaFunction{∎eq}}\<%
\\
%
\\[\AgdaEmptyExtraSkip]%
\>[0]\AgdaComment{--\ Preservation\ and\ Reflection\ of\ }\<%
\\
\>[0]\AgdaComment{--\ simultaneous\ substitution\ of\ Girard's\ translation}\<%
\\
\>[0]\AgdaFunction{pres-simultaneous-subst}\AgdaSpace{}%
\AgdaSymbol{:}\AgdaSpace{}%
\AgdaSymbol{∀}\AgdaSpace{}%
\AgdaSymbol{\{}\AgdaBound{Γ}\AgdaSpace{}%
\AgdaBound{Δ}\AgdaSpace{}%
\AgdaBound{A}\AgdaSymbol{\}}\<%
\\
\>[0][@{}l@{\AgdaIndent{0}}]%
\>[2]\AgdaSymbol{→}\AgdaSpace{}%
\AgdaSymbol{(}\AgdaBound{fb'}\AgdaSpace{}%
\AgdaSymbol{:}\AgdaSpace{}%
\AgdaSymbol{∀}\AgdaSpace{}%
\AgdaSymbol{\{}\AgdaBound{Ab}\AgdaSymbol{\}}\AgdaSpace{}%
\AgdaSymbol{→}\AgdaSpace{}%
\AgdaInductiveConstructor{□}\AgdaSpace{}%
\AgdaBound{Ab}\AgdaSpace{}%
\AgdaOperator{\AgdaDatatype{∈b}}\AgdaSpace{}%
\AgdaFunction{embedContext}\AgdaSpace{}%
\AgdaBound{Γ}\AgdaSpace{}%
\AgdaSymbol{→}\AgdaSpace{}%
\AgdaFunction{embedContext}\AgdaSpace{}%
\AgdaBound{Δ}\AgdaSpace{}%
\AgdaOperator{\AgdaDatatype{⊢b}}\AgdaSpace{}%
\AgdaBound{Ab}\AgdaSymbol{)}\<%
\\
%
\>[2]\AgdaSymbol{→}\AgdaSpace{}%
\AgdaSymbol{(}\AgdaBound{f'}\AgdaSpace{}%
\AgdaSymbol{:}\AgdaSpace{}%
\AgdaSymbol{∀}\AgdaSpace{}%
\AgdaSymbol{\{}\AgdaBound{A}\AgdaSymbol{\}}\AgdaSpace{}%
\AgdaSymbol{→}\AgdaSpace{}%
\AgdaBound{A}\AgdaSpace{}%
\AgdaOperator{\AgdaDatatype{∈}}\AgdaSpace{}%
\AgdaBound{Γ}\AgdaSpace{}%
\AgdaSymbol{→}\AgdaSpace{}%
\AgdaBound{Δ}\AgdaSpace{}%
\AgdaOperator{\AgdaDatatype{⊢}}\AgdaSpace{}%
\AgdaBound{A}\AgdaSymbol{)}\<%
\\
%
\>[2]\AgdaSymbol{→}\AgdaSpace{}%
\AgdaSymbol{(}\AgdaBound{h}\AgdaSpace{}%
\AgdaSymbol{:}\AgdaSpace{}%
\AgdaSymbol{∀}\AgdaSpace{}%
\AgdaSymbol{\{}\AgdaBound{A'}\AgdaSymbol{\}}\AgdaSpace{}%
\AgdaSymbol{→}\AgdaSpace{}%
\AgdaSymbol{(}\AgdaBound{x}\AgdaSpace{}%
\AgdaSymbol{:}\AgdaSpace{}%
\AgdaBound{A'}\AgdaSpace{}%
\AgdaOperator{\AgdaDatatype{∈}}\AgdaSpace{}%
\AgdaBound{Γ}\AgdaSymbol{)}\AgdaSpace{}%
\AgdaSymbol{→}\AgdaSpace{}%
\AgdaBound{fb'}\AgdaSpace{}%
\AgdaSymbol{(}\AgdaFunction{embedLookup}\AgdaSpace{}%
\AgdaBound{x}\AgdaSymbol{)}\AgdaSpace{}%
\AgdaOperator{\AgdaDatatype{≡}}\AgdaSpace{}%
\AgdaFunction{embedTerm}\AgdaSpace{}%
\AgdaSymbol{(}\AgdaBound{f'}\AgdaSpace{}%
\AgdaBound{x}\AgdaSymbol{)}\AgdaSpace{}%
\AgdaSymbol{)}\<%
\\
%
\>[2]\AgdaSymbol{→}\AgdaSpace{}%
\AgdaSymbol{(}\AgdaBound{M}\AgdaSpace{}%
\AgdaSymbol{:}\AgdaSpace{}%
\AgdaBound{Γ}\AgdaSpace{}%
\AgdaOperator{\AgdaDatatype{⊢}}\AgdaSpace{}%
\AgdaBound{A}\AgdaSymbol{)}\<%
\\
%
\>[2]\AgdaSymbol{→}\AgdaSpace{}%
\AgdaFunction{substb}\AgdaSpace{}%
\AgdaBound{fb'}\AgdaSpace{}%
\AgdaSymbol{(}\AgdaFunction{embedTerm}\AgdaSpace{}%
\AgdaBound{M}\AgdaSymbol{)}\AgdaSpace{}%
\AgdaOperator{\AgdaDatatype{≡}}\AgdaSpace{}%
\AgdaFunction{embedTerm}\AgdaSpace{}%
\AgdaSymbol{(}\AgdaFunction{subst}\AgdaSpace{}%
\AgdaBound{f'}\AgdaSpace{}%
\AgdaBound{M}\AgdaSymbol{)}\<%
\\
\>[0]\AgdaFunction{pres-simultaneous-subst}\AgdaSpace{}%
\AgdaBound{fb'}\AgdaSpace{}%
\AgdaBound{f'}\AgdaSpace{}%
\AgdaBound{h}\AgdaSpace{}%
\AgdaSymbol{(}\AgdaInductiveConstructor{`}\AgdaSpace{}%
\AgdaBound{x}\AgdaSymbol{)}\AgdaSpace{}%
\AgdaSymbol{=}\<%
\\
\>[0][@{}l@{\AgdaIndent{0}}]%
\>[2]\AgdaBound{h}\AgdaSpace{}%
\AgdaBound{x}\<%
\\
\>[0]\AgdaFunction{pres-simultaneous-subst}\AgdaSpace{}%
\AgdaBound{fb'}\AgdaSpace{}%
\AgdaBound{f'}\AgdaSpace{}%
\AgdaBound{h}\AgdaSpace{}%
\AgdaSymbol{(}\AgdaInductiveConstructor{ƛ}\AgdaSpace{}%
\AgdaBound{M}\AgdaSymbol{)}\AgdaSpace{}%
\AgdaSymbol{=}\<%
\\
\>[0][@{}l@{\AgdaIndent{0}}]%
\>[2]\AgdaFunction{cong}\<%
\\
\>[2][@{}l@{\AgdaIndent{0}}]%
\>[4]\AgdaInductiveConstructor{ƛb}\<%
\\
%
\>[4]\AgdaSymbol{(}\AgdaSpace{}%
\AgdaFunction{pres-simultaneous-subst}\AgdaSpace{}%
\AgdaSymbol{(}\AgdaFunction{extsb}\AgdaSpace{}%
\AgdaBound{fb'}\AgdaSymbol{)}\AgdaSpace{}%
\AgdaSymbol{(}\AgdaFunction{exts}\AgdaSpace{}%
\AgdaBound{f'}\AgdaSymbol{)}\AgdaSpace{}%
\AgdaSymbol{(}\AgdaFunction{pres-exts}\AgdaSpace{}%
\AgdaBound{fb'}\AgdaSpace{}%
\AgdaBound{f'}\AgdaSpace{}%
\AgdaBound{h}\AgdaSymbol{)}\AgdaSpace{}%
\AgdaBound{M}\AgdaSpace{}%
\AgdaSymbol{)}\<%
\\
\>[0]\AgdaFunction{pres-simultaneous-subst}\AgdaSpace{}%
\AgdaBound{fb'}\AgdaSpace{}%
\AgdaBound{f'}\AgdaSpace{}%
\AgdaBound{h}\AgdaSpace{}%
\AgdaSymbol{(}\AgdaBound{M}\AgdaSpace{}%
\AgdaOperator{\AgdaInductiveConstructor{·}}\AgdaSpace{}%
\AgdaBound{M₁}\AgdaSymbol{)}\AgdaSpace{}%
\AgdaSymbol{=}\<%
\\
\>[0][@{}l@{\AgdaIndent{0}}]%
\>[2]\AgdaFunction{cong₂}\<%
\\
\>[2][@{}l@{\AgdaIndent{0}}]%
\>[4]\AgdaOperator{\AgdaInductiveConstructor{\AgdaUnderscore{}·b\AgdaUnderscore{}}}\<%
\\
%
\>[4]\AgdaSymbol{(}\AgdaFunction{pres-simultaneous-subst}\AgdaSpace{}%
\AgdaBound{fb'}\AgdaSpace{}%
\AgdaBound{f'}\AgdaSpace{}%
\AgdaBound{h}\AgdaSpace{}%
\AgdaBound{M}\AgdaSymbol{)}\<%
\\
%
\>[4]\AgdaSymbol{(}\AgdaFunction{cong}\AgdaSpace{}%
\AgdaInductiveConstructor{box}\AgdaSpace{}%
\AgdaSymbol{(}\AgdaFunction{pres-simultaneous-subst}\AgdaSpace{}%
\AgdaBound{fb'}\AgdaSpace{}%
\AgdaBound{f'}\AgdaSpace{}%
\AgdaBound{h}\AgdaSpace{}%
\AgdaBound{M₁}\AgdaSymbol{))}\<%
\\
%
\\[\AgdaEmptyExtraSkip]%
%
\\[\AgdaEmptyExtraSkip]%
\>[0]\AgdaComment{--\ Proof\ of\ h}\<%
\\
\>[0]\AgdaFunction{pres-subst-func}\AgdaSpace{}%
\AgdaSymbol{:}\AgdaSpace{}%
\AgdaSymbol{∀}\AgdaSpace{}%
\AgdaSymbol{\{}\AgdaBound{A}\AgdaSpace{}%
\AgdaBound{C}\AgdaSpace{}%
\AgdaBound{Γ}\AgdaSymbol{\}}\<%
\\
\>[0][@{}l@{\AgdaIndent{0}}]%
\>[2]\AgdaSymbol{→}\AgdaSpace{}%
\AgdaSymbol{(}\AgdaBound{N}\AgdaSpace{}%
\AgdaSymbol{:}\AgdaSpace{}%
\AgdaBound{Γ}\AgdaSpace{}%
\AgdaOperator{\AgdaDatatype{⊢}}\AgdaSpace{}%
\AgdaBound{C}\AgdaSymbol{)}\<%
\\
%
\>[2]\AgdaSymbol{→}\AgdaSpace{}%
\AgdaSymbol{(}\AgdaBound{x}\AgdaSpace{}%
\AgdaSymbol{:}\AgdaSpace{}%
\AgdaBound{A}\AgdaSpace{}%
\AgdaOperator{\AgdaDatatype{∈}}\AgdaSpace{}%
\AgdaBound{Γ}\AgdaSpace{}%
\AgdaOperator{\AgdaInductiveConstructor{,}}\AgdaSpace{}%
\AgdaBound{C}\AgdaSymbol{)}\<%
\\
%
\>[2]\AgdaSymbol{→}\AgdaSpace{}%
\AgdaFunction{fb}\AgdaSpace{}%
\AgdaSymbol{(}\AgdaFunction{embedTerm}\AgdaSpace{}%
\AgdaBound{N}\AgdaSymbol{)}\AgdaSpace{}%
\AgdaSymbol{(}\AgdaFunction{embedLookup}\AgdaSpace{}%
\AgdaBound{x}\AgdaSymbol{)}\AgdaSpace{}%
\AgdaOperator{\AgdaDatatype{≡}}\AgdaSpace{}%
\AgdaFunction{embedTerm}\AgdaSpace{}%
\AgdaSymbol{((}\AgdaFunction{f}\AgdaSpace{}%
\AgdaBound{N}\AgdaSymbol{)}\AgdaSpace{}%
\AgdaBound{x}\AgdaSymbol{)}\<%
\\
\>[0]\AgdaFunction{pres-subst-func}\AgdaSpace{}%
\AgdaBound{N}\AgdaSpace{}%
\AgdaInductiveConstructor{Z}\AgdaSpace{}%
\AgdaSymbol{=}\AgdaSpace{}%
\AgdaInductiveConstructor{refl}\<%
\\
\>[0]\AgdaFunction{pres-subst-func}\AgdaSpace{}%
\AgdaBound{N}\AgdaSpace{}%
\AgdaSymbol{(}\AgdaInductiveConstructor{S}\AgdaSpace{}%
\AgdaBound{x}\AgdaSymbol{)}\AgdaSpace{}%
\AgdaSymbol{=}\AgdaSpace{}%
\AgdaInductiveConstructor{refl}\<%
\\
\>[0]\AgdaComment{--\ Reflection\ and\ preservation\ of\ single\ substitution\ }\<%
\\
\>[0]\AgdaComment{--\ of\ refined\ Girard's\ translation\ }\<%
\\
\>[0]\AgdaFunction{pres-subst}\AgdaSpace{}%
\AgdaSymbol{:}\AgdaSpace{}%
\AgdaSymbol{∀}\AgdaSpace{}%
\AgdaSymbol{\{}\AgdaBound{Γ}\AgdaSpace{}%
\AgdaBound{A}\AgdaSpace{}%
\AgdaBound{B}\AgdaSymbol{\}}\<%
\\
\>[0][@{}l@{\AgdaIndent{0}}]%
\>[2]\AgdaSymbol{→}\AgdaSpace{}%
\AgdaSymbol{(}\AgdaBound{M}\AgdaSpace{}%
\AgdaSymbol{:}\AgdaSpace{}%
\AgdaBound{Γ}\AgdaSpace{}%
\AgdaOperator{\AgdaInductiveConstructor{,}}\AgdaSpace{}%
\AgdaBound{B}\AgdaSpace{}%
\AgdaOperator{\AgdaDatatype{⊢}}\AgdaSpace{}%
\AgdaBound{A}\AgdaSymbol{)}\<%
\\
%
\>[2]\AgdaSymbol{→}\AgdaSpace{}%
\AgdaSymbol{(}\AgdaBound{N}\AgdaSpace{}%
\AgdaSymbol{:}\AgdaSpace{}%
\AgdaBound{Γ}\AgdaSpace{}%
\AgdaOperator{\AgdaDatatype{⊢}}\AgdaSpace{}%
\AgdaBound{B}\AgdaSymbol{)}\<%
\\
%
\>[2]\AgdaSymbol{→}\AgdaSpace{}%
\AgdaSymbol{(}\AgdaFunction{embedTerm}\AgdaSpace{}%
\AgdaBound{M}\AgdaSymbol{)}\AgdaSpace{}%
\AgdaOperator{\AgdaFunction{[}}\AgdaSpace{}%
\AgdaFunction{embedTerm}\AgdaSpace{}%
\AgdaBound{N}\AgdaSpace{}%
\AgdaOperator{\AgdaFunction{]b}}\AgdaSpace{}%
\AgdaOperator{\AgdaDatatype{≡}}\AgdaSpace{}%
\AgdaFunction{embedTerm}\AgdaSpace{}%
\AgdaSymbol{(}\AgdaBound{M}\AgdaSpace{}%
\AgdaOperator{\AgdaFunction{[}}\AgdaSpace{}%
\AgdaBound{N}\AgdaSpace{}%
\AgdaOperator{\AgdaFunction{]}}\AgdaSymbol{)}\<%
\\
\>[0]\AgdaFunction{pres-subst}\AgdaSpace{}%
\AgdaBound{M}\AgdaSpace{}%
\AgdaBound{N}\AgdaSpace{}%
\AgdaSymbol{=}\<%
\\
\>[0][@{}l@{\AgdaIndent{0}}]%
\>[2]\AgdaFunction{pres-simultaneous-subst}\AgdaSpace{}%
\AgdaSymbol{(}\AgdaFunction{fb}\AgdaSpace{}%
\AgdaSymbol{(}\AgdaFunction{embedTerm}\AgdaSpace{}%
\AgdaBound{N}\AgdaSymbol{))}\AgdaSpace{}%
\AgdaSymbol{(}\AgdaFunction{f}\AgdaSpace{}%
\AgdaBound{N}\AgdaSymbol{)}\AgdaSpace{}%
\AgdaSymbol{(}\AgdaFunction{pres-subst-func}\AgdaSpace{}%
\AgdaBound{N}\AgdaSymbol{)}\AgdaSpace{}%
\AgdaBound{M}\<%
\\
\>[0]\<%
\end{code}
\begin{code}%
\>[0]\AgdaComment{--\ Preservation\ of\ reduction\ of\ Girard's\ translation}\<%
\\
\>[0]\AgdaFunction{pres-red}\AgdaSpace{}%
\AgdaSymbol{:}\AgdaSpace{}%
\AgdaSymbol{∀}\AgdaSpace{}%
\AgdaSymbol{\{}\AgdaBound{Γ}\AgdaSpace{}%
\AgdaBound{A}\AgdaSymbol{\}}\<%
\\
\>[0][@{}l@{\AgdaIndent{0}}]%
\>[2]\AgdaSymbol{→}\AgdaSpace{}%
\AgdaSymbol{(}\AgdaBound{M}\AgdaSpace{}%
\AgdaSymbol{:}\AgdaSpace{}%
\AgdaBound{Γ}\AgdaSpace{}%
\AgdaOperator{\AgdaDatatype{⊢}}\AgdaSpace{}%
\AgdaBound{A}\AgdaSymbol{)}\<%
\\
%
\>[2]\AgdaSymbol{→}\AgdaSpace{}%
\AgdaSymbol{(}\AgdaBound{N}\AgdaSpace{}%
\AgdaSymbol{:}\AgdaSpace{}%
\AgdaBound{Γ}\AgdaSpace{}%
\AgdaOperator{\AgdaDatatype{⊢}}\AgdaSpace{}%
\AgdaBound{A}\AgdaSymbol{)}\<%
\\
%
\>[2]\AgdaSymbol{→}\AgdaSpace{}%
\AgdaSymbol{(}\AgdaBound{M}\AgdaSpace{}%
\AgdaOperator{\AgdaDatatype{↝βₙ}}\AgdaSpace{}%
\AgdaBound{N}\AgdaSymbol{)}\<%
\\
%
\>[2]\AgdaSymbol{→}\AgdaSpace{}%
\AgdaSymbol{(}\AgdaFunction{embedTerm}\AgdaSpace{}%
\AgdaBound{M}\AgdaSpace{}%
\AgdaOperator{\AgdaDatatype{↝βb}}\AgdaSpace{}%
\AgdaFunction{embedTerm}\AgdaSpace{}%
\AgdaBound{N}\AgdaSymbol{)}\<%
\\
\>[0]\AgdaFunction{pres-red}\AgdaSpace{}%
\AgdaSymbol{(}\AgdaInductiveConstructor{ƛ}\AgdaSpace{}%
\AgdaBound{M₁}\AgdaSpace{}%
\AgdaOperator{\AgdaInductiveConstructor{·}}\AgdaSpace{}%
\AgdaBound{M₂}\AgdaSymbol{)}\AgdaSpace{}%
\AgdaBound{N}\AgdaSpace{}%
\AgdaInductiveConstructor{βₙ}\AgdaSpace{}%
\AgdaSymbol{=}\<%
\\
\>[0][@{}l@{\AgdaIndent{0}}]%
\>[2]\AgdaFunction{Eq.subst}\AgdaSpace{}%
\AgdaSymbol{(λ}\AgdaSpace{}%
\AgdaBound{t}\AgdaSpace{}%
\AgdaSymbol{→}\AgdaSpace{}%
\AgdaInductiveConstructor{ƛb}\AgdaSpace{}%
\AgdaSymbol{(}\AgdaFunction{embedTerm}\AgdaSpace{}%
\AgdaBound{M₁}\AgdaSymbol{)}\AgdaSpace{}%
\AgdaOperator{\AgdaInductiveConstructor{·b}}\AgdaSpace{}%
\AgdaInductiveConstructor{box}\AgdaSpace{}%
\AgdaSymbol{(}\AgdaFunction{embedTerm}\AgdaSpace{}%
\AgdaBound{M₂}\AgdaSymbol{)}\AgdaSpace{}%
\AgdaOperator{\AgdaDatatype{↝βb}}\AgdaSpace{}%
\AgdaBound{t}\AgdaSymbol{)}\<%
\\
%
\>[2]\AgdaSymbol{(}\AgdaFunction{pres-subst}\AgdaSpace{}%
\AgdaBound{M₁}\AgdaSpace{}%
\AgdaBound{M₂}\AgdaSymbol{)}\<%
\\
%
\>[2]\AgdaInductiveConstructor{βb}\<%
\\
\>[0]\AgdaFunction{pres-red}\AgdaSpace{}%
\AgdaBound{M}\AgdaSpace{}%
\AgdaBound{N}\AgdaSpace{}%
\AgdaSymbol{(}\AgdaInductiveConstructor{μ}\AgdaSpace{}%
\AgdaSymbol{\{}\AgdaArgument{M}\AgdaSpace{}%
\AgdaSymbol{=}\AgdaSpace{}%
\AgdaBound{M₁}\AgdaSymbol{\}}\AgdaSpace{}%
\AgdaSymbol{\{}\AgdaArgument{M'}\AgdaSpace{}%
\AgdaSymbol{=}\AgdaSpace{}%
\AgdaBound{M'}\AgdaSymbol{\}}\AgdaSpace{}%
\AgdaBound{red}\AgdaSymbol{)}\AgdaSpace{}%
\AgdaSymbol{=}\<%
\\
\>[0][@{}l@{\AgdaIndent{0}}]%
\>[2]\AgdaInductiveConstructor{μ}\AgdaSpace{}%
\AgdaSymbol{(}\AgdaFunction{pres-red}\AgdaSpace{}%
\AgdaBound{M₁}\AgdaSpace{}%
\AgdaBound{M'}\AgdaSpace{}%
\AgdaBound{red}\AgdaSymbol{)}\<%
\\
\>[0]\AgdaFunction{pres-red}\AgdaSpace{}%
\AgdaBound{M}\AgdaSpace{}%
\AgdaBound{N}\AgdaSpace{}%
\AgdaSymbol{(}\AgdaInductiveConstructor{ν<}\AgdaSpace{}%
\AgdaSymbol{\{}\AgdaArgument{M}\AgdaSpace{}%
\AgdaSymbol{=}\AgdaSpace{}%
\AgdaBound{M₁}\AgdaSymbol{\}}\AgdaSpace{}%
\AgdaSymbol{\{}\AgdaArgument{N}\AgdaSpace{}%
\AgdaSymbol{=}\AgdaSpace{}%
\AgdaBound{N₁}\AgdaSymbol{\}}\AgdaSpace{}%
\AgdaSymbol{\{}\AgdaArgument{N'}\AgdaSpace{}%
\AgdaSymbol{=}\AgdaSpace{}%
\AgdaBound{N'}\AgdaSymbol{\}}\AgdaSpace{}%
\AgdaBound{x}\AgdaSpace{}%
\AgdaBound{red}\AgdaSymbol{)}\AgdaSpace{}%
\AgdaSymbol{=}\<%
\\
\>[0][@{}l@{\AgdaIndent{0}}]%
\>[2]\AgdaInductiveConstructor{ν}\AgdaSpace{}%
\AgdaSymbol{(}\AgdaInductiveConstructor{ζ}\AgdaSpace{}%
\AgdaSymbol{(}\AgdaFunction{pres-red}\AgdaSpace{}%
\AgdaBound{N₁}\AgdaSpace{}%
\AgdaBound{N'}\AgdaSpace{}%
\AgdaBound{red}\AgdaSymbol{))}\<%
\\
\>[0]\AgdaFunction{pres-red}\AgdaSpace{}%
\AgdaBound{M}\AgdaSpace{}%
\AgdaBound{N}\AgdaSpace{}%
\AgdaSymbol{(}\AgdaInductiveConstructor{ν}\AgdaSpace{}%
\AgdaSymbol{\{}\AgdaArgument{M}\AgdaSpace{}%
\AgdaSymbol{=}\AgdaSpace{}%
\AgdaBound{M₁}\AgdaSymbol{\}}\AgdaSpace{}%
\AgdaSymbol{\{}\AgdaArgument{N}\AgdaSpace{}%
\AgdaSymbol{=}\AgdaSpace{}%
\AgdaBound{N₁}\AgdaSymbol{\}}\AgdaSpace{}%
\AgdaSymbol{\{}\AgdaArgument{N'}\AgdaSpace{}%
\AgdaSymbol{=}\AgdaSpace{}%
\AgdaBound{N'}\AgdaSymbol{\}}\AgdaSpace{}%
\AgdaBound{red}\AgdaSymbol{)}\AgdaSpace{}%
\AgdaSymbol{=}\<%
\\
\>[0][@{}l@{\AgdaIndent{0}}]%
\>[2]\AgdaInductiveConstructor{ν}\AgdaSpace{}%
\AgdaSymbol{(}\AgdaInductiveConstructor{ζ}\AgdaSpace{}%
\AgdaSymbol{(}\AgdaFunction{pres-red}\AgdaSpace{}%
\AgdaBound{N₁}\AgdaSpace{}%
\AgdaBound{N'}\AgdaSpace{}%
\AgdaBound{red}\AgdaSymbol{))}\<%
\\
\>[0]\AgdaFunction{pres-red}\AgdaSpace{}%
\AgdaBound{M}\AgdaSpace{}%
\AgdaBound{N}\AgdaSpace{}%
\AgdaSymbol{(}\AgdaInductiveConstructor{ξ}\AgdaSpace{}%
\AgdaSymbol{\{}\AgdaArgument{M}\AgdaSpace{}%
\AgdaSymbol{=}\AgdaSpace{}%
\AgdaBound{M₁}\AgdaSymbol{\}}\AgdaSpace{}%
\AgdaSymbol{\{}\AgdaArgument{M'}\AgdaSpace{}%
\AgdaSymbol{=}\AgdaSpace{}%
\AgdaBound{M'}\AgdaSymbol{\}}\AgdaSpace{}%
\AgdaBound{red}\AgdaSymbol{)}\AgdaSpace{}%
\AgdaSymbol{=}\<%
\\
\>[0][@{}l@{\AgdaIndent{0}}]%
\>[2]\AgdaInductiveConstructor{ξ}\AgdaSpace{}%
\AgdaSymbol{(}\AgdaFunction{pres-red}\AgdaSpace{}%
\AgdaBound{M₁}\AgdaSpace{}%
\AgdaBound{M'}\AgdaSpace{}%
\AgdaBound{red}\AgdaSymbol{)}\<%
\end{code}
\begin{code}[hide]%
\>[0]\<%
\\
\>[0]\AgdaFunction{pres-eval}\AgdaSpace{}%
\AgdaSymbol{:}\AgdaSpace{}%
\AgdaSymbol{∀}\AgdaSpace{}%
\AgdaSymbol{\{}\AgdaBound{Γ}\AgdaSpace{}%
\AgdaBound{A}\AgdaSymbol{\}}\<%
\\
\>[0][@{}l@{\AgdaIndent{0}}]%
\>[2]\AgdaSymbol{→}\AgdaSpace{}%
\AgdaSymbol{(}\AgdaBound{M}\AgdaSpace{}%
\AgdaSymbol{:}\AgdaSpace{}%
\AgdaBound{Γ}\AgdaSpace{}%
\AgdaOperator{\AgdaDatatype{⊢}}\AgdaSpace{}%
\AgdaBound{A}\AgdaSymbol{)}\<%
\\
%
\>[2]\AgdaSymbol{→}\AgdaSpace{}%
\AgdaSymbol{(}\AgdaBound{N}\AgdaSpace{}%
\AgdaSymbol{:}\AgdaSpace{}%
\AgdaBound{Γ}\AgdaSpace{}%
\AgdaOperator{\AgdaDatatype{⊢}}\AgdaSpace{}%
\AgdaBound{A}\AgdaSymbol{)}\<%
\\
%
\>[2]\AgdaSymbol{→}\AgdaSpace{}%
\AgdaSymbol{(}\AgdaBound{M}\AgdaSpace{}%
\AgdaOperator{\AgdaDatatype{↝ₙ}}\AgdaSpace{}%
\AgdaBound{N}\AgdaSymbol{)}\<%
\\
%
\>[2]\AgdaSymbol{→}\AgdaSpace{}%
\AgdaSymbol{(}\AgdaFunction{embedTerm}\AgdaSpace{}%
\AgdaBound{M}\AgdaSpace{}%
\AgdaOperator{\AgdaDatatype{↝bₙ}}\AgdaSpace{}%
\AgdaFunction{embedTerm}\AgdaSpace{}%
\AgdaBound{N}\AgdaSymbol{)}\<%
\\
\>[0]\AgdaFunction{pres-eval}\AgdaSpace{}%
\AgdaSymbol{(}\AgdaInductiveConstructor{ƛ}\AgdaSpace{}%
\AgdaBound{M₁}\AgdaSpace{}%
\AgdaOperator{\AgdaInductiveConstructor{·}}\AgdaSpace{}%
\AgdaBound{M₂}\AgdaSymbol{)}\AgdaSpace{}%
\AgdaBound{N}\AgdaSpace{}%
\AgdaInductiveConstructor{βₙ}\AgdaSpace{}%
\AgdaSymbol{=}\<%
\\
\>[0][@{}l@{\AgdaIndent{0}}]%
\>[2]\AgdaFunction{Eq.subst}\<%
\\
\>[2][@{}l@{\AgdaIndent{0}}]%
\>[4]\AgdaSymbol{(λ}\AgdaSpace{}%
\AgdaBound{t}\AgdaSpace{}%
\AgdaSymbol{→}\AgdaSpace{}%
\AgdaInductiveConstructor{ƛb}\AgdaSpace{}%
\AgdaSymbol{(}\AgdaFunction{embedTerm}\AgdaSpace{}%
\AgdaBound{M₁}\AgdaSymbol{)}\AgdaSpace{}%
\AgdaOperator{\AgdaInductiveConstructor{·b}}\AgdaSpace{}%
\AgdaInductiveConstructor{box}\AgdaSpace{}%
\AgdaSymbol{(}\AgdaFunction{embedTerm}\AgdaSpace{}%
\AgdaBound{M₂}\AgdaSymbol{)}\AgdaSpace{}%
\AgdaOperator{\AgdaDatatype{↝bₙ}}\AgdaSpace{}%
\AgdaBound{t}\AgdaSymbol{)}\<%
\\
%
\>[4]\AgdaSymbol{(}\AgdaFunction{pres-subst}\AgdaSpace{}%
\AgdaBound{M₁}\AgdaSpace{}%
\AgdaBound{M₂}\AgdaSymbol{)}\<%
\\
%
\>[4]\AgdaInductiveConstructor{βb}\<%
\\
\>[0]\AgdaFunction{pres-eval}\AgdaSpace{}%
\AgdaBound{M}\AgdaSpace{}%
\AgdaBound{N}\AgdaSpace{}%
\AgdaSymbol{(}\AgdaInductiveConstructor{μ}\AgdaSpace{}%
\AgdaSymbol{\{}\AgdaArgument{M}\AgdaSpace{}%
\AgdaSymbol{=}\AgdaSpace{}%
\AgdaBound{M₁}\AgdaSymbol{\}}\AgdaSpace{}%
\AgdaSymbol{\{}\AgdaArgument{M'}\AgdaSpace{}%
\AgdaSymbol{=}\AgdaSpace{}%
\AgdaBound{M'}\AgdaSymbol{\}}\AgdaSpace{}%
\AgdaBound{red}\AgdaSymbol{)}\AgdaSpace{}%
\AgdaSymbol{=}\<%
\\
\>[0][@{}l@{\AgdaIndent{0}}]%
\>[2]\AgdaInductiveConstructor{μ}\AgdaSpace{}%
\AgdaSymbol{(}\AgdaFunction{pres-eval}\AgdaSpace{}%
\AgdaBound{M₁}\AgdaSpace{}%
\AgdaBound{M'}\AgdaSpace{}%
\AgdaBound{red}\AgdaSymbol{)}\<%
\end{code}
  \hfill $\qed$
\end{proofthm}

\begin{proofthm}
  \label{proof:girard-preservation-eval}
  \textbf{Preservation of evaluation} under Girard's translation.
  \small
  \begin{code}[hide]%
\>[0]\AgdaKeyword{module}\AgdaSpace{}%
\AgdaModule{nameProofs4}\AgdaSpace{}%
\AgdaKeyword{where}\<%
\\
%
\\[\AgdaEmptyExtraSkip]%
\>[0]\AgdaKeyword{open}\AgdaSpace{}%
\AgdaKeyword{import}\AgdaSpace{}%
\AgdaModule{boxCalc}\<%
\\
\>[0]\AgdaKeyword{open}\AgdaSpace{}%
\AgdaKeyword{import}\AgdaSpace{}%
\AgdaModule{calc}\<%
\\
\>[0]\AgdaKeyword{open}\AgdaSpace{}%
\AgdaKeyword{import}\AgdaSpace{}%
\AgdaModule{embedCBNIntoCBB}\<%
\\
\>[0]\AgdaKeyword{open}\AgdaSpace{}%
\AgdaKeyword{import}\AgdaSpace{}%
\AgdaModule{nameCalcRed}\<%
\\
\>[0]\AgdaKeyword{open}\AgdaSpace{}%
\AgdaKeyword{import}\AgdaSpace{}%
\AgdaModule{boxCalcRed}\<%
\\
%
\\[\AgdaEmptyExtraSkip]%
\>[0]\AgdaKeyword{open}\AgdaSpace{}%
\AgdaKeyword{import}\AgdaSpace{}%
\AgdaModule{Data.Empty}\AgdaSpace{}%
\AgdaKeyword{using}\AgdaSpace{}%
\AgdaSymbol{(}\AgdaFunction{⊥}\AgdaSymbol{;}\AgdaSpace{}%
\AgdaFunction{⊥-elim}\AgdaSymbol{)}\<%
\\
\>[0]\AgdaKeyword{import}\AgdaSpace{}%
\AgdaModule{Relation.Binary.PropositionalEquality}\AgdaSpace{}%
\AgdaSymbol{as}\AgdaSpace{}%
\AgdaModule{Eq}\<%
\\
\>[0]\AgdaKeyword{open}\AgdaSpace{}%
\AgdaModule{Eq}\AgdaSpace{}%
\AgdaKeyword{using}\AgdaSpace{}%
\AgdaSymbol{(}\AgdaOperator{\AgdaDatatype{\AgdaUnderscore{}≡\AgdaUnderscore{}}}\AgdaSymbol{;}\AgdaSpace{}%
\AgdaInductiveConstructor{refl}\AgdaSymbol{;}\AgdaSpace{}%
\AgdaFunction{cong}\AgdaSymbol{;}\AgdaSpace{}%
\AgdaFunction{cong₂}\AgdaSymbol{;}\AgdaSpace{}%
\AgdaFunction{sym}\AgdaSymbol{)}\<%
\\
%
\\[\AgdaEmptyExtraSkip]%
\>[0]\AgdaKeyword{open}\AgdaSpace{}%
\AgdaModule{Eq.≡-Reasoning}\AgdaSpace{}%
\AgdaKeyword{using}\AgdaSpace{}%
\AgdaSymbol{(}\AgdaOperator{\AgdaFunction{begin\AgdaUnderscore{}}}\AgdaSymbol{;}\AgdaSpace{}%
\AgdaFunction{step-≡-∣}\AgdaSymbol{;}\AgdaSpace{}%
\AgdaFunction{step-≡-⟩}\AgdaSymbol{;}\AgdaSpace{}%
\AgdaOperator{\AgdaFunction{\AgdaUnderscore{}∎}}\AgdaSymbol{)}\<%
\\
\>[0]\AgdaKeyword{open}\AgdaSpace{}%
\AgdaModule{Eq.≡-Reasoning}\AgdaSpace{}%
\AgdaKeyword{renaming}\AgdaSpace{}%
\AgdaSymbol{(}\AgdaOperator{\AgdaFunction{begin\AgdaUnderscore{}}}\AgdaSpace{}%
\AgdaSymbol{to}\AgdaSpace{}%
\AgdaOperator{\AgdaFunction{beginEq\AgdaUnderscore{}}}\AgdaSymbol{;}\AgdaSpace{}%
\AgdaOperator{\AgdaFunction{\AgdaUnderscore{}∎}}\AgdaSpace{}%
\AgdaSymbol{to}\AgdaSpace{}%
\AgdaOperator{\AgdaFunction{\AgdaUnderscore{}∎eq}}\AgdaSymbol{)}\<%
\\
\>[0]\AgdaKeyword{open}\AgdaSpace{}%
\AgdaKeyword{import}\AgdaSpace{}%
\AgdaModule{Function.Bundles}\AgdaSpace{}%
\AgdaKeyword{using}\AgdaSpace{}%
\AgdaSymbol{(}\AgdaOperator{\AgdaFunction{\AgdaUnderscore{}⇔\AgdaUnderscore{}}}\AgdaSymbol{;}\AgdaSpace{}%
\AgdaFunction{mk⇔}\AgdaSymbol{)}\<%
\\
%
\\[\AgdaEmptyExtraSkip]%
\>[0]\AgdaComment{--\ GIRARD'S\ embedding}\<%
\\
%
\\[\AgdaEmptyExtraSkip]%
\>[0]\AgdaComment{--\ Preservation\ of\ typing\ by\ Girard's\ translation}\<%
\\
\>[0]\AgdaFunction{pres-typing}\AgdaSpace{}%
\AgdaSymbol{:}\AgdaSpace{}%
\AgdaSymbol{∀}\AgdaSpace{}%
\AgdaSymbol{\{}\AgdaBound{Γ}\AgdaSpace{}%
\AgdaBound{A}\AgdaSymbol{\}}\AgdaSpace{}%
\AgdaSymbol{→}\AgdaSpace{}%
\AgdaBound{Γ}\AgdaSpace{}%
\AgdaOperator{\AgdaDatatype{⊢}}\AgdaSpace{}%
\AgdaBound{A}\AgdaSpace{}%
\AgdaSymbol{→}\AgdaSpace{}%
\AgdaFunction{embedContext}\AgdaSpace{}%
\AgdaBound{Γ}\AgdaSpace{}%
\AgdaOperator{\AgdaDatatype{⊢b}}\AgdaSpace{}%
\AgdaFunction{embedType}\AgdaSpace{}%
\AgdaBound{A}\<%
\\
\>[0]\AgdaFunction{pres-typing}\AgdaSpace{}%
\AgdaBound{x}\AgdaSpace{}%
\AgdaSymbol{=}\AgdaSpace{}%
\AgdaFunction{embedTerm}\AgdaSpace{}%
\AgdaBound{x}\<%
\\
%
\\[\AgdaEmptyExtraSkip]%
\>[0]\AgdaFunction{pres-ext}\AgdaSpace{}%
\AgdaSymbol{:}\AgdaSpace{}%
\AgdaSymbol{∀}\AgdaSpace{}%
\AgdaSymbol{\{}\AgdaBound{Γ}\AgdaSpace{}%
\AgdaBound{Δ}\AgdaSpace{}%
\AgdaBound{A'}\AgdaSpace{}%
\AgdaBound{A}\AgdaSymbol{\}}\<%
\\
\>[0][@{}l@{\AgdaIndent{0}}]%
\>[2]\AgdaSymbol{→}\AgdaSpace{}%
\AgdaSymbol{(}\AgdaBound{f'}\AgdaSpace{}%
\AgdaSymbol{:}\AgdaSpace{}%
\AgdaSymbol{∀}\AgdaSpace{}%
\AgdaSymbol{\{}\AgdaBound{A}\AgdaSymbol{\}}\AgdaSpace{}%
\AgdaSymbol{→}\AgdaSpace{}%
\AgdaBound{A}\AgdaSpace{}%
\AgdaOperator{\AgdaDatatype{∈}}\AgdaSpace{}%
\AgdaBound{Γ}\AgdaSpace{}%
\AgdaSymbol{→}\AgdaSpace{}%
\AgdaBound{A}\AgdaSpace{}%
\AgdaOperator{\AgdaDatatype{∈}}\AgdaSpace{}%
\AgdaBound{Δ}\AgdaSymbol{)}\<%
\\
%
\>[2]\AgdaSymbol{→}\AgdaSpace{}%
\AgdaSymbol{(}\AgdaBound{fb'}\AgdaSpace{}%
\AgdaSymbol{:}\AgdaSpace{}%
\AgdaSymbol{∀}\AgdaSpace{}%
\AgdaSymbol{\{}\AgdaBound{Ab}\AgdaSymbol{\}}\AgdaSpace{}%
\AgdaSymbol{→}\AgdaSpace{}%
\AgdaBound{Ab}\AgdaSpace{}%
\AgdaOperator{\AgdaDatatype{∈b}}\AgdaSpace{}%
\AgdaFunction{embedContext}\AgdaSpace{}%
\AgdaBound{Γ}\AgdaSpace{}%
\AgdaSymbol{→}\AgdaSpace{}%
\AgdaBound{Ab}\AgdaSpace{}%
\AgdaOperator{\AgdaDatatype{∈b}}\AgdaSpace{}%
\AgdaFunction{embedContext}\AgdaSpace{}%
\AgdaBound{Δ}\AgdaSymbol{)}\<%
\\
%
\>[2]\AgdaSymbol{→}\AgdaSpace{}%
\AgdaSymbol{(\{}\AgdaBound{A'}\AgdaSpace{}%
\AgdaSymbol{:}\AgdaSpace{}%
\AgdaDatatype{Type}\AgdaSymbol{\}}\AgdaSpace{}%
\AgdaSymbol{(}\AgdaBound{x}\AgdaSpace{}%
\AgdaSymbol{:}\AgdaSpace{}%
\AgdaBound{A'}\AgdaSpace{}%
\AgdaOperator{\AgdaDatatype{∈}}\AgdaSpace{}%
\AgdaBound{Γ}\AgdaSymbol{)}\AgdaSpace{}%
\AgdaSymbol{→}\AgdaSpace{}%
\AgdaBound{fb'}\AgdaSpace{}%
\AgdaSymbol{(}\AgdaFunction{embedLookup}\AgdaSpace{}%
\AgdaBound{x}\AgdaSymbol{)}\AgdaSpace{}%
\AgdaOperator{\AgdaDatatype{≡}}\AgdaSpace{}%
\AgdaFunction{embedLookup}\AgdaSpace{}%
\AgdaSymbol{(}\AgdaBound{f'}\AgdaSpace{}%
\AgdaBound{x}\AgdaSymbol{))}\<%
\\
%
\>[2]\AgdaSymbol{→}\AgdaSpace{}%
\AgdaSymbol{(}\AgdaBound{x}\AgdaSpace{}%
\AgdaSymbol{:}\AgdaSpace{}%
\AgdaBound{A'}\AgdaSpace{}%
\AgdaOperator{\AgdaDatatype{∈}}\AgdaSpace{}%
\AgdaBound{Γ}\AgdaSpace{}%
\AgdaOperator{\AgdaInductiveConstructor{,}}\AgdaSpace{}%
\AgdaBound{A}\AgdaSymbol{)}\<%
\\
%
\>[2]\AgdaSymbol{→}\AgdaSpace{}%
\AgdaFunction{extb}\AgdaSpace{}%
\AgdaBound{fb'}\AgdaSpace{}%
\AgdaSymbol{(}\AgdaFunction{embedLookup}\AgdaSpace{}%
\AgdaBound{x}\AgdaSymbol{)}\AgdaSpace{}%
\AgdaOperator{\AgdaDatatype{≡}}\AgdaSpace{}%
\AgdaFunction{embedLookup}\AgdaSpace{}%
\AgdaSymbol{(}\AgdaFunction{ext}\AgdaSpace{}%
\AgdaBound{f'}\AgdaSpace{}%
\AgdaBound{x}\AgdaSymbol{)}\<%
\\
\>[0]\AgdaFunction{pres-ext}\AgdaSpace{}%
\AgdaBound{f'}\AgdaSpace{}%
\AgdaBound{fb'}\AgdaSpace{}%
\AgdaBound{h}\AgdaSpace{}%
\AgdaInductiveConstructor{Z}\AgdaSpace{}%
\AgdaSymbol{=}\AgdaSpace{}%
\AgdaInductiveConstructor{refl}\<%
\\
\>[0]\AgdaFunction{pres-ext}\AgdaSpace{}%
\AgdaBound{f'}\AgdaSpace{}%
\AgdaBound{fb'}\AgdaSpace{}%
\AgdaBound{h}\AgdaSpace{}%
\AgdaSymbol{(}\AgdaInductiveConstructor{S}\AgdaSpace{}%
\AgdaBound{x}\AgdaSymbol{)}\AgdaSpace{}%
\AgdaSymbol{=}\AgdaSpace{}%
\AgdaFunction{cong}\AgdaSpace{}%
\AgdaInductiveConstructor{Sb}\AgdaSpace{}%
\AgdaSymbol{(}\AgdaBound{h}\AgdaSpace{}%
\AgdaBound{x}\AgdaSymbol{)}\<%
\\
%
\\[\AgdaEmptyExtraSkip]%
\>[0]\AgdaFunction{pres-rename}\AgdaSpace{}%
\AgdaSymbol{:}\AgdaSpace{}%
\AgdaSymbol{∀}\AgdaSpace{}%
\AgdaSymbol{\{}\AgdaBound{Γ}\AgdaSpace{}%
\AgdaBound{A}\AgdaSpace{}%
\AgdaBound{Δ}\AgdaSymbol{\}}\<%
\\
\>[0][@{}l@{\AgdaIndent{0}}]%
\>[2]\AgdaSymbol{→}\AgdaSpace{}%
\AgdaSymbol{(}\AgdaBound{f'}\AgdaSpace{}%
\AgdaSymbol{:}\AgdaSpace{}%
\AgdaSymbol{∀}\AgdaSpace{}%
\AgdaSymbol{\{}\AgdaBound{A}\AgdaSymbol{\}}\AgdaSpace{}%
\AgdaSymbol{→}\AgdaSpace{}%
\AgdaBound{A}\AgdaSpace{}%
\AgdaOperator{\AgdaDatatype{∈}}\AgdaSpace{}%
\AgdaBound{Γ}\AgdaSpace{}%
\AgdaSymbol{→}\AgdaSpace{}%
\AgdaBound{A}\AgdaSpace{}%
\AgdaOperator{\AgdaDatatype{∈}}\AgdaSpace{}%
\AgdaBound{Δ}\AgdaSymbol{)}\<%
\\
%
\>[2]\AgdaSymbol{→}\AgdaSpace{}%
\AgdaSymbol{(}\AgdaBound{fb'}\AgdaSpace{}%
\AgdaSymbol{:}\AgdaSpace{}%
\AgdaSymbol{∀}\AgdaSpace{}%
\AgdaSymbol{\{}\AgdaBound{Ab}\AgdaSymbol{\}}\AgdaSpace{}%
\AgdaSymbol{→}\AgdaSpace{}%
\AgdaBound{Ab}\AgdaSpace{}%
\AgdaOperator{\AgdaDatatype{∈b}}\AgdaSpace{}%
\AgdaFunction{embedContext}\AgdaSpace{}%
\AgdaBound{Γ}\AgdaSpace{}%
\AgdaSymbol{→}\AgdaSpace{}%
\AgdaBound{Ab}\AgdaSpace{}%
\AgdaOperator{\AgdaDatatype{∈b}}\AgdaSpace{}%
\AgdaFunction{embedContext}\AgdaSpace{}%
\AgdaBound{Δ}\AgdaSymbol{)}\<%
\\
%
\>[2]\AgdaSymbol{→}\AgdaSpace{}%
\AgdaSymbol{(\{}\AgdaBound{A'}\AgdaSpace{}%
\AgdaSymbol{:}\AgdaSpace{}%
\AgdaDatatype{Type}\AgdaSymbol{\}}\AgdaSpace{}%
\AgdaSymbol{(}\AgdaBound{x}\AgdaSpace{}%
\AgdaSymbol{:}\AgdaSpace{}%
\AgdaBound{A'}\AgdaSpace{}%
\AgdaOperator{\AgdaDatatype{∈}}\AgdaSpace{}%
\AgdaBound{Γ}\AgdaSymbol{)}\AgdaSpace{}%
\AgdaSymbol{→}\AgdaSpace{}%
\AgdaBound{fb'}\AgdaSpace{}%
\AgdaSymbol{(}\AgdaFunction{embedLookup}\AgdaSpace{}%
\AgdaBound{x}\AgdaSymbol{)}\AgdaSpace{}%
\AgdaOperator{\AgdaDatatype{≡}}\AgdaSpace{}%
\AgdaFunction{embedLookup}\AgdaSpace{}%
\AgdaSymbol{(}\AgdaBound{f'}\AgdaSpace{}%
\AgdaBound{x}\AgdaSymbol{))}\<%
\\
%
\>[2]\AgdaSymbol{→}\AgdaSpace{}%
\AgdaSymbol{(}\AgdaBound{t}\AgdaSpace{}%
\AgdaSymbol{:}\AgdaSpace{}%
\AgdaBound{Γ}\AgdaSpace{}%
\AgdaOperator{\AgdaDatatype{⊢}}\AgdaSpace{}%
\AgdaBound{A}\AgdaSymbol{)}\<%
\\
%
\>[2]\AgdaSymbol{→}\AgdaSpace{}%
\AgdaFunction{renameb}\AgdaSpace{}%
\AgdaBound{fb'}\AgdaSpace{}%
\AgdaSymbol{(}\AgdaFunction{embedTerm}\AgdaSpace{}%
\AgdaBound{t}\AgdaSymbol{)}\AgdaSpace{}%
\AgdaOperator{\AgdaDatatype{≡}}\AgdaSpace{}%
\AgdaFunction{embedTerm}\AgdaSpace{}%
\AgdaSymbol{(}\AgdaFunction{rename}\AgdaSpace{}%
\AgdaBound{f'}\AgdaSpace{}%
\AgdaBound{t}\AgdaSymbol{)}\<%
\\
\>[0]\AgdaFunction{pres-rename}\AgdaSpace{}%
\AgdaBound{f'}\AgdaSpace{}%
\AgdaBound{fb'}\AgdaSpace{}%
\AgdaBound{h}\AgdaSpace{}%
\AgdaSymbol{(}\AgdaInductiveConstructor{`}\AgdaSpace{}%
\AgdaBound{x}\AgdaSymbol{)}\AgdaSpace{}%
\AgdaSymbol{=}\<%
\\
\>[0][@{}l@{\AgdaIndent{0}}]%
\>[2]\AgdaFunction{cong}\AgdaSpace{}%
\AgdaInductiveConstructor{ε}\AgdaSpace{}%
\AgdaSymbol{(}\AgdaBound{h}\AgdaSpace{}%
\AgdaBound{x}\AgdaSymbol{)}\<%
\\
\>[0]\AgdaFunction{pres-rename}\AgdaSpace{}%
\AgdaBound{f'}\AgdaSpace{}%
\AgdaBound{fb'}\AgdaSpace{}%
\AgdaBound{h}\AgdaSpace{}%
\AgdaSymbol{(}\AgdaInductiveConstructor{ƛ}\AgdaSpace{}%
\AgdaBound{t}\AgdaSymbol{)}\AgdaSpace{}%
\AgdaSymbol{=}\<%
\\
\>[0][@{}l@{\AgdaIndent{0}}]%
\>[2]\AgdaFunction{cong}\AgdaSpace{}%
\AgdaInductiveConstructor{ƛb}\AgdaSpace{}%
\AgdaSymbol{(}\AgdaFunction{pres-rename}\AgdaSpace{}%
\AgdaSymbol{(}\AgdaFunction{calc.ext}\AgdaSpace{}%
\AgdaBound{f'}\AgdaSymbol{)}\AgdaSpace{}%
\AgdaSymbol{(}\AgdaFunction{extb}\AgdaSpace{}%
\AgdaBound{fb'}\AgdaSymbol{)}\AgdaSpace{}%
\AgdaSymbol{(}\AgdaFunction{pres-ext}\AgdaSpace{}%
\AgdaBound{f'}\AgdaSpace{}%
\AgdaBound{fb'}\AgdaSpace{}%
\AgdaBound{h}\AgdaSymbol{)}\AgdaSpace{}%
\AgdaBound{t}\AgdaSymbol{)}\<%
\\
\>[0]\AgdaFunction{pres-rename}\AgdaSpace{}%
\AgdaBound{f'}\AgdaSpace{}%
\AgdaBound{fb'}\AgdaSpace{}%
\AgdaBound{h}\AgdaSpace{}%
\AgdaSymbol{(}\AgdaBound{t}\AgdaSpace{}%
\AgdaOperator{\AgdaInductiveConstructor{·}}\AgdaSpace{}%
\AgdaBound{t₁}\AgdaSymbol{)}\AgdaSpace{}%
\AgdaSymbol{=}\<%
\\
\>[0][@{}l@{\AgdaIndent{0}}]%
\>[2]\AgdaFunction{cong₂}\AgdaSpace{}%
\AgdaOperator{\AgdaInductiveConstructor{\AgdaUnderscore{}·b\AgdaUnderscore{}}}\AgdaSpace{}%
\AgdaSymbol{(}\AgdaFunction{pres-rename}\AgdaSpace{}%
\AgdaBound{f'}\AgdaSpace{}%
\AgdaBound{fb'}\AgdaSpace{}%
\AgdaBound{h}\AgdaSpace{}%
\AgdaBound{t}\AgdaSymbol{)}\AgdaSpace{}%
\AgdaSymbol{(}\AgdaFunction{cong}\AgdaSpace{}%
\AgdaInductiveConstructor{box}\AgdaSpace{}%
\AgdaSymbol{(}\AgdaFunction{pres-rename}\AgdaSpace{}%
\AgdaBound{f'}\AgdaSpace{}%
\AgdaBound{fb'}\AgdaSpace{}%
\AgdaBound{h}\AgdaSpace{}%
\AgdaBound{t₁}\AgdaSymbol{))}\<%
\\
%
\\[\AgdaEmptyExtraSkip]%
%
\\[\AgdaEmptyExtraSkip]%
\>[0]\AgdaFunction{pres-exts}\AgdaSpace{}%
\AgdaSymbol{:}\AgdaSpace{}%
\AgdaSymbol{∀}\AgdaSpace{}%
\AgdaSymbol{\{}\AgdaBound{Γ}\AgdaSpace{}%
\AgdaBound{Δ}\AgdaSymbol{\}}\<%
\\
\>[0][@{}l@{\AgdaIndent{0}}]%
\>[2]\AgdaSymbol{→}\AgdaSpace{}%
\AgdaSymbol{(}\AgdaBound{fb'}\AgdaSpace{}%
\AgdaSymbol{:}\AgdaSpace{}%
\AgdaSymbol{∀}\AgdaSpace{}%
\AgdaSymbol{\{}\AgdaBound{Ab}\AgdaSymbol{\}}\AgdaSpace{}%
\AgdaSymbol{→}\AgdaSpace{}%
\AgdaInductiveConstructor{□}\AgdaSpace{}%
\AgdaBound{Ab}\AgdaSpace{}%
\AgdaOperator{\AgdaDatatype{∈b}}\AgdaSpace{}%
\AgdaFunction{embedContext}\AgdaSpace{}%
\AgdaBound{Γ}\AgdaSpace{}%
\AgdaSymbol{→}\AgdaSpace{}%
\AgdaFunction{embedContext}\AgdaSpace{}%
\AgdaBound{Δ}\AgdaSpace{}%
\AgdaOperator{\AgdaDatatype{⊢b}}\AgdaSpace{}%
\AgdaBound{Ab}\AgdaSymbol{)}\<%
\\
%
\>[2]\AgdaSymbol{→}\AgdaSpace{}%
\AgdaSymbol{(}\AgdaBound{f'}\AgdaSpace{}%
\AgdaSymbol{:}\AgdaSpace{}%
\AgdaSymbol{∀}\AgdaSpace{}%
\AgdaSymbol{\{}\AgdaBound{A}\AgdaSymbol{\}}\AgdaSpace{}%
\AgdaSymbol{→}\AgdaSpace{}%
\AgdaBound{A}\AgdaSpace{}%
\AgdaOperator{\AgdaDatatype{∈}}\AgdaSpace{}%
\AgdaBound{Γ}\AgdaSpace{}%
\AgdaSymbol{→}\AgdaSpace{}%
\AgdaBound{Δ}\AgdaSpace{}%
\AgdaOperator{\AgdaDatatype{⊢}}\AgdaSpace{}%
\AgdaBound{A}\AgdaSymbol{)}\<%
\\
%
\>[2]\AgdaSymbol{→}\AgdaSpace{}%
\AgdaSymbol{(\{}\AgdaBound{A'}\AgdaSpace{}%
\AgdaSymbol{:}\AgdaSpace{}%
\AgdaDatatype{Type}\AgdaSymbol{\}}\AgdaSpace{}%
\AgdaSymbol{(}\AgdaBound{x}\AgdaSpace{}%
\AgdaSymbol{:}\AgdaSpace{}%
\AgdaBound{A'}\AgdaSpace{}%
\AgdaOperator{\AgdaDatatype{∈}}\AgdaSpace{}%
\AgdaBound{Γ}\AgdaSymbol{)}\AgdaSpace{}%
\AgdaSymbol{→}\AgdaSpace{}%
\AgdaBound{fb'}\AgdaSpace{}%
\AgdaSymbol{(}\AgdaFunction{embedLookup}\AgdaSpace{}%
\AgdaBound{x}\AgdaSymbol{)}\AgdaSpace{}%
\AgdaOperator{\AgdaDatatype{≡}}\AgdaSpace{}%
\AgdaFunction{embedTerm}\AgdaSpace{}%
\AgdaSymbol{(}\AgdaBound{f'}\AgdaSpace{}%
\AgdaBound{x}\AgdaSymbol{))}\<%
\\
%
\>[2]\AgdaSymbol{→}\AgdaSpace{}%
\AgdaSymbol{\{}\AgdaBound{A'}\AgdaSpace{}%
\AgdaBound{A}\AgdaSpace{}%
\AgdaSymbol{:}\AgdaSpace{}%
\AgdaDatatype{Type}\AgdaSymbol{\}}\AgdaSpace{}%
\AgdaSymbol{(}\AgdaBound{x}\AgdaSpace{}%
\AgdaSymbol{:}\AgdaSpace{}%
\AgdaBound{A'}\AgdaSpace{}%
\AgdaOperator{\AgdaDatatype{∈}}\AgdaSpace{}%
\AgdaBound{Γ}\AgdaSpace{}%
\AgdaOperator{\AgdaInductiveConstructor{,}}\AgdaSpace{}%
\AgdaBound{A}\AgdaSymbol{)}\AgdaSpace{}%
\AgdaSymbol{→}\AgdaSpace{}%
\AgdaFunction{extsb}\AgdaSpace{}%
\AgdaBound{fb'}\AgdaSpace{}%
\AgdaSymbol{(}\AgdaFunction{embedLookup}\AgdaSpace{}%
\AgdaBound{x}\AgdaSymbol{)}\AgdaSpace{}%
\AgdaOperator{\AgdaDatatype{≡}}\AgdaSpace{}%
\AgdaFunction{embedTerm}\AgdaSpace{}%
\AgdaSymbol{(}\AgdaFunction{exts}\AgdaSpace{}%
\AgdaBound{f'}\AgdaSpace{}%
\AgdaBound{x}\AgdaSymbol{)}\<%
\\
\>[0]\AgdaFunction{pres-exts}\AgdaSpace{}%
\AgdaBound{fb'}\AgdaSpace{}%
\AgdaBound{f'}\AgdaSpace{}%
\AgdaBound{h}\AgdaSpace{}%
\AgdaInductiveConstructor{Z}\AgdaSpace{}%
\AgdaSymbol{=}\AgdaSpace{}%
\AgdaInductiveConstructor{refl}\<%
\\
\>[0]\AgdaFunction{pres-exts}\AgdaSpace{}%
\AgdaBound{fb'}\AgdaSpace{}%
\AgdaBound{f'}\AgdaSpace{}%
\AgdaBound{h}\AgdaSpace{}%
\AgdaSymbol{(}\AgdaInductiveConstructor{S}\AgdaSpace{}%
\AgdaBound{x}\AgdaSymbol{)}\AgdaSpace{}%
\AgdaSymbol{=}\<%
\\
\>[0][@{}l@{\AgdaIndent{0}}]%
\>[2]\AgdaOperator{\AgdaFunction{beginEq}}\<%
\\
\>[2][@{}l@{\AgdaIndent{0}}]%
\>[4]\AgdaFunction{renameb}\AgdaSpace{}%
\AgdaInductiveConstructor{Sb}\AgdaSpace{}%
\AgdaSymbol{(}\AgdaBound{fb'}\AgdaSpace{}%
\AgdaSymbol{(}\AgdaFunction{embedLookup}\AgdaSpace{}%
\AgdaBound{x}\AgdaSymbol{))}\<%
\\
%
\>[2]\AgdaFunction{≡⟨}\AgdaSpace{}%
\AgdaFunction{cong}\AgdaSpace{}%
\AgdaSymbol{(}\AgdaFunction{renameb}\AgdaSpace{}%
\AgdaInductiveConstructor{Sb}\AgdaSymbol{)}\AgdaSpace{}%
\AgdaSymbol{(}\AgdaBound{h}\AgdaSpace{}%
\AgdaBound{x}\AgdaSymbol{)}\AgdaSpace{}%
\AgdaFunction{⟩}\<%
\\
\>[2][@{}l@{\AgdaIndent{0}}]%
\>[4]\AgdaFunction{renameb}\AgdaSpace{}%
\AgdaInductiveConstructor{Sb}\AgdaSpace{}%
\AgdaSymbol{(}\AgdaFunction{embedTerm}\AgdaSpace{}%
\AgdaSymbol{(}\AgdaBound{f'}\AgdaSpace{}%
\AgdaBound{x}\AgdaSymbol{))}\<%
\\
%
\>[2]\AgdaFunction{≡⟨}\AgdaSpace{}%
\AgdaFunction{pres-rename}\AgdaSpace{}%
\AgdaInductiveConstructor{S}\AgdaSpace{}%
\AgdaInductiveConstructor{Sb}\AgdaSpace{}%
\AgdaSymbol{(λ}\AgdaSpace{}%
\AgdaBound{x₁}\AgdaSpace{}%
\AgdaSymbol{→}\AgdaSpace{}%
\AgdaInductiveConstructor{refl}\AgdaSymbol{)}\AgdaSpace{}%
\AgdaSymbol{(}\AgdaBound{f'}\AgdaSpace{}%
\AgdaBound{x}\AgdaSymbol{)}\AgdaSpace{}%
\AgdaFunction{⟩}\<%
\\
\>[2][@{}l@{\AgdaIndent{0}}]%
\>[4]\AgdaFunction{embedTerm}\AgdaSpace{}%
\AgdaSymbol{(}\AgdaFunction{rename}\AgdaSpace{}%
\AgdaInductiveConstructor{S}\AgdaSpace{}%
\AgdaSymbol{(}\AgdaBound{f'}\AgdaSpace{}%
\AgdaBound{x}\AgdaSymbol{))}\<%
\\
%
\>[2]\AgdaOperator{\AgdaFunction{∎eq}}\<%
\\
%
\\[\AgdaEmptyExtraSkip]%
\>[0]\AgdaComment{--\ Preservation\ and\ Reflection\ of\ }\<%
\\
\>[0]\AgdaComment{--\ simultaneous\ substitution\ of\ Girard's\ translation}\<%
\\
\>[0]\AgdaFunction{pres-simultaneous-subst}\AgdaSpace{}%
\AgdaSymbol{:}\AgdaSpace{}%
\AgdaSymbol{∀}\AgdaSpace{}%
\AgdaSymbol{\{}\AgdaBound{Γ}\AgdaSpace{}%
\AgdaBound{Δ}\AgdaSpace{}%
\AgdaBound{A}\AgdaSymbol{\}}\<%
\\
\>[0][@{}l@{\AgdaIndent{0}}]%
\>[2]\AgdaSymbol{→}\AgdaSpace{}%
\AgdaSymbol{(}\AgdaBound{fb'}\AgdaSpace{}%
\AgdaSymbol{:}\AgdaSpace{}%
\AgdaSymbol{∀}\AgdaSpace{}%
\AgdaSymbol{\{}\AgdaBound{Ab}\AgdaSymbol{\}}\AgdaSpace{}%
\AgdaSymbol{→}\AgdaSpace{}%
\AgdaInductiveConstructor{□}\AgdaSpace{}%
\AgdaBound{Ab}\AgdaSpace{}%
\AgdaOperator{\AgdaDatatype{∈b}}\AgdaSpace{}%
\AgdaFunction{embedContext}\AgdaSpace{}%
\AgdaBound{Γ}\AgdaSpace{}%
\AgdaSymbol{→}\AgdaSpace{}%
\AgdaFunction{embedContext}\AgdaSpace{}%
\AgdaBound{Δ}\AgdaSpace{}%
\AgdaOperator{\AgdaDatatype{⊢b}}\AgdaSpace{}%
\AgdaBound{Ab}\AgdaSymbol{)}\<%
\\
%
\>[2]\AgdaSymbol{→}\AgdaSpace{}%
\AgdaSymbol{(}\AgdaBound{f'}\AgdaSpace{}%
\AgdaSymbol{:}\AgdaSpace{}%
\AgdaSymbol{∀}\AgdaSpace{}%
\AgdaSymbol{\{}\AgdaBound{A}\AgdaSymbol{\}}\AgdaSpace{}%
\AgdaSymbol{→}\AgdaSpace{}%
\AgdaBound{A}\AgdaSpace{}%
\AgdaOperator{\AgdaDatatype{∈}}\AgdaSpace{}%
\AgdaBound{Γ}\AgdaSpace{}%
\AgdaSymbol{→}\AgdaSpace{}%
\AgdaBound{Δ}\AgdaSpace{}%
\AgdaOperator{\AgdaDatatype{⊢}}\AgdaSpace{}%
\AgdaBound{A}\AgdaSymbol{)}\<%
\\
%
\>[2]\AgdaSymbol{→}\AgdaSpace{}%
\AgdaSymbol{(}\AgdaBound{h}\AgdaSpace{}%
\AgdaSymbol{:}\AgdaSpace{}%
\AgdaSymbol{∀}\AgdaSpace{}%
\AgdaSymbol{\{}\AgdaBound{A'}\AgdaSymbol{\}}\AgdaSpace{}%
\AgdaSymbol{→}\AgdaSpace{}%
\AgdaSymbol{(}\AgdaBound{x}\AgdaSpace{}%
\AgdaSymbol{:}\AgdaSpace{}%
\AgdaBound{A'}\AgdaSpace{}%
\AgdaOperator{\AgdaDatatype{∈}}\AgdaSpace{}%
\AgdaBound{Γ}\AgdaSymbol{)}\AgdaSpace{}%
\AgdaSymbol{→}\AgdaSpace{}%
\AgdaBound{fb'}\AgdaSpace{}%
\AgdaSymbol{(}\AgdaFunction{embedLookup}\AgdaSpace{}%
\AgdaBound{x}\AgdaSymbol{)}\AgdaSpace{}%
\AgdaOperator{\AgdaDatatype{≡}}\AgdaSpace{}%
\AgdaFunction{embedTerm}\AgdaSpace{}%
\AgdaSymbol{(}\AgdaBound{f'}\AgdaSpace{}%
\AgdaBound{x}\AgdaSymbol{)}\AgdaSpace{}%
\AgdaSymbol{)}\<%
\\
%
\>[2]\AgdaSymbol{→}\AgdaSpace{}%
\AgdaSymbol{(}\AgdaBound{M}\AgdaSpace{}%
\AgdaSymbol{:}\AgdaSpace{}%
\AgdaBound{Γ}\AgdaSpace{}%
\AgdaOperator{\AgdaDatatype{⊢}}\AgdaSpace{}%
\AgdaBound{A}\AgdaSymbol{)}\<%
\\
%
\>[2]\AgdaSymbol{→}\AgdaSpace{}%
\AgdaFunction{substb}\AgdaSpace{}%
\AgdaBound{fb'}\AgdaSpace{}%
\AgdaSymbol{(}\AgdaFunction{embedTerm}\AgdaSpace{}%
\AgdaBound{M}\AgdaSymbol{)}\AgdaSpace{}%
\AgdaOperator{\AgdaDatatype{≡}}\AgdaSpace{}%
\AgdaFunction{embedTerm}\AgdaSpace{}%
\AgdaSymbol{(}\AgdaFunction{subst}\AgdaSpace{}%
\AgdaBound{f'}\AgdaSpace{}%
\AgdaBound{M}\AgdaSymbol{)}\<%
\\
\>[0]\AgdaFunction{pres-simultaneous-subst}\AgdaSpace{}%
\AgdaBound{fb'}\AgdaSpace{}%
\AgdaBound{f'}\AgdaSpace{}%
\AgdaBound{h}\AgdaSpace{}%
\AgdaSymbol{(}\AgdaInductiveConstructor{`}\AgdaSpace{}%
\AgdaBound{x}\AgdaSymbol{)}\AgdaSpace{}%
\AgdaSymbol{=}\<%
\\
\>[0][@{}l@{\AgdaIndent{0}}]%
\>[2]\AgdaBound{h}\AgdaSpace{}%
\AgdaBound{x}\<%
\\
\>[0]\AgdaFunction{pres-simultaneous-subst}\AgdaSpace{}%
\AgdaBound{fb'}\AgdaSpace{}%
\AgdaBound{f'}\AgdaSpace{}%
\AgdaBound{h}\AgdaSpace{}%
\AgdaSymbol{(}\AgdaInductiveConstructor{ƛ}\AgdaSpace{}%
\AgdaBound{M}\AgdaSymbol{)}\AgdaSpace{}%
\AgdaSymbol{=}\<%
\\
\>[0][@{}l@{\AgdaIndent{0}}]%
\>[2]\AgdaFunction{cong}\<%
\\
\>[2][@{}l@{\AgdaIndent{0}}]%
\>[4]\AgdaInductiveConstructor{ƛb}\<%
\\
%
\>[4]\AgdaSymbol{(}\AgdaSpace{}%
\AgdaFunction{pres-simultaneous-subst}\AgdaSpace{}%
\AgdaSymbol{(}\AgdaFunction{extsb}\AgdaSpace{}%
\AgdaBound{fb'}\AgdaSymbol{)}\AgdaSpace{}%
\AgdaSymbol{(}\AgdaFunction{exts}\AgdaSpace{}%
\AgdaBound{f'}\AgdaSymbol{)}\AgdaSpace{}%
\AgdaSymbol{(}\AgdaFunction{pres-exts}\AgdaSpace{}%
\AgdaBound{fb'}\AgdaSpace{}%
\AgdaBound{f'}\AgdaSpace{}%
\AgdaBound{h}\AgdaSymbol{)}\AgdaSpace{}%
\AgdaBound{M}\AgdaSpace{}%
\AgdaSymbol{)}\<%
\\
\>[0]\AgdaFunction{pres-simultaneous-subst}\AgdaSpace{}%
\AgdaBound{fb'}\AgdaSpace{}%
\AgdaBound{f'}\AgdaSpace{}%
\AgdaBound{h}\AgdaSpace{}%
\AgdaSymbol{(}\AgdaBound{M}\AgdaSpace{}%
\AgdaOperator{\AgdaInductiveConstructor{·}}\AgdaSpace{}%
\AgdaBound{M₁}\AgdaSymbol{)}\AgdaSpace{}%
\AgdaSymbol{=}\<%
\\
\>[0][@{}l@{\AgdaIndent{0}}]%
\>[2]\AgdaFunction{cong₂}\<%
\\
\>[2][@{}l@{\AgdaIndent{0}}]%
\>[4]\AgdaOperator{\AgdaInductiveConstructor{\AgdaUnderscore{}·b\AgdaUnderscore{}}}\<%
\\
%
\>[4]\AgdaSymbol{(}\AgdaFunction{pres-simultaneous-subst}\AgdaSpace{}%
\AgdaBound{fb'}\AgdaSpace{}%
\AgdaBound{f'}\AgdaSpace{}%
\AgdaBound{h}\AgdaSpace{}%
\AgdaBound{M}\AgdaSymbol{)}\<%
\\
%
\>[4]\AgdaSymbol{(}\AgdaFunction{cong}\AgdaSpace{}%
\AgdaInductiveConstructor{box}\AgdaSpace{}%
\AgdaSymbol{(}\AgdaFunction{pres-simultaneous-subst}\AgdaSpace{}%
\AgdaBound{fb'}\AgdaSpace{}%
\AgdaBound{f'}\AgdaSpace{}%
\AgdaBound{h}\AgdaSpace{}%
\AgdaBound{M₁}\AgdaSymbol{))}\<%
\\
%
\\[\AgdaEmptyExtraSkip]%
%
\\[\AgdaEmptyExtraSkip]%
\>[0]\AgdaComment{--\ Proof\ of\ h}\<%
\\
\>[0]\AgdaFunction{pres-subst-func}\AgdaSpace{}%
\AgdaSymbol{:}\AgdaSpace{}%
\AgdaSymbol{∀}\AgdaSpace{}%
\AgdaSymbol{\{}\AgdaBound{A}\AgdaSpace{}%
\AgdaBound{C}\AgdaSpace{}%
\AgdaBound{Γ}\AgdaSymbol{\}}\<%
\\
\>[0][@{}l@{\AgdaIndent{0}}]%
\>[2]\AgdaSymbol{→}\AgdaSpace{}%
\AgdaSymbol{(}\AgdaBound{N}\AgdaSpace{}%
\AgdaSymbol{:}\AgdaSpace{}%
\AgdaBound{Γ}\AgdaSpace{}%
\AgdaOperator{\AgdaDatatype{⊢}}\AgdaSpace{}%
\AgdaBound{C}\AgdaSymbol{)}\<%
\\
%
\>[2]\AgdaSymbol{→}\AgdaSpace{}%
\AgdaSymbol{(}\AgdaBound{x}\AgdaSpace{}%
\AgdaSymbol{:}\AgdaSpace{}%
\AgdaBound{A}\AgdaSpace{}%
\AgdaOperator{\AgdaDatatype{∈}}\AgdaSpace{}%
\AgdaBound{Γ}\AgdaSpace{}%
\AgdaOperator{\AgdaInductiveConstructor{,}}\AgdaSpace{}%
\AgdaBound{C}\AgdaSymbol{)}\<%
\\
%
\>[2]\AgdaSymbol{→}\AgdaSpace{}%
\AgdaFunction{fb}\AgdaSpace{}%
\AgdaSymbol{(}\AgdaFunction{embedTerm}\AgdaSpace{}%
\AgdaBound{N}\AgdaSymbol{)}\AgdaSpace{}%
\AgdaSymbol{(}\AgdaFunction{embedLookup}\AgdaSpace{}%
\AgdaBound{x}\AgdaSymbol{)}\AgdaSpace{}%
\AgdaOperator{\AgdaDatatype{≡}}\AgdaSpace{}%
\AgdaFunction{embedTerm}\AgdaSpace{}%
\AgdaSymbol{((}\AgdaFunction{f}\AgdaSpace{}%
\AgdaBound{N}\AgdaSymbol{)}\AgdaSpace{}%
\AgdaBound{x}\AgdaSymbol{)}\<%
\\
\>[0]\AgdaFunction{pres-subst-func}\AgdaSpace{}%
\AgdaBound{N}\AgdaSpace{}%
\AgdaInductiveConstructor{Z}\AgdaSpace{}%
\AgdaSymbol{=}\AgdaSpace{}%
\AgdaInductiveConstructor{refl}\<%
\\
\>[0]\AgdaFunction{pres-subst-func}\AgdaSpace{}%
\AgdaBound{N}\AgdaSpace{}%
\AgdaSymbol{(}\AgdaInductiveConstructor{S}\AgdaSpace{}%
\AgdaBound{x}\AgdaSymbol{)}\AgdaSpace{}%
\AgdaSymbol{=}\AgdaSpace{}%
\AgdaInductiveConstructor{refl}\<%
\\
\>[0]\AgdaComment{--\ Reflection\ and\ preservation\ of\ single\ substitution\ }\<%
\\
\>[0]\AgdaComment{--\ of\ refined\ Girard's\ translation\ }\<%
\\
\>[0]\AgdaFunction{pres-subst}\AgdaSpace{}%
\AgdaSymbol{:}\AgdaSpace{}%
\AgdaSymbol{∀}\AgdaSpace{}%
\AgdaSymbol{\{}\AgdaBound{Γ}\AgdaSpace{}%
\AgdaBound{A}\AgdaSpace{}%
\AgdaBound{B}\AgdaSymbol{\}}\<%
\\
\>[0][@{}l@{\AgdaIndent{0}}]%
\>[2]\AgdaSymbol{→}\AgdaSpace{}%
\AgdaSymbol{(}\AgdaBound{M}\AgdaSpace{}%
\AgdaSymbol{:}\AgdaSpace{}%
\AgdaBound{Γ}\AgdaSpace{}%
\AgdaOperator{\AgdaInductiveConstructor{,}}\AgdaSpace{}%
\AgdaBound{B}\AgdaSpace{}%
\AgdaOperator{\AgdaDatatype{⊢}}\AgdaSpace{}%
\AgdaBound{A}\AgdaSymbol{)}\<%
\\
%
\>[2]\AgdaSymbol{→}\AgdaSpace{}%
\AgdaSymbol{(}\AgdaBound{N}\AgdaSpace{}%
\AgdaSymbol{:}\AgdaSpace{}%
\AgdaBound{Γ}\AgdaSpace{}%
\AgdaOperator{\AgdaDatatype{⊢}}\AgdaSpace{}%
\AgdaBound{B}\AgdaSymbol{)}\<%
\\
%
\>[2]\AgdaSymbol{→}\AgdaSpace{}%
\AgdaSymbol{(}\AgdaFunction{embedTerm}\AgdaSpace{}%
\AgdaBound{M}\AgdaSymbol{)}\AgdaSpace{}%
\AgdaOperator{\AgdaFunction{[}}\AgdaSpace{}%
\AgdaFunction{embedTerm}\AgdaSpace{}%
\AgdaBound{N}\AgdaSpace{}%
\AgdaOperator{\AgdaFunction{]b}}\AgdaSpace{}%
\AgdaOperator{\AgdaDatatype{≡}}\AgdaSpace{}%
\AgdaFunction{embedTerm}\AgdaSpace{}%
\AgdaSymbol{(}\AgdaBound{M}\AgdaSpace{}%
\AgdaOperator{\AgdaFunction{[}}\AgdaSpace{}%
\AgdaBound{N}\AgdaSpace{}%
\AgdaOperator{\AgdaFunction{]}}\AgdaSymbol{)}\<%
\\
\>[0]\AgdaFunction{pres-subst}\AgdaSpace{}%
\AgdaBound{M}\AgdaSpace{}%
\AgdaBound{N}\AgdaSpace{}%
\AgdaSymbol{=}\<%
\\
\>[0][@{}l@{\AgdaIndent{0}}]%
\>[2]\AgdaFunction{pres-simultaneous-subst}\AgdaSpace{}%
\AgdaSymbol{(}\AgdaFunction{fb}\AgdaSpace{}%
\AgdaSymbol{(}\AgdaFunction{embedTerm}\AgdaSpace{}%
\AgdaBound{N}\AgdaSymbol{))}\AgdaSpace{}%
\AgdaSymbol{(}\AgdaFunction{f}\AgdaSpace{}%
\AgdaBound{N}\AgdaSymbol{)}\AgdaSpace{}%
\AgdaSymbol{(}\AgdaFunction{pres-subst-func}\AgdaSpace{}%
\AgdaBound{N}\AgdaSymbol{)}\AgdaSpace{}%
\AgdaBound{M}\<%
\\
%
\\[\AgdaEmptyExtraSkip]%
\>[0]\AgdaComment{--\ Preservation\ of\ reduction\ of\ Girard's\ translation}\<%
\\
\>[0]\AgdaFunction{pres-red}\AgdaSpace{}%
\AgdaSymbol{:}\AgdaSpace{}%
\AgdaSymbol{∀}\AgdaSpace{}%
\AgdaSymbol{\{}\AgdaBound{Γ}\AgdaSpace{}%
\AgdaBound{A}\AgdaSymbol{\}}\<%
\\
\>[0][@{}l@{\AgdaIndent{0}}]%
\>[2]\AgdaSymbol{→}\AgdaSpace{}%
\AgdaSymbol{(}\AgdaBound{M}\AgdaSpace{}%
\AgdaSymbol{:}\AgdaSpace{}%
\AgdaBound{Γ}\AgdaSpace{}%
\AgdaOperator{\AgdaDatatype{⊢}}\AgdaSpace{}%
\AgdaBound{A}\AgdaSymbol{)}\<%
\\
%
\>[2]\AgdaSymbol{→}\AgdaSpace{}%
\AgdaSymbol{(}\AgdaBound{N}\AgdaSpace{}%
\AgdaSymbol{:}\AgdaSpace{}%
\AgdaBound{Γ}\AgdaSpace{}%
\AgdaOperator{\AgdaDatatype{⊢}}\AgdaSpace{}%
\AgdaBound{A}\AgdaSymbol{)}\<%
\\
%
\>[2]\AgdaSymbol{→}\AgdaSpace{}%
\AgdaSymbol{(}\AgdaBound{M}\AgdaSpace{}%
\AgdaOperator{\AgdaDatatype{↝βₙ}}\AgdaSpace{}%
\AgdaBound{N}\AgdaSymbol{)}\<%
\\
%
\>[2]\AgdaSymbol{→}\AgdaSpace{}%
\AgdaSymbol{(}\AgdaFunction{embedTerm}\AgdaSpace{}%
\AgdaBound{M}\AgdaSpace{}%
\AgdaOperator{\AgdaDatatype{↝βb}}\AgdaSpace{}%
\AgdaFunction{embedTerm}\AgdaSpace{}%
\AgdaBound{N}\AgdaSymbol{)}\<%
\\
\>[0]\AgdaFunction{pres-red}\AgdaSpace{}%
\AgdaSymbol{(}\AgdaInductiveConstructor{ƛ}\AgdaSpace{}%
\AgdaBound{M₁}\AgdaSpace{}%
\AgdaOperator{\AgdaInductiveConstructor{·}}\AgdaSpace{}%
\AgdaBound{M₂}\AgdaSymbol{)}\AgdaSpace{}%
\AgdaBound{N}\AgdaSpace{}%
\AgdaInductiveConstructor{βₙ}\AgdaSpace{}%
\AgdaSymbol{=}\<%
\\
\>[0][@{}l@{\AgdaIndent{0}}]%
\>[2]\AgdaFunction{Eq.subst}\AgdaSpace{}%
\AgdaSymbol{(λ}\AgdaSpace{}%
\AgdaBound{t}\AgdaSpace{}%
\AgdaSymbol{→}\AgdaSpace{}%
\AgdaInductiveConstructor{ƛb}\AgdaSpace{}%
\AgdaSymbol{(}\AgdaFunction{embedTerm}\AgdaSpace{}%
\AgdaBound{M₁}\AgdaSymbol{)}\AgdaSpace{}%
\AgdaOperator{\AgdaInductiveConstructor{·b}}\AgdaSpace{}%
\AgdaInductiveConstructor{box}\AgdaSpace{}%
\AgdaSymbol{(}\AgdaFunction{embedTerm}\AgdaSpace{}%
\AgdaBound{M₂}\AgdaSymbol{)}\AgdaSpace{}%
\AgdaOperator{\AgdaDatatype{↝βb}}\AgdaSpace{}%
\AgdaBound{t}\AgdaSymbol{)}\<%
\\
%
\>[2]\AgdaSymbol{(}\AgdaFunction{pres-subst}\AgdaSpace{}%
\AgdaBound{M₁}\AgdaSpace{}%
\AgdaBound{M₂}\AgdaSymbol{)}\<%
\\
%
\>[2]\AgdaInductiveConstructor{βb}\<%
\\
\>[0]\AgdaFunction{pres-red}\AgdaSpace{}%
\AgdaBound{M}\AgdaSpace{}%
\AgdaBound{N}\AgdaSpace{}%
\AgdaSymbol{(}\AgdaInductiveConstructor{μ}\AgdaSpace{}%
\AgdaSymbol{\{}\AgdaArgument{M}\AgdaSpace{}%
\AgdaSymbol{=}\AgdaSpace{}%
\AgdaBound{M₁}\AgdaSymbol{\}}\AgdaSpace{}%
\AgdaSymbol{\{}\AgdaArgument{M'}\AgdaSpace{}%
\AgdaSymbol{=}\AgdaSpace{}%
\AgdaBound{M'}\AgdaSymbol{\}}\AgdaSpace{}%
\AgdaBound{red}\AgdaSymbol{)}\AgdaSpace{}%
\AgdaSymbol{=}\<%
\\
\>[0][@{}l@{\AgdaIndent{0}}]%
\>[2]\AgdaInductiveConstructor{μ}\AgdaSpace{}%
\AgdaSymbol{(}\AgdaFunction{pres-red}\AgdaSpace{}%
\AgdaBound{M₁}\AgdaSpace{}%
\AgdaBound{M'}\AgdaSpace{}%
\AgdaBound{red}\AgdaSymbol{)}\<%
\\
\>[0]\AgdaFunction{pres-red}\AgdaSpace{}%
\AgdaBound{M}\AgdaSpace{}%
\AgdaBound{N}\AgdaSpace{}%
\AgdaSymbol{(}\AgdaInductiveConstructor{ν<}\AgdaSpace{}%
\AgdaSymbol{\{}\AgdaArgument{M}\AgdaSpace{}%
\AgdaSymbol{=}\AgdaSpace{}%
\AgdaBound{M₁}\AgdaSymbol{\}}\AgdaSpace{}%
\AgdaSymbol{\{}\AgdaArgument{N}\AgdaSpace{}%
\AgdaSymbol{=}\AgdaSpace{}%
\AgdaBound{N₁}\AgdaSymbol{\}}\AgdaSpace{}%
\AgdaSymbol{\{}\AgdaArgument{N'}\AgdaSpace{}%
\AgdaSymbol{=}\AgdaSpace{}%
\AgdaBound{N'}\AgdaSymbol{\}}\AgdaSpace{}%
\AgdaBound{x}\AgdaSpace{}%
\AgdaBound{red}\AgdaSymbol{)}\AgdaSpace{}%
\AgdaSymbol{=}\<%
\\
\>[0][@{}l@{\AgdaIndent{0}}]%
\>[2]\AgdaInductiveConstructor{ν}\AgdaSpace{}%
\AgdaSymbol{(}\AgdaInductiveConstructor{ζ}\AgdaSpace{}%
\AgdaSymbol{(}\AgdaFunction{pres-red}\AgdaSpace{}%
\AgdaBound{N₁}\AgdaSpace{}%
\AgdaBound{N'}\AgdaSpace{}%
\AgdaBound{red}\AgdaSymbol{))}\<%
\\
\>[0]\AgdaFunction{pres-red}\AgdaSpace{}%
\AgdaBound{M}\AgdaSpace{}%
\AgdaBound{N}\AgdaSpace{}%
\AgdaSymbol{(}\AgdaInductiveConstructor{ν}\AgdaSpace{}%
\AgdaSymbol{\{}\AgdaArgument{M}\AgdaSpace{}%
\AgdaSymbol{=}\AgdaSpace{}%
\AgdaBound{M₁}\AgdaSymbol{\}}\AgdaSpace{}%
\AgdaSymbol{\{}\AgdaArgument{N}\AgdaSpace{}%
\AgdaSymbol{=}\AgdaSpace{}%
\AgdaBound{N₁}\AgdaSymbol{\}}\AgdaSpace{}%
\AgdaSymbol{\{}\AgdaArgument{N'}\AgdaSpace{}%
\AgdaSymbol{=}\AgdaSpace{}%
\AgdaBound{N'}\AgdaSymbol{\}}\AgdaSpace{}%
\AgdaBound{red}\AgdaSymbol{)}\AgdaSpace{}%
\AgdaSymbol{=}\<%
\\
\>[0][@{}l@{\AgdaIndent{0}}]%
\>[2]\AgdaInductiveConstructor{ν}\AgdaSpace{}%
\AgdaSymbol{(}\AgdaInductiveConstructor{ζ}\AgdaSpace{}%
\AgdaSymbol{(}\AgdaFunction{pres-red}\AgdaSpace{}%
\AgdaBound{N₁}\AgdaSpace{}%
\AgdaBound{N'}\AgdaSpace{}%
\AgdaBound{red}\AgdaSymbol{))}\<%
\\
\>[0]\AgdaFunction{pres-red}\AgdaSpace{}%
\AgdaBound{M}\AgdaSpace{}%
\AgdaBound{N}\AgdaSpace{}%
\AgdaSymbol{(}\AgdaInductiveConstructor{ξ}\AgdaSpace{}%
\AgdaSymbol{\{}\AgdaArgument{M}\AgdaSpace{}%
\AgdaSymbol{=}\AgdaSpace{}%
\AgdaBound{M₁}\AgdaSymbol{\}}\AgdaSpace{}%
\AgdaSymbol{\{}\AgdaArgument{M'}\AgdaSpace{}%
\AgdaSymbol{=}\AgdaSpace{}%
\AgdaBound{M'}\AgdaSymbol{\}}\AgdaSpace{}%
\AgdaBound{red}\AgdaSymbol{)}\AgdaSpace{}%
\AgdaSymbol{=}\<%
\\
\>[0][@{}l@{\AgdaIndent{0}}]%
\>[2]\AgdaInductiveConstructor{ξ}\AgdaSpace{}%
\AgdaSymbol{(}\AgdaFunction{pres-red}\AgdaSpace{}%
\AgdaBound{M₁}\AgdaSpace{}%
\AgdaBound{M'}\AgdaSpace{}%
\AgdaBound{red}\AgdaSymbol{)}\<%
\\
\>[0]\<%
\end{code}
\begin{code}%
\>[0]\AgdaComment{--\ Preservation\ of\ evaluation\ of\ Girard's\ translation}\<%
\\
\>[0]\AgdaFunction{pres-eval}\AgdaSpace{}%
\AgdaSymbol{:}\AgdaSpace{}%
\AgdaSymbol{∀}\AgdaSpace{}%
\AgdaSymbol{\{}\AgdaBound{Γ}\AgdaSpace{}%
\AgdaBound{A}\AgdaSymbol{\}}\<%
\\
\>[0][@{}l@{\AgdaIndent{0}}]%
\>[2]\AgdaSymbol{→}\AgdaSpace{}%
\AgdaSymbol{(}\AgdaBound{M}\AgdaSpace{}%
\AgdaSymbol{:}\AgdaSpace{}%
\AgdaBound{Γ}\AgdaSpace{}%
\AgdaOperator{\AgdaDatatype{⊢}}\AgdaSpace{}%
\AgdaBound{A}\AgdaSymbol{)}\<%
\\
%
\>[2]\AgdaSymbol{→}\AgdaSpace{}%
\AgdaSymbol{(}\AgdaBound{N}\AgdaSpace{}%
\AgdaSymbol{:}\AgdaSpace{}%
\AgdaBound{Γ}\AgdaSpace{}%
\AgdaOperator{\AgdaDatatype{⊢}}\AgdaSpace{}%
\AgdaBound{A}\AgdaSymbol{)}\<%
\\
%
\>[2]\AgdaSymbol{→}\AgdaSpace{}%
\AgdaSymbol{(}\AgdaBound{M}\AgdaSpace{}%
\AgdaOperator{\AgdaDatatype{↝ₙ}}\AgdaSpace{}%
\AgdaBound{N}\AgdaSymbol{)}\<%
\\
%
\>[2]\AgdaSymbol{→}\AgdaSpace{}%
\AgdaSymbol{(}\AgdaFunction{embedTerm}\AgdaSpace{}%
\AgdaBound{M}\AgdaSpace{}%
\AgdaOperator{\AgdaDatatype{↝bₙ}}\AgdaSpace{}%
\AgdaFunction{embedTerm}\AgdaSpace{}%
\AgdaBound{N}\AgdaSymbol{)}\<%
\\
\>[0]\AgdaFunction{pres-eval}\AgdaSpace{}%
\AgdaSymbol{(}\AgdaInductiveConstructor{ƛ}\AgdaSpace{}%
\AgdaBound{M₁}\AgdaSpace{}%
\AgdaOperator{\AgdaInductiveConstructor{·}}\AgdaSpace{}%
\AgdaBound{M₂}\AgdaSymbol{)}\AgdaSpace{}%
\AgdaBound{N}\AgdaSpace{}%
\AgdaInductiveConstructor{βₙ}\AgdaSpace{}%
\AgdaSymbol{=}\<%
\\
\>[0][@{}l@{\AgdaIndent{0}}]%
\>[2]\AgdaFunction{Eq.subst}\<%
\\
\>[2][@{}l@{\AgdaIndent{0}}]%
\>[4]\AgdaSymbol{(λ}\AgdaSpace{}%
\AgdaBound{t}\AgdaSpace{}%
\AgdaSymbol{→}\AgdaSpace{}%
\AgdaInductiveConstructor{ƛb}\AgdaSpace{}%
\AgdaSymbol{(}\AgdaFunction{embedTerm}\AgdaSpace{}%
\AgdaBound{M₁}\AgdaSymbol{)}\AgdaSpace{}%
\AgdaOperator{\AgdaInductiveConstructor{·b}}\AgdaSpace{}%
\AgdaInductiveConstructor{box}\AgdaSpace{}%
\AgdaSymbol{(}\AgdaFunction{embedTerm}\AgdaSpace{}%
\AgdaBound{M₂}\AgdaSymbol{)}\AgdaSpace{}%
\AgdaOperator{\AgdaDatatype{↝bₙ}}\AgdaSpace{}%
\AgdaBound{t}\AgdaSymbol{)}\<%
\\
%
\>[4]\AgdaSymbol{(}\AgdaFunction{pres-subst}\AgdaSpace{}%
\AgdaBound{M₁}\AgdaSpace{}%
\AgdaBound{M₂}\AgdaSymbol{)}\<%
\\
%
\>[4]\AgdaInductiveConstructor{βb}\<%
\\
\>[0]\AgdaFunction{pres-eval}\AgdaSpace{}%
\AgdaBound{M}\AgdaSpace{}%
\AgdaBound{N}\AgdaSpace{}%
\AgdaSymbol{(}\AgdaInductiveConstructor{μ}\AgdaSpace{}%
\AgdaSymbol{\{}\AgdaArgument{M}\AgdaSpace{}%
\AgdaSymbol{=}\AgdaSpace{}%
\AgdaBound{M₁}\AgdaSymbol{\}}\AgdaSpace{}%
\AgdaSymbol{\{}\AgdaArgument{M'}\AgdaSpace{}%
\AgdaSymbol{=}\AgdaSpace{}%
\AgdaBound{M'}\AgdaSymbol{\}}\AgdaSpace{}%
\AgdaBound{red}\AgdaSymbol{)}\AgdaSpace{}%
\AgdaSymbol{=}\<%
\\
\>[0][@{}l@{\AgdaIndent{0}}]%
\>[2]\AgdaInductiveConstructor{μ}\AgdaSpace{}%
\AgdaSymbol{(}\AgdaFunction{pres-eval}\AgdaSpace{}%
\AgdaBound{M₁}\AgdaSpace{}%
\AgdaBound{M'}\AgdaSpace{}%
\AgdaBound{red}\AgdaSymbol{)}\<%
\end{code}
  \hfill $\qed$
\end{proofthm}

Here, we give the proofs of \cref{prop:godel-preservation-red} and \cref{prop:godel-preservation-red2s}.
These proofs are more tricky than the proofs regarding Girard's translation because of the use of \textsf{renameb} inside \raiseembn.
Some helper lemmas that we needed are discussed in the \cref{sec:helper-lemmas}.

\begin{proofthm}
  \label{proof:gogel-preservation-2step-red}
  \textbf{Preservation of reduction} under Gödel's translation with $\redbboxts$ reduction rule.
  \small
  \begin{code}[hide]%
\>[0]\AgdaKeyword{module}\AgdaSpace{}%
\AgdaModule{valueProofs4}\AgdaSpace{}%
\AgdaKeyword{where}\<%
\\
\>[0]\<%
\\
\>[0]\AgdaKeyword{open}\AgdaSpace{}%
\AgdaKeyword{import}\AgdaSpace{}%
\AgdaModule{calc}\<%
\\
\>[0]\AgdaKeyword{open}\AgdaSpace{}%
\AgdaKeyword{import}\AgdaSpace{}%
\AgdaModule{boxCalc}\<%
\\
\>[0]\AgdaKeyword{open}\AgdaSpace{}%
\AgdaKeyword{import}\AgdaSpace{}%
\AgdaModule{embedCBVIntoCBB}\<%
\\
\>[0]\AgdaKeyword{open}\AgdaSpace{}%
\AgdaKeyword{import}\AgdaSpace{}%
\AgdaModule{valueCalcRed}\<%
\\
\>[0]\AgdaKeyword{open}\AgdaSpace{}%
\AgdaKeyword{import}\AgdaSpace{}%
\AgdaModule{boxCalcRed}\<%
\\
%
\\[\AgdaEmptyExtraSkip]%
\>[0]\AgdaKeyword{open}\AgdaSpace{}%
\AgdaKeyword{import}\AgdaSpace{}%
\AgdaModule{Data.Empty}\AgdaSpace{}%
\AgdaKeyword{using}\AgdaSpace{}%
\AgdaSymbol{(}\AgdaFunction{⊥}\AgdaSymbol{;}\AgdaSpace{}%
\AgdaFunction{⊥-elim}\AgdaSymbol{)}\<%
\\
\>[0]\AgdaKeyword{import}\AgdaSpace{}%
\AgdaModule{Relation.Binary.PropositionalEquality}\AgdaSpace{}%
\AgdaSymbol{as}\AgdaSpace{}%
\AgdaModule{Eq}\<%
\\
\>[0]\AgdaKeyword{open}\AgdaSpace{}%
\AgdaModule{Eq}\AgdaSpace{}%
\AgdaKeyword{using}\AgdaSpace{}%
\AgdaSymbol{(}\AgdaOperator{\AgdaDatatype{\AgdaUnderscore{}≡\AgdaUnderscore{}}}\AgdaSymbol{;}\AgdaSpace{}%
\AgdaInductiveConstructor{refl}\AgdaSymbol{;}\AgdaSpace{}%
\AgdaFunction{cong}\AgdaSymbol{;}\AgdaSpace{}%
\AgdaFunction{cong₂}\AgdaSymbol{;}\AgdaSpace{}%
\AgdaFunction{sym}\AgdaSymbol{)}\<%
\\
%
\\[\AgdaEmptyExtraSkip]%
\>[0]\AgdaKeyword{open}\AgdaSpace{}%
\AgdaModule{Eq.≡-Reasoning}\AgdaSpace{}%
\AgdaKeyword{using}\AgdaSpace{}%
\AgdaSymbol{(}\AgdaOperator{\AgdaFunction{begin\AgdaUnderscore{}}}\AgdaSymbol{;}\AgdaSpace{}%
\AgdaFunction{step-≡-∣}\AgdaSymbol{;}\AgdaSpace{}%
\AgdaFunction{step-≡-⟩}\AgdaSymbol{;}\AgdaSpace{}%
\AgdaOperator{\AgdaFunction{\AgdaUnderscore{}∎}}\AgdaSymbol{)}\<%
\\
\>[0]\AgdaKeyword{open}\AgdaSpace{}%
\AgdaModule{Eq.≡-Reasoning}\AgdaSpace{}%
\AgdaKeyword{renaming}\AgdaSpace{}%
\AgdaSymbol{(}\AgdaOperator{\AgdaFunction{begin\AgdaUnderscore{}}}\AgdaSpace{}%
\AgdaSymbol{to}\AgdaSpace{}%
\AgdaOperator{\AgdaFunction{beginEq\AgdaUnderscore{}}}\AgdaSymbol{;}\AgdaSpace{}%
\AgdaOperator{\AgdaFunction{\AgdaUnderscore{}∎}}\AgdaSpace{}%
\AgdaSymbol{to}\AgdaSpace{}%
\AgdaOperator{\AgdaFunction{\AgdaUnderscore{}∎eq}}\AgdaSymbol{)}\<%
\\
%
\\[\AgdaEmptyExtraSkip]%
\>[0]\AgdaComment{--\ Preservation\ of\ typing\ by\ Gödel's\ translation}\<%
\\
\>[0]\AgdaFunction{pres-typing}\AgdaSpace{}%
\AgdaSymbol{:}\AgdaSpace{}%
\AgdaSymbol{∀}\AgdaSpace{}%
\AgdaSymbol{\{}\AgdaBound{Γ}\AgdaSpace{}%
\AgdaBound{A}\AgdaSymbol{\}}\AgdaSpace{}%
\AgdaSymbol{→}\AgdaSpace{}%
\AgdaBound{Γ}\AgdaSpace{}%
\AgdaOperator{\AgdaDatatype{⊢}}\AgdaSpace{}%
\AgdaBound{A}\AgdaSpace{}%
\AgdaSymbol{→}\AgdaSpace{}%
\AgdaFunction{embedContext}\AgdaSpace{}%
\AgdaBound{Γ}\AgdaSpace{}%
\AgdaOperator{\AgdaDatatype{⊢b}}\AgdaSpace{}%
\AgdaFunction{embedType}\AgdaSpace{}%
\AgdaBound{A}\<%
\\
\>[0]\AgdaFunction{pres-typing}\AgdaSpace{}%
\AgdaBound{x}\AgdaSpace{}%
\AgdaSymbol{=}\AgdaSpace{}%
\AgdaFunction{embedTerm}\AgdaSpace{}%
\AgdaBound{x}\<%
\\
%
\\[\AgdaEmptyExtraSkip]%
\>[0]\AgdaComment{--\ Preservation\ of\ typing\ of\ Values\ by\ Gödel's\ translation}\<%
\\
\>[0]\AgdaFunction{pres-typing-val}\AgdaSpace{}%
\AgdaSymbol{:}\AgdaSpace{}%
\AgdaSymbol{∀}\AgdaSpace{}%
\AgdaSymbol{\{}\AgdaBound{Γ}\AgdaSpace{}%
\AgdaBound{A}\AgdaSymbol{\}}\AgdaSpace{}%
\AgdaSymbol{→}\AgdaSpace{}%
\AgdaSymbol{(}\AgdaBound{V}\AgdaSpace{}%
\AgdaSymbol{:}\AgdaSpace{}%
\AgdaBound{Γ}\AgdaSpace{}%
\AgdaOperator{\AgdaDatatype{⊢}}\AgdaSpace{}%
\AgdaBound{A}\AgdaSymbol{)}\AgdaSpace{}%
\AgdaSymbol{→}\AgdaSpace{}%
\AgdaSymbol{(}\AgdaDatatype{Value}\AgdaSpace{}%
\AgdaBound{V}\AgdaSymbol{)}\AgdaSpace{}%
\AgdaSymbol{→}\AgdaSpace{}%
\AgdaFunction{embedContext}\AgdaSpace{}%
\AgdaBound{Γ}\AgdaSpace{}%
\AgdaOperator{\AgdaDatatype{⊢b}}\AgdaSpace{}%
\AgdaFunction{embedTypeValue}\AgdaSpace{}%
\AgdaBound{A}\<%
\\
\>[0]\AgdaFunction{pres-typing-val}\AgdaSpace{}%
\AgdaBound{V}\AgdaSpace{}%
\AgdaBound{x}\AgdaSpace{}%
\AgdaSymbol{=}\AgdaSpace{}%
\AgdaFunction{embedValue}\AgdaSpace{}%
\AgdaBound{V}\AgdaSpace{}%
\AgdaBound{x}\<%
\\
%
\\[\AgdaEmptyExtraSkip]%
\>[0]\AgdaComment{--\ Helper\ lemma:\ applying\ embedTerm\ to\ a\ Value\ is\ equivalent\ to\ wrapping\ the\ result\ of\ embedValue\ in\ a\ box}\<%
\\
\>[0]\AgdaFunction{box-embedValue}\AgdaSpace{}%
\AgdaSymbol{:}\AgdaSpace{}%
\AgdaSymbol{∀}\AgdaSpace{}%
\AgdaSymbol{\{}\AgdaBound{Γ}\AgdaSpace{}%
\AgdaBound{A}\AgdaSymbol{\}}\AgdaSpace{}%
\AgdaSymbol{→}\AgdaSpace{}%
\AgdaSymbol{\{}\AgdaBound{V}\AgdaSpace{}%
\AgdaSymbol{:}\AgdaSpace{}%
\AgdaBound{Γ}\AgdaSpace{}%
\AgdaOperator{\AgdaDatatype{⊢}}\AgdaSpace{}%
\AgdaBound{A}\AgdaSymbol{\}}\AgdaSpace{}%
\AgdaSymbol{→}\AgdaSpace{}%
\AgdaSymbol{(}\AgdaBound{x}\AgdaSpace{}%
\AgdaSymbol{:}\AgdaSpace{}%
\AgdaDatatype{Value}\AgdaSpace{}%
\AgdaBound{V}\AgdaSymbol{)}\AgdaSpace{}%
\AgdaSymbol{→}\AgdaSpace{}%
\AgdaInductiveConstructor{box}\AgdaSpace{}%
\AgdaSymbol{(}\AgdaFunction{embedValue}\AgdaSpace{}%
\AgdaBound{V}\AgdaSpace{}%
\AgdaBound{x}\AgdaSymbol{)}\AgdaSpace{}%
\AgdaOperator{\AgdaDatatype{≡}}\AgdaSpace{}%
\AgdaFunction{embedTerm}\AgdaSpace{}%
\AgdaBound{V}\<%
\\
\>[0]\AgdaFunction{box-embedValue}\AgdaSpace{}%
\AgdaInductiveConstructor{V-ƛ}\AgdaSpace{}%
\AgdaSymbol{=}\AgdaSpace{}%
\AgdaInductiveConstructor{refl}\<%
\\
\>[0]\AgdaFunction{box-embedValue}\AgdaSpace{}%
\AgdaInductiveConstructor{V-`}\AgdaSpace{}%
\AgdaSymbol{=}\AgdaSpace{}%
\AgdaInductiveConstructor{refl}\<%
\\
%
\\[\AgdaEmptyExtraSkip]%
\>[0]\AgdaComment{--\ HELPER\ LEMMA'S\ SIMULTANIOUS\ SUBSTITUTION}\<%
\\
%
\\[\AgdaEmptyExtraSkip]%
\>[0]\AgdaComment{--\ Preservation\ and\ Reflection\ of\ }\<%
\\
\>[0]\AgdaComment{--\ ext\ of\ Gödel's\ translation}\<%
\\
\>[0]\AgdaFunction{pres-ext}\AgdaSpace{}%
\AgdaSymbol{:}\AgdaSpace{}%
\AgdaSymbol{∀}\AgdaSpace{}%
\AgdaSymbol{\{}\AgdaBound{Γ}\AgdaSpace{}%
\AgdaBound{Δ}\AgdaSpace{}%
\AgdaBound{A'}\AgdaSpace{}%
\AgdaBound{A}\AgdaSymbol{\}}\<%
\\
\>[0][@{}l@{\AgdaIndent{0}}]%
\>[2]\AgdaSymbol{→}\AgdaSpace{}%
\AgdaSymbol{(}\AgdaBound{f'}\AgdaSpace{}%
\AgdaSymbol{:}\AgdaSpace{}%
\AgdaSymbol{∀}\AgdaSpace{}%
\AgdaSymbol{\{}\AgdaBound{A}\AgdaSymbol{\}}\AgdaSpace{}%
\AgdaSymbol{→}\AgdaSpace{}%
\AgdaBound{A}\AgdaSpace{}%
\AgdaOperator{\AgdaDatatype{∈}}\AgdaSpace{}%
\AgdaBound{Γ}\AgdaSpace{}%
\AgdaSymbol{→}\AgdaSpace{}%
\AgdaBound{A}\AgdaSpace{}%
\AgdaOperator{\AgdaDatatype{∈}}\AgdaSpace{}%
\AgdaBound{Δ}\AgdaSymbol{)}\<%
\\
%
\>[2]\AgdaSymbol{→}\AgdaSpace{}%
\AgdaSymbol{(}\AgdaBound{fb'}\AgdaSpace{}%
\AgdaSymbol{:}\AgdaSpace{}%
\AgdaSymbol{∀}\AgdaSpace{}%
\AgdaSymbol{\{}\AgdaBound{Ab}\AgdaSymbol{\}}\AgdaSpace{}%
\AgdaSymbol{→}\AgdaSpace{}%
\AgdaBound{Ab}\AgdaSpace{}%
\AgdaOperator{\AgdaDatatype{∈b}}\AgdaSpace{}%
\AgdaFunction{embedContext}\AgdaSpace{}%
\AgdaBound{Γ}\AgdaSpace{}%
\AgdaSymbol{→}\AgdaSpace{}%
\AgdaBound{Ab}\AgdaSpace{}%
\AgdaOperator{\AgdaDatatype{∈b}}\AgdaSpace{}%
\AgdaFunction{embedContext}\AgdaSpace{}%
\AgdaBound{Δ}\AgdaSymbol{)}\<%
\\
%
\>[2]\AgdaSymbol{→}\AgdaSpace{}%
\AgdaSymbol{(\{}\AgdaBound{A'}\AgdaSpace{}%
\AgdaSymbol{:}\AgdaSpace{}%
\AgdaDatatype{Type}\AgdaSymbol{\}}\AgdaSpace{}%
\AgdaSymbol{(}\AgdaBound{x}\AgdaSpace{}%
\AgdaSymbol{:}\AgdaSpace{}%
\AgdaBound{A'}\AgdaSpace{}%
\AgdaOperator{\AgdaDatatype{∈}}\AgdaSpace{}%
\AgdaBound{Γ}\AgdaSymbol{)}\AgdaSpace{}%
\AgdaSymbol{→}\AgdaSpace{}%
\AgdaBound{fb'}\AgdaSpace{}%
\AgdaSymbol{(}\AgdaFunction{embedLookup}\AgdaSpace{}%
\AgdaBound{x}\AgdaSymbol{)}\AgdaSpace{}%
\AgdaOperator{\AgdaDatatype{≡}}\AgdaSpace{}%
\AgdaFunction{embedLookup}\AgdaSpace{}%
\AgdaSymbol{(}\AgdaBound{f'}\AgdaSpace{}%
\AgdaBound{x}\AgdaSymbol{))}\<%
\\
%
\>[2]\AgdaSymbol{→}\AgdaSpace{}%
\AgdaSymbol{(}\AgdaBound{x}\AgdaSpace{}%
\AgdaSymbol{:}\AgdaSpace{}%
\AgdaBound{A'}\AgdaSpace{}%
\AgdaOperator{\AgdaDatatype{∈}}\AgdaSpace{}%
\AgdaBound{Γ}\AgdaSpace{}%
\AgdaOperator{\AgdaInductiveConstructor{,}}\AgdaSpace{}%
\AgdaBound{A}\AgdaSymbol{)}\<%
\\
%
\>[2]\AgdaSymbol{→}\AgdaSpace{}%
\AgdaFunction{extb}\AgdaSpace{}%
\AgdaBound{fb'}\AgdaSpace{}%
\AgdaSymbol{(}\AgdaFunction{embedLookup}\AgdaSpace{}%
\AgdaBound{x}\AgdaSymbol{)}\AgdaSpace{}%
\AgdaOperator{\AgdaDatatype{≡}}\AgdaSpace{}%
\AgdaFunction{embedLookup}\AgdaSpace{}%
\AgdaSymbol{(}\AgdaFunction{ext}\AgdaSpace{}%
\AgdaBound{f'}\AgdaSpace{}%
\AgdaBound{x}\AgdaSymbol{)}\<%
\\
\>[0]\AgdaFunction{pres-ext}\AgdaSpace{}%
\AgdaBound{f'}\AgdaSpace{}%
\AgdaBound{fb'}\AgdaSpace{}%
\AgdaBound{h}\AgdaSpace{}%
\AgdaInductiveConstructor{Z}\AgdaSpace{}%
\AgdaSymbol{=}\AgdaSpace{}%
\AgdaInductiveConstructor{refl}\<%
\\
\>[0]\AgdaFunction{pres-ext}\AgdaSpace{}%
\AgdaBound{f'}\AgdaSpace{}%
\AgdaBound{fb'}\AgdaSpace{}%
\AgdaBound{h}\AgdaSpace{}%
\AgdaSymbol{(}\AgdaInductiveConstructor{S}\AgdaSpace{}%
\AgdaBound{x}\AgdaSymbol{)}\AgdaSpace{}%
\AgdaSymbol{=}\AgdaSpace{}%
\AgdaFunction{cong}\AgdaSpace{}%
\AgdaInductiveConstructor{Sb}\AgdaSpace{}%
\AgdaSymbol{(}\AgdaBound{h}\AgdaSpace{}%
\AgdaBound{x}\AgdaSymbol{)}\<%
\\
%
\\[\AgdaEmptyExtraSkip]%
\>[0]\AgdaComment{--\ Extb\ preserves\ pointwise\ equality\ }\<%
\\
\>[0]\AgdaComment{--\ of\ renaming\ functions}\<%
\\
\>[0]\AgdaFunction{extb-cong}\AgdaSpace{}%
\AgdaSymbol{:}\AgdaSpace{}%
\AgdaSymbol{∀}\AgdaSpace{}%
\AgdaSymbol{\{}\AgdaBound{Γ₁}\AgdaSpace{}%
\AgdaBound{Γ₂}\AgdaSpace{}%
\AgdaBound{A}\AgdaSpace{}%
\AgdaBound{A'}\AgdaSymbol{\}}\<%
\\
\>[0][@{}l@{\AgdaIndent{0}}]%
\>[2]\AgdaSymbol{→}\AgdaSpace{}%
\AgdaSymbol{\{}\AgdaBound{fb'}\AgdaSpace{}%
\AgdaSymbol{:}\AgdaSpace{}%
\AgdaSymbol{∀}\AgdaSpace{}%
\AgdaSymbol{\{}\AgdaBound{Ab}\AgdaSymbol{\}}\AgdaSpace{}%
\AgdaSymbol{→}\AgdaSpace{}%
\AgdaBound{Ab}\AgdaSpace{}%
\AgdaOperator{\AgdaDatatype{∈b}}\AgdaSpace{}%
\AgdaBound{Γ₁}\AgdaSpace{}%
\AgdaSymbol{→}\AgdaSpace{}%
\AgdaBound{Ab}\AgdaSpace{}%
\AgdaOperator{\AgdaDatatype{∈b}}\AgdaSpace{}%
\AgdaBound{Γ₂}\AgdaSymbol{\}}\<%
\\
%
\>[2]\AgdaSymbol{→}\AgdaSpace{}%
\AgdaSymbol{\{}\AgdaBound{fb''}\AgdaSpace{}%
\AgdaSymbol{:}\AgdaSpace{}%
\AgdaSymbol{∀}\AgdaSpace{}%
\AgdaSymbol{\{}\AgdaBound{Ab}\AgdaSymbol{\}}\AgdaSpace{}%
\AgdaSymbol{→}\AgdaSpace{}%
\AgdaBound{Ab}\AgdaSpace{}%
\AgdaOperator{\AgdaDatatype{∈b}}\AgdaSpace{}%
\AgdaBound{Γ₁}\AgdaSpace{}%
\AgdaSymbol{→}\AgdaSpace{}%
\AgdaBound{Ab}\AgdaSpace{}%
\AgdaOperator{\AgdaDatatype{∈b}}\AgdaSpace{}%
\AgdaBound{Γ₂}\AgdaSymbol{\}}\<%
\\
%
\>[2]\AgdaSymbol{→}\AgdaSpace{}%
\AgdaSymbol{(}\AgdaBound{h}\AgdaSpace{}%
\AgdaSymbol{:}\AgdaSpace{}%
\AgdaSymbol{∀}\AgdaSpace{}%
\AgdaSymbol{\{}\AgdaBound{Ab}\AgdaSymbol{\}}\AgdaSpace{}%
\AgdaSymbol{→}\AgdaSpace{}%
\AgdaSymbol{(}\AgdaBound{x}\AgdaSpace{}%
\AgdaSymbol{:}\AgdaSpace{}%
\AgdaBound{Ab}\AgdaSpace{}%
\AgdaOperator{\AgdaDatatype{∈b}}\AgdaSpace{}%
\AgdaBound{Γ₁}\AgdaSymbol{)}\AgdaSpace{}%
\AgdaSymbol{→}\AgdaSpace{}%
\AgdaBound{fb'}\AgdaSpace{}%
\AgdaBound{x}\AgdaSpace{}%
\AgdaOperator{\AgdaDatatype{≡}}\AgdaSpace{}%
\AgdaBound{fb''}\AgdaSpace{}%
\AgdaBound{x}\AgdaSymbol{)}\<%
\\
%
\>[2]\AgdaSymbol{→}\AgdaSpace{}%
\AgdaSymbol{(}\AgdaBound{x}\AgdaSpace{}%
\AgdaSymbol{:}\AgdaSpace{}%
\AgdaBound{A}\AgdaSpace{}%
\AgdaOperator{\AgdaDatatype{∈b}}\AgdaSpace{}%
\AgdaBound{Γ₁}\AgdaSpace{}%
\AgdaOperator{\AgdaInductiveConstructor{,b}}\AgdaSpace{}%
\AgdaInductiveConstructor{□}\AgdaSpace{}%
\AgdaBound{A'}\AgdaSymbol{)}\<%
\\
%
\>[2]\AgdaSymbol{→}\AgdaSpace{}%
\AgdaFunction{extb}\AgdaSpace{}%
\AgdaBound{fb'}\AgdaSpace{}%
\AgdaBound{x}\AgdaSpace{}%
\AgdaOperator{\AgdaDatatype{≡}}\AgdaSpace{}%
\AgdaFunction{extb}\AgdaSpace{}%
\AgdaBound{fb''}\AgdaSpace{}%
\AgdaBound{x}\<%
\\
\>[0]\AgdaFunction{extb-cong}\AgdaSpace{}%
\AgdaBound{h}\AgdaSpace{}%
\AgdaInductiveConstructor{Zb}\AgdaSpace{}%
\AgdaSymbol{=}\AgdaSpace{}%
\AgdaInductiveConstructor{refl}\<%
\\
\>[0]\AgdaFunction{extb-cong}\AgdaSpace{}%
\AgdaBound{h}\AgdaSpace{}%
\AgdaSymbol{(}\AgdaInductiveConstructor{Sb}\AgdaSpace{}%
\AgdaBound{x}\AgdaSymbol{)}\AgdaSpace{}%
\AgdaSymbol{=}\AgdaSpace{}%
\AgdaFunction{cong}\AgdaSpace{}%
\AgdaInductiveConstructor{Sb}\AgdaSpace{}%
\AgdaSymbol{(}\AgdaBound{h}\AgdaSpace{}%
\AgdaBound{x}\AgdaSymbol{)}\<%
\\
%
\\[\AgdaEmptyExtraSkip]%
\>[0]\AgdaComment{--\ Renameb\ preserves\ pointwise\ equality}\<%
\\
\>[0]\AgdaComment{--\ of\ renaming\ functions}\<%
\\
\>[0]\AgdaFunction{renameb-cong}\AgdaSpace{}%
\AgdaSymbol{:}\AgdaSpace{}%
\AgdaSymbol{∀}\AgdaSpace{}%
\AgdaSymbol{\{}\AgdaBound{Γ₁}\AgdaSpace{}%
\AgdaBound{Γ₂}\AgdaSpace{}%
\AgdaBound{A}\AgdaSymbol{\}}\<%
\\
\>[0][@{}l@{\AgdaIndent{0}}]%
\>[2]\AgdaSymbol{→}\AgdaSpace{}%
\AgdaSymbol{\{}\AgdaBound{fb'}\AgdaSpace{}%
\AgdaSymbol{:}\AgdaSpace{}%
\AgdaSymbol{∀}\AgdaSpace{}%
\AgdaSymbol{\{}\AgdaBound{Ab}\AgdaSymbol{\}}\AgdaSpace{}%
\AgdaSymbol{→}\AgdaSpace{}%
\AgdaBound{Ab}\AgdaSpace{}%
\AgdaOperator{\AgdaDatatype{∈b}}\AgdaSpace{}%
\AgdaBound{Γ₁}\AgdaSpace{}%
\AgdaSymbol{→}\AgdaSpace{}%
\AgdaBound{Ab}\AgdaSpace{}%
\AgdaOperator{\AgdaDatatype{∈b}}\AgdaSpace{}%
\AgdaBound{Γ₂}\AgdaSymbol{\}}\<%
\\
%
\>[2]\AgdaSymbol{→}\AgdaSpace{}%
\AgdaSymbol{\{}\AgdaBound{fb''}\AgdaSpace{}%
\AgdaSymbol{:}\AgdaSpace{}%
\AgdaSymbol{∀}\AgdaSpace{}%
\AgdaSymbol{\{}\AgdaBound{Ab}\AgdaSymbol{\}}\AgdaSpace{}%
\AgdaSymbol{→}\AgdaSpace{}%
\AgdaBound{Ab}\AgdaSpace{}%
\AgdaOperator{\AgdaDatatype{∈b}}\AgdaSpace{}%
\AgdaBound{Γ₁}\AgdaSpace{}%
\AgdaSymbol{→}\AgdaSpace{}%
\AgdaBound{Ab}\AgdaSpace{}%
\AgdaOperator{\AgdaDatatype{∈b}}\AgdaSpace{}%
\AgdaBound{Γ₂}\AgdaSymbol{\}}\<%
\\
%
\>[2]\AgdaSymbol{→}\AgdaSpace{}%
\AgdaSymbol{(}\AgdaBound{h}\AgdaSpace{}%
\AgdaSymbol{:}\AgdaSpace{}%
\AgdaSymbol{∀}\AgdaSpace{}%
\AgdaSymbol{\{}\AgdaBound{Ab}\AgdaSymbol{\}}\AgdaSpace{}%
\AgdaSymbol{→}\AgdaSpace{}%
\AgdaSymbol{(}\AgdaBound{x}\AgdaSpace{}%
\AgdaSymbol{:}\AgdaSpace{}%
\AgdaBound{Ab}\AgdaSpace{}%
\AgdaOperator{\AgdaDatatype{∈b}}\AgdaSpace{}%
\AgdaBound{Γ₁}\AgdaSymbol{)}\AgdaSpace{}%
\AgdaSymbol{→}\AgdaSpace{}%
\AgdaBound{fb'}\AgdaSpace{}%
\AgdaBound{x}\AgdaSpace{}%
\AgdaOperator{\AgdaDatatype{≡}}\AgdaSpace{}%
\AgdaBound{fb''}\AgdaSpace{}%
\AgdaBound{x}\AgdaSymbol{)}\<%
\\
%
\>[2]\AgdaSymbol{→}\AgdaSpace{}%
\AgdaSymbol{(}\AgdaBound{t}\AgdaSpace{}%
\AgdaSymbol{:}\AgdaSpace{}%
\AgdaBound{Γ₁}\AgdaSpace{}%
\AgdaOperator{\AgdaDatatype{⊢b}}\AgdaSpace{}%
\AgdaBound{A}\AgdaSymbol{)}\<%
\\
%
\>[2]\AgdaSymbol{→}\AgdaSpace{}%
\AgdaFunction{renameb}\AgdaSpace{}%
\AgdaBound{fb'}\AgdaSpace{}%
\AgdaBound{t}\AgdaSpace{}%
\AgdaOperator{\AgdaDatatype{≡}}\AgdaSpace{}%
\AgdaFunction{renameb}\AgdaSpace{}%
\AgdaBound{fb''}\AgdaSpace{}%
\AgdaBound{t}\<%
\\
\>[0]\AgdaFunction{renameb-cong}\AgdaSpace{}%
\AgdaBound{h}\AgdaSpace{}%
\AgdaSymbol{(}\AgdaInductiveConstructor{ε}\AgdaSpace{}%
\AgdaBound{x}\AgdaSymbol{)}\AgdaSpace{}%
\AgdaSymbol{=}\AgdaSpace{}%
\AgdaFunction{cong}\AgdaSpace{}%
\AgdaInductiveConstructor{ε}\AgdaSpace{}%
\AgdaSymbol{(}\AgdaBound{h}\AgdaSpace{}%
\AgdaBound{x}\AgdaSymbol{)}\<%
\\
\>[0]\AgdaFunction{renameb-cong}\AgdaSpace{}%
\AgdaBound{h}\AgdaSpace{}%
\AgdaSymbol{(}\AgdaInductiveConstructor{box}\AgdaSpace{}%
\AgdaBound{t}\AgdaSymbol{)}\AgdaSpace{}%
\AgdaSymbol{=}\AgdaSpace{}%
\AgdaFunction{cong}\AgdaSpace{}%
\AgdaInductiveConstructor{box}\AgdaSpace{}%
\AgdaSymbol{(}\AgdaFunction{renameb-cong}\AgdaSpace{}%
\AgdaBound{h}\AgdaSpace{}%
\AgdaBound{t}\AgdaSymbol{)}\<%
\\
\>[0]\AgdaFunction{renameb-cong}\AgdaSpace{}%
\AgdaBound{h}\AgdaSpace{}%
\AgdaSymbol{(}\AgdaInductiveConstructor{ƛb}\AgdaSpace{}%
\AgdaBound{t}\AgdaSymbol{)}\AgdaSpace{}%
\AgdaSymbol{=}\AgdaSpace{}%
\AgdaFunction{cong}\AgdaSpace{}%
\AgdaInductiveConstructor{ƛb}\AgdaSpace{}%
\AgdaSymbol{(}\AgdaFunction{renameb-cong}\AgdaSpace{}%
\AgdaSymbol{(}\AgdaFunction{extb-cong}\AgdaSpace{}%
\AgdaBound{h}\AgdaSymbol{)}\AgdaSpace{}%
\AgdaBound{t}\AgdaSymbol{)}\<%
\\
\>[0]\AgdaFunction{renameb-cong}\AgdaSpace{}%
\AgdaBound{h}\AgdaSpace{}%
\AgdaSymbol{(}\AgdaBound{t}\AgdaSpace{}%
\AgdaOperator{\AgdaInductiveConstructor{·b}}\AgdaSpace{}%
\AgdaBound{t₁}\AgdaSymbol{)}\AgdaSpace{}%
\AgdaSymbol{=}\AgdaSpace{}%
\AgdaFunction{cong₂}\AgdaSpace{}%
\AgdaOperator{\AgdaInductiveConstructor{\AgdaUnderscore{}·b\AgdaUnderscore{}}}\AgdaSpace{}%
\AgdaSymbol{(}\AgdaFunction{renameb-cong}\AgdaSpace{}%
\AgdaBound{h}\AgdaSpace{}%
\AgdaBound{t}\AgdaSymbol{)}\AgdaSpace{}%
\AgdaSymbol{(}\AgdaFunction{renameb-cong}\AgdaSpace{}%
\AgdaBound{h}\AgdaSpace{}%
\AgdaBound{t₁}\AgdaSymbol{)}\<%
\\
%
\\[\AgdaEmptyExtraSkip]%
\>[0]\AgdaComment{--\ extb\ distributes\ over\ composition:}\<%
\\
\>[0]\AgdaComment{--\ Extending\ first\ with\ fb',\ then\ with\ fb'',\ }\<%
\\
\>[0]\AgdaComment{--\ is\ the\ same\ as\ extending\ once\ with\ fb''\ ∘\ fb'.}\<%
\\
\>[0]\AgdaFunction{extb-comp}\AgdaSpace{}%
\AgdaSymbol{:}\AgdaSpace{}%
\AgdaSymbol{∀}\AgdaSpace{}%
\AgdaSymbol{\{}\AgdaBound{Γ₁}\AgdaSpace{}%
\AgdaBound{Γ₂}\AgdaSpace{}%
\AgdaBound{Γ₃}\AgdaSpace{}%
\AgdaBound{A}\AgdaSpace{}%
\AgdaBound{B}\AgdaSymbol{\}}\<%
\\
\>[0][@{}l@{\AgdaIndent{0}}]%
\>[2]\AgdaSymbol{→}\AgdaSpace{}%
\AgdaSymbol{(}\AgdaBound{fb'}\AgdaSpace{}%
\AgdaSymbol{:}\AgdaSpace{}%
\AgdaSymbol{∀}\AgdaSpace{}%
\AgdaSymbol{\{}\AgdaBound{Ab}\AgdaSymbol{\}}\AgdaSpace{}%
\AgdaSymbol{→}\AgdaSpace{}%
\AgdaBound{Ab}\AgdaSpace{}%
\AgdaOperator{\AgdaDatatype{∈b}}\AgdaSpace{}%
\AgdaBound{Γ₁}\AgdaSpace{}%
\AgdaSymbol{→}\AgdaSpace{}%
\AgdaBound{Ab}\AgdaSpace{}%
\AgdaOperator{\AgdaDatatype{∈b}}\AgdaSpace{}%
\AgdaBound{Γ₂}\AgdaSymbol{)}\<%
\\
%
\>[2]\AgdaSymbol{→}\AgdaSpace{}%
\AgdaSymbol{(}\AgdaBound{fb''}\AgdaSpace{}%
\AgdaSymbol{:}\AgdaSpace{}%
\AgdaSymbol{∀}\AgdaSpace{}%
\AgdaSymbol{\{}\AgdaBound{Ab}\AgdaSymbol{\}}\AgdaSpace{}%
\AgdaSymbol{→}\AgdaSpace{}%
\AgdaBound{Ab}\AgdaSpace{}%
\AgdaOperator{\AgdaDatatype{∈b}}\AgdaSpace{}%
\AgdaBound{Γ₂}\AgdaSpace{}%
\AgdaSymbol{→}\AgdaSpace{}%
\AgdaBound{Ab}\AgdaSpace{}%
\AgdaOperator{\AgdaDatatype{∈b}}\AgdaSpace{}%
\AgdaBound{Γ₃}\AgdaSymbol{)}\<%
\\
%
\>[2]\AgdaSymbol{→}\AgdaSpace{}%
\AgdaSymbol{(}\AgdaBound{x}\AgdaSpace{}%
\AgdaSymbol{:}\AgdaSpace{}%
\AgdaBound{A}\AgdaSpace{}%
\AgdaOperator{\AgdaDatatype{∈b}}\AgdaSpace{}%
\AgdaBound{Γ₁}\AgdaSpace{}%
\AgdaOperator{\AgdaInductiveConstructor{,b}}\AgdaSpace{}%
\AgdaBound{B}\AgdaSymbol{)}\<%
\\
%
\>[2]\AgdaSymbol{→}\AgdaSpace{}%
\AgdaFunction{extb}\AgdaSpace{}%
\AgdaBound{fb''}\AgdaSpace{}%
\AgdaSymbol{(}\AgdaFunction{extb}\AgdaSpace{}%
\AgdaBound{fb'}\AgdaSpace{}%
\AgdaBound{x}\AgdaSymbol{)}\AgdaSpace{}%
\AgdaOperator{\AgdaDatatype{≡}}\AgdaSpace{}%
\AgdaFunction{extb}\AgdaSpace{}%
\AgdaSymbol{(λ}\AgdaSpace{}%
\AgdaBound{x₁}\AgdaSpace{}%
\AgdaSymbol{→}\AgdaSpace{}%
\AgdaBound{fb''}\AgdaSpace{}%
\AgdaSymbol{(}\AgdaBound{fb'}\AgdaSpace{}%
\AgdaBound{x₁}\AgdaSymbol{))}\AgdaSpace{}%
\AgdaBound{x}\<%
\\
\>[0]\AgdaFunction{extb-comp}\AgdaSpace{}%
\AgdaBound{fb'}\AgdaSpace{}%
\AgdaBound{fb''}\AgdaSpace{}%
\AgdaInductiveConstructor{Zb}\AgdaSpace{}%
\AgdaSymbol{=}\AgdaSpace{}%
\AgdaInductiveConstructor{refl}\<%
\\
\>[0]\AgdaFunction{extb-comp}\AgdaSpace{}%
\AgdaBound{fb'}\AgdaSpace{}%
\AgdaBound{fb''}\AgdaSpace{}%
\AgdaSymbol{(}\AgdaInductiveConstructor{Sb}\AgdaSpace{}%
\AgdaBound{x}\AgdaSymbol{)}\AgdaSpace{}%
\AgdaSymbol{=}\AgdaSpace{}%
\AgdaInductiveConstructor{refl}\<%
\\
%
\\[\AgdaEmptyExtraSkip]%
\>[0]\AgdaComment{--\ renameb\ distributes\ over\ composition:}\<%
\\
\>[0]\AgdaComment{--\ Renaming\ first\ with\ fb',\ then\ with\ fb''}\<%
\\
\>[0]\AgdaComment{--\ is\ the\ same\ as\ renaming\ once\ with\ fb''\ ∘\ fb'.}\<%
\\
\>[0]\AgdaFunction{renameb-comp}\AgdaSpace{}%
\AgdaSymbol{:}\AgdaSpace{}%
\AgdaSymbol{∀}\AgdaSpace{}%
\AgdaSymbol{\{}\AgdaBound{Γ₁}\AgdaSpace{}%
\AgdaBound{Γ₂}\AgdaSpace{}%
\AgdaBound{Γ₃}\AgdaSpace{}%
\AgdaBound{A}\AgdaSymbol{\}}\<%
\\
\>[0][@{}l@{\AgdaIndent{0}}]%
\>[2]\AgdaSymbol{→}\AgdaSpace{}%
\AgdaSymbol{(}\AgdaBound{fb'}\AgdaSpace{}%
\AgdaSymbol{:}\AgdaSpace{}%
\AgdaSymbol{∀}\AgdaSpace{}%
\AgdaSymbol{\{}\AgdaBound{Ab}\AgdaSymbol{\}}\AgdaSpace{}%
\AgdaSymbol{→}\AgdaSpace{}%
\AgdaBound{Ab}\AgdaSpace{}%
\AgdaOperator{\AgdaDatatype{∈b}}\AgdaSpace{}%
\AgdaBound{Γ₁}\AgdaSpace{}%
\AgdaSymbol{→}\AgdaSpace{}%
\AgdaBound{Ab}\AgdaSpace{}%
\AgdaOperator{\AgdaDatatype{∈b}}\AgdaSpace{}%
\AgdaBound{Γ₂}\AgdaSymbol{)}\<%
\\
%
\>[2]\AgdaSymbol{→}\AgdaSpace{}%
\AgdaSymbol{(}\AgdaBound{fb''}\AgdaSpace{}%
\AgdaSymbol{:}\AgdaSpace{}%
\AgdaSymbol{∀}\AgdaSpace{}%
\AgdaSymbol{\{}\AgdaBound{Ab}\AgdaSymbol{\}}\AgdaSpace{}%
\AgdaSymbol{→}\AgdaSpace{}%
\AgdaBound{Ab}\AgdaSpace{}%
\AgdaOperator{\AgdaDatatype{∈b}}\AgdaSpace{}%
\AgdaBound{Γ₂}\AgdaSpace{}%
\AgdaSymbol{→}\AgdaSpace{}%
\AgdaBound{Ab}\AgdaSpace{}%
\AgdaOperator{\AgdaDatatype{∈b}}\AgdaSpace{}%
\AgdaBound{Γ₃}\AgdaSymbol{)}\<%
\\
%
\>[2]\AgdaSymbol{→}\AgdaSpace{}%
\AgdaSymbol{(}\AgdaBound{t}\AgdaSpace{}%
\AgdaSymbol{:}\AgdaSpace{}%
\AgdaBound{Γ₁}\AgdaSpace{}%
\AgdaOperator{\AgdaDatatype{⊢b}}\AgdaSpace{}%
\AgdaBound{A}\AgdaSymbol{)}\<%
\\
%
\>[2]\AgdaSymbol{→}\AgdaSpace{}%
\AgdaFunction{renameb}\AgdaSpace{}%
\AgdaBound{fb''}\AgdaSpace{}%
\AgdaSymbol{(}\AgdaFunction{renameb}\AgdaSpace{}%
\AgdaBound{fb'}\AgdaSpace{}%
\AgdaBound{t}\AgdaSymbol{)}\AgdaSpace{}%
\AgdaOperator{\AgdaDatatype{≡}}\AgdaSpace{}%
\AgdaFunction{renameb}\AgdaSpace{}%
\AgdaSymbol{(λ}\AgdaSpace{}%
\AgdaBound{x}\AgdaSpace{}%
\AgdaSymbol{→}\AgdaSpace{}%
\AgdaBound{fb''}\AgdaSpace{}%
\AgdaSymbol{(}\AgdaBound{fb'}\AgdaSpace{}%
\AgdaBound{x}\AgdaSymbol{))}\AgdaSpace{}%
\AgdaBound{t}\<%
\\
\>[0]\AgdaFunction{renameb-comp}\AgdaSpace{}%
\AgdaBound{fb'}\AgdaSpace{}%
\AgdaBound{fb''}\AgdaSpace{}%
\AgdaSymbol{(}\AgdaInductiveConstructor{ε}\AgdaSpace{}%
\AgdaBound{x}\AgdaSymbol{)}\AgdaSpace{}%
\AgdaSymbol{=}\AgdaSpace{}%
\AgdaInductiveConstructor{refl}\<%
\\
\>[0]\AgdaFunction{renameb-comp}\AgdaSpace{}%
\AgdaBound{fb'}\AgdaSpace{}%
\AgdaBound{fb''}\AgdaSpace{}%
\AgdaSymbol{(}\AgdaInductiveConstructor{box}\AgdaSpace{}%
\AgdaBound{t}\AgdaSymbol{)}\AgdaSpace{}%
\AgdaSymbol{=}\AgdaSpace{}%
\AgdaFunction{cong}\AgdaSpace{}%
\AgdaInductiveConstructor{box}\AgdaSpace{}%
\AgdaSymbol{(}\AgdaFunction{renameb-comp}\AgdaSpace{}%
\AgdaBound{fb'}\AgdaSpace{}%
\AgdaBound{fb''}\AgdaSpace{}%
\AgdaBound{t}\AgdaSymbol{)}\<%
\\
\>[0]\AgdaFunction{renameb-comp}\AgdaSpace{}%
\AgdaBound{fb'}\AgdaSpace{}%
\AgdaBound{fb''}\AgdaSpace{}%
\AgdaSymbol{(}\AgdaInductiveConstructor{ƛb}\AgdaSpace{}%
\AgdaBound{t}\AgdaSymbol{)}\AgdaSpace{}%
\AgdaSymbol{=}\<%
\\
\>[0][@{}l@{\AgdaIndent{0}}]%
\>[2]\AgdaOperator{\AgdaFunction{beginEq}}\<%
\\
\>[2][@{}l@{\AgdaIndent{0}}]%
\>[4]\AgdaInductiveConstructor{ƛb}\AgdaSpace{}%
\AgdaSymbol{(}\AgdaFunction{renameb}\AgdaSpace{}%
\AgdaSymbol{(}\AgdaFunction{extb}\AgdaSpace{}%
\AgdaBound{fb''}\AgdaSymbol{)}\AgdaSpace{}%
\AgdaSymbol{(}\AgdaFunction{renameb}\AgdaSpace{}%
\AgdaSymbol{(}\AgdaFunction{extb}\AgdaSpace{}%
\AgdaBound{fb'}\AgdaSymbol{)}\AgdaSpace{}%
\AgdaBound{t}\AgdaSymbol{))}\<%
\\
%
\>[2]\AgdaFunction{≡⟨}\AgdaSpace{}%
\AgdaFunction{cong}\AgdaSpace{}%
\AgdaInductiveConstructor{ƛb}\AgdaSpace{}%
\AgdaSymbol{(}\AgdaFunction{renameb-comp}\AgdaSpace{}%
\AgdaSymbol{(}\AgdaFunction{extb}\AgdaSpace{}%
\AgdaBound{fb'}\AgdaSymbol{)}\AgdaSpace{}%
\AgdaSymbol{(}\AgdaFunction{extb}\AgdaSpace{}%
\AgdaBound{fb''}\AgdaSymbol{)}\AgdaSpace{}%
\AgdaBound{t}\AgdaSymbol{)}\AgdaSpace{}%
\AgdaFunction{⟩}\<%
\\
\>[2][@{}l@{\AgdaIndent{0}}]%
\>[4]\AgdaInductiveConstructor{ƛb}\AgdaSpace{}%
\AgdaSymbol{(}\AgdaFunction{renameb}\AgdaSpace{}%
\AgdaSymbol{(λ}\AgdaSpace{}%
\AgdaBound{x}\AgdaSpace{}%
\AgdaSymbol{→}\AgdaSpace{}%
\AgdaFunction{extb}\AgdaSpace{}%
\AgdaBound{fb''}\AgdaSpace{}%
\AgdaSymbol{(}\AgdaFunction{extb}\AgdaSpace{}%
\AgdaBound{fb'}\AgdaSpace{}%
\AgdaBound{x}\AgdaSymbol{))}\AgdaSpace{}%
\AgdaBound{t}\AgdaSymbol{)}\<%
\\
%
\>[2]\AgdaFunction{≡⟨}\AgdaSpace{}%
\AgdaFunction{cong}\AgdaSpace{}%
\AgdaInductiveConstructor{ƛb}\AgdaSpace{}%
\AgdaSymbol{(}\AgdaFunction{renameb-cong}\AgdaSpace{}%
\AgdaSymbol{(}\AgdaFunction{extb-comp}\AgdaSpace{}%
\AgdaBound{fb'}\AgdaSpace{}%
\AgdaBound{fb''}\AgdaSymbol{)}\AgdaSpace{}%
\AgdaBound{t}\AgdaSymbol{)}\AgdaSpace{}%
\AgdaFunction{⟩}\<%
\\
\>[2][@{}l@{\AgdaIndent{0}}]%
\>[4]\AgdaInductiveConstructor{ƛb}\AgdaSpace{}%
\AgdaSymbol{(}\AgdaFunction{renameb}\AgdaSpace{}%
\AgdaSymbol{(}\AgdaFunction{extb}\AgdaSpace{}%
\AgdaSymbol{(λ}\AgdaSpace{}%
\AgdaBound{x}\AgdaSpace{}%
\AgdaSymbol{→}\AgdaSpace{}%
\AgdaBound{fb''}\AgdaSpace{}%
\AgdaSymbol{(}\AgdaBound{fb'}\AgdaSpace{}%
\AgdaBound{x}\AgdaSymbol{)))}\AgdaSpace{}%
\AgdaBound{t}\AgdaSymbol{)}\<%
\\
%
\>[2]\AgdaOperator{\AgdaFunction{∎eq}}\<%
\\
\>[0]\AgdaFunction{renameb-comp}\AgdaSpace{}%
\AgdaBound{fb'}\AgdaSpace{}%
\AgdaBound{fb''}\AgdaSpace{}%
\AgdaSymbol{(}\AgdaBound{t}\AgdaSpace{}%
\AgdaOperator{\AgdaInductiveConstructor{·b}}\AgdaSpace{}%
\AgdaBound{t₁}\AgdaSymbol{)}\AgdaSpace{}%
\AgdaSymbol{=}\AgdaSpace{}%
\AgdaFunction{cong₂}\AgdaSpace{}%
\AgdaOperator{\AgdaInductiveConstructor{\AgdaUnderscore{}·b\AgdaUnderscore{}}}\AgdaSpace{}%
\AgdaSymbol{(}\AgdaFunction{renameb-comp}\AgdaSpace{}%
\AgdaBound{fb'}\AgdaSpace{}%
\AgdaBound{fb''}\AgdaSpace{}%
\AgdaBound{t}\AgdaSymbol{)}\AgdaSpace{}%
\AgdaSymbol{(}\AgdaFunction{renameb-comp}\AgdaSpace{}%
\AgdaBound{fb'}\AgdaSpace{}%
\AgdaBound{fb''}\AgdaSpace{}%
\AgdaBound{t₁}\AgdaSymbol{)}\<%
\\
%
\\[\AgdaEmptyExtraSkip]%
\>[0]\AgdaComment{--\ extb\ and\ Sb\ commute\ under\ renameb.}\<%
\\
\>[0]\AgdaFunction{rename-extb-Sb-commute}\AgdaSpace{}%
\AgdaSymbol{:}\AgdaSpace{}%
\AgdaSymbol{∀}\AgdaSpace{}%
\AgdaSymbol{\{}\AgdaBound{Γ}\AgdaSpace{}%
\AgdaBound{A}\AgdaSpace{}%
\AgdaBound{B}\AgdaSpace{}%
\AgdaBound{Δ}\AgdaSymbol{\}}\<%
\\
\>[0][@{}l@{\AgdaIndent{0}}]%
\>[2]\AgdaSymbol{→}\AgdaSpace{}%
\AgdaSymbol{(}\AgdaBound{fb'}\AgdaSpace{}%
\AgdaSymbol{:}\AgdaSpace{}%
\AgdaSymbol{∀}\AgdaSpace{}%
\AgdaSymbol{\{}\AgdaBound{Ab}\AgdaSymbol{\}}\AgdaSpace{}%
\AgdaSymbol{→}\AgdaSpace{}%
\AgdaBound{Ab}\AgdaSpace{}%
\AgdaOperator{\AgdaDatatype{∈b}}\AgdaSpace{}%
\AgdaBound{Γ}\AgdaSpace{}%
\AgdaSymbol{→}\AgdaSpace{}%
\AgdaBound{Ab}\AgdaSpace{}%
\AgdaOperator{\AgdaDatatype{∈b}}\AgdaSpace{}%
\AgdaBound{Δ}\AgdaSymbol{)}\<%
\\
%
\>[2]\AgdaSymbol{→}\AgdaSpace{}%
\AgdaSymbol{(}\AgdaBound{t}\AgdaSpace{}%
\AgdaSymbol{:}\AgdaSpace{}%
\AgdaBound{Γ}\AgdaSpace{}%
\AgdaOperator{\AgdaDatatype{⊢b}}\AgdaSpace{}%
\AgdaBound{A}\AgdaSymbol{)}\<%
\\
%
\>[2]\AgdaSymbol{→}%
\>[579I]\AgdaFunction{renameb}\AgdaSpace{}%
\AgdaSymbol{(}\AgdaFunction{extb}\AgdaSpace{}%
\AgdaBound{fb'}\AgdaSpace{}%
\AgdaSymbol{\{}\AgdaArgument{B}\AgdaSpace{}%
\AgdaSymbol{=}\AgdaSpace{}%
\AgdaBound{B}\AgdaSymbol{\})}\AgdaSpace{}%
\AgdaSymbol{(}\AgdaFunction{renameb}\AgdaSpace{}%
\AgdaInductiveConstructor{Sb}\AgdaSpace{}%
\AgdaBound{t}\AgdaSymbol{)}\<%
\\
\>[.][@{}l@{}]\<[579I]%
\>[4]\AgdaOperator{\AgdaDatatype{≡}}\AgdaSpace{}%
\AgdaFunction{renameb}\AgdaSpace{}%
\AgdaInductiveConstructor{Sb}\AgdaSpace{}%
\AgdaSymbol{(}\AgdaFunction{renameb}\AgdaSpace{}%
\AgdaBound{fb'}\AgdaSpace{}%
\AgdaBound{t}\AgdaSymbol{)}\<%
\\
\>[0]\AgdaFunction{rename-extb-Sb-commute}\AgdaSpace{}%
\AgdaBound{fb'}\AgdaSpace{}%
\AgdaBound{t}\AgdaSpace{}%
\AgdaSymbol{=}\<%
\\
\>[0][@{}l@{\AgdaIndent{0}}]%
\>[2]\AgdaOperator{\AgdaFunction{beginEq}}\<%
\\
\>[2][@{}l@{\AgdaIndent{0}}]%
\>[4]\AgdaFunction{renameb}\AgdaSpace{}%
\AgdaSymbol{(}\AgdaFunction{extb}\AgdaSpace{}%
\AgdaBound{fb'}\AgdaSymbol{)}\AgdaSpace{}%
\AgdaSymbol{(}\AgdaFunction{renameb}\AgdaSpace{}%
\AgdaInductiveConstructor{Sb}\AgdaSpace{}%
\AgdaBound{t}\AgdaSymbol{)}\<%
\\
%
\>[2]\AgdaFunction{≡⟨}\AgdaSpace{}%
\AgdaFunction{renameb-comp}\AgdaSpace{}%
\AgdaInductiveConstructor{Sb}\AgdaSpace{}%
\AgdaSymbol{(}\AgdaFunction{extb}\AgdaSpace{}%
\AgdaBound{fb'}\AgdaSymbol{)}\AgdaSpace{}%
\AgdaBound{t}\AgdaSpace{}%
\AgdaFunction{⟩}\<%
\\
\>[2][@{}l@{\AgdaIndent{0}}]%
\>[4]\AgdaFunction{renameb}\AgdaSpace{}%
\AgdaSymbol{(λ}\AgdaSpace{}%
\AgdaBound{x}\AgdaSpace{}%
\AgdaSymbol{→}\AgdaSpace{}%
\AgdaFunction{extb}\AgdaSpace{}%
\AgdaBound{fb'}\AgdaSpace{}%
\AgdaSymbol{(}\AgdaInductiveConstructor{Sb}\AgdaSpace{}%
\AgdaBound{x}\AgdaSymbol{))}\AgdaSpace{}%
\AgdaBound{t}\<%
\\
%
\>[2]\AgdaFunction{≡⟨}\AgdaSpace{}%
\AgdaFunction{renameb-cong}\AgdaSpace{}%
\AgdaSymbol{(λ}\AgdaSpace{}%
\AgdaBound{x}\AgdaSpace{}%
\AgdaSymbol{→}\AgdaSpace{}%
\AgdaInductiveConstructor{refl}\AgdaSymbol{)}\AgdaSpace{}%
\AgdaBound{t}\AgdaSpace{}%
\AgdaFunction{⟩}\<%
\\
\>[2][@{}l@{\AgdaIndent{0}}]%
\>[4]\AgdaFunction{renameb}\AgdaSpace{}%
\AgdaSymbol{(λ}\AgdaSpace{}%
\AgdaBound{x}\AgdaSpace{}%
\AgdaSymbol{→}\AgdaSpace{}%
\AgdaInductiveConstructor{Sb}\AgdaSpace{}%
\AgdaSymbol{(}\AgdaBound{fb'}\AgdaSpace{}%
\AgdaBound{x}\AgdaSymbol{))}\AgdaSpace{}%
\AgdaBound{t}\<%
\\
%
\>[2]\AgdaFunction{≡⟨}\AgdaSpace{}%
\AgdaFunction{sym}\AgdaSpace{}%
\AgdaSymbol{(}\AgdaFunction{renameb-comp}\AgdaSpace{}%
\AgdaBound{fb'}\AgdaSpace{}%
\AgdaInductiveConstructor{Sb}\AgdaSpace{}%
\AgdaBound{t}\AgdaSymbol{)}\AgdaSpace{}%
\AgdaFunction{⟩}\<%
\\
\>[2][@{}l@{\AgdaIndent{0}}]%
\>[4]\AgdaFunction{renameb}\AgdaSpace{}%
\AgdaInductiveConstructor{Sb}\AgdaSpace{}%
\AgdaSymbol{(}\AgdaFunction{renameb}\AgdaSpace{}%
\AgdaBound{fb'}\AgdaSpace{}%
\AgdaBound{t}\AgdaSymbol{)}\<%
\\
%
\>[2]\AgdaOperator{\AgdaFunction{∎eq}}\<%
\\
%
\\[\AgdaEmptyExtraSkip]%
\>[0]\AgdaComment{--\ Preservation\ and\ Reflection\ of\ }\<%
\\
\>[0]\AgdaComment{--\ rename\ of\ Gödel's\ translation}\<%
\\
\>[0]\AgdaFunction{pres-rename}\AgdaSpace{}%
\AgdaSymbol{:}\AgdaSpace{}%
\AgdaSymbol{∀}\AgdaSpace{}%
\AgdaSymbol{\{}\AgdaBound{Γ}\AgdaSpace{}%
\AgdaBound{A}\AgdaSpace{}%
\AgdaBound{Δ}\AgdaSymbol{\}}\<%
\\
\>[0][@{}l@{\AgdaIndent{0}}]%
\>[2]\AgdaSymbol{→}\AgdaSpace{}%
\AgdaSymbol{(}\AgdaBound{f'}\AgdaSpace{}%
\AgdaSymbol{:}\AgdaSpace{}%
\AgdaSymbol{∀}\AgdaSpace{}%
\AgdaSymbol{\{}\AgdaBound{A}\AgdaSymbol{\}}\AgdaSpace{}%
\AgdaSymbol{→}\AgdaSpace{}%
\AgdaBound{A}\AgdaSpace{}%
\AgdaOperator{\AgdaDatatype{∈}}\AgdaSpace{}%
\AgdaBound{Γ}\AgdaSpace{}%
\AgdaSymbol{→}\AgdaSpace{}%
\AgdaBound{A}\AgdaSpace{}%
\AgdaOperator{\AgdaDatatype{∈}}\AgdaSpace{}%
\AgdaBound{Δ}\AgdaSymbol{)}\<%
\\
%
\>[2]\AgdaSymbol{→}\AgdaSpace{}%
\AgdaSymbol{(}\AgdaBound{fb'}\AgdaSpace{}%
\AgdaSymbol{:}\AgdaSpace{}%
\AgdaSymbol{∀}\AgdaSpace{}%
\AgdaSymbol{\{}\AgdaBound{Ab}\AgdaSymbol{\}}\AgdaSpace{}%
\AgdaSymbol{→}\AgdaSpace{}%
\AgdaBound{Ab}\AgdaSpace{}%
\AgdaOperator{\AgdaDatatype{∈b}}\AgdaSpace{}%
\AgdaFunction{embedContext}\AgdaSpace{}%
\AgdaBound{Γ}\AgdaSpace{}%
\AgdaSymbol{→}\AgdaSpace{}%
\AgdaBound{Ab}\AgdaSpace{}%
\AgdaOperator{\AgdaDatatype{∈b}}\AgdaSpace{}%
\AgdaFunction{embedContext}\AgdaSpace{}%
\AgdaBound{Δ}\AgdaSymbol{)}\<%
\\
%
\>[2]\AgdaSymbol{→}\AgdaSpace{}%
\AgdaSymbol{(\{}\AgdaBound{A'}\AgdaSpace{}%
\AgdaSymbol{:}\AgdaSpace{}%
\AgdaDatatype{Type}\AgdaSymbol{\}}\AgdaSpace{}%
\AgdaSymbol{(}\AgdaBound{x}\AgdaSpace{}%
\AgdaSymbol{:}\AgdaSpace{}%
\AgdaBound{A'}\AgdaSpace{}%
\AgdaOperator{\AgdaDatatype{∈}}\AgdaSpace{}%
\AgdaBound{Γ}\AgdaSymbol{)}\AgdaSpace{}%
\AgdaSymbol{→}\AgdaSpace{}%
\AgdaBound{fb'}\AgdaSpace{}%
\AgdaSymbol{(}\AgdaFunction{embedLookup}\AgdaSpace{}%
\AgdaBound{x}\AgdaSymbol{)}\AgdaSpace{}%
\AgdaOperator{\AgdaDatatype{≡}}\AgdaSpace{}%
\AgdaFunction{embedLookup}\AgdaSpace{}%
\AgdaSymbol{(}\AgdaBound{f'}\AgdaSpace{}%
\AgdaBound{x}\AgdaSymbol{))}\<%
\\
%
\>[2]\AgdaSymbol{→}\AgdaSpace{}%
\AgdaSymbol{(}\AgdaBound{t}\AgdaSpace{}%
\AgdaSymbol{:}\AgdaSpace{}%
\AgdaBound{Γ}\AgdaSpace{}%
\AgdaOperator{\AgdaDatatype{⊢}}\AgdaSpace{}%
\AgdaBound{A}\AgdaSymbol{)}\<%
\\
%
\>[2]\AgdaSymbol{→}\AgdaSpace{}%
\AgdaFunction{renameb}\AgdaSpace{}%
\AgdaBound{fb'}\AgdaSpace{}%
\AgdaSymbol{(}\AgdaFunction{embedTerm}\AgdaSpace{}%
\AgdaBound{t}\AgdaSymbol{)}\AgdaSpace{}%
\AgdaOperator{\AgdaDatatype{≡}}\AgdaSpace{}%
\AgdaFunction{embedTerm}\AgdaSpace{}%
\AgdaSymbol{(}\AgdaFunction{rename}\AgdaSpace{}%
\AgdaBound{f'}\AgdaSpace{}%
\AgdaBound{t}\AgdaSymbol{)}\<%
\\
\>[0]\AgdaFunction{pres-rename}\AgdaSpace{}%
\AgdaBound{f'}\AgdaSpace{}%
\AgdaBound{fb'}\AgdaSpace{}%
\AgdaBound{h}\AgdaSpace{}%
\AgdaSymbol{(}\AgdaInductiveConstructor{`}\AgdaSpace{}%
\AgdaBound{x}\AgdaSymbol{)}\AgdaSpace{}%
\AgdaSymbol{=}\AgdaSpace{}%
\AgdaFunction{cong}\AgdaSpace{}%
\AgdaInductiveConstructor{box}\AgdaSpace{}%
\AgdaSymbol{(}\AgdaFunction{cong}\AgdaSpace{}%
\AgdaInductiveConstructor{ε}\AgdaSpace{}%
\AgdaSymbol{(}\AgdaBound{h}\AgdaSpace{}%
\AgdaBound{x}\AgdaSymbol{))}\<%
\\
\>[0]\AgdaFunction{pres-rename}\AgdaSpace{}%
\AgdaBound{f'}\AgdaSpace{}%
\AgdaBound{fb'}\AgdaSpace{}%
\AgdaBound{h}\AgdaSpace{}%
\AgdaSymbol{(}\AgdaInductiveConstructor{ƛ}\AgdaSpace{}%
\AgdaBound{t}\AgdaSymbol{)}\AgdaSpace{}%
\AgdaSymbol{=}\AgdaSpace{}%
\AgdaFunction{cong}\AgdaSpace{}%
\AgdaInductiveConstructor{box}\AgdaSpace{}%
\AgdaSymbol{(}\AgdaFunction{cong}\AgdaSpace{}%
\AgdaInductiveConstructor{ƛb}\AgdaSpace{}%
\AgdaSymbol{(}\AgdaFunction{pres-rename}\AgdaSpace{}%
\AgdaSymbol{(}\AgdaFunction{calc.ext}\AgdaSpace{}%
\AgdaBound{f'}\AgdaSymbol{)}\AgdaSpace{}%
\AgdaSymbol{(}\AgdaFunction{boxCalc.extb}\AgdaSpace{}%
\AgdaBound{fb'}\AgdaSymbol{)}\AgdaSpace{}%
\AgdaSymbol{(}\AgdaFunction{pres-ext}\AgdaSpace{}%
\AgdaBound{f'}\AgdaSpace{}%
\AgdaBound{fb'}\AgdaSpace{}%
\AgdaBound{h}\AgdaSymbol{)}\AgdaSpace{}%
\AgdaBound{t}\AgdaSymbol{))}\<%
\\
\>[0]\AgdaFunction{pres-rename}\AgdaSpace{}%
\AgdaBound{f'}\AgdaSpace{}%
\AgdaBound{fb'}\AgdaSpace{}%
\AgdaBound{h}\AgdaSpace{}%
\AgdaSymbol{(}\AgdaBound{t}\AgdaSpace{}%
\AgdaOperator{\AgdaInductiveConstructor{·}}\AgdaSpace{}%
\AgdaBound{t₁}\AgdaSymbol{)}\AgdaSpace{}%
\AgdaSymbol{=}\<%
\\
\>[0][@{}l@{\AgdaIndent{0}}]%
\>[2]\AgdaOperator{\AgdaFunction{beginEq}}\<%
\\
\>[2][@{}l@{\AgdaIndent{0}}]%
\>[4]\AgdaInductiveConstructor{ƛb}\AgdaSpace{}%
\AgdaSymbol{(}\AgdaInductiveConstructor{ε}\AgdaSpace{}%
\AgdaInductiveConstructor{Zb}\AgdaSpace{}%
\AgdaOperator{\AgdaInductiveConstructor{·b}}\AgdaSpace{}%
\AgdaFunction{renameb}\AgdaSpace{}%
\AgdaSymbol{(}\AgdaFunction{extb}\AgdaSpace{}%
\AgdaBound{fb'}\AgdaSymbol{)}\AgdaSpace{}%
\AgdaSymbol{(}\AgdaFunction{renameb}\AgdaSpace{}%
\AgdaInductiveConstructor{Sb}\AgdaSpace{}%
\AgdaSymbol{(}\AgdaFunction{embedTerm}\AgdaSpace{}%
\AgdaBound{t₁}\AgdaSymbol{)))}\AgdaSpace{}%
\AgdaOperator{\AgdaInductiveConstructor{·b}}\AgdaSpace{}%
\AgdaFunction{renameb}\AgdaSpace{}%
\AgdaBound{fb'}\AgdaSpace{}%
\AgdaSymbol{(}\AgdaFunction{embedTerm}\AgdaSpace{}%
\AgdaBound{t}\AgdaSymbol{)}\<%
\\
%
\>[2]\AgdaFunction{≡⟨}\AgdaSpace{}%
\AgdaFunction{cong₂}\AgdaSpace{}%
\AgdaOperator{\AgdaInductiveConstructor{\AgdaUnderscore{}·b\AgdaUnderscore{}}}\AgdaSpace{}%
\AgdaInductiveConstructor{refl}\AgdaSpace{}%
\AgdaSymbol{(}\AgdaFunction{pres-rename}\AgdaSpace{}%
\AgdaBound{f'}\AgdaSpace{}%
\AgdaBound{fb'}\AgdaSpace{}%
\AgdaBound{h}\AgdaSpace{}%
\AgdaBound{t}\AgdaSymbol{)}\AgdaSpace{}%
\AgdaFunction{⟩}\<%
\\
\>[2][@{}l@{\AgdaIndent{0}}]%
\>[4]\AgdaInductiveConstructor{ƛb}\AgdaSpace{}%
\AgdaSymbol{(}\AgdaInductiveConstructor{ε}\AgdaSpace{}%
\AgdaInductiveConstructor{Zb}\AgdaSpace{}%
\AgdaOperator{\AgdaInductiveConstructor{·b}}\AgdaSpace{}%
\AgdaFunction{renameb}\AgdaSpace{}%
\AgdaSymbol{(}\AgdaFunction{extb}\AgdaSpace{}%
\AgdaBound{fb'}\AgdaSymbol{)}\AgdaSpace{}%
\AgdaSymbol{(}\AgdaFunction{renameb}\AgdaSpace{}%
\AgdaInductiveConstructor{Sb}\AgdaSpace{}%
\AgdaSymbol{(}\AgdaFunction{embedTerm}\AgdaSpace{}%
\AgdaBound{t₁}\AgdaSymbol{)))}\AgdaSpace{}%
\AgdaOperator{\AgdaInductiveConstructor{·b}}\AgdaSpace{}%
\AgdaFunction{embedTerm}\AgdaSpace{}%
\AgdaSymbol{(}\AgdaFunction{rename}\AgdaSpace{}%
\AgdaBound{f'}\AgdaSpace{}%
\AgdaBound{t}\AgdaSymbol{)}\<%
\\
%
\>[2]\AgdaFunction{≡⟨}\AgdaSpace{}%
\AgdaFunction{cong₂}\AgdaSpace{}%
\AgdaOperator{\AgdaInductiveConstructor{\AgdaUnderscore{}·b\AgdaUnderscore{}}}\AgdaSpace{}%
\AgdaSymbol{(}\AgdaFunction{cong}\AgdaSpace{}%
\AgdaInductiveConstructor{ƛb}\AgdaSpace{}%
\AgdaSymbol{(}\AgdaFunction{cong₂}\AgdaSpace{}%
\AgdaOperator{\AgdaInductiveConstructor{\AgdaUnderscore{}·b\AgdaUnderscore{}}}\AgdaSpace{}%
\AgdaInductiveConstructor{refl}\AgdaSpace{}%
\AgdaSymbol{(}\AgdaFunction{rename-extb-Sb-commute}\AgdaSpace{}%
\AgdaBound{fb'}\AgdaSpace{}%
\AgdaSymbol{(}\AgdaFunction{embedTerm}\AgdaSpace{}%
\AgdaBound{t₁}\AgdaSymbol{))))}\AgdaSpace{}%
\AgdaInductiveConstructor{refl}\AgdaSpace{}%
\AgdaFunction{⟩}\<%
\\
\>[2][@{}l@{\AgdaIndent{0}}]%
\>[4]\AgdaInductiveConstructor{ƛb}\AgdaSpace{}%
\AgdaSymbol{(}\AgdaInductiveConstructor{ε}\AgdaSpace{}%
\AgdaInductiveConstructor{Zb}\AgdaSpace{}%
\AgdaOperator{\AgdaInductiveConstructor{·b}}\AgdaSpace{}%
\AgdaFunction{renameb}\AgdaSpace{}%
\AgdaInductiveConstructor{Sb}\AgdaSpace{}%
\AgdaSymbol{(}\AgdaFunction{renameb}\AgdaSpace{}%
\AgdaBound{fb'}\AgdaSpace{}%
\AgdaSymbol{(}\AgdaFunction{embedTerm}\AgdaSpace{}%
\AgdaBound{t₁}\AgdaSymbol{)))}\AgdaSpace{}%
\AgdaOperator{\AgdaInductiveConstructor{·b}}\AgdaSpace{}%
\AgdaFunction{embedTerm}\AgdaSpace{}%
\AgdaSymbol{(}\AgdaFunction{rename}\AgdaSpace{}%
\AgdaBound{f'}\AgdaSpace{}%
\AgdaBound{t}\AgdaSymbol{)}\<%
\\
%
\>[2]\AgdaFunction{≡⟨}\AgdaSpace{}%
\AgdaFunction{cong₂}\AgdaSpace{}%
\AgdaOperator{\AgdaInductiveConstructor{\AgdaUnderscore{}·b\AgdaUnderscore{}}}\AgdaSpace{}%
\AgdaSymbol{(}\AgdaFunction{cong}\AgdaSpace{}%
\AgdaInductiveConstructor{ƛb}\AgdaSpace{}%
\AgdaSymbol{(}\AgdaFunction{cong₂}\AgdaSpace{}%
\AgdaOperator{\AgdaInductiveConstructor{\AgdaUnderscore{}·b\AgdaUnderscore{}}}\AgdaSpace{}%
\AgdaInductiveConstructor{refl}\AgdaSpace{}%
\AgdaSymbol{(}\AgdaFunction{cong}\AgdaSpace{}%
\AgdaSymbol{(}\AgdaFunction{renameb}\AgdaSpace{}%
\AgdaInductiveConstructor{Sb}\AgdaSymbol{)}\AgdaSpace{}%
\AgdaSymbol{(}\AgdaFunction{pres-rename}\AgdaSpace{}%
\AgdaBound{f'}\AgdaSpace{}%
\AgdaBound{fb'}\AgdaSpace{}%
\AgdaBound{h}\AgdaSpace{}%
\AgdaBound{t₁}\AgdaSymbol{))))}\AgdaSpace{}%
\AgdaInductiveConstructor{refl}\AgdaSpace{}%
\AgdaFunction{⟩}\<%
\\
\>[2][@{}l@{\AgdaIndent{0}}]%
\>[4]\AgdaInductiveConstructor{ƛb}\AgdaSpace{}%
\AgdaSymbol{(}\AgdaInductiveConstructor{ε}\AgdaSpace{}%
\AgdaInductiveConstructor{Zb}\AgdaSpace{}%
\AgdaOperator{\AgdaInductiveConstructor{·b}}\AgdaSpace{}%
\AgdaFunction{renameb}\AgdaSpace{}%
\AgdaInductiveConstructor{Sb}\AgdaSpace{}%
\AgdaSymbol{(}\AgdaFunction{embedTerm}\AgdaSpace{}%
\AgdaSymbol{(}\AgdaFunction{rename}\AgdaSpace{}%
\AgdaBound{f'}\AgdaSpace{}%
\AgdaBound{t₁}\AgdaSymbol{)))}\AgdaSpace{}%
\AgdaOperator{\AgdaInductiveConstructor{·b}}\AgdaSpace{}%
\AgdaFunction{embedTerm}\AgdaSpace{}%
\AgdaSymbol{(}\AgdaFunction{rename}\AgdaSpace{}%
\AgdaBound{f'}\AgdaSpace{}%
\AgdaBound{t}\AgdaSymbol{)}\<%
\\
%
\>[2]\AgdaOperator{\AgdaFunction{∎eq}}\<%
\\
%
\\[\AgdaEmptyExtraSkip]%
\>[0]\AgdaComment{--\ Preservation\ and\ Reflection\ of\ }\<%
\\
\>[0]\AgdaComment{--\ exts\ of\ Gödel's\ translation}\<%
\\
\>[0]\AgdaFunction{pres-exts}\AgdaSpace{}%
\AgdaSymbol{:}\AgdaSpace{}%
\AgdaSymbol{∀}\AgdaSpace{}%
\AgdaSymbol{\{}\AgdaBound{Γ}\AgdaSpace{}%
\AgdaBound{Δ}\AgdaSymbol{\}}\<%
\\
\>[0][@{}l@{\AgdaIndent{0}}]%
\>[2]\AgdaSymbol{→}\AgdaSpace{}%
\AgdaSymbol{(}\AgdaBound{fb'}\AgdaSpace{}%
\AgdaSymbol{:}\AgdaSpace{}%
\AgdaSymbol{∀}\AgdaSpace{}%
\AgdaSymbol{\{}\AgdaBound{Ab}\AgdaSymbol{\}}\AgdaSpace{}%
\AgdaSymbol{→}\AgdaSpace{}%
\AgdaInductiveConstructor{□}\AgdaSpace{}%
\AgdaBound{Ab}\AgdaSpace{}%
\AgdaOperator{\AgdaDatatype{∈b}}\AgdaSpace{}%
\AgdaFunction{embedContext}\AgdaSpace{}%
\AgdaBound{Γ}\AgdaSpace{}%
\AgdaSymbol{→}\AgdaSpace{}%
\AgdaFunction{embedContext}\AgdaSpace{}%
\AgdaBound{Δ}\AgdaSpace{}%
\AgdaOperator{\AgdaDatatype{⊢b}}\AgdaSpace{}%
\AgdaBound{Ab}\AgdaSymbol{)}\<%
\\
%
\>[2]\AgdaSymbol{→}\AgdaSpace{}%
\AgdaSymbol{(}\AgdaBound{f'}\AgdaSpace{}%
\AgdaSymbol{:}\AgdaSpace{}%
\AgdaSymbol{∀}\AgdaSpace{}%
\AgdaSymbol{\{}\AgdaBound{A}\AgdaSymbol{\}}\AgdaSpace{}%
\AgdaSymbol{→}\AgdaSpace{}%
\AgdaBound{A}\AgdaSpace{}%
\AgdaOperator{\AgdaDatatype{∈}}\AgdaSpace{}%
\AgdaBound{Γ}\AgdaSpace{}%
\AgdaSymbol{→}\AgdaSpace{}%
\AgdaBound{Δ}\AgdaSpace{}%
\AgdaOperator{\AgdaDatatype{⊢}}\AgdaSpace{}%
\AgdaBound{A}\AgdaSymbol{)}\<%
\\
%
\>[2]\AgdaSymbol{→}\AgdaSpace{}%
\AgdaSymbol{(\{}\AgdaBound{A'}\AgdaSpace{}%
\AgdaSymbol{:}\AgdaSpace{}%
\AgdaDatatype{Type}\AgdaSymbol{\}}\AgdaSpace{}%
\AgdaSymbol{(}\AgdaBound{x}\AgdaSpace{}%
\AgdaSymbol{:}\AgdaSpace{}%
\AgdaBound{A'}\AgdaSpace{}%
\AgdaOperator{\AgdaDatatype{∈}}\AgdaSpace{}%
\AgdaBound{Γ}\AgdaSymbol{)}\AgdaSpace{}%
\AgdaSymbol{→}\AgdaSpace{}%
\AgdaInductiveConstructor{box}\AgdaSpace{}%
\AgdaSymbol{(}\AgdaBound{fb'}\AgdaSpace{}%
\AgdaSymbol{(}\AgdaFunction{embedLookup}\AgdaSpace{}%
\AgdaBound{x}\AgdaSymbol{))}\AgdaSpace{}%
\AgdaOperator{\AgdaDatatype{≡}}\AgdaSpace{}%
\AgdaFunction{embedTerm}\AgdaSpace{}%
\AgdaSymbol{(}\AgdaBound{f'}\AgdaSpace{}%
\AgdaBound{x}\AgdaSymbol{))}\<%
\\
%
\>[2]\AgdaSymbol{→}\AgdaSpace{}%
\AgdaSymbol{\{}\AgdaBound{A'}\AgdaSpace{}%
\AgdaBound{A}\AgdaSpace{}%
\AgdaSymbol{:}\AgdaSpace{}%
\AgdaDatatype{Type}\AgdaSymbol{\}}\AgdaSpace{}%
\AgdaSymbol{(}\AgdaBound{x}\AgdaSpace{}%
\AgdaSymbol{:}\AgdaSpace{}%
\AgdaBound{A'}\AgdaSpace{}%
\AgdaOperator{\AgdaDatatype{∈}}\AgdaSpace{}%
\AgdaBound{Γ}\AgdaSpace{}%
\AgdaOperator{\AgdaInductiveConstructor{,}}\AgdaSpace{}%
\AgdaBound{A}\AgdaSymbol{)}\AgdaSpace{}%
\AgdaSymbol{→}\AgdaSpace{}%
\AgdaInductiveConstructor{box}\AgdaSpace{}%
\AgdaSymbol{(}\AgdaFunction{extsb}\AgdaSpace{}%
\AgdaBound{fb'}\AgdaSpace{}%
\AgdaSymbol{(}\AgdaFunction{embedLookup}\AgdaSpace{}%
\AgdaBound{x}\AgdaSymbol{))}\AgdaSpace{}%
\AgdaOperator{\AgdaDatatype{≡}}\AgdaSpace{}%
\AgdaFunction{embedTerm}\AgdaSpace{}%
\AgdaSymbol{(}\AgdaFunction{exts}\AgdaSpace{}%
\AgdaBound{f'}\AgdaSpace{}%
\AgdaBound{x}\AgdaSymbol{)}\<%
\\
\>[0]\AgdaFunction{pres-exts}\AgdaSpace{}%
\AgdaBound{fb'}\AgdaSpace{}%
\AgdaBound{f'}\AgdaSpace{}%
\AgdaBound{h}\AgdaSpace{}%
\AgdaInductiveConstructor{Z}\AgdaSpace{}%
\AgdaSymbol{=}\AgdaSpace{}%
\AgdaInductiveConstructor{refl}\<%
\\
\>[0]\AgdaFunction{pres-exts}\AgdaSpace{}%
\AgdaBound{fb'}\AgdaSpace{}%
\AgdaBound{f'}\AgdaSpace{}%
\AgdaBound{h}\AgdaSpace{}%
\AgdaSymbol{(}\AgdaInductiveConstructor{S}\AgdaSpace{}%
\AgdaBound{x}\AgdaSymbol{)}\AgdaSpace{}%
\AgdaSymbol{=}\<%
\\
\>[0][@{}l@{\AgdaIndent{0}}]%
\>[2]\AgdaOperator{\AgdaFunction{beginEq}}\<%
\\
\>[2][@{}l@{\AgdaIndent{0}}]%
\>[4]\AgdaInductiveConstructor{box}\AgdaSpace{}%
\AgdaSymbol{(}\AgdaFunction{renameb}\AgdaSpace{}%
\AgdaInductiveConstructor{Sb}\AgdaSpace{}%
\AgdaSymbol{(}\AgdaBound{fb'}\AgdaSpace{}%
\AgdaSymbol{(}\AgdaFunction{embedLookup}\AgdaSpace{}%
\AgdaBound{x}\AgdaSymbol{)))}\<%
\\
%
\>[2]\AgdaFunction{≡⟨}\AgdaSpace{}%
\AgdaInductiveConstructor{refl}\AgdaSpace{}%
\AgdaFunction{⟩}\<%
\\
\>[2][@{}l@{\AgdaIndent{0}}]%
\>[4]\AgdaFunction{renameb}\AgdaSpace{}%
\AgdaInductiveConstructor{Sb}\AgdaSpace{}%
\AgdaSymbol{(}\AgdaInductiveConstructor{box}\AgdaSpace{}%
\AgdaSymbol{(}\AgdaBound{fb'}\AgdaSpace{}%
\AgdaSymbol{(}\AgdaFunction{embedLookup}\AgdaSpace{}%
\AgdaBound{x}\AgdaSymbol{)))}\<%
\\
%
\>[2]\AgdaFunction{≡⟨}\AgdaSpace{}%
\AgdaFunction{cong}\AgdaSpace{}%
\AgdaSymbol{(}\AgdaFunction{renameb}\AgdaSpace{}%
\AgdaInductiveConstructor{Sb}\AgdaSymbol{)}\AgdaSpace{}%
\AgdaSymbol{(}\AgdaBound{h}\AgdaSpace{}%
\AgdaBound{x}\AgdaSymbol{)}\AgdaSpace{}%
\AgdaFunction{⟩}\<%
\\
\>[2][@{}l@{\AgdaIndent{0}}]%
\>[4]\AgdaFunction{renameb}\AgdaSpace{}%
\AgdaInductiveConstructor{Sb}\AgdaSpace{}%
\AgdaSymbol{(}\AgdaFunction{embedTerm}\AgdaSpace{}%
\AgdaSymbol{(}\AgdaBound{f'}\AgdaSpace{}%
\AgdaBound{x}\AgdaSymbol{))}\<%
\\
%
\>[2]\AgdaFunction{≡⟨}\AgdaSpace{}%
\AgdaFunction{pres-rename}\AgdaSpace{}%
\AgdaInductiveConstructor{S}\AgdaSpace{}%
\AgdaInductiveConstructor{Sb}\AgdaSpace{}%
\AgdaSymbol{(λ}\AgdaSpace{}%
\AgdaBound{x₁}\AgdaSpace{}%
\AgdaSymbol{→}\AgdaSpace{}%
\AgdaInductiveConstructor{refl}\AgdaSymbol{)}\AgdaSpace{}%
\AgdaSymbol{(}\AgdaBound{f'}\AgdaSpace{}%
\AgdaBound{x}\AgdaSymbol{)}\AgdaSpace{}%
\AgdaFunction{⟩}\<%
\\
\>[2][@{}l@{\AgdaIndent{0}}]%
\>[4]\AgdaFunction{embedTerm}\AgdaSpace{}%
\AgdaSymbol{(}\AgdaFunction{rename}\AgdaSpace{}%
\AgdaInductiveConstructor{S}\AgdaSpace{}%
\AgdaSymbol{(}\AgdaBound{f'}\AgdaSpace{}%
\AgdaBound{x}\AgdaSymbol{))}\<%
\\
%
\>[2]\AgdaOperator{\AgdaFunction{∎eq}}\<%
\\
%
\\[\AgdaEmptyExtraSkip]%
\>[0]\AgdaComment{--\ Extsb\ preserves\ pointwise\ equality}\<%
\\
\>[0]\AgdaComment{--\ of\ renaming\ functions}\<%
\\
\>[0]\AgdaFunction{extsb-cong}\AgdaSpace{}%
\AgdaSymbol{:}\AgdaSpace{}%
\AgdaSymbol{∀}\AgdaSpace{}%
\AgdaSymbol{\{}\AgdaBound{Γ₁}\AgdaSpace{}%
\AgdaBound{Γ₂}\AgdaSpace{}%
\AgdaBound{A}\AgdaSpace{}%
\AgdaBound{A'}\AgdaSymbol{\}}\<%
\\
\>[0][@{}l@{\AgdaIndent{0}}]%
\>[2]\AgdaSymbol{→}\AgdaSpace{}%
\AgdaSymbol{\{}\AgdaBound{fb'}\AgdaSpace{}%
\AgdaSymbol{:}\AgdaSpace{}%
\AgdaSymbol{∀}\AgdaSpace{}%
\AgdaSymbol{\{}\AgdaBound{Ab}\AgdaSymbol{\}}\AgdaSpace{}%
\AgdaSymbol{→}\AgdaSpace{}%
\AgdaInductiveConstructor{□}\AgdaSpace{}%
\AgdaBound{Ab}\AgdaSpace{}%
\AgdaOperator{\AgdaDatatype{∈b}}\AgdaSpace{}%
\AgdaBound{Γ₁}\AgdaSpace{}%
\AgdaSymbol{→}\AgdaSpace{}%
\AgdaBound{Γ₂}\AgdaSpace{}%
\AgdaOperator{\AgdaDatatype{⊢b}}\AgdaSpace{}%
\AgdaBound{Ab}\AgdaSymbol{\}}\<%
\\
%
\>[2]\AgdaSymbol{→}\AgdaSpace{}%
\AgdaSymbol{\{}\AgdaBound{fb''}\AgdaSpace{}%
\AgdaSymbol{:}\AgdaSpace{}%
\AgdaSymbol{∀}\AgdaSpace{}%
\AgdaSymbol{\{}\AgdaBound{Ab}\AgdaSymbol{\}}\AgdaSpace{}%
\AgdaSymbol{→}\AgdaSpace{}%
\AgdaInductiveConstructor{□}\AgdaSpace{}%
\AgdaBound{Ab}\AgdaSpace{}%
\AgdaOperator{\AgdaDatatype{∈b}}\AgdaSpace{}%
\AgdaBound{Γ₁}\AgdaSpace{}%
\AgdaSymbol{→}\AgdaSpace{}%
\AgdaBound{Γ₂}\AgdaSpace{}%
\AgdaOperator{\AgdaDatatype{⊢b}}\AgdaSpace{}%
\AgdaBound{Ab}\AgdaSymbol{\}}\<%
\\
%
\>[2]\AgdaSymbol{→}\AgdaSpace{}%
\AgdaSymbol{(}\AgdaBound{h}\AgdaSpace{}%
\AgdaSymbol{:}\AgdaSpace{}%
\AgdaSymbol{∀}\AgdaSpace{}%
\AgdaSymbol{\{}\AgdaBound{Ab}\AgdaSymbol{\}}\AgdaSpace{}%
\AgdaSymbol{→}\AgdaSpace{}%
\AgdaSymbol{(}\AgdaBound{x}\AgdaSpace{}%
\AgdaSymbol{:}\AgdaSpace{}%
\AgdaInductiveConstructor{□}\AgdaSpace{}%
\AgdaBound{Ab}\AgdaSpace{}%
\AgdaOperator{\AgdaDatatype{∈b}}\AgdaSpace{}%
\AgdaBound{Γ₁}\AgdaSymbol{)}\AgdaSpace{}%
\AgdaSymbol{→}\AgdaSpace{}%
\AgdaBound{fb'}\AgdaSpace{}%
\AgdaBound{x}\AgdaSpace{}%
\AgdaOperator{\AgdaDatatype{≡}}\AgdaSpace{}%
\AgdaBound{fb''}\AgdaSpace{}%
\AgdaBound{x}\AgdaSymbol{)}\<%
\\
%
\>[2]\AgdaSymbol{→}\AgdaSpace{}%
\AgdaSymbol{(}\AgdaBound{x}\AgdaSpace{}%
\AgdaSymbol{:}\AgdaSpace{}%
\AgdaInductiveConstructor{□}\AgdaSpace{}%
\AgdaBound{A}\AgdaSpace{}%
\AgdaOperator{\AgdaDatatype{∈b}}\AgdaSpace{}%
\AgdaBound{Γ₁}\AgdaSpace{}%
\AgdaOperator{\AgdaInductiveConstructor{,b}}\AgdaSpace{}%
\AgdaInductiveConstructor{□}\AgdaSpace{}%
\AgdaBound{A'}\AgdaSymbol{)}\<%
\\
%
\>[2]\AgdaSymbol{→}\AgdaSpace{}%
\AgdaFunction{extsb}\AgdaSpace{}%
\AgdaBound{fb'}\AgdaSpace{}%
\AgdaBound{x}\AgdaSpace{}%
\AgdaOperator{\AgdaDatatype{≡}}\AgdaSpace{}%
\AgdaFunction{extsb}\AgdaSpace{}%
\AgdaBound{fb''}\AgdaSpace{}%
\AgdaBound{x}\<%
\\
\>[0]\AgdaFunction{extsb-cong}\AgdaSpace{}%
\AgdaBound{h}\AgdaSpace{}%
\AgdaInductiveConstructor{Zb}\AgdaSpace{}%
\AgdaSymbol{=}\AgdaSpace{}%
\AgdaInductiveConstructor{refl}\<%
\\
\>[0]\AgdaFunction{extsb-cong}\AgdaSpace{}%
\AgdaBound{h}\AgdaSpace{}%
\AgdaSymbol{(}\AgdaInductiveConstructor{Sb}\AgdaSpace{}%
\AgdaBound{x}\AgdaSymbol{)}\AgdaSpace{}%
\AgdaSymbol{=}\AgdaSpace{}%
\AgdaFunction{cong}\AgdaSpace{}%
\AgdaSymbol{(}\AgdaFunction{renameb}\AgdaSpace{}%
\AgdaInductiveConstructor{Sb}\AgdaSymbol{)}\AgdaSpace{}%
\AgdaSymbol{(}\AgdaBound{h}\AgdaSpace{}%
\AgdaBound{x}\AgdaSymbol{)}\<%
\\
%
\\[\AgdaEmptyExtraSkip]%
\>[0]\AgdaComment{--\ substb\ preserves\ pointwise\ equality}\<%
\\
\>[0]\AgdaComment{--\ of\ renaming\ functions}\<%
\\
\>[0]\AgdaFunction{substb-cong}\AgdaSpace{}%
\AgdaSymbol{:}\AgdaSpace{}%
\AgdaSymbol{∀}\AgdaSpace{}%
\AgdaSymbol{\{}\AgdaBound{Γ₁}\AgdaSpace{}%
\AgdaBound{Γ₂}\AgdaSpace{}%
\AgdaBound{A}\AgdaSymbol{\}}\<%
\\
\>[0][@{}l@{\AgdaIndent{0}}]%
\>[2]\AgdaSymbol{→}\AgdaSpace{}%
\AgdaSymbol{\{}\AgdaBound{fb'}\AgdaSpace{}%
\AgdaSymbol{:}\AgdaSpace{}%
\AgdaSymbol{∀}\AgdaSpace{}%
\AgdaSymbol{\{}\AgdaBound{Ab}\AgdaSymbol{\}}\AgdaSpace{}%
\AgdaSymbol{→}\AgdaSpace{}%
\AgdaInductiveConstructor{□}\AgdaSpace{}%
\AgdaBound{Ab}\AgdaSpace{}%
\AgdaOperator{\AgdaDatatype{∈b}}\AgdaSpace{}%
\AgdaBound{Γ₁}\AgdaSpace{}%
\AgdaSymbol{→}\AgdaSpace{}%
\AgdaBound{Γ₂}\AgdaSpace{}%
\AgdaOperator{\AgdaDatatype{⊢b}}\AgdaSpace{}%
\AgdaBound{Ab}\AgdaSymbol{\}}\<%
\\
%
\>[2]\AgdaSymbol{→}\AgdaSpace{}%
\AgdaSymbol{\{}\AgdaBound{fb''}\AgdaSpace{}%
\AgdaSymbol{:}\AgdaSpace{}%
\AgdaSymbol{∀}\AgdaSpace{}%
\AgdaSymbol{\{}\AgdaBound{Ab}\AgdaSymbol{\}}\AgdaSpace{}%
\AgdaSymbol{→}\AgdaSpace{}%
\AgdaInductiveConstructor{□}\AgdaSpace{}%
\AgdaBound{Ab}\AgdaSpace{}%
\AgdaOperator{\AgdaDatatype{∈b}}\AgdaSpace{}%
\AgdaBound{Γ₁}\AgdaSpace{}%
\AgdaSymbol{→}\AgdaSpace{}%
\AgdaBound{Γ₂}\AgdaSpace{}%
\AgdaOperator{\AgdaDatatype{⊢b}}\AgdaSpace{}%
\AgdaBound{Ab}\AgdaSymbol{\}}\<%
\\
%
\>[2]\AgdaSymbol{→}\AgdaSpace{}%
\AgdaSymbol{(}\AgdaBound{h}\AgdaSpace{}%
\AgdaSymbol{:}\AgdaSpace{}%
\AgdaSymbol{∀}\AgdaSpace{}%
\AgdaSymbol{\{}\AgdaBound{Ab}\AgdaSymbol{\}}\AgdaSpace{}%
\AgdaSymbol{→}\AgdaSpace{}%
\AgdaSymbol{(}\AgdaBound{x}\AgdaSpace{}%
\AgdaSymbol{:}\AgdaSpace{}%
\AgdaInductiveConstructor{□}\AgdaSpace{}%
\AgdaBound{Ab}\AgdaSpace{}%
\AgdaOperator{\AgdaDatatype{∈b}}\AgdaSpace{}%
\AgdaBound{Γ₁}\AgdaSymbol{)}\AgdaSpace{}%
\AgdaSymbol{→}\AgdaSpace{}%
\AgdaBound{fb'}\AgdaSpace{}%
\AgdaBound{x}\AgdaSpace{}%
\AgdaOperator{\AgdaDatatype{≡}}\AgdaSpace{}%
\AgdaBound{fb''}\AgdaSpace{}%
\AgdaBound{x}\AgdaSymbol{)}\<%
\\
%
\>[2]\AgdaSymbol{→}\AgdaSpace{}%
\AgdaSymbol{(}\AgdaBound{t}\AgdaSpace{}%
\AgdaSymbol{:}\AgdaSpace{}%
\AgdaBound{Γ₁}\AgdaSpace{}%
\AgdaOperator{\AgdaDatatype{⊢b}}\AgdaSpace{}%
\AgdaBound{A}\AgdaSymbol{)}\<%
\\
%
\>[2]\AgdaSymbol{→}\AgdaSpace{}%
\AgdaFunction{substb}\AgdaSpace{}%
\AgdaBound{fb'}\AgdaSpace{}%
\AgdaBound{t}\AgdaSpace{}%
\AgdaOperator{\AgdaDatatype{≡}}\AgdaSpace{}%
\AgdaFunction{substb}\AgdaSpace{}%
\AgdaBound{fb''}\AgdaSpace{}%
\AgdaBound{t}\<%
\\
\>[0]\AgdaFunction{substb-cong}\AgdaSpace{}%
\AgdaBound{h}\AgdaSpace{}%
\AgdaSymbol{(}\AgdaInductiveConstructor{ε}\AgdaSpace{}%
\AgdaBound{x}\AgdaSymbol{)}\AgdaSpace{}%
\AgdaSymbol{=}\AgdaSpace{}%
\AgdaBound{h}\AgdaSpace{}%
\AgdaBound{x}\<%
\\
\>[0]\AgdaFunction{substb-cong}\AgdaSpace{}%
\AgdaBound{h}\AgdaSpace{}%
\AgdaSymbol{(}\AgdaInductiveConstructor{box}\AgdaSpace{}%
\AgdaBound{t}\AgdaSymbol{)}\AgdaSpace{}%
\AgdaSymbol{=}\AgdaSpace{}%
\AgdaFunction{cong}\AgdaSpace{}%
\AgdaInductiveConstructor{box}\AgdaSpace{}%
\AgdaSymbol{(}\AgdaFunction{substb-cong}\AgdaSpace{}%
\AgdaBound{h}\AgdaSpace{}%
\AgdaBound{t}\AgdaSymbol{)}\<%
\\
\>[0]\AgdaFunction{substb-cong}\AgdaSpace{}%
\AgdaBound{h}\AgdaSpace{}%
\AgdaSymbol{(}\AgdaInductiveConstructor{ƛb}\AgdaSpace{}%
\AgdaBound{t}\AgdaSymbol{)}\AgdaSpace{}%
\AgdaSymbol{=}\AgdaSpace{}%
\AgdaFunction{cong}\AgdaSpace{}%
\AgdaInductiveConstructor{ƛb}\AgdaSpace{}%
\AgdaSymbol{(}\AgdaFunction{substb-cong}\AgdaSpace{}%
\AgdaSymbol{(}\AgdaFunction{extsb-cong}\AgdaSpace{}%
\AgdaBound{h}\AgdaSymbol{)}\AgdaSpace{}%
\AgdaBound{t}\AgdaSymbol{)}\<%
\\
\>[0]\AgdaFunction{substb-cong}\AgdaSpace{}%
\AgdaBound{h}\AgdaSpace{}%
\AgdaSymbol{(}\AgdaBound{t}\AgdaSpace{}%
\AgdaOperator{\AgdaInductiveConstructor{·b}}\AgdaSpace{}%
\AgdaBound{t₁}\AgdaSymbol{)}\AgdaSpace{}%
\AgdaSymbol{=}\AgdaSpace{}%
\AgdaFunction{cong₂}\AgdaSpace{}%
\AgdaOperator{\AgdaInductiveConstructor{\AgdaUnderscore{}·b\AgdaUnderscore{}}}\AgdaSpace{}%
\AgdaSymbol{(}\AgdaFunction{substb-cong}\AgdaSpace{}%
\AgdaBound{h}\AgdaSpace{}%
\AgdaBound{t}\AgdaSymbol{)}\AgdaSpace{}%
\AgdaSymbol{(}\AgdaFunction{substb-cong}\AgdaSpace{}%
\AgdaBound{h}\AgdaSpace{}%
\AgdaBound{t₁}\AgdaSymbol{)}\<%
\\
%
\\[\AgdaEmptyExtraSkip]%
\>[0]\AgdaComment{--\ Commutativity\ of\ renameb\ and\ extsb:}\<%
\\
\>[0]\AgdaFunction{extsb-rename-commute}\AgdaSpace{}%
\AgdaSymbol{:}\AgdaSpace{}%
\AgdaSymbol{∀}\AgdaSpace{}%
\AgdaSymbol{\{}\AgdaBound{Γ₁}\AgdaSpace{}%
\AgdaBound{Γ₂}\AgdaSpace{}%
\AgdaBound{Γ₃}\AgdaSpace{}%
\AgdaBound{A}\AgdaSpace{}%
\AgdaBound{B}\AgdaSymbol{\}}\<%
\\
\>[0][@{}l@{\AgdaIndent{0}}]%
\>[2]\AgdaSymbol{→}\AgdaSpace{}%
\AgdaSymbol{(}\AgdaBound{fb'}\AgdaSpace{}%
\AgdaSymbol{:}\AgdaSpace{}%
\AgdaSymbol{∀}\AgdaSpace{}%
\AgdaSymbol{\{}\AgdaBound{Ab}\AgdaSymbol{\}}\AgdaSpace{}%
\AgdaSymbol{→}\AgdaSpace{}%
\AgdaBound{Ab}\AgdaSpace{}%
\AgdaOperator{\AgdaDatatype{∈b}}\AgdaSpace{}%
\AgdaBound{Γ₂}\AgdaSpace{}%
\AgdaSymbol{→}\AgdaSpace{}%
\AgdaBound{Ab}\AgdaSpace{}%
\AgdaOperator{\AgdaDatatype{∈b}}\AgdaSpace{}%
\AgdaBound{Γ₃}\AgdaSymbol{)}\<%
\\
%
\>[2]\AgdaSymbol{→}\AgdaSpace{}%
\AgdaSymbol{(}\AgdaBound{fb''}\AgdaSpace{}%
\AgdaSymbol{:}\AgdaSpace{}%
\AgdaSymbol{∀}\AgdaSpace{}%
\AgdaSymbol{\{}\AgdaBound{Ab}\AgdaSymbol{\}}\AgdaSpace{}%
\AgdaSymbol{→}\AgdaSpace{}%
\AgdaInductiveConstructor{□}\AgdaSpace{}%
\AgdaBound{Ab}\AgdaSpace{}%
\AgdaOperator{\AgdaDatatype{∈b}}\AgdaSpace{}%
\AgdaBound{Γ₁}\AgdaSpace{}%
\AgdaSymbol{→}\AgdaSpace{}%
\AgdaBound{Γ₂}\AgdaSpace{}%
\AgdaOperator{\AgdaDatatype{⊢b}}\AgdaSpace{}%
\AgdaBound{Ab}\AgdaSymbol{)}\<%
\\
%
\>[2]\AgdaSymbol{→}\AgdaSpace{}%
\AgdaSymbol{(}\AgdaBound{x}\AgdaSpace{}%
\AgdaSymbol{:}\AgdaSpace{}%
\AgdaInductiveConstructor{□}\AgdaSpace{}%
\AgdaBound{A}\AgdaSpace{}%
\AgdaOperator{\AgdaDatatype{∈b}}\AgdaSpace{}%
\AgdaBound{Γ₁}\AgdaSpace{}%
\AgdaOperator{\AgdaInductiveConstructor{,b}}\AgdaSpace{}%
\AgdaBound{B}\AgdaSymbol{)}\<%
\\
%
\>[2]\AgdaSymbol{→}%
\>[1170I]\AgdaFunction{extsb}\AgdaSpace{}%
\AgdaSymbol{(λ}\AgdaSpace{}%
\AgdaBound{x₁}\AgdaSpace{}%
\AgdaSymbol{→}\AgdaSpace{}%
\AgdaFunction{renameb}\AgdaSpace{}%
\AgdaBound{fb'}\AgdaSpace{}%
\AgdaSymbol{(}\AgdaBound{fb''}\AgdaSpace{}%
\AgdaBound{x₁}\AgdaSymbol{))}\AgdaSpace{}%
\AgdaBound{x}\AgdaSpace{}%
\AgdaOperator{\AgdaDatatype{≡}}\<%
\\
\>[.][@{}l@{}]\<[1170I]%
\>[4]\AgdaFunction{renameb}\AgdaSpace{}%
\AgdaSymbol{(}\AgdaFunction{extb}\AgdaSpace{}%
\AgdaBound{fb'}\AgdaSymbol{)}\AgdaSpace{}%
\AgdaSymbol{(}\AgdaFunction{extsb}\AgdaSpace{}%
\AgdaBound{fb''}\AgdaSpace{}%
\AgdaBound{x}\AgdaSymbol{)}\<%
\\
\>[0]\AgdaFunction{extsb-rename-commute}\AgdaSpace{}%
\AgdaBound{fb'}\AgdaSpace{}%
\AgdaBound{fb''}\AgdaSpace{}%
\AgdaInductiveConstructor{Zb}\AgdaSpace{}%
\AgdaSymbol{=}\AgdaSpace{}%
\AgdaInductiveConstructor{refl}\<%
\\
\>[0]\AgdaFunction{extsb-rename-commute}\AgdaSpace{}%
\AgdaBound{fb'}\AgdaSpace{}%
\AgdaBound{fb''}\AgdaSpace{}%
\AgdaSymbol{(}\AgdaInductiveConstructor{Sb}\AgdaSpace{}%
\AgdaBound{x}\AgdaSymbol{)}\AgdaSpace{}%
\AgdaSymbol{=}\AgdaSpace{}%
\AgdaFunction{sym}\AgdaSpace{}%
\AgdaSymbol{(}\AgdaFunction{rename-extb-Sb-commute}\AgdaSpace{}%
\AgdaBound{fb'}\AgdaSpace{}%
\AgdaSymbol{(}\AgdaBound{fb''}\AgdaSpace{}%
\AgdaBound{x}\AgdaSymbol{))}\<%
\\
%
\\[\AgdaEmptyExtraSkip]%
\>[0]\AgdaComment{--\ Commutativity\ of\ substb\ and\ renameb}\<%
\\
\>[0]\AgdaFunction{substb-rename-commute}\AgdaSpace{}%
\AgdaSymbol{:}\AgdaSpace{}%
\AgdaSymbol{∀}\AgdaSpace{}%
\AgdaSymbol{\{}\AgdaBound{Γ₁}\AgdaSpace{}%
\AgdaBound{Γ₂}\AgdaSpace{}%
\AgdaBound{Γ₃}\AgdaSpace{}%
\AgdaBound{A}\AgdaSymbol{\}}\<%
\\
\>[0][@{}l@{\AgdaIndent{0}}]%
\>[2]\AgdaSymbol{→}\AgdaSpace{}%
\AgdaSymbol{(}\AgdaBound{fb'}\AgdaSpace{}%
\AgdaSymbol{:}\AgdaSpace{}%
\AgdaSymbol{∀}\AgdaSpace{}%
\AgdaSymbol{\{}\AgdaBound{Ab}\AgdaSymbol{\}}\AgdaSpace{}%
\AgdaSymbol{→}\AgdaSpace{}%
\AgdaBound{Ab}\AgdaSpace{}%
\AgdaOperator{\AgdaDatatype{∈b}}\AgdaSpace{}%
\AgdaBound{Γ₂}\AgdaSpace{}%
\AgdaSymbol{→}\AgdaSpace{}%
\AgdaBound{Ab}\AgdaSpace{}%
\AgdaOperator{\AgdaDatatype{∈b}}\AgdaSpace{}%
\AgdaBound{Γ₃}\AgdaSymbol{)}\<%
\\
%
\>[2]\AgdaSymbol{→}\AgdaSpace{}%
\AgdaSymbol{(}\AgdaBound{fb''}\AgdaSpace{}%
\AgdaSymbol{:}\AgdaSpace{}%
\AgdaSymbol{∀}\AgdaSpace{}%
\AgdaSymbol{\{}\AgdaBound{Ab}\AgdaSymbol{\}}\AgdaSpace{}%
\AgdaSymbol{→}\AgdaSpace{}%
\AgdaInductiveConstructor{□}\AgdaSpace{}%
\AgdaBound{Ab}\AgdaSpace{}%
\AgdaOperator{\AgdaDatatype{∈b}}\AgdaSpace{}%
\AgdaBound{Γ₁}\AgdaSpace{}%
\AgdaSymbol{→}\AgdaSpace{}%
\AgdaBound{Γ₂}\AgdaSpace{}%
\AgdaOperator{\AgdaDatatype{⊢b}}\AgdaSpace{}%
\AgdaBound{Ab}\AgdaSymbol{)}\<%
\\
%
\>[2]\AgdaSymbol{→}\AgdaSpace{}%
\AgdaSymbol{(}\AgdaBound{t}\AgdaSpace{}%
\AgdaSymbol{:}\AgdaSpace{}%
\AgdaBound{Γ₁}\AgdaSpace{}%
\AgdaOperator{\AgdaDatatype{⊢b}}\AgdaSpace{}%
\AgdaBound{A}\AgdaSymbol{)}\<%
\\
%
\>[2]\AgdaSymbol{→}\AgdaSpace{}%
\AgdaFunction{substb}\AgdaSpace{}%
\AgdaSymbol{(λ}\AgdaSpace{}%
\AgdaBound{x}\AgdaSpace{}%
\AgdaSymbol{→}\AgdaSpace{}%
\AgdaFunction{renameb}\AgdaSpace{}%
\AgdaBound{fb'}\AgdaSpace{}%
\AgdaSymbol{(}\AgdaBound{fb''}\AgdaSpace{}%
\AgdaBound{x}\AgdaSymbol{))}\AgdaSpace{}%
\AgdaBound{t}\AgdaSpace{}%
\AgdaOperator{\AgdaDatatype{≡}}\AgdaSpace{}%
\AgdaFunction{renameb}\AgdaSpace{}%
\AgdaBound{fb'}\AgdaSpace{}%
\AgdaSymbol{(}\AgdaFunction{substb}\AgdaSpace{}%
\AgdaBound{fb''}\AgdaSpace{}%
\AgdaBound{t}\AgdaSymbol{)}\<%
\\
\>[0]\AgdaFunction{substb-rename-commute}\AgdaSpace{}%
\AgdaBound{fb'}\AgdaSpace{}%
\AgdaBound{fb''}\AgdaSpace{}%
\AgdaSymbol{(}\AgdaInductiveConstructor{ε}\AgdaSpace{}%
\AgdaBound{x}\AgdaSymbol{)}\AgdaSpace{}%
\AgdaSymbol{=}\AgdaSpace{}%
\AgdaInductiveConstructor{refl}\<%
\\
\>[0]\AgdaFunction{substb-rename-commute}\AgdaSpace{}%
\AgdaBound{fb'}\AgdaSpace{}%
\AgdaBound{fb''}\AgdaSpace{}%
\AgdaSymbol{(}\AgdaInductiveConstructor{box}\AgdaSpace{}%
\AgdaBound{t}\AgdaSymbol{)}\AgdaSpace{}%
\AgdaSymbol{=}\<%
\\
\>[0][@{}l@{\AgdaIndent{0}}]%
\>[2]\AgdaFunction{cong}\<%
\\
\>[2][@{}l@{\AgdaIndent{0}}]%
\>[4]\AgdaInductiveConstructor{box}\<%
\\
%
\>[4]\AgdaSymbol{(}\AgdaFunction{substb-rename-commute}\AgdaSpace{}%
\AgdaBound{fb'}\AgdaSpace{}%
\AgdaBound{fb''}\AgdaSpace{}%
\AgdaBound{t}\AgdaSymbol{)}\<%
\\
\>[0]\AgdaFunction{substb-rename-commute}\AgdaSpace{}%
\AgdaBound{fb'}\AgdaSpace{}%
\AgdaBound{fb''}\AgdaSpace{}%
\AgdaSymbol{(}\AgdaInductiveConstructor{ƛb}\AgdaSpace{}%
\AgdaBound{t}\AgdaSymbol{)}\AgdaSpace{}%
\AgdaSymbol{=}\<%
\\
\>[0][@{}l@{\AgdaIndent{0}}]%
\>[2]\AgdaOperator{\AgdaFunction{beginEq}}\<%
\\
\>[2][@{}l@{\AgdaIndent{0}}]%
\>[4]\AgdaInductiveConstructor{ƛb}\AgdaSpace{}%
\AgdaSymbol{(}\AgdaFunction{substb}\AgdaSpace{}%
\AgdaSymbol{(}\AgdaFunction{extsb}\AgdaSpace{}%
\AgdaSymbol{(λ}\AgdaSpace{}%
\AgdaBound{x}\AgdaSpace{}%
\AgdaSymbol{→}\AgdaSpace{}%
\AgdaFunction{renameb}\AgdaSpace{}%
\AgdaBound{fb'}\AgdaSpace{}%
\AgdaSymbol{(}\AgdaBound{fb''}\AgdaSpace{}%
\AgdaBound{x}\AgdaSymbol{)))}\AgdaSpace{}%
\AgdaBound{t}\AgdaSymbol{)}\<%
\\
%
\>[2]\AgdaFunction{≡⟨}\AgdaSpace{}%
\AgdaFunction{cong}\AgdaSpace{}%
\AgdaInductiveConstructor{ƛb}\AgdaSpace{}%
\AgdaSymbol{(}\AgdaFunction{substb-cong}\AgdaSpace{}%
\AgdaSymbol{(λ}\AgdaSpace{}%
\AgdaBound{x}\AgdaSpace{}%
\AgdaSymbol{→}\AgdaSpace{}%
\AgdaFunction{extsb-rename-commute}\AgdaSpace{}%
\AgdaBound{fb'}\AgdaSpace{}%
\AgdaBound{fb''}\AgdaSpace{}%
\AgdaBound{x}\AgdaSymbol{)}\AgdaSpace{}%
\AgdaBound{t}\AgdaSymbol{)}\AgdaSpace{}%
\AgdaFunction{⟩}\<%
\\
\>[2][@{}l@{\AgdaIndent{0}}]%
\>[4]\AgdaInductiveConstructor{ƛb}\AgdaSpace{}%
\AgdaSymbol{(}\AgdaFunction{substb}\AgdaSpace{}%
\AgdaSymbol{(λ}\AgdaSpace{}%
\AgdaBound{x}\AgdaSpace{}%
\AgdaSymbol{→}\AgdaSpace{}%
\AgdaFunction{renameb}\AgdaSpace{}%
\AgdaSymbol{(}\AgdaFunction{extb}\AgdaSpace{}%
\AgdaBound{fb'}\AgdaSymbol{)}\AgdaSpace{}%
\AgdaSymbol{((}\AgdaFunction{extsb}\AgdaSpace{}%
\AgdaBound{fb''}\AgdaSymbol{)}\AgdaSpace{}%
\AgdaBound{x}\AgdaSymbol{))}\AgdaSpace{}%
\AgdaBound{t}\AgdaSymbol{)}\<%
\\
%
\>[2]\AgdaFunction{≡⟨}\AgdaSpace{}%
\AgdaFunction{cong}\AgdaSpace{}%
\AgdaInductiveConstructor{ƛb}\AgdaSpace{}%
\AgdaSymbol{(}\AgdaFunction{substb-rename-commute}\AgdaSpace{}%
\AgdaSymbol{(}\AgdaFunction{extb}\AgdaSpace{}%
\AgdaBound{fb'}\AgdaSymbol{)}\AgdaSpace{}%
\AgdaSymbol{(}\AgdaFunction{extsb}\AgdaSpace{}%
\AgdaBound{fb''}\AgdaSymbol{)}\AgdaSpace{}%
\AgdaBound{t}\AgdaSymbol{)}\AgdaSpace{}%
\AgdaFunction{⟩}\<%
\\
\>[2][@{}l@{\AgdaIndent{0}}]%
\>[4]\AgdaInductiveConstructor{ƛb}\AgdaSpace{}%
\AgdaSymbol{(}\AgdaFunction{renameb}\AgdaSpace{}%
\AgdaSymbol{(}\AgdaFunction{extb}\AgdaSpace{}%
\AgdaBound{fb'}\AgdaSymbol{)}\AgdaSpace{}%
\AgdaSymbol{(}\AgdaFunction{substb}\AgdaSpace{}%
\AgdaSymbol{(}\AgdaFunction{extsb}\AgdaSpace{}%
\AgdaBound{fb''}\AgdaSymbol{)}\AgdaSpace{}%
\AgdaBound{t}\AgdaSymbol{))}\<%
\\
%
\>[2]\AgdaOperator{\AgdaFunction{∎eq}}\<%
\\
\>[0]\AgdaFunction{substb-rename-commute}\AgdaSpace{}%
\AgdaBound{fb'}\AgdaSpace{}%
\AgdaBound{fb''}\AgdaSpace{}%
\AgdaSymbol{(}\AgdaBound{t}\AgdaSpace{}%
\AgdaOperator{\AgdaInductiveConstructor{·b}}\AgdaSpace{}%
\AgdaBound{t₁}\AgdaSymbol{)}\AgdaSpace{}%
\AgdaSymbol{=}\<%
\\
\>[0][@{}l@{\AgdaIndent{0}}]%
\>[2]\AgdaFunction{cong₂}\<%
\\
\>[2][@{}l@{\AgdaIndent{0}}]%
\>[4]\AgdaOperator{\AgdaInductiveConstructor{\AgdaUnderscore{}·b\AgdaUnderscore{}}}\<%
\\
%
\>[4]\AgdaSymbol{(}\AgdaFunction{substb-rename-commute}\AgdaSpace{}%
\AgdaBound{fb'}\AgdaSpace{}%
\AgdaBound{fb''}\AgdaSpace{}%
\AgdaBound{t}\AgdaSymbol{)}\<%
\\
%
\>[4]\AgdaSymbol{(}\AgdaFunction{substb-rename-commute}\AgdaSpace{}%
\AgdaBound{fb'}\AgdaSpace{}%
\AgdaBound{fb''}\AgdaSpace{}%
\AgdaBound{t₁}\AgdaSymbol{)}\<%
\\
%
\\[\AgdaEmptyExtraSkip]%
\>[0]\AgdaComment{--\ extsb\ distributes\ over\ composition:}\<%
\\
\>[0]\AgdaComment{--\ Executing\ extsb\ first\ with\ fb',\ then\ with\ fb''}\<%
\\
\>[0]\AgdaComment{--\ is\ the\ same\ as\ doing\ extsb\ once\ with\ fb''\ ∘\ fb'.}\<%
\\
\>[0]\AgdaFunction{extsb-comp}\AgdaSpace{}%
\AgdaSymbol{:}\AgdaSpace{}%
\AgdaSymbol{∀}\AgdaSpace{}%
\AgdaSymbol{\{}\AgdaBound{Γ₁}\AgdaSpace{}%
\AgdaBound{Γ₂}\AgdaSpace{}%
\AgdaBound{Γ₃}\AgdaSpace{}%
\AgdaBound{A}\AgdaSpace{}%
\AgdaBound{B}\AgdaSymbol{\}}\<%
\\
\>[0][@{}l@{\AgdaIndent{0}}]%
\>[2]\AgdaSymbol{→}\AgdaSpace{}%
\AgdaSymbol{(}\AgdaBound{fb'}\AgdaSpace{}%
\AgdaSymbol{:}\AgdaSpace{}%
\AgdaSymbol{∀}\AgdaSpace{}%
\AgdaSymbol{\{}\AgdaBound{Ab}\AgdaSymbol{\}}\AgdaSpace{}%
\AgdaSymbol{→}\AgdaSpace{}%
\AgdaBound{Ab}\AgdaSpace{}%
\AgdaOperator{\AgdaDatatype{∈b}}\AgdaSpace{}%
\AgdaBound{Γ₁}\AgdaSpace{}%
\AgdaSymbol{→}\AgdaSpace{}%
\AgdaBound{Ab}\AgdaSpace{}%
\AgdaOperator{\AgdaDatatype{∈b}}\AgdaSpace{}%
\AgdaBound{Γ₂}\AgdaSymbol{)}\<%
\\
%
\>[2]\AgdaSymbol{→}\AgdaSpace{}%
\AgdaSymbol{(}\AgdaBound{fb''}\AgdaSpace{}%
\AgdaSymbol{:}\AgdaSpace{}%
\AgdaSymbol{∀}\AgdaSpace{}%
\AgdaSymbol{\{}\AgdaBound{Ab}\AgdaSymbol{\}}\AgdaSpace{}%
\AgdaSymbol{→}\AgdaSpace{}%
\AgdaInductiveConstructor{□}\AgdaSpace{}%
\AgdaBound{Ab}\AgdaSpace{}%
\AgdaOperator{\AgdaDatatype{∈b}}\AgdaSpace{}%
\AgdaBound{Γ₂}\AgdaSpace{}%
\AgdaSymbol{→}\AgdaSpace{}%
\AgdaBound{Γ₃}\AgdaSpace{}%
\AgdaOperator{\AgdaDatatype{⊢b}}\AgdaSpace{}%
\AgdaBound{Ab}\AgdaSymbol{)}\<%
\\
%
\>[2]\AgdaSymbol{→}\AgdaSpace{}%
\AgdaSymbol{(}\AgdaBound{x}\AgdaSpace{}%
\AgdaSymbol{:}\AgdaSpace{}%
\AgdaInductiveConstructor{□}\AgdaSpace{}%
\AgdaBound{A}\AgdaSpace{}%
\AgdaOperator{\AgdaDatatype{∈b}}\AgdaSpace{}%
\AgdaBound{Γ₁}\AgdaSpace{}%
\AgdaOperator{\AgdaInductiveConstructor{,b}}\AgdaSpace{}%
\AgdaBound{B}\AgdaSymbol{)}\<%
\\
%
\>[2]\AgdaSymbol{→}\AgdaSpace{}%
\AgdaFunction{extsb}\AgdaSpace{}%
\AgdaBound{fb''}\AgdaSpace{}%
\AgdaSymbol{(}\AgdaFunction{extb}\AgdaSpace{}%
\AgdaBound{fb'}\AgdaSpace{}%
\AgdaBound{x}\AgdaSymbol{)}\AgdaSpace{}%
\AgdaOperator{\AgdaDatatype{≡}}\AgdaSpace{}%
\AgdaFunction{extsb}\AgdaSpace{}%
\AgdaSymbol{(λ}\AgdaSpace{}%
\AgdaBound{x₁}\AgdaSpace{}%
\AgdaSymbol{→}\AgdaSpace{}%
\AgdaBound{fb''}\AgdaSpace{}%
\AgdaSymbol{(}\AgdaBound{fb'}\AgdaSpace{}%
\AgdaBound{x₁}\AgdaSymbol{))}\AgdaSpace{}%
\AgdaBound{x}\<%
\\
\>[0]\AgdaFunction{extsb-comp}\AgdaSpace{}%
\AgdaBound{fb'}\AgdaSpace{}%
\AgdaBound{fb''}\AgdaSpace{}%
\AgdaInductiveConstructor{Zb}\AgdaSpace{}%
\AgdaSymbol{=}\AgdaSpace{}%
\AgdaInductiveConstructor{refl}\<%
\\
\>[0]\AgdaFunction{extsb-comp}\AgdaSpace{}%
\AgdaBound{fb'}\AgdaSpace{}%
\AgdaBound{fb''}\AgdaSpace{}%
\AgdaSymbol{(}\AgdaInductiveConstructor{Sb}\AgdaSpace{}%
\AgdaBound{x}\AgdaSymbol{)}\AgdaSpace{}%
\AgdaSymbol{=}\AgdaSpace{}%
\AgdaInductiveConstructor{refl}\<%
\\
%
\\[\AgdaEmptyExtraSkip]%
\>[0]\AgdaComment{--\ substb\ and\ rename\ distribute\ over\ composition:}\<%
\\
\>[0]\AgdaComment{--\ Renaming\ first\ with\ fb'\ and\ substituting\ with\ fb''}\<%
\\
\>[0]\AgdaComment{--\ is\ the\ same\ as\ substituting\ with\ fb''\ ∘\ fb'}\<%
\\
\>[0]\AgdaFunction{subst-renameb-comp}\AgdaSpace{}%
\AgdaSymbol{:}\AgdaSpace{}%
\AgdaSymbol{∀}\AgdaSpace{}%
\AgdaSymbol{\{}\AgdaBound{Γ₁}\AgdaSpace{}%
\AgdaBound{Γ₂}\AgdaSpace{}%
\AgdaBound{Γ₃}\AgdaSpace{}%
\AgdaBound{A}\AgdaSymbol{\}}\<%
\\
\>[0][@{}l@{\AgdaIndent{0}}]%
\>[2]\AgdaSymbol{→}\AgdaSpace{}%
\AgdaSymbol{(}\AgdaBound{fb'}\AgdaSpace{}%
\AgdaSymbol{:}\AgdaSpace{}%
\AgdaSymbol{∀}\AgdaSpace{}%
\AgdaSymbol{\{}\AgdaBound{Ab}\AgdaSymbol{\}}\AgdaSpace{}%
\AgdaSymbol{→}\AgdaSpace{}%
\AgdaBound{Ab}\AgdaSpace{}%
\AgdaOperator{\AgdaDatatype{∈b}}\AgdaSpace{}%
\AgdaBound{Γ₁}\AgdaSpace{}%
\AgdaSymbol{→}\AgdaSpace{}%
\AgdaBound{Ab}\AgdaSpace{}%
\AgdaOperator{\AgdaDatatype{∈b}}\AgdaSpace{}%
\AgdaBound{Γ₂}\AgdaSymbol{)}\<%
\\
%
\>[2]\AgdaSymbol{→}\AgdaSpace{}%
\AgdaSymbol{(}\AgdaBound{fb''}\AgdaSpace{}%
\AgdaSymbol{:}\AgdaSpace{}%
\AgdaSymbol{∀}\AgdaSpace{}%
\AgdaSymbol{\{}\AgdaBound{Ab}\AgdaSymbol{\}}\AgdaSpace{}%
\AgdaSymbol{→}\AgdaSpace{}%
\AgdaInductiveConstructor{□}\AgdaSpace{}%
\AgdaBound{Ab}\AgdaSpace{}%
\AgdaOperator{\AgdaDatatype{∈b}}\AgdaSpace{}%
\AgdaBound{Γ₂}\AgdaSpace{}%
\AgdaSymbol{→}\AgdaSpace{}%
\AgdaBound{Γ₃}\AgdaSpace{}%
\AgdaOperator{\AgdaDatatype{⊢b}}\AgdaSpace{}%
\AgdaBound{Ab}\AgdaSymbol{)}\<%
\\
%
\>[2]\AgdaSymbol{→}\AgdaSpace{}%
\AgdaSymbol{(}\AgdaBound{t}\AgdaSpace{}%
\AgdaSymbol{:}\AgdaSpace{}%
\AgdaBound{Γ₁}\AgdaSpace{}%
\AgdaOperator{\AgdaDatatype{⊢b}}\AgdaSpace{}%
\AgdaBound{A}\AgdaSymbol{)}\<%
\\
%
\>[2]\AgdaSymbol{→}\AgdaSpace{}%
\AgdaFunction{substb}\AgdaSpace{}%
\AgdaBound{fb''}\AgdaSpace{}%
\AgdaSymbol{(}\AgdaFunction{renameb}\AgdaSpace{}%
\AgdaBound{fb'}\AgdaSpace{}%
\AgdaBound{t}\AgdaSymbol{)}\AgdaSpace{}%
\AgdaOperator{\AgdaDatatype{≡}}\AgdaSpace{}%
\AgdaFunction{substb}\AgdaSpace{}%
\AgdaSymbol{(λ}\AgdaSpace{}%
\AgdaBound{x}\AgdaSpace{}%
\AgdaSymbol{→}\AgdaSpace{}%
\AgdaBound{fb''}\AgdaSpace{}%
\AgdaSymbol{(}\AgdaBound{fb'}\AgdaSpace{}%
\AgdaBound{x}\AgdaSymbol{))}\AgdaSpace{}%
\AgdaBound{t}\<%
\\
\>[0]\AgdaFunction{subst-renameb-comp}\AgdaSpace{}%
\AgdaBound{fb'}\AgdaSpace{}%
\AgdaBound{fb''}\AgdaSpace{}%
\AgdaSymbol{(}\AgdaInductiveConstructor{ε}\AgdaSpace{}%
\AgdaBound{x}\AgdaSymbol{)}\AgdaSpace{}%
\AgdaSymbol{=}\AgdaSpace{}%
\AgdaInductiveConstructor{refl}\<%
\\
\>[0]\AgdaFunction{subst-renameb-comp}\AgdaSpace{}%
\AgdaBound{fb'}\AgdaSpace{}%
\AgdaBound{fb''}\AgdaSpace{}%
\AgdaSymbol{(}\AgdaInductiveConstructor{box}\AgdaSpace{}%
\AgdaBound{t}\AgdaSymbol{)}\AgdaSpace{}%
\AgdaSymbol{=}\AgdaSpace{}%
\AgdaFunction{cong}\AgdaSpace{}%
\AgdaInductiveConstructor{box}\AgdaSpace{}%
\AgdaSymbol{(}\AgdaFunction{subst-renameb-comp}\AgdaSpace{}%
\AgdaBound{fb'}\AgdaSpace{}%
\AgdaBound{fb''}\AgdaSpace{}%
\AgdaBound{t}\AgdaSymbol{)}\<%
\\
\>[0]\AgdaFunction{subst-renameb-comp}\AgdaSpace{}%
\AgdaBound{fb'}\AgdaSpace{}%
\AgdaBound{fb''}\AgdaSpace{}%
\AgdaSymbol{(}\AgdaInductiveConstructor{ƛb}\AgdaSpace{}%
\AgdaBound{t}\AgdaSymbol{)}\AgdaSpace{}%
\AgdaSymbol{=}\<%
\\
\>[0][@{}l@{\AgdaIndent{0}}]%
\>[2]\AgdaOperator{\AgdaFunction{beginEq}}\<%
\\
\>[2][@{}l@{\AgdaIndent{0}}]%
\>[4]\AgdaInductiveConstructor{ƛb}\AgdaSpace{}%
\AgdaSymbol{(}\AgdaFunction{substb}\AgdaSpace{}%
\AgdaSymbol{(}\AgdaFunction{extsb}\AgdaSpace{}%
\AgdaBound{fb''}\AgdaSymbol{)}\AgdaSpace{}%
\AgdaSymbol{(}\AgdaFunction{renameb}\AgdaSpace{}%
\AgdaSymbol{(}\AgdaFunction{extb}\AgdaSpace{}%
\AgdaBound{fb'}\AgdaSymbol{)}\AgdaSpace{}%
\AgdaBound{t}\AgdaSymbol{))}\<%
\\
%
\>[2]\AgdaFunction{≡⟨}\AgdaSpace{}%
\AgdaFunction{cong}\AgdaSpace{}%
\AgdaInductiveConstructor{ƛb}\AgdaSpace{}%
\AgdaSymbol{(}\AgdaFunction{subst-renameb-comp}\AgdaSpace{}%
\AgdaSymbol{(}\AgdaFunction{extb}\AgdaSpace{}%
\AgdaBound{fb'}\AgdaSymbol{)}\AgdaSpace{}%
\AgdaSymbol{(}\AgdaFunction{extsb}\AgdaSpace{}%
\AgdaBound{fb''}\AgdaSymbol{)}\AgdaSpace{}%
\AgdaBound{t}\AgdaSymbol{)}\AgdaSpace{}%
\AgdaFunction{⟩}\<%
\\
\>[2][@{}l@{\AgdaIndent{0}}]%
\>[4]\AgdaInductiveConstructor{ƛb}\AgdaSpace{}%
\AgdaSymbol{(}\AgdaFunction{substb}\AgdaSpace{}%
\AgdaSymbol{(λ}\AgdaSpace{}%
\AgdaBound{x}\AgdaSpace{}%
\AgdaSymbol{→}\AgdaSpace{}%
\AgdaSymbol{(}\AgdaFunction{extsb}\AgdaSpace{}%
\AgdaBound{fb''}\AgdaSymbol{)}\AgdaSpace{}%
\AgdaSymbol{((}\AgdaFunction{extb}\AgdaSpace{}%
\AgdaBound{fb'}\AgdaSymbol{)}\AgdaSpace{}%
\AgdaBound{x}\AgdaSymbol{))}\AgdaSpace{}%
\AgdaBound{t}\AgdaSymbol{)}\<%
\\
%
\>[2]\AgdaFunction{≡⟨}\AgdaSpace{}%
\AgdaFunction{cong}\AgdaSpace{}%
\AgdaInductiveConstructor{ƛb}\AgdaSpace{}%
\AgdaSymbol{(}\AgdaFunction{substb-cong}\AgdaSpace{}%
\AgdaSymbol{(λ}\AgdaSpace{}%
\AgdaBound{x}\AgdaSpace{}%
\AgdaSymbol{→}\AgdaSpace{}%
\AgdaFunction{extsb-comp}\AgdaSpace{}%
\AgdaBound{fb'}\AgdaSpace{}%
\AgdaBound{fb''}\AgdaSpace{}%
\AgdaBound{x}\AgdaSymbol{)}\AgdaSpace{}%
\AgdaBound{t}\AgdaSymbol{)}\AgdaSpace{}%
\AgdaFunction{⟩}\<%
\\
\>[2][@{}l@{\AgdaIndent{0}}]%
\>[4]\AgdaInductiveConstructor{ƛb}\AgdaSpace{}%
\AgdaSymbol{(}\AgdaFunction{substb}\AgdaSpace{}%
\AgdaSymbol{(}\AgdaFunction{extsb}\AgdaSpace{}%
\AgdaSymbol{(λ}\AgdaSpace{}%
\AgdaBound{x}\AgdaSpace{}%
\AgdaSymbol{→}\AgdaSpace{}%
\AgdaBound{fb''}\AgdaSpace{}%
\AgdaSymbol{(}\AgdaBound{fb'}\AgdaSpace{}%
\AgdaBound{x}\AgdaSymbol{)))}\AgdaSpace{}%
\AgdaBound{t}\AgdaSymbol{)}\<%
\\
%
\>[2]\AgdaOperator{\AgdaFunction{∎eq}}\<%
\\
\>[0]\AgdaFunction{subst-renameb-comp}\AgdaSpace{}%
\AgdaBound{fb'}\AgdaSpace{}%
\AgdaBound{fb''}\AgdaSpace{}%
\AgdaSymbol{(}\AgdaBound{t}\AgdaSpace{}%
\AgdaOperator{\AgdaInductiveConstructor{·b}}\AgdaSpace{}%
\AgdaBound{t₁}\AgdaSymbol{)}\AgdaSpace{}%
\AgdaSymbol{=}\AgdaSpace{}%
\AgdaFunction{cong₂}\AgdaSpace{}%
\AgdaOperator{\AgdaInductiveConstructor{\AgdaUnderscore{}·b\AgdaUnderscore{}}}\AgdaSpace{}%
\AgdaSymbol{(}\AgdaFunction{subst-renameb-comp}\AgdaSpace{}%
\AgdaBound{fb'}\AgdaSpace{}%
\AgdaBound{fb''}\AgdaSpace{}%
\AgdaBound{t}\AgdaSymbol{)}\AgdaSpace{}%
\AgdaSymbol{(}\AgdaFunction{subst-renameb-comp}\AgdaSpace{}%
\AgdaBound{fb'}\AgdaSpace{}%
\AgdaBound{fb''}\AgdaSpace{}%
\AgdaBound{t₁}\AgdaSymbol{)}\<%
\\
\>[0]\AgdaComment{--\ Renaming\ followed\ by\ substitution\ with\ an\ extended\ context\ (extsb)}\<%
\\
\>[0]\AgdaComment{--\ commutes\ with\ substitution\ followed\ by\ renaming.}\<%
\\
\>[0]\AgdaFunction{substb-exts-rename-Sb-commute}\AgdaSpace{}%
\AgdaSymbol{:}\AgdaSpace{}%
\AgdaSymbol{∀}\AgdaSpace{}%
\AgdaSymbol{\{}\AgdaBound{Γ}\AgdaSpace{}%
\AgdaBound{A}\AgdaSpace{}%
\AgdaBound{B}\AgdaSpace{}%
\AgdaBound{Δ}\AgdaSymbol{\}}\<%
\\
\>[0][@{}l@{\AgdaIndent{0}}]%
\>[2]\AgdaSymbol{→}\AgdaSpace{}%
\AgdaSymbol{(}\AgdaBound{fb'}\AgdaSpace{}%
\AgdaSymbol{:}\AgdaSpace{}%
\AgdaSymbol{∀}\AgdaSpace{}%
\AgdaSymbol{\{}\AgdaBound{Ab}\AgdaSymbol{\}}\AgdaSpace{}%
\AgdaSymbol{→}\AgdaSpace{}%
\AgdaInductiveConstructor{□}\AgdaSpace{}%
\AgdaBound{Ab}\AgdaSpace{}%
\AgdaOperator{\AgdaDatatype{∈b}}\AgdaSpace{}%
\AgdaBound{Γ}\AgdaSpace{}%
\AgdaSymbol{→}\AgdaSpace{}%
\AgdaBound{Δ}\AgdaSpace{}%
\AgdaOperator{\AgdaDatatype{⊢b}}\AgdaSpace{}%
\AgdaBound{Ab}\AgdaSymbol{)}\<%
\\
%
\>[2]\AgdaSymbol{→}\AgdaSpace{}%
\AgdaSymbol{(}\AgdaBound{t}\AgdaSpace{}%
\AgdaSymbol{:}\AgdaSpace{}%
\AgdaBound{Γ}\AgdaSpace{}%
\AgdaOperator{\AgdaDatatype{⊢b}}\AgdaSpace{}%
\AgdaBound{A}\AgdaSymbol{)}\<%
\\
%
\>[2]\AgdaSymbol{→}%
\>[1555I]\AgdaFunction{substb}\AgdaSpace{}%
\AgdaSymbol{(}\AgdaFunction{extsb}\AgdaSpace{}%
\AgdaBound{fb'}\AgdaSpace{}%
\AgdaSymbol{\{}\AgdaArgument{B}\AgdaSpace{}%
\AgdaSymbol{=}\AgdaSpace{}%
\AgdaBound{B}\AgdaSymbol{\})}\AgdaSpace{}%
\AgdaSymbol{(}\AgdaFunction{renameb}\AgdaSpace{}%
\AgdaInductiveConstructor{Sb}\AgdaSpace{}%
\AgdaBound{t}\AgdaSymbol{)}\<%
\\
\>[.][@{}l@{}]\<[1555I]%
\>[4]\AgdaOperator{\AgdaDatatype{≡}}\AgdaSpace{}%
\AgdaFunction{renameb}\AgdaSpace{}%
\AgdaInductiveConstructor{Sb}\AgdaSpace{}%
\AgdaSymbol{(}\AgdaFunction{substb}\AgdaSpace{}%
\AgdaBound{fb'}\AgdaSpace{}%
\AgdaBound{t}\AgdaSymbol{)}\<%
\\
\>[0]\AgdaFunction{substb-exts-rename-Sb-commute}\AgdaSpace{}%
\AgdaBound{fb'}\AgdaSpace{}%
\AgdaBound{t}\AgdaSpace{}%
\AgdaSymbol{=}\<%
\\
\>[0][@{}l@{\AgdaIndent{0}}]%
\>[2]\AgdaOperator{\AgdaFunction{beginEq}}\<%
\\
\>[2][@{}l@{\AgdaIndent{0}}]%
\>[4]\AgdaFunction{substb}\AgdaSpace{}%
\AgdaSymbol{(}\AgdaFunction{extsb}\AgdaSpace{}%
\AgdaBound{fb'}\AgdaSymbol{)}\AgdaSpace{}%
\AgdaSymbol{(}\AgdaFunction{renameb}\AgdaSpace{}%
\AgdaInductiveConstructor{Sb}\AgdaSpace{}%
\AgdaBound{t}\AgdaSymbol{)}\<%
\\
%
\>[2]\AgdaFunction{≡⟨}\AgdaSpace{}%
\AgdaFunction{subst-renameb-comp}\AgdaSpace{}%
\AgdaInductiveConstructor{Sb}\AgdaSpace{}%
\AgdaSymbol{(}\AgdaFunction{extsb}\AgdaSpace{}%
\AgdaBound{fb'}\AgdaSymbol{)}\AgdaSpace{}%
\AgdaBound{t}\AgdaSpace{}%
\AgdaFunction{⟩}\<%
\\
\>[2][@{}l@{\AgdaIndent{0}}]%
\>[4]\AgdaFunction{substb}\AgdaSpace{}%
\AgdaSymbol{(λ}\AgdaSpace{}%
\AgdaBound{x}\AgdaSpace{}%
\AgdaSymbol{→}\AgdaSpace{}%
\AgdaFunction{renameb}\AgdaSpace{}%
\AgdaInductiveConstructor{Sb}\AgdaSpace{}%
\AgdaSymbol{(}\AgdaBound{fb'}\AgdaSpace{}%
\AgdaBound{x}\AgdaSymbol{))}\AgdaSpace{}%
\AgdaBound{t}\<%
\\
%
\>[2]\AgdaFunction{≡⟨}\AgdaSpace{}%
\AgdaFunction{substb-rename-commute}\AgdaSpace{}%
\AgdaInductiveConstructor{Sb}\AgdaSpace{}%
\AgdaBound{fb'}\AgdaSpace{}%
\AgdaBound{t}\AgdaSpace{}%
\AgdaFunction{⟩}\<%
\\
\>[2][@{}l@{\AgdaIndent{0}}]%
\>[4]\AgdaFunction{renameb}\AgdaSpace{}%
\AgdaInductiveConstructor{Sb}\AgdaSpace{}%
\AgdaSymbol{(}\AgdaFunction{substb}\AgdaSpace{}%
\AgdaBound{fb'}\AgdaSpace{}%
\AgdaBound{t}\AgdaSymbol{)}\<%
\\
%
\>[2]\AgdaOperator{\AgdaFunction{∎eq}}\<%
\\
\>[0]\AgdaComment{--\ Preservation\ and\ Reflection\ of\ }\<%
\\
\>[0]\AgdaComment{--\ simultaneous\ substitution\ of\ Gödel's\ translation}\<%
\\
\>[0]\AgdaFunction{pres-simultaneous-subst}\AgdaSpace{}%
\AgdaSymbol{:}\AgdaSpace{}%
\AgdaSymbol{∀}\AgdaSpace{}%
\AgdaSymbol{\{}\AgdaBound{Γ}\AgdaSpace{}%
\AgdaBound{Δ}\AgdaSpace{}%
\AgdaBound{A}\AgdaSymbol{\}}\<%
\\
\>[0][@{}l@{\AgdaIndent{0}}]%
\>[2]\AgdaSymbol{→}\AgdaSpace{}%
\AgdaSymbol{(}\AgdaBound{fb'}\AgdaSpace{}%
\AgdaSymbol{:}\AgdaSpace{}%
\AgdaSymbol{∀}\AgdaSpace{}%
\AgdaSymbol{\{}\AgdaBound{Ab}\AgdaSymbol{\}}\AgdaSpace{}%
\AgdaSymbol{→}\AgdaSpace{}%
\AgdaInductiveConstructor{□}\AgdaSpace{}%
\AgdaBound{Ab}\AgdaSpace{}%
\AgdaOperator{\AgdaDatatype{∈b}}\AgdaSpace{}%
\AgdaFunction{embedContext}\AgdaSpace{}%
\AgdaBound{Γ}\AgdaSpace{}%
\AgdaSymbol{→}\AgdaSpace{}%
\AgdaFunction{embedContext}\AgdaSpace{}%
\AgdaBound{Δ}\AgdaSpace{}%
\AgdaOperator{\AgdaDatatype{⊢b}}\AgdaSpace{}%
\AgdaBound{Ab}\AgdaSymbol{)}\<%
\\
%
\>[2]\AgdaSymbol{→}\AgdaSpace{}%
\AgdaSymbol{(}\AgdaBound{f'}\AgdaSpace{}%
\AgdaSymbol{:}\AgdaSpace{}%
\AgdaSymbol{∀}\AgdaSpace{}%
\AgdaSymbol{\{}\AgdaBound{A}\AgdaSymbol{\}}\AgdaSpace{}%
\AgdaSymbol{→}\AgdaSpace{}%
\AgdaBound{A}\AgdaSpace{}%
\AgdaOperator{\AgdaDatatype{∈}}\AgdaSpace{}%
\AgdaBound{Γ}\AgdaSpace{}%
\AgdaSymbol{→}\AgdaSpace{}%
\AgdaBound{Δ}\AgdaSpace{}%
\AgdaOperator{\AgdaDatatype{⊢}}\AgdaSpace{}%
\AgdaBound{A}\AgdaSymbol{)}\<%
\\
%
\>[2]\AgdaSymbol{→}\AgdaSpace{}%
\AgdaSymbol{(}\AgdaBound{h}\AgdaSpace{}%
\AgdaSymbol{:}\AgdaSpace{}%
\AgdaSymbol{∀}\AgdaSpace{}%
\AgdaSymbol{\{}\AgdaBound{A'}\AgdaSymbol{\}}\AgdaSpace{}%
\AgdaSymbol{→}\AgdaSpace{}%
\AgdaSymbol{(}\AgdaBound{x}\AgdaSpace{}%
\AgdaSymbol{:}\AgdaSpace{}%
\AgdaBound{A'}\AgdaSpace{}%
\AgdaOperator{\AgdaDatatype{∈}}\AgdaSpace{}%
\AgdaBound{Γ}\AgdaSymbol{)}\AgdaSpace{}%
\AgdaSymbol{→}\AgdaSpace{}%
\AgdaInductiveConstructor{box}\AgdaSpace{}%
\AgdaSymbol{(}\AgdaBound{fb'}\AgdaSpace{}%
\AgdaSymbol{(}\AgdaFunction{embedLookup}\AgdaSpace{}%
\AgdaBound{x}\AgdaSymbol{))}\AgdaSpace{}%
\AgdaOperator{\AgdaDatatype{≡}}\AgdaSpace{}%
\AgdaFunction{embedTerm}\AgdaSpace{}%
\AgdaSymbol{(}\AgdaBound{f'}\AgdaSpace{}%
\AgdaBound{x}\AgdaSymbol{)}\AgdaSpace{}%
\AgdaSymbol{)}\<%
\\
%
\>[2]\AgdaSymbol{→}\AgdaSpace{}%
\AgdaSymbol{(}\AgdaBound{M}\AgdaSpace{}%
\AgdaSymbol{:}\AgdaSpace{}%
\AgdaBound{Γ}\AgdaSpace{}%
\AgdaOperator{\AgdaDatatype{⊢}}\AgdaSpace{}%
\AgdaBound{A}\AgdaSymbol{)}\<%
\\
%
\>[2]\AgdaSymbol{→}\AgdaSpace{}%
\AgdaFunction{substb}\AgdaSpace{}%
\AgdaBound{fb'}\AgdaSpace{}%
\AgdaSymbol{(}\AgdaFunction{embedTerm}\AgdaSpace{}%
\AgdaBound{M}\AgdaSymbol{)}\AgdaSpace{}%
\AgdaOperator{\AgdaDatatype{≡}}\AgdaSpace{}%
\AgdaFunction{embedTerm}\AgdaSpace{}%
\AgdaSymbol{(}\AgdaFunction{subst}\AgdaSpace{}%
\AgdaBound{f'}\AgdaSpace{}%
\AgdaBound{M}\AgdaSymbol{)}\<%
\\
\>[0]\AgdaFunction{pres-simultaneous-subst}\AgdaSpace{}%
\AgdaBound{fb'}\AgdaSpace{}%
\AgdaBound{f'}\AgdaSpace{}%
\AgdaBound{h}\AgdaSpace{}%
\AgdaSymbol{(}\AgdaInductiveConstructor{`}\AgdaSpace{}%
\AgdaBound{x}\AgdaSymbol{)}\AgdaSpace{}%
\AgdaSymbol{=}\<%
\\
\>[0][@{}l@{\AgdaIndent{0}}]%
\>[2]\AgdaBound{h}\AgdaSpace{}%
\AgdaBound{x}\<%
\\
\>[0]\AgdaFunction{pres-simultaneous-subst}\AgdaSpace{}%
\AgdaBound{fb'}\AgdaSpace{}%
\AgdaBound{f'}\AgdaSpace{}%
\AgdaBound{h}\AgdaSpace{}%
\AgdaSymbol{(}\AgdaInductiveConstructor{ƛ}\AgdaSpace{}%
\AgdaBound{M}\AgdaSymbol{)}\AgdaSpace{}%
\AgdaSymbol{=}\<%
\\
\>[0][@{}l@{\AgdaIndent{0}}]%
\>[2]\AgdaFunction{cong}\AgdaSpace{}%
\AgdaInductiveConstructor{box}\AgdaSpace{}%
\AgdaSymbol{(}\AgdaFunction{cong}\AgdaSpace{}%
\AgdaInductiveConstructor{ƛb}\AgdaSpace{}%
\AgdaSymbol{(}\AgdaFunction{pres-simultaneous-subst}\AgdaSpace{}%
\AgdaSymbol{(}\AgdaFunction{extsb}\AgdaSpace{}%
\AgdaBound{fb'}\AgdaSymbol{)}\AgdaSpace{}%
\AgdaSymbol{(}\AgdaFunction{exts}\AgdaSpace{}%
\AgdaBound{f'}\AgdaSymbol{)}\AgdaSpace{}%
\AgdaSymbol{(}\AgdaFunction{pres-exts}\AgdaSpace{}%
\AgdaBound{fb'}\AgdaSpace{}%
\AgdaBound{f'}\AgdaSpace{}%
\AgdaBound{h}\AgdaSymbol{)}\AgdaSpace{}%
\AgdaBound{M}\AgdaSymbol{))}\<%
\\
\>[0]\AgdaFunction{pres-simultaneous-subst}\AgdaSpace{}%
\AgdaBound{fb'}\AgdaSpace{}%
\AgdaBound{f'}\AgdaSpace{}%
\AgdaBound{h}\AgdaSpace{}%
\AgdaSymbol{(}\AgdaBound{M}\AgdaSpace{}%
\AgdaOperator{\AgdaInductiveConstructor{·}}\AgdaSpace{}%
\AgdaBound{M₁}\AgdaSymbol{)}\AgdaSpace{}%
\AgdaSymbol{=}\<%
\\
\>[0][@{}l@{\AgdaIndent{0}}]%
\>[2]\AgdaOperator{\AgdaFunction{beginEq}}\<%
\\
\>[2][@{}l@{\AgdaIndent{0}}]%
\>[4]\AgdaInductiveConstructor{ƛb}\AgdaSpace{}%
\AgdaSymbol{(}\AgdaInductiveConstructor{ε}\AgdaSpace{}%
\AgdaInductiveConstructor{Zb}\AgdaSpace{}%
\AgdaOperator{\AgdaInductiveConstructor{·b}}\AgdaSpace{}%
\AgdaFunction{substb}\AgdaSpace{}%
\AgdaSymbol{(}\AgdaFunction{extsb}\AgdaSpace{}%
\AgdaBound{fb'}\AgdaSymbol{)}\AgdaSpace{}%
\AgdaSymbol{(}\AgdaFunction{renameb}\AgdaSpace{}%
\AgdaInductiveConstructor{Sb}\AgdaSpace{}%
\AgdaSymbol{(}\AgdaFunction{embedTerm}\AgdaSpace{}%
\AgdaBound{M₁}\AgdaSymbol{)))}\AgdaSpace{}%
\AgdaOperator{\AgdaInductiveConstructor{·b}}\AgdaSpace{}%
\AgdaFunction{substb}\AgdaSpace{}%
\AgdaBound{fb'}\AgdaSpace{}%
\AgdaSymbol{(}\AgdaFunction{embedTerm}\AgdaSpace{}%
\AgdaBound{M}\AgdaSymbol{)}\<%
\\
%
\>[2]\AgdaFunction{≡⟨}\AgdaSpace{}%
\AgdaFunction{cong₂}\AgdaSpace{}%
\AgdaOperator{\AgdaInductiveConstructor{\AgdaUnderscore{}·b\AgdaUnderscore{}}}\AgdaSpace{}%
\AgdaInductiveConstructor{refl}\AgdaSpace{}%
\AgdaSymbol{(}\AgdaFunction{pres-simultaneous-subst}\AgdaSpace{}%
\AgdaBound{fb'}\AgdaSpace{}%
\AgdaBound{f'}\AgdaSpace{}%
\AgdaBound{h}\AgdaSpace{}%
\AgdaBound{M}\AgdaSymbol{)}\AgdaSpace{}%
\AgdaFunction{⟩}\<%
\\
\>[2][@{}l@{\AgdaIndent{0}}]%
\>[4]\AgdaInductiveConstructor{ƛb}\AgdaSpace{}%
\AgdaSymbol{(}\AgdaInductiveConstructor{ε}\AgdaSpace{}%
\AgdaInductiveConstructor{Zb}\AgdaSpace{}%
\AgdaOperator{\AgdaInductiveConstructor{·b}}\AgdaSpace{}%
\AgdaFunction{substb}\AgdaSpace{}%
\AgdaSymbol{(}\AgdaFunction{extsb}\AgdaSpace{}%
\AgdaBound{fb'}\AgdaSymbol{)}\AgdaSpace{}%
\AgdaSymbol{(}\AgdaFunction{renameb}\AgdaSpace{}%
\AgdaInductiveConstructor{Sb}\AgdaSpace{}%
\AgdaSymbol{(}\AgdaFunction{embedTerm}\AgdaSpace{}%
\AgdaBound{M₁}\AgdaSymbol{)))}\AgdaSpace{}%
\AgdaOperator{\AgdaInductiveConstructor{·b}}\AgdaSpace{}%
\AgdaFunction{embedTerm}\AgdaSpace{}%
\AgdaSymbol{(}\AgdaFunction{subst}\AgdaSpace{}%
\AgdaBound{f'}\AgdaSpace{}%
\AgdaBound{M}\AgdaSymbol{)}\<%
\\
%
\>[2]\AgdaFunction{≡⟨}\AgdaSpace{}%
\AgdaFunction{cong₂}\AgdaSpace{}%
\AgdaOperator{\AgdaInductiveConstructor{\AgdaUnderscore{}·b\AgdaUnderscore{}}}\AgdaSpace{}%
\AgdaSymbol{(}\AgdaFunction{cong}\AgdaSpace{}%
\AgdaInductiveConstructor{ƛb}\AgdaSpace{}%
\AgdaSymbol{(}\AgdaFunction{cong₂}\AgdaSpace{}%
\AgdaOperator{\AgdaInductiveConstructor{\AgdaUnderscore{}·b\AgdaUnderscore{}}}\AgdaSpace{}%
\AgdaInductiveConstructor{refl}\AgdaSpace{}%
\AgdaSymbol{(}\AgdaFunction{substb-exts-rename-Sb-commute}\AgdaSpace{}%
\AgdaBound{fb'}\AgdaSpace{}%
\AgdaSymbol{(}\AgdaFunction{embedTerm}\AgdaSpace{}%
\AgdaBound{M₁}\AgdaSymbol{))))}\AgdaSpace{}%
\AgdaInductiveConstructor{refl}\AgdaSpace{}%
\AgdaFunction{⟩}\<%
\\
\>[2][@{}l@{\AgdaIndent{0}}]%
\>[4]\AgdaInductiveConstructor{ƛb}\AgdaSpace{}%
\AgdaSymbol{(}\AgdaInductiveConstructor{ε}\AgdaSpace{}%
\AgdaInductiveConstructor{Zb}\AgdaSpace{}%
\AgdaOperator{\AgdaInductiveConstructor{·b}}\AgdaSpace{}%
\AgdaFunction{renameb}\AgdaSpace{}%
\AgdaInductiveConstructor{Sb}\AgdaSpace{}%
\AgdaSymbol{(}\AgdaFunction{substb}\AgdaSpace{}%
\AgdaBound{fb'}\AgdaSpace{}%
\AgdaSymbol{(}\AgdaFunction{embedTerm}\AgdaSpace{}%
\AgdaBound{M₁}\AgdaSymbol{)))}\AgdaSpace{}%
\AgdaOperator{\AgdaInductiveConstructor{·b}}\AgdaSpace{}%
\AgdaFunction{embedTerm}\AgdaSpace{}%
\AgdaSymbol{(}\AgdaFunction{subst}\AgdaSpace{}%
\AgdaBound{f'}\AgdaSpace{}%
\AgdaBound{M}\AgdaSymbol{)}\<%
\\
%
\>[2]\AgdaFunction{≡⟨}\AgdaSpace{}%
\AgdaFunction{cong₂}\AgdaSpace{}%
\AgdaOperator{\AgdaInductiveConstructor{\AgdaUnderscore{}·b\AgdaUnderscore{}}}\AgdaSpace{}%
\AgdaSymbol{(}\AgdaFunction{cong}\AgdaSpace{}%
\AgdaInductiveConstructor{ƛb}\AgdaSpace{}%
\AgdaSymbol{(}\AgdaFunction{cong₂}\AgdaSpace{}%
\AgdaOperator{\AgdaInductiveConstructor{\AgdaUnderscore{}·b\AgdaUnderscore{}}}\AgdaSpace{}%
\AgdaInductiveConstructor{refl}\AgdaSpace{}%
\AgdaSymbol{(}\AgdaFunction{cong}\AgdaSpace{}%
\AgdaSymbol{(}\AgdaFunction{renameb}\AgdaSpace{}%
\AgdaInductiveConstructor{Sb}\AgdaSymbol{)}\AgdaSpace{}%
\AgdaSymbol{(}\AgdaFunction{pres-simultaneous-subst}\AgdaSpace{}%
\AgdaBound{fb'}\AgdaSpace{}%
\AgdaBound{f'}\AgdaSpace{}%
\AgdaBound{h}\AgdaSpace{}%
\AgdaBound{M₁}\AgdaSymbol{))))}\AgdaSpace{}%
\AgdaInductiveConstructor{refl}\AgdaSpace{}%
\AgdaFunction{⟩}\<%
\\
\>[2][@{}l@{\AgdaIndent{0}}]%
\>[4]\AgdaInductiveConstructor{ƛb}\AgdaSpace{}%
\AgdaSymbol{(}\AgdaInductiveConstructor{ε}\AgdaSpace{}%
\AgdaInductiveConstructor{Zb}\AgdaSpace{}%
\AgdaOperator{\AgdaInductiveConstructor{·b}}\AgdaSpace{}%
\AgdaFunction{renameb}\AgdaSpace{}%
\AgdaInductiveConstructor{Sb}\AgdaSpace{}%
\AgdaSymbol{(}\AgdaFunction{embedTerm}\AgdaSpace{}%
\AgdaSymbol{(}\AgdaFunction{subst}\AgdaSpace{}%
\AgdaBound{f'}\AgdaSpace{}%
\AgdaBound{M₁}\AgdaSymbol{)))}\AgdaSpace{}%
\AgdaOperator{\AgdaInductiveConstructor{·b}}\AgdaSpace{}%
\AgdaFunction{embedTerm}\AgdaSpace{}%
\AgdaSymbol{(}\AgdaFunction{subst}\AgdaSpace{}%
\AgdaBound{f'}\AgdaSpace{}%
\AgdaBound{M}\AgdaSymbol{)}\<%
\\
%
\>[2]\AgdaOperator{\AgdaFunction{∎eq}}\<%
\\
%
\\[\AgdaEmptyExtraSkip]%
\>[0]\AgdaComment{--\ HELPER\ LEMMA'S\ TO\ PROVE\ h}\<%
\\
%
\\[\AgdaEmptyExtraSkip]%
\>[0]\AgdaComment{--\ Lemma\ for\ assumption\ h\ in\ lemma}\<%
\\
\>[0]\AgdaComment{--\ pres-simultaneous-subst}\<%
\\
\>[0]\AgdaFunction{pres-subst-func}\AgdaSpace{}%
\AgdaSymbol{:}\AgdaSpace{}%
\AgdaSymbol{∀}\AgdaSpace{}%
\AgdaSymbol{\{}\AgdaBound{A}\AgdaSpace{}%
\AgdaBound{C}\AgdaSpace{}%
\AgdaBound{Γ}\AgdaSymbol{\}}\<%
\\
\>[0][@{}l@{\AgdaIndent{0}}]%
\>[2]\AgdaSymbol{→}\AgdaSpace{}%
\AgdaSymbol{(}\AgdaBound{V}\AgdaSpace{}%
\AgdaSymbol{:}\AgdaSpace{}%
\AgdaBound{Γ}\AgdaSpace{}%
\AgdaOperator{\AgdaDatatype{⊢}}\AgdaSpace{}%
\AgdaBound{C}\AgdaSymbol{)}\<%
\\
%
\>[2]\AgdaSymbol{→}\AgdaSpace{}%
\AgdaSymbol{(}\AgdaBound{y}\AgdaSpace{}%
\AgdaSymbol{:}\AgdaSpace{}%
\AgdaDatatype{Value}\AgdaSpace{}%
\AgdaBound{V}\AgdaSymbol{)}\<%
\\
%
\>[2]\AgdaSymbol{→}\AgdaSpace{}%
\AgdaSymbol{(}\AgdaBound{x}\AgdaSpace{}%
\AgdaSymbol{:}\AgdaSpace{}%
\AgdaBound{A}\AgdaSpace{}%
\AgdaOperator{\AgdaDatatype{∈}}\AgdaSpace{}%
\AgdaBound{Γ}\AgdaSpace{}%
\AgdaOperator{\AgdaInductiveConstructor{,}}\AgdaSpace{}%
\AgdaBound{C}\AgdaSymbol{)}\<%
\\
%
\>[2]\AgdaSymbol{→}\AgdaSpace{}%
\AgdaInductiveConstructor{box}\AgdaSpace{}%
\AgdaSymbol{(}\AgdaFunction{fb}\AgdaSpace{}%
\AgdaSymbol{(}\AgdaFunction{embedValue}\AgdaSpace{}%
\AgdaBound{V}\AgdaSpace{}%
\AgdaBound{y}\AgdaSymbol{)}\AgdaSpace{}%
\AgdaSymbol{(}\AgdaFunction{embedLookup}\AgdaSpace{}%
\AgdaBound{x}\AgdaSymbol{))}\AgdaSpace{}%
\AgdaOperator{\AgdaDatatype{≡}}\AgdaSpace{}%
\AgdaFunction{embedTerm}\AgdaSpace{}%
\AgdaSymbol{((}\AgdaFunction{f}\AgdaSpace{}%
\AgdaBound{V}\AgdaSymbol{)}\AgdaSpace{}%
\AgdaBound{x}\AgdaSymbol{)}\<%
\\
\>[0]\AgdaFunction{pres-subst-func}\AgdaSpace{}%
\AgdaBound{V}\AgdaSpace{}%
\AgdaBound{y}\AgdaSpace{}%
\AgdaInductiveConstructor{Z}\AgdaSpace{}%
\AgdaSymbol{=}\AgdaSpace{}%
\AgdaFunction{box-embedValue}\AgdaSpace{}%
\AgdaBound{y}\<%
\\
\>[0]\AgdaFunction{pres-subst-func}\AgdaSpace{}%
\AgdaBound{V}\AgdaSpace{}%
\AgdaBound{y}\AgdaSpace{}%
\AgdaSymbol{(}\AgdaInductiveConstructor{S}\AgdaSpace{}%
\AgdaBound{x}\AgdaSymbol{)}\AgdaSpace{}%
\AgdaSymbol{=}\AgdaSpace{}%
\AgdaInductiveConstructor{refl}\<%
\\
%
\\[\AgdaEmptyExtraSkip]%
\>[0]\AgdaComment{--\ exts\ lemma\ for\ assumption\ h\ in\ lemma}\<%
\\
\>[0]\AgdaFunction{exts-pres-subst-func}\AgdaSpace{}%
\AgdaSymbol{:}\AgdaSpace{}%
\AgdaSymbol{∀}\AgdaSpace{}%
\AgdaSymbol{\{}\AgdaBound{A}\AgdaSpace{}%
\AgdaBound{C}\AgdaSpace{}%
\AgdaBound{Γ}\AgdaSpace{}%
\AgdaBound{B}\AgdaSymbol{\}}\<%
\\
\>[0][@{}l@{\AgdaIndent{0}}]%
\>[2]\AgdaSymbol{→}\AgdaSpace{}%
\AgdaSymbol{(}\AgdaBound{V}\AgdaSpace{}%
\AgdaSymbol{:}\AgdaSpace{}%
\AgdaBound{Γ}\AgdaSpace{}%
\AgdaOperator{\AgdaDatatype{⊢}}\AgdaSpace{}%
\AgdaBound{C}\AgdaSymbol{)}\<%
\\
%
\>[2]\AgdaSymbol{→}\AgdaSpace{}%
\AgdaSymbol{(}\AgdaBound{y}\AgdaSpace{}%
\AgdaSymbol{:}\AgdaSpace{}%
\AgdaDatatype{Value}\AgdaSpace{}%
\AgdaBound{V}\AgdaSymbol{)}\<%
\\
%
\>[2]\AgdaSymbol{→}\AgdaSpace{}%
\AgdaSymbol{(}\AgdaBound{x}\AgdaSpace{}%
\AgdaSymbol{:}\AgdaSpace{}%
\AgdaBound{A}\AgdaSpace{}%
\AgdaOperator{\AgdaDatatype{∈}}\AgdaSpace{}%
\AgdaBound{Γ}\AgdaSpace{}%
\AgdaOperator{\AgdaInductiveConstructor{,}}\AgdaSpace{}%
\AgdaBound{C}\AgdaSpace{}%
\AgdaOperator{\AgdaInductiveConstructor{,}}\AgdaSpace{}%
\AgdaBound{B}\AgdaSymbol{)}\<%
\\
%
\>[2]\AgdaSymbol{→}\AgdaSpace{}%
\AgdaInductiveConstructor{box}\AgdaSpace{}%
\AgdaSymbol{(}\AgdaFunction{extsb}\AgdaSpace{}%
\AgdaSymbol{(}\AgdaFunction{fb}\AgdaSpace{}%
\AgdaSymbol{(}\AgdaFunction{embedValue}\AgdaSpace{}%
\AgdaBound{V}\AgdaSpace{}%
\AgdaBound{y}\AgdaSymbol{))}\AgdaSpace{}%
\AgdaSymbol{(}\AgdaFunction{embedLookup}\AgdaSpace{}%
\AgdaBound{x}\AgdaSymbol{))}\AgdaSpace{}%
\AgdaOperator{\AgdaDatatype{≡}}\AgdaSpace{}%
\AgdaFunction{embedTerm}\AgdaSpace{}%
\AgdaSymbol{(}\AgdaFunction{exts}\AgdaSpace{}%
\AgdaSymbol{(}\AgdaFunction{f}\AgdaSpace{}%
\AgdaBound{V}\AgdaSymbol{)}\AgdaSpace{}%
\AgdaBound{x}\AgdaSymbol{)}\<%
\\
\>[0]\AgdaFunction{exts-pres-subst-func}\AgdaSpace{}%
\AgdaBound{V}\AgdaSpace{}%
\AgdaBound{y}\AgdaSpace{}%
\AgdaInductiveConstructor{Z}\AgdaSpace{}%
\AgdaSymbol{=}\AgdaSpace{}%
\AgdaInductiveConstructor{refl}\<%
\\
\>[0]\AgdaFunction{exts-pres-subst-func}\AgdaSpace{}%
\AgdaBound{V}\AgdaSpace{}%
\AgdaBound{y}\AgdaSpace{}%
\AgdaSymbol{(}\AgdaInductiveConstructor{S}\AgdaSpace{}%
\AgdaBound{x}\AgdaSymbol{)}\AgdaSpace{}%
\AgdaSymbol{=}\<%
\\
\>[0][@{}l@{\AgdaIndent{0}}]%
\>[2]\AgdaOperator{\AgdaFunction{beginEq}}\<%
\\
\>[2][@{}l@{\AgdaIndent{0}}]%
\>[4]\AgdaInductiveConstructor{box}\AgdaSpace{}%
\AgdaSymbol{(}\AgdaFunction{renameb}\AgdaSpace{}%
\AgdaInductiveConstructor{Sb}\AgdaSpace{}%
\AgdaSymbol{(}\AgdaFunction{fb}\AgdaSpace{}%
\AgdaSymbol{(}\AgdaFunction{embedValue}\AgdaSpace{}%
\AgdaBound{V}\AgdaSpace{}%
\AgdaBound{y}\AgdaSymbol{)}\AgdaSpace{}%
\AgdaSymbol{(}\AgdaFunction{embedLookup}\AgdaSpace{}%
\AgdaBound{x}\AgdaSymbol{)))}\<%
\\
%
\>[2]\AgdaFunction{≡⟨}\AgdaSpace{}%
\AgdaInductiveConstructor{refl}\AgdaSpace{}%
\AgdaFunction{⟩}\<%
\\
\>[2][@{}l@{\AgdaIndent{0}}]%
\>[4]\AgdaFunction{renameb}\AgdaSpace{}%
\AgdaInductiveConstructor{Sb}\AgdaSpace{}%
\AgdaSymbol{(}\AgdaInductiveConstructor{box}\AgdaSpace{}%
\AgdaSymbol{(}\AgdaFunction{fb}\AgdaSpace{}%
\AgdaSymbol{(}\AgdaFunction{embedValue}\AgdaSpace{}%
\AgdaBound{V}\AgdaSpace{}%
\AgdaBound{y}\AgdaSymbol{)}\AgdaSpace{}%
\AgdaSymbol{(}\AgdaFunction{embedLookup}\AgdaSpace{}%
\AgdaBound{x}\AgdaSymbol{)))}\<%
\\
%
\>[2]\AgdaFunction{≡⟨}\AgdaSpace{}%
\AgdaFunction{cong}\AgdaSpace{}%
\AgdaSymbol{(}\AgdaFunction{renameb}\AgdaSpace{}%
\AgdaInductiveConstructor{Sb}\AgdaSymbol{)}\AgdaSpace{}%
\AgdaSymbol{(}\AgdaFunction{pres-subst-func}\AgdaSpace{}%
\AgdaBound{V}\AgdaSpace{}%
\AgdaBound{y}\AgdaSpace{}%
\AgdaBound{x}\AgdaSymbol{)}\AgdaSpace{}%
\AgdaFunction{⟩}\<%
\\
\>[2][@{}l@{\AgdaIndent{0}}]%
\>[4]\AgdaFunction{renameb}\AgdaSpace{}%
\AgdaInductiveConstructor{Sb}\AgdaSpace{}%
\AgdaSymbol{(}\AgdaFunction{embedTerm}\AgdaSpace{}%
\AgdaSymbol{(}\AgdaSpace{}%
\AgdaSymbol{(}\AgdaFunction{f}\AgdaSpace{}%
\AgdaBound{V}\AgdaSymbol{)}\AgdaSpace{}%
\AgdaBound{x}\AgdaSpace{}%
\AgdaSymbol{))}\<%
\\
%
\>[2]\AgdaFunction{≡⟨}\AgdaSpace{}%
\AgdaFunction{pres-rename}\AgdaSpace{}%
\AgdaInductiveConstructor{S}\AgdaSpace{}%
\AgdaInductiveConstructor{Sb}\AgdaSpace{}%
\AgdaSymbol{(λ}\AgdaSpace{}%
\AgdaBound{x₁}\AgdaSpace{}%
\AgdaSymbol{→}\AgdaSpace{}%
\AgdaInductiveConstructor{refl}\AgdaSymbol{)}\AgdaSpace{}%
\AgdaSymbol{(}\AgdaFunction{f}\AgdaSpace{}%
\AgdaBound{V}\AgdaSpace{}%
\AgdaBound{x}\AgdaSymbol{)}\AgdaSpace{}%
\AgdaFunction{⟩}\<%
\\
\>[2][@{}l@{\AgdaIndent{0}}]%
\>[4]\AgdaFunction{embedTerm}\AgdaSpace{}%
\AgdaSymbol{(}\AgdaFunction{rename}\AgdaSpace{}%
\AgdaInductiveConstructor{S}\AgdaSpace{}%
\AgdaSymbol{(}\AgdaFunction{f}\AgdaSpace{}%
\AgdaBound{V}\AgdaSpace{}%
\AgdaBound{x}\AgdaSymbol{))}\<%
\\
%
\>[2]\AgdaOperator{\AgdaFunction{∎eq}}\<%
\\
%
\\[\AgdaEmptyExtraSkip]%
\>[0]\AgdaComment{--\ Preservation\ and\ Reflection\ of\ }\<%
\\
\>[0]\AgdaComment{--\ single\ substitution\ of\ Gödel's\ translation}\<%
\\
\>[0]\AgdaFunction{pres-subst}\AgdaSpace{}%
\AgdaSymbol{:}\AgdaSpace{}%
\AgdaSymbol{∀}\AgdaSpace{}%
\AgdaSymbol{\{}\AgdaBound{Γ}\AgdaSpace{}%
\AgdaBound{A}\AgdaSpace{}%
\AgdaBound{B}\AgdaSymbol{\}}\AgdaSpace{}%
\AgdaSymbol{→}\AgdaSpace{}%
\AgdaSymbol{(}\AgdaBound{M}\AgdaSpace{}%
\AgdaSymbol{:}\AgdaSpace{}%
\AgdaBound{Γ}\AgdaSpace{}%
\AgdaOperator{\AgdaInductiveConstructor{,}}\AgdaSpace{}%
\AgdaBound{B}\AgdaSpace{}%
\AgdaOperator{\AgdaDatatype{⊢}}\AgdaSpace{}%
\AgdaBound{A}\AgdaSymbol{)}\AgdaSpace{}%
\AgdaSymbol{→}\AgdaSpace{}%
\AgdaSymbol{(}\AgdaBound{V}\AgdaSpace{}%
\AgdaSymbol{:}\AgdaSpace{}%
\AgdaBound{Γ}\AgdaSpace{}%
\AgdaOperator{\AgdaDatatype{⊢}}\AgdaSpace{}%
\AgdaBound{B}\AgdaSymbol{)}\AgdaSpace{}%
\AgdaSymbol{→}\AgdaSpace{}%
\AgdaSymbol{(}\AgdaBound{x}\AgdaSpace{}%
\AgdaSymbol{:}\AgdaSpace{}%
\AgdaDatatype{Value}\AgdaSpace{}%
\AgdaBound{V}\AgdaSymbol{)}\AgdaSpace{}%
\AgdaSymbol{→}\AgdaSpace{}%
\AgdaFunction{embedTerm}\AgdaSpace{}%
\AgdaBound{M}\AgdaSpace{}%
\AgdaOperator{\AgdaFunction{[}}\AgdaSpace{}%
\AgdaFunction{embedValue}\AgdaSpace{}%
\AgdaBound{V}\AgdaSpace{}%
\AgdaBound{x}\AgdaSpace{}%
\AgdaOperator{\AgdaFunction{]b}}\AgdaSpace{}%
\AgdaOperator{\AgdaDatatype{≡}}\AgdaSpace{}%
\AgdaFunction{embedTerm}\AgdaSpace{}%
\AgdaSymbol{(}\AgdaBound{M}\AgdaSpace{}%
\AgdaOperator{\AgdaFunction{[}}\AgdaSpace{}%
\AgdaBound{V}\AgdaSpace{}%
\AgdaOperator{\AgdaFunction{]}}\AgdaSymbol{)}\<%
\\
\>[0]\AgdaFunction{pres-subst}\AgdaSpace{}%
\AgdaBound{M}\AgdaSpace{}%
\AgdaBound{V}\AgdaSpace{}%
\AgdaBound{x}\AgdaSpace{}%
\AgdaSymbol{=}\AgdaSpace{}%
\AgdaFunction{pres-simultaneous-subst}\AgdaSpace{}%
\AgdaSymbol{(}\AgdaFunction{fb}\AgdaSpace{}%
\AgdaSymbol{(}\AgdaFunction{embedValue}\AgdaSpace{}%
\AgdaBound{V}\AgdaSpace{}%
\AgdaBound{x}\AgdaSymbol{))}\AgdaSpace{}%
\AgdaSymbol{(}\AgdaFunction{f}\AgdaSpace{}%
\AgdaBound{V}\AgdaSymbol{)}\AgdaSpace{}%
\AgdaSymbol{(}\AgdaFunction{pres-subst-func}\AgdaSpace{}%
\AgdaBound{V}\AgdaSpace{}%
\AgdaBound{x}\AgdaSymbol{)}\AgdaSpace{}%
\AgdaBound{M}\<%
\\
%
\\[\AgdaEmptyExtraSkip]%
\>[0]\AgdaComment{--\ Helper\ lemma:\ if\ a\ value\ is\ substituted\ in\ a\ value,\ the\ resulting\ lambda\ term\ will\ also\ be\ a\ value}\<%
\\
\>[0]\AgdaFunction{val-after-subst}\AgdaSpace{}%
\AgdaSymbol{:}\AgdaSpace{}%
\AgdaSymbol{∀}\AgdaSpace{}%
\AgdaSymbol{\{}\AgdaBound{Γ}\AgdaSpace{}%
\AgdaBound{A}\AgdaSpace{}%
\AgdaBound{B}\AgdaSymbol{\}}\AgdaSpace{}%
\AgdaSymbol{→}\AgdaSpace{}%
\AgdaSymbol{(}\AgdaBound{M}\AgdaSpace{}%
\AgdaSymbol{:}\AgdaSpace{}%
\AgdaBound{Γ}\AgdaSpace{}%
\AgdaOperator{\AgdaInductiveConstructor{,}}\AgdaSpace{}%
\AgdaBound{B}\AgdaSpace{}%
\AgdaOperator{\AgdaDatatype{⊢}}\AgdaSpace{}%
\AgdaBound{A}\AgdaSymbol{)}\AgdaSpace{}%
\AgdaSymbol{→}\AgdaSpace{}%
\AgdaSymbol{(}\AgdaBound{N}\AgdaSpace{}%
\AgdaSymbol{:}\AgdaSpace{}%
\AgdaBound{Γ}\AgdaSpace{}%
\AgdaOperator{\AgdaDatatype{⊢}}\AgdaSpace{}%
\AgdaBound{B}\AgdaSymbol{)}\AgdaSpace{}%
\AgdaSymbol{→}\AgdaSpace{}%
\AgdaSymbol{(}\AgdaBound{x}\AgdaSpace{}%
\AgdaSymbol{:}\AgdaSpace{}%
\AgdaDatatype{Value}\AgdaSpace{}%
\AgdaBound{M}\AgdaSymbol{)}\AgdaSpace{}%
\AgdaSymbol{→}\AgdaSpace{}%
\AgdaSymbol{(}\AgdaBound{y}\AgdaSpace{}%
\AgdaSymbol{:}\AgdaSpace{}%
\AgdaDatatype{Value}\AgdaSpace{}%
\AgdaBound{N}\AgdaSymbol{)}\AgdaSpace{}%
\AgdaSymbol{→}\AgdaSpace{}%
\AgdaDatatype{Value}\AgdaSpace{}%
\AgdaSymbol{(}\AgdaBound{M}\AgdaSpace{}%
\AgdaOperator{\AgdaFunction{[}}\AgdaSpace{}%
\AgdaBound{N}\AgdaSpace{}%
\AgdaOperator{\AgdaFunction{]}}\AgdaSymbol{)}\<%
\\
\>[0]\AgdaFunction{val-after-subst}\AgdaSpace{}%
\AgdaSymbol{(}\AgdaInductiveConstructor{`}\AgdaSpace{}%
\AgdaInductiveConstructor{Z}\AgdaSymbol{)}\AgdaSpace{}%
\AgdaBound{N}\AgdaSpace{}%
\AgdaBound{x}\AgdaSpace{}%
\AgdaBound{y}\AgdaSpace{}%
\AgdaSymbol{=}\AgdaSpace{}%
\AgdaBound{y}\<%
\\
\>[0]\AgdaFunction{val-after-subst}\AgdaSpace{}%
\AgdaSymbol{(}\AgdaInductiveConstructor{`}\AgdaSpace{}%
\AgdaSymbol{(}\AgdaInductiveConstructor{S}\AgdaSpace{}%
\AgdaBound{x₁}\AgdaSymbol{))}\AgdaSpace{}%
\AgdaBound{N}\AgdaSpace{}%
\AgdaBound{x}\AgdaSpace{}%
\AgdaBound{y}\AgdaSpace{}%
\AgdaSymbol{=}\AgdaSpace{}%
\AgdaInductiveConstructor{V-`}\<%
\\
\>[0]\AgdaFunction{val-after-subst}\AgdaSpace{}%
\AgdaSymbol{(}\AgdaInductiveConstructor{ƛ}\AgdaSpace{}%
\AgdaBound{M}\AgdaSymbol{)}\AgdaSpace{}%
\AgdaBound{N}\AgdaSpace{}%
\AgdaBound{x}\AgdaSpace{}%
\AgdaBound{y}\AgdaSpace{}%
\AgdaSymbol{=}\AgdaSpace{}%
\AgdaInductiveConstructor{V-ƛ}\<%
\\
%
\\[\AgdaEmptyExtraSkip]%
\>[0]\AgdaComment{--\ Preservation\ and\ Reflection\ of\ }\<%
\\
\>[0]\AgdaComment{--\ single\ substitution\ in\ values\ of\ Gödel's\ translation\ }\<%
\\
\>[0]\AgdaFunction{pres-subst-val}\AgdaSpace{}%
\AgdaSymbol{:}\AgdaSpace{}%
\AgdaSymbol{∀}\AgdaSpace{}%
\AgdaSymbol{\{}\AgdaBound{Γ}\AgdaSpace{}%
\AgdaBound{A}\AgdaSpace{}%
\AgdaBound{B}\AgdaSymbol{\}}\<%
\\
\>[0][@{}l@{\AgdaIndent{0}}]%
\>[2]\AgdaSymbol{→}\AgdaSpace{}%
\AgdaSymbol{(}\AgdaBound{M}\AgdaSpace{}%
\AgdaSymbol{:}\AgdaSpace{}%
\AgdaBound{Γ}\AgdaSpace{}%
\AgdaOperator{\AgdaInductiveConstructor{,}}\AgdaSpace{}%
\AgdaBound{B}\AgdaSpace{}%
\AgdaOperator{\AgdaDatatype{⊢}}\AgdaSpace{}%
\AgdaBound{A}\AgdaSymbol{)}\<%
\\
%
\>[2]\AgdaSymbol{→}\AgdaSpace{}%
\AgdaSymbol{(}\AgdaBound{V}\AgdaSpace{}%
\AgdaSymbol{:}\AgdaSpace{}%
\AgdaBound{Γ}\AgdaSpace{}%
\AgdaOperator{\AgdaDatatype{⊢}}\AgdaSpace{}%
\AgdaBound{B}\AgdaSymbol{)}\<%
\\
%
\>[2]\AgdaSymbol{→}\AgdaSpace{}%
\AgdaSymbol{(}\AgdaBound{x}\AgdaSpace{}%
\AgdaSymbol{:}\AgdaSpace{}%
\AgdaDatatype{Value}\AgdaSpace{}%
\AgdaBound{M}\AgdaSymbol{)}\<%
\\
%
\>[2]\AgdaSymbol{→}\AgdaSpace{}%
\AgdaSymbol{(}\AgdaBound{y}\AgdaSpace{}%
\AgdaSymbol{:}\AgdaSpace{}%
\AgdaDatatype{Value}\AgdaSpace{}%
\AgdaBound{V}\AgdaSymbol{)}\<%
\\
%
\>[2]\AgdaSymbol{→}\AgdaSpace{}%
\AgdaFunction{embedValue}\AgdaSpace{}%
\AgdaBound{M}\AgdaSpace{}%
\AgdaBound{x}\AgdaSpace{}%
\AgdaOperator{\AgdaFunction{[}}\AgdaSpace{}%
\AgdaFunction{embedValue}\AgdaSpace{}%
\AgdaBound{V}\AgdaSpace{}%
\AgdaBound{y}\AgdaSpace{}%
\AgdaOperator{\AgdaFunction{]b}}\AgdaSpace{}%
\AgdaOperator{\AgdaDatatype{≡}}\AgdaSpace{}%
\AgdaFunction{embedValue}\AgdaSpace{}%
\AgdaSymbol{(}\AgdaBound{M}\AgdaSpace{}%
\AgdaOperator{\AgdaFunction{[}}\AgdaSpace{}%
\AgdaBound{V}\AgdaSpace{}%
\AgdaOperator{\AgdaFunction{]}}\AgdaSymbol{)}\AgdaSpace{}%
\AgdaSymbol{(}\AgdaFunction{val-after-subst}\AgdaSpace{}%
\AgdaBound{M}\AgdaSpace{}%
\AgdaBound{V}\AgdaSpace{}%
\AgdaBound{x}\AgdaSpace{}%
\AgdaBound{y}\AgdaSymbol{)}\<%
\\
\>[0]\AgdaFunction{pres-subst-val}\AgdaSpace{}%
\AgdaSymbol{(}\AgdaInductiveConstructor{`}\AgdaSpace{}%
\AgdaInductiveConstructor{Z}\AgdaSymbol{)}\AgdaSpace{}%
\AgdaBound{V}\AgdaSpace{}%
\AgdaBound{x}\AgdaSpace{}%
\AgdaBound{y}\AgdaSpace{}%
\AgdaSymbol{=}\AgdaSpace{}%
\AgdaInductiveConstructor{refl}\<%
\\
\>[0]\AgdaFunction{pres-subst-val}\AgdaSpace{}%
\AgdaSymbol{(}\AgdaInductiveConstructor{`}\AgdaSpace{}%
\AgdaSymbol{(}\AgdaInductiveConstructor{S}\AgdaSpace{}%
\AgdaBound{x₁}\AgdaSymbol{))}\AgdaSpace{}%
\AgdaBound{V}\AgdaSpace{}%
\AgdaBound{x}\AgdaSpace{}%
\AgdaBound{y}\AgdaSpace{}%
\AgdaSymbol{=}\AgdaSpace{}%
\AgdaInductiveConstructor{refl}\<%
\\
\>[0]\AgdaFunction{pres-subst-val}\AgdaSpace{}%
\AgdaSymbol{(}\AgdaInductiveConstructor{ƛ}\AgdaSpace{}%
\AgdaBound{M}\AgdaSymbol{)}\AgdaSpace{}%
\AgdaBound{V}\AgdaSpace{}%
\AgdaBound{x}\AgdaSpace{}%
\AgdaBound{y}\AgdaSpace{}%
\AgdaSymbol{=}\<%
\\
\>[0][@{}l@{\AgdaIndent{0}}]%
\>[2]\AgdaOperator{\AgdaFunction{beginEq}}\<%
\\
\>[2][@{}l@{\AgdaIndent{0}}]%
\>[4]\AgdaFunction{embedValue}\AgdaSpace{}%
\AgdaSymbol{(}\AgdaInductiveConstructor{ƛ}\AgdaSpace{}%
\AgdaBound{M}\AgdaSymbol{)}\AgdaSpace{}%
\AgdaBound{x}\AgdaSpace{}%
\AgdaOperator{\AgdaFunction{[}}\AgdaSpace{}%
\AgdaFunction{embedValue}\AgdaSpace{}%
\AgdaBound{V}\AgdaSpace{}%
\AgdaBound{y}\AgdaSpace{}%
\AgdaOperator{\AgdaFunction{]b}}\<%
\\
%
\>[2]\AgdaFunction{≡⟨}\AgdaSpace{}%
\AgdaInductiveConstructor{refl}\AgdaSpace{}%
\AgdaFunction{⟩}\<%
\\
\>[2][@{}l@{\AgdaIndent{0}}]%
\>[4]\AgdaInductiveConstructor{ƛb}\AgdaSpace{}%
\AgdaSymbol{(}\AgdaFunction{substb}\AgdaSpace{}%
\AgdaSymbol{(}\AgdaFunction{extsb}\AgdaSpace{}%
\AgdaSymbol{(}\AgdaFunction{fb}\AgdaSpace{}%
\AgdaSymbol{(}\AgdaFunction{embedValue}\AgdaSpace{}%
\AgdaBound{V}\AgdaSpace{}%
\AgdaBound{y}\AgdaSymbol{)))}\AgdaSpace{}%
\AgdaSymbol{(}\AgdaFunction{embedTerm}\AgdaSpace{}%
\AgdaBound{M}\AgdaSymbol{))}\<%
\\
%
\>[2]\AgdaFunction{≡⟨}\AgdaSpace{}%
\AgdaFunction{cong}\AgdaSpace{}%
\AgdaInductiveConstructor{ƛb}\AgdaSpace{}%
\AgdaSymbol{(}\AgdaFunction{pres-simultaneous-subst}\AgdaSpace{}%
\AgdaSymbol{(}\AgdaFunction{extsb}\AgdaSpace{}%
\AgdaSymbol{(}\AgdaFunction{fb}\AgdaSpace{}%
\AgdaSymbol{(}\AgdaFunction{embedValue}\AgdaSpace{}%
\AgdaBound{V}\AgdaSpace{}%
\AgdaBound{y}\AgdaSymbol{)))}\AgdaSpace{}%
\AgdaSymbol{(}\AgdaFunction{exts}\AgdaSpace{}%
\AgdaSymbol{(}\AgdaFunction{f}\AgdaSpace{}%
\AgdaBound{V}\AgdaSymbol{))}\AgdaSpace{}%
\AgdaSymbol{(}\AgdaFunction{exts-pres-subst-func}\AgdaSpace{}%
\AgdaBound{V}\AgdaSpace{}%
\AgdaBound{y}\AgdaSymbol{)}\AgdaSpace{}%
\AgdaBound{M}\AgdaSymbol{)}\AgdaSpace{}%
\AgdaFunction{⟩}\<%
\\
\>[2][@{}l@{\AgdaIndent{0}}]%
\>[4]\AgdaInductiveConstructor{ƛb}\AgdaSpace{}%
\AgdaSymbol{(}\AgdaFunction{embedTerm}\AgdaSpace{}%
\AgdaSymbol{(}\AgdaFunction{subst}\AgdaSpace{}%
\AgdaSymbol{(}\AgdaFunction{exts}\AgdaSpace{}%
\AgdaSymbol{(}\AgdaFunction{f}\AgdaSpace{}%
\AgdaBound{V}\AgdaSymbol{))}\AgdaSpace{}%
\AgdaBound{M}\AgdaSymbol{))}\<%
\\
%
\>[2]\AgdaFunction{≡⟨}\AgdaSpace{}%
\AgdaInductiveConstructor{refl}\AgdaSpace{}%
\AgdaFunction{⟩}\<%
\\
\>[2][@{}l@{\AgdaIndent{0}}]%
\>[4]\AgdaFunction{embedValue}\AgdaSpace{}%
\AgdaSymbol{((}\AgdaInductiveConstructor{ƛ}\AgdaSpace{}%
\AgdaBound{M}\AgdaSymbol{)}\AgdaSpace{}%
\AgdaOperator{\AgdaFunction{[}}\AgdaSpace{}%
\AgdaBound{V}\AgdaSpace{}%
\AgdaOperator{\AgdaFunction{]}}\AgdaSymbol{)}\AgdaSpace{}%
\AgdaSymbol{(}\AgdaFunction{val-after-subst}\AgdaSpace{}%
\AgdaSymbol{(}\AgdaInductiveConstructor{ƛ}\AgdaSpace{}%
\AgdaBound{M}\AgdaSymbol{)}\AgdaSpace{}%
\AgdaBound{V}\AgdaSpace{}%
\AgdaBound{x}\AgdaSpace{}%
\AgdaBound{y}\AgdaSymbol{)}\<%
\\
%
\>[2]\AgdaOperator{\AgdaFunction{∎eq}}\<%
\\
%
\\[\AgdaEmptyExtraSkip]%
\>[0]\AgdaComment{--\ fb\ and\ renameb\ commute\ under\ extended\ context.}\<%
\\
\>[0]\AgdaFunction{fb-renameb-extb-commute}\AgdaSpace{}%
\AgdaSymbol{:}\AgdaSpace{}%
\AgdaSymbol{∀}\AgdaSpace{}%
\AgdaSymbol{\{}\AgdaBound{Γ}\AgdaSpace{}%
\AgdaBound{A}\AgdaSpace{}%
\AgdaBound{B}\AgdaSpace{}%
\AgdaBound{Δ}\AgdaSymbol{\}}\<%
\\
\>[0][@{}l@{\AgdaIndent{0}}]%
\>[2]\AgdaSymbol{→}\AgdaSpace{}%
\AgdaSymbol{(}\AgdaBound{fb'}\AgdaSpace{}%
\AgdaSymbol{:}\AgdaSpace{}%
\AgdaSymbol{∀}\AgdaSpace{}%
\AgdaSymbol{\{}\AgdaBound{Ab}\AgdaSymbol{\}}\AgdaSpace{}%
\AgdaSymbol{→}\AgdaSpace{}%
\AgdaBound{Ab}\AgdaSpace{}%
\AgdaOperator{\AgdaDatatype{∈b}}\AgdaSpace{}%
\AgdaBound{Γ}\AgdaSpace{}%
\AgdaSymbol{→}\AgdaSpace{}%
\AgdaBound{Ab}\AgdaSpace{}%
\AgdaOperator{\AgdaDatatype{∈b}}\AgdaSpace{}%
\AgdaBound{Δ}\AgdaSymbol{)}\<%
\\
%
\>[2]\AgdaSymbol{→}\AgdaSpace{}%
\AgdaSymbol{(}\AgdaBound{W}\AgdaSpace{}%
\AgdaSymbol{:}\AgdaSpace{}%
\AgdaBound{Γ}\AgdaSpace{}%
\AgdaOperator{\AgdaDatatype{⊢b}}\AgdaSpace{}%
\AgdaBound{B}\AgdaSymbol{)}\<%
\\
%
\>[2]\AgdaSymbol{→}\AgdaSpace{}%
\AgdaSymbol{(}\AgdaBound{x}\AgdaSpace{}%
\AgdaSymbol{:}\AgdaSpace{}%
\AgdaInductiveConstructor{□}\AgdaSpace{}%
\AgdaBound{A}\AgdaSpace{}%
\AgdaOperator{\AgdaDatatype{∈b}}\AgdaSpace{}%
\AgdaBound{Γ}\AgdaSpace{}%
\AgdaOperator{\AgdaInductiveConstructor{,b}}\AgdaSpace{}%
\AgdaInductiveConstructor{□}\AgdaSpace{}%
\AgdaBound{B}\AgdaSymbol{)}\<%
\\
%
\>[2]\AgdaSymbol{→}\AgdaSpace{}%
\AgdaFunction{fb}\AgdaSpace{}%
\AgdaSymbol{(}\AgdaFunction{renameb}\AgdaSpace{}%
\AgdaBound{fb'}\AgdaSpace{}%
\AgdaBound{W}\AgdaSymbol{)}\AgdaSpace{}%
\AgdaSymbol{(}\AgdaFunction{extb}\AgdaSpace{}%
\AgdaBound{fb'}\AgdaSpace{}%
\AgdaBound{x}\AgdaSymbol{)}\AgdaSpace{}%
\AgdaOperator{\AgdaDatatype{≡}}\AgdaSpace{}%
\AgdaFunction{renameb}\AgdaSpace{}%
\AgdaBound{fb'}\AgdaSpace{}%
\AgdaSymbol{(}\AgdaFunction{fb}\AgdaSpace{}%
\AgdaBound{W}\AgdaSpace{}%
\AgdaBound{x}\AgdaSymbol{)}\<%
\\
\>[0]\AgdaFunction{fb-renameb-extb-commute}\AgdaSpace{}%
\AgdaBound{fb'}\AgdaSpace{}%
\AgdaBound{W}\AgdaSpace{}%
\AgdaInductiveConstructor{Zb}\AgdaSpace{}%
\AgdaSymbol{=}\AgdaSpace{}%
\AgdaInductiveConstructor{refl}\<%
\\
\>[0]\AgdaFunction{fb-renameb-extb-commute}\AgdaSpace{}%
\AgdaBound{fb'}\AgdaSpace{}%
\AgdaBound{W}\AgdaSpace{}%
\AgdaSymbol{(}\AgdaInductiveConstructor{Sb}\AgdaSpace{}%
\AgdaBound{x}\AgdaSymbol{)}\AgdaSpace{}%
\AgdaSymbol{=}\AgdaSpace{}%
\AgdaInductiveConstructor{refl}\<%
\\
%
\\[\AgdaEmptyExtraSkip]%
\>[0]\AgdaComment{--\ Preservation\ and\ reflection\ of\ }\<%
\\
\>[0]\AgdaComment{--\ single\ substitution\ under\ renaming}\<%
\\
\>[0]\AgdaFunction{pres-subst-func-rename}\AgdaSpace{}%
\AgdaSymbol{:}\AgdaSpace{}%
\AgdaSymbol{∀}\AgdaSpace{}%
\AgdaSymbol{\{}\AgdaBound{Γ}\AgdaSpace{}%
\AgdaBound{B}\AgdaSpace{}%
\AgdaBound{A}\AgdaSpace{}%
\AgdaBound{Δ}\AgdaSymbol{\}}\<%
\\
\>[0][@{}l@{\AgdaIndent{0}}]%
\>[2]\AgdaSymbol{→}\AgdaSpace{}%
\AgdaSymbol{(}\AgdaBound{fb'}\AgdaSpace{}%
\AgdaSymbol{:}\AgdaSpace{}%
\AgdaSymbol{∀}\AgdaSpace{}%
\AgdaSymbol{\{}\AgdaBound{Ab}\AgdaSymbol{\}}\AgdaSpace{}%
\AgdaSymbol{→}\AgdaSpace{}%
\AgdaBound{Ab}\AgdaSpace{}%
\AgdaOperator{\AgdaDatatype{∈b}}\AgdaSpace{}%
\AgdaBound{Γ}\AgdaSpace{}%
\AgdaSymbol{→}\AgdaSpace{}%
\AgdaBound{Ab}\AgdaSpace{}%
\AgdaOperator{\AgdaDatatype{∈b}}\AgdaSpace{}%
\AgdaBound{Δ}\AgdaSymbol{)}\<%
\\
%
\>[2]\AgdaSymbol{→}\AgdaSpace{}%
\AgdaSymbol{(}\AgdaBound{M}\AgdaSpace{}%
\AgdaSymbol{:}\AgdaSpace{}%
\AgdaBound{Γ}\AgdaSpace{}%
\AgdaOperator{\AgdaInductiveConstructor{,b}}\AgdaSpace{}%
\AgdaInductiveConstructor{□}\AgdaSpace{}%
\AgdaBound{B}\AgdaSpace{}%
\AgdaOperator{\AgdaDatatype{⊢b}}\AgdaSpace{}%
\AgdaBound{A}\AgdaSymbol{)}\<%
\\
%
\>[2]\AgdaSymbol{→}\AgdaSpace{}%
\AgdaSymbol{(}\AgdaBound{W}\AgdaSpace{}%
\AgdaSymbol{:}\AgdaSpace{}%
\AgdaBound{Γ}\AgdaSpace{}%
\AgdaOperator{\AgdaDatatype{⊢b}}\AgdaSpace{}%
\AgdaBound{B}\AgdaSymbol{)}\<%
\\
%
\>[2]\AgdaSymbol{→}%
\>[2253I]\AgdaFunction{renameb}\AgdaSpace{}%
\AgdaSymbol{(}\AgdaFunction{extb}\AgdaSpace{}%
\AgdaBound{fb'}\AgdaSymbol{)}\AgdaSpace{}%
\AgdaBound{M}\AgdaSpace{}%
\AgdaOperator{\AgdaFunction{[}}\AgdaSpace{}%
\AgdaFunction{renameb}\AgdaSpace{}%
\AgdaBound{fb'}\AgdaSpace{}%
\AgdaBound{W}\AgdaSpace{}%
\AgdaOperator{\AgdaFunction{]b}}\<%
\\
\>[.][@{}l@{}]\<[2253I]%
\>[4]\AgdaOperator{\AgdaDatatype{≡}}\AgdaSpace{}%
\AgdaFunction{renameb}\AgdaSpace{}%
\AgdaBound{fb'}\AgdaSpace{}%
\AgdaSymbol{(}\AgdaBound{M}\AgdaSpace{}%
\AgdaOperator{\AgdaFunction{[}}\AgdaSpace{}%
\AgdaBound{W}\AgdaSpace{}%
\AgdaOperator{\AgdaFunction{]b}}\AgdaSymbol{)}\<%
\\
\>[0]\AgdaFunction{pres-subst-func-rename}\AgdaSpace{}%
\AgdaBound{fb'}\AgdaSpace{}%
\AgdaBound{M}\AgdaSpace{}%
\AgdaBound{W}\AgdaSpace{}%
\AgdaSymbol{=}\<%
\\
\>[0][@{}l@{\AgdaIndent{0}}]%
\>[2]\AgdaOperator{\AgdaFunction{beginEq}}\<%
\\
\>[2][@{}l@{\AgdaIndent{0}}]%
\>[4]\AgdaFunction{substb}\AgdaSpace{}%
\AgdaSymbol{(}\AgdaFunction{fb}\AgdaSpace{}%
\AgdaSymbol{(}\AgdaFunction{renameb}\AgdaSpace{}%
\AgdaBound{fb'}\AgdaSpace{}%
\AgdaBound{W}\AgdaSymbol{))}\AgdaSpace{}%
\AgdaSymbol{(}\AgdaFunction{renameb}\AgdaSpace{}%
\AgdaSymbol{(}\AgdaFunction{extb}\AgdaSpace{}%
\AgdaBound{fb'}\AgdaSymbol{)}\AgdaSpace{}%
\AgdaBound{M}\AgdaSymbol{)}\<%
\\
%
\>[2]\AgdaFunction{≡⟨}\AgdaSpace{}%
\AgdaFunction{subst-renameb-comp}\AgdaSpace{}%
\AgdaSymbol{(}\AgdaFunction{extb}\AgdaSpace{}%
\AgdaBound{fb'}\AgdaSymbol{)}\AgdaSpace{}%
\AgdaSymbol{(}\AgdaFunction{fb}\AgdaSpace{}%
\AgdaSymbol{(}\AgdaFunction{renameb}\AgdaSpace{}%
\AgdaBound{fb'}\AgdaSpace{}%
\AgdaBound{W}\AgdaSymbol{))}\AgdaSpace{}%
\AgdaBound{M}\AgdaSpace{}%
\AgdaFunction{⟩}\<%
\\
\>[2][@{}l@{\AgdaIndent{0}}]%
\>[4]\AgdaFunction{substb}\AgdaSpace{}%
\AgdaSymbol{(λ}\AgdaSpace{}%
\AgdaBound{x}\AgdaSpace{}%
\AgdaSymbol{→}\AgdaSpace{}%
\AgdaFunction{fb}\AgdaSpace{}%
\AgdaSymbol{(}\AgdaFunction{renameb}\AgdaSpace{}%
\AgdaBound{fb'}\AgdaSpace{}%
\AgdaBound{W}\AgdaSymbol{)}\AgdaSpace{}%
\AgdaSymbol{(}\AgdaFunction{extb}\AgdaSpace{}%
\AgdaBound{fb'}\AgdaSpace{}%
\AgdaBound{x}\AgdaSymbol{))}\AgdaSpace{}%
\AgdaBound{M}\<%
\\
%
\>[2]\AgdaFunction{≡⟨}\AgdaSpace{}%
\AgdaFunction{substb-cong}\AgdaSpace{}%
\AgdaSymbol{(}\AgdaFunction{fb-renameb-extb-commute}\AgdaSpace{}%
\AgdaBound{fb'}\AgdaSpace{}%
\AgdaBound{W}\AgdaSymbol{)}\AgdaSpace{}%
\AgdaBound{M}\AgdaSpace{}%
\AgdaFunction{⟩}\<%
\\
\>[2][@{}l@{\AgdaIndent{0}}]%
\>[4]\AgdaFunction{substb}\AgdaSpace{}%
\AgdaSymbol{(λ}\AgdaSpace{}%
\AgdaBound{x}\AgdaSpace{}%
\AgdaSymbol{→}\AgdaSpace{}%
\AgdaFunction{renameb}\AgdaSpace{}%
\AgdaBound{fb'}\AgdaSpace{}%
\AgdaSymbol{(}\AgdaFunction{fb}\AgdaSpace{}%
\AgdaBound{W}\AgdaSpace{}%
\AgdaBound{x}\AgdaSymbol{))}\AgdaSpace{}%
\AgdaBound{M}\<%
\\
%
\>[2]\AgdaFunction{≡⟨}\AgdaSpace{}%
\AgdaFunction{substb-rename-commute}\AgdaSpace{}%
\AgdaBound{fb'}\AgdaSpace{}%
\AgdaSymbol{(}\AgdaFunction{fb}\AgdaSpace{}%
\AgdaBound{W}\AgdaSymbol{)}\AgdaSpace{}%
\AgdaBound{M}\AgdaSpace{}%
\AgdaFunction{⟩}\<%
\\
\>[2][@{}l@{\AgdaIndent{0}}]%
\>[4]\AgdaFunction{renameb}\AgdaSpace{}%
\AgdaBound{fb'}\AgdaSpace{}%
\AgdaSymbol{(}\AgdaFunction{substb}\AgdaSpace{}%
\AgdaSymbol{(}\AgdaFunction{fb}\AgdaSpace{}%
\AgdaBound{W}\AgdaSymbol{)}\AgdaSpace{}%
\AgdaBound{M}\AgdaSymbol{)}\<%
\\
%
\>[2]\AgdaOperator{\AgdaFunction{∎eq}}\<%
\\
%
\\[\AgdaEmptyExtraSkip]%
\>[0]\AgdaComment{--\ Preservation\ of\ reduction\ under\ renaming\ \ \ }\<%
\\
\>[0]\AgdaFunction{pres-rename-red}\AgdaSpace{}%
\AgdaSymbol{:}\AgdaSpace{}%
\AgdaSymbol{∀}\AgdaSpace{}%
\AgdaSymbol{\{}\AgdaBound{Γ}\AgdaSpace{}%
\AgdaBound{A}\AgdaSpace{}%
\AgdaBound{Δ}\AgdaSymbol{\}}\<%
\\
\>[0][@{}l@{\AgdaIndent{0}}]%
\>[2]\AgdaSymbol{→}\AgdaSpace{}%
\AgdaSymbol{(}\AgdaBound{fb'}\AgdaSpace{}%
\AgdaSymbol{:}\AgdaSpace{}%
\AgdaSymbol{∀}\AgdaSpace{}%
\AgdaSymbol{\{}\AgdaBound{Ab}\AgdaSymbol{\}}\AgdaSpace{}%
\AgdaSymbol{→}\AgdaSpace{}%
\AgdaBound{Ab}\AgdaSpace{}%
\AgdaOperator{\AgdaDatatype{∈b}}\AgdaSpace{}%
\AgdaBound{Γ}\AgdaSpace{}%
\AgdaSymbol{→}\AgdaSpace{}%
\AgdaBound{Ab}\AgdaSpace{}%
\AgdaOperator{\AgdaDatatype{∈b}}\AgdaSpace{}%
\AgdaBound{Δ}\AgdaSymbol{)}\<%
\\
%
\>[2]\AgdaSymbol{→}\AgdaSpace{}%
\AgdaSymbol{(}\AgdaBound{M}\AgdaSpace{}%
\AgdaSymbol{:}\AgdaSpace{}%
\AgdaBound{Γ}\AgdaSpace{}%
\AgdaOperator{\AgdaDatatype{⊢b}}\AgdaSpace{}%
\AgdaBound{A}\AgdaSymbol{)}\<%
\\
%
\>[2]\AgdaSymbol{→}\AgdaSpace{}%
\AgdaSymbol{(}\AgdaBound{N}\AgdaSpace{}%
\AgdaSymbol{:}\AgdaSpace{}%
\AgdaBound{Γ}\AgdaSpace{}%
\AgdaOperator{\AgdaDatatype{⊢b}}\AgdaSpace{}%
\AgdaBound{A}\AgdaSymbol{)}\<%
\\
%
\>[2]\AgdaSymbol{→}\AgdaSpace{}%
\AgdaSymbol{(}\AgdaBound{M}\AgdaSpace{}%
\AgdaOperator{\AgdaDatatype{↝βb}}\AgdaSpace{}%
\AgdaBound{N}\AgdaSymbol{)}\<%
\\
%
\>[2]\AgdaSymbol{→}\AgdaSpace{}%
\AgdaFunction{renameb}\AgdaSpace{}%
\AgdaBound{fb'}\AgdaSpace{}%
\AgdaBound{M}\AgdaSpace{}%
\AgdaOperator{\AgdaDatatype{↝βb}}\AgdaSpace{}%
\AgdaFunction{renameb}\AgdaSpace{}%
\AgdaBound{fb'}\AgdaSpace{}%
\AgdaBound{N}\<%
\\
\>[0]\AgdaFunction{pres-rename-red}\AgdaSpace{}%
\AgdaBound{fb'}\AgdaSpace{}%
\AgdaSymbol{(}\AgdaInductiveConstructor{ƛb}\AgdaSpace{}%
\AgdaBound{M₁}\AgdaSpace{}%
\AgdaOperator{\AgdaInductiveConstructor{·b}}\AgdaSpace{}%
\AgdaInductiveConstructor{box}\AgdaSpace{}%
\AgdaBound{M₂}\AgdaSymbol{)}\AgdaSpace{}%
\AgdaBound{N}\AgdaSpace{}%
\AgdaInductiveConstructor{βb}\AgdaSpace{}%
\AgdaSymbol{=}\AgdaSpace{}%
\AgdaFunction{Eq.subst}\AgdaSpace{}%
\AgdaSymbol{(λ}\AgdaSpace{}%
\AgdaBound{x}\AgdaSpace{}%
\AgdaSymbol{→}\AgdaSpace{}%
\AgdaInductiveConstructor{ƛb}\AgdaSpace{}%
\AgdaSymbol{(}\AgdaFunction{renameb}\AgdaSpace{}%
\AgdaSymbol{(}\AgdaFunction{extb}\AgdaSpace{}%
\AgdaBound{fb'}\AgdaSymbol{)}\AgdaSpace{}%
\AgdaBound{M₁}\AgdaSymbol{)}\AgdaSpace{}%
\AgdaOperator{\AgdaInductiveConstructor{·b}}\AgdaSpace{}%
\AgdaInductiveConstructor{box}\AgdaSpace{}%
\AgdaSymbol{(}\AgdaFunction{renameb}\AgdaSpace{}%
\AgdaBound{fb'}\AgdaSpace{}%
\AgdaBound{M₂}\AgdaSymbol{)}\AgdaSpace{}%
\AgdaOperator{\AgdaDatatype{↝βb}}\AgdaSpace{}%
\AgdaBound{x}\AgdaSymbol{)}\AgdaSpace{}%
\AgdaSymbol{(}\AgdaFunction{pres-subst-func-rename}\AgdaSpace{}%
\AgdaBound{fb'}\AgdaSpace{}%
\AgdaBound{M₁}\AgdaSpace{}%
\AgdaBound{M₂}\AgdaSymbol{)}\AgdaSpace{}%
\AgdaInductiveConstructor{βb}\<%
\\
\>[0]\AgdaFunction{pres-rename-red}\AgdaSpace{}%
\AgdaBound{fb'}\AgdaSpace{}%
\AgdaBound{M}\AgdaSpace{}%
\AgdaBound{N}\AgdaSpace{}%
\AgdaSymbol{(}\AgdaInductiveConstructor{μ}\AgdaSpace{}%
\AgdaSymbol{\{}\AgdaArgument{M}\AgdaSpace{}%
\AgdaSymbol{=}\AgdaSpace{}%
\AgdaBound{M₁}\AgdaSymbol{\}}\AgdaSpace{}%
\AgdaSymbol{\{}\AgdaArgument{M'}\AgdaSpace{}%
\AgdaSymbol{=}\AgdaSpace{}%
\AgdaBound{M'}\AgdaSymbol{\}}\AgdaSpace{}%
\AgdaBound{red}\AgdaSymbol{)}\AgdaSpace{}%
\AgdaSymbol{=}\AgdaSpace{}%
\AgdaInductiveConstructor{μ}\AgdaSpace{}%
\AgdaSymbol{(}\AgdaFunction{pres-rename-red}\AgdaSpace{}%
\AgdaBound{fb'}\AgdaSpace{}%
\AgdaBound{M₁}\AgdaSpace{}%
\AgdaBound{M'}\AgdaSpace{}%
\AgdaBound{red}\AgdaSymbol{)}\<%
\\
\>[0]\AgdaFunction{pres-rename-red}\AgdaSpace{}%
\AgdaBound{fb'}\AgdaSpace{}%
\AgdaBound{M}\AgdaSpace{}%
\AgdaBound{N}\AgdaSpace{}%
\AgdaSymbol{(}\AgdaInductiveConstructor{ν}\AgdaSpace{}%
\AgdaSymbol{\{}\AgdaArgument{M}\AgdaSpace{}%
\AgdaSymbol{=}\AgdaSpace{}%
\AgdaBound{M₁}\AgdaSymbol{\}}\AgdaSpace{}%
\AgdaSymbol{\{}\AgdaArgument{N}\AgdaSpace{}%
\AgdaSymbol{=}\AgdaSpace{}%
\AgdaBound{N₁}\AgdaSymbol{\}}\AgdaSpace{}%
\AgdaSymbol{\{}\AgdaArgument{N'}\AgdaSpace{}%
\AgdaSymbol{=}\AgdaSpace{}%
\AgdaBound{N'}\AgdaSymbol{\}}\AgdaSpace{}%
\AgdaBound{red}\AgdaSymbol{)}\AgdaSpace{}%
\AgdaSymbol{=}\AgdaSpace{}%
\AgdaInductiveConstructor{ν}\AgdaSpace{}%
\AgdaSymbol{(}\AgdaFunction{pres-rename-red}\AgdaSpace{}%
\AgdaBound{fb'}\AgdaSpace{}%
\AgdaBound{N₁}\AgdaSpace{}%
\AgdaBound{N'}\AgdaSpace{}%
\AgdaBound{red}\AgdaSymbol{)}\<%
\\
\>[0]\AgdaFunction{pres-rename-red}\AgdaSpace{}%
\AgdaBound{fb'}\AgdaSpace{}%
\AgdaBound{M}\AgdaSpace{}%
\AgdaBound{N}\AgdaSpace{}%
\AgdaSymbol{(}\AgdaInductiveConstructor{ξ}\AgdaSpace{}%
\AgdaSymbol{\{}\AgdaArgument{M}\AgdaSpace{}%
\AgdaSymbol{=}\AgdaSpace{}%
\AgdaBound{M₁}\AgdaSymbol{\}}\AgdaSpace{}%
\AgdaSymbol{\{}\AgdaArgument{M'}\AgdaSpace{}%
\AgdaSymbol{=}\AgdaSpace{}%
\AgdaBound{M'}\AgdaSymbol{\}}\AgdaSpace{}%
\AgdaBound{red}\AgdaSymbol{)}\AgdaSpace{}%
\AgdaSymbol{=}\AgdaSpace{}%
\AgdaInductiveConstructor{ξ}\AgdaSpace{}%
\AgdaSymbol{(}\AgdaFunction{pres-rename-red}\AgdaSpace{}%
\AgdaSymbol{(}\AgdaFunction{extb}\AgdaSpace{}%
\AgdaBound{fb'}\AgdaSymbol{)}\AgdaSpace{}%
\AgdaBound{M₁}\AgdaSpace{}%
\AgdaBound{M'}\AgdaSpace{}%
\AgdaBound{red}\AgdaSymbol{)}\<%
\\
\>[0]\AgdaFunction{pres-rename-red}\AgdaSpace{}%
\AgdaBound{fb'}\AgdaSpace{}%
\AgdaBound{M}\AgdaSpace{}%
\AgdaBound{N}\AgdaSpace{}%
\AgdaSymbol{(}\AgdaInductiveConstructor{ζ}\AgdaSpace{}%
\AgdaSymbol{\{}\AgdaArgument{M}\AgdaSpace{}%
\AgdaSymbol{=}\AgdaSpace{}%
\AgdaBound{M₁}\AgdaSymbol{\}}\AgdaSpace{}%
\AgdaSymbol{\{}\AgdaArgument{M'}\AgdaSpace{}%
\AgdaSymbol{=}\AgdaSpace{}%
\AgdaBound{M'}\AgdaSymbol{\}}\AgdaSpace{}%
\AgdaBound{red}\AgdaSymbol{)}\AgdaSpace{}%
\AgdaSymbol{=}\AgdaSpace{}%
\AgdaInductiveConstructor{ζ}\AgdaSpace{}%
\AgdaSymbol{(}\AgdaFunction{pres-rename-red}\AgdaSpace{}%
\AgdaBound{fb'}\AgdaSpace{}%
\AgdaBound{M₁}\AgdaSpace{}%
\AgdaBound{M'}\AgdaSpace{}%
\AgdaBound{red}\AgdaSymbol{)}\<%
\\
%
\\[\AgdaEmptyExtraSkip]%
\>[0]\AgdaComment{--\ Preservation\ of\ 2\ step\ reduction}\<%
\\
\>[0]\AgdaComment{--\ under\ renaming}\<%
\\
\>[0]\AgdaFunction{pres-rename-red-2}\AgdaSpace{}%
\AgdaSymbol{:}\AgdaSpace{}%
\AgdaSymbol{∀}\AgdaSpace{}%
\AgdaSymbol{\{}\AgdaBound{Γ}\AgdaSpace{}%
\AgdaBound{A}\AgdaSpace{}%
\AgdaBound{Δ}\AgdaSymbol{\}}\<%
\\
\>[0][@{}l@{\AgdaIndent{0}}]%
\>[2]\AgdaSymbol{→}\AgdaSpace{}%
\AgdaSymbol{(}\AgdaBound{fb'}\AgdaSpace{}%
\AgdaSymbol{:}\AgdaSpace{}%
\AgdaSymbol{∀}\AgdaSpace{}%
\AgdaSymbol{\{}\AgdaBound{Ab}\AgdaSymbol{\}}\AgdaSpace{}%
\AgdaSymbol{→}\AgdaSpace{}%
\AgdaBound{Ab}\AgdaSpace{}%
\AgdaOperator{\AgdaDatatype{∈b}}\AgdaSpace{}%
\AgdaBound{Γ}\AgdaSpace{}%
\AgdaSymbol{→}\AgdaSpace{}%
\AgdaBound{Ab}\AgdaSpace{}%
\AgdaOperator{\AgdaDatatype{∈b}}\AgdaSpace{}%
\AgdaBound{Δ}\AgdaSymbol{)}\<%
\\
%
\>[2]\AgdaSymbol{→}\AgdaSpace{}%
\AgdaSymbol{(}\AgdaBound{M}\AgdaSpace{}%
\AgdaSymbol{:}\AgdaSpace{}%
\AgdaBound{Γ}\AgdaSpace{}%
\AgdaOperator{\AgdaDatatype{⊢b}}\AgdaSpace{}%
\AgdaBound{A}\AgdaSymbol{)}\<%
\\
%
\>[2]\AgdaSymbol{→}\AgdaSpace{}%
\AgdaSymbol{(}\AgdaBound{N}\AgdaSpace{}%
\AgdaSymbol{:}\AgdaSpace{}%
\AgdaBound{Γ}\AgdaSpace{}%
\AgdaOperator{\AgdaDatatype{⊢b}}\AgdaSpace{}%
\AgdaBound{A}\AgdaSymbol{)}\<%
\\
%
\>[2]\AgdaSymbol{→}\AgdaSpace{}%
\AgdaSymbol{(}\AgdaBound{M}\AgdaSpace{}%
\AgdaOperator{\AgdaDatatype{↝βb²}}\AgdaSpace{}%
\AgdaBound{N}\AgdaSymbol{)}\<%
\\
%
\>[2]\AgdaSymbol{→}\AgdaSpace{}%
\AgdaFunction{renameb}\AgdaSpace{}%
\AgdaBound{fb'}\AgdaSpace{}%
\AgdaBound{M}\AgdaSpace{}%
\AgdaOperator{\AgdaDatatype{↝βb²}}\AgdaSpace{}%
\AgdaFunction{renameb}\AgdaSpace{}%
\AgdaBound{fb'}\AgdaSpace{}%
\AgdaBound{N}\<%
\\
\>[0]\AgdaFunction{pres-rename-red-2}\AgdaSpace{}%
\AgdaBound{fb'}\AgdaSpace{}%
\AgdaBound{M}\AgdaSpace{}%
\AgdaBound{N}\AgdaSpace{}%
\AgdaSymbol{(}\AgdaInductiveConstructor{βb²}\AgdaSpace{}%
\AgdaDottedPattern{\AgdaSymbol{.}}\AgdaDottedPattern{\AgdaBound{M}}\AgdaSpace{}%
\AgdaBound{t₂}\AgdaSpace{}%
\AgdaDottedPattern{\AgdaSymbol{.}}\AgdaDottedPattern{\AgdaBound{N}}\AgdaSpace{}%
\AgdaBound{x}\AgdaSpace{}%
\AgdaBound{x₁}\AgdaSymbol{)}\AgdaSpace{}%
\AgdaSymbol{=}\<%
\\
\>[0][@{}l@{\AgdaIndent{0}}]%
\>[2]\AgdaInductiveConstructor{βb²}\<%
\\
\>[2][@{}l@{\AgdaIndent{0}}]%
\>[4]\AgdaSymbol{(}\AgdaFunction{renameb}\AgdaSpace{}%
\AgdaBound{fb'}\AgdaSpace{}%
\AgdaBound{M}\AgdaSymbol{)}\<%
\\
%
\>[4]\AgdaSymbol{(}\AgdaFunction{renameb}\AgdaSpace{}%
\AgdaBound{fb'}\AgdaSpace{}%
\AgdaBound{t₂}\AgdaSymbol{)}\<%
\\
%
\>[4]\AgdaSymbol{(}\AgdaFunction{renameb}\AgdaSpace{}%
\AgdaBound{fb'}\AgdaSpace{}%
\AgdaBound{N}\AgdaSymbol{)}\<%
\\
%
\>[4]\AgdaSymbol{(}\AgdaFunction{pres-rename-red}\AgdaSpace{}%
\AgdaBound{fb'}\AgdaSpace{}%
\AgdaBound{M}\AgdaSpace{}%
\AgdaBound{t₂}\AgdaSpace{}%
\AgdaBound{x}\AgdaSymbol{)}\<%
\\
%
\>[4]\AgdaSymbol{(}\AgdaFunction{pres-rename-red}\AgdaSpace{}%
\AgdaBound{fb'}\AgdaSpace{}%
\AgdaBound{t₂}\AgdaSpace{}%
\AgdaBound{N}\AgdaSpace{}%
\AgdaBound{x₁}\AgdaSymbol{)}\<%
\\
\>[0]\AgdaFunction{pres-rename-red-2}\AgdaSpace{}%
\AgdaBound{fb'}\AgdaSpace{}%
\AgdaBound{M}\AgdaSpace{}%
\AgdaBound{N}\AgdaSpace{}%
\AgdaSymbol{(}\AgdaInductiveConstructor{μ}\AgdaSpace{}%
\AgdaSymbol{\{}\AgdaArgument{M}\AgdaSpace{}%
\AgdaSymbol{=}\AgdaSpace{}%
\AgdaBound{M₁}\AgdaSymbol{\}}\AgdaSpace{}%
\AgdaSymbol{\{}\AgdaArgument{M'}\AgdaSpace{}%
\AgdaSymbol{=}\AgdaSpace{}%
\AgdaBound{M'}\AgdaSymbol{\}}\AgdaSpace{}%
\AgdaBound{red}\AgdaSymbol{)}\AgdaSpace{}%
\AgdaSymbol{=}\<%
\\
\>[0][@{}l@{\AgdaIndent{0}}]%
\>[2]\AgdaInductiveConstructor{μ}\AgdaSpace{}%
\AgdaSymbol{(}\AgdaFunction{pres-rename-red-2}\AgdaSpace{}%
\AgdaBound{fb'}\AgdaSpace{}%
\AgdaBound{M₁}\AgdaSpace{}%
\AgdaBound{M'}\AgdaSpace{}%
\AgdaBound{red}\AgdaSymbol{)}\<%
\\
\>[0]\AgdaFunction{pres-rename-red-2}\AgdaSpace{}%
\AgdaBound{fb'}\AgdaSpace{}%
\AgdaBound{M}\AgdaSpace{}%
\AgdaBound{N}\AgdaSpace{}%
\AgdaSymbol{(}\AgdaInductiveConstructor{ν}\AgdaSpace{}%
\AgdaSymbol{\{}\AgdaArgument{M}\AgdaSpace{}%
\AgdaSymbol{=}\AgdaSpace{}%
\AgdaBound{M₁}\AgdaSymbol{\}}\AgdaSpace{}%
\AgdaSymbol{\{}\AgdaArgument{N}\AgdaSpace{}%
\AgdaSymbol{=}\AgdaSpace{}%
\AgdaBound{N₁}\AgdaSymbol{\}}\AgdaSpace{}%
\AgdaSymbol{\{}\AgdaArgument{N'}\AgdaSpace{}%
\AgdaSymbol{=}\AgdaSpace{}%
\AgdaBound{N'}\AgdaSymbol{\}}\AgdaSpace{}%
\AgdaBound{red}\AgdaSymbol{)}\AgdaSpace{}%
\AgdaSymbol{=}\<%
\\
\>[0][@{}l@{\AgdaIndent{0}}]%
\>[2]\AgdaInductiveConstructor{ν}\AgdaSpace{}%
\AgdaSymbol{(}\AgdaFunction{pres-rename-red-2}\AgdaSpace{}%
\AgdaBound{fb'}\AgdaSpace{}%
\AgdaBound{N₁}\AgdaSpace{}%
\AgdaBound{N'}\AgdaSpace{}%
\AgdaBound{red}\AgdaSymbol{)}\<%
\\
\>[0]\AgdaFunction{pres-rename-red-2}\AgdaSpace{}%
\AgdaBound{fb'}\AgdaSpace{}%
\AgdaBound{M}\AgdaSpace{}%
\AgdaBound{N}\AgdaSpace{}%
\AgdaSymbol{(}\AgdaInductiveConstructor{ξ}\AgdaSpace{}%
\AgdaSymbol{\{}\AgdaArgument{M}\AgdaSpace{}%
\AgdaSymbol{=}\AgdaSpace{}%
\AgdaBound{M₁}\AgdaSymbol{\}}\AgdaSpace{}%
\AgdaSymbol{\{}\AgdaArgument{M'}\AgdaSpace{}%
\AgdaSymbol{=}\AgdaSpace{}%
\AgdaBound{M'}\AgdaSymbol{\}}\AgdaSpace{}%
\AgdaBound{red}\AgdaSymbol{)}\AgdaSpace{}%
\AgdaSymbol{=}\<%
\\
\>[0][@{}l@{\AgdaIndent{0}}]%
\>[2]\AgdaInductiveConstructor{ξ}\AgdaSpace{}%
\AgdaSymbol{(}\AgdaFunction{pres-rename-red-2}\AgdaSpace{}%
\AgdaSymbol{(}\AgdaFunction{extb}\AgdaSpace{}%
\AgdaBound{fb'}\AgdaSymbol{)}\AgdaSpace{}%
\AgdaBound{M₁}\AgdaSpace{}%
\AgdaBound{M'}\AgdaSpace{}%
\AgdaBound{red}\AgdaSymbol{)}\<%
\\
\>[0]\AgdaFunction{pres-rename-red-2}\AgdaSpace{}%
\AgdaBound{fb'}\AgdaSpace{}%
\AgdaBound{M}\AgdaSpace{}%
\AgdaBound{N}\AgdaSpace{}%
\AgdaSymbol{(}\AgdaInductiveConstructor{ζ}\AgdaSpace{}%
\AgdaSymbol{\{}\AgdaArgument{M}\AgdaSpace{}%
\AgdaSymbol{=}\AgdaSpace{}%
\AgdaBound{M₁}\AgdaSymbol{\}}\AgdaSpace{}%
\AgdaSymbol{\{}\AgdaArgument{M'}\AgdaSpace{}%
\AgdaSymbol{=}\AgdaSpace{}%
\AgdaBound{M'}\AgdaSymbol{\}}\AgdaSpace{}%
\AgdaBound{red}\AgdaSymbol{)}\AgdaSpace{}%
\AgdaSymbol{=}\<%
\\
\>[0][@{}l@{\AgdaIndent{0}}]%
\>[2]\AgdaInductiveConstructor{ζ}\AgdaSpace{}%
\AgdaSymbol{(}\AgdaFunction{pres-rename-red-2}\AgdaSpace{}%
\AgdaBound{fb'}\AgdaSpace{}%
\AgdaBound{M₁}\AgdaSpace{}%
\AgdaBound{M'}\AgdaSpace{}%
\AgdaBound{red}\AgdaSymbol{)}\<%
\\
%
\\[\AgdaEmptyExtraSkip]%
\>[0]\AgdaComment{--\ extsb\ of\ ε\ is\ ε}\<%
\\
\>[0]\AgdaFunction{ε-exts}\AgdaSpace{}%
\AgdaSymbol{:}\AgdaSpace{}%
\AgdaSymbol{∀}\AgdaSpace{}%
\AgdaSymbol{\{}\AgdaBound{Γ}\AgdaSpace{}%
\AgdaBound{A}\AgdaSpace{}%
\AgdaBound{B}\AgdaSymbol{\}}\<%
\\
\>[0][@{}l@{\AgdaIndent{0}}]%
\>[2]\AgdaSymbol{→}\AgdaSpace{}%
\AgdaSymbol{(}\AgdaBound{x}\AgdaSpace{}%
\AgdaSymbol{:}\AgdaSpace{}%
\AgdaInductiveConstructor{□}\AgdaSpace{}%
\AgdaBound{A}\AgdaSpace{}%
\AgdaOperator{\AgdaDatatype{∈b}}\AgdaSpace{}%
\AgdaBound{Γ}\AgdaSpace{}%
\AgdaOperator{\AgdaInductiveConstructor{,b}}\AgdaSpace{}%
\AgdaBound{B}\AgdaSymbol{)}\<%
\\
%
\>[2]\AgdaSymbol{→}\AgdaSpace{}%
\AgdaFunction{extsb}\AgdaSpace{}%
\AgdaInductiveConstructor{ε}\AgdaSpace{}%
\AgdaBound{x}\AgdaSpace{}%
\AgdaOperator{\AgdaDatatype{≡}}\AgdaSpace{}%
\AgdaInductiveConstructor{ε}\AgdaSpace{}%
\AgdaBound{x}\<%
\\
\>[0]\AgdaFunction{ε-exts}\AgdaSpace{}%
\AgdaInductiveConstructor{Zb}\AgdaSpace{}%
\AgdaSymbol{=}\AgdaSpace{}%
\AgdaInductiveConstructor{refl}\<%
\\
\>[0]\AgdaFunction{ε-exts}\AgdaSpace{}%
\AgdaSymbol{(}\AgdaInductiveConstructor{Sb}\AgdaSpace{}%
\AgdaBound{x}\AgdaSymbol{)}\AgdaSpace{}%
\AgdaSymbol{=}\AgdaSpace{}%
\AgdaInductiveConstructor{refl}\<%
\\
%
\\[\AgdaEmptyExtraSkip]%
\>[0]\AgdaComment{--\ Simultaneous\ substitution\ with\ ε\ }\<%
\\
\>[0]\AgdaComment{--\ leaves\ the\ term\ unchanged\ \ }\<%
\\
\>[0]\AgdaFunction{subst-with-id}\AgdaSpace{}%
\AgdaSymbol{:}\AgdaSpace{}%
\AgdaSymbol{∀}\AgdaSpace{}%
\AgdaSymbol{\{}\AgdaBound{Γ}\AgdaSpace{}%
\AgdaBound{A}\AgdaSymbol{\}}\<%
\\
\>[0][@{}l@{\AgdaIndent{0}}]%
\>[2]\AgdaSymbol{→}\AgdaSpace{}%
\AgdaSymbol{(}\AgdaBound{M}\AgdaSpace{}%
\AgdaSymbol{:}\AgdaSpace{}%
\AgdaBound{Γ}\AgdaSpace{}%
\AgdaOperator{\AgdaDatatype{⊢b}}\AgdaSpace{}%
\AgdaBound{A}\AgdaSymbol{)}\<%
\\
%
\>[2]\AgdaSymbol{→}\AgdaSpace{}%
\AgdaFunction{substb}\AgdaSpace{}%
\AgdaSymbol{(λ}\AgdaSpace{}%
\AgdaBound{x₁}\AgdaSpace{}%
\AgdaSymbol{→}\AgdaSpace{}%
\AgdaInductiveConstructor{ε}\AgdaSpace{}%
\AgdaBound{x₁}\AgdaSymbol{)}\AgdaSpace{}%
\AgdaBound{M}\AgdaSpace{}%
\AgdaOperator{\AgdaDatatype{≡}}\AgdaSpace{}%
\AgdaBound{M}\<%
\\
\>[0]\AgdaFunction{subst-with-id}\AgdaSpace{}%
\AgdaSymbol{(}\AgdaInductiveConstructor{ε}\AgdaSpace{}%
\AgdaBound{x}\AgdaSymbol{)}\AgdaSpace{}%
\AgdaSymbol{=}\AgdaSpace{}%
\AgdaInductiveConstructor{refl}\<%
\\
\>[0]\AgdaFunction{subst-with-id}\AgdaSpace{}%
\AgdaSymbol{(}\AgdaInductiveConstructor{box}\AgdaSpace{}%
\AgdaBound{M}\AgdaSymbol{)}\AgdaSpace{}%
\AgdaSymbol{=}\AgdaSpace{}%
\AgdaFunction{cong}\AgdaSpace{}%
\AgdaInductiveConstructor{box}\AgdaSpace{}%
\AgdaSymbol{(}\AgdaFunction{subst-with-id}\AgdaSpace{}%
\AgdaBound{M}\AgdaSymbol{)}\<%
\\
\>[0]\AgdaFunction{subst-with-id}\AgdaSpace{}%
\AgdaSymbol{(}\AgdaInductiveConstructor{ƛb}\AgdaSpace{}%
\AgdaBound{M}\AgdaSymbol{)}\AgdaSpace{}%
\AgdaSymbol{=}\<%
\\
\>[0][@{}l@{\AgdaIndent{0}}]%
\>[2]\AgdaOperator{\AgdaFunction{beginEq}}\<%
\\
\>[2][@{}l@{\AgdaIndent{0}}]%
\>[4]\AgdaInductiveConstructor{ƛb}\AgdaSpace{}%
\AgdaSymbol{(}\AgdaFunction{substb}\AgdaSpace{}%
\AgdaSymbol{(}\AgdaFunction{extsb}\AgdaSpace{}%
\AgdaInductiveConstructor{ε}\AgdaSymbol{)}\AgdaSpace{}%
\AgdaBound{M}\AgdaSymbol{)}\<%
\\
%
\>[2]\AgdaFunction{≡⟨}\AgdaSpace{}%
\AgdaFunction{cong}\AgdaSpace{}%
\AgdaInductiveConstructor{ƛb}\AgdaSpace{}%
\AgdaSymbol{(}\AgdaFunction{substb-cong}\AgdaSpace{}%
\AgdaSymbol{(λ}\AgdaSpace{}%
\AgdaBound{x}\AgdaSpace{}%
\AgdaSymbol{→}\AgdaSpace{}%
\AgdaFunction{ε-exts}\AgdaSpace{}%
\AgdaBound{x}\AgdaSymbol{)}\AgdaSpace{}%
\AgdaBound{M}\AgdaSymbol{)}\AgdaSpace{}%
\AgdaFunction{⟩}\<%
\\
\>[2][@{}l@{\AgdaIndent{0}}]%
\>[4]\AgdaInductiveConstructor{ƛb}\AgdaSpace{}%
\AgdaSymbol{(}\AgdaFunction{substb}\AgdaSpace{}%
\AgdaInductiveConstructor{ε}\AgdaSpace{}%
\AgdaBound{M}\AgdaSymbol{)}\<%
\\
%
\>[2]\AgdaFunction{≡⟨}\AgdaSpace{}%
\AgdaFunction{cong}\AgdaSpace{}%
\AgdaInductiveConstructor{ƛb}\AgdaSpace{}%
\AgdaSymbol{(}\AgdaFunction{subst-with-id}\AgdaSpace{}%
\AgdaBound{M}\AgdaSymbol{)}\AgdaSpace{}%
\AgdaFunction{⟩}\<%
\\
\>[2][@{}l@{\AgdaIndent{0}}]%
\>[4]\AgdaInductiveConstructor{ƛb}\AgdaSpace{}%
\AgdaBound{M}\<%
\\
%
\>[2]\AgdaOperator{\AgdaFunction{∎eq}}\<%
\\
\>[0]\AgdaFunction{subst-with-id}\AgdaSpace{}%
\AgdaSymbol{(}\AgdaBound{M}\AgdaSpace{}%
\AgdaOperator{\AgdaInductiveConstructor{·b}}\AgdaSpace{}%
\AgdaBound{M₁}\AgdaSymbol{)}\AgdaSpace{}%
\AgdaSymbol{=}\AgdaSpace{}%
\AgdaFunction{cong₂}\AgdaSpace{}%
\AgdaOperator{\AgdaInductiveConstructor{\AgdaUnderscore{}·b\AgdaUnderscore{}}}\AgdaSpace{}%
\AgdaSymbol{(}\AgdaFunction{subst-with-id}\AgdaSpace{}%
\AgdaBound{M}\AgdaSymbol{)}\AgdaSpace{}%
\AgdaSymbol{(}\AgdaFunction{subst-with-id}\AgdaSpace{}%
\AgdaBound{M₁}\AgdaSymbol{)}\<%
\\
%
\\[\AgdaEmptyExtraSkip]%
\>[0]\AgdaComment{--\ Substitution\ composed\ with\ renaming\ with\ Sb\ }\<%
\\
\>[0]\AgdaComment{--\ is\ the\ same\ as\ embedding\ that\ term.}\<%
\\
\>[0]\AgdaFunction{subst-rename-fb-comp}\AgdaSpace{}%
\AgdaSymbol{:}\AgdaSpace{}%
\AgdaSymbol{∀}\AgdaSpace{}%
\AgdaSymbol{\{}\AgdaBound{Γ}\AgdaSpace{}%
\AgdaBound{A}\AgdaSpace{}%
\AgdaBound{B}\AgdaSymbol{\}}\<%
\\
\>[0][@{}l@{\AgdaIndent{0}}]%
\>[2]\AgdaSymbol{→}\AgdaSpace{}%
\AgdaSymbol{(}\AgdaBound{M₁}\AgdaSpace{}%
\AgdaSymbol{:}\AgdaSpace{}%
\AgdaBound{Γ}\AgdaSpace{}%
\AgdaOperator{\AgdaInductiveConstructor{,}}\AgdaSpace{}%
\AgdaBound{B}\AgdaSpace{}%
\AgdaOperator{\AgdaDatatype{⊢}}\AgdaSpace{}%
\AgdaBound{A}\AgdaSymbol{)}\<%
\\
%
\>[2]\AgdaSymbol{→}\AgdaSpace{}%
\AgdaSymbol{(}\AgdaBound{M₂}\AgdaSpace{}%
\AgdaSymbol{:}\AgdaSpace{}%
\AgdaBound{Γ}\AgdaSpace{}%
\AgdaOperator{\AgdaDatatype{⊢}}\AgdaSpace{}%
\AgdaBound{B}\AgdaSymbol{)}\<%
\\
%
\>[2]\AgdaSymbol{→}\AgdaSpace{}%
\AgdaSymbol{(}\AgdaBound{x}\AgdaSpace{}%
\AgdaSymbol{:}\AgdaSpace{}%
\AgdaDatatype{Value}\AgdaSpace{}%
\AgdaBound{M₂}\AgdaSpace{}%
\AgdaSymbol{)}\<%
\\
%
\>[2]\AgdaSymbol{→}%
\>[2715I]\AgdaFunction{substb}\AgdaSpace{}%
\AgdaSymbol{(}\AgdaFunction{fb}\AgdaSpace{}%
\AgdaSymbol{(}\AgdaInductiveConstructor{ƛb}\AgdaSpace{}%
\AgdaSymbol{(}\AgdaFunction{embedTerm}\AgdaSpace{}%
\AgdaBound{M₁}\AgdaSymbol{)))}\AgdaSpace{}%
\AgdaSymbol{(}\AgdaFunction{renameb}\AgdaSpace{}%
\AgdaInductiveConstructor{Sb}\AgdaSpace{}%
\AgdaSymbol{(}\AgdaFunction{embedTerm}\AgdaSpace{}%
\AgdaBound{M₂}\AgdaSymbol{))}\<%
\\
\>[.][@{}l@{}]\<[2715I]%
\>[4]\AgdaOperator{\AgdaDatatype{≡}}\AgdaSpace{}%
\AgdaInductiveConstructor{box}\AgdaSpace{}%
\AgdaSymbol{(}\AgdaFunction{embedValue}\AgdaSpace{}%
\AgdaBound{M₂}\AgdaSpace{}%
\AgdaBound{x}\AgdaSymbol{)}\<%
\\
\>[0]\AgdaFunction{subst-rename-fb-comp}\AgdaSpace{}%
\AgdaBound{M₁}\AgdaSpace{}%
\AgdaBound{M₂}\AgdaSpace{}%
\AgdaBound{x}\AgdaSpace{}%
\AgdaSymbol{=}\<%
\\
\>[0][@{}l@{\AgdaIndent{0}}]%
\>[2]\AgdaOperator{\AgdaFunction{beginEq}}\<%
\\
\>[2][@{}l@{\AgdaIndent{0}}]%
\>[4]\AgdaFunction{substb}\AgdaSpace{}%
\AgdaSymbol{(}\AgdaFunction{fb}\AgdaSpace{}%
\AgdaSymbol{(}\AgdaInductiveConstructor{ƛb}\AgdaSpace{}%
\AgdaSymbol{(}\AgdaFunction{embedTerm}\AgdaSpace{}%
\AgdaBound{M₁}\AgdaSymbol{)))}\AgdaSpace{}%
\AgdaSymbol{(}\AgdaFunction{renameb}\AgdaSpace{}%
\AgdaInductiveConstructor{Sb}\AgdaSpace{}%
\AgdaSymbol{(}\AgdaFunction{embedTerm}\AgdaSpace{}%
\AgdaBound{M₂}\AgdaSymbol{))}\<%
\\
%
\>[2]\AgdaFunction{≡⟨}\AgdaSpace{}%
\AgdaFunction{subst-renameb-comp}\AgdaSpace{}%
\AgdaInductiveConstructor{Sb}\AgdaSpace{}%
\AgdaSymbol{(}\AgdaFunction{fb}\AgdaSpace{}%
\AgdaSymbol{(}\AgdaInductiveConstructor{ƛb}\AgdaSpace{}%
\AgdaSymbol{(}\AgdaFunction{embedTerm}\AgdaSpace{}%
\AgdaBound{M₁}\AgdaSymbol{)))}\AgdaSpace{}%
\AgdaSymbol{(}\AgdaFunction{embedTerm}\AgdaSpace{}%
\AgdaBound{M₂}\AgdaSymbol{)}\AgdaSpace{}%
\AgdaFunction{⟩}\<%
\\
\>[2][@{}l@{\AgdaIndent{0}}]%
\>[4]\AgdaFunction{substb}\AgdaSpace{}%
\AgdaSymbol{(λ}\AgdaSpace{}%
\AgdaBound{x₁}\AgdaSpace{}%
\AgdaSymbol{→}\AgdaSpace{}%
\AgdaInductiveConstructor{ε}\AgdaSpace{}%
\AgdaBound{x₁}\AgdaSymbol{)}\AgdaSpace{}%
\AgdaSymbol{(}\AgdaFunction{embedTerm}\AgdaSpace{}%
\AgdaBound{M₂}\AgdaSymbol{)}\<%
\\
%
\>[2]\AgdaFunction{≡⟨}\AgdaSpace{}%
\AgdaFunction{subst-with-id}\AgdaSpace{}%
\AgdaSymbol{(}\AgdaFunction{embedTerm}\AgdaSpace{}%
\AgdaBound{M₂}\AgdaSymbol{)}\AgdaSpace{}%
\AgdaFunction{⟩}\<%
\\
\>[2][@{}l@{\AgdaIndent{0}}]%
\>[4]\AgdaFunction{embedTerm}\AgdaSpace{}%
\AgdaBound{M₂}\<%
\\
%
\>[2]\AgdaFunction{≡⟨}\AgdaSpace{}%
\AgdaFunction{sym}\AgdaSpace{}%
\AgdaSymbol{(}\AgdaFunction{box-embedValue}\AgdaSpace{}%
\AgdaBound{x}\AgdaSymbol{)}\AgdaSpace{}%
\AgdaFunction{⟩}\<%
\\
\>[2][@{}l@{\AgdaIndent{0}}]%
\>[4]\AgdaInductiveConstructor{box}\AgdaSpace{}%
\AgdaSymbol{(}\AgdaFunction{embedValue}\AgdaSpace{}%
\AgdaBound{M₂}\AgdaSpace{}%
\AgdaBound{x}\AgdaSymbol{)}\<%
\\
%
\>[2]\AgdaOperator{\AgdaFunction{∎eq}}\<%
\\
\>[0]\<%
\end{code}
\begin{code}%
\>[0]\AgdaComment{--\ Preservation\ of\ 2\ step\ reduction\ of\ Gödel's\ translation}\<%
\\
\>[0]\AgdaFunction{pres-red}\AgdaSpace{}%
\AgdaSymbol{:}\AgdaSpace{}%
\AgdaSymbol{∀}\AgdaSpace{}%
\AgdaSymbol{\{}\AgdaBound{Γ}\AgdaSpace{}%
\AgdaBound{A}\AgdaSymbol{\}}\<%
\\
\>[0][@{}l@{\AgdaIndent{0}}]%
\>[2]\AgdaSymbol{→}\AgdaSpace{}%
\AgdaSymbol{(}\AgdaBound{M}\AgdaSpace{}%
\AgdaSymbol{:}\AgdaSpace{}%
\AgdaBound{Γ}\AgdaSpace{}%
\AgdaOperator{\AgdaDatatype{⊢}}\AgdaSpace{}%
\AgdaBound{A}\AgdaSymbol{)}\<%
\\
%
\>[2]\AgdaSymbol{→}\AgdaSpace{}%
\AgdaSymbol{(}\AgdaBound{N}\AgdaSpace{}%
\AgdaSymbol{:}\AgdaSpace{}%
\AgdaBound{Γ}\AgdaSpace{}%
\AgdaOperator{\AgdaDatatype{⊢}}\AgdaSpace{}%
\AgdaBound{A}\AgdaSymbol{)}\<%
\\
%
\>[2]\AgdaSymbol{→}\AgdaSpace{}%
\AgdaSymbol{(}\AgdaBound{M}\AgdaSpace{}%
\AgdaOperator{\AgdaDatatype{↝βᵥ}}\AgdaSpace{}%
\AgdaBound{N}\AgdaSymbol{)}\<%
\\
%
\>[2]\AgdaSymbol{→}\AgdaSpace{}%
\AgdaSymbol{(}\AgdaFunction{embedTerm}\AgdaSpace{}%
\AgdaBound{M}\AgdaSpace{}%
\AgdaOperator{\AgdaDatatype{↝βb²}}\AgdaSpace{}%
\AgdaFunction{embedTerm}\AgdaSpace{}%
\AgdaBound{N}\AgdaSymbol{)}\<%
\end{code}
\begin{code}[hide]%
\>[0]\AgdaFunction{pres-red}\AgdaSpace{}%
\AgdaSymbol{(}\AgdaInductiveConstructor{ƛ}\AgdaSpace{}%
\AgdaBound{M₁}\AgdaSpace{}%
\AgdaOperator{\AgdaInductiveConstructor{·}}\AgdaSpace{}%
\AgdaBound{M₂}\AgdaSymbol{)}\AgdaSpace{}%
\AgdaBound{N}\AgdaSpace{}%
\AgdaSymbol{(}\AgdaInductiveConstructor{βᵥ}\AgdaSpace{}%
\AgdaBound{x}\AgdaSymbol{)}\AgdaSpace{}%
\AgdaSymbol{=}\<%
\\
\>[0][@{}l@{\AgdaIndent{0}}]%
\>[2]\AgdaInductiveConstructor{βb²}\<%
\\
%
\>[2]\AgdaSymbol{(}\AgdaFunction{raise}\AgdaSpace{}%
\AgdaSymbol{(}\AgdaFunction{embedTerm}\AgdaSpace{}%
\AgdaBound{M₂}\AgdaSymbol{)}\AgdaSpace{}%
\AgdaOperator{\AgdaInductiveConstructor{·b}}\AgdaSpace{}%
\AgdaInductiveConstructor{box}\AgdaSpace{}%
\AgdaSymbol{(}\AgdaInductiveConstructor{ƛb}\AgdaSpace{}%
\AgdaSymbol{(}\AgdaFunction{embedTerm}\AgdaSpace{}%
\AgdaBound{M₁}\AgdaSymbol{)))}\<%
\\
%
\>[2]\AgdaSymbol{(}\AgdaInductiveConstructor{ƛb}\AgdaSpace{}%
\AgdaSymbol{(}\AgdaFunction{embedTerm}\AgdaSpace{}%
\AgdaBound{M₁}\AgdaSymbol{)}\AgdaSpace{}%
\AgdaOperator{\AgdaInductiveConstructor{·b}}\AgdaSpace{}%
\AgdaInductiveConstructor{box}\AgdaSpace{}%
\AgdaSymbol{(}\AgdaFunction{embedValue}\AgdaSpace{}%
\AgdaBound{M₂}\AgdaSpace{}%
\AgdaBound{x}\AgdaSymbol{))}\<%
\\
%
\>[2]\AgdaSymbol{(}\AgdaFunction{embedTerm}\AgdaSpace{}%
\AgdaSymbol{(}\AgdaBound{M₁}\AgdaSpace{}%
\AgdaOperator{\AgdaFunction{[}}\AgdaSpace{}%
\AgdaBound{M₂}\AgdaSpace{}%
\AgdaOperator{\AgdaFunction{]}}\AgdaSymbol{))}\<%
\\
%
\>[2]\AgdaSymbol{(}\AgdaFunction{Eq.subst}\<%
\\
\>[2][@{}l@{\AgdaIndent{0}}]%
\>[4]\AgdaSymbol{(λ}\AgdaSpace{}%
\AgdaBound{x}\AgdaSpace{}%
\AgdaSymbol{→}\<%
\\
\>[4][@{}l@{\AgdaIndent{0}}]%
\>[6]\AgdaInductiveConstructor{ƛb}\AgdaSpace{}%
\AgdaSymbol{(}\AgdaInductiveConstructor{ε}\AgdaSpace{}%
\AgdaInductiveConstructor{Zb}\AgdaSpace{}%
\AgdaOperator{\AgdaInductiveConstructor{·b}}\AgdaSpace{}%
\AgdaFunction{renameb}\AgdaSpace{}%
\AgdaInductiveConstructor{Sb}\AgdaSpace{}%
\AgdaSymbol{(}\AgdaFunction{embedTerm}\AgdaSpace{}%
\AgdaBound{M₂}\AgdaSymbol{))}\AgdaSpace{}%
\AgdaOperator{\AgdaInductiveConstructor{·b}}\AgdaSpace{}%
\AgdaInductiveConstructor{box}\AgdaSpace{}%
\AgdaSymbol{(}\AgdaInductiveConstructor{ƛb}\AgdaSpace{}%
\AgdaSymbol{(}\AgdaFunction{embedTerm}\AgdaSpace{}%
\AgdaBound{M₁}\AgdaSymbol{))}\<%
\\
%
\>[6]\AgdaOperator{\AgdaDatatype{↝βb}}\<%
\\
%
\>[6]\AgdaBound{x}\AgdaSymbol{)}\<%
\\
%
\>[4]\AgdaSymbol{(}\AgdaFunction{cong₂}\AgdaSpace{}%
\AgdaOperator{\AgdaInductiveConstructor{\AgdaUnderscore{}·b\AgdaUnderscore{}}}\AgdaSpace{}%
\AgdaInductiveConstructor{refl}\AgdaSpace{}%
\AgdaSymbol{(}\AgdaFunction{subst-rename-fb-comp}\AgdaSpace{}%
\AgdaBound{M₁}\AgdaSpace{}%
\AgdaBound{M₂}\AgdaSpace{}%
\AgdaBound{x}\AgdaSymbol{))}\<%
\\
%
\>[4]\AgdaInductiveConstructor{βb}\AgdaSymbol{)}\<%
\\
%
\>[2]\AgdaSymbol{(}\AgdaFunction{Eq.subst}\<%
\\
\>[2][@{}l@{\AgdaIndent{0}}]%
\>[4]\AgdaSymbol{(λ}\AgdaSpace{}%
\AgdaBound{t}\AgdaSpace{}%
\AgdaSymbol{→}\AgdaSpace{}%
\AgdaInductiveConstructor{ƛb}\AgdaSpace{}%
\AgdaSymbol{(}\AgdaFunction{embedTerm}\AgdaSpace{}%
\AgdaBound{M₁}\AgdaSymbol{)}\AgdaSpace{}%
\AgdaOperator{\AgdaInductiveConstructor{·b}}\AgdaSpace{}%
\AgdaInductiveConstructor{box}\AgdaSpace{}%
\AgdaSymbol{(}\AgdaFunction{embedValue}\AgdaSpace{}%
\AgdaBound{M₂}\AgdaSpace{}%
\AgdaBound{x}\AgdaSymbol{)}\AgdaSpace{}%
\AgdaOperator{\AgdaDatatype{↝βb}}\AgdaSpace{}%
\AgdaBound{t}\AgdaSymbol{)}\<%
\\
%
\>[4]\AgdaSymbol{(}\AgdaFunction{pres-subst}\AgdaSpace{}%
\AgdaBound{M₁}\AgdaSpace{}%
\AgdaBound{M₂}\AgdaSpace{}%
\AgdaBound{x}\AgdaSymbol{)}\<%
\\
%
\>[4]\AgdaInductiveConstructor{βb}\AgdaSymbol{)}\<%
\\
\>[0]\AgdaFunction{pres-red}\AgdaSpace{}%
\AgdaBound{M}\AgdaSpace{}%
\AgdaBound{N}\AgdaSpace{}%
\AgdaSymbol{(}\AgdaInductiveConstructor{μ}\AgdaSpace{}%
\AgdaSymbol{\{}\AgdaArgument{M}\AgdaSpace{}%
\AgdaSymbol{=}\AgdaSpace{}%
\AgdaBound{M₁}\AgdaSymbol{\}}\AgdaSpace{}%
\AgdaSymbol{\{}\AgdaArgument{M'}\AgdaSpace{}%
\AgdaSymbol{=}\AgdaSpace{}%
\AgdaBound{M'}\AgdaSymbol{\}}\AgdaSpace{}%
\AgdaBound{red}\AgdaSymbol{)}\AgdaSpace{}%
\AgdaSymbol{=}\<%
\\
\>[0][@{}l@{\AgdaIndent{0}}]%
\>[2]\AgdaInductiveConstructor{ν}\AgdaSpace{}%
\AgdaSymbol{(}\AgdaFunction{pres-red}\AgdaSpace{}%
\AgdaBound{M₁}\AgdaSpace{}%
\AgdaBound{M'}\AgdaSpace{}%
\AgdaBound{red}\AgdaSymbol{)}\<%
\\
\>[0]\AgdaFunction{pres-red}\AgdaSpace{}%
\AgdaBound{M}\AgdaSpace{}%
\AgdaBound{N}\AgdaSpace{}%
\AgdaSymbol{(}\AgdaInductiveConstructor{ν<}\AgdaSpace{}%
\AgdaSymbol{\{}\AgdaArgument{M}\AgdaSpace{}%
\AgdaSymbol{=}\AgdaSpace{}%
\AgdaBound{M₁}\AgdaSymbol{\}}\AgdaSpace{}%
\AgdaSymbol{\{}\AgdaArgument{N}\AgdaSpace{}%
\AgdaSymbol{=}\AgdaSpace{}%
\AgdaBound{N₁}\AgdaSymbol{\}}\AgdaSpace{}%
\AgdaSymbol{\{}\AgdaArgument{N'}\AgdaSpace{}%
\AgdaSymbol{=}\AgdaSpace{}%
\AgdaBound{N'}\AgdaSymbol{\}}\AgdaSpace{}%
\AgdaBound{x}\AgdaSpace{}%
\AgdaBound{red}\AgdaSymbol{)}\AgdaSpace{}%
\AgdaSymbol{=}\<%
\\
\>[0][@{}l@{\AgdaIndent{0}}]%
\>[2]\AgdaInductiveConstructor{μ}\AgdaSpace{}%
\AgdaSymbol{(}\AgdaInductiveConstructor{ξ}\AgdaSpace{}%
\AgdaSymbol{(}\AgdaInductiveConstructor{ν}%
\>[2883I]\AgdaSymbol{(}\AgdaFunction{pres-rename-red-2}\<%
\\
\>[2883I][@{}l@{\AgdaIndent{0}}]%
\>[12]\AgdaInductiveConstructor{Sb}\<%
\\
%
\>[12]\AgdaSymbol{(}\AgdaFunction{embedTerm}\AgdaSpace{}%
\AgdaBound{N₁}\AgdaSymbol{)}\<%
\\
%
\>[12]\AgdaSymbol{(}\AgdaFunction{embedTerm}\AgdaSpace{}%
\AgdaBound{N'}\AgdaSymbol{)}\<%
\\
%
\>[12]\AgdaSymbol{(}\AgdaFunction{pres-red}\AgdaSpace{}%
\AgdaBound{N₁}\AgdaSpace{}%
\AgdaBound{N'}\AgdaSpace{}%
\AgdaBound{red}\AgdaSymbol{))))}\<%
\\
\>[0]\AgdaFunction{pres-red}\AgdaSpace{}%
\AgdaBound{M}\AgdaSpace{}%
\AgdaBound{N}\AgdaSpace{}%
\AgdaSymbol{(}\AgdaInductiveConstructor{ν}\AgdaSpace{}%
\AgdaSymbol{\{}\AgdaArgument{M}\AgdaSpace{}%
\AgdaSymbol{=}\AgdaSpace{}%
\AgdaBound{M₁}\AgdaSymbol{\}}\AgdaSpace{}%
\AgdaSymbol{\{}\AgdaArgument{N}\AgdaSpace{}%
\AgdaSymbol{=}\AgdaSpace{}%
\AgdaBound{N₁}\AgdaSymbol{\}}\AgdaSpace{}%
\AgdaSymbol{\{}\AgdaArgument{N'}\AgdaSpace{}%
\AgdaSymbol{=}\AgdaSpace{}%
\AgdaBound{N'}\AgdaSymbol{\}}\AgdaSpace{}%
\AgdaBound{red}\AgdaSymbol{)}\AgdaSpace{}%
\AgdaSymbol{=}\<%
\\
\>[0][@{}l@{\AgdaIndent{0}}]%
\>[2]\AgdaInductiveConstructor{μ}\AgdaSpace{}%
\AgdaSymbol{(}\AgdaInductiveConstructor{ξ}\AgdaSpace{}%
\AgdaSymbol{(}\AgdaInductiveConstructor{ν}%
\>[2905I]\AgdaSymbol{(}\AgdaFunction{pres-rename-red-2}\<%
\\
\>[2905I][@{}l@{\AgdaIndent{0}}]%
\>[12]\AgdaInductiveConstructor{Sb}\<%
\\
%
\>[12]\AgdaSymbol{(}\AgdaFunction{embedTerm}\AgdaSpace{}%
\AgdaBound{N₁}\AgdaSymbol{)}\<%
\\
%
\>[12]\AgdaSymbol{(}\AgdaFunction{embedTerm}\AgdaSpace{}%
\AgdaBound{N'}\AgdaSymbol{)}\<%
\\
%
\>[12]\AgdaSymbol{(}\AgdaFunction{pres-red}\AgdaSpace{}%
\AgdaBound{N₁}\AgdaSpace{}%
\AgdaBound{N'}\AgdaSpace{}%
\AgdaBound{red}\AgdaSymbol{))))}\<%
\\
\>[0]\AgdaFunction{pres-red}\AgdaSpace{}%
\AgdaBound{M}\AgdaSpace{}%
\AgdaBound{N}\AgdaSpace{}%
\AgdaSymbol{(}\AgdaInductiveConstructor{ξ}\AgdaSpace{}%
\AgdaSymbol{\{}\AgdaArgument{M}\AgdaSpace{}%
\AgdaSymbol{=}\AgdaSpace{}%
\AgdaBound{M₁}\AgdaSymbol{\}}\AgdaSpace{}%
\AgdaSymbol{\{}\AgdaArgument{M'}\AgdaSpace{}%
\AgdaSymbol{=}\AgdaSpace{}%
\AgdaBound{M'}\AgdaSymbol{\}}\AgdaSpace{}%
\AgdaBound{red}\AgdaSymbol{)}\AgdaSpace{}%
\AgdaSymbol{=}\<%
\\
\>[0][@{}l@{\AgdaIndent{0}}]%
\>[2]\AgdaInductiveConstructor{ζ}\AgdaSpace{}%
\AgdaSymbol{(}\AgdaInductiveConstructor{ξ}\AgdaSpace{}%
\AgdaSymbol{(}\AgdaFunction{pres-red}\AgdaSpace{}%
\AgdaBound{M₁}\AgdaSpace{}%
\AgdaBound{M'}\AgdaSpace{}%
\AgdaBound{red}\AgdaSymbol{))}\<%
\\
%
\\[\AgdaEmptyExtraSkip]%
\>[0]\AgdaComment{--\ Preservation\ of\ the\ 2\ step\ reduction\ rule\ under\ rename}\<%
\\
\>[0]\AgdaFunction{pres-rename-red-↝βb²}\AgdaSpace{}%
\AgdaSymbol{:}\AgdaSpace{}%
\AgdaSymbol{∀}\AgdaSpace{}%
\AgdaSymbol{\{}\AgdaBound{Γ}\AgdaSpace{}%
\AgdaBound{A}\AgdaSpace{}%
\AgdaBound{Δ}\AgdaSymbol{\}}\<%
\\
\>[0][@{}l@{\AgdaIndent{0}}]%
\>[2]\AgdaSymbol{→}\AgdaSpace{}%
\AgdaSymbol{(}\AgdaBound{fb'}\AgdaSpace{}%
\AgdaSymbol{:}\AgdaSpace{}%
\AgdaSymbol{∀}\AgdaSpace{}%
\AgdaSymbol{\{}\AgdaBound{Ab}\AgdaSymbol{\}}\AgdaSpace{}%
\AgdaSymbol{→}\AgdaSpace{}%
\AgdaBound{Ab}\AgdaSpace{}%
\AgdaOperator{\AgdaDatatype{∈b}}\AgdaSpace{}%
\AgdaBound{Γ}\AgdaSpace{}%
\AgdaSymbol{→}\AgdaSpace{}%
\AgdaBound{Ab}\AgdaSpace{}%
\AgdaOperator{\AgdaDatatype{∈b}}\AgdaSpace{}%
\AgdaBound{Δ}\AgdaSymbol{)}\<%
\\
%
\>[2]\AgdaSymbol{→}\AgdaSpace{}%
\AgdaSymbol{(}\AgdaBound{M}\AgdaSpace{}%
\AgdaSymbol{:}\AgdaSpace{}%
\AgdaBound{Γ}\AgdaSpace{}%
\AgdaOperator{\AgdaDatatype{⊢b}}\AgdaSpace{}%
\AgdaBound{A}\AgdaSymbol{)}\<%
\\
%
\>[2]\AgdaSymbol{→}\AgdaSpace{}%
\AgdaSymbol{(}\AgdaBound{N}\AgdaSpace{}%
\AgdaSymbol{:}\AgdaSpace{}%
\AgdaBound{Γ}\AgdaSpace{}%
\AgdaOperator{\AgdaDatatype{⊢b}}\AgdaSpace{}%
\AgdaBound{A}\AgdaSymbol{)}\<%
\\
%
\>[2]\AgdaSymbol{→}\AgdaSpace{}%
\AgdaSymbol{(}\AgdaBound{M}\AgdaSpace{}%
\AgdaOperator{\AgdaDatatype{↝βb₂}}\AgdaSpace{}%
\AgdaBound{N}\AgdaSymbol{)}\<%
\\
%
\>[2]\AgdaSymbol{→}\AgdaSpace{}%
\AgdaFunction{renameb}\AgdaSpace{}%
\AgdaBound{fb'}\AgdaSpace{}%
\AgdaBound{M}\AgdaSpace{}%
\AgdaOperator{\AgdaDatatype{↝βb₂}}\AgdaSpace{}%
\AgdaFunction{renameb}\AgdaSpace{}%
\AgdaBound{fb'}\AgdaSpace{}%
\AgdaBound{N}\<%
\\
\>[0]\AgdaFunction{pres-rename-red-↝βb²}\AgdaSpace{}%
\AgdaBound{fb'}\AgdaSpace{}%
\AgdaBound{M}\AgdaSpace{}%
\AgdaBound{N}\AgdaSpace{}%
\AgdaSymbol{(}\AgdaInductiveConstructor{βb₂}\AgdaSpace{}%
\AgdaSymbol{\{}\AgdaArgument{M}\AgdaSpace{}%
\AgdaSymbol{=}\AgdaSpace{}%
\AgdaBound{M'}\AgdaSymbol{\}}\AgdaSpace{}%
\AgdaSymbol{\{}\AgdaArgument{W}\AgdaSpace{}%
\AgdaSymbol{=}\AgdaSpace{}%
\AgdaBound{W'}\AgdaSymbol{\})}\AgdaSpace{}%
\AgdaSymbol{=}\<%
\\
\>[0][@{}l@{\AgdaIndent{0}}]%
\>[2]\AgdaFunction{Eq.subst}\<%
\\
\>[2][@{}l@{\AgdaIndent{0}}]%
\>[4]\AgdaSymbol{(λ}\AgdaSpace{}%
\AgdaBound{x}\AgdaSpace{}%
\AgdaSymbol{→}\AgdaSpace{}%
\AgdaBound{x}\AgdaSpace{}%
\AgdaOperator{\AgdaDatatype{↝βb₂}}\AgdaSpace{}%
\AgdaFunction{renameb}\AgdaSpace{}%
\AgdaBound{fb'}\<%
\\
%
\>[4]\AgdaSymbol{(}\AgdaBound{M'}\AgdaSpace{}%
\AgdaOperator{\AgdaFunction{[}}\AgdaSpace{}%
\AgdaBound{W'}\AgdaSpace{}%
\AgdaOperator{\AgdaFunction{]b}}\AgdaSymbol{))}\<%
\\
%
\>[4]\AgdaSymbol{(}\AgdaFunction{cong₂}\AgdaSpace{}%
\AgdaOperator{\AgdaInductiveConstructor{\AgdaUnderscore{}·b\AgdaUnderscore{}}}\AgdaSpace{}%
\AgdaSymbol{(}\AgdaFunction{cong}\AgdaSpace{}%
\AgdaInductiveConstructor{ƛb}\AgdaSpace{}%
\AgdaSymbol{(}\AgdaFunction{cong₂}\AgdaSpace{}%
\AgdaOperator{\AgdaInductiveConstructor{\AgdaUnderscore{}·b\AgdaUnderscore{}}}\AgdaSpace{}%
\AgdaInductiveConstructor{refl}\AgdaSpace{}%
\AgdaSymbol{(}\AgdaFunction{cong}\AgdaSpace{}%
\AgdaInductiveConstructor{box}\AgdaSpace{}%
\AgdaSymbol{(}\AgdaFunction{sym}\AgdaSpace{}%
\AgdaSymbol{(}\AgdaFunction{rename-extb-Sb-commute}\AgdaSpace{}%
\AgdaBound{fb'}\AgdaSpace{}%
\AgdaBound{W'}\AgdaSymbol{)))))}\AgdaSpace{}%
\AgdaInductiveConstructor{refl}\AgdaSymbol{)}\<%
\\
%
\>[4]\AgdaSymbol{(}\AgdaFunction{Eq.subst}\<%
\\
\>[4][@{}l@{\AgdaIndent{0}}]%
\>[6]\AgdaSymbol{(λ}\AgdaSpace{}%
\AgdaBound{x}\AgdaSpace{}%
\AgdaSymbol{→}\AgdaSpace{}%
\AgdaSymbol{(}\AgdaFunction{raise}%
\>[21]\AgdaSymbol{(}\AgdaInductiveConstructor{box}\AgdaSpace{}%
\AgdaSymbol{(}\AgdaFunction{renameb}\AgdaSpace{}%
\AgdaBound{fb'}\AgdaSpace{}%
\AgdaBound{W'}\AgdaSymbol{)))}\AgdaSpace{}%
\AgdaOperator{\AgdaInductiveConstructor{·b}}\AgdaSpace{}%
\AgdaInductiveConstructor{box}\AgdaSpace{}%
\AgdaSymbol{(}\AgdaInductiveConstructor{ƛb}\AgdaSpace{}%
\AgdaSymbol{(}\AgdaFunction{renameb}\AgdaSpace{}%
\AgdaSymbol{(}\AgdaFunction{extb}\AgdaSpace{}%
\AgdaBound{fb'}\AgdaSymbol{)}\AgdaSpace{}%
\AgdaBound{M'}\AgdaSymbol{))}\AgdaSpace{}%
\AgdaOperator{\AgdaDatatype{↝βb₂}}\AgdaSpace{}%
\AgdaBound{x}\AgdaSymbol{)}\<%
\\
%
\>[6]\AgdaSymbol{(}\AgdaFunction{pres-subst-func-rename}\AgdaSpace{}%
\AgdaBound{fb'}\AgdaSpace{}%
\AgdaBound{M'}\AgdaSpace{}%
\AgdaBound{W'}\AgdaSymbol{)}\<%
\\
%
\>[6]\AgdaInductiveConstructor{βb₂}\<%
\\
%
\>[4]\AgdaSymbol{)}\<%
\\
\>[0]\AgdaFunction{pres-rename-red-↝βb²}\AgdaSpace{}%
\AgdaBound{fb'}\AgdaSpace{}%
\AgdaBound{M}\AgdaSpace{}%
\AgdaBound{N}\AgdaSpace{}%
\AgdaSymbol{(}\AgdaInductiveConstructor{μ}\AgdaSpace{}%
\AgdaSymbol{\{}\AgdaArgument{M}\AgdaSpace{}%
\AgdaSymbol{=}\AgdaSpace{}%
\AgdaBound{M₁}\AgdaSymbol{\}}\AgdaSpace{}%
\AgdaSymbol{\{}\AgdaArgument{M'}\AgdaSpace{}%
\AgdaSymbol{=}\AgdaSpace{}%
\AgdaBound{M'}\AgdaSymbol{\}}\AgdaSpace{}%
\AgdaBound{red}\AgdaSymbol{)}\AgdaSpace{}%
\AgdaSymbol{=}\AgdaSpace{}%
\AgdaInductiveConstructor{μ}\AgdaSpace{}%
\AgdaSymbol{(}\AgdaFunction{pres-rename-red-↝βb²}\AgdaSpace{}%
\AgdaBound{fb'}\AgdaSpace{}%
\AgdaBound{M₁}\AgdaSpace{}%
\AgdaBound{M'}\AgdaSpace{}%
\AgdaBound{red}\AgdaSymbol{)}\<%
\\
\>[0]\AgdaFunction{pres-rename-red-↝βb²}\AgdaSpace{}%
\AgdaBound{fb'}\AgdaSpace{}%
\AgdaBound{M}\AgdaSpace{}%
\AgdaBound{N}\AgdaSpace{}%
\AgdaSymbol{(}\AgdaInductiveConstructor{ν}\AgdaSpace{}%
\AgdaSymbol{\{}\AgdaArgument{N}\AgdaSpace{}%
\AgdaSymbol{=}\AgdaSpace{}%
\AgdaBound{N₁}\AgdaSymbol{\}}\AgdaSpace{}%
\AgdaSymbol{\{}\AgdaArgument{N'}\AgdaSpace{}%
\AgdaSymbol{=}\AgdaSpace{}%
\AgdaBound{N'}\AgdaSymbol{\}}\AgdaSpace{}%
\AgdaBound{red}\AgdaSymbol{)}\AgdaSpace{}%
\AgdaSymbol{=}\AgdaSpace{}%
\AgdaInductiveConstructor{ν}\AgdaSpace{}%
\AgdaSymbol{(}\AgdaFunction{pres-rename-red-↝βb²}\AgdaSpace{}%
\AgdaBound{fb'}\AgdaSpace{}%
\AgdaBound{N₁}\AgdaSpace{}%
\AgdaBound{N'}\AgdaSpace{}%
\AgdaBound{red}\AgdaSymbol{)}\<%
\\
\>[0]\AgdaFunction{pres-rename-red-↝βb²}\AgdaSpace{}%
\AgdaBound{fb'}\AgdaSpace{}%
\AgdaBound{M}\AgdaSpace{}%
\AgdaBound{N}\AgdaSpace{}%
\AgdaSymbol{(}\AgdaInductiveConstructor{ξ}\AgdaSpace{}%
\AgdaSymbol{\{}\AgdaArgument{M}\AgdaSpace{}%
\AgdaSymbol{=}\AgdaSpace{}%
\AgdaBound{M₁}\AgdaSymbol{\}}\AgdaSpace{}%
\AgdaSymbol{\{}\AgdaArgument{M'}\AgdaSpace{}%
\AgdaSymbol{=}\AgdaSpace{}%
\AgdaBound{M'}\AgdaSymbol{\}}\AgdaSpace{}%
\AgdaBound{red}\AgdaSymbol{)}\AgdaSpace{}%
\AgdaSymbol{=}\AgdaSpace{}%
\AgdaInductiveConstructor{ξ}\AgdaSpace{}%
\AgdaSymbol{(}\AgdaFunction{pres-rename-red-↝βb²}\AgdaSpace{}%
\AgdaSymbol{(}\AgdaFunction{extb}\AgdaSpace{}%
\AgdaBound{fb'}\AgdaSymbol{)}\AgdaSpace{}%
\AgdaBound{M₁}\AgdaSpace{}%
\AgdaBound{M'}\AgdaSpace{}%
\AgdaBound{red}\AgdaSymbol{)}\<%
\\
\>[0]\AgdaFunction{pres-rename-red-↝βb²}\AgdaSpace{}%
\AgdaBound{fb'}\AgdaSpace{}%
\AgdaBound{M}\AgdaSpace{}%
\AgdaBound{N}\AgdaSpace{}%
\AgdaSymbol{(}\AgdaInductiveConstructor{ζ}\AgdaSpace{}%
\AgdaSymbol{\{}\AgdaArgument{M}\AgdaSpace{}%
\AgdaSymbol{=}\AgdaSpace{}%
\AgdaBound{M₁}\AgdaSymbol{\}}\AgdaSpace{}%
\AgdaSymbol{\{}\AgdaArgument{M'}\AgdaSpace{}%
\AgdaSymbol{=}\AgdaSpace{}%
\AgdaBound{M'}\AgdaSymbol{\}}\AgdaSpace{}%
\AgdaBound{red}\AgdaSymbol{)}\AgdaSpace{}%
\AgdaSymbol{=}\AgdaSpace{}%
\AgdaInductiveConstructor{ζ}\AgdaSpace{}%
\AgdaSymbol{(}\AgdaFunction{pres-rename-red-↝βb²}\AgdaSpace{}%
\AgdaBound{fb'}\AgdaSpace{}%
\AgdaBound{M₁}\AgdaSpace{}%
\AgdaBound{M'}\AgdaSpace{}%
\AgdaBound{red}\AgdaSymbol{)}\<%
\\
%
\\[\AgdaEmptyExtraSkip]%
\>[0]\AgdaComment{--\ Preservation\ of\ reduction\ by\ Gödel's\ translation}\<%
\\
\>[0]\AgdaComment{--\ with\ the\ two\ step\ reduction\ rule}\<%
\\
\>[0]\AgdaFunction{pres-red↝βb²}\AgdaSpace{}%
\AgdaSymbol{:}\AgdaSpace{}%
\AgdaSymbol{∀}\AgdaSpace{}%
\AgdaSymbol{\{}\AgdaBound{Γ}\AgdaSpace{}%
\AgdaBound{A}\AgdaSymbol{\}}\<%
\\
\>[0][@{}l@{\AgdaIndent{0}}]%
\>[2]\AgdaSymbol{→}\AgdaSpace{}%
\AgdaSymbol{(}\AgdaBound{M}\AgdaSpace{}%
\AgdaSymbol{:}\AgdaSpace{}%
\AgdaBound{Γ}\AgdaSpace{}%
\AgdaOperator{\AgdaDatatype{⊢}}\AgdaSpace{}%
\AgdaBound{A}\AgdaSymbol{)}\<%
\\
%
\>[2]\AgdaSymbol{→}\AgdaSpace{}%
\AgdaSymbol{(}\AgdaBound{N}\AgdaSpace{}%
\AgdaSymbol{:}\AgdaSpace{}%
\AgdaBound{Γ}\AgdaSpace{}%
\AgdaOperator{\AgdaDatatype{⊢}}\AgdaSpace{}%
\AgdaBound{A}\AgdaSymbol{)}\<%
\\
%
\>[2]\AgdaSymbol{→}\AgdaSpace{}%
\AgdaSymbol{(}\AgdaBound{M}\AgdaSpace{}%
\AgdaOperator{\AgdaDatatype{↝βᵥ}}\AgdaSpace{}%
\AgdaBound{N}\AgdaSymbol{)}\<%
\\
%
\>[2]\AgdaSymbol{→}\AgdaSpace{}%
\AgdaSymbol{(}\AgdaFunction{embedTerm}\AgdaSpace{}%
\AgdaBound{M}\AgdaSpace{}%
\AgdaOperator{\AgdaDatatype{↝βb₂}}\AgdaSpace{}%
\AgdaFunction{embedTerm}\AgdaSpace{}%
\AgdaBound{N}\AgdaSymbol{)}\<%
\\
\>[0]\AgdaFunction{pres-red↝βb²}\AgdaSpace{}%
\AgdaSymbol{(}\AgdaInductiveConstructor{ƛ}\AgdaSpace{}%
\AgdaBound{M₁}\AgdaSpace{}%
\AgdaOperator{\AgdaInductiveConstructor{·}}\AgdaSpace{}%
\AgdaBound{M₂}\AgdaSymbol{)}\AgdaSpace{}%
\AgdaBound{N}\AgdaSpace{}%
\AgdaSymbol{(}\AgdaInductiveConstructor{βᵥ}\AgdaSpace{}%
\AgdaBound{x}\AgdaSymbol{)}\AgdaSpace{}%
\AgdaSymbol{=}\<%
\\
\>[0][@{}l@{\AgdaIndent{0}}]%
\>[2]\AgdaFunction{Eq.subst}\<%
\\
\>[2][@{}l@{\AgdaIndent{0}}]%
\>[4]\AgdaSymbol{(λ}\AgdaSpace{}%
\AgdaBound{t}\AgdaSpace{}%
\AgdaSymbol{→}\AgdaSpace{}%
\AgdaBound{t}\AgdaSpace{}%
\AgdaOperator{\AgdaDatatype{↝βb₂}}\AgdaSpace{}%
\AgdaFunction{embedTerm}\AgdaSpace{}%
\AgdaSymbol{(}\AgdaBound{M₁}\AgdaSpace{}%
\AgdaOperator{\AgdaFunction{[}}\AgdaSpace{}%
\AgdaBound{M₂}\AgdaSpace{}%
\AgdaOperator{\AgdaFunction{]}}\AgdaSymbol{))}\<%
\\
%
\>[4]\AgdaSymbol{(}\AgdaFunction{cong₂}\AgdaSpace{}%
\AgdaOperator{\AgdaInductiveConstructor{\AgdaUnderscore{}·b\AgdaUnderscore{}}}\AgdaSpace{}%
\AgdaSymbol{(}\AgdaFunction{cong}\AgdaSpace{}%
\AgdaInductiveConstructor{ƛb}\AgdaSpace{}%
\AgdaSymbol{(}\AgdaFunction{cong₂}\AgdaSpace{}%
\AgdaOperator{\AgdaInductiveConstructor{\AgdaUnderscore{}·b\AgdaUnderscore{}}}\AgdaSpace{}%
\AgdaInductiveConstructor{refl}\AgdaSpace{}%
\AgdaSymbol{(}\AgdaFunction{cong}\AgdaSpace{}%
\AgdaSymbol{(}\AgdaFunction{renameb}\AgdaSpace{}%
\AgdaInductiveConstructor{Sb}\AgdaSymbol{)}\AgdaSpace{}%
\AgdaSymbol{(}\AgdaFunction{box-embedValue}\AgdaSpace{}%
\AgdaBound{x}\AgdaSymbol{))))}\AgdaSpace{}%
\AgdaInductiveConstructor{refl}\AgdaSymbol{)}\<%
\\
%
\>[4]\AgdaSymbol{(}\AgdaFunction{Eq.subst}\<%
\\
\>[4][@{}l@{\AgdaIndent{0}}]%
\>[6]\AgdaSymbol{(λ}\AgdaSpace{}%
\AgdaBound{t}\AgdaSpace{}%
\AgdaSymbol{→}\AgdaSpace{}%
\AgdaFunction{raise}\AgdaSpace{}%
\AgdaSymbol{(}\AgdaSpace{}%
\AgdaInductiveConstructor{box}\AgdaSpace{}%
\AgdaSymbol{(}\AgdaFunction{embedValue}\AgdaSpace{}%
\AgdaBound{M₂}\AgdaSpace{}%
\AgdaBound{x}\AgdaSymbol{)}\AgdaSpace{}%
\AgdaSymbol{)}\AgdaSpace{}%
\AgdaOperator{\AgdaInductiveConstructor{·b}}\AgdaSpace{}%
\AgdaInductiveConstructor{box}\AgdaSpace{}%
\AgdaSymbol{(}\AgdaInductiveConstructor{ƛb}\AgdaSpace{}%
\AgdaSymbol{(}\AgdaFunction{embedTerm}\AgdaSpace{}%
\AgdaBound{M₁}\AgdaSymbol{))}\AgdaSpace{}%
\AgdaOperator{\AgdaDatatype{↝βb₂}}\AgdaSpace{}%
\AgdaBound{t}\AgdaSymbol{)}\<%
\\
%
\>[6]\AgdaSymbol{(}\AgdaFunction{pres-subst}\AgdaSpace{}%
\AgdaBound{M₁}\AgdaSpace{}%
\AgdaBound{M₂}\AgdaSpace{}%
\AgdaBound{x}\AgdaSymbol{)}\<%
\\
%
\>[6]\AgdaInductiveConstructor{βb₂}\<%
\\
%
\>[4]\AgdaSymbol{)}\<%
\\
\>[0]\AgdaFunction{pres-red↝βb²}\AgdaSpace{}%
\AgdaBound{M}\AgdaSpace{}%
\AgdaBound{N}\AgdaSpace{}%
\AgdaSymbol{(}\AgdaInductiveConstructor{μ}\AgdaSpace{}%
\AgdaSymbol{\{}\AgdaArgument{M}\AgdaSpace{}%
\AgdaSymbol{=}\AgdaSpace{}%
\AgdaBound{M₁}\AgdaSymbol{\}}\AgdaSpace{}%
\AgdaSymbol{\{}\AgdaArgument{M'}\AgdaSpace{}%
\AgdaSymbol{=}\AgdaSpace{}%
\AgdaBound{M'}\AgdaSymbol{\}}\AgdaSpace{}%
\AgdaBound{red}\AgdaSymbol{)}\AgdaSpace{}%
\AgdaSymbol{=}\AgdaSpace{}%
\AgdaInductiveConstructor{ν}\AgdaSpace{}%
\AgdaSymbol{(}\AgdaFunction{pres-red↝βb²}\AgdaSpace{}%
\AgdaBound{M₁}\AgdaSpace{}%
\AgdaBound{M'}\AgdaSpace{}%
\AgdaBound{red}\AgdaSymbol{)}\<%
\\
\>[0]\AgdaFunction{pres-red↝βb²}\AgdaSpace{}%
\AgdaBound{M}\AgdaSpace{}%
\AgdaBound{N}\AgdaSpace{}%
\AgdaSymbol{(}\AgdaInductiveConstructor{ν<}\AgdaSpace{}%
\AgdaSymbol{\{}\AgdaArgument{N}\AgdaSpace{}%
\AgdaSymbol{=}\AgdaSpace{}%
\AgdaBound{N₁}\AgdaSymbol{\}}\AgdaSpace{}%
\AgdaSymbol{\{}\AgdaArgument{N'}\AgdaSpace{}%
\AgdaSymbol{=}\AgdaSpace{}%
\AgdaBound{N'}\AgdaSymbol{\}}\AgdaSpace{}%
\AgdaBound{x}\AgdaSpace{}%
\AgdaBound{red}\AgdaSymbol{)}\AgdaSpace{}%
\AgdaSymbol{=}\AgdaSpace{}%
\AgdaInductiveConstructor{μ}\AgdaSpace{}%
\AgdaSymbol{(}\AgdaInductiveConstructor{ξ}\AgdaSpace{}%
\AgdaSymbol{(}\AgdaInductiveConstructor{ν}\AgdaSpace{}%
\AgdaSymbol{(}\AgdaFunction{pres-rename-red-↝βb²}\AgdaSpace{}%
\AgdaInductiveConstructor{Sb}\AgdaSpace{}%
\AgdaSymbol{(}\AgdaFunction{embedTerm}\AgdaSpace{}%
\AgdaBound{N₁}\AgdaSymbol{)}\AgdaSpace{}%
\AgdaSymbol{(}\AgdaFunction{embedTerm}\AgdaSpace{}%
\AgdaBound{N'}\AgdaSymbol{)}\AgdaSpace{}%
\AgdaSymbol{(}\AgdaFunction{pres-red↝βb²}\AgdaSpace{}%
\AgdaBound{N₁}\AgdaSpace{}%
\AgdaBound{N'}\AgdaSpace{}%
\AgdaBound{red}\AgdaSymbol{))))}\<%
\\
\>[0]\AgdaFunction{pres-red↝βb²}\AgdaSpace{}%
\AgdaBound{M}\AgdaSpace{}%
\AgdaBound{N}\AgdaSpace{}%
\AgdaSymbol{(}\AgdaInductiveConstructor{ν}\AgdaSpace{}%
\AgdaSymbol{\{}\AgdaArgument{N}\AgdaSpace{}%
\AgdaSymbol{=}\AgdaSpace{}%
\AgdaBound{N₁}\AgdaSymbol{\}}\AgdaSpace{}%
\AgdaSymbol{\{}\AgdaArgument{N'}\AgdaSpace{}%
\AgdaSymbol{=}\AgdaSpace{}%
\AgdaBound{N'}\AgdaSymbol{\}}\AgdaSpace{}%
\AgdaBound{red}\AgdaSymbol{)}\AgdaSpace{}%
\AgdaSymbol{=}\AgdaSpace{}%
\AgdaInductiveConstructor{μ}\AgdaSpace{}%
\AgdaSymbol{(}\AgdaInductiveConstructor{ξ}\AgdaSpace{}%
\AgdaSymbol{(}\AgdaInductiveConstructor{ν}\AgdaSpace{}%
\AgdaSymbol{(}\AgdaFunction{pres-rename-red-↝βb²}\AgdaSpace{}%
\AgdaInductiveConstructor{Sb}\AgdaSpace{}%
\AgdaSymbol{(}\AgdaFunction{embedTerm}\AgdaSpace{}%
\AgdaBound{N₁}\AgdaSymbol{)}\AgdaSpace{}%
\AgdaSymbol{(}\AgdaFunction{embedTerm}\AgdaSpace{}%
\AgdaBound{N'}\AgdaSymbol{)}\AgdaSpace{}%
\AgdaSymbol{(}\AgdaFunction{pres-red↝βb²}\AgdaSpace{}%
\AgdaBound{N₁}\AgdaSpace{}%
\AgdaBound{N'}\AgdaSpace{}%
\AgdaBound{red}\AgdaSymbol{))))}\<%
\\
\>[0]\AgdaFunction{pres-red↝βb²}\AgdaSpace{}%
\AgdaBound{M}\AgdaSpace{}%
\AgdaBound{N}\AgdaSpace{}%
\AgdaSymbol{(}\AgdaInductiveConstructor{ξ}\AgdaSpace{}%
\AgdaSymbol{\{}\AgdaArgument{M}\AgdaSpace{}%
\AgdaSymbol{=}\AgdaSpace{}%
\AgdaBound{M₁}\AgdaSymbol{\}}\AgdaSpace{}%
\AgdaSymbol{\{}\AgdaArgument{M'}\AgdaSpace{}%
\AgdaSymbol{=}\AgdaSpace{}%
\AgdaBound{M'}\AgdaSymbol{\}}\AgdaSpace{}%
\AgdaBound{red}\AgdaSymbol{)}\AgdaSpace{}%
\AgdaSymbol{=}\AgdaSpace{}%
\AgdaInductiveConstructor{ζ}\AgdaSpace{}%
\AgdaSymbol{(}\AgdaInductiveConstructor{ξ}\AgdaSpace{}%
\AgdaSymbol{(}\AgdaFunction{pres-red↝βb²}\AgdaSpace{}%
\AgdaBound{M₁}\AgdaSpace{}%
\AgdaBound{M'}\AgdaSpace{}%
\AgdaBound{red}\AgdaSymbol{))}\<%
\end{code}  
  \hfill $\qed$
\end{proofthm}

\begin{proofthm}
  \label{proof:godel-preservation-red}
  \textbf{Preservation of reduction} under Gödel's translation with $\redbboxt$ reduction rule.
  \small
  \begin{code}[hide]%
\>[0]\AgdaKeyword{module}\AgdaSpace{}%
\AgdaModule{valueProofs5}\AgdaSpace{}%
\AgdaKeyword{where}\<%
\\
\>[0]\<%
\\
\>[0]\AgdaKeyword{open}\AgdaSpace{}%
\AgdaKeyword{import}\AgdaSpace{}%
\AgdaModule{calc}\<%
\\
\>[0]\AgdaKeyword{open}\AgdaSpace{}%
\AgdaKeyword{import}\AgdaSpace{}%
\AgdaModule{boxCalc}\<%
\\
\>[0]\AgdaKeyword{open}\AgdaSpace{}%
\AgdaKeyword{import}\AgdaSpace{}%
\AgdaModule{embedCBVIntoCBB}\<%
\\
\>[0]\AgdaKeyword{open}\AgdaSpace{}%
\AgdaKeyword{import}\AgdaSpace{}%
\AgdaModule{valueCalcRed}\<%
\\
\>[0]\AgdaKeyword{open}\AgdaSpace{}%
\AgdaKeyword{import}\AgdaSpace{}%
\AgdaModule{boxCalcRed}\<%
\\
%
\\[\AgdaEmptyExtraSkip]%
\>[0]\AgdaKeyword{open}\AgdaSpace{}%
\AgdaKeyword{import}\AgdaSpace{}%
\AgdaModule{Data.Empty}\AgdaSpace{}%
\AgdaKeyword{using}\AgdaSpace{}%
\AgdaSymbol{(}\AgdaFunction{⊥}\AgdaSymbol{;}\AgdaSpace{}%
\AgdaFunction{⊥-elim}\AgdaSymbol{)}\<%
\\
\>[0]\AgdaKeyword{import}\AgdaSpace{}%
\AgdaModule{Relation.Binary.PropositionalEquality}\AgdaSpace{}%
\AgdaSymbol{as}\AgdaSpace{}%
\AgdaModule{Eq}\<%
\\
\>[0]\AgdaKeyword{open}\AgdaSpace{}%
\AgdaModule{Eq}\AgdaSpace{}%
\AgdaKeyword{using}\AgdaSpace{}%
\AgdaSymbol{(}\AgdaOperator{\AgdaDatatype{\AgdaUnderscore{}≡\AgdaUnderscore{}}}\AgdaSymbol{;}\AgdaSpace{}%
\AgdaInductiveConstructor{refl}\AgdaSymbol{;}\AgdaSpace{}%
\AgdaFunction{cong}\AgdaSymbol{;}\AgdaSpace{}%
\AgdaFunction{cong₂}\AgdaSymbol{;}\AgdaSpace{}%
\AgdaFunction{sym}\AgdaSymbol{)}\<%
\\
%
\\[\AgdaEmptyExtraSkip]%
\>[0]\AgdaKeyword{open}\AgdaSpace{}%
\AgdaModule{Eq.≡-Reasoning}\AgdaSpace{}%
\AgdaKeyword{using}\AgdaSpace{}%
\AgdaSymbol{(}\AgdaOperator{\AgdaFunction{begin\AgdaUnderscore{}}}\AgdaSymbol{;}\AgdaSpace{}%
\AgdaFunction{step-≡-∣}\AgdaSymbol{;}\AgdaSpace{}%
\AgdaFunction{step-≡-⟩}\AgdaSymbol{;}\AgdaSpace{}%
\AgdaOperator{\AgdaFunction{\AgdaUnderscore{}∎}}\AgdaSymbol{)}\<%
\\
\>[0]\AgdaKeyword{open}\AgdaSpace{}%
\AgdaModule{Eq.≡-Reasoning}\AgdaSpace{}%
\AgdaKeyword{renaming}\AgdaSpace{}%
\AgdaSymbol{(}\AgdaOperator{\AgdaFunction{begin\AgdaUnderscore{}}}\AgdaSpace{}%
\AgdaSymbol{to}\AgdaSpace{}%
\AgdaOperator{\AgdaFunction{beginEq\AgdaUnderscore{}}}\AgdaSymbol{;}\AgdaSpace{}%
\AgdaOperator{\AgdaFunction{\AgdaUnderscore{}∎}}\AgdaSpace{}%
\AgdaSymbol{to}\AgdaSpace{}%
\AgdaOperator{\AgdaFunction{\AgdaUnderscore{}∎eq}}\AgdaSymbol{)}\<%
\\
%
\\[\AgdaEmptyExtraSkip]%
\>[0]\AgdaComment{--\ Preservation\ of\ typing\ by\ Gödel's\ translation}\<%
\\
\>[0]\AgdaFunction{pres-typing}\AgdaSpace{}%
\AgdaSymbol{:}\AgdaSpace{}%
\AgdaSymbol{∀}\AgdaSpace{}%
\AgdaSymbol{\{}\AgdaBound{Γ}\AgdaSpace{}%
\AgdaBound{A}\AgdaSymbol{\}}\AgdaSpace{}%
\AgdaSymbol{→}\AgdaSpace{}%
\AgdaBound{Γ}\AgdaSpace{}%
\AgdaOperator{\AgdaDatatype{⊢}}\AgdaSpace{}%
\AgdaBound{A}\AgdaSpace{}%
\AgdaSymbol{→}\AgdaSpace{}%
\AgdaFunction{embedContext}\AgdaSpace{}%
\AgdaBound{Γ}\AgdaSpace{}%
\AgdaOperator{\AgdaDatatype{⊢b}}\AgdaSpace{}%
\AgdaFunction{embedType}\AgdaSpace{}%
\AgdaBound{A}\<%
\\
\>[0]\AgdaFunction{pres-typing}\AgdaSpace{}%
\AgdaBound{x}\AgdaSpace{}%
\AgdaSymbol{=}\AgdaSpace{}%
\AgdaFunction{embedTerm}\AgdaSpace{}%
\AgdaBound{x}\<%
\\
%
\\[\AgdaEmptyExtraSkip]%
\>[0]\AgdaComment{--\ Preservation\ of\ typing\ of\ Values\ by\ Gödel's\ translation}\<%
\\
\>[0]\AgdaFunction{pres-typing-val}\AgdaSpace{}%
\AgdaSymbol{:}\AgdaSpace{}%
\AgdaSymbol{∀}\AgdaSpace{}%
\AgdaSymbol{\{}\AgdaBound{Γ}\AgdaSpace{}%
\AgdaBound{A}\AgdaSymbol{\}}\AgdaSpace{}%
\AgdaSymbol{→}\AgdaSpace{}%
\AgdaSymbol{(}\AgdaBound{V}\AgdaSpace{}%
\AgdaSymbol{:}\AgdaSpace{}%
\AgdaBound{Γ}\AgdaSpace{}%
\AgdaOperator{\AgdaDatatype{⊢}}\AgdaSpace{}%
\AgdaBound{A}\AgdaSymbol{)}\AgdaSpace{}%
\AgdaSymbol{→}\AgdaSpace{}%
\AgdaSymbol{(}\AgdaDatatype{Value}\AgdaSpace{}%
\AgdaBound{V}\AgdaSymbol{)}\AgdaSpace{}%
\AgdaSymbol{→}\AgdaSpace{}%
\AgdaFunction{embedContext}\AgdaSpace{}%
\AgdaBound{Γ}\AgdaSpace{}%
\AgdaOperator{\AgdaDatatype{⊢b}}\AgdaSpace{}%
\AgdaFunction{embedTypeValue}\AgdaSpace{}%
\AgdaBound{A}\<%
\\
\>[0]\AgdaFunction{pres-typing-val}\AgdaSpace{}%
\AgdaBound{V}\AgdaSpace{}%
\AgdaBound{x}\AgdaSpace{}%
\AgdaSymbol{=}\AgdaSpace{}%
\AgdaFunction{embedValue}\AgdaSpace{}%
\AgdaBound{V}\AgdaSpace{}%
\AgdaBound{x}\<%
\\
%
\\[\AgdaEmptyExtraSkip]%
\>[0]\AgdaComment{--\ Helper\ lemma:\ applying\ embedTerm\ to\ a\ Value\ is\ equivalent\ to\ wrapping\ the\ result\ of\ embedValue\ in\ a\ box}\<%
\\
\>[0]\AgdaFunction{box-embedValue}\AgdaSpace{}%
\AgdaSymbol{:}\AgdaSpace{}%
\AgdaSymbol{∀}\AgdaSpace{}%
\AgdaSymbol{\{}\AgdaBound{Γ}\AgdaSpace{}%
\AgdaBound{A}\AgdaSymbol{\}}\AgdaSpace{}%
\AgdaSymbol{→}\AgdaSpace{}%
\AgdaSymbol{\{}\AgdaBound{V}\AgdaSpace{}%
\AgdaSymbol{:}\AgdaSpace{}%
\AgdaBound{Γ}\AgdaSpace{}%
\AgdaOperator{\AgdaDatatype{⊢}}\AgdaSpace{}%
\AgdaBound{A}\AgdaSymbol{\}}\AgdaSpace{}%
\AgdaSymbol{→}\AgdaSpace{}%
\AgdaSymbol{(}\AgdaBound{x}\AgdaSpace{}%
\AgdaSymbol{:}\AgdaSpace{}%
\AgdaDatatype{Value}\AgdaSpace{}%
\AgdaBound{V}\AgdaSymbol{)}\AgdaSpace{}%
\AgdaSymbol{→}\AgdaSpace{}%
\AgdaInductiveConstructor{box}\AgdaSpace{}%
\AgdaSymbol{(}\AgdaFunction{embedValue}\AgdaSpace{}%
\AgdaBound{V}\AgdaSpace{}%
\AgdaBound{x}\AgdaSymbol{)}\AgdaSpace{}%
\AgdaOperator{\AgdaDatatype{≡}}\AgdaSpace{}%
\AgdaFunction{embedTerm}\AgdaSpace{}%
\AgdaBound{V}\<%
\\
\>[0]\AgdaFunction{box-embedValue}\AgdaSpace{}%
\AgdaInductiveConstructor{V-ƛ}\AgdaSpace{}%
\AgdaSymbol{=}\AgdaSpace{}%
\AgdaInductiveConstructor{refl}\<%
\\
\>[0]\AgdaFunction{box-embedValue}\AgdaSpace{}%
\AgdaInductiveConstructor{V-`}\AgdaSpace{}%
\AgdaSymbol{=}\AgdaSpace{}%
\AgdaInductiveConstructor{refl}\<%
\\
%
\\[\AgdaEmptyExtraSkip]%
\>[0]\AgdaComment{--\ HELPER\ LEMMA'S\ SIMULTANIOUS\ SUBSTITUTION}\<%
\\
%
\\[\AgdaEmptyExtraSkip]%
\>[0]\AgdaComment{--\ Preservation\ and\ Reflection\ of\ }\<%
\\
\>[0]\AgdaComment{--\ ext\ of\ Gödel's\ translation}\<%
\\
\>[0]\AgdaFunction{pres-ext}\AgdaSpace{}%
\AgdaSymbol{:}\AgdaSpace{}%
\AgdaSymbol{∀}\AgdaSpace{}%
\AgdaSymbol{\{}\AgdaBound{Γ}\AgdaSpace{}%
\AgdaBound{Δ}\AgdaSpace{}%
\AgdaBound{A'}\AgdaSpace{}%
\AgdaBound{A}\AgdaSymbol{\}}\<%
\\
\>[0][@{}l@{\AgdaIndent{0}}]%
\>[2]\AgdaSymbol{→}\AgdaSpace{}%
\AgdaSymbol{(}\AgdaBound{f'}\AgdaSpace{}%
\AgdaSymbol{:}\AgdaSpace{}%
\AgdaSymbol{∀}\AgdaSpace{}%
\AgdaSymbol{\{}\AgdaBound{A}\AgdaSymbol{\}}\AgdaSpace{}%
\AgdaSymbol{→}\AgdaSpace{}%
\AgdaBound{A}\AgdaSpace{}%
\AgdaOperator{\AgdaDatatype{∈}}\AgdaSpace{}%
\AgdaBound{Γ}\AgdaSpace{}%
\AgdaSymbol{→}\AgdaSpace{}%
\AgdaBound{A}\AgdaSpace{}%
\AgdaOperator{\AgdaDatatype{∈}}\AgdaSpace{}%
\AgdaBound{Δ}\AgdaSymbol{)}\<%
\\
%
\>[2]\AgdaSymbol{→}\AgdaSpace{}%
\AgdaSymbol{(}\AgdaBound{fb'}\AgdaSpace{}%
\AgdaSymbol{:}\AgdaSpace{}%
\AgdaSymbol{∀}\AgdaSpace{}%
\AgdaSymbol{\{}\AgdaBound{Ab}\AgdaSymbol{\}}\AgdaSpace{}%
\AgdaSymbol{→}\AgdaSpace{}%
\AgdaBound{Ab}\AgdaSpace{}%
\AgdaOperator{\AgdaDatatype{∈b}}\AgdaSpace{}%
\AgdaFunction{embedContext}\AgdaSpace{}%
\AgdaBound{Γ}\AgdaSpace{}%
\AgdaSymbol{→}\AgdaSpace{}%
\AgdaBound{Ab}\AgdaSpace{}%
\AgdaOperator{\AgdaDatatype{∈b}}\AgdaSpace{}%
\AgdaFunction{embedContext}\AgdaSpace{}%
\AgdaBound{Δ}\AgdaSymbol{)}\<%
\\
%
\>[2]\AgdaSymbol{→}\AgdaSpace{}%
\AgdaSymbol{(\{}\AgdaBound{A'}\AgdaSpace{}%
\AgdaSymbol{:}\AgdaSpace{}%
\AgdaDatatype{Type}\AgdaSymbol{\}}\AgdaSpace{}%
\AgdaSymbol{(}\AgdaBound{x}\AgdaSpace{}%
\AgdaSymbol{:}\AgdaSpace{}%
\AgdaBound{A'}\AgdaSpace{}%
\AgdaOperator{\AgdaDatatype{∈}}\AgdaSpace{}%
\AgdaBound{Γ}\AgdaSymbol{)}\AgdaSpace{}%
\AgdaSymbol{→}\AgdaSpace{}%
\AgdaBound{fb'}\AgdaSpace{}%
\AgdaSymbol{(}\AgdaFunction{embedLookup}\AgdaSpace{}%
\AgdaBound{x}\AgdaSymbol{)}\AgdaSpace{}%
\AgdaOperator{\AgdaDatatype{≡}}\AgdaSpace{}%
\AgdaFunction{embedLookup}\AgdaSpace{}%
\AgdaSymbol{(}\AgdaBound{f'}\AgdaSpace{}%
\AgdaBound{x}\AgdaSymbol{))}\<%
\\
%
\>[2]\AgdaSymbol{→}\AgdaSpace{}%
\AgdaSymbol{(}\AgdaBound{x}\AgdaSpace{}%
\AgdaSymbol{:}\AgdaSpace{}%
\AgdaBound{A'}\AgdaSpace{}%
\AgdaOperator{\AgdaDatatype{∈}}\AgdaSpace{}%
\AgdaBound{Γ}\AgdaSpace{}%
\AgdaOperator{\AgdaInductiveConstructor{,}}\AgdaSpace{}%
\AgdaBound{A}\AgdaSymbol{)}\<%
\\
%
\>[2]\AgdaSymbol{→}\AgdaSpace{}%
\AgdaFunction{extb}\AgdaSpace{}%
\AgdaBound{fb'}\AgdaSpace{}%
\AgdaSymbol{(}\AgdaFunction{embedLookup}\AgdaSpace{}%
\AgdaBound{x}\AgdaSymbol{)}\AgdaSpace{}%
\AgdaOperator{\AgdaDatatype{≡}}\AgdaSpace{}%
\AgdaFunction{embedLookup}\AgdaSpace{}%
\AgdaSymbol{(}\AgdaFunction{ext}\AgdaSpace{}%
\AgdaBound{f'}\AgdaSpace{}%
\AgdaBound{x}\AgdaSymbol{)}\<%
\\
\>[0]\AgdaFunction{pres-ext}\AgdaSpace{}%
\AgdaBound{f'}\AgdaSpace{}%
\AgdaBound{fb'}\AgdaSpace{}%
\AgdaBound{h}\AgdaSpace{}%
\AgdaInductiveConstructor{Z}\AgdaSpace{}%
\AgdaSymbol{=}\AgdaSpace{}%
\AgdaInductiveConstructor{refl}\<%
\\
\>[0]\AgdaFunction{pres-ext}\AgdaSpace{}%
\AgdaBound{f'}\AgdaSpace{}%
\AgdaBound{fb'}\AgdaSpace{}%
\AgdaBound{h}\AgdaSpace{}%
\AgdaSymbol{(}\AgdaInductiveConstructor{S}\AgdaSpace{}%
\AgdaBound{x}\AgdaSymbol{)}\AgdaSpace{}%
\AgdaSymbol{=}\AgdaSpace{}%
\AgdaFunction{cong}\AgdaSpace{}%
\AgdaInductiveConstructor{Sb}\AgdaSpace{}%
\AgdaSymbol{(}\AgdaBound{h}\AgdaSpace{}%
\AgdaBound{x}\AgdaSymbol{)}\<%
\\
%
\\[\AgdaEmptyExtraSkip]%
\>[0]\AgdaComment{--\ Extb\ preserves\ pointwise\ equality\ }\<%
\\
\>[0]\AgdaComment{--\ of\ renaming\ functions}\<%
\\
\>[0]\AgdaFunction{extb-cong}\AgdaSpace{}%
\AgdaSymbol{:}\AgdaSpace{}%
\AgdaSymbol{∀}\AgdaSpace{}%
\AgdaSymbol{\{}\AgdaBound{Γ₁}\AgdaSpace{}%
\AgdaBound{Γ₂}\AgdaSpace{}%
\AgdaBound{A}\AgdaSpace{}%
\AgdaBound{A'}\AgdaSymbol{\}}\<%
\\
\>[0][@{}l@{\AgdaIndent{0}}]%
\>[2]\AgdaSymbol{→}\AgdaSpace{}%
\AgdaSymbol{\{}\AgdaBound{fb'}\AgdaSpace{}%
\AgdaSymbol{:}\AgdaSpace{}%
\AgdaSymbol{∀}\AgdaSpace{}%
\AgdaSymbol{\{}\AgdaBound{Ab}\AgdaSymbol{\}}\AgdaSpace{}%
\AgdaSymbol{→}\AgdaSpace{}%
\AgdaBound{Ab}\AgdaSpace{}%
\AgdaOperator{\AgdaDatatype{∈b}}\AgdaSpace{}%
\AgdaBound{Γ₁}\AgdaSpace{}%
\AgdaSymbol{→}\AgdaSpace{}%
\AgdaBound{Ab}\AgdaSpace{}%
\AgdaOperator{\AgdaDatatype{∈b}}\AgdaSpace{}%
\AgdaBound{Γ₂}\AgdaSymbol{\}}\<%
\\
%
\>[2]\AgdaSymbol{→}\AgdaSpace{}%
\AgdaSymbol{\{}\AgdaBound{fb''}\AgdaSpace{}%
\AgdaSymbol{:}\AgdaSpace{}%
\AgdaSymbol{∀}\AgdaSpace{}%
\AgdaSymbol{\{}\AgdaBound{Ab}\AgdaSymbol{\}}\AgdaSpace{}%
\AgdaSymbol{→}\AgdaSpace{}%
\AgdaBound{Ab}\AgdaSpace{}%
\AgdaOperator{\AgdaDatatype{∈b}}\AgdaSpace{}%
\AgdaBound{Γ₁}\AgdaSpace{}%
\AgdaSymbol{→}\AgdaSpace{}%
\AgdaBound{Ab}\AgdaSpace{}%
\AgdaOperator{\AgdaDatatype{∈b}}\AgdaSpace{}%
\AgdaBound{Γ₂}\AgdaSymbol{\}}\<%
\\
%
\>[2]\AgdaSymbol{→}\AgdaSpace{}%
\AgdaSymbol{(}\AgdaBound{h}\AgdaSpace{}%
\AgdaSymbol{:}\AgdaSpace{}%
\AgdaSymbol{∀}\AgdaSpace{}%
\AgdaSymbol{\{}\AgdaBound{Ab}\AgdaSymbol{\}}\AgdaSpace{}%
\AgdaSymbol{→}\AgdaSpace{}%
\AgdaSymbol{(}\AgdaBound{x}\AgdaSpace{}%
\AgdaSymbol{:}\AgdaSpace{}%
\AgdaBound{Ab}\AgdaSpace{}%
\AgdaOperator{\AgdaDatatype{∈b}}\AgdaSpace{}%
\AgdaBound{Γ₁}\AgdaSymbol{)}\AgdaSpace{}%
\AgdaSymbol{→}\AgdaSpace{}%
\AgdaBound{fb'}\AgdaSpace{}%
\AgdaBound{x}\AgdaSpace{}%
\AgdaOperator{\AgdaDatatype{≡}}\AgdaSpace{}%
\AgdaBound{fb''}\AgdaSpace{}%
\AgdaBound{x}\AgdaSymbol{)}\<%
\\
%
\>[2]\AgdaSymbol{→}\AgdaSpace{}%
\AgdaSymbol{(}\AgdaBound{x}\AgdaSpace{}%
\AgdaSymbol{:}\AgdaSpace{}%
\AgdaBound{A}\AgdaSpace{}%
\AgdaOperator{\AgdaDatatype{∈b}}\AgdaSpace{}%
\AgdaBound{Γ₁}\AgdaSpace{}%
\AgdaOperator{\AgdaInductiveConstructor{,b}}\AgdaSpace{}%
\AgdaInductiveConstructor{□}\AgdaSpace{}%
\AgdaBound{A'}\AgdaSymbol{)}\<%
\\
%
\>[2]\AgdaSymbol{→}\AgdaSpace{}%
\AgdaFunction{extb}\AgdaSpace{}%
\AgdaBound{fb'}\AgdaSpace{}%
\AgdaBound{x}\AgdaSpace{}%
\AgdaOperator{\AgdaDatatype{≡}}\AgdaSpace{}%
\AgdaFunction{extb}\AgdaSpace{}%
\AgdaBound{fb''}\AgdaSpace{}%
\AgdaBound{x}\<%
\\
\>[0]\AgdaFunction{extb-cong}\AgdaSpace{}%
\AgdaBound{h}\AgdaSpace{}%
\AgdaInductiveConstructor{Zb}\AgdaSpace{}%
\AgdaSymbol{=}\AgdaSpace{}%
\AgdaInductiveConstructor{refl}\<%
\\
\>[0]\AgdaFunction{extb-cong}\AgdaSpace{}%
\AgdaBound{h}\AgdaSpace{}%
\AgdaSymbol{(}\AgdaInductiveConstructor{Sb}\AgdaSpace{}%
\AgdaBound{x}\AgdaSymbol{)}\AgdaSpace{}%
\AgdaSymbol{=}\AgdaSpace{}%
\AgdaFunction{cong}\AgdaSpace{}%
\AgdaInductiveConstructor{Sb}\AgdaSpace{}%
\AgdaSymbol{(}\AgdaBound{h}\AgdaSpace{}%
\AgdaBound{x}\AgdaSymbol{)}\<%
\\
%
\\[\AgdaEmptyExtraSkip]%
\>[0]\AgdaComment{--\ Renameb\ preserves\ pointwise\ equality}\<%
\\
\>[0]\AgdaComment{--\ of\ renaming\ functions}\<%
\\
\>[0]\AgdaFunction{renameb-cong}\AgdaSpace{}%
\AgdaSymbol{:}\AgdaSpace{}%
\AgdaSymbol{∀}\AgdaSpace{}%
\AgdaSymbol{\{}\AgdaBound{Γ₁}\AgdaSpace{}%
\AgdaBound{Γ₂}\AgdaSpace{}%
\AgdaBound{A}\AgdaSymbol{\}}\<%
\\
\>[0][@{}l@{\AgdaIndent{0}}]%
\>[2]\AgdaSymbol{→}\AgdaSpace{}%
\AgdaSymbol{\{}\AgdaBound{fb'}\AgdaSpace{}%
\AgdaSymbol{:}\AgdaSpace{}%
\AgdaSymbol{∀}\AgdaSpace{}%
\AgdaSymbol{\{}\AgdaBound{Ab}\AgdaSymbol{\}}\AgdaSpace{}%
\AgdaSymbol{→}\AgdaSpace{}%
\AgdaBound{Ab}\AgdaSpace{}%
\AgdaOperator{\AgdaDatatype{∈b}}\AgdaSpace{}%
\AgdaBound{Γ₁}\AgdaSpace{}%
\AgdaSymbol{→}\AgdaSpace{}%
\AgdaBound{Ab}\AgdaSpace{}%
\AgdaOperator{\AgdaDatatype{∈b}}\AgdaSpace{}%
\AgdaBound{Γ₂}\AgdaSymbol{\}}\<%
\\
%
\>[2]\AgdaSymbol{→}\AgdaSpace{}%
\AgdaSymbol{\{}\AgdaBound{fb''}\AgdaSpace{}%
\AgdaSymbol{:}\AgdaSpace{}%
\AgdaSymbol{∀}\AgdaSpace{}%
\AgdaSymbol{\{}\AgdaBound{Ab}\AgdaSymbol{\}}\AgdaSpace{}%
\AgdaSymbol{→}\AgdaSpace{}%
\AgdaBound{Ab}\AgdaSpace{}%
\AgdaOperator{\AgdaDatatype{∈b}}\AgdaSpace{}%
\AgdaBound{Γ₁}\AgdaSpace{}%
\AgdaSymbol{→}\AgdaSpace{}%
\AgdaBound{Ab}\AgdaSpace{}%
\AgdaOperator{\AgdaDatatype{∈b}}\AgdaSpace{}%
\AgdaBound{Γ₂}\AgdaSymbol{\}}\<%
\\
%
\>[2]\AgdaSymbol{→}\AgdaSpace{}%
\AgdaSymbol{(}\AgdaBound{h}\AgdaSpace{}%
\AgdaSymbol{:}\AgdaSpace{}%
\AgdaSymbol{∀}\AgdaSpace{}%
\AgdaSymbol{\{}\AgdaBound{Ab}\AgdaSymbol{\}}\AgdaSpace{}%
\AgdaSymbol{→}\AgdaSpace{}%
\AgdaSymbol{(}\AgdaBound{x}\AgdaSpace{}%
\AgdaSymbol{:}\AgdaSpace{}%
\AgdaBound{Ab}\AgdaSpace{}%
\AgdaOperator{\AgdaDatatype{∈b}}\AgdaSpace{}%
\AgdaBound{Γ₁}\AgdaSymbol{)}\AgdaSpace{}%
\AgdaSymbol{→}\AgdaSpace{}%
\AgdaBound{fb'}\AgdaSpace{}%
\AgdaBound{x}\AgdaSpace{}%
\AgdaOperator{\AgdaDatatype{≡}}\AgdaSpace{}%
\AgdaBound{fb''}\AgdaSpace{}%
\AgdaBound{x}\AgdaSymbol{)}\<%
\\
%
\>[2]\AgdaSymbol{→}\AgdaSpace{}%
\AgdaSymbol{(}\AgdaBound{t}\AgdaSpace{}%
\AgdaSymbol{:}\AgdaSpace{}%
\AgdaBound{Γ₁}\AgdaSpace{}%
\AgdaOperator{\AgdaDatatype{⊢b}}\AgdaSpace{}%
\AgdaBound{A}\AgdaSymbol{)}\<%
\\
%
\>[2]\AgdaSymbol{→}\AgdaSpace{}%
\AgdaFunction{renameb}\AgdaSpace{}%
\AgdaBound{fb'}\AgdaSpace{}%
\AgdaBound{t}\AgdaSpace{}%
\AgdaOperator{\AgdaDatatype{≡}}\AgdaSpace{}%
\AgdaFunction{renameb}\AgdaSpace{}%
\AgdaBound{fb''}\AgdaSpace{}%
\AgdaBound{t}\<%
\\
\>[0]\AgdaFunction{renameb-cong}\AgdaSpace{}%
\AgdaBound{h}\AgdaSpace{}%
\AgdaSymbol{(}\AgdaInductiveConstructor{ε}\AgdaSpace{}%
\AgdaBound{x}\AgdaSymbol{)}\AgdaSpace{}%
\AgdaSymbol{=}\AgdaSpace{}%
\AgdaFunction{cong}\AgdaSpace{}%
\AgdaInductiveConstructor{ε}\AgdaSpace{}%
\AgdaSymbol{(}\AgdaBound{h}\AgdaSpace{}%
\AgdaBound{x}\AgdaSymbol{)}\<%
\\
\>[0]\AgdaFunction{renameb-cong}\AgdaSpace{}%
\AgdaBound{h}\AgdaSpace{}%
\AgdaSymbol{(}\AgdaInductiveConstructor{box}\AgdaSpace{}%
\AgdaBound{t}\AgdaSymbol{)}\AgdaSpace{}%
\AgdaSymbol{=}\AgdaSpace{}%
\AgdaFunction{cong}\AgdaSpace{}%
\AgdaInductiveConstructor{box}\AgdaSpace{}%
\AgdaSymbol{(}\AgdaFunction{renameb-cong}\AgdaSpace{}%
\AgdaBound{h}\AgdaSpace{}%
\AgdaBound{t}\AgdaSymbol{)}\<%
\\
\>[0]\AgdaFunction{renameb-cong}\AgdaSpace{}%
\AgdaBound{h}\AgdaSpace{}%
\AgdaSymbol{(}\AgdaInductiveConstructor{ƛb}\AgdaSpace{}%
\AgdaBound{t}\AgdaSymbol{)}\AgdaSpace{}%
\AgdaSymbol{=}\AgdaSpace{}%
\AgdaFunction{cong}\AgdaSpace{}%
\AgdaInductiveConstructor{ƛb}\AgdaSpace{}%
\AgdaSymbol{(}\AgdaFunction{renameb-cong}\AgdaSpace{}%
\AgdaSymbol{(}\AgdaFunction{extb-cong}\AgdaSpace{}%
\AgdaBound{h}\AgdaSymbol{)}\AgdaSpace{}%
\AgdaBound{t}\AgdaSymbol{)}\<%
\\
\>[0]\AgdaFunction{renameb-cong}\AgdaSpace{}%
\AgdaBound{h}\AgdaSpace{}%
\AgdaSymbol{(}\AgdaBound{t}\AgdaSpace{}%
\AgdaOperator{\AgdaInductiveConstructor{·b}}\AgdaSpace{}%
\AgdaBound{t₁}\AgdaSymbol{)}\AgdaSpace{}%
\AgdaSymbol{=}\AgdaSpace{}%
\AgdaFunction{cong₂}\AgdaSpace{}%
\AgdaOperator{\AgdaInductiveConstructor{\AgdaUnderscore{}·b\AgdaUnderscore{}}}\AgdaSpace{}%
\AgdaSymbol{(}\AgdaFunction{renameb-cong}\AgdaSpace{}%
\AgdaBound{h}\AgdaSpace{}%
\AgdaBound{t}\AgdaSymbol{)}\AgdaSpace{}%
\AgdaSymbol{(}\AgdaFunction{renameb-cong}\AgdaSpace{}%
\AgdaBound{h}\AgdaSpace{}%
\AgdaBound{t₁}\AgdaSymbol{)}\<%
\\
%
\\[\AgdaEmptyExtraSkip]%
\>[0]\AgdaComment{--\ extb\ distributes\ over\ composition:}\<%
\\
\>[0]\AgdaComment{--\ Extending\ first\ with\ fb',\ then\ with\ fb'',\ }\<%
\\
\>[0]\AgdaComment{--\ is\ the\ same\ as\ extending\ once\ with\ fb''\ ∘\ fb'.}\<%
\\
\>[0]\AgdaFunction{extb-comp}\AgdaSpace{}%
\AgdaSymbol{:}\AgdaSpace{}%
\AgdaSymbol{∀}\AgdaSpace{}%
\AgdaSymbol{\{}\AgdaBound{Γ₁}\AgdaSpace{}%
\AgdaBound{Γ₂}\AgdaSpace{}%
\AgdaBound{Γ₃}\AgdaSpace{}%
\AgdaBound{A}\AgdaSpace{}%
\AgdaBound{B}\AgdaSymbol{\}}\<%
\\
\>[0][@{}l@{\AgdaIndent{0}}]%
\>[2]\AgdaSymbol{→}\AgdaSpace{}%
\AgdaSymbol{(}\AgdaBound{fb'}\AgdaSpace{}%
\AgdaSymbol{:}\AgdaSpace{}%
\AgdaSymbol{∀}\AgdaSpace{}%
\AgdaSymbol{\{}\AgdaBound{Ab}\AgdaSymbol{\}}\AgdaSpace{}%
\AgdaSymbol{→}\AgdaSpace{}%
\AgdaBound{Ab}\AgdaSpace{}%
\AgdaOperator{\AgdaDatatype{∈b}}\AgdaSpace{}%
\AgdaBound{Γ₁}\AgdaSpace{}%
\AgdaSymbol{→}\AgdaSpace{}%
\AgdaBound{Ab}\AgdaSpace{}%
\AgdaOperator{\AgdaDatatype{∈b}}\AgdaSpace{}%
\AgdaBound{Γ₂}\AgdaSymbol{)}\<%
\\
%
\>[2]\AgdaSymbol{→}\AgdaSpace{}%
\AgdaSymbol{(}\AgdaBound{fb''}\AgdaSpace{}%
\AgdaSymbol{:}\AgdaSpace{}%
\AgdaSymbol{∀}\AgdaSpace{}%
\AgdaSymbol{\{}\AgdaBound{Ab}\AgdaSymbol{\}}\AgdaSpace{}%
\AgdaSymbol{→}\AgdaSpace{}%
\AgdaBound{Ab}\AgdaSpace{}%
\AgdaOperator{\AgdaDatatype{∈b}}\AgdaSpace{}%
\AgdaBound{Γ₂}\AgdaSpace{}%
\AgdaSymbol{→}\AgdaSpace{}%
\AgdaBound{Ab}\AgdaSpace{}%
\AgdaOperator{\AgdaDatatype{∈b}}\AgdaSpace{}%
\AgdaBound{Γ₃}\AgdaSymbol{)}\<%
\\
%
\>[2]\AgdaSymbol{→}\AgdaSpace{}%
\AgdaSymbol{(}\AgdaBound{x}\AgdaSpace{}%
\AgdaSymbol{:}\AgdaSpace{}%
\AgdaBound{A}\AgdaSpace{}%
\AgdaOperator{\AgdaDatatype{∈b}}\AgdaSpace{}%
\AgdaBound{Γ₁}\AgdaSpace{}%
\AgdaOperator{\AgdaInductiveConstructor{,b}}\AgdaSpace{}%
\AgdaBound{B}\AgdaSymbol{)}\<%
\\
%
\>[2]\AgdaSymbol{→}\AgdaSpace{}%
\AgdaFunction{extb}\AgdaSpace{}%
\AgdaBound{fb''}\AgdaSpace{}%
\AgdaSymbol{(}\AgdaFunction{extb}\AgdaSpace{}%
\AgdaBound{fb'}\AgdaSpace{}%
\AgdaBound{x}\AgdaSymbol{)}\AgdaSpace{}%
\AgdaOperator{\AgdaDatatype{≡}}\AgdaSpace{}%
\AgdaFunction{extb}\AgdaSpace{}%
\AgdaSymbol{(λ}\AgdaSpace{}%
\AgdaBound{x₁}\AgdaSpace{}%
\AgdaSymbol{→}\AgdaSpace{}%
\AgdaBound{fb''}\AgdaSpace{}%
\AgdaSymbol{(}\AgdaBound{fb'}\AgdaSpace{}%
\AgdaBound{x₁}\AgdaSymbol{))}\AgdaSpace{}%
\AgdaBound{x}\<%
\\
\>[0]\AgdaFunction{extb-comp}\AgdaSpace{}%
\AgdaBound{fb'}\AgdaSpace{}%
\AgdaBound{fb''}\AgdaSpace{}%
\AgdaInductiveConstructor{Zb}\AgdaSpace{}%
\AgdaSymbol{=}\AgdaSpace{}%
\AgdaInductiveConstructor{refl}\<%
\\
\>[0]\AgdaFunction{extb-comp}\AgdaSpace{}%
\AgdaBound{fb'}\AgdaSpace{}%
\AgdaBound{fb''}\AgdaSpace{}%
\AgdaSymbol{(}\AgdaInductiveConstructor{Sb}\AgdaSpace{}%
\AgdaBound{x}\AgdaSymbol{)}\AgdaSpace{}%
\AgdaSymbol{=}\AgdaSpace{}%
\AgdaInductiveConstructor{refl}\<%
\\
%
\\[\AgdaEmptyExtraSkip]%
\>[0]\AgdaComment{--\ renameb\ distributes\ over\ composition:}\<%
\\
\>[0]\AgdaComment{--\ Renaming\ first\ with\ fb',\ then\ with\ fb''}\<%
\\
\>[0]\AgdaComment{--\ is\ the\ same\ as\ renaming\ once\ with\ fb''\ ∘\ fb'.}\<%
\\
\>[0]\AgdaFunction{renameb-comp}\AgdaSpace{}%
\AgdaSymbol{:}\AgdaSpace{}%
\AgdaSymbol{∀}\AgdaSpace{}%
\AgdaSymbol{\{}\AgdaBound{Γ₁}\AgdaSpace{}%
\AgdaBound{Γ₂}\AgdaSpace{}%
\AgdaBound{Γ₃}\AgdaSpace{}%
\AgdaBound{A}\AgdaSymbol{\}}\<%
\\
\>[0][@{}l@{\AgdaIndent{0}}]%
\>[2]\AgdaSymbol{→}\AgdaSpace{}%
\AgdaSymbol{(}\AgdaBound{fb'}\AgdaSpace{}%
\AgdaSymbol{:}\AgdaSpace{}%
\AgdaSymbol{∀}\AgdaSpace{}%
\AgdaSymbol{\{}\AgdaBound{Ab}\AgdaSymbol{\}}\AgdaSpace{}%
\AgdaSymbol{→}\AgdaSpace{}%
\AgdaBound{Ab}\AgdaSpace{}%
\AgdaOperator{\AgdaDatatype{∈b}}\AgdaSpace{}%
\AgdaBound{Γ₁}\AgdaSpace{}%
\AgdaSymbol{→}\AgdaSpace{}%
\AgdaBound{Ab}\AgdaSpace{}%
\AgdaOperator{\AgdaDatatype{∈b}}\AgdaSpace{}%
\AgdaBound{Γ₂}\AgdaSymbol{)}\<%
\\
%
\>[2]\AgdaSymbol{→}\AgdaSpace{}%
\AgdaSymbol{(}\AgdaBound{fb''}\AgdaSpace{}%
\AgdaSymbol{:}\AgdaSpace{}%
\AgdaSymbol{∀}\AgdaSpace{}%
\AgdaSymbol{\{}\AgdaBound{Ab}\AgdaSymbol{\}}\AgdaSpace{}%
\AgdaSymbol{→}\AgdaSpace{}%
\AgdaBound{Ab}\AgdaSpace{}%
\AgdaOperator{\AgdaDatatype{∈b}}\AgdaSpace{}%
\AgdaBound{Γ₂}\AgdaSpace{}%
\AgdaSymbol{→}\AgdaSpace{}%
\AgdaBound{Ab}\AgdaSpace{}%
\AgdaOperator{\AgdaDatatype{∈b}}\AgdaSpace{}%
\AgdaBound{Γ₃}\AgdaSymbol{)}\<%
\\
%
\>[2]\AgdaSymbol{→}\AgdaSpace{}%
\AgdaSymbol{(}\AgdaBound{t}\AgdaSpace{}%
\AgdaSymbol{:}\AgdaSpace{}%
\AgdaBound{Γ₁}\AgdaSpace{}%
\AgdaOperator{\AgdaDatatype{⊢b}}\AgdaSpace{}%
\AgdaBound{A}\AgdaSymbol{)}\<%
\\
%
\>[2]\AgdaSymbol{→}\AgdaSpace{}%
\AgdaFunction{renameb}\AgdaSpace{}%
\AgdaBound{fb''}\AgdaSpace{}%
\AgdaSymbol{(}\AgdaFunction{renameb}\AgdaSpace{}%
\AgdaBound{fb'}\AgdaSpace{}%
\AgdaBound{t}\AgdaSymbol{)}\AgdaSpace{}%
\AgdaOperator{\AgdaDatatype{≡}}\AgdaSpace{}%
\AgdaFunction{renameb}\AgdaSpace{}%
\AgdaSymbol{(λ}\AgdaSpace{}%
\AgdaBound{x}\AgdaSpace{}%
\AgdaSymbol{→}\AgdaSpace{}%
\AgdaBound{fb''}\AgdaSpace{}%
\AgdaSymbol{(}\AgdaBound{fb'}\AgdaSpace{}%
\AgdaBound{x}\AgdaSymbol{))}\AgdaSpace{}%
\AgdaBound{t}\<%
\\
\>[0]\AgdaFunction{renameb-comp}\AgdaSpace{}%
\AgdaBound{fb'}\AgdaSpace{}%
\AgdaBound{fb''}\AgdaSpace{}%
\AgdaSymbol{(}\AgdaInductiveConstructor{ε}\AgdaSpace{}%
\AgdaBound{x}\AgdaSymbol{)}\AgdaSpace{}%
\AgdaSymbol{=}\AgdaSpace{}%
\AgdaInductiveConstructor{refl}\<%
\\
\>[0]\AgdaFunction{renameb-comp}\AgdaSpace{}%
\AgdaBound{fb'}\AgdaSpace{}%
\AgdaBound{fb''}\AgdaSpace{}%
\AgdaSymbol{(}\AgdaInductiveConstructor{box}\AgdaSpace{}%
\AgdaBound{t}\AgdaSymbol{)}\AgdaSpace{}%
\AgdaSymbol{=}\AgdaSpace{}%
\AgdaFunction{cong}\AgdaSpace{}%
\AgdaInductiveConstructor{box}\AgdaSpace{}%
\AgdaSymbol{(}\AgdaFunction{renameb-comp}\AgdaSpace{}%
\AgdaBound{fb'}\AgdaSpace{}%
\AgdaBound{fb''}\AgdaSpace{}%
\AgdaBound{t}\AgdaSymbol{)}\<%
\\
\>[0]\AgdaFunction{renameb-comp}\AgdaSpace{}%
\AgdaBound{fb'}\AgdaSpace{}%
\AgdaBound{fb''}\AgdaSpace{}%
\AgdaSymbol{(}\AgdaInductiveConstructor{ƛb}\AgdaSpace{}%
\AgdaBound{t}\AgdaSymbol{)}\AgdaSpace{}%
\AgdaSymbol{=}\<%
\\
\>[0][@{}l@{\AgdaIndent{0}}]%
\>[2]\AgdaOperator{\AgdaFunction{beginEq}}\<%
\\
\>[2][@{}l@{\AgdaIndent{0}}]%
\>[4]\AgdaInductiveConstructor{ƛb}\AgdaSpace{}%
\AgdaSymbol{(}\AgdaFunction{renameb}\AgdaSpace{}%
\AgdaSymbol{(}\AgdaFunction{extb}\AgdaSpace{}%
\AgdaBound{fb''}\AgdaSymbol{)}\AgdaSpace{}%
\AgdaSymbol{(}\AgdaFunction{renameb}\AgdaSpace{}%
\AgdaSymbol{(}\AgdaFunction{extb}\AgdaSpace{}%
\AgdaBound{fb'}\AgdaSymbol{)}\AgdaSpace{}%
\AgdaBound{t}\AgdaSymbol{))}\<%
\\
%
\>[2]\AgdaFunction{≡⟨}\AgdaSpace{}%
\AgdaFunction{cong}\AgdaSpace{}%
\AgdaInductiveConstructor{ƛb}\AgdaSpace{}%
\AgdaSymbol{(}\AgdaFunction{renameb-comp}\AgdaSpace{}%
\AgdaSymbol{(}\AgdaFunction{extb}\AgdaSpace{}%
\AgdaBound{fb'}\AgdaSymbol{)}\AgdaSpace{}%
\AgdaSymbol{(}\AgdaFunction{extb}\AgdaSpace{}%
\AgdaBound{fb''}\AgdaSymbol{)}\AgdaSpace{}%
\AgdaBound{t}\AgdaSymbol{)}\AgdaSpace{}%
\AgdaFunction{⟩}\<%
\\
\>[2][@{}l@{\AgdaIndent{0}}]%
\>[4]\AgdaInductiveConstructor{ƛb}\AgdaSpace{}%
\AgdaSymbol{(}\AgdaFunction{renameb}\AgdaSpace{}%
\AgdaSymbol{(λ}\AgdaSpace{}%
\AgdaBound{x}\AgdaSpace{}%
\AgdaSymbol{→}\AgdaSpace{}%
\AgdaFunction{extb}\AgdaSpace{}%
\AgdaBound{fb''}\AgdaSpace{}%
\AgdaSymbol{(}\AgdaFunction{extb}\AgdaSpace{}%
\AgdaBound{fb'}\AgdaSpace{}%
\AgdaBound{x}\AgdaSymbol{))}\AgdaSpace{}%
\AgdaBound{t}\AgdaSymbol{)}\<%
\\
%
\>[2]\AgdaFunction{≡⟨}\AgdaSpace{}%
\AgdaFunction{cong}\AgdaSpace{}%
\AgdaInductiveConstructor{ƛb}\AgdaSpace{}%
\AgdaSymbol{(}\AgdaFunction{renameb-cong}\AgdaSpace{}%
\AgdaSymbol{(}\AgdaFunction{extb-comp}\AgdaSpace{}%
\AgdaBound{fb'}\AgdaSpace{}%
\AgdaBound{fb''}\AgdaSymbol{)}\AgdaSpace{}%
\AgdaBound{t}\AgdaSymbol{)}\AgdaSpace{}%
\AgdaFunction{⟩}\<%
\\
\>[2][@{}l@{\AgdaIndent{0}}]%
\>[4]\AgdaInductiveConstructor{ƛb}\AgdaSpace{}%
\AgdaSymbol{(}\AgdaFunction{renameb}\AgdaSpace{}%
\AgdaSymbol{(}\AgdaFunction{extb}\AgdaSpace{}%
\AgdaSymbol{(λ}\AgdaSpace{}%
\AgdaBound{x}\AgdaSpace{}%
\AgdaSymbol{→}\AgdaSpace{}%
\AgdaBound{fb''}\AgdaSpace{}%
\AgdaSymbol{(}\AgdaBound{fb'}\AgdaSpace{}%
\AgdaBound{x}\AgdaSymbol{)))}\AgdaSpace{}%
\AgdaBound{t}\AgdaSymbol{)}\<%
\\
%
\>[2]\AgdaOperator{\AgdaFunction{∎eq}}\<%
\\
\>[0]\AgdaFunction{renameb-comp}\AgdaSpace{}%
\AgdaBound{fb'}\AgdaSpace{}%
\AgdaBound{fb''}\AgdaSpace{}%
\AgdaSymbol{(}\AgdaBound{t}\AgdaSpace{}%
\AgdaOperator{\AgdaInductiveConstructor{·b}}\AgdaSpace{}%
\AgdaBound{t₁}\AgdaSymbol{)}\AgdaSpace{}%
\AgdaSymbol{=}\AgdaSpace{}%
\AgdaFunction{cong₂}\AgdaSpace{}%
\AgdaOperator{\AgdaInductiveConstructor{\AgdaUnderscore{}·b\AgdaUnderscore{}}}\AgdaSpace{}%
\AgdaSymbol{(}\AgdaFunction{renameb-comp}\AgdaSpace{}%
\AgdaBound{fb'}\AgdaSpace{}%
\AgdaBound{fb''}\AgdaSpace{}%
\AgdaBound{t}\AgdaSymbol{)}\AgdaSpace{}%
\AgdaSymbol{(}\AgdaFunction{renameb-comp}\AgdaSpace{}%
\AgdaBound{fb'}\AgdaSpace{}%
\AgdaBound{fb''}\AgdaSpace{}%
\AgdaBound{t₁}\AgdaSymbol{)}\<%
\\
%
\\[\AgdaEmptyExtraSkip]%
\>[0]\AgdaComment{--\ extb\ and\ Sb\ commute\ under\ renameb.}\<%
\\
\>[0]\AgdaFunction{rename-extb-Sb-commute}\AgdaSpace{}%
\AgdaSymbol{:}\AgdaSpace{}%
\AgdaSymbol{∀}\AgdaSpace{}%
\AgdaSymbol{\{}\AgdaBound{Γ}\AgdaSpace{}%
\AgdaBound{A}\AgdaSpace{}%
\AgdaBound{B}\AgdaSpace{}%
\AgdaBound{Δ}\AgdaSymbol{\}}\<%
\\
\>[0][@{}l@{\AgdaIndent{0}}]%
\>[2]\AgdaSymbol{→}\AgdaSpace{}%
\AgdaSymbol{(}\AgdaBound{fb'}\AgdaSpace{}%
\AgdaSymbol{:}\AgdaSpace{}%
\AgdaSymbol{∀}\AgdaSpace{}%
\AgdaSymbol{\{}\AgdaBound{Ab}\AgdaSymbol{\}}\AgdaSpace{}%
\AgdaSymbol{→}\AgdaSpace{}%
\AgdaBound{Ab}\AgdaSpace{}%
\AgdaOperator{\AgdaDatatype{∈b}}\AgdaSpace{}%
\AgdaBound{Γ}\AgdaSpace{}%
\AgdaSymbol{→}\AgdaSpace{}%
\AgdaBound{Ab}\AgdaSpace{}%
\AgdaOperator{\AgdaDatatype{∈b}}\AgdaSpace{}%
\AgdaBound{Δ}\AgdaSymbol{)}\<%
\\
%
\>[2]\AgdaSymbol{→}\AgdaSpace{}%
\AgdaSymbol{(}\AgdaBound{t}\AgdaSpace{}%
\AgdaSymbol{:}\AgdaSpace{}%
\AgdaBound{Γ}\AgdaSpace{}%
\AgdaOperator{\AgdaDatatype{⊢b}}\AgdaSpace{}%
\AgdaBound{A}\AgdaSymbol{)}\<%
\\
%
\>[2]\AgdaSymbol{→}%
\>[579I]\AgdaFunction{renameb}\AgdaSpace{}%
\AgdaSymbol{(}\AgdaFunction{extb}\AgdaSpace{}%
\AgdaBound{fb'}\AgdaSpace{}%
\AgdaSymbol{\{}\AgdaArgument{B}\AgdaSpace{}%
\AgdaSymbol{=}\AgdaSpace{}%
\AgdaBound{B}\AgdaSymbol{\})}\AgdaSpace{}%
\AgdaSymbol{(}\AgdaFunction{renameb}\AgdaSpace{}%
\AgdaInductiveConstructor{Sb}\AgdaSpace{}%
\AgdaBound{t}\AgdaSymbol{)}\<%
\\
\>[.][@{}l@{}]\<[579I]%
\>[4]\AgdaOperator{\AgdaDatatype{≡}}\AgdaSpace{}%
\AgdaFunction{renameb}\AgdaSpace{}%
\AgdaInductiveConstructor{Sb}\AgdaSpace{}%
\AgdaSymbol{(}\AgdaFunction{renameb}\AgdaSpace{}%
\AgdaBound{fb'}\AgdaSpace{}%
\AgdaBound{t}\AgdaSymbol{)}\<%
\\
\>[0]\AgdaFunction{rename-extb-Sb-commute}\AgdaSpace{}%
\AgdaBound{fb'}\AgdaSpace{}%
\AgdaBound{t}\AgdaSpace{}%
\AgdaSymbol{=}\<%
\\
\>[0][@{}l@{\AgdaIndent{0}}]%
\>[2]\AgdaOperator{\AgdaFunction{beginEq}}\<%
\\
\>[2][@{}l@{\AgdaIndent{0}}]%
\>[4]\AgdaFunction{renameb}\AgdaSpace{}%
\AgdaSymbol{(}\AgdaFunction{extb}\AgdaSpace{}%
\AgdaBound{fb'}\AgdaSymbol{)}\AgdaSpace{}%
\AgdaSymbol{(}\AgdaFunction{renameb}\AgdaSpace{}%
\AgdaInductiveConstructor{Sb}\AgdaSpace{}%
\AgdaBound{t}\AgdaSymbol{)}\<%
\\
%
\>[2]\AgdaFunction{≡⟨}\AgdaSpace{}%
\AgdaFunction{renameb-comp}\AgdaSpace{}%
\AgdaInductiveConstructor{Sb}\AgdaSpace{}%
\AgdaSymbol{(}\AgdaFunction{extb}\AgdaSpace{}%
\AgdaBound{fb'}\AgdaSymbol{)}\AgdaSpace{}%
\AgdaBound{t}\AgdaSpace{}%
\AgdaFunction{⟩}\<%
\\
\>[2][@{}l@{\AgdaIndent{0}}]%
\>[4]\AgdaFunction{renameb}\AgdaSpace{}%
\AgdaSymbol{(λ}\AgdaSpace{}%
\AgdaBound{x}\AgdaSpace{}%
\AgdaSymbol{→}\AgdaSpace{}%
\AgdaFunction{extb}\AgdaSpace{}%
\AgdaBound{fb'}\AgdaSpace{}%
\AgdaSymbol{(}\AgdaInductiveConstructor{Sb}\AgdaSpace{}%
\AgdaBound{x}\AgdaSymbol{))}\AgdaSpace{}%
\AgdaBound{t}\<%
\\
%
\>[2]\AgdaFunction{≡⟨}\AgdaSpace{}%
\AgdaFunction{renameb-cong}\AgdaSpace{}%
\AgdaSymbol{(λ}\AgdaSpace{}%
\AgdaBound{x}\AgdaSpace{}%
\AgdaSymbol{→}\AgdaSpace{}%
\AgdaInductiveConstructor{refl}\AgdaSymbol{)}\AgdaSpace{}%
\AgdaBound{t}\AgdaSpace{}%
\AgdaFunction{⟩}\<%
\\
\>[2][@{}l@{\AgdaIndent{0}}]%
\>[4]\AgdaFunction{renameb}\AgdaSpace{}%
\AgdaSymbol{(λ}\AgdaSpace{}%
\AgdaBound{x}\AgdaSpace{}%
\AgdaSymbol{→}\AgdaSpace{}%
\AgdaInductiveConstructor{Sb}\AgdaSpace{}%
\AgdaSymbol{(}\AgdaBound{fb'}\AgdaSpace{}%
\AgdaBound{x}\AgdaSymbol{))}\AgdaSpace{}%
\AgdaBound{t}\<%
\\
%
\>[2]\AgdaFunction{≡⟨}\AgdaSpace{}%
\AgdaFunction{sym}\AgdaSpace{}%
\AgdaSymbol{(}\AgdaFunction{renameb-comp}\AgdaSpace{}%
\AgdaBound{fb'}\AgdaSpace{}%
\AgdaInductiveConstructor{Sb}\AgdaSpace{}%
\AgdaBound{t}\AgdaSymbol{)}\AgdaSpace{}%
\AgdaFunction{⟩}\<%
\\
\>[2][@{}l@{\AgdaIndent{0}}]%
\>[4]\AgdaFunction{renameb}\AgdaSpace{}%
\AgdaInductiveConstructor{Sb}\AgdaSpace{}%
\AgdaSymbol{(}\AgdaFunction{renameb}\AgdaSpace{}%
\AgdaBound{fb'}\AgdaSpace{}%
\AgdaBound{t}\AgdaSymbol{)}\<%
\\
%
\>[2]\AgdaOperator{\AgdaFunction{∎eq}}\<%
\\
%
\\[\AgdaEmptyExtraSkip]%
\>[0]\AgdaComment{--\ Preservation\ and\ Reflection\ of\ }\<%
\\
\>[0]\AgdaComment{--\ rename\ of\ Gödel's\ translation}\<%
\\
\>[0]\AgdaFunction{pres-rename}\AgdaSpace{}%
\AgdaSymbol{:}\AgdaSpace{}%
\AgdaSymbol{∀}\AgdaSpace{}%
\AgdaSymbol{\{}\AgdaBound{Γ}\AgdaSpace{}%
\AgdaBound{A}\AgdaSpace{}%
\AgdaBound{Δ}\AgdaSymbol{\}}\<%
\\
\>[0][@{}l@{\AgdaIndent{0}}]%
\>[2]\AgdaSymbol{→}\AgdaSpace{}%
\AgdaSymbol{(}\AgdaBound{f'}\AgdaSpace{}%
\AgdaSymbol{:}\AgdaSpace{}%
\AgdaSymbol{∀}\AgdaSpace{}%
\AgdaSymbol{\{}\AgdaBound{A}\AgdaSymbol{\}}\AgdaSpace{}%
\AgdaSymbol{→}\AgdaSpace{}%
\AgdaBound{A}\AgdaSpace{}%
\AgdaOperator{\AgdaDatatype{∈}}\AgdaSpace{}%
\AgdaBound{Γ}\AgdaSpace{}%
\AgdaSymbol{→}\AgdaSpace{}%
\AgdaBound{A}\AgdaSpace{}%
\AgdaOperator{\AgdaDatatype{∈}}\AgdaSpace{}%
\AgdaBound{Δ}\AgdaSymbol{)}\<%
\\
%
\>[2]\AgdaSymbol{→}\AgdaSpace{}%
\AgdaSymbol{(}\AgdaBound{fb'}\AgdaSpace{}%
\AgdaSymbol{:}\AgdaSpace{}%
\AgdaSymbol{∀}\AgdaSpace{}%
\AgdaSymbol{\{}\AgdaBound{Ab}\AgdaSymbol{\}}\AgdaSpace{}%
\AgdaSymbol{→}\AgdaSpace{}%
\AgdaBound{Ab}\AgdaSpace{}%
\AgdaOperator{\AgdaDatatype{∈b}}\AgdaSpace{}%
\AgdaFunction{embedContext}\AgdaSpace{}%
\AgdaBound{Γ}\AgdaSpace{}%
\AgdaSymbol{→}\AgdaSpace{}%
\AgdaBound{Ab}\AgdaSpace{}%
\AgdaOperator{\AgdaDatatype{∈b}}\AgdaSpace{}%
\AgdaFunction{embedContext}\AgdaSpace{}%
\AgdaBound{Δ}\AgdaSymbol{)}\<%
\\
%
\>[2]\AgdaSymbol{→}\AgdaSpace{}%
\AgdaSymbol{(\{}\AgdaBound{A'}\AgdaSpace{}%
\AgdaSymbol{:}\AgdaSpace{}%
\AgdaDatatype{Type}\AgdaSymbol{\}}\AgdaSpace{}%
\AgdaSymbol{(}\AgdaBound{x}\AgdaSpace{}%
\AgdaSymbol{:}\AgdaSpace{}%
\AgdaBound{A'}\AgdaSpace{}%
\AgdaOperator{\AgdaDatatype{∈}}\AgdaSpace{}%
\AgdaBound{Γ}\AgdaSymbol{)}\AgdaSpace{}%
\AgdaSymbol{→}\AgdaSpace{}%
\AgdaBound{fb'}\AgdaSpace{}%
\AgdaSymbol{(}\AgdaFunction{embedLookup}\AgdaSpace{}%
\AgdaBound{x}\AgdaSymbol{)}\AgdaSpace{}%
\AgdaOperator{\AgdaDatatype{≡}}\AgdaSpace{}%
\AgdaFunction{embedLookup}\AgdaSpace{}%
\AgdaSymbol{(}\AgdaBound{f'}\AgdaSpace{}%
\AgdaBound{x}\AgdaSymbol{))}\<%
\\
%
\>[2]\AgdaSymbol{→}\AgdaSpace{}%
\AgdaSymbol{(}\AgdaBound{t}\AgdaSpace{}%
\AgdaSymbol{:}\AgdaSpace{}%
\AgdaBound{Γ}\AgdaSpace{}%
\AgdaOperator{\AgdaDatatype{⊢}}\AgdaSpace{}%
\AgdaBound{A}\AgdaSymbol{)}\<%
\\
%
\>[2]\AgdaSymbol{→}\AgdaSpace{}%
\AgdaFunction{renameb}\AgdaSpace{}%
\AgdaBound{fb'}\AgdaSpace{}%
\AgdaSymbol{(}\AgdaFunction{embedTerm}\AgdaSpace{}%
\AgdaBound{t}\AgdaSymbol{)}\AgdaSpace{}%
\AgdaOperator{\AgdaDatatype{≡}}\AgdaSpace{}%
\AgdaFunction{embedTerm}\AgdaSpace{}%
\AgdaSymbol{(}\AgdaFunction{rename}\AgdaSpace{}%
\AgdaBound{f'}\AgdaSpace{}%
\AgdaBound{t}\AgdaSymbol{)}\<%
\\
\>[0]\AgdaFunction{pres-rename}\AgdaSpace{}%
\AgdaBound{f'}\AgdaSpace{}%
\AgdaBound{fb'}\AgdaSpace{}%
\AgdaBound{h}\AgdaSpace{}%
\AgdaSymbol{(}\AgdaInductiveConstructor{`}\AgdaSpace{}%
\AgdaBound{x}\AgdaSymbol{)}\AgdaSpace{}%
\AgdaSymbol{=}\AgdaSpace{}%
\AgdaFunction{cong}\AgdaSpace{}%
\AgdaInductiveConstructor{box}\AgdaSpace{}%
\AgdaSymbol{(}\AgdaFunction{cong}\AgdaSpace{}%
\AgdaInductiveConstructor{ε}\AgdaSpace{}%
\AgdaSymbol{(}\AgdaBound{h}\AgdaSpace{}%
\AgdaBound{x}\AgdaSymbol{))}\<%
\\
\>[0]\AgdaFunction{pres-rename}\AgdaSpace{}%
\AgdaBound{f'}\AgdaSpace{}%
\AgdaBound{fb'}\AgdaSpace{}%
\AgdaBound{h}\AgdaSpace{}%
\AgdaSymbol{(}\AgdaInductiveConstructor{ƛ}\AgdaSpace{}%
\AgdaBound{t}\AgdaSymbol{)}\AgdaSpace{}%
\AgdaSymbol{=}\AgdaSpace{}%
\AgdaFunction{cong}\AgdaSpace{}%
\AgdaInductiveConstructor{box}\AgdaSpace{}%
\AgdaSymbol{(}\AgdaFunction{cong}\AgdaSpace{}%
\AgdaInductiveConstructor{ƛb}\AgdaSpace{}%
\AgdaSymbol{(}\AgdaFunction{pres-rename}\AgdaSpace{}%
\AgdaSymbol{(}\AgdaFunction{calc.ext}\AgdaSpace{}%
\AgdaBound{f'}\AgdaSymbol{)}\AgdaSpace{}%
\AgdaSymbol{(}\AgdaFunction{boxCalc.extb}\AgdaSpace{}%
\AgdaBound{fb'}\AgdaSymbol{)}\AgdaSpace{}%
\AgdaSymbol{(}\AgdaFunction{pres-ext}\AgdaSpace{}%
\AgdaBound{f'}\AgdaSpace{}%
\AgdaBound{fb'}\AgdaSpace{}%
\AgdaBound{h}\AgdaSymbol{)}\AgdaSpace{}%
\AgdaBound{t}\AgdaSymbol{))}\<%
\\
\>[0]\AgdaFunction{pres-rename}\AgdaSpace{}%
\AgdaBound{f'}\AgdaSpace{}%
\AgdaBound{fb'}\AgdaSpace{}%
\AgdaBound{h}\AgdaSpace{}%
\AgdaSymbol{(}\AgdaBound{t}\AgdaSpace{}%
\AgdaOperator{\AgdaInductiveConstructor{·}}\AgdaSpace{}%
\AgdaBound{t₁}\AgdaSymbol{)}\AgdaSpace{}%
\AgdaSymbol{=}\<%
\\
\>[0][@{}l@{\AgdaIndent{0}}]%
\>[2]\AgdaOperator{\AgdaFunction{beginEq}}\<%
\\
\>[2][@{}l@{\AgdaIndent{0}}]%
\>[4]\AgdaInductiveConstructor{ƛb}\AgdaSpace{}%
\AgdaSymbol{(}\AgdaInductiveConstructor{ε}\AgdaSpace{}%
\AgdaInductiveConstructor{Zb}\AgdaSpace{}%
\AgdaOperator{\AgdaInductiveConstructor{·b}}\AgdaSpace{}%
\AgdaFunction{renameb}\AgdaSpace{}%
\AgdaSymbol{(}\AgdaFunction{extb}\AgdaSpace{}%
\AgdaBound{fb'}\AgdaSymbol{)}\AgdaSpace{}%
\AgdaSymbol{(}\AgdaFunction{renameb}\AgdaSpace{}%
\AgdaInductiveConstructor{Sb}\AgdaSpace{}%
\AgdaSymbol{(}\AgdaFunction{embedTerm}\AgdaSpace{}%
\AgdaBound{t₁}\AgdaSymbol{)))}\AgdaSpace{}%
\AgdaOperator{\AgdaInductiveConstructor{·b}}\AgdaSpace{}%
\AgdaFunction{renameb}\AgdaSpace{}%
\AgdaBound{fb'}\AgdaSpace{}%
\AgdaSymbol{(}\AgdaFunction{embedTerm}\AgdaSpace{}%
\AgdaBound{t}\AgdaSymbol{)}\<%
\\
%
\>[2]\AgdaFunction{≡⟨}\AgdaSpace{}%
\AgdaFunction{cong₂}\AgdaSpace{}%
\AgdaOperator{\AgdaInductiveConstructor{\AgdaUnderscore{}·b\AgdaUnderscore{}}}\AgdaSpace{}%
\AgdaInductiveConstructor{refl}\AgdaSpace{}%
\AgdaSymbol{(}\AgdaFunction{pres-rename}\AgdaSpace{}%
\AgdaBound{f'}\AgdaSpace{}%
\AgdaBound{fb'}\AgdaSpace{}%
\AgdaBound{h}\AgdaSpace{}%
\AgdaBound{t}\AgdaSymbol{)}\AgdaSpace{}%
\AgdaFunction{⟩}\<%
\\
\>[2][@{}l@{\AgdaIndent{0}}]%
\>[4]\AgdaInductiveConstructor{ƛb}\AgdaSpace{}%
\AgdaSymbol{(}\AgdaInductiveConstructor{ε}\AgdaSpace{}%
\AgdaInductiveConstructor{Zb}\AgdaSpace{}%
\AgdaOperator{\AgdaInductiveConstructor{·b}}\AgdaSpace{}%
\AgdaFunction{renameb}\AgdaSpace{}%
\AgdaSymbol{(}\AgdaFunction{extb}\AgdaSpace{}%
\AgdaBound{fb'}\AgdaSymbol{)}\AgdaSpace{}%
\AgdaSymbol{(}\AgdaFunction{renameb}\AgdaSpace{}%
\AgdaInductiveConstructor{Sb}\AgdaSpace{}%
\AgdaSymbol{(}\AgdaFunction{embedTerm}\AgdaSpace{}%
\AgdaBound{t₁}\AgdaSymbol{)))}\AgdaSpace{}%
\AgdaOperator{\AgdaInductiveConstructor{·b}}\AgdaSpace{}%
\AgdaFunction{embedTerm}\AgdaSpace{}%
\AgdaSymbol{(}\AgdaFunction{rename}\AgdaSpace{}%
\AgdaBound{f'}\AgdaSpace{}%
\AgdaBound{t}\AgdaSymbol{)}\<%
\\
%
\>[2]\AgdaFunction{≡⟨}\AgdaSpace{}%
\AgdaFunction{cong₂}\AgdaSpace{}%
\AgdaOperator{\AgdaInductiveConstructor{\AgdaUnderscore{}·b\AgdaUnderscore{}}}\AgdaSpace{}%
\AgdaSymbol{(}\AgdaFunction{cong}\AgdaSpace{}%
\AgdaInductiveConstructor{ƛb}\AgdaSpace{}%
\AgdaSymbol{(}\AgdaFunction{cong₂}\AgdaSpace{}%
\AgdaOperator{\AgdaInductiveConstructor{\AgdaUnderscore{}·b\AgdaUnderscore{}}}\AgdaSpace{}%
\AgdaInductiveConstructor{refl}\AgdaSpace{}%
\AgdaSymbol{(}\AgdaFunction{rename-extb-Sb-commute}\AgdaSpace{}%
\AgdaBound{fb'}\AgdaSpace{}%
\AgdaSymbol{(}\AgdaFunction{embedTerm}\AgdaSpace{}%
\AgdaBound{t₁}\AgdaSymbol{))))}\AgdaSpace{}%
\AgdaInductiveConstructor{refl}\AgdaSpace{}%
\AgdaFunction{⟩}\<%
\\
\>[2][@{}l@{\AgdaIndent{0}}]%
\>[4]\AgdaInductiveConstructor{ƛb}\AgdaSpace{}%
\AgdaSymbol{(}\AgdaInductiveConstructor{ε}\AgdaSpace{}%
\AgdaInductiveConstructor{Zb}\AgdaSpace{}%
\AgdaOperator{\AgdaInductiveConstructor{·b}}\AgdaSpace{}%
\AgdaFunction{renameb}\AgdaSpace{}%
\AgdaInductiveConstructor{Sb}\AgdaSpace{}%
\AgdaSymbol{(}\AgdaFunction{renameb}\AgdaSpace{}%
\AgdaBound{fb'}\AgdaSpace{}%
\AgdaSymbol{(}\AgdaFunction{embedTerm}\AgdaSpace{}%
\AgdaBound{t₁}\AgdaSymbol{)))}\AgdaSpace{}%
\AgdaOperator{\AgdaInductiveConstructor{·b}}\AgdaSpace{}%
\AgdaFunction{embedTerm}\AgdaSpace{}%
\AgdaSymbol{(}\AgdaFunction{rename}\AgdaSpace{}%
\AgdaBound{f'}\AgdaSpace{}%
\AgdaBound{t}\AgdaSymbol{)}\<%
\\
%
\>[2]\AgdaFunction{≡⟨}\AgdaSpace{}%
\AgdaFunction{cong₂}\AgdaSpace{}%
\AgdaOperator{\AgdaInductiveConstructor{\AgdaUnderscore{}·b\AgdaUnderscore{}}}\AgdaSpace{}%
\AgdaSymbol{(}\AgdaFunction{cong}\AgdaSpace{}%
\AgdaInductiveConstructor{ƛb}\AgdaSpace{}%
\AgdaSymbol{(}\AgdaFunction{cong₂}\AgdaSpace{}%
\AgdaOperator{\AgdaInductiveConstructor{\AgdaUnderscore{}·b\AgdaUnderscore{}}}\AgdaSpace{}%
\AgdaInductiveConstructor{refl}\AgdaSpace{}%
\AgdaSymbol{(}\AgdaFunction{cong}\AgdaSpace{}%
\AgdaSymbol{(}\AgdaFunction{renameb}\AgdaSpace{}%
\AgdaInductiveConstructor{Sb}\AgdaSymbol{)}\AgdaSpace{}%
\AgdaSymbol{(}\AgdaFunction{pres-rename}\AgdaSpace{}%
\AgdaBound{f'}\AgdaSpace{}%
\AgdaBound{fb'}\AgdaSpace{}%
\AgdaBound{h}\AgdaSpace{}%
\AgdaBound{t₁}\AgdaSymbol{))))}\AgdaSpace{}%
\AgdaInductiveConstructor{refl}\AgdaSpace{}%
\AgdaFunction{⟩}\<%
\\
\>[2][@{}l@{\AgdaIndent{0}}]%
\>[4]\AgdaInductiveConstructor{ƛb}\AgdaSpace{}%
\AgdaSymbol{(}\AgdaInductiveConstructor{ε}\AgdaSpace{}%
\AgdaInductiveConstructor{Zb}\AgdaSpace{}%
\AgdaOperator{\AgdaInductiveConstructor{·b}}\AgdaSpace{}%
\AgdaFunction{renameb}\AgdaSpace{}%
\AgdaInductiveConstructor{Sb}\AgdaSpace{}%
\AgdaSymbol{(}\AgdaFunction{embedTerm}\AgdaSpace{}%
\AgdaSymbol{(}\AgdaFunction{rename}\AgdaSpace{}%
\AgdaBound{f'}\AgdaSpace{}%
\AgdaBound{t₁}\AgdaSymbol{)))}\AgdaSpace{}%
\AgdaOperator{\AgdaInductiveConstructor{·b}}\AgdaSpace{}%
\AgdaFunction{embedTerm}\AgdaSpace{}%
\AgdaSymbol{(}\AgdaFunction{rename}\AgdaSpace{}%
\AgdaBound{f'}\AgdaSpace{}%
\AgdaBound{t}\AgdaSymbol{)}\<%
\\
%
\>[2]\AgdaOperator{\AgdaFunction{∎eq}}\<%
\\
%
\\[\AgdaEmptyExtraSkip]%
\>[0]\AgdaComment{--\ Preservation\ and\ Reflection\ of\ }\<%
\\
\>[0]\AgdaComment{--\ exts\ of\ Gödel's\ translation}\<%
\\
\>[0]\AgdaFunction{pres-exts}\AgdaSpace{}%
\AgdaSymbol{:}\AgdaSpace{}%
\AgdaSymbol{∀}\AgdaSpace{}%
\AgdaSymbol{\{}\AgdaBound{Γ}\AgdaSpace{}%
\AgdaBound{Δ}\AgdaSymbol{\}}\<%
\\
\>[0][@{}l@{\AgdaIndent{0}}]%
\>[2]\AgdaSymbol{→}\AgdaSpace{}%
\AgdaSymbol{(}\AgdaBound{fb'}\AgdaSpace{}%
\AgdaSymbol{:}\AgdaSpace{}%
\AgdaSymbol{∀}\AgdaSpace{}%
\AgdaSymbol{\{}\AgdaBound{Ab}\AgdaSymbol{\}}\AgdaSpace{}%
\AgdaSymbol{→}\AgdaSpace{}%
\AgdaInductiveConstructor{□}\AgdaSpace{}%
\AgdaBound{Ab}\AgdaSpace{}%
\AgdaOperator{\AgdaDatatype{∈b}}\AgdaSpace{}%
\AgdaFunction{embedContext}\AgdaSpace{}%
\AgdaBound{Γ}\AgdaSpace{}%
\AgdaSymbol{→}\AgdaSpace{}%
\AgdaFunction{embedContext}\AgdaSpace{}%
\AgdaBound{Δ}\AgdaSpace{}%
\AgdaOperator{\AgdaDatatype{⊢b}}\AgdaSpace{}%
\AgdaBound{Ab}\AgdaSymbol{)}\<%
\\
%
\>[2]\AgdaSymbol{→}\AgdaSpace{}%
\AgdaSymbol{(}\AgdaBound{f'}\AgdaSpace{}%
\AgdaSymbol{:}\AgdaSpace{}%
\AgdaSymbol{∀}\AgdaSpace{}%
\AgdaSymbol{\{}\AgdaBound{A}\AgdaSymbol{\}}\AgdaSpace{}%
\AgdaSymbol{→}\AgdaSpace{}%
\AgdaBound{A}\AgdaSpace{}%
\AgdaOperator{\AgdaDatatype{∈}}\AgdaSpace{}%
\AgdaBound{Γ}\AgdaSpace{}%
\AgdaSymbol{→}\AgdaSpace{}%
\AgdaBound{Δ}\AgdaSpace{}%
\AgdaOperator{\AgdaDatatype{⊢}}\AgdaSpace{}%
\AgdaBound{A}\AgdaSymbol{)}\<%
\\
%
\>[2]\AgdaSymbol{→}\AgdaSpace{}%
\AgdaSymbol{(\{}\AgdaBound{A'}\AgdaSpace{}%
\AgdaSymbol{:}\AgdaSpace{}%
\AgdaDatatype{Type}\AgdaSymbol{\}}\AgdaSpace{}%
\AgdaSymbol{(}\AgdaBound{x}\AgdaSpace{}%
\AgdaSymbol{:}\AgdaSpace{}%
\AgdaBound{A'}\AgdaSpace{}%
\AgdaOperator{\AgdaDatatype{∈}}\AgdaSpace{}%
\AgdaBound{Γ}\AgdaSymbol{)}\AgdaSpace{}%
\AgdaSymbol{→}\AgdaSpace{}%
\AgdaInductiveConstructor{box}\AgdaSpace{}%
\AgdaSymbol{(}\AgdaBound{fb'}\AgdaSpace{}%
\AgdaSymbol{(}\AgdaFunction{embedLookup}\AgdaSpace{}%
\AgdaBound{x}\AgdaSymbol{))}\AgdaSpace{}%
\AgdaOperator{\AgdaDatatype{≡}}\AgdaSpace{}%
\AgdaFunction{embedTerm}\AgdaSpace{}%
\AgdaSymbol{(}\AgdaBound{f'}\AgdaSpace{}%
\AgdaBound{x}\AgdaSymbol{))}\<%
\\
%
\>[2]\AgdaSymbol{→}\AgdaSpace{}%
\AgdaSymbol{\{}\AgdaBound{A'}\AgdaSpace{}%
\AgdaBound{A}\AgdaSpace{}%
\AgdaSymbol{:}\AgdaSpace{}%
\AgdaDatatype{Type}\AgdaSymbol{\}}\AgdaSpace{}%
\AgdaSymbol{(}\AgdaBound{x}\AgdaSpace{}%
\AgdaSymbol{:}\AgdaSpace{}%
\AgdaBound{A'}\AgdaSpace{}%
\AgdaOperator{\AgdaDatatype{∈}}\AgdaSpace{}%
\AgdaBound{Γ}\AgdaSpace{}%
\AgdaOperator{\AgdaInductiveConstructor{,}}\AgdaSpace{}%
\AgdaBound{A}\AgdaSymbol{)}\AgdaSpace{}%
\AgdaSymbol{→}\AgdaSpace{}%
\AgdaInductiveConstructor{box}\AgdaSpace{}%
\AgdaSymbol{(}\AgdaFunction{extsb}\AgdaSpace{}%
\AgdaBound{fb'}\AgdaSpace{}%
\AgdaSymbol{(}\AgdaFunction{embedLookup}\AgdaSpace{}%
\AgdaBound{x}\AgdaSymbol{))}\AgdaSpace{}%
\AgdaOperator{\AgdaDatatype{≡}}\AgdaSpace{}%
\AgdaFunction{embedTerm}\AgdaSpace{}%
\AgdaSymbol{(}\AgdaFunction{exts}\AgdaSpace{}%
\AgdaBound{f'}\AgdaSpace{}%
\AgdaBound{x}\AgdaSymbol{)}\<%
\\
\>[0]\AgdaFunction{pres-exts}\AgdaSpace{}%
\AgdaBound{fb'}\AgdaSpace{}%
\AgdaBound{f'}\AgdaSpace{}%
\AgdaBound{h}\AgdaSpace{}%
\AgdaInductiveConstructor{Z}\AgdaSpace{}%
\AgdaSymbol{=}\AgdaSpace{}%
\AgdaInductiveConstructor{refl}\<%
\\
\>[0]\AgdaFunction{pres-exts}\AgdaSpace{}%
\AgdaBound{fb'}\AgdaSpace{}%
\AgdaBound{f'}\AgdaSpace{}%
\AgdaBound{h}\AgdaSpace{}%
\AgdaSymbol{(}\AgdaInductiveConstructor{S}\AgdaSpace{}%
\AgdaBound{x}\AgdaSymbol{)}\AgdaSpace{}%
\AgdaSymbol{=}\<%
\\
\>[0][@{}l@{\AgdaIndent{0}}]%
\>[2]\AgdaOperator{\AgdaFunction{beginEq}}\<%
\\
\>[2][@{}l@{\AgdaIndent{0}}]%
\>[4]\AgdaInductiveConstructor{box}\AgdaSpace{}%
\AgdaSymbol{(}\AgdaFunction{renameb}\AgdaSpace{}%
\AgdaInductiveConstructor{Sb}\AgdaSpace{}%
\AgdaSymbol{(}\AgdaBound{fb'}\AgdaSpace{}%
\AgdaSymbol{(}\AgdaFunction{embedLookup}\AgdaSpace{}%
\AgdaBound{x}\AgdaSymbol{)))}\<%
\\
%
\>[2]\AgdaFunction{≡⟨}\AgdaSpace{}%
\AgdaInductiveConstructor{refl}\AgdaSpace{}%
\AgdaFunction{⟩}\<%
\\
\>[2][@{}l@{\AgdaIndent{0}}]%
\>[4]\AgdaFunction{renameb}\AgdaSpace{}%
\AgdaInductiveConstructor{Sb}\AgdaSpace{}%
\AgdaSymbol{(}\AgdaInductiveConstructor{box}\AgdaSpace{}%
\AgdaSymbol{(}\AgdaBound{fb'}\AgdaSpace{}%
\AgdaSymbol{(}\AgdaFunction{embedLookup}\AgdaSpace{}%
\AgdaBound{x}\AgdaSymbol{)))}\<%
\\
%
\>[2]\AgdaFunction{≡⟨}\AgdaSpace{}%
\AgdaFunction{cong}\AgdaSpace{}%
\AgdaSymbol{(}\AgdaFunction{renameb}\AgdaSpace{}%
\AgdaInductiveConstructor{Sb}\AgdaSymbol{)}\AgdaSpace{}%
\AgdaSymbol{(}\AgdaBound{h}\AgdaSpace{}%
\AgdaBound{x}\AgdaSymbol{)}\AgdaSpace{}%
\AgdaFunction{⟩}\<%
\\
\>[2][@{}l@{\AgdaIndent{0}}]%
\>[4]\AgdaFunction{renameb}\AgdaSpace{}%
\AgdaInductiveConstructor{Sb}\AgdaSpace{}%
\AgdaSymbol{(}\AgdaFunction{embedTerm}\AgdaSpace{}%
\AgdaSymbol{(}\AgdaBound{f'}\AgdaSpace{}%
\AgdaBound{x}\AgdaSymbol{))}\<%
\\
%
\>[2]\AgdaFunction{≡⟨}\AgdaSpace{}%
\AgdaFunction{pres-rename}\AgdaSpace{}%
\AgdaInductiveConstructor{S}\AgdaSpace{}%
\AgdaInductiveConstructor{Sb}\AgdaSpace{}%
\AgdaSymbol{(λ}\AgdaSpace{}%
\AgdaBound{x₁}\AgdaSpace{}%
\AgdaSymbol{→}\AgdaSpace{}%
\AgdaInductiveConstructor{refl}\AgdaSymbol{)}\AgdaSpace{}%
\AgdaSymbol{(}\AgdaBound{f'}\AgdaSpace{}%
\AgdaBound{x}\AgdaSymbol{)}\AgdaSpace{}%
\AgdaFunction{⟩}\<%
\\
\>[2][@{}l@{\AgdaIndent{0}}]%
\>[4]\AgdaFunction{embedTerm}\AgdaSpace{}%
\AgdaSymbol{(}\AgdaFunction{rename}\AgdaSpace{}%
\AgdaInductiveConstructor{S}\AgdaSpace{}%
\AgdaSymbol{(}\AgdaBound{f'}\AgdaSpace{}%
\AgdaBound{x}\AgdaSymbol{))}\<%
\\
%
\>[2]\AgdaOperator{\AgdaFunction{∎eq}}\<%
\\
%
\\[\AgdaEmptyExtraSkip]%
\>[0]\AgdaComment{--\ Extsb\ preserves\ pointwise\ equality}\<%
\\
\>[0]\AgdaComment{--\ of\ renaming\ functions}\<%
\\
\>[0]\AgdaFunction{extsb-cong}\AgdaSpace{}%
\AgdaSymbol{:}\AgdaSpace{}%
\AgdaSymbol{∀}\AgdaSpace{}%
\AgdaSymbol{\{}\AgdaBound{Γ₁}\AgdaSpace{}%
\AgdaBound{Γ₂}\AgdaSpace{}%
\AgdaBound{A}\AgdaSpace{}%
\AgdaBound{A'}\AgdaSymbol{\}}\<%
\\
\>[0][@{}l@{\AgdaIndent{0}}]%
\>[2]\AgdaSymbol{→}\AgdaSpace{}%
\AgdaSymbol{\{}\AgdaBound{fb'}\AgdaSpace{}%
\AgdaSymbol{:}\AgdaSpace{}%
\AgdaSymbol{∀}\AgdaSpace{}%
\AgdaSymbol{\{}\AgdaBound{Ab}\AgdaSymbol{\}}\AgdaSpace{}%
\AgdaSymbol{→}\AgdaSpace{}%
\AgdaInductiveConstructor{□}\AgdaSpace{}%
\AgdaBound{Ab}\AgdaSpace{}%
\AgdaOperator{\AgdaDatatype{∈b}}\AgdaSpace{}%
\AgdaBound{Γ₁}\AgdaSpace{}%
\AgdaSymbol{→}\AgdaSpace{}%
\AgdaBound{Γ₂}\AgdaSpace{}%
\AgdaOperator{\AgdaDatatype{⊢b}}\AgdaSpace{}%
\AgdaBound{Ab}\AgdaSymbol{\}}\<%
\\
%
\>[2]\AgdaSymbol{→}\AgdaSpace{}%
\AgdaSymbol{\{}\AgdaBound{fb''}\AgdaSpace{}%
\AgdaSymbol{:}\AgdaSpace{}%
\AgdaSymbol{∀}\AgdaSpace{}%
\AgdaSymbol{\{}\AgdaBound{Ab}\AgdaSymbol{\}}\AgdaSpace{}%
\AgdaSymbol{→}\AgdaSpace{}%
\AgdaInductiveConstructor{□}\AgdaSpace{}%
\AgdaBound{Ab}\AgdaSpace{}%
\AgdaOperator{\AgdaDatatype{∈b}}\AgdaSpace{}%
\AgdaBound{Γ₁}\AgdaSpace{}%
\AgdaSymbol{→}\AgdaSpace{}%
\AgdaBound{Γ₂}\AgdaSpace{}%
\AgdaOperator{\AgdaDatatype{⊢b}}\AgdaSpace{}%
\AgdaBound{Ab}\AgdaSymbol{\}}\<%
\\
%
\>[2]\AgdaSymbol{→}\AgdaSpace{}%
\AgdaSymbol{(}\AgdaBound{h}\AgdaSpace{}%
\AgdaSymbol{:}\AgdaSpace{}%
\AgdaSymbol{∀}\AgdaSpace{}%
\AgdaSymbol{\{}\AgdaBound{Ab}\AgdaSymbol{\}}\AgdaSpace{}%
\AgdaSymbol{→}\AgdaSpace{}%
\AgdaSymbol{(}\AgdaBound{x}\AgdaSpace{}%
\AgdaSymbol{:}\AgdaSpace{}%
\AgdaInductiveConstructor{□}\AgdaSpace{}%
\AgdaBound{Ab}\AgdaSpace{}%
\AgdaOperator{\AgdaDatatype{∈b}}\AgdaSpace{}%
\AgdaBound{Γ₁}\AgdaSymbol{)}\AgdaSpace{}%
\AgdaSymbol{→}\AgdaSpace{}%
\AgdaBound{fb'}\AgdaSpace{}%
\AgdaBound{x}\AgdaSpace{}%
\AgdaOperator{\AgdaDatatype{≡}}\AgdaSpace{}%
\AgdaBound{fb''}\AgdaSpace{}%
\AgdaBound{x}\AgdaSymbol{)}\<%
\\
%
\>[2]\AgdaSymbol{→}\AgdaSpace{}%
\AgdaSymbol{(}\AgdaBound{x}\AgdaSpace{}%
\AgdaSymbol{:}\AgdaSpace{}%
\AgdaInductiveConstructor{□}\AgdaSpace{}%
\AgdaBound{A}\AgdaSpace{}%
\AgdaOperator{\AgdaDatatype{∈b}}\AgdaSpace{}%
\AgdaBound{Γ₁}\AgdaSpace{}%
\AgdaOperator{\AgdaInductiveConstructor{,b}}\AgdaSpace{}%
\AgdaInductiveConstructor{□}\AgdaSpace{}%
\AgdaBound{A'}\AgdaSymbol{)}\<%
\\
%
\>[2]\AgdaSymbol{→}\AgdaSpace{}%
\AgdaFunction{extsb}\AgdaSpace{}%
\AgdaBound{fb'}\AgdaSpace{}%
\AgdaBound{x}\AgdaSpace{}%
\AgdaOperator{\AgdaDatatype{≡}}\AgdaSpace{}%
\AgdaFunction{extsb}\AgdaSpace{}%
\AgdaBound{fb''}\AgdaSpace{}%
\AgdaBound{x}\<%
\\
\>[0]\AgdaFunction{extsb-cong}\AgdaSpace{}%
\AgdaBound{h}\AgdaSpace{}%
\AgdaInductiveConstructor{Zb}\AgdaSpace{}%
\AgdaSymbol{=}\AgdaSpace{}%
\AgdaInductiveConstructor{refl}\<%
\\
\>[0]\AgdaFunction{extsb-cong}\AgdaSpace{}%
\AgdaBound{h}\AgdaSpace{}%
\AgdaSymbol{(}\AgdaInductiveConstructor{Sb}\AgdaSpace{}%
\AgdaBound{x}\AgdaSymbol{)}\AgdaSpace{}%
\AgdaSymbol{=}\AgdaSpace{}%
\AgdaFunction{cong}\AgdaSpace{}%
\AgdaSymbol{(}\AgdaFunction{renameb}\AgdaSpace{}%
\AgdaInductiveConstructor{Sb}\AgdaSymbol{)}\AgdaSpace{}%
\AgdaSymbol{(}\AgdaBound{h}\AgdaSpace{}%
\AgdaBound{x}\AgdaSymbol{)}\<%
\\
%
\\[\AgdaEmptyExtraSkip]%
\>[0]\AgdaComment{--\ substb\ preserves\ pointwise\ equality}\<%
\\
\>[0]\AgdaComment{--\ of\ renaming\ functions}\<%
\\
\>[0]\AgdaFunction{substb-cong}\AgdaSpace{}%
\AgdaSymbol{:}\AgdaSpace{}%
\AgdaSymbol{∀}\AgdaSpace{}%
\AgdaSymbol{\{}\AgdaBound{Γ₁}\AgdaSpace{}%
\AgdaBound{Γ₂}\AgdaSpace{}%
\AgdaBound{A}\AgdaSymbol{\}}\<%
\\
\>[0][@{}l@{\AgdaIndent{0}}]%
\>[2]\AgdaSymbol{→}\AgdaSpace{}%
\AgdaSymbol{\{}\AgdaBound{fb'}\AgdaSpace{}%
\AgdaSymbol{:}\AgdaSpace{}%
\AgdaSymbol{∀}\AgdaSpace{}%
\AgdaSymbol{\{}\AgdaBound{Ab}\AgdaSymbol{\}}\AgdaSpace{}%
\AgdaSymbol{→}\AgdaSpace{}%
\AgdaInductiveConstructor{□}\AgdaSpace{}%
\AgdaBound{Ab}\AgdaSpace{}%
\AgdaOperator{\AgdaDatatype{∈b}}\AgdaSpace{}%
\AgdaBound{Γ₁}\AgdaSpace{}%
\AgdaSymbol{→}\AgdaSpace{}%
\AgdaBound{Γ₂}\AgdaSpace{}%
\AgdaOperator{\AgdaDatatype{⊢b}}\AgdaSpace{}%
\AgdaBound{Ab}\AgdaSymbol{\}}\<%
\\
%
\>[2]\AgdaSymbol{→}\AgdaSpace{}%
\AgdaSymbol{\{}\AgdaBound{fb''}\AgdaSpace{}%
\AgdaSymbol{:}\AgdaSpace{}%
\AgdaSymbol{∀}\AgdaSpace{}%
\AgdaSymbol{\{}\AgdaBound{Ab}\AgdaSymbol{\}}\AgdaSpace{}%
\AgdaSymbol{→}\AgdaSpace{}%
\AgdaInductiveConstructor{□}\AgdaSpace{}%
\AgdaBound{Ab}\AgdaSpace{}%
\AgdaOperator{\AgdaDatatype{∈b}}\AgdaSpace{}%
\AgdaBound{Γ₁}\AgdaSpace{}%
\AgdaSymbol{→}\AgdaSpace{}%
\AgdaBound{Γ₂}\AgdaSpace{}%
\AgdaOperator{\AgdaDatatype{⊢b}}\AgdaSpace{}%
\AgdaBound{Ab}\AgdaSymbol{\}}\<%
\\
%
\>[2]\AgdaSymbol{→}\AgdaSpace{}%
\AgdaSymbol{(}\AgdaBound{h}\AgdaSpace{}%
\AgdaSymbol{:}\AgdaSpace{}%
\AgdaSymbol{∀}\AgdaSpace{}%
\AgdaSymbol{\{}\AgdaBound{Ab}\AgdaSymbol{\}}\AgdaSpace{}%
\AgdaSymbol{→}\AgdaSpace{}%
\AgdaSymbol{(}\AgdaBound{x}\AgdaSpace{}%
\AgdaSymbol{:}\AgdaSpace{}%
\AgdaInductiveConstructor{□}\AgdaSpace{}%
\AgdaBound{Ab}\AgdaSpace{}%
\AgdaOperator{\AgdaDatatype{∈b}}\AgdaSpace{}%
\AgdaBound{Γ₁}\AgdaSymbol{)}\AgdaSpace{}%
\AgdaSymbol{→}\AgdaSpace{}%
\AgdaBound{fb'}\AgdaSpace{}%
\AgdaBound{x}\AgdaSpace{}%
\AgdaOperator{\AgdaDatatype{≡}}\AgdaSpace{}%
\AgdaBound{fb''}\AgdaSpace{}%
\AgdaBound{x}\AgdaSymbol{)}\<%
\\
%
\>[2]\AgdaSymbol{→}\AgdaSpace{}%
\AgdaSymbol{(}\AgdaBound{t}\AgdaSpace{}%
\AgdaSymbol{:}\AgdaSpace{}%
\AgdaBound{Γ₁}\AgdaSpace{}%
\AgdaOperator{\AgdaDatatype{⊢b}}\AgdaSpace{}%
\AgdaBound{A}\AgdaSymbol{)}\<%
\\
%
\>[2]\AgdaSymbol{→}\AgdaSpace{}%
\AgdaFunction{substb}\AgdaSpace{}%
\AgdaBound{fb'}\AgdaSpace{}%
\AgdaBound{t}\AgdaSpace{}%
\AgdaOperator{\AgdaDatatype{≡}}\AgdaSpace{}%
\AgdaFunction{substb}\AgdaSpace{}%
\AgdaBound{fb''}\AgdaSpace{}%
\AgdaBound{t}\<%
\\
\>[0]\AgdaFunction{substb-cong}\AgdaSpace{}%
\AgdaBound{h}\AgdaSpace{}%
\AgdaSymbol{(}\AgdaInductiveConstructor{ε}\AgdaSpace{}%
\AgdaBound{x}\AgdaSymbol{)}\AgdaSpace{}%
\AgdaSymbol{=}\AgdaSpace{}%
\AgdaBound{h}\AgdaSpace{}%
\AgdaBound{x}\<%
\\
\>[0]\AgdaFunction{substb-cong}\AgdaSpace{}%
\AgdaBound{h}\AgdaSpace{}%
\AgdaSymbol{(}\AgdaInductiveConstructor{box}\AgdaSpace{}%
\AgdaBound{t}\AgdaSymbol{)}\AgdaSpace{}%
\AgdaSymbol{=}\AgdaSpace{}%
\AgdaFunction{cong}\AgdaSpace{}%
\AgdaInductiveConstructor{box}\AgdaSpace{}%
\AgdaSymbol{(}\AgdaFunction{substb-cong}\AgdaSpace{}%
\AgdaBound{h}\AgdaSpace{}%
\AgdaBound{t}\AgdaSymbol{)}\<%
\\
\>[0]\AgdaFunction{substb-cong}\AgdaSpace{}%
\AgdaBound{h}\AgdaSpace{}%
\AgdaSymbol{(}\AgdaInductiveConstructor{ƛb}\AgdaSpace{}%
\AgdaBound{t}\AgdaSymbol{)}\AgdaSpace{}%
\AgdaSymbol{=}\AgdaSpace{}%
\AgdaFunction{cong}\AgdaSpace{}%
\AgdaInductiveConstructor{ƛb}\AgdaSpace{}%
\AgdaSymbol{(}\AgdaFunction{substb-cong}\AgdaSpace{}%
\AgdaSymbol{(}\AgdaFunction{extsb-cong}\AgdaSpace{}%
\AgdaBound{h}\AgdaSymbol{)}\AgdaSpace{}%
\AgdaBound{t}\AgdaSymbol{)}\<%
\\
\>[0]\AgdaFunction{substb-cong}\AgdaSpace{}%
\AgdaBound{h}\AgdaSpace{}%
\AgdaSymbol{(}\AgdaBound{t}\AgdaSpace{}%
\AgdaOperator{\AgdaInductiveConstructor{·b}}\AgdaSpace{}%
\AgdaBound{t₁}\AgdaSymbol{)}\AgdaSpace{}%
\AgdaSymbol{=}\AgdaSpace{}%
\AgdaFunction{cong₂}\AgdaSpace{}%
\AgdaOperator{\AgdaInductiveConstructor{\AgdaUnderscore{}·b\AgdaUnderscore{}}}\AgdaSpace{}%
\AgdaSymbol{(}\AgdaFunction{substb-cong}\AgdaSpace{}%
\AgdaBound{h}\AgdaSpace{}%
\AgdaBound{t}\AgdaSymbol{)}\AgdaSpace{}%
\AgdaSymbol{(}\AgdaFunction{substb-cong}\AgdaSpace{}%
\AgdaBound{h}\AgdaSpace{}%
\AgdaBound{t₁}\AgdaSymbol{)}\<%
\\
%
\\[\AgdaEmptyExtraSkip]%
\>[0]\AgdaComment{--\ Commutativity\ of\ renameb\ and\ extsb:}\<%
\\
\>[0]\AgdaFunction{extsb-rename-commute}\AgdaSpace{}%
\AgdaSymbol{:}\AgdaSpace{}%
\AgdaSymbol{∀}\AgdaSpace{}%
\AgdaSymbol{\{}\AgdaBound{Γ₁}\AgdaSpace{}%
\AgdaBound{Γ₂}\AgdaSpace{}%
\AgdaBound{Γ₃}\AgdaSpace{}%
\AgdaBound{A}\AgdaSpace{}%
\AgdaBound{B}\AgdaSymbol{\}}\<%
\\
\>[0][@{}l@{\AgdaIndent{0}}]%
\>[2]\AgdaSymbol{→}\AgdaSpace{}%
\AgdaSymbol{(}\AgdaBound{fb'}\AgdaSpace{}%
\AgdaSymbol{:}\AgdaSpace{}%
\AgdaSymbol{∀}\AgdaSpace{}%
\AgdaSymbol{\{}\AgdaBound{Ab}\AgdaSymbol{\}}\AgdaSpace{}%
\AgdaSymbol{→}\AgdaSpace{}%
\AgdaBound{Ab}\AgdaSpace{}%
\AgdaOperator{\AgdaDatatype{∈b}}\AgdaSpace{}%
\AgdaBound{Γ₂}\AgdaSpace{}%
\AgdaSymbol{→}\AgdaSpace{}%
\AgdaBound{Ab}\AgdaSpace{}%
\AgdaOperator{\AgdaDatatype{∈b}}\AgdaSpace{}%
\AgdaBound{Γ₃}\AgdaSymbol{)}\<%
\\
%
\>[2]\AgdaSymbol{→}\AgdaSpace{}%
\AgdaSymbol{(}\AgdaBound{fb''}\AgdaSpace{}%
\AgdaSymbol{:}\AgdaSpace{}%
\AgdaSymbol{∀}\AgdaSpace{}%
\AgdaSymbol{\{}\AgdaBound{Ab}\AgdaSymbol{\}}\AgdaSpace{}%
\AgdaSymbol{→}\AgdaSpace{}%
\AgdaInductiveConstructor{□}\AgdaSpace{}%
\AgdaBound{Ab}\AgdaSpace{}%
\AgdaOperator{\AgdaDatatype{∈b}}\AgdaSpace{}%
\AgdaBound{Γ₁}\AgdaSpace{}%
\AgdaSymbol{→}\AgdaSpace{}%
\AgdaBound{Γ₂}\AgdaSpace{}%
\AgdaOperator{\AgdaDatatype{⊢b}}\AgdaSpace{}%
\AgdaBound{Ab}\AgdaSymbol{)}\<%
\\
%
\>[2]\AgdaSymbol{→}\AgdaSpace{}%
\AgdaSymbol{(}\AgdaBound{x}\AgdaSpace{}%
\AgdaSymbol{:}\AgdaSpace{}%
\AgdaInductiveConstructor{□}\AgdaSpace{}%
\AgdaBound{A}\AgdaSpace{}%
\AgdaOperator{\AgdaDatatype{∈b}}\AgdaSpace{}%
\AgdaBound{Γ₁}\AgdaSpace{}%
\AgdaOperator{\AgdaInductiveConstructor{,b}}\AgdaSpace{}%
\AgdaBound{B}\AgdaSymbol{)}\<%
\\
%
\>[2]\AgdaSymbol{→}%
\>[1170I]\AgdaFunction{extsb}\AgdaSpace{}%
\AgdaSymbol{(λ}\AgdaSpace{}%
\AgdaBound{x₁}\AgdaSpace{}%
\AgdaSymbol{→}\AgdaSpace{}%
\AgdaFunction{renameb}\AgdaSpace{}%
\AgdaBound{fb'}\AgdaSpace{}%
\AgdaSymbol{(}\AgdaBound{fb''}\AgdaSpace{}%
\AgdaBound{x₁}\AgdaSymbol{))}\AgdaSpace{}%
\AgdaBound{x}\AgdaSpace{}%
\AgdaOperator{\AgdaDatatype{≡}}\<%
\\
\>[.][@{}l@{}]\<[1170I]%
\>[4]\AgdaFunction{renameb}\AgdaSpace{}%
\AgdaSymbol{(}\AgdaFunction{extb}\AgdaSpace{}%
\AgdaBound{fb'}\AgdaSymbol{)}\AgdaSpace{}%
\AgdaSymbol{(}\AgdaFunction{extsb}\AgdaSpace{}%
\AgdaBound{fb''}\AgdaSpace{}%
\AgdaBound{x}\AgdaSymbol{)}\<%
\\
\>[0]\AgdaFunction{extsb-rename-commute}\AgdaSpace{}%
\AgdaBound{fb'}\AgdaSpace{}%
\AgdaBound{fb''}\AgdaSpace{}%
\AgdaInductiveConstructor{Zb}\AgdaSpace{}%
\AgdaSymbol{=}\AgdaSpace{}%
\AgdaInductiveConstructor{refl}\<%
\\
\>[0]\AgdaFunction{extsb-rename-commute}\AgdaSpace{}%
\AgdaBound{fb'}\AgdaSpace{}%
\AgdaBound{fb''}\AgdaSpace{}%
\AgdaSymbol{(}\AgdaInductiveConstructor{Sb}\AgdaSpace{}%
\AgdaBound{x}\AgdaSymbol{)}\AgdaSpace{}%
\AgdaSymbol{=}\AgdaSpace{}%
\AgdaFunction{sym}\AgdaSpace{}%
\AgdaSymbol{(}\AgdaFunction{rename-extb-Sb-commute}\AgdaSpace{}%
\AgdaBound{fb'}\AgdaSpace{}%
\AgdaSymbol{(}\AgdaBound{fb''}\AgdaSpace{}%
\AgdaBound{x}\AgdaSymbol{))}\<%
\\
%
\\[\AgdaEmptyExtraSkip]%
\>[0]\AgdaComment{--\ Commutativity\ of\ substb\ and\ renameb}\<%
\\
\>[0]\AgdaFunction{substb-rename-commute}\AgdaSpace{}%
\AgdaSymbol{:}\AgdaSpace{}%
\AgdaSymbol{∀}\AgdaSpace{}%
\AgdaSymbol{\{}\AgdaBound{Γ₁}\AgdaSpace{}%
\AgdaBound{Γ₂}\AgdaSpace{}%
\AgdaBound{Γ₃}\AgdaSpace{}%
\AgdaBound{A}\AgdaSymbol{\}}\<%
\\
\>[0][@{}l@{\AgdaIndent{0}}]%
\>[2]\AgdaSymbol{→}\AgdaSpace{}%
\AgdaSymbol{(}\AgdaBound{fb'}\AgdaSpace{}%
\AgdaSymbol{:}\AgdaSpace{}%
\AgdaSymbol{∀}\AgdaSpace{}%
\AgdaSymbol{\{}\AgdaBound{Ab}\AgdaSymbol{\}}\AgdaSpace{}%
\AgdaSymbol{→}\AgdaSpace{}%
\AgdaBound{Ab}\AgdaSpace{}%
\AgdaOperator{\AgdaDatatype{∈b}}\AgdaSpace{}%
\AgdaBound{Γ₂}\AgdaSpace{}%
\AgdaSymbol{→}\AgdaSpace{}%
\AgdaBound{Ab}\AgdaSpace{}%
\AgdaOperator{\AgdaDatatype{∈b}}\AgdaSpace{}%
\AgdaBound{Γ₃}\AgdaSymbol{)}\<%
\\
%
\>[2]\AgdaSymbol{→}\AgdaSpace{}%
\AgdaSymbol{(}\AgdaBound{fb''}\AgdaSpace{}%
\AgdaSymbol{:}\AgdaSpace{}%
\AgdaSymbol{∀}\AgdaSpace{}%
\AgdaSymbol{\{}\AgdaBound{Ab}\AgdaSymbol{\}}\AgdaSpace{}%
\AgdaSymbol{→}\AgdaSpace{}%
\AgdaInductiveConstructor{□}\AgdaSpace{}%
\AgdaBound{Ab}\AgdaSpace{}%
\AgdaOperator{\AgdaDatatype{∈b}}\AgdaSpace{}%
\AgdaBound{Γ₁}\AgdaSpace{}%
\AgdaSymbol{→}\AgdaSpace{}%
\AgdaBound{Γ₂}\AgdaSpace{}%
\AgdaOperator{\AgdaDatatype{⊢b}}\AgdaSpace{}%
\AgdaBound{Ab}\AgdaSymbol{)}\<%
\\
%
\>[2]\AgdaSymbol{→}\AgdaSpace{}%
\AgdaSymbol{(}\AgdaBound{t}\AgdaSpace{}%
\AgdaSymbol{:}\AgdaSpace{}%
\AgdaBound{Γ₁}\AgdaSpace{}%
\AgdaOperator{\AgdaDatatype{⊢b}}\AgdaSpace{}%
\AgdaBound{A}\AgdaSymbol{)}\<%
\\
%
\>[2]\AgdaSymbol{→}\AgdaSpace{}%
\AgdaFunction{substb}\AgdaSpace{}%
\AgdaSymbol{(λ}\AgdaSpace{}%
\AgdaBound{x}\AgdaSpace{}%
\AgdaSymbol{→}\AgdaSpace{}%
\AgdaFunction{renameb}\AgdaSpace{}%
\AgdaBound{fb'}\AgdaSpace{}%
\AgdaSymbol{(}\AgdaBound{fb''}\AgdaSpace{}%
\AgdaBound{x}\AgdaSymbol{))}\AgdaSpace{}%
\AgdaBound{t}\AgdaSpace{}%
\AgdaOperator{\AgdaDatatype{≡}}\AgdaSpace{}%
\AgdaFunction{renameb}\AgdaSpace{}%
\AgdaBound{fb'}\AgdaSpace{}%
\AgdaSymbol{(}\AgdaFunction{substb}\AgdaSpace{}%
\AgdaBound{fb''}\AgdaSpace{}%
\AgdaBound{t}\AgdaSymbol{)}\<%
\\
\>[0]\AgdaFunction{substb-rename-commute}\AgdaSpace{}%
\AgdaBound{fb'}\AgdaSpace{}%
\AgdaBound{fb''}\AgdaSpace{}%
\AgdaSymbol{(}\AgdaInductiveConstructor{ε}\AgdaSpace{}%
\AgdaBound{x}\AgdaSymbol{)}\AgdaSpace{}%
\AgdaSymbol{=}\AgdaSpace{}%
\AgdaInductiveConstructor{refl}\<%
\\
\>[0]\AgdaFunction{substb-rename-commute}\AgdaSpace{}%
\AgdaBound{fb'}\AgdaSpace{}%
\AgdaBound{fb''}\AgdaSpace{}%
\AgdaSymbol{(}\AgdaInductiveConstructor{box}\AgdaSpace{}%
\AgdaBound{t}\AgdaSymbol{)}\AgdaSpace{}%
\AgdaSymbol{=}\<%
\\
\>[0][@{}l@{\AgdaIndent{0}}]%
\>[2]\AgdaFunction{cong}\<%
\\
\>[2][@{}l@{\AgdaIndent{0}}]%
\>[4]\AgdaInductiveConstructor{box}\<%
\\
%
\>[4]\AgdaSymbol{(}\AgdaFunction{substb-rename-commute}\AgdaSpace{}%
\AgdaBound{fb'}\AgdaSpace{}%
\AgdaBound{fb''}\AgdaSpace{}%
\AgdaBound{t}\AgdaSymbol{)}\<%
\\
\>[0]\AgdaFunction{substb-rename-commute}\AgdaSpace{}%
\AgdaBound{fb'}\AgdaSpace{}%
\AgdaBound{fb''}\AgdaSpace{}%
\AgdaSymbol{(}\AgdaInductiveConstructor{ƛb}\AgdaSpace{}%
\AgdaBound{t}\AgdaSymbol{)}\AgdaSpace{}%
\AgdaSymbol{=}\<%
\\
\>[0][@{}l@{\AgdaIndent{0}}]%
\>[2]\AgdaOperator{\AgdaFunction{beginEq}}\<%
\\
\>[2][@{}l@{\AgdaIndent{0}}]%
\>[4]\AgdaInductiveConstructor{ƛb}\AgdaSpace{}%
\AgdaSymbol{(}\AgdaFunction{substb}\AgdaSpace{}%
\AgdaSymbol{(}\AgdaFunction{extsb}\AgdaSpace{}%
\AgdaSymbol{(λ}\AgdaSpace{}%
\AgdaBound{x}\AgdaSpace{}%
\AgdaSymbol{→}\AgdaSpace{}%
\AgdaFunction{renameb}\AgdaSpace{}%
\AgdaBound{fb'}\AgdaSpace{}%
\AgdaSymbol{(}\AgdaBound{fb''}\AgdaSpace{}%
\AgdaBound{x}\AgdaSymbol{)))}\AgdaSpace{}%
\AgdaBound{t}\AgdaSymbol{)}\<%
\\
%
\>[2]\AgdaFunction{≡⟨}\AgdaSpace{}%
\AgdaFunction{cong}\AgdaSpace{}%
\AgdaInductiveConstructor{ƛb}\AgdaSpace{}%
\AgdaSymbol{(}\AgdaFunction{substb-cong}\AgdaSpace{}%
\AgdaSymbol{(λ}\AgdaSpace{}%
\AgdaBound{x}\AgdaSpace{}%
\AgdaSymbol{→}\AgdaSpace{}%
\AgdaFunction{extsb-rename-commute}\AgdaSpace{}%
\AgdaBound{fb'}\AgdaSpace{}%
\AgdaBound{fb''}\AgdaSpace{}%
\AgdaBound{x}\AgdaSymbol{)}\AgdaSpace{}%
\AgdaBound{t}\AgdaSymbol{)}\AgdaSpace{}%
\AgdaFunction{⟩}\<%
\\
\>[2][@{}l@{\AgdaIndent{0}}]%
\>[4]\AgdaInductiveConstructor{ƛb}\AgdaSpace{}%
\AgdaSymbol{(}\AgdaFunction{substb}\AgdaSpace{}%
\AgdaSymbol{(λ}\AgdaSpace{}%
\AgdaBound{x}\AgdaSpace{}%
\AgdaSymbol{→}\AgdaSpace{}%
\AgdaFunction{renameb}\AgdaSpace{}%
\AgdaSymbol{(}\AgdaFunction{extb}\AgdaSpace{}%
\AgdaBound{fb'}\AgdaSymbol{)}\AgdaSpace{}%
\AgdaSymbol{((}\AgdaFunction{extsb}\AgdaSpace{}%
\AgdaBound{fb''}\AgdaSymbol{)}\AgdaSpace{}%
\AgdaBound{x}\AgdaSymbol{))}\AgdaSpace{}%
\AgdaBound{t}\AgdaSymbol{)}\<%
\\
%
\>[2]\AgdaFunction{≡⟨}\AgdaSpace{}%
\AgdaFunction{cong}\AgdaSpace{}%
\AgdaInductiveConstructor{ƛb}\AgdaSpace{}%
\AgdaSymbol{(}\AgdaFunction{substb-rename-commute}\AgdaSpace{}%
\AgdaSymbol{(}\AgdaFunction{extb}\AgdaSpace{}%
\AgdaBound{fb'}\AgdaSymbol{)}\AgdaSpace{}%
\AgdaSymbol{(}\AgdaFunction{extsb}\AgdaSpace{}%
\AgdaBound{fb''}\AgdaSymbol{)}\AgdaSpace{}%
\AgdaBound{t}\AgdaSymbol{)}\AgdaSpace{}%
\AgdaFunction{⟩}\<%
\\
\>[2][@{}l@{\AgdaIndent{0}}]%
\>[4]\AgdaInductiveConstructor{ƛb}\AgdaSpace{}%
\AgdaSymbol{(}\AgdaFunction{renameb}\AgdaSpace{}%
\AgdaSymbol{(}\AgdaFunction{extb}\AgdaSpace{}%
\AgdaBound{fb'}\AgdaSymbol{)}\AgdaSpace{}%
\AgdaSymbol{(}\AgdaFunction{substb}\AgdaSpace{}%
\AgdaSymbol{(}\AgdaFunction{extsb}\AgdaSpace{}%
\AgdaBound{fb''}\AgdaSymbol{)}\AgdaSpace{}%
\AgdaBound{t}\AgdaSymbol{))}\<%
\\
%
\>[2]\AgdaOperator{\AgdaFunction{∎eq}}\<%
\\
\>[0]\AgdaFunction{substb-rename-commute}\AgdaSpace{}%
\AgdaBound{fb'}\AgdaSpace{}%
\AgdaBound{fb''}\AgdaSpace{}%
\AgdaSymbol{(}\AgdaBound{t}\AgdaSpace{}%
\AgdaOperator{\AgdaInductiveConstructor{·b}}\AgdaSpace{}%
\AgdaBound{t₁}\AgdaSymbol{)}\AgdaSpace{}%
\AgdaSymbol{=}\<%
\\
\>[0][@{}l@{\AgdaIndent{0}}]%
\>[2]\AgdaFunction{cong₂}\<%
\\
\>[2][@{}l@{\AgdaIndent{0}}]%
\>[4]\AgdaOperator{\AgdaInductiveConstructor{\AgdaUnderscore{}·b\AgdaUnderscore{}}}\<%
\\
%
\>[4]\AgdaSymbol{(}\AgdaFunction{substb-rename-commute}\AgdaSpace{}%
\AgdaBound{fb'}\AgdaSpace{}%
\AgdaBound{fb''}\AgdaSpace{}%
\AgdaBound{t}\AgdaSymbol{)}\<%
\\
%
\>[4]\AgdaSymbol{(}\AgdaFunction{substb-rename-commute}\AgdaSpace{}%
\AgdaBound{fb'}\AgdaSpace{}%
\AgdaBound{fb''}\AgdaSpace{}%
\AgdaBound{t₁}\AgdaSymbol{)}\<%
\\
%
\\[\AgdaEmptyExtraSkip]%
\>[0]\AgdaComment{--\ extsb\ distributes\ over\ composition:}\<%
\\
\>[0]\AgdaComment{--\ Executing\ extsb\ first\ with\ fb',\ then\ with\ fb''}\<%
\\
\>[0]\AgdaComment{--\ is\ the\ same\ as\ doing\ extsb\ once\ with\ fb''\ ∘\ fb'.}\<%
\\
\>[0]\AgdaFunction{extsb-comp}\AgdaSpace{}%
\AgdaSymbol{:}\AgdaSpace{}%
\AgdaSymbol{∀}\AgdaSpace{}%
\AgdaSymbol{\{}\AgdaBound{Γ₁}\AgdaSpace{}%
\AgdaBound{Γ₂}\AgdaSpace{}%
\AgdaBound{Γ₃}\AgdaSpace{}%
\AgdaBound{A}\AgdaSpace{}%
\AgdaBound{B}\AgdaSymbol{\}}\<%
\\
\>[0][@{}l@{\AgdaIndent{0}}]%
\>[2]\AgdaSymbol{→}\AgdaSpace{}%
\AgdaSymbol{(}\AgdaBound{fb'}\AgdaSpace{}%
\AgdaSymbol{:}\AgdaSpace{}%
\AgdaSymbol{∀}\AgdaSpace{}%
\AgdaSymbol{\{}\AgdaBound{Ab}\AgdaSymbol{\}}\AgdaSpace{}%
\AgdaSymbol{→}\AgdaSpace{}%
\AgdaBound{Ab}\AgdaSpace{}%
\AgdaOperator{\AgdaDatatype{∈b}}\AgdaSpace{}%
\AgdaBound{Γ₁}\AgdaSpace{}%
\AgdaSymbol{→}\AgdaSpace{}%
\AgdaBound{Ab}\AgdaSpace{}%
\AgdaOperator{\AgdaDatatype{∈b}}\AgdaSpace{}%
\AgdaBound{Γ₂}\AgdaSymbol{)}\<%
\\
%
\>[2]\AgdaSymbol{→}\AgdaSpace{}%
\AgdaSymbol{(}\AgdaBound{fb''}\AgdaSpace{}%
\AgdaSymbol{:}\AgdaSpace{}%
\AgdaSymbol{∀}\AgdaSpace{}%
\AgdaSymbol{\{}\AgdaBound{Ab}\AgdaSymbol{\}}\AgdaSpace{}%
\AgdaSymbol{→}\AgdaSpace{}%
\AgdaInductiveConstructor{□}\AgdaSpace{}%
\AgdaBound{Ab}\AgdaSpace{}%
\AgdaOperator{\AgdaDatatype{∈b}}\AgdaSpace{}%
\AgdaBound{Γ₂}\AgdaSpace{}%
\AgdaSymbol{→}\AgdaSpace{}%
\AgdaBound{Γ₃}\AgdaSpace{}%
\AgdaOperator{\AgdaDatatype{⊢b}}\AgdaSpace{}%
\AgdaBound{Ab}\AgdaSymbol{)}\<%
\\
%
\>[2]\AgdaSymbol{→}\AgdaSpace{}%
\AgdaSymbol{(}\AgdaBound{x}\AgdaSpace{}%
\AgdaSymbol{:}\AgdaSpace{}%
\AgdaInductiveConstructor{□}\AgdaSpace{}%
\AgdaBound{A}\AgdaSpace{}%
\AgdaOperator{\AgdaDatatype{∈b}}\AgdaSpace{}%
\AgdaBound{Γ₁}\AgdaSpace{}%
\AgdaOperator{\AgdaInductiveConstructor{,b}}\AgdaSpace{}%
\AgdaBound{B}\AgdaSymbol{)}\<%
\\
%
\>[2]\AgdaSymbol{→}\AgdaSpace{}%
\AgdaFunction{extsb}\AgdaSpace{}%
\AgdaBound{fb''}\AgdaSpace{}%
\AgdaSymbol{(}\AgdaFunction{extb}\AgdaSpace{}%
\AgdaBound{fb'}\AgdaSpace{}%
\AgdaBound{x}\AgdaSymbol{)}\AgdaSpace{}%
\AgdaOperator{\AgdaDatatype{≡}}\AgdaSpace{}%
\AgdaFunction{extsb}\AgdaSpace{}%
\AgdaSymbol{(λ}\AgdaSpace{}%
\AgdaBound{x₁}\AgdaSpace{}%
\AgdaSymbol{→}\AgdaSpace{}%
\AgdaBound{fb''}\AgdaSpace{}%
\AgdaSymbol{(}\AgdaBound{fb'}\AgdaSpace{}%
\AgdaBound{x₁}\AgdaSymbol{))}\AgdaSpace{}%
\AgdaBound{x}\<%
\\
\>[0]\AgdaFunction{extsb-comp}\AgdaSpace{}%
\AgdaBound{fb'}\AgdaSpace{}%
\AgdaBound{fb''}\AgdaSpace{}%
\AgdaInductiveConstructor{Zb}\AgdaSpace{}%
\AgdaSymbol{=}\AgdaSpace{}%
\AgdaInductiveConstructor{refl}\<%
\\
\>[0]\AgdaFunction{extsb-comp}\AgdaSpace{}%
\AgdaBound{fb'}\AgdaSpace{}%
\AgdaBound{fb''}\AgdaSpace{}%
\AgdaSymbol{(}\AgdaInductiveConstructor{Sb}\AgdaSpace{}%
\AgdaBound{x}\AgdaSymbol{)}\AgdaSpace{}%
\AgdaSymbol{=}\AgdaSpace{}%
\AgdaInductiveConstructor{refl}\<%
\\
%
\\[\AgdaEmptyExtraSkip]%
\>[0]\AgdaComment{--\ substb\ and\ rename\ distribute\ over\ composition:}\<%
\\
\>[0]\AgdaComment{--\ Renaming\ first\ with\ fb'\ and\ substituting\ with\ fb''}\<%
\\
\>[0]\AgdaComment{--\ is\ the\ same\ as\ substituting\ with\ fb''\ ∘\ fb'}\<%
\\
\>[0]\AgdaFunction{subst-renameb-comp}\AgdaSpace{}%
\AgdaSymbol{:}\AgdaSpace{}%
\AgdaSymbol{∀}\AgdaSpace{}%
\AgdaSymbol{\{}\AgdaBound{Γ₁}\AgdaSpace{}%
\AgdaBound{Γ₂}\AgdaSpace{}%
\AgdaBound{Γ₃}\AgdaSpace{}%
\AgdaBound{A}\AgdaSymbol{\}}\<%
\\
\>[0][@{}l@{\AgdaIndent{0}}]%
\>[2]\AgdaSymbol{→}\AgdaSpace{}%
\AgdaSymbol{(}\AgdaBound{fb'}\AgdaSpace{}%
\AgdaSymbol{:}\AgdaSpace{}%
\AgdaSymbol{∀}\AgdaSpace{}%
\AgdaSymbol{\{}\AgdaBound{Ab}\AgdaSymbol{\}}\AgdaSpace{}%
\AgdaSymbol{→}\AgdaSpace{}%
\AgdaBound{Ab}\AgdaSpace{}%
\AgdaOperator{\AgdaDatatype{∈b}}\AgdaSpace{}%
\AgdaBound{Γ₁}\AgdaSpace{}%
\AgdaSymbol{→}\AgdaSpace{}%
\AgdaBound{Ab}\AgdaSpace{}%
\AgdaOperator{\AgdaDatatype{∈b}}\AgdaSpace{}%
\AgdaBound{Γ₂}\AgdaSymbol{)}\<%
\\
%
\>[2]\AgdaSymbol{→}\AgdaSpace{}%
\AgdaSymbol{(}\AgdaBound{fb''}\AgdaSpace{}%
\AgdaSymbol{:}\AgdaSpace{}%
\AgdaSymbol{∀}\AgdaSpace{}%
\AgdaSymbol{\{}\AgdaBound{Ab}\AgdaSymbol{\}}\AgdaSpace{}%
\AgdaSymbol{→}\AgdaSpace{}%
\AgdaInductiveConstructor{□}\AgdaSpace{}%
\AgdaBound{Ab}\AgdaSpace{}%
\AgdaOperator{\AgdaDatatype{∈b}}\AgdaSpace{}%
\AgdaBound{Γ₂}\AgdaSpace{}%
\AgdaSymbol{→}\AgdaSpace{}%
\AgdaBound{Γ₃}\AgdaSpace{}%
\AgdaOperator{\AgdaDatatype{⊢b}}\AgdaSpace{}%
\AgdaBound{Ab}\AgdaSymbol{)}\<%
\\
%
\>[2]\AgdaSymbol{→}\AgdaSpace{}%
\AgdaSymbol{(}\AgdaBound{t}\AgdaSpace{}%
\AgdaSymbol{:}\AgdaSpace{}%
\AgdaBound{Γ₁}\AgdaSpace{}%
\AgdaOperator{\AgdaDatatype{⊢b}}\AgdaSpace{}%
\AgdaBound{A}\AgdaSymbol{)}\<%
\\
%
\>[2]\AgdaSymbol{→}\AgdaSpace{}%
\AgdaFunction{substb}\AgdaSpace{}%
\AgdaBound{fb''}\AgdaSpace{}%
\AgdaSymbol{(}\AgdaFunction{renameb}\AgdaSpace{}%
\AgdaBound{fb'}\AgdaSpace{}%
\AgdaBound{t}\AgdaSymbol{)}\AgdaSpace{}%
\AgdaOperator{\AgdaDatatype{≡}}\AgdaSpace{}%
\AgdaFunction{substb}\AgdaSpace{}%
\AgdaSymbol{(λ}\AgdaSpace{}%
\AgdaBound{x}\AgdaSpace{}%
\AgdaSymbol{→}\AgdaSpace{}%
\AgdaBound{fb''}\AgdaSpace{}%
\AgdaSymbol{(}\AgdaBound{fb'}\AgdaSpace{}%
\AgdaBound{x}\AgdaSymbol{))}\AgdaSpace{}%
\AgdaBound{t}\<%
\\
\>[0]\AgdaFunction{subst-renameb-comp}\AgdaSpace{}%
\AgdaBound{fb'}\AgdaSpace{}%
\AgdaBound{fb''}\AgdaSpace{}%
\AgdaSymbol{(}\AgdaInductiveConstructor{ε}\AgdaSpace{}%
\AgdaBound{x}\AgdaSymbol{)}\AgdaSpace{}%
\AgdaSymbol{=}\AgdaSpace{}%
\AgdaInductiveConstructor{refl}\<%
\\
\>[0]\AgdaFunction{subst-renameb-comp}\AgdaSpace{}%
\AgdaBound{fb'}\AgdaSpace{}%
\AgdaBound{fb''}\AgdaSpace{}%
\AgdaSymbol{(}\AgdaInductiveConstructor{box}\AgdaSpace{}%
\AgdaBound{t}\AgdaSymbol{)}\AgdaSpace{}%
\AgdaSymbol{=}\AgdaSpace{}%
\AgdaFunction{cong}\AgdaSpace{}%
\AgdaInductiveConstructor{box}\AgdaSpace{}%
\AgdaSymbol{(}\AgdaFunction{subst-renameb-comp}\AgdaSpace{}%
\AgdaBound{fb'}\AgdaSpace{}%
\AgdaBound{fb''}\AgdaSpace{}%
\AgdaBound{t}\AgdaSymbol{)}\<%
\\
\>[0]\AgdaFunction{subst-renameb-comp}\AgdaSpace{}%
\AgdaBound{fb'}\AgdaSpace{}%
\AgdaBound{fb''}\AgdaSpace{}%
\AgdaSymbol{(}\AgdaInductiveConstructor{ƛb}\AgdaSpace{}%
\AgdaBound{t}\AgdaSymbol{)}\AgdaSpace{}%
\AgdaSymbol{=}\<%
\\
\>[0][@{}l@{\AgdaIndent{0}}]%
\>[2]\AgdaOperator{\AgdaFunction{beginEq}}\<%
\\
\>[2][@{}l@{\AgdaIndent{0}}]%
\>[4]\AgdaInductiveConstructor{ƛb}\AgdaSpace{}%
\AgdaSymbol{(}\AgdaFunction{substb}\AgdaSpace{}%
\AgdaSymbol{(}\AgdaFunction{extsb}\AgdaSpace{}%
\AgdaBound{fb''}\AgdaSymbol{)}\AgdaSpace{}%
\AgdaSymbol{(}\AgdaFunction{renameb}\AgdaSpace{}%
\AgdaSymbol{(}\AgdaFunction{extb}\AgdaSpace{}%
\AgdaBound{fb'}\AgdaSymbol{)}\AgdaSpace{}%
\AgdaBound{t}\AgdaSymbol{))}\<%
\\
%
\>[2]\AgdaFunction{≡⟨}\AgdaSpace{}%
\AgdaFunction{cong}\AgdaSpace{}%
\AgdaInductiveConstructor{ƛb}\AgdaSpace{}%
\AgdaSymbol{(}\AgdaFunction{subst-renameb-comp}\AgdaSpace{}%
\AgdaSymbol{(}\AgdaFunction{extb}\AgdaSpace{}%
\AgdaBound{fb'}\AgdaSymbol{)}\AgdaSpace{}%
\AgdaSymbol{(}\AgdaFunction{extsb}\AgdaSpace{}%
\AgdaBound{fb''}\AgdaSymbol{)}\AgdaSpace{}%
\AgdaBound{t}\AgdaSymbol{)}\AgdaSpace{}%
\AgdaFunction{⟩}\<%
\\
\>[2][@{}l@{\AgdaIndent{0}}]%
\>[4]\AgdaInductiveConstructor{ƛb}\AgdaSpace{}%
\AgdaSymbol{(}\AgdaFunction{substb}\AgdaSpace{}%
\AgdaSymbol{(λ}\AgdaSpace{}%
\AgdaBound{x}\AgdaSpace{}%
\AgdaSymbol{→}\AgdaSpace{}%
\AgdaSymbol{(}\AgdaFunction{extsb}\AgdaSpace{}%
\AgdaBound{fb''}\AgdaSymbol{)}\AgdaSpace{}%
\AgdaSymbol{((}\AgdaFunction{extb}\AgdaSpace{}%
\AgdaBound{fb'}\AgdaSymbol{)}\AgdaSpace{}%
\AgdaBound{x}\AgdaSymbol{))}\AgdaSpace{}%
\AgdaBound{t}\AgdaSymbol{)}\<%
\\
%
\>[2]\AgdaFunction{≡⟨}\AgdaSpace{}%
\AgdaFunction{cong}\AgdaSpace{}%
\AgdaInductiveConstructor{ƛb}\AgdaSpace{}%
\AgdaSymbol{(}\AgdaFunction{substb-cong}\AgdaSpace{}%
\AgdaSymbol{(λ}\AgdaSpace{}%
\AgdaBound{x}\AgdaSpace{}%
\AgdaSymbol{→}\AgdaSpace{}%
\AgdaFunction{extsb-comp}\AgdaSpace{}%
\AgdaBound{fb'}\AgdaSpace{}%
\AgdaBound{fb''}\AgdaSpace{}%
\AgdaBound{x}\AgdaSymbol{)}\AgdaSpace{}%
\AgdaBound{t}\AgdaSymbol{)}\AgdaSpace{}%
\AgdaFunction{⟩}\<%
\\
\>[2][@{}l@{\AgdaIndent{0}}]%
\>[4]\AgdaInductiveConstructor{ƛb}\AgdaSpace{}%
\AgdaSymbol{(}\AgdaFunction{substb}\AgdaSpace{}%
\AgdaSymbol{(}\AgdaFunction{extsb}\AgdaSpace{}%
\AgdaSymbol{(λ}\AgdaSpace{}%
\AgdaBound{x}\AgdaSpace{}%
\AgdaSymbol{→}\AgdaSpace{}%
\AgdaBound{fb''}\AgdaSpace{}%
\AgdaSymbol{(}\AgdaBound{fb'}\AgdaSpace{}%
\AgdaBound{x}\AgdaSymbol{)))}\AgdaSpace{}%
\AgdaBound{t}\AgdaSymbol{)}\<%
\\
%
\>[2]\AgdaOperator{\AgdaFunction{∎eq}}\<%
\\
\>[0]\AgdaFunction{subst-renameb-comp}\AgdaSpace{}%
\AgdaBound{fb'}\AgdaSpace{}%
\AgdaBound{fb''}\AgdaSpace{}%
\AgdaSymbol{(}\AgdaBound{t}\AgdaSpace{}%
\AgdaOperator{\AgdaInductiveConstructor{·b}}\AgdaSpace{}%
\AgdaBound{t₁}\AgdaSymbol{)}\AgdaSpace{}%
\AgdaSymbol{=}\AgdaSpace{}%
\AgdaFunction{cong₂}\AgdaSpace{}%
\AgdaOperator{\AgdaInductiveConstructor{\AgdaUnderscore{}·b\AgdaUnderscore{}}}\AgdaSpace{}%
\AgdaSymbol{(}\AgdaFunction{subst-renameb-comp}\AgdaSpace{}%
\AgdaBound{fb'}\AgdaSpace{}%
\AgdaBound{fb''}\AgdaSpace{}%
\AgdaBound{t}\AgdaSymbol{)}\AgdaSpace{}%
\AgdaSymbol{(}\AgdaFunction{subst-renameb-comp}\AgdaSpace{}%
\AgdaBound{fb'}\AgdaSpace{}%
\AgdaBound{fb''}\AgdaSpace{}%
\AgdaBound{t₁}\AgdaSymbol{)}\<%
\\
\>[0]\AgdaComment{--\ Renaming\ followed\ by\ substitution\ with\ an\ extended\ context\ (extsb)}\<%
\\
\>[0]\AgdaComment{--\ commutes\ with\ substitution\ followed\ by\ renaming.}\<%
\\
\>[0]\AgdaFunction{substb-exts-rename-Sb-commute}\AgdaSpace{}%
\AgdaSymbol{:}\AgdaSpace{}%
\AgdaSymbol{∀}\AgdaSpace{}%
\AgdaSymbol{\{}\AgdaBound{Γ}\AgdaSpace{}%
\AgdaBound{A}\AgdaSpace{}%
\AgdaBound{B}\AgdaSpace{}%
\AgdaBound{Δ}\AgdaSymbol{\}}\<%
\\
\>[0][@{}l@{\AgdaIndent{0}}]%
\>[2]\AgdaSymbol{→}\AgdaSpace{}%
\AgdaSymbol{(}\AgdaBound{fb'}\AgdaSpace{}%
\AgdaSymbol{:}\AgdaSpace{}%
\AgdaSymbol{∀}\AgdaSpace{}%
\AgdaSymbol{\{}\AgdaBound{Ab}\AgdaSymbol{\}}\AgdaSpace{}%
\AgdaSymbol{→}\AgdaSpace{}%
\AgdaInductiveConstructor{□}\AgdaSpace{}%
\AgdaBound{Ab}\AgdaSpace{}%
\AgdaOperator{\AgdaDatatype{∈b}}\AgdaSpace{}%
\AgdaBound{Γ}\AgdaSpace{}%
\AgdaSymbol{→}\AgdaSpace{}%
\AgdaBound{Δ}\AgdaSpace{}%
\AgdaOperator{\AgdaDatatype{⊢b}}\AgdaSpace{}%
\AgdaBound{Ab}\AgdaSymbol{)}\<%
\\
%
\>[2]\AgdaSymbol{→}\AgdaSpace{}%
\AgdaSymbol{(}\AgdaBound{t}\AgdaSpace{}%
\AgdaSymbol{:}\AgdaSpace{}%
\AgdaBound{Γ}\AgdaSpace{}%
\AgdaOperator{\AgdaDatatype{⊢b}}\AgdaSpace{}%
\AgdaBound{A}\AgdaSymbol{)}\<%
\\
%
\>[2]\AgdaSymbol{→}%
\>[1555I]\AgdaFunction{substb}\AgdaSpace{}%
\AgdaSymbol{(}\AgdaFunction{extsb}\AgdaSpace{}%
\AgdaBound{fb'}\AgdaSpace{}%
\AgdaSymbol{\{}\AgdaArgument{B}\AgdaSpace{}%
\AgdaSymbol{=}\AgdaSpace{}%
\AgdaBound{B}\AgdaSymbol{\})}\AgdaSpace{}%
\AgdaSymbol{(}\AgdaFunction{renameb}\AgdaSpace{}%
\AgdaInductiveConstructor{Sb}\AgdaSpace{}%
\AgdaBound{t}\AgdaSymbol{)}\<%
\\
\>[.][@{}l@{}]\<[1555I]%
\>[4]\AgdaOperator{\AgdaDatatype{≡}}\AgdaSpace{}%
\AgdaFunction{renameb}\AgdaSpace{}%
\AgdaInductiveConstructor{Sb}\AgdaSpace{}%
\AgdaSymbol{(}\AgdaFunction{substb}\AgdaSpace{}%
\AgdaBound{fb'}\AgdaSpace{}%
\AgdaBound{t}\AgdaSymbol{)}\<%
\\
\>[0]\AgdaFunction{substb-exts-rename-Sb-commute}\AgdaSpace{}%
\AgdaBound{fb'}\AgdaSpace{}%
\AgdaBound{t}\AgdaSpace{}%
\AgdaSymbol{=}\<%
\\
\>[0][@{}l@{\AgdaIndent{0}}]%
\>[2]\AgdaOperator{\AgdaFunction{beginEq}}\<%
\\
\>[2][@{}l@{\AgdaIndent{0}}]%
\>[4]\AgdaFunction{substb}\AgdaSpace{}%
\AgdaSymbol{(}\AgdaFunction{extsb}\AgdaSpace{}%
\AgdaBound{fb'}\AgdaSymbol{)}\AgdaSpace{}%
\AgdaSymbol{(}\AgdaFunction{renameb}\AgdaSpace{}%
\AgdaInductiveConstructor{Sb}\AgdaSpace{}%
\AgdaBound{t}\AgdaSymbol{)}\<%
\\
%
\>[2]\AgdaFunction{≡⟨}\AgdaSpace{}%
\AgdaFunction{subst-renameb-comp}\AgdaSpace{}%
\AgdaInductiveConstructor{Sb}\AgdaSpace{}%
\AgdaSymbol{(}\AgdaFunction{extsb}\AgdaSpace{}%
\AgdaBound{fb'}\AgdaSymbol{)}\AgdaSpace{}%
\AgdaBound{t}\AgdaSpace{}%
\AgdaFunction{⟩}\<%
\\
\>[2][@{}l@{\AgdaIndent{0}}]%
\>[4]\AgdaFunction{substb}\AgdaSpace{}%
\AgdaSymbol{(λ}\AgdaSpace{}%
\AgdaBound{x}\AgdaSpace{}%
\AgdaSymbol{→}\AgdaSpace{}%
\AgdaFunction{renameb}\AgdaSpace{}%
\AgdaInductiveConstructor{Sb}\AgdaSpace{}%
\AgdaSymbol{(}\AgdaBound{fb'}\AgdaSpace{}%
\AgdaBound{x}\AgdaSymbol{))}\AgdaSpace{}%
\AgdaBound{t}\<%
\\
%
\>[2]\AgdaFunction{≡⟨}\AgdaSpace{}%
\AgdaFunction{substb-rename-commute}\AgdaSpace{}%
\AgdaInductiveConstructor{Sb}\AgdaSpace{}%
\AgdaBound{fb'}\AgdaSpace{}%
\AgdaBound{t}\AgdaSpace{}%
\AgdaFunction{⟩}\<%
\\
\>[2][@{}l@{\AgdaIndent{0}}]%
\>[4]\AgdaFunction{renameb}\AgdaSpace{}%
\AgdaInductiveConstructor{Sb}\AgdaSpace{}%
\AgdaSymbol{(}\AgdaFunction{substb}\AgdaSpace{}%
\AgdaBound{fb'}\AgdaSpace{}%
\AgdaBound{t}\AgdaSymbol{)}\<%
\\
%
\>[2]\AgdaOperator{\AgdaFunction{∎eq}}\<%
\\
\>[0]\AgdaComment{--\ Preservation\ and\ Reflection\ of\ }\<%
\\
\>[0]\AgdaComment{--\ simultaneous\ substitution\ of\ Gödel's\ translation}\<%
\\
\>[0]\AgdaFunction{pres-simultaneous-subst}\AgdaSpace{}%
\AgdaSymbol{:}\AgdaSpace{}%
\AgdaSymbol{∀}\AgdaSpace{}%
\AgdaSymbol{\{}\AgdaBound{Γ}\AgdaSpace{}%
\AgdaBound{Δ}\AgdaSpace{}%
\AgdaBound{A}\AgdaSymbol{\}}\<%
\\
\>[0][@{}l@{\AgdaIndent{0}}]%
\>[2]\AgdaSymbol{→}\AgdaSpace{}%
\AgdaSymbol{(}\AgdaBound{fb'}\AgdaSpace{}%
\AgdaSymbol{:}\AgdaSpace{}%
\AgdaSymbol{∀}\AgdaSpace{}%
\AgdaSymbol{\{}\AgdaBound{Ab}\AgdaSymbol{\}}\AgdaSpace{}%
\AgdaSymbol{→}\AgdaSpace{}%
\AgdaInductiveConstructor{□}\AgdaSpace{}%
\AgdaBound{Ab}\AgdaSpace{}%
\AgdaOperator{\AgdaDatatype{∈b}}\AgdaSpace{}%
\AgdaFunction{embedContext}\AgdaSpace{}%
\AgdaBound{Γ}\AgdaSpace{}%
\AgdaSymbol{→}\AgdaSpace{}%
\AgdaFunction{embedContext}\AgdaSpace{}%
\AgdaBound{Δ}\AgdaSpace{}%
\AgdaOperator{\AgdaDatatype{⊢b}}\AgdaSpace{}%
\AgdaBound{Ab}\AgdaSymbol{)}\<%
\\
%
\>[2]\AgdaSymbol{→}\AgdaSpace{}%
\AgdaSymbol{(}\AgdaBound{f'}\AgdaSpace{}%
\AgdaSymbol{:}\AgdaSpace{}%
\AgdaSymbol{∀}\AgdaSpace{}%
\AgdaSymbol{\{}\AgdaBound{A}\AgdaSymbol{\}}\AgdaSpace{}%
\AgdaSymbol{→}\AgdaSpace{}%
\AgdaBound{A}\AgdaSpace{}%
\AgdaOperator{\AgdaDatatype{∈}}\AgdaSpace{}%
\AgdaBound{Γ}\AgdaSpace{}%
\AgdaSymbol{→}\AgdaSpace{}%
\AgdaBound{Δ}\AgdaSpace{}%
\AgdaOperator{\AgdaDatatype{⊢}}\AgdaSpace{}%
\AgdaBound{A}\AgdaSymbol{)}\<%
\\
%
\>[2]\AgdaSymbol{→}\AgdaSpace{}%
\AgdaSymbol{(}\AgdaBound{h}\AgdaSpace{}%
\AgdaSymbol{:}\AgdaSpace{}%
\AgdaSymbol{∀}\AgdaSpace{}%
\AgdaSymbol{\{}\AgdaBound{A'}\AgdaSymbol{\}}\AgdaSpace{}%
\AgdaSymbol{→}\AgdaSpace{}%
\AgdaSymbol{(}\AgdaBound{x}\AgdaSpace{}%
\AgdaSymbol{:}\AgdaSpace{}%
\AgdaBound{A'}\AgdaSpace{}%
\AgdaOperator{\AgdaDatatype{∈}}\AgdaSpace{}%
\AgdaBound{Γ}\AgdaSymbol{)}\AgdaSpace{}%
\AgdaSymbol{→}\AgdaSpace{}%
\AgdaInductiveConstructor{box}\AgdaSpace{}%
\AgdaSymbol{(}\AgdaBound{fb'}\AgdaSpace{}%
\AgdaSymbol{(}\AgdaFunction{embedLookup}\AgdaSpace{}%
\AgdaBound{x}\AgdaSymbol{))}\AgdaSpace{}%
\AgdaOperator{\AgdaDatatype{≡}}\AgdaSpace{}%
\AgdaFunction{embedTerm}\AgdaSpace{}%
\AgdaSymbol{(}\AgdaBound{f'}\AgdaSpace{}%
\AgdaBound{x}\AgdaSymbol{)}\AgdaSpace{}%
\AgdaSymbol{)}\<%
\\
%
\>[2]\AgdaSymbol{→}\AgdaSpace{}%
\AgdaSymbol{(}\AgdaBound{M}\AgdaSpace{}%
\AgdaSymbol{:}\AgdaSpace{}%
\AgdaBound{Γ}\AgdaSpace{}%
\AgdaOperator{\AgdaDatatype{⊢}}\AgdaSpace{}%
\AgdaBound{A}\AgdaSymbol{)}\<%
\\
%
\>[2]\AgdaSymbol{→}\AgdaSpace{}%
\AgdaFunction{substb}\AgdaSpace{}%
\AgdaBound{fb'}\AgdaSpace{}%
\AgdaSymbol{(}\AgdaFunction{embedTerm}\AgdaSpace{}%
\AgdaBound{M}\AgdaSymbol{)}\AgdaSpace{}%
\AgdaOperator{\AgdaDatatype{≡}}\AgdaSpace{}%
\AgdaFunction{embedTerm}\AgdaSpace{}%
\AgdaSymbol{(}\AgdaFunction{subst}\AgdaSpace{}%
\AgdaBound{f'}\AgdaSpace{}%
\AgdaBound{M}\AgdaSymbol{)}\<%
\\
\>[0]\AgdaFunction{pres-simultaneous-subst}\AgdaSpace{}%
\AgdaBound{fb'}\AgdaSpace{}%
\AgdaBound{f'}\AgdaSpace{}%
\AgdaBound{h}\AgdaSpace{}%
\AgdaSymbol{(}\AgdaInductiveConstructor{`}\AgdaSpace{}%
\AgdaBound{x}\AgdaSymbol{)}\AgdaSpace{}%
\AgdaSymbol{=}\<%
\\
\>[0][@{}l@{\AgdaIndent{0}}]%
\>[2]\AgdaBound{h}\AgdaSpace{}%
\AgdaBound{x}\<%
\\
\>[0]\AgdaFunction{pres-simultaneous-subst}\AgdaSpace{}%
\AgdaBound{fb'}\AgdaSpace{}%
\AgdaBound{f'}\AgdaSpace{}%
\AgdaBound{h}\AgdaSpace{}%
\AgdaSymbol{(}\AgdaInductiveConstructor{ƛ}\AgdaSpace{}%
\AgdaBound{M}\AgdaSymbol{)}\AgdaSpace{}%
\AgdaSymbol{=}\<%
\\
\>[0][@{}l@{\AgdaIndent{0}}]%
\>[2]\AgdaFunction{cong}\AgdaSpace{}%
\AgdaInductiveConstructor{box}\AgdaSpace{}%
\AgdaSymbol{(}\AgdaFunction{cong}\AgdaSpace{}%
\AgdaInductiveConstructor{ƛb}\AgdaSpace{}%
\AgdaSymbol{(}\AgdaFunction{pres-simultaneous-subst}\AgdaSpace{}%
\AgdaSymbol{(}\AgdaFunction{extsb}\AgdaSpace{}%
\AgdaBound{fb'}\AgdaSymbol{)}\AgdaSpace{}%
\AgdaSymbol{(}\AgdaFunction{exts}\AgdaSpace{}%
\AgdaBound{f'}\AgdaSymbol{)}\AgdaSpace{}%
\AgdaSymbol{(}\AgdaFunction{pres-exts}\AgdaSpace{}%
\AgdaBound{fb'}\AgdaSpace{}%
\AgdaBound{f'}\AgdaSpace{}%
\AgdaBound{h}\AgdaSymbol{)}\AgdaSpace{}%
\AgdaBound{M}\AgdaSymbol{))}\<%
\\
\>[0]\AgdaFunction{pres-simultaneous-subst}\AgdaSpace{}%
\AgdaBound{fb'}\AgdaSpace{}%
\AgdaBound{f'}\AgdaSpace{}%
\AgdaBound{h}\AgdaSpace{}%
\AgdaSymbol{(}\AgdaBound{M}\AgdaSpace{}%
\AgdaOperator{\AgdaInductiveConstructor{·}}\AgdaSpace{}%
\AgdaBound{M₁}\AgdaSymbol{)}\AgdaSpace{}%
\AgdaSymbol{=}\<%
\\
\>[0][@{}l@{\AgdaIndent{0}}]%
\>[2]\AgdaOperator{\AgdaFunction{beginEq}}\<%
\\
\>[2][@{}l@{\AgdaIndent{0}}]%
\>[4]\AgdaInductiveConstructor{ƛb}\AgdaSpace{}%
\AgdaSymbol{(}\AgdaInductiveConstructor{ε}\AgdaSpace{}%
\AgdaInductiveConstructor{Zb}\AgdaSpace{}%
\AgdaOperator{\AgdaInductiveConstructor{·b}}\AgdaSpace{}%
\AgdaFunction{substb}\AgdaSpace{}%
\AgdaSymbol{(}\AgdaFunction{extsb}\AgdaSpace{}%
\AgdaBound{fb'}\AgdaSymbol{)}\AgdaSpace{}%
\AgdaSymbol{(}\AgdaFunction{renameb}\AgdaSpace{}%
\AgdaInductiveConstructor{Sb}\AgdaSpace{}%
\AgdaSymbol{(}\AgdaFunction{embedTerm}\AgdaSpace{}%
\AgdaBound{M₁}\AgdaSymbol{)))}\AgdaSpace{}%
\AgdaOperator{\AgdaInductiveConstructor{·b}}\AgdaSpace{}%
\AgdaFunction{substb}\AgdaSpace{}%
\AgdaBound{fb'}\AgdaSpace{}%
\AgdaSymbol{(}\AgdaFunction{embedTerm}\AgdaSpace{}%
\AgdaBound{M}\AgdaSymbol{)}\<%
\\
%
\>[2]\AgdaFunction{≡⟨}\AgdaSpace{}%
\AgdaFunction{cong₂}\AgdaSpace{}%
\AgdaOperator{\AgdaInductiveConstructor{\AgdaUnderscore{}·b\AgdaUnderscore{}}}\AgdaSpace{}%
\AgdaInductiveConstructor{refl}\AgdaSpace{}%
\AgdaSymbol{(}\AgdaFunction{pres-simultaneous-subst}\AgdaSpace{}%
\AgdaBound{fb'}\AgdaSpace{}%
\AgdaBound{f'}\AgdaSpace{}%
\AgdaBound{h}\AgdaSpace{}%
\AgdaBound{M}\AgdaSymbol{)}\AgdaSpace{}%
\AgdaFunction{⟩}\<%
\\
\>[2][@{}l@{\AgdaIndent{0}}]%
\>[4]\AgdaInductiveConstructor{ƛb}\AgdaSpace{}%
\AgdaSymbol{(}\AgdaInductiveConstructor{ε}\AgdaSpace{}%
\AgdaInductiveConstructor{Zb}\AgdaSpace{}%
\AgdaOperator{\AgdaInductiveConstructor{·b}}\AgdaSpace{}%
\AgdaFunction{substb}\AgdaSpace{}%
\AgdaSymbol{(}\AgdaFunction{extsb}\AgdaSpace{}%
\AgdaBound{fb'}\AgdaSymbol{)}\AgdaSpace{}%
\AgdaSymbol{(}\AgdaFunction{renameb}\AgdaSpace{}%
\AgdaInductiveConstructor{Sb}\AgdaSpace{}%
\AgdaSymbol{(}\AgdaFunction{embedTerm}\AgdaSpace{}%
\AgdaBound{M₁}\AgdaSymbol{)))}\AgdaSpace{}%
\AgdaOperator{\AgdaInductiveConstructor{·b}}\AgdaSpace{}%
\AgdaFunction{embedTerm}\AgdaSpace{}%
\AgdaSymbol{(}\AgdaFunction{subst}\AgdaSpace{}%
\AgdaBound{f'}\AgdaSpace{}%
\AgdaBound{M}\AgdaSymbol{)}\<%
\\
%
\>[2]\AgdaFunction{≡⟨}\AgdaSpace{}%
\AgdaFunction{cong₂}\AgdaSpace{}%
\AgdaOperator{\AgdaInductiveConstructor{\AgdaUnderscore{}·b\AgdaUnderscore{}}}\AgdaSpace{}%
\AgdaSymbol{(}\AgdaFunction{cong}\AgdaSpace{}%
\AgdaInductiveConstructor{ƛb}\AgdaSpace{}%
\AgdaSymbol{(}\AgdaFunction{cong₂}\AgdaSpace{}%
\AgdaOperator{\AgdaInductiveConstructor{\AgdaUnderscore{}·b\AgdaUnderscore{}}}\AgdaSpace{}%
\AgdaInductiveConstructor{refl}\AgdaSpace{}%
\AgdaSymbol{(}\AgdaFunction{substb-exts-rename-Sb-commute}\AgdaSpace{}%
\AgdaBound{fb'}\AgdaSpace{}%
\AgdaSymbol{(}\AgdaFunction{embedTerm}\AgdaSpace{}%
\AgdaBound{M₁}\AgdaSymbol{))))}\AgdaSpace{}%
\AgdaInductiveConstructor{refl}\AgdaSpace{}%
\AgdaFunction{⟩}\<%
\\
\>[2][@{}l@{\AgdaIndent{0}}]%
\>[4]\AgdaInductiveConstructor{ƛb}\AgdaSpace{}%
\AgdaSymbol{(}\AgdaInductiveConstructor{ε}\AgdaSpace{}%
\AgdaInductiveConstructor{Zb}\AgdaSpace{}%
\AgdaOperator{\AgdaInductiveConstructor{·b}}\AgdaSpace{}%
\AgdaFunction{renameb}\AgdaSpace{}%
\AgdaInductiveConstructor{Sb}\AgdaSpace{}%
\AgdaSymbol{(}\AgdaFunction{substb}\AgdaSpace{}%
\AgdaBound{fb'}\AgdaSpace{}%
\AgdaSymbol{(}\AgdaFunction{embedTerm}\AgdaSpace{}%
\AgdaBound{M₁}\AgdaSymbol{)))}\AgdaSpace{}%
\AgdaOperator{\AgdaInductiveConstructor{·b}}\AgdaSpace{}%
\AgdaFunction{embedTerm}\AgdaSpace{}%
\AgdaSymbol{(}\AgdaFunction{subst}\AgdaSpace{}%
\AgdaBound{f'}\AgdaSpace{}%
\AgdaBound{M}\AgdaSymbol{)}\<%
\\
%
\>[2]\AgdaFunction{≡⟨}\AgdaSpace{}%
\AgdaFunction{cong₂}\AgdaSpace{}%
\AgdaOperator{\AgdaInductiveConstructor{\AgdaUnderscore{}·b\AgdaUnderscore{}}}\AgdaSpace{}%
\AgdaSymbol{(}\AgdaFunction{cong}\AgdaSpace{}%
\AgdaInductiveConstructor{ƛb}\AgdaSpace{}%
\AgdaSymbol{(}\AgdaFunction{cong₂}\AgdaSpace{}%
\AgdaOperator{\AgdaInductiveConstructor{\AgdaUnderscore{}·b\AgdaUnderscore{}}}\AgdaSpace{}%
\AgdaInductiveConstructor{refl}\AgdaSpace{}%
\AgdaSymbol{(}\AgdaFunction{cong}\AgdaSpace{}%
\AgdaSymbol{(}\AgdaFunction{renameb}\AgdaSpace{}%
\AgdaInductiveConstructor{Sb}\AgdaSymbol{)}\AgdaSpace{}%
\AgdaSymbol{(}\AgdaFunction{pres-simultaneous-subst}\AgdaSpace{}%
\AgdaBound{fb'}\AgdaSpace{}%
\AgdaBound{f'}\AgdaSpace{}%
\AgdaBound{h}\AgdaSpace{}%
\AgdaBound{M₁}\AgdaSymbol{))))}\AgdaSpace{}%
\AgdaInductiveConstructor{refl}\AgdaSpace{}%
\AgdaFunction{⟩}\<%
\\
\>[2][@{}l@{\AgdaIndent{0}}]%
\>[4]\AgdaInductiveConstructor{ƛb}\AgdaSpace{}%
\AgdaSymbol{(}\AgdaInductiveConstructor{ε}\AgdaSpace{}%
\AgdaInductiveConstructor{Zb}\AgdaSpace{}%
\AgdaOperator{\AgdaInductiveConstructor{·b}}\AgdaSpace{}%
\AgdaFunction{renameb}\AgdaSpace{}%
\AgdaInductiveConstructor{Sb}\AgdaSpace{}%
\AgdaSymbol{(}\AgdaFunction{embedTerm}\AgdaSpace{}%
\AgdaSymbol{(}\AgdaFunction{subst}\AgdaSpace{}%
\AgdaBound{f'}\AgdaSpace{}%
\AgdaBound{M₁}\AgdaSymbol{)))}\AgdaSpace{}%
\AgdaOperator{\AgdaInductiveConstructor{·b}}\AgdaSpace{}%
\AgdaFunction{embedTerm}\AgdaSpace{}%
\AgdaSymbol{(}\AgdaFunction{subst}\AgdaSpace{}%
\AgdaBound{f'}\AgdaSpace{}%
\AgdaBound{M}\AgdaSymbol{)}\<%
\\
%
\>[2]\AgdaOperator{\AgdaFunction{∎eq}}\<%
\\
%
\\[\AgdaEmptyExtraSkip]%
\>[0]\AgdaComment{--\ HELPER\ LEMMA'S\ TO\ PROVE\ h}\<%
\\
%
\\[\AgdaEmptyExtraSkip]%
\>[0]\AgdaComment{--\ Lemma\ for\ assumption\ h\ in\ lemma}\<%
\\
\>[0]\AgdaComment{--\ pres-simultaneous-subst}\<%
\\
\>[0]\AgdaFunction{pres-subst-func}\AgdaSpace{}%
\AgdaSymbol{:}\AgdaSpace{}%
\AgdaSymbol{∀}\AgdaSpace{}%
\AgdaSymbol{\{}\AgdaBound{A}\AgdaSpace{}%
\AgdaBound{C}\AgdaSpace{}%
\AgdaBound{Γ}\AgdaSymbol{\}}\<%
\\
\>[0][@{}l@{\AgdaIndent{0}}]%
\>[2]\AgdaSymbol{→}\AgdaSpace{}%
\AgdaSymbol{(}\AgdaBound{V}\AgdaSpace{}%
\AgdaSymbol{:}\AgdaSpace{}%
\AgdaBound{Γ}\AgdaSpace{}%
\AgdaOperator{\AgdaDatatype{⊢}}\AgdaSpace{}%
\AgdaBound{C}\AgdaSymbol{)}\<%
\\
%
\>[2]\AgdaSymbol{→}\AgdaSpace{}%
\AgdaSymbol{(}\AgdaBound{y}\AgdaSpace{}%
\AgdaSymbol{:}\AgdaSpace{}%
\AgdaDatatype{Value}\AgdaSpace{}%
\AgdaBound{V}\AgdaSymbol{)}\<%
\\
%
\>[2]\AgdaSymbol{→}\AgdaSpace{}%
\AgdaSymbol{(}\AgdaBound{x}\AgdaSpace{}%
\AgdaSymbol{:}\AgdaSpace{}%
\AgdaBound{A}\AgdaSpace{}%
\AgdaOperator{\AgdaDatatype{∈}}\AgdaSpace{}%
\AgdaBound{Γ}\AgdaSpace{}%
\AgdaOperator{\AgdaInductiveConstructor{,}}\AgdaSpace{}%
\AgdaBound{C}\AgdaSymbol{)}\<%
\\
%
\>[2]\AgdaSymbol{→}\AgdaSpace{}%
\AgdaInductiveConstructor{box}\AgdaSpace{}%
\AgdaSymbol{(}\AgdaFunction{fb}\AgdaSpace{}%
\AgdaSymbol{(}\AgdaFunction{embedValue}\AgdaSpace{}%
\AgdaBound{V}\AgdaSpace{}%
\AgdaBound{y}\AgdaSymbol{)}\AgdaSpace{}%
\AgdaSymbol{(}\AgdaFunction{embedLookup}\AgdaSpace{}%
\AgdaBound{x}\AgdaSymbol{))}\AgdaSpace{}%
\AgdaOperator{\AgdaDatatype{≡}}\AgdaSpace{}%
\AgdaFunction{embedTerm}\AgdaSpace{}%
\AgdaSymbol{((}\AgdaFunction{f}\AgdaSpace{}%
\AgdaBound{V}\AgdaSymbol{)}\AgdaSpace{}%
\AgdaBound{x}\AgdaSymbol{)}\<%
\\
\>[0]\AgdaFunction{pres-subst-func}\AgdaSpace{}%
\AgdaBound{V}\AgdaSpace{}%
\AgdaBound{y}\AgdaSpace{}%
\AgdaInductiveConstructor{Z}\AgdaSpace{}%
\AgdaSymbol{=}\AgdaSpace{}%
\AgdaFunction{box-embedValue}\AgdaSpace{}%
\AgdaBound{y}\<%
\\
\>[0]\AgdaFunction{pres-subst-func}\AgdaSpace{}%
\AgdaBound{V}\AgdaSpace{}%
\AgdaBound{y}\AgdaSpace{}%
\AgdaSymbol{(}\AgdaInductiveConstructor{S}\AgdaSpace{}%
\AgdaBound{x}\AgdaSymbol{)}\AgdaSpace{}%
\AgdaSymbol{=}\AgdaSpace{}%
\AgdaInductiveConstructor{refl}\<%
\\
%
\\[\AgdaEmptyExtraSkip]%
\>[0]\AgdaComment{--\ exts\ lemma\ for\ assumption\ h\ in\ lemma}\<%
\\
\>[0]\AgdaFunction{exts-pres-subst-func}\AgdaSpace{}%
\AgdaSymbol{:}\AgdaSpace{}%
\AgdaSymbol{∀}\AgdaSpace{}%
\AgdaSymbol{\{}\AgdaBound{A}\AgdaSpace{}%
\AgdaBound{C}\AgdaSpace{}%
\AgdaBound{Γ}\AgdaSpace{}%
\AgdaBound{B}\AgdaSymbol{\}}\<%
\\
\>[0][@{}l@{\AgdaIndent{0}}]%
\>[2]\AgdaSymbol{→}\AgdaSpace{}%
\AgdaSymbol{(}\AgdaBound{V}\AgdaSpace{}%
\AgdaSymbol{:}\AgdaSpace{}%
\AgdaBound{Γ}\AgdaSpace{}%
\AgdaOperator{\AgdaDatatype{⊢}}\AgdaSpace{}%
\AgdaBound{C}\AgdaSymbol{)}\<%
\\
%
\>[2]\AgdaSymbol{→}\AgdaSpace{}%
\AgdaSymbol{(}\AgdaBound{y}\AgdaSpace{}%
\AgdaSymbol{:}\AgdaSpace{}%
\AgdaDatatype{Value}\AgdaSpace{}%
\AgdaBound{V}\AgdaSymbol{)}\<%
\\
%
\>[2]\AgdaSymbol{→}\AgdaSpace{}%
\AgdaSymbol{(}\AgdaBound{x}\AgdaSpace{}%
\AgdaSymbol{:}\AgdaSpace{}%
\AgdaBound{A}\AgdaSpace{}%
\AgdaOperator{\AgdaDatatype{∈}}\AgdaSpace{}%
\AgdaBound{Γ}\AgdaSpace{}%
\AgdaOperator{\AgdaInductiveConstructor{,}}\AgdaSpace{}%
\AgdaBound{C}\AgdaSpace{}%
\AgdaOperator{\AgdaInductiveConstructor{,}}\AgdaSpace{}%
\AgdaBound{B}\AgdaSymbol{)}\<%
\\
%
\>[2]\AgdaSymbol{→}\AgdaSpace{}%
\AgdaInductiveConstructor{box}\AgdaSpace{}%
\AgdaSymbol{(}\AgdaFunction{extsb}\AgdaSpace{}%
\AgdaSymbol{(}\AgdaFunction{fb}\AgdaSpace{}%
\AgdaSymbol{(}\AgdaFunction{embedValue}\AgdaSpace{}%
\AgdaBound{V}\AgdaSpace{}%
\AgdaBound{y}\AgdaSymbol{))}\AgdaSpace{}%
\AgdaSymbol{(}\AgdaFunction{embedLookup}\AgdaSpace{}%
\AgdaBound{x}\AgdaSymbol{))}\AgdaSpace{}%
\AgdaOperator{\AgdaDatatype{≡}}\AgdaSpace{}%
\AgdaFunction{embedTerm}\AgdaSpace{}%
\AgdaSymbol{(}\AgdaFunction{exts}\AgdaSpace{}%
\AgdaSymbol{(}\AgdaFunction{f}\AgdaSpace{}%
\AgdaBound{V}\AgdaSymbol{)}\AgdaSpace{}%
\AgdaBound{x}\AgdaSymbol{)}\<%
\\
\>[0]\AgdaFunction{exts-pres-subst-func}\AgdaSpace{}%
\AgdaBound{V}\AgdaSpace{}%
\AgdaBound{y}\AgdaSpace{}%
\AgdaInductiveConstructor{Z}\AgdaSpace{}%
\AgdaSymbol{=}\AgdaSpace{}%
\AgdaInductiveConstructor{refl}\<%
\\
\>[0]\AgdaFunction{exts-pres-subst-func}\AgdaSpace{}%
\AgdaBound{V}\AgdaSpace{}%
\AgdaBound{y}\AgdaSpace{}%
\AgdaSymbol{(}\AgdaInductiveConstructor{S}\AgdaSpace{}%
\AgdaBound{x}\AgdaSymbol{)}\AgdaSpace{}%
\AgdaSymbol{=}\<%
\\
\>[0][@{}l@{\AgdaIndent{0}}]%
\>[2]\AgdaOperator{\AgdaFunction{beginEq}}\<%
\\
\>[2][@{}l@{\AgdaIndent{0}}]%
\>[4]\AgdaInductiveConstructor{box}\AgdaSpace{}%
\AgdaSymbol{(}\AgdaFunction{renameb}\AgdaSpace{}%
\AgdaInductiveConstructor{Sb}\AgdaSpace{}%
\AgdaSymbol{(}\AgdaFunction{fb}\AgdaSpace{}%
\AgdaSymbol{(}\AgdaFunction{embedValue}\AgdaSpace{}%
\AgdaBound{V}\AgdaSpace{}%
\AgdaBound{y}\AgdaSymbol{)}\AgdaSpace{}%
\AgdaSymbol{(}\AgdaFunction{embedLookup}\AgdaSpace{}%
\AgdaBound{x}\AgdaSymbol{)))}\<%
\\
%
\>[2]\AgdaFunction{≡⟨}\AgdaSpace{}%
\AgdaInductiveConstructor{refl}\AgdaSpace{}%
\AgdaFunction{⟩}\<%
\\
\>[2][@{}l@{\AgdaIndent{0}}]%
\>[4]\AgdaFunction{renameb}\AgdaSpace{}%
\AgdaInductiveConstructor{Sb}\AgdaSpace{}%
\AgdaSymbol{(}\AgdaInductiveConstructor{box}\AgdaSpace{}%
\AgdaSymbol{(}\AgdaFunction{fb}\AgdaSpace{}%
\AgdaSymbol{(}\AgdaFunction{embedValue}\AgdaSpace{}%
\AgdaBound{V}\AgdaSpace{}%
\AgdaBound{y}\AgdaSymbol{)}\AgdaSpace{}%
\AgdaSymbol{(}\AgdaFunction{embedLookup}\AgdaSpace{}%
\AgdaBound{x}\AgdaSymbol{)))}\<%
\\
%
\>[2]\AgdaFunction{≡⟨}\AgdaSpace{}%
\AgdaFunction{cong}\AgdaSpace{}%
\AgdaSymbol{(}\AgdaFunction{renameb}\AgdaSpace{}%
\AgdaInductiveConstructor{Sb}\AgdaSymbol{)}\AgdaSpace{}%
\AgdaSymbol{(}\AgdaFunction{pres-subst-func}\AgdaSpace{}%
\AgdaBound{V}\AgdaSpace{}%
\AgdaBound{y}\AgdaSpace{}%
\AgdaBound{x}\AgdaSymbol{)}\AgdaSpace{}%
\AgdaFunction{⟩}\<%
\\
\>[2][@{}l@{\AgdaIndent{0}}]%
\>[4]\AgdaFunction{renameb}\AgdaSpace{}%
\AgdaInductiveConstructor{Sb}\AgdaSpace{}%
\AgdaSymbol{(}\AgdaFunction{embedTerm}\AgdaSpace{}%
\AgdaSymbol{(}\AgdaSpace{}%
\AgdaSymbol{(}\AgdaFunction{f}\AgdaSpace{}%
\AgdaBound{V}\AgdaSymbol{)}\AgdaSpace{}%
\AgdaBound{x}\AgdaSpace{}%
\AgdaSymbol{))}\<%
\\
%
\>[2]\AgdaFunction{≡⟨}\AgdaSpace{}%
\AgdaFunction{pres-rename}\AgdaSpace{}%
\AgdaInductiveConstructor{S}\AgdaSpace{}%
\AgdaInductiveConstructor{Sb}\AgdaSpace{}%
\AgdaSymbol{(λ}\AgdaSpace{}%
\AgdaBound{x₁}\AgdaSpace{}%
\AgdaSymbol{→}\AgdaSpace{}%
\AgdaInductiveConstructor{refl}\AgdaSymbol{)}\AgdaSpace{}%
\AgdaSymbol{(}\AgdaFunction{f}\AgdaSpace{}%
\AgdaBound{V}\AgdaSpace{}%
\AgdaBound{x}\AgdaSymbol{)}\AgdaSpace{}%
\AgdaFunction{⟩}\<%
\\
\>[2][@{}l@{\AgdaIndent{0}}]%
\>[4]\AgdaFunction{embedTerm}\AgdaSpace{}%
\AgdaSymbol{(}\AgdaFunction{rename}\AgdaSpace{}%
\AgdaInductiveConstructor{S}\AgdaSpace{}%
\AgdaSymbol{(}\AgdaFunction{f}\AgdaSpace{}%
\AgdaBound{V}\AgdaSpace{}%
\AgdaBound{x}\AgdaSymbol{))}\<%
\\
%
\>[2]\AgdaOperator{\AgdaFunction{∎eq}}\<%
\\
%
\\[\AgdaEmptyExtraSkip]%
\>[0]\AgdaComment{--\ Preservation\ and\ Reflection\ of\ }\<%
\\
\>[0]\AgdaComment{--\ single\ substitution\ of\ Gödel's\ translation}\<%
\\
\>[0]\AgdaFunction{pres-subst}\AgdaSpace{}%
\AgdaSymbol{:}\AgdaSpace{}%
\AgdaSymbol{∀}\AgdaSpace{}%
\AgdaSymbol{\{}\AgdaBound{Γ}\AgdaSpace{}%
\AgdaBound{A}\AgdaSpace{}%
\AgdaBound{B}\AgdaSymbol{\}}\AgdaSpace{}%
\AgdaSymbol{→}\AgdaSpace{}%
\AgdaSymbol{(}\AgdaBound{M}\AgdaSpace{}%
\AgdaSymbol{:}\AgdaSpace{}%
\AgdaBound{Γ}\AgdaSpace{}%
\AgdaOperator{\AgdaInductiveConstructor{,}}\AgdaSpace{}%
\AgdaBound{B}\AgdaSpace{}%
\AgdaOperator{\AgdaDatatype{⊢}}\AgdaSpace{}%
\AgdaBound{A}\AgdaSymbol{)}\AgdaSpace{}%
\AgdaSymbol{→}\AgdaSpace{}%
\AgdaSymbol{(}\AgdaBound{V}\AgdaSpace{}%
\AgdaSymbol{:}\AgdaSpace{}%
\AgdaBound{Γ}\AgdaSpace{}%
\AgdaOperator{\AgdaDatatype{⊢}}\AgdaSpace{}%
\AgdaBound{B}\AgdaSymbol{)}\AgdaSpace{}%
\AgdaSymbol{→}\AgdaSpace{}%
\AgdaSymbol{(}\AgdaBound{x}\AgdaSpace{}%
\AgdaSymbol{:}\AgdaSpace{}%
\AgdaDatatype{Value}\AgdaSpace{}%
\AgdaBound{V}\AgdaSymbol{)}\AgdaSpace{}%
\AgdaSymbol{→}\AgdaSpace{}%
\AgdaFunction{embedTerm}\AgdaSpace{}%
\AgdaBound{M}\AgdaSpace{}%
\AgdaOperator{\AgdaFunction{[}}\AgdaSpace{}%
\AgdaFunction{embedValue}\AgdaSpace{}%
\AgdaBound{V}\AgdaSpace{}%
\AgdaBound{x}\AgdaSpace{}%
\AgdaOperator{\AgdaFunction{]b}}\AgdaSpace{}%
\AgdaOperator{\AgdaDatatype{≡}}\AgdaSpace{}%
\AgdaFunction{embedTerm}\AgdaSpace{}%
\AgdaSymbol{(}\AgdaBound{M}\AgdaSpace{}%
\AgdaOperator{\AgdaFunction{[}}\AgdaSpace{}%
\AgdaBound{V}\AgdaSpace{}%
\AgdaOperator{\AgdaFunction{]}}\AgdaSymbol{)}\<%
\\
\>[0]\AgdaFunction{pres-subst}\AgdaSpace{}%
\AgdaBound{M}\AgdaSpace{}%
\AgdaBound{V}\AgdaSpace{}%
\AgdaBound{x}\AgdaSpace{}%
\AgdaSymbol{=}\AgdaSpace{}%
\AgdaFunction{pres-simultaneous-subst}\AgdaSpace{}%
\AgdaSymbol{(}\AgdaFunction{fb}\AgdaSpace{}%
\AgdaSymbol{(}\AgdaFunction{embedValue}\AgdaSpace{}%
\AgdaBound{V}\AgdaSpace{}%
\AgdaBound{x}\AgdaSymbol{))}\AgdaSpace{}%
\AgdaSymbol{(}\AgdaFunction{f}\AgdaSpace{}%
\AgdaBound{V}\AgdaSymbol{)}\AgdaSpace{}%
\AgdaSymbol{(}\AgdaFunction{pres-subst-func}\AgdaSpace{}%
\AgdaBound{V}\AgdaSpace{}%
\AgdaBound{x}\AgdaSymbol{)}\AgdaSpace{}%
\AgdaBound{M}\<%
\\
%
\\[\AgdaEmptyExtraSkip]%
\>[0]\AgdaComment{--\ Helper\ lemma:\ if\ a\ value\ is\ substituted\ in\ a\ value,\ the\ resulting\ lambda\ term\ will\ also\ be\ a\ value}\<%
\\
\>[0]\AgdaFunction{val-after-subst}\AgdaSpace{}%
\AgdaSymbol{:}\AgdaSpace{}%
\AgdaSymbol{∀}\AgdaSpace{}%
\AgdaSymbol{\{}\AgdaBound{Γ}\AgdaSpace{}%
\AgdaBound{A}\AgdaSpace{}%
\AgdaBound{B}\AgdaSymbol{\}}\AgdaSpace{}%
\AgdaSymbol{→}\AgdaSpace{}%
\AgdaSymbol{(}\AgdaBound{M}\AgdaSpace{}%
\AgdaSymbol{:}\AgdaSpace{}%
\AgdaBound{Γ}\AgdaSpace{}%
\AgdaOperator{\AgdaInductiveConstructor{,}}\AgdaSpace{}%
\AgdaBound{B}\AgdaSpace{}%
\AgdaOperator{\AgdaDatatype{⊢}}\AgdaSpace{}%
\AgdaBound{A}\AgdaSymbol{)}\AgdaSpace{}%
\AgdaSymbol{→}\AgdaSpace{}%
\AgdaSymbol{(}\AgdaBound{N}\AgdaSpace{}%
\AgdaSymbol{:}\AgdaSpace{}%
\AgdaBound{Γ}\AgdaSpace{}%
\AgdaOperator{\AgdaDatatype{⊢}}\AgdaSpace{}%
\AgdaBound{B}\AgdaSymbol{)}\AgdaSpace{}%
\AgdaSymbol{→}\AgdaSpace{}%
\AgdaSymbol{(}\AgdaBound{x}\AgdaSpace{}%
\AgdaSymbol{:}\AgdaSpace{}%
\AgdaDatatype{Value}\AgdaSpace{}%
\AgdaBound{M}\AgdaSymbol{)}\AgdaSpace{}%
\AgdaSymbol{→}\AgdaSpace{}%
\AgdaSymbol{(}\AgdaBound{y}\AgdaSpace{}%
\AgdaSymbol{:}\AgdaSpace{}%
\AgdaDatatype{Value}\AgdaSpace{}%
\AgdaBound{N}\AgdaSymbol{)}\AgdaSpace{}%
\AgdaSymbol{→}\AgdaSpace{}%
\AgdaDatatype{Value}\AgdaSpace{}%
\AgdaSymbol{(}\AgdaBound{M}\AgdaSpace{}%
\AgdaOperator{\AgdaFunction{[}}\AgdaSpace{}%
\AgdaBound{N}\AgdaSpace{}%
\AgdaOperator{\AgdaFunction{]}}\AgdaSymbol{)}\<%
\\
\>[0]\AgdaFunction{val-after-subst}\AgdaSpace{}%
\AgdaSymbol{(}\AgdaInductiveConstructor{`}\AgdaSpace{}%
\AgdaInductiveConstructor{Z}\AgdaSymbol{)}\AgdaSpace{}%
\AgdaBound{N}\AgdaSpace{}%
\AgdaBound{x}\AgdaSpace{}%
\AgdaBound{y}\AgdaSpace{}%
\AgdaSymbol{=}\AgdaSpace{}%
\AgdaBound{y}\<%
\\
\>[0]\AgdaFunction{val-after-subst}\AgdaSpace{}%
\AgdaSymbol{(}\AgdaInductiveConstructor{`}\AgdaSpace{}%
\AgdaSymbol{(}\AgdaInductiveConstructor{S}\AgdaSpace{}%
\AgdaBound{x₁}\AgdaSymbol{))}\AgdaSpace{}%
\AgdaBound{N}\AgdaSpace{}%
\AgdaBound{x}\AgdaSpace{}%
\AgdaBound{y}\AgdaSpace{}%
\AgdaSymbol{=}\AgdaSpace{}%
\AgdaInductiveConstructor{V-`}\<%
\\
\>[0]\AgdaFunction{val-after-subst}\AgdaSpace{}%
\AgdaSymbol{(}\AgdaInductiveConstructor{ƛ}\AgdaSpace{}%
\AgdaBound{M}\AgdaSymbol{)}\AgdaSpace{}%
\AgdaBound{N}\AgdaSpace{}%
\AgdaBound{x}\AgdaSpace{}%
\AgdaBound{y}\AgdaSpace{}%
\AgdaSymbol{=}\AgdaSpace{}%
\AgdaInductiveConstructor{V-ƛ}\<%
\\
%
\\[\AgdaEmptyExtraSkip]%
\>[0]\AgdaComment{--\ Preservation\ and\ Reflection\ of\ }\<%
\\
\>[0]\AgdaComment{--\ single\ substitution\ in\ values\ of\ Gödel's\ translation\ }\<%
\\
\>[0]\AgdaFunction{pres-subst-val}\AgdaSpace{}%
\AgdaSymbol{:}\AgdaSpace{}%
\AgdaSymbol{∀}\AgdaSpace{}%
\AgdaSymbol{\{}\AgdaBound{Γ}\AgdaSpace{}%
\AgdaBound{A}\AgdaSpace{}%
\AgdaBound{B}\AgdaSymbol{\}}\<%
\\
\>[0][@{}l@{\AgdaIndent{0}}]%
\>[2]\AgdaSymbol{→}\AgdaSpace{}%
\AgdaSymbol{(}\AgdaBound{M}\AgdaSpace{}%
\AgdaSymbol{:}\AgdaSpace{}%
\AgdaBound{Γ}\AgdaSpace{}%
\AgdaOperator{\AgdaInductiveConstructor{,}}\AgdaSpace{}%
\AgdaBound{B}\AgdaSpace{}%
\AgdaOperator{\AgdaDatatype{⊢}}\AgdaSpace{}%
\AgdaBound{A}\AgdaSymbol{)}\<%
\\
%
\>[2]\AgdaSymbol{→}\AgdaSpace{}%
\AgdaSymbol{(}\AgdaBound{V}\AgdaSpace{}%
\AgdaSymbol{:}\AgdaSpace{}%
\AgdaBound{Γ}\AgdaSpace{}%
\AgdaOperator{\AgdaDatatype{⊢}}\AgdaSpace{}%
\AgdaBound{B}\AgdaSymbol{)}\<%
\\
%
\>[2]\AgdaSymbol{→}\AgdaSpace{}%
\AgdaSymbol{(}\AgdaBound{x}\AgdaSpace{}%
\AgdaSymbol{:}\AgdaSpace{}%
\AgdaDatatype{Value}\AgdaSpace{}%
\AgdaBound{M}\AgdaSymbol{)}\<%
\\
%
\>[2]\AgdaSymbol{→}\AgdaSpace{}%
\AgdaSymbol{(}\AgdaBound{y}\AgdaSpace{}%
\AgdaSymbol{:}\AgdaSpace{}%
\AgdaDatatype{Value}\AgdaSpace{}%
\AgdaBound{V}\AgdaSymbol{)}\<%
\\
%
\>[2]\AgdaSymbol{→}\AgdaSpace{}%
\AgdaFunction{embedValue}\AgdaSpace{}%
\AgdaBound{M}\AgdaSpace{}%
\AgdaBound{x}\AgdaSpace{}%
\AgdaOperator{\AgdaFunction{[}}\AgdaSpace{}%
\AgdaFunction{embedValue}\AgdaSpace{}%
\AgdaBound{V}\AgdaSpace{}%
\AgdaBound{y}\AgdaSpace{}%
\AgdaOperator{\AgdaFunction{]b}}\AgdaSpace{}%
\AgdaOperator{\AgdaDatatype{≡}}\AgdaSpace{}%
\AgdaFunction{embedValue}\AgdaSpace{}%
\AgdaSymbol{(}\AgdaBound{M}\AgdaSpace{}%
\AgdaOperator{\AgdaFunction{[}}\AgdaSpace{}%
\AgdaBound{V}\AgdaSpace{}%
\AgdaOperator{\AgdaFunction{]}}\AgdaSymbol{)}\AgdaSpace{}%
\AgdaSymbol{(}\AgdaFunction{val-after-subst}\AgdaSpace{}%
\AgdaBound{M}\AgdaSpace{}%
\AgdaBound{V}\AgdaSpace{}%
\AgdaBound{x}\AgdaSpace{}%
\AgdaBound{y}\AgdaSymbol{)}\<%
\\
\>[0]\AgdaFunction{pres-subst-val}\AgdaSpace{}%
\AgdaSymbol{(}\AgdaInductiveConstructor{`}\AgdaSpace{}%
\AgdaInductiveConstructor{Z}\AgdaSymbol{)}\AgdaSpace{}%
\AgdaBound{V}\AgdaSpace{}%
\AgdaBound{x}\AgdaSpace{}%
\AgdaBound{y}\AgdaSpace{}%
\AgdaSymbol{=}\AgdaSpace{}%
\AgdaInductiveConstructor{refl}\<%
\\
\>[0]\AgdaFunction{pres-subst-val}\AgdaSpace{}%
\AgdaSymbol{(}\AgdaInductiveConstructor{`}\AgdaSpace{}%
\AgdaSymbol{(}\AgdaInductiveConstructor{S}\AgdaSpace{}%
\AgdaBound{x₁}\AgdaSymbol{))}\AgdaSpace{}%
\AgdaBound{V}\AgdaSpace{}%
\AgdaBound{x}\AgdaSpace{}%
\AgdaBound{y}\AgdaSpace{}%
\AgdaSymbol{=}\AgdaSpace{}%
\AgdaInductiveConstructor{refl}\<%
\\
\>[0]\AgdaFunction{pres-subst-val}\AgdaSpace{}%
\AgdaSymbol{(}\AgdaInductiveConstructor{ƛ}\AgdaSpace{}%
\AgdaBound{M}\AgdaSymbol{)}\AgdaSpace{}%
\AgdaBound{V}\AgdaSpace{}%
\AgdaBound{x}\AgdaSpace{}%
\AgdaBound{y}\AgdaSpace{}%
\AgdaSymbol{=}\<%
\\
\>[0][@{}l@{\AgdaIndent{0}}]%
\>[2]\AgdaOperator{\AgdaFunction{beginEq}}\<%
\\
\>[2][@{}l@{\AgdaIndent{0}}]%
\>[4]\AgdaFunction{embedValue}\AgdaSpace{}%
\AgdaSymbol{(}\AgdaInductiveConstructor{ƛ}\AgdaSpace{}%
\AgdaBound{M}\AgdaSymbol{)}\AgdaSpace{}%
\AgdaBound{x}\AgdaSpace{}%
\AgdaOperator{\AgdaFunction{[}}\AgdaSpace{}%
\AgdaFunction{embedValue}\AgdaSpace{}%
\AgdaBound{V}\AgdaSpace{}%
\AgdaBound{y}\AgdaSpace{}%
\AgdaOperator{\AgdaFunction{]b}}\<%
\\
%
\>[2]\AgdaFunction{≡⟨}\AgdaSpace{}%
\AgdaInductiveConstructor{refl}\AgdaSpace{}%
\AgdaFunction{⟩}\<%
\\
\>[2][@{}l@{\AgdaIndent{0}}]%
\>[4]\AgdaInductiveConstructor{ƛb}\AgdaSpace{}%
\AgdaSymbol{(}\AgdaFunction{substb}\AgdaSpace{}%
\AgdaSymbol{(}\AgdaFunction{extsb}\AgdaSpace{}%
\AgdaSymbol{(}\AgdaFunction{fb}\AgdaSpace{}%
\AgdaSymbol{(}\AgdaFunction{embedValue}\AgdaSpace{}%
\AgdaBound{V}\AgdaSpace{}%
\AgdaBound{y}\AgdaSymbol{)))}\AgdaSpace{}%
\AgdaSymbol{(}\AgdaFunction{embedTerm}\AgdaSpace{}%
\AgdaBound{M}\AgdaSymbol{))}\<%
\\
%
\>[2]\AgdaFunction{≡⟨}\AgdaSpace{}%
\AgdaFunction{cong}\AgdaSpace{}%
\AgdaInductiveConstructor{ƛb}\AgdaSpace{}%
\AgdaSymbol{(}\AgdaFunction{pres-simultaneous-subst}\AgdaSpace{}%
\AgdaSymbol{(}\AgdaFunction{extsb}\AgdaSpace{}%
\AgdaSymbol{(}\AgdaFunction{fb}\AgdaSpace{}%
\AgdaSymbol{(}\AgdaFunction{embedValue}\AgdaSpace{}%
\AgdaBound{V}\AgdaSpace{}%
\AgdaBound{y}\AgdaSymbol{)))}\AgdaSpace{}%
\AgdaSymbol{(}\AgdaFunction{exts}\AgdaSpace{}%
\AgdaSymbol{(}\AgdaFunction{f}\AgdaSpace{}%
\AgdaBound{V}\AgdaSymbol{))}\AgdaSpace{}%
\AgdaSymbol{(}\AgdaFunction{exts-pres-subst-func}\AgdaSpace{}%
\AgdaBound{V}\AgdaSpace{}%
\AgdaBound{y}\AgdaSymbol{)}\AgdaSpace{}%
\AgdaBound{M}\AgdaSymbol{)}\AgdaSpace{}%
\AgdaFunction{⟩}\<%
\\
\>[2][@{}l@{\AgdaIndent{0}}]%
\>[4]\AgdaInductiveConstructor{ƛb}\AgdaSpace{}%
\AgdaSymbol{(}\AgdaFunction{embedTerm}\AgdaSpace{}%
\AgdaSymbol{(}\AgdaFunction{subst}\AgdaSpace{}%
\AgdaSymbol{(}\AgdaFunction{exts}\AgdaSpace{}%
\AgdaSymbol{(}\AgdaFunction{f}\AgdaSpace{}%
\AgdaBound{V}\AgdaSymbol{))}\AgdaSpace{}%
\AgdaBound{M}\AgdaSymbol{))}\<%
\\
%
\>[2]\AgdaFunction{≡⟨}\AgdaSpace{}%
\AgdaInductiveConstructor{refl}\AgdaSpace{}%
\AgdaFunction{⟩}\<%
\\
\>[2][@{}l@{\AgdaIndent{0}}]%
\>[4]\AgdaFunction{embedValue}\AgdaSpace{}%
\AgdaSymbol{((}\AgdaInductiveConstructor{ƛ}\AgdaSpace{}%
\AgdaBound{M}\AgdaSymbol{)}\AgdaSpace{}%
\AgdaOperator{\AgdaFunction{[}}\AgdaSpace{}%
\AgdaBound{V}\AgdaSpace{}%
\AgdaOperator{\AgdaFunction{]}}\AgdaSymbol{)}\AgdaSpace{}%
\AgdaSymbol{(}\AgdaFunction{val-after-subst}\AgdaSpace{}%
\AgdaSymbol{(}\AgdaInductiveConstructor{ƛ}\AgdaSpace{}%
\AgdaBound{M}\AgdaSymbol{)}\AgdaSpace{}%
\AgdaBound{V}\AgdaSpace{}%
\AgdaBound{x}\AgdaSpace{}%
\AgdaBound{y}\AgdaSymbol{)}\<%
\\
%
\>[2]\AgdaOperator{\AgdaFunction{∎eq}}\<%
\\
%
\\[\AgdaEmptyExtraSkip]%
\>[0]\AgdaComment{--\ fb\ and\ renameb\ commute\ under\ extended\ context.}\<%
\\
\>[0]\AgdaFunction{fb-renameb-extb-commute}\AgdaSpace{}%
\AgdaSymbol{:}\AgdaSpace{}%
\AgdaSymbol{∀}\AgdaSpace{}%
\AgdaSymbol{\{}\AgdaBound{Γ}\AgdaSpace{}%
\AgdaBound{A}\AgdaSpace{}%
\AgdaBound{B}\AgdaSpace{}%
\AgdaBound{Δ}\AgdaSymbol{\}}\<%
\\
\>[0][@{}l@{\AgdaIndent{0}}]%
\>[2]\AgdaSymbol{→}\AgdaSpace{}%
\AgdaSymbol{(}\AgdaBound{fb'}\AgdaSpace{}%
\AgdaSymbol{:}\AgdaSpace{}%
\AgdaSymbol{∀}\AgdaSpace{}%
\AgdaSymbol{\{}\AgdaBound{Ab}\AgdaSymbol{\}}\AgdaSpace{}%
\AgdaSymbol{→}\AgdaSpace{}%
\AgdaBound{Ab}\AgdaSpace{}%
\AgdaOperator{\AgdaDatatype{∈b}}\AgdaSpace{}%
\AgdaBound{Γ}\AgdaSpace{}%
\AgdaSymbol{→}\AgdaSpace{}%
\AgdaBound{Ab}\AgdaSpace{}%
\AgdaOperator{\AgdaDatatype{∈b}}\AgdaSpace{}%
\AgdaBound{Δ}\AgdaSymbol{)}\<%
\\
%
\>[2]\AgdaSymbol{→}\AgdaSpace{}%
\AgdaSymbol{(}\AgdaBound{W}\AgdaSpace{}%
\AgdaSymbol{:}\AgdaSpace{}%
\AgdaBound{Γ}\AgdaSpace{}%
\AgdaOperator{\AgdaDatatype{⊢b}}\AgdaSpace{}%
\AgdaBound{B}\AgdaSymbol{)}\<%
\\
%
\>[2]\AgdaSymbol{→}\AgdaSpace{}%
\AgdaSymbol{(}\AgdaBound{x}\AgdaSpace{}%
\AgdaSymbol{:}\AgdaSpace{}%
\AgdaInductiveConstructor{□}\AgdaSpace{}%
\AgdaBound{A}\AgdaSpace{}%
\AgdaOperator{\AgdaDatatype{∈b}}\AgdaSpace{}%
\AgdaBound{Γ}\AgdaSpace{}%
\AgdaOperator{\AgdaInductiveConstructor{,b}}\AgdaSpace{}%
\AgdaInductiveConstructor{□}\AgdaSpace{}%
\AgdaBound{B}\AgdaSymbol{)}\<%
\\
%
\>[2]\AgdaSymbol{→}\AgdaSpace{}%
\AgdaFunction{fb}\AgdaSpace{}%
\AgdaSymbol{(}\AgdaFunction{renameb}\AgdaSpace{}%
\AgdaBound{fb'}\AgdaSpace{}%
\AgdaBound{W}\AgdaSymbol{)}\AgdaSpace{}%
\AgdaSymbol{(}\AgdaFunction{extb}\AgdaSpace{}%
\AgdaBound{fb'}\AgdaSpace{}%
\AgdaBound{x}\AgdaSymbol{)}\AgdaSpace{}%
\AgdaOperator{\AgdaDatatype{≡}}\AgdaSpace{}%
\AgdaFunction{renameb}\AgdaSpace{}%
\AgdaBound{fb'}\AgdaSpace{}%
\AgdaSymbol{(}\AgdaFunction{fb}\AgdaSpace{}%
\AgdaBound{W}\AgdaSpace{}%
\AgdaBound{x}\AgdaSymbol{)}\<%
\\
\>[0]\AgdaFunction{fb-renameb-extb-commute}\AgdaSpace{}%
\AgdaBound{fb'}\AgdaSpace{}%
\AgdaBound{W}\AgdaSpace{}%
\AgdaInductiveConstructor{Zb}\AgdaSpace{}%
\AgdaSymbol{=}\AgdaSpace{}%
\AgdaInductiveConstructor{refl}\<%
\\
\>[0]\AgdaFunction{fb-renameb-extb-commute}\AgdaSpace{}%
\AgdaBound{fb'}\AgdaSpace{}%
\AgdaBound{W}\AgdaSpace{}%
\AgdaSymbol{(}\AgdaInductiveConstructor{Sb}\AgdaSpace{}%
\AgdaBound{x}\AgdaSymbol{)}\AgdaSpace{}%
\AgdaSymbol{=}\AgdaSpace{}%
\AgdaInductiveConstructor{refl}\<%
\\
%
\\[\AgdaEmptyExtraSkip]%
\>[0]\AgdaComment{--\ Preservation\ and\ reflection\ of\ }\<%
\\
\>[0]\AgdaComment{--\ single\ substitution\ under\ renaming}\<%
\\
\>[0]\AgdaFunction{pres-subst-func-rename}\AgdaSpace{}%
\AgdaSymbol{:}\AgdaSpace{}%
\AgdaSymbol{∀}\AgdaSpace{}%
\AgdaSymbol{\{}\AgdaBound{Γ}\AgdaSpace{}%
\AgdaBound{B}\AgdaSpace{}%
\AgdaBound{A}\AgdaSpace{}%
\AgdaBound{Δ}\AgdaSymbol{\}}\<%
\\
\>[0][@{}l@{\AgdaIndent{0}}]%
\>[2]\AgdaSymbol{→}\AgdaSpace{}%
\AgdaSymbol{(}\AgdaBound{fb'}\AgdaSpace{}%
\AgdaSymbol{:}\AgdaSpace{}%
\AgdaSymbol{∀}\AgdaSpace{}%
\AgdaSymbol{\{}\AgdaBound{Ab}\AgdaSymbol{\}}\AgdaSpace{}%
\AgdaSymbol{→}\AgdaSpace{}%
\AgdaBound{Ab}\AgdaSpace{}%
\AgdaOperator{\AgdaDatatype{∈b}}\AgdaSpace{}%
\AgdaBound{Γ}\AgdaSpace{}%
\AgdaSymbol{→}\AgdaSpace{}%
\AgdaBound{Ab}\AgdaSpace{}%
\AgdaOperator{\AgdaDatatype{∈b}}\AgdaSpace{}%
\AgdaBound{Δ}\AgdaSymbol{)}\<%
\\
%
\>[2]\AgdaSymbol{→}\AgdaSpace{}%
\AgdaSymbol{(}\AgdaBound{M}\AgdaSpace{}%
\AgdaSymbol{:}\AgdaSpace{}%
\AgdaBound{Γ}\AgdaSpace{}%
\AgdaOperator{\AgdaInductiveConstructor{,b}}\AgdaSpace{}%
\AgdaInductiveConstructor{□}\AgdaSpace{}%
\AgdaBound{B}\AgdaSpace{}%
\AgdaOperator{\AgdaDatatype{⊢b}}\AgdaSpace{}%
\AgdaBound{A}\AgdaSymbol{)}\<%
\\
%
\>[2]\AgdaSymbol{→}\AgdaSpace{}%
\AgdaSymbol{(}\AgdaBound{W}\AgdaSpace{}%
\AgdaSymbol{:}\AgdaSpace{}%
\AgdaBound{Γ}\AgdaSpace{}%
\AgdaOperator{\AgdaDatatype{⊢b}}\AgdaSpace{}%
\AgdaBound{B}\AgdaSymbol{)}\<%
\\
%
\>[2]\AgdaSymbol{→}%
\>[2253I]\AgdaFunction{renameb}\AgdaSpace{}%
\AgdaSymbol{(}\AgdaFunction{extb}\AgdaSpace{}%
\AgdaBound{fb'}\AgdaSymbol{)}\AgdaSpace{}%
\AgdaBound{M}\AgdaSpace{}%
\AgdaOperator{\AgdaFunction{[}}\AgdaSpace{}%
\AgdaFunction{renameb}\AgdaSpace{}%
\AgdaBound{fb'}\AgdaSpace{}%
\AgdaBound{W}\AgdaSpace{}%
\AgdaOperator{\AgdaFunction{]b}}\<%
\\
\>[.][@{}l@{}]\<[2253I]%
\>[4]\AgdaOperator{\AgdaDatatype{≡}}\AgdaSpace{}%
\AgdaFunction{renameb}\AgdaSpace{}%
\AgdaBound{fb'}\AgdaSpace{}%
\AgdaSymbol{(}\AgdaBound{M}\AgdaSpace{}%
\AgdaOperator{\AgdaFunction{[}}\AgdaSpace{}%
\AgdaBound{W}\AgdaSpace{}%
\AgdaOperator{\AgdaFunction{]b}}\AgdaSymbol{)}\<%
\\
\>[0]\AgdaFunction{pres-subst-func-rename}\AgdaSpace{}%
\AgdaBound{fb'}\AgdaSpace{}%
\AgdaBound{M}\AgdaSpace{}%
\AgdaBound{W}\AgdaSpace{}%
\AgdaSymbol{=}\<%
\\
\>[0][@{}l@{\AgdaIndent{0}}]%
\>[2]\AgdaOperator{\AgdaFunction{beginEq}}\<%
\\
\>[2][@{}l@{\AgdaIndent{0}}]%
\>[4]\AgdaFunction{substb}\AgdaSpace{}%
\AgdaSymbol{(}\AgdaFunction{fb}\AgdaSpace{}%
\AgdaSymbol{(}\AgdaFunction{renameb}\AgdaSpace{}%
\AgdaBound{fb'}\AgdaSpace{}%
\AgdaBound{W}\AgdaSymbol{))}\AgdaSpace{}%
\AgdaSymbol{(}\AgdaFunction{renameb}\AgdaSpace{}%
\AgdaSymbol{(}\AgdaFunction{extb}\AgdaSpace{}%
\AgdaBound{fb'}\AgdaSymbol{)}\AgdaSpace{}%
\AgdaBound{M}\AgdaSymbol{)}\<%
\\
%
\>[2]\AgdaFunction{≡⟨}\AgdaSpace{}%
\AgdaFunction{subst-renameb-comp}\AgdaSpace{}%
\AgdaSymbol{(}\AgdaFunction{extb}\AgdaSpace{}%
\AgdaBound{fb'}\AgdaSymbol{)}\AgdaSpace{}%
\AgdaSymbol{(}\AgdaFunction{fb}\AgdaSpace{}%
\AgdaSymbol{(}\AgdaFunction{renameb}\AgdaSpace{}%
\AgdaBound{fb'}\AgdaSpace{}%
\AgdaBound{W}\AgdaSymbol{))}\AgdaSpace{}%
\AgdaBound{M}\AgdaSpace{}%
\AgdaFunction{⟩}\<%
\\
\>[2][@{}l@{\AgdaIndent{0}}]%
\>[4]\AgdaFunction{substb}\AgdaSpace{}%
\AgdaSymbol{(λ}\AgdaSpace{}%
\AgdaBound{x}\AgdaSpace{}%
\AgdaSymbol{→}\AgdaSpace{}%
\AgdaFunction{fb}\AgdaSpace{}%
\AgdaSymbol{(}\AgdaFunction{renameb}\AgdaSpace{}%
\AgdaBound{fb'}\AgdaSpace{}%
\AgdaBound{W}\AgdaSymbol{)}\AgdaSpace{}%
\AgdaSymbol{(}\AgdaFunction{extb}\AgdaSpace{}%
\AgdaBound{fb'}\AgdaSpace{}%
\AgdaBound{x}\AgdaSymbol{))}\AgdaSpace{}%
\AgdaBound{M}\<%
\\
%
\>[2]\AgdaFunction{≡⟨}\AgdaSpace{}%
\AgdaFunction{substb-cong}\AgdaSpace{}%
\AgdaSymbol{(}\AgdaFunction{fb-renameb-extb-commute}\AgdaSpace{}%
\AgdaBound{fb'}\AgdaSpace{}%
\AgdaBound{W}\AgdaSymbol{)}\AgdaSpace{}%
\AgdaBound{M}\AgdaSpace{}%
\AgdaFunction{⟩}\<%
\\
\>[2][@{}l@{\AgdaIndent{0}}]%
\>[4]\AgdaFunction{substb}\AgdaSpace{}%
\AgdaSymbol{(λ}\AgdaSpace{}%
\AgdaBound{x}\AgdaSpace{}%
\AgdaSymbol{→}\AgdaSpace{}%
\AgdaFunction{renameb}\AgdaSpace{}%
\AgdaBound{fb'}\AgdaSpace{}%
\AgdaSymbol{(}\AgdaFunction{fb}\AgdaSpace{}%
\AgdaBound{W}\AgdaSpace{}%
\AgdaBound{x}\AgdaSymbol{))}\AgdaSpace{}%
\AgdaBound{M}\<%
\\
%
\>[2]\AgdaFunction{≡⟨}\AgdaSpace{}%
\AgdaFunction{substb-rename-commute}\AgdaSpace{}%
\AgdaBound{fb'}\AgdaSpace{}%
\AgdaSymbol{(}\AgdaFunction{fb}\AgdaSpace{}%
\AgdaBound{W}\AgdaSymbol{)}\AgdaSpace{}%
\AgdaBound{M}\AgdaSpace{}%
\AgdaFunction{⟩}\<%
\\
\>[2][@{}l@{\AgdaIndent{0}}]%
\>[4]\AgdaFunction{renameb}\AgdaSpace{}%
\AgdaBound{fb'}\AgdaSpace{}%
\AgdaSymbol{(}\AgdaFunction{substb}\AgdaSpace{}%
\AgdaSymbol{(}\AgdaFunction{fb}\AgdaSpace{}%
\AgdaBound{W}\AgdaSymbol{)}\AgdaSpace{}%
\AgdaBound{M}\AgdaSymbol{)}\<%
\\
%
\>[2]\AgdaOperator{\AgdaFunction{∎eq}}\<%
\\
%
\\[\AgdaEmptyExtraSkip]%
\>[0]\AgdaComment{--\ Preservation\ of\ reduction\ under\ renaming\ \ \ }\<%
\\
\>[0]\AgdaFunction{pres-rename-red}\AgdaSpace{}%
\AgdaSymbol{:}\AgdaSpace{}%
\AgdaSymbol{∀}\AgdaSpace{}%
\AgdaSymbol{\{}\AgdaBound{Γ}\AgdaSpace{}%
\AgdaBound{A}\AgdaSpace{}%
\AgdaBound{Δ}\AgdaSymbol{\}}\<%
\\
\>[0][@{}l@{\AgdaIndent{0}}]%
\>[2]\AgdaSymbol{→}\AgdaSpace{}%
\AgdaSymbol{(}\AgdaBound{fb'}\AgdaSpace{}%
\AgdaSymbol{:}\AgdaSpace{}%
\AgdaSymbol{∀}\AgdaSpace{}%
\AgdaSymbol{\{}\AgdaBound{Ab}\AgdaSymbol{\}}\AgdaSpace{}%
\AgdaSymbol{→}\AgdaSpace{}%
\AgdaBound{Ab}\AgdaSpace{}%
\AgdaOperator{\AgdaDatatype{∈b}}\AgdaSpace{}%
\AgdaBound{Γ}\AgdaSpace{}%
\AgdaSymbol{→}\AgdaSpace{}%
\AgdaBound{Ab}\AgdaSpace{}%
\AgdaOperator{\AgdaDatatype{∈b}}\AgdaSpace{}%
\AgdaBound{Δ}\AgdaSymbol{)}\<%
\\
%
\>[2]\AgdaSymbol{→}\AgdaSpace{}%
\AgdaSymbol{(}\AgdaBound{M}\AgdaSpace{}%
\AgdaSymbol{:}\AgdaSpace{}%
\AgdaBound{Γ}\AgdaSpace{}%
\AgdaOperator{\AgdaDatatype{⊢b}}\AgdaSpace{}%
\AgdaBound{A}\AgdaSymbol{)}\<%
\\
%
\>[2]\AgdaSymbol{→}\AgdaSpace{}%
\AgdaSymbol{(}\AgdaBound{N}\AgdaSpace{}%
\AgdaSymbol{:}\AgdaSpace{}%
\AgdaBound{Γ}\AgdaSpace{}%
\AgdaOperator{\AgdaDatatype{⊢b}}\AgdaSpace{}%
\AgdaBound{A}\AgdaSymbol{)}\<%
\\
%
\>[2]\AgdaSymbol{→}\AgdaSpace{}%
\AgdaSymbol{(}\AgdaBound{M}\AgdaSpace{}%
\AgdaOperator{\AgdaDatatype{↝βb}}\AgdaSpace{}%
\AgdaBound{N}\AgdaSymbol{)}\<%
\\
%
\>[2]\AgdaSymbol{→}\AgdaSpace{}%
\AgdaFunction{renameb}\AgdaSpace{}%
\AgdaBound{fb'}\AgdaSpace{}%
\AgdaBound{M}\AgdaSpace{}%
\AgdaOperator{\AgdaDatatype{↝βb}}\AgdaSpace{}%
\AgdaFunction{renameb}\AgdaSpace{}%
\AgdaBound{fb'}\AgdaSpace{}%
\AgdaBound{N}\<%
\\
\>[0]\AgdaFunction{pres-rename-red}\AgdaSpace{}%
\AgdaBound{fb'}\AgdaSpace{}%
\AgdaSymbol{(}\AgdaInductiveConstructor{ƛb}\AgdaSpace{}%
\AgdaBound{M₁}\AgdaSpace{}%
\AgdaOperator{\AgdaInductiveConstructor{·b}}\AgdaSpace{}%
\AgdaInductiveConstructor{box}\AgdaSpace{}%
\AgdaBound{M₂}\AgdaSymbol{)}\AgdaSpace{}%
\AgdaBound{N}\AgdaSpace{}%
\AgdaInductiveConstructor{βb}\AgdaSpace{}%
\AgdaSymbol{=}\AgdaSpace{}%
\AgdaFunction{Eq.subst}\AgdaSpace{}%
\AgdaSymbol{(λ}\AgdaSpace{}%
\AgdaBound{x}\AgdaSpace{}%
\AgdaSymbol{→}\AgdaSpace{}%
\AgdaInductiveConstructor{ƛb}\AgdaSpace{}%
\AgdaSymbol{(}\AgdaFunction{renameb}\AgdaSpace{}%
\AgdaSymbol{(}\AgdaFunction{extb}\AgdaSpace{}%
\AgdaBound{fb'}\AgdaSymbol{)}\AgdaSpace{}%
\AgdaBound{M₁}\AgdaSymbol{)}\AgdaSpace{}%
\AgdaOperator{\AgdaInductiveConstructor{·b}}\AgdaSpace{}%
\AgdaInductiveConstructor{box}\AgdaSpace{}%
\AgdaSymbol{(}\AgdaFunction{renameb}\AgdaSpace{}%
\AgdaBound{fb'}\AgdaSpace{}%
\AgdaBound{M₂}\AgdaSymbol{)}\AgdaSpace{}%
\AgdaOperator{\AgdaDatatype{↝βb}}\AgdaSpace{}%
\AgdaBound{x}\AgdaSymbol{)}\AgdaSpace{}%
\AgdaSymbol{(}\AgdaFunction{pres-subst-func-rename}\AgdaSpace{}%
\AgdaBound{fb'}\AgdaSpace{}%
\AgdaBound{M₁}\AgdaSpace{}%
\AgdaBound{M₂}\AgdaSymbol{)}\AgdaSpace{}%
\AgdaInductiveConstructor{βb}\<%
\\
\>[0]\AgdaFunction{pres-rename-red}\AgdaSpace{}%
\AgdaBound{fb'}\AgdaSpace{}%
\AgdaBound{M}\AgdaSpace{}%
\AgdaBound{N}\AgdaSpace{}%
\AgdaSymbol{(}\AgdaInductiveConstructor{μ}\AgdaSpace{}%
\AgdaSymbol{\{}\AgdaArgument{M}\AgdaSpace{}%
\AgdaSymbol{=}\AgdaSpace{}%
\AgdaBound{M₁}\AgdaSymbol{\}}\AgdaSpace{}%
\AgdaSymbol{\{}\AgdaArgument{M'}\AgdaSpace{}%
\AgdaSymbol{=}\AgdaSpace{}%
\AgdaBound{M'}\AgdaSymbol{\}}\AgdaSpace{}%
\AgdaBound{red}\AgdaSymbol{)}\AgdaSpace{}%
\AgdaSymbol{=}\AgdaSpace{}%
\AgdaInductiveConstructor{μ}\AgdaSpace{}%
\AgdaSymbol{(}\AgdaFunction{pres-rename-red}\AgdaSpace{}%
\AgdaBound{fb'}\AgdaSpace{}%
\AgdaBound{M₁}\AgdaSpace{}%
\AgdaBound{M'}\AgdaSpace{}%
\AgdaBound{red}\AgdaSymbol{)}\<%
\\
\>[0]\AgdaFunction{pres-rename-red}\AgdaSpace{}%
\AgdaBound{fb'}\AgdaSpace{}%
\AgdaBound{M}\AgdaSpace{}%
\AgdaBound{N}\AgdaSpace{}%
\AgdaSymbol{(}\AgdaInductiveConstructor{ν}\AgdaSpace{}%
\AgdaSymbol{\{}\AgdaArgument{M}\AgdaSpace{}%
\AgdaSymbol{=}\AgdaSpace{}%
\AgdaBound{M₁}\AgdaSymbol{\}}\AgdaSpace{}%
\AgdaSymbol{\{}\AgdaArgument{N}\AgdaSpace{}%
\AgdaSymbol{=}\AgdaSpace{}%
\AgdaBound{N₁}\AgdaSymbol{\}}\AgdaSpace{}%
\AgdaSymbol{\{}\AgdaArgument{N'}\AgdaSpace{}%
\AgdaSymbol{=}\AgdaSpace{}%
\AgdaBound{N'}\AgdaSymbol{\}}\AgdaSpace{}%
\AgdaBound{red}\AgdaSymbol{)}\AgdaSpace{}%
\AgdaSymbol{=}\AgdaSpace{}%
\AgdaInductiveConstructor{ν}\AgdaSpace{}%
\AgdaSymbol{(}\AgdaFunction{pres-rename-red}\AgdaSpace{}%
\AgdaBound{fb'}\AgdaSpace{}%
\AgdaBound{N₁}\AgdaSpace{}%
\AgdaBound{N'}\AgdaSpace{}%
\AgdaBound{red}\AgdaSymbol{)}\<%
\\
\>[0]\AgdaFunction{pres-rename-red}\AgdaSpace{}%
\AgdaBound{fb'}\AgdaSpace{}%
\AgdaBound{M}\AgdaSpace{}%
\AgdaBound{N}\AgdaSpace{}%
\AgdaSymbol{(}\AgdaInductiveConstructor{ξ}\AgdaSpace{}%
\AgdaSymbol{\{}\AgdaArgument{M}\AgdaSpace{}%
\AgdaSymbol{=}\AgdaSpace{}%
\AgdaBound{M₁}\AgdaSymbol{\}}\AgdaSpace{}%
\AgdaSymbol{\{}\AgdaArgument{M'}\AgdaSpace{}%
\AgdaSymbol{=}\AgdaSpace{}%
\AgdaBound{M'}\AgdaSymbol{\}}\AgdaSpace{}%
\AgdaBound{red}\AgdaSymbol{)}\AgdaSpace{}%
\AgdaSymbol{=}\AgdaSpace{}%
\AgdaInductiveConstructor{ξ}\AgdaSpace{}%
\AgdaSymbol{(}\AgdaFunction{pres-rename-red}\AgdaSpace{}%
\AgdaSymbol{(}\AgdaFunction{extb}\AgdaSpace{}%
\AgdaBound{fb'}\AgdaSymbol{)}\AgdaSpace{}%
\AgdaBound{M₁}\AgdaSpace{}%
\AgdaBound{M'}\AgdaSpace{}%
\AgdaBound{red}\AgdaSymbol{)}\<%
\\
\>[0]\AgdaFunction{pres-rename-red}\AgdaSpace{}%
\AgdaBound{fb'}\AgdaSpace{}%
\AgdaBound{M}\AgdaSpace{}%
\AgdaBound{N}\AgdaSpace{}%
\AgdaSymbol{(}\AgdaInductiveConstructor{ζ}\AgdaSpace{}%
\AgdaSymbol{\{}\AgdaArgument{M}\AgdaSpace{}%
\AgdaSymbol{=}\AgdaSpace{}%
\AgdaBound{M₁}\AgdaSymbol{\}}\AgdaSpace{}%
\AgdaSymbol{\{}\AgdaArgument{M'}\AgdaSpace{}%
\AgdaSymbol{=}\AgdaSpace{}%
\AgdaBound{M'}\AgdaSymbol{\}}\AgdaSpace{}%
\AgdaBound{red}\AgdaSymbol{)}\AgdaSpace{}%
\AgdaSymbol{=}\AgdaSpace{}%
\AgdaInductiveConstructor{ζ}\AgdaSpace{}%
\AgdaSymbol{(}\AgdaFunction{pres-rename-red}\AgdaSpace{}%
\AgdaBound{fb'}\AgdaSpace{}%
\AgdaBound{M₁}\AgdaSpace{}%
\AgdaBound{M'}\AgdaSpace{}%
\AgdaBound{red}\AgdaSymbol{)}\<%
\\
\>[0]\AgdaComment{--\ Preservation\ of\ 2\ step\ reduction}\<%
\\
\>[0]\AgdaComment{--\ under\ renaming}\<%
\\
\>[0]\AgdaFunction{pres-rename-red-2}\AgdaSpace{}%
\AgdaSymbol{:}\AgdaSpace{}%
\AgdaSymbol{∀}\AgdaSpace{}%
\AgdaSymbol{\{}\AgdaBound{Γ}\AgdaSpace{}%
\AgdaBound{A}\AgdaSpace{}%
\AgdaBound{Δ}\AgdaSymbol{\}}\<%
\\
\>[0][@{}l@{\AgdaIndent{0}}]%
\>[2]\AgdaSymbol{→}\AgdaSpace{}%
\AgdaSymbol{(}\AgdaBound{fb'}\AgdaSpace{}%
\AgdaSymbol{:}\AgdaSpace{}%
\AgdaSymbol{∀}\AgdaSpace{}%
\AgdaSymbol{\{}\AgdaBound{Ab}\AgdaSymbol{\}}\AgdaSpace{}%
\AgdaSymbol{→}\AgdaSpace{}%
\AgdaBound{Ab}\AgdaSpace{}%
\AgdaOperator{\AgdaDatatype{∈b}}\AgdaSpace{}%
\AgdaBound{Γ}\AgdaSpace{}%
\AgdaSymbol{→}\AgdaSpace{}%
\AgdaBound{Ab}\AgdaSpace{}%
\AgdaOperator{\AgdaDatatype{∈b}}\AgdaSpace{}%
\AgdaBound{Δ}\AgdaSymbol{)}\<%
\\
%
\>[2]\AgdaSymbol{→}\AgdaSpace{}%
\AgdaSymbol{(}\AgdaBound{M}\AgdaSpace{}%
\AgdaSymbol{:}\AgdaSpace{}%
\AgdaBound{Γ}\AgdaSpace{}%
\AgdaOperator{\AgdaDatatype{⊢b}}\AgdaSpace{}%
\AgdaBound{A}\AgdaSymbol{)}\<%
\\
%
\>[2]\AgdaSymbol{→}\AgdaSpace{}%
\AgdaSymbol{(}\AgdaBound{N}\AgdaSpace{}%
\AgdaSymbol{:}\AgdaSpace{}%
\AgdaBound{Γ}\AgdaSpace{}%
\AgdaOperator{\AgdaDatatype{⊢b}}\AgdaSpace{}%
\AgdaBound{A}\AgdaSymbol{)}\<%
\\
%
\>[2]\AgdaSymbol{→}\AgdaSpace{}%
\AgdaSymbol{(}\AgdaBound{M}\AgdaSpace{}%
\AgdaOperator{\AgdaDatatype{↝βb²}}\AgdaSpace{}%
\AgdaBound{N}\AgdaSymbol{)}\<%
\\
%
\>[2]\AgdaSymbol{→}\AgdaSpace{}%
\AgdaFunction{renameb}\AgdaSpace{}%
\AgdaBound{fb'}\AgdaSpace{}%
\AgdaBound{M}\AgdaSpace{}%
\AgdaOperator{\AgdaDatatype{↝βb²}}\AgdaSpace{}%
\AgdaFunction{renameb}\AgdaSpace{}%
\AgdaBound{fb'}\AgdaSpace{}%
\AgdaBound{N}\<%
\\
\>[0]\AgdaFunction{pres-rename-red-2}\AgdaSpace{}%
\AgdaBound{fb'}\AgdaSpace{}%
\AgdaBound{M}\AgdaSpace{}%
\AgdaBound{N}\AgdaSpace{}%
\AgdaSymbol{(}\AgdaInductiveConstructor{βb²}\AgdaSpace{}%
\AgdaDottedPattern{\AgdaSymbol{.}}\AgdaDottedPattern{\AgdaBound{M}}\AgdaSpace{}%
\AgdaBound{t₂}\AgdaSpace{}%
\AgdaDottedPattern{\AgdaSymbol{.}}\AgdaDottedPattern{\AgdaBound{N}}\AgdaSpace{}%
\AgdaBound{x}\AgdaSpace{}%
\AgdaBound{x₁}\AgdaSymbol{)}\AgdaSpace{}%
\AgdaSymbol{=}\<%
\\
\>[0][@{}l@{\AgdaIndent{0}}]%
\>[2]\AgdaInductiveConstructor{βb²}\<%
\\
\>[2][@{}l@{\AgdaIndent{0}}]%
\>[4]\AgdaSymbol{(}\AgdaFunction{renameb}\AgdaSpace{}%
\AgdaBound{fb'}\AgdaSpace{}%
\AgdaBound{M}\AgdaSymbol{)}\<%
\\
%
\>[4]\AgdaSymbol{(}\AgdaFunction{renameb}\AgdaSpace{}%
\AgdaBound{fb'}\AgdaSpace{}%
\AgdaBound{t₂}\AgdaSymbol{)}\<%
\\
%
\>[4]\AgdaSymbol{(}\AgdaFunction{renameb}\AgdaSpace{}%
\AgdaBound{fb'}\AgdaSpace{}%
\AgdaBound{N}\AgdaSymbol{)}\<%
\\
%
\>[4]\AgdaSymbol{(}\AgdaFunction{pres-rename-red}\AgdaSpace{}%
\AgdaBound{fb'}\AgdaSpace{}%
\AgdaBound{M}\AgdaSpace{}%
\AgdaBound{t₂}\AgdaSpace{}%
\AgdaBound{x}\AgdaSymbol{)}\<%
\\
%
\>[4]\AgdaSymbol{(}\AgdaFunction{pres-rename-red}\AgdaSpace{}%
\AgdaBound{fb'}\AgdaSpace{}%
\AgdaBound{t₂}\AgdaSpace{}%
\AgdaBound{N}\AgdaSpace{}%
\AgdaBound{x₁}\AgdaSymbol{)}\<%
\\
\>[0]\AgdaFunction{pres-rename-red-2}\AgdaSpace{}%
\AgdaBound{fb'}\AgdaSpace{}%
\AgdaBound{M}\AgdaSpace{}%
\AgdaBound{N}\AgdaSpace{}%
\AgdaSymbol{(}\AgdaInductiveConstructor{μ}\AgdaSpace{}%
\AgdaSymbol{\{}\AgdaArgument{M}\AgdaSpace{}%
\AgdaSymbol{=}\AgdaSpace{}%
\AgdaBound{M₁}\AgdaSymbol{\}}\AgdaSpace{}%
\AgdaSymbol{\{}\AgdaArgument{M'}\AgdaSpace{}%
\AgdaSymbol{=}\AgdaSpace{}%
\AgdaBound{M'}\AgdaSymbol{\}}\AgdaSpace{}%
\AgdaBound{red}\AgdaSymbol{)}\AgdaSpace{}%
\AgdaSymbol{=}\<%
\\
\>[0][@{}l@{\AgdaIndent{0}}]%
\>[2]\AgdaInductiveConstructor{μ}\AgdaSpace{}%
\AgdaSymbol{(}\AgdaFunction{pres-rename-red-2}\AgdaSpace{}%
\AgdaBound{fb'}\AgdaSpace{}%
\AgdaBound{M₁}\AgdaSpace{}%
\AgdaBound{M'}\AgdaSpace{}%
\AgdaBound{red}\AgdaSymbol{)}\<%
\\
\>[0]\AgdaFunction{pres-rename-red-2}\AgdaSpace{}%
\AgdaBound{fb'}\AgdaSpace{}%
\AgdaBound{M}\AgdaSpace{}%
\AgdaBound{N}\AgdaSpace{}%
\AgdaSymbol{(}\AgdaInductiveConstructor{ν}\AgdaSpace{}%
\AgdaSymbol{\{}\AgdaArgument{M}\AgdaSpace{}%
\AgdaSymbol{=}\AgdaSpace{}%
\AgdaBound{M₁}\AgdaSymbol{\}}\AgdaSpace{}%
\AgdaSymbol{\{}\AgdaArgument{N}\AgdaSpace{}%
\AgdaSymbol{=}\AgdaSpace{}%
\AgdaBound{N₁}\AgdaSymbol{\}}\AgdaSpace{}%
\AgdaSymbol{\{}\AgdaArgument{N'}\AgdaSpace{}%
\AgdaSymbol{=}\AgdaSpace{}%
\AgdaBound{N'}\AgdaSymbol{\}}\AgdaSpace{}%
\AgdaBound{red}\AgdaSymbol{)}\AgdaSpace{}%
\AgdaSymbol{=}\<%
\\
\>[0][@{}l@{\AgdaIndent{0}}]%
\>[2]\AgdaInductiveConstructor{ν}\AgdaSpace{}%
\AgdaSymbol{(}\AgdaFunction{pres-rename-red-2}\AgdaSpace{}%
\AgdaBound{fb'}\AgdaSpace{}%
\AgdaBound{N₁}\AgdaSpace{}%
\AgdaBound{N'}\AgdaSpace{}%
\AgdaBound{red}\AgdaSymbol{)}\<%
\\
\>[0]\AgdaFunction{pres-rename-red-2}\AgdaSpace{}%
\AgdaBound{fb'}\AgdaSpace{}%
\AgdaBound{M}\AgdaSpace{}%
\AgdaBound{N}\AgdaSpace{}%
\AgdaSymbol{(}\AgdaInductiveConstructor{ξ}\AgdaSpace{}%
\AgdaSymbol{\{}\AgdaArgument{M}\AgdaSpace{}%
\AgdaSymbol{=}\AgdaSpace{}%
\AgdaBound{M₁}\AgdaSymbol{\}}\AgdaSpace{}%
\AgdaSymbol{\{}\AgdaArgument{M'}\AgdaSpace{}%
\AgdaSymbol{=}\AgdaSpace{}%
\AgdaBound{M'}\AgdaSymbol{\}}\AgdaSpace{}%
\AgdaBound{red}\AgdaSymbol{)}\AgdaSpace{}%
\AgdaSymbol{=}\<%
\\
\>[0][@{}l@{\AgdaIndent{0}}]%
\>[2]\AgdaInductiveConstructor{ξ}\AgdaSpace{}%
\AgdaSymbol{(}\AgdaFunction{pres-rename-red-2}\AgdaSpace{}%
\AgdaSymbol{(}\AgdaFunction{extb}\AgdaSpace{}%
\AgdaBound{fb'}\AgdaSymbol{)}\AgdaSpace{}%
\AgdaBound{M₁}\AgdaSpace{}%
\AgdaBound{M'}\AgdaSpace{}%
\AgdaBound{red}\AgdaSymbol{)}\<%
\\
\>[0]\AgdaFunction{pres-rename-red-2}\AgdaSpace{}%
\AgdaBound{fb'}\AgdaSpace{}%
\AgdaBound{M}\AgdaSpace{}%
\AgdaBound{N}\AgdaSpace{}%
\AgdaSymbol{(}\AgdaInductiveConstructor{ζ}\AgdaSpace{}%
\AgdaSymbol{\{}\AgdaArgument{M}\AgdaSpace{}%
\AgdaSymbol{=}\AgdaSpace{}%
\AgdaBound{M₁}\AgdaSymbol{\}}\AgdaSpace{}%
\AgdaSymbol{\{}\AgdaArgument{M'}\AgdaSpace{}%
\AgdaSymbol{=}\AgdaSpace{}%
\AgdaBound{M'}\AgdaSymbol{\}}\AgdaSpace{}%
\AgdaBound{red}\AgdaSymbol{)}\AgdaSpace{}%
\AgdaSymbol{=}\<%
\\
\>[0][@{}l@{\AgdaIndent{0}}]%
\>[2]\AgdaInductiveConstructor{ζ}\AgdaSpace{}%
\AgdaSymbol{(}\AgdaFunction{pres-rename-red-2}\AgdaSpace{}%
\AgdaBound{fb'}\AgdaSpace{}%
\AgdaBound{M₁}\AgdaSpace{}%
\AgdaBound{M'}\AgdaSpace{}%
\AgdaBound{red}\AgdaSymbol{)}\<%
\\
%
\\[\AgdaEmptyExtraSkip]%
\>[0]\AgdaComment{--\ extsb\ of\ ε\ is\ ε}\<%
\\
\>[0]\AgdaFunction{ε-exts}\AgdaSpace{}%
\AgdaSymbol{:}\AgdaSpace{}%
\AgdaSymbol{∀}\AgdaSpace{}%
\AgdaSymbol{\{}\AgdaBound{Γ}\AgdaSpace{}%
\AgdaBound{A}\AgdaSpace{}%
\AgdaBound{B}\AgdaSymbol{\}}\<%
\\
\>[0][@{}l@{\AgdaIndent{0}}]%
\>[2]\AgdaSymbol{→}\AgdaSpace{}%
\AgdaSymbol{(}\AgdaBound{x}\AgdaSpace{}%
\AgdaSymbol{:}\AgdaSpace{}%
\AgdaInductiveConstructor{□}\AgdaSpace{}%
\AgdaBound{A}\AgdaSpace{}%
\AgdaOperator{\AgdaDatatype{∈b}}\AgdaSpace{}%
\AgdaBound{Γ}\AgdaSpace{}%
\AgdaOperator{\AgdaInductiveConstructor{,b}}\AgdaSpace{}%
\AgdaBound{B}\AgdaSymbol{)}\<%
\\
%
\>[2]\AgdaSymbol{→}\AgdaSpace{}%
\AgdaFunction{extsb}\AgdaSpace{}%
\AgdaInductiveConstructor{ε}\AgdaSpace{}%
\AgdaBound{x}\AgdaSpace{}%
\AgdaOperator{\AgdaDatatype{≡}}\AgdaSpace{}%
\AgdaInductiveConstructor{ε}\AgdaSpace{}%
\AgdaBound{x}\<%
\\
\>[0]\AgdaFunction{ε-exts}\AgdaSpace{}%
\AgdaInductiveConstructor{Zb}\AgdaSpace{}%
\AgdaSymbol{=}\AgdaSpace{}%
\AgdaInductiveConstructor{refl}\<%
\\
\>[0]\AgdaFunction{ε-exts}\AgdaSpace{}%
\AgdaSymbol{(}\AgdaInductiveConstructor{Sb}\AgdaSpace{}%
\AgdaBound{x}\AgdaSymbol{)}\AgdaSpace{}%
\AgdaSymbol{=}\AgdaSpace{}%
\AgdaInductiveConstructor{refl}\<%
\\
%
\\[\AgdaEmptyExtraSkip]%
\>[0]\AgdaComment{--\ Simultaneous\ substitution\ with\ ε\ }\<%
\\
\>[0]\AgdaComment{--\ leaves\ the\ term\ unchanged\ \ }\<%
\\
\>[0]\AgdaFunction{subst-with-id}\AgdaSpace{}%
\AgdaSymbol{:}\AgdaSpace{}%
\AgdaSymbol{∀}\AgdaSpace{}%
\AgdaSymbol{\{}\AgdaBound{Γ}\AgdaSpace{}%
\AgdaBound{A}\AgdaSymbol{\}}\<%
\\
\>[0][@{}l@{\AgdaIndent{0}}]%
\>[2]\AgdaSymbol{→}\AgdaSpace{}%
\AgdaSymbol{(}\AgdaBound{M}\AgdaSpace{}%
\AgdaSymbol{:}\AgdaSpace{}%
\AgdaBound{Γ}\AgdaSpace{}%
\AgdaOperator{\AgdaDatatype{⊢b}}\AgdaSpace{}%
\AgdaBound{A}\AgdaSymbol{)}\<%
\\
%
\>[2]\AgdaSymbol{→}\AgdaSpace{}%
\AgdaFunction{substb}\AgdaSpace{}%
\AgdaSymbol{(λ}\AgdaSpace{}%
\AgdaBound{x₁}\AgdaSpace{}%
\AgdaSymbol{→}\AgdaSpace{}%
\AgdaInductiveConstructor{ε}\AgdaSpace{}%
\AgdaBound{x₁}\AgdaSymbol{)}\AgdaSpace{}%
\AgdaBound{M}\AgdaSpace{}%
\AgdaOperator{\AgdaDatatype{≡}}\AgdaSpace{}%
\AgdaBound{M}\<%
\\
\>[0]\AgdaFunction{subst-with-id}\AgdaSpace{}%
\AgdaSymbol{(}\AgdaInductiveConstructor{ε}\AgdaSpace{}%
\AgdaBound{x}\AgdaSymbol{)}\AgdaSpace{}%
\AgdaSymbol{=}\AgdaSpace{}%
\AgdaInductiveConstructor{refl}\<%
\\
\>[0]\AgdaFunction{subst-with-id}\AgdaSpace{}%
\AgdaSymbol{(}\AgdaInductiveConstructor{box}\AgdaSpace{}%
\AgdaBound{M}\AgdaSymbol{)}\AgdaSpace{}%
\AgdaSymbol{=}\AgdaSpace{}%
\AgdaFunction{cong}\AgdaSpace{}%
\AgdaInductiveConstructor{box}\AgdaSpace{}%
\AgdaSymbol{(}\AgdaFunction{subst-with-id}\AgdaSpace{}%
\AgdaBound{M}\AgdaSymbol{)}\<%
\\
\>[0]\AgdaFunction{subst-with-id}\AgdaSpace{}%
\AgdaSymbol{(}\AgdaInductiveConstructor{ƛb}\AgdaSpace{}%
\AgdaBound{M}\AgdaSymbol{)}\AgdaSpace{}%
\AgdaSymbol{=}\<%
\\
\>[0][@{}l@{\AgdaIndent{0}}]%
\>[2]\AgdaOperator{\AgdaFunction{beginEq}}\<%
\\
\>[2][@{}l@{\AgdaIndent{0}}]%
\>[4]\AgdaInductiveConstructor{ƛb}\AgdaSpace{}%
\AgdaSymbol{(}\AgdaFunction{substb}\AgdaSpace{}%
\AgdaSymbol{(}\AgdaFunction{extsb}\AgdaSpace{}%
\AgdaInductiveConstructor{ε}\AgdaSymbol{)}\AgdaSpace{}%
\AgdaBound{M}\AgdaSymbol{)}\<%
\\
%
\>[2]\AgdaFunction{≡⟨}\AgdaSpace{}%
\AgdaFunction{cong}\AgdaSpace{}%
\AgdaInductiveConstructor{ƛb}\AgdaSpace{}%
\AgdaSymbol{(}\AgdaFunction{substb-cong}\AgdaSpace{}%
\AgdaSymbol{(λ}\AgdaSpace{}%
\AgdaBound{x}\AgdaSpace{}%
\AgdaSymbol{→}\AgdaSpace{}%
\AgdaFunction{ε-exts}\AgdaSpace{}%
\AgdaBound{x}\AgdaSymbol{)}\AgdaSpace{}%
\AgdaBound{M}\AgdaSymbol{)}\AgdaSpace{}%
\AgdaFunction{⟩}\<%
\\
\>[2][@{}l@{\AgdaIndent{0}}]%
\>[4]\AgdaInductiveConstructor{ƛb}\AgdaSpace{}%
\AgdaSymbol{(}\AgdaFunction{substb}\AgdaSpace{}%
\AgdaInductiveConstructor{ε}\AgdaSpace{}%
\AgdaBound{M}\AgdaSymbol{)}\<%
\\
%
\>[2]\AgdaFunction{≡⟨}\AgdaSpace{}%
\AgdaFunction{cong}\AgdaSpace{}%
\AgdaInductiveConstructor{ƛb}\AgdaSpace{}%
\AgdaSymbol{(}\AgdaFunction{subst-with-id}\AgdaSpace{}%
\AgdaBound{M}\AgdaSymbol{)}\AgdaSpace{}%
\AgdaFunction{⟩}\<%
\\
\>[2][@{}l@{\AgdaIndent{0}}]%
\>[4]\AgdaInductiveConstructor{ƛb}\AgdaSpace{}%
\AgdaBound{M}\<%
\\
%
\>[2]\AgdaOperator{\AgdaFunction{∎eq}}\<%
\\
\>[0]\AgdaFunction{subst-with-id}\AgdaSpace{}%
\AgdaSymbol{(}\AgdaBound{M}\AgdaSpace{}%
\AgdaOperator{\AgdaInductiveConstructor{·b}}\AgdaSpace{}%
\AgdaBound{M₁}\AgdaSymbol{)}\AgdaSpace{}%
\AgdaSymbol{=}\AgdaSpace{}%
\AgdaFunction{cong₂}\AgdaSpace{}%
\AgdaOperator{\AgdaInductiveConstructor{\AgdaUnderscore{}·b\AgdaUnderscore{}}}\AgdaSpace{}%
\AgdaSymbol{(}\AgdaFunction{subst-with-id}\AgdaSpace{}%
\AgdaBound{M}\AgdaSymbol{)}\AgdaSpace{}%
\AgdaSymbol{(}\AgdaFunction{subst-with-id}\AgdaSpace{}%
\AgdaBound{M₁}\AgdaSymbol{)}\<%
\\
%
\\[\AgdaEmptyExtraSkip]%
\>[0]\AgdaComment{--\ Substitution\ composed\ with\ renaming\ with\ Sb\ }\<%
\\
\>[0]\AgdaComment{--\ is\ the\ same\ as\ embedding\ that\ term.}\<%
\\
\>[0]\AgdaFunction{subst-rename-fb-comp}\AgdaSpace{}%
\AgdaSymbol{:}\AgdaSpace{}%
\AgdaSymbol{∀}\AgdaSpace{}%
\AgdaSymbol{\{}\AgdaBound{Γ}\AgdaSpace{}%
\AgdaBound{A}\AgdaSpace{}%
\AgdaBound{B}\AgdaSymbol{\}}\<%
\\
\>[0][@{}l@{\AgdaIndent{0}}]%
\>[2]\AgdaSymbol{→}\AgdaSpace{}%
\AgdaSymbol{(}\AgdaBound{M₁}\AgdaSpace{}%
\AgdaSymbol{:}\AgdaSpace{}%
\AgdaBound{Γ}\AgdaSpace{}%
\AgdaOperator{\AgdaInductiveConstructor{,}}\AgdaSpace{}%
\AgdaBound{B}\AgdaSpace{}%
\AgdaOperator{\AgdaDatatype{⊢}}\AgdaSpace{}%
\AgdaBound{A}\AgdaSymbol{)}\<%
\\
%
\>[2]\AgdaSymbol{→}\AgdaSpace{}%
\AgdaSymbol{(}\AgdaBound{M₂}\AgdaSpace{}%
\AgdaSymbol{:}\AgdaSpace{}%
\AgdaBound{Γ}\AgdaSpace{}%
\AgdaOperator{\AgdaDatatype{⊢}}\AgdaSpace{}%
\AgdaBound{B}\AgdaSymbol{)}\<%
\\
%
\>[2]\AgdaSymbol{→}\AgdaSpace{}%
\AgdaSymbol{(}\AgdaBound{x}\AgdaSpace{}%
\AgdaSymbol{:}\AgdaSpace{}%
\AgdaDatatype{Value}\AgdaSpace{}%
\AgdaBound{M₂}\AgdaSpace{}%
\AgdaSymbol{)}\<%
\\
%
\>[2]\AgdaSymbol{→}%
\>[2715I]\AgdaFunction{substb}\AgdaSpace{}%
\AgdaSymbol{(}\AgdaFunction{fb}\AgdaSpace{}%
\AgdaSymbol{(}\AgdaInductiveConstructor{ƛb}\AgdaSpace{}%
\AgdaSymbol{(}\AgdaFunction{embedTerm}\AgdaSpace{}%
\AgdaBound{M₁}\AgdaSymbol{)))}\AgdaSpace{}%
\AgdaSymbol{(}\AgdaFunction{renameb}\AgdaSpace{}%
\AgdaInductiveConstructor{Sb}\AgdaSpace{}%
\AgdaSymbol{(}\AgdaFunction{embedTerm}\AgdaSpace{}%
\AgdaBound{M₂}\AgdaSymbol{))}\<%
\\
\>[.][@{}l@{}]\<[2715I]%
\>[4]\AgdaOperator{\AgdaDatatype{≡}}\AgdaSpace{}%
\AgdaInductiveConstructor{box}\AgdaSpace{}%
\AgdaSymbol{(}\AgdaFunction{embedValue}\AgdaSpace{}%
\AgdaBound{M₂}\AgdaSpace{}%
\AgdaBound{x}\AgdaSymbol{)}\<%
\\
\>[0]\AgdaFunction{subst-rename-fb-comp}\AgdaSpace{}%
\AgdaBound{M₁}\AgdaSpace{}%
\AgdaBound{M₂}\AgdaSpace{}%
\AgdaBound{x}\AgdaSpace{}%
\AgdaSymbol{=}\<%
\\
\>[0][@{}l@{\AgdaIndent{0}}]%
\>[2]\AgdaOperator{\AgdaFunction{beginEq}}\<%
\\
\>[2][@{}l@{\AgdaIndent{0}}]%
\>[4]\AgdaFunction{substb}\AgdaSpace{}%
\AgdaSymbol{(}\AgdaFunction{fb}\AgdaSpace{}%
\AgdaSymbol{(}\AgdaInductiveConstructor{ƛb}\AgdaSpace{}%
\AgdaSymbol{(}\AgdaFunction{embedTerm}\AgdaSpace{}%
\AgdaBound{M₁}\AgdaSymbol{)))}\AgdaSpace{}%
\AgdaSymbol{(}\AgdaFunction{renameb}\AgdaSpace{}%
\AgdaInductiveConstructor{Sb}\AgdaSpace{}%
\AgdaSymbol{(}\AgdaFunction{embedTerm}\AgdaSpace{}%
\AgdaBound{M₂}\AgdaSymbol{))}\<%
\\
%
\>[2]\AgdaFunction{≡⟨}\AgdaSpace{}%
\AgdaFunction{subst-renameb-comp}\AgdaSpace{}%
\AgdaInductiveConstructor{Sb}\AgdaSpace{}%
\AgdaSymbol{(}\AgdaFunction{fb}\AgdaSpace{}%
\AgdaSymbol{(}\AgdaInductiveConstructor{ƛb}\AgdaSpace{}%
\AgdaSymbol{(}\AgdaFunction{embedTerm}\AgdaSpace{}%
\AgdaBound{M₁}\AgdaSymbol{)))}\AgdaSpace{}%
\AgdaSymbol{(}\AgdaFunction{embedTerm}\AgdaSpace{}%
\AgdaBound{M₂}\AgdaSymbol{)}\AgdaSpace{}%
\AgdaFunction{⟩}\<%
\\
\>[2][@{}l@{\AgdaIndent{0}}]%
\>[4]\AgdaFunction{substb}\AgdaSpace{}%
\AgdaSymbol{(λ}\AgdaSpace{}%
\AgdaBound{x₁}\AgdaSpace{}%
\AgdaSymbol{→}\AgdaSpace{}%
\AgdaInductiveConstructor{ε}\AgdaSpace{}%
\AgdaBound{x₁}\AgdaSymbol{)}\AgdaSpace{}%
\AgdaSymbol{(}\AgdaFunction{embedTerm}\AgdaSpace{}%
\AgdaBound{M₂}\AgdaSymbol{)}\<%
\\
%
\>[2]\AgdaFunction{≡⟨}\AgdaSpace{}%
\AgdaFunction{subst-with-id}\AgdaSpace{}%
\AgdaSymbol{(}\AgdaFunction{embedTerm}\AgdaSpace{}%
\AgdaBound{M₂}\AgdaSymbol{)}\AgdaSpace{}%
\AgdaFunction{⟩}\<%
\\
\>[2][@{}l@{\AgdaIndent{0}}]%
\>[4]\AgdaFunction{embedTerm}\AgdaSpace{}%
\AgdaBound{M₂}\<%
\\
%
\>[2]\AgdaFunction{≡⟨}\AgdaSpace{}%
\AgdaFunction{sym}\AgdaSpace{}%
\AgdaSymbol{(}\AgdaFunction{box-embedValue}\AgdaSpace{}%
\AgdaBound{x}\AgdaSymbol{)}\AgdaSpace{}%
\AgdaFunction{⟩}\<%
\\
\>[2][@{}l@{\AgdaIndent{0}}]%
\>[4]\AgdaInductiveConstructor{box}\AgdaSpace{}%
\AgdaSymbol{(}\AgdaFunction{embedValue}\AgdaSpace{}%
\AgdaBound{M₂}\AgdaSpace{}%
\AgdaBound{x}\AgdaSymbol{)}\<%
\\
%
\>[2]\AgdaOperator{\AgdaFunction{∎eq}}\<%
\\
%
\\[\AgdaEmptyExtraSkip]%
\>[0]\AgdaComment{--\ Preservation\ of\ 2\ step\ reduction\ of\ Gödel's\ translation}\<%
\\
\>[0]\AgdaFunction{pres-red}\AgdaSpace{}%
\AgdaSymbol{:}\AgdaSpace{}%
\AgdaSymbol{∀}\AgdaSpace{}%
\AgdaSymbol{\{}\AgdaBound{Γ}\AgdaSpace{}%
\AgdaBound{A}\AgdaSymbol{\}}\<%
\\
\>[0][@{}l@{\AgdaIndent{0}}]%
\>[2]\AgdaSymbol{→}\AgdaSpace{}%
\AgdaSymbol{(}\AgdaBound{M}\AgdaSpace{}%
\AgdaSymbol{:}\AgdaSpace{}%
\AgdaBound{Γ}\AgdaSpace{}%
\AgdaOperator{\AgdaDatatype{⊢}}\AgdaSpace{}%
\AgdaBound{A}\AgdaSymbol{)}\<%
\\
%
\>[2]\AgdaSymbol{→}\AgdaSpace{}%
\AgdaSymbol{(}\AgdaBound{N}\AgdaSpace{}%
\AgdaSymbol{:}\AgdaSpace{}%
\AgdaBound{Γ}\AgdaSpace{}%
\AgdaOperator{\AgdaDatatype{⊢}}\AgdaSpace{}%
\AgdaBound{A}\AgdaSymbol{)}\<%
\\
%
\>[2]\AgdaSymbol{→}\AgdaSpace{}%
\AgdaSymbol{(}\AgdaBound{M}\AgdaSpace{}%
\AgdaOperator{\AgdaDatatype{↝βᵥ}}\AgdaSpace{}%
\AgdaBound{N}\AgdaSymbol{)}\<%
\\
%
\>[2]\AgdaSymbol{→}\AgdaSpace{}%
\AgdaSymbol{(}\AgdaFunction{embedTerm}\AgdaSpace{}%
\AgdaBound{M}\AgdaSpace{}%
\AgdaOperator{\AgdaDatatype{↝βb²}}\AgdaSpace{}%
\AgdaFunction{embedTerm}\AgdaSpace{}%
\AgdaBound{N}\AgdaSymbol{)}\<%
\\
\>[0]\AgdaFunction{pres-red}\AgdaSpace{}%
\AgdaSymbol{(}\AgdaInductiveConstructor{ƛ}\AgdaSpace{}%
\AgdaBound{M₁}\AgdaSpace{}%
\AgdaOperator{\AgdaInductiveConstructor{·}}\AgdaSpace{}%
\AgdaBound{M₂}\AgdaSymbol{)}\AgdaSpace{}%
\AgdaBound{N}\AgdaSpace{}%
\AgdaSymbol{(}\AgdaInductiveConstructor{βᵥ}\AgdaSpace{}%
\AgdaBound{x}\AgdaSymbol{)}\AgdaSpace{}%
\AgdaSymbol{=}\<%
\\
\>[0][@{}l@{\AgdaIndent{0}}]%
\>[2]\AgdaInductiveConstructor{βb²}\<%
\\
%
\>[2]\AgdaSymbol{(}\AgdaFunction{raise}\AgdaSpace{}%
\AgdaSymbol{(}\AgdaFunction{embedTerm}\AgdaSpace{}%
\AgdaBound{M₂}\AgdaSymbol{)}\AgdaSpace{}%
\AgdaOperator{\AgdaInductiveConstructor{·b}}\AgdaSpace{}%
\AgdaInductiveConstructor{box}\AgdaSpace{}%
\AgdaSymbol{(}\AgdaInductiveConstructor{ƛb}\AgdaSpace{}%
\AgdaSymbol{(}\AgdaFunction{embedTerm}\AgdaSpace{}%
\AgdaBound{M₁}\AgdaSymbol{)))}\<%
\\
%
\>[2]\AgdaSymbol{(}\AgdaInductiveConstructor{ƛb}\AgdaSpace{}%
\AgdaSymbol{(}\AgdaFunction{embedTerm}\AgdaSpace{}%
\AgdaBound{M₁}\AgdaSymbol{)}\AgdaSpace{}%
\AgdaOperator{\AgdaInductiveConstructor{·b}}\AgdaSpace{}%
\AgdaInductiveConstructor{box}\AgdaSpace{}%
\AgdaSymbol{(}\AgdaFunction{embedValue}\AgdaSpace{}%
\AgdaBound{M₂}\AgdaSpace{}%
\AgdaBound{x}\AgdaSymbol{))}\<%
\\
%
\>[2]\AgdaSymbol{(}\AgdaFunction{embedTerm}\AgdaSpace{}%
\AgdaSymbol{(}\AgdaBound{M₁}\AgdaSpace{}%
\AgdaOperator{\AgdaFunction{[}}\AgdaSpace{}%
\AgdaBound{M₂}\AgdaSpace{}%
\AgdaOperator{\AgdaFunction{]}}\AgdaSymbol{))}\<%
\\
%
\>[2]\AgdaSymbol{(}\AgdaFunction{Eq.subst}\<%
\\
\>[2][@{}l@{\AgdaIndent{0}}]%
\>[4]\AgdaSymbol{(λ}\AgdaSpace{}%
\AgdaBound{x}\AgdaSpace{}%
\AgdaSymbol{→}\AgdaSpace{}%
\AgdaInductiveConstructor{ƛb}\AgdaSpace{}%
\AgdaSymbol{(}\AgdaInductiveConstructor{ε}\AgdaSpace{}%
\AgdaInductiveConstructor{Zb}\AgdaSpace{}%
\AgdaOperator{\AgdaInductiveConstructor{·b}}\AgdaSpace{}%
\AgdaFunction{renameb}\AgdaSpace{}%
\AgdaInductiveConstructor{Sb}\AgdaSpace{}%
\AgdaSymbol{(}\AgdaFunction{embedTerm}\AgdaSpace{}%
\AgdaBound{M₂}\AgdaSymbol{))}\AgdaSpace{}%
\AgdaOperator{\AgdaInductiveConstructor{·b}}\AgdaSpace{}%
\AgdaInductiveConstructor{box}\AgdaSpace{}%
\AgdaSymbol{(}\AgdaInductiveConstructor{ƛb}\AgdaSpace{}%
\AgdaSymbol{(}\AgdaFunction{embedTerm}\AgdaSpace{}%
\AgdaBound{M₁}\AgdaSymbol{))}\AgdaSpace{}%
\AgdaOperator{\AgdaDatatype{↝βb}}\AgdaSpace{}%
\AgdaBound{x}\AgdaSymbol{)}\<%
\\
%
\>[4]\AgdaSymbol{(}\AgdaFunction{cong₂}\AgdaSpace{}%
\AgdaOperator{\AgdaInductiveConstructor{\AgdaUnderscore{}·b\AgdaUnderscore{}}}\AgdaSpace{}%
\AgdaInductiveConstructor{refl}\AgdaSpace{}%
\AgdaSymbol{(}\AgdaFunction{subst-rename-fb-comp}\AgdaSpace{}%
\AgdaBound{M₁}\AgdaSpace{}%
\AgdaBound{M₂}\AgdaSpace{}%
\AgdaBound{x}\AgdaSymbol{))}\<%
\\
%
\>[4]\AgdaInductiveConstructor{βb}\AgdaSymbol{)}\<%
\\
%
\>[2]\AgdaSymbol{(}\AgdaFunction{Eq.subst}\AgdaSpace{}%
\AgdaSymbol{(λ}\AgdaSpace{}%
\AgdaBound{t}\AgdaSpace{}%
\AgdaSymbol{→}\AgdaSpace{}%
\AgdaInductiveConstructor{ƛb}\AgdaSpace{}%
\AgdaSymbol{(}\AgdaFunction{embedTerm}\AgdaSpace{}%
\AgdaBound{M₁}\AgdaSymbol{)}\AgdaSpace{}%
\AgdaOperator{\AgdaInductiveConstructor{·b}}\AgdaSpace{}%
\AgdaInductiveConstructor{box}\AgdaSpace{}%
\AgdaSymbol{(}\AgdaFunction{embedValue}\AgdaSpace{}%
\AgdaBound{M₂}\AgdaSpace{}%
\AgdaBound{x}\AgdaSymbol{)}\AgdaSpace{}%
\AgdaOperator{\AgdaDatatype{↝βb}}\AgdaSpace{}%
\AgdaBound{t}\AgdaSymbol{)}\AgdaSpace{}%
\AgdaSymbol{(}\AgdaFunction{pres-subst}\AgdaSpace{}%
\AgdaBound{M₁}\AgdaSpace{}%
\AgdaBound{M₂}\AgdaSpace{}%
\AgdaBound{x}\AgdaSymbol{)}\AgdaSpace{}%
\AgdaInductiveConstructor{βb}\AgdaSymbol{)}\<%
\\
\>[0]\AgdaFunction{pres-red}\AgdaSpace{}%
\AgdaBound{M}\AgdaSpace{}%
\AgdaBound{N}\AgdaSpace{}%
\AgdaSymbol{(}\AgdaInductiveConstructor{μ}\AgdaSpace{}%
\AgdaSymbol{\{}\AgdaArgument{M}\AgdaSpace{}%
\AgdaSymbol{=}\AgdaSpace{}%
\AgdaBound{M₁}\AgdaSymbol{\}}\AgdaSpace{}%
\AgdaSymbol{\{}\AgdaArgument{M'}\AgdaSpace{}%
\AgdaSymbol{=}\AgdaSpace{}%
\AgdaBound{M'}\AgdaSymbol{\}}\AgdaSpace{}%
\AgdaBound{red}\AgdaSymbol{)}\AgdaSpace{}%
\AgdaSymbol{=}\AgdaSpace{}%
\AgdaInductiveConstructor{ν}\AgdaSpace{}%
\AgdaSymbol{(}\AgdaFunction{pres-red}\AgdaSpace{}%
\AgdaBound{M₁}\AgdaSpace{}%
\AgdaBound{M'}\AgdaSpace{}%
\AgdaBound{red}\AgdaSymbol{)}\<%
\\
\>[0]\AgdaFunction{pres-red}\AgdaSpace{}%
\AgdaBound{M}\AgdaSpace{}%
\AgdaBound{N}\AgdaSpace{}%
\AgdaSymbol{(}\AgdaInductiveConstructor{ν<}\AgdaSpace{}%
\AgdaSymbol{\{}\AgdaArgument{M}\AgdaSpace{}%
\AgdaSymbol{=}\AgdaSpace{}%
\AgdaBound{M₁}\AgdaSymbol{\}}\AgdaSpace{}%
\AgdaSymbol{\{}\AgdaArgument{N}\AgdaSpace{}%
\AgdaSymbol{=}\AgdaSpace{}%
\AgdaBound{N₁}\AgdaSymbol{\}}\AgdaSpace{}%
\AgdaSymbol{\{}\AgdaArgument{N'}\AgdaSpace{}%
\AgdaSymbol{=}\AgdaSpace{}%
\AgdaBound{N'}\AgdaSymbol{\}}\AgdaSpace{}%
\AgdaBound{x}\AgdaSpace{}%
\AgdaBound{red}\AgdaSymbol{)}\AgdaSpace{}%
\AgdaSymbol{=}\AgdaSpace{}%
\AgdaInductiveConstructor{μ}\AgdaSpace{}%
\AgdaSymbol{(}\AgdaInductiveConstructor{ξ}\AgdaSpace{}%
\AgdaSymbol{(}\AgdaInductiveConstructor{ν}\AgdaSpace{}%
\AgdaSymbol{(}\AgdaFunction{pres-rename-red-2}\AgdaSpace{}%
\AgdaInductiveConstructor{Sb}\AgdaSpace{}%
\AgdaSymbol{(}\AgdaFunction{embedTerm}\AgdaSpace{}%
\AgdaBound{N₁}\AgdaSymbol{)}\AgdaSpace{}%
\AgdaSymbol{(}\AgdaFunction{embedTerm}\AgdaSpace{}%
\AgdaBound{N'}\AgdaSymbol{)}\AgdaSpace{}%
\AgdaSymbol{(}\AgdaFunction{pres-red}\AgdaSpace{}%
\AgdaBound{N₁}\AgdaSpace{}%
\AgdaBound{N'}\AgdaSpace{}%
\AgdaBound{red}\AgdaSymbol{))))}\<%
\\
\>[0]\AgdaFunction{pres-red}\AgdaSpace{}%
\AgdaBound{M}\AgdaSpace{}%
\AgdaBound{N}\AgdaSpace{}%
\AgdaSymbol{(}\AgdaInductiveConstructor{ν}\AgdaSpace{}%
\AgdaSymbol{\{}\AgdaArgument{M}\AgdaSpace{}%
\AgdaSymbol{=}\AgdaSpace{}%
\AgdaBound{M₁}\AgdaSymbol{\}}\AgdaSpace{}%
\AgdaSymbol{\{}\AgdaArgument{N}\AgdaSpace{}%
\AgdaSymbol{=}\AgdaSpace{}%
\AgdaBound{N₁}\AgdaSymbol{\}}\AgdaSpace{}%
\AgdaSymbol{\{}\AgdaArgument{N'}\AgdaSpace{}%
\AgdaSymbol{=}\AgdaSpace{}%
\AgdaBound{N'}\AgdaSymbol{\}}\AgdaSpace{}%
\AgdaBound{red}\AgdaSymbol{)}\AgdaSpace{}%
\AgdaSymbol{=}\AgdaSpace{}%
\AgdaInductiveConstructor{μ}\AgdaSpace{}%
\AgdaSymbol{(}\AgdaInductiveConstructor{ξ}\AgdaSpace{}%
\AgdaSymbol{(}\AgdaInductiveConstructor{ν}\AgdaSpace{}%
\AgdaSymbol{(}\AgdaFunction{pres-rename-red-2}\AgdaSpace{}%
\AgdaInductiveConstructor{Sb}\AgdaSpace{}%
\AgdaSymbol{(}\AgdaFunction{embedTerm}\AgdaSpace{}%
\AgdaBound{N₁}\AgdaSymbol{)}\AgdaSpace{}%
\AgdaSymbol{(}\AgdaFunction{embedTerm}\AgdaSpace{}%
\AgdaBound{N'}\AgdaSymbol{)}\AgdaSpace{}%
\AgdaSymbol{(}\AgdaFunction{pres-red}\AgdaSpace{}%
\AgdaBound{N₁}\AgdaSpace{}%
\AgdaBound{N'}\AgdaSpace{}%
\AgdaBound{red}\AgdaSymbol{))))}\<%
\\
\>[0]\AgdaFunction{pres-red}\AgdaSpace{}%
\AgdaBound{M}\AgdaSpace{}%
\AgdaBound{N}\AgdaSpace{}%
\AgdaSymbol{(}\AgdaInductiveConstructor{ξ}\AgdaSpace{}%
\AgdaSymbol{\{}\AgdaArgument{M}\AgdaSpace{}%
\AgdaSymbol{=}\AgdaSpace{}%
\AgdaBound{M₁}\AgdaSymbol{\}}\AgdaSpace{}%
\AgdaSymbol{\{}\AgdaArgument{M'}\AgdaSpace{}%
\AgdaSymbol{=}\AgdaSpace{}%
\AgdaBound{M'}\AgdaSymbol{\}}\AgdaSpace{}%
\AgdaBound{red}\AgdaSymbol{)}\AgdaSpace{}%
\AgdaSymbol{=}\AgdaSpace{}%
\AgdaInductiveConstructor{ζ}\AgdaSpace{}%
\AgdaSymbol{(}\AgdaInductiveConstructor{ξ}\AgdaSpace{}%
\AgdaSymbol{(}\AgdaFunction{pres-red}\AgdaSpace{}%
\AgdaBound{M₁}\AgdaSpace{}%
\AgdaBound{M'}\AgdaSpace{}%
\AgdaBound{red}\AgdaSymbol{))}\<%
\\
%
\\[\AgdaEmptyExtraSkip]%
\>[0]\AgdaComment{--\ Preservation\ of\ the\ 2\ step\ reduction\ rule\ under\ rename}\<%
\\
\>[0]\AgdaFunction{pres-rename-red-↝βb₂}\AgdaSpace{}%
\AgdaSymbol{:}\AgdaSpace{}%
\AgdaSymbol{∀}\AgdaSpace{}%
\AgdaSymbol{\{}\AgdaBound{Γ}\AgdaSpace{}%
\AgdaBound{A}\AgdaSpace{}%
\AgdaBound{Δ}\AgdaSymbol{\}}\<%
\\
\>[0][@{}l@{\AgdaIndent{0}}]%
\>[2]\AgdaSymbol{→}\AgdaSpace{}%
\AgdaSymbol{(}\AgdaBound{fb'}\AgdaSpace{}%
\AgdaSymbol{:}\AgdaSpace{}%
\AgdaSymbol{∀}\AgdaSpace{}%
\AgdaSymbol{\{}\AgdaBound{Ab}\AgdaSymbol{\}}\AgdaSpace{}%
\AgdaSymbol{→}\AgdaSpace{}%
\AgdaBound{Ab}\AgdaSpace{}%
\AgdaOperator{\AgdaDatatype{∈b}}\AgdaSpace{}%
\AgdaBound{Γ}\AgdaSpace{}%
\AgdaSymbol{→}\AgdaSpace{}%
\AgdaBound{Ab}\AgdaSpace{}%
\AgdaOperator{\AgdaDatatype{∈b}}\AgdaSpace{}%
\AgdaBound{Δ}\AgdaSymbol{)}\<%
\\
%
\>[2]\AgdaSymbol{→}\AgdaSpace{}%
\AgdaSymbol{(}\AgdaBound{M}\AgdaSpace{}%
\AgdaSymbol{:}\AgdaSpace{}%
\AgdaBound{Γ}\AgdaSpace{}%
\AgdaOperator{\AgdaDatatype{⊢b}}\AgdaSpace{}%
\AgdaBound{A}\AgdaSymbol{)}\<%
\\
%
\>[2]\AgdaSymbol{→}\AgdaSpace{}%
\AgdaSymbol{(}\AgdaBound{N}\AgdaSpace{}%
\AgdaSymbol{:}\AgdaSpace{}%
\AgdaBound{Γ}\AgdaSpace{}%
\AgdaOperator{\AgdaDatatype{⊢b}}\AgdaSpace{}%
\AgdaBound{A}\AgdaSymbol{)}\<%
\\
%
\>[2]\AgdaSymbol{→}\AgdaSpace{}%
\AgdaSymbol{(}\AgdaBound{M}\AgdaSpace{}%
\AgdaOperator{\AgdaDatatype{↝βb₂}}\AgdaSpace{}%
\AgdaBound{N}\AgdaSymbol{)}\<%
\\
%
\>[2]\AgdaSymbol{→}\AgdaSpace{}%
\AgdaFunction{renameb}\AgdaSpace{}%
\AgdaBound{fb'}\AgdaSpace{}%
\AgdaBound{M}\AgdaSpace{}%
\AgdaOperator{\AgdaDatatype{↝βb₂}}\AgdaSpace{}%
\AgdaFunction{renameb}\AgdaSpace{}%
\AgdaBound{fb'}\AgdaSpace{}%
\AgdaBound{N}\<%
\\
\>[0]\AgdaFunction{pres-rename-red-↝βb₂}\AgdaSpace{}%
\AgdaBound{fb'}\AgdaSpace{}%
\AgdaBound{M}\AgdaSpace{}%
\AgdaBound{N}\AgdaSpace{}%
\AgdaSymbol{(}\AgdaInductiveConstructor{βb₂}\AgdaSpace{}%
\AgdaSymbol{\{}\AgdaArgument{M}\AgdaSpace{}%
\AgdaSymbol{=}\AgdaSpace{}%
\AgdaBound{M'}\AgdaSymbol{\}}\AgdaSpace{}%
\AgdaSymbol{\{}\AgdaArgument{W}\AgdaSpace{}%
\AgdaSymbol{=}\AgdaSpace{}%
\AgdaBound{W'}\AgdaSymbol{\})}\AgdaSpace{}%
\AgdaSymbol{=}\<%
\\
\>[0][@{}l@{\AgdaIndent{0}}]%
\>[2]\AgdaFunction{Eq.subst}\<%
\\
\>[2][@{}l@{\AgdaIndent{0}}]%
\>[4]\AgdaSymbol{(λ}\AgdaSpace{}%
\AgdaBound{x}\AgdaSpace{}%
\AgdaSymbol{→}\AgdaSpace{}%
\AgdaBound{x}\AgdaSpace{}%
\AgdaOperator{\AgdaDatatype{↝βb₂}}\AgdaSpace{}%
\AgdaFunction{renameb}\AgdaSpace{}%
\AgdaBound{fb'}\<%
\\
%
\>[4]\AgdaSymbol{(}\AgdaBound{M'}\AgdaSpace{}%
\AgdaOperator{\AgdaFunction{[}}\AgdaSpace{}%
\AgdaBound{W'}\AgdaSpace{}%
\AgdaOperator{\AgdaFunction{]b}}\AgdaSymbol{))}\<%
\\
%
\>[4]\AgdaSymbol{(}\AgdaFunction{cong₂}\AgdaSpace{}%
\AgdaOperator{\AgdaInductiveConstructor{\AgdaUnderscore{}·b\AgdaUnderscore{}}}\AgdaSpace{}%
\AgdaSymbol{(}\AgdaFunction{cong}\AgdaSpace{}%
\AgdaInductiveConstructor{ƛb}\AgdaSpace{}%
\AgdaSymbol{(}\AgdaFunction{cong₂}\AgdaSpace{}%
\AgdaOperator{\AgdaInductiveConstructor{\AgdaUnderscore{}·b\AgdaUnderscore{}}}\AgdaSpace{}%
\AgdaInductiveConstructor{refl}\AgdaSpace{}%
\AgdaSymbol{(}\AgdaFunction{cong}\AgdaSpace{}%
\AgdaInductiveConstructor{box}\AgdaSpace{}%
\AgdaSymbol{(}\AgdaFunction{sym}\AgdaSpace{}%
\AgdaSymbol{(}\AgdaFunction{rename-extb-Sb-commute}\AgdaSpace{}%
\AgdaBound{fb'}\AgdaSpace{}%
\AgdaBound{W'}\AgdaSymbol{)))))}\AgdaSpace{}%
\AgdaInductiveConstructor{refl}\AgdaSymbol{)}\<%
\\
%
\>[4]\AgdaSymbol{(}\AgdaFunction{Eq.subst}\<%
\\
\>[4][@{}l@{\AgdaIndent{0}}]%
\>[6]\AgdaSymbol{(λ}\AgdaSpace{}%
\AgdaBound{x}\AgdaSpace{}%
\AgdaSymbol{→}\AgdaSpace{}%
\AgdaSymbol{(}\AgdaFunction{raise}%
\>[21]\AgdaSymbol{(}\AgdaInductiveConstructor{box}\AgdaSpace{}%
\AgdaSymbol{(}\AgdaFunction{renameb}\AgdaSpace{}%
\AgdaBound{fb'}\AgdaSpace{}%
\AgdaBound{W'}\AgdaSymbol{)))}\AgdaSpace{}%
\AgdaOperator{\AgdaInductiveConstructor{·b}}\AgdaSpace{}%
\AgdaInductiveConstructor{box}\AgdaSpace{}%
\AgdaSymbol{(}\AgdaInductiveConstructor{ƛb}\AgdaSpace{}%
\AgdaSymbol{(}\AgdaFunction{renameb}\AgdaSpace{}%
\AgdaSymbol{(}\AgdaFunction{extb}\AgdaSpace{}%
\AgdaBound{fb'}\AgdaSymbol{)}\AgdaSpace{}%
\AgdaBound{M'}\AgdaSymbol{))}\AgdaSpace{}%
\AgdaOperator{\AgdaDatatype{↝βb₂}}\AgdaSpace{}%
\AgdaBound{x}\AgdaSymbol{)}\<%
\\
%
\>[6]\AgdaSymbol{(}\AgdaFunction{pres-subst-func-rename}\AgdaSpace{}%
\AgdaBound{fb'}\AgdaSpace{}%
\AgdaBound{M'}\AgdaSpace{}%
\AgdaBound{W'}\AgdaSymbol{)}\<%
\\
%
\>[6]\AgdaInductiveConstructor{βb₂}\<%
\\
%
\>[4]\AgdaSymbol{)}\<%
\\
\>[0]\AgdaFunction{pres-rename-red-↝βb₂}\AgdaSpace{}%
\AgdaBound{fb'}\AgdaSpace{}%
\AgdaBound{M}\AgdaSpace{}%
\AgdaBound{N}\AgdaSpace{}%
\AgdaSymbol{(}\AgdaInductiveConstructor{μ}\AgdaSpace{}%
\AgdaSymbol{\{}\AgdaArgument{M}\AgdaSpace{}%
\AgdaSymbol{=}\AgdaSpace{}%
\AgdaBound{M₁}\AgdaSymbol{\}}\AgdaSpace{}%
\AgdaSymbol{\{}\AgdaArgument{M'}\AgdaSpace{}%
\AgdaSymbol{=}\AgdaSpace{}%
\AgdaBound{M'}\AgdaSymbol{\}}\AgdaSpace{}%
\AgdaBound{red}\AgdaSymbol{)}\AgdaSpace{}%
\AgdaSymbol{=}\AgdaSpace{}%
\AgdaInductiveConstructor{μ}\AgdaSpace{}%
\AgdaSymbol{(}\AgdaFunction{pres-rename-red-↝βb₂}\AgdaSpace{}%
\AgdaBound{fb'}\AgdaSpace{}%
\AgdaBound{M₁}\AgdaSpace{}%
\AgdaBound{M'}\AgdaSpace{}%
\AgdaBound{red}\AgdaSymbol{)}\<%
\\
\>[0]\AgdaFunction{pres-rename-red-↝βb₂}\AgdaSpace{}%
\AgdaBound{fb'}\AgdaSpace{}%
\AgdaBound{M}\AgdaSpace{}%
\AgdaBound{N}\AgdaSpace{}%
\AgdaSymbol{(}\AgdaInductiveConstructor{ν}\AgdaSpace{}%
\AgdaSymbol{\{}\AgdaArgument{N}\AgdaSpace{}%
\AgdaSymbol{=}\AgdaSpace{}%
\AgdaBound{N₁}\AgdaSymbol{\}}\AgdaSpace{}%
\AgdaSymbol{\{}\AgdaArgument{N'}\AgdaSpace{}%
\AgdaSymbol{=}\AgdaSpace{}%
\AgdaBound{N'}\AgdaSymbol{\}}\AgdaSpace{}%
\AgdaBound{red}\AgdaSymbol{)}\AgdaSpace{}%
\AgdaSymbol{=}\AgdaSpace{}%
\AgdaInductiveConstructor{ν}\AgdaSpace{}%
\AgdaSymbol{(}\AgdaFunction{pres-rename-red-↝βb₂}\AgdaSpace{}%
\AgdaBound{fb'}\AgdaSpace{}%
\AgdaBound{N₁}\AgdaSpace{}%
\AgdaBound{N'}\AgdaSpace{}%
\AgdaBound{red}\AgdaSymbol{)}\<%
\\
\>[0]\AgdaFunction{pres-rename-red-↝βb₂}\AgdaSpace{}%
\AgdaBound{fb'}\AgdaSpace{}%
\AgdaBound{M}\AgdaSpace{}%
\AgdaBound{N}\AgdaSpace{}%
\AgdaSymbol{(}\AgdaInductiveConstructor{ξ}\AgdaSpace{}%
\AgdaSymbol{\{}\AgdaArgument{M}\AgdaSpace{}%
\AgdaSymbol{=}\AgdaSpace{}%
\AgdaBound{M₁}\AgdaSymbol{\}}\AgdaSpace{}%
\AgdaSymbol{\{}\AgdaArgument{M'}\AgdaSpace{}%
\AgdaSymbol{=}\AgdaSpace{}%
\AgdaBound{M'}\AgdaSymbol{\}}\AgdaSpace{}%
\AgdaBound{red}\AgdaSymbol{)}\AgdaSpace{}%
\AgdaSymbol{=}\AgdaSpace{}%
\AgdaInductiveConstructor{ξ}\AgdaSpace{}%
\AgdaSymbol{(}\AgdaFunction{pres-rename-red-↝βb₂}\AgdaSpace{}%
\AgdaSymbol{(}\AgdaFunction{extb}\AgdaSpace{}%
\AgdaBound{fb'}\AgdaSymbol{)}\AgdaSpace{}%
\AgdaBound{M₁}\AgdaSpace{}%
\AgdaBound{M'}\AgdaSpace{}%
\AgdaBound{red}\AgdaSymbol{)}\<%
\\
\>[0]\AgdaFunction{pres-rename-red-↝βb₂}\AgdaSpace{}%
\AgdaBound{fb'}\AgdaSpace{}%
\AgdaBound{M}\AgdaSpace{}%
\AgdaBound{N}\AgdaSpace{}%
\AgdaSymbol{(}\AgdaInductiveConstructor{ζ}\AgdaSpace{}%
\AgdaSymbol{\{}\AgdaArgument{M}\AgdaSpace{}%
\AgdaSymbol{=}\AgdaSpace{}%
\AgdaBound{M₁}\AgdaSymbol{\}}\AgdaSpace{}%
\AgdaSymbol{\{}\AgdaArgument{M'}\AgdaSpace{}%
\AgdaSymbol{=}\AgdaSpace{}%
\AgdaBound{M'}\AgdaSymbol{\}}\AgdaSpace{}%
\AgdaBound{red}\AgdaSymbol{)}\AgdaSpace{}%
\AgdaSymbol{=}\AgdaSpace{}%
\AgdaInductiveConstructor{ζ}\AgdaSpace{}%
\AgdaSymbol{(}\AgdaFunction{pres-rename-red-↝βb₂}\AgdaSpace{}%
\AgdaBound{fb'}\AgdaSpace{}%
\AgdaBound{M₁}\AgdaSpace{}%
\AgdaBound{M'}\AgdaSpace{}%
\AgdaBound{red}\AgdaSymbol{)}\<%
\\
\>[0]\<%
\end{code}
\begin{code}%
\>[0]\AgdaComment{--\ Preservation\ of\ reduction\ by\ Gödel's\ translation}\<%
\\
\>[0]\AgdaComment{--\ with\ the\ two\ step\ reduction\ rule}\<%
\\
\>[0]\AgdaFunction{pres-red↝βb₂}\AgdaSpace{}%
\AgdaSymbol{:}\AgdaSpace{}%
\AgdaSymbol{∀}\AgdaSpace{}%
\AgdaSymbol{\{}\AgdaBound{Γ}\AgdaSpace{}%
\AgdaBound{A}\AgdaSymbol{\}}\<%
\\
\>[0][@{}l@{\AgdaIndent{0}}]%
\>[2]\AgdaSymbol{→}\AgdaSpace{}%
\AgdaSymbol{(}\AgdaBound{M}\AgdaSpace{}%
\AgdaSymbol{:}\AgdaSpace{}%
\AgdaBound{Γ}\AgdaSpace{}%
\AgdaOperator{\AgdaDatatype{⊢}}\AgdaSpace{}%
\AgdaBound{A}\AgdaSymbol{)}\<%
\\
%
\>[2]\AgdaSymbol{→}\AgdaSpace{}%
\AgdaSymbol{(}\AgdaBound{N}\AgdaSpace{}%
\AgdaSymbol{:}\AgdaSpace{}%
\AgdaBound{Γ}\AgdaSpace{}%
\AgdaOperator{\AgdaDatatype{⊢}}\AgdaSpace{}%
\AgdaBound{A}\AgdaSymbol{)}\<%
\\
%
\>[2]\AgdaSymbol{→}\AgdaSpace{}%
\AgdaSymbol{(}\AgdaBound{M}\AgdaSpace{}%
\AgdaOperator{\AgdaDatatype{↝βᵥ}}\AgdaSpace{}%
\AgdaBound{N}\AgdaSymbol{)}\<%
\\
%
\>[2]\AgdaSymbol{→}\AgdaSpace{}%
\AgdaSymbol{(}\AgdaFunction{embedTerm}\AgdaSpace{}%
\AgdaBound{M}\AgdaSpace{}%
\AgdaOperator{\AgdaDatatype{↝βb₂}}\AgdaSpace{}%
\AgdaFunction{embedTerm}\AgdaSpace{}%
\AgdaBound{N}\AgdaSymbol{)}\<%
\end{code}
\begin{code}[hide]%
\>[0]\AgdaFunction{pres-red↝βb₂}\AgdaSpace{}%
\AgdaSymbol{(}\AgdaInductiveConstructor{ƛ}\AgdaSpace{}%
\AgdaBound{M₁}\AgdaSpace{}%
\AgdaOperator{\AgdaInductiveConstructor{·}}\AgdaSpace{}%
\AgdaBound{M₂}\AgdaSymbol{)}\AgdaSpace{}%
\AgdaBound{N}\AgdaSpace{}%
\AgdaSymbol{(}\AgdaInductiveConstructor{βᵥ}\AgdaSpace{}%
\AgdaBound{x}\AgdaSymbol{)}\AgdaSpace{}%
\AgdaSymbol{=}\<%
\\
\>[0][@{}l@{\AgdaIndent{0}}]%
\>[2]\AgdaFunction{Eq.subst}\<%
\\
\>[2][@{}l@{\AgdaIndent{0}}]%
\>[4]\AgdaSymbol{(λ}\AgdaSpace{}%
\AgdaBound{t}\AgdaSpace{}%
\AgdaSymbol{→}\AgdaSpace{}%
\AgdaBound{t}\AgdaSpace{}%
\AgdaOperator{\AgdaDatatype{↝βb₂}}\AgdaSpace{}%
\AgdaFunction{embedTerm}\AgdaSpace{}%
\AgdaSymbol{(}\AgdaBound{M₁}\AgdaSpace{}%
\AgdaOperator{\AgdaFunction{[}}\AgdaSpace{}%
\AgdaBound{M₂}\AgdaSpace{}%
\AgdaOperator{\AgdaFunction{]}}\AgdaSymbol{))}\<%
\\
%
\>[4]\AgdaSymbol{(}\AgdaFunction{cong₂}\<%
\\
\>[4][@{}l@{\AgdaIndent{0}}]%
\>[6]\AgdaOperator{\AgdaInductiveConstructor{\AgdaUnderscore{}·b\AgdaUnderscore{}}}\<%
\\
%
\>[6]\AgdaSymbol{(}\AgdaFunction{cong}\<%
\\
\>[6][@{}l@{\AgdaIndent{0}}]%
\>[8]\AgdaInductiveConstructor{ƛb}\<%
\\
%
\>[8]\AgdaSymbol{(}\AgdaFunction{cong₂}\AgdaSpace{}%
\AgdaOperator{\AgdaInductiveConstructor{\AgdaUnderscore{}·b\AgdaUnderscore{}}}\AgdaSpace{}%
\AgdaInductiveConstructor{refl}\AgdaSpace{}%
\AgdaSymbol{(}\AgdaFunction{cong}\AgdaSpace{}%
\AgdaSymbol{(}\AgdaFunction{renameb}\AgdaSpace{}%
\AgdaInductiveConstructor{Sb}\AgdaSymbol{)}\AgdaSpace{}%
\AgdaSymbol{(}\AgdaFunction{box-embedValue}\AgdaSpace{}%
\AgdaBound{x}\AgdaSymbol{))))}\<%
\\
%
\>[6]\AgdaInductiveConstructor{refl}\AgdaSymbol{)}\<%
\\
%
\>[4]\AgdaSymbol{(}\AgdaFunction{Eq.subst}\<%
\\
\>[4][@{}l@{\AgdaIndent{0}}]%
\>[6]\AgdaSymbol{(λ}\AgdaSpace{}%
\AgdaBound{t}\AgdaSpace{}%
\AgdaSymbol{→}\<%
\\
\>[6][@{}l@{\AgdaIndent{0}}]%
\>[8]\AgdaFunction{raise}\AgdaSpace{}%
\AgdaSymbol{(}\AgdaSpace{}%
\AgdaInductiveConstructor{box}\AgdaSpace{}%
\AgdaSymbol{(}\AgdaFunction{embedValue}\AgdaSpace{}%
\AgdaBound{M₂}\AgdaSpace{}%
\AgdaBound{x}\AgdaSymbol{)}\AgdaSpace{}%
\AgdaSymbol{)}\AgdaSpace{}%
\AgdaOperator{\AgdaInductiveConstructor{·b}}\AgdaSpace{}%
\AgdaInductiveConstructor{box}\AgdaSpace{}%
\AgdaSymbol{(}\AgdaInductiveConstructor{ƛb}\AgdaSpace{}%
\AgdaSymbol{(}\AgdaFunction{embedTerm}\AgdaSpace{}%
\AgdaBound{M₁}\AgdaSymbol{))}\AgdaSpace{}%
\AgdaOperator{\AgdaDatatype{↝βb₂}}\AgdaSpace{}%
\AgdaBound{t}\AgdaSymbol{)}\<%
\\
%
\>[6]\AgdaSymbol{(}\AgdaFunction{pres-subst}\AgdaSpace{}%
\AgdaBound{M₁}\AgdaSpace{}%
\AgdaBound{M₂}\AgdaSpace{}%
\AgdaBound{x}\AgdaSymbol{)}\<%
\\
%
\>[6]\AgdaInductiveConstructor{βb₂}\<%
\\
%
\>[4]\AgdaSymbol{)}\<%
\\
\>[0]\AgdaFunction{pres-red↝βb₂}\AgdaSpace{}%
\AgdaBound{M}\AgdaSpace{}%
\AgdaBound{N}\AgdaSpace{}%
\AgdaSymbol{(}\AgdaInductiveConstructor{μ}\AgdaSpace{}%
\AgdaSymbol{\{}\AgdaArgument{M}\AgdaSpace{}%
\AgdaSymbol{=}\AgdaSpace{}%
\AgdaBound{M₁}\AgdaSymbol{\}}\AgdaSpace{}%
\AgdaSymbol{\{}\AgdaArgument{M'}\AgdaSpace{}%
\AgdaSymbol{=}\AgdaSpace{}%
\AgdaBound{M'}\AgdaSymbol{\}}\AgdaSpace{}%
\AgdaBound{red}\AgdaSymbol{)}\AgdaSpace{}%
\AgdaSymbol{=}\<%
\\
\>[0][@{}l@{\AgdaIndent{0}}]%
\>[2]\AgdaInductiveConstructor{ν}\AgdaSpace{}%
\AgdaSymbol{(}\AgdaFunction{pres-red↝βb₂}\AgdaSpace{}%
\AgdaBound{M₁}\AgdaSpace{}%
\AgdaBound{M'}\AgdaSpace{}%
\AgdaBound{red}\AgdaSymbol{)}\<%
\\
\>[0]\AgdaFunction{pres-red↝βb₂}\AgdaSpace{}%
\AgdaBound{M}\AgdaSpace{}%
\AgdaBound{N}\AgdaSpace{}%
\AgdaSymbol{(}\AgdaInductiveConstructor{ν<}\AgdaSpace{}%
\AgdaSymbol{\{}\AgdaArgument{N}\AgdaSpace{}%
\AgdaSymbol{=}\AgdaSpace{}%
\AgdaBound{N₁}\AgdaSymbol{\}}\AgdaSpace{}%
\AgdaSymbol{\{}\AgdaArgument{N'}\AgdaSpace{}%
\AgdaSymbol{=}\AgdaSpace{}%
\AgdaBound{N'}\AgdaSymbol{\}}\AgdaSpace{}%
\AgdaBound{x}\AgdaSpace{}%
\AgdaBound{red}\AgdaSymbol{)}\AgdaSpace{}%
\AgdaSymbol{=}\<%
\\
\>[0][@{}l@{\AgdaIndent{0}}]%
\>[2]\AgdaInductiveConstructor{μ}\AgdaSpace{}%
\AgdaSymbol{(}\AgdaInductiveConstructor{ξ}\AgdaSpace{}%
\AgdaSymbol{(}\AgdaInductiveConstructor{ν}%
\>[3199I]\AgdaSymbol{(}\AgdaFunction{pres-rename-red-↝βb₂}\<%
\\
\>[3199I][@{}l@{\AgdaIndent{0}}]%
\>[12]\AgdaInductiveConstructor{Sb}\<%
\\
%
\>[12]\AgdaSymbol{(}\AgdaFunction{embedTerm}\AgdaSpace{}%
\AgdaBound{N₁}\AgdaSymbol{)}\<%
\\
%
\>[12]\AgdaSymbol{(}\AgdaFunction{embedTerm}\AgdaSpace{}%
\AgdaBound{N'}\AgdaSymbol{)}\<%
\\
%
\>[12]\AgdaSymbol{(}\AgdaFunction{pres-red↝βb₂}\AgdaSpace{}%
\AgdaBound{N₁}\AgdaSpace{}%
\AgdaBound{N'}\AgdaSpace{}%
\AgdaBound{red}\AgdaSymbol{))))}\<%
\\
\>[0]\AgdaFunction{pres-red↝βb₂}\AgdaSpace{}%
\AgdaBound{M}\AgdaSpace{}%
\AgdaBound{N}\AgdaSpace{}%
\AgdaSymbol{(}\AgdaInductiveConstructor{ν}\AgdaSpace{}%
\AgdaSymbol{\{}\AgdaArgument{N}\AgdaSpace{}%
\AgdaSymbol{=}\AgdaSpace{}%
\AgdaBound{N₁}\AgdaSymbol{\}}\AgdaSpace{}%
\AgdaSymbol{\{}\AgdaArgument{N'}\AgdaSpace{}%
\AgdaSymbol{=}\AgdaSpace{}%
\AgdaBound{N'}\AgdaSymbol{\}}\AgdaSpace{}%
\AgdaBound{red}\AgdaSymbol{)}\AgdaSpace{}%
\AgdaSymbol{=}\<%
\\
\>[0][@{}l@{\AgdaIndent{0}}]%
\>[2]\AgdaInductiveConstructor{μ}\AgdaSpace{}%
\AgdaSymbol{(}\AgdaInductiveConstructor{ξ}\AgdaSpace{}%
\AgdaSymbol{(}\AgdaInductiveConstructor{ν}%
\>[3218I]\AgdaSymbol{(}\AgdaFunction{pres-rename-red-↝βb₂}\<%
\\
\>[3218I][@{}l@{\AgdaIndent{0}}]%
\>[12]\AgdaInductiveConstructor{Sb}\<%
\\
%
\>[12]\AgdaSymbol{(}\AgdaFunction{embedTerm}\AgdaSpace{}%
\AgdaBound{N₁}\AgdaSymbol{)}\<%
\\
%
\>[12]\AgdaSymbol{(}\AgdaFunction{embedTerm}\AgdaSpace{}%
\AgdaBound{N'}\AgdaSymbol{)}\<%
\\
%
\>[12]\AgdaSymbol{(}\AgdaFunction{pres-red↝βb₂}\AgdaSpace{}%
\AgdaBound{N₁}\AgdaSpace{}%
\AgdaBound{N'}\AgdaSpace{}%
\AgdaBound{red}\AgdaSymbol{))))}\<%
\\
\>[0]\AgdaFunction{pres-red↝βb₂}\AgdaSpace{}%
\AgdaBound{M}\AgdaSpace{}%
\AgdaBound{N}\AgdaSpace{}%
\AgdaSymbol{(}\AgdaInductiveConstructor{ξ}\AgdaSpace{}%
\AgdaSymbol{\{}\AgdaArgument{M}\AgdaSpace{}%
\AgdaSymbol{=}\AgdaSpace{}%
\AgdaBound{M₁}\AgdaSymbol{\}}\AgdaSpace{}%
\AgdaSymbol{\{}\AgdaArgument{M'}\AgdaSpace{}%
\AgdaSymbol{=}\AgdaSpace{}%
\AgdaBound{M'}\AgdaSymbol{\}}\AgdaSpace{}%
\AgdaBound{red}\AgdaSymbol{)}\AgdaSpace{}%
\AgdaSymbol{=}\<%
\\
\>[0][@{}l@{\AgdaIndent{0}}]%
\>[2]\AgdaInductiveConstructor{ζ}\AgdaSpace{}%
\AgdaSymbol{(}\AgdaInductiveConstructor{ξ}\AgdaSpace{}%
\AgdaSymbol{(}\AgdaFunction{pres-red↝βb₂}\AgdaSpace{}%
\AgdaBound{M₁}\AgdaSpace{}%
\AgdaBound{M'}\AgdaSpace{}%
\AgdaBound{red}\AgdaSymbol{))}\<%
\end{code}
  \hfill $\qed$
\end{proofthm}

% wat is er en hoe verhoudt het zich tot de scriptie qua resultaten